\documentclass[11pt]{article}
\usepackage[T1]{fontenc}\usepackage{lmodern}\usepackage[sc]{mathpazo}
\usepackage{amsmath,amssymb,amsthm,mathtools,booktabs,array}
\usepackage[paperwidth=7.65in,paperheight=10.4in,inner=.88in,outer=.78in,top=.8in,bottom=.83in]{geometry}
\usepackage{xcolor,fancyhdr}\usepackage[unicode,pdfencoding=auto]{hyperref}
\definecolor{ink}{HTML}{163745}\definecolor{accent}{HTML}{246064}
\hypersetup{colorlinks=true,linkcolor=accent,urlcolor=accent,pdftitle={Critical zeros of automorphic and elliptic-curve L-functions},pdfauthor={Mathematics project}}
\pagestyle{fancy}\fancyhf{}\fancyhead[L]{\small\itshape Critical zeros of automorphic and elliptic-curve $L$-functions}\fancyhead[R]{\small\thepage}\renewcommand{\headrulewidth}{.25pt}\setlength{\headheight}{15pt}
\setlength{\parindent}{1.18em}\setlength{\parskip}{.12em}\setlength{\emergencystretch}{2em}
\newtheorem{theorem}{Theorem}[section]\newtheorem{lemma}[theorem]{Lemma}\newtheorem{corollary}[theorem]{Corollary}
\newtheorem{proposition}[theorem]{Proposition}\theoremstyle{definition}\newtheorem{arinput}[theorem]{Input}
\newtheorem{definition}[theorem]{Definition}\newtheorem{remark}[theorem]{Remark}\newtheorem{example}[theorem]{Example}
\newcommand{\R}{\mathbb R}\newcommand{\C}{\mathbb C}\newcommand{\Z}{\mathbb Z}\newcommand{\Q}{\mathbb Q}\newcommand{\F}{\mathbb F}
\newcommand{\ord}{\operatorname{ord}}\newcommand{\re}{\operatorname{Re}}\newcommand{\im}{\operatorname{Im}}\newcommand{\pos}[1]{\left(#1\right)_{\!+}}
\newcommand{\src}[1]{\hyperlink{src-#1}{\textsf{\scriptsize[#1]}}}
\newcommand{\lead}[1]{\par\smallskip\noindent\textbf{#1}\enspace}
\newcommand{\chapterstart}{\clearpage}\setcounter{tocdepth}{1}
\makeatletter
\renewcommand*\l@section{\@dottedtocline{1}{.6em}{2.1em}}
\renewcommand*\l@part[2]{\addpenalty{-100}\vskip .70em
  \noindent{\color{accent}\sffamily\bfseries #1}\par\nobreak\vskip .18em}
\makeatother
\begin{document}\thispagestyle{empty}
{\color{accent}\small\sffamily MATHEMATICS PROJECT\hfill NOVEMBER 2026}\par
\vspace*{.65in}
{\raggedright\color{ink}\fontsize{29}{34}\selectfont Critical zeros of automorphic and elliptic-curve $L$-functions\par}
\vspace{.22in}{\Large\itshape Transfer theorems, obstructions, and certified zeros\par}
\vspace{.35in}\noindent\rule{\linewidth}{.6pt}\par\bigskip
\noindent For an elliptic curve over a number field the arithmetic critical line is $\Re s=1$. For a standard automorphic function in unitary normalization it is $\Re s=1/2$. Over global function fields the full assertion follows from Lafforgue and Deligne. The all-height assertion over number fields remains open in this investigation, already for the conductor-$11$ elliptic curve.\par\medskip
\noindent We express the reciprocal of a rational elliptic $L$-function as a continued average of quadratic Dirichlet products. A projective family gives half the exponent in a sufficient two-scale estimate. We also characterize the abelian varieties whose first Frobenius coefficients such a rational-parameter quadratic family can represent.\par\medskip
\noindent Multiplication by translates of one fixed Dirichlet completion gives another exact comparison. For a fixed automorphic target of absolute degree at least two, the least nonnegative heat time for an entirely critical divisor, multiplied by the logarithm of the shift, tends to its original horizontal zero displacement. The limit recovers that displacement even when no zero attains it.\par\medskip
\noindent For the undeformed functions we give explicit all-height exclusions, common narrow strips in growing intervals, a uniform finite-height Li implication and individual certificates. A positive kernel detects a hypothetical off-line zero in ordinary Frobenius traces on every sufficiently distant prime window of a suitable fixed ratio.\par
\vfill\noindent{\small Companion files contain the source, interval verifiers, exact certificate data, and a record of the proof attempts.}\clearpage
\section*{Principal results}
\pdfbookmark[1]{Principal results}{reader-guide}
We write $c=1/2$ in the unitary variable and $c=1$ in the arithmetic elliptic variable. In either case $\Delta$ is the supremum of the horizontal distances of the nontrivial zeros from the critical line. ERH for an individual function is the assertion $\Delta=0$. The number-field investigation below leaves that assertion open even for the rational curve of conductor eleven.

\subsection*{Three proofs to read first}
\noindent\textbf{A direct transfer from quadratic Dirichlet characters.} Theorem~\ref{qcm:inverse} writes $1/L(E,1+s)$, up to a holomorphic nonzero factor, as a continued character average. The projective two-scale construction proves an error $O(N^{-2u+\epsilon})$; $O(N^{-4u})$ would exclude all unitary zeros to the right of $1-u$ for $0<u\le1/4$ (Theorem~\ref{pc:main}). Theorems~\ref{albs:rigid} and~\ref{albs:hyper} characterize which abelian varieties can occur in a rational-parameter quadratic-character average. These proofs are collected in \emph{Quadratic character transfer}.

\medskip\noindent\textbf{How a Dirichlet translate records a full automorphic divisor.} Form the completed product with $\Lambda(s+i\tau,\chi_{-4})$. At Gaussian heat time $a/\log\tau$ its entire zero strip tends in width to $(\Delta-a)_+$, measured from its center. The rescaled time at which every zero is on the line tends to $\Delta$ along every sequence $\tau\to\infty$. On fixed off-line disks, a quotient converges to the translated target function and preserves multiplicities. Theorem~\ref{hol:profile} proves all three assertions. Corollaries~\ref{oldheat:localvector} and~\ref{hol:abelianfield} identify the associated automorphic vectors and ground fields. The short paper \emph{Recovering the critical zero strip under heat deformation} contains the proofs.

\medskip\noindent\textbf{An ordinary prime-counting criterion for a fixed curve.} For any prescribed $H(Y)\to\infty$, one-sided cancellation of the raw elliptic prime sums throughout arbitrarily distant intervals $[Y,YH(Y)]$ is equivalent to ERH. This remains true in terms of the product of the group orders $\#E(\F_p)/p$. A single computable set of $\sim2\sqrt{x}/\log x$ primes up to $x$ retains the full right-hand zero divisors, with multiplicities, of every rational elliptic curve. See Theorem~\ref{rawfilter:main} and Corollary~\ref{sample:divisor} or the paper \emph{Compact prime windows}.

\subsection*{The original zeros and the finite-height information}
\noindent For the conductor-eleven curve, all $3\,848\,682$ zeros through $|\gamma|\le524\,288$ are simple and lie on the line. Above height $262\,144$ the explicit exclusion $1/(13.518\log|\gamma|)<\Re\rho<1-1/(13.518\log|\gamma|)$ holds in the unitary coordinate (Section~\ref{exstrip:chapter}). For every fixed non-CM rational elliptic curve the rank-one input sharpens the denominator to $12.335\log|\gamma|$ at sufficiently large heights (Corollary~\ref{rqstrip:rankone}). For conductor eleven the $12.335$ bound holds for $\log|\gamma|\ge300000$. The same calculation yields explicit exclusions for every elliptic curve over every real quadratic field. For non-CM modular curves the effective formula extends to totally real fields of any degree (Theorem~\ref{rqstrip:effective} and Corollary~\ref{rqstrip:arbitraryScalar}). A sixth symmetric-power argument sharpens the eventual coefficient to $12.326\,[K:\Q]$ for the latter curves (Theorem~\ref{rqstrip:six}). Over arbitrary number fields the degree-two theorem allows arbitrary central character, with no global choice of a square root (Theorem~\ref{rqstrip:central}). Theorem~\ref{rqstrip:arbitraryrank} gives explicit asymptotic coefficients in every rank, including $10.919$ times the absolute degree in the essentially self-dual case. Potential automorphy and Brauer induction give high-ordinate zero- and pole-free regions for every elliptic curve over a totally real or CM field, even before modularity over the original field is known (Theorem~\ref{rqstrip:potential}). The index also lists all $38\,042$ rational isogeny classes of conductor below $10\,000$, with exact central orders and $4\,256\,834$ distinct positive zeros through height $64$. Higher-rank finite verifications include a non-CM fifth symmetric power and a mixed Rankin--Selberg tensor.

\medskip\noindent For standard and Rankin--Selberg functions over number fields, verification through $T$ forces strict positivity of all Li coefficients through $T^2$ (Theorem~\ref{li:everyheight}); the unknown high zeros enter the proof through a positive tail estimate. The conductor-eleven example has at least $10\,035\,901\,614\,218$ positive initial coefficients. The short paper \emph{Uniform Li criterion} gives the proof. Section~\ref{hlc:section} proves that an entire locally narrow strip for the original zeros becomes a line in a smaller window at a heat time on the scale of the square of its width. For every fixed entire standard completion of absolute degree $d$, Theorem~\ref{hdp:fulldisk} places all zeros of its $j$th ordinary derivative exactly on the line through height $e^{Hj}$ for all large $j$ and every $0<H<4/d$.

\medskip\noindent For CM elliptic curves, Hu's Hecke-character input gives $5/8\le\Re s\le11/8$ in the arithmetic variable at every height (Section~\ref{recentcm:section}). For global function fields the full critical-line theorem follows from Lafforgue and Deligne (Sections~\ref{ff:global}--\ref{ff:elliptic}). Section~\ref{close:main} identifies the remaining characteristic-zero estimates.

\clearpage
\section*{Numerical results at a glance}
\pdfbookmark[1]{Numerical results at a glance}{numerical-index}
Each height \(T\) in the table means the entire closed band \(|\Im s|\le T\). Counts include both signs of the ordinate and the center, with multiplicity. A row marked \emph{simple} has every zero in the band on the critical line, each of order one. Degree means the degree of the \(L\)-function over \(\Q\).
\begin{center}\small\renewcommand{\arraystretch}{1.18}\begin{tabular}{@{}lrrrr@{}}\toprule Function & Degree & $T$ & Zeros & Multiplicity\\\midrule
$11.\mathrm{a1}$ & 2 & $524\,288$ & $3\,848\,682$ & simple\\
$32.\mathrm{a}$ & 2 & $131\,072$ & $891\,046$ & simple\\
Ramanujan's $\Delta$ & 2 & $131\,072$ & $746\,456$ & simple\\
$\mathrm{Sym}^{2}(43.\mathrm{a1})$ & 3 & $256$ & $1\,276$ & simple\\
$\mathrm{Sym}^{3}(43.\mathrm{a1})$ & 4 & $256$ & $1\,803$ & simple\\
$\mathrm{Sym}^{4}(43.\mathrm{a1})$ & 5 & $32$ & $188$ & simple\\
$\mathrm{Sym}^5(11.\mathrm{a1})$ & 6 & $132$ & $1\,023$ & simple\\
$11.\mathrm{a1}\times15.\mathrm{a1}$ & 4 & $256$ & $1\,716$ & simple\\
$11.\mathrm{a1}\otimes\chi_5$ & 2 & $4\,096$ & $21\,613$ & simple\\
$11.\mathrm{a1}/\Q(\zeta_5)$ & 8 & $4\,096$ & $82\,256$ & simple\\
$234446.\mathrm{a1}$ & 2 & $2\,048$ & $14\,302$ & four in disk\\
\bottomrule\end{tabular}\end{center}
For \(234446.\mathrm{a1}\), the \(14\,298\) zeros outside the central disk are simple and on \(\Re s=1\). The four inside it have total multiplicity four. Their exact positions remain undecided. See Section~\ref{r4:chapter}.
\paragraph{Complete conductor ranges.} Each isogeny class contributes a single elliptic \(L\)-function, regardless of how many rational Weierstrass models represent it. The centers are counted with their exact orders.
\begin{center}\small\begin{tabular}{@{}rrrrr@{}}\toprule Maximum $N$ & Classes & $T$ & Positive simple & Total zeros\\\midrule
$1\,000$ & $2\,463$ & $256$ & $1\,163\,668$ & $2\,328\,496$\\
$9\,999$ & $38\,042$ & $64$ & $4\,256\,834$ & $8\,537\,231$\\
\bottomrule\end{tabular}\end{center}
\paragraph{How the counts become proofs.} A closed contour gives the number of zeros with multiplicity. An equal number of disjoint sign brackets on the line accounts for all of them, proving simplicity. All $4\,256\,834$ positive ordinates in the conductor-$9999$ census are also mutually distinct across the whole family. A multiple central zero calls for an additional vanishing argument. Section~\ref{cens:chapter} gives those arguments for the census. Section~\ref{gauss:chapter} gives a second counting method and tests isolated bands at much higher ordinates.
\clearpage\tableofcontents

\clearpage
\phantomsection\addcontentsline{toc}{part}{I. Prime sums and common height windows}
\noindent{\small\sffamily\color{accent}I. Prime sums and common height windows}\par
\section{The extension problem}\label{intro:target}
Let $E$ be an elliptic curve over $\Q$, and write $N$ for its conductor. At a good prime $p$ put $a_p=p+1-\#E(\F_p)$. The requested assertion is that the nontrivial zeros of
\[
 L(E,s)=\prod_{p\nmid N}(1-a_pp^{-s}+p^{1-2s})^{-1}
          \prod_{p\mid N}L_p(E,s)
\]
have real part one. The more general assertion concerns the standard $L$-function of a primitive unitary cuspidal representation of $\mathrm{GL}_n(\mathbb A_F)$ for a number field $F$. Its unitary critical line has real part one half. Modularity over $\Q$ identifies the two conventions by $L(s,\pi_E)=L(E,s+1/2)$. The extra real half in the elliptic equation is the weight of each Frobenius eigenvalue at a prime. Extending elliptic curves to other number fields requires care about analytic continuation; we state the hypotheses in the relevant results.

The starting theorem of Alp\"oge, checked in November 2026, gives $L(\chi,s)\ne0$ for $\Re s>1/2$ for every rational Dirichlet character, excluding the pole of the principal character. Its quantifier is visible in the exported declaration \texttt{LFunction\_ne\_zero\_of\_half\_lt\_re}. We invoke that theorem in the results that factor into Dirichlet $L$-functions. Andrew Sutherland's \href{https://math.mit.edu/~drew/grh/}{independent verification record} gives its exact statement and the \href{https://leanprover-community.github.io/blog/posts/riemann-hypothesis/}{Lean and Mathlib joint statement} describes the kernel checks. A source comparison for this project matched all 175,187 original Lean source files against a pinned mirror. The comparison did not repeat the full kernel build.

\paragraph{The missing implication.}
The input theorem restricts weighted prime sums in which a prime is seen through its congruence class. The coefficient needed here is $a_p/\sqrt p$. Even the fixed CM curve $y^2=x^3-x$ assigns this coefficient from the angle of a Gaussian prime. A finite congruence class gives no such angle. For non-CM curves the distribution of this coefficient on a fixed abelian class has positive variance. Replacing the label by a growing torsion field recovers the trace, but the required error must then be uniform as the field grows with the prime cutoff. Section~\ref{geo:bridge} gives an explicit inequality sufficient for this step. Sections~\ref{aw:section} and~\ref{ep:primes} develop equivalent prime-sum tests under their stated hypotheses. There is no proof here of the required cancellation for a fixed arbitrary elliptic curve at unbounded scales.

For an elliptic curve over $\Q$ an exact form of the remaining estimate is especially short. Choose any fixed $H(Y)\to\infty$, however slowly, for example $H(Y)=\log\log(3Y)$. For every $\varepsilon>0$ there must be a $C_\varepsilon$ and arbitrarily large $Y$ such that
\[
  \sum_{\substack{p\le x\\p\nmid N}}(p+1-\#E(\F_p))\ge -C_\varepsilon x^{1+\varepsilon}
        \qquad(Y\le x\le YH(Y)).
\]
Theorem~\ref{rawfilter:main} proves the equivalence. Corollary~\ref{rawfilter:pointproduct} gives an equivalent upper bound on a product of group orders. A positive average of ordinary prime cutoffs detects a hypothetical off-line zero in every sufficiently distant interval of a suitable fixed ratio. A single computable set with $\sim2\sqrt x/\log x$ primes up to $x$ retains the whole right-half zero divisor for every rational elliptic curve (Corollary~\ref{sample:divisor}). It is selected by balancing an increasing list of Frobenius traces. The dependence of $Y$ on the \emph{fixed} curve is essential.

\paragraph{Unconditional conclusions.}
The three kinds of zero statements in the paper require different evidence. An argument-principle count and an equal number of certified sign changes locate \emph{every} zero up to a finite height for concrete characteristic-zero functions. A root-number calculation can force central zeros at \emph{every} degree of an infinite tower, with multiplicities bounded below, without locating the other zeros. A product identity valid at every place can transfer the critical-line statement at \emph{all} heights, provided each factor has the required zero theorem. Over a global function field, purity supplies the all-height result in every cuspidal rank. None of the finite-height calculations is used to infer a statement at unbounded height.

There is an all-height approximation to ERH in characteristic zero at every fixed absolute degree $d$. Under the real heat deformation of the completion, its least physical time $\lambda_\pi$ with all zeros on the line obeys $0\le\lambda_\pi\ll_d1/\log(3+\mathcal Q_*)$, where $\mathcal Q_*$ is the analytic conductor minimized over norm twists (Theorem~\ref{hc:automorphic}). Thus every fixed positive time suffices at every height for every sufficiently large conductor. ERH still requires time zero for the individual function.

An unconditional statement for the undeformed functions applies to every cuspidal $\mathrm{GL}_2/\Q$ function, including all rational elliptic curves. For almost every $v\in[1,2]$, the height intervals centered at $v2^j$ of radius $\exp(\sqrt{\log 2^j})$ eventually contain only zeros at distance at most $(1+o(1))/\sqrt{\log 2^j}$ from the critical line, for each fixed function. Each interval contains an increasing number of zeros. Sections~\ref{common:section} and~\ref{ang:chapter} derive this simultaneous conclusion by zero density and prove near-line behavior at polynomially growing heights in CM weight families.

\paragraph{How to read the proofs.}
For an individual rational elliptic curve start with Section~\ref{num:chapter}, or Section~\ref{r2:chapter} for a curve with a double central zero. The analytic bridge from prime counts to possible off-axis zeros can be read independently of the Fourier computations. Section~\ref{zf:chapter} uses symmetric powers to obtain an unconditional bound at every height, including real zeros near the right edge. Section~\ref{ebc:chapter} distinguishes passage to a nonabelian field from replacing a representation by its abelianization. The closing geometric result explains where a weight theorem over a finite field completes exactly the step still missing over a number field. We cite established inputs at the place they enter. An explicit specialization or a certificate newly calculated for this project carries no claim of historical priority.

\subsection*{Four zero patterns in the arithmetic coordinate}
\begin{center}\setlength{\unitlength}{.96pt}\begin{picture}(415,94)
\put(248,14){\color{black!24}\line(0,1){68}}
\put(5,77){\makebox(75,0)[l]{$11\mathrm a1$}}
\put(98,77){\color{black!20}\line(1,0){300}}
\put(152.624,77){\color{accent}\circle*{3.5}}
\put(118.889,77){\color{accent}\circle*{3.5}}
\put(343.376,77){\color{accent}\circle*{3.5}}
\put(377.111,77){\color{accent}\circle*{3.5}}
\put(5,59){\makebox(75,0)[l]{$37\mathrm a1$}}
\put(98,59){\color{black!20}\line(1,0){300}}
\put(172.953,59){\color{accent}\circle*{3.5}}
\put(144.945,59){\color{accent}\circle*{3.5}}
\put(127.785,59){\color{accent}\circle*{3.5}}
\put(99.004,59){\color{accent}\circle*{3.5}}
\put(323.047,59){\color{accent}\circle*{3.5}}
\put(351.055,59){\color{accent}\circle*{3.5}}
\put(368.215,59){\color{accent}\circle*{3.5}}
\put(396.996,59){\color{accent}\circle*{3.5}}
\put(248,59){\color{accent}\circle*{3.5}}
\put(5,41){\makebox(75,0)[l]{$389\mathrm a1$}}
\put(98,41){\color{black!20}\line(1,0){300}}
\put(204.858,41){\color{accent}\circle*{3.5}}
\put(181.874,41){\color{accent}\circle*{3.5}}
\put(161.151,41){\color{accent}\circle*{3.5}}
\put(143.066,41){\color{accent}\circle*{3.5}}
\put(135.907,41){\color{accent}\circle*{3.5}}
\put(118.575,41){\color{accent}\circle*{3.5}}
\put(103.503,41){\color{accent}\circle*{3.5}}
\put(291.142,41){\color{accent}\circle*{3.5}}
\put(314.126,41){\color{accent}\circle*{3.5}}
\put(334.849,41){\color{accent}\circle*{3.5}}
\put(352.934,41){\color{accent}\circle*{3.5}}
\put(360.093,41){\color{accent}\circle*{3.5}}
\put(377.425,41){\color{accent}\circle*{3.5}}
\put(392.497,41){\color{accent}\circle*{3.5}}
\put(248,41){\color{accent}\circle*{3.5}}
\put(248,41){\color{accent}\circle{7.5}}
\put(5,23){\makebox(75,0)[l]{$11\mathrm a1\otimes\chi_5$}}
\put(98,23){\color{black!20}\line(1,0){300}}
\put(217.192,23){\color{accent}\circle*{3.5}}
\put(184.559,23){\color{accent}\circle*{3.5}}
\put(166.253,23){\color{accent}\circle*{3.5}}
\put(149.746,23){\color{accent}\circle*{3.5}}
\put(142.773,23){\color{accent}\circle*{3.5}}
\put(124.849,23){\color{accent}\circle*{3.5}}
\put(114.661,23){\color{accent}\circle*{3.5}}
\put(265.216,23){\color{accent}\circle*{3.5}}
\put(292.237,23){\color{accent}\circle*{3.5}}
\put(307.492,23){\color{accent}\circle*{3.5}}
\put(321.438,23){\color{accent}\circle*{3.5}}
\put(347.589,23){\color{accent}\circle*{3.5}}
\put(365.458,23){\color{accent}\circle*{3.5}}
\put(386.814,23){\color{accent}\circle*{3.5}}
\put(98,8){\makebox(0,0){\footnotesize \(-10\)}}
\put(173,8){\makebox(0,0){\footnotesize \(-5\)}}
\put(248,8){\makebox(0,0){\footnotesize \(0\)}}
\put(323,8){\makebox(0,0){\footnotesize \(5\)}}
\put(398,8){\makebox(0,0){\footnotesize \(10\)}}
\end{picture}\end{center}
Each horizontal position is the ordinate $t$ of a zero at $s=1+it$, on the range $-10<t<10$. Dots come from certified intervals; the ring marks a double zero. The last row twists the conductor-$11$ form by the quartic character with $\chi_5(2)=i$. Its two halves differ. Its first positive zero lies between $1.1447753$ and $1.1447754$, while the nearest negative zero lies between $-2.0532709$ and $-2.0532708$. The picture is an index to the proofs and their brackets.

\paragraph{Where a new input changes the answer.}
A unitary Hecke character over a CM field can have an infinite-order angular part. Section~\ref{recentcm:section} examines recent zero-free half-planes for that actual character, including the formal evidence for the boundary $7/8$. It translates that input to an arithmetic strip $5/8\le\Re s\le 11/8$ for a CM elliptic curve, at every height. For a non-CM curve the connected image is different. Section~\ref{fc:section} proves an optimal obstruction to its recovery by bounding residual character modes separately. Section~\ref{it:section} uses the original Dirichlet theorem positively to construct twists with prescribed small Euler data and forced central zeros.

Two arithmetic consequences explain why a weaker zero result can still answer a sharp question about individual curves. If $E_1,E_2/\Q$ are nonisogenous, then their point counts differ at a common good prime $p\ll[\log(2\operatorname{rad}(N_1N_2))]^8$, with an absolute constant (Section~\ref{ld:section}). For one non-CM curve, the number of primes up to $x$ having a prescribed nonzero trace is $O_E(x^{19/20}/(\log x)^{3/5})$ (Section~\ref{conc:section}). Both statements use the verified rank-one boundary in number fields. Neither requires ERH for the elliptic curve. The finite-computation chapters go in a different direction by locating individual zeros of genuinely cuspidal non-CM functions of ranks two through six, including a mixed elliptic tensor.

Three public expositions help identify the arithmetic step in that comparison. Rubinstein's working account of the zeta specialization reduces its main inequality to thirteen further inputs. Slijep\v{c}evi\'c traces the Dirichlet argument through weighted prime sums, Hermitian lattices from theta lifts, and bounds on Tate--Shafarevich groups of Tate twists. His final section specifies the gaps in that exposition. Jaswal follows the declarations in the formal archive. For the established Dirichlet input we cite the verified theorem itself. Section~\ref{pp:section} tests the prime-power identity on a non-CM elliptic curve, where the distinction between an exact signed identity and an acceptable analytic error is decisive.
\href{https://www.math.uwaterloo.ca/~mrubinst/alpoge_claude_grh/}{Rubinstein's working guide.}
\href{https://hexagonmath.org/pdf/2611.00812}{Slijep\v{c}evi\'c's outline.}
\href{https://hexagonmath.org/2611.00799v1}{Jaswal's guide to the archive.}

\clearpage
\section{A square-root window for a standard automorphic function}\label{aw:section}

Let $F$ be a number field and $\pi$ a unitary cuspidal automorphic representation of $\mathrm{GL}_n(\mathbb A_F)$. We first take $n\ge2$. Write $q=\mathrm N\mathfrak p$, the absolute norm of a prime ideal. Let $\alpha_j(\mathfrak p)$ be the inverse roots of the local Euler polynomial in the unitary normalization, with zeros appended to make a list of length $n$ at a ramified prime. Put
\[
 a(\mathfrak p,k)=\sum_{j=1}^n\alpha_j(\mathfrak p)^k,\qquad
 S_\pi(x)=\sum_{q^k\le x}a(\mathfrak p,k)\log q\,(x/q^k-1).
\]
Every sum over $q^k$ here runs over prime ideals and integers $k\ge1$, retaining multiplicities when distinct ideals have the same norm. For an elliptic curve with good reduction, $a(\mathfrak p,1)=(q+1-\#E(\F_q))/\sqrt q$. We will show exactly how little of this local information remains in an equivalent form of GRH.

Fix $b>1$, for example $b=2$. Use the continuous taper
\[
 w_b(v)=\begin{cases}
 b-1,&0<v\le1,\\
 b/v-1,&1<v<b,\\
 0,&v\ge b.
 \end{cases}
\]
Our observable is the real quantity
\begin{equation}
 T_{\pi,b}(x)=\sum_{\mathfrak p,k}
  \operatorname{Re}a(\mathfrak p,k)\log q\;w_b(q^k/x).
 \label{aw:taper}
\end{equation}
It can be evaluated from finite Euler data up to $bx$. It assigns a constant weight to all the terms up to $x$ and softens the endpoint between $x$ and $bx$. For complex $\pi$, taking real parts preserves all the zeros that must be detected.

\begin{theorem}[One-sided automorphic criterion]\label{aw:criterion}
Fix $\delta>0$ and $b>1$. The following statements are equivalent for $\pi$.
\begin{enumerate}
\item Every zero of the completed standard $L$-function has real part $1/2$.
\item $|T_{\pi,b}(x)|\ll_{\pi,b}\sqrt x$ for every $x\ge2$.
\item For every $\varepsilon>0$ there is a constant $C_\varepsilon$ and there are arbitrarily large $Y$ for which $T_{\pi,b}(x)\ge-C_\varepsilon x^{1/2+\varepsilon}$ for all $Y\le x\le Y^{1+\delta}$.
\item For every $\eta>0$ there are arbitrarily large $Y$ for which
\[
 \int_Y^{Y^{1+\delta}}
       \max\{-T_{\pi,b}(x),0\}\frac{dx}{x^{3/2}}\le Y^\eta.
\]
\end{enumerate}
The same four equivalences hold with $-T_{\pi,b}(x)$ replaced throughout by
\[
 H_\pi(x)=\operatorname{Re}\{S_\pi(x)+x(L'/L)(1,\pi)\},
\]
using upper bounds for $H_\pi$. The constants may depend on the representation and the field. A single fixed $b$ and a single fixed $\delta$ suffice.
\end{theorem}

We give a proof over $F$ directly. In particular the zeros tested belong to the function over $F$ even when no automorphic induction from $F$ to $\Q$ has been established. The following standard inputs suffice. Godement--Jacquet theory gives the continuation, the entire completion of order one in rank at least two, and the functional equation relating $\pi$ and its contragredient $\widetilde\pi$. Jacquet--Shalika nonvanishing gives regularity and nonvanishing on $\operatorname{Re}s=1$. Unitarity gives
$L(s,\widetilde\pi)=\overline{L(\bar s,\pi)}$. Write
\[
 \Lambda(s,\pi)=Q_\pi^{s/2}L(s,\pi)
        \prod_{v\mid\infty}\prod_{j=1}^n
             \Gamma_v(s+\mu_j(v)),
 \quad Q_\pi=|D_F|^n\mathrm N\mathfrak q_\pi ,
\]
where $D_F$ is the discriminant, $\mathfrak q_\pi$ the conductor, $\Gamma_{\R}(s)=\pi^{-s/2}\Gamma(s/2)$ and $\Gamma_{\C}(s)=2(2\pi)^{-s}\Gamma(s)$. For $v$ real or complex take the corresponding gamma function. The bounds of Luo--Rudnick--Sarnak and their local extensions give
\begin{equation}
 |\alpha_j(\mathfrak p)|\le q^{\vartheta_n},\qquad
 \operatorname{Re}\mu_j(v)\ge-\vartheta_n,\qquad
 \vartheta_n=\tfrac12-\frac1{n^2+1}.
 \label{aw:lrs}
\end{equation}
These bounds also hold at ramified places for the indicated parameters. The choices $7/64$, $5/14$, $9/22$ are available in degrees two, three and four, respectively. A twist by a unitary power of the norm translates the imaginary parameter in the analytic formulas and preserves these bounds. See Humphries--Thorner, \emph{Zeros of Rankin--Selberg $L$-functions in families}, \S3, equations (3.1)--(3.4), for the local and global normalizations and references. \href{https://arxiv.org/pdf/2103.05634}{Humphries--Thorner.}

\subsection{The meromorphic transform}

For $x\ge1$ the sums defining $S_\pi$ are continuous, with a possible change of derivative when $x$ passes an ideal norm. They grow at most as $O_{F,n}(x^2)$ by \eqref{aw:lrs}. Indeed, if $d=[F:\Q]$, the sum of $\log q$ for prime ideal powers of norm at most $x$ is at most $d\sum_{p^m\le x}\log p\ll_F x$: for a fixed rational power $p^m$ the sum of residue degrees $f_{\mathfrak p}$ over $\mathfrak p\mid p$ such that $f_{\mathfrak p}\mid m$ is at most $d$. Here we used Chebyshev's elementary bound. Partial summation now gives the claimed growth.

For $\operatorname{Re}s>2$, denote Mellin transformation on $[1,\infty)$ against $x^{-s-1}dx$ by $\mathcal M$. Termwise integration gives
\[
 \mathcal M S_\pi(s)=\frac{-(L'/L)(s,\pi)}{s(s-1)}.
\]
In fact the integral of $(x/m-1)$ on $x\ge m$ is $m^{-s}/[s(s-1)]$. Taking the average of this meromorphic expression for $\pi$ and $\widetilde\pi$, and adding $\operatorname{Re}(L'/L)(1,\pi)/(s-1)$, gives $\mathcal M H_\pi$. The pole at $1$ is removable. Only at this step have we used the boundary nonvanishing. At any zero $\rho$ of the completed function or its dual with $\operatorname{Re}\rho>\vartheta_n$, the residue is
\begin{equation}
 \frac{m_\pi(\rho)+m_{\widetilde\pi}(\rho)}
              {2\rho(1-\rho)},\qquad
 m_\pi(\rho)=\operatorname{ord}_\rho\Lambda(s,\pi).
 \label{aw:residue}
\end{equation}
The numerator is positive there. In this half-plane the gamma factors are regular and have no zeros. Farther left, a pole of a gamma factor contributes a trivial zero to the finite $L$-function, sometimes with positive real part. The cut at $\vartheta_n$ avoids dropping such a zero from the transform.

To find the transform of the finite observable, the identity
\begin{equation}
 -T_{\pi,b}(x)=bH_\pi(x)-H_\pi(bx)
  =\operatorname{Re}\{bS_\pi(x)-S_\pi(bx)\}
 \label{aw:difference}
\end{equation}
follows by separating $q^k\le x$ and $x<q^k\le bx$. Thus $\mathcal M(-T_{\pi,b})$ equals $(b-b^s)\mathcal M H_\pi(s)$ plus the entire function $b^s\int_1^b H_\pi(y)y^{-s-1}dy$. At every pole in the open strip the multiplier is nonzero, since $|b^s|<b$. At a real pole there its value is positive, as is \eqref{aw:residue}. These two sign computations explain the direction of the one-sided prime inequality. For a nonreal representation the orders in \eqref{aw:residue} add at the same $s$. There is no opportunity for the conjugate functions to cancel a pole.

\subsection{Why selected windows detect every pole}

For completeness we isolate the complex-analytic implication. It applies with either of the transforms above. Let $a<\theta<1$. Suppose $g(u)$ is real, locally integrable on $u\ge0$, and $|g(u)|\ll e^{A u}$ for some $A>1$. Its Laplace transform $P$ continues meromorphically to $\operatorname{Re}s>a$, with only simple poles there, all in $\operatorname{Re}s<1$. Suppose all real poles in this half-plane have positive residues. If for every $\eta>0$ arbitrarily large $U$ satisfy
\begin{equation}
 \int_U^{(1+\delta)U}\max\{g(u),0\}e^{-\theta u}du\le e^{\eta U},
 \label{aw:abstractwindow}
\end{equation}
then $P$ has no pole in $\operatorname{Re}s>\theta$. The following proof also permits the windows for different $\eta$ to be unrelated.

For $\sigma>A$, set $B_k(z)=(-1)^kP^{(k)}(z)/k!$ and $K_k(u)=u^k e^{-\sigma u}/k!$. The transform on its domain of integration gives
\[
 |B_k(\sigma+it)|+B_k(\sigma)
   \le 2\int_0^\infty \max\{g(u),0\}K_k(u)du .
\]
Put $c=1+\delta/2$ and $k=\lfloor c(\sigma-A)U\rfloor$. Within the window $K_k(u)e^{\theta u}\le(\sigma-\theta)^{-k}$ by the exponential series. Outside, use the growth bound for $g$. The full integral of $K_k(u)e^{Au}$ equals $(\sigma-A)^{-k-1}$, and at most an $O_\delta(e^{-d_\delta k})$ fraction of this integral lies outside the window, for some $d_\delta>0$. To see this, substitute $v=(\sigma-A)u/k$. The density is proportional to $(ve^{1-v})^k$, whose base is uniquely maximized at $1$. The window endpoints tend to $1/c$ and $(1+\delta)/c$. Off the window the base is bounded by $h<1$ and on a fixed subinterval about $1$ it is bounded below by $h'>h$. On the unbounded tail integrate one factor $ve^{1-v}$ and bound the other $k-1$ by $h^{k-1}$. This proves the claimed fraction. Choose $\sigma$ large enough that
$\log((\sigma-\theta)/(\sigma-A))<d_\delta$ and $c(\sigma-A)>1$. With $E=\sigma-\theta$, the windows in \eqref{aw:abstractwindow} give arbitrarily large $k$, with $U\le k$, for which
\begin{equation}
 |B_k(\sigma+it)|+B_k(\sigma)\ll e^{\eta k}E^{-k}
 \quad\text{for all real }t.
 \label{aw:coeff}
\end{equation}

Suppose a pole has real part greater than $\theta$. Let $D$ be the lesser of $E$ and the distance from $\sigma$ to the nearest pole in the closed disk of radius $E$, with $D=E$ if this disk is empty. If some pole lies at horizontal distance less than $D$ from the line through $\sigma$, choose $t$ so $z=\sigma+it$ has a unique nearest pole $\rho$ at distance $d<D$. Such a choice follows by starting at the height of the given pole and perturbing $t$ slightly. Only finitely many poles enter consideration, and equal distances to a distinct pair impose at most one value of $t$, since all poles lie to the left of the line. Choose $d<h<D$ such that the disk of radius $h$ around $z$ contains only $\rho$. Taylor expansion of its simple principal part and Cauchy's estimate for the remainder give
\[
 |B_k(z)|\ge |\operatorname{Res}_\rho P|d^{-k-1}-O(h^{-k}),
 \qquad |B_k(\sigma)|\ll h^{-k}.
\]
This contradicts \eqref{aw:coeff} for $0<\eta<\log(E/d)$.

Otherwise every pole has horizontal distance at least $D$ from the chosen line. The assumed pole gives $D<E$. A nearest pole to $\sigma$ at distance $D$ must then have zero vertical displacement, hence be the real point $\sigma-D$. Its residue $r$ is positive. An $h$ with $D<h<E$ can be chosen so this is the sole pole in the closed disk about $\sigma$ of radius $h$. Now $B_k(\sigma)=rD^{-k-1}+O(h^{-k})$. Set $t=0$ in \eqref{aw:coeff} and choose $0<\eta<\log(E/D)$ to obtain the contradiction. All disks used in this argument lie strictly in the domain of meromorphy because their left edges are at or to the right of $\theta>a$.

Take $g(u)=-T_{\pi,b}(e^u)$ or $H_\pi(e^u)$, $\theta=1/2$ and $a$ strictly between $\vartheta_n$ and $1/2$. The criterion just proved and \eqref{aw:residue} exclude all completed zeros in $\operatorname{Re}s>1/2$ for both representations. The functional equation excludes the left half of the strip. The completed functions have no zeros outside the open strip. This proves (4)$\Rightarrow$(1). To see (3)$\Rightarrow$(4), take $\varepsilon<\eta/(2(1+\delta))$. The integral is bounded by $C_\varepsilon\delta(\log Y)Y^{(1+\delta)\varepsilon}\le Y^\eta$ on the indicated sequence after discarding a finite initial segment.

\subsection{The converse and the archimedean terms}

The converse follows from an explicit formula useful even when the zeros are off the line. In what follows the sums over completed zeros list multiplicities. For a fixed $\pi$ let $a$ be any number with $\vartheta_n<a<1/2$. Then for $x\ge2$,
\begin{equation}
 S_\pi(x)+x(L'/L)(1,\pi)
     =\sum_{\Lambda(\rho,\pi)=0}
         \frac{x^\rho}{\rho(1-\rho)}+O_{\pi,a}(x^a).
 \label{aw:explicit}
\end{equation}
The series is absolutely convergent for each $x$, and the error estimate holds without a hypothesis on the zeros. In particular it is legitimate to separate the rightmost and leftmost zero contributions even when there is no empty vertical zone between them.

Here are details of the shift. The initial line $\operatorname{Re}s=3$ gives
\[
 S_\pi(x)=\frac1{2\pi i}\int_{3-i\infty}^{3+i\infty}
       \frac{-(L'/L)(s,\pi)x^s}{s(s-1)}ds .
\]
Absolute convergence gives inversion term by term. The kernel has inverse $(y-1)_+$, where the subscript means the positive part, including the value zero at $y=1$. This identity follows by moving the scalar integration line across $1$ and $0$ when $y>1$, or taking the line to the right when $y<1$, with the endpoint given by continuity.

There are $O_\pi(\log(2+|t|))$ completed zeros with imaginary part in $[t-1,t+1]$. For example, the logarithmic derivative of the Hadamard product for $\Lambda$ at $3+it$ has real part
\[
 C_\pi+\sum_\rho
 \frac{3-\operatorname{Re}\rho}
      {(3-\operatorname{Re}\rho)^2+(t-\operatorname{Im}\rho)^2}.
\]
The constant is finite since all zeros are in $0<\operatorname{Re}s<1$ and the canonical product has genus at most one. The Euler series for $L'/L$ and Stirling's estimate for the gamma factors bound the left side by $O_\pi(\log(2+|t|))$. Each indicated nearby zero contributes at least a positive absolute constant. This proves the local zero count and, by summing over unit intervals, the convergence of $\sum_\rho |\rho(1-\rho)|^{-1}$. The compact part of the sum is finite by boundary nonvanishing.

Choose $h$ strictly between $-2$ and $-1$, with $\operatorname{Re}s=h$ avoiding the gamma poles. The possible excluded real parts form a locally finite set. On this line the functional equation and absolute convergence at $1-h-it$ give $(L'/L)(h+it,\pi)\ll_{\pi,h}\log(2+|t|)$. At large positive and negative heights one may join the lines $h$ and $3$ on horizontal segments where $L'/L\ll_{\pi,h}\log^2(2+|t|)$. In fact the local zero count allows selection of a height in each sufficiently large unit interval at distance $\gg_\pi1/\log(2+|t|)$ from all zero ordinates. In the difference of the Hadamard logarithmic derivatives at $s$ and at $3+it$, nearby zeros each contribute $O_\pi(\log(2+|t|))$. At distance between $j$ and $j+1$, $j\ge1$, each contributes $O_h(j^{-2})$ by subtraction of the two fractions. The sum of these tails is $O_{\pi,h}(\log(2+|t|))$ by the same count. Stirling's estimate is uniform on these segments after the finitely many imaginary parts of the gamma poles are passed. Thus the horizontal integrals vanish at infinity.

Crossing $s=1$ contributes $-x(L'/L)(1,\pi)$. At a nonzero finite zero the sign of the residue gives $x^\rho/[\rho(1-\rho)]$ per unit of multiplicity. The order of a finite zero equals the sum of its completed zero order and its archimedean pole order, including at a coincident point. In the strip $h\le\operatorname{Re}s\le3$ the gamma poles form a finite set; their real parts are at most $\vartheta_n$. The point zero contributes $O_\pi(1+\log x)$ since the rational kernel and $L'/L$ have at most simple poles there. Every other gamma contribution is $O_\pi(x^{\vartheta_n})$ for $x\ge2$. The remaining vertical integral is $O_{\pi,h}(x^h)$. These facts prove \eqref{aw:explicit} with the asserted error.

Under GRH the series in \eqref{aw:explicit} is bounded in modulus by $\sqrt x\sum_\rho(1/4+(\operatorname{Im}\rho)^2)^{-1}$. Equation \eqref{aw:difference} proves (1)$\Rightarrow$(2), and (2)$\Rightarrow$(3) follows at once. The same proof works for $H_\pi$. Theorem~\ref{aw:criterion} follows.

\subsection{What a rightmost zero would force on shorter intervals}

The power-length windows used in Theorem~\ref{aw:criterion} allow completely arbitrary zero configurations. For a narrower class of potential configurations we can quantify the supply of large positive values on much shorter intervals.

\begin{theorem}[An attained right edge]\label{aw:edge}
Let $\beta=\sup_\rho\operatorname{Re}\rho$ for the zeros of $\Lambda(s,\pi)$, and suppose some zero has real part $\beta$. There are constants $c>0$, $B>1$ such that every sufficiently large $[X,BX]$ contains an $x$ with $H_\pi(x)\ge c x^\beta$. There are also such constants, depending on $b$, for which some $x\in[X,BX]$ has $T_{\pi,b}(x)\le-c x^\beta$.
\end{theorem}
\begin{proof}
The real parts of the zero set are symmetric about $1/2$ by the functional equation and conjugacy, so $1/2\le\beta<1$. Put $u=\log x$. Divide \eqref{aw:explicit} by $x^\beta$, average with its conjugate equation and let $x\to\infty$. Each zero with smaller real part gives a term tending to zero, dominated in modulus by $1/|\rho(1-\rho)|$. The summable bound proves
\[
 e^{-\beta u}H_\pi(e^u)=P(u)+o(1),\quad
 P(u)=\sum_{\substack{\operatorname{Re}\rho=\beta\\
                   m_\pi(\rho)+m_{\widetilde\pi}(\rho)>0}}
        \frac{m_\pi(\rho)+m_{\widetilde\pi}(\rho)}
        {2\rho(1-\rho)}e^{i(\operatorname{Im}\rho)u}.
\]
Here the displayed sum is over distinct points and includes their indicated weights. $P$ is a real uniformly absolutely convergent Fourier series. Its value at $0$ is positive since
$\operatorname{Re}(\rho(1-\rho))=\beta(1-\beta)+(\operatorname{Im}\rho)^2>0$. Its large positive values recur in every interval in $u$ of a fixed sufficiently large length. Indeed, truncate the series to uniform accuracy less than $P(0)/8$. The closure of the real one-parameter group of its finitely many phases in a torus is a compact group. Finitely many translates by times $s_1,\dots,s_m$ of any neighborhood of the identity cover this group, by density and compactness. For any time $t$, at least one of $t-s_j$ then gives phases inside that neighborhood. Choosing the neighborhood small makes the finite series close to its value at zero. This proves the recurrence with a bounded gap and yields the first assertion.

For $-T_{\pi,b}$ the edge series has each coefficient multiplied by $b-b^\rho$ by \eqref{aw:difference}. All of these factors are nonzero and its constant coefficient, if present, is positive. This real series has some positive value. In fact its mean over $[0,U]$ tends to its constant coefficient by uniform convergence. With a positive constant coefficient the assertion follows immediately. If the constant coefficient is zero and the series were everywhere nonpositive, it would still be strictly negative somewhere, since a nonzero Fourier coefficient is recovered by averaging against its character. Uniform continuity and the recurrence just proved would give fixed-length intervals of negative values bounded away from zero with bounded gaps, contradicting mean zero. Recur at one positive value and absorb the $o(1)$ as above. Exponentiating the bounded gaps gives the claimed $B$.
\end{proof}

It follows that an upper bound for $H_\pi$ of order $x^{1/2+o(1)}$ on arbitrarily distant intervals $[X,M_X X]$ with $M_X\to\infty$ rules out an attained rightmost off-line zero. No uniform distance between that zero and the zeros to its left was assumed. When real parts tend up to a supremum which no zero attains, every individual term in the normalized edge series tends to zero and this shorter-interval argument gives no exclusion. The power-length criterion still applies.

\subsection{Removing higher prime powers without a zero hypothesis}

The tapered sum has no linear centering term. Consequently deleting a set of Euler data costs at most $(b-1)$ times the sum of their absolute coefficients up to $bx$. The original kernel $S_\pi$ requires an additional linear term. For any selected set $\mathcal D$ of pairs $(\mathfrak p,k)$, put
\[
 A_{\mathcal D}(y)=\sum_{\substack{(\mathfrak p,k)\in\mathcal D\\q^k\le y}}
            |a(\mathfrak p,k)|\log q.
\]
If $A_{\mathcal D}(y)\ll y^\tau$ with $0<\tau<1$ then
$c_{\mathcal D}=\sum_{\mathcal D}a(\mathfrak p,k)\log q/q^k$ converges absolutely and the portion $S_{\mathcal D}(x)$ of $S_\pi(x)$ satisfies
\begin{equation}
 S_{\mathcal D}(x)-xc_{\mathcal D}
 =-\sum_{\substack{\mathcal D\\q^k\le x}}a(\mathfrak p,k)\log q
  -x\sum_{\substack{\mathcal D\\q^k>x}}\frac{a(\mathfrak p,k)\log q}{q^k}
 =O(x^\tau).
 \label{aw:delete}
\end{equation}
Indeed the first partial sum has the asserted bound and dyadically splitting the tail gives $O(x^{\tau-1})$ before multiplication by $x$. When changing from the full $S_\pi$ to its retained portion in $H_\pi$, one must \emph{add} $c_{\mathcal D}$ to $(L'/L)(1,\pi)$. This follows from \eqref{aw:delete} and fixes the sign independently of any convention for Frobenius.

For every rank and integer $h\ge2$,
\[
 A_{k\ge h}(y)\ll_{F,n,h} y^{\vartheta_n+1/h},\qquad y\ge2 .
\]
In fact $|a(\mathfrak p,k)|\le n y^{\vartheta_n}$ here. For the purely prime-ideal weight, $k=h$ contributes $O_F(y^{1/h})$ by Chebyshev. The other $k$ contribute $O_F((\log y)y^{1/(h+1)})=O_{F,h}(y^{1/h})$. Thus all $k\ge n^2+1$ can be deleted at cost $O(\sqrt y)$. At a finite fixed set of prime ideals, \eqref{aw:lrs} and a geometric sum give $O_{\pi,n}(y^{\vartheta_n})$ when $\vartheta_n>0$, or $O_\pi(\log(2y))$ if the modulus bound is one. These facts alone leave only finitely many exponents to control at any rank.

\begin{theorem}[Reduction to first powers]\label{aw:prime}
Write $T^{(1)}_{\pi,b}$ for \eqref{aw:taper} restricted to unramified prime ideals and $k=1$. Uniformly for $x\ge2$, the following are unconditional bounds.
\begin{enumerate}
\item If $n=2$, $3$ or $4$, then $T_{\pi,b}(x)-T^{(1)}_{\pi,b}(x)\ll_{\pi,b}\sqrt x$ in absolute value.
\item In every rank satisfying the second-order mean bound \eqref{aw:brumley} below, the difference in absolute value is $O_{\pi,b,\varepsilon}(x^{1/2+\varepsilon})$ for every $\varepsilon>0$. The sharper $O_{\pi,b}(\sqrt x)$ holds if each exterior-power transfer through degree $\lfloor n/2\rfloor$ is an isobaric sum of unitary cusp forms.
\item If $\pi$ is tempered at all unramified finite places, every rank has the bound $O_{\pi,b}(\sqrt x)$. In that case one may additionally delete the unramified prime ideals with residue degree larger than one over $\Q$ at the same cost.
\end{enumerate}
Consequently condition (3) or (4) of Theorem~\ref{aw:criterion} is equivalent to GRH with $T^{(1)}_{\pi,b}$ in place of $T_{\pi,b}$ in each indicated case. Condition (2) makes the same substitution whenever the error is $O(\sqrt x)$.
\end{theorem}
\begin{proof}
We establish the rank two and three estimate using only Rankin--Selberg theory for the given representation. At an unramified prime let $M=\max_j|\alpha_j|$. The multiset is stable under $\alpha\mapsto1/\bar\alpha$. In ranks two and three, if $M>1$, there is exactly one off-circle pair $M e^{it}, M^{-1}e^{it}$. In rank three the one remaining root has modulus one. Hence $M\le|a(\mathfrak p,1)|+1$ in those cases, and for $M=1$ this bound persists. For every $k\ge2$ this gives
\begin{equation}
 |a(\mathfrak p,k)|\le 2n(q^{\vartheta_n})^{k-2}
                       (1+|a(\mathfrak p,1)|^2).
 \label{aw:powertr}
\end{equation}
For unramified primes
\[
 \sum_{q\le z}(1+|a(\mathfrak p,1)|^2)\log q\ll_\pi z
 \qquad(z\ge2).
\]
The first term is Chebyshev's estimate. For the second, the prime coefficient of $-L'/L(s,\pi\times\widetilde\pi)$ at an unramified prime is $|a(\mathfrak p,1)|^2\log q$, and the higher unramified coefficients are nonnegative. The Rankin--Selberg prime number theorem supplies this upper bound. Equivalently apply Wiener--Ikehara to the logarithmic derivative with the ramified factors removed. The incomplete product has a simple pole at $1$, no other zero or pole on the boundary, and nonnegative logarithmic-derivative coefficients with the minus sign. The removed factors have no poles on that boundary since their reciprocal roots have modulus at most $q^{2\vartheta_n}<q$. See Humphries--Thorner, \emph{Towards a $\mathrm{GL}_n$ variant of the Hoheisel phenomenon}, Theorem 2.1 and its fixed-representation specialization. \href{https://arxiv.org/pdf/2004.09488}{Rankin--Selberg estimates.}

The sum in \eqref{aw:powertr} for a fixed $k$ over $q^k\le y$, weighted by $\log q$, is therefore
\[
 \ll_\pi y^{\vartheta_n+(1-2\vartheta_n)/k}.
\]
The constant is independent of $k$. At $k=2$ the exponent is exactly $1/2$. For $k\ge3$ it is at most $\vartheta_n+(1-2\vartheta_n)/3<1/2$ and there are $O(\log(2y))$ possible exponents. The strict gap absorbs this factor. The finitely many ramified primes satisfy the earlier bound. This proves the first assertion in ranks two and three.

In rank four there may be two large roots whose phases cancel in the first trace. Use the trace of the exterior square to control this case. More generally, at an unramified prime let $e_j$ be the elementary symmetric functions of the $n$ roots. The central character is unitary, so $|e_n|=1$, and the unitary multiset identity gives $e_{n-j}=e_n\overline{e_j}$. The characteristic polynomial implies
\[
 M^2\le4\left(1+\sum_{j=1}^{\lfloor n/2\rfloor}|e_j|^{2/j}\right).
\]
Indeed, if a root had magnitude greater than $2\max(1,\max_{j}|e_j|^{1/j})$, dividing the polynomial by its leading term would bound $1$ by a sum strictly below $\sum_{j=1}^{n}2^{-j}<1$. The symmetry just established removes the indices above $n/2$. This is the local estimate of Brumley, Proposition~5.

Suppose that for $1\le j\le\lfloor n/2\rfloor$ the corresponding exterior-power representation agrees outside finitely many primes with an isobaric sum of unitary cuspidal representations. Then $e_j$ is the sum of their ordinary traces at those primes. The fixed-representation Rankin--Selberg bound used above, Cauchy--Schwarz in the finitely many constituents, and $|t|^{2/j}\le1+|t|^2$ show that $\sum_{q\le z}M(\mathfrak p)^2\log q\ll_\pi z$ outside these finitely many places. The sum at a fixed $k\ge2$ is bounded as before, now using $|a(\mathfrak p,k)|\le n q^{(k-2)\vartheta_n}M(\mathfrak p)^2$. This yields $O_\pi(\sqrt y)$ for all deleted powers. In rank four only the exterior square in addition to $\pi$ is needed. Kim's exterior-square transfer is a special case of the split $\mathrm{GSpin}_6$ transfer. Asgari--Shahidi's description of its image supplies the unitary cuspidal constituents. Agreement at the unramified places outside a finite set suffices, since all powers of any exceptional finite prime have already been bounded. This proves the remaining rank-four assertion. \href{https://math.okstate.edu/people/asgari/Files/gspin21.pdf}{Asgari--Shahidi, \emph{Image of functoriality for general spin groups}, Theorem 4.26 and \S5.3.}

In the second assertion the input is
\begin{equation}
 \sum_{\mathfrak p\text{ unramified}}
            \frac{M(\mathfrak p)^2}{q^{1+u}}<\infty
               \qquad (u>0).
 \label{aw:brumley}
\end{equation}
Brumley's Theorem 1 proves this for $n\le4$, using Kim's exterior-square transfer in rank four. In higher rank its hypotheses give further individual examples when the indicated exterior-power lifts are available. To apply the displayed sum, take $s=1/2+\varepsilon$ real. Geometric summation gives
\[
 \sum_{\mathfrak p\text{ unramified}}\sum_{k\ge2}
   \frac{|a(\mathfrak p,k)|\log q}{q^{ks}}
 \le n\sum_{\mathfrak p\text{ unramified}}
    \frac{M(\mathfrak p)^2\log q}{q^{2s}}
        \frac1{1-M(\mathfrak p)q^{-s}}<\infty .
\]
The denominators on the right are bounded below by $1-2^{\vartheta_n-s}>0$, and $\log q\ll_\varepsilon q^{\varepsilon}$. Bounding each term up to $y$ by $y^s$ times its weighted counterpart gives $A_{k\ge2}(y)\ll_{\pi,\varepsilon}y^{1/2+\varepsilon}$. For the consequence concerning the integral in (4), choose $\varepsilon$ strictly smaller than $\eta/(2(1+\delta))$ and absorb the logarithmic factor; the positive part changes by at most the modulus of the deleted contribution. \href{https://doi.org/10.1007/s00013-005-1632-3}{Brumley, \emph{Second order average estimates on local data of cusp forms}.}

For the last assertion, $|a(\mathfrak p,k)|\le n$ at the indicated primes. The earlier prime-power count with $h=2$ is $O_F(\sqrt y)$. At $k=1$, prime ideals of residue degree $f\ge2$ necessarily lie above rational primes $p\le\sqrt y$. For any fixed such $p$ the sum of the $\log q$ is at most $[F:\Q]\log p$. Chebyshev gives $O_F(\sqrt y)$ for their aggregate weight. This proves the claim. The comparison of the two versions of the criterion follows from the triangle inequality.
\end{proof}

\section{A criterion using only finite-field point counts}\label{aw:pointsection}

A curve can be tested by the ordinary sum of its Frobenius traces. We first eliminate the taper, then the factor $\log q$, and finally the square-root normalization of each trace. Removing $\log q$ changes a pole of a logarithmic derivative into a logarithmic singularity. Its leading Taylor coefficient still has the sign needed for a one-sided estimate.

For $\ell\ge0$ define real cumulative sums over prime ideal powers by
\[
 \Theta_{\pi,\ell}(x)=\sum_{q^k\le x}\operatorname{Re}a(\mathfrak p,k)
            q^{k\ell}\log q,\qquad
 B_{\pi,\ell}(x)=\sum_{q^k\le x}
            \frac{\operatorname{Re}a(\mathfrak p,k)}{k}q^{k\ell}.
\]
In \(B_{\pi,0}\) a first-power contribution is just the real part of a Satake trace. In $B_{\pi,1/2}$ it is the usual trace of Frobenius of an elliptic curve in classical normalization.

\begin{theorem}[Removing the weights]\label{aw:raw}
Fix $\ell\ge0$ and $\delta>0$. For either $A=\Theta_{\pi,\ell}$ or $A=B_{\pi,\ell}$, GRH for the standard function of $\pi$ is equivalent to the following condition. For every $\varepsilon>0$, a constant $C_\varepsilon$ and arbitrarily large $Y$ exist such that
\begin{equation}
 A(x)\ge-C_\varepsilon x^{1/2+\ell+\varepsilon}
     \qquad(Y\le x\le Y^{1+\delta}).
 \label{aw:rawwindow}
\end{equation}
It is also equivalent to the condition that for every $\eta>0$ there are arbitrarily large $Y$ with
\[
 \int_Y^{Y^{1+\delta}}\max\{-A(x),0\}
             \frac{dx}{x^{3/2+\ell}}\le Y^\eta .
\]
Under GRH, for all $x\ge2$ one has
\[
 |\Theta_{\pi,\ell}(x)|\ll_{\pi,\ell}x^{1/2+\ell}\log^2(2x),
 \qquad |B_{\pi,\ell}(x)|\ll_{\pi,\ell} x^{1/2+\ell}\log(2x).
\]
\end{theorem}

\begin{proof}
We first establish the displayed bounds under GRH. Retain the more precise contour shift in the proof of \eqref{aw:explicit}, keeping each individual trivial-zero term, the polynomial of degree at most one in $\log x$ at zero, and a remainder $O_\pi(x^h)$ with $-2<h<-1$. Let $m\ge2$ be an integer, $t=1/(2m)$ and $b=1+t$. No possible ideal norm lies strictly between $m$ and $bm=m+1/2$. Consequently $T_{\pi,b}(m)/t=\Theta_{\pi,0}(m)$. Apply the identity \eqref{aw:difference} to the more precise formula for $H_\pi$ at $m$ and $bm$, and divide by $t$. Its linear terms cancel. For each $\rho=1/2+i\gamma$ the multiplier satisfies
\[
 \frac{|b-b^\rho|}{t}\ll \min\{1+|\gamma|,1/t\}.
\]
For the first bound integrate the derivative of $v-v^\rho$ from $1$ to $b$ and use $1\le v\le2$; the second follows directly by the triangle inequality. Hence the sum of the absolute values of the normalized zero terms is at most a constant times
\[
 \sqrt m\sum_\rho
 \frac{\min\{1+|\gamma|,1/t\}}{1/4+\gamma^2}
       \ll_\pi \sqrt m\log^2(2m).
\]
Indeed the unit ordinates up to $1/t$ give a harmonic sum weighted by $\log(2+j)$, and the longer tail gives
$O_\pi(t^{-1}\sum_{j>1/t}\log(2+j)/j^2)=O_\pi(\log(2/t))$.
Each of the finitely many nonzero gamma terms at $\alpha$ contributes $O_\pi(m^{\operatorname{Re}\alpha})$ since $(b-b^\alpha)/t$ stays bounded for fixed $\alpha$. The polynomial at zero contributes $O_\pi(1+\log m)$. Bounding the two negative-line remainders separately costs only $O_\pi(m^{h+1})$. This proves the bound for $\Theta_{\pi,0}(m)$, and constancy between integer norms gives it for real $x$. Partial summation against $y^\ell$ then gives the bound for $\Theta_{\pi,\ell}$. Since $k\log q=\log(q^k)$, partial summation against $1/\log y$ on $y\ge2$ gives the stated single logarithm for $B_{\pi,\ell}$. These bounds give \eqref{aw:rawwindow}, which gives the integrated condition by taking $\varepsilon$ small, as in Theorem~\ref{aw:criterion}.

For the reverse implication form $g(u)=-e^{-\ell u}A(e^u)$ and its Laplace transform $P$. These functions have the crude exponential growth needed for \eqref{aw:coeff}, by \eqref{aw:lrs}. For $A=\Theta_{\pi,\ell}$ direct integration gives
\[
 P(s)=\frac{1}{2(s+\ell)}
 \left\{\frac{L'}{L}(s,\pi)+\frac{L'}{L}(s,\widetilde\pi)\right\}.
\]
In $\operatorname{Re}s>\vartheta_n$ the only singularities are simple poles at the completed zeros of the pair. At a real zero $r>\vartheta_n$ the residue equals $(m_\pi(r)+m_{\widetilde\pi}(r))/(2(r+\ell))>0$. The preceding Laplace argument with $\theta=1/2$ applies and yields GRH.

For $A=B_{\pi,\ell}$, start in $\operatorname{Re}s>2$ and choose the branches of the two logarithms specified by the Euler products. Their expansions in $q^{-s}$ give
\begin{equation}
 P(s)=-\frac{\log L(s,\pi)+\log L(s,\widetilde\pi)}
                    {2(s+\ell)} .
 \label{aw:logtransform}
\end{equation}
The integral moment estimate \eqref{aw:coeff} did not use meromorphic continuation and applies unchanged. The singularities in $\operatorname{Re}s>\vartheta_n$ are now logarithms. We verify the change in the nearest-singularity step explicitly.

Let $z$ be a center on the line $\operatorname{Re}z=\sigma$ with a unique nearest zero $\rho$ of the pair at distance $d<\sigma-1/2$. Take $d<h<\sigma-1/2$ with no other zero in the closed disk. By shrinking $h$ slightly if needed we may use an open neighborhood of that disk. Factoring out $(s-\rho)$ to the exact order from each function leaves analytic nonvanishing factors with analytic logarithms there. Up to a holomorphic remainder, the transform around $z$ thus equals $-m\log(s-\rho)/(s+\ell)$ where $m=(m_\pi(\rho)+m_{\widetilde\pi}(\rho))/2>0$. Branch constants add holomorphic terms. If $B_k(z)$ denotes the coefficient of $(z-s)^k$ in $P(s)$, then
\begin{equation}
 B_k(z)=\frac{m(z-\rho)^{-k}}{k(\rho+\ell)}
       +O\left(\frac{d^{-k}}{k^2}\right)+O(h^{-k}).
 \label{aw:logcoeff}
\end{equation}
In fact write $r=(z-\rho)/(z+\ell)$, whose modulus is strictly less than one. Multiply the convergent expansions of $1/(s+\ell)$ and $-\log(1-(z-s)/(z-\rho))$ about $z$. The coefficient from these two factors equals
\[
 \frac{(z-\rho)^{-k}}{z+\ell}
     \sum_{j=0}^{k-1}\frac{r^j}{k-j}.
\]
For $j\le k/2$, $|(k-j)^{-1}-k^{-1}|\le2j/k^2$. The complementary indices have geometrically small total weight. The sum is therefore $1/[k(1-r)]+O_r(k^{-2})$, which proves \eqref{aw:logcoeff} after multiplying by $m$. The omitted constant of the logarithm has a Taylor coefficient $O(|z+\ell|^{-k})$ and the analytic remainder contributes $O(h^{-k})$.

Use the same disks and two cases as in the pole-exclusion argument following \eqref{aw:coeff}. A unique nearest nonreal zero gives
$|B_k(z)|\gg d^{-k}/k$ for all large $k$. The disk about $\sigma$ of radius $h$ in the first case is zero-free and the chosen branch there is analytic, whence $|B_k(\sigma)|\ll h^{-k}$. In the second case the unique nearest zero is real, $\rho=\sigma-D>1/2$, and \eqref{aw:logcoeff} gives $B_k(\sigma)\gg D^{-k}/k$ with a positive leading term. The factor $1/k$ leaves both exponential contradictions unchanged. Every zero in the right half of the strip is excluded. The functional equation completes the proof.
\end{proof}

The deletions in Theorem~\ref{aw:prime} also apply to the raw window criteria. On norms at most $x$, the additional factor $q^{k\ell}$ costs at most $x^\ell$, and passing from $\Theta$ to $B$ multiplies a coefficient by $1/\log(q^k)\le1/\log2$. Thus rank at most four already permits a test using only real parts of ordinary unramified Satake traces.\par

The elementary identity
\[
 T_{\pi,b}(x)=b x\int_x^{b x}\Theta_{\pi,0}(u)\,u^{-2}du
\]
gives an independent check on the sign in the ordinary cutoff criterion. A first-power trace belonging to an ideal of norm at most $x$ appears in the integral with the positive coefficient $b-1$.

\subsection{Elliptic curves over a number field}

For an elliptic curve $E/F$, let $G_1(E/F)$ be the prime ideals of good reduction with residue degree one over $\Q$. For $\mathfrak p\in G_1(E/F)$ write $p=\mathrm N\mathfrak p$ and $k_{\mathfrak p}=\mathcal O_F/\mathfrak p$. Define
\begin{equation}
 D_{E/F}(x)=\sum_{\substack{\mathfrak p\in G_1(E/F)\\
                                     p\le x}}
       \bigl(p+1-\#\widetilde E_{\mathfrak p}(k_{\mathfrak p})\bigr).
 \label{aw:points}
\end{equation}
This is an integer, evaluated solely by counting points over prime fields, with one summand for each prime ideal present. Its sign depends on the curve. Elliptic prime races and their relation to zeros have been studied by Sarnak, Fiorilli and others. The condition below concerns a bound for the size of the deficit for one fixed curve on arbitrarily far power-length windows. \href{https://arxiv.org/abs/1304.8011}{Fiorilli, \emph{Elliptic curves of unbounded rank and Chebyshev's bias}.}

\begin{theorem}[An elliptic point-count criterion]\label{aw:elliptic}
Suppose the unitary Hasse--Weil function $L_E^u(s)=L(E/F,s+1/2)$ is the standard function of a unitary cuspidal representation of $\mathrm{GL}_2(\mathbb A_F)$, or is a product of two entire standard functions of unitary Hecke characters of $F$ with the usual completion. Fix $\delta>0$. All nontrivial zeros of $L(E/F,s)$ have real part $1$ if and only if, for every $\varepsilon>0$, there is a constant $C_{E,F,\varepsilon}$ and there are arbitrarily large $Y$ for which
\[
 D_{E/F}(x)\ge -C_{E,F,\varepsilon} x^{1+\varepsilon}
          \qquad(Y\le x\le Y^{1+\delta}).
\]
Equivalently, for every $\eta>0$ and arbitrarily large $Y$ the integral of $\max\{-D_{E/F}(x),0\}/x^2$ over this interval is at most $Y^\eta$. Under the asserted zero location one has $|D_{E/F}(x)|\ll_{E,F} x\log(2x)$ on the full half-line.
\end{theorem}
\begin{proof}
At a good prime ideal the local polynomial in classical normalization is
$1-a_{\mathfrak p}q^{-s}+q^{1-2s}$, with $a_{\mathfrak p}=q+1-\#\widetilde E_{\mathfrak p}(k_{\mathfrak p})$. Hasse's theorem puts its two reciprocal eigenvalues on the circle of radius $\sqrt q$. At multiplicative primes the single classical inverse root is $1$ or $-1$, and at additive primes the local polynomial is $1$. Thus for all primes the normalized inverse roots have modulus at most one, and for the good primes $a(\mathfrak p,1)q^{1/2}=a_{\mathfrak p}$. These descriptions hold over any number field. \href{https://people.math.harvard.edu/~gross/preprints/ell2.pdf}{B. Gross, \emph{The arithmetic of elliptic curves---An update}, \S1.}

Use Theorem~\ref{aw:raw} with $\ell=1/2$ and $A=B_{\pi,1/2}$. The finite sums of the curve and this automorphic function agree, with its true local factors at all primes. If it is the indicated Hecke product, the proof of Theorem~\ref{aw:raw} applies to the sum of the two functions' logarithms and their conjugates. Each zero contributes its nonnegative order and the argument detecting the product is unchanged. The entireness assumption excludes possible poles from both factors.

Deleting $k\ge2$ and the prime ideals of residue degree larger than one changes $B_{\pi,1/2}(x)$ by $O_F(x)$: in the proof of Theorem~\ref{aw:prime}(3) the total unweighted count with weights $\log q$ is $O_F(\sqrt x)$, the factor $1/(k\log q)$ is bounded, and $q^{k/2}\le\sqrt x$. The finitely many bad first powers add $O_{E,F}(1)$. The retained sum is exactly \eqref{aw:points}. This error is admissible for both window conditions, and the bound on the full half-line follows from the bound for $B_{\pi,1/2}$. Finally $s$ in $L_E^u$ corresponds to $s+1/2$ in the classical Hasse--Weil function. Thus the respective critical lines are $1/2$ and $1$ and the classical functional equation relates $s$ to $2-s$.
\end{proof}

This theorem applies, for example, to every elliptic curve over $\Q$, over a real quadratic field, and over a totally real cubic field by the modularity theorems of Breuil--Conrad--Diamond--Taylor, Freitas--Le Hung--Siksek, and Derickx--Najman--Siksek, respectively. When the CM endomorphisms are defined over the base, the two Hecke factors supply the second alternative by complex multiplication theory. Under the first alternative the linear centering in $H_\pi$ would use $(L'(E/F,3/2)/L(E/F,3/2))x$, whereas \eqref{aw:points} requires only the indicated finite-field counts. \href{https://arxiv.org/abs/1310.7088}{Real quadratic modularity.} \href{https://arxiv.org/abs/1901.03436}{Totally real cubic modularity.}

If $F/\Q$ is Galois, each rational prime unramified in $F$ that has one degree-one prime above it splits completely. The finitely many primes ramified in $F$ and those below bad primes of $E$ can be dropped at bounded cost in \eqref{aw:points}. Hence only the completely split rational primes are needed. Over a nongalois field one keeps each degree-one prime even if the same rational prime has other primes above it of higher degree. For instance a factorization with residue degrees $1,2$ contributes precisely its degree-one prime in a cubic field.

The estimate still required for an unconditional zero theorem is the lower bound for \eqref{aw:points} over its selected windows. Hasse bounds an individual deficit by $2\sqrt p$. Applying this separately to the degree-one primes gives $O_F(x^{3/2}/\log(2x))$ for the sum of absolute values, too large by a factor close to $\sqrt x$. Rankin--Selberg estimates on squared traces likewise give information on their magnitudes. The window condition asks for cancellation of the signed traces for this fixed curve. All the discarded local terms satisfy their bounds independently of this cancellation.

\subsection{Bounds attained in families of twists}

The missing fixed-curve estimate has a finite family version that can be proved directly. It is important here that the length of the interval need not be small compared with its lower endpoint. The distinction is which curve is tested.

\begin{theorem}[Finite windows for quadratic twists]\label{aw:twists}
Let $E/F$ be any elliptic curve, $\varepsilon>0$, and $2\le Y\le X<\infty$. At all good degree-one primes of odd residue characteristic and norm at most $X$, independently prescribe the unramified quadratic twist sign to be $+1$ or $-1$. Keep the local twist trivial at the primes of bad reduction of $E$ and at all primes above $2$. Every resulting pattern is realized by infinitely many quadratic twisting classes in $F^\times/(F^\times)^2$. For $Y$ sufficiently large in terms of $F$ and $\varepsilon$, all but a fraction
\[
 O_{F,\varepsilon}\!\left(\frac{Y^{-2\varepsilon}}{\log(2Y)}\right)
\]
of the uniformly counted patterns satisfy
$|D_{E^{(d)}/F}(x)|\le x^{1+\varepsilon}$ simultaneously for all $Y\le x\le X$, for every twist realizing the stated local data. This assertion uses no modularity theorem. When indexing the patterns, omit primes already prescribed to have trivial twist.
\end{theorem}
\begin{proof}
At each odd prime ideal in question choose a nonzero residue of the prescribed quadratic character. A unit with this residue determines an unramified or trivial local quadratic twist and multiplies the Frobenius trace by the prescribed sign. To keep the twist trivial at a place $v$ above $2$ or of bad reduction, prescribe the element to be congruent to $1$ modulo a sufficiently high power of $v$ so that it is a local square, by Hensel's lemma. These conditions concern coprime ideals and the Chinese remainder theorem supplies $d\in\mathcal O_F$ satisfying them all. To realize infinitely many distinct square classes, at each stage choose a new prime ideal outside this set and outside the support of the previously chosen elements, and prescribe valuation exactly one there using its square as a modulus. A representative with this valuation has a different square class from every earlier representative.

Uniformly over the sign patterns, the good odd contributions to \eqref{aw:points} are independent centered random variables $Z_{\mathfrak p}$, each of variance $a_{\mathfrak p}^2\le4p$. Every good degree-one prime of odd characteristic of norm at most $X$ was included. Primes originally bad stay bad under the prescribed local square, and the contribution at $2$ is bounded by a constant depending on $[F:\Q]$. For $V\ge2$ the total variance of the random variables with norm at most $\min(2V,X)$ is
\[
 \sum_{p\le2V}\sum_{\mathfrak p\mid p}\operatorname{Var}(Z_{\mathfrak p})
       \ll_F \sum_{p\le2V}p\ll_F V^2/\log(2V).
\]
Order the variables by norm, with any fixed order within each norm. Kolmogorov's maximal inequality bounds the fraction for which any partial sum up to $\min(2V,X)$ exceeds $V^{1+\varepsilon}/2$ in modulus by $O_F(V^{-2\varepsilon}/\log(2V))$. For completeness, stop at the first such crossing. Independence gives zero expected cross term between the stopped partial sum and its future increments on each first-crossing event. The second moment of the final sum is therefore at least the square threshold times the probability of crossing. This proves the inequality and also handles equal norms, where only a subset of the tested partial sums occurs in \eqref{aw:points}.

Use $V=2^jY$, summing over $j\ge0$ with $V\le X$. The geometric factor $2^{-2\varepsilon j}$ bounds the total exceptional fraction by the assertion. For large $Y$ the deterministic contribution over $2$ is at most $Y^{1+\varepsilon}/2$. Adding it completes the proof.
\end{proof}

For $F=\Q$ these proportions are also relative natural densities among positive prime twisting parameters lying in any fixed admissible progression that makes the twist a local square at $2$ and at the bad primes. To see this, combine the fixed progression with congruences at the finitely many remaining primes through $X$. On the units at each odd such prime, precisely half of the classes realize each sign. The prime number theorem in arithmetic progressions for this fixed modulus gives equal relative densities to the resulting sign patterns. The finitely many prime parameters at most $X$ have density zero. The modulus in this argument grows with $X$.

For each finite $X$ the theorem permits one to impose arbitrarily long continuous windows on many twisting classes. Taking $X$ to infinity changes the finite set of prescribed signs. The Chinese remainder argument constructs a new element for each finite system, and it gives no element of $F^\times$ realizing an assigned infinite system. In particular, its conclusion does not imply \eqref{aw:rawwindow} for the original curve or for a single fixed twist on infinitely many windows.

\begin{theorem}[Twists of bounded parameter over a fixed field]\label{aw:sieve}
Fix $E/F$, $\varepsilon>0$ and $0<c\le1$. Let $D\ge2$ and $D^c\le Y\le X\le D$. Among odd positive squarefree integers $d\le D$, all but a proportion $O_{E,F,c,\varepsilon}(Y^{-\varepsilon})$ satisfy
\begin{equation}
 |D_{E^{(d)}/F}(x)|\le x^{1+\varepsilon}
        \qquad\text{simultaneously for }Y\le x\le X.
 \label{aw:actualfamily}
\end{equation}
Here the twists of the curve over $F$ are taken by rational integers viewed as elements of $F$. For the shorter range $2\le Y\le X\le\sqrt D$, the exceptional proportion is $O_{E,F,\varepsilon}(Y^{-2\varepsilon}\log(2Y))$ without a lower power-size restriction on $Y$. In particular, for any fixed $\delta>0$, a proportion tending to one of odd squarefree parameters $d\le Y^{1+\delta}$ obeys \eqref{aw:actualfamily} on $[Y,Y^{1+\delta}]$.
\end{theorem}
\begin{proof}
Let $S$ contain $2$, the rational primes ramified in $F$, and the rational primes under the bad places of $E$. For each rational prime $p\notin S$, let
\[
 c_p=\sum_{\substack{\mathfrak p\mid p\\ f_{\mathfrak p}=1}}
         (p+1-\#\widetilde E_{\mathfrak p}(\F_p)),\qquad
                    |c_p|\le2[F:\Q]\sqrt p,
\]
and set $c_p=0$ on $S$. If $p\nmid d$ the trace at every one of these degree-one primes is multiplied by the same Legendre symbol $(d/p)$. If $p\mid d$, all these places have odd valuation of $d$ since $p$ is unramified in $F$ and $d$ is squarefree. The good unramified two-dimensional Tate representation acquires nontrivial inertia by the quadratic twist, so the twisted curve has bad reduction there. Thus those places give zero to the restricted sum in either description. At primes in $S$ Hasse bounds any contribution of the twisted curve independently of $d$. We have obtained the uniform identity with bounded error
\[
 D_{E^{(d)}/F}(u)=\sum_{p\le u}c_p(d/p)+O_{E,F}(1).
\]
Chebyshev and Hasse give $\sum_{p\le N}|c_p|^2\ll_F N^2/\log(2N)$.

The coefficient form of the quadratic large sieve in the proof of Theorem~\ref{ep:sieve} gives, for any coefficients $v_p$ supported on odd primes up to $N\le D$ and any $\kappa>0$,
\[
 \frac1D\sum_{\substack{d\le D\\d\text{ odd squarefree}}}
 \left|\sum_{p\le N}v_p(d/p)\right|^2
 \ll_\kappa (DN)^\kappa\sum_{p\le N}|v_p|^2 .
\]
Here $d$ is the positive odd integer parameter itself, so only the split modulo $4$ in quadratic reciprocity is needed. As in that proof the maximum of partial sums costs a factor $O(\log^2(2N))$.

The same bound with factor $1$ in place of $(DN)^\kappa$ holds if $N\le\sqrt D$. To prove it, extend the sum of squares to all integers $d\le D$ and expand. The diagonal terms are at most $D\sum_p |v_p|^2$. For distinct odd primes $p,r$, P\'olya--Vinogradov for the primitive quadratic character of conductor $pr$ bounds $|\sum_{d\le D}(d/p)(d/r)|\ll\sqrt{pr}\log(2N)$. The total off-diagonal contribution is therefore at most
\[
 O\left(\log(2N)\Big(\sum_{p\le N}|v_p|\sqrt p\Big)^2\right)
  \ll N^2\sum_{p\le N}|v_p|^2,
\]
by Cauchy--Schwarz and $\sum_{p\le N}p\ll N^2/\log(2N)$. Division by $D$ proves the claimed estimate.

The dyadic maximal argument in Theorem~\ref{ep:sieve} applies to this shorter-range bound as well, since it holds for every coefficient vector.

Use $V=2^jY$ for $j\ge0$ with $V\le X$, and apply the maximal bound with $N=\min(2V,X)$ and $v_p=c_p$. The relative number of $d$ giving a partial sum exceeding $V^{1+\varepsilon}/2$ in modulus is
\[
 \ll_{F,\kappa}D^{2\kappa}\frac{\log(2V)}{V^{2\varepsilon}}
\]
in the first range, or the same expression without $D^{2\kappa}$ in the shorter range. The normalization by $D$ is comparable to normalization by the number of odd squarefree $d\le D$, by their positive natural density. Summing the geometric progression with its logarithmic factor costs $O_\varepsilon(Y^{-2\varepsilon}\log(2Y))$. Choose $\kappa=c\varepsilon/4$ in the first range so $D^{2\kappa}\le Y^{\varepsilon/2}$, and absorb the remaining logarithm. For large $Y$ the fixed contribution of $S$ is at most $Y^{1+\varepsilon}/2$. Enlarging the constant handles bounded $Y$. The last statement follows on taking $X=D=Y^{1+\delta}$ and $c=1/(1+\delta)$.
\end{proof}

The parameter $d$ is bounded here as an integer. Its relation to the conductor of the twist depends on $F$ and on the local conductors. The exceptional-set estimates count parameters at each scale. They allow any fixed parameter to lie in every exceptional set. In particular they do not prove even a single fixed-curve instance of the infinite-window assertion in Theorem~\ref{aw:elliptic}.

\begin{corollary}[Quadratic characters of the number field]\label{aw:hec}
Fix $F$, a unitary cuspidal $\pi$ of rank $n\ge2$, and $\ell\ge0$. For a quadratic Fisher--Friedberg Hecke family $\chi_{\mathfrak a}$ over $F$, order its squarefree ideal parameters by $\mathrm N\mathfrak a\le D$. In the range $D^c\le Y\le X\le D$ of Theorem~\ref{aw:sieve}, a proportion $1-O_{\pi,F,c,\ell,\varepsilon}(Y^{-\varepsilon})$ of these characters satisfies the following statement on $Y\le x\le X$. The sum of $q^\ell$ times the real first Satake trace of $\pi\otimes\chi_{\mathfrak a}$ over its unramified primes of norm at most $x$ is $O_{\pi,F,\ell,\varepsilon}(x^{1/2+\ell+\varepsilon})$ in modulus. If $n\le4$ the same conclusion holds for the full raw sum $B_{\pi\otimes\chi_{\mathfrak a},\ell}(x)$. For every elliptic curve $E/F$, irrespective of modularity, the corresponding quadratic twists obey $|D_{E^{(\mathfrak a)}/F}(x)|\ll_{E,F,\varepsilon}x^{1+\varepsilon}$ throughout the same window for the same proportion of parameters, with the implied exceptional-set constant depending on $E$. Here the norm of the finite conductor of $\chi_{\mathfrak a}$ is comparable to $\mathrm N\mathfrak a$, with constants fixed by the family.
\end{corollary}
\begin{proof}
Take the quadratic family with fixed auxiliary integral ideal $\mathfrak c$ constructed by Fisher--Friedberg. Goldmakher--Louvel, \S2, give its finite conductor as $\mathfrak a\mathfrak c_{\mathfrak a}$ for squarefree $\mathfrak a$ prime to $\mathfrak c$, where $\mathfrak c_{\mathfrak a}\mid\mathfrak c$. In particular distinct such parameters give distinct primitive characters. Their Theorem~1.1 yields
\[
 \sum_{\mathrm N\mathfrak a\le D}^{*}
 \left|\sum_{\mathrm N\mathfrak b\le N}^{*}
                      v_{\mathfrak b}\chi_{\mathfrak a}(\mathfrak b)\right|^2
 \ll_{F,\mathfrak c,\kappa}(D+N)(DN)^\kappa
                           \sum_{\mathrm N\mathfrak b\le N}^{*}|v_{\mathfrak b}|^2 .
\]
The stars denote squarefree integral ideals prime to $\mathfrak c$, and evaluations are extended by zero at common factors. The paper's displayed character indices are transposed. Applying that inequality with the two bounds exchanged and taking the adjoint of its finite matrix gives exactly this orientation since the quadratic characters are real. The number of parameters is $\asymp_{F,\mathfrak c}D$ by the usual squarefree ideal count (equivalently apply Wiener--Ikehara to $\zeta_F(s)/\zeta_F(2s)$ with the finitely many local factors omitted). \href{https://arxiv.org/pdf/1112.1642}{Goldmakher--Louvel, \emph{A quadratic large sieve inequality over number fields}.}

Take $v_{\mathfrak p}=q^\ell\operatorname{Re}a(\mathfrak p,1)$ away from $\mathfrak c$ and the conductor of $\pi$, and set all remaining coefficients to zero. The fixed-representation Rankin--Selberg estimate from Theorem~\ref{aw:prime} gives
$\sum_{q\le N}|v_{\mathfrak p}|^2\ll_\pi N^{1+2\ell}/\log(2N)$. For this bound separate $q\le\sqrt N$ before removing the logarithm in the prime coefficient. Outside the fixed exceptional set the coefficient of the twist at an unramified prime is precisely $a(\mathfrak p,1)\chi_{\mathfrak a}(\mathfrak p)$. Where $\mathfrak p\mid\mathfrak a$ the originally unramified representation has trivial inertia before the twist and nontrivial scalar inertia after it, hence local Euler factor one. At the fixed exceptional primes the individual twisted roots are uniformly bounded by \eqref{aw:lrs}. The dyadic maximal estimate and the sum over scales in Theorem~\ref{aw:sieve} now apply to these coefficients with squared threshold $V^{1+2\ell+2\varepsilon}/4$. Ordering the ideals by norm breaks ties arbitrarily. There are $O_F(N)$ eligible degree-unrestricted prime ideals, so the logarithmic cost of the maximal estimate is the same.

For $n\le4$ the higher coefficients outside the fixed exceptional set have absolute values bounded term by term by those of the fixed $\pi$, since the twist either multiplies them by a sign or deletes them. Their total is $O_\pi(x^{1/2+\ell})$ by the proof of Theorem~\ref{aw:prime}, uniformly in the twisting character. All powers of the fixed exceptional primes cost $O_{\pi,F}(x^{\vartheta_n+\ell})$. This proves the second statement. For the elliptic statement take instead $v_{\mathfrak p}=q+1-\#\widetilde E_{\mathfrak p}(k_{\mathfrak p})$ at good degree-one primes away from $\mathfrak c$, zero elsewhere. Hasse gives squared coefficient mass $O_F(N^2/\log(2N))$ directly. Quadratic class field theory identifies each nontrivial family character with a quadratic twist of $E/F$. The same good, ramified and fixed exceptional place analysis proves the claimed estimate.
\end{proof}

\subsection{Character factors and a translated pole}

A unitary Hecke character with entire completed function satisfies the same theorems as in rank at least two. Its unramified parameters have modulus one. A unitary norm character has a different main term. To specify it, write its function as $L(s)=\zeta_F(s+i\tau)$ with $\tau\in\R$ and $s_0=1-i\tau$. Its pole at $s_0$ is simple. For $\tau\ne0$ the centered expression replacing the complex expression inside $H_\pi$ is
\[
 S_\pi(x)+x(L'/L)(1)-\frac{x^{s_0}}{s_0(s_0-1)}.
\]
The last term cancels the residue of $-(L'/L)(s)/[s(s-1)]$ at $s_0$. For $\tau=0$ the kernel pole and zeta pole coincide and the centered expression is
$S_\pi(x)-x\log x+(1+\gamma_F)x$, where $\gamma_F$ is the ratio of constant and polar coefficients of $\zeta_F$ at $1$. In the latter case \eqref{aw:difference} becomes $-T_{\pi,b}(x)+bx\log b$ after centering.

For the raw sums the corrections are especially simple. From $\Theta_{\pi,\ell}$ subtract $\operatorname{Re}\int_2^x y^{\ell-i\tau}dy$ and from $B_{\pi,\ell}$ subtract $\operatorname{Re}\int_2^x y^{\ell-i\tau}dy/\log y$, taking each correction to be zero for $x<2$. The equivalences of Theorem~\ref{aw:raw} then hold with these corrected sums. Indeed, the contribution to the negative-sum transform from adding the first integral back, before conjugate averaging, is $2^{-(s-s_0)}/[(s+\ell)(s-s_0)]$, which cancels the principal pole of $(L'/L)(s)/(s+\ell)$ at $s_0$. For the second integral the corresponding term is
\[
 \frac{1}{s+\ell}\int_{\log2}^\infty e^{-(s-s_0)v}\frac{dv}{v}.
\]
The derivative of the integral with respect to $s$ is $-2^{-(s-s_0)}/(s-s_0)$, so the integral differs from $-\log(s-s_0)$ by an entire function after continuation. It cancels the singular part of $-\log L(s)/(s+\ell)$ at the translated pole. The disk argument now detects only zeros. In the converse direction the classical prime ideal explicit formula under GRH gives $\sum_{q^k\le x}\log q=x+O_F(\sqrt x\log^2(2x))$. It also follows by the preceding contour argument with the displayed principal term at $1$. Partial summation against $y^{\ell-i\tau}$ and then $1/\log y$ gives the asserted two bounds after the indicated subtractions. The horizontal translation preserves the real parts of the nontrivial zeros.

\clearpage
\section{Prime sums on intervals of nearly linear length}\label{apw:section}
In the equivalent criterion of Section~\ref{aw:section} the required prime intervals extend from \(Y\) to \(Y^{1+\delta}\). We prove a one-sided estimate for the very same prime statistic on intervals of length \(Y^{1-\eta}\), for every fixed \(\eta>0\), subject to an explicit condition on a possible real zero. A single fixed unitary norm twist removes that condition for any cuspidal representation. For the positive eta products, including the conductor-\(11\) elliptic curve, no twist is needed.

Let \(n\ge2\), let \(\pi\) be a fixed unitary cuspidal representation of \(\mathrm{GL}_n(\mathbb A_F)\), and let
\(\Theta=\sup\{\Re\rho:\Lambda(\rho,\pi)=0\}\). Fix \(b>1\) and put \(P(x)=T_{\pi,b}(x)\) as in \eqref{aw:taper}; this is a real prime-power sum.

\begin{theorem}\label{apw:intervals}
If \(\Theta=1/2\), then \(|P(x)|\ll_\pi\sqrt x\) at every scale, with \(b\) fixed. Suppose \(\Theta>1/2\) and the real number \(\Theta\) itself is not a zero. For each \(0<\eta<1\) and each \(\varepsilon>0\) there are arbitrarily large \(Y\) for which
\[
          P(x)>x^{\Theta-\varepsilon}\qquad
                     (Y\le x\le Y+Y^{1-\eta}).
\]
There are also arbitrarily large \(Y\) on which \(P(x)<-x^{\Theta-\varepsilon}\) throughout the same interval. In particular, under either alternative there are arbitrarily far intervals of the displayed length on which the lower bound in Theorem~\ref{aw:criterion} holds.

If in the second alternative there is some zero \(\Theta+i\gamma\), there are constants \(c,\delta_0>0\) such that \(P(x)>c x^\Theta\) throughout intervals \([Y,e^{\delta_0}Y]\), with \(\log Y\) relatively dense in a half-line. The same is true for \(P(x)<-c x^\Theta\), with possibly smaller constants.
\end{theorem}
\begin{proof}
The first alternative is Theorem~\ref{aw:criterion}. For the others, write
\[
          \Psi(x)=\sum_{q^j\le x}\Re a(\mathfrak p,j)\log q .
\]
The finite-difference argument in Section~\ref{aw:section} gives, with no hypothesis on local temperedness,
\begin{equation}\label{apw:psi}
                  \Psi(x)\ll_\pi x^\Theta\log^2(2x).
\end{equation}
Indeed at the integer \(m=\lfloor x\rfloor\) the two cutoffs \(m,m+1/2\) separate no ideal norms. The exact smoothed explicit formula with \(b'=1+1/(2m)\) expresses the sharp sum using the multiplier \((b'-b'^\rho)/(b'-1)\). Its absolute value is \(O(\min(1+|\Im\rho|,m))\) for \(0\le\Re\rho\le1\). Replace \(|m^\rho|\) by \(m^\Theta\) in the zero sum there. The finitely many bounded zeros cost \(O_\pi(m^\Theta)\), with all denominators nonzero by boundary nonvanishing. For larger ordinates the bound is \(O_\pi(m^\Theta\log^2(2m))\) by the unit-interval zero count and the factor \(1/|\rho(1-\rho)|\). The archimedean contributions and the negative-line remainder are the same as in that argument: their exponents are strictly below \(1/2\). This proves \eqref{apw:psi}.

At points away from the cutoffs, differentiation of the taper gives
\[
        P'(x)=b\sum_{x<q^j<bx}\frac{\Re a(\mathfrak p,j)\log q}{q^j}
                   \ll_{\pi,b}x^{\Theta-1}\log^2(2x)
\]
by partial summation in \eqref{apw:psi}. The finite, piecewise smooth sum \(P\) is continuous and locally absolutely continuous, so this derivative bound integrates across its corners.

The Mellin transform of \(P\) has a pole at every nontrivial zero of \(\pi\) with real part \(>1/2\), including those approaching \(\Theta\). This follows from \eqref{aw:residue} and \eqref{aw:difference}; the factor \(b-b^\rho\) never vanishes in the open strip, and multiplicities of \(\pi\) and its dual add. The same transform is regular on the real axis near \(\Theta\) and to its right, by the present hypothesis and the removable pole at one. Real zeros in \([1/2,1]\) are finite by compactness and isolatedness. Choose \(a=\Theta-\xi>1/2\) above the real zeros, with \(0<\xi<\min(\eta/2,\varepsilon)\).

If \(P(x)\le C x^a\) eventually for some constant \(C\), the eventually nonnegative function \(Cx^a-P(x)\) has a Mellin transform with a real singularity at its abscissa of convergence, by Landau's nonnegativity theorem. Omitting an initial compact interval adds an entire function. The polynomial growth of the prime sum gives a finite upper bound for the abscissa, and the nonreal poles with real parts above \(a\) give a lower bound strictly greater than \(a\). Yet the meromorphic transform is regular at every real point strictly greater than \(a\). This contradiction applies as well to \(-P\). Thus each sign exceeds \(2Y^{\Theta-\xi}\) at arbitrarily large \(Y\). On the interval of the theorem its variation is \(O_{\pi,b}(Y^{\Theta-\eta}\log^2 Y)=o(Y^{\Theta-\xi})\). Because \(\xi<\varepsilon\), the asserted strict inequalities follow.

When the edge is attained, the proof of Theorem~\ref{aw:edge} gives \(e^{-\Theta u}P(e^u)=Q(u)+o(1)\) as \(u\to+\infty\), where \(Q\) is the absolutely convergent Fourier series of the edge zeros. Its constant coefficient is zero because no edge zero is real. At least one other coefficient survives by the nonzero multiplier and positive zero multiplicities. A nonzero continuous uniformly almost periodic function of mean zero takes both signs. Indeed a strictly positive value gives a fixed positively bounded interval; its relatively dense almost periods give a positive mean if the function is nonnegative everywhere. Apply the same argument to its negative. Uniform continuity and those almost periods now translate an interval on which \(Q\) is bounded away from zero with either sign to relatively dense locations, with a smaller common margin. The \(o(1)\) preserves that margin far out on the half-line.
\end{proof}

\subsection*{The original unsmoothed elliptic trace}
For a rational elliptic curve the one-sided condition in the introduction involves a sharp cutoff and ordinary integral traces. The following result uses precisely that observable.

\begin{theorem}\label{apw:raw}
Put \(A_E(x)=\sum_{p\le x,\ p\nmid N} a_p\), and let \(\Theta\) be the supremum of real parts in the \emph{unitary} zero set. Suppose \(\Theta>1/2\) and it is not a real zero. For each \(\eta\in(0,1)\), \(\varepsilon>0\) there are arbitrarily large \(Y\) such that
\[
      A_E(x)>x^{1/2+\Theta-\varepsilon}
            \qquad (Y\le x\le Y+Y^{1-\eta}).
\]
There are arbitrarily large intervals of the same form throughout which \(A_E(x)<-x^{1/2+\Theta-\varepsilon}\). If the edge is attained, there are fixed \(c,h>0\) such that \(A_E(x)>c x^{1/2+\Theta}/\log x\) throughout intervals \([Y,e^hY]\) with \(\log Y\) relatively dense, and likewise for the negative sign. Consequently every elliptic curve in the complete critical-line census, unconditionally, has arbitrarily distant intervals \([Y,Y+Y^{1-\eta}]\) on which \(A_E(x)\ge-C_{E,\varepsilon}x^{1+\varepsilon}\).
\end{theorem}
\begin{proof}
When \(\Theta>1/2\), choose \(1/2<\sigma<\Theta\). The fixed-form density estimate of Section~\ref{common:section} gives \(N_\pi(\sigma,V)\ll_\pi V^{1-\delta}\log(2V)\) with \(\delta=(\sigma-1/2)/8>0\), for both signs of the ordinate. In particular,
\begin{equation}\label{apw:highzero}
   \sum_{\beta\ge\sigma,\ |\gamma|>K}\frac1{|\gamma|}
                \ll_{\pi,\sigma} K^{-\delta}\log(2K)\quad(K\ge2)
\end{equation}
by partial summation. We now show that the sharp prime-power statistic has an absolutely convergent leading sum,
\begin{equation}\label{apw:sharp}
 \Psi(m)=-\Re\sum_{\beta\ge\sigma}\frac{m^\rho}{\rho}
    +O_{\pi,\sigma}\bigl(m^\sigma\log^2(2m)+m^{\Theta-\delta}\log(2m)\bigr)
                      \qquad(m\in\Z,\ m\ge2).
\end{equation}
In the integer-cutoff calculation of \eqref{apw:psi}, put \(t=1/(2m)\). For each zero, counted with multiplicity, its coefficient after extracting \(-m^\rho\) is
\[
   \frac{1+t-(1+t)^\rho}{t\rho(1-\rho)}
       =\frac1\rho+O(t).
\]
Taylor's integral remainder proves this uniformly in \(\rho\), since the second derivative of \((1+v)^\rho\) has modulus at most \(|\rho(1-\rho)|\) on \(0\le v\le t\). For \(\beta\ge\sigma,\ |\gamma|\le m\), the total error from this replacement is \(O_\pi(m^{\Theta-\delta}\log(2m))\) by the density estimate. At larger ordinates, both the exact multiplier and \(1/\rho\) are \(O(1/|\gamma|)\), so \eqref{apw:highzero} gives the same bound. For all zeros with \(\beta<\sigma\) retain the original multiplier \(O_\pi(\min(1+|\gamma|,m)/(1+\gamma^2))\). As in \eqref{apw:psi}, their sum and the remaining explicit-formula terms cost \(O_\pi(m^\sigma\log^2(2m))\). This proves \eqref{apw:sharp}.

Let \(X\) be large, \(0\le v\le X^{1-\eta}\), and \(K=X^\eta\). In differences of the leading sum up to ordinate \(K\), integrate \(\rho y^{\rho-1}\) between the two integer cutoffs. After the division by \(\rho\), the bound per zero is \(O(X^\Theta/K)\), including the possible extra unit from the floor function. There are \(O_\pi(K^{1-\delta}\log(2K))\) such zeros. For the remainder use \eqref{apw:highzero} without taking a difference. It follows uniformly in \(v\) that
\[
 |\Psi(X+v)-\Psi(X)|\ll_{\pi,\sigma}
     X^\sigma\log^2(2X)+X^{\Theta-\eta\delta}\log(2X).
\]
Deleting powers \(p^j\) with \(j\ge2\) and the bad primes adds \(O_E(\sqrt X\log^2(2X))\), even when estimated on the full interval up to \(2X\), by the normalized Hasse bound. The remaining weight is \((a_p/\sqrt p)\log p\). Partial summation on the short interval against \(\sqrt y/\log y\) now bounds \(|A_E(X+v)-A_E(X)|\) by \(O_E(\sqrt X)\) times the preceding bounds. In particular it is \(O_{E,\eta}(X^{1/2+\Theta-\kappa})\) for some \(\kappa>0\); take any \(\kappa<\min(\Theta-\sigma,\eta\delta)\).

Finally the Mellin transform of \(A_E\), multiplied by \(s\), is \(\sum_{p\nmid N}a_p p^{-s}\) initially far to the right. It differs from \(\log L(E,s)\) by a function holomorphic on \(\Re s>1\): there the terms of order at least two in the logarithmic Euler product converge locally absolutely by Hasse, and the omitted bad factors have holomorphic logarithms. Hence at each zero \(\rho\) of the unitary completion with \(\Re\rho>1/2\), its derivative has a pole at \(s=\rho+1/2\). On the real axis above the rightmost real zero it has a holomorphic local branch, continued from the right. Apply the nonnegativity argument used in Theorem~\ref{apw:intervals}, now to the sharp sum and either sign. For arbitrarily small \(\xi>0\), both signs exceed \(2X^{1/2+\Theta-\xi}\) along sequences. Taking \(\xi<\min(\kappa,\varepsilon)\) and using the local variation estimate proves the strict assertions.

When the edge is attained, divide \eqref{apw:sharp} by \(m^\Theta\). Absolute convergence from \eqref{apw:highzero} allows dominated convergence on all zeros of smaller real part. In the retained edge series one can replace \(\lfloor x\rfloor\) by \(x\) after normalization: each individual exponential changes by \(o(1)\), and the common series of absolute majorants converges. The same limit holds after deleting the higher prime powers since \(\Theta>1/2\). Partial summation against \(\sqrt x/\log x\) gives
\[
 \frac{\log x}{x^{\Theta+1/2}}\,A_E(x)
    =-\Re\sum_{\beta=\Theta}\frac{e^{i\gamma\log x}}{\rho+1/2}+o(1).
\]
To check this calculation, apply partial summation first to finitely many terms \(-y^\rho/\rho\). Substitution \(y=xe^{-w}\) in the integral of \(y^{\rho-1/2}/\log y\), or integration by parts, gives the endpoint coefficient \(-1/\rho+1/(2\rho(\rho+1/2))=-1/(\rho+1/2)\). The integral with an extra factor \(1/\log y\) is of lower order. The errors and the omitted series tails can be passed to the limit because partial summation sends a bound \(\epsilon y^\Theta\), above any fixed starting point, to \(O_\Theta(\epsilon x^{\Theta+1/2}/\log x)\). This is again a nonzero mean-zero uniformly almost periodic function plus a vanishing error. The sign intervals follow as in Theorem~\ref{apw:intervals}.

For every curve in the census the only possible real zero is central. The desired one-sided estimate follows in the \(\Theta>1/2\) case just proved. In the other case ERH holds for this curve, and Theorem~\ref{aw:elliptic} or its forward explicit-formula proof supplies the bound at all scales. This proves the unconditional disjunction needed for the final conclusion.
\end{proof}

For any fixed \(\pi\), the forbidden norm twists \(\pi\otimes|\det|^{i\tau}\) for which an edge zero is real are at most countable. A twist translates zero ordinates, leaving \(\Theta\) and ERH unchanged. Thus for almost every *fixed* \(\tau\), the one-sided bound holds on the near-linear additive intervals above. At \(\tau=0\) itself the hypothesis holds for every form in Section~\ref{num:chapter}: its integral against the positive real Mellin kernel excludes all real zeros of its completed function. The bound controls exactly the taper of the equivalence, but its intervals are shorter. In the variable \(u=\log x\) their lengths are \(\log(1+Y^{-\eta})\), while those required there have length \(\delta\log Y\).

\begin{corollary}\label{apw:census}
For every rational elliptic isogeny class in the complete critical-line census of Section~\ref{cens:chapter}, with no norm twist, every \(\eta,\varepsilon>0\) with \(\eta<1\) admits arbitrarily large \(Y\) on which its normalized prime-power statistic satisfies
\[
        T_{\pi_E,b}(x)\ge-C_{E,b,\varepsilon}x^{1/2+\varepsilon}
                  \quad (Y\le x\le Y+Y^{1-\eta}).
\]
The conclusion concerns unbounded prime scales.
\end{corollary}
\begin{proof}
Every real zero has ordinate zero. The closed rectangle certified for the curve therefore already contains its entire real zero set in the critical strip. Each such zero, of whatever central multiplicity, lies on the line. If \(\Theta>1/2\) the real number \(\Theta\) cannot be a zero, and the positive intervals of Theorem~\ref{apw:intervals} apply. If \(\Theta=1/2\) the claimed bound holds for every \(x\ge2\).
\end{proof}

\subsection{Sharp Frobenius sums in abelian extensions}\label{ao:split}
Let $E/\Q$ be an elliptic curve, let $K/\Q$ be finite abelian with group $G$, and fix $g\in G$. Omit the rational primes of bad reduction and the primes ramified in $K$. With arithmetic Frobenius at the remaining primes, put
\[
 A_g(x)=\sum_{\substack{p\le x\\ \operatorname{Frob}_p=g}}a_p(E).
\]
For $\chi\in\widehat G$, take the primitive unitary elliptic twist with good Euler parameters $\chi(\operatorname{Frob}_p)\alpha_p,\chi(\operatorname{Frob}_p)\beta_p$, where $\alpha_p+\beta_p=a_p(E)/\sqrt p$. Write $m_\chi(\rho)$ for its multiplicity at a nontrivial zero, with zero multiplicity at regular nonzeros, and define
\[
 r_g(\rho)=\frac{1}{|G|}\sum_{\chi\in\widehat G}\overline{\chi(g)}m_\chi(\rho),\qquad
 \Theta_g=\sup\bigl(\{1/2\}\cup\{\Re\rho:r_g(\rho)\ne0\}\bigr).
\]
Conjugating the twists gives $r_g(\bar\rho)=\overline{r_g(\rho)}$. In the completely split class $r_1(\rho)$ is the multiplicity of the base-change zero divided by $[K:\Q]$, by Theorem~\ref{ebc:divisor}. Other classes can have cancellations.

\begin{corollary}\label{ao:class}
If $\Theta_g=1/2$, then $A_g(x)\ll_{E,K}x\log(2x)$. If $\Theta_g>1/2$ and $r_g(\Theta_g)=0$, the two sign conclusions of Theorem~\ref{apw:raw} hold for $A_g$ with $\Theta_g$ in the exponent. Its attained-edge conclusion holds when some $\Theta_g+i\gamma$ has $r_g(\Theta_g+i\gamma)\ne0$. In particular every invertible residue class modulo a fixed integer $m$ is covered on taking $K=\Q(\zeta_m)$.
\end{corollary}
\begin{proof}
Use the complex logarithmic derivatives of the individual twists. The proof of the tapered explicit formula applies before taking real parts, by the same contour shift. Summing with weights $\overline{\chi(g)}/|G|$ and using character orthogonality selects precisely the powers with $\operatorname{Frob}_p^j=g$ at good unramified primes. The finitely many other primes contribute $O_{E,K}(\log(2x))$ to the sharp logarithmically weighted sum, since all the unitary local roots have modulus at most one. Denote the resulting good-prime-power sum by $\Psi_g(x)$.

At the integer cutoff $m_0=\lfloor x\rfloor$, the exact multiplier in the proof of \eqref{apw:sharp} depends only on $\rho$ and $m_0$. Thus it multiplies $r_g(\rho)$, including its phase, after coincident zeros have been combined. The bound $|r_g(\rho)|\le |G|^{-1}\sum_\chi m_\chi(\rho)$ and the two-sided density estimate of Section~\ref{common:section} give \eqref{apw:highzero} for the weighted absolute tail at every $\sigma>1/2$. Those zeros with $r_g(\rho)\ne0$ have real part at most $\Theta_g$. Consequently \eqref{apw:sharp} and its short-increment bound hold with $\Theta_g$ and with the weighted leading sum $-\sum_{\Re\rho\ge\sigma}r_g(\rho)m_0^\rho/\rho$. In the case $\Theta_g=1/2$, simply retain the original multiplier at all zeros as in \eqref{apw:psi}. This gives $\Psi_g(x)\ll_{E,K}\sqrt x\log^2(2x)$. The powers with $j\ge2$ cost $O_{E}(\sqrt x\log^2(2x))$ under either hypothesis. Partial summation against $\sqrt x/\log x$ proves the asserted bound when $\Theta_g=1/2$, and transfers the short-increment estimate otherwise.

To check the sign argument, the derivative of $\sum_{\operatorname{Frob}_p=g}a_p(E)p^{-s}$ differs on $\Re s>1$ from
\[
 \frac1{|G|}\sum_\chi\overline{\chi(g)}
   \frac{L'}{L}(s-1/2,\pi_E\otimes\chi)
\]
by a holomorphic function. This follows by differentiating the logarithmic Euler products and summing the absolutely convergent higher powers. Thus each residue at $s=1/2+\rho$ equals $r_g(\rho)$. On the real axis to the right of the last effective real zero the continued derivative admits a local primitive agreeing with the prime series continued from the right. Landau's argument in Theorem~\ref{apw:raw} uses precisely these nonreal singularities and real regularity. At an attained edge the limiting Fourier series for $\log x\,A_g(x)/x^{\Theta_g+1/2}$ has coefficients $-r_g(\rho)/(\rho+1/2)$. It is real, has a nonzero coefficient, and has mean zero by hypothesis. The same almost-period argument completes the proof.
\end{proof}

One can select individual conjugacy classes in certain nonabelian extensions as well. Suppose $K/\Q$ is finite Galois, with group $G$, and that the Weil restriction of $E_K$ is $\Q$-isogenous to a product of rational elliptic curves $E_j$ with multiplicities. For a conjugacy class $C$ and $c\in C$, the prime sum is $A_C(x)=\sum_{p\le x,\,\operatorname{Frob}_p\in C}a_p(E)$, with the same finite omissions. Fix one embedding $\overline{\Q}\hookrightarrow\overline{\Q}_\ell$, realize every irreducible $G$-representation $\tau$ over that algebraically closed coefficient field, and let $V_E$ be the elliptic Tate representation there. The regular-representation identity gives
\[
      V_E\otimes\overline{\Q}_\ell[G]
                \simeq\bigoplus_j V_{E_j}^{\,m_j}.
\]
Each $V_{E_j}$ is absolutely irreducible over the rational Galois group, by semisimplicity and the endomorphism theorem of Faltings with $\operatorname{End}_{\Q}(E_j)=\Z$. In the CM case complex conjugation acts nontrivially on the imaginary quadratic endomorphism field, giving the same rational endomorphism ring. Thus the direct summand $V_E\otimes\tau$ has the form $\bigoplus_j V_{E_j}^{\,m_{\tau,j}}$ with nonnegative multiplicities. The corresponding unitarily normalized products of rational elliptic functions have precisely the twisted Frobenius parameters at every common good prime.

Character orthogonality gives
\[
  {\bf1}_{C}(h)=\frac{|C|}{|G|}\sum_\tau
                         \overline{\operatorname{tr}\tau(c)}\operatorname{tr}\tau(h).
\]
Adding the logarithmic derivatives using these weights selects $\operatorname{Frob}_p^k\in C$ at the $k$th prime power. In Corollary~\ref{ao:class} replace $r_g$ by
\[
        r_C(\rho)=\frac{|C|}{|G|}\sum_\tau
               \overline{\operatorname{tr}\tau(c)}
                    \sum_j m_{\tau,j}\operatorname{ord}_\rho\Lambda(s,\pi_{E_j}).
\]
The prime sum is real, so the residues of its differentiated prime series at conjugate points are conjugate. The finite products give the required power-saving density estimate for the absolute weighted zero count, and the proof of the corollary gives all its conclusions for $A_C$. A full rectangle certificate for each $E_j$ again excludes effective off-center real zeros for every $C$. This argument also applies if the displayed decomposition of the regular Tate representation has been established directly.

There is a parallel splitting assertion beyond abelian extensions. Let $K/\Q$ be any finite extension such that the completed elliptic function over $K$ has a fixed product decomposition into positive integral powers of unitary cuspidal $\mathrm{GL}_2/\Q$ functions, with the arithmetic argument shifted by $1/2$. Write $u_K(p)$ for the number of degree-one primes above an unramified $p$ in $K$. Then Theorem~\ref{apw:raw}, including its $\Theta=1/2$ bound from the proof, holds for
\[
                 A_{E,K}(x)=\sum_{p\le x} u_K(p)\,a_p(E)
\]
after omitting the finite exceptional set and using the edge of this product. Indeed the $\Psi$ formula for the product follows by adding the factors with their multiplicities, and the required density bound follows from the same finite sum. Over $K$, for each remaining term the unitary local parameters still have modulus one. For $q^j\le x$ with $q=\mathrm N\mathfrak p$, the contributions with $j\ge2$ or with relative residue degree $f(\mathfrak p/p)\ge2$ have total absolute logarithmic weight $O_K(\sqrt x\log^2(2x))$. For such terms $p\le\sqrt x$, there are at most $[K:\Q]$ places above $p$ and $O(\log(2x))$ exponents at each. The surviving logarithmic weight at an unramified $p$ is $u_K(p)a_p(E)\log p/\sqrt p$. This also proves convergence of the deleted logarithmic Euler terms to the right of unitary abscissa $1/2$. The sharp proof now applies. If $K/\Q$ is Galois, $u_K(p)=[K:\Q]$ for completely split primes and is zero for all others. Dividing by this fixed degree preserves the assertions, on choosing a slightly smaller $\varepsilon$ in the strict power bounds.

A full rectangle certificate at any positive height excludes every off-center \emph{real} zero at every scale, since real zeros have ordinate zero. For each of the $1791$ multiquadratic pairs in Section~\ref{cens:chapter}, the certificate therefore implies, unconditionally, that every $g$ has arbitrarily far intervals $[Y,Y+Y^{1-\eta}]$ on which $A_g(x)\ge-C_{E,K,\varepsilon}x^{1+\varepsilon}$. Each constituent of the positive product is certified, so every $r_g$ vanishes on the real interval $(1/2,1]$. For instance this applies separately to every invertible residue class modulo $24$ for the base change of $36\mathrm a$ over $\Q(\zeta_{24})$. This example illustrates why effective multiplicities enter. For the four invertible classes that are $2$ modulo $3$, the CM self-twist by the quadratic character of $\Q(\sqrt{-3})$ pairs the summands of $r_g$ with opposite signs, and $a_p(E)=0$ at every such good prime. For each of the $331$ nonabelian normal-closure pairs of Theorem~\ref{cmr:finite}, the proof of Theorem~\ref{cmr:tower} gives the regular Tate decomposition with two copies of every rational factor. Thus the unconditional one-sided conclusion holds in every individual conjugacy class, including the completely split primes. The same conclusion holds for the degree-one weighted sum in each of the $331$ nonnormal pairs. One of the normal fields is $\Q(189^{1/6},2^{1/3},\sqrt{-3})$ of degree $36$, with the curve $27\mathrm a1$. The large opposite-sign conclusions in each case remain subject to an edge strictly to the right of $1/2$.

Choosing a rational representative for each catalogued Galois field/isogeny pair gives $3976$ individual class cases from the multiquadratic list and $1696$ from the normal multiradical list. The first number is $2\cdot1604+4\cdot182+8\cdot5$. For the second, if $G=A\rtimes C_2$ with inversion action, there are $(|A|+|A[2]|)/2$ classes in $A$ and $|A/2A|=|A[2]|$ in the other coset. Multiplying these class totals by the eight pair counts in Section~\ref{cmr:chapter} gives $1696$. These counts distinguish indexed cases; in particular several of the prime sums can vanish identically.

\subsection*{Scope in higher absolute degree}
Let $\pi$ be fixed unitary cuspidal on $\mathrm{GL}_n(\mathbb A_F)$ with entire completion, absolute degree $d=n[F:\Q]\ge2$, and $\ell\ge0$. In unitary normalization form the sharp sum
\[
 A_{\pi,\ell}(x)=\sum_{\substack{\mathrm N\mathfrak p\le x,\ \mathfrak p\ {\rm unramified}\\ f(\mathfrak p/p)=1}}
       (\mathrm N\mathfrak p)^\ell\,\Re a(\mathfrak p,1).
\]
If $\Theta>1-1/d$ and $\Theta$ is not a real zero, then for every $\eta\in(0,1)$ and $\varepsilon>0$ there are arbitrarily far intervals $[Y,Y+Y^{1-\eta}]$ throughout which $A_{\pi,\ell}(x)>x^{\Theta+\ell-\varepsilon}$, and others throughout which $A_{\pi,\ell}(x)<-x^{\Theta+\ell-\varepsilon}$. If the edge is attained the respective bounds improve to $\pm c x^{\Theta+\ell}/\log x$ on intervals $[Y,e^hY]$ whose starts are relatively dense in $\log Y$. Here $c,h>0$ may depend on the fixed data. In the level-one $\mathrm{GL}_3/\Q$ case of Theorem~\ref{dw:windows}, $\Theta>5/8$ suffices.

For the proof take $1-1/d<\sigma<\Theta$ (or $5/8<\sigma<\Theta$ in the last example). Proposition~\ref{dw:density} supplies $N_\pi(\sigma,T)\ll_{\pi,\sigma}T^{1-\delta}$ for some $\delta>0$. This supplies the absolute high-zero tail and the short-increment bound of Theorem~\ref{apw:raw}. We verify passage to the specified prime sum even without local temperedness. At an unramified place, the coefficient of the negative logarithmic derivative of the Rankin--Selberg product at $\mathfrak p^k$ is exactly $|a(\mathfrak p,k)|^2\log q$, with $q=\mathrm N\mathfrak p$. The Rankin--Selberg prime number theorem used in the proof of Theorem~\ref{aw:prime}, applied with the finitely many ramified Euler factors removed, therefore gives
\[
 \sum_{\substack{q^k\le x\\\mathfrak p\ {\rm unramified}}}
                |a(\mathfrak p,k)|^2\log q\ll_\pi x.
\]
If $\mathcal D$ comprises $k\ge2$ and the first powers of residue degree at least two, then $\sum_{\mathcal D,\ q^k\le x}\log q\ll_F\sqrt x$. For $k=2$, use the prime-ideal Chebyshev bound through $\sqrt x$; the terms $k\ge3$ cost $O_F(x^{1/3}\log(2x))$. For the first powers of higher residue degree, the underlying rational primes satisfy $p\le\sqrt x$ and the total of their residue degrees above each $p$ is at most $[F:\Q]$. Cauchy--Schwarz on the measure with atom $\log q$ yields
\[
       \sum_{\mathcal D,\ q^k\le x,\ \mathfrak p\ {\rm unramified}}
                  |a(\mathfrak p,k)|\log q\ll_\pi x^{3/4}.
\]
The finitely many ramified places cost $O_\pi(\sqrt x\log(2x))$ by the local unitary bound. For $d\ge4$ one has $3/4<\sigma$. For $d=2,3$ the remaining cases are $\mathrm{GL}_2$ and $\mathrm{GL}_3$ over $\Q$, where Theorem~\ref{aw:prime} gives $O_{\pi,\xi}(x^{1/2+\xi})$ deletion, and unitary Hecke characters over fields of degree two or three, where the same estimate follows by counting with modulus-one coefficients. The assumption of an entire completion excludes polar norm characters.

These same absolute bounds give local uniform convergence of the deleted logarithmic Euler series and its derivative to the right of $\sigma$, by partial summation. For the retained real coefficients the two logarithmic derivatives of $\pi$ and its dual are averaged, so a nontrivial zero has positive combined multiplicity and cannot cancel. Applying the real-axis and nonreal-singularity argument after the shift by $\ell$ proves the sign peaks. Partial summation against $x^\ell/\log x$ preserves the strict power saving of the short increments. At an attained edge the averaging of the dual factors gives the real part of the same absolutely convergent Fourier series as before, with coefficient $-1/(\rho+\ell)$ per zero of $\pi$ before taking the real part. This proves the asserted extension.

\subsection{A nonattained zero supremum with rapid sign changes}\label{ao:section}
The sign conclusion for an attained supremum cannot be extended using the absolutely convergent zero expansion, its exact residues and the usual zero-counting bounds alone. The following model uses one zero set for every fixed taper parameter $b>1$.

\begin{proposition}\label{ao:counter}
Fix $1/2<\Theta<1$, $A>0$ and $B\in\R$. There is a real entire function $Z$ of order one with $Z(1-s)=Z(s)$. Its zeros belong to $0<\Re s<1$, none is real, and the supremum $\Theta$ of their real parts is never attained. Their one-sided count is
\[
        N^+(T)=AT\log T+BT+O(\log(2+T)),
\]
and every unit interval of ordinates contains $O(\log(2+|T|))$ zeros. Define, with multiplicities,
\begin{equation}\label{ao:series}
 \begin{gathered}
 H(x)=\sum_{Z(\rho)=0}\frac{x^\rho}{\rho(1-\rho)},\qquad
 P_b(x)=bH(x)-H(bx)=\sum_\rho c_b(\rho)x^\rho,\\
 c_b(\rho)=\frac{b-b^\rho}{\rho(1-\rho)} .
 \end{gathered}
\end{equation}
This uses the sign of $-T_{\pi,b}$ in \eqref{aw:difference}. For each fixed $b>1$, $P_b(x)=o_b(x^\Theta)$ and $P_b'(x)\ll_b x^{\Theta-1}$ for $x\ge e$; implied constants may also depend on $Z$. There exist $c,C>0$ depending on the fixed data such that, for all sufficiently large $X$, the interval
\[
       \left[X,\ X+X\exp\left\{-\frac{c\log X}{(\log\log X)^2}\right\}\right]
\]
contains points $x_+,x_-$ with $\pm P_b(x_\pm)\ge X^{\Theta-C/\log\log X}$. In particular arbitrarily late intervals of a fixed positive multiplicative width cannot have just one sign in this model.
\end{proposition}

\begin{proof}
Choose $a>0$ so small that $\beta_1>1/2$, and for $n\ge1$ set
\[
          \beta_n=\Theta-\frac{a}{\log(n+1)},\qquad
          \gamma_n=\pi\,3^n .
\]
Use all four zeros $\beta_n\pm i\gamma_n$, $1-\beta_n\pm i\gamma_n$. Add a central-line background as follows. For all sufficiently large integers $j$, let $t_j$ be the unique sufficiently large solution of $A t_j\log t_j+B t_j=j$, and include $1/2\pm i t_j$. The counts on the line follow by differencing this smooth increasing function, while the quadruples contribute $O(\log(2+T))$ to the one-sided count and $O(1)$ per unit interval. The quadruples never touch the central line and place zeros arbitrarily close to the specified supremum.

To obtain $Z$, put $w=s-1/2$ and $z_n=\beta_n-1/2+i\gamma_n$. Take the product of the factors $1+w^2/t_j^2$ and the pairs $(1-w^2/z_n^2)(1-w^2/\bar z_n^2)$. All products converge locally uniformly and have precisely the indicated zeros. They are even in $w$ and have real coefficients. The zero count and summation of $\log(1+R^2/|z|^2)$ by parts give $\log\max_{|w|\le R}|Z(1/2+w)|=O(R\log(2+R))$. Conversely the exponent of convergence of the background zero set is one, since $t_j\asymp j/\log j$. This proves the order assertion.

All series in \eqref{ao:series} converge absolutely and locally uniformly for $x>0$. Since $\sum_\rho|c_b(\rho)|<\infty$ and every $\Re\rho<\Theta$, dominated convergence gives $P_b(x)=o(x^\Theta)$. Its Mellin transform on $[1,\infty)$ initially converges absolutely for $\Re s>\Theta$ and equals $\sum_\rho c_b(\rho)/(s-\rho)$. This sum continues meromorphically to the plane, with precisely the displayed residues. The same assertion holds for $H$ with numerator one in its coefficients.

For $\rho=\sigma+i\gamma_n$ from either upper member of a quadruple one has, uniformly in $n$,
\[
  |c_b(\rho)|\le C_b9^{-n},\qquad
  \Re c_b(\rho)\ge c_b^*9^{-n}\quad(n\ge n_b),
\]
where $c_b^*>0$. Indeed $\Re(b-b^\rho)\ge b-b^\Theta>0$, $|b-b^\rho|\le 2b$, and
$1/(\rho(1-\rho))=\gamma_n^{-2}(1+O(1/\gamma_n))$. This works for every fixed $b$ on the same zero set.

Write $V=\log X$ and choose $m=\lfloor \kappa V/(\log V)^2\rfloor$ with fixed $\kappa>0$ sufficiently small. At each grid point $u=k/3^m$ in $[V,V+1]$, for every $n\ge m$ one has $e^{i\gamma_nu}=(-1)^k$ since $3^{n-m}$ is odd. After pairing conjugates, the entire tail therefore has sign $(-1)^k$ and its size is at least $2c_b^*9^{-m}e^{\beta_m u}$. The reflected real parts $1-\beta_n$ reinforce the same sign in this tail. In earlier quadruples their size is bounded by that of the corresponding right-hand exponent. The absolute contribution of all $j<m$, in units of $9^{-m}e^{\beta_m u}$, is at most
\[
 4C_b\sum_{j<m}
  \exp\left\{(m-j)\left(\log9-
                   \frac{a u}{(m+1)\log^2(m+1)}\right)\right\}.
\]
Here the derivative of $\Theta-a/\log(v+1)$ is positive and decreasing, so it gives the stated lower bound for $\beta_m-\beta_j$. By reducing $\kappa$ this geometric sum is smaller than $c_b^*/2$ uniformly for $u\ge V$ and $V$ large. The background contributes $O_b(e^{u/2})$, which is smaller still, because $m=o(V)$ and $\beta_m\to\Theta>1/2$. We have proved the asserted sign at every indicated grid point, with magnitude at least $c_b^*9^{-m}e^{\beta_m u}$. Since $\log(m+1)\sim\log V$, this exceeds $X^{\Theta-C/\log V}$ on increasing $C$ if needed. For $0<c<\kappa\log3$, the length in $u$ corresponding to the interval in the proposition is $\log(1+\exp(-cV/(\log V)^2))>2\cdot3^{-m}$ eventually. It contains two consecutive grid points with opposite signs.

We finish by checking that sharp changes of sign have not been put in by dispensing with the derivative estimate. For the quadruples the differentiated series in $u$ is normally convergent on compact sets, with total magnitude $O_b(e^{\Theta u}\sum_n 3^{-n})$. On the background $\rho(1-\rho)=1/4+t_j^2$. Apart from terms whose absolute sum is $O_b(e^{u/2})$, its $u$-derivative reduces to $e^{u/2}$ times linear combinations of
\[
                 \sum_j\frac{t_j e^{it_j v}}{1/4+t_j^2},
                  \qquad  v=u\ \text{or}\ v=u+\log b ,
\]
and their conjugates. For $v\ge1$ these sums are $O((1+v)^2)$. Up to $t_j\le e^{D v}$ this follows by taking absolute values and partial summation of the zero count. Choose $D$ large depending only on $A,B$. On a dyadic interval $R<t_j\le2R$ with $R\ge e^{Dv}$, the increments $v(t_{j+1}-t_j)$ are decreasing and between $c_Av/\log R$ and $\pi$. Every partial sum of $e^{it_jv}$ there is $O(\log R/v)$: write $e^{it_jv}=(e^{it_jv}-e^{it_{j+1}v})/(1-e^{iv(t_{j+1}-t_j)})$ and sum by parts. The reciprocal denominators have bounded variation $O(\log R/v)$ because their real parts are constant and their imaginary parts are monotone. A second partial summation with $t/(1/4+t^2)$ gives $O(\log R/(vR))$ for the weighted interval. Summing the dyadic tails proves the claim, locally uniformly for $v\ge1$, and justifies differentiation. Thus the background $u$-derivative is $O_b(e^{u/2}(1+u)^2)$, and conversion from $u$ to $x$ gives the asserted bound.
\end{proof}

The endpoint $\Theta=1$ also admits a counterexample with the same ordinates and $\beta_n=1-a/\log(n+1)$. Now $b-b^{\beta_n}\gg_b1/\log(n+1)$, so the real-part lower bound at layer $n$ becomes $c_b^*9^{-n}/\log(n+1)$, still valid after the $O_b(3^{-n})$ relative denominator correction. Take $m=\lfloor V/(\log V)^3\rfloor$. Then $au/((m+1)\log^2(m+1))\gg_a\log V$, and the earlier geometric sum is $o(1/\log(m+1))$. The same proof gives both signs throughout intervals of the stated form with denominator $(\log\log X)^3$ in place of $(\log\log X)^2$, and amplitudes at least $X^{1-C/\log\log X}$. This zero set also satisfies a classical zero-free region $\Re\rho\le1-c_0/\log(2+|\Im\rho|)$ for some $c_0>0$.

The proposition concerns an entire function with the prescribed symmetry and tapered zero sum. Its construction supplies no Euler product or automorphic representation. A one-sign lower bound throughout an interval of fixed logarithmic width for an actual automorphic prime sum in the nonattained case would require input that distinguishes its zero divisor from this model.

\clearpage
\section{A fixed multiplicative interval detects an off-line zero}\label{compact:section}
There is a shorter interval on which a one-sided prime estimate would settle elliptic ERH. Let $H(Y)$ tend to infinity as slowly as desired. For a fixed rational elliptic curve it is enough to find the desired estimate on arbitrarily distant intervals $[Y,YH(Y)]$. We prove this first for the automorphic taper and then for the ordinary elliptic point-count sum. The proof uses positivity to make the converse local in $\log x$.

In the classical setting, Kaczorowski and Pintz deduce oscillations of size $x^{\Theta-\varepsilon}$ in every sufficiently large power interval and a lower bound for the total number of changes of sign from the singularities of a Mellin transform. Ingham's fixed-ratio conclusion for the prime number theorem when a rightmost zero exists is discussed in Montgomery and Vaughan. We use the fixed-form zero-density bound to construct a positive average with bounded support, including when there is no rightmost zero. See \href{https://personal.math.ubc.ca/~gerg/teaching/592-Fall2018/papers/1987.Kaczorowski_2.pdf}{Kaczorowski--Pintz, \emph{Oscillatory properties of arithmetical functions. II}, Theorem 1}, and \href{https://personal.science.psu.edu/rcv4/personal/Publications/MNTI/19.0_pp_463_485_Oscillations_of_error_terms.pdf}{Montgomery--Vaughan, \emph{Multiplicative Number Theory I}, Theorem 15.8}.

\begin{lemma}[A positive measure that removes prescribed frequencies]\label{compact:filter}
Let $Z$ be a locally finite subset of $0<\Re s<1$ with $\Im s>0$ and $\sum_{\rho\in Z}1/\Im\rho<\infty$. There is a probability measure $\mu$ supported in $[0,C]$, where $C=\pi\sum_{\rho\in Z}1/\Im\rho$, such that $M(s)=\int e^{sv}\,d\mu(v)$ vanishes on $Z$ and its conjugate. It is strictly positive on the real axis. Given a point $s_*=\beta_*+i\gamma_*$ with $\gamma_*>0$, $M(s_*)\ne0$ provided no member of $Z$ has real part $\beta_*$ and ordinate $\gamma_*/(2k+1)$ for an integer $k\ge0$.
\end{lemma}
\begin{proof}
For $\rho=\beta+i\gamma$ set $h=\pi/\gamma$ and let $X_\rho$ be a random variable taking $0,h$ with probabilities $e^{\beta h}/(1+e^{\beta h}),1/(1+e^{\beta h})$, respectively. Take the variables independent. Their sum lies in $[0,C]$ for every outcome. Its distribution has transform
\[
 M(s)=\prod_{\rho\in Z}\frac{e^{\beta h}+e^{sh}}{1+e^{\beta h}} .
\]
The difference of a factor from one is bounded by a constant times $h$ on every compact $s$-set, since all $h\le C$. Hence the product converges locally uniformly, is equal to the transform by bounded convergence, and is nonzero wherever every individual factor is nonzero. The zeros of a factor have real part $\beta$ and imaginary part an odd multiple of $\gamma$. For $Z$ empty take $\mu=\delta_0$.
\end{proof}

\begin{theorem}\label{compact:taper}
Let $\pi$ be unitary cuspidal on $\mathrm{GL}_2/\Q$ and let $P=T_{\pi,b}$ be the real prime-power statistic in \eqref{aw:taper}, with $b>1$ fixed. Write $\Theta=\sup_\rho\Re\rho$ in the unitary normalization.
\begin{enumerate}
\item Suppose $\Theta>1/2$. For every $a<\Theta$ there is a finite $H_a>0$ such that for every sufficiently large $Y$ some $x\in[Y,e^{H_a}Y]$ has $P(x)<-x^a$. If the real number $\Theta$ is not itself a zero, there is also an $x$ in every such interval with $P(x)>x^a$.
\item If no zero attains $\Theta>1/2$, the conclusion for both signs holds for \emph{every} prescribed $H_a>0$. For each $a<\Theta$ the minimum number of alternating excursions across $\pm x^a$ inside a dyadic interval $[Y,2Y]$ tends to infinity as $Y\to\infty$.
\item For any fixed function $H(Y)\to\infty$, ERH for $\pi$ is equivalent to the following statement. For every $\varepsilon>0$ there is a constant $C_\varepsilon$ and arbitrarily large $Y$ such that
\[
            P(x)\ge-C_\varepsilon x^{1/2+\varepsilon}
                 \qquad(Y\le x\le YH(Y)).
\]
\end{enumerate}
\end{theorem}
\begin{proof}
Choose $1/2<\sigma<\Theta$. The fixed-form estimate of Beckwith, Liu, Thorner and Zaharescu in Section~\ref{common:section}, applied also to the contragredient, gives
\[
 \sum_{\substack{\beta\ge\sigma\\ |\gamma|\ge1}}\frac1{|\gamma|}<\infty .
\]
Indeed the count through $T$ is $O_{\pi,\sigma}(T^{1-\delta}\log(2T))$ for some $\delta>0$, so dyadic summation suffices. The finitely many nonreal zeros below height one cause no problem. Group identical points and conjugate pairs in \eqref{aw:explicit} and \eqref{aw:difference}. We can write
\[
 P(e^u)=-\sum_{\beta\ge\sigma} c_\rho e^{\rho u}+O_{\pi,b,\sigma}(e^{\sigma u}),
 \qquad
 c_\rho=\frac{m_\pi(\rho)+m_{\widetilde\pi}(\rho)}2
                  \frac{b-b^\rho}{\rho(1-\rho)} .
\]
The real-part convention makes $c_{\bar\rho}=\overline{c_\rho}$. Every indicated coefficient is nonzero, since $0<\beta<1$ and $b>1$. Real coefficients at real zeros are positive. The discarded zeros are bounded by absolute convergence of the coefficients. The series is absolutely convergent locally uniformly in $u$. All subsequent convolutions are on bounded intervals and may be performed term by term.

First assume the real number $\Theta$ is not a zero. The finitely many real zeros have real parts bounded away from $\Theta$. Take $\sigma>\max(1/2,a)$ above all of them. Choose some real part $\beta_*>\sigma$ occurring among the zeros. Among the zeros at that real part choose the least positive ordinate $\gamma_*$ in the combined divisor. In Lemma~\ref{compact:filter} take $Z$ to be all the other distinct points with $\beta\ge\sigma$, $\gamma>0$. No excluded point gives an accidental zero at $\rho_*=\beta_*+i\gamma_*$. We obtain
\[
 \int P(e^{u+v})\,d\mu(v)
   =-2\Re\{c_{\rho_*} M(\rho_*) e^{\rho_*u}\}+O(e^{\sigma u}).
\]
For every $U$ there are points in $[U,U+2\pi/\gamma_*]$ at which the cosine in this expression is $+1$ and $-1$. As $U\to\infty$, the error is smaller than half the magnitude of the corresponding main term. A positive average places a value of the same sign and at least that magnitude somewhere in its support. Taking $H_a=C+2\pi/\gamma_*$ and using $\beta_*>a$ proves the two assertions, absorbing fixed multiplicative constants by increasing the threshold for $Y$.

If $\Theta$ is itself a real zero, keep all real zeros with $\beta\ge\sigma>\max(a,1/2)$ and remove all nonreal pairs in this range. The leading term of the average is $-c_\Theta M(\Theta)e^{\Theta u}$ with a strictly negative coefficient, and all other real terms have the same sign. This gives the negative excursion on an interval of fixed multiplicative length even in this case.

If $\Theta$ is not attained, let $\sigma$ increase to $\Theta$. The sum of $1/\gamma$ over the remaining positive ordinates tends to zero by dominated convergence, and their least possible ordinate tends to infinity by local finiteness. Hence $C+2\pi/\gamma_*$ can be made smaller than any prescribed positive number. To force any fixed number of alternating excursions, subdivide $[\log Y,\log Y+\log2]$ into that many disjoint short intervals separated by gaps and use the two signs in each. This holds eventually for every fixed number of intervals.

Finally ERH gives $P(x)=O(\sqrt x)$ by Theorem~\ref{aw:criterion}. If it fails, choose $1/2<1/2+\varepsilon<a<\Theta$. The negative excursion in each fixed-ratio interval contradicts the stated lower bound once $H(Y)>e^{H_a}$ and $Y$ is sufficiently large.
\end{proof}

The proof uses only an absolutely convergent tapered explicit formula and the convergence of the reciprocal ordinates in the chosen right-hand strip. It therefore applies above any line $\Re s=\theta\ge1/2$ for which every $\sigma>\theta$ has that convergence. In particular the higher-degree density bounds of Section~\ref{dw:section} give an equivalence between $\Theta\le\theta$ and the preceding condition with $1/2$ replaced by $\theta$, whenever their right-half zero count is sublinear at every $\sigma>\theta$. This assertion concerns each fixed representation.

\clearpage
\section{A positive cutoff for ordinary Frobenius traces and products of point counts}\label{rawfilter:section}

Let \(E/\Q\) be an elliptic curve of conductor \(N\). At a good prime write
\(\lambda_p=(p+1-\#E(\F_p))/\sqrt p\). Put
\[
       B_\kappa(x)=\sum_{\substack{p\le x\\p\nmid N}}p^\kappa\lambda_p,
       \qquad -\tfrac12\le\kappa\le\tfrac12 .
\]
In particular \(B_{1/2}\) is the ordinary sum of the integers \(a_p\). Write
\(\Theta=\sup_\rho\Re\rho\) for zeros of the unitary completion. All assertions
below concern a fixed curve. Neither the interval ratio nor the starting point
in the oscillation assertion is uniform in the curve.

\begin{theorem}\label{rawfilter:main}
If \(\Theta>1/2\), then for each real \(a<\Theta+\kappa\) there is \(D>1\)
such that every sufficiently large interval \([X,DX]\) contains an \(x\) with
\[
                         B_\kappa(x)<-x^a .
\]
If the real number \(\Theta\) is not a zero, the same intervals contain a
point with \(B_\kappa(x)>x^a\). If no zero has real part \(\Theta\), both
conclusions hold for every prescribed \(D>1\).

In particular, let \(H(Y)\to\infty\) be any fixed function. ERH for \(E\) is
equivalent to the following assertion for the single choice \(\kappa=1/2\).
For each \(\epsilon>0\) there are arbitrarily large \(Y\) such that
\[
   \sum_{\substack{p\le x\\p\nmid N}}a_p
              \ge -C_{E,\epsilon}x^{1+\epsilon}
                \qquad(Y\le x\le YH(Y)).
\]
\end{theorem}

The logarithmic derivative supplies the prime weight \(\log p\).
Partial summation over a bounded ratio of cutoffs introduces a term from
primes below the first cutoff. We construct a positive average of ordinary
sharp sums whose explicit formula retains the desired zeros exactly.

\begin{lemma}\label{rawfilter:untaper}
Let \(\mu\) be a probability measure supported on \([0,c]\), with
\(M(z)=\int e^{zv}\,d\mu(v)\). For \(t>0\) set
\[
            k(t)=\int_0^\infty \frac{e^{-yt}}{M(y)}\,dy .
\]
If \(U>c+1\), there is a nonincreasing weight \(W_U(t)\), equal to one
for \(t<U\) and zero for \(t>U+c\), such that for \(t>c\)
\[
       \frac{W_U(t)}t
          =\int_{v\ge t-U} k(t-v)\,d\mu(v).
\]
If \(\mu\) has a smooth density, \(W_U\) is continuous on \(t>0\)
and is the tail probability of a measure on \([U,U+c]\).
Moreover,
\[
      \frac1{t+c}\le k(t)\le\frac1t,\qquad |k'(t)|\le t^{-2}.
\]
\end{lemma}
\begin{proof}
The estimates follow from \(1\le M(y)\le e^{cy}\), including differentiation
under the integral. At \(t>c\), Fubini and the change of variable \(z=ty\)
give
\[
 t\int_{v\ge t-U} k(t-v)\,d\mu(v)
    =\int_0^\infty e^{-z}\,
       \nu_{z/t}([t-U,\infty))\,dz,\qquad
 d\nu_y(v)=\frac{e^{yv}}{M(y)}d\mu(v).
\]
For every fixed upper interval \(J\), the derivative of \(\nu_y(J)\)
in \(y\) is \(\operatorname{Cov}_{\nu_y}(\mathbf1_J(v),v)\ge0\).
Indeed, twice this covariance is the mean of
\((\mathbf1_J(v)-\mathbf1_J(w))(v-w)\) over two independent samples.
As \(t\) increases the tilt parameter decreases and the upper interval
shrinks. Thus the integral decreases. It equals one below \(U\) and zero
above \(U+c\). Extend it by one on \(0<t\le c\). An atomless \(\mu\)
gives continuity by dominated convergence. The negative Stieltjes
derivative of \(W_U\) is then the required probability measure.
\end{proof}

\begin{proof}[Proof of Theorem~\ref{rawfilter:main}]
Choose a real \(b\) satisfying \(1/2<b<\Theta\) and \(b+\kappa>a\), so
\(b+\kappa>0\). If \(\Theta\) is not itself a real zero, increase \(b\)
beyond all real zeros. This is possible because the compact real portion
of the critical strip has only finitely many zeros. The
\href{https://doi.org/10.1017/S0305004122000445}{density estimate of Beckwith,
Liu, Thorner and Zaharescu, Theorem~1.1}, gives
\[
     \sum_{\substack{\rho=\beta+i\gamma\\\beta\ge b,\ \gamma>0}}
                         \frac1{\gamma}<\infty .
\]
Indeed there are finitely many terms at bounded height, and the count up to
\(T\) is \(O_{E,b}(T^{1-(b-1/2)/8}\log(2T))\). In the nonreal case choose
a real part \(\beta_*>b\) occurring among the zeros, and let \(\gamma_*>0\)
be the smallest positive ordinate at that real part. Write
\(z_*=\beta_*+\kappa+i\gamma_*\).

For each distinct other conjugate pair on \(\beta\ge b\), with \(\gamma>0\),
put \(h=\pi/\gamma\) and take the probability measure
\[
       \frac{e^{(\beta+\kappa)h}\delta_0+\delta_h}
            {1+e^{(\beta+\kappa)h}} .
\]
Convolve these measures, omitting the pair at \(z_*\). Since \(\sum h<\infty\),
the sum of independent variables with these laws converges and lies in
\([0,\sum h]\). Its moment generating function is the product of the
individual ones, by bounded convergence. Its value at \(\rho+\kappa\)
vanishes for each discarded pair, regardless of multiplicity. No
individual factor vanishes at \(z_*\): a zero would require
\(\beta=\beta_*\) and \(\gamma_*/\gamma\) an odd integer, whereas each
remaining positive ordinate there exceeds \(\gamma_*\).
The product at \(z_*\) is nonzero because the deviations from one are
summable, using \(|e^{z_*h}-1|\ll_{z_*}h\) for small \(h\).
Convolve once more with a nonnegative smooth probability density on
\([0,d]\), with \(d>0\) small enough that its transform at \(z_*\)
is nonzero. For example \(d\gamma_*<\pi/2\) suffices. Denote the
result by \(\mu\), supported on \([0,c]\). Its entire transform satisfies
\[
       M(\sigma+it)\ll_{E,b,\mu,J,(a_1,a_2)}(1+|t|)^{-J}
                       \qquad(a_1\le\sigma\le a_2)
\]
for every fixed bounded strip and every \(J\), by integration by parts
in the smooth factor. If \(\Theta\) is a real zero, instead include
\emph{all} the nonreal pairs on \(\beta\ge b\), with the same smoothing.
Then \(M(r)>0\) at every real \(r\).

Put \(L_\kappa(s)=L(E,s+1/2-\kappa)\) in arithmetic normalization. Its
nontrivial zeros are precisely the \(z=\rho+\kappa\). The \(k\)-th
prime-power coefficient of \(-L_\kappa'/L_\kappa\) is
\(p^{k\kappa}a(p,k)\log p\), where \(a(p,k)\) is the sum of the \(k\)-th
powers of the unitary inverse roots. Thus \(a(p,1)=\lambda_p\) at good
primes and \(|a(p,k)|\le2\) at every prime. To avoid any endpoint issue
choose a fixed \(A>c+1\) and \(U>A+c\), and use the smooth log-norm test
\[
           w_U=\mu*\bigl(k\,\mathbf1_{[A,U]}\bigr),\qquad
           \widehat w_U(s)=\int_{\R}w_U(t)e^{st}\,dt
                       =M(s)\int_A^U e^{st}k(t)\,dt .
\]
Above \(A+c\) the preceding lemma gives \(w_U(t)=W_U(t)/t\).
Below \(A+c\) its bounded contribution to any prime sum is \(O_{E,\mu,A}(1)\).
Consequently
\begin{equation}\label{rawfilter:average}
 \sum_{p\nmid N}p^\kappa\lambda_p W_U(\log p)
   =\sum_{p,j\ge1}p^{j\kappa}a(p,j)\log p\,w_U(j\log p)
       +O_{E,\mu,\epsilon}\big(e^{(\kappa+1/2+\epsilon)U}\big).
\end{equation}
Here and below a fixed \(A\) can be included in the constants. To see the
error bound set \(X=e^{U+c}\). The number of prime powers with exponent
\(j\ge2\) and size at most \(x\) is \(O(\sqrt{x}\log(2x))\), by summing
\(x^{1/j}\) over \(j\le\log_2 x\). In each dyadic interval for \(p^j\)
the absolute weight after division by \(j\log p\) is at most
\(2(p^j)^\kappa/j\). Summing these bounds gives
\(O_\epsilon(X^{\kappa+1/2+\epsilon})\) throughout the specified
\(\kappa\)-range, including \(\kappa=-1/2\). The first powers at bad
primes contribute only a fixed constant. The left side of
\eqref{rawfilter:average} is \(\int B_\kappa(e^t)(-dW_U(t))\), a
probability average of sharp prefixes with \(t\in[U,U+c]\).

For completeness, Mellin inversion on \(s=2+\kappa+it\) writes the
prime-power sum in \eqref{rawfilter:average} as the integral of
\((-L_\kappa'/L_\kappa)(s)\widehat w_U(s)\,ds/(2\pi i)\).
Shift to \(\Re s=\kappa-1/4\). There are no crossed trivial zeros,
since the classical gamma factor is \(\Gamma(s+1/2-\kappa)\).
On the new line the functional equation relates the log derivative to
one in the absolutely convergent unitary half-plane \(\Re s_{\rm unit}=5/4\)
and to gamma derivatives of size \(O_E(\log(2+|t|))\).
Along suitable horizontal segments the standard local partial fractions
and the unit-interval zero count give \(O_E(\log^2(2+|t|))\) for the log
derivative. They follow as in Section~\ref{ep:primes}, comparing the
Hadamard derivatives with a point in absolute convergence. At each fixed
\(U\), the rapid decay of \(M\) allows passage to the limit. The integral
on the new line is \(O_{E,\mu}((1+U)e^{\max(0,\kappa-1/4)U})\). The
remaining terms are \(-\sum_\rho\widehat w_U(\rho+\kappa)\), with
multiplicity. The contribution of all zeros with \(\Re\rho<b\) is
\(O_{E,\mu,b}(e^{(b+\kappa)U})\): use rapid decay of \(M\), the
unit-interval count, and
\(\int_A^U e^{(b+\kappa)t}k(t)\,dt\ll e^{(b+\kappa)U}\).
Take \(0<\epsilon<b-1/2\) in \eqref{rawfilter:average}.

In the nonreal case this leaves exactly one conjugate pair above \(b\).
For a fixed \(z\) with \(\Re z=r>0\), integration by parts, the derivative
bound on \(k\), and \(k(U)=U^{-1}+O_c(U^{-2})\) give
\[
        \int_A^U e^{zt}k(t)\,dt
              =\frac{e^{zU}}{zU}+O_{z,c,A}(e^{rU}/U^2).
\]
If \(m_*\) is the multiplicity, the average is thus
\[
   -\frac{2m_*e^{(\beta_*+\kappa)U}}{U}
       \Re\!\left(\frac{M(z_*)}{z_*}e^{i\gamma_* U}\right)
          +o\!\left(e^{(\beta_*+\kappa)U}/U\right).
\]
The amplitude is nonzero. For either sign and every \(X\) large enough
there is an appropriate phase \(U\in[\log X,\log X+2\pi/\gamma_*]\).
A sharp prefix in the averaging interval has the indicated sign and
magnitude at least a constant times \(X^{\beta_*+\kappa}/\log X\).
This implies the asserted bound with \(D=\exp(c+2\pi/\gamma_*)\).
In the real-edge case the retained real zeros give the leading term
\(-m\,M(\Theta+\kappa)e^{(\Theta+\kappa)U}/((\Theta+\kappa)U)\),
strictly negative, and all other retained real zeros are smaller.
The negative conclusion follows with \(D=e^c\).

When the edge is unattained, fix \(b_0\in(1/2,\Theta)\).
As \(b\) increases to \(\Theta\), dominated convergence in the
reciprocal-ordinate sum for \(b_0\) shows \(\sum_{\beta\ge b,\gamma>0}
\pi/\gamma\to0\). The smallest positive ordinate on \(\beta\ge b\)
tends to infinity by discreteness. Hence the support length, smoothing
length and period can all be taken with sum below \(\log D\) for any
prescribed \(D>1\).

Under ERH the ordinary prime estimate in Theorem~\ref{ep:equiv} and
partial summation give \(B_{1/2}(x)\ll_E x\log(2x)\). Conversely, if
\(\Theta>1/2\), take \(1+\epsilon<a<\Theta+1/2\). For large \(Y\),
\([Y,DY]\subseteq[Y,YH(Y)]\) and the negative excursion contradicts
the required lower bound. This proves the equivalence.
\end{proof}

\begin{corollary}\label{rawfilter:pointproduct}
For the same arbitrary \(H(Y)\to\infty\) put
\[
            P_E(x)=\prod_{\substack{p\le x\\p\nmid N}}\frac{\#E(\F_p)}p .
\]
ERH for the fixed curve is equivalent to the assertion that for every
\(\epsilon>0\) there are arbitrarily large \(Y\) such that
\[
          \log P_E(x)\le C_{E,\epsilon}x^\epsilon
                    \qquad(Y\le x\le YH(Y)).
\]
If ERH fails, for every \(a<\Theta-1/2\) there exists \(D>1\) such
that every sufficiently large \([X,DX]\) contains an \(x\) with
\(\log P_E(x)>x^a\). If the edge is unattained, every fixed \(D>1\)
has this property.
\end{corollary}
\begin{proof}
At every sufficiently large good prime, Hasse's bound and the Taylor
estimate \(\log(1+z)=z+O(z^2)\), applied to \(z=(1-a_p)/p\), give
\[
                \log P_E(x)=-B_{-1/2}(x)+O_E(\log\log(3x)).
\]
The error uses \(\sum_{p\le x}1/p\ll\log\log(3x)\); the finitely
many small point counts are positive. Under ERH partial summation of
\(B_{1/2}(x)\ll_E x\log(2x)\) against \(1/x\) gives
\(\log P_E(x)\ll_E\log^2(2x)\) in absolute value. When ERH fails,
apply the negative excursion theorem at an intermediate exponent
strictly between \(a\) and \(\Theta-1/2\), chosen positive, to absorb
the error.
\end{proof}

This criterion concerns an upper bound on a product of positive rational
numbers. Equivalently it bounds the product of the integral group orders
by \(\prod_{p\le x,\ p\nmid N}p\) times \(\exp(C_{E,\epsilon}x^\epsilon)\).
It can also be stated as a single count over a finite ring. If
\(m_x=\prod_{p\le x,\ p\nmid N}p\) and \(\mathcal E/\Z[1/N]\) is the
smooth elliptic model, the Chinese remainder theorem gives
\[
     P_E(x)=\frac{\#\mathcal E(\Z/m_x\Z)}{m_x}.
\]
Indeed the spectrum of this finite product of fields is the disjoint
union of their spectra, so its points with values in \(\mathcal E\)
are independent tuples of points in the good fibers.
The classical partial-product conjecture predicts the much finer
asymptotic \(P_E(x)\sim C(\log x)^r\). Goldfeld deduced ERH from that
asymptotic. Conrad treats the critical-line Euler-product problem in
greater generality, and Sheth obtains the predicted asymptotic outside
a set of finite logarithmic measure under ERH. Our estimate uses only an
upper bound on specified intervals and makes no assertion about \(r\).
\href{https://doi.org/10.4153/CJM-2005-012-6}{K.~Conrad,
\emph{Partial Euler products on the critical line}, Theorem~1.1
(quoting Goldfeld)};
\href{https://doi.org/10.1093/imrn/rnaf214}{A.~Sheth,
\emph{Euler Product Asymptotics for \(L\)-functions of Elliptic Curves}}.

\clearpage
\section{How many prime cutoffs detect an off-line zero?}\label{mass:section}

Keep the notation of Theorem~\ref{rawfilter:main}, with
$-1/2\le\kappa\le1/2$. Set $e=\Theta+\kappa$ and
$F_+(x)=\max(B_\kappa(x),0)$,
$F_-(x)=\max(-B_\kappa(x),0)$. Assume $\Theta>1/2$, so $e>0$.
The negative sign is always available. The positive sign is also
available if the real number $\Theta$ is not a zero.

\begin{theorem}\label{mass:block}
For each $a<e$ and $0<\eta<1$ there is $D>1$ with the following
properties in every sufficiently large $[Y,DY]$, for each available sign.
\begin{enumerate}
\item There is $r>\max(a,e-\eta)$ and a constant $c>0$, independent
of $Y$, such that
\[
 \int_Y^{DY}F_\pm(x)\frac{dx}{x}\ge c\frac{Y^r}{\log Y},
 \qquad
 \int_{\substack{Y\le x\le DY\\ F_\pm(x)\ge cY^r/\log Y}}
                \frac{dx}{x}\ge cY^{r-e}.
\]
\item There is a whole interval of ordinary length $Y^{1-\eta}$
inside $[Y,DY]$ on which $F_\pm(x)>x^a$.
At least $c_{E,a,\eta,\kappa}Y^{1-\eta}/\log Y$ good primes
$p\in[Y,DY]$ satisfy $F_\pm(p)>p^a$.
\end{enumerate}
If the edge is unattained, either assertion holds for every
prescribed $D>1$, with constants allowed to depend on $D$. If the edge is attained, one may take $r=e$ for every available sign.
\end{theorem}

\begin{proof}
The smooth probability measure of Lemma~\ref{rawfilter:untaper}
has density $f$ of finite supremum $K$ and support $[0,C]$.
The resulting probability on sharp cutoffs $t\in[U,U+C]$ has a
density bounded uniformly for $U>2C+1$. To see this, use the
tilted survival probability formula for $W_U$. Its negative derivative is
\[
 \int_0^\infty e^{-z}\left[
    \frac{e^{(z/t)(t-U)} f(t-U)}{M(z/t)}
     +\frac{z}{t^2}\operatorname{Cov}_{\nu_{z/t}}
            (\mathbf1_{[t-U,\infty)}(v),v)
                         \right] dz .
\]
The first term inside brackets is at most $K e^{zC/t}$ and the
second at most $zC/t^2$. Both are nonnegative and allow differentiation
under the integral. Thus $-W_U'(t)\le 2K+C/U^2$. Continuity
at the two endpoints rules out endpoint atoms.

Select a phase and a zero in the proof of Theorem~\ref{rawfilter:main}
with $r=\beta_*+\kappa>\max(a,e-\eta)$, or $r=e$ for the real edge.
The signed average at that phase is $\gg Y^r/\log Y$.
Dropping its negative part and using the density bound proves the
first integral inequality. We also have the global estimate
\[
                    |B_\kappa(x)|\ll_{E,\kappa}x^e/\log(2x).
\]
Indeed, \eqref{apw:sharp} and \eqref{apw:highzero} imply
$\Psi(x)\ll_E x^\Theta$ at integer cutoffs and hence at every cutoff.
Removing higher powers and bad primes changes the normalized
log-weighted statistic by $O_E(\sqrt x\log^2(2x))$. Partial summation
against $x^\kappa/\log x$ now gives the claim since $e>0$.
The portion of the first integral with $F_\pm<c_1Y^r/\log Y$ costs
at most $c_1(\log D)Y^r/\log Y$. Choosing $c_1$ small retains half
the lower bound. The global upper estimate on the remainder implies
the asserted logarithmic measure, after reducing the common $c$.

For the second claim it is important to place a whole interval
around the excursion. Fix $0<\eta'<\min(\eta,1/5)$.
The high- and low-ordinate difference estimate proved in
Theorem~\ref{apw:raw} gives some $d>0$ such that, for
$0\le h\le X^{1-\eta'}$,
\[
 \left|\sum_{X<p\le X+h,\ p\nmid N}\lambda_p\log p\right|
                 \ll_{E,\eta'} X^{\Theta-d}.
\]
Here $d$ may be decreased to absorb logarithms and the higher-power
bound. This estimate holds at every intermediate endpoint in the interval.
Partial summation there against $t^\kappa/\log t$ gives
$|B_\kappa(X+h)-B_\kappa(X)|\ll_{E,\eta',\kappa} X^{e-d}$,
with a smaller $d>0$ if needed. Apply Theorem~\ref{rawfilter:main}
at an exponent strictly between $\max(a,e-d)$ and $e$.
The witness $X$ lies in $[Y,D_0Y]$ and its signed value dominates both
the variation and $x^a$ throughout $[X,X+X^{1-\eta'}]$.
Choose $D>D_0$ to leave space for this interval. If $D$ was prescribed
in the unattained case, first take $1<D_0<D$.
A subinterval has length $Y^{1-\eta}$.
The longer interval has $\gg X^{1-\eta'}/\log X$ good primes
by the classical short-interval prime number theorem, since
$1-\eta'>4/5>7/12$. This proves the stated count.
See M.~N.~Huxley, \emph{The Distribution of Prime Numbers}, Oxford
(1972), or the statement of the classical theorem in the introduction to
\href{https://doi.org/10.1017/S1446788709000020}{D.~Bazzanella, \emph{Prime numbers in intervals starting at a fixed
power of the integers}}.
\end{proof}

\begin{corollary}\label{mass:criterion}
Let $A_E(x)=\sum_{p\le x,\ p\nmid N} a_p(E)$.
Fix any function $H(Y)\to\infty$ and any $0<\delta<1$.
ERH for the fixed curve is equivalent to each of the following.
\begin{enumerate}
\item For every $\epsilon>0$ there are arbitrarily large $Y$ for which
\[
  \int_Y^{YH(Y)}
      \frac{\max(-A_E(x),0)}{x^{1+\epsilon}}\,\frac{dx}{x}
                         \le 1 .
\]
\item For every $\epsilon>0$ there is a constant $C_{E,\epsilon}$
and arbitrarily large $Y$ such that all but at most
$C_{E,\epsilon}Y^{1-\delta}$ good primes $p\in[Y,YH(Y)]$ satisfy
$A_E(p)\ge-C_{E,\epsilon}p^{1+\epsilon}$.
\end{enumerate}
\end{corollary}
\begin{proof}
Under ERH, Theorem~\ref{ep:equiv} and partial summation give
$|A_E(x)|\ll_E x\log(2x)$. The integral is bounded by
$O_{E,\epsilon}(Y^{-\epsilon}\log(2Y))$ even upon extending its
upper limit to infinity. Every sufficiently large prime satisfies the
second bound. Conversely, choose $1+\epsilon<a<\Theta+1/2$ when
ERH fails. Take $\kappa=1/2$. In the first part of
Theorem~\ref{mass:block}, $r>a$, so the integral over just
$[Y,DY]$ is $\gg Y^{r-1-\epsilon}/\log Y$, tending to infinity.
For the second test take $0<\eta<\delta$. There are
$\gg Y^{1-\eta}/\log Y$ good prime indices at which
$A_E(p)<-p^a$. This exceeds the allowed exception count for
large $Y$, and $p^a>C_{E,\epsilon}p^{1+\epsilon}$. In both cases
$H(Y)>D$ eventually.
\end{proof}

If the edge is unattained, both signs of the ordinary sum cross their
prescribed thresholds arbitrarily many times inside every sufficiently
large dyadic block, by subdividing into any fixed number of disjoint
fixed-ratio intervals. At each passage between opposite signs there
is a prime cutoff $p$ with $|A_E(p)|\le 2\sqrt p$. Indeed, at the
first prime reaching the other closed half-line, the change is the
single Hasse trace $a_p$. Thus these blocks also contain arbitrarily
many distinct returns to this absolute range.

\begin{corollary}[Prime cutoffs selected by a second curve]\label{mass:sample}
Fix any elliptic curve $E/\Q$, any non-CM elliptic curve $E'/\Q$, and
a fixed interval $I\subset[-1,1]$ of positive length. The prime test of
Corollary~\ref{mass:criterion} remains equivalent to ERH for $E$ if
one tests only the good primes of both curves for which
\[
                  \frac{a_p(E')}{2\sqrt p}\in I .
\]
The exception allowance is still $C_{E,E',I,\epsilon}Y^{1-\delta}$.
When $E$ is non-CM one may take $E'=E$.
\end{corollary}
\begin{proof}
For the fixed interval choose a continuous nonnegative bump supported
inside $I$, at most one and of positive Sato--Tate integral.
Uniform polynomial approximation followed by subtraction of its error
gives a polynomial minorant of $\mathbf1_I$ with positive integral.
Its expansion in the $U_j$ basis therefore has positive constant
coefficient. The normalized modular form of $E'$ has trivial central
character and satisfies the temperedness hypothesis. Non-CM gives genuine
type, and Newton--Thorne supplies cuspidal automorphy for all its
positive symmetric powers. Hence \href{https://arxiv.org/abs/1505.03122}
{Lemke Oliver--Thorner, \emph{Effective log-free zero density estimates
for automorphic $L$-functions and the Sato--Tate conjecture},
Theorem~1.8} gives a $d_I>0$ for which
\[
     \sum_{X<p\le X+h,\ a_p(E')/(2\sqrt p)\in I}\log p
                          \gg_{E',I} h
                 \qquad (X^{1-d_I}\le h\le X)
\]
for sufficiently large $X$, with bad primes omitted.
In the interval proof of Theorem~\ref{mass:block} choose
$0<\eta'<\min(\eta,d_I,1/5)$. Its witness interval contains
$\gg_{E',I} X^{1-\eta'}/\log X$ selected prime cutoffs, all with the
same asserted excursion of the cumulative sum for $E$. Choosing
$\eta<\delta$ again exceeds the allowed exceptions. The forward
implication continues to hold at all good cutoffs.
\end{proof}

\clearpage
\section{Two-sided elliptic tests on shrinking relative intervals}\label{absolute:section}

For the ordinary point-count sum, an absolute-value estimate can be tested
on intervals even shorter than a fixed multiple of their starting point.
Keep $B_\kappa$ from Theorem~\ref{rawfilter:main}, with
$-1/2\le\kappa\le1/2$. For $\alpha>0$ set
\[
             h_\alpha(Y)=(\log Y)^{-\alpha},\qquad
             I_\alpha(Y)=[Y,Y\exp(h_\alpha(Y))].
\]
\begin{theorem}\label{absolute:excursion}
Suppose $\Theta>1/2$, and take
\[
           0<\alpha<\frac{\Theta-1/2}{2(1-\Theta)},
\]
with the upper endpoint understood as $+\infty$ when $\Theta=1$.
For every $a<\Theta+\kappa$ and all sufficiently large $Y$,
\[
           \int_{I_\alpha(Y)} |B_\kappa(x)|\frac{dx}{x}>Y^a.
\]
For every $0<\eta<1$ each such interval also contains a subinterval of
length $Y^{1-\eta}$ throughout which $|B_\kappa(x)|>x^a$.
These conclusions include an attained real or nonreal edge.
\end{theorem}

\begin{proof}
Write $u=\log Y$. Choose $1/2<b<\Theta$ so close to $\Theta$ that
$\alpha<\delta/(1-\delta)$, where
$\delta=(b-1/2)/(3/2-b)$. By Theorem~\ref{ang:density} the count
in this right strip through height $T$ is
$O_{E,b}(T^{1-\delta}(\log(2T))^C)$, so the reciprocal-ordinate
tail above $K$ is $O_{E,b}(K^{-\delta}(\log(2K))^C)$.
For the sharp expansion it is sufficient to keep the independently
established saving $\omega=(b-1/2)/8$ in \eqref{apw:sharp}.
Let $s_0=\max(b,\Theta-\omega)<\Theta$. For the sharp log-weighted sum
$Q_\kappa(t)=\sum_{\log p\le t,\ p\nmid N}p^\kappa\lambda_p\log p$,
\[
 Q_\kappa(t)=
      -\sum_{\Re\rho\ge b}\frac{e^{(\rho+\kappa)t}}{\rho+\kappa}
        +O_{E,b,\kappa}\bigl((1+t)^2e^{(s_0+\kappa)t}\bigr).
 \tag{*}
\]
The zeros are counted with multiplicity. To verify this for real cutoffs,
begin with \eqref{apw:sharp}. In its leading sum, moving from
$\lfloor x\rfloor$ to $x$ costs $O(x^{\Theta-\omega}\log(2x))$:
each term with $|\gamma|\le x$ costs $O(x^{\Theta-1})$, while the
rest is bounded by \eqref{apw:highzero}. Deleting higher powers and
bad prime terms costs $O_E(\sqrt x\log^2(2x))$. Partial summation
against $x^\kappa$ sends $-x^\rho/\rho$ to
$-x^{\rho+\kappa}/(\rho+\kappa)$ up to its fixed lower-end term.
All these integrations are justified by absolute convergence in the
chosen strip. Also $s_0+\kappa>0$, so the displayed error follows.

Select a single distinct zero with $z_*=\rho_*+\kappa$ and
$\Re\rho_*>\max(s_0,a-\kappa)$. Its ordinate and its multiplicity
$m_*$ are fixed throughout. We construct a possibly complex measure
$\nu_K$ of total variation at most one, with
\[
 \int d\nu_K=0,\qquad
 \int e^{zv}d\nu_K(v)=0
       \quad(z=\rho+\kappa\ne z_*,\ \Re\rho\ge b).
\]
For each distinct conjugate pair with $\Re\rho\ge b$,
$\gamma>K>2+2|\Im z_*|$, use the positive probability pulse
of Theorem~\ref{rawfilter:main}, of length $\pi/\gamma$ and real
exponent $\Re\rho+\kappa$. Their total support length is
$O(K^{-\delta}(\log(2K))^C)$ by the improved tail estimate.
At $z_*$ their joint transform differs from one by
$O_{z_*}(K^{-\delta}(\log(2K))^C)$ since it is an expectation of
$e^{z_*v}$ over this support.

For the remaining distinct $z\ne z_*$ with $|\Im z|\le K$, and also
for $z=0$, take the signed pulse
\[
        \frac{\delta_d-e^{zd}\delta_0}{1+|e^{zd}|},
                       \qquad d=\frac1{4K}.
\]
At the test exponent $z$ its transform vanishes. We include the
conjugate of $z_*$ when it is distinct. Each measure has variation
at most one. At the retained exponent, its transform has modulus
at least $c_*/K$ with $c_*>0$ independent of $K$. Indeed
$|(z-z_*)d|\le1/2$ for $K$ sufficiently large, and
$|e^w-1|\gg |w|$ there. Discreteness of the zeros in a bounded
neighborhood of $z_*$ gives a strictly positive minimum separation
from all distinct $z$ under consideration, including zero.
There are $O(K^{1-\delta}(\log(2K))^C)$ such pulses.
Finally convolve with the uniform probability on $[0,d]$.
Its transform at $z_*$ tends to one; the resulting density has absolute
value at most $1/d$ because the preceding convolution has variation
at most one. This proves
\[
 \begin{split}
  \operatorname{supp}\nu_K&\subset[0,H_K],\qquad
       H_K\ll K^{-\delta}(\log(2K))^C,\\
  \left|\int e^{z_*v}d\nu_K(v)\right|
       &\ge \exp\{-C_* K^{1-\delta}(\log(2K))^{C+1}\}.
 \end{split}
\]
Taking $K=u^v$ with $\alpha/\delta<v<1/(1-\delta)$ gives
$H_K=o(u^{-\alpha})$ and a retained transform at least $e^{-o(u)}$.
Integrate $(*)$ at $t=u+v'$ against $d\nu_K(v')$. Its error is bounded
using total variation at most one. Its single surviving term is
$-m_*e^{z_*u}\int e^{z_*v'}d\nu_K(v')/z_*$. As
$\Re z_*>s_0+\kappa$, this proves
\[
   \left|\int Q_\kappa(u+v')d\nu_K(v')\right|
                              \ge \exp\{(\Re z_*)u-o(u)\}.
\]

Put $A(t)=B_\kappa(e^t)$. Local Abel summation, including primes at
endpoints by the right-continuous convention, gives
\[
 Q_\kappa(u+v')=Q_\kappa(u)-uA(u)
        +(u+v')A(u+v')-\int_0^{v'}A(u+w)\,dw .
\]
The two quantities inherited from before the interval disappear because
$\nu_K$ has mass zero. Its bounded density and its variation bound then give
\[
  \left|\int Q_\kappa(u+v')d\nu_K(v')\right|
       \le [4K(u+H_K)+1]\int_0^{H_K}|A(u+w)|\,dw .
\]
The bracket is polynomial in $u$. Since $\Re z_*>a$, this proves the
integral assertion, already on the shorter log interval $[u,u+H_K]$.

For a subinterval use the short-increment saving $d_0>0$ in the proof of
Theorem~\ref{mass:block} at a parameter $0<\eta'<\eta$.
Apply the argument just proved with an exponent strictly between
$\max(a,\Theta+\kappa-d_0)$ and $\Theta+\kappa$.
It gives a peak on $[Y,Ye^{H_K}]$ that dominates the variation throughout
the following interval of length $Y^{1-\eta'}$. That interval lies
inside $I_\alpha(Y)$ eventually, since $H_K=o(u^{-\alpha})$.
This proves the second assertion.
\end{proof}

\begin{corollary}\label{absolute:criteria}
Fix any elliptic curve $E/\Q$ and use its ordinary raw sum $A_E$.
For any prescribed function $g(\epsilon)>0$ with
$g(\epsilon)\to0$ as $\epsilon\downarrow0$, ERH is equivalent to the
following condition: for each $\epsilon>0$ there are arbitrarily large
$Y$ with
\[
     \int_{I_{g(\epsilon)}(Y)}
                 \frac{|A_E(x)|}{x^{1+\epsilon}}\frac{dx}{x}\le1 .
\]
The same equivalence holds if the integral condition is replaced by
$|A_E(x)|\le C_{E,\epsilon}x^{1+\epsilon}$ throughout the interval.
One may instead use the same interval $[Y,Y\exp(h(Y))]$ for every
$\epsilon$, whenever $0<h(Y)\le1$ and
$\log(1/h(Y))=o(\log\log Y)$.
\end{corollary}
\begin{proof}
Under ERH, $|A_E(x)|\ll_E x\log(2x)$ on the entire half-line.
This proves both assertions, including the integral bound even if one
extends its upper limit to infinity. If ERH fails, choose $\epsilon>0$
small enough that $1+\epsilon<\Theta+1/2$ and $g(\epsilon)$ is below
the positive threshold of Theorem~\ref{absolute:excursion}. Choose
$1+\epsilon<a<\Theta+1/2$ and apply that theorem with $\kappa=1/2$.
All arguments in the interval are bounded by $eY$, so the normalized
absolute integral tends to infinity. Its subinterval assertion
contradicts the pointwise condition. For a common $h(Y)$ as specified,
the interval eventually contains $I_\alpha(Y)$ for every fixed
$\alpha>0$, and the same argument applies.
\end{proof}

For the pointwise criterion one may again test only the prime cutoffs
of Corollary~\ref{mass:sample}, allowing $O(Y^{1-\delta_0})$
exceptions for any fixed $0<\delta_0<1$. Choose
$\eta'<\min(\delta_0,d_I,1/5)$ in the subinterval proof; the published
short-interval estimate then supplies more selected prime cutoffs
than exceptions. The one-sided argument of Theorem~\ref{rawfilter:main}
relies on positivity of its filter.

At a fixed prescribed $\alpha$, the same implication has a numerical
strip interpretation. If the displayed test holds at $I_\alpha(Y)$
for every $\epsilon>0$, then
\[
                  \Theta\le\theta(\alpha):=\frac{1/2+2\alpha}{1+2\alpha}.
\]
Otherwise the strict inequality on $\alpha$ in
Theorem~\ref{absolute:excursion} holds and gives a contradiction.
The bound is strictly below one for every finite $\alpha$.

\paragraph{Additional information near real part one.}
An independent published fixed-degree log-free density theorem also gives a strip for every fixed $\alpha$, with an unspecified absolute constant. There is an absolute $A\ge2$ such that the same
two-sided test with any fixed $\alpha>0$ implies
\[
                  \Theta\le 1-\frac{1}{A(1+\alpha)}.
\]
In particular, if $\Theta=1$, the two conclusions of
Theorem~\ref{absolute:excursion} hold for \emph{every} fixed
$\alpha>0$. To prove this, use the fixed elliptic specialization of
\href{https://arxiv.org/abs/1505.03122}{Lemke Oliver--Thorner,
Corollary~1.2 or Theorem~1.5 with a trivial degree-one factor}:
for each fixed $1/2<b<1$, $N_E(b,T)\ll_{E,b} T^{A(1-b)}$.
Suppose $\Theta>1-1/[A(1+\alpha)]$ and choose $b<\Theta$ so close
to it that $\tau=A(1-b)<1/(1+\alpha)$ and $b>1/2$.
Retain the same sharp expansion $(*)$ using its original strictly
sub-edge error. The compressed pulse count is now $O(K^\tau)$ and
the tail support is $O(K^{\tau-1})$ by partial summation.
A choice $\alpha/(1-\tau)<v<1/\tau$ leaves enough support and
preserves the target exponential exactly as before. This proves
both conclusions for the specified $\alpha$. The good-window
condition contradicts the resulting excursion when
$1+\epsilon<a<\Theta+1/2$, giving the stated strip.

\paragraph{Attainment and the size of a return.}
An attained nonreal edge supplies signs, but crossing zero gives no
comparable size for a two-sided window. Even the single-pair profile
$G(e^u)=e^{eu}\cos(\gamma u)/u$, with $e>1$ and $\gamma\ne0$,
has
\[
  \inf_{U\le v\le U+1}\sup_{v\le t\le v+h}
                   |G(e^t)|\gg_{e,\gamma} e^{eU}h/U
       \qquad (0<h<\min(1,1/|\gamma|),\ U\ge2).
\]
Indeed on any interval of phases of length $|\gamma|h$ the largest
absolute cosine is at least $\sin(|\gamma|h/2)$; equality is attained
when the interval is centered at a cosine zero. Although this profile
has exact zeros at arbitrarily large arguments, relative log widths
$h=Y^{-o(1)}$, $Y=e^U$, keep its supremum at $Y^{e-o(1)}$.
This elementary model demonstrates the limitation of an inference
from the leading profile alone.

There is a corresponding fixed-form deduction if some strip
$\Re\rho\ge b>1/2$ contains only finitely many distinct zeros and
contains an edge zero. Let $z_*$ be an edge zero shifted by $\kappa$.
Use only the finitely many signed pulses in the proof above, one for
each distinct other shifted zero in this strip and one for zero;
write $n\ge1$ for their number. Take their delay $d=Y^{-c}/(n+1)$
and omit the final smoothing. Their variation is at most one, support
is shorter than $Y^{-c}$, and the retained transform is $\gg d^n$.
In $(*)$ the remainder has exponent $s_0+\kappa<\Theta+\kappa$.
The local Abel identity now bounds the averaged input by
$(2u+2)\sup_{0\le t\le Y^{-c}}|B_\kappa(e^{u+t})|$ for $u=\log Y$.
It follows that, whenever $0<nc<\Theta-s_0$, every sufficiently
large such window contains a point with
\[
         |B_\kappa(x)|\gg_{E,b,c,\kappa}
             \frac{Y^{\Theta+\kappa-nc}}{\log Y}.
\]
Thus, under this specific spectral gap, even power-length intervals
can be too wide to sustain a two-sided near-critical bound.

\subsection{An affirmative bound on long intervals around crossings}
\begin{proposition}\label{absolute:crossings}
Suppose $0<D\le1/2$, $\Theta\le1/2+D$, and either $\Theta=1/2$
or the real number $\Theta$ is not a zero.
For each fixed $0<\eta<1/4$ and every $\xi>0$ there is $C>1$ such
that every sufficiently large $[Y,CY]$ contains a subinterval of
length $Y^{1-\eta}$ throughout which the ordinary raw sum satisfies
\[
       |A_E(x)|\ll_{E,D,\eta,\xi}
               Y^{1+D-\eta D/(1-D)+\xi}.
\]
For $E:y^2=x^3-x$, the CM bound of Section~\ref{recentcm:section}
gives $D=3/8$, and the certified real-axis result excludes a noncentral
real zero. In particular every sufficiently large fixed-ratio block
with the indicated $C$ contains an interval of length $Y^{7/8}$ with
$|A_E(x)|\ll_{E,\xi}Y^{13/10+\xi}$. At this length the direct Hasse
variation gives exponent $11/8$.
\end{proposition}
\begin{proof}
We first bound increments at arbitrary integer endpoints near $X$.
For fixed $b_0>1/2$ the sharp prime-power formula can be written
at each integer $m\ge2$ as
\[
 \Psi(m)=-\sum_{\beta\ge b_0}m^\rho/\rho
                                  +O_{E,b_0}(m^{b_0}\log^2(2m)).
\]
Here the error can be brought down to $b_0$ even when the edge is far
to its right. In the proof of \eqref{apw:sharp} take the taper parameter
$t=1/(4m^L)$ with $L$ a sufficiently large fixed integer, so the two
cutoffs still separate no integer. Replace the zero coefficient by
$1/\rho$ through height $m^L$ as there, with error
$O(m^{\Theta-L\delta_0}(\log(2m))^C)$, where
$\delta_0=(b_0-1/2)/8$. Above that height use the absolute reciprocal
tail for the same error. The lower zeros cost $O_L(m^{b_0}\log^2(2m))$
by the unit-interval count and the exact taper coefficient.
For the negative-line integrals in the underlying twice-integrated
formula take a fixed line $h<-L-2$ clear of poles. Bounding the two
integrals separately and dividing by $t$ costs $O_{E,L}(m^{h+L})$.
The finitely many further trivial-zero residues have negative real
part, and each fixed difference quotient stays bounded after its
factor $m^\rho$ is removed. This proves the sharper formula.

Choose $b_0=1/2+v_0$ with $v_0>0$ small. In a fixed slice
$v\le\beta-1/2\le v+\Delta$, $v\ge v_0$, difference the leading
terms at integers $m,m'$ near $X$ with $0\le m'-m\le 2X^{1-\eta}$.
For each summand the cost is
$O(X^{1/2+v+\Delta}\min(X^{-\eta},1/(1+|\gamma|)))$.
Theorem~\ref{ang:density}, split at $X^\eta$ and summed by parts
above it, bounds the whole slice by
\[
      O_{E,v}\bigl(X^{1/2+v+\Delta-\eta v/(1-v)}
                                (\log X)^C\bigr).
\]
Use finitely many slices through $D$, with a last slice ending there.
The derivative of $v-\eta v/(1-v)$ is $1-\eta/(1-v)^2>0$ for
$0\le v\le D\le1/2$. Choose $v_0$ small compared with
$D-\eta D/(1-D)>0$ and $\Delta$ small compared with $\xi$.
Thus the exponent for the entire increment is at most
$1/2+D-\eta D/(1-D)+\xi/2$, after absorbing logarithms.
If the initial slice already lies past the edge, only the error at
$b_0$ is required. Deleting higher powers and bad primes costs
$O_E(\sqrt X\log^2(2X))$ at each endpoint. The same estimate holds
at real cutoffs by taking floors. Partial summation on the short
interval against $\sqrt x/\log x$ adds $1/2$ to the exponent.

Under ERH the claimed absolute bound follows on every interval.
Otherwise Theorem~\ref{rawfilter:main} supplies both signs in every
large $[Y,C_0Y]$. At a passage between opposite signs the ordinary
sum has magnitude at most $2\sqrt p$ at a prime cutoff $p$ by Hasse.
The increment just proved bounds its absolute value throughout the
following length $Y^{1-\eta}$. Taking $C>C_0$ completes the proof.
\end{proof}

\clearpage
\section{Reading the rightmost real part from ordinary prime counts}\label{cinv:section}
Let $E/\Q$ be fixed, write $A(x)=\sum_{p\le x,\ p\nmid N}a_p(E)$ and let $\Theta$ be the supremum in unitary coordinates. Two questions about the zero divisor can be read from the size and duration of the negative values of $A$. The first asks how far the divisor extends. The second asks if a zero realizes this supremum.

\begin{proposition}\label{cinv:exponent}
Fix any positive $H(Y)\to\infty$. Then
\[
\Theta+\tfrac12=\inf\left\{a\ge1:\ \begin{array}{l}\text{for arbitrarily large $Y$, } A(x)\ge -x^a\\[-2pt]\text{for every }Y\le x\le YH(Y)\end{array}\right\}.
\]
\end{proposition}
\begin{proof}
If $1\le a<\Theta+1/2$, choose $a<b<\Theta+1/2$ and apply Theorem~\ref{rawfilter:main} at exponent $b$. For every sufficiently large $Y$ its fixed-ratio negative excursion lies in the interval being tested. To prove the upper bound on the infimum, under ERH use $A(x)\ll_E x\log(2x)$. Otherwise the absolutely convergent right-half part in the sharp explicit formula of Section~\ref{apw:section} gives $A(x)\ll_{E,\epsilon}x^{\Theta+1/2+\epsilon}$ for every $\epsilon>0$. More explicitly, choose $1/2<\sigma<\Theta$ in \eqref{apw:sharp}: the coefficients $m_\rho/\rho$ in the retained sum are absolutely summable and its terms are bounded by $x^\Theta$. Partial summation against $\sqrt x/\log x$ and deletion of the higher prime powers give the stated bound. In either case every $a>\Theta+1/2$ belongs to the set.
\end{proof}

Suppose $\Theta>1/2$, fix $1\le a<\Theta+1/2$ and suppose the real point $\Theta$ is not a zero. Define the length of the longest unbroken interval satisfying the lower bound arbitrarily far out by
\[
 G_E(a)=\lim_{V\to\infty}\sup\bigl(\{0\}\cup\{h>0:\ \exists\, U\ge V,
  \ A(e^u)\ge -e^{au}\ \text{for all }U\le u\le U+h\}\bigr).
\]
The limit exists in $[0,\infty]$ by monotonicity. The upper bound below holds for every $V$ sufficiently large.

\begin{theorem}\label{cinv:attained}
Under these hypotheses $G_E(a)$ is finite. It is positive exactly when some zero has real part $\Theta$. If this happens, both signs have fixed-length intervals on which $\pm A(x)\gg_E x^{\Theta+1/2}/\log x$, and the starts of these intervals have bounded gaps in $\log x$.
\end{theorem}
\begin{proof}
Theorem~\ref{rawfilter:main} at any exponent between $a$ and $\Theta+1/2$ bounds $G_E(a)$ by the fixed logarithmic length for a negative excursion. If the supremum is unattained it bounds $G_E(a)$ by every positive real number. Conversely, if the edge is attained, the proof of Theorem~\ref{apw:raw} gives
\[
 ue^{-(\Theta+1/2)u}A(e^u)=F(u)+o(1),\qquad
 F(u)=-\sum_{\Re\rho=\Theta}\frac{m_\rho e^{i\Im(\rho)u}}{\rho+1/2}.
\]
The series is uniformly absolutely convergent by zero density. It defines a real uniformly almost periodic function with mean zero and at least one nonzero Fourier coefficient. It has positive and negative values: a continuous nonzero function of a single sign with relatively dense near-recurrences has nonzero mean, by integrating on disjoint intervals around the near-recurrences. Choose a small closed interval of strict positivity and one of strict negativity. Uniform almost periodicity gives their translates, with bounded gaps, on which the absolute value still exceeds a fixed constant of the corresponding sign. The displayed $o(1)$ is uniform for $u$ sufficiently large. This proves the last assertion and shows $G_E(a)>0$ using the positive intervals.
\end{proof}

\begin{proposition}[Effect of a finite height certificate]\label{cinv:height}
Let $1/2<\sigma<1$, $B\ge2$, and suppose no zero with $\Re\rho\ge\sigma$ has $|\Im\rho|\le B$. Suppose the number $N^+(\sigma,t)$ of such zeros with $0<\Im\rho\le t$, counted with multiplicity, satisfies $N^+(\sigma,t)\le K t^{1-\delta}\log(2t)$ for $t\ge B$, with $K,\delta>0$. Put
\[
 W=\frac{2\pi}{B}+\pi K B^{-\delta}
        \left(\frac{\log(2B)}\delta+\frac1{\delta^2}\right).
\]
If $\Theta>\sigma$, then every sufficiently far interval $[Y,e^hY]$, for each $h>W$, contains both signs of the excursions in Theorem~\ref{rawfilter:main} at every fixed $a<\Theta+\kappa$. The same holds for Theorem~\ref{compact:taper} at $a<\Theta$ for an elliptic form. For a general cuspidal $\mathrm{GL}_2/\Q$ taper take the combined conjugate divisor in the two hypotheses.
\end{proposition}
\begin{proof}
Choose the right-hand cutoff above $\sigma$ and close enough to $\Theta$ to accommodate $a$. No real zero is retained. All the positive ordinates retained are greater than $B$. Tonelli's theorem, counting multiplicities as an upper bound on distinct points, gives
\[
 \sum_{\substack{\beta\ge\sigma\\\gamma>B}}\frac1\gamma
 =\int_B^\infty\frac{N^+(\sigma,t)}{t^2}\,dt
 \le K B^{-\delta}\left(\frac{\log(2B)}\delta+\frac1{\delta^2}\right).
\]
The equality uses $N^+(\sigma,B)=0$. In the pulse construction the support costs at most $\pi$ times this sum, and one phase period costs less than $2\pi/B$. In the raw sum the terminal smooth probability can have arbitrarily small positive support. Thus $h>W$ accommodates both signs in the stated proofs. The threshold on $Y$ can still depend on the chosen zero and the value of the transform there.
\end{proof}

In particular, arbitrarily far intervals of this \emph{fixed} ratio on which $A(x)\ge -C_\epsilon x^{\sigma+1/2+\epsilon}$ for each $\epsilon>0$ would certify the full half-plane $\Re\rho\le\sigma$ under the two stated inputs. The density constant and exponent are needed to turn the ratio into a numerical value. The result supplies no upper bound for the first scale at which a hypothetical zero must become visible.

\clearpage
\section{One sparse set of primes for every rational elliptic curve}\label{sample:section}
What is the cost of evaluating the point-count criterion after choosing the primes once for all curves? There is a single computable set $\mathcal P$ with
\[
             \#\{p\le x:p\in\mathcal P\}\sim\frac{2\sqrt x}{\log x}.
\]
The group orders at these primes suffice for each of the rational elliptic curves. The selection uses Frobenius data of an increasing finite list of curves when it is constructed. After construction the same primes and the same integer exponents serve all curves. The argument gives a representation of the remaining cancellation problem and does not itself establish the cancellation.

\begin{theorem}\label{sample:elliptic}
Let $g(x)$ be any nondecreasing function tending to infinity. There is a set $\mathcal P$ of the size displayed above such that, for every $E/\Q$, simultaneously for all $x\ge3$,
\begin{align}
 \log\prod_{\substack{p\in\mathcal P,\ p\le x\\p\nmid N}}
       \left(\frac{\#E(\F_p)}p\right)^{\lfloor\sqrt p\rfloor}
   &= \log P_E(x)+O_E(1+g(x)),\label{sample:product}\\
 \sum_{\substack{p\in\mathcal P,\ p\le x\\p\nmid N}}a_p(E)
   &=\sum_{\substack{p\le x\\p\nmid N}}\frac{a_p(E)}{\sqrt p}
          +O_E\big(\sqrt x(1+g(x))\big).\label{sample:raw}
\end{align}
One may always replace $g(x)$ by $\min(g(x),\log\log(3x))$. In particular one fixed computable choice of $\mathcal P$ satisfies both estimates with $1+\log\log(3x)$. Fix any $H(Y)\to\infty$. Each of the following tests is equivalent to ERH for the individual $E$. For each $\epsilon>0$ there should be arbitrarily large $Y$ on which, throughout $[Y,YH(Y)]$, either the logarithm on the left of \eqref{sample:product} is $\le C_{E,\epsilon}x^\epsilon$, or, for the second test, the sum on the left of \eqref{sample:raw} is $\ge-C_{E,\epsilon}x^{1/2+\epsilon}$. Each test alone suffices. Each may be checked only at cutoffs $x$ that belong to $\mathcal P$.
\end{theorem}

\begin{lemma}[Rounding with control at every prefix]\label{sample:round}
Let $v_1,\ldots,v_k\in[-B,B]^d$ and $q_i\in[0,1]$. There are $e_i\in\{0,1\}$ with
\[
     \left\|\sum_{i\le h}(e_i-q_i)v_i\right\|_\infty\le dB
                              \quad(0\le h\le k).
\]
For rational data this can be done by rational arithmetic.
\end{lemma}
\begin{proof}
This is the floating-variable argument of finite-dimensional vector balancing. Process the positions in increasing order. Insert the current coordinate with value $q_i$. Keep coordinates already equal to zero or one fixed. If more than $d$ coordinates are fractional, take a nontrivial linear dependence of their vectors, and move their scalar values along that dependence until some coordinate reaches zero or one, keeping all of them in $[0,1]$. Repeat if needed. After processing $h$, the prefix vector sum equals its original fractional target at that time. Of the first $h$ coordinates at most $d$ can subsequently change, each by at most one. After the last step round the remaining fractional coordinates arbitrarily. This gives the bound simultaneously for all prefixes. Deterministic pivot and tie conventions give a finite algorithm. For the classical vector-partition setting see \href{https://www.renyi.hu/~barany/cikkek/103.pdf}{I.~B\'ar\'any and B.~Doerr, \emph{Balanced partitions of vector sequences}}, and the earlier B\'ar\'any--Grinberg theorem cited there.
\end{proof}

\begin{proof}[Proof of Theorem~\ref{sample:elliptic}]
Enumerate rational elliptic curves by their nonsingular integral Weierstrass equations; repetitions cause no change. Set $k_p=\lfloor\sqrt p\rfloor$, $q_p=1/k_p$. For a good prime use the two coordinates
\[
 \lambda_{E,p}=a_p(E)/\sqrt p,\qquad v_{E,p}=k_p\log(\#E(\F_p)/p),
\]
and at the finitely many bad primes set both to zero. The coordinates have an absolute bound. For the second one this follows for large $p$ from Hasse and $\log(1+h)=O(h)$ for $|h|\le1/2$; for the remaining primes only finitely many positive orders in the Hasse interval occur. Choose increasing integer bounds $X_j\to\infty$ such that $g(X_j)\ge j^2$ and, if necessary, replace $g$ by its indicated minimum. On each block $[X_j,X_{j+1})$, apply Lemma~\ref{sample:round} to vectors whose first coordinate is one, and whose other coordinates are the two just given for each of the first $j$ curves. Use dyadic rational approximations within $p^{-2}$ of each real coordinate. For each individual coordinate the cost of replacing all these approximants by their true values in any partial sum is bounded by $\sum_{n\ge2}n^{-2}$. Include the prime in $\mathcal P$ if its rounded scalar is one. The finitely many primes below $X_1$ can be assigned arbitrarily.

For each coordinate belonging to a fixed $E$, omission in the first finitely many blocks costs $O_E(1)$. Through any prefix of block $j$ the total rounding discrepancy is $O_E(1+\sum_{i\le j}(2i+1))=O_E(1+g(x))$. The count coordinate has the same bound. Consequently
\[
 \#(\mathcal P\cap[2,x])=\sum_{p\le x}q_p+O(1+g(x))
            \sim 2\sqrt x/\log x
\]
by the prime number theorem and partial summation, using $q_p=p^{-1/2}+O(1/p)$. For the computable choice $g(x)=\log\log(3x)$ one may take $X_j=2^{2^{2j^2+2}}$. Each block is finite, its coordinates can be computed, and the rational algorithm terminates. This effectiveness gives no practical size bound on the initial blocks needed for a given curve.

The target in the $v_{E,p}$ coordinate is exactly $q_pv_{E,p}=\log(\#E(\F_p)/p)$. This proves \eqref{sample:product}. For \eqref{sample:raw}, partial summation in the $\lambda_{E,p}$ coordinate gives the weighted discrepancy $O_E(\sqrt x(1+g(x)))$ after multiplication by $\sqrt p$. Meanwhile $q_p\sqrt p=1+O(p^{-1/2})$, so replacing the fractional target by $\sum_{p\le x}\lambda_{E,p}$ costs $O(\sqrt x)$ by Hasse. Corollary~\ref{rawfilter:pointproduct} gives the first equivalence. Under ERH the sharp prime bound in Theorem~\ref{ep:equiv} gives $\sum_{p\le x}\lambda_{E,p}\ll_E\sqrt x\log(2x)$ by partial summation. Under ERH failure Theorem~\ref{rawfilter:main} with $\kappa=0$ gives negative excursions with any exponent below $\Theta$. An exponent strictly between $1/2+\epsilon$ and $\Theta$ absorbs the rounding error. This proves the second equivalence.

Finally both sparse expressions are constant between selected primes, including any omitted bad prime. The counting asymptotic puts a selected prime in $[Y,2Y]$ for all sufficiently large $Y$. A witness in $[2Y,2DY]$ for a fixed excursion ratio $D$ has a last selected cutoff in $[Y,2DY]$. At a negative lower-bound violation the smaller cutoff only tightens a positive-power tolerance; this also holds for an upper-bound violation. By using an intermediate exponent we absorb any fixed constants. Eventually $2D<H(Y)$. Thus testing just the selected cutoffs suffices in both cases.
\end{proof}

\begin{corollary}[The zero multiplicities are retained]\label{sample:divisor}
For this same computable set of primes and any $E/\Q$, the function initially given on $\Re s>1/2$ by
\[
         D_E(s)=\exp\!\left(\sum_{p\in\mathcal P,\ p\nmid N}
                         \frac{a_p(E)}{\sqrt p}\,p^{-s}\right)
\]
continues holomorphically to $\Re s>0$. In this half-plane its zeros, with their exact multiplicities, are those of the arithmetic function $L(E,1+s)$. In particular elliptic ERH is equivalent to the absence of zeros of $D_E$ on this half-plane.
\end{corollary}
\begin{proof}
The initial series converges absolutely by the counting bound. The discrepancy gives a single-valued holomorphic function for $\Re s>0$ by partial summation of $\sum_p(\mathbf1_{\mathcal P}(p)-q_p)\lambda_{E,p}p^{-s}$, locally uniformly on compact subsets. Replacing $q_p$ by $p^{-1/2}$ adds an absolutely convergent holomorphic function there. On $\Re s>1/2$ the logarithm of $L(E,1+s)$ has first good-prime terms $\sum_p\lambda_{E,p}p^{-1/2-s}$. Its higher powers converge absolutely locally uniformly on $\Re s>0$ by Hasse; the finitely many bad local logarithms are regular there. Thus the difference of the defining logarithms extends as a single-valued holomorphic $G(s)$ to this half-plane. The identity $D_E(s)=L(E,1+s)e^{G(s)}$ gives the continuation with every order preserved. The functional equation gives the final equivalence.
\end{proof}

\subsection{A shrinking window using only the selected cutoffs}\label{sample:shrinking}

Take the computable set $\mathcal P$ of Theorem~\ref{sample:elliptic}
with $g(x)=\log\log(3x)$. Put
\[
 S_E(x)=\sum_{\substack{p\in\mathcal P,\ p\le x\\p\nmid N}}a_p(E).
\]
Fix $0<\delta_0<1/2$ and any $g_0(\epsilon)>0$ tending to zero
with $\epsilon$. For each $E/\Q$, ERH is equivalent to the following.
For every $\epsilon>0$ there are arbitrarily large $Y$ such that
all but at most $C_{E,\epsilon}Y^{1/2-\delta_0}$ good selected primes
$p\in\mathcal P\cap I_{g_0(\epsilon)}(Y)$ satisfy
$|S_E(p)|\le C_{E,\epsilon}p^{1/2+\epsilon}$.
Here $I_\alpha(Y)=[Y,Y\exp((\log Y)^{-\alpha})]$.
The same statement holds for any common width $h(Y)$ allowed in
Corollary~\ref{absolute:criteria}.

Indeed the constant coordinate in the rounding construction gives
\[
     \#\{p\in\mathcal P:p\le x\}
       -\sum_{p\le x}\frac1{\lfloor\sqrt p\rfloor}
                              =O(1+\log\log(3x)).
\]
Subtract this estimate at $X$ and $X+L$, where $L=X^{1-\eta}$ and
$0<\eta<1/5$. The ordinary short-interval prime number theorem puts
$\gg L/\log X$ primes there, and each fractional weight is
$\gg X^{-1/2}$. Thus $\gg X^{1/2-\eta}/\log X$ selected primes
lie there, even after deleting the finitely many bad primes of $E$.
Under ERH, $B_0(x)\ll_E\sqrt x\log(2x)$ and
\eqref{sample:raw} supplies the asserted estimate at every selected
cutoff. Otherwise choose $\epsilon$ and $g_0(\epsilon)$ within the
range of Theorem~\ref{absolute:excursion} with $1/2+\epsilon<a<\Theta$.
In every large test window the proof gives an interval
$[X,X+X^{1-\eta}]$, $X\asymp Y$, with
$|B_0(x)|>x^a$ throughout, where
$0<\eta<\min(\delta_0,1/5)$. In \eqref{sample:raw} the approximation
error for $S_E$ is only $O_E(\sqrt x\log\log(3x))$. Consequently
each selected cutoff just counted violates the proposed bound.
Their number exceeds the allowed exceptions. The proof for a common
width uses the same eventual containment as Corollary~\ref{absolute:criteria}.

\subsection{Why the approximation constants depend on the curve}\label{sample:cost}
The size of a test set alone cannot make its accuracy uniform over all quadratic twists. Here the original Dirichlet theorem bounds the conductors needed to exhibit the obstruction. Fix a non-CM $E/\Q$ and any set $\mathcal P$ of primes with $\#\mathcal P(x)=o(x/\log x)$. Write $P_{E}(X)$ for the full good-prime product of Corollary~\ref{rawfilter:pointproduct}.

\begin{proposition}\label{sample:opposite}
For every sufficiently large $X$ there are two positive prime discriminants $d_1,d_2>X$ for which $\log d_i\ll_E X/\log X$ and whose twists have identical local Euler factors at every selected good prime $p\le X$. The twists may also be taken locally equal at the bad primes of $E$. Nevertheless
\[
 \log P_{E^{d_1}}(X)-\log P_{E^{d_2}}(X)\gg_E \sqrt X/\log X.
\]
Thus any approximation computed solely from the selected local data through $X$ must have absolute error at least half this difference for one of the two curves, unless it uses additional information such as the twist parameter.
\end{proposition}
\begin{proof}
For the two twists prescribe identical quadratic symbols at every selected prime and at the primes of bad reduction. At all omitted odd good primes $p\le X$, prescribe the signs so that the two twisted traces are respectively $-|a_p|$ and $|a_p|$. A zero trace imposes no restriction. At the prime two use $d_i\equiv1\pmod8$. We recall the elementary effective existence argument to make the dependence on $X$ explicit.

For $m$ prescribed odd Legendre symbols, average their indicator over prime candidates $\ell\equiv1\pmod8$. The expanded indicator, on primes coprime to the combined modulus $Q=8\prod p$, is the sum of $4\cdot2^m$ Dirichlet characters with coefficients of absolute value $1/(4\cdot2^m)$. Only the principal term has nonzero main term, by independence on the factors of the Chinese remainder decomposition. On $[z/2,z]$ the Dirichlet critical-line theorem and its standard explicit formula bound a nonprincipal prime sum weighted by $\log\ell$ by $O(\sqrt z\log^2(Qz))$. Removing prime powers costs $O(\sqrt z\log^2z)$, and the principal term contributes $z/(8\cdot2^m)$ with the same error allowance. Take $z=\exp(C X/\log X)$ with an absolute sufficiently large $C$. Here $m\le\pi(X)+O_E(1)$ and $\log Q=O_E(X)$, so the main term exceeds the total error. This produces each required prime $d_i\in[z/2,z]$. A square unit twist is locally trivial at an odd bad prime, and the congruence modulo eight gives the same property at two.

At each omitted odd good prime the difference of the two logarithms is
\[
      \log\frac{p+1+|a_p|}{p+1-|a_p|}\ \ge\ \frac{2|a_p|}{p+1}.
\]
The other good primes contribute zero. By the Sato--Tate theorem for $E$, a positive proportion of primes in $[X/2,X]$ have $|a_p|\ge\sqrt p$. Only $o(X/\log X)$ of these can have been selected. Summing the displayed lower bounds proves the claim. The two different symbols at any one such prime make the two resulting discriminants distinct.
\end{proof}

In a finite illustration of Theorem~\ref{sample:elliptic}, seven coordinates encode the count and the two observables of conductors $11,32,37$. Applied to all $6\,542$ primes through $65\,536$, the rational algorithm retains $61$ primes. An independent final pass checks all prefixes with certified absolute discrepancy at most $4$ for the count coordinate, $4$ for both conductor-$11$ coordinates, $5$ for the conductor-$32$ coordinates and $6$ for conductor $37$. The finite block can be prescribed at the start of any universal construction above. The program and the exact list are archived in \texttt{sparse\_prime\_sampler.py} and \texttt{sparse-prime-sampler.json}.

\clearpage
\section{Sparse Euler products that preserve the automorphic divisor}\label{spauto:section}
The preceding construction can keep the actual local Euler polynomial at each selected prime and raise it to an integer power. The locations and orders of all zeros and poles strictly to the right of the critical line survive. The sparse Euler product receives its continuation from the original function. Consequently the assertion here is exact compression of arithmetic data, with no new zero-free region for an individual standard function.

\begin{theorem}\label{spauto:theorem}
Fix $0\le\vartheta<1/2$ and put $c=1/2-\vartheta$. Let $\mathcal F$ be a collection of standard unitary automorphic functions, over possibly different number fields, countable modulo unitary norm twists, whose local inverse roots satisfy $|\alpha(\mathfrak p)|\le (\mathrm N\mathfrak p)^{\vartheta}$ outside a finite set for each function. For each rational prime $p$, let $L_{\pi,p}$ denote the product of the finite Euler factors at places over $p$. There is one set of rational primes $\mathcal P$ such that
\[
 \#(\mathcal P\cap[2,x])\sim\frac{x^{1-c}}{(1-c)\log x},\qquad
 \mathcal C_\pi(s)=\prod_{p\in\mathcal P} L_{\pi,p}(s)^{\lfloor p^c\rfloor}
\]
satisfies $\mathcal C_\pi(s)=L(s,\pi)\exp(G_\pi(s))$ with $G_\pi$ holomorphic on $\Re s>1/2$ for every $\pi\in\mathcal F$. Here $L(s,\pi)$ denotes its finite, uncompleted function and the product is initially taken far to the right. At the finitely many exceptional primes it suffices that the local inverse roots have modulus at most $(\mathrm N\mathfrak p)^{1/2}$; this is automatic for the standard unitary functions in question. Thus all right-half zero and pole multiplicities agree, also for non-self-dual representatives and their duals.
\end{theorem}
\begin{proof}
For $\Re s$ sufficiently large write the canonical local Euler logarithm as
\[
 \log L_{\pi,p}(s)=\sum_{k\ge1} A_{\pi,p,k}\,p^{-ks},\qquad
 A_{\pi,p,k}=\sum_{\substack{v\mid p\\ f_v\mid k}}
       \frac{f_v}{k}\sum_{j=1}^{n}\alpha_j(v)^{k/f_v} .
\]
Here $f_v$ is the residue degree, and zero inverse roots can be appended at ramified places. The bound $|A_{\pi,p,k}|\le n[F:\Q]p^{k\vartheta}$ holds at good primes for every $k$. Set the coordinates to zero at the finitely many excluded primes. Use $q_p=1/\lfloor p^c\rfloor$ in Lemma~\ref{sample:round}. Include the count coordinate and, for a countable set of representatives and all $k\ge1$, the real and imaginary coordinates of
\[
            A_{\pi,p,k}/(n[F:\Q]p^{k\vartheta}).
\]
On successive blocks use an increasing finite initial segment as in Theorem~\ref{sample:elliptic}. We can take each fixed-coordinate discrepancy up to $x$ to be $O_{\pi,k}(1+\log\log(3x))$. The count discrepancy has the same size. The prime number theorem with partial summation gives the cardinality, since $q_p=p^{-c}(1+O(p^{-c}))$ at large $p$. For a norm twist, $A_{\pi,p,k}$ is multiplied by $p^{ik\tau}$ for some real $\tau$ of fixed sign convention. On any fixed dyadic block its variation is bounded by $O(1+k|\tau|)$; the dyadic discrepancy bounds below therefore also hold for every such twist.

Let $\delta_p=\mathbf1_{\mathcal P}(p)$ and $b_p=\lfloor p^c\rfloor$. For fixed $k$ the logarithm of the quotient formally contributes $\sum_p(\delta_p-q_p)b_p A_{\pi,p,k}p^{-ks}$. On a dyadic block $X\le p<2X$ the absolute value of its sum is bounded by
\[
 O_{\pi,k,K}\!\left((1+\log\log(3X))X^{c-k(\Re s-\vartheta)}\right)
 \qquad(s\in K),
\]
where $K$ is a compact subset of $\Re s>1/2$. Indeed on that block $b_p$ is monotone of total variation $O(X^c)$, and the variation of $t^{-k(s-\vartheta)}$ is $O_{k,K}(X^{-k(\Re s-\vartheta)})$. Partial summation against the bounded prefix discrepancies proves the claim. The dyadic majorant is summable locally uniformly because $k(\Re s-\vartheta)>c$. To sum over $k$, choose $k_0$ so large that $k_0\inf_{s\in K}(\Re s-\vartheta)>c+2$. The tail for $k\ge k_0$ is absolutely and locally uniformly convergent by comparison with $\sum_p\sum_{k\ge k_0}(1+p^c)p^{-k\inf_K(\Re s-\vartheta)}$. Only finitely many individual $k$ need the discrepancy estimate. At each excluded prime its full local Euler logarithm is holomorphic for $\Re s>1/2$ by the local hypothesis. Adding these terms gives a holomorphic $G_\pi$. Equality in absolute convergence followed by continuation proves the theorem.
\end{proof}

For all fixed cuspidal ranks one may apply this theorem to the whole unitary cuspidal spectrum over every number field. There are countably many number fields. For a fixed field its unitary central characters modulo norm twists form a countable set: the norm-one id\`ele class group is compact second countable and has countable character group. For each central character the cuspidal spectrum is discrete in a separable Hilbert space. In rank one this also handles the pole of a norm character, with its multiplicity. The standard local bounds in \eqref{aw:lrs} apply in all higher ranks, so in rank $n$ the density exponent $1-c=1/2+\vartheta$ can always be taken to be $1-1/(n^2+1)$. In ranks two, three and four the stated improved local estimates give $39/64$, $6/7$, and $10/11$, respectively. For tempered forms in any rank it is $1/2$. Over general spectra this is an existence theorem; the computable elliptic construction of Section~\ref{sample:section} relies on exact point-count algorithms.

A single list of rational primes can even contain a suitable marked subsequence for each rank, with any prescribed subpower saving in its cardinality. Let $w(x)\to\infty$ be nondecreasing with $w(x)=x^{o(1)}$. For each $n$, choose $X_n$ so large that the entire rank-$n$ set just constructed has count at most $2^{-n}x/(w(x)\log x)$ at every $x\ge X_n$. Retain from it only primes larger than $X_n$ and mark them by $n$; multiple marks are allowed. The union has count at most $x/(w(x)\log x)$ by summing the estimates. Deleting finitely many primes for any fixed rank changes its quotient by a holomorphic nonvanishing factor to the right of the line. The list with its marks therefore suffices for every rank. It remains necessary to prove nonvanishing of the individual continuations there.

\clearpage
\section{A bounded discrepancy for countably many prime observables}\label{countsample:chapter}
The computable construction in Theorem~\ref{sample:elliptic} makes the error grow arbitrarily slowly. By allowing an existence construction one can bound it independently of the cutoff for each curve. We prove a general ordered rounding lemma, since a fixed-dimensional rounding constant does not immediately give a bound for infinitely many rows.

\begin{theorem}[Coordinatewise bounded prefix rounding]\label{countsample:round}
Let $x_{i,n}\in[-1,1]$ for $i,n\ge1$, and let $q_n\in[0,1]$. There is a sequence $e_n\in\{0,1\}$ such that
\[
            \sup_{h\ge1}\left|\sum_{n=1}^h(e_n-q_n)x_{i,n}\right|\le C_i
                     \qquad(i\ge1).
\]
The finite constants $C_i$ may be chosen in advance, depending only on the row index. No arithmetic independence between the rows is needed.
\end{theorem}

\begin{proof}
We first sign any finite ordered list of columns. Allow countably many auxiliary integer rows, numbered by $s\ge1$, whose entries in row $s$ have absolute value at most the positive integer $D_s$. Process rows in increasing order. Start with singleton groups, each carrying a sign pattern determined up to a common flip. At level $s$ the available groups have size $2^{s-1}$. Take them in order of their largest original position. Classify a group by the absolute value of its signed sum in row $s$, an integer between $0$ and $2^{s-1}D_s$. Pair consecutive groups in each class, and choose their relative orientation to make the combined signed sum in that row zero. A combined group retains its common flip and enters the next level. An unpaired final group in a class keeps its existing sign pattern, with arbitrary orientation. This procedure terminates on the finite list. Every completed level-$s$ group has sum zero in all integer rows through $s$.

For a fixed original prefix, consider the level-$j$ groups completed by that prefix, meaning that their largest position belongs to it. At stage $j$, in each class at most one available group completed before the cutoff fails to belong to a pair completed before the cutoff. This counts the group even when its partner is completed later. The new loss in covered original positions is therefore at most $2^{j-1}(1+2^{j-1}D_j)$. Iteration shows that at most
\[
       A_s=\sum_{j=1}^s2^{j-1}(1+2^{j-1}D_j)
                      \le 2\cdot4^s\max_{j\le s}D_j
\]
positions in the prefix remain outside a completed level-$s$ group. Hence its final signed sum in integer row $s$ has absolute value at most $D_s A_s$. This also applies beyond the last completed level: the same count covers the entire prefix when no group survives.

We encode the real rows into these integer rows. For real row $i$ use radix $B_i=2^{4\cdot2^i}$. Expand $|x_{i,n}|$ in that radix, and put its sign on every digit; for absolute value one use the expansion with every digit equal to $B_i-1$. Assign the successive digits of row $i$ to the levels
\[
          \{s\ge4^i:v_2(s)=i-1\},
\]
in increasing order. These schedules are disjoint. Unassigned levels may carry a zero row with $D_s=1$. At an assigned level use $D_s=B_i$. If $s_{i,k}$ is the level of the $k$th digit of row $i$, then
\[
       s_{i,k}\le4^i+2^i k,\qquad
                 \max_{j\le s}D_j\le2^{4\sqrt{s}} .
\]
The delayed start at $4^i$ is what ensures the second estimate even when a new row uses a very large radix.

Multiply each digit-row discrepancy by its place value. For the real row $i$ the resulting bound is
\[
 C_i' = 2 B_i\sum_{k\ge1}
           B_i^{-k}4^{\,4^i+2^i k}\,2^{4\sqrt{\,4^i+2^i k\,}} <\infty .
\]
Putting $r=2^i$ and using $\sqrt{r^2+rk}\le r+k/2$ gives the explicit estimate
\[
             C_i'\le\frac{2^{2r^2+8r+1}}{2^{2r-2}-1}.
\]
Here $r\ge2$, so the geometric series converges. Thus the finite ordered sign choice controls every prefix of every real row by $C_i'$, uniformly in the list length. The same assertion applies to any subsequence of columns in their inherited order.

To round a finite list of fractional $q_n$, first approximate them by numbers with a common dyadic denominator $2^b$ so that their total absolute approximation error is at most one. For $l=b,b-1,\ldots,1$, the coordinates with odd numerator at denominator $2^l$ may be changed by $+2^{-l}$ or $-2^{-l}$, staying in $[0,1]$ and leaving denominators at most $2^{l-1}$. Choose the signs by the preceding construction on this ordered subsequence. Each adjustment changes the discrepancy in every prefix by at most $2^{-l}C_i'$ in row $i$. The terminal coordinates are zero or one, and the bound from the original target is $C_i=C_i'+1$.

Finally each row and cutoff specifies a closed cylinder constraint in the compact space $\{0,1\}^{\mathbb N}$. Every finite set of constraints is satisfied by the finite-list argument up to its largest cutoff. Compactness supplies one infinite choice satisfying all of them.
\end{proof}

\begin{corollary}\label{countsample:elliptic}
One set of primes $\mathcal P$ satisfies, for every $E/\Q$,
\begin{align*}
 \#\{p\le x:p\in\mathcal P\}
   &=\sum_{p\le x}\frac1{\lfloor\sqrt p\rfloor}+O(1)
                \sim\frac{2\sqrt x}{\log x},\\
 \log\prod_{\substack{p\le x,\ p\in\mathcal P\\p\nmid N_E}}
       \left(\frac{\#E(\F_p)}p\right)^{\lfloor\sqrt p\rfloor}
   &=\log P_E(x)+O_E(1),\\
 \sum_{\substack{p\le x,\ p\in\mathcal P\\p\nmid N_E}}a_p(E)
   &=\sum_{\substack{p\le x\\p\nmid N_E}}\frac{a_p(E)}{\sqrt p}
                                  +O_E(\sqrt x).
\end{align*}
The point-product ratio is bounded above and away from zero for each fixed curve. For the elliptic rows the prime set may be chosen limit-computably, as shown below.
\end{corollary}
\begin{proof}
Take columns indexed by the primes and $q_p=1/\lfloor\sqrt p\rfloor$. The rows are the constant one and the two bounded observables of Theorem~\ref{sample:elliptic} for each curve, scaled by a fixed constant into $[-1,1]$. The preceding theorem gives bounded errors simultaneously in all prefix sums. In the normalized-trace row, partial summation with the factor $\sqrt p$ then costs $O_E(\sqrt x)$. Replacing the fractional coefficient $q_p\sqrt p$ by one costs $O(\sum_{p\le x}p^{-1/2})=O(\sqrt x)$. The count and product rows give the other two formulas directly.
\end{proof}

\begin{corollary}[One sampler for tempered standard functions]\label{countsample:global}
One may require the same $\mathcal P$ to have the following additional property. Fix in advance a countable list of finite-degree Euler products $L_\nu(z)$, primitive or imprimitive, in unitary normalization, with a meromorphic continuation to $\Re z>1/2$. Assume at all rational primes outside a finite set for each $\nu$ that their local inverse factors are polynomials
\[
       L_{\nu,p}(z)^{-1}=\prod_{l=1}^{d_{\nu,p}}
                        (1-\alpha_{\nu,p,l}p^{-z}),\qquad
       |\alpha_{\nu,p,l}|\le1,\quad d_{\nu,p}\le d_\nu.
\]
Omit the exceptional local factors from $L_\nu$; equivalently, include them if they have neither a zero nor a pole on $\Re z>1/2$. Write $b_\nu(p)=\sum_l\alpha_{\nu,p,l}$ at the retained primes and set $b_\nu(p)=0$ at the excluded primes. For every $\nu$ and $\tau\in\R$ the function
\[
       D_{\nu,\tau}(s)=\exp\left(\sum_{p\in\mathcal P}
                   b_\nu(p)p^{-s-i\tau}\right),\qquad \Re s>1/2,
\]
has a meromorphic continuation throughout $\Re s>0$, with the same divisor, including multiplicities, as $L_\nu(1/2+s+i\tau)$ there.

In particular the same prime indices work for every symmetric power of every non-CM rational elliptic curve and every finite-order Dirichlet twist of these powers. A single choice also works for any prescribed countable family of tempered standard functions over different number fields.
\end{corollary}

\begin{proof}
Adjoin $\Re b_\nu(p)/d_\nu$ and $\Im b_\nu(p)/d_\nu$ to the rows before applying Theorem~\ref{countsample:round}. With $\tau=0$, the discrepancy series
\[
              \sum_p(\mathbf1_{\mathcal P}(p)-q_p)b_\nu(p)p^{-s}
\]
is holomorphic for $\Re s>0$ by locally uniform partial summation. So is the correction formed with $q_p-p^{-1/2}=O(p^{-1})$. The terms of order $k\ge2$ in $\log L_\nu(1/2+s)$ converge absolutely locally uniformly in the same half-plane, since the $k$th local trace is bounded by $d_\nu$. In the common initial half-plane the difference of logarithms consequently has a single-valued holomorphic extension $G_\nu$ to $\Re s>0$. There
\[
                  D_{\nu,0}(s)=e^{G_\nu(s)}L_\nu(1/2+s).
\]
Translation $s\mapsto s+i\tau$ gives every real twist, with no additional rounding conditions.

For a number field of degree $n$, rewrite the local factors above a rational prime $p$ using $z\mapsto p^{-z}$. A factor $1-\alpha(p^{-z})^f$ at a place of residue degree $f$ splits into $f$ linear factors with roots of modulus at most one in the indicated convention. Their first trace is zero for $f>1$. The rational local degree is bounded by $n$ times the degree over the number field. Thus only the traces at residue-degree-one places enter $b_\nu(p)$. For the rational symmetric powers the good parameters have modulus one by Hasse, and automorphy and continuation are supplied by the theorems cited in Section~\ref{amp:chapter}. The finitely many omitted elliptic local factors are invertible on the required open half-plane by their tempered bounds.
\end{proof}

\subsection*{Effective information in the bounded-discrepancy statement}

\begin{proposition}\label{countsample:effective}
For uniformly computable $x_{i,n}$ and $q_n$, a choice in Theorem~\ref{countsample:round} is uniformly limit-computable. There is a computable sequence of guesses for which the first $h$ final membership bits collectively change at most $2^h-1$ times.

Conversely, for the specific probabilities $q_n=1/\lfloor\sqrt{p_n}\rfloor$ at the $n$th prime, there is a uniformly computable rational array $x_{i,n}\in[-1,1]$ such that no computable choice of bits has bounded prefix discrepancy in every coordinate, even with unspecified coordinate constants.
\end{proposition}
\begin{proof}
In the first direction take the integer thresholds $C_i=1+2^{2\cdot4^i+8\cdot2^i+1}$ justified by the proof of Theorem~\ref{countsample:round}. A finite bit string specifying a prefix that exceeds one of these thresholds in one of the computable rows can be enumerated as forbidden by rational interval approximations. Strict violation is an open condition, and every violation is eventually enumerated. At stage $t$ impose the first $t$ computation steps of this enumeration. The complement of the finitely many forbidden cylinders has a lexicographically least infinite binary sequence, computable by extending through the largest forbidden length and appending zeros. This complement is nonempty by Theorem~\ref{countsample:round}. Each stage shrinks the complement, so each length-$h$ prefix of the least point is nondecreasing in the finite lexicographic order. There are at most $2^h-1$ changes. Its limit belongs to every stage's closed complement; it is a solution.

For the converse enumerate the partial computable functions from natural numbers to natural numbers. Assign two private rows to each machine $e$. Its entries are either zero in both rows, or they form a rational vector $v\in[-1,1]^2$ with $\|v\|_\infty=1$. Start after $q_n\le1/8$ holds permanently and set earlier entries zero. Maintain the rational discrepancy vector $S$ from the machine's already completed nonzero blocks; initially $S=0$. At the next unused index start a fresh consecutive block. Use the constant vector
\[
          v=\begin{cases}
               (1,0),&S=0,\\
               (-S_2,S_1)/\max(|S_1|,|S_2|),&S\ne0
             \end{cases}
\]
through this block. End it when its cumulative target mass $w=\sum q_n$ first reaches $3/8$. It is finite because each unused tail has divergent mass, and $3/8\le w\le 1/2$.

Afterward write zeros while waiting to see the outputs of machine $e$ on the just-completed block. At index $n$ one may simulate those finitely many computations for at most $n$ steps apiece. If an output outside $\{0,1\}$ is detected, leave both rows zero forever. Once all outputs have been obtained as bits, let $b$ be their integer sum, replace $S$ by $S+(b-w)v$, and begin the next block at the following index. All decisions for entries through any fixed $n$ require only finite rational arithmetic and bounded computations, uniformly in $e,n$. In particular the array is computable even for a machine that never answers.

If the machine is a total bit procedure it completes infinitely many blocks. Its own discrepancy through each completed block is exactly the new $S$. Orthogonality gives
\[
       \|S+(b-w)v\|_2^2
         =\|S\|_2^2+(b-w)^2\|v\|_2^2
         \ge \|S\|_2^2+(3/8)^2.
\]
At least one of its two rows therefore has unbounded prefix discrepancy. Every computable candidate appears in the list. This proves the obstruction for general input arrays. It leaves open whether the special elliptic-curve rows admit a computable sampler with the strengthened errors.
\end{proof}

The constants in Theorem~\ref{countsample:round} grow with the row index. For example, take $q_n=1/2$ and the rows of a $d\times d$ Sylvester Hadamard matrix $H$, where $d$ is a power of two. At the final prefix, for the sign vector $\varepsilon=2e-1$, the squared discrepancies sum to $\|H\varepsilon/2\|_2^2=d^2/4$. Some row has discrepancy at least $\sqrt d/2$. These columns have Euclidean norm $\sqrt d$. The Koml\'os problem imposes unit Euclidean column norm and asks for one constant independent of dimension. Coordinate-dependent bounds of the form obtained here give no such uniform estimate.

\clearpage
\section{An unweighted Euler-product test on a curve-independent random set}\label{blind:section}
The rounding construction reads elliptic traces when it selects the primes. A second construction uses independent coin tosses whose probabilities depend only on $p$. It retains more primes, still of density zero. Along the resulting set even ordinary convergence of the group-order product is an exact ERH test. The product uses each selected prime once.

\begin{theorem}\label{blind:product}
For each prime $p$ retain it independently with probability $Q(\log p)=\exp(-\sqrt{\log p})$. With probability one the resulting set $\mathcal S$ has the following properties simultaneously for every rational elliptic curve $E$. The full elliptic ERH for $E$ is equivalent to convergence to a finite strictly positive real number of
\[
       P_{E,\mathcal S}(x)=\prod_{\substack{p\le x,\ p\in\mathcal S\\p\nmid N}}
                            \frac{\#E(\F_p)}p .
\]
It is also equivalent, for any fixed $H(Y)\to\infty$, to the existence of a constant $C_E$ and arbitrarily large $Y$ on which $P_{E,\mathcal S}(x)\le C_E$ throughout $[Y,YH(Y)]$. Moreover $\sum_{p\in\mathcal S}1/p<\infty$. More generally, for any prescribed nondecreasing $D(x)\to\infty$ with $D(x)=x^{o(1)}$, a retention law independent of all elliptic curves gives the same assertions together with $\#\{p\le x:p\in\mathcal S\}=O(x/(D(x)\log x))$ almost surely.
\end{theorem}

\begin{proof}
We give conditions on $Q$ that cover both constructions. Take $Q(t)$ to be the Laplace transform of a finite positive measure of mass at most one on $[0,\infty)$ whose support has infimum zero. Require $Q(t)\ll(1+t)^{-4}$. Then $Q(t)=\exp(-o(t))$, and $-Q'(t)/Q(t)\to0$ as $t\to\infty$. For the latter assertion the probability proportional to $e^{-ty}d\nu(y)$ concentrates at zero: for every $a>0$ the integral over $y>a$ of $y e^{-ty}$ is exponentially smaller than the mass on $[0,a/4]$, for $t$ large. The example $Q(t)=e^{-\sqrt t}$ is the Laplace transform of $(2\sqrt\pi)^{-1}y^{-3/2}\exp(-1/(4y))\,dy$; the Gaussian integral verifies the identity.

Write $\xi_p=\mathbf1_{\mathcal S}(p)$ and $\lambda_p=a_p/\sqrt p$ at good primes, zero at bad ones. The independent centered series
\[
       \sum_p(\xi_p-Q(\log p))\lambda_p/\sqrt p
\]
has finite total variance, at most $4\sum_p Q(\log p)/p<\infty$. It converges almost surely by the square-integrable martingale convergence theorem. The same holds for every rational $E$ simultaneously by countability, and $\sum_{p\in\mathcal S}1/p<\infty$ almost surely by its finite mean. The Taylor estimate for the logarithm of a local group order now gives
\[
 \log P_{E,\mathcal S}(x)=-\sum_{p\le x}Q(\log p)\frac{a_p}{p}+O_E(1),
\]
where the error in fact tends to a finite limit. Under ERH, the full sum $\sum_{p\le x}a_p/p$ is $O_E(\log^2(2x))$ by Corollary~\ref{rawfilter:pointproduct} and its proof. Partial summation against $Q$ gives convergence of the displayed weighted sum: $t^2Q(t)\to0$ and $\int_1^\infty t^2(-Q'(t))dt<\infty$ by integration by parts. This proves the positive finite limit, and hence the asserted upper bound.

For the converse we extend the positive cutoff in Lemma~\ref{rawfilter:untaper}. Set $r(t)=Q(t)/t$, itself the Laplace transform of a positive measure $d\omega(y)$, by convolution of $d\nu$ with Lebesgue measure on the positive half-line. Use the same smooth compact probability $\mu$ chosen in Theorem~\ref{rawfilter:main} for $\kappa=-1/2$, supported on $[0,c]$ with transform $M$. Put
\[
 k_r(t)=\int_0^\infty\frac{e^{-ty}}{M(y)}\,d\omega(y),\qquad
 W_U(t)=\frac{1}{r(t)}\int_{v\ge t-U}k_r(t-v)\,d\mu(v)\qquad(t>c).
\]
There is again $W_U=1$ below $U$, and $W_U=0$ above $U+c$. To see that $W_U$ decreases, let $\nu_y$ denote the $e^{yv}$ tilt of $\mu$. Then $W_U(t)$ is the mean of $\nu_y([t-U,\infty))$ with respect to $e^{-ty}d\omega(y)/r(t)$. The integrand increases in $y$, by the covariance argument of Lemma~\ref{rawfilter:untaper}. As $t$ increases its upper interval shrinks and the mixing distribution moves towards smaller $y$. The latter follows from the same covariance calculation, with $-y$ in the derivative. Thus $W_U$ decreases, and it is continuous by the smoothness of $\mu$. It defines a positive average of sharp prefixes precisely as before.

The inequalities $r(t+c)\le k_r(t)\le r(t)$ show that $k_r(t)\sim r(t)$, since $-r'(t)/r(t)\to0$. For each $z$ with $\Re z>0$ and fixed $A>c+1$,
\[
            \int_A^U e^{zt}k_r(t)\,dt\sim e^{zU}r(U)/z.
\]
Indeed substitute $w=U-t$: for each fixed $w$ the ratio $k_r(U-w)/r(U)$ tends to one, and for $U-w$ sufficiently large it is at most $e^{\epsilon w}$ for any fixed $0<\epsilon<\Re z$. The bounded initial range is negligible since $r(U)=e^{-o(U)}$. Exponential domination gives the asserted asymptotic, including its complex phase.

Apply the Mellin inversion proof of Theorem~\ref{rawfilter:main} with $k_r$ in place of $k$. Above $A+c$, convolving $\mu$ with $k_r\mathbf1_{[A,U]}$ gives $r(t)W_U(t)$. Thus the first powers in the logarithmic derivative give precisely the weighted sharp sums with factor $Q(\log p)$. All earlier lower-zero and contour estimates still apply, since $k_r(t)\le 1/t$; the deletion of prime powers is bounded by the previous estimate as $0<Q\le1$. At the selected zero $z_*=\rho_*-1/2$ the main term has magnitude of order $e^{(\Re\rho_*-1/2)U}Q(U)/U$ at a selected phase. Since $Q(U)=e^{-o(U)}$, its negative excursion dominates $x^a$ for each $a<\Re\rho_*-1/2$. In the real-edge case the leading coefficient has the same strict negative sign as before. Consequently ERH failure gives a fixed $D_0>1$ and a positive $a$ such that every sufficiently large $[Y,D_0Y]$ contains an $x$ with $\log P_{E,\mathcal S}(x)>x^a$. This rules out the claimed window bound and any finite positive limit.

For the adjustable sparsity pick $W(t)=(1+t)^4D(e^{2t})^2$ for $t\ge1$. Then $\log W(t)=o(t)$. For $j\ge1$ set $C_j=1+\sup_{t\ge1}e^{-t/j}W(t)<\infty$ and take
\[
                  Q(t)=\sum_{j\ge1}2^{-j}C_j^{-1}e^{-t/j}.
\]
This is an admissible Laplace transform and $Q(t)\le1/W(t)$ for $t\ge1$. Its expected count through $x$ is at most $\sqrt x+\pi(x)/D(x)^2$, by splitting at $\sqrt x$. The centered count through $x$ is $O(x^{3/4})$ almost surely: the sum of Bernoulli variances through $2^{j+1}$ is at most $2^{j+1}$, and the maximal $L^2$ inequality gives a summable exceptional probability $O(2^{-j/2})$ at threshold $2^{3j/4}$. The Borel--Cantelli lemma and $D(x)=x^{o(1)}$ prove the count claimed. For $Q(t)=e^{-\sqrt t}$ the same argument gives density zero directly. The full-measure event for the fluctuations does not depend on $H$, so the window quantifier is simultaneous as stated.
\end{proof}

There is a shorter way to see why actual convergence detects ERH. Bounded partial sums of $\sum Q(\log p)a_p/p$ imply convergence of $\sum (a_p/p)p^{-s}$ on $\Re s>0$ after partial summation with $e^{-st}/Q(t)$, whose variation is integrable there by $-Q'/Q\to0$ and $Q(t)=e^{-o(t)}$. The terms of higher prime powers in $\log L(E,1+s)$ converge absolutely in that half-plane. Exponentiation gives nonvanishing. The positive cutoff above is needed to deduce the stronger conclusion from upper bounds confined to arbitrarily distant intervals.

\subsection{The zero boundary measured by polynomial thinning}\label{blind:polynomial}

Let $Q$ be any one of the admissible completely monotone functions in
Theorem~\ref{blind:product}. Choose independent variables $V_p$
uniform on $[0,1]$, once for all curves, and set
\[
  \mathcal S_a=\{p:V_p\le p^{-a}Q(\log p)\},\qquad a\ge0.
\]
For each fixed $a$ the following statements hold almost surely for
every rational elliptic curve $E$. They hold simultaneously for all
nonnegative rational $a$ on one event of probability one.
\begin{enumerate}
\item The ordinary unweighted product of $\#E(\F_p)/p$ over good
$p\in\mathcal S_a$ converges to a finite positive limit if and only if
$\Theta_E\le1/2+a$.
For any $H(Y)\to\infty$ the same boundary is equivalent to a fixed
upper bound for the product throughout arbitrarily distant
$[Y,YH(Y)]$.
\item The product converges absolutely for $a\ge1/2$.
For $0\le a<1/2$ its logarithmic series fails to converge absolutely,
including when it converges conditionally. For $0\le a<1$,
\[
       \#\{p\le x:p\in\mathcal S_a\}
          \sim\frac{x^{1-a}Q(\log x)}{(1-a)\log x},
\]
while $\mathcal S_a$ is finite almost surely for $a\ge1$.
\end{enumerate}
In particular this nested random family recovers the exact rightmost
real-part supremum simultaneously for every rational elliptic curve:
\[
 \Theta_E=\tfrac12+\inf\{a\in\Q_{\ge0}:
              P_{E,\mathcal S_a}(x)\text{ has a positive finite limit}\}.
\]

\begin{proof}
Write $t=\log p$. The variance and local Taylor estimates of
Theorem~\ref{blind:product} still apply, since $p^{-a}Q(t)\le Q(t)$.
Thus $\log P_{E,\mathcal S_a}(x)$ differs by a term with a finite
limit from $-\sum_{p\le x,\ p\nmid N} Q(\log p)a_p/p^{1+a}$.
If $\Theta_E\le1/2+a$, the sharp explicit formula and partial summation
give $A_E(x)\ll_E x^{1+a}\log(2x)$; for $\Theta_E=1/2$ use
Theorem~\ref{ep:equiv}, and for $\Theta_E>1/2$ use the global estimate
in Theorem~\ref{mass:block}. Hence
$\sum_{p\le x,\ p\nmid N}a_p/p^{1+a}=O_{E,a}(\log^2(2x))$.
Partial summation against $Q$ proves convergence, including at equality
of the zero boundary.

For the other direction suppose $\Theta_E>1/2+a$. Choose
$1/2+a<b<\Theta_E$, above all real zeros in the nonreal-edge case.
Use the same probability filter $\mu$ as in the proof of
Theorem~\ref{blind:product}, for $z=\rho-1/2$. Now replace $r(t)$ by
$e^{-at}Q(t)/t$. Its representing positive measure is the original
one translated by $a$, so its support has infimum $a$. The monotone
upper-tail proof and the exact averaging identity are unchanged.
In the quotient $k_r(t)/r(t)$ the exponentially tilted distribution
of the representing variable concentrates at $a$. As $1/M$ is
continuous and bounded by one on the positive axis,
$k_r(t)/r(t)\to1/M(a)$. For each fixed selected $z$ with
$\Re z>a$, substitution $w=U-t$ and exponential domination as in
Theorem~\ref{blind:product} give
\[
   \int_A^U e^{zt}k_r(t)\,dt
       \sim\frac{e^{(z-a)U}Q(U)}{M(a)(z-a)U}.
\]
Here $e^{at}r(t)=Q(t)/t$ has logarithmic derivative tending to zero.
The lower zeros in the smoothed formula cost
$O(e^{(b-1/2-a)U})$, using $k_r(t)\le e^{-at}/t$ and rapid decay of
$M$. The shifted contour is $O(1+U)$, and higher prime powers for
$a>0$ cost $O(1)$ by absolute convergence (the case $a=0$ was
already proved). Thus the selected conjugate pair gives a negative
weighted sharp excursion after choosing its phase, or the real leading
zero gives a negative term directly. The factor $1/M(a)$ is positive.
The resulting logarithm of the point product exceeds $x^d$ somewhere
in every sufficiently large fixed-ratio interval, for some $d>0$,
since $Q(U)=e^{-o(U)}$. This excludes both proposed bounds.

For absolute convergence use
$|\log(\#E(\F_p)/p)|\le 2/\sqrt p+O(1/p)$.
At $a\ge1/2$ the expected absolute sum is finite, including the
endpoint by $\sum_p Q(\log p)/p<\infty$. For $a<1/2$, the fixed
Rankin--Selberg prime number theorem gives
$\sum_{x<p\le2x}|\lambda_p|^2\log p\sim x$ at good primes;
the higher powers are negligible by Hasse. Since $|\lambda_p|\le2$,
it follows that $\sum_{x<p\le2x}|\lambda_p|\gg_E x/\log x$.
The expected absolute logarithmic sum on this dyadic block is
$\gg x^{1/2-a-o(1)}/\log x$, by the reverse local Taylor inequality
and monotonicity of $Q$. The centered series of these absolute
logarithms again has finite total variance, so the absolute sum
diverges almost surely.

Finally, partial summation of the prime number theorem against the
decreasing weight $x^{-a}Q(\log x)$ gives the stated asymptotic for
the expected count when $a<1$. Its comparison integral is
$\int_2^x y^{-a}Q(\log y)\,dy/\log y$; differentiating the displayed
asymptotic equivalent gives the integrand times $1+o(1)$ because
$Q'/Q\to0$. If $M_p$ denotes the expected count through $p$, then
$\sum_p\operatorname{Var}(\mathbf1_{p\in\mathcal S_a})/(1+M_p)^2
\le\int_0^\infty(1+t)^{-2}dt$. Martingale convergence and
Kronecker's lemma give an actual count asymptotic to its divergent
expectation. At $a\ge1$ the expectation is finite. Taking the
countable intersection over rational $a$ and curves proves the joint
claim and the infimum formula.
\end{proof}

\begin{theorem}[Pure power sampling]\label{blind:pure}
Fix $a>0$ and now retain $p$ with probability exactly $p^{-a}$.
Almost surely the following hold for all rational elliptic curves.
The unweighted product has a finite positive limit precisely when
$L(E,1+s)$ has no zero in $\Re s>a$ and no zero at the real point
$s=a$. Equivalently, $\Theta_E\le1/2+a$ and $1/2+a$ is not a real
unitary zero. The fixed upper-bound test on arbitrarily distant
$[Y,YH(Y)]$ has exactly the same criterion.
If $\Theta_E=1/2+a$ and $1/2+a$ is a real zero of multiplicity $m$,
then for some almost surely finite positive $C_{E,a}$ one has
\[
                   P_{E,\mathcal S_a}(x)
                      \sim C_{E,a}(\log x)^m .
\]
For $a=1/2$ the product converges conditionally for every $E$.
Absolute convergence holds exactly for $a>1/2$.
\end{theorem}

\begin{proof}
The centered series again converges almost surely because
$\sum_p p^{-1-a}<\infty$, and the random finite harmonic sum
controls the logarithmic Taylor remainder. Thus the deterministic
quantity to analyze is $-\sum_{p\le x,\ p\nmid N}a_p/p^{1+a}$.
For $\Theta_E<1/2+a$ convergence follows by partial summation of
the global prime bound at a strictly smaller exponent. For
$\Theta_E>1/2+a$ the preceding compact positive filter proof applies
with $Q=1$ and $r(t)=e^{-at}/t$: its translated Laplace measure is
Lebesgue measure on $[a,\infty)$, and the higher prime powers are
absolutely summable because $a>0$. Thus there are unbounded positive
power excursions in every sufficiently late fixed-ratio block.

It remains to examine the boundary. Take $b\in(1/2,\Theta_E)$ in
the sharp expansion $(*)$ of Theorem~\ref{absolute:excursion}, with
$\kappa=-1/2$. Write $t=\log x$, $w=\rho-1/2$ and take a fixed
$T>1$. The error in the log-weighted sharp sum there is
$O((1+t)^2e^{(a-\delta)t})$ for some $\delta>0$.
Stieltjes partial summation against $e^{-at}/t$ therefore makes its
contribution converge to a constant. Each leading term
$-e^{wt}/w$ contributes
\[
                -\int_T^{\log x}\frac{e^{(w-a)t}}t\,dt
\]
to the deterministic sum of the traces, up to the fixed lower-cutoff
convention. To justify termwise integration first apply ordinary
integration by parts on a compact $t$ interval, where the leading
series with coefficient $1/w$ converges absolutely uniformly.
For each nonreal zero, $\Re(w-a)\le0$, and integration by parts in
the last integral bounds its tail on $[R,S]$, $S\ge R\ge T$, by
$3/(R|\Im\rho|)$. These bounds are summable over all the nonreal
zeros in the indicated strip. The finitely many real terms with
$w<a$ converge. A real term with $w=a$ gives $-\log\log x+O(1)$
per multiplicity. Recalling the negative sign in the logarithm of
the point product proves both alternatives, with an actual convergent
constant. A real zero at unitary abscissa 1 is excluded
by boundary nonvanishing, giving convergence at $a=1/2$.

For absolute convergence use the same variance and
Rankin--Selberg prime estimates as above. At $a<1/2$ the dyadic
expected absolute sums grow without bound. At $a=1/2$ their lower
bound at $X=2^j$ is $\gg_E1/j$ after the summable local Taylor error.
Thus their sum diverges there too. For $a>1/2$ the expectation is
finite. This proves the sharp absolute threshold.
\end{proof}

For pure powers with $0<a<1$ the counting asymptotic becomes
$x^{1-a}/((1-a)\log x)$. At $a=1$ there are almost surely
$(1+o(1))\log\log x$ selected primes through $x$, by the prime
harmonic sum and the same martingale strong law. Above $a=1$ only
finitely many primes are selected. Independent uniforms give a joint
coupling for all positive rational exponents, also in this pure case.

\clearpage
\section{One coefficient-independent random prime set in every cuspidal rank}\label{blindall:section}
For elliptic curves the local factors at the critical point are ordinary point counts. For a general standard function it is useful to retain the logarithmic coefficients of the higher prime powers separately. The set of rational primes can be the same for all ranks, all fields, and every imaginary norm twist, selected without referring to an automorphic coefficient.

\begin{theorem}\label{blindall:main}
Choose any of the independent prime laws constructed in Theorem~\ref{blind:product}, and keep the resulting $\mathcal S$. With probability one the following holds simultaneously for every unitary cuspidal standard function $L(s,\pi)$ over every number field whose completion is entire. Use the rational local logarithmic coefficients $A_{\pi,p,k}$ of Theorem~\ref{spauto:theorem}. The ERH for this individual function is equivalent to convergence to a finite complex limit, as $x\to\infty$, of the sums
\[
 \sum_{\substack{p\in\mathcal S,\ k\ge1\\p^k\le x}}
       \frac{Q(k\log p)}{Q(\log p)}\,A_{\pi,p,k}\,p^{-k/2} .
\]
If all unramified inverse roots of $\pi$ are tempered, the same equivalence holds for convergence to a finite nonzero limit of the ordinary partial Euler product $\prod_{p\le x,\,p\in\mathcal S}L_{\pi,p}(1/2)$. This also holds for every unitary cuspidal $\mathrm{GL}_2/\Q$ function without a temperedness assumption. Each selected prime occurs once in that product, with all its local factors. Finite changes at ramified primes are immaterial.
\end{theorem}

\begin{proof}
Write $n=p^k$ and $A(n)=A_{\pi,p,k}$, with zero at other positive integers. All bounds in the proof are for a fixed function unless a compact set of translates is indicated. We use the standard local bound $|A(p^k)|\ll_{\pi}p^{k\vartheta}$ for some $\vartheta<1/2$ and the estimate
\begin{equation}\label{blindall:second}
 \sum_{p,k}\frac{|A(p^k)|^2}{p^{k(1+v)}}\ll_\pi 1+\log(1/v)
                      \qquad(0<v\le1/2).
\end{equation}
To prove the latter, first remove the finitely many ramified rational primes, which contribute boundedly by the local bound. For $v_0\mid p$ of residue degree $f$, the term in $A(p^k)$ when $f\mid k$ is $a(v_0,j)/j$ with $j=k/f$. Cauchy--Schwarz costs at most $[F:\Q]$ when summing over places above $p$. Since $1/j^2\le1/j$, the resulting nonnegative series is bounded by a constant times the unramified Euler logarithm of $L(1+v,\pi\times\widetilde\pi)$: its coefficient at each $(v_0,j)$ is exactly $|a(v_0,j)|^2/j$. This incomplete Rankin--Selberg function has a simple pole at one, giving \eqref{blindall:second}. Equivalently one can apply the standard unramified Rankin--Selberg prime-coefficient bound. The local root bounds and these analytic properties are recalled in Section~\ref{aw:section} and \href{https://arxiv.org/abs/2103.05634}{Humphries--Thorner, \S3}.

It follows that $\sum_{n\le e^T}|A(n)|^2/n\ll_\pi1+\log(2T)$, on taking $v=1/T$. Partial summation and $Q(t)\ll(1+t)^{-4}$ imply
\[
       \sum_{n\ge2}\frac{|A(n)|^2}{n}(1+\log^2 n)Q(\log n)<\infty.
\]
Set $q_p=Q(\log p)$ and $\xi_p=\mathbf1_{\mathcal S}(p)$. For each fixed $k$, the independent centered sum over the rational primes with terms
\[
      (\xi_p-q_p)\frac{Q(k\log p)}{q_p}A(p^k)p^{-k/2-ik\tau}
\]
has finite total variance, uniformly also after one derivative in $\tau$. Indeed $Q(k\log p)\le q_p$, so its variance with the derivative included is bounded by a fixed subseries of the preceding display. More precisely on each fixed interval $[-m,m]$ the sum of expected squared $H^1([-m,m])$ norms is finite. The partial sums form an $L^2$-bounded Hilbert-space martingale; they converge almost surely in $H^1$, and the one-dimensional Sobolev inequality gives uniform convergence in $\tau$ there. Countably many $m$ allow every real translate. For $k$ sufficiently large, the whole double tail is even absolutely convergent, deterministically, since $Q(k\log p)/q_p\le1$ and $\sum_p\sum_{k\ge k_0}p^{-k(1/2-\vartheta)}<\infty$. Thus the difference between the sampled sum and
\[
                      \sum_{n\le x} Q(\log n)A(n)n^{-1/2-i\tau}
\]
converges for each real $\tau$. The order $p^k\le x$ causes no rearrangement problem: only finitely many fixed-$k$ subseries remain after the absolutely convergent tail. Countability of automorphic functions modulo norm twists, established in Section~\ref{spauto:section}, gives a single full-measure event for all the families.

It remains to identify the deterministic weighted convergence. Under ERH,
\[
              \sum_{n\le x}A(n)\log n\ll_\pi \sqrt{x}\log^2(2x).
\]
For clarity, apply the complex explicit formula \eqref{aw:explicit} at an integer $m$ and at $m+1/2$. The difference taper of \eqref{aw:difference} before taking real parts, divided by $1/(2m)$, is exactly the sharp sum at $m$. Each zero on the line costs $O(\sqrt m\min((1+|\gamma|)^{-1},m/(1+\gamma^2)))$. The unit-interval zero count gives $O(\sqrt m\log^2(2m))$ on summing. The finitely many archimedean terms cost $O_\pi(m^{\vartheta}+\log(2m))$ after differencing and the negative contour may be taken below $-2$, as in Section~\ref{aw:pointsection}. All zeros are counted with multiplicity. For each fixed translate repeat the argument on that function. Partial summation against $x^{-1/2}/\log x$ gives $\sum_{n\le x}A(n)n^{-1/2}=O_\pi(\log^2(2x))$. Damping by $Q$ now gives a convergent series exactly as in Theorem~\ref{blind:product}.

Conversely, bounded partial sums of the deterministic damped series give local uniform convergence of $\sum_n A(n)n^{-s}$ for $\Re s>1/2$, by partial summation against $e^{-(s-1/2)t}/Q(t)$. Its derivative in $t$ is integrable uniformly on compact $s$-sets in this region because $Q(t)=e^{-o(t)}$ and $-Q'(t)/Q(t)\to0$. Exponentiating continues the ordinary Euler logarithm and excludes all zeros there. Reflection of the individual unitary completion across the critical line excludes any on the other side. This proves the equivalence.

For tempered factors, $|A(p^k)|\ll_\pi1$ for all good $p,k$. The expected sum of absolute values of the selected higher powers at $1/2$ is $O_\pi(\sum_p q_p/p)<\infty$. Thus their ordinary sum is absolutely convergent almost surely. For $\mathrm{GL}_2/\Q$ the good second-power trace is $\lambda(p)^2-2\omega(p)$, where $|\omega(p)|=1$. Its contribution is absolutely summable in expectation by the preceding Rankin--Selberg estimate at first powers. For $k\ge3$ the bound $7/64$ gives $3(1/2-7/64)>1$ and hence absolute convergence even before selection. These absolute bounds are unchanged by imaginary twists and again hold for all representatives by countability. The damped higher terms and the ordinary higher terms consequently both converge; their difference is inessential. Each canonical logarithm of a local product tends to zero as $p\to\infty$. Hence convergence of the ordinary product to a nonzero complex number is equivalent to convergence of the sum of these logarithms: eventually the partial products stay in a small disk avoiding zero and one can use its single logarithm branch, since each successive logarithmic increment tends to zero.
\end{proof}

For a norm character with a pole, one can instead subtract at every $p^k$ the known coefficients of the rational zeta function with the identical imaginary shift. Its simple pole cancels the one over the number field. The Dirichlet theorem on rational zeta places its nontrivial zeros on the line. The same proof with this subtraction gives the right-half test for the pole-removed quotient: the explicit prime-sum main terms cancel, and the second moment bound survives the additional degree-one coefficients. This paragraph is the only place in this section requiring the previously verified rational Dirichlet result.

\clearpage
\section{Which Frobenius classes detect a zero?}\label{rfc:section}

For a rational elliptic curve \(E\), ordinary point counts at just the
completely split primes can detect all off-line zeros over any abelian
extension of \(\Q\). At other Frobenius classes there can be exact
cancellation. We also give a criterion in which the quadratic field can
change at every prime cutoff.

\subsection{Fixed fields and cancellation}
Use the conventions of Section~\ref{ao:split}. If \(K/\Q\) is finite
abelian with group \(G\) and \(g\in G\), put
\[
 B_{\kappa,g}(x)=\sum_{\substack{p\le x\\\operatorname{Frob}_p=g}}
                    p^\kappa\lambda_p,\qquad
 r_g(\rho)=\frac1{|G|}\sum_{\chi\in\widehat G}
               \overline{\chi(g)}\,m_\chi(\rho),
 \qquad -1/2\le\kappa\le1/2 .
\]
All bad or ramified primes are omitted. The multiplicity \(m_\chi\) is
taken in the completed unitary twist of \(E\). Set
\(\Theta_g=\sup(\{1/2\}\cup\{\Re\rho:r_g(\rho)\ne0\})\).

\begin{theorem}\label{rfc:class}
Suppose \(\Theta_g>1/2\) and \(a<\Theta_g+\kappa\).
There is a \(D>1\) such that each sufficiently large \([X,DX]\)
contains the following excursions.
\begin{enumerate}
\item If \(r_g(\Theta_g)=0\), both \(B_{\kappa,g}(x)>x^a\) and
\(B_{\kappa,g}(x)<-x^a\) occur. If the supremum is unattained by
the effective zeros, every prescribed \(D>1\) works.
\item If \(r_g(\Theta_g)\ne0\), one has a value of magnitude \(>x^a\)
with sign opposite to \(r_g(\Theta_g)\).
\end{enumerate}
For \(g=1\) the one-sided criterion in Theorem~\ref{rawfilter:main}
with \(B_{1/2,1}\) is equivalent to ERH for \(E/K\).
The upper point-product criterion in Corollary~\ref{rawfilter:pointproduct}
is likewise equivalent if its product is taken over just the completely
split good primes.
\end{theorem}
\begin{proof}
The logarithmic derivative for the finite linear combination of twists
has coefficient at the unramified \(p^j\) given by character orthogonality,
with selector \(\mathbf1_{\operatorname{Frob}_p^j=g}\). Its zero
contributions after the shift by \(\kappa\) are
\(-r_g(\rho)M(\rho+\kappa)\int_A^Ue^{(\rho+\kappa)t}k(t)\,dt\).
Conjugate points have conjugate residues because the prime coefficients
are real. The absolute sum of inverse ordinates on every
\(\Re\rho\ge b>1/2\), counted over the finite collection of twists, is
finite by the same rational degree-two density theorem. The finite
prime omissions and the prime powers obey the bounds in
\eqref{rawfilter:average}. In particular one applies that formula to
the combination before deleting the higher powers. The fixed twists
have the usual weight-two gamma factors and functional equations.

Use the positive measure of Theorem~\ref{rawfilter:main} on the set of
effective zeros with \(r_g(\rho)\ne0\). In the nonreal case the
retained amplitude is multiplied by the nonzero complex number
\(r_g(\rho_*)\), which only changes its size and phase. In the real
case annihilate all effective nonreal pairs. The surviving real
coefficient \(-r_g(\Theta_g)M(\Theta_g+\kappa)/(\Theta_g+\kappa)\)
has the asserted sign. In the split class every residue is a
nonnegative real number. Its value is the multiplicity of the
base-change zero divided by \(|G|\), since the completed base-change
function is the product of the twists. Thus every effective real
edge forces the negative excursion needed in the one-sided criterion.
When the entire divisor is on the line, the sharp bound and the
point-product estimate follow by adding the finite set of explicit
formulas as in Section~\ref{ao:split}, followed by partial summation.
The identity with the logarithm of the restricted point product
has error \(O_E(\log\log(3x))\), by the same local Taylor
estimate irrespective of which primes are selected.
\end{proof}

The splitting assertion has the following extension. Suppose
\(K/\Q\) is finite and the completed base-change function of \(E/K\)
is a product of fixed cuspidal unitary \(\mathrm{GL}_2/\Q\) functions
with positive integral multiplicities, with unitary variable used
in all factors. Put \(u_K(p)=\#\{\mathfrak p\mid p:f(\mathfrak p/p)=1\}\).
Then the same two criteria apply with the weight \(u_K(p)\) in the
sharp sum and as an exponent in the point product. Indeed adding the
explicit formulas gives exactly the zero multiplicities of the product.
In the prime powers with \(f(\mathfrak p/p)j\ge2\), group by
\(p^{fj}\). The number of places over a given \(p\) is bounded by
\([K:\Q]\), and for each exponent there are at most \([K:\Q]\)
choices for \(f\). Their total error after unweighting is still
\(O_{E,K,\epsilon}(X^{\kappa+1/2+\epsilon})\). The surviving primes
give precisely \(u_K(p)\lambda_p\). Equivalently this comparison
can be made first on the common good Euler factors, using
Section~\ref{ao:split}; the finitely many exceptional local terms have
the same harmless bound by elliptic temperedness, or, for an auxiliary
fixed rational unitary \(\mathrm{GL}_2\) factor, the strict local
exponent bound \(<1/2\). For general archimedean parameters one
uses the smooth explicit formula of Section~\ref{aw:section}, or
shifts the Mellin integral to a line below \(1/2\) and above all the
finitely many rightmost trivial zeros. Those gamma-factor exponents
are strictly less than \(1/2\) in the unitary normalization. When
\(K/\Q\) is Galois, \(u_K(p)=[K:\Q]\) at split primes and is zero
at the others.

For a nonidentity conjugacy class \(C\) of a nonabelian Galois
extension, suppose in addition that each \(V_E\otimes\tau\) for an
irreducible finite-group representation \(\tau\) is realized by a
fixed positive product of cuspidal rational degree-two functions,
with unitary local parameters at all common good primes.
This holds, for example, in the rational elliptic Weil restriction
decompositions of Section~\ref{ao:split}. Replace \(r_g\) above by
\[
       r_C(\rho)=\frac{|C|}{|G|}\sum_\tau
                  \overline{\operatorname{tr}\tau(c)}\,m_\tau(\rho),
                      \qquad c\in C ,
\]
where \(m_\tau\) counts the zeros of the indicated product. Finite
group character orthogonality proves Theorem~\ref{rfc:class} for this
class and residue. A splitting-only factorization of the full product
is sufficient for the preceding \(u_K\) assertion.

\begin{proposition}\label{rfc:cyclic}
Let \(g\) have order \(n>1\) in the abelian group \(G\).
Fix a primitive \(n\)-th root \(\zeta\) and put
\[
       M_j=\sum_{\chi(g)=\zeta^j}m_\chi(\rho),\qquad
       P(T)=\sum_{j=0}^{n-1}M_jT^j.
\]
A zero is invisible in the \(g\)-class precisely if \(\Phi_n(T)\)
divides \(P(T)\) over \(\Q\). If \(n=\ell^v\) for a prime \(\ell\),
this says that for each \(0\le j<\ell^{v-1}\) the \(\ell\) entries
\(M_{j+k\ell^{v-1}}\), \(0\le k<\ell\), agree.
In a cyclic extension of prime degree \(\ell\), any one nonsplit
class sees the zero unless \emph{all} the twists have identical
multiplicity there. In particular every zero of base-change
multiplicity strictly between \(0\) and \(\ell\) is seen by every
nonsplit class.
\end{proposition}
\begin{proof}
The residue equals \(P(\zeta^{-1})/|G|\). The minimal polynomial of
\(\zeta^{-1}\) is \(\Phi_n\). For \(n=\ell^v\) this polynomial is
\(\sum_{k=0}^{\ell-1}T^{k\ell^{v-1}}\), and the quotient must have
degree less than \(\ell^{v-1}\), so there is no overlap among the
copies of its coefficient vector in the product. In the prime-degree
case the nonidentity \(g\) is a generator and every fiber contains
one character. The base-change multiplicity is the sum of all
the entries.
\end{proof}

Exact cancellation does occur in the elliptic setting. Let \(E\) have
CM by an imaginary quadratic field \(K\) (over \(\overline{\Q}\)) and
take that field for the quadratic extension. Its quadratic character
is a self-twist of the induced two-dimensional parameter, since its
restriction to the inducing subgroup is trivial. Thus the two
multiplicities are equal at \emph{every} point, and all the inert
traces vanish. For \(E:y^2=x^3-x\) this is the elementary cancellation
at \(p\equiv3\bmod4\) in Section~\ref{ip:section}. One prescribed
inert class can therefore carry no standard trace information at all.

\subsection{A fixed off-line zero is visible at almost every quadratic inert class}
Let \(\mathcal D\) be the positive squarefree \(d\equiv1\bmod4\)
prime to \(2N\). Each \(d>1\) is a fundamental discriminant. For a
compact set \(K_0\subset\{\Re s>1/2\}\) we need the following
conductor-aspect consequence of the quadratic large sieve.

\begin{lemma}\label{rfc:compact}
For a fixed \(E\) and \(K_0\) there is \(\delta>0\) such that
\[
 \#\{d\in\mathcal D:D<d\le2D,\ L(s,\pi_E\otimes\chi_d)
                              \text{ vanishes somewhere in }K_0\}
                         \ll_{E,K_0} D^{1-\delta}.
\]
\end{lemma}
\begin{proof}
Nonvanishing on \(\Re s>1\) allows us to cover the relevant compact
set by the interior of a bounded rectangle \(\mathcal R\) with
\(1/2<\sigma_0\le\Re s\le\sigma_1<5/4\), taking \(\sigma_0<1\).
The distance of \(K_0\cap\{\Re s\le1\}\) to its boundary is positive;
an empty intersection needs no estimate. Use the family norm
\(\|b_d\|_D=(D^{-1}\sum_{d\in\mathcal D,\ D<d\le2D}|b_d|^2)^{1/2}\).
We first record the two uniform mean-square bounds entering the proof.

\href{https://doi.org/10.4064/aa-72-3-235-275}{Heath-Brown's quadratic
large sieve, Theorem~1}, bounds the squared, unnormalized family norm of
\(\sum_{m\le Q}b_m\chi_d(m)\) over squarefree odd \(m\) by
\((D+Q)(DQ)^\eta\sum |b_m|^2\), up to a constant depending on
\(\eta>0\). One may allow even squarefree \(m\) by separating a factor
two. For coefficients \(|b_n|\ll_\eta n^\eta\), \(Q\le D\), and
\(\sigma_0\le\Re s\le\sigma_1\), it follows that
\begin{equation}\label{rfc:tail}
  \left\|\sum_{R<n\le Q} b_n\chi_d(n)n^{-s}\right\|_D
                         \ll_\eta D^\eta R^{1/2-\sigma_0}
                    \qquad (1\le R\le Q),
\end{equation}
where the coefficients satisfy their stated bound for arbitrarily
small exponents and constants can also depend on \(\mathcal R\).
Indeed write \(n=h^2m\) with \(m\) squarefree. In the \(h\)-term
the additional factor \(\mathbf1_{(h,d)=1}\) can be dropped in the
upper bound for its norm. Minkowski's inequality and the sieve
bound the total, at the cost of \(D^\eta\), by
\[
 \sum_{h\le\sqrt Q}h^{-2\sigma_0}
    \left(\sum_{m>R/h^2}m^{-2\sigma_0}\right)^{1/2}
 \ll_{\sigma_0} R^{1/2-\sigma_0}\log(2R).
\]
For \(h\le\sqrt R\) the summand is \(O(R^{1/2-\sigma_0}/h)\);
the remaining tail is \(O(R^{1/2-\sigma_0})\). Use the individual
coefficient and sieve estimates with a smaller exponent to absorb
the logarithm.

For \(L_d=L(s,\pi_E\otimes\chi_d)\) one also has
\begin{equation}\label{rfc:central}
          \|L_d(1/2+it)\|_D\ll_{E,\eta}D^\eta(2+|t|)^{C_E}.
\end{equation}
This is the quadratic large-sieve second-moment bound. We indicate
the fixed-level details. For \(d\) coprime to \(2N\) the primitive
twist has conductor \(Nd^2\), the same weight-two gamma factor, and
coefficients \(\lambda(n)\chi_d(n)\) with
\(|\lambda(n)|\le \tau(n)\), the divisor function. The usual two
sums in the approximate functional equation have length
\(\ll_E D(2+|t|)\) with rapidly decreasing smooth weights.
The common weights at scale \(d\) can be written by Mellin inversion
at arbitrarily small fixed positive real part against
\((d/n)^{iv}\), with an arbitrarily fast decaying tail in \(v\)
and costs polynomial in \(2+|t|\). Equivalently one may use partial
summation and a maximal sieve bound on each dyadic block, obtained
from the sieve on its binary subintervals at a cost of
\(O(\log^2(2Q))\). For \(n\le Q\), the \(h^2m\) decomposition above
at \(\Re s=1/2\) bounds the norm by
\((1+Q/D)^{1/2}(DQ)^\eta\), since the two harmonic sums cost
only logarithms. Dyadic blocks beyond \(D(2+|t|)\) can be summed
with arbitrary inverse powers of their relative length, and the
twist root number has modulus one. This gives \eqref{rfc:central}.
For the approximate functional equation and the same large-sieve
argument see \href{https://doi.org/10.4171/JEMS/224}{Soundararajan
and Young, \emph{The second moment of quadratic twists of modular
\(L\)-functions}, Lemma~2.1 and Corollary~2.5 in their full-level
formulation}.

Set \(X=D^{1/4}\), and write
\[
 S_X(s,d)=\sum_n\lambda(n)\chi_d(n)e^{-n/X}n^{-s},\qquad
 J_R(s,d)=\sum_{n\le R}\nu(n)\chi_d(n)n^{-s},
\]
where \(\nu\) are the inverse Dirichlet coefficients of \(L(s,\pi_E)\).
The reciprocal local polynomials give \(|\nu(n)|\le\tau(n)\).
For \(s\in\mathcal R\), Mellin inversion and shifting the gamma
integral to \(\Re z=1/2-\Re s\) cross only the pole at \(z=0\).
By \eqref{rfc:central}, Minkowski and the exponential decay of
\(\Gamma(z)\) there,
\[
       \|S_X(s,d)-L_d(s)\|_D\ll_{E,\mathcal R,\eta}
                               D^\eta X^{1/2-\sigma_0}.
\]
The trivial bound on \(J_R\) is \(O_\eta(R^{1-\sigma_0+\eta})\).
In \(S_X\) truncate at \(D^{1/3}\), with an error smaller than any
inverse power of \(D\). In the product \(S_XJ_R\), the coefficient
of \(\chi_d(n)n^{-s}\) for \(n\le R\), after subtracting the
constant \(1\), is \(O_\eta(n^{1+\eta}/X)\). This follows from
\(\lambda*\nu=\mathbf1_{n=1}\) and
\(|1-e^{-b/X}|\le b/X\) for each divisor \(b\mid n\).
Above \(R\) the coefficients are \(O_\eta(n^\eta)\).
If \(R=D^r\) with \(0<r<1/12\), the resulting length is \(<D\),
and \eqref{rfc:tail} applies. Uniformly on \(\mathcal R\) we
therefore obtain
\[
 \|L_d(s)J_R(s,d)-1\|_D\ll_{E,\mathcal R,\eta}D^\eta
 \bigl(X^{1/2-\sigma_0}R^{1-\sigma_0}
          +R^{2-\sigma_0}/X+R^{1/2-\sigma_0}\bigr).
\]
Choose \(0<r<\min(1/12,(\sigma_0-1/2)/8)\), then \(\eta>0\)
small enough. Each term saves a fixed power of \(D\). If \(L_d\)
has a zero in \(K_0\), the holomorphic function \(L_dJ_R-1\)
equals \(-1\) there; its squared absolute value consequently has
integral at least a positive fixed constant over \(\mathcal R\)
by the disk submean inequality. Integrate the preceding squared
norm bound over the rectangle to count such \(d\).
\end{proof}

\begin{corollary}\label{rfc:detection}
Let \(E/\Q\) have an off-line unitary zero \(\rho\) of real part
\(\beta>1/2\), and fix \(0<\epsilon<\beta-1/2\). Except for
\(O_{E,\rho,\epsilon}(D^{1-\delta})\) discriminants in
\(\mathcal D\cap(D,2D]\), its inert-prime sum
\[
            I_d(x)=\sum_{\substack{p\le x,\ p\nmid N d\\\chi_d(p)=-1}}a_p
\]
has arbitrarily large negative excursions exceeding \(x^{1+\epsilon}\)
in size in every sufficiently large fixed-ratio interval of some
width depending on \(E,d,\rho,\epsilon\). The zero \(\rho\) itself
has inert-class residue \(m_E(\rho)/2>0\) for these discriminants.

In particular to deduce ERH for a fixed \(E\), it suffices, for every
\(\epsilon>0\), to have a set of positive upper relative density
among \(\mathcal D\), or among its prime members, for which each
\(I_d(x)\ge-C_{d,\epsilon}x^{1+\epsilon}\) throughout arbitrarily
distant intervals \([Y,YH_d(Y)]\), with \(H_d(Y)\to\infty\).
The set and the intervals may depend on \(\epsilon\), and the
intervals for distinct \(d\) may be disjoint.
\end{corollary}
\begin{proof}
Choose \(1/2<b<\beta\) and apply Lemma~\ref{rfc:compact} with
\(K_0=\{\rho\}\cup[b,1]\). For each remaining \(d\),
\[
          r_{-}(\rho)=\tfrac12(m_E(\rho)-m_{E_d}(\rho))
                        =\tfrac12m_E(\rho)>0 .
\]
Every effective real zero with real part at least \(b\) has positive
residue as well. Thus \(\Theta_-\ge\beta\), and in both possible
real-edge cases Theorem~\ref{rfc:class} supplies negative excursions
with any exponent between \(1+\epsilon\) and \(\beta+1/2\).
This strict gap also absorbs any fixed \(C_{d,\epsilon}\).
Summing the exceptional estimate over dyadic intervals gives a
power-saving count up to \(D\). There are \(\gg_E D\) discriminants
in \(\mathcal D\) by elementary squarefree counting, and
\(\asymp D/\log D\) prime members by the prime number theorem modulo
four. The exceptional set has zero relative upper density in both.
\end{proof}

\subsection{An auxiliary prime discriminant that can change at each cutoff}

\begin{theorem}\label{rfc:moving}
Fix \(H(Y)\to\infty\). For \(x\ge2\) use only prime discriminants
\(d\equiv1\bmod4\) in \(x<d\le2x\), prime to \(N\).
Each of the following is equivalent to ERH for the fixed \(E/\Q\).
In each assertion, for every \(0<\epsilon<1/4\) there are a constant
\(C_{E,\epsilon}\) and arbitrarily large \(Y\) such that at \emph{each}
\(x\in[Y,YH(Y)]\) at least \(x^{1-\epsilon}\) such prime
discriminants satisfy the stated inequality:
\begin{enumerate}
\item \(I_d(x)\ge-C_{E,\epsilon}x^{1+\epsilon}\).
\item
\(\displaystyle
 \log\prod_{\substack{p\le x,\ p\nmid N\\\chi_d(p)=-1}}
         \frac{\#E(\F_p)}p \ \le C_{E,\epsilon}x^\epsilon\).
\end{enumerate}
The discriminants are allowed to differ for every real \(x\).
Under ERH each inequality in fact holds for all but
\(O_{E,\epsilon}(x^{1-\epsilon})\) of the candidates at that cutoff.
\end{theorem}
\begin{proof}
Write \(S(x)=\sum_{p\le x,\ p\nmid N}a_p\) and
\(Q_d(x)=\sum_{p\le x,\ p\nmid N}a_p\chi_d(p)\).
Since \(d>x\), no prime in either sum is ramified in the quadratic
field and \(2I_d(x)=S(x)-Q_d(x)\) exactly.
For arbitrary coefficients \(u_p\) supported on good primes up to
\(x\), the quadratic large sieve gives
\[
  \sum_{x<d\le2x}^{*}
       \left|\sum_p u_p\chi_d(p)\right|^2
            \ll_\eta x^{1+\eta}\sum_p|u_p|^2 .
\]
The star may include all \(d\in\mathcal D\), hence in particular
all candidate primes; separate \(p=2\) and use Cauchy--Schwarz.
Hasse gives \(\sum a_p^2\ll x^2\), so
\(\#\{d:|Q_d(x)|>x^{1+2\epsilon}\}
\ll_\epsilon x^{1-3\epsilon}\). Likewise, put
\(v_p=\log(\#E(\F_p)/p)\), so \(|v_p|\ll_E p^{-1/2}\).
Then \(\sum v_p^2\ll_E\log\log(3x)\), and
\(\#\{d:|\sum v_p\chi_d(p)|>x^{2\epsilon}\}
\ll_{E,\epsilon}x^{1-3\epsilon}\), using a smaller \(\eta\).
The exact analogous identity is
\[
 2\sum_{\chi_d(p)=-1}v_p=\log P_E(x)-\sum_p v_p\chi_d(p),
\]
with the same prime cutoffs and omissions on both sides.
If at least \(x^{1-\epsilon}\) candidates obey the assumed bound,
at least one is outside the corresponding exceptional set. At each
\(x\) this gives \(S(x)\gg_{E,\epsilon}-x^{1+2\epsilon}\) or
\(\log P_E(x)\ll_{E,\epsilon}x^{2\epsilon}\), respectively.
Use Theorem~\ref{rawfilter:main} or
Corollary~\ref{rawfilter:pointproduct} with \(\epsilon\) arbitrarily
small. For the forward implication use their global ERH estimates
and apply the same sieve with thresholds \(x^{1+\epsilon}\) and
\(x^\epsilon\). The resulting exceptional sets have size
\(O_{E,\epsilon}(x^{1-\epsilon})\). By the prime number theorem for
\(1\bmod4\) the candidate list has \(\asymp x/\log x\) members.
This proves every assertion.
\end{proof}

\clearpage
\section{An inverse elliptic function from quadratic Dirichlet characters}\label{qcm:section}

Fix an integral equation \(E:y^2=f(t)=t^3+At+B\) with nonzero
discriminant. Let \(S\) contain the primes of bad reduction of this
equation and the primes at most \(7\). Throughout this section a
superscript \(S\) removes these Euler factors. For \(f(n)\ne0\) set
\[
     \xi_n(p)=\left(\frac{f(n)}p\right)\quad(p\notin S),
 \qquad
     Q_n(s)=\prod_{p\notin S}(1+\xi_n(p)p^{-s})
                     \qquad(\Re s>1).
\]
The finite-place character is imprimitive when the polynomial value
has square factors. Explicitly take the primitive quadratic
character of \(\Q(\sqrt{f(n)})\) and delete the Euler factors at
\emph{all} primes dividing \(f(n)\). For square values of \(f(n)\)
the primitive character would be trivial. We omit those \(n\) and
the roots of \(f\). This omits finitely many integers by
\href{https://doi.org/10.1007/978-0-387-09494-6_IX}{Siegel's
integral-points theorem, Silverman, Chapter~IX, Theorem~3.1}.
The prime on \(\sum_n'\) indicates these fixed omissions.

\begin{theorem}\label{qcm:inverse}
The limit of \(N^{-1}\sum_{N<n\le2N}'Q_n(s)\) on \(\Re s>1\) has the
meromorphic continuation to \(\Re s>0\)
\[
       F_E(s)=\frac{C_E(s)}{L^S(E,s+1)},\qquad
       C_E(s)=\prod_{p\notin S}
       \frac{1-a_p p^{-1-s}}{1-a_p p^{-1-s}+p^{-1-2s}} .
\]
The factor \(C_E\) is holomorphic and never zero there. Consequently
ERH for the fixed curve is equivalent to holomorphy of this continued
quadratic-character mean on \(\Re s>0\). Assuming the Dirichlet
critical-line theorem, the finite means themselves converge locally
uniformly to \(F_E\) throughout \(\Re s>1/2\).
\end{theorem}
\begin{proof}
For each \(p\notin S\), counting the affine points and the one point
at infinity gives \(\sum_{u\bmod p}(f(u)/p)=-a_p\).
For a squarefree \(m\) prime to \(S\), the residue-row coefficient
\(\xi_n(m)=\prod_{p\mid m}\xi_n(p)\) is periodic in \(n\) modulo
\(m\). The Chinese remainder theorem gives its average over a
complete period as \(\prod_{p\mid m}(-a_p/p)\).
For \(\Re s>1\) the coefficients of \(Q_n\) are bounded by one;
absolute convergence therefore allows passage to the mean limit.
Its series is the squarefree coefficient series of
\(\prod_{p\notin S}(1-a_p p^{-1-s})\).
This latter product is already absolutely convergent for \(\Re s>1/2\)
by Hasse. In arithmetic normalization the reciprocal Euler factor
of \(L(E,s+1)\) is \(1-a_pp^{-1-s}+p^{-1-2s}\), which gives the
displayed identity. For \(\Re s>0\), neither the numerator nor
the denominator of a local correction vanishes: use
\(|a_p|/p\le2/\sqrt p<1\) and the two Frobenius eigenvalues of
modulus \(\sqrt p\). Uniformly on each \(\Re s\ge\sigma_0>0\) the
factor differs from one by \(O_{\sigma_0}(p^{-1-2\sigma_0})\).
The infinite product has the claimed properties. The finitely
removed elliptic Euler factors have no divisor at \(\Re(s+1)>1\).
Modularity and the functional equation now give the equivalence.

For the assertion about actual means write \(L^S(s,\xi_n)\) for
the primitive quadratic Dirichlet function with the described
additional deletions. Local multiplication gives
\begin{equation}\label{qcm:dir}
 Q_n(s)=\frac{L^S(s,\xi_n)}{\zeta^S(2s)}
           \prod_{\substack{p\mid f(n)\\p\notin S}}(1-p^{-2s})^{-1}.
\end{equation}
In particular \(Q_n\) is holomorphic on \(\Re s>1/2\).
On Dirichlet GRH, for a primitive nonprincipal character of
conductor \(q\), every \(\delta,\eta>0\) gives
\[
 |L(\sigma+it,\chi)|+|L(\sigma+it,\chi)|^{-1}
       \ll_{\delta,\eta} (q(2+|t|))^\eta
                         \qquad(\sigma\ge1/2+\delta).
\]
For the standard short derivation, start the zero-free logarithm
at \(2+it\). Convexity bounds its real part by
\(O_\delta(\log(q(2+|t|)))\) on the concentric disk of radius
\(3/2-\delta/8\), taking \(0<\delta<1/2\). Borel--Carathéodory
bounds its absolute value by the same order on radius
\(3/2-\delta/4\). On radius \(1/2\) the logarithm is bounded
absolutely by the Euler product. Hadamard's three-circles theorem
at radius \(3/2-\delta\) bounds its modulus by a fixed power
strictly less than one of \(\log(3q(2+|t|))\). Exponentiation and
absolute convergence further to the right give the estimate.
The conductor is at most \(4|f(n)|\), and the finite imprimitive
Euler products in \eqref{qcm:dir} contribute at most
\(O_{\delta,\eta}((2+|f(n)|)^\eta)\). This follows by comparing
each sufficiently large local factor with \(p^\eta\) and treating
the finitely many smaller primes separately.

Fix a compact subset of \(1/2<\Re s\le 1\) and take
\(0<\delta<\inf(\Re s-1/2)\). In the squarefree Dirichlet series
for \(Q_n(s)\), introduce the weight \(e^{-m/M}\), \(M=N^{1/4}\).
Mellin inversion, followed by a shift to
\(\Re(s+z)=1/2+\delta/2\), gives an error
\(O_{\delta,\eta}(N^\eta M^{-\delta/2})\), uniformly for
\(N<n\le2N\) and \(s\) in the compact set. Indeed the only
crossed pole of \(\Gamma(z)\) is at zero, and \eqref{qcm:dir}
and the preceding bounds control the vertical integral. The same
calculation applies across \(\Re s=1\) on compact subrectangles
whose right edge is below \(5/4\). Farther right the original
dominated convergence is already locally uniform.

In a dyadic integer interval of length \(N\), the error from
replacing the average of the individual \(m\)-th coefficient by
its complete-period mean is \(O(m/N)\), uniformly in \(m\).
Thus the smoothed series incurs an error
\(\ll N^{-1}\sum_{m\ge1}m^{1-\Re s}e^{-m/M}=o(1)\).
For large \(N\) there are no exceptional polynomial values in the
interval. The limiting smoothed series converges to the absolutely
convergent Euler product for \(F_E(s)\), locally uniformly on
\(\Re s>1/2\). Taking \(\eta\) small proves the claim.
\end{proof}

Thus the actual finite means already reproduce the full analytic
function in the half-plane of absolute convergence of its
elliptic-side product. On the other side of \(\Re s=1/2\) their
positive-conductor growth creates a specific obstruction.

\begin{proposition}\label{qcm:growth}
Assume Dirichlet GRH. For every fixed \(0<\sigma<1/2\) and \(\eta>0\),
\[
       \frac1N\sum_{N<n\le2N}' Q_n(\sigma)
                      \le -c_{E,\sigma,\eta}\,
                                     N^{3(1/2-\sigma)-\eta}
\]
for all sufficiently large \(N\).
\end{proposition}
\begin{proof}
The alternating eta series gives \(\zeta(2\sigma)<0\). All the
additional real Euler factors in \eqref{qcm:dir} are positive, and
the real primitive quadratic \(L(\sigma,\chi)\) is positive. Indeed
its sign is positive to the right of one by the Euler product,
remains positive on the interior real interval \((1/2,1]\) by
Dirichlet GRH and boundary nonvanishing, and on \((0,1/2)\) follows
from its functional equation with root number \(+1\) and positive
real gamma factors. Thus all terms in the mean are negative.

We need many polynomial values whose primitive conductor has full
cubic size. The following use of a prime-input squarefree sieve
suffices. For each \(\ell=2,3\) there is a residue \(r_\ell\bmod
\ell^2\) with \(\ell^2\nmid f(r_\ell)\), since the third finite
difference of the monic cubic is \(6\), divisible by neither \(4\)
nor \(9\). Choose \(r\bmod36\) satisfying these conditions. Put
\(g(u)=f(36u+r)\). It is a cubic with distinct roots. For every
prime \(\ell\ge5\) at most three residues modulo \(\ell\) are
roots, so some unit residue is a nonroot. At \(2\) and \(3\) the
progression already rules out square divisibility. All local
factors in
\href{https://doi.org/10.1007/s11511-014-0117-2}{Helfgott's
squarefree-values theorem for cubic polynomials at prime arguments}
are therefore positive. For primes outside the fixed discriminant
the excluded proportion modulo \(\ell^2\) is \(O(1/\ell^2)\)
by simple-root lifting, so their product is positive. Helfgott's
asymptotic supplies \(\gg_E N/\log N\) indices \(n=36p+r\)
in \((N,2N]\) with \(p\) prime and \(f(n)\) squarefree.

For these indices \(f(n)>1\), the conductor \(q\) is \(f(n)\) or
\(4f(n)\), and \(q\asymp_E N^3\). There are no extra odd
imprimitive deletions outside \(S\). The primitive functional
equation gives
\[
       L(\sigma,\chi)=c_\sigma q^{1/2-\sigma}
                                    L(1-\sigma,\chi),
       \qquad
       c_\sigma=\pi^{\sigma-1/2}
                    \frac{\Gamma((1-\sigma)/2)}{\Gamma(\sigma/2)}>0
\]
for these positive discriminants. The reciprocal bound in the
preceding proof yields \(L(1-\sigma,\chi)\gg_{\sigma,\epsilon}q^{-\epsilon}\).
The primes of the fixed set \(S\) cost at most a fixed positive
factor, and the product over \(p\mid f(n)\) in \eqref{qcm:dir}
has magnitude at least one. Sum over these indices, use the common
sign on the remainder, and absorb the logarithm by decreasing
\(\epsilon\).
\end{proof}

\subsection{A polynomial-conductor Dirichlet statistic for the prime window}
There is also a finite real formulation that avoids continuation of
any averaged functions. Write \(H(Y)\to\infty\) and let \(x\ge3\).
Take \(M=\lceil x\rceil\), \(k=\lceil2\log_2 x\rceil\), and let
\(n=U_1+\cdots+U_k\) where the \(U_j\) are independent uniform
integers in \(\{1,\ldots,M\}\). Define \(\mathbb E_x\) using these
explicit nonnegative rational weights. For every sufficiently large
\(x\), all sampled polynomial values are nonzero and nonsquare since
\(n\ge k\to\infty\). The conductors are at most
\(O_E(x^3\log^3(2x))\).

\begin{theorem}\label{qcm:window}
For this family one has, uniformly for all sufficiently large \(x\),
\[
 \mathbb E_x \sum_{\substack{p\le x\\p\notin S}}\xi_n(p)
     =-\sum_{\substack{p\le x\\p\notin S}}\frac{a_p}{p}+O(x^{-1})
     =\log P_E(x)+O_E(\log\log(3x)).
\]
In particular ERH for \(E\) is equivalent to an upper bound
\(\mathbb E_x\sum_{p\le x,\ p\notin S}\xi_n(p)\le C_{E,\epsilon}x^\epsilon\)
throughout arbitrarily distant \([Y,YH(Y)]\), for each \(\epsilon>0\).
\end{theorem}
\begin{proof}
For \(1\le h<p\) put \(\|h\|_p=\min(h,p-h)\). A finite geometric
series and \(\sin(\pi\|h\|_p/p)\ge2\|h\|_p/p\) give
\[
 \left|\mathbb E_x e^{2\pi ihn/p}\right|
       \le\left(\frac{p}{2M\|h\|_p}\right)^k
       \le(2\|h\|_p)^{-k}.
\]
The normalized finite Fourier coefficients of \(u\mapsto(f(u)/p)\)
have modulus at most one. Its mean against the sample therefore
differs from its complete mean by at most \(2^{1-k}\zeta(k)\ll x^{-2}\).
Sum over primes \(p\le x\), and apply
Corollary~\ref{rawfilter:pointproduct}. The finitely many extra good
primes in \(S\) affect only its bounded error.
\end{proof}

On Dirichlet GRH an individual character prime sum here is
\(O_E(\sqrt x\log^2(2x))\), by the usual sharp explicit formula and
partial summation, with at most \(O_E(\log(2x))\) from the omitted
ramified Euler factors. The required gain in the average is the
remaining fixed-form cancellation. The exact identity identifies a
specific family to which a relative trace or conductor moment
estimate could be applied.

This gain cannot be obtained by replacing the cubic with a fixed
collection of other nonsquare polynomials sampled uniformly at each
prime. For squarefree \(F_j\bmod p\) of positive degree, the
multiplicative
\href{https://doi.org/10.1017/CBO9780511525926}{Weil bound,
Lidl and Niederreiter, \emph{Finite Fields}, Theorem~5.41}, gives
\[
 \left|\sum_j c_j\,\frac1p\sum_{u\bmod p}\left(\frac{F_j(u)}p\right)\right|
             \le p^{-1/2}\sum_j|c_j|(\deg F_j-1).
\]
For all sufficiently large \(X\), some good \(p\in(X,2X]\) has
\(|\lambda_p|>1/2\). This follows, for example, by differencing
the fixed-curve Rankin--Selberg asymptotic for the second moment
at \(2X\) and at \(X\). Thus approximation of \(\lambda_p\)
to within \(1/4\) at every such prime, using polynomials satisfying
the stated smoothness there, requires
\(\sum_j|c_j|(\deg F_j-1)\ge\sqrt X/4\).
The single cubic recovers the trace exactly with coefficient
\(-\sqrt p\). The half-power cost persists for complete uniform
polynomial sampling.

\subsection{The exact saving required of a centered Dirichlet moment}
For the following test take \(E:y^2=x^3-x\), \(f(n)=n^3-n\).
Let \(S_0=\{2,3,5,7,11\}\). Fix the progression
\(n\equiv2\bmod4\), \(n\equiv p\bmod p^2\) for each odd \(p\in S_0\),
of modulus \(B=4\prod_{p\in S_0,\ p>2}p^2\).
Write \(\mathcal A\) for its positive members with \(f(n)\) squarefree.
For large \(n\in\mathcal A\) the primitive character has positive
discriminant \(q_n=4f(n)\), ramified at every member of \(S_0\).
For a smooth weight \(V\) supported in \((1,2)\), put
\[
 A_N(w;V)=\frac1N\sum_{n\in\mathcal A}V(n/N)L(w,\chi_{q_n}),
 \qquad W_z(t)=W(t)t^z,
\]
where \(W\ge0\) is one fixed smooth weight of integral one with this
support. Use the full primitive Dirichlet functions in \(A_N\).

\begin{theorem}\label{qcm:centered}
Assume Dirichlet GRH, fix \(0<u\le1/4\), and let \(w=1/2+u+it\).
Suppose that for some fixed \(C\), uniformly in \(t\in\R\) and \(N\ge2\),
\begin{equation}\label{qcm:target}
 |A_{2N}(w;W_{3(w-1/2)})-A_N(w;W_{3(w-1/2)})|
                        \ll_{u,W} (2+|t|)^C N^{-3u}.
\end{equation}
Then \(L(E,s+1/2)\) has no zero in \(\Re s>1-u\).
The power in \eqref{qcm:target} is sharp for such a uniform
bound: for each fixed \(C\ge0\) and all sufficiently large \(N\),
\[
 \sup_{t\in\R}\frac{|A_{2N}(w;W_{3(w-1/2)})-A_N(w;W_{3(w-1/2)})|}
                         {(2+|t|)^C}\gg_{u,W,C}N^{-3u}.
\]
\end{theorem}
\begin{proof}
First compute the mean for \(\Re s>1\). At \(p\notin S_0\) the
three roots of \(f\bmod p\) are simple. The density of
\(p^2\nmid f(n)\) is \(\Delta_p=1-3/p^2\). Conditional on this
event, the expectation of \(\chi_{q_n}(p)\) is
\(b_p=-a_p/(p\Delta_p)\), since the removed residues are roots and
have character zero already. Its square has mean
\(v_p=(1-3/p)/\Delta_p\). The conditions at distinct primes give
independent local means by the Chinese remainder theorem. Passing
to the infinite squarefree condition is valid: outside \(S_0\),
if \(p^2\mid f(n)\), then \(p^2\) divides one of \(n-1,n,n+1\).
The number of \(n\in[N,2N]\) excluded by some \(p>Z\) is therefore
\(O(N/Z+\sqrt N)\) by summing \(O(N/p^2+1)\) to \(p\le\sqrt{2N+1}\).
Consequently, setting \(d_0=B^{-1}\prod_{p\notin S_0}\Delta_p>0\),
one obtains by dominated convergence
\[
  \lim_{N\to\infty}A_N(s;W)
       =d_0 Z(s),\qquad
  Z(s)=\prod_{p\notin S_0}
        \frac{1+b_p p^{-s}+(v_p-1)p^{-2s}}{1-p^{-2s}}.
\]
This is the average of the ordinary primitive \(L\), which accounts
for the even powers in the denominator. The same formula holds with
\(W\) replaced by \(V\) and multiplied by \(\int V\).

In \(\Re s>0\) the continuation has the form
\[
       Z(s)=\frac{\zeta^{S_0}(2s)D(s)}{L^{S_0}(E,s+1)}
\]
with \(D\) holomorphic and nonzero. To check this, divide each
local numerator by \(1-a_p p^{-s-1}+p^{-1-2s}\). The difference
in the linear coefficients is \(O(p^{-5/2})\), while that in the
quadratic coefficients is \(O(1/p)\). The product of ratios thus
converges normally. Its numerator cannot vanish since, for
\(p\ge13\),
\(|b_p|+|v_p-1|\le(2/\sqrt p+3/p)/(1-3/p^2)<1\).

Set \(v(s)=3(1/2-s)\) and \(T_N(s)=A_{2N}(s;W)-2^{v(s)}A_N(s;W)\).
These are entire functions. For \(\Re s>1\) the limit is
\((1-2^{v(s)})d_0Z(s)\). At \(s=1-w\) the primitive even functional
equation gives the conductor factor \(q_n^{w-1/2}\) and multiplier
\(c(s)=\pi^{s-1/2}\Gamma((1-s)/2)/\Gamma(s/2)\).
Since \(q_n=4n^3(1-n^{-2})\), uniformly for real \(t\)
\[
 \begin{split}
 T_N(1-w)={}&c(1-w)4^{w-1/2}(2N)^{3(w-1/2)}
       \big[A_{2N}(w;W_{3(w-1/2)})-
                      A_N(w;W_{3(w-1/2)})\big]\\
       &+O_{u,W,\eta}\big((2+|t|)^{C_1} N^{3u-2+\eta}\big).
 \end{split}
\]
Indeed \((1-n^{-2})^{w-1/2}=1+O_u((1+|t|)n^{-2})\), Dirichlet GRH
gives \(|L(w,\chi_{q_n})|\ll_\eta N^\eta(2+|t|)^\eta\), and the
gamma ratio has polynomial growth in \(t\) by Stirling. This holds
for any \(\eta>0\). It yields a bound on \(T_N\) polynomial in
height and independent of \(N\) on the left boundary of
\(1/2-u\le\Re s\le2\), under \eqref{qcm:target}.
On the right boundary there is an absolute bound by the Euler
product. Divide by a sufficiently large fixed power of \(3+s\);
the three-lines theorem then gives local boundedness uniformly in
\(N\) on the interior. For each fixed \(N\) the required growth
condition for three-lines follows from the finite sum of primitive
Dirichlet functions and their standard vertical-strip bounds.
By Montel's theorem and uniqueness on \(\Re s>1\), the limit there
continues holomorphically to the whole interior.

In \(1/2-u<\Re s<1/2\), neither \(1-2^{v(s)}\) nor
\(\zeta^{S_0}(2s)\) vanishes, the latter by Dirichlet GRH because
\(\Re(2s)>1/2\). Thus \(L(E,s+1)\) has no zero there. The rest of
the claimed region follows from absolute convergence and the
known boundary nonvanishing in unitary normalization.

For sharpness the right-hand limit at the real point \(s=3/2\)
is nonzero: \(W\) has integral one, \(d_0>0\), and every quadratic
\(L\) in the absolutely convergent calculation has positive
local means. For any fixed sufficiently large \(m\), the supremum
of \(|T_N(s)/(3+s)^m|\) on the left boundary is consequently bounded
below by a positive constant for all large \(N\), by three-lines
and the uniform upper bound on the right boundary. Given \(C\)
choose \(m\) large relative to \(C,C_1\) and the growth of
\(c(1-w)\). The displayed functional-equation identity then yields
the stated lower bound, since \(3u-2<0\).
\end{proof}

The elementary moment estimate available here has a different
exponent. For every \(0<u\le1/4\) and \(\eta>0\), the left side
of \eqref{qcm:target} is
\begin{equation}\label{qcm:provedmoment}
             O_{u,W,\eta}\big((2+|t|)^{C_\eta}N^{-u/2+\eta}\big).
\end{equation}
Here are the details of the squarefree sieve and of the truncation.
Take \(L=N^{1/4-\varepsilon}\) and truncate
\(\mu^2(f(n))=\sum_{e^2\mid f(n)}\mu(e)\) at \(e\le L\), within the
fixed progression. The mean absolute error is
\(O_\varepsilon(N^\varepsilon(L^{-1}+N^{-1/4}))\), with an
arbitrarily small power loss. To see this, if some \(n+i\),
\(i=-1,0,1\), has a square divisor \(h^2\) with \(h>N^{1/4}\),
there are \(O(N^{3/4}+\sqrt N)\) possible \(n\), and
\(O_\varepsilon(N^\varepsilon)\) divisor terms for each.
For all remaining \(n\), the three integers are pairwise coprime
on this progression, so any squarefree \(e\) in the omitted
tail is at most \(N^{3/4}\). For fixed \(e\) prime to \(S_0\)
there are at most \(3^{\omega(e)}\) permitted classes modulo
\(e^2\), by simple-root lifting. The total further count is
\(\ll_\varepsilon N^\varepsilon(N/L+N^{3/4})\). A divisor \(e\)
meeting \(S_0\) gives no terms. This proves the error bound.
For the limiting density the corresponding tail is
\(\ll_\varepsilon L^{-1+\varepsilon}\), by the absolutely convergent
sum over \(3^{\omega(e)}/e^2\).

For each Dirichlet series index \(m\), \(\chi_{q_n}(m)\) on the
squarefree family equals \((f(n)/m)\) at odd \(m\) and is zero
at even \(m\). Use that same periodic sequence also inside the
truncated sieve; when \(m\) meets \(S_0\) the coefficient vanishes
identically. Each truncated condition now has period at most
\(Bm e^2\). For any periodic coefficients bounded by one of period
\(P\le N\), Poisson summation for a smooth weight \(V(n/N)\)
replaces the normalized sum by its complete mean times \(\int V\)
with error \(O_{J,V}((P/N)^J)\) for any \(J>1\). This follows
by finite Fourier expansion, bounding its normalized coefficients
by one and using the rapid decay of \(\widehat V\).
For \(m\le N^{1/2-\varepsilon}\) and \(e\le L\), use \(J\)
sufficiently large to absorb the number of sieve terms. For
\(V=W_{3(w-1/2)}\), its finitely many needed seminorms grow only
polynomially with \(2+|t|\). We thus know each averaged coefficient
in this range with error
\(O_\varepsilon((2+|t|)^{C_\varepsilon}N^{-1/4+O(\varepsilon)})\).

Approximate the individual \(L(w,\chi_{q_n})\) by its series with
weight \(e^{-m/M}\), \(M=N^{1/2-2\varepsilon}\). The Mellin shift
under Dirichlet GRH used above gives error
\(O_{u,\varepsilon}((2+|t|)^{C_\varepsilon}N^\varepsilon
M^{-u+\varepsilon})\); the series beyond \(N^{1/2-\varepsilon}\)
is negligible by exponential decay. Summing the coefficient errors
costs \(N^{-1/4+O(\varepsilon)}M^{1/2-u}\) times the same polynomial
in height. The limiting coefficients of \(Z\) converge absolutely
with the weight \(m^{-1/2-\varepsilon}\), since \(b_p=O(p^{-1/2})\)
and the coefficients at higher powers have absolute value at most
one. Removing the smooth weight from this limiting series costs
\(O_{u,\varepsilon}(M^{-u+\varepsilon})\). The main term
\(d_0 Z(w)\int W_{3(w-1/2)}\) is independent of \(N\) and cancels
between scales. Choosing \(\varepsilon\) sufficiently small in
terms of \(u,\eta\) proves \eqref{qcm:provedmoment}.
Thus this calculation gives exponent \(u/2\), while the
fixed-form zero-free implication needs \(3u\). A uniform estimate
at \(u>1/8\) in \eqref{qcm:target} would improve the \(7/8\)
all-height CM boundary from Section~\ref{recentcm:section}.

\subsection{An Eisenstein coefficient at every cubic value}
There is a version for every rational elliptic curve that also gives
a larger proved saving. It uses the local factors in the Fourier
coefficients of a half-integral weight Eisenstein series. These
factors give a functional equation with the full polynomial value
in its conductor and preserve periodic coefficients without
discarding nonsquarefree values. The correction itself and its
functional equation are classical. See
\href{https://chinta.ccny.cuny.edu/publ/cfh2.pdf}{Chinta, Friedberg
and Hoffstein, \emph{Multiple Dirichlet series and automorphic
forms}, \S3.4.1, equations (3.15)--(3.17)}.

Return to the smooth integral cubic \(f\) of
Theorem~\ref{qcm:inverse}. Fix the set \(S\) from that theorem.
On one fixed arithmetic progression \(n\equiv b\bmod B\),
arrange that \(p^2\nmid f(n)\) for \(p\in S\), that \(v_p(f(n))\)
is fixed there, and that the local quadratic characters at those
primes are fixed. For odd \(p\) this uses a choice modulo \(p^2\)
with \(p^2\nmid f(b)\); such a choice always exists, since a fixed
square divisor would divide the third finite difference \(6\).
At \(2\) choose a class modulo \(8\) with \(4\nmid f(b)\). If its
value is odd then its unit residue modulo \(8\) is fixed. If it is
even then the quadratic discriminant will be even. The Chinese
remainder theorem gives all conditions at once.

For all sufficiently large members of this progression, write
\[
      f(n)=d_0e^2,\qquad d_0>1\text{ squarefree},\qquad q_0=c\,d_0
\]
for the primitive quadratic conductor, with \(c=1\) or \(4\)
fixed by the progression. In particular \((e,\prod_{p\in S}p)=1\).
For \(p^r\Vert e\) define
\[
 P_{p,r}(s)=\sum_{j=0}^r(p^{1-2s})^j
          -\chi_{q_0}(p)p^{-s}\sum_{j=0}^{r-1}(p^{1-2s})^j,\qquad
       D_n(s)= L(s,\chi_{q_0})\prod_{p^r\Vert e}P_{p,r}(s).
\]
The finite omitted exceptional \(n\) are treated as in
Theorem~\ref{qcm:inverse}. On writing \(\widetilde A_N(s;V)\) for
the smoothed \(N^{-1}\) average of these \(D_n(s)\) over the
progression, the following version of the centered criterion holds.

\begin{theorem}\label{qcm:eisenstein}
For every fixed \(E/\Q\), Theorem~\ref{qcm:centered}, including its
sharpness assertion, holds with \(A_N\) replaced by
\(\widetilde A_N\), using the indicated progression.
For this version the established bound on the two-scale difference is
\[
       O_{E,u,W,\eta}\big((2+|t|)^{C_\eta}N^{-u+\eta}\big)
               \qquad (0<u\le1/4,\ \eta>0).
\]
In both versions the sufficient estimate \eqref{qcm:target} need
only hold at \(|t|\le 10\log\log(10N)\), with the same polynomial
weight in \(t\). The supremum in the sharpness assertion can also be
restricted to that height range.
\end{theorem}
\begin{proof}
Writing \(z=p^{-s}\) shows by reversing both finite sums that
\(P_{p,r}(s)=p^{r(1-2s)}P_{p,r}(1-s)\).
Consequently \(D_n(s)=c(s)(c f(n))^{1/2-s}D_n(1-s)\).
For every fixed \(\sigma\ge1/2+\delta\), each local polynomial has
size at most \((1+p^{-\sigma})/(1-p^{-2\delta})\). Their product is
\(O_{\delta,\eta}(f(n)^\eta)\), by treating finitely many small
primes separately. The same Dirichlet GRH bound used in the
previous contour shifts now holds for \(D_n\). The finitely many
small-prime factors of \(L(s,\chi_{q_0})\) are fixed and have no
zeros or poles in \(\Re s>0\).

We compute the full local mean outside \(S\). If \(p^r\Vert e\)
(allow \(r=0\)) and \(\chi=\chi_{q_0}(p)\), the local generating
series of \(D_n\) is
\[
        \sum_{j=0}^{r-1}p^jz^{2j}+
                      \frac{p^rz^{2r}}{1-\chi z}.
\]
Let \(v=v_p(f(n))\). For even \(k=2j\), its coefficient at \(z^k\)
is \(p^j\) if \(v\ge2j\), is \(p^h\) if \(v=2h<2j\), and is zero
in the other cases. For odd \(k=2j+1\), the coefficient is nonzero
only at \(v=2h\le2j\), when it equals
\(p^h(f(n)/p^{2h}\mid p)\), where the last symbol is Legendre's.
Thus the coefficient at \(z^k\) is determined by \(f(n)\bmod p^k\)
and has modulus at most \(p^{k/2}\). The coefficients at primes
of \(S\) are fixed. It follows that the coefficient at any integer
\(m\) is periodic in the progression parameter with period dividing
\(m\), and has absolute value at most \(\sqrt m\).

Let \(h_p\) be the number of roots of \(f\bmod p\). Simple lifting
gives \(\Pr(v\ge j)=h_p/p^j\) for \(j\ge1\). At a given even
valuation \(2h\ge2\), the unit \(f(n)/p^{2h}\bmod p\) runs equally
often through the nonzero residue classes: each lifted root is
simple and its next digit is uniform. Thus all these strata have
zero mean in an odd coefficient. The remaining stratum \(v=0\)
contributes \(-a_p/p\). For each even \(k=2j\ge2\) the mean is
\[
       \frac{h_p}{p^j}+
       1-\frac{h_p}{p}+
       h_p(1-p^{-1})\sum_{h=1}^{j-1}p^{-h}=1.
\]
It follows by independence of the complete means that the averaged
series is
\[
    Z_*(s)=\prod_{p\in S}(1-\chi_{q_0}(p)p^{-s})^{-1}
             \,\zeta^S(2s)\prod_{p\notin S}(1-a_p p^{-1-s}).
\]
The fixed small-prime values in this expression are understood.
In particular the correction leaves exactly the squarefree
numerator of Theorem~\ref{qcm:inverse} at every other prime.

On \(\Re s>3/2\) the individual coefficients are dominated by
\(\sqrt m\), so the smoothed right-hand limit is
\(B^{-1}Z_*(s)\int V\) by dominated convergence. The mean
coefficients themselves are absolutely summable with
\(m^{-1/2-\delta}\) for each \(\delta>0\), as is seen directly
from the last Euler product. For \(w=1/2+u+it\) approximate
\(D_n(w)\) by its series with \(e^{-m/M}\),
\(M=N^{1-2\varepsilon}\). The right-half-plane bound and Mellin
shift give error
\(O_{E,u,\varepsilon}((2+|t|)^{C_\varepsilon}N^\varepsilon
M^{-u+\varepsilon})\). Beyond \(m>N^{1-\varepsilon}\) the
smoothed sum is negligible even using \(\sqrt m\). Below it,
Poisson summation with period at most \(Bm\) replaces the
average of each coefficient by its complete mean with error
\(O_{J,E,V}(\sqrt m(Bm/N)^J)\) for any \(J>1\).
Take \(J\) sufficiently large in terms of \(\varepsilon\).
The cost of removing \(e^{-m/M}\) from the limiting series is
\(O_{E,u,\varepsilon}(M^{-u+\varepsilon})\).
For \(V=W_{3(w-1/2)}\) all seminorms used grow polynomially in
height. The two dyadic limits agree, and these estimates prove
the asserted \(N^{-u+\eta}\) bound.

The conductor factor in the functional equation is now
\(c f(n)=c n^3(1+O_E(n^{-2}))\). Thus the two-scale identity
in the proof of Theorem~\ref{qcm:centered} applies with \(4\)
replaced by \(c\), \(\widetilde A_N\) in place of \(A_N\), and
an error of the same order. At \(\Re s=2\) the individual bound
by \(\sum\sqrt m/m^2\) is uniform. For the sharpness argument
take the interior real point \(s=7/4\), where the limiting
\((1-2^{v(s)})Z_*(s)/B\) is nonzero. Three-lines proves both
the zero-free implication and the full-height sharpness exactly
as above.

Finally let \(a=1/2-u\), \(b'=2\), and apply the Poisson inequality
for the subharmonic function \(\log|T_N(s)/(3+s)^m|\) on the strip
\(a<\Re s<b'\), with the corresponding centered functions in
either proof. For an interior point \(s_0\), the harmonic measure
density on each boundary line at ordinate \(t\) is
\(O_{s_0}(\exp(-\pi|t-\Im s_0|/(b'-a)))\). This follows by
mapping the strip exponentially onto a half-plane and using its
Poisson kernel; the constants are locally uniform in \(s_0\).
Choose \(m\) large. The functional equation and the GRH subpower
bound on the dual show that the positive logarithm on the
uncontrolled left-boundary tail is \(O_{E,u,W}(\log(2N))\).
The total contribution above \(10\log\log(10N)\) in absolute
ordinate is therefore \(o(1)\) locally uniformly, because
\(b'-a\le7/4\). On the rest of the left boundary use
\eqref{qcm:target}, and on the right use absolute convergence.
This proves local boundedness and hence the implication. For
sharpness, if the normalized left boundary in this restricted
range were arbitrarily small along some \(N\to\infty\), the
positive fixed harmonic measure of that range, the same \(o(1)\)
tail bound, and the bounded right boundary would make the
nonzero interior limit impossible. Applying the two-scale
identity with \(m\) large relative to the specified \(C\) gives
the restricted supremum assertion.
\end{proof}

\clearpage
\section{A projective quadratic family and its exact finite Fourier transform}
\label{pc:section}

The integer representation of an elliptic curve by quadratic
characters in Theorem~\ref{qcm:eisenstein} can be enlarged by
sampling its rational abscissas. The extra variable changes both
the conductor and the attainable accuracy of the mean. Its effect
comes from an exact Fourier identity on primitive vectors modulo
each prime power.

Fix a nonsingular integral short model
\(E:y^2=x^3+Ax+B\), write
\[
 F(a,b)=b(a^3+Aab^2+Bb^3),\qquad J(h,k)=(-k,h).
\]
Choose a fixed finite set \(S\) containing \(2,3\), all primes of
bad reduction of the model, and all primes at most \(11\). Choose
an integer \(B_0\) supported on \(S\) and a primitive residue
\(v_0=(a_0,1)\bmod B_0\) such that, when \(v\equiv v_0\),
the valuations of \(F(v)\) at \(p\in S\) are fixed and at most one,
the local quadratic characters there are fixed, and the fundamental
discriminant belonging to a positive nonsquare \(F(v)=d e^2\),
\(d\) squarefree, is \(c d\) with \(c\in\{1,4\}\) fixed.
The progression of Theorem~\ref{qcm:eisenstein}, with \(b\equiv1\),
gives precisely these conditions.

Let \(W\ge0\) be a smooth function of integral one supported in
a compact subset of \(\{(a,b):1<b<2,\ F(a,b)>0\}\).
All test functions below have support in this same compact set.
For a primitive vector \(v=(a,b)\) in the progression with positive
nonsquare \(F(v)\), let
\[
 D_v(s)=L(s,\chi_{cd})\prod_{p^j\Vert e}P_{p,j}(s)
\]
use the reciprocal polynomials of Theorem~\ref{qcm:eisenstein}.
Write
\[
 A_N(s;V)=N^{-2}\!\sum_{\substack{v\in\Z^2,\ (a,b)=1,\ v\equiv v_0\ (B_0)\\
                                    F(v)>0,\ F(v)\ {\rm nonsquare}}}
                  V(v/N)D_v(s),
 \qquad W^{w}(v)=W(v)F(v)^{w-1/2}.
\]
Powers of positive numbers here and below use the real logarithm.

\begin{theorem}\label{pc:main}
Assume Dirichlet GRH. For every rational elliptic curve, every
fixed \(0<u\le1/4\), \(w=1/2+u+it\), and every \(\varepsilon>0\),
\[
 |A_{2N}(w;W^{w})-A_N(w;W^{w})|
    \ll_{E,W,u,\varepsilon}(2+|t|)^{C_\varepsilon}
                                  N^{-2u+\varepsilon}.
\]
If this difference is \(O_{E,W,u}((2+|t|)^C N^{-4u})\)
for \(|t|\le10\log\log(10N)\), then \(L(E,s+1/2)\) has no zero
in \(\Re s>1-u\). For any fixed \(C\ge0\), the supremum
of the difference divided by \((2+|t|)^C\) on this same height range
is \(\gg_{E,W,u,C}N^{-4u}\) for all sufficiently large \(N\).
\end{theorem}

\subsection{Projective fibres}
For a good prime \(p\notin S\) and a vector primitive modulo \(p\),
write \(b_k(v)\) for the coefficient of \(z^k\) in the corrected
local Euler factor, \(z=p^{-s}\), and \(b_0=1\).
The valuation formulas in Theorem~\ref{qcm:eisenstein} apply
also to negative values of \(F\) and, by setting the valuation
equal to infinity, to zero. They depend only on the projective
class modulo \(p^k\), since unit dilations change \(F\) by a fourth
power. Put \(c_p=a_p(E)/(p+1)\), and define on all vectors
modulo \(p^k\)
\[
 t_p(v)=b_1(v)+c_p,\qquad t_{p^k}(v)=b_k(v)-b_{k-2}(v)\quad(k\ge2)
\]
when \(v\) is primitive modulo \(p\), and put all these functions
equal to zero otherwise. Extend multiplicatively in the modulus:
\(t_m(v)=\prod_{p^k\Vert m}t_{p^k}(v)\) for \((m,B_0)=1\).

The projective mean of \(b_1\) is \(-c_p\): the affine
character sum is \(-a_p(E)\) and infinity contributes zero.
For \(k\ge2\) the sum of \(t_{p^k}\) on every fibre
of \(\mathbf P^1(\Z/p^k)\to\mathbf P^1(\Z/p^{k-1})\)
vanishes. Indeed, for \(k=2j\) the increment is
\[
  p^{j-1}{\bf1}_{p^{k-1}\mid F(v)}
                (p{\bf1}_{p^k\mid F(v)}-1).
\]
For \(k=2j+1\ge3\) it vanishes except where \(p^{k-1}\mid
F(v)\), when it equals \(p^j(F(v)/p^{k-1}\mid p)\).
Each projective root is simple, including infinity from
the factor \(b\). In a chart about such a root the next digit
of \(F\) runs through all of \(\F_p\). These formulas prove the
claim in both parities.

\begin{lemma}\label{pc:poisson}
For the normalized Fourier transform on \((\Z/m)^2\),
\[
  m^{-2}\sum_{v\bmod m}t_m(v)e^{-2\pi i r\cdot v/m}
                                          =m^{-1}t_m(Jr).
\]
For \(m>1\) its zero-frequency term vanishes.
\end{lemma}
\begin{proof}
For \(m=p^k\), group primitive vectors by projective line. The
sum over unit generators of the line generated by \(v\) is
\(p^k{\bf1}_{p^k\mid r\cdot v}-p^{k-1}{\bf1}_{p^{k-1}\mid
r\cdot v}\), the prime-power Ramanujan sum. For primitive \(r\)
the first indicator selects the line generated by \(Jr\). The
second gives a complete fibre, or all of \(\mathbf P^1(\F_p)\)
if \(k=1\), and sums to zero. When \(p\mid r\), both indicators
are constant on these fibres and both sums vanish. CRT combines
the prime powers; the dual unit scalings in CRT preserve every
local function. This also proves the last assertion.
\end{proof}

\subsection{The dual series and its exceptional directions}
Let \(r\ne0\), \(l\ge1\), and let \(\eta\) be a Dirichlet character
modulo \(B_0\). Consider
\[
 G_{r,l,\eta}(s)=\sum_{(m,l B_0)=1} t_m(r)\eta(m)m^{-s}.
\]
Unless \(F(r)=0\) or \(F(r)\) belongs to one of finitely many
nonzero rational squareclasses depending only on \(B_0\),
this series extends holomorphically to \(\Re s>1/2\) and
satisfies there, for fixed \(\delta,\epsilon>0\),
\begin{equation}\label{pc:dualbound}
 |G_{r,l,\eta}(s)|\ll_{E,B_0,\delta,\epsilon}
       ((2+\|r\|)l(2+|\Im s|))^\epsilon
                 \qquad(\Re s\ge1/2+\delta).
\end{equation}
To see this, put \(\chi=\chi_{\operatorname{disc}(\Q(\sqrt{F(r)}))}\)
at nonsquares and use the trivial character at squares.
At \(p\nmid B_0 l\gcd(r_1,r_2)\), put \(z=\eta(p)p^{-s}\).
The local generating series is
\[
 (1-z^2)\frac{P_{p,h}(z)}{1-\chi(p)z}+c_p z,
            \qquad h=\lfloor v_p(F(r))/2\rfloor.
\]
Here \(P_{p,h}(z)\) means the polynomial in \(z\) with coefficients
involving powers of \(p\). On dividing by the local factor of
\(L(s,\chi\eta)\), the result is
\((1-z^2)P_{p,h}(z)+c_p z(1-\chi(p)z)\).
Away from primes dividing \(B_0lF(r)\) this differs from one by
\(O(p^{-2\Re s}+p^{-1/2-\Re s})\).
At the remaining primes the bound
\( |P_{p,h}(z)|\le(1+p^{-\Re s})/(1-p^{1-2\Re s})\)
is independent of \(h\); omission of a whole Euler factor contributes
at most \(1+p^{-\Re s}\). Their total cost is
\(O_{\delta,\epsilon}((l |F(r)|)^{\epsilon})\).
Finally \(|F(r)|\ll_E(1+\|r\|)^4\) and Dirichlet GRH gives
the subpower bound for the primitive Dirichlet function away
from \(\Re s=1/2\). The only possible pole is the principal
character, which requires the conductor of the quadratic character
to divide \(B_0\). This proves \eqref{pc:dualbound}.

For every \(r\), including exceptional ones, the elementary
estimates \(|t_p|\ll1\) and
\(|t_{p^k}|\le2p^{\lfloor k/2\rfloor}\) give
\begin{equation}\label{pc:absolute}
     \sum_{m\ge1,(m,B_0)=1}|t_m(r)|m^{-1-\epsilon}
                           \ll_\epsilon 1 .
\end{equation}
Indeed the \(k=1,2\) terms are \(O(p^{-1-\epsilon})\)
and the remaining prime-power terms form summable geometric tails.

There are \(O_{E,B_0}(T)\) exceptional integral vectors with
\(0<\|r\|\le T\), including all their multiples. For \(F(r)=0\)
there are only finitely many rational lines. For a fixed nonzero
squareclass \(d\), a primitive slope \(a/b\) with \(F(a,b)\in
d\Q^{\times2}\) is the abscissa on \(dy^2=x^3+Ax+B\).
The logarithm of \(\max(|a|,|b|)\) is twice the canonical height
up to a bounded error. By Mordell--Weil and positive definiteness
on the free group, only \(O_d((1+\log(2T))^{r_d/2})\) such
primitive slopes have height at most \(T\), for fixed rank \(r_d\).
Thus the sum of the inverse naive heights over these primitive
vectors converges. Summing the multiples of each proves the claim.
For the height and lattice assertions see
\href{https://doi.org/10.1007/978-0-387-09494-6}{Silverman,
\emph{The Arithmetic of Elliptic Curves}, Chapter~VIII}.

\subsection{Poisson summation on a dyadic block}
Fix a smooth function \(V_1\) supported in a compact subinterval
of \((0,\infty)\), a smooth compactly supported function \(V_2\)
in the cone used above, and \(w=1/2+u+it\). For \(M\ge1\)
consider the block, before imposing primitivity,
\[
 Y^{-2}\sum_{v\equiv v_1\ (B_0)}V_2(v/Y)
       \sum_{d>1,(d,lB_0)=1} d^{-w}V_1(d/M)t_d(v).
\]
For \(Y\gg_E1\) its size is
\begin{equation}\label{pc:unprimitive}
 \ll (MlY)^\epsilon(2+|t|)^C M^{-u}
                  \left(\frac M{Y^2}+\frac{\sqrt M}{Y}\right),
\end{equation}
with constants uniform in the residue class \(v_1\).
Indeed, lattice Poisson summation, CRT, and Lemma~\ref{pc:poisson}
give the exact expression
\[
 B_0^{-2}\sum_{r\ne0}\sum_{d>1,(d,lB_0)=1}
   d^{-1-w}V_1(d/M)t_d(Jr)
   e^{2\pi i(r\cdot v_1)\bar d/B_0}
   \widehat V_2(Yr/(B_0d)).
\]
The Fourier kernel for the displayed convention has negative sign.
The finite congruence weight on the units modulo \(B_0\) expands
in the finitely many Dirichlet characters of that modulus.
For nonexceptional \(Jr\), Mellin inversion and
\eqref{pc:dualbound}, shifting the Dirichlet-series abscissa to
\(1/2+\epsilon\), bound the inner sum by
\[
   (Ml(2+\|r\|))^{O(\epsilon)}(2+|t|)^C
               M^{-1-u}(1+Y\|r\|/M)^{-J_0}.
\]
Here \(J_0\) can be arbitrarily large: apply repeated integration
by parts to the Mellin transform of
\(V_1(x)\widehat V_2(Yr/(B_0Mx))\). Its derivatives in \(x\)
obey the same arbitrary power decay, with a changed constant.
For exceptional \(Jr\), \eqref{pc:absolute} gives the corresponding
bound \(M^{-1/2-u+\epsilon}(1+Y\|r\|/M)^{-J_0}\).
The total counts of these frequencies to radius \(T\) are
\(O(T^2)\) and \(O(T)\), respectively. Summing the two decays
proves \eqref{pc:unprimitive}; for \(M/Y<1\) one can absorb any
additional fixed power of \(M/Y\) by increasing \(J_0\). If
\(M\) is bounded where the cutoff might include \(d=1\), treat
the finitely many \(d>1\) individually using the same Poisson
formula and rapid Fourier decay.

\begin{lemma}\label{pc:block}
For \(1\le M\le N^2\), the same block with \(Y=N\), \(v_1=v_0\),
restricted to primitive vectors with \(F(v)\) nonsquare, satisfies
\[
   \ll (MlN)^\epsilon(2+|t|)^C M^{-u}\frac{\sqrt M}{N}.
\]
The constants in both block estimates use finitely many test
function seminorms.
\end{lemma}
\begin{proof}
First keep all values of \(F(v)\). Detect primitivity by
\(\sum_{h\mid\gcd(v)}\mu(h)\). Only \(h\) coprime to \(B_0\)
contribute. If \((h,d)>1\) then \(t_d(v)=0\); for \((h,d)=1\)
we have \(t_d(hy)=t_d(y)\). Put \(H=N/\sqrt M\).
For \(h\le H\), apply \eqref{pc:unprimitive} with \(Y=N/h\),
\(l\) replaced by \(lh\), and the progression \(h^{-1}v_0\).
The extra factor is \(h^{-2}\), and the two contributions sum to
\[
   \ll (MlN)^{O(\epsilon)}(2+|t|)^C M^{-u}
          (MH/N^2+\sqrt M\log(2H)/N).
\]
For \(h>H\) work with the original vector \(v\). If \(F(v)\)
is nonexceptional, \eqref{pc:dualbound} and Mellin inversion
bound its \(d\)-sum by
\((MlN)^\epsilon(2+|t|)^C M^{-u}\). There are
\(O(N^2/H)\) pairs \((h,v)\) with \(h>H\), \(h\mid\gcd(v)\),
since \(h\ll N\) on the given compact support and
\(\sum_{H<h\ll N}(N/h)^2\ll N^2/H\).
Only \(O(N)\) vectors are exceptional, and each has
\(O_\epsilon(N^\epsilon)\) divisors of its gcd. Their individual
\(d\)-sum costs at most \(O(M^{1/2-u+\epsilon})\)
by \eqref{pc:absolute}. Dividing by \(N^2\) proves the assertion
for all primitive vectors. Subtract primitive square-value
vectors using the same exceptional estimate. Shrinking the
auxiliary epsilon absorbs the logarithm and fixed powers.
\end{proof}

\subsection{Completion of the transfer}
On primitive vectors the Euler factor of
\(D_v^S(s)/\zeta^S(2s)\) is
\[
   \left(\sum_{k\ge0}t_{p^k}(v)p^{-ks}\right)-c_p p^{-s},
             \qquad t_1=1 .
\]
Here \(D_v^S\) omits the factors at \(S\). In a coefficient
of index \(m\) one therefore sums over factorizations \(ed=m\),
\((e,d)=1\), \(e\) squarefree and supported outside \(S\), with
weight \(c'_e=\prod_{p\mid e}(-c_p)\) and remaining coefficient
\(t_d(v)\). Since \(|c_p|\le2/\sqrt p\), we have
\(|c'_e|\ll_\epsilon e^{-1/2+\epsilon}\).
At \(d=1\) the vector average is the density of primitive vectors
times \(\int V_2\), up to \(O_{V_2,\epsilon}(N^{-1+\epsilon})\),
also with the square-value exclusion. More precisely this density is
\[
    \kappa=B_0^{-2}\prod_{p\nmid B_0}(1-p^{-2})>0.
\]
Möbius inversion and the smooth lattice sum with mesh \(B_0h/N\)
give this estimate directly: for \(h\ll N\) the unnormalized
lattice error is \(O_{V_2}(1+N/h)\), and
\(\sum_{h\gg N}h^{-2}\ll N^{-1}\).
The square exceptions cost \(O(N^{-1})\) by the preceding height
count.

Consequently in a weighted dyadic block \(m\asymp M\le N^2\),
the vector average of \(D_v^S/\zeta^S(2\cdot)\) differs from
the block of \(\kappa\int V_2\prod_{p\notin S}(1-c_p p^{-s})\)
by
\[
    \ll_{\epsilon} N^{-1+\epsilon}M^{1/2-u+\epsilon}(2+|t|)^C.
\]
Indeed use Lemma~\ref{pc:block} on \(d\asymp M/e\) with \(l=e\),
multiply by \(|c'_e|e^{-1/2-u}\), and sum the resulting
convergent \(\sum_e e^{-3/2+O(\epsilon)}\). Bounded blocks of
\(d>1\) are included by changing fixed cutoff constants. The
\(d=1\) error is \(O(N^{-1+\epsilon}M^{-u+\epsilon})\).

Smooth the Dirichlet series by \(e^{-m/X}\),
\(X=N^{2-2\epsilon}\), taking the auxiliary \(\epsilon>0\) small
compared with \(u\). Dirichlet GRH and the uniform estimate on
the local correcting polynomials give a bound
\(O_{E,u,\epsilon}(N^\epsilon(2+|t|)^C)\) on every fixed
line to the right of \(1/2\), after choosing the subpower exponents
sufficiently small. This also holds after division by
\(\zeta^S(2s)\), since its Euler product there converges absolutely.
The usual gamma Mellin shift to \(\Re(w+z)=1/2+\epsilon\)
thus costs \(O(N^\epsilon(2+|t|)^C X^{-u+\epsilon})\).
Terms with \(m>N^{2-\epsilon}\) are negligible: each individual
absolute coefficient after division is \(O_\epsilon(m^{1/2+\epsilon})\).
Inside this range use a smooth dyadic partition, absorbing
\(e^{-m/X}\) into the test function with uniformly bounded
seminorms. The sum of the dyadic discrepancies is
\(O(N^{-2u+O(\epsilon)}(2+|t|)^C)\). In the limiting
Euler product the absolute coefficients are summable on each
\(\Re s>1/2\), so undoing its smoothing costs
\(O(X^{-u+\epsilon})\). Restoring the fixed Euler factors
\(L_S(s)\) and \(\zeta^S(2s)\) gives, on the line of \(w\),
\begin{align*}
 A_N(w;V_2)&=\kappa\left(\int V_2\right) Z_P(w)
             +O(N^{-2u+O(\epsilon)}(2+|t|)^C),\\
 Z_P(s)&=L_S(s)\zeta^S(2s)\prod_{p\notin S}(1-c_p p^{-s}).
\end{align*}
The projective means of the corrected odd prime-power coefficients
are all \(-c_p\), and those of the even coefficients are all one,
also directly from the vanishing fibre sums.
Taking \(V_2=W^w\), whose test seminorms grow polynomially in its
height, proves the first assertion of Theorem~\ref{pc:main}.

For the last two assertions write \(v(s)=4(1/2-s)\).
The individual functional equation and exact homogeneity of \(F\)
give, with \(g(s)=\pi^{s-1/2}\Gamma((1-s)/2)/\Gamma(s/2)\),
\[
\begin{split}
 A_{2N}(1-w;W)-2^{v(1-w)}A_N(1-w;W)
  ={}&g(1-w)c^{w-1/2}(2N)^{4(w-1/2)}\\
    &\times[A_{2N}(w;W^w)-A_N(w;W^w)].
\end{split}
\]
The functions on the left are entire; square values have been
removed throughout. For \(\Re s>3/2\), dominated convergence
gives their limit \((1-2^{v(s)})\kappa Z_P(s)\int W\)
with \(s=1-w\). This follows from the coefficient means already
computed and the uniform \(\sqrt m\) bound on the coefficients
of \(D_v\). In \(\Re s>0\),
\[
   Z_P(s)=\frac{L_S(s)\zeta^S(2s)H(s)}
                            {L^S(E,s+1)},\qquad
   H(s)=\prod_{p\notin S}
           \frac{1-c_p p^{-s}}{1-a_p p^{-s-1}+p^{-1-2s}} .
\]
Here \(H\) is holomorphic and nowhere zero. Indeed
\(|a_p/p-c_p|\ll p^{-3/2}\);
the remaining quadratic discrepancy is \(p^{-1-2s}\), and
the local factors on both sides are nonzero throughout that
half-plane. Thus the quotient product converges normally
on compact sets. To apply the three-lines and finite-height
rectangle proof of Theorem~\ref{qcm:eisenstein}, the right boundary
is \(\Re s=2\), the nonzero limiting anchor is \(s=7/4\),
and the reflection exponent is four. All other hypotheses are
identical. On the uncontrolled part of the left boundary the
trivial GRH bound for the two right-hand moments, multiplied
by the conductor factor, costs \(N^{4u+\epsilon}(2+|t|)^C\),
which suffices for the same rectangle argument. For
\(\Re s>1/2-u\), \(u\le1/4\), Dirichlet GRH excludes zeros
of \(\zeta(2s)\), the elementary factor \(1-2^{v(s)}\)
can vanish only at \(\Re s=1/2\), and \(L(E,s+1)\) has no
zero on that line (the classical zero-free line at unitary real
part one). This proves both remaining assertions of the theorem.
The argument establishes half of the exponent required by its
sufficient criterion, for a family with conductors of order \(N^4\).

\subsection{The generic frequencies and the elliptic factor they share}
\begin{proposition}\label{pc:budget}
Deleting every exceptional direction in the Poisson sum and
omitting the cost of the primitivity sieve leave the exponent
\(2u\) in the optimization of the separate bounds used above.
This holds for every nonzero fixed smooth spatial weight.
\end{proposition}
\begin{proof}
Let \(V\in C_c^\infty(\R^2)\) be nonzero. As \(R\to\infty\),
\[
 \sum_{r:\ F(Jr)\ {\rm nonexceptional}}|\widehat V(r/R)|
                     \asymp_V R^2.
\]
For the lower bound take a ball where \(|\widehat V|\) is bounded
below. Of its \(\gg R^2\) lattice samples, at most \(O(R)\) are
exceptional. Schwartz decay gives the upper bound. In a block
\(m\asymp M\) the generic number of frequencies is thus of order
\((M/N)^2\) at the relevant scale. Combining this frequency count
with the individual square-root Dirichlet bound yields a budget
\(M^{1-u}/N^2\), even after all exceptional terms have been
discarded. If the truncation length is \(X\), the two separate
majorants at its endpoint and for its tail have respective scales
\(X^{1-u}/N^2\) and \(X^{-u}\). For \(0<u<1\) the minimum of
their maximum is exactly \(N^{-2u}\), at \(X=N^2\).
At real \(w=1/2+u\) the obstruction also holds uniformly
for moving nonnegative weights of mass one supported in the
same compact set. Indeed write \(V=W F^u\), whose positive
mass stays bounded above and below, and normalize its mass
to one. For the two-scale difference its Fourier weight is
\(\widehat V(2\xi)-\widehat V(\xi)\). If \(m=\int vV(v)\,dv\),
this is \(-2\pi i\, m\cdot\xi+O(\|\xi\|^2)\), with a common
Taylor constant from the fixed compact support. The second
component of \(m\) is at least one and the first is uniformly
bounded. A sufficiently small fixed ball about a point on the
positive second coordinate axis therefore gives a uniform
positive lower bound for the modulus. Its generic lattice
samples again number \(\gg R^2\). This computes the best exponent from these
majorants. Additional cancellation between distinct directions
could change the actual mean.
\end{proof}

There is a second reason for distinguishing individual dual
series from an average of them. An elliptic divisor is already
present in each individual frequency at its first Euler level.

\begin{proposition}\label{pc:divisor}
For any fixed \(v\) with \(F(v)\ne0\), let \(\chi\) be its
quadratic character (trivial for a square).
On \(\Re s>0\) the function
\(\mathcal G_v(s)=\zeta^S(2s)G_{v,1,1}(s)\) continues
meromorphically. For each \(\delta>0\), choose an integer
\(K\ge1\) with \((K+1)\delta>1/2\).
After deleting finitely many local factors depending on
\(v,E,\delta\), on \(\Re s\ge\delta\) its factorization is
\[
 L(s,\chi)\prod_{j=1}^{K}
     L(E\otimes\chi^{j-1},1+js)^{(-1)^{j-1}}\mathcal H_v(s),
\]
where all displayed \(L\)-functions have the same local deletions
and \(\mathcal H_v\) is holomorphic and nowhere zero on a
neighborhood of this half-plane. The original finitely many
local factors of \(\mathcal G_v\) are holomorphic in \(\Re s>0\).

In particular, if \(L(E,1+s_0)=0\) and \(\Re s_0>1/4\),
then \(\mathcal G_v(s_0)=0\) for every \(v\) with \(F(v)\ne0\).
If any quadratic twist of a fixed curve violates ERH, some
twist \(E'\) and some \(s_0\), \(0<\Re s_0<1/2\), have this
same common-zero property for all the functions \(\mathcal G_v\)
formed using \(E'\), without any fixed lower bound on the size
of the hypothetical violation.
\end{proposition}
\begin{proof}
For \(p\nmid B_0F(v)\) one has \(b_k(v)=\chi(p)^k\) and
\(\chi(p)^2=1\). Thus every \(t_{p^k}(v)\), \(k\ge2\), is
zero. The local factor of \(G_{v,1,1}\) is exactly
\(1+(\chi(p)+c_p)z\), \(z=p^{-s}\).
After including \(\zeta_p(2s)\) and dividing by \(L_p(s,\chi)\)
one obtains \(1+c_p z/(1+\chi(p)z)\).
For large enough \(p\) this has a single-valued Euler logarithm
on \(\Re s\ge\delta\). There
\[
 \log\left(1+\frac{c_p z}{1+\chi(p)z}\right)
  =\sum_{j=1}^{K}(-\chi(p))^{j-1}c_p z^j
      +O_{K,\delta}(p^{-1-2\delta}+p^{-1/2-(K+1)\delta}).
\]
On the same half-plane the logarithm of the elliptic local
factor of \(L(E\otimes\chi^{j-1},1+js)\) is
\((a_p/p)\chi(p)^{j-1}z^j+O(p^{-1-2j\delta})\).
Since \(c_p-a_p/p=O(p^{-3/2})\), the difference of the
logarithms after accounting for all \(K\) factors sums normally.
The strict choice of \(K\) allows a small enlargement of the
half-plane. At the deleted primes the original factors have at
most denominators \(1-\chi(p)p^{-s}\) and \(1-p^{-2s}\);
these do not vanish on \(\Re s>0\). Elliptic and Dirichlet
finite Euler factors in the displayed arguments likewise have
no zeros or poles there. This proves the continuation and the
first zero assertion by setting \(K=1\) near the zero in question.

For the last statement let \(b>1/2\) be the supremum of the
unitary real parts of the zeros of all quadratic twists in this
family. Choose one of these twists \(E'\) with a zero of unitary
real part \(1/2+\Re s_0\) satisfying
\(2\Re s_0>b-1/2\). This is possible by the definition of \(b\).
For every \(j\ge2\), every elliptic factor in the displayed
factorization now has unitary real part \(1/2+j\Re s_0>b\).
All belong to the same twist family, so none can cancel the
zero of the \(j=1\) factor. The Dirichlet factor has no pole
at \(s_0\). This works simultaneously for every nonsingular
frequency. Neither assertion here requires Dirichlet GRH.
\end{proof}

\subsection{Odd-degree hyperelliptic curves}
\begin{corollary}\label{pc:hyperelliptic}
Let \(C:y^2=f(x)\) with \(f\in\Z[x]\) monic, squarefree and
of degree \(2g+1\), \(g\ge1\). Put \(D=2g+2\).
For \(F(a,b)=b^D f(a/b)\), replace the exponent four in
the definition of the centered comparison by \(D\).
Under Dirichlet GRH the projective moment again satisfies the
\(N^{-2u+\varepsilon}\) estimate of Theorem~\ref{pc:main}.
The necessary lower bound of that theorem becomes \(N^{-Du}\).

If the Hasse--Weil function of \(H^1(C)\) is automorphic with
the standard weight-one normalization, the estimate
\(O((2+|t|)^C N^{-Du})\) up to \(10\log\log(10N)\) implies
that its unitary standard function has no zeros in \(\Re s>1-u\).
The same implication applies given a meromorphic continuation
and nonvanishing on the unitary line \(\Re s=1\).
\end{corollary}
\begin{proof}
Enlarge \(S\) to contain \(2\), the primes dividing the discriminant,
and enough initial primes in terms of \(g\). One may impose
\(a\equiv1,\ b\equiv p\bmod p^2\) for each \(p\in S\);
the simple factor at infinity gives \(v_p(F)=1\) there.
The even squarefree part has positive discriminant \(4d\)
on the chosen positive support. For \(p\notin S\), all projective
roots are simple and \(\sum_{\mathbf P^1(\F_p)}(F\mid p)
=\#C(\F_p)-(p+1)=-a_p\). Here
\(a_p=\sum_{i=1}^{2g}\alpha_{i,p}\) and
\(|\alpha_{i,p}|=\sqrt p\) by Weil's theorem for curves.
Every fibre, Fourier identity and Dirichlet-series estimate above
is unchanged with \(c_p=a_p/(p+1)\). Unit dilations preserve
the character because \(D\) is even. The conductor of any
dual character is \(O_C((1+\|r\|)^D)\); a fixed power only
changes the auxiliary subpower exponents in \eqref{pc:dualbound}.
For \(g\ge2\), each exceptional nonzero squareclass lies on a
fixed quadratic twist of \(C\), which has finitely many rational
points by \href{https://geodesic.mathdoc.fr/item/IM_1983__73_143051/}
{Faltings, \emph{Endlichkeitssätze für abelsche Varietäten über
Zahlkörpern}, Invent. Math. 73 (1983)}. These exceptional vectors
thus still number \(O_{C,B_0}(T)\) up to radius \(T\).
No effective estimate for the finitely many rational points
is required. This proves the same block estimate and convergence rate.

The inverse local Hasse--Weil factor at \(1+s\) has linear term
\(-a_p p^{-1-s}\). Its terms of degree at least two have total
size \(O_{g,\sigma}(p^{-1-2\sigma})\) on \(\Re s\ge\sigma>0\).
This proves the same nonvanishing correction product in \(\Re s>0\)
as in Theorem~\ref{pc:main}, after the indicated finite exclusions.
The primitive Dirichlet functional equations have conductor factor
\((4F(v))^{1/2-s}\); their rescaling therefore introduces \(D\)
in place of four. Three-lines with the same nonzero anchor proves
both the necessary lower bound and the stated sufficient estimate.
Under the alternative continuation assumption, the usual absolute
Euler product gives nonvanishing for unitary real part greater
than one.
\end{proof}

The number of available primitive arguments is itself constrained
by the corrected conductor. For any fixed integral binary form
of degree \(D\ge3\) and nonzero discriminant, one has
\[
   \#\{(a,b)\in\Z^2:0<|F(a,b)|\le Q\}\ll_F Q^{2/D}.
\]
This is the Mahler--Thunder lattice point estimate, including
reducible forms; see
\href{https://arxiv.org/abs/1605.03427}{Stewart and Xiao,
\emph{On the representation of integers by binary forms},
Lemma~2.1}. The projective family above attains this order
in a fixed positive cone, by its positive primitive density.
Thus even allowing unbounded-height points near real roots
cannot increase the power of the number of distinct primitive
slopes with corrected conductor \(c|F(a,b)|\le Q\).
The primitive fundamental discriminant after deleting the
square correction is a different size parameter.

\subsection{The remaining possible cancellations in a full-strip criterion}
The exact homogeneous reflection also gives a version up to the
elliptic critical line. One must account for the factor
\(\zeta^S(2s)\) in the limiting mean.

\begin{corollary}\label{pc:fullstrip}
In Theorem~\ref{pc:main} and Corollary~\ref{pc:hyperelliptic},
the proved discrepancy bound and the necessary lower bound hold
for every fixed \(0<u<1/2\). Assume, for a sequence of such
\(u\to1/2\), the sufficient estimate with exponent \(Du\) and
the same restricted heights (with constants permitted to depend on
\(u\)). Write
\[
   T_N(s)=A_{2N}(s;W)-2^{D(1/2-s)}A_N(s;W).
\]
Then, under the analytic hypotheses of the respective theorem,
all zeros right of the unitary critical line must have the form
\(1/2+\rho/2\) for a nontrivial zeta zero \(\rho\). Their
multiplicity is at most that of \(\rho\). Under the same
assumptions, ERH holds if and only if
\[
    \lim_{N\to\infty} T_N^{(j)}(\rho/2)=0
             \quad\text{for every }\rho
             \quad\text{and }0\le j<\operatorname{ord}_{\rho}\zeta .
\]
In particular, the additional tests concern only a countable set of points
on \(\Re s=1/4\), and impose no
little-oh bound on a full vertical line.
\end{corollary}
\begin{proof}
In the dyadic calculation the exponent \(1/2-u\) is still
positive; all conductor subpower bounds and the exact reflection
are unchanged. The strip in the harmonic measure argument now
has width less than two, for which the same height cutoff makes
its uncontrolled tails \(o(1)\). Thus under the hypotheses
\(T_N\) and all its derivatives converge locally uniformly on
\(0<\Re s<2\) to the holomorphic continuation of
\((1-2^{D(1/2-s)})\kappa Z_P(s)\). The preliminary convergence
with limiting \(Z_P\) at \(\Re s>3/2\) fixes this continuation
uniquely. In \(0<\Re s<1/2\), all factors except
\(\zeta^S(2s)\) and the inverse Hasse--Weil function are
holomorphic and nonvanishing. Orders therefore give both the
claimed restriction and exactly the derivative tests. The
boundary at unitary real part one is handled by the assumed
nonvanishing or the automorphic theorem. Dirichlet GRH places
the zeta zeros on \(\Re\rho=1/2\).
\end{proof}

\subsection{Failure of meromorphic continuation at a random projective direction}
The distinction between integer directions and the complete local
average can also be seen at a single real point. Independently for
each \(p\notin S\) choose \(v_p\) with the invariant probability
distribution on \(\mathbf P^1(\Z_p)\), so every finite projective
reduction is uniform. Form the centered local product
\[
       G_{\boldsymbol v}(s)=\prod_{p\notin S}
                     \left(1+\sum_{k\ge1}t_{p^k}(v_p)p^{-ks}\right).
\]

\begin{proposition}\label{pc:random}
For each real \(\sigma>1/2\), the product is positive almost
surely and has expectation one. With probability one,
\[
    \lim_{\sigma\downarrow1/2}
          \frac{\log G_{\boldsymbol v}(\sigma)}{\log(2\sigma-1)}
                      =\frac12.
\]
In particular a random projective direction almost surely gives
no meromorphic extension across the real point \(1/2\), and the
products are not uniformly integrable as \(\sigma\downarrow1/2\).
Both conclusions hold in every hyperelliptic genus above.
\end{proposition}
\begin{proof}
Write \(x_p=t_p(v_p)\) and \(e_p=v_p(F(v_p))\), where the
valuation of the homogeneous polynomial is taken on primitive
coordinates. The independent bounded variables \(x_p\) have
mean zero. The projective root-lifting rule gives
\(\Pr(e_p\ge j)\le D/((p+1)p^{j-1})\) for \(j\ge1\).
Thus almost surely only finitely many \(p\) have \(e_p\ge2\),
and \(\sum_{e_p\ge1}1/p<\infty\) almost surely, by summing its
expectation. Also \(e_p<\infty\) simultaneously at every prime
almost surely.

Set \(z=p^{-\sigma}\). At \(e_p=0\) the local factor is exactly
\(1+x_p z\). At \(e_p=1\) it is \(1+x_p z-z^2\).
Every local factor is positive, including at \(\sigma=1/2\):
from its uncentered formula it is at least
\(1-z-|c_p|z>0\) after the fixed exclusions in \(S\).
Indeed the uncentered local series equals \(1/(1-\chi z)\)
at \(\lfloor e_p/2\rfloor=0\), and is at least one otherwise.
Consequently the finitely many higher-valuation factors contribute
\(O_{\boldsymbol v}(1)\) to the real logarithm near this point.
Taylor expansion of the other factors, with
\(\sum_p p^{-3/2}<\infty\), gives
\[
 \log G_{\boldsymbol v}(\sigma)
     =\sum_p x_pp^{-\sigma}-\tfrac12\sum_p p^{-2\sigma}
                                      +O_{\boldsymbol v}(1).
\]
Here \(x_p^2=1+O_D(p^{-1/2})\) when \(e_p=0\), and the total
error at \(e_p\ge1\) is bounded by the summable exceptional
series just proved. All prime sums omit \(S\).

For completeness,
\(\sum_{p\ge 5}\operatorname{Var}(x_p)/(p(\log\log p)^2)<\infty\)
by the prime harmonic sum. Kolmogorov's convergence theorem and
Kronecker's lemma give
\(\sum_{p\le X}x_p/\sqrt p=o(\log\log X)\) almost surely.
Partial summation now gives
\(\sum_p x_p p^{-1/2-h}=o(\log(1/h))\) as \(h\downarrow0\).
Indeed, after putting \(v=\log X\), the Abel kernel
\(he^{-hv}\,dv\) integrates \(\log(2+v)\) to
\(\log(1/h)+O(1)\). Meanwhile
\(\sum_p p^{-1-2h}=\log(1/(2h))+O(1)\) by the simple pole
of \(\zeta\). This proves the limiting ratio simultaneously
as \(h\downarrow0\). A nonzero meromorphic function has integral
order at a point, proving the continuation claim.

Finally the projective fibre sums make the \(t_{p^k}\) at
different levels mutually orthogonal in real \(L^2\).
Their second moments are \(O_D(1)\) uniformly in \(p,k\):
at even \(k\) the two possible nonzero sizes are
\(O(p^{k/2})\) on \(e_p\ge k\) and \(p^{k/2-1}\) on
\(e_p=k-1\), and at odd \(k\ge3\) the size
\(p^{(k-1)/2}\) requires \(e_p=k-1\). The case \(k=1\)
is bounded directly. The local factor thus has mean one and
variance \(O_D(\sum_{k\ge1}p^{-2k\sigma})\). Independence
makes the partial products a mean-one martingale bounded in
\(L^2\) for each fixed \(\sigma>1/2\). Its expectation passes
to the limit. Uniform integrability along \(\sigma\downarrow1/2\)
would instead make this expectation converge to zero, since
the almost sure limit is zero. This proves the final claim.
\end{proof}

\begin{corollary}\label{pc:natural}
With probability one, \(G_{\boldsymbol v}(s)\) is holomorphic
on \(\Re s>1/2\) and admits no meromorphic continuation across
any point of its boundary line.
\end{corollary}
\begin{proof}
The almost sure bound on \(\sum_{p\le X}x_p/\sqrt p\) gives
locally uniform convergence of the linear Dirichlet series on
\(\Re s>1/2\) by partial summation. The remaining Euler logarithms
converge absolutely there, after removing the almost surely finite
set with \(e_p\ge2\). Their local factors have no poles on
\(\Re s>0\), hence the full function is holomorphic on the
claimed half-plane. Removing these finite factors cannot affect
whether meromorphic continuation is possible at a point. Denote the
product after their removal by \(G^*\).

Fix \(t\ne0\) and put \(s=1/2+h+it\), \(h>0\).
As in Proposition~\ref{pc:random}, Taylor expansion gives
\[
 \log|G^*(s)|=\sum_p x_p\cos(t\log p)p^{-1/2-h}
                                      +O_{\boldsymbol v,t}(1).
\]
Indeed the real part of the quadratic sum
\(\sum_p p^{-1-2h-2it}\) is bounded as \(h\downarrow0\),
by nonvanishing and regularity of \(\zeta(1+2it)\) and absolute
convergence of the remaining Euler logarithm. The finitely many
omitted linear terms are bounded. The same Kronecker argument
with \(x_p\cos(t\log p)\) gives
\(\log|G^*(s)|=o(|\log h|)\) almost surely. Its linear sum has
variance
\[
 V_h=\sum_p(1+O_D(p^{-1}))\cos^2(t\log p)p^{-1-2h}
                        =\tfrac12\log(1/h)+O_t(1),
\]
since \(\operatorname{Var}(x_p)=1-\Pr(e_p\ge1)-c_p^2\).
Dividing that sum by \(\sqrt{V_h}\) gives convergence in law
to a standard normal variable. For example, expand the logarithm
of its characteristic function to second order: the sum of the
absolute third moments before normalization is uniformly bounded
by \(O_D(\sum_pp^{-3/2})\), and the summands are independent
and bounded. The bounded almost sure remainder divided by
\(\sqrt{V_h}\) tends to zero. In particular, for every \(A>0\),
\(\Pr(|\log|G^*(s)||\le A)\to0\). Along \(h=1/n\) the event
that this logarithm is eventually bounded has probability zero,
by Fatou's lemma and a countable union over integral \(A\).

A meromorphic zero or pole would violate the almost sure
little-oh estimate, and a regular nonzero value would give
an eventually bounded logarithm. This excludes meromorphic
extension almost surely at the given point. Intersect the
probability-one events over nonzero rational \(t\); their density
excludes extension through any neighborhood of any boundary point.
For comparison, random Euler products with independent unbiased
Rademacher coefficients are classical; see
\href{https://www.math.mcgill.ca/radziwill/helsonRev4.pdf}
{Harper, Nikeghbali and Radziwiłł, \emph{A note on Helson's
conjecture on moments of random multiplicative functions}, \S1}.
The local variables here have curve-dependent bias before
centering, and higher terms from projective ramification.
\end{proof}

The same proof also permits a separable integral binary branch
form of even degree with no rational branch point. Choose a
primitive integral base vector with nonzero value; a sufficiently
fine fixed congruence keeps all valuations, squarefree unit
classes and quadratic character values at \(S\) fixed. Omit
the fixed polynomials \(P_{p,h}\) at these primes from the
definition of each entire \(D_v\). Its corrected conductor is
then the old conductor divided by the fixed square
\(\prod_{p\in S}p^{2\lfloor v_p(F)/2\rfloor}\), hence is still
a fixed multiple of \(|F|\). Choose a compact cone of constant
sign. For a negative value the character is odd and the gamma
ratio in the primitive functional equation is replaced by its
odd-parity version, with the same polynomial vertical bounds.
All good-prime projective fibre and trace formulas hold in either
chart including infinity, and no other part of the proof changes.

\paragraph{Amplification and quadratic ranks.}
If a hyperelliptic Jacobian is \(\Q\)-isogenous to \(E^m\times B\),
its Hasse--Weil function contains \(L(E,s)^m\). Whenever
\(L(B,s)\) is holomorphic at the point in question, the
multiplicity allowance in Corollary~\ref{pc:fullstrip} for the
elliptic factor becomes \(\lfloor\operatorname{ord}_\rho\zeta/m\rfloor\).
In particular \(m=2\) suffices at every simple zeta zero, under
the hypothesized moment control for this curve.

Arbitrarily large such multiplicities for one fixed \(E/\Q\)
would already imply unbounded Mordell--Weil ranks among its
rational quadratic twists. Indeed write the hyperelliptic
curve as \(C:y^2=f(t)\), with involution \(\iota\).
The morphism \(P\mapsto[P-\iota P]\) into \(J_C\) is defined
over \(\Q\) even in the absence of a rational base point.
Composing it with the \(m\) independent rational homomorphisms
to \(E\) gives \(m\) anti-invariant points of \(E(\Q(C))\).
They are independent modulo torsion: on Jacobians the difference
map induces \(1-\iota_*=2\), while a constant map induces zero.
These are rational points of the quadratic twist over \(\Q(t)\),
and specialization preserves their independence outside bounded
height by the specialization theorem (the nonconstant quadratic
twist has zero constant trace). In the other direction, independent
points on this function-field twist give independent homomorphisms
from \(J_C\) to \(E\), since anti-invariant constants are
two-torsion. Classical examples therefore supply finite amplification:
for \(j(E)\ne0,1728\), Mestre constructs a quadratic twist over
\(\Q(t)\) of rank at least two, and Rubin and Silverberg give
explicit further families of rank at least three. See
\href{https://arxiv.org/abs/math/0010056}{Rubin and Silverberg,
\emph{Rank frequencies for quadratic twists of elliptic curves},
introduction and \S2}. This rank implication does not supply
the associated centered moment estimate.

\clearpage
\section{Which quadratic-character averages can represent an elliptic curve?}\label{albs:section}
The binary quartic in Section~\ref{pc:section} encodes the trace of a curve in the average of rational Dirichlet characters. One might try to gain a power in its analytic estimate by starting with a smooth polynomial in more independent variables. There is a geometric necessity for every such construction, including singular polynomials and rational functions.

Let $K=\Q(\mathbf P^n)$, $n\ge1$, and let $f\in K^\times$ be a nonsquare even over $\overline\Q(\mathbf P^n)$. Choose a dense open $U\subset\mathbf P^n$ where $f$ is regular and invertible. For sufficiently large primes define the projective average
\[
        M_f(p)=\frac1{\#\mathbf P^n(\F_p)}
                    \sum_{x\in U(\F_p)}\left(\frac{f(x)}p\right).
\]
Different fixed choices of $U$ change the answer by $O_f(p^{-1})$. Let $X$ be a smooth projective model of the geometrically integral extension $K(\sqrt f)$ and $A=\operatorname{Alb}_{X/\Q}$ its Albanese variety. For an abelian variety $A$ let $t_p(A)$ denote the trace on its first cohomology at a good prime. All $L$-functions in this section are in arithmetic normalization.

\begin{theorem}\label{albs:rigid}
For every such family,
\[
                         M_f(p)=-\frac{t_p(A)}p+O_f(p^{-1}).
\]
For every elliptic curve $E/\Q$ the following conditions are equivalent:
\begin{enumerate}
\item $M_f(p)=-a_p(E)/p+O_{f,E}(p^{-1})$ at all sufficiently large primes;
\item $A$ is isogenous over $\Q$ to $E$.
\end{enumerate}
It suffices in the first condition to use a set of primes of density one. When these conditions hold there is a separable quartic $g(z)\in\Q[z]$, a nonconstant $r\in K$ and $h\in K^\times$ such that
\[
                  f=h^2 g(r) \quad\text{in }K.
\]
The genus-one curve $y^2=g(z)$ has Jacobian $\Q$-isogenous to $E$. It may be a torsor without a rational point. Thus adding parameters in such a representation of the trace always yields a rational pullback of a one-parameter genus-one double cover. The distribution of the values of the rational map $r$ can still affect an analytic error estimate.
\end{theorem}

\begin{proof}
Spread the varieties, the maps on dense opens and their complements out over $\Z[1/S]$. After shrinking the opens, the double cover over $U$ is finite \'etale and agrees there with an open of $X$. Each fibre over a rational point $x$ has $1+(f(x)\mid p)$ points. Each fixed complement has dimension at most $n-1$ and $O(p^{n-1})$ points, uniformly outside $S$. Thus the character sum is $\#X(\F_p)-\#\mathbf P^n(\F_p)+O(p^{n-1})$. The top cohomology contributes $p^n$ for either variety. By hard Lefschetz for an ample class, the next contribution on $X$ is $-p^{n-1}t_p(A)$: the first cohomology of $X$ is the first cohomology of its Albanese, and cup product $n-1$ times supplies the Tate twist. The trace formula and Deligne's weight bound put the other terms in $O_X(p^{n-1})$. Division by $1+p+\cdots+p^n$ proves the first formula since $|t_p(A)|\le2\dim(A)\sqrt p$. The cited cohomological inputs are Grothendieck's trace formula and Deligne, \href{https://doi.org/10.1007/BF02684780}{\emph{La conjecture de Weil II}}; see also \href{https://kskedlaya.org/weil-cohom/}{Kedlaya's notes on Weil cohomology}.

We give the step from bounded traces to equal traces, since dropping it would lose the arithmetic content of the equivalence. Choose $\ell$ and let $G$ be the compact image of the joint $\ell$-adic first-cohomology representation of $A$ and $E$. The continuous conjugation-invariant trace difference $D$ takes values in a fixed finite set of integers on the good Frobenius classes in the matching set, by the first formula. Their conjugates are dense in $G$ by Chebotarev, also after deletion of a density-zero set. Thus $D$ takes values in this finite set everywhere. In particular for any good Frobenius $g$ with residue characteristic different from $\ell$, the integers $D(g^j)$ are bounded for every $j\ge1$. In a complex embedding, group its cohomology eigenvalues $\alpha$ with their signed multiplicities $n_\alpha$. The series on the left is holomorphic in the unit disk and satisfies, initially near zero,
\[
       \sum_{j\ge1}D(g^j)z^j
                         =\sum_\alpha n_\alpha\frac{\alpha z}{1-\alpha z}.
\]
The distinct poles on the right have modulus $p^{-1/2}<1$ by weight one. Holomorphy forces every $n_\alpha=0$, in particular $D(g)=0$. Chebotarev and the character criterion for semisimple representations now identify the two semisimplified modules. Faltings's semisimplicity and isogeny theorems give $A\sim_{\Q} E$. Conversely an isogeny gives identical good traces. For the isogeny input see \href{https://doi.org/10.1007/BF01388432}{Faltings, \emph{Endlichkeitss\"atze f\"ur abelsche Variet\"aten \"uber Zahlk\"orpern}}.

Choose an equivariant smooth projective resolution for the double-cover involution $\iota$ in characteristic zero, or use its birational action on the Albanese torsor. The canonical morphism $X\to C=\operatorname{Alb}_{X/\Q}^{1}$ is defined over $\Q$ without a choice of rational point and is functorial in $X$; see \href{https://arxiv.org/abs/math/0611284}{Wittenberg, \emph{On Albanese torsors and the elementary obstruction}, \S2}. Since $A$ has dimension one and the Albanese image generates it geometrically up to translation, the map is dominant. The induced involution on $C$ has linear part $+1$ or $-1$ on its Jacobian. If its linear part were $+1$ it would be a translation of order at most two, and $C/\langle\iota\rangle$ would have genus one. This would give a dominant rational map from $\mathbf P^n$ to a curve of genus one. Such a map is impossible: over $\overline\Q$ restrict it to a line through two points with different images and use properness and Riemann--Hurwitz. The linear part is therefore $-1$. Its four geometric fixed points give a genus-zero quotient $B$ with a double cover $C\to B$ branched in four distinct points.

The dominant rational map $\mathbf P^n\dashrightarrow B$ evaluates to a rational point on its nonempty domain, so $B\simeq\mathbf P^1$ over $\Q$. Choose a rational coordinate with infinity off the branch set. The quadratic extension $\Q(C)/\Q(B)$ then has the form $y^2=g(z)$ with $g$ a separable quartic, including its rational scalar twist. The pullback of $z$ is a nonconstant rational function $r\in K$. The element $y\in\Q(C)\subset\Q(X)$ is anti-invariant, hence $K(y)=K(\sqrt f)$. Equality of these nontrivial quadratic extensions is precisely $f/g(r)\in K^{\times2}$. This proves the assertion over $\Q$.
\end{proof}

\begin{corollary}\label{albs:euler}
After omitting a sufficiently large finite set $S$, the product with first coefficient $M_f(p)$ has the identity
\[
            \prod_{p\notin S}(1+M_f(p)p^{-s})
                        =\frac{H(s)}{L^S(A,1+s)} ,
\]
initially for $\Re s>1/2$, where $H$ is holomorphic and nowhere zero in $\Re s>0$. Wherever the abelian $L$-function continues meromorphically in that region the identity continues with it. If the branch divisor is a smooth hypersurface of even degree in $\mathbf P^n$ with $n\ge2$, then $A=0$. Such a smooth higher-dimensional family cannot satisfy the elliptic matching condition.
\end{corollary}
\begin{proof}
Form the ratio to the reciprocal of the local Hasse--Weil Euler factor. The first-order coefficient is $M_f(p)+t_p(A)/p=O_f(p^{-1})$. On each $\Re s\ge\epsilon>0$ the higher logarithmic terms are $O_{f,\epsilon}(p^{-1-2\epsilon})$ by weight one. The local logarithms therefore sum normally to a holomorphic function, proving the first assertion after exponentiating. For a smooth branch hypersurface of degree $2e$, write its equation homogeneously as $F=0$. Over $\overline\Q$ the ordinary hypersurface $z^{2e}=F(x_0,\ldots,x_n)$ in $\mathbf P^{n+1}$ is smooth and dominates the double cover via $y=z^e$. Weak Lefschetz gives $H^1=0$ on this hypersurface for $n\ge2$. The double cover itself is smooth, and pullback on rational first cohomology is injective under a generically finite surjective map, by trace and degree. Its Albanese is therefore zero. Birational invariance gives the same conclusion for the chosen smooth projective model.
\end{proof}

\subsection*{The classification in every weight-one degree}
The preceding obstruction extends to abelian varieties of arbitrary dimension and to elliptic and hyperelliptic Jacobians over every number field. For the geometric statement work over a number field $k$ and use the same symbols with base $k$. A genus-one double cover is included when we refer to curves of genus $g\ge1$ with a degree-two map to the projective line.

\begin{theorem}\label{albs:hyper}
If the double cover $X$ has $g=\dim\operatorname{Alb}_X>0$, then there is a smooth geometrically connected curve $C/k$ of genus $g$ with a degree-two map to $\mathbf P^1_k$. Its Jacobian is $k$-isomorphic, as an abelian variety, to $\operatorname{Alb}_X$. There is a dominant rational map $r:\mathbf P^n\dashrightarrow\mathbf P^1$ with geometrically connected generic fibre such that
\[
   k(X)= k(\mathbf P^n)\, k(C),\qquad f=h^2g_0(r),
\]
where $h\in k(\mathbf P^n)^\times$ and $g_0\in k[z]$ is separable of degree $2g+2$. Conversely, for any such $C$, taking $r$ a coordinate on an affine chart realizes its Jacobian by this quadratic-character average. Consequently an abelian variety $B/k$ of positive dimension has a rational multivariable quadratic-character representation of its first Frobenius trace to error $O(1/\mathrm N v)$ in the normalized average, at a density-one set of primes, exactly when $B$ is $k$-isogenous to the Jacobian of such a curve.
\end{theorem}

\begin{proof}
Use an equivariant smooth projective model and first pass to $\bar k$. For $j\ge1$, any involution-invariant regular $j$-form on $X$ descends rationally to $\mathbf P^n$, by separability and Galois descent for differentials of function fields. Its descent is regular in codimension one. Here are local details of the latter assertion. Over a general point of a prime divisor of $\mathbf P^n$, the normalized finite double cover is regular; the proper smooth model gives the same pulled-back form there by the valuative criterion, since its birational map to $X$ is defined at every codimension-one point. At an unramified divisor a pole persists under pullback. At a ramified divisor take a local transverse parameter $x$ on the base and, after an \'etale localization at its generic point, a parameter $z$ with $x=z^2$ upstairs; the other parameters are unchanged. A tangential component with a pole of order $j_0\ge1$ acquires order $2j_0$ upstairs. A component involving $dx$ acquires order $2j_0-1\ge1$ because $dx=2z\,dz$. These components lie in independent slots in a basis of differentials. No pole can disappear. Thus the descent is regular everywhere by smoothness of the base, and vanishes since $H^0(\mathbf P^n,\Omega^j)=0$.

In particular all global regular one-forms on $X$ are anti-invariant. The wedge of any two is invariant and therefore zero. This is the Castelnuovo--de Franchis argument in the present setting: the rank of the differential of the Albanese map is at most one at the generic point, since its cotangent forms are the pullbacks of invariant one-forms on its Albanese torsor. In characteristic zero this rank is the dimension of its image. The map is nonconstant as $g>0$, so its image is a curve. The image and Stein factorization of the Albanese morphism are defined over $k$; write $X\to C\to\operatorname{Alb}^1_X$ for this factorization onto the image, with normal projective curve $C$ and geometrically connected generic fibre. In characteristic zero $C$ is smooth. Global regular one-forms on $C$ inject into those on $X$; conversely every regular one-form on $X$ is pulled back from the Albanese torsor and hence from $C$. Therefore the genus of $C$ is exactly $g$. Functoriality gives homomorphisms $\operatorname{Alb}_X\to J_C\to\operatorname{Alb}_X$ whose composite is the identity by the defining property of the Albanese map. The dimensions agree, so both maps are isomorphisms. No assertion about a particular principal polarization is needed.

The covering involution descends to $C$ by functoriality of the Stein factorization. It is nonidentity: otherwise $C$, of positive genus, would be dominated by $\mathbf P^n$. Its quotient is dominated by $\mathbf P^n$, so has genus zero by restriction to a geometric line. As before it has a $k$-point and is $\mathbf P^1_k$. Riemann--Hurwitz gives $2g+2$ branch points. Choose a rational coordinate with infinity off the branch locus and form the corresponding squarefree polynomial $g_0$ over $k$, retaining its scalar twist. The same degree-two field argument as in Theorem~\ref{albs:rigid} gives $k(X)=k(\mathbf P^n)k(C)$ and the stated square class. Over $\bar k$, the algebraic closure of $\bar k(C)$ in $\bar k(X)$ is $\bar k(C)$ by the Stein factorization. If an element of $\bar k(\mathbf P^n)$ is algebraic over $\bar k(r)$ it thus lies in $\bar k(C)$, and is fixed by the involution. It belongs to $\bar k(r)$. This proves geometric connectedness for the rational map $r$ as well.

For the converse, $f=g_0(r)$ with $r$ a coordinate gives a model birational to $C\times\mathbf P^{n-1}$, whose Albanese is $J_C$. Finally the trace formula and bounded-trace proof of Theorem~\ref{albs:rigid} apply to any abelian $B$ over $k$, with $p$ replaced by the residue cardinality. The good Frobenius eigenvalues again have the same positive weight. This gives the claimed precise characterization. It concerns the first Euler coefficients and the normally convergent correction of Corollary~\ref{albs:euler}; zero-free bounds on Dirichlet functions do not themselves bound the resulting projected mean at the needed scale.
\end{proof}

\begin{corollary}\label{albs:density}
In the trace classification of Theorem~\ref{albs:hyper}, the scaled error $(\mathrm N v)M_f(v)+t_v(B)$ need only be bounded in prime density. Precisely, it suffices that for each $\epsilon>0$ some finite $M$ makes the set where its absolute value exceeds $M$ have upper density at most $\epsilon$. In fact if $A$ and $B$ are abelian varieties over the same number field, then, for every fixed $M\ge0$, the prime densities of $|t_v(A)-t_v(B)|\le M$ and $t_v(A)=t_v(B)$ exist and are equal. In particular a prescribed nonzero bounded difference of these two integral traces occurs only on a set of density zero.
\end{corollary}
\begin{proof}
Take the algebraic monodromy group of the joint first-cohomology representation in characteristic zero. The algebraic Chebotarev theorem gives the density of primes when Frobenius is constrained to a conjugation-stable Zariski-closed set: it is the fraction of geometric components lying wholly in that set. Here components are counted in the Zariski closure of the actual Galois image; the standard compact-group proof uses that a proper algebraic subset has Haar measure zero on a component of this image unless it contains the component. See \href{https://doi.org/10.1007/BF02698692}{Serre, \emph{Quelques applications du th\'eor\`eme de densit\'e de Chebotarev}}. Apply the result to the trace-difference hypersurface $D=h$ for an integer $h$. Suppose one component is contained in it. The conjugacy orbit of this component has positive density, so it contains a good Frobenius $g$ up to conjugacy. Let $k$ be the order of the finite component group. Every $g^{1+ki}$ is in the same component as $g$. Thus $D(g^{1+ki})=h$ for all $i\ge0$. This constant sequence is a finite sum of complex exponentials in $i$ with bases $\alpha^k$ of modulus $(\mathrm N v)^{k/2}>1$. The Vandermonde determinant for the distinct bases together with $1$ forces $h=0$. Every nonzero integer level therefore has density zero. Taking the finite union of levels proves the equality of densities. If the errors are bounded in density, apply the first formula of Theorem~\ref{albs:rigid}: the trace difference is bounded in density as well. Its zero level must have density at least $1-\epsilon$ for each $\epsilon$, hence contains every component. The two characters agree identically and Faltings gives the isogeny over the original field.
\end{proof}

\clearpage
\section{Minimum branch degree of an arithmetic quadratic sampler}\label{cbd:section}

Retain the double cover of $\mathbf P^n$ over a number field $k$ and its smooth projective model $X$ from Theorem~\ref{albs:hyper}. Suppose its Albanese has dimension $g\ge1$. Use the smooth curve $C$ of genus $g$ obtained from the Albanese fibration, the double map $C\to\mathbf P^1$ and the induced dominant rational map
\[
                    r:\mathbf P^n\dashrightarrow\mathbf P^1 .
\]
Write $r=P/Q$ with relatively prime homogeneous forms of the same degree $e\ge1$. Let $B$ denote the reduced geometric branch divisor on $\mathbf P^n$, including the hyperplane at infinity when it occurs with odd valuation. A generic line meets $B$ in $\deg B$ distinct points.

\begin{theorem}\label{cbd:minimum}
The branch divisor satisfies
\[
                        \deg B\ge 4+2e(g-1)\ge 2g+2 .
\]
At the least possible degree $2g+2$, the map $r$ is a ratio of independent linear forms. The branch divisor consists geometrically of $2g+2$ members of one hyperplane pencil, and $X$ is birational over $k$ to $C\times\mathbf P^{n-1}$. Conversely every separable binary form of degree $2g+2$ evaluated on such a linear pair gives this type of cover and has Albanese isomorphic to the Jacobian of its curve. In particular the four hyperplanes at minimum degree for an elliptic trace must be in a single pencil. The curve occurring in the binary quartic may have no $k$-point.
\end{theorem}

\begin{proof}
We use the primitive choice of $r$ supplied by the Stein factor of the Albanese map. For clarity, after passage to $\bar k$ the field $\bar k(C)$ is relatively algebraically closed in $\bar k(X)$. If an element of $\bar k(\mathbf P^n)$ is algebraic over $\bar k(r)$, it is also algebraic over $\bar k(C)$, hence belongs to $\bar k(C)$. Taking the deck invariants then puts it in $\bar k(r)$. Thus the generic fibre of $r$ is geometrically integral in characteristic zero. In the genus-one case this primitive choice also follows directly from the Albanese torsor map. Its Stein curve would be finite over the genus-one torsor, hence have genus at least one. Pullback of differentials bounds its genus by $\dim\operatorname{Alb}(X)=1$. Over $\bar k$ its map from $X$ therefore factors back through the Albanese torsor. The resulting section of the finite curve map makes it an isomorphism.

Work over $\bar k$. The branch is nonempty: an even divisor everywhere for the square class would make the class a square because $\operatorname{Pic}(\mathbf P^n)$ has no two-torsion and constants are squares. A general line $\ell$ avoids the base locus of $P,Q$ and meets $B$ transversely at its smooth points. For $n=1$ use the whole $\mathbf P^1$. The normalized double cover $C_\ell$ of the line is connected and has genus $(\deg B-2)/2$, by Riemann--Hurwitz. The map $r|_\ell$ has degree $e$. Since the quadratic extension of the line's function field is the pullback of $\bar k(C)/\bar k(r)$, the induced map $C_\ell\to C$ also has degree $e$, by the tower formula. Applying Riemann--Hurwitz to this map gives
\[
             \deg B-4=e(2g-2)+\deg\operatorname{Ram}(C_\ell/C).
\]
This is the asserted lower bound. For $g\ge2$ and $\deg B=2g+2$ it gives $e=1$ immediately.

Now let $g=1$ and $\deg B=4$. For every such general line the map $C_\ell\to C$ between genus-one curves is étale. Every critical value of $r|_\ell$ must therefore be a branch value of $C\to\mathbf P^1$. In fact away from its four branch values the base change is étale over the line. At a point where $r|_\ell$ ramifies over any other value, its ramification index is inherited by $C_\ell\to C$, a contradiction.

For $n=1$ relative algebraic closedness already gives $e=1$. For $n\ge2$, suppose $e\ge2$. Outside finitely many values, the closures $H_t$ of the fibres of $r$ are irreducible degree-$e$ hypersurfaces by geometric integrality of the generic fibre and openness of geometric integrality. No component of a member of the pencil can be wholly contained in its base locus, which has codimension at least two. Take $t$ outside the four branch values and in this open set. At a general smooth point $x\in H_t$ outside the base locus, $dr_x\ne0$. The lines through $x$ in its tangent hyperplane have dimension $n-2$. A dense open set of these lines is not contained in $H_t$, since otherwise the tangent hyperplane itself would be contained in the irreducible nonlinear hypersurface $H_t$.

Thus the incidence of $(t,x,\ell)$ with a noncontained tangent line has dimension $1+(n-1)+(n-2)=2n-2$. Its projection to the Grassmannian of lines has finite fibres: on every such line $r$ is nonconstant and separable, and the admissible points $x$ are critical points of its restriction, a finite set, even if one removes base points elsewhere on that line. Consequently the image is dense in the $(2n-2)$-dimensional Grassmannian and contains a nonempty open there. A general line would have a critical value outside the four allowed values. This contradiction proves $e=1$.

At degree one the two forms $P,Q$ are independent and $k(\mathbf P^n)$ is purely transcendental of degree $n-1$ over $k(P/Q)$. Extending scalars to $k(C)$ gives $k(X)=k(C)(u_2,\ldots,u_n)$. The square class of the pulled-back separable binary form branches exactly along the $2g+2$ distinct hyperplanes of this pencil, each with multiplicity one. This also proves the converse and the assertion about the Albanese by birational invariance.
\end{proof}

\begin{corollary}[The least degree of a nonlinear pencil]\label{cbd:next}
For a primitive nonlinear pencil $e\ge2$ one has $n\ge2$ and
\[
                  \deg B\ge 6+2e(g-1)\ge 4g+2 .
\]
For every $n\ge2$ and $g\ge1$ the last bound is attained. In particular the least branch degree of an elliptic sampler coming from a nonlinear primitive pencil is six.
\end{corollary}
\begin{proof}
Apply the same tangent incidence argument with all the $2g+2$ branch values of $C\to\mathbf P^1$ excluded. For a general line at least one critical point of $r|_\ell$ then has value outside this set. It lifts to two distinct points on $C_\ell$, each ramified in $C_\ell\to C$. Thus the ramification term in the displayed Riemann--Hurwitz equality is at least two.

For sharpness let $H(t)$ be a squarefree polynomial of degree $2g+1$ over $k$ with nonzero roots and take
\[
             r=\frac{x_0x_1}{x_2^2},\qquad f=H(r).
\]
The finite branch components over $\bar k$ are the $2g+1$ distinct irreducible quadrics $x_0x_1-a x_2^2=0$ for the roots $a$ of $H$. The pole divisor above infinity has even multiplicity and gives no further component. On $x_2\ne0$, write $u=x_0/x_2$, so $x_1/x_2=r/u$. The remaining affine coordinates are free. Thus $r$ is primitive of degree two and $k(\mathbf P^n)$ is purely transcendental over $k(r)$. Its double cover is birational to the product of $y^2=H(t)$ with $\mathbf P^{n-1}$ and realizes branch degree $4g+2$. In genus one, any given elliptic curve over $k$ has a separable cubic model; translate its coordinate to make the roots nonzero. The example can therefore realize the trace of each individual elliptic curve.
\end{proof}

At every odd residue field $\mathbb F_q$ of good reduction the nonlinear example has an exact finite-field sum. Write $t_q(C)=q+1-\#C(\mathbb F_q)$ and extend the quadratic character $\chi$ by $\chi(0)=0$. On the chart $x_2=1$ its sum is
\[
 \sum_{\mathbf x\in\mathbb A^n(\mathbb F_q)}\chi(H(x_0x_1))
       =q^{n-2}\bigl(-(q-1)t_q(C)+q\chi(H(0))\bigr).
\]
Indeed $uv=t$ has $q-1$ ordered solutions for $t\ne0$ and $2q-1$ for $t=0$. The odd-degree model has one rational point at infinity, so $\sum_t\chi(H(t))=-t_q(C)$. The complementary projective hyperplane is the pole divisor; deleting it is a valid choice of domain for the sampler. This computes its order-$q^{-1}$ correction for that choice.

\begin{corollary}[Arithmetic of the nonlinear minimum]\label{cbd:weierstrass}
Fix the curve $C/k$ in Theorem~\ref{albs:hyper}, with $g\ge2$, split rational quotient, and $n\ge2$. Among nonlinear primitive pencils whose double covers have this Stein curve, the minimum branch degree is $4g+2$ if $C$ has a $k$-rational Weierstrass point, and $4g+4$ otherwise. In degree $4g+2$ every pencil has the shape $P/L^2$ over $k$ after a change of coordinate on its rational target, with $P$ a quadratic form and $L$ linear.
\end{corollary}
\begin{proof}
Equality at $4g+2$ forces $e=2$ by Corollary~\ref{cbd:next}, since $g-1>0$. In a quadratic pencil each of the $2g+2$ branch fibres contributes degree two to $B$ unless the quadratic is a double hyperplane, in which case it contributes zero. Distinct fibre divisors have no common component by coprimality. At most one fibre in the pencil can be a double hyperplane: two such members would exhibit a fractional-linear transform of $r$ as a square of a ratio of linear forms over $\bar k$, contradicting primitivity. At degree $4g+2$ precisely one branch value has a double hyperplane fibre. Galois invariance makes that value and the hyperplane $k$-rational. The unique point of $C$ above it is a $k$-rational Weierstrass point. Taking that value as infinity gives the required shape.

If $C$ has a rational Weierstrass point, placing it at infinity gives the odd-degree model used in Corollary~\ref{cbd:next}, with the scalar quadratic twist retained. Otherwise the degree is even and must be at least $4g+4$. Choose a rational point of the quotient away from the branch locus as infinity, and translate the remaining coordinate to make all the finite branch values nonzero. Then $C$ has an equation $y^2=H(t)$ with $H$ squarefree of degree $2g+2$. Taking the same pencil $x_0x_1/x_2^2$ gives $2g+2$ distinct branch quadrics, and the same purely transcendental extension shows that the Stein curve is indeed $C$.
\end{proof}

\begin{corollary}\label{cbd:singular}
For $n\ge2$, a geometrically integral double cover with positive-dimensional Albanese on its smooth projective model has a nonrational singularity on the normal double cover of $\mathbf P^n$. In particular a reduced branch divisor with simple normal crossings cannot occur.
\end{corollary}
\begin{proof}
Over $\bar k$ let $2b=\deg B$ and let $Y\to\mathbf P^n$ be the normal double cover. Locally it has equation $z^2=F_B$ for a local equation of the reduced branch. These hypersurfaces are normal by regularity in codimension one and the Cohen--Macaulay property. Globally the invariant and anti-invariant summands of its finite direct image are $\mathcal O$ and $\mathcal O(-b)$. This follows by gluing that local quadratic equation with the square root line bundle $\mathcal O(b)$, since $\operatorname{Pic}(\mathbf P^n)=\Z$. Hence $H^1(Y,\mathcal O_Y)=0$ by the cohomology of line bundles on projective space for $n\ge2$. If $Y$ had rational singularities, resolution $\mu:X'\to Y$ would give $\mu_*\mathcal O_{X'}=\mathcal O_Y$ and $R^i\mu_*\mathcal O_{X'}=0$ for $i>0$. It would follow that $H^1(X',\mathcal O_{X'})=0$, contradicting $\dim\operatorname{Alb}(X')=g$.

At a point on $j$ simple normal-crossing branch components the equation is étale locally $z^2=x_1\cdots x_j$, after extracting the square root of a unit. This is the quotient of a smooth coordinate space by the finite group of sign changes on $u_1,\ldots,u_j$ with product of signs $+1$. Its invariants are generated by $x_i=u_i^2$ and $z=\prod_i u_i$, with the remaining coordinates fixed. Finite quotient singularities in characteristic zero are rational by \href{https://doi.org/10.1007/BF01405091}{Boutot's theorem on rational singularities of reductive quotients}. This proves the exclusion.
\end{proof}

\clearpage
\section{Frobenius traces of a fixed elliptic curve and of its twists}\label{ep:primes}

For an elliptic curve $E/\Q$ of conductor $N$, write
\[
 \lambda_E(p)=\frac{p+1-\#E(\F_p)}{\sqrt p},\qquad
 B_E(x)=\sum_{\substack{p\le x\\p\nmid N}}\lambda_E(p)\log p .
\]
The first expression concerns good primes only. Each summand is accessible by counting points in a finite field. The following criterion concerns their signed sum over rational primes. The size $\sqrt{x}$ here is a cancellation estimate across different primes. Hasse's theorem gives the distinct local estimate $|\lambda_E(p)|\le2$.

\begin{theorem}[The sharp prime criterion]\label{ep:equiv}
The following conditions on a fixed $E/\Q$ are equivalent.
\begin{enumerate}
\item Every nontrivial zero of the arithmetic $L(E,s)$ has real part $1$.
\item For every $\epsilon>0$, $B_E(x)=O_{E,\epsilon}(x^{1/2+\epsilon})$ for $x\ge2$.
\item For every $\epsilon>0$,
$\sum_{p\le x,\ p\nmid N}(p+1-\#E(\F_p))\log p
=O_{E,\epsilon}(x^{1+\epsilon})$ for $x\ge2$.
\end{enumerate}
Under the first condition there is an absolute implied constant in
\begin{equation}\label{ep:sharp}
 |B_E(x)|\ll \sqrt x\log(2x)\log(2Nx).
\end{equation}
Let $w\in C_c^\infty((0,1))$ be fixed. If the first condition holds, $x\ge2$ and $0<h\le x$, then
\begin{equation}\label{ep:smooth}
 \left|\sum_{p\nmid N}\lambda_E(p)\log p\,w((p-x)/h)\right|
 \ll_w\sqrt x\log(2Nx/h).
\end{equation}
\end{theorem}

\subsection*{The logarithmic derivative and its prime powers}
Use the unitary normalization $L_u(s)=L(E,s+1/2)$. At a good prime the two inverse roots $\alpha_p,\beta_p$ have absolute value one, product one, and sum $\lambda_E(p)$. At a multiplicative prime the single nonzero inverse root is $\pm p^{-1/2}$. At an additive prime there are none. Sum the $k$th powers of the local inverse roots to form $b(p,k)$, with empty sum zero. For $\operatorname{Re}s>1$,
\[
 -L'_u(s)/L_u(s)=\sum_{p,k\ge1} b(p,k)\log p\ p^{-ks}.
\]
Here and below the double sum runs over primes $p$ and positive integers $k$. The bound $|b(p,k)|\le2$ gives absolute convergence. By modularity, $L_u$ is entire and its completion is of order one and satisfies the functional equation with symmetry $s\leftrightarrow1-s$. The gamma factor in this normalization is a constant times $(2\pi)^{-s}\Gamma(s+1/2)$. The completed zeros are precisely the nontrivial zeros.

Put $A_E(x)=\sum_{p^k\le x}b(p,k)\log p$. One has, uniformly in the curve,
\begin{equation}\label{ep:powers}
 A_E(x)=B_E(x)+O(\sqrt x).
\end{equation}
Indeed, the multiplicative primes with $k=1$ contribute at most $\sum_{p\le x}\log p/\sqrt p\ll\sqrt x$, by the upper bound $\theta(t)\ll t$ and partial summation. The squares cost $O(\sqrt x)$ by the same upper bound. Each remaining exponent costs $O(x^{1/3})$ and at most $\log x/\log 2$ exponents occur. These estimates also bound the total absolute value of all omitted terms up to $2x$. This observation will permit us to restrict to good primes even in the short interval.

\subsection*{Two cutoffs in the explicit formula}
For completeness we establish both uniform estimates. We use the standard unit-interval zero count and local partial-fraction formula for an entire primitive cuspidal $L$-function of fixed degree and fixed archimedean type. For this degree-two function, if $t$ is real then
\[
 \#\{\rho:\ |\operatorname{Im}\rho-t|\le1\}
       \ll\log(2N(|t|+2)).
\]
In $1/4\le\operatorname{Re}s\le2$ its logarithmic derivative equals the sum of $1/(s-\rho)$ over these zeros plus $O(\log(2N(|t|+2)))$, away from the poles of the two sides. These standard forms follow by logarithmically differentiating the Hadamard product, comparing at $2+it$, and applying Stirling's formula. The absolute convergence at $2+it$ is uniform by $|b(p,k)|\le2$. The positive real parts of $1/(2+it-\rho)$ majorize a constant times $(1+|t-\operatorname{Im}\rho|^2)^{-1}$, which gives the stated local count. Subtracting the two partial fractions makes the contribution of each distant zero $O((1+|t-\operatorname{Im}\rho|^2)^{-1})$. For the general analytic input see Iwaniec and Kowalski, \href{https://bookstore.ams.org/coll-53}{\emph{Analytic Number Theory}}, Chapter~5.

Suppose henceforth in this subsection that all completed zeros have real part $1/2$. The line $\operatorname{Re}s=1/4$ stays at distance at least $1/4$ from them, so $|L'_u/L_u|\ll\log(2N(|t|+2))$ there. Put $y=h/x$, and define
\[
 v_y(t)=w((t-1)/y),\qquad
 V_y(s)=y\int_0^1w(a)(1+ya)^{s-1}\,da .
\]
For the sharp cutoff choose a nonnegative smooth $w_0$ of integral one supported in $(0,1)$ and put
\[
 b_y(t)=\int_0^1w_0(a){\bf1}_{t\le1+ya}\,da,\qquad
 H_y(s)=s^{-1}\int_0^1w_0(a)(1+ya)^s\,da .
\]
The latter two transforms $V_y$ and $H_y$ are the Mellin transforms of $v_y$ and $b_y$, respectively. Throughout the closed strip $1/4\le\sigma\le2$, and for every integer $j\ge0$, they satisfy
\begin{align*}
 |V_y(\sigma+it)|&\ll_{j,w} y(1+y|t|)^{-j},\\
 |H_y(\sigma+it)|&\ll_{j,w_0}(1+|t|)^{-1}(1+y|t|)^{-j}.
\end{align*}
For $y|t|\le1$ these follow from the defining integrals. Otherwise integrate $j$ times by parts against $\exp(it\log(1+ya))$, each time gaining $(y|t|)^{-1}$ since $1\le1+ya\le2$ and the differentiated functions vanish at the endpoints.

For either pair $(v,V)$, inversion at $\operatorname{Re}s=2$ and a contour shift give the exact identity
\begin{equation}\label{ep:explicit}
 \sum_{p,k}b(p,k)\log p\,v(p^k/x)
   =-\sum_{\rho}x^\rho V(\rho)
    +\frac1{2\pi i}\int_{(1/4)}
       -\frac{L'_u}{L_u}(s)x^s V(s)\,ds .
\end{equation}
Both terms on the right converge absolutely. Indeed the unit-interval zero count permits the horizontal segments to be placed at positive and negative heights of size $T$ whose distance from all ordinates is $\gg1/\log(2N(T+2))$. The local partial fractions bound the logarithmic derivative by $O(\log^2(2N(T+2)))$ there. Arbitrary power decay of the transforms at fixed $y>0$ makes these horizontal contributions tend to zero. The residues have the stated negative sign. There is no pole of $L_u$ and there are no gamma-factor poles in the strip of integration.

Let $U=1/y$. Below $|t|=U$, the sums over zeros of the bounds for $V_y$ and $H_y$ are at most $O_w(\log(2NU))$ and $O_{w_0}(\log(2U)\log(2NU))$, respectively. Here $V_y$ contributes a factor $U^{-1}$, while $H_y$ costs the harmonic sum up to $U$. Above $U$, divide the ordinates into bands $2^jU<|t|\le2^{j+1}U$. Taking the transform bounds with exponent $2$ makes the tail for either kernel $O(\sum_{j\ge0}2^{-j}(\log(2NU)+j+1))$. Integrating the corresponding majorants against $\log(2N(|t|+2))\,dt$ gives the same two estimates. The zero term in \eqref{ep:explicit} multiplies them by $\sqrt x$, the vertical integral by $x^{1/4}$.

For $v_y$ this gives \eqref{ep:smooth} after applying the absolute omission bound following \eqref{ep:powers}. For the sharp sum set $y=1/x$ in $b_y$. The smoothed sum differs from $A_E(x)$ only on $(x,x+1]$, at at most one integer, which costs $O(\log(2x))$. Equation~\eqref{ep:powers} proves \eqref{ep:sharp}.

\subsection*{The reverse implication and a secondary term}
If $B_E(x)\ll_{E,\epsilon}x^{1/2+\epsilon}$ for every $\epsilon>0$, then the same statement holds for $A_E$ by \eqref{ep:powers}. Partial summation on $\operatorname{Re}s>1$ yields
\[
 -L'_u(s)/L_u(s)=s\int_1^\infty A_E(t)t^{-s-1}\,dt .
\]
On any compact subset of $\operatorname{Re}s>1/2$, choose $\epsilon$ smaller than its distance to the boundary. The integral converges locally uniformly there and is holomorphic. The identity theorem for meromorphic functions rules out a zero of $L_u$ in this open half-plane, since its logarithmic derivative would have a pole with nonzero residue. The functional equation rules out completed zeros in the opposite half of the strip. Finally partial summation first with the factor $\sqrt p$, then with $1/\sqrt p$, proves the equivalence of the second and third conditions of Theorem~\ref{ep:equiv}. Specifically the two integrals are bounded by constants times $\int_2^x t^{\epsilon}\,dt$ and $\int_2^x t^{-1/2+\epsilon}\,dt$, respectively.

For a fixed non-CM curve one can determine the leading term in the correction \eqref{ep:powers} unconditionally. The Sato--Tate theorem gives limiting trace distribution $(2\pi)^{-1}\sqrt{4-t^2}\,dt$ on $[-2,2]$, with first moment $0$ and second moment $1$. The prime number theorem and partial summation give $B_E(x)=o_E(x)$ and $\sum_{p\le y,\ p\nmid N}\lambda_E(p)^2\log p\sim y$. At a good prime, $b(p,2)=\lambda_E(p)^2-2$. At the fixed finitely many bad primes the total of all powers is $O_E(\log x)$, and the good powers of exponent at least three cost $O(x^{1/3}\log x)$. Consequently
\begin{equation}\label{ep:bias}
 A_E(x)=B_E(x)-\sqrt x+o_E(\sqrt x).
\end{equation}
This signed term has exactly the power of $x$ in the zero criterion. It is present before any conjecture on zeros is invoked. For Sato--Tate in the stated generality see Barnet-Lamb, Geraghty, Harris and Taylor, \href{https://doi.org/10.2977/PRIMS/31}{\emph{A family of Calabi--Yau varieties and potential automorphy~II}}.

The qualitative trace law by itself imposes no square-root estimate. To see the logical distinction, fix $1/2<\beta<1$ and take $X_j=4^j$. In each interval $[X_j,2X_j]$, Sato--Tate and the prime number theorem supply $\gg X_j/\log X_j$ primes having $\lambda_E(p)<0$. At $\lfloor X_j^\beta/\log X_j\rfloor$ such primes replace the trace by $2$ and keep every other trace. The altered sequence has the same limiting distribution. In fact the modified primes up to $x$ number $O(x^\beta)$, which is $o(\pi(x))$. Its sum up to $2X_j$ exceeds the old sum by $\gg X_j^\beta$, since all changes have positive sign. The old and modified sequences cannot both satisfy condition~2 when $\epsilon<\beta-1/2$. This is a construction of bounded sequences; it makes no assertion that the altered sequence is attached to a curve.

\subsection*{A maximal estimate across quadratic twists}
Let $d$ range over signed fundamental discriminants with $|d|\le D$, let $E_d$ be the quadratic twist by $\chi_d=(d/\cdot)$, and put
\[
 T_d(x)=\sum_{\substack{p\le x\\p\nmid2N}}\lambda_E(p)\chi_d(p)\log p .
\]
At $p\nmid2N$ with $p\nmid d$, the two curves' traces differ by the factor $\chi_d(p)$. When $p\nmid2N$ divides $d$, the twist of the unramified two-dimensional local representation by the ramified quadratic character has no inertia invariants. Its local Euler factor is one. Hence $T_d(x)$ differs from $B_{E_d}(x)$ only at primes dividing $2N$, with total difference at most $2\log(2N)$.

\begin{theorem}[Unconditional family estimate]\label{ep:sieve}
For $D,X\ge2$ and every $\epsilon>0$,
\[
 \sum_{|d|\le D}^{\rm fund}\ \,\max_{2\le y\le X}|T_d(y)|^2
 \ll_\epsilon (D+X)(DX)^\epsilon X\log^3(2X).
\]
In particular, fix $0<\alpha,\eta\le1$ and an elliptic curve $E$. Apart from at most $O_{\alpha,\eta}(D^{1-\alpha\eta/2})$ fundamental discriminants of absolute value at most $D$, one has, simultaneously for $D^\alpha\le y\le D$, $y\ge2$,
\[
 |B_{E_d}(y)|\ll_E y^{1/2+\eta}.
\]
\end{theorem}
\begin{proof}
Heath-Brown's quadratic large sieve for positive odd squarefree integers states that
\[
 \sum_{m\le M}^{*}\left|\sum_{n\le H}^{*}
        c_n\left(\frac{n}{m}\right)\right|^2
 \ll_\epsilon (MH)^\epsilon(M+H)\sum_{n\le H}^{*}|c_n|^2 .
\]
This is Theorem~1 in \href{https://doi.org/10.4064/aa-72-3-235-275}{\emph{A mean value estimate for real character sums}}. The same bound applies to the characters $(d/n)$ above. Indeed write $d=u m$ with $m$ odd positive squarefree and $u$ one of $\pm1,\pm4,\pm8$. Split $m$ by residue modulo $4$. Quadratic reciprocity converts $(m/n)$ to $(n/m)$ times a factor depending on $n$ only within each group. The factor $(u/n)$ is likewise absorbed into $c_n$. Dropping restrictions on $m$ then increases the sum of squares. There are boundedly many groups.

Take $c_p=\lambda_E(p)\log p$ on allowed primes and zero elsewhere. By Chebyshev, $\sum_{p\le X}|c_p|^2\ll X\log(2X)$. For the maximum, embed $[1,X]$ in $[1,2^J]$ with $2^{J-1}<X\le2^J$, and extend the coefficients by zero. Each initial interval decomposes into at most $J+1$ standard dyadic blocks, at most one of each length. Cauchy--Schwarz bounds its squared sum by $(J+1)$ times the sum of squared block sums over all blocks. Apply the large sieve to each block while keeping $H=2^J$. At each fixed length the blocks partition the support. Their squared coefficient masses add to the original mass. The $J+1$ levels therefore cost a second factor $J+1$. This proves the maximal estimate.

Cover $[D^\alpha,D]$ by $O(\log(2D))$ dyadic intervals $(X/2,X]$, with $D^\alpha\le X\le2D$. On any interval containing some $y\ge D^\alpha$ for which $|T_d(y)|>y^{1/2+\eta}$, the maximum exceeds $\sqrt{X/2}\,D^{\alpha\eta}$. Markov's inequality and the first assertion bound the number of exceptions on that interval by $O_\epsilon(D^{1+2\epsilon-2\alpha\eta}\log^3(2D))$. Take $\epsilon=\alpha\eta/4$ and absorb the logarithms after summing the intervals. This implies the displayed, weaker exceptional-set exponent. Adding back the primes over $2N$ proves the result.
\end{proof}

For $X\le D$ the root mean square of the maxima in the first assertion is $\ll_\epsilon \sqrt X(DX)^{\epsilon/2}\log^{3/2}(2X)$, since there are $\asymp D$ fundamental discriminants. The curve and the observation window vary together in the second assertion. A fixed discriminant need not occur outside the exceptional set for any unbounded sequence of $D$. In addition, for fixed $d$ the lower end $D^\alpha$ goes to infinity. The quantifiers in Theorem~\ref{ep:equiv} would require estimates at every scale for the same curve.

\subsection*{Zeros unconditionally on the line}
There are two distinct unconditional statements about actual zeros. First, for every $E/\Q$ the proportion of positive-height zeros on the critical line, in the sense of the lower limit as the height bound goes to infinity, is greater than $0.0696$. Indeed its form is a holomorphic Hecke newform on $\Gamma_0(N)$ with trivial nebentypus. Apply Theorem~1.1 of Andersen and Thorner in the unitary normalization and translate by $1/2$. This statement uses their unconditional percentage. The stronger percentage in the same theorem assumes the generalized Ramanujan conjecture for auxiliary forms occurring in their method. See \href{https://arxiv.org/abs/2004.03581}{Andersen--Thorner, \emph{Zeros of $\mathrm{GL}_2$ $L$-functions on the critical line}}.

Second, positive proportions of quadratic twists of any rational elliptic curve with a rational $3$-isogeny have analytic rank zero and analytic rank one, respectively. The two proportions are relative to the set of fundamental discriminants ordered by absolute value. This is Theorem~1.5 of Kriz and Li, \href{https://doi.org/10.1017/fms.2019.9}{\emph{Goldfeld's conjecture and congruences between Heegner points}}. For a non-CM example take $y^2+xy+y=x^3$. The tangent $y=0$ at $(0,0)$ meets the curve there with multiplicity three, so the point generates a rational group of order three. The discriminant is $-26$, and $j=12167/26$ is nonintegral, which excludes complex multiplication. Hence a positive proportion of this family's $L$-functions vanish simply at the arithmetic center and a positive proportion are nonzero there. The theorem concerns the order at the single point $s=1$. Neither this multiplicity nor the critical-line proportion specifies the real part of every other zero.

\begin{corollary}[Distances to the line for fixed degree-two functions]\label{ep:density}
Let $\pi$ be a fixed cuspidal representation of $\mathrm{GL}_2/\Q$ with unitary central character. For its nontrivial zeros $\beta+i\gamma$ in unitary normalization, put $v=|\beta-1/2|$. For every real $k>0$ and $T\ge3$ one has unconditionally
\[
 \sum_{0<\gamma\le T} v^k\ll_{\pi,k} T(\log T)^{1-k}.
\]
For $w\ge0$, the proportion of positive-ordinate zeros through $T$ with $v\ge w/\log T$ is $O_\pi(e^{-w/8})$ as $T\to\infty$, uniformly in $w$. For elliptic curves in arithmetic coordinates use $v=|\beta-1|$. All counts include multiplicity.
\end{corollary}
\begin{proof}
Take $c=1/8$ in the Selberg-type density estimate of \href{https://arxiv.org/pdf/2103.01956}{Beckwith--Liu--Thorner--Zaharescu, Theorem 1.1}. This choice is admissible unconditionally since $1/8<1/4-(7/64)/2$. For $u\ge0$ it gives
\[
 \#\{\rho:\ |\gamma|\le T,\ \beta\ge1/2+u\}
       \ll_\pi T\log T\,e^{-u\log T/8}.
\]
The functional equation and complex conjugation of the dual identify zeros of $\pi$ at $\beta+i\gamma$ and $1-\beta+i\gamma$ with equal multiplicity. Thus twice the displayed bound controls both sides at positive height. Divide by the standard positive-ordinate count $N_\pi^+(T)\sim(T/\pi)\log T$ to obtain the tail bound. Integrating the same inequality against $k u^{k-1}\,du$ gives the moment estimate since the integral from zero to infinity is $\Gamma(k+1)(8/\log T)^k$. The implied constants may depend on the fixed representation and do not assert uniformity as its conductor varies.
\end{proof}

\section{A Mellin certificate for the odd-sign curve of conductor 37}\label{ep:37}
Take $E:y^2+y=x^3-x$, the curve $37.\mathrm{a}1$. Write
\[
 \Lambda_E(s)=\left(\frac{\sqrt{37}}{2\pi}\right)^s
       \Gamma(s)L(E,s)
\]
in arithmetic normalization. Here is an individual unconditional result, for a non-CM degree-two function with an actual zero. All real zeros of the completion are covered, irrespective of their real part.

\begin{theorem}\label{ep:certificate}
The sole real zero of $\Lambda_E$ is $s=1$, of multiplicity one. The same point is the only zero, counting multiplicity, in the closed rectangle
\[
 \frac12\le\operatorname{Re}s\le\frac32,\qquad
 |\operatorname{Im}s|\le\frac92 .
\]
\end{theorem}

The argument uses one Fourier expansion and the odd functional equation. Positivity of its real Mellin kernel will handle the real axis. To count the complex zeros, we bound the sign of the imaginary part throughout the rectangle, including its vertical edges. The rational polynomial used to certify this sign has small enough degree that all its coefficients in a positive basis can be checked individually. For identification the curve and its modular form are recorded at \href{https://www.lmfdb.org/EllipticCurve/Q/37/a/1}{LMFDB, $37.\mathrm{a}1$}. The numerical zeros displayed by databases are not needed in the certificate.

\subsection*{The kernel and the real-axis theorem}
The invariants are $\Delta=37$, $c_4=48$, $c_6=-216$. This integral equation is minimal since no twelfth power of a prime divides the discriminant. Thus there is multiplicative reduction at $37$ and good reduction elsewhere. Since $(216/37)=(6/37)=-1$, the multiplicative reduction is nonsplit. The finite root number at $37$ is $+1$, and the real place contributes $-1$. The conductor is $37$ and the global root number is $-1$. Also $j=110592/37$ is nonintegral, so $E$ has no complex multiplication.

Let $f$ denote the weight-two newform supplied by modularity. Put
\[
 a=2\pi/\sqrt{37},\quad r=e^{-a},\quad
 F(u)=f(iu/\sqrt{37})=\sum_{n\ge1}a_n e^{-anu}.
\]
The coefficients $a_n$ are the arithmetic coefficients, with $a_1=1$ and $a_p=p+1-\#\widetilde E(\F_p)$, where the point count for the singular cubic at $37$ includes the singular point and the point at infinity. At a good prime the prime-power recursion
$a_{p^k}=a_pa_{p^{k-1}}-pa_{p^{k-2}}$ applies for $k\ge2$. At $37$ it reads $a_{37^k}=a_{37}^k$. Together with multiplicativity these determine every $a_n$. Hasse's theorem in each good Euler factor, followed by multiplication, gives
\[
 |a_n|\le d(n)\sqrt n\le2n,
\]
since divisors pair across $\sqrt n$. The Fricke involution with this root number gives $F(1/u)=-u^2F(u)$ and $F(1)=0$. Mellin integration first in the half-plane of absolute convergence, followed by inversion $u\mapsto1/u$ below one, gives the entire identity
\begin{equation}\label{ep:mellin}
 \Lambda_E(1+z)=2\int_1^\infty F(u)\sinh(z\log u)\,du .
\end{equation}
Exponential decay for $u\ge1$ justifies analytic continuation and differentiation on every compact subset of the $z$-plane.

\begin{lemma}\label{ep:positive}
For all real $u>1$ one has $0<F(u)\le e^{-au}$.
\end{lemma}
\begin{proof}
Write $H(t)=\sum_{n\ge1}a_nt^{n-1}$. Then $F(u)=e^{-au}H(e^{-au})$, $H(r)=0$ and $H(0)=1$. Direct counting and the Hecke recursion give
\[
 (a_1,\ldots,a_{10})=(1,-2,-3,2,-2,6,-1,0,6,4).
\]
We have $a>1$ and $r<3/8$. Write $r_0=3/8$. For $0\le t\le r$,
\[
 H'(t)\le -2-6t+t\left[
 \sum_{n=4}^{10}(n-1)\max(a_n,0)r_0^{\,n-3}
 +\frac{220r_0^8}{1-(12/10)r_0}\right].
\]
For the tail $n\ge11$ use $2n(n-1)r_0^{n-3}$ inside the bracket. Its consecutive ratios are at most $(12/10)r_0$. The finite sum and the geometric bound in brackets add to less than $5$, a rational inequality. Hence $H'(t)\le-2$ and $0<H(t)\le1$ for $0\le t<r$. This proves the claim.
\end{proof}
In \eqref{ep:mellin} the sign for nonzero real $z$ is now the sign of $z$. At zero the first derivative is $2\int_1^\infty F(u)\log u\,du>0$. This proves the assertion about every real zero, independently of the computation below. In particular it proves analytic rank one for this curve without a rank computation.

\subsection*{From a rectangle to finitely many inequalities}
Put $T=9/2$, $z=\sigma+i\tau$, with $|\sigma|\le1/2$ and $0<\tau\le T$. By \eqref{ep:mellin},
\begin{equation}\label{ep:imag}
 \frac{\operatorname{Im}\Lambda_E(1+z)}{2\tau}
 =\int_1^\infty F(u)\cosh(\sigma\log u)
                  \frac{\sin(\tau\log u)}{\tau}\,du .
\end{equation}
This extends continuously to $\tau=0$. Introduce $Y=(\tau/T)^2$, $V=(2\sigma)^2$ and moments
\[
 D_m=\int_1^\infty F(u)\frac{(10\log u)^{2m+1}}{(2m+1)!}\,du.
\]
On the unit square the right side of \eqref{ep:imag} differs uniformly from the polynomial $P$ below by at most $1575/68719476736$:
\begin{align*}
 P(Y,V)&=\sum_{j+k\le21}c_{jk}Y^jV^k,\\
 c_{jk}&=(-1)^j\binom{2j+2k+1}{2j+1}
 \frac{T^{2j}2^{-2k}}{10^{2j+2k+1}}D_{j+k}.
\end{align*}
To see this, expand the sine and the hyperbolic cosine at zero and multiply. For every $m$ the sum of the absolute bounds on the terms with $j+k=m$, before integration against $F(u)$, is at most
\[
 \frac1T\frac{((T+1/2)\log u)^{2m+1}}{(2m+1)!}.
\]
This follows from the binomial formula after majorizing $|\tau|$ by $T$ and $|\sigma|$ by $1/2$. For degree $m\ge22$ extract from each term $((T+1/2)/10)^{2m+1}\le2^{-45}$. The remaining terms sum to at most $u^{10}$. Because $F(u)\le e^{-au}\le e^{-u}$, their integral is at most $10!$. Thus the discarded tail is bounded by
\begin{equation}\label{ep:series-tail}
 \frac{10!}{T\,2^{45}}=\frac{1575}{68719476736}<10^{-7}.
\end{equation}
The same majorant justifies integration of the entire series. Notice the error is absolute on the full square, even as $\tau$ approaches zero.

For $0\le i\le m$ put $\mathcal B_{i,m}(x)=\binom mi x^i(1-x)^{m-i}$. These polynomials are nonnegative on $[0,1]$ and sum to one. The binomial identity
\[
 x^j=\sum_{i=j}^m\frac{\binom{i}{j}}{\binom{m}{j}}\mathcal B_{i,m}(x)
\]
follows by factoring out $x^j$ and applying the binomial theorem of degree $m-j$. Apply it separately to $Y$ and $V$ with $m=21$. Then
\begin{align}
 P(Y,V)&=\sum_{r,s=0}^{21} b_{rs}\mathcal B_{r,21}(Y)
                           \mathcal B_{s,21}(V),\nonumber\\
 b_{rs}&=\sum_{\substack{j\le r,\ k\le s\\j+k\le21}}c_{jk}
 \frac{\binom{r}{j}\binom{s}{k}}{\binom{21}{j}\binom{21}{k}}.
 \label{ep:bernstein}
\end{align}
It suffices to enclose the 22 moments and check 484 rational lower endpoints. The enclosure method and the computed minimum appear next.

\subsection*{Validated integration of the moments}
Only coefficients through $a_{60}$ are needed. Set $F_{60}(u)=\sum_{n\le60}a_ne^{-anu}$. Truncate its moment integrals to $[1,65]$. There are three analytic remainders.
\begin{enumerate}
\item For $u\ge1$, $e^{-au}\le r_0$ and $e^{-au}\le e^{-u}$. Summing $2n(e^{-au})^{n-1}$ over $n\ge61$ shows $|F(u)-F_{60}(u)|\le 200r_0^{60}e^{-u}$. Since $|(10\log u)^k/k!|\le u^{10}$, the total replacement error on the finite interval is at most $200r_0^{60}10!<2^{-48}$.
\item The original integral from $65$ to infinity costs at most
\[
 \int_{65}^\infty e^{-u}u^{10}\,du
 \le \frac{65^{10}e^{-65}}{1-10/65}<2^{-31}.
\]
The first bound follows by writing $u=65+v$ and using $(1+v/65)^{10}\le e^{10v/65}$. The last inequality can be checked with the first 170 terms of the exponential series at $65$.
\item Divide $[1,65]$ into panels of length $h=1/32$. On each panel interpolate the normalized integrand
\[
 G_m(w)=F_{60}(w)\frac{(10\log w)^{2m+1}}{(2m+1)!}
\]
at 28 equally spaced nodes including its endpoints, and integrate the interpolation polynomial exactly. For each $u\in[1,65]$, the disk $|w-u|\le1/2$ is contained in the right half-plane, with $|\arg w|\le\pi/6$, $1/2\le |w|\le65.5$, and $|\log w|<5$ for the principal logarithm. There
\[
 |F_{60}(w)|\le\sum_{n\ge1}2ne^{-an/2}<12,
 \qquad |G_m(w)|<12e^{50}<2^{77}.
\]
Here $e^{-a/2}<2/3$ since $a>1$. Cauchy's derivative bound gives $|G_m^{(28)}(u)|/28!\le2^{77}2^{28}$. The real Lagrange remainder is at most this bound times the product of the 28 distances from the nodes. Each distance is at most $h$, so integration over all panels costs at most $64\cdot2^{77}/16^{28}<2^{-28}$.
\end{enumerate}
The three analytic error bounds sum to less than $2^{-27}$, independently of $m\le21$.

The remaining arithmetic is recorded so that the check can be reproduced without a zero-finding library. Endpoints of intervals are integers in units of $10^{-36}$. Products use all sign-relevant endpoints, rounding down for the lower end and up for the upper end. The quadrature weights are rational. On the node set $j/27$, $0\le j\le27$, the weight of node $j$ is
\[
 \frac1{27}\int_0^{27}\prod_{0\le k\le27,\ k\ne j}
                   \frac{t-k}{j-k}\,dt .
\]
The checker verifies each monomial integral of degree at most 27 exactly. Consecutive panels share their endpoint. This gives $64\cdot32\cdot27+1=55297$ distinct nodes.

Evaluate $\pi$ with Machin's identity $\pi=16\arctan(1/5)-4\arctan(1/239)$ and the alternating-series error after 48 terms. The square root of $37$ is enclosed by two consecutive integers after scaling. At each node compute $q=e^{-au}$ and $\log u$ by correctly rounded decimal exponentials and logarithms with 74 digits of precision, with each endpoint widened by a further unit of $10^{-36}$. The arguments to these functions are exact decimals for the interval endpoints. Their outputs here have magnitude below $6$, so a unit at that scale dominates the rounding errors. Horner evaluation gives $F_{60}=q(a_1+q(a_2+\cdots+qa_{60}))$. For the normalized power at the next $m$, multiply by $(10\log u)^2/((2m+2)(2m+3))$ with outward rounding. Sum all weighted node intervals, allowing for negative weights, then enlarge each resulting interval by $2^{-27}$ in both directions. Decimal implementations used for this audit specify correctly rounded exponentials and logarithms under round-to-nearest with ties to even. See the \href{https://docs.python.org/3/library/decimal.html}{Python decimal specification}.

For every $c_{jk}$ use the lower endpoint of $D_{j+k}$ when $j$ is even and its upper endpoint when $j$ is odd. All remaining factors in \eqref{ep:bernstein} are positive rationals. The exact rational check, with complete moment intervals in the accompanying files, gives
\[
 \min_{0\le r,s\le21}\ \underline{b}_{rs}
   >\frac{218956871}{10^{12}}>\frac1{5000},
 \qquad \text{minimum attained at }(r,s)=(21,21).
\]
Underlining denotes the computed lower bound. The program \texttt{check\_37a\_moments.py} regenerates the finite-field coefficients, the integration bounds and this minimum from the Weierstrass equation. Its saved intervals are supplementary data and are not accepted as inputs by the checker.

Every value of $P$ is a convex combination of its Bernstein coefficients. Subtracting \eqref{ep:series-tail} leaves a strictly positive lower bound throughout the square. Equation~\eqref{ep:imag} therefore gives $\operatorname{Im}\Lambda_E(1+\sigma+i\tau)>0$ whenever $0<\tau\le9/2$ and $|\sigma|\le1/2$. Complex conjugation gives strict negativity below the axis. Lemma~\ref{ep:positive} and \eqref{ep:mellin} handle the remaining axis, proving Theorem~\ref{ep:certificate}.

At larger heights the sine in \eqref{ep:imag} oscillates more rapidly. The positivity test is specific to the stated $T$; it supplies an all-points assertion on a finite region of the critical strip and the separate real-axis assertion at all real parameters.

\subsection*{Counting four further conjugate pairs}
A second component of the same integral reaches beyond the first noncentral zero. On the imaginary axis this component is the derivative of the function whose sign changed. It therefore tests both the position and the multiplicity of a zero.

\begin{theorem}\label{ep:ten}
For $E:y^2+y=x^3-x$, the completed function $\Lambda_E$ has exactly nine zeros, counted with multiplicity, in
\[
 1/2\le\operatorname{Re}s\le3/2,\qquad |\operatorname{Im}s|\le10 .
\]
All are simple and have real part $1$. They consist of the central zero, one zero at $1+i\tau$ in each of the four following ordinate intervals, and their conjugates:
\[
 (5.0031,5.0032),\quad(6.8703,6.8704),\quad
 (8.0143,8.0144),\quad(9.9328,9.9334).
\]
All displayed decimal endpoints are exact rationals.
\end{theorem}
\begin{proof}
In addition to $I=\operatorname{Im}\Lambda_E(1+\sigma+i\tau)/(2\tau)$ in \eqref{ep:imag}, use
\[
 R(\sigma,\tau)=\frac{\operatorname{Re}\Lambda_E(1+\sigma+i\tau)}{2\sigma}
 =\int_1^\infty F(u)\frac{\sinh(\sigma\log u)}{\sigma}\cos(\tau\log u)\,du ,
\]
taking the continuous extension at $\sigma=0$. Real and imaginary parts of the integral prove both identities. A nonzero sign of $I$ excludes zeros wherever $\tau>0$. A nonzero sign of $R$ excludes zeros wherever $\sigma\ne0$. Moreover, if $J(\tau)=\operatorname{Im}\Lambda_E(1+i\tau)$ then
\begin{equation}\label{ep:derivative}
 J'(\tau)=2R(0,\tau).
\end{equation}
Thus on a band where $R$ has strict sign, $J$ is strictly monotone and any zero it has is simple: the derivative along the axis is already nonzero. A sign change brackets its unique zero. Along that axis the real part vanishes by the odd Fricke relation.

Set $T=10$, $Y=(\tau/T)^2$, $V=(2\sigma)^2$. For $I$ use the previous polynomial with total degree cutoff $49$. For $R$ only its binomial coefficient changes. For $j+k\le49$ the coefficients of the two polynomials are
\[
 c^{(C)}_{jk}=(-1)^j
 \begin{cases}
 \binom{2j+2k+1}{2j+1},&C=I,\\[2pt]
 \binom{2j+2k+1}{2j},&C=R
 \end{cases}
 \frac{10^{2j}2^{-2k}}{10^{2j+2k+1}}D_{j+k}.
\]
For $R$ the factor $(\log u)^{2k+1}/(2k+1)!$ comes from $\sinh(\sigma\log u)/\sigma$ instead of the cosine. Thus the same 50 moments $D_0,\ldots,D_{49}$ serve both polynomials. Majorize the omitted terms by the binomial theorem as before. For this cutoff take the exponential majorant $u^{30}$. The ratio of the two power-series parameters is $(T+1/2)/30=21/60$, and the lowest discarded ordinary degree is $101$. The analogues of the prefactor $1/T$ are $1/T$ for $I$ and $2$ for $R$. Hence both remainders on the full rectangle are bounded by
\[
 \mathcal E=2\cdot30!\left(\frac{21}{60}\right)^{101}<10^{-13}.
\]

For these moments use 90 Fourier coefficients, integrate up to $96$, and put 36 evenly spaced nodes in each panel of length $1/32$. The coefficient replacement error is at most $300(3/8)^{90}10!<2^{-85}$ by summing $2n(3/8)^{n-1}$ from $n=91$. The original integral past $96$ costs at most
$96^{10}e^{-96}/(1-10/96)<2^{-70}$. This last inequality follows, for example, from the first 230 terms of the series for $e^{96}$. On $|w-u|\le1/2$, $1\le u\le96$, the bound on the principal logarithm becomes $|\log w|<5.2$. The other factor still satisfies $|F_{90}(w)|<12$. We therefore have $|G_m(w)|<12e^{52}<2^{80}$, uniformly in $m$. The Lagrange error of the 36-node rule over the whole interval is at most $95\cdot2^{80}/16^{36}<2^{-57}$. All analytic errors together are strictly less than $2^{-56}$ per moment. The same outward-rounded fixed-point quadrature now uses $95\cdot32\cdot35+1=106401$ distinct nodes.

Convert coefficients to the Bernstein basis of degree 49 in each variable by the binomial identity used in \eqref{ep:bernstein}. To bound each polynomial on a closed height band $[a,b]$, restrict its $Y$ interval to $[(a/10)^2,(b/10)^2]$; the $V$ axis still covers $[0,1]$. The restriction operation itself can be checked as follows. For a univariate polynomial expressed with Bernstein coefficients $c_0,\ldots,c_n$, set $c^{(0)}_i=c_i$ and iterate
\[
 c^{(j+1)}_i=(1-v)c^{(j)}_i+v c^{(j)}_{i+1}.
\]
The coefficient at index $j$ of its restriction to $[0,v]$ is $c^{(j)}_0$. The coefficient at index $j$ of its restriction to $[v,1]$ is $c^{(n-j)}_j$. For the first formula, expand a polynomial at $vx$: sampling $k$ times from $n$ with probability $x$ and then retaining each selected trial with probability $v$ gives a binomial distribution of parameter $vx$. This proves the Bernstein composition formula by conditioning on $k$. The second formula follows by reversing the coefficient sequence and replacing $x$ by $1-x$. The relations hold by linearity for arbitrary real coefficients. Applying two such rational subdivisions in $Y$ gives exactly the requested restriction. Each subdivision uses positive weights and therefore preserves rigorous enclosures. The least signed lower endpoint of the resulting control coefficients, diminished by $\mathcal E$, bounds the signed value from below throughout the band.

Here are the resulting strict lower bounds for the signed components over the indicated closed bands and the entire horizontal interval $|\sigma|\le1/2$. The right column already subtracts $\mathcal E$.
\begin{center}
\begin{tabular}{@{}lll@{}}
\toprule
Height interval & Signed component & Lower bound\\
\midrule
$[0,4.5]$ & $ I$ & $2189578180\cdot10^{-13}$\\
$[4.5,5.2]$ & $-R$ & $11150547269\cdot10^{-13}$\\
$[5.2,6.25]$ & $-I$ & $259498611\cdot10^{-13}$\\
$[6.25,7]$ & $ R$ & $563385832\cdot10^{-13}$\\
$[7,7.65]$ & $ I$ & $19713095\cdot10^{-13}$\\
$[7.65,8.25]$ & $-R$ & $151352849\cdot10^{-13}$\\
$[8.25,9.25]$ & $-I$ & $2269603\cdot10^{-13}$\\
$[9.25,10]$ & $ R$ & $3753921\cdot10^{-13}$\\
\bottomrule
\end{tabular}
\end{center}

On the axis $V=0$, only its zeroth Bernstein coefficient survives. For each pair of endpoints in the theorem, repeat the subdivision to a point in $Y$ there. The signed lower endpoints after deducting $\mathcal E$ are in the last two columns below, in units of $10^{-15}$. Signs are those of $I(0,\tau)$ from left to right.
\begin{center}
\begin{tabular}{@{}lcrr@{}}
\toprule
Ordinate interval & Signs & Left bound & Right bound\\
\midrule
$(5.0031,5.0032)$ & $+,-$ & $61694534$ & $26417969$\\
$(6.8703,6.8704)$ & $-,+$ & $2603092$ & $250551$\\
$(8.0143,8.0144)$ & $+,-$ & $142571$ & $320236$\\
$(9.9328,9.9334)$ & $-,+$ & $13391$ & $13761$\\
\bottomrule
\end{tabular}
\end{center}

Each pair belongs to its own $R$ band. By \eqref{ep:derivative} these four bands contribute exactly one simple zero each in the upper half of the rectangle, all on its central line. The intervening $I$ bands contain none at positive height. Complex conjugation and Lemma~\ref{ep:positive} supply the remaining five zeros and complete the proof. The file \texttt{check\_37a\_strips.py} regenerates both tables and all 50 moment intervals.
\end{proof}

\clearpage
\section{Growing height intervals shared by all rational degree-two forms}\label{common:section}
A zero-density theorem bounds the number of zeros far from the critical line. It also locates \emph{every} zero in certain long intervals of height. The intervals can be chosen simultaneously for all rational elliptic curves, and even for all cuspidal standard functions of $\mathrm{GL}_2/\Q$. The eventual threshold depends on the function.

\begin{theorem}\label{common:main}
Put $T_j=2^j$, $H_j=\exp(\sqrt{\log T_j})$ and $u_j=24/\sqrt{\log T_j}$. For Lebesgue-almost every $v\in[1,2]$ the following assertion holds. For every cuspidal representation $\pi$ of $\mathrm{GL}_2(\mathbb A_\Q)$ with unitary central character, there is $j_0(\pi,v)$ such that, for all $j\ge j_0(\pi,v)$, every nontrivial zero $\beta+i\gamma$ of its unitary standard function satisfying
\[
                       |\gamma-vT_j|\le H_j
             \quad\text{obeys}\quad |\beta-1/2|\le u_j .
\]
For each fixed $\pi$, each such interval contains
\[
                      (2/\pi+o_\pi(1)) H_j\log T_j
\]
zeros counted with multiplicity. In particular the same choices of $v$ give, for every rational elliptic curve, growing height intervals on which all zeros have arithmetic real part within $u_j$ of one.
\end{theorem}

\begin{proof}
We use \href{https://arxiv.org/abs/2103.01956}{Beckwith, Liu, Thorner and Zaharescu, Theorem 1.1}, as in Corollary~\ref{ep:density}. For fixed $\pi$, its right-half zero count at distance $u$ from the line, up to absolute height $3T$, is
\[
                 \ll_\pi T^{1-u/8}\log T,\qquad 0\le u\le1/2.
\]
Functional-equation symmetry gives the same count on the left. For $1\le H\le T/4$ consider $v\in[1,2]$ for which a zero with $|\beta-1/2|>u$ has $|\gamma-vT|\le2H$. Each possible zero has $|\gamma|\le3T$ and excludes a set of $v$ of length at most $4H/T$. Thus the exceptional measure is
\[
                    \ll_\pi H T^{-u/8}\log T .
\]
With the parameters in the theorem this is $O_\pi(\log T_j\,e^{-2\sqrt{\log T_j}})$. The series in $j$ converges, by comparison with the integral after substituting $y=\sqrt j$. The first Borel--Cantelli lemma now gives an eventual result for width $2H_j$ and this fixed $\pi$. No independence assumption on the ordinates occurs.

There is a countable set of representatives for the indicated cuspidal representations modulo the twists $|\det|^{i\tau}$, $\tau\in\R$. In fact a unitary central character over $\Q$ is a finite-order Dirichlet character times an imaginary norm power. Twisting removes the norm power. For each remaining central character the cuspidal Hilbert space is separable and its cuspidal summands occur discretely, so there are at most countably many; there are countably many finite-order central characters. These are the classical discreteness and central-character facts for the cuspidal spectrum. Intersect the measure-one sets just obtained over the representatives. For any fixed real twist the zero ordinates translate by a fixed real number. Since $H_j\to\infty$, an interval of width $H_j$ for that twist is contained, after translation, in the width-$2H_j$ interval for its representative, for all sufficiently large $j$. This proves the simultaneous assertion including all real twists.

Finally the one-sided Riemann--von Mangoldt formula for fixed $\pi$ gives
$N_\pi^+(t)=\pi^{-1}t\log t+c_\pi t+O_\pi(\log(2+t))$ for $t\ge2$. Subtract at $vT_j\pm H_j$. The leading difference is $(2/\pi)H_j\log T_j+O(H_j)$, uniformly for $1\le v\le2$. The linear and remainder terms are $O_\pi(H_j+\log T_j)$, which proves the count. Including endpoints has the same error by the unit-interval zero count. Modularity gives the final statement for elliptic curves.
\end{proof}

For a single fixed form the same computation permits other choices of width. At scale $T$ the exceptional fraction is $O_\pi(A^{-1})$ if
\[
 u=\frac{8}{\log T}\log(AH\log T),\qquad 1\le H\le T/4,\quad 0<u\le1/2 .
\]
For example, at any fixed $0<u<1/2$ one may take $H=T^{u/16}$ and obtain a fraction tending to zero. The narrowing strips above allow zeros to stay slightly off the line and leave gaps between successive height intervals. A zero outside the critical line at a fixed height would not contradict the theorem.

\clearpage
\section{Near-line windows in the height and CM weight aspects}\label{ang:chapter}

There are two ways to obtain a large family from a degree-two function. One can translate the height on one fixed function. For a CM elliptic curve one can also vary the angular power of its Hecke character. In the first case the good intervals can be shared by every fixed cusp form. In the second they start at height zero, and even a growing interval of consecutive angular powers can simultaneously be good. We prove both statements with the logarithmic losses displayed separately from the power saving. These losses matter when the strip about the line shrinks.

Put $d(u)=u/(1-u)$ for $0<u\le1/2$. For a completed cusp function in unitary normalization let $N(\sigma,H)$ count its zeros with $\Re\rho\ge\sigma$ and $|\Im\rho|\le H$, with their orders. The implied constants and the positive constant $C$ below may depend on the fixed automorphic data. They do not depend on $u$, the large family parameter or the height.

\begin{theorem}\label{ang:density}
For each fixed unitary cuspidal $\pi$ of $\mathrm{GL}_2/\mathbb Q$ and $R\ge3$,
\begin{equation}\label{ang:vertical}
 N_\pi(1/2+u,R)\ll_\pi
       (1+u^{-1})R^{1-d(u)}(\log(2R))^C .
\end{equation}
Let $F$ be an imaginary quadratic field and let $\psi$ be a fixed unitary Hecke character of nonzero integral angular type, with its norm twist removed. Fix a unitary Hecke character $\eta$ of $F$. Write $L_m(s)=L_F(s,\eta\psi^m)$. For $M\ge3$ and $1\le H\le M/4$ one has
\begin{equation}\label{ang:angular}
 \sum_{M\le m\le2M} N_m(1/2+u,H)
 \ll_{F,\psi,\eta}(1+u^{-1})(H+1)M^{1-d(u)}(\log(2M))^C .
\end{equation}
In this sum the finitely many indices of zero total angular type may be omitted. Both estimates hold also for leftward deviations by the functional equation.
\end{theorem}

For fixed $\sigma=1/2+u$, the exponent $1-d(u)=4(1-\sigma)/(3-2\sigma)$ is classical in the vertical aspect. Ivi\'c's bound and its proof by \href{https://arxiv.org/abs/1310.0765}{Yashiro, Theorem~1.1} give this exponent with $R^\varepsilon$ for level-one holomorphic forms up to $\sigma=3/4$, and a better exponent further right. The single logarithm of \href{https://arxiv.org/abs/2103.01956}{Beckwith--Liu--Thorner--Zaharescu} retains information still closer to $1/2$. We shall use \eqref{ang:vertical} where a fixed power of $\log R$ is harmless, including for $u\asymp1/\sqrt{\log R}$. For Gaussian angular characters Ricci's estimate, in the form of \href{https://www.math.tau.ac.il/~rudnick/papers/aa190331-23-7.pdf}{Huang--Liu--Rudnick, Theorem~3}, is $\ll H M^{(10/3)(1-\sigma)}(\log M)^C$ when $H=o(M)$. The exponent in \eqref{ang:angular} is smaller for $1/2<\sigma<9/10$. The Gaussian angular sieve and its central mean square appear in \href{https://doi.org/10.1112/S0025579300014388}{Harman--Lewis, Lemmas~1 and~2}.

\subsection*{Two polynomial mean squares}
For the vertical family let $L(s)=\sum\lambda(n)n^{-s}$ and $1/L(s)=\sum\mu(n)n^{-s}$ initially on $\Re s>1$. The multiplicative majorant $h=|\lambda|*|\mu|$ satisfies
\begin{equation}\label{ang:coeff}
 \sum_{n\le 2B} h(n)^2/n\ll_\pi (\log(2B))^C .
\end{equation}
For an elliptic curve one may simply use $h(n)\le d_4(n)$ from its local roots. In general, at a good prime $h(p)=2|\lambda(p)|$. The Rankin--Selberg Euler logarithm at $1+1/\log(2B)$ gives $\sum_{p\le 2B}|\lambda(p)|^2/p\ll_\pi\log\log(3B)$. Indeed its good prime-power coefficients are nonnegative, and the Rankin--Selberg function has a simple pole at one. For $j\ge2$, the Satake bounds give $h(p^j)\ll (j+1)^3p^{7j/64}$. The resulting sum of $h(p^j)^2/p^j$ over $p,j\ge2$ is bounded. At each of the finitely many bad primes the local standard roots have modulus strictly less than $p^{1/2}$ by unitary generic classification, so the corresponding local series also converges. Taking the partial Euler product proves \eqref{ang:coeff}. We use the degree-two bound of Kim--Sarnak and its number-field formulation in \href{https://doi.org/10.4007/annals.2011.174.1.18}{Blomer--Brumley}.

For $B\ge1$, coefficients $|c_n|\le h(n)$ on $B<n\le2B$, and any subinterval of heights of length at most $5R$, the ordinary Dirichlet polynomial mean-value theorem yields
\begin{equation}\label{ang:vls}
 \int \left|\sum c_n n^{-it}\right|^2dt
          \ll (R+B)B(\log(2B))^C .
\end{equation}
The same estimate, with an increased logarithmic power, holds for the squared maximum over prefixes of this sum. To see this, dyadically dissect the integer positions in $(B,2B]$. A prefix is a union of $O(\log(2B))$ disjoint dyadic intervals. Apply Cauchy--Schwarz and sum the mean-value bound over the intervals at each scale; their coefficient square sums are disjoint. This is also the proof we use below if a bounded scalar weight depends on the member of the family. Partial summation only costs the supremum and the total variation of the scalar weight.

Here is the analogue for ideals. Take a common residue class of $m$ modulo an integer $q$ such that $\Theta=\psi^q$ is unramified at all finite places. Such a $q$ kills the finite image on the local unit groups at the finitely many ramified places. For $m=m_0+qk$, write $\eta\psi^m=\xi\Theta^k$ with $\xi=\eta\psi^{m_0}$ fixed. These primitive characters have the same finite conductor. On principal ideals,
$\Theta((z))=(z/|z|)^b$ for a fixed nonzero integer $b$, with one of the two sign conventions for infinity type. For each ideal class choose an integral ideal $\mathfrak b$ in its inverse class. An integral ideal $\mathfrak a$ in the specified class has $\mathfrak a\mathfrak b=(z)$ with $z\in\mathfrak b$. Choose one argument sector for representatives modulo units and split it into a fixed number of smaller sectors of width at most $\pi/|b|$. Within a smaller sector two different rays of norm $N\mathfrak a\le2B$ have angles separated by $\gg_F 1/B$. In fact a nonzero determinant of two vectors of the fixed lattice $\mathfrak b\subset\mathbb C$ has absolute value at least its covolume, and both vector lengths are $O_{\mathfrak b}(\sqrt B)$. The phases $b\arg z$ have the same separation modulo $2\pi$. The factor $\Theta(\mathfrak b)^{-k}$ has absolute value one throughout the class.

On a single ray write $z=d z_0$, with $z_0$ primitive in the lattice $\mathfrak b$ and $d\ge1$ integral. Then $\mathfrak a=(d)\mathfrak a_0$ for $\mathfrak a_0=(z_0)\mathfrak b^{-1}$ integral. Let $\tau_F$ be the ideal-divisor function. Aggregating coefficients bounded by $\tau_F(\mathfrak a)$ over this ray costs at most
\[
 \tau_F(\mathfrak a_0)
 \sum_{d\le\sqrt{2B/N\mathfrak a_0}}\tau(d)^2
 \ll \tau_F(\mathfrak a_0)\sqrt{B/N\mathfrak a_0}\, (\log(2B))^3 .
\]
Here $\tau_F((d))\le\tau(d)^2$ follows prime by prime in a quadratic field. Summing the squares over rays costs $O_F(B(\log(2B))^{10})$, since $\tau_F^2\le d_{4,F}$ and the partial harmonic sums of $d_{4,F}$ are $O_F((\log(2B))^4)$ by the simple pole of $\zeta_F$. Applying the classical large sieve for separated points of the circle in its dual form, and then summing over the finitely many classes and sectors, gives
\begin{equation}\label{ang:als}
 \sum_{M\le m\le2M}
     \left|\sum_{B<N\mathfrak a\le2B}
         c_{\mathfrak a}(\eta\psi^m)(\mathfrak a)\right|^2
            \ll (M+B)B(\log(2B))^C
\end{equation}
for $|c_{\mathfrak a}|\le\tau_F(\mathfrak a)$. Characters vanish on ideals outside their primitive conductor domain. The maximum over norm prefixes obeys the same estimate with another fixed logarithmic power: apply the dyadic dissection just used for~\eqref{ang:vls}. For each ray, the sum of squares of absolute aggregated coefficients over disjoint subintervals is at most the square of their total on the ray. Thus their summed cost at any one dissection scale obeys the same bound. Either mean-square estimate still applies when $c$ is multiplied by $n^{-it}$ for fixed $t$ in the angular case, or when a fixed power of $\log n$ is included with the corresponding logarithmic cost.

\subsection*{The central moment and the detector}
The two central estimates required below are
\begin{align}
 \int_{-4R}^{4R}|L_\pi(1/2+it)|^2dt&\ll R(\log(2R))^C,\\
 \sum_{M\le m\le2M}|L_m(1/2+it)|^2&\ll M(\log(2M))^C
                                   \qquad(|t|\le M).
 \label{ang:means}
\end{align}
For clarity, apply the smooth approximate functional equation with test function $e^{w^2}$; see \href{https://bookstore.ams.org/coll-53}{Iwaniec--Kowalski, Theorem~5.3}. Its two sums have coefficients $\lambda(n)n^{-1/2-it}$ and the dual coefficients with conjugate height. The multiplier on the second sum has modulus one. The analogous ideal sums have coefficients $(\eta\psi^m)(\mathfrak a)(N\mathfrak a)^{-1/2-it}$ and their duals. The archimedean factor is one $\Gamma_{\mathbb C}(s+|b_m|/2+i t_0)$, where $|b_m|\asymp M$ and $t_0$ is fixed. On a norm shell $(B,2B]$ the scalar weight and its variation are $O_A(\min(1,(M/B)^A))$ for any fixed $A>0$, uniformly for $|t|\le M$. For the vertical case first restrict to $R_0\le |t|\le2R_0$; the same assertion holds with $R_0$ in place of $M$, with constants depending on the fixed representation. A bounded initial range costs a fixed amount. These weight bounds follow by moving the gamma-ratio contour to $\Re w=A$ for large $B$. For small $B$ move to a fixed small negative real part, crossing only the residue at zero; the derivative of the constant residue vanishes. Stirling's formula gives the indicated scale, and $e^{w^2}$ absorbs the tails in $\Im w$. In the vertical case take the negative shift small enough to avoid all poles of the real gamma factors. Applying the prefix form of \eqref{ang:vls} or \eqref{ang:als} on each shell, dividing by $\sqrt B$ and taking the triangle inequality in mean square gives \eqref{ang:means}. There are $O(\log(2M))$ shells below $M$, and the rest are geometrically summable; the same argument uses $R_0$ vertically. Summing the dyadic vertical bounds gives the first assertion.

We give the detection details, including the cost of choosing a zero height after the coefficients are known. Write $S=R$ in the vertical family and $S=M$ in the angular family, set $X=S$ and $Y=S^{1/(1-u)}$. Let $M_X(s)$ be the Dirichlet polynomial of inverse Euler coefficients up to $X$, indexed by integers or ideals as appropriate. The coefficients of $L(s)M_X(s)$ up to $X$ are $1$ at the unit ideal or integer and $0$ elsewhere. Above $X$ their moduli are bounded by $h(n)$ vertically. In the angular case they are $(\eta\psi^m)(\mathfrak a)b_X(\mathfrak a)$ with $b_X(\mathfrak a)=\sum_{\mathfrak d\mid\mathfrak a,\,N\mathfrak d\le X}\mu_F(\mathfrak d)$ on ideals prime to the conductor, and $|b_X|\le\tau_F$. Formulae \eqref{ang:vls} and~\eqref{ang:als}, applied to the inverse coefficients themselves, also give the two estimates \eqref{ang:means} for $M_X$.

At an offending zero $\rho=\beta+i\gamma$ put $v=\beta-1/2\ge u$. Mellin inversion of $e^{-n/Y}$ and shifting from $\Re w=2$ to $\Re w=-v$ give
\begin{equation}\label{ang:detect}
 e^{-1/Y}+\sum_{n>X} b(n)n^{-\rho}e^{-n/Y}
 =\frac{1}{2\pi}\int_{\mathbb R}
      (LM_X)(1/2+i\gamma+it)Y^{-v+it}\Gamma(-v+it)\,dt .
\end{equation}
Here and below $n$ denotes the norm in an ideal-indexed sum, with one summand for each ideal, and $b$ denotes the actual convolution coefficient. The only pole crossed is the Mellin pole at zero, whose residue vanishes because $L(\rho)=0$. The shift is justified by fixed-strip polynomial growth and gamma decay. For $u\le v\le1/2$,
$|\Gamma(-v+it)|\ll (1+u^{-1})e^{-c|t|}$ with an absolute $c>0$.

Divide the relevant heights into unit intervals and count occupied pairs of a function and an interval, taking one offending zero in each. Each occupied pair accounts for $O(\log(2S))$ zeros with multiplicity by the standard local zero count in conductor-height aspect; this follows also from the logarithmic derivative of the completion at $2+it$. For the integral in~\eqref{ang:detect}, the exponential majorant for a zero in a given interval is bounded by a constant times $e^{-c|t'-j|}$ in the physical height $t'$, where $j$ is the interval's integer endpoint. Summing these translates over $j$ has bounded overlap. Cauchy--Schwarz and~\eqref{ang:means} thus bound the number of occupied pairs for which the right side has modulus at least $e^{-1/Y}/2$ by
\begin{equation}\label{ang:typeii}
 \ll (1+u^{-1})Y^{-u}S(\log(2S))^C
 \begin{cases}1&\text{vertically},\\H+1&\text{angularly}.\end{cases}
\end{equation}
For the vertical estimate use the central means over $|t'|\le4R$ and exponential decay beyond that range. In the angular estimate integrate the kernel over $|t'|\le M$: the sum of interval kernels there has integral $O(H+1)$. Beyond that range the factor $e^{-c(M-H)}$ absorbs the fixed polynomial conductor and height bounds for $L_mM_X$. The analogous vertical tails are bounded in the same way. These standard polynomial bounds follow directly from the functional equation and Stirling by convexity. If $u^{-1}$ exceeds an arbitrarily large fixed power of $S$, the stated theorem follows already from the total zero count; one can therefore assume $u^{-1}\le S^{10}$ in these tail estimates.

It remains to bound the number of occupied pairs for which the nonconstant sum over $n>X$ on the left of~\eqref{ang:detect} has modulus at least $e^{-1/Y}/2$. Truncate it at $Y(\log(2S))^2$; the omitted tail is smaller than a fixed small constant for large $S$, by exponential decay and the coefficient bounds above. There are $O(\log(2S))$ remaining norm shells $(B,2B]$ starting at $X$. One shell at the chosen zero has modulus $\gg1/\log(2S)$. On it partial summation removes $n^{-\beta}e^{-n/Y}$ at a cost $O(B^{-1/2-u})$ times the maximum of the unweighted prefix sums in $n^{-i\gamma}$. To control the maximum over a unit height interval use
\[
 \sup_{j\le t\le j+1}|P(t)|^2
 \le 2\int_j^{j+1}(|P(t)|^2+|P'(t)|^2)\,dt .
\]
Indeed choose a point where $|P|^2$ is no larger than its average and apply the fundamental theorem of calculus and Cauchy--Schwarz. Apply this to every dyadic piece of the prefix decomposition. Differentiation multiplies each coefficient by $-i\log n$. Equations~\eqref{ang:vls} and~\eqref{ang:als} then bound the contributed count, summed over shells, by
\begin{equation}\label{ang:typei}
 \ll (\log(2S))^C\bigl(SX^{-2u}+Y^{1-2u}\bigr)
 \begin{cases}1&\text{vertically},\\H+1&\text{angularly}.\end{cases}
\end{equation}
In fact the individual shell costs $(S+B)B^{-2u}$ times those logarithmic and height factors. The cutoff $B\le Y(\log(2S))^2$ adds only a fixed logarithmic power. The same estimates cover the finitely many bounded parameters after increasing the constants.

Now $S Y^{-u}=Y^{1-2u}=S^{1-d(u)}$ and $S X^{-2u}\le S^{1-d(u)}$ for $u\le1/2$. Combining~\eqref{ang:typeii} and~\eqref{ang:typei} and including the local zero count proves Theorem~\ref{ang:density}. The unitary functional equation and conjugation give the reflection $\rho\mapsto1-\bar\rho$ on the zeros of each individual function.

\subsection*{Longer intervals for each fixed degree-two function}
Let $R_j=2^j$. For a positive radius $B_j\le R_j/4$, remove those $v\in[1,2]$ for which some ordinate $|\gamma-vR_j|\le2B_j$ has $|\beta-1/2|\ge u_j$. Each offending zero removes length at most $4B_j/R_j$. By \eqref{ang:vertical}, the removed measure is at most
\[
    O_\pi\bigl((1+u_j^{-1})B_j R_j^{-d(u_j)}(\log(2R_j))^C\bigr).
\]
The two signs of the height can be included simultaneously by doubling this bound. Whenever the series of these bounds converges for every fixed $C$, Borel--Cantelli gives eventual avoidance. The countable intersection over the cuspidal representations modulo unitary norm twist has full measure. The reserve factor two in the radius absorbs each fixed real twist for all large $j$, provided $B_j\to\infty$, exactly as in Theorem~\ref{common:main}.

\begin{corollary}\label{ang:long}
For almost every $v\in[1,2]$, each fixed cuspidal standard function of $\mathrm{GL}_2/\mathbb Q$ eventually has all its zeros in both windows
\[
            |\gamma\mp vR_j|\le B_j
\]
within distance $u_j$ of the unitary critical line, for either of the following choices:
\begin{align*}
 B_j&=\exp\sqrt{\log R_j},&
 u_j&=\frac1{\sqrt{\log R_j}}+
       \frac{(\log\log R_j)^2}{\log R_j};\\
 B_j&=R_j^{\alpha},&
 u_j&=u,\qquad 0<\alpha<d(u),\quad 0<u<1/2 .
\end{align*}
Each individual window contains $(2/\pi+o_\pi(1))B_j\log R_j$ zeros with multiplicity. For rational elliptic curves in arithmetic normalization the reference line is $\Re s=1$.
\end{corollary}
\begin{proof}
In the first case $d(u_j)\ge u_j$ and the gain left after multiplication by $B_j$ is at least $\exp(- (\log\log R_j)^2)$, which absorbs every fixed power of $\log R_j$ and is summable in $j$. In the second there is a fixed power saving. Rational parameter pairs and monotonicity give simultaneous validity for every admissible real pair $(\alpha,u)$. The count follows by subtraction of the usual degree-two Riemann--von Mangoldt formula, as in Theorem~\ref{common:main}.
\end{proof}

\subsection*{Centered windows for a growing collection of CM motives}
Let $E/\mathbb Q$ be an elliptic curve with complex multiplication by an order in $F$. Write $\psi$ for its unitary Hecke character over $F$, so that $\psi^c=\psi^{-1}$ and $\pi_E=\mathrm{AI}_{F/\mathbb Q}\psi$. Put $\Pi_m=\mathrm{AI}_{F/\mathbb Q}\psi^m$ for $m\ge1$. This is a genuine cuspidal representation of degree two, since the two angular types $m$ and $-m$ differ. The $L_m$ in this specialization is its standard $L$-function. It is also a primitive factor of symmetric powers of $E$ in their unitary normalization.

\begin{corollary}\label{ang:cmblock}
Let $0<u<1/2$, $0\le D\le M/4$, $1\le H\le M/8$. Outside at most
\begin{equation}\label{ang:blockexc}
 O_E\bigl((1+u^{-1})(D+1)(H+1)M^{1-d(u)}(\log(2M))^C\bigr)
\end{equation}
integers $m\in[M,2M]$, every integer $0\le d\le D$ has the following property. All zeros of the completed standard tensor function for
$\Pi_m\otimes\operatorname{Sym}^d\pi_E$, throughout $|\Im s|\le H$, satisfy $|\Re s-1/2|<u$. The degree of this tensor function is $2(d+1)$. When $H\to\infty$, the number of its zeros in the window is $(2/\pi+o_E(1))(d+1)H\log M$, uniformly in the indicated $m,d$. The same clean-window assertion holds after any fixed abelian base change of $E$ and all these CM factors.
\end{corollary}
\begin{proof}
Induction of the actual CM Weil representation gives, at every place, the identity
\[
 (\operatorname{Ind}\psi^m)\otimes\operatorname{Sym}^d(\operatorname{Ind}\psi)
       \simeq \bigoplus_{r=0}^d\operatorname{Ind}(\psi^{m+d-2r}).
\]
Indeed restriction to $F$ decomposes the symmetric power into characters $\psi^{d-2r}$, and tensoring commutes with induction in this formula. Since $m>d$, all factors on the right have nonzero angular type; induction gives the corresponding primitive automorphic functions and their complete local factors. A bad constituent index belongs to $[M-D,2M+D]$. This range is covered by finitely many dyadic ranges comparable with $M$. It can spoil at most $2\lfloor D\rfloor+1$ choices of $m$. Applying \eqref{ang:angular} and reflection proves~\eqref{ang:blockexc}.

The conductor and parity possibilities of the primitive powers form a fixed finite set at finite places. Their archimedean factor is $\Gamma_{\mathbb C}(s+(m+d-2r)/2)$. The argument principle in the symmetric window gives the change of argument of this gamma factor divided by $\pi$, a linear conductor term and an error $O_E(\log(2M))$. The usual error bound is uniform here: on a disk about $2\pm iH$ of fixed radius the finite $L$-function has a polynomial conductor bound by its functional equation, whereas the Euler product is uniformly bounded above and away from zero at $2\pm iH$. The usual Jensen argument for the argument variation therefore costs $O_E(\log M)$; endpoints at a zero are handled by the corresponding local zero count. For $|t|\le H$, the real part of the logarithmic gamma derivative is $\log M+O_E(1)$, by Stirling, since $m+d-2r\asymp M$. Thus each factor contributes $(2H/\pi)\log M+O_E(H+\log M)$, proving the count after summation. Over a fixed abelian extension of $\mathbb Q$ the base-change function factors into twists by the finite character group. After induction these are the fixed twists $(\chi\circ N_{F/\mathbb Q})\psi^{m+d-2r}$. Applying the same exceptional-set estimate finitely many times proves the last assertion.
\end{proof}

For example, take $u=1/10$ and $D=H=M^{1/20}$. Then $1/20+1/20<1/9=d(1/10)$, so a proportion tending to one of all $m$ in the dyadic interval works simultaneously for all the displayed tensors. All their zeros up to polynomially growing height lie in $2/5<\Re s<3/5$. A shrinking-strip example takes
\[
 D=H=\exp\bigl(\tfrac12\sqrt{\log M}\bigr),\qquad
 u=(\log M)^{-1/2}+(\log\log M)^2/\log M .
\]
Again the exceptional fraction tends to zero. The windows are centered at the real axis; the individual weights and the number of simultaneous primitive factors grow. This does not remove a possible off-line zero of any particular fixed primitive factor.

\subsection*{Actual quadratic twists, including curves without CM}
For a rational elliptic curve $E$ and a nonzero squarefree integer $a$, let $E^{(a)}$ denote the quadratic twist. Here the large parameter changes the curve itself. The following conclusions apply in particular when $E$ has no complex multiplication.

\begin{theorem}\label{ang:quadratic}
Fix $E/\mathbb Q$, $0<u<1/2$, and $0<\alpha<2u/(3-4u)$. There is a $\delta>0$ such that, with at most $O_{E,u,\alpha}(D^{1-\delta})$ exceptions among squarefree $a$ with $D\le |a|\le2D$, every nontrivial zero $w$ of $L(E^{(a)},w)$ with $|\Im w|\le |a|^\alpha$ satisfies
\[
                         |\Re w-1|<u.
\]
Consequently the same holds for a proportion tending to one of the positive prime twisting parameters, including within every fixed reduced progression of primes. It holds simultaneously after base change to any fixed finite collection of abelian number fields.
\end{theorem}

\begin{proof}
First take a fixed normalized holomorphic newform $f$ of weight two and primitive quadratic characters $\chi$ of conductor at most $2D$ coprime to its level. Its primitive twist has coefficients $\lambda_f(n)\chi(n)$, including at the ramified twisting primes. Its conductor is the level of $f$ times $\operatorname{cond}(\chi)^2$. Indeed at a new ramified prime the original two-dimensional representation is unramified and twisting gives conductor exponent twice that of $\chi$ and no inertia invariants. At its original bad primes the twist is unramified. The archimedean types belong to a fixed finite set. The same reasoning applies to fixed finite-order twists of $f$.

\href{https://ora.ox.ac.uk/objects/uuid:b188b365-d8d1-4267-8548-49f01e6cb6a7}{Heath--Brown, \emph{A mean value estimate for real character sums}, Corollary~2} gives on a shell $B<n\le2B$ the estimate
\[
 \sum_\chi\left|\sum c_n\chi(n)\right|^2
 \ll_{\varepsilon}(DB)^\varepsilon(D+B)
       \sum_{n_1n_2=\square}|c_{n_1}c_{n_2}|.
\]
For $|c_n|\le d_4(n)$ the last sum is $\ll_\varepsilon B^{1+\varepsilon}$. Write $n_i=h b_i^2$ with the same squarefree $h$. There are $O(B\sum_{h\le2B}1/h)$ such pairs and $d_4(n)\ll_\varepsilon n^\varepsilon$. This also bounds the sum over a disjoint decomposition at each scale, so the usual maximal-prefix version holds with its logarithms absorbed in $\varepsilon$. Both the twist coefficients and the inverse-convolution detector coefficients are bounded by $d_4$ before multiplication by $\chi$.

In the central approximate functional equation the length is now $O_f(D(2+|t|))$. The dyadic-shell proof of~\eqref{ang:means} gives
\[
 \sum_\chi |L(1/2+it,f\otimes\chi)|^2
       \ll_{f,\varepsilon} D(2+|t|)\,[D(2+|t|)]^\varepsilon .
\]
For the inverse polynomial of length $X=D$ the corresponding bound is $\ll_{f,\varepsilon}D^{1+\varepsilon}$. The smoothing and its total variation are uniform on each normal length scale. One may split the primitive conductors into dyadic intervals, or use the upper length bound together with each character's own smooth weight.

Repeat~\eqref{ang:detect} for the primitive twist. For $H\ge1$ the large-integral count, including the possible unit-interval multiplicity and the gamma singularity as $u\downarrow0$, is
\[
 \ll_{f,\varepsilon}
   (1+u^{-1})(DHY)^\varepsilon D(H+1)^{3/2}Y^{-u}.
\]
Indeed Cauchy--Schwarz at height $t'$ costs $D(2+|t'|)^{1/2}$ times the epsilon loss. Integrating against the sum of interval kernels for $|j|\le H+1$ costs $O((H+1)^{3/2})$, including the exponentially decreasing tails. The large-polynomial count is
\[
 \ll_{f,\varepsilon}(DHY)^\varepsilon (H+1)
                \{D X^{-2u}+Y^{1-2u}\}.
\]
This uses the real-character large sieve in each unit interval, including once with coefficients multiplied by $\log n$. The truncation, Sobolev bound and prefix dissection are those in the proof of~\eqref{ang:typei}. Choose
$Y=(D\sqrt{H+1})^{1/(1-u)}$. Since $d(u)\le2u$ and $Y\le[D(H+1)]^2$, decreasing the initial epsilon gives
\begin{equation}\label{ang:quaddensity}
 \sum_\chi N_{f\otimes\chi}(1/2+u,H)
 \ll_{f,\varepsilon}
 (1+u^{-1})[D(H+1)]^\varepsilon D^{1-d(u)}
       (H+1)^{(3-4u)/(2(1-u))}.
\end{equation}
At $H=(2D)^\alpha$ the exponent of $D$ is strictly below one for sufficiently small fixed $\varepsilon>0$. Reflection handles the other side.

To include all squarefree twisting parameters, let $s$ be the product of the positive odd prime factors of $|a|$ away from the conductor of $E$. Take $b=s$ or $-s$, choosing $b\equiv1\pmod4$. Apart from $s=1$, $b$ is an odd fundamental discriminant. The integer $a/b$ ranges over a fixed finite set of signed divisors of $2\operatorname{rad}(N_E)$. The character of discriminant $b$ is prime to the conductor of the fixed curve $E^{(a/b)}$, so that~\eqref{ang:quaddensity} applies to each of these finitely many base curves. The finitely many parameters with $s=1$ affect no asymptotic. The number of primes in a fixed reduced progression in $[D,2D]$ is asymptotic to a positive constant times $D/\log D$, which exceeds the exceptional count by a factor tending to infinity.

For an abelian base field take all its Dirichlet characters and factor the complete base-change function accordingly. Include their finite conductors in the fixed set of bad primes in the preceding decomposition of $a$. The constituent forms are fixed finite-order twists of the finitely many base curves before twisting by $b$. The same proof and a finite union of exceptional sets give the simultaneous statement.
\end{proof}

In particular, for each fixed $E$, almost every prime $p$ gives an elliptic curve $E^{(p)}$ all of whose nontrivial zeros through $p^{1/16}$ have arithmetic real part between $9/10$ and $11/10$. This gives longer centered windows as the conductor of the rational elliptic curve grows. For any prescribed $H(D)\ge1$ with $H(D)=D^{o(1)}$, the largest deviation from the line in that window tends to zero in density over the squarefree twisting parameters or over the primes. Indeed each fixed positive deviation threshold has exceptional density zero by the theorem. This also gives a subset of relative natural density one along which the largest deviation tends to zero, by taking successive deviation thresholds $1/3,1/4,\ldots$ on sufficiently late consecutive dyadic ranges.

\subsection*{Dirichlet twists with a common prime conductor}
A sharper hybrid conclusion is available for the full character group. Fix a unitary cuspidal $\pi$ of $\mathrm{GL}_2/\mathbb Q$. For every prime $q$ outside its conductor, all nonprincipal characters modulo $q$ give primitive twists of conductor $N_\pi q^2$. For every $H\ge1$,
\begin{equation}\label{ang:primetwists}
 \sum_{\substack{\chi\bmod q\\\chi\ne1}}N_{\pi\otimes\chi}(1/2+u,H)
 \ll_\pi (1+u^{-1})[q(H+1)]^{1-d(u)}
                          (\log(2q(H+1)))^C .
\end{equation}
To prove this, observe that for a Dirichlet polynomial on $B<n\le2B$,
\[
 \sum_{\chi\bmod q}\int_{-H}^{H}\left|\sum_n c_n\chi(n)n^{-it}\right|^2dt
             \ll (q(H+1)+B)\sum_n|c_n|^2 .
\]
Orthogonality first separates the coprime residue classes. In one class consecutive logarithmic frequencies are separated by $\gg\min(1,q/B)$, or there is a single term. The separated-frequency mean-value theorem gives the displayed estimate upon multiplication by $q-1$. This is also the single-modulus hybrid large sieve in \href{https://doi.org/10.1112/plms.12546}{Klurman--Mangerel--Ter\"av\"ainen}. The same prefix dissection handles smooth weights. With $S=q(H+1)$ the central approximate functional equation has length $O_\pi(S)$. The hybrid bound and~\eqref{ang:coeff} give the integrated central mean $O_\pi(S(\log(2S))^C)$ for $L$, and the same for $M_X$ of length $X=S$. For the exponentially weighted integral outside $[-H-1,H+1]$ apply the corresponding bounds successively up to $H+l+1$ and sum against $e^{-cl}$. The type-I and type-II estimates are now exactly~\eqref{ang:typei} and~\eqref{ang:typeii} with unit height factor and $S=q(H+1)$. This proves~\eqref{ang:primetwists}.

For fixed $0<u<1/2$, almost every character of each large prime modulus $q$ consequently has a good centered window of height $q^\alpha$ for every $\alpha<u/(1-2u)$. This means the proportion among the nonprincipal characters tends to one as $q\to\infty$. One can also prescribe open arcs of positive length for $\chi(p)$ at any fixed finitely many primes $p$. The admissible characters have positive limiting proportion by finite-prime equidistribution, while the exceptional characters have proportion tending to zero. Thus the same conclusion holds in every such local box.

\subsection*{The limiting law and a test for one fixed CM function}
The weight-aspect argument also clarifies the step required to infer an all-height conclusion for one member. Let $\mathcal H=\{\Re s>1/2\}$ and use the compact-open topology on holomorphic functions. Fix a unitary ideal character $\xi$ of $F$ of nonzero angular type, and an unramified unitary character $\Theta$ of nonzero angular type without norm twist. Replace $\Theta$ by its square if necessary so that $\Theta(\mathfrak p)=1$ at the primes over ramified rational primes. This is possible since their squares are principal ideals generated by positive rational numbers. Consider the primitive Hecke functions
\[
                     F_k(s)=L_F(s,\xi\Theta^k),\qquad k\ge1 .
\]
At an inert rational prime the local factor is fixed since $\Theta((p))=1$. At a split rational prime choose one prime $\mathfrak p$ above $p$ and an independent uniform unit complex number $U_p$. In the two factors replace $\Theta(\mathfrak p)^k$ by $U_p$ and $\Theta(\bar{\mathfrak p})^k$ by $\bar U_p$. Keep all the factors $\xi(\mathfrak p)$, with value zero where the primitive character is ramified. Denote the resulting random Euler function by $\mathcal L$.

\begin{theorem}\label{ang:bagchi}
On the range $K\le k\le2K$, with the uniform probability measure, $F_k$ converges in distribution as $K\to\infty$ to $\mathcal L$ on $\mathcal H$. Almost surely $\mathcal L$ is holomorphic and has no zeros there. The following conditions on $F_0(s)=L_F(s,\xi)$ are equivalent:
\begin{enumerate}
\item All its nontrivial zeros are on the unitary critical line.
\item For every compact $J\Subset\mathcal H$ and every $\epsilon>0$ the integers $k$ satisfying $\sup_J|F_k-F_0|<\epsilon$ have positive lower density in the dyadic ranges.
\end{enumerate}
Under the first condition one may simultaneously require, at finitely many split primes, that the phases $\Theta(\mathfrak p)^k$ lie in prescribed open neighborhoods of $1$. These statements apply with $\xi$ the normalized elliptic CM character and $\Theta$ a sufficiently divisible power of it.
\end{theorem}

\begin{proof}
For distinct split primes the numbers $\Theta(\mathfrak p)$ are multiplicatively independent. Indeed a multiplicative relation, raised to the class number of $F$, gives a principal fractional ideal $(z)=\prod\mathfrak p^{h n_p}$ with $(z/|z|)^b=1$, where $b\ne0$ is the angular type of $\Theta$. Thus $z/\bar z$ is a root of unity in $F$, so $(z)=(\bar z)$. Unique ideal factorization now makes all $n_p$ vanish. Also $\Theta(\mathfrak p)\Theta(\bar{\mathfrak p})=\Theta((p))=1$. The circle equidistribution criterion applied to each nonconstant monomial implies joint convergence on any finite set of primes to the independent uniform $U_p$.

The random Euler logarithm converges locally uniformly almost surely in $\mathcal H$. Its split first powers
\[
 [\xi(\mathfrak p)U_p+\xi(\bar{\mathfrak p})\bar U_p]p^{-s}
\]
are independent and centered. The sum of their squared norms in $L^2(V,dA)$ is finite on each bounded open rectangle $V$ with closure in $\mathcal H$, by $\sum_p p^{-2\inf_V\Re s}<\infty$. Hilbert-space martingale convergence gives convergence in this area norm almost surely and in mean square. The disk mean-value bound gives local uniform convergence on the interior, and a countable exhaustion handles the half-plane. The inert factors and split higher powers converge absolutely there; the finitely many ramified rational primes contribute fixed nonvanishing factors. Exponentiating proves the assertion for $\mathcal L$.

We explain why this Euler law is the limit of the actual primitive functions. Let $P_{k,X}$ be their ordinary Dirichlet sum through ideal norm $X$. If $J\Subset\mathcal H$ and $0<a<\inf_J\Re s-1/2$, the angular large sieve gives
\begin{equation}\label{ang:lawtail}
 \limsup_{K\to\infty}\left(\frac1K\sum_{K\le k\le2K}
          \sup_J|F_k-P_{k,X}|^2\right)^{1/2}
 \ll_{J,a} X^{-a}(\log(2X))^C.
\end{equation}
Here the logarithm depends on $X$, which is fixed before the limit in $K$. To verify the estimate, use the smooth approximate functional equation on a bounded open neighborhood of $J$ contained in $\Re s>1/2+a$. On a shell $B<n\le2B$ the root mean square of its first sum is bounded by
\[
 (1+B/K)^{1/2} B^{1/2-\sigma}(\log(2B))^C
                       \min(1,(K/B)^A),\qquad \sigma=\Re s.
\]
This follows from~\eqref{ang:als} with its norm-prefix form. The implicit constant is uniform on the neighborhood. For each fixed $X$, the weight at $n\le X$ tends uniformly to one, by the negative shift of its Mellin integral. The displayed bound summed over $B>X$ has limsup $O_{J,a}(X^{-a}(\log(2X))^C)$; the portion $B>K$ is $O_{J,a}(K^{-a}(\log(2K))^C)$ for $A$ sufficiently large. For the dual sum replace $B^{1/2-\sigma}$ by $B^{\sigma-1/2}$ and multiply by the gamma-ratio modulus $O_J(K^{1-2\sigma})$. Its root mean square tends to zero after summing the shells. Finally integrate the squared estimates on the neighborhood and apply the disk mean-value bound, obtaining~\eqref{ang:lawtail}.

For fixed $X$ the polynomial $P_{k,X}$ therefore has as its limit in law the finite ideal sum $P_{U,X}$ obtained from the random phases. To see that $P_{U,X}$ tends to $\mathcal L$, take $Y\ge X$, use finite-phase equidistribution on $P_{k,Y}-P_{k,X}$, and pass the bound~\eqref{ang:lawtail} to the mean-square sup norm on $J$. These random polynomials are Cauchy in that norm as $X\to\infty$. A single subsequence converges locally uniformly almost surely on an exhaustion. Its holomorphic limit equals $\mathcal L$ on $\Re s>1$ by absolute convergence and hence on $\mathcal H$ by the identity theorem. The finite-polynomial approximation and~\eqref{ang:lawtail} now give weak convergence, for example by testing bounded Lipschitz functions of the compact-open metric and first keeping finitely many terms of that metric. The proof gives joint convergence with any fixed finite set of the prime phases.

If $F_0$ has an off-line zero, the functional equation puts one in $\mathcal H$. Choose a disk about it whose boundary is free of zeros. Approximation sufficiently close on that boundary forces a zero of $F_k$ in the disk by Rouch\'e's theorem. Equation~\eqref{ang:angular} makes the proportion of such $k$ tend to zero. This proves necessity.

Conversely, if the zeros of $F_0$ are all on the line, the standard explicit formula for a fixed nonprincipal Hecke character gives
\[
 \sum_{N\mathfrak a\le x}\Lambda_F(\mathfrak a)\xi(\mathfrak a)
                         \ll_\xi \sqrt x\,(\log(2x))^2 .
\]
This is the usual GRH prime-ideal estimate, here used only in the forward implication. It follows by truncating the explicit formula at height $x$ and summing reciprocal zero ordinates using the degree-two zero count. Partial summation gives locally uniform convergence of the Euler logarithm of $F_0$ on $\mathcal H$ in increasing norm order. Using instead all powers of the rational primes up to a cutoff changes no limit, since all higher powers and all inert factors are absolutely summable on these compacts.

Given $J$ and $\epsilon$, choose a finite initial segment of primes so the tail of the target logarithm is small. Enlarge that segment so the random first-power logarithm beyond it has small mean-square supremum on $J$, and the tail of higher powers is uniformly small. Independently of the initial phases, the random tail is then small with positive probability. Constraining the finitely many initial split phases to sufficiently short neighborhoods of $1$ also has positive probability and makes the finite Euler factor close to the target's factor. Exponentiation on the bounded neighborhood of the target logarithm now shows that the open approximation event for $\mathcal L$ has positive probability. The lower bound for open sets under weak convergence and the joint finite-phase assertion prove the required positive lower density.
\end{proof}

\clearpage
\section{An all-height recurrence test using prime twists of the same elliptic curve}\label{etr:chapter}

Fix an elliptic curve $E/\mathbb Q$ and let $G(s)=L(E,s+1/2)$, with central line $\Re s=1/2$. For primes $p\equiv1\pmod{8N_E}$ write $G_p(s)=L(E^{(p)},s+1/2)$. These are primitive functions of actual rational elliptic curves of the same $j$-invariant. The available near-line theorem for almost all of them does not require ERH for any individual twist. This makes the following recurrence test independent of Dirichlet GRH as well. We give a simultaneous formulation because all the curves can use the same twisting primes.

\begin{theorem}\label{etr:main}
Let $E_1,\ldots,E_r$ be any finite collection of rational elliptic curves. Fix any integer $Q$ divisible by $8\prod_j N_{E_j}$. For primes $p\equiv1\pmod Q$ the following are equivalent.
\begin{enumerate}
\item Every $G_j(s)=L(E_j,s+1/2)$ satisfies ERH.
\item For every compact $J\Subset\{\Re s>1/2\}$ and every $\epsilon>0$, a positive lower proportion of the primes $D\le p\le2D$ in the progression satisfy
\[
          \max_{1\le j\le r}\sup_{s\in J}|G_{j,p}(s)-G_j(s)|<\epsilon .
\]
\end{enumerate}
Under the first condition the same approximation may also require any finite collection of rational primes away from $Q$ to split in $\mathbb Q(\sqrt p)$.
\end{theorem}

For one direction choose a compact circle around an off-line zero, lying in the right half-plane and containing no boundary zero. Rouch\'e's theorem and Theorem~\ref{ang:quadratic} make the proportion of approximating twists tend to zero. For the converse and a sharper description of the domain of approximation we determine a probability law. Results in related families include \href{https://www.impan.pl/shop/publication/transaction/download/product/81899}{Mishou--Nagoshi on rank-one prime discriminants} and \href{https://people.math.ethz.ch/~kowalski/universality.pdf}{Kowalski on level-aspect cusp forms}. Here a single isomorphism class over $\overline{\mathbb Q}$ is fixed in each coordinate.

\subsection*{The joint limiting Euler functions}
At every rational prime $\ell\nmid Q$ choose independent signs $X_\ell\in\{-1,1\}$, each with probability $1/2$. Put $X_\ell=1$ for $\ell\mid Q$. To obtain $\mathcal G_j$, replace every normalized reciprocal root of the local Euler factor of $E_j$ at $\ell$ by $X_\ell$ times that root. The signs are shared by all coordinates, including when the curves are not isogenous. The logarithm defines almost surely a holomorphic function on $\mathcal H=\{\Re s>1/2\}$. In fact the random first powers $\lambda_j(\ell)X_\ell\ell^{-s}$ are independent and centered away from $Q$, with $|\lambda_j(\ell)|\le2$. Their squared norms on each bounded rectangle compactly contained in $\mathcal H$ are summable. Martingale convergence in the area norm and the disk mean-value inequality give local uniform convergence, exactly as in Theorem~\ref{ang:bagchi}. All higher prime powers converge absolutely and the finitely many deterministic factors do not vanish in $\mathcal H$. Thus each $\mathcal G_j$ is holomorphic and nowhere zero there.

\begin{proposition}\label{etr:law}
On primes $p\equiv1\pmod Q$, $D\le p\le2D$, the joint distribution of $(G_{j,p})_{j\le r}$ in the compact-open topology on $\mathcal H$ converges to that of $(\mathcal G_j)_{j\le r}$. The convergence holds jointly with the quadratic signs at any prescribed finite set of primes away from $Q$.
\end{proposition}
\begin{proof}
For each fixed set of primes the assertion about the signs follows from quadratic reciprocity, the Chinese remainder theorem and the prime number theorem in fixed progressions. In particular, the limit law of every finite vector of Dirichlet polynomials is the law obtained by replacing their local signs. We prove the compact approximation for a single coordinate. A finite union or a sum of squares then suffices.

Let $P_{p,X}(s)=\sum_{n\le X}\lambda_E(n)(n/p)n^{-s}$. Write $\mathbb E_D$ for the normalized mean on the indicated primes. If $J\Subset\mathcal H$ and $0<a<\inf_J\Re s-1/2$, we claim
\begin{equation}\label{etr:tail}
 \limsup_{D\to\infty}
    \left(\mathbb E_D\sup_J|G_p-P_{p,X}|^2\right)^{1/2}
           \ll_{E,Q,J,a} X^{-a}(\log(2X))^C .
\end{equation}
A loss $D^\varepsilon$ at every length would not give this statement, since $X$ is fixed before taking the limsup. We therefore use three length ranges. On $B<n\le2B$, let $|c_n|\le d_4(n)$. All the estimates below allow a maximum over norm prefixes with a fixed further power of $\log(2B)$, by the same dyadic dissection as in~\eqref{ang:vls}.

First, for any fixed finite set of positive squareclasses $\mathcal S$ there is the uniform bound
\begin{equation}\label{etr:squares}
 \sum_{\operatorname{sf}(n_1n_2)\in\mathcal S}|c_{n_1}c_{n_2}|
                       \ll_{\mathcal S} B(\log(2B))^C ,
\end{equation}
where $\operatorname{sf}$ denotes squarefree kernel. Indeed group $n=h b^2$ by its squarefree part $h$. On this group the sum of absolute coefficients is at most
$d_4(h)\sqrt{2B/h}(\log(2B))^C$. One uses $d_4(b^2)\le\tau(b)^6\le d_{64}(b)$ and the elementary divisor mean. The sum of group squares is $O(B(\log(2B))^C)$ by $\sum_{h\le2B}d_4(h)^2/h\ll(\log(2B))^{16}$. Multiplication by a fixed squareclass permutes all squareclasses. Cauchy--Schwarz now proves~\eqref{etr:squares} for each of them.

Take any fixed $A>0$. If $B\le(\log D)^A$, expand
$\mathbb E_D|\sum c_n(n/p)|^2$ by pairs. When the squarefree kernel of $n_1n_2$ involves a prime outside $Q$, its quadratic character in $p$ is nonprincipal on the progression. The relevant modulus divides $4Qn_1n_2\le16QB^2$. The Siegel--Walfisz theorem, or its character-sum form after detecting the fixed progression, makes the normalized prime average $O_{Q,A,C_0}((\log D)^{-C_0})$ for any fixed $C_0$. Primes dividing $n_1n_2$ do not enter the dyadic prime range for large $D$. The remaining squareclasses are divisors of $\operatorname{rad}(Q)$ and are bounded by~\eqref{etr:squares}. Taking $C_0$ larger than the fixed powers from the number and sizes of the coefficients gives
\begin{equation}\label{etr:short}
       \mathbb E_D\left|\sum c_n(n/p)\right|^2
                   \ll_{Q,A} B(\log(2B))^C .
\end{equation}

For $(\log D)^A<B\le D^{1/3}$, majorize the nonnegative square by summing over \emph{all} positive odd integers $z\equiv1\pmod Q$ in the range, using the Jacobi symbols $(n/z)$. Expand as before. The periodic sums for the nonprincipal terms on that progression are $O_Q(B^2)$ by splitting into complete periods of length at most $16QB^2$. This remains true for imprimitive symbols: an odd prime in the squarefree kernel outside $Q$ contributes a nonprincipal local character, while zeros from other prime factors are periodic. The other terms cost $O_Q(DB(\log(2B))^C)$. Before normalization the contribution from incomplete periods is
$O_Q(B^2(\sum |c_n|)^2)\ll_Q B^4(\log(2B))^C$.
Since $B^3\le D$ and the number of sampled primes is $\asymp_Q D/\log D$, this gives
\begin{equation}\label{etr:mid}
       \mathbb E_D\left|\sum c_n(n/p)\right|^2
               \ll_Q (\log D) B(\log(2B))^C .
\end{equation}
For a disjoint partition at one level of the norm-prefix dissection the squareclass bounds add to at most their full positive majorant. The squares of the incomplete-period coefficient sums do so as well. Thus the asserted prefix estimates follow without a factor equal to the number of intervals.

For $B>D^{1/3}$ use all primitive real characters of conductor at most $2D$. Heath--Brown's Corollary~2 and~\eqref{etr:squares} give
\begin{equation}\label{etr:large}
       \mathbb E_D\left|\sum c_n(n/p)\right|^2
          \ll_{\varepsilon,Q}(DB)^\varepsilon(1+B/D)B(\log(2DB))^C .
\end{equation}
The same estimate is valid at all lengths when the epsilon loss is harmless.

Insert these estimates in the smooth approximate functional equation at $s$ in a bounded open neighborhood of $J$ with real part $>1/2+a$. Its main weight and variation on a shell are $O_L(\min(1,(D/B)^L))$ for any fixed $L$, and its weight at fixed $n$ tends to one. Its dual multiplier is $O(D^{1-2\Re s})$, uniformly in the neighborhood; the two sums have the same shell lengths up to fixed factors. These statements follow from the gamma-ratio integral, since the conductor is $N_Ep^2$ with $p\asymp D$ and both the fixed weight and the height range are bounded. Work first on the strip $1/2<\Re s<5/4$. There both numerator gamma factors in the weight-two functional equation have positive real part on a sufficiently small negative shift of the weight contour, uniformly on the fixed neighborhood. The second domain $\Re s>1$ will use absolute convergence.

More explicitly on this strip, apply~\eqref{etr:short} to the primary shells $B>X$. After the factor $B^{-\Re s}$ their summed root mean square is $O(X^{-a}(\log(2X))^C)$. In~\eqref{etr:mid} the additional $\sqrt{\log D}$ is absorbed at the lower bound $B>(\log D)^A$ by taking $A>3/a$; the geometric sum of those shells tends to zero. In~\eqref{etr:large} first take $\varepsilon$ sufficiently small in terms of $a$. The shells from $D^{1/3}$ to $D$ then give $o(1)$ after the $B^{-\Re s}$ factor, and those above $D$ do so by the weight decay. The dual contribution is $o(1)$ by~\eqref{etr:large} at every length: through $D$ the shell factors are $O(D^{1-2\Re s+\varepsilon}B^{\Re s-1/2+\varepsilon})$ after logarithms are absorbed, and the longer shells have arbitrary decay. The disk mean-value bound gives~\eqref{etr:tail} on the strip.

On a compact inside $\Re s>1$ one may also see the same limiting estimate~\eqref{etr:tail} directly. For fixed $X<Y$, finite-phase equidistribution makes the limiting squared average of the polynomial from $X$ to $Y$ its random-model second moment. Equations~\eqref{etr:squares} and the dyadic decomposition bound its root mean square by $O(X^{-a}(\log(2X))^C)$ uniformly in $Y$, using area means on an enlarged compact. The ordinary series tails beyond $Y$ tend uniformly to zero in $p$ by absolute convergence. A finite compact cover inside the two overlapping domains now proves~\eqref{etr:tail} as stated.

For fixed $X$ the finite polynomial also converges jointly with finitely many prescribed signs. The random versions of the polynomials are Cauchy in mean square uniformly on compacta by~\eqref{etr:tail} and finite-phase equidistribution, as in Theorem~\ref{ang:bagchi}. Their limit agrees with the random Euler products on $\Re s>1$ and therefore throughout $\mathcal H$. Testing bounded Lipschitz functions and using~\eqref{etr:tail} proves the proposition.
\end{proof}

\begin{proof}[Completion of the proof of Theorem~\ref{etr:main}]
Under ERH for each fixed $G_j$, the explicit formula gives $\sum_{n\le x}\Lambda_{G_j}(n)\ll_{E_j}\sqrt{x}\,(\log(2x))^2$ for its normalized logarithmic-derivative coefficients. Partial summation shows that the ordinary Euler logarithm of each target converges uniformly on compact subsets of $\mathcal H$, as in Theorem~\ref{ang:bagchi}. Fix a large finite initial segment of primes to make all its target tails small, and enlarge it to make the random log tails simultaneously small with probability greater than $1/2$. This last claim follows from the squared area norm of the random first powers and absolute convergence of higher powers. Putting $X_\ell=1$ for each of the finitely many initial primes has positive probability, independently of these tails. The same operation enforces any requested extra splitting conditions. The joint weak limit and the open-set lower bound give the positive dyadic proportions in the statement.
\end{proof}

\subsection*{What is approximated without assuming the target hypothesis}
For each target remove the portion in $\mathcal H$ of the leftward horizontal ray from each of its zeros in $\mathcal H$. Let $\Omega$ be the intersection of the resulting open domains. Thus at a point of $\Omega$, every rightward horizontal translate avoids every target zero.

\begin{theorem}\label{etr:slit}
On every compact $J\Subset\Omega$, for every tolerance, a positive lower proportion of the prime twists in Proposition~\ref{etr:law} approximate all the targets simultaneously to that tolerance. Finitely many prescribed primes may be required to split in the twisting fields. There is a sequence of such twisting primes tending to infinity for which the exact simultaneous normality locus of the tuple $(G_{j,p})_{j\le r}$ in $\mathcal H$ is $\Omega$. This assertion is unconditional, including for non-CM elliptic curves.
\end{theorem}
\begin{proof}
For each target divide off the fixed Euler factors at the primes dividing $Q$; these factors are zero-free in $\mathcal H$ and are identical for every twist in the limiting model. Let $F_j$ be the logarithm of the remaining target on $\Omega$, normalized by the ordinary Euler logarithm on the right. The rightward translation argument in Lemma~\ref{st:taylor} gives these logarithms. Its Taylor construction applies simultaneously to the finitely many $F_j$: use the same tube and take the largest of the finitely many sup norms in choosing its degrees. It gives a common scalar weight $W(y)$ in $[0,1]$, arbitrarily close to one on any prescribed bounded interval of $y$, with the absolutely converging weighted prime-power logarithms approximating every $F_j$ on $J$.

Fix a large initial cutoff $X$, including the requested split primes away from $Q$. For larger finitely many primes $\ell\le Y$ use auxiliary independent signs $S_\ell$, with probability $(1+W(\log\ell))/2$ of $+1$. Use $S_\ell=1$ for $\ell\le X$. Their expectations at first prime powers give exactly the Taylor weight beyond $X$. On the initial primes the difference is arbitrarily small by the finite-coordinate condition on $W$. At powers of order at least two the expectation is either $1$ or $W(\log\ell)$. Replacing it by $W(\log\ell^k)$ has arbitrarily small total effect: the tail $\sum_{\ell,k\ge2}\ell^{-k\sigma_*}/k$ is summable for any fixed $\sigma_*>1/2$ below the compact and the weights are arbitrarily close to one on the finitely many remaining arguments. For fixed $W$ its expected first-power tail after $Y$ tends absolutely to zero because $W(\log\ell)=\ell^{-B}$ times a fixed polynomial in $\log\ell$ for a sufficiently large positive $B$.

The variance of the centered finite random logarithms from $\ell>X$ on a bounded open neighborhood of $J$ in $\mathcal H$ is $O(\sum_{\ell>X}\ell^{-2\sigma_*})$ after integrating their squared modulus. This estimate is uniform in the Taylor degrees and in $Y$. Applying the disk bound and summing over the finitely many coordinates gives a deterministic choice of signs through $Y$ whose finite Euler logarithms approximate all the targets on $J$. Choose $X$ first for the variance error, then the Taylor parameters, then $Y$ for the expected tail and for the small variance of the remaining unbiased random signs. The chosen finite pattern has positive probability under the \emph{unbiased} limiting law of Proposition~\ref{etr:law}. Independently of it, the unbiased Euler tails beyond $Y$ are simultaneously small with positive probability. Exponentiation and the weak-limit open-set bound prove the first assertion.

At each stage along a compact exhaustion of $\Omega$ the approximation event has positive lower prime density. The primes for which one of the twists has a zero in a prescribed compact subset of $\mathcal H$ have density zero by Theorem~\ref{ang:quadratic}. Thus one can choose an increasing sequence of approximating primes so that every coordinate is eventually zero-free on each compact of $\mathcal H$. The sequence converges locally uniformly to the targets on $\Omega$.

If a point $s_0=x+i\gamma$ is on a removed ray, one target, say $G_j$, has a zero at or to its right on that ordinate. Simultaneous normality near $s_0$ would imply local convergence to $G_j$: every subsequential meromorphic limit agrees with $G_j$ on a neighboring open subset of $\Omega$ and hence on the disk. If $G_j(s_0)=0$, this contradicts Hurwitz. Otherwise take a short rectangle with left edge through $s_0$ in this disk and right edge past one. Its top and bottom avoid all the finitely many removed ordinates in the rectangle. The $j$th component would converge to $G_j$ on the entire boundary, but is eventually zero-free inside while $G_j$ has a zero there. Rouch\'e gives the contradiction. This proves exactness of the stated simultaneous normality locus.
\end{proof}

\subsection*{A shrinking strip for almost every prime twist}

\begin{theorem}\label{etr:strip}
Suppose $E$ has odd conductor. Write $\mathbb E_D$ for the normalized mean over primes $p\in[D,2D]$, $p\equiv1\bmod Q$, and let $N_p(\sigma,H)$ count zeros of $G_p$ with $\Re\rho\ge\sigma$ and $|\Im\rho|\le H$. There is a constant $C$ such that for $D\ge3$, $1\le H\le D$ and $0<u<1/2$,
\[
 \mathbb E_D N_p(1/2+u,H)
 \ll_{E,Q}(1+u^{-1})(H+1)^{5/2}D^{-u/2}
                  [\log(2D(H+1))]^C .
\]
The assertion concerns dyadic ranges with at least one prime in the progression. For any fixed $A,B>0$, all but $O_{E,Q,A,B}((\log D)^{-A})$ of the prime proportion have
\[
 |\Re\rho-1/2|\le C_{E,Q,A,B}\frac{\log\log D}{\log D}
              \qquad (|\Im\rho|\le(\log D)^B).
\]
For fixed $c>0$, the height may instead be $\exp(c\sqrt{\log D})$ with width
$5c/\sqrt{\log D}+(\log\log D)^2/\log D$ and exceptional proportion smaller than every fixed negative power of $\log D$.
\end{theorem}
\begin{proof}
We use the smooth quadratic-twist estimate of Xiannan Li in the fixed modular-form version stated as \href{https://doi.org/10.1093/qmath/haae028}{Kumar--Mallesham--Sharma--Singh, Proposition~1}; see also the \href{https://arxiv.org/pdf/2303.16864}{displayed proposition and partition of unity}. For the fixed newform $f$ of $E$ and their particular smooth dyadic partition function $g$, its bound on the sum of squares of
\[
       \sum_n\lambda_f(n)(m/n)n^{-1/2-it}g(n/B)
\]
over odd squarefree $D\le|m|\le2D$ is
$\ll_f (1+|t|)^2[D+B\log(2+B/D)]$. Restricting this nonnegative sum to the sampled primes is permissible. For those primes $(p/n)=(n/p)$ also at powers of two. The auxiliary estimate used here is for the fixed form coefficients. The inverse-coefficient polynomials below are estimated separately.

We first obtain, uniformly for real $t$,
\begin{equation}\label{etr:central}
 \sum_{\substack{D\le p\le2D\\p\equiv1\bmod Q}}
     |G_p(1/2+it)|^2
      \ll_{E,Q}D(2+|t|)^3[\log(2D(2+|t|))]^C .
\end{equation}
Insert the smooth approximate functional equation with test function $e^{w^2}$ and partition both sums by $g(n/B)$. On the support of this factor write the additional scalar weight as a Fourier integral in $\log(n/B)$, inserting a fixed compactly supported smooth cutoff equal to one on $\operatorname{supp}(g)$. The Fourier coefficient at $v$ has absolute value
\[
 \ll_{E,L}(1+|v|)^{-4}
       \min\{1,(D(2+|t|)/B)^L\}
\]
for every fixed $L>0$, uniformly in $p$. Indeed this follows by four integrations by parts. Each of the first four logarithmic derivatives of the approximate-functional-equation weight satisfies the same length bound by the gamma-ratio contour and Stirling estimate used in~\eqref{ang:means}. A small negative shift handles short lengths and a positive shift handles long lengths. The completed gamma factor is $\Gamma(s+1/2)$, so at the center the former shift may be chosen strictly between $-1/2$ and zero. Apply Minkowski's inequality to the vector indexed by $p$, taking the supremum of the Fourier coefficient in $p$ before using the quoted estimate at $t-v$. Each root mean square before normalization costs at most a constant times
\[
  (2+|t|)\sqrt{D+B\log(2+B/D)}
                      \min\{1,(D(2+|t|)/B)^L\}.
\]
The dual sum has multiplier of modulus one. Summing the dyadic lengths through $D(2+|t|)$ and their geometrically decreasing tail gives~\eqref{etr:central}.

For $B\le\sqrt D$ we also have the estimate
\begin{equation}\label{etr:polya}
 \sum_{\substack{D\le p\le2D\\p\equiv1\bmod Q}}
    \left|\sum_{B<n\le2B} c_n(n/p)\right|^2
      \ll_Q D B[\log(2D)]^C,\qquad |c_n|\le d_4(n).
\end{equation}
As for~\eqref{etr:mid}, sum the nonnegative square over all denominators $z\equiv1\bmod Q$ in the interval. The frequencies $m=n_1n_2$ whose squarefree kernels divide $\operatorname{rad} Q$ cost $O_Q(DB(\log(2B))^C)$ by~\eqref{etr:squares}. For any other frequency detect the progression with characters modulo $Q$. The product character is nonprincipal, with primitive conductor dividing $4Qm$. P\'olya--Vinogradov for its primitive character, followed by M\"obius exclusion of the extra primes of the imprimitive modulus, bounds the interval sum by $O_Q(\tau(m)\sqrt m\log(2Qm))$. Summed over the two indices, these errors cost
\[
   \ll_Q B\log(2B)\left(\sum_{B<n\le2B}d_4(n)\tau(n)\right)^2
         \ll_Q B^3(\log(2B))^C .
\]
Here $d_4(n)\tau(n)\le d_8(n)$, as follows prime by prime. This proves~\eqref{etr:polya}. Its prefix maximum costs an additional fixed logarithmic power by dyadic dissection. On disjoint pieces at the same scale the absolute pair sums for both estimates add to at most the original full positive majorant.

Put $Y=X=\sqrt D/[\log(2D)]^4$ for sufficiently large $D$. For the polynomial $M_X$ of inverse Euler coefficients,~\eqref{etr:polya} and dyadic summation give $\sum_p|M_X(1/2+it)|^2\ll D(\log(2D))^C$ uniformly in $t$. Its coefficients and the convolution coefficients in~\eqref{ang:detect} after the quadratic sign is removed are bounded in absolute value by $d_4(n)$. Indeed both the form and its inverse have coefficients bounded by $d_2(n)$, and at the twisting prime the ordinary sign convention sets every positive local exponent to zero.

Use~\eqref{ang:detect}, selecting one offending zero per occupied pair of a prime and a unit height interval. The convolution sum truncates at $Y(\log(2D))^2$ with error $o(1)$ uniformly in the zero. Each remaining shell has $X/2\le B\le\sqrt D$. Partial summation and the unit-interval supremum bound for a prefix and its derivative, as in~\eqref{ang:typei}, give total large-polynomial count
$\ll D(H+1)X^{-2u}[\log(2D(H+1))]^C$.
For the integral side, Cauchy--Schwarz at physical height $t'$ uses~\eqref{etr:central} and the inverse-polynomial bound, at a cost $\ll D(2+|t'|)^{3/2}[\log(2D(2+|t'|))]^C$. The translated gamma majorants over the unit intervals centered at $|j|\le H+1$ integrate this to
\[
 \ll (1+u^{-1})Y^{-u}D(H+1)^{5/2}[\log(2D(H+1))]^C .
\]
Both counts absorb the unit-interval multiplicity $O_E(\log(2D(H+1)))$. Divide by $\asymp_Q D/\log D$ and use $Y^{-u}=D^{-u/2}[\log(2D)]^{4u}$. This proves the first assertion. Horizontal reflection and substitution give the two exceptional-proportion statements.
\end{proof}

\subsection*{Recovering an elliptic zero divisor from elliptic zero ordinates}

For $x\in(1/2,1)$ put $b=x-1/2$ and $P(b,y)=b/(b^2+y^2)$. For each prime twist record the multiset of signed imaginary parts $\gamma_p$ of all its nontrivial zeros, including multiplicities, and form
\[
 D_{p,x}(t)=\sum_{\gamma_p}P(b,t-\gamma_p)
       -\log\frac{\sqrt{N_E}\,p}{2\pi}
       -\Re\frac{\Gamma'}{\Gamma}(x+1/2+it).
\]
The sum converges absolutely and locally uniformly in $t$.

\begin{theorem}\label{etr:ordinates}
For every rational elliptic curve $E$ of odd conductor, the primes $p_j$ in Theorem~\ref{etr:slit} can be chosen so that the following holds at every real ordinate $\gamma$ and every $x\in(1/2,1)$ with $G(x+i\gamma)\ne0$:
\[
 \sum_{\substack{G(\beta+i\gamma)=0\\ \beta>x}}
        \operatorname{ord}_{\beta+i\gamma}G
   =\frac1{2\pi}\lim_{h\downarrow0}\lim_{j\to\infty}
                  \int_{\gamma-h}^{\gamma+h}D_{p_j,x}(t)\,dt .
\]
In the sum each distinct zero is indexed once. The target masses can therefore be recovered from the signed zero ordinates of prime quadratic twists of the same curve, their conductors and the common gamma factor. A simultaneous choice works for every fixed finite tuple of odd-conductor curves.
\end{theorem}
\begin{proof}
At each positive-proportion approximation stage in Theorem~\ref{etr:slit}, discard the vanishing proportion of primes for which the conclusion of Theorem~\ref{etr:strip} fails with $H=(\log(2D))^2$ and width $\eta=(\log\log D)^2/\log D$. This is still possible along arbitrarily late dyadic intervals. For a finite tuple discard the union of the exceptional sets. The resulting selected primes retain the convergence on $\Omega$ and eventual zero-freeness on all compact subsets of $\mathcal H$.

For fixed $x$ and a bounded interval of $t$, the real Hadamard identity for the completed function gives $\Re(G_p'/G_p)(x+it)$ by the displayed expression for $D_{p,x}$ with each $P(b,t-\gamma_p)$ replaced by $P(x-\beta_p,t-\gamma_p)$. This identity follows by taking real parts in the logarithmic derivative of the order-one Hadamard product. The remaining real constant for its symmetric zero sum is zero by the functional equation at a nonzero point on the critical line. The gamma and conductor terms displayed above correspond to the primitive complete function.

Pair $\beta_p+i\gamma_p$ with $1-\beta_p+i\gamma_p$. Through $H$, their first-order displacements cancel in the kernel. For large $p$ we have $\eta<b/2$, and the second derivative in the first argument is $\ll_b (1+y^2)^{-1}$ for $b/2\le a\le3b/2$. The kernel replacement thus costs a total $O_{E,x,I}(\eta^2\log(2p))$ on a fixed $t$-interval $I$, by the unit-interval zero bound. For $|\gamma_p|>H$ both kernels are $O_{x,I}((1+\gamma_p^2)^{-1})$ as $H\to\infty$, even if $x-\beta_p$ is negative. The tail costs $O_{E,x,I}(\log(2pH)/H)$. Both errors tend to zero.

Choose $h$ sufficiently small so the left segment $x+i[\gamma-h,\gamma+h]$ has no target zero and the rectangle ending at real part $2$ encloses exactly the target mass in the theorem. Its other three sides lie in $\Omega$ for one curve. For a finite tuple take $h$ small enough to avoid the finitely many other ordinates of all targets in this region. The top, right and bottom logarithmic derivatives of the selected twists then converge uniformly to those of the target, and the twists eventually have no zero in the rectangle. The counterclockwise argument principle, with the left edge traversed downwards, gives
\[
 \lim_{j\to\infty}\int_{\gamma-h}^{\gamma+h}D_{p_j,x}(t)\,dt
   =2\pi\sum_{\beta>x}\operatorname{ord}_{\beta+i\gamma}G
      +\int_{\gamma-h}^{\gamma+h}\Re(G'/G)(x+it)\,dt ,
\]
where the sum is restricted to the target zeros at this ordinate. The last integral tends to zero with $h$ because its integrand is regular there. Varying $x$ yields the cumulative masses and hence the real parts and orders of all right-half-plane target zeros.
\end{proof}

\clearpage
\section{A common interval for every elliptic conductor up to a growing bound}
\label{uw:chapter}

In a countable-family theorem the height after which a particular curve has good windows may depend arbitrarily on the curve. Here the conductor cutoff grows with the height and every curve below that cutoff is included simultaneously. We also control the zeros of each entire symmetric power of a fixed CM elliptic curve as its degree grows.

\begin{theorem}\label{uw:hybrid}
Let $f$ be a normalized primitive holomorphic cusp form of weight $k\ge1$, level $N$ and arbitrary nebentypus, in unitary normalization. Put
\[
 D(u)=\frac{u}{1-u},\qquad
 S=\sqrt{H\{H+\sqrt N(k+H)\}},\qquad \ell=\log(2N(k+H)).
\]
For $H\ge2$, $0<u<1/2$, the number of its nontrivial zeros with
$|\Im\rho|\le H$ and $\Re\rho\ge1/2+u$ satisfies
\[
       N_f(1/2+u,H)\ll(1+u^{-1})S^{\,1-D(u)}\ell^C ,
\]
with absolute implied constant and absolute $C$. Horizontal reflection gives the same bound on the other side.
\end{theorem}
\begin{proof}
The finite roots have modulus at most one, by holomorphic temperedness and the weight-one finite-image theorem; see the sources in Theorem~\ref{li:holodyadic}. Thus the coefficients of $L_f$, its inverse and their truncated convolution are bounded uniformly by $d_2,d_2,d_4$ respectively. In particular \eqref{ang:coeff} and the prefix polynomial mean square \eqref{ang:vls} hold with absolute logarithmic powers and constants.

For a height interval $[-V,V]$, $V\ge2$, the smooth central approximate functional equation has effective dyadic length $\sqrt N(k+V)$: the integration variables in its single complex gamma ratio have real part $k/2$ before shift. For the scalar weight and its first logarithmic derivative on a shell $(P,2P]$ the bound is
\[
        O_A\!\left(\min(1,(\sqrt N(k+V)/P)^A)\right)
\]
uniformly for $|t|\le V$. To see uniformity at small lengths shift a fixed distance less than $1/4$ to the left; the gamma numerator then has positive real part even at $k=1$. For greater lengths shift to $A$ on the right. Stirling's gamma-ratio bound with the test $e^{w^2}$ handles all $k$ and $t$, as in~\eqref{ang:means}. Partial summation with the prefix mean bound on each shell therefore gives
\[
     \int_{-V}^V |L_f(1/2+it)|^2\,dt
       \ll (V+\sqrt N(k+V))[\log(2N(k+V))]^C .
\]
The multiplier of the dual approximate sum has modulus one. For the inverse polynomial $M_H$ the same integral is
$\ll(V+H)[\log(2NH(k+V))]^C$, by the ordinary polynomial mean alone.

In \eqref{ang:detect} use $X=H$ and $Y=S^{1/(1-u)}$. The count from the left nonconstant polynomial is at most
\[
      \ell^C\{H^{1-2u}+Y^{1-2u}\},
\]
by the prefix and unit-height supremum proof of~\eqref{ang:typei}. Its dyadic lengths stop at $Y\ell^2$; since $Y\le S^2$, all their logarithms are $O(\ell)$. The zero multiplicity per unit interval is $O(\ell)$ in the full initial height range, uniformly in the form by the real Hadamard formula at $2+it$, the gamma bound, and the bounded logarithmic Euler derivative there.

For the right integral use Cauchy--Schwarz on $[-H-j-2,H+j+2]$ and sum the outer intervals with exponentially decreasing weight $O(e^{-c j})$, as in~\eqref{ang:typeii}. The two preceding means at $V=H+j+2$ have product of square roots $\ll S(1+j)^C\ell^C$. Thus this count is $\ll(1+u^{-1})S Y^{-u}\ell^C$, including the unit-interval zero multiplicity. The identities $SY^{-u}=Y^{1-2u}=S^{1-D(u)}$ and $H^{1-2u}\le S^{1-D(u)}$ finish the proof.
\end{proof}

\begin{corollary}\label{uw:all}
Fix $0<u<1/2$ and $\kappa,\alpha>0$ such that
\[
       \alpha+\kappa\,\frac{9-D(u)}4<D(u).
\]
For almost every $v\in[1,2]$, at every sufficiently large dyadic $T$, \emph{every} rational elliptic curve of conductor $N\le T^\kappa$ has every nontrivial zero with $\big||\Im\rho|-vT\big|\le T^\alpha$ in $|\Re\rho-1/2|<u$ in unitary normalization. For example $N\le T^{1/40}$ and windows of radius $T^{1/20}$ work simultaneously in the arithmetic strip $0.9<\Re s<1.1$.

More generally, let $a,c>0$, put $L=\log T$, and take conductors $N\le\exp(c\sqrt L)$ and window radius $\exp(a\sqrt L)$. The common normalized width can be taken to be
\[
       \frac{a+9c/4}{\sqrt L}+\frac{(\log L)^2}{L}.
\]
Both conclusions include all primitive weight-two newforms of trivial character at the specified levels, irrespective of their coefficient fields.
\end{corollary}
\begin{proof}
By the dimension formula the total number of normalized weight-two eigenforms of trivial nebentypus over levels $N\le Y$ is $\ll Y^2[\log(2Y)]^C$. Indeed each dimension is $O([\mathrm{SL}_2(\mathbb Z):\Gamma_0(N)])=O(N\prod_{p\mid N}(1+1/p))$; use the ordinary dimension formula, for instance \href{https://wstein.org/books/modform/modform/dimension_formulas.html}{Stein, \emph{Modular Forms}, dimension formulas}. Modularity includes one such eigenform for every rational elliptic isogeny class. Apply Theorem~\ref{uw:hybrid} at $H=3T$. Its factor $S$ is $O(TN^{1/4})$. Each offending signed zero excludes a measurable set of $v$ of length at most twice the radius divided by $T$. If the radius is $W\le T/2$, the union bound over all curves and zeros gives
\[
 \ll(1+u^{-1})\,W\,T^{-D(u)}\,
          Y^{\,2+(1-D(u))/4}[\log(2TY)]^C .
\]
This is summable over dyadic $T$ in the first parameter range. With the second choices and the displayed $u$, its logarithm is at most $-(\log L)^2+O_{a,c}(\log L)$: use $D(u)\ge u$ and $2+(1-D(u))/4\le9/4$. Borel--Cantelli proves both assertions.
\end{proof}

The same argument includes all trivial-character holomorphic newforms with $2\le k\le\exp(e\sqrt L)$ and $N\le\exp(c\sqrt L)$, with the coefficient of $1/\sqrt L$ enlarged to $a+9c/4+2e$. The aggregate dimension through weight $K$ costs an additional factor $O(K^2)$ and $k=o(T)$ keeps the same central length. For any fixed included weight and conductor the zero count in each of the two growing windows is asymptotic to $(2/\pi)W\log T$ by the usual degree-two zero count. The hypothesis on every zero is thus a statement on populated windows.

\subsection*{Each complete symmetric power of a CM curve}
\begin{theorem}\label{uw:sym}
Let $E/\mathbb Q$ have CM and put $Z_n(\sigma,H)$ for the right-side zero count of its \emph{full} unitary $\operatorname{Sym}^n$ function, with multiplicity, $n\ge1$. For $H\ge2$ put $R=\max(n,H)$. Then unconditionally
\[
 Z_n(1/2+u,H)\ll_E (1+u^{-1})\,nH\, R^{-D(u)}
                                  [\log(2nH)]^C
             \qquad(0<u<1/2).
\]
In particular along \emph{every} sequence with $H\to\infty$, $R\to\infty$, all but $O_{E,A}((\log R)^{-A})$ of the zero proportion are in
\[
        |\Re\rho-1/2|\le C_{E,A}\log\log R/\log R
\]
for each fixed $A>0$. The degree $n+1$ can grow without any relation to $H$.
\end{theorem}
\begin{proof}
By the CM decomposition in Theorem~\ref{recentcm:main}, the degree-two factors are $L_F(s,\psi^j)$ for positive $j=n,n-2,\ldots$, where $\psi$ is the fixed angular character of the elliptic curve. Their associated primitive holomorphic forms have weights $j+1$ and uniformly bounded levels: the finite inertia characters of powers of a fixed Hecke character have bounded conductor, and the quadratic induction multiplies the norm conductor by the fixed field discriminant. For even $n$ there is also one of two fixed degree-one middle characters.

On a dyadic range $M\le j\le2M$ with $H\le M/4$, use~\eqref{ang:angular} and enlarge the range to all integers since zero counts are nonnegative. The resulting bound is
\[
  \ll_E(1+u^{-1}) M H M^{-D(u)}[\log(2nH)]^C .
\] If $M<4H$, apply Theorem~\ref{uw:hybrid} separately at the $O(M)$ weights. The levels are bounded and $j+1\ll H$, whence $S\ll_E H$, yielding the same bound with saving $H^{-D(u)}$ instead. Sum the dyadic ranges. For $M\le 4H$ their lengths sum to $O(\min(n,H))$. For $M>4H$ the sum of $M^{1-D(u)}$ is bounded by $O(n^{1-D(u)}\log(2n))$, uniformly even if $D(u)$ approaches one. This gives the theorem for the degree-two constituents, both when $n\ge H$ and when $H>n$.

For the fixed middle character its degree-one central approximate functional equation at height $H$ has length $O_E(\sqrt H)$, or one can use the Dirichlet polynomial mean for its finite coefficients. The detector with $X=H$, $Y=H^{1/(1-u)}$ yields $O_E((1+u^{-1})H^{1-D(u)}(\log(2H))^C)$. For the principal character the shift in~\eqref{ang:detect} crosses the pole of $\zeta$ at $1$. Its extra term is $Y^{1-\rho}\Gamma(1-\rho)M_H(1)$, with $|M_H(1)|\ll\log(2H)$. If $|\Im\rho|\ge C_1\log(2H)$ this is smaller than $1/10$ for large fixed $C_1$ by gamma decay and $Y\le H^2$; the other heights contain only $O(\log(2H)\log\log(3H))$ zeros by the fixed zeta zero count. These are absorbed by the logarithmic power uniformly in $u$. In both ranges of $n,H$ the degree-one bound is at most the right side of the theorem. Multiplicities, including common ordinates, add in this complete factorization.

For each bounded-level weight $k$, the standard uniform argument-principle zero count and Stirling give
\[
   N_f(0,H)=\tfrac{2H}{\pi}\log(k+H)
                     +O_E(H+\log(k+H)).
\]
One may obtain the error directly from the usual argument variation bound on the horizontal segments, $O(\log(N(k+H)))$, by Jensen on a fixed disk about $2\pm iH$, using the bounded Euler logarithm there and uniform convexity as in the proof of Theorem~\ref{fl:global}. Along the right edge the Euler argument has bounded variation between its endpoints on its Euler branch. The gamma phase increment is the integral of its real digamma derivative and equals $2H\log(k+H)+O(H)$, since the average of $\log(k+|t|)$ there differs from $\log(k+H)$ by $O(1)$. This also holds by limiting past endpoint ordinates and the local zero count. The degree-one factor has the corresponding leading coefficient $H/\pi$ and scale $H$. Summing over the parity class of weights gives
\[
 Z_n(0,H)=\frac{(n+1)H}{\pi}\log(n+H)
                 +O_E\bigl(nH+n\log(n+H)\bigr).
\]
Indeed averaging $\log(j+1+H)$ loses only a constant from its top-end value, by integral comparison. For the possible middle term the discrepancy $H\log(1+n/H)$ is absorbed as well. Divide the off-line count by this total, use reflection and take $u=C_{E,A}\log\log R/\log R$ with sufficiently large constant.
\end{proof}

\clearpage
\section{The spacing scale of zeros in prime elliptic twists}\label{ms:chapter}
The shrinking strip controls real parts. A count in every interval of imaginary parts gives further information. For almost all prime twists of a fixed odd-conductor elliptic curve we obtain the expected count uniformly down to a multiple of $\log\log p$ mean spacings. Thus the line is close to actual zeros throughout the indicated height range, as well as the zeros being close to the line.

\begin{theorem}\label{ms:local}
Let $E/\mathbb Q$ have odd conductor and retain primes $p\equiv1\bmod Q$, where $8N_E\mid Q$. Fix $A,B>0$ and $0<\varepsilon<1$. Write $L=\log D$. There exists $K=K(E,Q,A,B,\varepsilon)$ such that for all sufficiently large $D$, outside a proportion $O_{E,Q,A,B,\varepsilon}(L^{-A})$ of these primes $D\le p\le2D$, the following holds for \emph{every} interval $I\subset[-L^B,L^B]$ of length at least $K\log L/L$:
\[
 (1-\varepsilon)\,\frac{|I|\log p}{\pi}
 \ \le\ \#\{\rho:G_p(\rho)=0,\ \Im\rho\in I\}
 \ \le\ (1+\varepsilon)\,\frac{|I|\log p}{\pi}.
\]
Only nontrivial zeros are counted, with multiplicity. The endpoints of $I$ may be included or excluded. Every such zero may simultaneously be required to have horizontal distance $O_{E,Q,A,B}(\log L/L)$ from the critical line.
\end{theorem}
The largest gap among these ordinates is $O_{E,Q,A,B}(\log L/L)$. Every point $1/2+it$, $|t|\le L^B$, is within complex distance $O_{E,Q,A,B}(\log L/L)$ of a nontrivial zero of the primitive elliptic function. At an endpoint take the entire testing interval on the available side.

\subsection*{A high moment on the twisting primes}
Put $Z=D^{1/16}$. Upper-bound Selberg weights for the progression $z\equiv1\bmod Q$ give a nonnegative function
\[
 {\cal V}(z)=\left|\sum_{d\mid z}\lambda_d\right|^2,\qquad
 \lambda_1=1,\quad\lambda_d=0\quad(d>Z\ \text{or }(d,Q)>1),
 \qquad |\lambda_d|\ll_Q\log(2Z)
\]
for which the sum over $D\le z\le2D$, $z\equiv1\bmod Q$, is $O_Q(D/\log D)$. Indeed, take the squarefree Selberg coefficients with local densities $1/\ell$ at primes $\ell\nmid Q$ and $d\le Z$. Their optimized quadratic form is the reciprocal of
$\sum_{r\le Z,\,(r,Q)=1}\mu^2(r)/\varphi(r)\gg_Q\log Z$ and the coefficients have modulus at most $d/\varphi(d)$. The $O_Q(1)$ counting error per common-multiple progression costs $O_Q(Z^2\log^2(2Z))$ in all. See \href{https://dms.umontreal.ca/~koukoulo/documents/publications/primes.pdf}{Koukoulopoulos, \emph{The Distribution of Prime Numbers}, Chapter~21}. Each sampled prime has weight one.

Let $\sum_{\ell\le U,\,\ell\nmid Q} a_\ell(\ell/p)$ be a sum supported on primes, with $\sum |a_\ell|^2\le C_0$ and $|a_\ell|\le C_0$, $C_0\ge1$. For each integer $k\ge1$ with $(2C_0 U)^{4k}\le D^{1/4}$,
\begin{equation}\label{ms:moment}
 \mathbb E_D\left|\sum_\ell a_\ell(\ell/p)\right|^{2k}
                \ll_Q (2C_0k)^k+D^{-1/2}.
\end{equation}
Here $\mathbb E_D$ is normalized on the dyadic primes in the progression. Majorize the nonnegative moment by the corresponding sum over the weighted integer Jacobi denominators $z$. Every square frequency contributes in absolute value at most the total sieve weight. The absolute coefficient products for those frequencies sum to at most $(2k-1)!!(\sum|a_\ell|^2)^k$, since every square product permits a pairing of equal primes. For a nonsquare frequency $m$, its squarefree kernel has a prime outside $Q$. On expanding the sieve square put $z=[d,e]w$. If $([d,e],m)>1$ the Jacobi symbol vanishes. Otherwise the sum over $w$ is on one reduced progression modulo $Q$; the symbol is nonprincipal there, of period dividing $4Qm$. Reciprocity and the Chinese remainder theorem give sum zero on every complete period. This remains true for imprimitive symbols because an outside prime contributes a nonprincipal local character. All $d,e$ are odd since $2\mid Q$. Thus the incomplete sum costs $O_Q(m)$. Combining the absolute errors gives
\[
 \ll_Q Z^2\log^2(2Z)\,U^{2k}(C_0 U)^{2k}.
\]
Divide by $\asymp_Q D/\log D$ to prove~\eqref{ms:moment}. The square-frequency contribution has the correct prime-scale normalization.

\subsection*{The explicit formula and the actual ordinates}
Take $h=C_1\log L/L$ with $C_1\ge1$ to be fixed in terms of the accuracy. A real test $\phi$ will be Schwartz with smooth compactly supported Fourier transform, or the nonnegative function
\[
 \phi_0(y)=\operatorname{sinc}^2(y)
       +\operatorname{sinc}^2(y-1/2)+\operatorname{sinc}^2(y+1/2),
 \qquad \operatorname{sinc}(y)=\sin(\pi y)/(\pi y),
\]
with the removable values filled in. Our Fourier transform uses $e^{-2\pi ixy}$. These functions are entire and satisfy
\[
 |\phi^{(j)}(y+iv)|\ll_\phi
      (1+y^2)^{-1}\exp(C_\phi|v|)\qquad(j=0,1,2).
\]
For the Schwartz tests this follows by two integrations by parts in Fourier inversion, absorbing the resulting powers of $|v|$ into the exponential. For $\phi_0$ it follows from the three displayed terms. Moreover $\phi_0(y)\gg(1+y^2)^{-1}$ on $\mathbb R$, since a sine and its half-period translate cannot be simultaneously small.

Write $v_\rho=\Re\rho-1/2$ and $\gamma_\rho=\Im\rho$. The explicit formula with test $\phi((\,\cdot-t)/h)$ at $|t|\le2L^B+2$ gives
\begin{equation}\label{ms:explicit}
 \Re\sum_\rho\phi\big((\gamma_\rho-t-i v_\rho)/h\big)
  =\frac{h\log p}{\pi}\int_{\mathbb R}\phi
       +O_{E,\phi}(h\log(3+|t|)+h)+O_\phi(1)+P_{p,t,\phi}.
\end{equation}
This is the Guinand--Weil formula for the weight-two completion, with gamma $\Gamma(s+1/2)$, conductor $N_Ep^2$ and no poles. See \href{https://bookstore.ams.org/coll-53}{Iwaniec--Kowalski, Theorem~5.12}. For the test $\phi_0$, which decays quadratically, first multiply by $\omega(y/R)$ with $\omega$ real Schwartz, $\omega(0)=1$ and smooth compact Fourier support. Apply the formula at fixed $p,t,h$, then let $R\to\infty$. At the zeros the imaginary displacement is bounded for fixed $h$, so the quadratic decay and the unit-interval zero count justify dominated convergence. The archimedean term has the same integrable majorant and the prime side remains supported on a fixed finite set. Its archimedean remainder follows by integrating the bound $|\Re(\Gamma'/\Gamma)(1+iw)|\ll\log(2+|w|)$ against the test. The term $P_{p,t,\phi}$ is a real part of a prime sum as in~\eqref{ms:moment}, with
\[
 U\le e^{C_\phi/h},\qquad
 |a_\ell|\ll_{E,\phi}h\frac{\log\ell}{\sqrt\ell},\qquad
 \sum_{\ell\le U}|a_\ell|^2\ll_{E,\phi}h^2(\log(2U))^2\ll_{E,\phi}1.
\]
Indeed the two Fourier terms at a prime have bounded transform and phases of modulus one. Fixed primes dividing $Q$ cost $O_{E,\phi}(h)$. The second prime powers cost $O_\phi(h\sum_{\ell\le U}\log\ell/\ell)=O_\phi(1)$ in absolute value; higher powers cost $O_\phi(h)$. The squared-coefficient bound follows from Chebyshev and partial summation. The twisting prime exceeds the Fourier-support cutoff for large $D$.

For any finite collection of fixed tests take $k=\lfloor c_1hL\rfloor$, where $c_1>0$ is sufficiently small for~\eqref{ms:moment}. At each $t$ and every fixed $\delta>0$, Markov's inequality gives
\[
 \operatorname{Prob}\{|P_{p,t,\phi}|>\delta hL\}
     \ll_{\delta,\phi,Q}\exp[-c_2hL\log(hL)] .
\]
This is smaller than every fixed inverse power of $L$. In particular outside a proportion $O_{A,B,\phi,\delta}(L^{-A})$ the bound holds simultaneously at $O(L^{B+2})$ grid points for all the tests.

We also need control of the real parts at distant heights. Outside a proportion $O_{E,Q,A}(L^{-A})$, Theorem~\ref{etr:strip} gives
\begin{equation}\label{ms:taper}
 |v_\rho|\le K_0\frac{\log L+\log(2+|\gamma_\rho|)}{L}
                       \qquad(|\gamma_\rho|\le D).
\end{equation}
Apply that mean estimate at all $H=2^j\le D$ and at $H=D$, with $u=K'_0(\log L+\log(2+H))/L$. For $u\ge1/2$ the ordinary critical strip suffices. Otherwise $D^{-u/2}$ absorbs the height power, $1+u^{-1}$, all logarithmic losses and the $O(L)$ applications. A union bound and fixed enlargement from the dyadic endpoints give~\eqref{ms:taper}.

Put $M_p(t)=\sum_\rho[1+((\gamma_\rho-t)/h)^2]^{-1}$. At $|\gamma_\rho|\le J=L^{B+5}$, \eqref{ms:taper} gives $|v_\rho|/h\le C_{E,Q,A,B}/C_1$. Once this is at most one, Taylor's formula bounds the total error in replacing the real part of the test by its value on the real axis by
\[
       O_\phi\big((C_{E,Q,A,B}/C_1)^2 M_p(t)\big).
\]
The first-order term is purely imaginary. At $J<|\gamma_\rho|\le D$, the factor $\exp(C_\phi|v_\rho|/h)$ is $O_\phi((2+|\gamma_\rho|)^{1/4})$ for sufficiently large $D$, since $hL=C_1\log L$. The term $\log L$ in~\eqref{ms:taper} contributes a bounded factor. At $|\gamma_\rho|>D$ use $\exp(C_\phi/(2h))\le D^{1/4}$. The unit-interval zero count is $O_E(L+\log(2+|\gamma_\rho|))$. Since $|\gamma_\rho-t|\gg|\gamma_\rho|$ in the two distant ranges, the combined real and complex tails there cost
\[
 O_{E,\phi}\big(h^2(L+\log J)J^{-3/4}+h^2L D^{-3/4}\big)=o(hL).
\]
The same tail estimate applies to the real weight defining $M_p(t)$.

Apply~\eqref{ms:explicit} first to $\phi_0$ on the grid. Its left-hand real sum is $O(hL)$ by the prime-moment estimate. Its sum on the real ordinates is at least a positive absolute multiple of $M_p(t)$. Taking $C_1$ sufficiently large absorbs the displacement error and gives $M_p(t)\ll hL$ on the entire grid. For every other fixed test, a further enlargement of $C_1$ and small fixed choice of the moment tolerance therefore give
\begin{equation}\label{ms:real}
 \left|\frac1{h\log p}\sum_\rho\phi((\gamma_\rho-t)/h)
           -\frac1\pi\int_{\mathbb R}\phi\right|\le \delta'
\end{equation}
for any prescribed $\delta'>0$, throughout the grid. We first fix the finite list of tests and all accuracies, then $C_1$, then sufficiently large $D$.

\begin{proof}[Completion of the proof of Theorem~\ref{ms:local}]
Choose smooth nonnegative compactly supported functions $g_-$ and $g_+$ with $g_-\le\mathbf1_{(-1,1)}\le\mathbf1_{[-1,1]}\le g_+$ and integrals arbitrarily close to $2$. Each can be approximated by a real Schwartz test $\phi_\pm$ with smooth compact Fourier support in the weighted norm $\sup_y(1+y^2)|g_\pm(y)-\phi_\pm(y)|$. Indeed multiply the Fourier transform of $g_\pm$ by smooth compact cutoffs tending to one; convergence holds in the Schwartz topology. The error in the ordinate sum is at most this norm times $M_p(t)$, and the integral error at most $\pi$ times the norm. Thus~\eqref{ms:real} proves the upper count for a closed bin $[t-h,t+h]$ and the lower count for its interior, with any fixed relative accuracy.

Take a grid of adjacent bins of length $2h$ covering $[-L^B-1,L^B+1]$. There are $O(L^{B+2})$ centers. Sum open-bin lower bounds over the bins contained in $I$ and closed-bin upper bounds over all bins that meet $I$. Only a bounded number of boundary bins can affect the enclosing lengths, giving relative error $O(h/|I|)$. Take the bin error smaller than $\varepsilon/4$ and $K$ sufficiently large relative to $C_1/\varepsilon$. This proves the assertion with every endpoint convention. The horizontal conclusion follows from~\eqref{ms:taper} on the same primes.
\end{proof}

\subsection*{Angular characters over any imaginary quadratic field}
\begin{theorem}\label{ms:angular}
Let $F,\eta,\psi$ be as in Theorem~\ref{ang:density}, with nonzero angular type for $\psi$ and zero norm component. For $M\le m\le2M$ the conclusions of Theorem~\ref{ms:local} hold for $L_F(s,\eta\psi^m)$ with $\log m$ in the main term, $L=\log M$, and exceptional proportion $O_{A,B,\varepsilon,F,\eta,\psi}(L^{-A})$. They hold also relative to each prescribed fixed progression of $m$. No assumption on Hecke zeros enters this assertion.
\end{theorem}
\begin{proof}
Refine the progression into the finitely many classes $m=m_0+qj$ in which $\Theta=\psi^q$ is unramified at every finite place. Put $\xi=\eta\psi^{m_0}$. On principal ideals $\Theta((z))=(z/|z|)^b$ for a fixed nonzero integer $b$. Inert rational primes have norm $\ell^2$, so their contribution to the explicit formula is $O_\phi(h)$ times $\sum_{\ell\le U}\log\ell/\ell=O_\phi(1/h)$, just as for second powers. The finitely many ramified primes cause bounded terms. At a split rational prime with chosen prime ideal $\mathfrak l$ put $z_\ell=\Theta(\mathfrak l)$. Then $\Theta(\bar{\mathfrak l})=z_\ell^{-1}$, since $\Theta((\ell))=1$. The varying first primes therefore contribute a real part of
\[
      \sum_{\ell\le U,\,\ell\ {\rm split}}(a_\ell z_\ell^j+b_\ell z_\ell^{-j}),
      \qquad \sum_\ell(|a_\ell|^2+|b_\ell|^2)\ll_{\phi,F}1,
      \qquad U\le e^{C_\phi/h}.
\]
The coefficients absorb the fixed parts $\xi$ and the oscillations in $t$; the same coefficient estimates as before apply.

Here a direct separation estimate replaces the Selberg sieve. In the $2k$th moment, a frequency is a product $\prod z_\ell^{e_\ell}$ with $\sum|e_\ell|\le2k$. If any $e_\ell$ is nonzero, form the integral ideal $\mathfrak a$ using $\mathfrak l^{e_\ell}$ at positive exponents and $\bar{\mathfrak l}^{-e_\ell}$ at negative ones. Its norm is at most $U^{2k}$ and its character value is the required frequency. With $h_F$ the class number, write $\mathfrak a^{h_F}=(w)$, $w\in\mathcal O_F$. Then $|w|\le U^{k h_F}$. The algebraic integer $w^{|b|}$ is nonreal: equality with its conjugate would give equality of the ideals $\mathfrak a^{h_F|b|}$ and $\bar{\mathfrak a}^{h_F|b|}$, impossible at a split prime with nonzero imbalance. Every nonzero imaginary part of an element of $\mathcal O_F$ is bounded away from zero by a fixed field constant. Consequently
\[
     |1-\Theta(\mathfrak a)|\ \ge\
         |1-\Theta(\mathfrak a)^{h_F}|/h_F
          \gg_{F,b} U^{-k h_F|b|}.
\]
Averaging the geometric progression in $j$ now costs $O_{F,b,q}(M^{-1}U^{k h_F|b|})$ for each nonconstant frequency. Their total coefficient mass is at most $(C_\phi U)^{2k}$. For $k=\lfloor c h\log M\rfloor$ and sufficiently small $c$, all such frequencies cost $O(M^{-1/2})$. A constant frequency requires every prime to occur an even number of times. Summing the absolute phase choices first gives at most
$(2k-1)!![\sum(|a_\ell|+|b_\ell|)^2]^k$, by the same pairing argument. This proves the high-moment bound used in~\eqref{ms:explicit}.

The $m$th gamma factor is a single complex gamma with real shift $|b_m|/2\asymp M$ and a fixed imaginary shift. Its real logarithmic derivative contributes the main $h\log m\,\int\phi/\pi$ at $|t|\le2L^B+2$, with error $O_{F,\eta,\psi,\phi}(h(1+\log(2+|t|)))$. Indeed Stirling gives $\log m+O(1+\log(1+|w|/M))$ at integration height $w$, which is integrable with the same test bound. The finite conductor is bounded on the classes under consideration. The taper~\eqref{ms:taper} through $M/4$ now follows from~\eqref{ang:angular}, using $D(u)\ge u$. Beyond $M/4$ use the whole critical strip in the distant-tail estimate. All remaining steps in Theorem~\ref{ms:local} use only these inputs, including the unit-interval bound $O_{F,\eta,\psi}(\log(M+|\gamma|+2))$. This proves the result on each prescribed class.
\end{proof}

\begin{corollary}\label{ms:sym}
Fix a CM elliptic curve $E/\mathbb Q$. Given $A,B>0$ and $0<\varepsilon<1$, there are constants $K,C$ such that for every sufficiently large $n$, in \emph{every} interval $I\subset[-(\log n)^B,(\log n)^B]$ with $|I|\ge K\log\log n/\log n$, the complete $\operatorname{Sym}^n E$ has between $1-\varepsilon$ and $1+\varepsilon$ times
\[
                  \frac{(n+1)|I|\log n}{2\pi}
\]
nontrivial zeros with signed ordinate in $I$. At most a proportion $O_{E,A,B,\varepsilon}((\log n)^{-A})$ of the zeros in \emph{any one} of those intervals can have $|\Re\rho-1/2|>C\log\log n/\log n$. Multiplicities are included in both statements.
\end{corollary}
\begin{proof}
Write $L=\log n$. Among the positive constituent weights $j\equiv n\bmod2$ in Theorem~\ref{uw:sym}, first set aside $j\le\sqrt n$ and the possible fixed middle character. Their unit-interval zero counts are $O_E(L)$ in the stated region; they contribute $O_E(\sqrt n(1+|I|)L)$ zeros in any $I$. Cover the remaining indices by dyadic ranges $M\le j\le2M$ with $\log M\ge L/2-O(1)$. Apply Theorem~\ref{ms:angular} there with height exponent $B+1$, tolerance $\varepsilon/4$, and exceptional power $A+4$. Truncating the last dyadic interval can only reduce its absolute exceptional count. Summing the exceptional counts costs $O(nL^{-A-4})$ weights, with total contribution $O_E(nL^{-A-4}(1+|I|)L)$. One fixed increase in $K,C$ covers all the dyadic ranges.

For the remaining weights each interval has count within $\varepsilon/4$ relative error of $|I|\log j/\pi$ and all its zeros in the requested strip. Summing $\log j$ over the parity class gives $(n+1)\log n/2+O(n)$ by integral comparison. Removing the small or exceptional weights changes this sum by a negligible proportion. Since $|I|\gg\log L/L$, both upper bounds for excluded zeros are $O_{E,A}(L^{-A}n|I|L)$. These estimates prove both assertions, with no exceptional degree or interval.
\end{proof}

There is a limitation on the length of intervals in a statement of this form. For $E_d:y^2=x^3-dx$, $d$ odd squarefree, \href{https://doi.org/10.1112/mtk.70067}{David--Devin--Waxman, Lemma~2.5} show that exactly two odd classes modulo $8$ have root number $-1$ for the primitive angular powers, and all even classes have root number $+1$. For odd $n$ the complete symmetric power contains every odd primitive exponent through $n$ once. Its analytic central multiplicity is therefore at least $n/4-O(1)$, unconditionally. A uniform relative main term as in Corollary~\ref{ms:sym} cannot hold at all intervals of length $o(1/\log n)$, since one such interval can contain this forced center. The proved scale differs from this obstruction by a factor of $\log\log n$.

\subsection*{An interval shared by cusp forms of growing weight and level}
\begin{theorem}\label{ms:vertical}
Fix $A,B,C>0$ and $0<\varepsilon<1$. Put $L=\log T$. Apart from a set of centers $t_0\in[T,2T]$ of length $O_{A,B,C,\varepsilon}(T L^{-A})$, the following holds simultaneously for every primitive holomorphic cusp form over $\mathbb Q$ with $k,N\le L^C$, of any nebentypus. On both windows $|t-t_0|\le L^B$ and $|t+t_0|\le L^B$, all the nontrivial zeros have $|\Re\rho-1/2|\le C_0\log L/L$. Every subinterval $I$ of either window with $|I|\ge K\log L/L$ contains between $1-\varepsilon$ and $1+\varepsilon$ times $|I|\log T/\pi$ of their signed ordinates with multiplicities. The constants $C_0,K$ are independent of the form.
\end{theorem}
\begin{proof}
The argument of Theorem~\ref{ms:local} is applied with random $t_0$ under normalized Lebesgue measure. At heights up to $3T$, Theorem~\ref{uw:hybrid} gives
\[
     N_f(1/2+u,3T)\ll (1+u^{-1})T^{1-D(u)}L^{C_2}
                     \qquad(k,N\le L^C)
\]
with $C_2$ depending only on $C$. For a fixed zero, the fraction of centers at distance at most $H$ from its signed ordinate, or its negative, is $O(H/T)$. Thus the taper~\eqref{ms:taper} has a vertical analogue, on both signs, with $|\gamma_\rho|$ on its right replaced by the distance $|\gamma_\rho\mp t_0|$. It holds out to distance $T/4$, with exceptional center proportion $O(L^{-A'})$ for arbitrary fixed $A'$. Indeed apply the density bound at each dyadic radius $H\le T/4$, with $u=K_0(\log L+\log(2+H))/L$, using the trivial strip once $u\ge1/2$. The same computation includes the last radius and both reflections.

For a fixed offset $t$ in $[-2L^B-2,2L^B+2]$ and a fixed test, the short prime terms in the explicit formula centered at $t_0+t$ are the real part of $\sum_{\ell\le U} a_\ell\ell^{-it_0}$, with precisely the same length and coefficient-square bounds as before. The possible pair of conjugate terms from the Fourier transform reduces to this real part, including for a nonreal form because the reflected term uses the dual coefficients. Holomorphic temperedness gives the coefficient bounds with absolute constants, including at bad primes. At $-t_0+t$ the phases have opposite sign, which gives the same mean. Put $r=\lfloor c hL\rfloor$, with $c$ sufficiently small that $U^r\le T^{1/4}$. Write the $r$th power as an ordinary Dirichlet polynomial with frequencies $n\le U^r$. The Montgomery--Vaughan mean-value theorem and unique factorization give
\[
 \frac1T\int_T^{2T}\left|\sum_{\ell\le U}a_\ell\ell^{-it_0}\right|^{2r}dt_0
       \ll (1+U^r/T)\,r!\left(\sum_\ell|a_\ell|^2\right)^r .
\]
Indeed, two ordered lists of $r$ primes with equal product are permutations of one another. The large-deviation estimate for the entire growing grid follows as in~\eqref{ms:moment}.

The gamma factor is one complex gamma with real shift $(k-1)/2$. At a signed center of absolute height $\asymp T$, its logarithmic derivative contributes the main $h\log T\,\int\phi/\pi$ and a uniform $O_\phi(h)$ error, up to the finite-conductor contribution $O(h\log N)=O_C(h\log L)$. To see the gamma estimate, split its real integral at distance $T/2$ from the center. In the nearby range Stirling gives real part $\log T+O(1)$ uniformly for $k\le L^C$. In the far range the test is $O_\phi(h^2/|w\mp t_0-t|^2)$ and Stirling, with the uniform bound on bounded gamma arguments for $k\ge1$, makes its contribution $O_\phi(h^2\log(2T)/T)$. The unit-interval zero bound is $O_C(\log(2T+|w|))$ at height $w$ for the included parameters. Thus the estimates following~\eqref{ms:taper} apply with relative distance to the signed center and ordinary critical-strip control beyond $T/4$. This gives the local counts from the grid exactly as before.

Finally the dimension formula at $\Gamma_1(N)$ bounds the total number of such forms by $O(L^{C_3})$, with $C_3$ depending only on $C$. Indeed the index is $O(N^2)$, the dimension at weight $k\ge2$ is $O(kN^2)$, and multiplication by the level-one Eisenstein series of weight four embeds every weight-one space in its corresponding weight-five space. Choose $A'>A+C_3+2$ and sum the exceptional measures. This gives a single set of good centers for all the forms and both signs.
\end{proof}

For each fixed unitary Maass cusp form over $\mathbb Q$ the same statement holds with constants allowed to depend on the form, by Theorem~\ref{ang:density} for its vertical density. The prime coefficients in the short sum still satisfy the bounded-square condition. Indeed the good-prime Rankin--Selberg logarithmic derivative has nonnegative coefficients, and its prime-power sum through $x$ is $O_\pi(x)$ by the fixed Rankin--Selberg prime number theorem. Therefore $\sum_{\ell\le U}|\lambda_\pi(\ell)|^2(\log\ell)^2/\ell\ll_\pi(\log(2U))^2$. For the second powers use $|\alpha_\ell^2+\beta_\ell^2|\le|\lambda_\pi(\ell)|^2+2$ and the same estimate by partial summation. At powers of order at least three the Kim--Sarnak bound gives an absolutely convergent sum since $3(1/2-7/64)>1$. Finitely many bad primes contribute $O_\pi(h)$ by the strict local bound $|\alpha|<\sqrt\ell$. The fixed two real gamma factors have the same combined leading term $h\log T/\pi$ and are regular in the real gamma integral on the critical line. All other steps, including the prime polynomial mean, carry over.

Taking $A>1$ and applying Borel--Cantelli on dyadic heights, almost every $v\in[1,2]$ therefore gives such windows centered at $\pm v2^j$ at every sufficiently large $j$. One may do this simultaneously for all fixed unitary cusp forms on $\mathrm{GL}_2/\mathbb Q$, using countably many representatives modulo norm twists. A fixed translation of the ordinates is absorbed by first increasing the polylogarithmic window exponent, and changes neither local counts nor normalized real parts. Rational choices of the positive exponents and accuracies suffice to cover all fixed choices by inclusion.

\subsection*{Exact critical-line zeros inside the common windows}
\begin{proposition}\label{ms:dyadic-line}
For the newform of a fixed rational elliptic curve and any fixed $0\le a<1$, every interval $J\subset[T,2T]$ of length $H\ge T/(\log T)^a$ contains at least $0.0696$ of its zeros exactly on the critical line, as a fraction of its total zero count, for all sufficiently large $T$. More precisely, a sufficiently large fixed $C_E$ permits every $H\ge C_E T(\log\log T)^4/\log T$.
\end{proposition}
\begin{proof}
We give the localization of the published mollified moment, since subtracting two lower bounds for the cumulative count would not prove the assertion. Use \href{https://arxiv.org/pdf/2004.03581}{Andersen--Thorner, Theorem~3.1 and Section~6} with its large parameter equal to $2T$. Put $L_1=\log(2T)$, $\sigma_0=1/2-R/L_1$, and use their fixed admissible mollifier $\psi$ and real polynomial
\[
 Q(x)=c_0+\sum_j d_j(1-2x)^{2j-1},\qquad Q(0)=1.
\]
Their parameter optimization is strict: one can fix $R,\nu,P,Q$ within the permitted lengths with $1-\log c(P,Q,2R,\nu/2)/(2R)>0.0696$. If necessary approximate the optimal smooth mollifier polynomial in the integral defining $c$ while preserving its endpoints and strict margin. Here $c_0\ne0$: otherwise $Q(1)=-1$, and their displayed lower bound for the optimized moment in Section~6, at their $R>9$, already exceeds the value required for a positive proportion.

Take a smooth majorant of the indicator of $J$, supported on its enlargement by $O(T/L_1)$ at both ends, bounded by one, and with derivatives $O_j((L_1/T)^j)$. It is an admissible weight in their Theorem~3.1 and its integral is $H+O(T/L_1)$. With $V=Q(-(2L_1)^{-1}d/ds)L_f(s)$, the same normalized shift differentiations as in their proof of Theorem~2.1 give
\[
  \frac1H\int_J |V(\sigma_0+it)\psi(\sigma_0+it)|^2dt
        \le c(P,Q,2R,\nu/2)+o(1).
\]
Indeed the smoothed error is $O(T(\log L_1)^4/L_1)=o(H)$, uniformly in the endpoints. Their differentiated estimates also give normalized moments $O(1)$ when $V$ is replaced by $L_f^{(j)}/L_1^j$ for every fixed needed $j$. This follows directly by Cauchy's formula in the shifts at radius comparable to $1/L_1$ in Theorem~3.1; opposite shifts can use unequal radii, keeping their sum away from zero on the contours. These differentiations cost a fixed power of $L_1$ precisely cancelled by the derivative normalization.

For clarity write the phased completion as $\Phi=X L_f$, real on its critical line, with $X$ the gamma-conductor factor including the fixed phase. Put
\[
 G=X^{-1}\left(c_0\Phi+\sum_j d_j\,\Phi^{(2j-1)}/L_1^{\,2j-1}\right).
\]
On the line $\Re(XG)=c_0\Phi$, because the odd derivatives are purely imaginary. In a fixed horizontal strip at $T\le t\le2T$ one has $X'/X=L_1+O(1)$ by Stirling, with bounded higher logarithmic derivatives. Leibniz's formula therefore gives $G=V+O(L_1^{-1})\sum_{j\le\deg Q}|L_f^{(j)}|/L_1^j$ as an error bound. The preceding moments and Cauchy--Schwarz show that the same mollified second-moment upper bound holds for $G\psi$ on $J$.

Apply the local Levinson argument on the rectangle with horizontal sides at the endpoints of $J$; compare \href{https://arxiv.org/pdf/1105.1177}{Conrey--Iwaniec--Soundararajan, Appendix, equations (A.16)--(A.25)}. We record the endpoint accounting. Each variation of $\pi$ in the argument of $XG$ on the critical edge requires a zero of its real part, with common zeros handled by the usual indentations to the right. The gamma contribution divided by $\pi$ is $N_f(J)+O_f(\log T)$ by the argument principle for the completion. The variation of $\arg G$ on the remaining three sides gives at worst twice its number of right-side zeros, apart from $O_{f,S}(\log T)$, on taking a fixed sufficiently large right edge $S$. Indeed $G=1+O_{f,Q}(L_1^{-1}+2^{-S/2})$ on this edge by absolute convergence and the same gamma derivative estimate. The two horizontal argument variations are $O_{f,Q,S}(\log T)$ by Jensen on a fixed disk about the right endpoint, using the polynomial bound in $T$ and the nonzero real part there. The same argument applies to the mollifier, since $\psi=1+O_{f,P}(2^{-S/2})$ on that edge. Endpoints can be displaced infinitesimally to avoid zeros; the local count of the target is $O_f(\log T)$ on any unit endpoint interval.

Thus $N_{f,0}(J)\ge N_f(J)-2N_G+O_{f,S}(\log T)$, where $N_G$ counts the zeros of $G$ between the critical line and $S$ on the rectangle. Littlewood's formula for $G\psi$ from $\sigma_0$ to $S$ gives
\[
 2\pi(1/2-\sigma_0)N_G
 \le \int_J\log |G\psi(\sigma_0+it)|dt
       +O_{f,Q,P}\big(H(L_1^{-1}+2^{-S/2})\big)
       +O_{f,Q,S}(\log T).
\]
In this use of Littlewood the extra mollifier zeros only increase the weighted count. Jensen bounds the integral by $(H/2)\log(H^{-1}\int_J|G\psi|^2)$. As $N_f(J)\sim HL_1/\pi$ uniformly in the stated range, division gives the numerical proportion of Andersen--Thorner with additional error $O_{f,Q,P,R}(2^{-S/2})+o_{f,Q,P,R,S}(1)$. Fix $S$ sufficiently far to the right for the first error to be less than half the strict numerical margin, and then take large $T$. In the more precise range the smoothed moment error divided by $H$ is $O_{f,P,Q,R}(1/C_E)$, and $T/L_1=o(H)$. Taking $C_E$ sufficiently large absorbs this additional error inside the same strict margin. This proves the proposition.
\end{proof}

\begin{corollary}\label{ms:online}
Fix a rational elliptic curve $E$ and $0\le a<1$. In every interval $J\subset[T,2T]$ of length at least $T/(\log T)^a$, at least $4\%$ of centers $t_0\in J$ satisfy the simultaneous conclusion of Theorem~\ref{ms:vertical} at relative tolerance $1/100$ and, in addition, at least $3\%$ of the zeros of this fixed curve in $[t_0-h,t_0+h]$ lie \emph{exactly} on the critical line. Here $T$ is sufficiently large, $h=K\log\log T/\log T$ with fixed sufficiently large $K$, and $B,C>0$ may be prescribed. The more precise lower length of Proposition~\ref{ms:dyadic-line} also suffices here.
\end{corollary}
\begin{proof}
Integrate the on-line zero count in $[t_0-h,t_0+h]$ over $t_0\in J$. Each interior zero contributes $2h$ and the two ends cost $O_E(h\log T)$. By Proposition~\ref{ms:dyadic-line} the integral is at least $(0.0696-o(1))\,2h|J|\log T/\pi$. The excluded centers of Theorem~\ref{ms:vertical}, with $A=4$, cost at most $O_E(T(\log T)^{-3})$ by the unit-interval count, negligible even for the smallest $J$. On the remaining centers the total local count is at most $1.01\,(2h\log T/\pi)$. If a proportion $s$ of the centers in $J$ is good and has the asserted $3\%$ on the line, the integral is consequently at most $(1.01)(s+0.03(1-s))+o(1)$ in units of $2h|J|\log T/\pi$. Thus $s\ge(0.0696/1.01-0.03)/0.97-o(1)>0.04$ for sufficiently large $T$.
\end{proof}

The lower proportion in this corollary does not determine an on-line proportion for almost every center. Even a configuration with equally spaced ordinates in pairs can concentrate its on-line pairs on a fixed fraction of many consecutive intermediate-length blocks, placing every other pair at small opposite horizontal displacements. Such a configuration can have the asserted near-line counts and a positive proportion on the line in every long interval. To obtain a density-one local exact-line conclusion would require further information on the local fluctuations of the on-line count or on sign changes of the Hardy function.

\clearpage
\section{Li coefficients as the conductor increases}\label{fl:chapter}
The real Li coefficients involve the entire nontrivial zero divisor. A growing initial interval can nevertheless be asymptotically determined by the conductor alone. A common all-height bound gives a statement for every member of a family. A distribution theorem on compact subsets of the right half-plane gives a different statement for prime twists of any fixed elliptic curve.

\subsection*{A deterministic consequence of a global zero-free half-plane}
Write the completion of an entire primitive standard function in unitary normalization as
\[
 \Phi(s)=N^{s/2}\prod_{j=1}^d\Gamma_{\mathbb R}(s+\mu_j)L(s),
       \qquad \mathcal C=N\prod_{j=1}^d(3+|\mu_j|).
\]
Assume fixed absolute degree $d$, tempered finite roots, $\Re\mu_j\ge0$, the standard functional equation and an entire completion of order one. These hypotheses include non-norm Hecke characters and the primitive factors of CM elliptic curves. Products with bounded total degree and the same hypotheses are also allowed.

\begin{theorem}\label{fl:global}
Suppose every nontrivial zero satisfies $\Re\rho\le\theta$ for a fixed $1/2\le\theta<1$. Define
\[
 c(\theta)=\sup_{0<r<(1-\theta)/\theta}
 \frac{(1+r)^{-1}-\theta}{(1-\theta)\log(1/r)}.
\]
For every $0<c<c(\theta)$ there is $\delta>0$ such that, uniformly over these functions as $\mathcal C\to\infty$,
\[
 \Re\lambda_m=\frac m2\log\mathcal C
       +O_{\theta,c,d}\bigl(m(\log\mathcal C)^{1-\delta}\bigr)
                 \qquad(1\le m\le c\log\log\mathcal C).
\]
With the all-height $7/8$ Hecke theorem reviewed in Section~\ref{recentcm:audit}, every CM elliptic curve over a field of bounded degree and sufficiently large absolute analytic conductor therefore has positive real Li coefficients through $0.218\log\log\mathcal C$. This conclusion requires no individual zero-height computation. It also applies to all non-norm unitary Hecke characters of bounded absolute degree.
\end{theorem}
\begin{proof}
On $\Re s\ge1+e$ with $e>0$ fixed, grouping by rational primes bounds the absolute Euler logarithm by $d\log\zeta(1+e)$. The functional equation and gamma ratios give the usual convexity bound
\[
       \log|L(s)|\ll_{d,e}\log(\mathcal C(3+|\Im s|)^d)
               \qquad (-e\le\Re s\le 3).
\]
Indeed on $\Re s=-e$ the conductor factor is $N^{1/2+e}$ and Stirling bounds the archimedean ratio by fixed powers of the quantities $3+|\mu_j|+|\Im s|$. An inverse gamma pole causes a zero. Take $0<e<1/8$, so the dual gamma numerator has positive real part there. The rest of the bound follows by Phragm\'en--Lindel\"of and finite order.

The Euler logarithm extends throughout $\Re s>\theta$. Apply Borel--Carath\'eodory to it in disks centered at $2+it$, with outer radius $2-\theta-e/3$ and inner radius $2-\theta-e$. Its value at the center is bounded by the Euler series, and its real part is bounded above by convexity. This gives
\[
       |\log L(\theta+e+it)|\ll_{d,e,\theta}
                      \log(\mathcal C(3+|t|)^d).
\]
Between $\theta+e$ and $1+e$ apply the three-lines theorem to $\log L(s)\exp((s-it_0)^2)$ for fixed bounded $t_0$. The Gaussian absorbs the height dependence on the left, and the right line is bounded. Thus for any fixed $\theta<\sigma_0<1$, sufficiently small $e>0$, and bounded $s$ with $\Re s\ge\sigma_0$,
\[
   |\log L(s)|\ll (\log\mathcal C)^{(1+e-\sigma_0)/(1-\theta)}.
\]
Points to the right of $1+e$ use absolute convergence.

Put $s(z)=1/(1-z)$. On $|z|\le r$ the least real part of $s(z)$ is $\sigma_r=1/(1+r)$. With the logarithm vanishing at $z=0$ one has
\[
 \log\frac{\prod_j\Gamma_{\mathbb R}(s(z)+\mu_j)}
              {\prod_j\Gamma_{\mathbb R}(1+\mu_j)}
      =\frac{s(z)-1}{2}\sum_j\log(3+|\mu_j|)+O_{d,r}(1).
\]
In fact integrate the digamma derivative from $1$ to $s(z)$. Stirling gives the stated leading derivative with bounded complex error for large $|\mu_j|$, since the argument of $s+\mu_j$ is bounded. Bounded $\mu_j$ also give bounded error, since the gamma arguments stay separated from their poles on this compact.

The real part of $m[z^m]\log(\Phi(s(z))/\Phi(1))$ is $\Re\lambda_m$. In the order-one Hadamard product a zero at $\rho$ contributes $1-(\rho/(\rho-1))^m+m/\rho$ to this coefficient after multiplication by $m$. The exponential factor contributes $mB$, where $\Re B+\sum_\rho\Re(1/\rho)=0$ by the functional equation, as in the real Hadamard identity. The remaining real parts give the usual Li sum by horizontal reflection. These sums of real parts converge absolutely, so the real coefficient identity also holds for non-self-dual functions.

Choose $r$ for which $c\log(1/r)<(\sigma_r-\theta)/(1-\theta)$ and then $e>0$ small enough to preserve this strict inequality after subtracting $e/(1-\theta)$ on the right. Cauchy's bound for the gamma remainder and the Euler logarithm gives
\[
 \left|\Re\lambda_m-\tfrac m2\log\mathcal C\right|
    \ll m r^{-m}\left(1+(\log\mathcal C)^{(1+e-\sigma_r)/(1-\theta)}\right).
\]
This proves the power saving. For $\theta=7/8$ one may take $r=1/30$, for which the quotient in the supremum exceeds $0.2181$. The CM factorization in Theorem~\ref{recentcm:main} has only non-norm factors for the first elliptic power. Equality of the completed factors gives the same conductor up to a bounded factor for bounded absolute degree. This proves the applications.
\end{proof}

\subsection*{Prime quadratic twists of any fixed elliptic curve}
\begin{theorem}\label{fl:prime}
Let $E/\mathbb Q$ be any elliptic curve and take primes $p\equiv1\bmod Q$ as in Theorem~\ref{etr:main}. For all fixed $A,\varepsilon>0$,
\[
 \lim_{D\to\infty}\operatorname{Prob}_{D\le p\le2D}
 \left\{\max_{1\le m\le A\log\log D}
       \left|\frac{\Re\lambda_m(E^{(p)})}{m\log p}-1\right|>\varepsilon
 \right\}=0 .
\]
If $N_E$ is odd, for every $B>0$ there exists $b=b(E,Q,B)>0$ such that, outside a proportion $O_{E,Q,B}((\log D)^{-B})$, every real Li coefficient of the twist through $b\log p$ is strictly positive.
\end{theorem}
\begin{proof}
Fix $A$ and choose $0<r<r'<1$ so that $A\log(1/r)<1/2$. The limiting function in Proposition~\ref{etr:law} is almost surely zero-free on $s(\{|z|\le r'\})$, a compact subset of $\mathcal H$. With probability arbitrarily close to one it is bounded above there by a fixed constant and bounded away from zero. By convergence in distribution and Cauchy's formula, the logarithmic derivative of $G_p(s(z))$ is bounded in probability on $|z|\le r$: the map $g\mapsto(g\circ s)'/(g\circ s)$ there is continuous under uniform convergence on $|z|\le r'$ at every zero-free limit. Integrating this derivative from $0$ makes the normalized local logarithm bounded in probability. The finitely many successive estimates require no value of that logarithm on a member having a zero in the larger disk, since the probability of that event tends to zero by Theorem~\ref{ang:quadratic}. Since $N=N_Ep^2$ and the gamma factor is fixed, the same Cauchy bound makes the maximum relative error bounded in probability by a constant times $r^{-A\log\log D}/\log D$. This factor tends to zero.

For odd $N_E$ take $H=(\log D)^2$ and $\eta=C\log\log D/\log D$, choosing $C$ by Theorem~\ref{etr:strip} for the stated exceptional proportion. Every included twist has $|\beta-1/2|\le\eta$ through $H$. Write $w_\rho=\log|1-1/\rho|$ and $\theta_\rho=\arg(1-1/\rho)$. Horizontal reflection gives the positive phase sum $B_m=\sum_\rho(1-\cos(m\theta_\rho))$ and a possible loss no greater than $\sum_\rho(\cosh(mw_\rho)-1)$. The weight-two bound in the proof of Theorem~\ref{li:holodyadic} gives an absolute constant $c_0>0$ for which
\[
     B_m\ge c_0\,\frac{m^2}{\max(m,1605)}\log D
\]
for all sufficiently large $D$. Indeed the bracket there is $4b_0/5+\log(2\max(m,1605))-5$ with $b_0=\log(N_Ep^2)/2$.

At $|\gamma|\le H$ put $v=\beta-1/2$. Differentiating
$w=\tfrac12\log[((1/2-v)^2+\gamma^2)/((1/2+v)^2+\gamma^2)]$ in $v$ gives $|w_\rho|\le5\eta/(1+\gamma^2)$ for sufficiently small $\eta$. Each derivative term $a/(a^2+\gamma^2)$, $a=1/2\pm v$, is at most $(1/a)/(1+\gamma^2)$ since $a\le1$. The unit-interval zero count is $O_E(\log(2p(2+|\gamma|)))$. Thus $\cosh t-1\le t^2e^{|t|}$ bounds the low-zero loss by $O_E((m\eta)^2e^{5m\eta}\log D)$. At $|\gamma|>H$, $|w_\rho|\le1/(2\gamma^2)$ by the ratio of the squared distances. For $m\le\log D$ summing the unit intervals bounds the high-zero loss by $O_E(m^2 H^{-3}\log(2pH))$. Hence for all $m\le b\log(2D)$ the ratio of the total loss to the positive lower bound is
\[
   \ll_E (\log(2D)+1605)\{\eta^2 e^{5b\eta\log(2D)}+H^{-3}\}.
\]
Choose $b<1/2$ with $5bC<1/2$. This expression tends to zero, proving strict positivity simultaneously.
\end{proof}

\subsection*{Every complete symmetric power in a CM tower}
\begin{theorem}\label{fl:sym}
Fix a CM elliptic curve $E/\mathbb Q$. With the all-height unitary Hecke bound $7/8$ of Section~\ref{recentcm:audit}, there is $b_E>0$ such that for every sufficiently large $n$ all real Li coefficients of the pole-removed complete function of $\operatorname{Sym}^n E$ are strictly positive through $b_E\log n$. In particular there is no exceptional degree in this conclusion.
\end{theorem}
\begin{proof}
Write $L=\log n$, $H=L^2$ and $u_0=K_0\log L/L$, where the fixed $K_0$ will be chosen large. The middle degree-one pole, when present, is removed as in the usual Li zero divisor. All the constituent zero divisors therefore add, including the real Li sums and their positive phase masses. For at least a fixed positive proportion of the parity exponents $j$ one has $n/2\le j\le n$. These degree-two factors have bounded levels and weights $j+1$, so the general-weight positive phase estimate in the proof of Theorem~\ref{li:holodyadic} gives, for every index $m\ge1$,
\[
 B_m(\operatorname{Sym}^n E)\gg_E
                  n L\,\frac{m^2}{\max(m,1605)}.
\]
Its bracket is $0.825H_0-3.85$, where
\[
 H_0=(\log N_j)/2+\log(j/2+2+\max(m,1605))\gg L .
\] All other phase terms are nonnegative.

By Section~\ref{recentcm:transfer} every normalized zero of the full product has $v=|\Re\rho-1/2|\le3/8$. Its possible Li loss is at most $\cosh(mw)-1$, where $w=\log|1-1/\rho|$. The same derivative calculation as in Theorem~\ref{fl:prime} gives
\[
                 |w|\le 16v/(1+\gamma^2).
\]
In fact the two derivative terms have $1/8\le a\le7/8$. Each is at most $8/(1+\gamma^2)$. The unit-interval count of the whole tower is $O_E(n\log(n+|\gamma|+2))$ by its constituent bounds. Zeros at $|\gamma|>H$ thus give total loss $O_E(nm^2H^{-3}\log(n+H))$ when $m\le L$, using $|w|\le1/(2\gamma^2)$. Among the remaining zeros those with $v\le u_0$ cost at most
\[
             O_E\big(nL(m u_0)^2 e^{16m u_0}\big).
\]

For $u_0\le v\le3/8$, Theorem~\ref{uw:sym} and horizontal reflection bound the number with deviation at least $v$ by $O_E(n L^{C_1}e^{-vL})$ at this $H$, absorbing $1+1/u_0$ into a power of $L$. If $m\le b_E L$ with $16b_E\le1/2$, partial summation bounds their total possible loss by
\[
 \ll_E n L^{C_1}
 \left(e^{-(L-16m)u_0}
             +16m\int_{u_0}^{3/8}e^{-(L-16m)v}\,dv\right)
 \ll_E n L^{C_1-K_0/2}.
\]
Take $K_0>2(C_1+4)$ and decrease $b_E$ so that $16b_E K_0<1/2$. The last loss is then $O_E(nL^{-4})$. Relative to the positive lower bound the low-deviation loss is
$O_E((L+1605)u_0^2 L^{16b_E K_0})=o(1)$, and the high-ordinate ratio is $O_E((L+1605)H^{-3})=o(1)$. The full real sums are consequently strictly positive, uniformly for all the stated indices and degrees.
\end{proof}

\clearpage
\section{Common height windows in every absolute degree}\label{dw:section}

Fix a number field $F$ and a unitary cuspidal representation $\pi$ of
$\mathrm{GL}_n(\mathbb A_F)$. Write $d=n[F:\Q]$ and use unitary normalization.
At absolute degree two the choice below puts every zero in a shrinking neighborhood of the critical line. At absolute degree three it puts every zero within a shrinking neighborhood of $[1/3,2/3]$. For cusp forms on $\mathrm{SL}_3(\Z)$ the interval improves to $[3/8,5/8]$.

\begin{theorem}\label{dw:windows}
Let $T_j=2^j$. Fix any positive sequence $H_j\to\infty$ such that $\log H_j=o(\log T_j)$, for instance $H_j=\exp(\sqrt{\log T_j})$. For almost every $v\in[1,2]$ the following holds simultaneously for every $F,n,\pi$ of absolute degree $d\ge2$. Every nontrivial zero $\beta+i\gamma$ with $\bigl||\gamma|-vT_j\bigr|\le H_j$ satisfies
\[
           \operatorname{dist}\bigl(\beta,[1/d,1-1/d]\bigr)=o_{\pi,v}(1)
           \qquad(j\to\infty),
\]
uniformly over the zeros in these two windows. For a Hecke cusp form on $\mathrm{SL}_3(\Z)$ as in the theorem of Dasgupta, Leung and Young, and for each of its unitary norm twists, one can use $[3/8,5/8]$ in the same assertion. Each of the two windows contains $(d/\pi+o_\pi(1))H_j\log T_j$ zeros with multiplicity, where the denominator is the circle constant.
\end{theorem}

Our input for the improved interval is \href{https://arxiv.org/pdf/2407.06962}{Dasgupta--Leung--Young, Theorem 1.1}, which bounds the integral from $-T$ to $T$ for a fixed level-one form. The passage from moments to density is classical. On $\Q$ it is treated by \href{https://doi.org/10.1016/j.jnt.2013.05.012}{Ye--Zhang} and, for degree three, by \href{https://doi.org/10.1093/imrn/rnaf234}{Pal}. In the latter setting Pal, \S1.1.4, records the Ye--Zhang exponent $2(1+m)(1-\sigma)/(3-2\sigma)$ from a second moment of size $T^{m+\epsilon}$. This exponent is less than one exactly when $\sigma>1-1/(2m)$, so the improved level-one boundary follows directly from those published estimates. We give a direct argument to keep track of the field degree, ramification and the quantifiers over twists. It requires no conjectural local bounds.

\subsection{Coefficients over a number field}
Write the rationally indexed Dirichlet series of $L$ and its reciprocal as
\[
 L(s,\pi)=\sum_{k\ge1}a(k)k^{-s},\qquad
 L(s,\pi)^{-1}=\sum_{k\ge1}b(k)k^{-s}\qquad(\Re s>1).
\]
For each $\epsilon>0$ one has
\begin{equation}\label{dw:coeff}
   \sum_{k\ge1}\frac{|a(k)|^2+|b(k)|^2}{k^{1+\epsilon}}<\infty.
\end{equation}
Indeed, at an unramified prime ideal let $\alpha$ be the multiset of Satake parameters. Unitarity identifies the dual multiset with $\bar\alpha$. The degree-$h$ coefficient of the Rankin--Selberg factor for $\pi\times\widetilde\pi$ is
\[
       \sum_{|\lambda|=h}|s_\lambda(\alpha)|^2
\]
by the Cauchy identity for Schur polynomials. The ordinary coefficient at that prime power is $s_{(h)}(\alpha)$. The reciprocal coefficient is $(-1)^h s_{(1^h)}(\alpha)$ for $h\le n$, and zero for $h>n$. Consequently the squares of both sets of good-ideal coefficients are bounded coefficientwise by those of the incomplete Rankin--Selberg product, which converges absolutely to the right of one. At a fixed ramified prime of norm $q$ the standard reciprocal roots satisfy $|\alpha_i|\le q^{1/2}$. Its ordinary degree-$h$ coefficient is therefore $O_n((1+h)^n q^{h/2})$, and the reciprocal polynomial has finite support. This gives a convergent local series of squares at every $1+\epsilon$. Restoring the finitely many such primes proves the square-summability assertion at the level of ideals. For the analytic and local inputs see \href{https://arxiv.org/pdf/2103.05634}{Humphries--Thorner, \S3}.

The number of ideals of norm $k$ is at most the $[F:\Q]$-fold divisor function. At a rational prime $p$ the exponent vectors $(e_i)$ contributing to $p^h$ satisfy $\sum_i f_i e_i=h$, where $f_i$ are residue degrees. The vectors $(f_i e_i)$, padded with zeros, form a subset of the weak compositions of $h$ with $[F:\Q]$ coordinates. Multiplying at rational primes gives the bound. Group by norm, apply Cauchy--Schwarz and absorb this divisor function into $k^{\epsilon/2}$ to obtain \eqref{dw:coeff}. Cauchy--Schwarz also gives absolute convergence for both rational series when $\Re s>1$.

For $R\ge2$ let $M_R(s)=\sum_{h\le R}b(h)h^{-s}$ and let $c_R(k)$ be the coefficients of $K_R(s)=L(s,\pi)M_R(s)$. Dirichlet inversion gives $c_R(1)=1$ and $c_R(k)=0$ for $1<k\le R$. For all $R$,
\begin{equation}\label{dw:cm}
 \sum_{k\le Z}|c_R(k)|^2\ll_{\pi,\epsilon} Z^{1+\epsilon},
 \qquad \sup_u|M_R(1/2+iu)|^2\ll_{\pi,\epsilon}R^{1+\epsilon}.
\end{equation}
For the first inequality use $|c_R(k)|^2\le\tau(k)\sum_{h\ell=k}|b(h)|^2|a(\ell)|^2$ and \eqref{dw:coeff}. For the second use Cauchy--Schwarz and \eqref{dw:coeff}, reserving an arbitrarily small power for $\sum_{h\le R}1/h$.

\subsection{A central mean square}
Suppose $m\ge1$ is such that
\begin{equation}\label{dw:moment}
 \int_{T\le |t|\le 2T}|L(1/2+it,\pi)|^2\,dt\ll_{\pi,\epsilon} T^{m+\epsilon}
 \qquad(T\ge2)
\end{equation}
for every $\epsilon>0$. All $\pi$ considered here admit $m=d/2$. We recall the details that justify this exponent for arbitrary $F$.

In the balanced smooth approximate functional equation on $\Re s=1/2$, the two Dirichlet sums have coefficients $a(k)$ and the coefficients of $\widetilde\pi$, respectively, each divided by $\sqrt{k}$. The latter sum has a scalar of modulus one. The analytic conductor for $|t|\asymp T$ is $\asymp_\pi T^{d}$, since a complex gamma factor counts twice. Given $\kappa>0$, the sums truncate at $A=T^{d/2+\kappa}$ with error $O_{\kappa,B,\pi}(T^{-B})$ for any fixed $B$. The weights $V_t(x)$ on $1\le x\le A$ satisfy
\[
                  |V_t(x)|+|xV'_t(x)|\ll_{\pi,\kappa} T^\kappa.
\]
These are the usual weight bounds in the \href{https://bookstore.ams.org/coll-53}{smooth approximate functional equation of Iwaniec and Kowalski, Theorem 5.3}. Explicitly one may insert $e^{z^2}/z$ in its gamma-ratio integral. On $\Re z=c>0$ fixed, Stirling's formula and the Gaussian give $O_{c,\pi}((T^{d/2}/x)^c)$ for the weight and its first logarithmic derivative; use small $c$ for the displayed bound and large $c$ for the tail together with \eqref{dw:coeff}. On $|\Im z|>T/2$ the Gaussian dominates the reciprocal gamma factors. Possible polar terms for a norm translate of $\zeta_F$ are rapidly decreasing here for the same reason.

Let $S_y(t)=\sum_{k\le y}a(k)k^{-1/2-it}$. Partial summation bounds the first smoothed sum by
$O_{\pi,\kappa}(T^\kappa)(|S_A(t)|+\int_1^A |S_y(t)|\,dy/y)$. The Dirichlet polynomial mean-value theorem gives
\[
       \int_{T\le |t|\le 2T}|S_y(t)|^2\,dt
           \ll (T+y)\sum_{k\le y}|a(k)|^2/k
           \ll_{\pi,\kappa}(T+A)A^\kappa .
\]
Take $L^2$ norms by the integral triangle inequality, and do the same for the dual. Since $d\ge2$, choosing $\kappa$ sufficiently small for the prescribed $\epsilon$ gives \eqref{dw:moment} with $m=d/2$. For the indicated level-one degree-three forms, Dasgupta--Leung--Young give $m=4/3$.

\subsection{From a central mean square to a density bound}
\begin{proposition}\label{dw:density}
Assume \eqref{dw:moment}, put $b_m=1-1/(2m)$ and let $b_m<\sigma<1$. Write $N_\pi(\sigma,T)$ for the number, with multiplicity, of nontrivial zeros with $\Re\rho\ge\sigma$, $|\Im\rho|\le T$. Then
\[
 N_\pi(\sigma,T)\ll_{\pi,m,\sigma} T^{1-\delta},
 \qquad
 \delta=\frac{m(\sigma-b_m)(2b_m-1+\sigma-b_m)}4>0.
\]
\end{proposition}

\begin{proof}
Put $g=\sigma-b_m$, $a=b_m+g/2$, $r=mg/2$, $Y=T^m$ and $R=T^r$. Choose $0<h<\min(\sigma-a,1/4)$ and allow $a\le u\le1+h$. On $T-1\le |t|\le 2T+1$ we approximate $K_R(u+it)$ by
\[
              P(u+it)=\sum_{k\ge1}c_R(k)e^{-k/Y}k^{-u-it}.
\]
Mellin inversion and a shift from $\Re z=2$ to $\Re z=1/2-u$ show that $P-K_R$ is the shifted integral of $K_R(u+it+z)\Gamma(z)Y^z/(2\pi i)$, up to polar terms. Only the gamma pole at zero is crossed. A pole of $L$ at $1+i\tau$ contributes a term involving $\Gamma(1+i\tau-u-it)$, exponentially small in $T$ for fixed $\pi$. The other factors there grow at most as a fixed power of $T$. The new real parts $1/2-u$ stay in a compact subset of $(-1,0)$.

By Stirling's formula the absolute kernel in the new ordinate $w=\Im z$ is bounded by $C_{a,h}Y^{1/2-u}e^{-c|w|}$. For $|w|\le T/4$, the shifted ordinates $t+w$ lie in a fixed union of dyadic intervals with $|t+w|\asymp T$. The $L^2$ triangle inequality, \eqref{dw:cm} and \eqref{dw:moment} imply, uniformly in $u$,
\begin{equation}\label{dw:shift}
 \int_{T-1\le|t|\le2T+1}|P(u+it)-K_R(u+it)|^2\,dt
      \ll_{\pi,a,h,\epsilon}Y^{1-2a}R^{1+\epsilon}T^{m+\epsilon}+O_\pi(T^{-10}).
\end{equation}
Here $|w|>T/4$ contributes less than any fixed negative power by exponential decay and polynomial growth on the central line. The norm translate of $\zeta_F$ has its pole to the right of that line.

The nonconstant coefficients of $P$ start at $k>R$. The weighted Dirichlet polynomial mean-value inequality
$\int_I|\sum d_k k^{-it}|^2dt\ll\sum (|I|+k)|d_k|^2$
is a consequence of \href{https://doi.org/10.1112/jlms/s2-8.1.73}{Montgomery and Vaughan's Hilbert inequality}. It applies to the exponential series by truncation. By \eqref{dw:cm}, partial summation yields
\begin{align}
 \int_{T-1\le|t|\le2T+1}|P(u+it)-1|^2dt
  &\ll_{\pi,a,\epsilon}
     T Y^{-2}+T R^{1-2a+\epsilon}+Y^{2-2a+\epsilon}.
                                                        \label{dw:poly}
\end{align}
The three terms come from the constant $e^{-1/Y}-1$, the $|I|$ part, and the $k$ part, respectively. Take $\epsilon<2a-1$ in the second term. In the third term $a<1$ ensures that dyadic summation up to $Y$ is bounded by its largest scale, after which the exponential decreases.

Before the arbitrarily small powers of $T$ are included, the last two terms of \eqref{dw:poly} have exponents $1-r(2a-1)$ and $1-mg$. The main term of \eqref{dw:shift} has exponent $1-mg+r=1-r$. Since $0<2a-1<1$, all three gains are at least $r(2a-1)=2\delta$. Fix $\epsilon$ small enough in terms of $m,\sigma$ that these estimates, including their sum, leave a gain larger than $\delta$. This order of choices permits the implied constants to depend on the fixed distance from $b_m$.

At each zero $\rho$ with $\Re\rho\ge\sigma$ and $T\le|\Im\rho|\le2T$, the value of $K_R-1$ is $-1$. The disk submean inequality bounds one by the integral of $|K_R-1|^2$ over the disk of radius $h$ about $\rho$, divided by $\pi h^2$. The disks lie in the rectangle ranges used above. Count disks with multiplicity. Their overlap is $O_\pi(\log T)$ by the usual unit-interval zero count for a fixed completed standard function, which also holds on negative ordinates. Integrating in $u$, summing the disk inequalities and absorbing the logarithm proves the stated bound on each dyadic height range. Summing the ranges proves the proposition.
\end{proof}

\subsection{Choosing the same ordinates}
\begin{proof}[Proof of Theorem~\ref{dw:windows}]
Fix one $\pi$, a permitted $m$ and $\sigma>b_m$. The zeros of its completed unitary function are symmetric under $\rho\mapsto1-\bar\rho$, including multiplicity. Thus Proposition~\ref{dw:density} also bounds the number on the left of $1-\sigma$. A zero of ordinate $\gamma$ with $|\gamma|\le3T_j$ excludes a set of $v\in[1,2]$ of length at most $4H_j/T_j$ if we seek the wider windows $\bigl||\gamma|-vT_j\bigr|\le2H_j$. Hence the exceptional measure is
\[
       O_{\pi,\sigma,m}(H_j T_j^{-\delta})=2^{-j\delta+o(j)}.
\]
Its sum in $j$ converges. Borel--Cantelli gives the conclusion for this $\pi,\sigma$ eventually on both windows. Intersect over a sequence of rational $\sigma\downarrow b_m$.

There are countably many number fields up to isomorphism and countably many ranks. For each pair, choose representatives for unitary cusp forms modulo norm twists $|\det|^{i\tau}$. The representatives can be taken countable. Indeed the norm-one idele class group is compact metrizable with countable character group. The extension of each of its characters to the full idele class group is unique up to imaginary powers of the norm, and a central twist by $|\det|^{i\tau}$ runs through all these powers on the center. At fixed central character the discrete cuspidal summands occur in a separable Hilbert space. Countable intersection over these representatives preserves full measure. For the improved degree-three assertion intersect also over the countable level-one representatives with their improved boundary. A fixed norm twist translates zero ordinates by a fixed real amount. For all sufficiently large $j$ the original width $H_j$ of the twist fits inside the width $2H_j$ of its representative. This proves the simultaneous assertions without requiring a uniform conductor estimate.

Finally the fixed-function one-sided Riemann--von Mangoldt formula on either sign of the height is $\frac{d}{2\pi}t\log t+c_{\pi,\pm}t+O_\pi(\log(2+t))$. Subtraction at $vT_j\pm H_j$ gives $(d/\pi)H_j\log T_j+O_\pi(H_j+\log T_j)$, which is the asserted asymptotic since $H_j\to\infty$.
\end{proof}

\begin{corollary}\label{dw:haus}
Fix $m,\sigma$ from Proposition~\ref{dw:density}, with its exponent $\delta$. Let $0\le\kappa<\delta$ and fix $H_j\to\infty$ with $H_j=T_j^{\kappa+o(1)}$. The set of $v\in[1,2]$ for which at least one fixed representation satisfying \eqref{dw:moment} has infinitely many windows meeting a zero outside $[1-\sigma,\sigma]$ has Hausdorff dimension at most
\[
                           \frac{1-\delta}{1-\kappa}.
\]
This includes the union over all number fields and unitary norm twists in this class.
\end{corollary}
\begin{proof}
For one fixed representative consider the doubled windows. At scale $T_j$ there are $O_{\pi,m,\sigma}(T_j^{1-\delta})$ possible zero ordinates, each excluding an interval of length at most $4H_j/T_j$. For each real $s>(1-\delta)/(1-\kappa)$ the sum of the $s$-th powers of these lengths over all sufficiently large scales is bounded by the tail of
\[
                   O_{\pi,m,\sigma,s}(1)\sum_j2^{j(1-\delta-s(1-\kappa)+o(1))}.
\]
The tail tends to zero. The defining limsup set therefore has zero $s$-dimensional Hausdorff measure. The moment bound is invariant under a fixed shift of height, by covering its enlarged dyadic intervals by finitely many dyads. The union over all twists of any representative lies inside that representative's doubled-window limsup set. Finally Hausdorff dimension is stable under countable unions.\end{proof}

The choice $m=d/2$ gives $1-b_m=1/d$, while $m=4/3$ gives $1-b_m=3/8$. Each bound describes all zeros in the selected windows. The intervals between successive windows, and the locations of the individual zeros inside the limiting band, remain unrestricted by this argument.

\clearpage
\section{One collection of intervals carrying almost every zero}\label{blocks:chapter}
Following one choice of normalized height in each dyadic range produces a sparse sequence of intervals. A second selection retains almost the entire height axis. The same intervals then contain a proportion tending to one of the actual zeros of each fixed standard function, while controlling \emph{every} zero they contain.

\begin{theorem}\label{blocks:main}
There is a set $G\subset[2,\infty)$, a union of closed intervals, and an absolute $c>0$ with the following properties. On $T\le t\le2T$, $T=2^j$, its generating intervals have length between $e^{\sqrt{\log T}}$ and $2e^{\sqrt{\log T}}$. For $X\to\infty$,
\[
 |[2,X]\setminus G|\ll Xe^{-c\sqrt{\log X}}.
\]
For every fixed unitary cuspidal standard function over a number field, of absolute degree $d\ge2$, let $N_\pi(X)$ count zeros with $|\Im\rho|\le X$ with multiplicities. Then
\[
 \#\{\rho:2\le|\Im\rho|\le X,\ |\Im\rho|\notin G\}
       \ll_\pi N_\pi(X)e^{-c\sqrt{\log X}}.
\]
Within $G$, every zero of this function at sufficiently large height has real part at distance $o_\pi(1)$ from $[1/d,1-1/d]$. For the level-one degree-three forms of Theorem~\ref{dw:windows}, $[3/8,5/8]$ may be used. For each cuspidal $\mathrm{GL}_2/\Q$ standard function the stronger bound
\[
             |\Re\rho-1/2|\le 2/\sqrt{\log T}
                  \qquad(T\le|\Im\rho|\le2T)
\]
holds on every retained sufficiently high dyadic block. Fixed unitary norm twists are included. Both $G$ and its generating blocks can be chosen once for all these functions. The construction is an existence statement; a uniform threshold in the conductor is not asserted.
\end{theorem}

\begin{proof}
Put $H=e^{\sqrt{\log T}}$ and partition $[T,2T]$ into $\lfloor T/H\rfloor$ equal intervals of length $h$, so $H\le h\le2H$ on all sufficiently large scales. Retain no cells on the finitely many initial scales where this fails. A test consists of a fixed representative modulo unitary norm twist and a specified real-part condition. Mark a cell if it comes within $2h$ of the absolute ordinate of a zero violating that condition. One ordinate marks at most six cells, up to an immaterial endpoint convention; using eight also suffices. A zero marking a cell has $|\gamma|\le3T$.

For a rational degree-two cuspidal representative impose $|\beta-1/2|\le u=2/\sqrt{\log T}$. Theorem~\ref{ang:density} and horizontal functional-equation symmetry show that the fraction of marked cells is bounded by
\[
  C_\pi(1+u^{-1})(\log T)^{C_\pi} H T^{-u/(1-u)}
                     =o_\pi(e^{-\frac12\sqrt{\log T}}).
\]
For every representative of absolute degree $d\ge2$, impose in turn all rational enlargements of $[1/d,1-1/d]$ by a fixed positive amount. Proposition~\ref{dw:density} gives $O_\pi(T^{1-\delta})$ violating zeros on each fixed test, with $\delta>0$ depending on the enlargement and degree. Its marked fraction is $O_\pi(HT^{-\delta})$, again $o_\pi(e^{-\frac12\sqrt{\log T}})$. Include the corresponding enlargements of the improved interval when the extra central moment estimate of Theorem~\ref{dw:windows} applies. Conditions whose enlarged interval contains the strip require no marks.

Enumerate the countably many tests (Theorem~\ref{dw:windows} gives countability of the representatives). On scale $j$ put $\epsilon_j=e^{-\sqrt{\log T}/2}$. Keep the largest initial list of at most $j$ tests for which the sum of their individually marked fractions is at most $\epsilon_j$, allowing the empty list. Its length tends to infinity because every fixed finite list has total marked fraction $o(\epsilon_j)$. Retain every unmarked cell for the chosen list, and take $G$ to be their union on all scales. The omitted length on a dyadic range is at most $\epsilon_jT$. Summing geometrically over lower dyadic ranges gives the asserted estimate for any $X$, for instance with $c=1/4$ after changing the constant. All sufficiently high retained blocks obey each fixed test. A fixed norm twist translates ordinates by a constant, eventually smaller than the $2h$ reserve. This proves the pointwise assertions on $G$, including both signs of the ordinate.

For each fixed function the usual unit-interval zero bound is $O_\pi(\log(2T))$ in the dyadic range and on its negative. The number of zeros in any marked cell is therefore $O_\pi(h\log(2T))$. Since at most $\epsilon_jT/h$ cells are discarded, the omitted zeros number $O_\pi(\epsilon_jT\log(2T))$ per scale. Summing as above and using $N_\pi(X)\asymp_\pi X\log(2X)$ proves the relative zero count. Subtracting the one-sided Riemann--von Mangoldt formulas at the endpoints shows in addition that every retained cell itself contains $(d/(2\pi)+o_\pi(1))h\log T$ zeros on each sign of the height.
\end{proof}

\clearpage
\section{Class numbers in the prime traces of two fixed forms}\label{cl:section}

The trace of a Hecke operator is a sum of class numbers of imaginary quadratic orders. Usually it combines many eigenforms. Two one-dimensional cusp spaces allow us to extract a fixed function without a projector. The first is $S_2(\Gamma_0(11))$, belonging to the non-CM elliptic curve
\[
 E=X_0(11):\quad y^2+y=x^3-x^2-10x-20 .
\]
The second is $S_{12}(\mathrm{SL}_2(\Z))$, spanned by $\Delta(z)=\sum_{n\ge1}\tau(n)e^{2\pi inz}$. The model of the first curve, with discriminant $-11^5$ and nonintegral $j$, fixes its isogeny class. Both forms have trivial character. We will express their unitary prime coefficients by finite quadratic character sums and identify a part for which Dirichlet GRH proves the required cancellation. \href{https://nms.kcl.ac.uk/martin.dickson/bb/notes/cremona_notes.pdf}{Cremona, \emph{Modular curves, forms, elliptic curves, and symbols}.}

Write $H(D)$ for the Hurwitz class number of positive $D$, with value zero unless $-D\equiv0,1\bmod4$. Each class of positive definite integral binary quadratic forms of discriminant $-D$, including imprimitive forms, has weight $2/|\operatorname{Aut}_{\mathrm{SL}_2(\Z)}|$. Thus multiples of the square and hexagonal forms have weights $1/2$ and $1/3$. Let $h_w(\delta)$ be the corresponding primitive weighted class number at the negative discriminant $\delta$. Division by the content gives
\[
 H(D)=\sum_{\substack{m^2\mid D\\-D/m^2\equiv0,1\bmod4}}h_w(-D/m^2).
\]
For a prime $p$ and an integer $|t|<2\sqrt p$ put $D=4p-t^2$. It is positive, so no convention for $H(0)$ is needed. In this section $(b/n)$ denotes the Kronecker symbol, with its usual conventions at even denominators.

\begin{theorem}[The level factor at eleven]\label{cl:level}
At every prime $p\ne11$ the elliptic trace $a_p=p+1-\#E(\F_p)$ equals
\[
 -1-\frac12\sum_{|t|<2\sqrt p}
 \left\{\left(\frac{t^2-4p}{11}\right)H(D)
         +11\,\mathbf1_{121\mid D}\,H(D/121)\right\}.
\]
For the level-one weight-twelve form, at every prime $p$ one has
\[
 \tau(p)=-1-\frac12\sum_{|t|<2\sqrt p}
                 p^5 U_{10}(t/(2\sqrt p))H(D),
\]
where $U_j(\cos u)=\sin((j+1)u)/\sin u$.
\end{theorem}
\begin{proof}
Apply the Eichler--Selberg formula to the trivial character at prime level $11$, weight $2$, with Hecke index $p\ne11$. The scalar term vanishes since $p$ is nonsquare. The exceptional weight-two correction is $\sigma_1(p)=p+1$, while the hyperbolic contribution is $-2$: the divisors $1,p$ each have minimum $1$ and each of the two divisors of $11$ contributes $\varphi(1)=1$, with prefactor $-1/2$. The remaining elliptic term is $-\tfrac12\sum_{t,m}\mu(t,m,p)h_w(-D/m^2)$.

For this formula the local multiplier is $\psi(11)/\psi(11/g)$ times the number of invertible classes $c\bmod 11$ which lift to roots of $c^2-tc+p$ modulo $11g$, where $g=(11,m)$ and $\psi(11)=12$, $\psi(1)=1$. Every root is automatically a unit since $p\ne11$. If $11\nmid m$, the multiplier is $1+((t^2-4p)/11)$. If $11\mid m$, eligibility forces $121\mid D$. The discriminant is zero modulo $121$, and there is exactly one root class modulo $11$ which lifts, namely $t/2\bmod 11$. It has lifts since $2$ is invertible. The multiplier is $12$. Thus in every case
\[
 \mu(t,m,p)-1=((t^2-4p)/11)+11\,\mathbf1_{11\mid m}.
\]
The unit term has sum $\sum_t H(4p-t^2)=2p$ by the classical Kronecker--Hurwitz relation. With the prefactor it cancels $p$ out of $p-1$. The part with $11\mid m$ has sum $H(D/121)$ when $121\mid D$, by putting $m=11m'$. The space has dimension one, so the trace is the eigenvalue of the indicated curve. This proves the first identity, including at $2$ and $3$, since the cited Hecke trace formula assumes only $p\nmid11$. \href{https://www.ams.org/books/surv/133/surv133-endmatter.pdf}{Knightly--Li, \emph{Traces of Hecke operators}, final statement, pp.\ 370--371.} \href{https://arxiv.org/abs/1208.4769}{Brown, Calkin, Flowers, James, Smith and Stout, Hurwitz class number relations.}

For level one the local multiplier is $1$, the weight-two exceptional correction is absent and the two divisors of $p$ contribute $-1$. The elliptic polynomial is the coefficient of $z^{10}$ in $(1-tz+pz^2)^{-1}$, which is $p^5U_{10}(t/(2\sqrt p))$. The eigenvalue on the one-dimensional space is $\tau(p)$. \href{https://people.mpim-bonn.mpg.de/zagier/files/tex/UtrechtLectures/UtBook.pdf}{Zagier, \emph{Elliptic modular forms and their applications}, \S6.3.}
\end{proof}

The term at $121$ is needed only when $p$ is a square modulo $11$. At $p=97$ the two lifts are $t=\pm5$, with $4p-t^2=363=121\cdot3$. The first sum in braces totals $14/3$ and the scalar-lift sum contributes $22/3$ because $H(3)=1/3$. This gives the integer $a_{97}=-7$. Direct enumeration of reduced binary quadratic forms and independent point counts verified the formula at every good $p\le499$, including nineteen primes with nonzero correction. The weight-twelve formula was independently tested against the coefficients of $q\prod_{n\ge1}(1-q^n)^{24}$ at all 95 primes in that range. The algorithm and exact-integer checks accompany the manuscript.

\subsection{One primitive character per trace parameter}

Write $-D=f^2 d$ with $d<0$ a fundamental discriminant, $f\ge1$, and put $\chi_d(n)=(d/n)$. This character is odd and primitive of conductor $q=|d|$. Define
\[
 G(f,d)=\prod_{\ell^e\parallel f}
       \left(1+(1-\chi_d(\ell))\frac{1-\ell^{-e}}{\ell-1}\right),
 \qquad G(1,d)=1 .
\]
Then
\begin{equation}
 H(D)=\frac{\sqrt D}{\pi}G(f,d)L(1,\chi_d).
 \label{cl:conductors}
\end{equation}
Indeed the admissible $m$ in the primitive class number sum are precisely the divisors of $f$, since every discriminant of an order in $\Q(\sqrt d)$ equals $d$ times the square of its conductor. The analytic and order class number formulas give
\[
 h_w((f/m)^2d)=\frac{(f/m)\sqrt q}{\pi}L(1,\chi_d)
          \prod_{\ell\mid f/m}(1-\chi_d(\ell)/\ell).
\]
For the maximal order the class number formula contains its number of units. The order class number formula divides by the unit index between the maximal order and the given order, and the factor $2/|\mathcal O^\times|$ in $h_w$ cancels the remaining units. This includes $d=-3,-4$. Dividing the sum over $m$ by $\sqrt D/\pi$ gives a multiplicative divisor sum whose factor at $\ell^e$ is
$(1-\chi_d(\ell)/\ell)\sum_{j=0}^{e-1}\ell^{-j}+\ell^{-e}$, as claimed. In particular $1\le G(f,d)\ll_\varepsilon f^\varepsilon$: bound each factor by $1+2/(\ell-1)$ and absorb it into $\ell^\varepsilon$ for all sufficiently large $\ell$.

For $u=t/(2\sqrt p)$ define the two signed weights
\[
 J_E(p,t)=((-D)/11)G(f,d)+\mathbf1_{11\mid f}\,G(f/11,d),
 \qquad J_\Delta(p,t)=U_{10}(u)G(f,d),
\]
where the second term for $E$ is omitted when $11\nmid f$. In particular $121\mid D$ is equivalent to $11\mid f$, since $11$ is odd and $d$ fundamental. Substitution into Theorem~\ref{cl:level} yields, with $g=E$ or $\Delta$,
\begin{equation}
 \lambda_g(p)=-c_g(p)
   -\frac1{2\pi}\sum_{|t|<2\sqrt p}\sqrt{D/p}\,
                           J_g(p,t)L(1,\chi_d).
 \label{cl:lvalues}
\end{equation}
Here $\lambda_E(p)=a_p/\sqrt p$, $c_E(p)=p^{-1/2}$, $\lambda_\Delta(p)=\tau(p)/p^{11/2}$ and $c_\Delta(p)=p^{-11/2}$. Thus the conductor-divisor variable has disappeared. Both weights are $O_\varepsilon(p^\varepsilon)$, by $|U_{10}(u)|\le11$ on $[-1,1]$. The character still varies with $p,t$ even when these two variables lie in fixed small residue classes.

\subsection{A finite elementary sum of square-root length}

For a negative primitive quadratic character of conductor $q$ put
\[
 V_q(n)=\frac{e^{-\pi n^2/q}}{n}
           +\frac{\pi}{\sqrt q}\operatorname{erfc}\!\left(n\sqrt{\pi/q}\right),
 \qquad
 \operatorname{erfc}(v)=\frac2{\sqrt\pi}\int_v^\infty e^{-y^2}\,dy .
\]
The exact identity is
\begin{equation}
 L(1,\chi_d)=\sum_{n\ge1}\chi_d(n)V_q(n).
 \label{cl:afe}
\end{equation}
To see this, use the classical odd primitive theta functional equation for $\Theta(v)=\sum_{n\ge1}n\chi_d(n)e^{-\pi n^2v/q}$, namely $\Theta(v)=v^{-3/2}\Theta(1/v)$. Its sign is $+1$ by the primitive quadratic Gauss sum $\tau(\chi_d)=i\sqrt q$. The Mellin transform is $(q/\pi)^{(s+1)/2}\Gamma((s+1)/2)L(s,\chi_d)$. At $s=1$ split at $v=1$ and invert the lower part to get $(q/\pi)L(1,\chi_d)=\int_1^\infty(1+v^{-1/2})\Theta(v)\,dv$. Termwise integration gives precisely \eqref{cl:afe}. The root number, including even conductors $q$, is the one in Rohrlich's quadratic Gauss-sum theorem. \href{https://math.bu.edu/people/rohrlich/pcmi.pdf}{Rohrlich, \emph{Root numbers}, Lecture 1.}

For $v>0$, $\operatorname{erfc}(v)\le e^{-v^2}/(\sqrt\pi v)$ by integration, so $V_q(n)\le 2e^{-\pi n^2/q}/n$. Given $A>0$ and $x\ge2$ use the cutoff
\[
 N(q,x)=\left\lceil\sqrt{(A+5)q\log(2x)/\pi}\right\rceil .
\]
The omitted positive majorant is at most $\int_{N(q,x)}^\infty 2e^{-\pi y^2/q}dy/y\le q e^{-\pi N(q,x)^2/q}/(\pi N(q,x)^2)$. Since $q\le D\le4p\le4x$, all lengths are $O_A(\sqrt{x\log(2x)})$.

\begin{theorem}[A quadratic-character test for each fixed form]\label{cl:triple}
For $p\le x$ restrict to good primes of $g$. Form
\[
 \mathcal C_g(x)=\sum_{p\le x}\log p
       \sum_{|t|<2\sqrt p}\sqrt{D/p}\,J_g(p,t)
           \sum_{n\le N(q,x)} (d/n)V_q(n).
\]
For every fixed $A>0$ one has
\[
 \sum_{p\le x}\lambda_g(p)\log p
       =-\sum_{p\le x}c_g(p)\log p
              -\mathcal C_g(x)/(2\pi)+O_{g,A}(x^{-A}).
\]
In particular fix $\delta>0$. GRH for the standard function of $g$ is equivalent to $\mathcal C_g(x)\le C_\varepsilon x^{1/2+\varepsilon}$ throughout arbitrarily far intervals $[Y,Y^{1+\delta}]$, for each $\varepsilon>0$ with its own constant. The underlying classical zero lines are $1$ for $E$ and $6$ for $\Delta$.
\end{theorem}
\begin{proof}
There are $O(\sqrt p)$ integers $t$ for each $p$. Bound $|J_g|\ll_g p$ crudely, $\sqrt{D/p}\le2$ and $\sum_{p\le x}\log p\ll x$. The uniform truncation error just obtained then contributes $O_{g,A}(x^{-A})$ to \eqref{cl:lvalues} after summation. The elliptic $c_g$ contributes $O(\sqrt x)$ by Chebyshev and partial summation, and the weight-twelve $c_g$ contributes $O(1)$. Apply Theorem~\ref{aw:raw} with $\ell=0$, its logarithmically weighted cutoff and the prime-power deletion in Theorem~\ref{aw:prime}. The sign changes because the class-number sum enters with a minus sign. For $\Delta$ unitary normalization shifts the classical variable by $11/2$, giving real part $6$ at the center. Neither the truncation nor the identity assumes Dirichlet GRH.
\end{proof}

A one-sided truncation of the Dirichlet series has a different cost. Under Dirichlet GRH, shifting the Mellin integral for $\sum_n\chi_d(n)e^{-n/Z}/n$ from positive real part to $-1/2+\epsilon$ gives
\[
 L(1,\chi_d)=\sum_{n\ge1}\frac{\chi_d(n)}{n}e^{-n/Z}
       +O_\epsilon((qZ)^\epsilon Z^{-1/2}).
\]
The gamma factor has one crossed pole, at zero. The usual hybrid Lindel\"of bound on the shifted line follows from Dirichlet GRH, and gamma decay makes the integral absolutely convergent. Bounding every such error separately in \eqref{cl:lvalues} costs $O_\epsilon(x^{3/2+\epsilon}Z^{-1/2+\epsilon})$ in the prime sum. This particular triangle-inequality argument would require $Z$ roughly $x^2$ to preserve a square-root test. The second term of $V_q$ uses the functional equation and achieves the smaller cutoff $\sqrt{x\log x}$ unconditionally. These power counts do not exclude an improvement through cancellation of the character tails across distinct $(p,t)$.

\subsection{The vanishing mean at the level prime}

There is cancellation inside the level-eleven expression before one attempts to sum over the remaining discriminants. In order to keep the arithmetic at $11$ independent of the archimedean weight, return temporarily to the imprimitive Dirichlet series in the order class number formula. Write $L=11$ in this subsection only, and set $\Delta_t=t^2-4p$. Factor an eligible order divisor as $m=L^e m'$ with $(m',L)=1$. Remove the Euler factor at $L$ from each imprimitive quadratic $L(1)$, and write its remaining Dirichlet series with $n'$ coprime to $L$. At those $n'$ the value of the quadratic character is unchanged by division of the discriminant by $L^{2e}$. Eligibility modulo $4$ is unchanged too. All dependence on $e$ and on powers of $L$ in the Dirichlet series therefore combines in the single function
\[
 K(b)=\sum_{\substack{e\ge0\\L^{2e}\mid b}}
 \frac{(b/L)+L\,\mathbf1_{e\ge1}}{L^e}\,
                 \frac{1}{1-((b/L^{2e})/L)/L},
 \qquad b=\Delta_t .
\]
In this formula $(b/L)$ and $((b/L^{2e})/L)$ denote Legendre symbols. The extra $m'^2$ in a discriminant leaves these symbols unchanged since $m'$ is a unit at $L$. Grouping is justified directly by the finite product identity
$L(1,\chi)=(1-\chi(L)/L)^{-1}\sum_{(n,L)=1}\chi(n)/n$ for each nonprincipal quadratic character. The series after omitting an Euler factor converges by periodicity and partial summation. Extending $K$ by its limiting value at zero gives a bounded continuous function on $\Z_L$: the tail for $e\ge E$ is uniformly $O(L^{1-E})$, and each summand is locally constant.

\begin{lemma}[Exact local cancellation]\label{cl:local}
For every prime $p\ne11$, with probability Haar measure on $\Z_{11}$, one has
\[
 \int_{\Z_{11}}K(t^2-4p)\,dt=0.
\]
If $p$ is a nonsquare modulo $11$, the integrand already depends only on $t\bmod11$.
\end{lemma}
\begin{proof}
Let $r=0$ or $2$ be the number of roots of $t^2=4p$ modulo $L$. The character sum $\sum_{t\bmod L}((t^2-4p)/L)=-1$ follows by counting $(t-y)(t+y)=4p$ over the finite field. Thus the numbers of residues at which the symbol is $+1,-1$ respectively are $(L-r-1)/2$ and $(L-r+1)/2$. The contribution from $e=0$ to the integral is
\[
 \frac{L-r-1}{2(L-1)}-\frac{L-r+1}{2(L+1)}
          =-\frac r{L^2-1}.
\]
When $r=0$ there is no other contribution. When $r=2$, each root is simple. The set $L^{2e}\mid t^2-4p$ has measure $rL^{-2e}$, and on each of its root balls $(t^2-4p)/L^{2e}$ is uniformly distributed on $\Z_L$. This follows by Hensel lifting one digit at a time using the unit derivative $2t$. For a uniformly distributed $y$ the mean of $(1-(y/L)/L)^{-1}$ is
\[
 A_L=\frac1L+\frac{L-1}{2L}
           \left(\frac1{1-1/L}+\frac1{1+1/L}\right)
       =\frac{L^2+L+1}{L(L+1)} .
\]
Hence the total of the $e\ge1$ terms is
$rA_L\sum_{e\ge1}L^{1-3e}=r/(L^2-1)$.
It cancels the $e=0$ term exactly. Without the $121$ lifting term in Theorem~\ref{cl:level} the cancellation would fail for quadratic-residue primes.
\end{proof}

For $J\ge1$ let $K_J(t,p)$ be the conditional mean of $K(u^2-4p)$ on $u\equiv t\bmod11^J$, with $p$ fixed. It is bounded uniformly in $J$, has mean zero over $t\bmod11^J$, and depends on $p\bmod11^J$. For the last assertion, if the residue of $t^2-4p$ is nonzero modulo $11^J$, the valuations and unit digits used in $K$ are determined by this residue. Otherwise $t$ is a unit and the normalized variable $(u^2-4p)/11^J$ is uniform on $\Z_{11}$ by the same Hensel argument. This finite projection converges uniformly to $K(t^2-4p)$ as $J$ increases, by uniform continuity of $K$ on $\Z_{11}$. For nonsquare $p\bmod 11$, already $K_1=K$.

\subsection{Square-root cancellation for the finite congruence terms}

Use $e(y)=e^{2\pi i y}$ and
$V_j(u)=\sqrt{1-u^2}\,U_j(u)$ for $|u|\le1$, extended by zero outside. Only $j=0$ and $10$ will be needed. The next estimate concerns a smooth sum on one interval $[P,2P]$. Its constants and its dependence on the period are stated in order to make explicit which part of the trace formula it can control.

\begin{theorem}[Periodic coefficients in the trace formula]\label{cl:modes}
Assume GRH for primitive Dirichlet functions, including the Riemann zeta function. Let $w$ be smooth and compactly supported in $(1,2)$, $Q\ge1$, and $b(t,p)$ a complex function of the two residues modulo $Q$, bounded by $1$. Suppose $j=10$, or suppose $j=0$ and $\sum_{t\bmod Q}b(t,a)=0$ for every unit $a\bmod Q$. Then for $P\ge2$
\begin{equation}
 \left|\sum_{p\nmid Q}\log p\,w(p/P)
       \sum_{t\in\Z}V_j(t/(2\sqrt p))\,b(t,p)\right|
 \ll_{w,j} Q^2\sqrt P\,\log^2(2QP).
 \label{cl:periodic}
\end{equation}
In particular periods bounded by a fixed power of $\log P$ have the square-root exponent.
\end{theorem}
\begin{proof}
First, for a fixed smooth $v$ supported in $(1,2)$, a unit class $a\bmod Q$ and $B\ge1$ one has under the stated GRH
\begin{equation}
 \left|\sum_{p\equiv a\bmod Q}\log p\,v(p/P)e(\pm B\sqrt{p/P})\right|
 \ll_v P/B+\sqrt{PB}\log(2Q(B+2)).
 \label{cl:primephase}
\end{equation}
To see this, expand the congruence indicator in Dirichlet characters. At primitive conductor $q\mid Q$ take the Mellin transform
$W_B(s)=\int v(y)e(\pm B\sqrt y)y^{s-1}dy$. In the variable $z=\log y$ the phase at $s=\sigma+i\gamma$ is $\pm2\pi B e^{z/2}+\gamma z$. Its second derivative is of constant sign and comparable in absolute value to $B$ on the support. The second-derivative estimate gives $|W_B(\sigma+i\gamma)|\ll_v B^{-1/2}$ for $-1/4\le\sigma\le2$. When $|\gamma|>C B$ the first derivative is bounded below by $|\gamma|/2$, and repeated integrations by parts yield $O_{v,N}(|\gamma|^{-N})$ for every fixed $N$. Also $W_B(1)\ll_v B^{-1}$ by integration by parts with $\gamma=0$.

Mellin inversion with $-L'/L$ and shift to $\operatorname{Re}s=-1/4$ crosses the possible principal pole at $1$, the nontrivial zeros, and at most the trivial zero at $0$. The unit-interval zero count is $O(\log(2q(|\gamma|+2)))$. Thus GRH and the preceding transform estimates bound the sum of nontrivial residues by $O_v(\sqrt{PB}\log(2q(B+2)))$. The principal residue is $O_v(P/B)$. The functional equation and the Euler product on $\operatorname{Re}s=5/4$ give $|L'/L(-1/4+i\gamma)|\ll\log(2q(|\gamma|+2))$, uniformly for primitive characters, including conductor one. The left integral is therefore bounded by the same zero-term bound, with $P^{-1/4}$ in place of $\sqrt P$. The bounded trivial-zero term is also admissible. The standard choice of horizontal contour away from ordinates is valid by the unit-interval count and the rapid decay beyond $CB$. Omitting the prime powers $p^k$, $k\ge2$, costs $O(\sqrt P)$ by Chebyshev. Passing from a primitive to an induced character changes only the primes dividing $Q$, whose power contributions in $[P,2P]$ total $O(\log(2Q))$, since each such prime has at most two powers there. Averaging the resulting uniform character estimates proves \eqref{cl:primephase}.

Return to the sum in the theorem. If $Q>\sqrt P$ its elementary bound $O_{w,j}(P^{3/2})$ suffices. Otherwise take $a=p\bmod Q$ and $\widehat b_a(h)=\sum_{t\bmod Q}b(t,a)e(ht/Q)$, of modulus at most $Q$. Poisson summation on the residue classes gives
\[
 \sum_{t\in\Z} V_j(t/(2\sqrt p))b(t,a)
   =\frac{2\sqrt p}{Q}\sum_{h\in\Z}
            \widehat b_a(h)\widehat V_j(2h\sqrt p/Q),
\]
using $\widehat V_j(y)=\int V_j(u)e(-yu)du$.
The $h=0$ term vanishes. For $j=10$ this follows from orthogonality of $U_{10}$ to the constant polynomial against the semicircle weight. For $j=0$ it is the hypothesis on $b$. The identity
\[
 \widehat V_j(y)
     =\pi(j+1)(-i)^j\,\frac{J_{j+1}(2\pi y)}{2\pi y}
\]
follows by putting $u=\cos\theta$, multiplying the two sines, and using the integral representation and recurrence for the Bessel function $J_{j+1}$. Consequently $|\widehat V_j(y)|\ll_j(1+|y|)^{-3/2}$, which justifies the absolutely convergent Poisson series by periodization. For positive $y\ge1$ its expansion is a fixed linear combination of $y^{-3/2}e(y)$ and $y^{-3/2}e(-y)$ with error $O_j(y^{-5/2})$. The negative range gives the same bounds. \href{https://dlmf.nist.gov/10.9}{Bessel integral representations.} \href{https://dlmf.nist.gov/10.17}{Large-argument expansion.}

Truncating the Poisson series at $|h|\le P$ costs $O_{w,j}(Q^{3/2}P^{1/4})$ after summing over primes: per prime the tail is $O_j(Q^{3/2}p^{-1/4}P^{-1/2})$. The Bessel expansion errors inside this cutoff total $O_{w,j}(Q^{5/2}P^{1/4})$, using $\sum_{h\ge1}h^{-5/2}<\infty$. For each nonzero retained $h$ put $d_h(a)=\widehat b_a(h)/Q$, so $|d_h(a)|\le1$. Expand $d_h$ in multiplicative characters on $(\Z/Q\Z)^\times$. By Parseval the sum of squares of its Fourier coefficients is at most one, hence their sum in absolute value is at most $\sqrt{\varphi(Q)}$. Its principal coefficient has modulus at most one. In the derivation of \eqref{cl:primephase} the $P/B$ term occurs only for the principal character. Therefore the prime phase sum with $d_h(p)$ inserted is bounded by $O_w(P/B+\sqrt{QPB}\log(2Q(B+2)))$.

The coefficient of this sum is $O_j(Q^{3/2}P^{-1/4}|h|^{-3/2})$. Its weight is $v(y)=y^{-1/4}w(y)$ and $B=2|h|\sqrt P/Q\ge1$. Consequently the zero term is $O_{w,j}(Q^{3/2}\sqrt P\,|h|^{-1}\log(2QP))$. Summing over $h$ costs one more logarithm. The principal term costs $O_{w,j}(Q^{5/2}P^{1/4}|h|^{-5/2})$. This and the Bessel error are $O_{w,j}(Q^2\sqrt P)$ because $Q\le\sqrt P$. This proves the theorem.
\end{proof}


\begin{corollary}[Sharp cutoffs for the same terms]\label{cl:sharp}
Under the hypotheses of Theorem~\ref{cl:modes}, for $2\le y\le x$ with $x/2\le y$ the left side of \eqref{cl:periodic}, summed over $y<p\le x$ with no factor $w$, is $O_j(Q^2\sqrt x\log^3(2Qx))$. By summing dyadic intervals the same bound holds through $p\le x$.
\end{corollary}
\begin{proof}
Put $P=x/2$ and ignore bounded $x$ by increasing the constant. Let $m,n$ be the smallest and largest integers of $(y,x]$. If the interval contains none the sum is zero. Convolve its enlarged rescaled interval $[(m-1/2)/P,(n+1/2)/P]$ with a smooth nonnegative mollifier of integral one supported in $[-1/(4P),1/(4P)]$. The result takes value exactly one at $k/P$ for $m\le k\le n$ and vanishes at every other integer point $k/P$. Its support stays in a fixed compact subset of $(0,\infty)$. The function has supremum at most one, first derivative of bounded $L^1$ norm independent of the endpoints, and second derivative of $L^1$ norm $O(P)$. These assertions follow by writing its first derivative as a difference of translates of the mollifier. Multiplication by $u^{-1/4}$ and change to $\log u$ preserve those bounds up to absolute constants on this support.

In the proof of \eqref{cl:primephase}, the second-derivative phase estimate costs only the supremum and total variation of the amplitude. Hence the bound $O(B^{-1/2})$ still applies, as does $O(B^{-1})$ at the principal pole. For $|\gamma|>CB$, one or two integrations by parts now give
\[
 |W_B(\sigma+i\gamma)|\ll
        \min\{|\gamma|^{-1},P|\gamma|^{-2}\}
                 \quad(-1/4\le\sigma\le2).
\]
Indeed the first phase derivative has size at least $|\gamma|/2$, subsequent phase derivatives are $O(B)$ there, and the first two amplitude derivatives have the stated integral bounds. Sum the first bound over unit intervals through $|\gamma|\le P$ and the second in the remaining tail. The unit zero count, or the analogous logarithmic majorant of the left integral, gives $O(\log^2(2Q(P+B+2)))$ for this far range. Thus the sharp version of \eqref{cl:primephase} holds with $\log(2Q(B+2))$ replaced by $\log^2(2Q(P+B+2))$. For each fixed $P$ the mollifier is smooth, so the contour shift remains justified. Cut off the Poisson series at $|h|\le P$ exactly as before, so $B\le 2P^{3/2}$. The same proof yields the stated bound with the third logarithm. Restricting to $p\nmid Q$ causes no difficulty even for the global version: the total absolute contribution from the finitely many $p\mid Q$, if they are restored with bounded coefficients, is $O_j(\sqrt x\log(2Q))$ since there are $O(\sqrt p)$ terms per such prime.
\end{proof}

The exponent two of the period in the general theorem cannot be decreased by a fixed amount, even at a smooth cutoff. Let $Q$ tend through integers greater than two, put $P=Q^2$, $j=0$, and choose $b(t,p)=\cos(2\pi t/Q)$, which has zero mean. The terms $h=\pm1$ in Poisson summation together give $2\sqrt p\,\widehat V_0(2\sqrt p/Q)$. All other nonzero modes are congruent to $\pm1\bmod Q$ and cost $O(Q^{-1/2})$ per prime by the $|y|^{-3/2}$ transform bound. Choose a fixed nonnegative $w$, not identically zero, supported where $\widehat V_0(2\sqrt y)$ has fixed nonzero sign; such an interval exists in $(1,2)$ by real analyticity and the nonzero value of the transform at zero. The prime number theorem gives a sum of size $\gg P^{3/2}$. An upper bound $Q^{2-\nu}\sqrt P(\log(2QP))^{O(1)}$ with any fixed $\nu>0$ would contradict this example. The example tests arbitrary periodic coefficients and does not assert optimality for the specific class-number coefficients.

To place the result inside the class-number expansion for $\Delta$, fix $m,n$ and take the coefficient
\[
 b(t,p)=\mathbf1_{m^2\mid 4p-t^2}
       \mathbf1_{(t^2-4p)/m^2\equiv0,1\bmod4}
             \left(\frac{(t^2-4p)/m^2}{n}\right),
\]
defined as zero when the divisibility fails. This is periodic in $t,p$ modulo $8m^2n$; this modulus determines both the order condition and the Kronecker symbol, including at $2$. It is exactly the mode from using the primitive-order divisor $m$ and the $n$th term of its imprimitive $L(1)$ series. The archimedean factor in \eqref{cl:lvalues} before primitivizing is $2V_{10}(t/(2\sqrt p))$. For $E$, keep $m,n$ prime to $11$, take the same coefficient times $K_J(t,p)$, and use $2V_0$. Its period divides $Q=8m^2n\,11^J$. By the Chinese remainder theorem and Lemma~\ref{cl:local} it has exact zero mean in $t$ for each $p$ prime to $Q$. On nonsquare classes $p\bmod11$ the coefficient with $J=1$ is the full local factor. A restriction to any such class is itself periodic and preserves that mean.

It follows that the complete sharp prime sum of these terms, with $m\le M$ and $n\le N$ and their actual weights $1/(mn)$, is bounded under Dirichlet GRH by
\begin{equation}
 \ll_g \sqrt x\,\log^3(2xMN\,11^J)\,M^4N^2\,11^{2J}.
 \label{cl:lowpower}
\end{equation}
In the weight-twelve case omit $11^J$. To see this, use Corollary~\ref{cl:sharp} with $Q\le8m^2n\,11^J$. Summing the period cost gives $\sum_{m\le M,n\le N}m^4n^2/(mn)\ll M^4N^2$. The exceptional $p\mid Q$ satisfy the bound at the end of that proof, or can be omitted and put back directly. For $E$ the result concerns the finite conditional projection at $11$ over all primes, or the exact level factor on nonsquare classes. The estimate is for genuine finite terms of the imprimitive order-class-number series. Its weights are independent of the conductor-dependent Gaussians used for the separate short formula in Theorem~\ref{cl:triple}.

In particular, for each $\varepsilon>0$ put $\theta=\min(1,\varepsilon)/10$ and take $M,N$ between $1$ and $x^\theta$, and $11^J\le\max(11,x^\theta)$ if needed. The bound in \eqref{cl:lowpower} is $O_{g,\varepsilon}(x^{1/2+\varepsilon})$ at \emph{every} sharp cutoff. This gives an explicit positive power of the order and character ranges that can already be removed at the target error by Dirichlet GRH. More precisely, let $\mathcal L_g(x)$ be the resulting rectangular sum with the archimedean factor $2V_j$, the period coefficient just described and weights $1/(mn)$. One may take $M=N=\max(1,\lfloor x^\theta\rfloor)$ and take $J$ the largest allowed positive integer. The pointwise bound applies with this choice even though the cutoffs change with $x$.

Under Dirichlet GRH the fixed-form GRH in Theorem~\ref{cl:triple} is therefore equivalent to the same one-sided sparse-window bound for the entirely finite difference $\mathcal C_g(x)-\mathcal L_g(x)$, with the above $\theta$ chosen for the tested $\varepsilon$. In the forward direction bound $\mathcal C_g$ by the theorem and bound $\mathcal L_g$ in absolute value by \eqref{cl:lowpower}. In the reverse direction add the same absolute majorant back. The difference represents, up to the controlled Gaussian truncation error, the remaining part of the order-class-number series. It retains growing character indices and discriminant conductors, and for $E$ at square residue classes modulo $11$ also the unretained higher $11$-adic digits. The loss from the square of the period prevents \eqref{cl:lowpower} from reaching the full range. The signed cancellation of this remaining explicit difference is the unproved input, separately for each fixed form.

\subsection{Removing unbounded square parts of the character indices}

One need not restrict the whole index $n$ of the imprimitive series. In each term take $n=r b^2$ with $r$ squarefree. For the character $\chi_\delta=(\delta/\cdot)$ the exact sum over all $b\ge1$ is
\begin{equation}
 \sum_{b\ge1}\frac{\chi_\delta(rb^2)}{rb^2}
 =\frac{\chi_\delta(r)}r\,\zeta(2)
       \prod_{\ell\mid\delta}(1-\ell^{-2})
 =\frac{\chi_\delta(r)}r\,\zeta(2)\sum_{h\mid\delta}\frac{\mu(h)}{h^2}.
 \label{cl:squares}
\end{equation}
Here $\mu$ denotes the M\"obius function. The identity is absolute for each $r$. To check that grouping the full series by $r$ is legitimate, first work at real part $s>1$, where both products are absolutely convergent. At the boundary the partial sums of $\mu^2(r)\chi_\delta(r)$ are $O_\delta(\sqrt y)$ up to $y$, by the identity $\mu^2(r)=\sum_{v^2\mid r}\mu(v)$ and bounded partial sums of the nonprincipal periodic character. Hence the squarefree series itself converges at $1$ by partial summation, giving the same boundary value. For the elliptic expansion outside $11$ impose $(r b,11)=1$. In \eqref{cl:squares} this replaces $\zeta(2)$ by $\zeta(2)(1-11^{-2})$ and restricts $h$ to $(h,11)=1$.

\begin{corollary}[Full square parts, bounded squarefree kernels]\label{cl:squarefree}
Assume Dirichlet GRH. In the sharp class-number prime sum retain order indices $m\le M$, squarefree kernels $r\le R$, and the \emph{entire} sum over $b$ in \eqref{cl:squares}. For $E$ use the normalized local projection $K_J$ at eleven as above. The resulting contribution satisfies
\[
 \ll_g \sqrt x\, M^2 R^{3/2}\,11^{2J}
              \log^7(2xMR\,11^J).
\]
Omit $11^J$ for $\Delta$. Thus one can allow $M,R,11^J$ up to $x^{\min(1,\varepsilon)/6}$ (with $11$ as a constant lower cutoff if needed) and still have $O_{g,\varepsilon}(x^{1/2+\varepsilon})$ at every sharp cutoff. The equivalent one-sided test in Theorem~\ref{cl:triple} also holds upon subtracting this enlarged finite contribution, evaluated by the final divisor sum in \eqref{cl:squares}.
\end{corollary}
\begin{proof}
Fix $m,r,h$, with $r,h$ squarefree. The condition $h\mid\delta$ for $\delta=(t^2-4p)/m^2$ restricts $t$ to the roots of $t^2\equiv4p\bmod m^2h$. Write $D_0=m^2h$. For $p\nmid2D_0$ the number of roots modulo $D_0$ is $O(\tau(mh))$, where $\tau$ here is the divisor-counting function. At each odd prime dividing $D_0$ there are at most two simple roots by Hensel's lemma. At $2$ there are at most eight roots: at large exponents divide by four and use the at most four odd roots of a unit square modulo a power of two, allowing the extra factor two from lifting back. The lower exponents satisfy the same bound. Apply the Chinese remainder theorem. The full periodic coefficient, including its order condition, symbol $(\delta/r)$ and any level projection, has a period dividing $Q=8D_0r\,11^J$, where the last factor is omitted at level one. The zero-frequency term in Poisson summation vanishes by $U_{10}$ or by the independent local mean at eleven.

Put $r_0=r/(r,22mh)$. For each odd prime $\ell\mid r_0$ the $\ell$-component of the periodic coefficient is $((t^2-4a)/\ell)$ at a unit $a=p\bmod Q$, because $m^2$ is a nonzero square modulo $\ell$. For $k\not\equiv0\bmod\ell$ its unnormalized additive transform is
\[
 \sum_{t\bmod\ell}\left(\frac{t^2-4a}{\ell}\right)e(kt/\ell)
   =\sum_{u\bmod\ell}^{*}e\bigl(k(u+a/u)/\ell\bigr).
\]
Indeed the number of $y$ with $y^2=t^2-4a$ is one plus the Legendre symbol. Parametrizing the resulting pair by $u=(t+y)/2$ gives $t=u+a/u$. The constant term sums to zero at nonzero frequency. At zero frequency the original transform equals $-1$ by the same count. Weil's Kloosterman bound therefore bounds all the normalized transforms by $2/\sqrt\ell$. \href{https://www.renyi.hu/~gharcos/weil.pdf}{Weil's bound, Theorem 1 in Harcos's notes.}

The Chinese remainder theorem factors the normalized additive transform into those at the primes $\ell\mid r_0$ and the complementary modulus $Q/r_0$, with rescaled frequencies. The complementary coefficient is bounded and supported on the roots modulo $D_0$ already counted. Thus at every frequency $k$, including nonzero multiples of $Q$, and at every unit $a$,
\[
 |d_k(a)|\ll A_{m,r,h},\qquad
 A_{m,r,h}=\frac{\tau(mh)\tau(r)}{m^2h\sqrt{r_0}}.
\]
For $Q\le\sqrt P$ all estimates in the proof of Corollary~\ref{cl:sharp} are linear in this uniform bound, from the Bessel remainder through the multiplicative Parseval estimate. For $Q>\sqrt P$ absolute Poisson summation gives the bound
\[
 \sqrt p A_{m,r,h}\sum_{k\in\Z}
       (1+2|k|\sqrt p/Q)^{-3/2}\ll A_{m,r,h}Q
\]
at each prime $p\asymp P$ of unit class. Summing with $\log p$ costs $O(P)$, and $QP\le Q^2\sqrt P$. Thus for every period the sharp dyadic bound at unit primes is $O_g(A_{m,r,h}Q^2\sqrt P\log^3(2QP))$. The exceptional primes dividing $Q$ contribute at most $O_g(\sqrt P\log(2Q))$ without imposing a support or Fourier bound. Summing dyadic intervals gives the same estimates with $x$ in place of $P$.

Multiply by the weights $1/(mrh^2)$. Since $h\mid\delta$ one has $h\le4x$. After taking out $\sqrt x\,11^{2J}\log^3(2xMR\,11^J)$ it suffices to sum
\[
  \sum_{m\le M}\sum_{h\le4x}\frac{m\tau(mh)}{h}
      \sum_{r\le R}\tau(r)\sqrt r\,(r,22mh)^{1/2}.
\]
To bound the inner sum put $B(v)=\sum_{d\mid v}\tau(d)/\sqrt d$, a submultiplicative function, by distributing each divisor of a product among its two factors and using the submultiplicativity of $\tau$. Majorize $(r,v)^{1/2}$ by $\sum_{d\mid(r,v)}\sqrt d$, write $r=de$ and use $\sum_{e\le z}\tau(e)\ll z\log(2z)$. The inner sum is $O(R^{3/2}\log(2R)B(22mh))$. Also
\[
 \sum_{v\le z}\tau(v)B(v)\ll z\log(2z)
                    \sum_{d\ge1}\frac{\tau(d)^2}{d^{3/2}}
             \ll z\log(2z),
\]
by writing $v=de$ and applying the same divisor estimate. Submultiplicativity at $22mh$ and partial summation now bound the $m$-sum by $O(M^2\log(2M))$ and the $h$-sum by $O(\log^2(2x))$. Including the logarithms outside gives the claimed seventh power. The exceptional-prime bound is absolutely summable in $h$ with weight $1/h^2$ and fits the same majorant. The stated common cutoff uses $2+3/2+2=11/2<6$, leaving room for all logarithms. Subtraction and reinsertion of a term uniformly bounded at the tested exponent preserve the one-sided equivalence.
\end{proof}

There is also an elementary estimate at the opposite end of the order range. Return to the complete imprimitive order expansion, with its exact level factor $K$ for $E$. For $1\le M\le 2\sqrt x$ the portion from $m>M$, including \emph{all} its character indices, has absolute value at most
\begin{equation}
 \ll_g\sqrt x\,(1+x/M^2)\log^3(2x).
 \label{cl:upperorder}
\end{equation}
Indeed each eligible triple satisfies $4p=t^2+m^2k$ with $1\le k\le4x/m^2$ and $|t|\le2\sqrt x$. For fixed $t,m,k$ this determines $p$. The sharp prime logarithm is at most $\log x$, both the archimedean factors and $K$ are bounded, and each imprimitive $L(1,\chi_\delta)$, with the factor at $11$ deleted if necessary, is $O(\log(2x))$. The last assertion follows by summing the harmonic series through a period of length $O(x)$ and bounding the tail by partial summation with character sums bounded by the period. The negative discriminant ensures a nonprincipal character, and deleting the Euler factor changes the bound by a constant. The remaining weight is $1/m$. Hence the contribution is $O_g(\sqrt x\log^2(2x)\sum_{M<m\le2\sqrt x}(1+x/m^2)/m)$, giving \eqref{cl:upperorder}. For example all $m>x^{1/2-\min(1,\varepsilon)/4}$ meet $O_{g,\varepsilon}(x^{1/2+\varepsilon})$. This estimate is unconditional.

The same counting gives a quantitative cost for approximating the full level function. With $L=11$, the difference $K(t^2-4p)-K_J(t,p)$ vanishes unless $L^J\mid t^2-4p$. On this set its modulus is $O(L^{1-J/2})$, because every summand with $2e<J$ in the local definition is fixed by the residue and the remaining geometric tail has this bound. For all $m$ at a fixed $(p,t)$ the order weights satisfy $\sum_{m^2\mid 4p-t^2}1/m\ll\log(2x)$. Counting $4p=t^2+L^Jk$ for each $t$, with $1\le k\le4x/L^J$, bounds the error in the entire sharp prime sum by
\[
 \ll \sqrt x\,L^{1-J/2}(1+x/L^J)\log^3(2x).
\]
Thus an $L^J\ge x^{2/3}$ projection already recovers the full level factor to the square-root exponent, with all other order and character components included. Corollary~\ref{cl:squarefree} controls the lower joint range of moduli. The ranges between these estimates still require cancellation across varying squarefree characters.

\subsection{Summing every level-eleven digit}

The period $11^J$ can in fact be removed from the final estimate for the elliptic curve. We need a second cancellation condition, this time in the prime residue. Write
\[
 d_k(a)=Q^{-1}\sum_{t\bmod Q} b(t,a)e(kt/Q),\qquad
 A=\max_{k,\,a\in(\Z/Q\Z)^\times}|d_k(a)|.
\]
Assume the Poisson zero term vanishes as in Theorem~\ref{cl:modes}. Assume also that each $d_k$ has sum zero over the unit classes $a$. This condition holds, for example, if the sum of $b(t,a)$ over the unit $a$ vanishes for every $t$. Then the sharp dyadic bound at primes $p\nmid Q$ improves to
\begin{equation}
 \ll_j A Q^{3/2}\sqrt P\log^3(2QP)
 \label{cl:noprincipal}
\end{equation}
for \emph{all} $Q\ge1$. In particular the estimate retains the cancellation in $b$ when $Q>\sqrt P$.

To prove this, the multiplicative expansion of $d_k$ has sum of absolute coefficients at most $A\sqrt Q$ and has zero principal coefficient. We use the exact Hankel decomposition of the Bessel function: for $y\ge1/2$ the transform $\widehat V_j(y)$ is a sum of $e(y)a_+(y)$ and $e(-y)a_-(y)$, with $|y^v a_\pm^{(v)}(y)|\ll_j y^{-3/2}$ for $v=0,1,2$. Indeed $J_{j+1}$ is half the sum of the two Hankel functions. Their differentiable large-argument expansions after removing the exponentials give these bounds, and compactness extends them down to $1/2$. The same assertion follows by using the expansions of the functions and their first two derivatives to sufficiently many terms. \href{https://dlmf.nist.gov/10.17}{Hankel function and derivative expansions, DLMF 10.17.}

Put $B=2|k|\sqrt P/Q$. For $B\ge1$ use this exact decomposition in the prime sum. After extracting $\sqrt P B^{-3/2}$ the weight, before the cutoff, has two bounded derivatives in the scaled prime variable. The proof of Corollary~\ref{cl:sharp} therefore applies uniformly, with no principal residue and no Bessel expansion remainder. The contribution per $k$ is bounded by
\[
 \sqrt P B^{-3/2}\, A\sqrt Q\sqrt{PB}
     \log^2(2Q(P+B+2))
 \ll A Q^{3/2}\sqrt P\,|k|^{-1}\log^2(2Q(P+B+2)).
\]
If $0<B<1$ instead retain $\widehat V_j(B\sqrt y)$ directly as a smooth amplitude. Its first two derivatives are uniformly bounded on the compact support; the zero-time Mellin argument for nonprincipal characters gives $O(\sqrt P\log^2(2QP))$ before inserting the Fourier and $\sqrt P$ factors. There are $O(Q/\sqrt P)$ such nonzero indices. Together they cost $O(AQ^{3/2}\sqrt P\log^2(2QP))$. The sharp cutoff is the same exact half-integer mollification. Truncate the remaining modes at $|k|\le \lceil QP^2\rceil$. Beyond this point the absolute Poisson tail over the primes is $O_j(AQ^{3/2}P^{3/4}(QP^2)^{-1/2})$. Up to this point $B\ll P^{5/2}$ and the harmonic sum costs $O(\log(2QP))$. This proves \eqref{cl:noprincipal}.

For the local function set $L=11$ and write $D_j(t,p)=K_j(t,p)-K_{j-1}(t,p)$ for $j\ge2$. At unit primes its support lies in $L^{j-1}\mid t^2-4p$ because $K$ is constant on every nonzero residue ball modulo $L^{j-1}$. There $t$ is a unit. Terms indexed by $2e<j-1$ in $K$ are fixed by this residue, so $|D_j|\ll L^{-j/2}$ with a fixed $L$-dependent constant. By the simple-root count the support occupies at most $2L^{1-j}$ of the $t$ residues. Thus every normalized additive Fourier coefficient in $t$ of $D_j$ is $O(L^{-3j/2})$.
Furthermore
\[
  \sum_{a\bmod L^j}^{*} D_j(t,a)=0\qquad(t\bmod L^j).
\]
For $t$ a nonunit the support is empty. Otherwise, on its support, $u\mapsto u^2-4a$ maps the ball $t+L^i\Z_L$ bijectively and measure-preservingly onto $t^2-4a+L^i\Z_L$ for $i\ge1$. Thus $K_i(t,a)$ there is the conditional mean of $K$ on the corresponding residue ball of its argument. With $t$ fixed, $a\mapsto t^2-4a$ ranges once through the subballs of the supported ball $L^{j-1}\Z_L$ as $a$ ranges through the supported unit residues modulo $L^j$. Subtracting the coarser conditional mean has sum zero.

\begin{theorem}[All level digits with bounded squarefree kernels]\label{cl:alllevel}
Under Dirichlet GRH, the complete contribution in Corollary~\ref{cl:squarefree} for $X_0(11)$, with the full level factor $K$ and all square parts $b$, is
\[
  O\!\left(\sqrt x\,M^2R^{3/2}\log^8(2xMR)\right).
\]
For $\Delta$ the same estimate already holds with seventh logarithmic power. In particular $m,r\le x^{\min(1,\varepsilon)/4}$ are allowed simultaneously for each form with no remaining local truncation.
\end{theorem}
\begin{proof}
With the notation of Corollary~\ref{cl:squarefree}, $Q_0=8m^2hr$ is prime to $11$ for the elliptic sum. The component without the level factor has normalized additive Fourier coefficients bounded by a constant times $A_{m,r,h}$. Multiplying by $D_j$ yields a period $Q=Q_0 L^j$. By the Chinese remainder theorem its Fourier coefficients are $O(A_{m,r,h}L^{-3j/2})$ and its principal coefficient in the prime residue is zero at every frequency. Its zero Poisson term vanishes by the zero $t$-mean of $D_j$. Apply \eqref{cl:noprincipal}. The bound at unit primes becomes $O(A_{m,r,h}Q_0^{3/2}\sqrt x\log^3(2Qx))$, uniformly in the level depth $j$. The level-$1$ projection $K_1$ costs $O(A_{m,r,h}Q_0^2\sqrt x\log^3(2Q_0x))$ by the earlier estimate and the fixed factor $L$. Between them these powers are at most $Q_0^2$. For the eligible $p,t$, $1\le4p-t^2\le4x$, so $D_j$ vanishes identically when $L^{j-1}>4x$, and the sum of the remaining $D_j$ and $K_1$ is exactly $K$. There are $O(\log(2x))$ such depths. For $p\mid Q$ the exception bound decays by $L^{-j/2}$ and $p=L$ is omitted throughout. The weighted sum over $m,r,h$ is now precisely the divisor estimate of Corollary~\ref{cl:squarefree} with its fixed level modulus and one additional logarithm. Finally $2+3/2<4$ absorbs the logarithms in the asserted growing cutoff.
\end{proof}

\subsection{What changes after base extension}

For a fixed Dirichlet character $\xi$ one obtains the corresponding statement for the twist of $g$ by inserting $\operatorname{Re}\xi(p)$ throughout the signed prime sums and deleting the finitely many new bad primes. Indeed $\lambda_g(p)$ is real, so this product is the real part of the twisted prime trace. Theorem~\ref{aw:raw} tests all zeros of the twist even if $\xi$ is nonreal. The additional multiplier is periodic, depends on $p$ alone and preserves the zero-mean conditions in Theorem~\ref{cl:modes}. It changes the period to at most $Q\,\operatorname{cond}(\xi)$, a fixed additional factor. Thus Dirichlet GRH gives every such fixed twist the same sharp bound for its low terms.

There is a direct version for the curves and forms considered over larger base fields. For any number field $F$ put $r_F(p)=\#\{\mathfrak p\mid p:f_{\mathfrak p}=1\}$. Outside the primes ramified in $F$ and the fixed bad primes of $g$, the sum of the normalized degree-one traces over $p$ is exactly $r_F(p)\lambda_g(p)$. This is unconditional equality of the good local data. For $E$ it follows immediately by reduction over the same prime field. For $\Delta$ it follows by restriction of its two-dimensional Galois representation. It follows that the identical finite sum with multiplier $r_F(p)$ gives the zero criterion over $F$ whenever the resulting standard function has the analytic continuation, functional equation and boundary nonvanishing used above. Finite-place purity again bounds the higher powers and primes of higher residue degree by $O_F(\sqrt x)$ in the weighted unitary sharp cutoff.

If $F/\Q$ is abelian, Kronecker--Weber gives a fixed modulus $M_F$ such that $r_F(p)$ is $[F:\Q]$ times a splitting indicator periodic modulo $M_F$ at the eligible primes. The preceding low-term estimate therefore follows from Dirichlet GRH alone. In this case the analytic input is known: the base-change function factors into the cusp-form twists indexed by the characters of $\operatorname{Gal}(F/\Q)$, by finite-dimensional induction and Artin formalism for the associated Galois representation. Equivalently one can use the abelian base-change theorem. The factors are entire and have the usual functional equations. Their zero orders add, so the proof of the one-sided test for this finite product is unchanged. For elliptic curves the local induction identity, including bad primes, is detailed in Theorem~\ref{ebc:divisor}.

Even for nonabelian $F/\Q$, the low-term estimate does not require a conjecture about zeros of twists of $g$. A sufficient additional hypothesis, beyond Dirichlet GRH over $\Q$, is GRH for the Dedekind zeta functions of $F(\zeta_Q)$ for the positive integers $Q$ occurring as periods. To see this, replace the prime sum with Dirichlet character $\chi$ in the proof of \eqref{cl:primephase} by the logarithmic derivative of the finite-order Hecke character $\chi\circ\mathrm N_{F/\Q}$. It gives the coefficient $r_F(p)\chi(p)\log p$ at norm $p$ outside the excluded primes. The total of higher prime powers and higher residue-degree terms on $[P,2P]$ is $O_F(\sqrt P)$ by Chebyshev, as in Theorem~\ref{aw:prime}. Primes above the finite excluded set have the same harmless correction as in the rational argument. The finite conductor of the norm-pulled character divides $Q\mathcal O_F$ because a local unit congruent to one modulo $Q$ has norm congruent to one modulo $Q$; its archimedean factors have fixed parity types. The primitive Hecke functions therefore have unit-interval zero counts $O_F(\log(2Q(|\gamma|+2)))$. Their nontrivial zeros lie on the line under the stated hypothesis, since $F(\zeta_Q)/F$ is abelian and its zeta function factors into the entire nontrivial Hecke factors and $\zeta_F$. The primitive characters underlying these norm-pulled characters are among these factors.

The number of Dirichlet characters modulo $Q$ whose pullback is the trivial Hecke character is at most $[F:\Q]$. In fact it is the index of $\operatorname{Gal}(F(\zeta_Q)/F)$ in $(\Z/Q\Z)^\times$, which is $[F\cap\Q(\zeta_Q):\Q]$. Thus only boundedly many principal residues enter where the proof used the unique principal Dirichlet character. The functional equation, contour argument and Parseval factor give exactly \eqref{cl:periodic} and its sharp version with the inserted factor $r_F(p)$ and constants also depending on $F$. This describes the precise additional degree-one zero information for this argument over a general field. It leaves the same correlation in the larger character and conductor ranges.


The full local result of Theorem~\ref{cl:alllevel} also persists with these fixed base and twist factors. For a fixed periodic multiplier in $p$, choose $j$ beyond one plus the $11$-adic valuation of its modulus. At fixed residues modulo the other primes and $11^{j-1}$, the mean of $D_j$ in the last prime-residue digit still vanishes. This removes its principal projection as before, and the finitely many smaller depths have fixed cost. Over general $F$ the characters whose pullback is principal factor through the fixed abelian field $F\cap\Q^{\mathrm{ab}}$. Their conductors therefore divide a fixed integer by Kronecker--Weber. Their projections of $D_j$ vanish by the same final-digit calculation for all larger $j$, so \eqref{cl:noprincipal} applies there under the stated Dedekind hypothesis.

\clearpage
\section{Prime generators of quadratic orders in the conductor-eleven trace}
\label{op:section}

The level factor in the class-number formula lets us count the large order divisors without truncating any Dirichlet series. The primes to be counted are generators in varying imaginary quadratic orders. This brings the Hecke input of Section~\ref{recentcm:section} directly into a sum for the fixed non-CM elliptic curve. We use $L=11$ throughout, and suppose that all unitary Hecke functions over imaginary quadratic fields have no zeros in $\operatorname{Re}s>\theta$, with $1/2<\theta<1$. The endpoint $\theta=7/8$ is the input audited in that section. Finite-order characters suffice for an order discriminant cutoff. The order-divisor cutoff also uses characters of nonzero angular type.

Let $\delta<0$ always denote an order discriminant, so $\delta\equiv0,1\bmod4$. Let $h_\delta$ and $w_\delta$ be the proper ideal class number and unit count of $\mathcal O_\delta$, and put $h_w(\delta)=2h_\delta/w_\delta$. Put $\eta_\delta=(\delta/L)$ and, in the notation of Lemma~\ref{cl:local}, define
\[
 B(\delta)=(1-\eta_\delta/L)K(\delta),\qquad
 c(\delta)=B(\delta)\left(1-\frac1{L-\eta_\delta}\right).
\]
If $4p=t^2+m^2|\delta|$ with $m>0$ and $L\nmid m$, then $K(t^2-4p)=K(\delta)$ because $m$ is an $L$-adic unit and its square preserves every Legendre symbol and valuation in $K$. The order class-number formula now gives the following \emph{exact} finite sum:
\begin{equation}
 S(x)=\sum_{\substack{p\le x\\p\ne L}}\frac{\log p}{\sqrt p}
       \sum_{\substack{t\in\Z,\,m\ge1,\,\delta<0\\
                    4p=t^2+m^2|\delta|,\ L\nmid m}}
                      h_w(\delta)B(\delta).
 \label{op:sum}
\end{equation}
Specifically, $\sum_{p\le x,p\ne L}\lambda_E(p)\log p=-\sum_{p\le x,p\ne L}\log p/\sqrt p-S(x)/2$. To check the normalization, the term in $2V_0/m$ equals $\sqrt{|\delta|/p}$. Its remaining character sum is $L(1,\chi_\delta)(1-\eta_\delta/L)$ and $\sqrt{|\delta|}L(1,\chi_\delta)/\pi=h_w(\delta)$. This is the imprimitive series regrouped at $L$ in Section~\ref{cl:section}. Each term of \eqref{op:sum} is computable by reduced binary quadratic forms.

\subsection{The main density and its cancellation}

Set
\[
 A_\delta(p)=\#\{(t,m)\in\Z\times\Z_{>0}:4p=t^2+m^2|\delta|\},
 \qquad
 \Psi_\delta(x)=h_w(\delta)\sum_{\substack{p\le x\\p\ne L}}
                          A_\delta(p)\frac{\log p}{\sqrt p}.
\]
The pairs correspond to the elements $(t+m\sqrt\delta)/2\in\mathcal O_\delta$ of norm $p$ with positive imaginary part. The parity condition for membership is automatic from the equation modulo $4$. The real elements have square norm, so the sign choice takes exactly half of all generators. Away from primes dividing $2L\delta$, extension and contraction give a norm-preserving bijection of invertible ideals with ideals prime to the conductor in the maximal order. Consequently $A_\delta(p)$ is $w_\delta/2$ times the number of principal prime ideals of norm $p$ in the order. This includes the extra generators at $\delta=-3,-4$. The ring class group is a quotient of the ray class group of modulus $f\mathcal O_{\Q(\sqrt\delta)}$, where $\delta=f^2d$ with $d$ fundamental. Its Hecke characters therefore have finite analytic conductor at most $|d|f^2=|\delta|$. \href{https://www.math.utoronto.ca/~ila/Cox-Primes_of_the_form_x2%2Bny2.pdf}{Cox, \emph{Primes of the Form $x^2+ny^2$}, §§7--9.}

The explicit formula and character orthogonality give, uniformly for $x\ge2$,
\begin{equation}
 \Psi_\delta(x)=2\sqrt x+
   O_\theta\!\left(h_\delta x^{\theta-1/2}\log^4(2x|\delta|)\right).
 \label{op:one}
\end{equation}
Here is a verification including the sharp cutoff. For $P\le y\le2P$ detect norms in $(P,y]$ by moving each endpoint to an adjacent half-integer and smoothing over less than $1/4$ in the norm variable. At integer norms the values are unchanged. On the scale $P$ the resulting cutoff has bounded total variation and its second derivative has $L^1$ norm $O(P)$. Its Mellin transform is $O(\min(1,1/|v|,P/|v|^2))$ on a fixed strip, by two integrations by parts. At archimedean angular index $h$ the complex gamma factor of a primitive character is proportional to $\Gamma(s+|h|/2)$. The unit-interval zero count is $O(\log(2|\delta|(|h|+|v|+2)))$, by the logarithmic-derivative argument of Section~\ref{ld:section} in degree two. Sum the Mellin bound over zeros up to height $P$ and beyond. The pole yields the usual integral for a trivial character; the zeros and left-line integral cost $O(P^\theta\log^2(2P|\delta|(2+|h|)))$. The horizontal segments are taken away from zeros as in \eqref{ld:smooth}. The smoothing changes the principal integral by $O(1)$. Higher prime powers and primes of residue degree two cost $O(\sqrt P\log(2P))$. Deleted prime ideals over $2L\delta$ cost $O(\log(2|\delta|))$ on this dyadic interval. These estimates are uniform in $h$. Ordinary partial summation allows an extra weight $p^{-1/2}$ and then dyadic summation gives the error in \eqref{op:one}, with $\log^4(2x|\delta|(2+|h|))$ when needed. The dyadic powers are summable since $\theta>1/2$.

For $h=0$ the principal class indicator averages the characters with coefficient $1/h_\delta$, and multiplication by $h_w(\delta)w_\delta/2=h_\delta$ gives the displayed main term with coefficient one. The discarded primes in the direct definition of $A_\delta$ contribute $O(h_\delta\log(2|\delta|))$, since $A_\delta(p)\ll\sqrt p$ by counting the possible traces. The integral $\int_2^x dy/\sqrt y=2\sqrt x+O(1)$ completes \eqref{op:one}.

Pairs with $L\mid m$ are exactly the pairs for the suborder of discriminant $L^2\delta$, by dividing $m$ by $L$. The order class-number formula gives
\[
       h_w(L^2\delta)/h_w(\delta)=L-\eta_\delta .
\]
Hence the contribution to \eqref{op:sum} at fixed $\delta$ is
\begin{equation}
 B(\delta)\left(\Psi_\delta(x)-
               \frac{\Psi_{L^2\delta}(x)}{L-\eta_\delta}\right).
 \label{op:scalar}
\end{equation}
The leading term is $2\sqrt x\,c(\delta)$ and the error in \eqref{op:one} retains its size by the same ratio and the boundedness of $B$.

\begin{lemma}\label{op:mean}
The incomplete sums of $c(\delta)$ over $0<|\delta|\le D$, $\delta\equiv0,1\bmod4$, are uniformly $O_L(1)$.
\end{lemma}
\begin{proof}
On the unit residues $\eta=\pm1$, one has $B=\eta$. Their total contribution to the Haar mean on $\Z_L$ is
\[
 \frac{L-1}{2L}
 \left[1-\frac1{L-1}-\left(1-\frac1{L+1}\right)\right]
       =-\frac1{L(L+1)} .
\]
On multiples of $L$ the factor outside $K$ in $c$ is $1-1/L$. For each $e\ge1$ its term in $K$ is $L^{1-e}$ times the indicator of $L^{2e}\mid\delta$ times $(1-(\frac{\delta/L^{2e}}{L})/L)^{-1}$. As computed in Lemma~\ref{cl:local}, the average of this last factor at an unrestricted $L$-adic argument is $A_L=(L^2+L+1)/(L(L+1))$. Thus the contribution is
\[
 (1-1/L) A_L \sum_{e\ge1} L^{1-3e}=1/[L(L+1)].
\]
The two terms cancel. On either of the two progressions modulo four, unit residues incur $O_L(1)$ discrepancy from the mean. In the $e$th term, divide first by $L^{2e}$: the remaining progression is still coprime in step to $L$, and a bounded function of one residue modulo $L$ has $O_L(1)$ discrepancy on any interval. Including the endpoint error and the coefficient gives $O_L(L^{1-e})$ independently of $D$. This argument also covers $L^{2e}>D$, when the actual contribution is zero. Sum the geometric series.
\end{proof}

\begin{theorem}[All character indices at small order discriminants]\label{op:disc}
Let $S_{\le D}(x)$ retain only $|\delta|\le D$ in \eqref{op:sum}. Then
\[
 S_{\le D}(x)\ll_\theta \sqrt x+
             x^{\theta-1/2}D^{3/2}\log^5(2xD).
\]
For $\theta=7/8$ the full portion $|\delta|\le x^{1/12}$ satisfies the square-root estimate up to logarithms, at every sharp prime cutoff.
\end{theorem}
\begin{proof}
Sum \eqref{op:scalar} and apply Lemma~\ref{op:mean}. The nonprincipal errors use $h_\delta\ll\sqrt{|\delta|}\log(2|\delta|)$ from the class number formula and the elementary harmonic bound for an imprimitive nonprincipal character at one. Thus their sum is as stated. All squarefree and square character indices at the indicated discriminants have been retained.
\end{proof}

\subsection{A divisor range with all moving characters}

\begin{theorem}\label{op:taper}
Fix a smooth nondecreasing function $W$ on $[0,\infty)$ that vanishes on $[0,1]$ and equals one on $[2,\infty)$. Insert $W(m/M)$ in \eqref{op:sum}, and write $D=4x/M^2$. If $D\ge1$, the resulting contribution satisfies
\[
 S_W(x,M)\ll_{W,\theta}\sqrt x+
                x^{\theta-1/2}D^{3/2}\log^5(2xD).
\]
If $D<1$ the sum vanishes. In particular take $M=2x^{(1+2\theta)/6}$. The square-root test for $E$ is equivalent to its restriction with factor $1-W(m/M)$. For the input $\theta=7/8$ the residual order divisors are supported on $m<4x^{11/24}$. This cutoff applies to the full class numbers with no restriction on the other character variables.
\end{theorem}
\begin{proof}
Use the angle $v=\arg((t+m\sqrt\delta)/2)$ of a prime generator. The weight is $W(A|\sin v|)$ for $A=2\sqrt p/(M\sqrt{|\delta|})$. It is unchanged at $-\alpha$ and $\bar\alpha$, so taking all generators with the factor $1/2$ is valid. Subtracting the index-$L$ suborder still uses precisely the same angular weight: both descriptions concern the same complex number $\alpha$. In the suborder the old coordinate is $m=Lm'$. The parameter in $W$ is the old $m/M$.

On $P<p\le2P$ write $y=p/P$ and $A_0=2\sqrt P/(M\sqrt{|\delta|})$. Let $c_h(y)$ denote the angular Fourier coefficients of $W(A_0\sqrt y\,|\sin v|)$. For the function and its $y$ derivative the supremum is $O_W(1)$, since $zW'(z)$ is bounded. The first two angular derivatives are $O_W((1+A_0)^2)$, including after one $y$ derivative. The absolute value at the zeros of the sine causes no singularity because $W$ vanishes in a neighborhood of zero. Two integrations by parts give
\[
 |c_h(y)|+|c'_h(y)|\ll_W
        \min(1,(1+A_0)^2/h^2) \qquad(h\ne0).
\]
For $|\delta|\le D$, summing with the conductor weights gives
\[
 \sum_{h\ne0} (|c_h(y)|+|c'_h(y)|)\log^4(2P|\delta|(2+|h|))
          \ll_W (1+A_0)\log^4(2xD).
\]

For a frequency $h$ the sum over the generators of a principal ideal vanishes unless $w_\delta\mid h$. Otherwise the value $(\alpha/|\alpha|)^h$ defines a character on the prime-to-conductor principal ideals of the order, independent of the generating unit. It extends to all invertible ideals prime to the conductor: choose roots for the values on cyclic generators of the finite quotient class group. This defines a Hecke character of angular type $h$ and finite conductor dividing $f\mathcal O_{\Q(\sqrt\delta)}$, since on maximal-order principal ideals with generator congruent to $1\bmod f$ it already has the prescribed archimedean value. Multiply its extension by the class characters to impose principality. For $h\ne0$ none has a pole, since the angular type is nonzero. Apply the sharp partial-sum estimate proved above with the coefficient $c_h(p/P)p^{-1/2}$ by partial summation. Summing the nonzero frequencies gives an error $O_{W,\theta}(h_\delta P^{\theta-1/2}(1+A_0)\log^4(2xD))$. The same calculation for the suborder includes only its own unit-compatible frequencies and the same average $c_0$. The zero frequency has the corresponding error and contributes the main density
\[
 c(\delta)\int_P^{\min(2P,x)}
    \left(\frac1{2\pi}\int_0^{2\pi}
          W\left(\frac{2\sqrt y\,|\sin v|}{M\sqrt{|\delta|}}\right),dv\right)
                       \frac{dy}{\sqrt y}.
\]
For each fixed $y$, the angular average is between zero and one and is nonincreasing in $|\delta|$. Abel summation and Lemma~\ref{op:mean} give $O(1)$ for its signed sum over the discriminants. The integrated main term is $O(\sqrt x)$. For the errors use
\[
 \sum_{|\delta|\le D} h_\delta(1+\sqrt{D/|\delta|})
                  \ll D^{3/2}\log(2D).
\]
Dyadic summation in $P$ and restoration of the finitely excluded primes at each order give the theorem. At the stated $M$, $D=x^{2(1-\theta)/3}$, making the second term $O_{W,\theta}(\sqrt x\log^5(2x))$. The square-root equivalence is preserved when this term is removed and restored at every sharp cutoff, by the same one-sided argument as in Theorem~\ref{cl:triple}.
\end{proof}

\subsection{A sharp order-divisor cutoff}

A moving sector can also be treated without a taper. The sharp sum here refers to the exact inequality $m>M$ at every prime and order. Its estimate is
\begin{theorem}\label{op:sharp}
Under the same Hecke hypothesis, for $M\ge1$, $D=4x/M^2\ge1$ and every integer $H\ge1$,
\[
 S_{\mathbf1_{m>M}}(x,M)
 \ll_\theta \sqrt x(1+D/H)
        +x^{\theta-1/2}D^{3/2}\sqrt H\log^6(2xDH).
\]
Taking $H=\lceil D\rceil$ in particular permits a sharp truncation to $m\le2 x^{(1+\theta)/4}$ in the square-root test. At $\theta=7/8$ the exponent is $15/32$.
\end{theorem}

\begin{proof}
We give the oscillatory estimate used for the moving endpoints. On $P<p\le2P$ put $z=\log(p/P)$ and
\[
 \phi(z)=\arcsin\bigl(\min(1,b e^{-z/2})\bigr),\qquad b\ge0.
\]
On a fixed enlargement of $[0,\log 2]$, its total variation is at most $\pi/2$ and its $1/2$-Hölder constant is bounded independently of $b$. For the second assertion the argument of the arcsine is uniformly Lipschitz, since its derivative where it is nonconstant has modulus at most $1/2$; the arcsine is $1/2$-Hölder. Consider a nonzero angular Hecke frequency $h$ in a ring class expansion and insert the phase $e^{\pm ih\phi(z)}$ in a prime sum over any subinterval of $(P,2P]$. Uniformly for $1\le |h|\le H$, its error before the order-class factor and before $p^{-1/2}$ is
\begin{equation}
  O\bigl(P^\theta \sqrt{1+|h|}\log^3(2P|\delta|(H+2))\bigr).
 \label{op:phase}
\end{equation}
The same holds with an extra fixed smooth bounded weight of $p/P$ on this interval, including $(p/P)^{-1/2}$.

Here are the details for the sharp endpoints. Replace the endpoints of the interval of norms by half-integers just outside its first and last integers. Let $F(z)$ be its resulting log-scale indicator times the phase and any stated smooth factor. It has bounded supremum, compact support and total variation $O(1+|h|)$. Smooth by convolution in $z$ with a nonnegative unit-mass bump of radius $\xi=(10P(H+2))^{-4}$. The indicator is constant on a $\xi$-neighborhood of each relevant integer norm point. The phase changes there by $O((1+|h|)\sqrt\xi)$, contributing $O(1)$ to the prime sum by Chebyshev's estimate in degree two. The Mellin transform of this convolution is the transform of $F$ times that of the bump; the latter is $O((1+\xi|\gamma|)^{-2})$ on a bounded strip.

For $G(z)=z^j e^{\beta z}F(z)$, $j=0,1,2$, with $\beta$ in a fixed bounded interval, the two variation bounds imply
\[
 \int_{\mathbb R}|G(z+t)-G(z)|^2 dz\ll (1+|h|)|t|
\]
by the $L^1$ translation bound for a function of bounded variation and its supremum norm. Plancherel with $t=1/V$ gives
\[
 \int_{V\le|\gamma|\le2V}|\widehat G(\gamma)|^2d\gamma
                   \ll (1+|h|)/V\qquad(V\ge1)
\]
in the angular Fourier convention $e^{i\gamma z}$, since $|e^{i\gamma/V}-1|$ stays bounded below on this band. The same estimate on bands enlarged by one unit follows by changing the endpoints by fixed factors or treating bounded $V$ directly. Write $\mathcal F(\beta,\gamma)=\int e^{(\beta+i\gamma)z}F(z)\,dz$. The above estimate applies also with either or both of one $\beta$ derivative and one $\gamma$ derivative. On each unit $\gamma$ interval and a fixed $\beta$ interval, the squared supremum of $\mathcal F$ is bounded by a constant times the sum of the squared $L^2$ norms of these four functions. This is the one-variable fundamental theorem of calculus and Cauchy--Schwarz applied in both coordinates. Summing the squared suprema on the band gives $O((1+|h|)/V)$. Cauchy--Schwarz over its $O(V)$ intervals now bounds the sum of the suprema by $O(\sqrt{1+|h|})$.

Apply the explicit formula in the proof of \eqref{op:one}. The zero multiplicity per unit interval costs only $O(\log(2|\delta|(H+|\gamma|+2)))$. Summing dyadic bands with the bump factor costs at most $O(\log^2(2P|\delta|(H+2)))$. The same reasoning, or its integral version on a single vertical line, bounds the left contour. There is no principal pole at this nonzero angular frequency. Higher powers, omitted primes and the endpoint replacement are within \eqref{op:phase}. This proves the phase estimate.

At fixed $\delta,p$, the allowed angles consist, for $p>M^2|\delta|/4$, of the two open intervals
\[
 (\phi,\pi-\phi),\qquad(\pi+\phi,2\pi-\phi),
 \qquad b=M\sqrt{|\delta|}/(2\sqrt P).
\]
For smaller $p$ the count is zero. Majorize and minorize each interval by the degree-$H$ Selberg trigonometric polynomials. Their constant coefficients equal its normalized length up to $O(1/H)$. At index $h\ne0$ their coefficients are linear combinations of exponentials of $ih$ times the two endpoints, with coefficients bounded by $O(1/|h|)$ independently of the interval. These follow from the translated one-sided approximations to the periodic sawtooth function; the majorant covers endpoints and the minorant applies to an open interval. \href{https://math.tsukuba.ac.jp/~akiyama/papers/Salem8.pdf}{Akiyama and Tanigawa, Selberg interval polynomials.}
Thus the extra phase in each nonzero angular term is of precisely the form in \eqref{op:phase}, up to a constant of modulus one. Summing $\sqrt{1+|h|}/|h|$ through $H$ costs $O(\sqrt H)$. The zero-frequency main density is the angular length as a function of $|\delta|$, a bounded nonincreasing function at fixed $p$. Its integral has bounded signed sum after Lemma~\ref{op:mean}. The constant-coefficient errors of the majorant and minorant cost $O(\sqrt P/H)$ per discriminant in the principal integral; their nonprincipal errors are allowed by \eqref{op:phase}. The exact zero-frequency length has bounded variation in $p$ as well, so its nonprincipal error follows by partial summation in \eqref{op:one}.

First apply these two inequalities to each of the nonnegative generator counts for $\mathcal O_\delta$ and $\mathcal O_{L^2\delta}$. Their angular regions concern the same elements, with the same angular endpoints, just as in Theorem~\ref{op:taper}. Then subtract the two main terms with their exact weights and add the error bounds in absolute value. Sum over the order discriminants and dyadic $P$ using $\sum h_\delta\ll D^{3/2}\log(2D)$. This proves the estimate, including both sharp cutoffs.
\end{proof}

\subsection{The complementary low squarefree modes}

The short-modulus projection below improves Theorem~\ref{cl:alllevel}. It uses the special dependence on $t^2-4p$ of the class-number coefficients and keeps every level digit and every square part.

\begin{theorem}\label{op:low}
Under Dirichlet GRH the full contribution with $m\le M$ and squarefree character kernel $r\le R$ satisfies
\[
          O_g\bigl(\sqrt x\, M R\log^{10}(2xMR)\bigr)
\]
for $g=X_0(11)$ or $\Delta$, at all sharp prime cutoffs. For $X_0(11)$ use the complete $K$ in the imprimitive order expansion away from eleven. In particular the square-root target holds in every rectangle $MR\le x^{\varepsilon/2}$ for a tested exponent $\varepsilon>0$.
\end{theorem}
\begin{proof}
Recall the divisor $h$ of \eqref{cl:squares}, the condition $m^2h\mid t^2-4p$, and the actual weight $1/(mrh^2)$. At level eleven all of $m,r,h$ are prime to $11$. Treat the projection $K_1$ first; at level one omit this factor. Put $D_0=m^2h$, $r_0=r/(r,22mh)$, and use period $Q=8D_0rL$ with $L=11$ (omit $L$ for $\Delta$). On every unit prime residue $a$, its coefficient $b(t,a)$ depends only on $t^2-4a$ modulo $Q$. In particular $K_1(t,a)$ has this property: the value on a zero residue of $t^2-4a\bmod L$ is the mean of $K$ on $L\Z_L$, by the unit-derivative substitution of Lemma~\ref{cl:local}; on the other residues $K$ is already determined.

Let $\bar b(t)$ be its average over the unit residues $a\bmod Q$. For $t'\equiv t\bmod q_0$, $q_0=4\operatorname{rad}(Q)$, the integer $(t'^2-t^2)/4$ is divisible by each prime of $Q$. Translation by this integer permutes the unit residues and leaves $t^2-4a$ unchanged when $t$ is replaced by $t'$. Thus $\bar b$ has period $q_0\ll mhr$. For a fixed $t$ there are at most four possible $a\bmod D_0$ making $4a\equiv t^2\bmod D_0$. The fraction among the units is therefore $O(1/\varphi(D_0))$. The Chinese remainder theorem separates each $\ell\mid r_0$ from these conditions and the other factors; its average over $a\bmod\ell$ coprime to $\ell$ is
\[
 \frac1{\ell-1}\sum_{a\ne0}
        \left(\frac{t^2-4a}{\ell}\right)
          =-\frac1{\ell-1}\left(\frac{t^2}{\ell}\right).
\]
Division by $m^2$ does not change this symbol. The remaining factors are bounded. Hence
\[
          |\bar b(t)|\ll \frac{\log^2(2Q)}{D_0 r_0},
\]
using $\prod_{\ell\mid v}(1-1/\ell)^{-1}\ll\log(2v)$. The mean of $\bar b$ in $t$ still vanishes when required; this follows by interchanging the two finite averages. Apply Corollary~\ref{cl:sharp} to $\bar b(t)$ with modulus $q_0$. Its contribution before the weights costs $O(\sqrt x\,h r(r,22mh)\log^5(2xQ))$. After the weights the sum left is bounded by
\[
 \sum_{m\le M,h\le4x}\frac1{mh}\sum_{r\le R}(r,22mh)
       \ll R\log^2(2M)\log^2(2x),
\]
since $\sum_{r\le R}(r,v)\le R\tau(v)$ and $\sum_{j\le y}\tau(j)/j\ll\log^2(2y)$. Distinct prime divisors of $q_0$ and $Q$ are the same.

The remainder $b-\bar b$ has zero average in $a$ for each $t$. All its normalized additive Fourier coefficients in $t$ are $O(A_{m,r,h})$ with $A_{m,r,h}=\tau(mh)\tau(r)/(D_0\sqrt{r_0})$ from Corollary~\ref{cl:squarefree}: averaging a Fourier coefficient over $a$ preserves this bound. Thus \eqref{cl:noprincipal} applies with period power $Q^{3/2}$. For each subsequent level digit the same bound applies with the factor $L^{-3j/2}$ against the period gain $L^{3j/2}$, as in Theorem~\ref{cl:alllevel}. There are $O(\log(2x))$ such digits; the usual exceptional-prime bound from that theorem is also admissible. After taking out the prime-cutoff factor and its logarithms, the remaining sum is
\[
 \sum_{m\le M,h\le4x}\frac{\tau(mh)}{h^{3/2}}
                 \sum_{r\le R}\tau(r)(r,22mh)^{1/2}.
\]
With $B(v)=\sum_{d\mid v}\tau(d)/\sqrt d$ as in Corollary~\ref{cl:squarefree}, the inner sum is $O(R\log(2R)B(22mh))$. The previously proved bound $\sum_{j\le z}\tau(j)B(j)\ll z\log(2z)$ handles the $m$-sum and gives absolute convergence for the $h$-sum with weight $h^{-3/2}$. Combining the projections gives the assertion, with the stated logarithmic power allowing all the endpoint and divisor factors.
\end{proof}

For a particular conductor-sized Gaussian range the interest of this estimate is its dependence on $m$ and $r$. The formal scale $mr\asymp\sqrt x$ lies at the $x$ exponent of this estimate. The Gaussian weights in Theorem~\ref{cl:triple} vary with the conductor; that expansion does not follow by substituting this scale into Theorem~\ref{op:low}. In particular the estimate does not give a square-root bound throughout the full transition between the small-character and small-order-discriminant ranges.

\subsection{Two failed absolute-value bounds}

The order of local summation matters. Before the powers of eleven are grouped as in \eqref{op:sum}, keep all primitive-order terms of Theorem~\ref{cl:level} with $|\delta|\le D$. Its multiplier at the order divisor $m$ is $\eta_\delta$ if $L\nmid m$, and $L$ if $L\mid m$. Consequently the contribution at $\delta$ is exactly
\[
   \eta_\delta\Psi_\delta(x)+\Psi_{L^2\delta}(x).
\]
Its main coefficient is $1+\eta_\delta$. Over the negative discriminants through $D$, the constant has sum $D/2+O(1)$ and the symbol has bounded sum by periodicity and the Chinese remainder theorem. Thus the \emph{uncollected} cutoff is
\[
 D\sqrt x+
 O_\theta\bigl(\sqrt x+x^{\theta-1/2}D^{3/2}\log^5(2xD)\bigr).
\]
For $D=x^b$ with $0<b<2(1-\theta)$ it is asymptotic to $D\sqrt x$. Both order expansions give the same sum when complete. A bound at the square-root exponent for this particular truncated version would be false.

There is a different obstruction directly at the moving squarefree indices. Expand all of \eqref{op:sum} back in the imprimitive Dirichlet series away from eleven. Write the character index as $rb^2$, $r$ squarefree, and sum over all square parts $b$ by \eqref{cl:squares}. Let $U_r(x)$ be the resulting real coefficient, \emph{including} its factor $1/r$ and every prime $p$, trace $t$, order divisor $m$ and level contribution. The conditional sum over $r$ equals $S(x)$; convergence follows for each fixed $x$ from the finite original set of orders and the squarefree boundary argument after \eqref{cl:squares}.

\begin{theorem}[An absolute obstruction for the actual squarefree coefficients]\label{op:barrier}
Assume Dirichlet GRH. There are absolute positive constants $C,c$ such that, for all sufficiently large $x$ and $R\ge Cx\log^8(2x)$,
\[
   \sum_{\substack{R<r\le2R\\r\ {\rm prime}}}|U_r(x)|
                   \ge c\,x/\log(2R).
\]
In particular absolute summation of these coefficients diverges for every sufficiently large fixed $x$. At $R=Cx\log^8(2x)$ an absolute dyadic bound of square-root size in $x$ fails even after all the other variables have been summed.
\end{theorem}
\begin{proof}
Write $\delta=f^2 d$ with $d$ fundamental. For all primes $r>\max(4x,L)$ there is an exact finite decomposition
\[
   rU_r(x)=\sum_d A_d(x)\chi_d(r).
\]
Indeed the fixed $(p,t,m)$ term in $A_d(x)$ is
\[
   \frac{\log p}{\pi\sqrt p}\sqrt{|\delta|}\,K(\delta)\,
       \zeta(2)(1-L^{-2})
        \prod_{\substack{\ell\mid\delta\\\ell\ne L}}(1-\ell^{-2}).
\]
The primitive value at $r$ agrees with the imprimitive one because $r$ is larger than every discriminant here. The product over $\ell$ is bounded both above and below by positive constants. Since $\sqrt{|\delta|/p}\le2/m$, $|K|\ll1$, there are $O(\sqrt p)$ possible traces and $\sum_{m\le 2\sqrt p}1/m\ll\log(2p)$, Chebyshev's bound gives $\sum_d |A_d(x)|\ll x^{3/2}\log(2x)$.

In the particular square class $d=-7$, the value $(-7/L)=1$ implies that $K(-7f^2)$ is positive and bounded below uniformly in $f$. If $L\nmid f$, only the $e=0$ term in $K$ contributes and equals $1/(1-1/L)$. If $L\mid f$, its $e=1$ term alone is at least $(1+1/L)^{-1}$ and all other terms are nonnegative. In fact
\begin{equation}
                   A_{-7}(x)\gg x .
 \label{op:largecomponent}
\end{equation}
It suffices to use only the terms $m=1$. The maximal order of $\Q(\sqrt{-7})$ has class number one and two units, as follows also by enumerating reduced forms of discriminant $-7$. For every nonexceptional split prime, take its generators $\alpha=(t+j\sqrt{-7})/2$ with $j>0$. Then $t^2-4p=-7j^2$ is exactly the order discriminant in this term, and $\sqrt7j/\sqrt p=2|\sin(\arg\alpha)|$. The sum of $\sin^2(\arg\alpha)\log p$ over these generators and $p\le x$ is $x/2+o(x)$: average $1/2-(1/2)\cos(2\arg\alpha)$ over degree-one prime ideals. The ordinary prime ideal theorem gives $x$ for the constant, and the prime ideal theorem for the nontrivial fixed angular Hecke character gives $o(x)$ for the cosine. Both are classical qualitative prime ideal theorems; the stated Hecke strip also gives the result. The choice $j>0$ takes one of the two associates per ideal and does not affect the cosine. Since $|\sin v|\ge\sin^2 v$, all these terms have the same positive sign and prove \eqref{op:largecomponent}.

Project $rU_r(x)$ onto $\chi_{-7}(r)$ along the primes in $(R,2R]$ using weight $\log r$. At every $d\ne-7$ the product character is nonprincipal with primitive conductor at most $7|d|\le28x$; distinct fundamental discriminants give distinct primitive characters. By the standard sharp explicit formula under Dirichlet GRH its sum is $O(\sqrt R\log^2(2Rx))$. The finitely many omitted bad Euler factors do not occur at these primes. At $d=-7$ the coefficient is $R+O(\sqrt R\log^2(2R))$. Thus
\[
 \sum_{R<r\le2R\atop r\ {\rm prime}}
      (\log r)\chi_{-7}(r)\,r U_r(x)
   =R A_{-7}(x)+O(x^{3/2}\log(2x)\sqrt R\log^2(2Rx)).
\]
The error is less than half the main term for the stated range, increasing $C$ if necessary. This remains valid for all larger $R$ by eventual monotonicity of $\sqrt R/\log^2(2Rx)$. Apply the triangle inequality and $r\log r\le2R\log(2R)$ to obtain the bound. The harmonic series over dyadic exponents proves the divergence assertion.
\end{proof}

This obstruction concerns the true conditionally convergent series on squarefree kernels. The finite Gaussian series of Theorem~\ref{cl:triple} contains no indices as large as $Cx\log^8(2x)$ at a prime cutoff $x$. The theorem gives no lower bound in that shorter interval. A square-root proof may exploit cancellation between the prime index $r$ and the varying primitive characters or use the two components of the finite Gaussian kernel together.

\subsection{The main term from one primitive quadratic field}

The short Gaussian expression has a different obstruction of its own. For each fundamental discriminant $d<0$ let $\Phi_d(x)$ be the contribution to $S(x)$ of pairs $(p,t)$ for which $t^2-4p=f^2d$, in the primitive expansion of \eqref{cl:lvalues}. Thus its weight is
\[
       \frac{\log p}{\pi}\sqrt{(4p-t^2)/p}\,
              J_E(p,t)L(1,\chi_d).
\]
This grouping is by the field, so it includes all conductors compatible with its pairs. The principal term for the field survives averaging over primes.

\begin{theorem}\label{op:primitive}
Uniformly for $d$ fundamental, $0<|d|\le4x$, one has under the stated Hecke hypothesis
\[
 \Phi_d(x)=C(d)x+O_\theta\!\left(x^\theta\log^9(2x|d|)\right).
\]
Writing $\eta=(d/L)$, the coefficient is
\begin{align*}
 C(d)&=\frac{4}{\pi\sqrt{|d|}}\,F_L(\eta)
  \prod_{\ell\ne L}\left(1+
         \frac{1-\chi_d(\ell)}
                 {(1-\chi_d(\ell)/\ell)(\ell^2-1)}\right),\\
 F_L(\eta)&=\eta+\frac{(1-\eta)L}{(1-\eta/L)(L^2-1)}.
\end{align*}
In particular $C(-7)>0$. There is an absolute constant $c>0$ such that for every fixed $0<b<2(1-\theta)$,
\[
  \sum_{\substack{|d|\le x^b\\d<0\ {\rm fundamental}}}\Phi_d(x)
                  \sim c\,x^{1+b/2}.
\]
At $\theta=7/8$ this holds for all $0<b<1/4$. The same asymptotic holds for the corresponding portion of the finite Gaussian expansion, in the normalization of $S(x)$.
\end{theorem}
\begin{proof}
At $p,t$ as above, $J_E(p,t)$ depends only on $d,f$. Write its divisor expansion as $J_E=\sum_{k\mid f}a_d(k)$. At primes $\ell\ne L$ the factors for $\ell^v\parallel k$ are $(1-\chi_d(\ell))/\ell^v$ for $v\ge1$, obtained by differencing $G(\ell^v,d)$ from Section~\ref{cl:section}; the factor is one for $v=0$. At $L$ the corresponding coefficients are
\[
              b_0=\eta,\qquad b_v=(1-\eta)/L^{v-1}\quad(v\ge1).
\]
In fact when $L\nmid f$, the local factor of $J_E$ is $\eta$. At $L^e\parallel f$, $e\ge1$, it is $G(L^{e-1},d)$, because the discriminant symbol in the first term of $J_E$ vanishes, leaving the scalar-lift term. First differences give the asserted $b_v$. The product of these rules, including the factor $b_0$ also when $L\nmid k$, defines $a_d(k)$. Uniformly $|a_d(k)|\ll_L \tau(k)/k$.

A pair with $k\mid f$ is exactly a prime generator in the order $\mathcal O_{k^2d}$, with coordinate $f/k$. Also $\sqrt{(4p-t^2)/p}=2|\sin(\arg\alpha)|$. The angular weight $|\sin v|$ has mean $2/\pi$ and Fourier coefficients $O((1+|h|)^{-2})$, by integration by parts on the two smooth half circles. Its logarithmically weighted prime sum over those generators with positive imaginary part, excluding $p=L$, is therefore
\[
 \frac{2}{\pi}\frac{w_{k^2d}}{2h_{k^2d}}\,x
                  +O_\theta(x^\theta\log^4(2x k|d|)).
\]
This is the angular Hecke argument in Theorem~\ref{op:taper} before division by $\sqrt p$ or multiplication by $h_w$. The class indicator has total absolute character coefficient one, and $w_{k^2d}\le6$. At primes dividing $2Lkd$, the order's generators are a subset of maximal-order elements of norm $p$, at most twelve by the ideal factorization in degree two. Restoring them is covered by the error.

Sum this last formula with weights $2L(1,\chi_d)a_d(k)/\pi$ through $k\le K=2\sqrt{x/|d|}$. The accumulated error is $O_\theta(x^\theta\log^8(2x|d|))$, since $L(1,\chi_d)\ll\log(2|d|)$ and $\sum_{k\le K}\tau(k)/k\ll\log^2(2K)$. In the main coefficient use
\[
 \frac{2h_{k^2d}}{w_{k^2d}}
     =\frac{k\sqrt{|d|}}{\pi}L(1,\chi_d)
                         \prod_{\ell\mid k}(1-\chi_d(\ell)/\ell).
\]
The primitive $L$-value cancels. Extending the main coefficient to all $k$ costs
\[
 \ll \frac{x}{\sqrt{|d|}}\sum_{k>K}\frac{\tau(k)\log(2k)}{k^2}
      \ll \sqrt x\,\log^2(2x),
\]
also when $K$ is bounded, by elementary divisor summation. Now multiply the absolutely convergent local sums over $k$. Away from $L$ the geometric series over $v\ge1$ gives the indicated Euler factor. At $L$ the sum of the $b_v$ divided by $L^v(1-\eta/L)$ for $v\ge1$, with $b_0$ added, gives $F_L$. This proves the uniform first assertion.

For its mean, the Euler product away from $L$ is between $1$ and an absolute constant and has a uniform multiplicative tail $1+O(\sum_{\ell>Y}\ell^{-2})$ on truncation to $\ell\le Y$. Among negative fundamental discriminants, the values $\eta=\pm1$ each have limiting probability $L/(2(L+1))$, and $\eta=0$ has probability $1/(L+1)$. They are independent of the symbols at any finite set of other primes. One elementary verification writes a fundamental discriminant as a squarefree integer in the prescribed classes modulo four, or four times a squarefree integer in the complementary classes. At each odd prime $\ell$, among the squarefree admissible classes modulo $\ell^2$ there are $\ell-1$ classes divisible by $\ell$ and $\ell(\ell-1)$ unit classes, the latter evenly split between the two signs. Multiplication by four preserves the odd symbols. The Chinese remainder theorem gives independence at the finite set, including any fixed $2$-adic conditions. Excluding squared prime factors at primes outside this set has a limiting independent density: truncate the squarefree sieve, then bound the number of integers divisible by some $q^2$ with $q>Y$ by $O(D/Y)$ through $D$. In particular the total density of negative fundamental discriminants is positive.

Under this distribution the mean of $\eta$ is zero; the second term of $F_L$ is nonnegative and strictly positive for $\eta=0,-1$. At two the three symbol values have probabilities $1/3$ each. Indeed the odd negative fundamental discriminants contribute density $2/\pi^2$, split evenly between their two classes modulo eight, and the even ones contribute $1/\pi^2$ by the squarefree characterization. Thus the total density is $3/\pi^2$ and the same probabilities $\ell/(2(\ell+1)),1/(\ell+1),\ell/(2(\ell+1))$ apply at every prime. Put
\[
 u_\ell=\frac{\ell^2+\ell+1}{(\ell^2-1)(\ell+1)^2}.
\]
The mean of the Euler factor away from $L$ is $1+u_\ell$, and that of $F_L$ is $L u_L$. Independence and uniform convergence therefore give the explicit constant
\[
       c=\frac{24}{\pi^3}L u_L\prod_{\ell\ne L}(1+u_\ell)>0 .
\]

One obtains the quantitative estimate $\sum_{|d|\le D}C(d)=c\sqrt D+O(\log(2D))$. To see this, expand the away-from-$L$ product over squarefree $v$. Its $v$th summand, as a function of the squarefree integer underlying the field discriminant, has period dividing $8Lv$ and supremum $\ll_L C_0^{\omega(v)}/v^2$ for an absolute $C_0$. Here $\omega(v)$ counts the distinct prime factors. Include the bounded $L$ factor and the fixed congruences modulo eight in these bounds. For any bounded periodic $V$ of period $q$, squaring out the squarefree indicator gives
\[
 \sum_{u\le X}\mu^2(u)V(u)
       =X c_V+O(\sqrt{Xq}\,\|V\|_\infty).
\]
Indeed for the terms $t^2\mid u$ with $t\le Z$ the inner progression sum has error $O(q\|V\|_\infty)$ from its periodic mean, and the discarded actual and mean tails have total $O(X\|V\|_\infty/Z)$. Take $Z$ of order $\sqrt{X/q}$ when $q\le X$; otherwise the trivial bound suffices. Summing the errors in $v$ costs $O_L(\sqrt D)$ because $\sum_v C_0^{\omega(v)}/v^{3/2}<\infty$. Apply this separately to the two fundamental-discriminant parametrizations. Partial summation with $|d|^{-1/2}$ proves the estimate. In particular uniformly for $1\le D\le4x$,
\[
 \sum_{\substack{|d|\le D\\d<0\ {\rm fundamental}}}\Phi_d(x)
   =c x\sqrt D+
       O_\theta\left(x\log(2D)+x^\theta D\log^9(2xD)\right).
\]
At $D=x^b$ both errors are $o(x\sqrt D)$ in the stated open range. Finally the Gaussian tail in Theorem~\ref{cl:triple} was bounded absolutely before summation, so it also bounds the error on this subset of primitive discriminants.
\end{proof}

Even the single field $\mathbb Q(\sqrt{-7})$ therefore contributes a positive term of order $x$ in the actual finite formula. The sum over fields $|d|\le x^b$ contributes the positive superlinear term above. Since Hasse's bound gives $S(x)=O(x)$, the complementary fields must supply its opposite to leading order. A proof using this particular decomposition has to account for that signed cancellation. Cutting off the order discriminant $\delta$ in Theorem~\ref{op:disc} and cutting off the primitive field discriminant $d$ select different terms.

For comparison, Ditchen's average theorem over fundamental discriminants $|d|\le D$, $d\not\equiv0\bmod8$, bounds the sum of the maximal unweighted class-prime counting errors by $O_{A,\epsilon}(D^{1/2}x/\log^A x)$ when $D^{20/3+\epsilon}\le x/\log^B x$. On a dyadic prime interval the elementary insertion of $h_d\ll\sqrt D\log(2D)$ and $(\log p)/\sqrt p$ into this absolute bound costs $O_{A,\epsilon}(\sqrt x\,D/\log^{A-2}x)$, which does not reach the square-root exponent for $D$ a fixed power of $x$. Its fundamental-discriminant range also omits the suborders in \eqref{op:scalar}. Bilinear Kloosterman estimates at a fixed prime modulus treat a different averaging variable from the changing squarefree modulus $r$ here. The projection in Theorem~\ref{op:barrier} and the main term in Theorem~\ref{op:primitive} give concrete obstructions to two absolute summations; neither rules out estimates that use cancellation between those moving variables. \href{https://arxiv.org/abs/1312.1502}{Ditchen, Theorems 1.1--1.2.} \href{https://doi.org/10.4007/annals.2017.186.2.2}{Kowalski, Michel and Sawin, bilinear forms at a prime modulus.}

\subsection{The exact abscissa for absolute summation over fields}

There is an elementary lower bound over the full range of primitive discriminants. It uses only the level-eleven formula, Hasse and primes in fixed progressions, and has no zero-strip hypothesis. Write $\Phi_{d,p}$ for the contribution at $p$ to $\Phi_d$, with the prime logarithm removed. If $(d/L)=1$, then $J_E(p,t)\ge0$ at every eligible pair. Indeed it equals $G(f,d)$ if $L\nmid f$, and $G(f/L,d)$ if $L\mid f$. Thus $\Phi_{d,p}\ge0$ even after regrouping within each such field.

For a prime $p$ that is a quadratic nonresidue modulo $L$, the symbol $(\frac{t^2-4p}{L})$ never vanishes. There are no terms with $L\mid f$ or any scalar lift. Let $A$ and $B$ be the sums of $H(4p-t^2)$ over the two symbol signs $+1,-1$. Kronecker--Hurwitz and Theorem~\ref{cl:level} give
\[
  A+B=2p,\qquad A-B=-2(a_p+1),\qquad A=p-a_p-1 .
\]
For these primes the field signs $(d/L)$ equal the discriminant symbols. Consequently the total from the positive family is exactly $(p-a_p-1)/\sqrt p$ at this prime. Summing over the nonresidue primes through $x$, the prime number theorem in fixed progressions and Hasse give $(1/3+o(1))x^{3/2}$ after logarithmic weighting. The other primes contribute nonnegative amounts to that family. Since $S(x)=O(x)$, one obtains
\[
 (2/3-o(1))x^{3/2}
       \ \le\sum_{d<0\ {\rm fundamental}}|\Phi_d(x)|\ll x^{3/2}.
\]
For the upper bound, the absolute mass at $p$ is at most $(1+L)\sum_t H(4p-t^2)/\sqrt p=2(1+L)\sqrt p$. Indeed $H(D/L^2)$, when present, sums a subset of the primitive-order terms of $H(D)$. This proves the sharp power of $x$ for the loss after summation over all primes and all conductors within each fixed field.

In particular, put $P_d(s)=\sum_{p\ne L}\Phi_{d,p}(\log p)p^{-s}$. For fixed $d$ this series converges absolutely in $\operatorname{Re}s>1$: the norm-$p$ generators in the maximal order are bounded in number and $G(f,d)\ll_\epsilon f^\epsilon$. The sum over all fields and primes converges absolutely in $\operatorname{Re}s>3/2$ by the preceding upper bound at each prime. At every real $\sigma$ with $1<\sigma\le3/2$ even $\sum_d|P_d(\sigma)|$ diverges. In fact the nonnegative fields with $(d/L)=1$ alone contain the subseries
\[
 \sum_{\substack{p\ \mathrm{prime}\\(p/L)=-1}}
          \frac{p-a_p-1}{\sqrt p}\frac{\log p}{p^\sigma},
\]
which diverges throughout this range. Thus $3/2$ is the exact absolute abscissa of this field regrouping, although the original elliptic prime series is absolutely convergent for $\operatorname{Re}s>1$.

There is a matching normal-convergence obstruction for the continuations. Under the Hecke hypothesis, partial summation of Theorem~\ref{op:primitive} extends each $P_d(s)$ meromorphically to $\operatorname{Re}s>\theta$, with only a possible pole at one and residue $C(d)$. For each $0<r<1-\theta$, the residue theorem gives
\[
 \max_{|s-1|=r}\left|\sum_{|d|\le D}P_d(s)\right|
       \ \ge\ \frac1r\left|\sum_{|d|\le D} C(d)\right|
       =\frac{c\sqrt D+O(\log(2D))}{r}.
\]
The finite field truncations therefore cannot converge uniformly on such a circle. These conclusions concern this specific order of summation and these continued functions.

\subsection{The six one-dimensional cusp spaces at level one}

The field calculation has a uniform extension in weight. Let $k$ be one of $12,16,18,20,22,26$ and $g$ the normalized eigenform spanning $S_k(\mathrm{SL}_2(\Z))$. These are the one-dimensional cusp spaces by the dimension formula, equivalently by the graded ring $\C[E_4,E_6]$ and the cusp ideal generated by $\Delta$. The level-one Eichler--Selberg formula at a prime gives
\[
 \frac{a_p(g)}{p^{(k-1)/2}}
  =-p^{-(k-1)/2}
       -\frac1{2\sqrt p}\sum_{|t|<2\sqrt p}
             U_{k-2}(t/(2\sqrt p))H(4p-t^2).
\]
The weight-twelve instance was given in Theorem~\ref{cl:level}; the same trace formula supplies the displayed elliptic polynomial in weight $k$. For $d$ fundamental, let $\Phi_{g,d}$ denote its primitive-field component in the logarithmically weighted class-number sum before the factor $-1/2$.

\begin{corollary}\label{op:levelone}
For every such $g$ and every fixed imaginary quadratic field discriminant $d$ the coefficient of $x$ in $\Phi_{g,d}(x)$ is strictly positive. More precisely, uniformly for $|d|\le4x$,
\[
 \Phi_{g,d}(x)=
 \frac{4x}{\pi(k-1)\sqrt{|d|}}
 \prod_\ell\left(1+
      \frac{1-\chi_d(\ell)}{(1-\chi_d(\ell)/\ell)(\ell^2-1)}\right)
   +O_{k,\theta}\!\left(x^\theta\log^9(2x|d|)\right).
\]
The portion $|d|\le x^b$ has asymptotic $c_k x^{1+b/2}$ for $0<b<2(1-\theta)$, with $c_k=24\prod_\ell(1+u_\ell)/(\pi^3(k-1))>0$.
\end{corollary}
\begin{proof}
Here the order divisor expansion is that of $G(f,d)$ at every prime; there is no distinguished local factor. The generator angle weight becomes $|\sin v|U_{k-2}(\cos v)$. Its Fourier coefficients are $O_k((1+|h|)^{-2})$, by piecewise integration by parts and continuity. As $k$ is even its mean equals
\[
 \frac1\pi\int_0^\pi \sin((k-1)v)\,dv=\frac{2}{\pi(k-1)}.
\]
Use this mean and the ordinary local divisor coefficients in the proof of Theorem~\ref{op:primitive}. The same unit, conductor and tail estimates apply. The elementary average of the positive Euler product now includes all primes and gives the asserted constant.
\end{proof}

In each of these six cases the signed sum over all quadratic fields has size $O_k(x)$ by the Ramanujan bound for the fixed holomorphic form. The complement of the fields $|d|\le x^b$ therefore has a negative aggregate of order $x^{1+b/2}$, although the leading term at every fixed individual field is positive. This identifies a failure of interchange of the field and prime limits within the actual fixed-form trace.

The cancellation cannot be confined to a sparse exceptional set of fields. For both the conductor-$11$ and the level-one expressions there is an elementary bound
\begin{equation}
 |\Phi_{g,d}(x)|\ll_g\frac{x\log^4(2x)}{\sqrt{|d|}},\qquad |d|\le4x.
 \label{op:lattice}
\end{equation}
Indeed $1\le f\le2\sqrt{x/|d|}$ and $|t|\le2\sqrt x$ in $4p=t^2+|d|f^2$, and $(t,f)$ fixes the prime. Each summand is $O_g(\log^4(2x))$: use $\sqrt{(4p-t^2)/p}\le2$, $L(1,\chi_d)\ll\log(2|d|)$, and
$G(f,d)\le\prod_{\ell\mid f}(1+2/(\ell-1))\ll\log^2(2f)$.
For $E$, the same bound holds for both terms of $J_E$.

It follows unconditionally that, for all sufficiently large $x$, each of the signs of $\Phi_d(x)$ for the conductor-$11$ curve occurs at least $\gg x/\log^8(2x)$ times among $|d|\gg x/\log^8(2x)$. In fact both signs carry total mass $\gg x^{3/2}$ as proved above. The total absolute contribution of $|d|\le y$ is $O(x\sqrt y\log^4(2x))$ by \eqref{op:lattice}. Take $y=c_1x/\log^8(2x)$ with $c_1>0$ small enough that this is less than half of each sign's mass. Every field left contributes at most $O(\sqrt x\log^8(2x))$, which proves the count.

For the six level-one eigenforms the signs visibly vary with the order of the limits. Take $D=x^{2(1-\theta)}/\log^{20}(2x)$. Uniformly through $|d|\le D$, their positive main term dominates the individual error in Corollary~\ref{op:levelone}, and their combined contribution is $\gg_k x\sqrt D$. Since the full signed sum is $O_k(x)$, negative contributions among $|d|>D$ must have total magnitude $\gg_k x\sqrt D$. Each such field has contribution at most $O_k(x\log^4(2x)/\sqrt D)$ by \eqref{op:lattice}. Thus at every sufficiently large $x$ there are at least
\[
     \gg_{k,\theta}x^{2(1-\theta)}/\log^{24}(2x)
\]
distinct fields with negative contribution although each fixed field eventually has positive contribution. At the input $\theta=7/8$ the numerator is $x^{1/4}$.

\clearpage
\section{Trying to put a weight-two eigenform into the geometric prime count}\label{geo:bridge}

Let $E/\Q$ be an elliptic curve with conductor $N$. For $p\nmid N$, put
\begin{align*}
 a_p&=p+1-\#E(\F_p), &t_p&=a_p/\sqrt p,\\
 U_p&=\operatorname{diag}(e^{i\theta_p},e^{-i\theta_p}),&
 2\cos\theta_p&=t_p,\quad 0\le\theta_p\le\pi .
\end{align*}
Hasse's bound makes this choice possible. These matrices represent the unitary Satake classes of the modular form attached to $E$. In this section $L_u(s,E)=L(s+1/2,E)$, so the target critical line in $s$ is $\operatorname{Re}s=1/2$. Its unramified logarithmic derivative has coefficients
\begin{equation}
 -\frac{(L_u^S)'(s,E)}{L_u^S(s,E)}
   =\sum_{p\notin S}\sum_{j\ge1}
       \operatorname{tr}(U_p^j)(\log p)p^{-js},
       \qquad \operatorname{Re}s>1 ,
 \label{geo:newformld}
\end{equation}
where $S$ contains the bad primes. Both eigenvalues have absolute value one. The coefficient at $p^2$ is $t_p^2-2$, whereas the normalized Hecke eigenvalue at $p^2$ is $t_p^2-1$. Putting $t_p^2$ or the latter Hecke eigenvalue in the logarithmic derivative would already change the Euler factor.

We will prove three assertions about the proposed geometric extension. At a fixed level the earlier prime count has zero projection on the elliptic trace in a natural prime mean. A growing torsion field does recover this trace exactly on a finite interval. In that field the recovering function has no one-dimensional Fourier coefficient, and its higher-dimensional Fourier coefficients have quantitatively large total cost. Together the assertions locate the missing arithmetic estimate. None of them settles the location of an individual elliptic-curve zero.

\subsection{The exact point where the two problems diverge}

The original geometric row and the prime weight it produces are different objects. Here is a precise account of the data passed between the two steps.

\begin{enumerate}
\item A frame specifies a finite \emph{commutative} group $A$, a cyclic subgroup of it, and a projection to an abelian cyclotomic layer. The field used by a host comes with both an identification of its Galois group with a product of $A$ and two sign groups, and an embedding into a cyclotomic field. Its Artin map on reduced residue classes modulo a single integer gives the associated Frobenius element. These restrictions appear at the definitions of the frame and host. \href{https://github.com/Deicyde/riemann-hypothesis/blob/ba2a065656ec7c638baa601994de3e0245ff6665/src/Formalism/SharpRestriction/RamificationFrobenius/TotientGaloisClasses/ZsmulGeneratorEqZeroOfProjection.lean#L15}{Frame.} \href{https://github.com/Deicyde/riemann-hypothesis/blob/ba2a065656ec7c638baa601994de3e0245ff6665/src/Formalism/SharpRestriction/RamificationFrobenius/TotientGaloisClasses/ConjEMemHostFIff.lean#L12}{Host, cyclotomic embedding and Artin map.}
\item At a prime of that field the count of partners after $j$ steps is read from translation by $j\operatorname{Frob}_p$ in that finite group, with averaging over inertia when needed. Finite Fourier inversion on the group therefore reads, on the $\chi$-mode, the scalar $\chi(\operatorname{Frob}_p)^j$ at every unramified prime. The resulting Dirichlet character is identified with its primitive character when the Euler factors at bad primes are accounted for. \href{https://github.com/Deicyde/riemann-hypothesis/blob/ba2a065656ec7c638baa601994de3e0245ff6665/src/Development/AutomorphicArithmetic/PrimitiveCosets/PrimRead.lean#L48}{Prime-sum identification.}
\item The global spectrum is used to obtain a lower bound for the dimension of a space on which the auxiliary restriction estimates hold. Its local labels are polynomials over the two auxiliary residue fields. The formal row-count identity sums local idempotent weights for the label tuples, and the surviving space counts rows passing the two tests. It places no occurrence of $t_p$ in the partner count of the preceding step. \href{https://github.com/Deicyde/riemann-hypothesis/blob/ba2a065656ec7c638baa601994de3e0245ff6665/src/Formalism/SharpRestriction/LevelLetterCocycles/ArithmeticCylinders.lean#L244}{Fields of the spectral row data.}
\item In the retained linear combination the index ranges over primitive Dirichlet characters. In particular the definition of the cyclic level takes the subgroup generated by one such character. The meromorphic functions that enter the pole argument are fixed at this step, even though the auxiliary rows varied earlier. \href{https://github.com/Deicyde/riemann-hypothesis/blob/ba2a065656ec7c638baa601994de3e0245ff6665/src/Formalism/SharpRestriction/RamificationFrobenius/TotientGaloisClasses/ExplicitFormulaErrorOfHostConditions.lean#L38}{The centered sum.} \href{https://github.com/Deicyde/riemann-hypothesis/blob/ba2a065656ec7c638baa601994de3e0245ff6665/src/Formalism/SharpRestriction/RamificationFrobenius/TotientGaloisClasses/ExplicitFormulaErrorOfHostConditions.lean#L121}{The cyclic level.}
\end{enumerate}

For example, consider
\[
  E:\quad y^2+y=x^3-x^2-10x-20.
\]
This is the conductor-$11$ curve with Cremona label $11\mathrm{a}1$, corresponding to the newform $11.2.\mathrm{a}.\mathrm{a}$. The primes $23,67,89$ are all $1$ modulo $11$. Their point counts are $25,75,75$, respectively, and their Frobenius traces are $-1,-7,15$. These counts can be computed without a modular-symbol calculation. For odd $p$ the number of affine points above $x$ is $1+(\frac{1+4(x^3-x^2-10x-20)}p)$, to which one adds the point at infinity after summing over $x$. Thus retaining every character modulo the level of the newform still assigns identical first-prime character values to these three distinct normalized Hecke eigenvalues.
\href{https://www.lmfdb.org/EllipticCurve/Q/11/a/2}{The curve $11.\mathrm{a}2$.}
\href{https://www.lmfdb.org/ModularForm/GL2/Q/holomorphic/11/2/a/a/}{Its weight-two newform.}

\begin{theorem}\label{geo:fixed}
For every elliptic curve $E/\Q$ and every finite Galois extension $K/\Q$, the numbers $t_p$ take infinitely many different values among good primes splitting completely in $K$. In particular no fixed finite Frobenius class function can give these values at all but finitely many primes.
\end{theorem}
\begin{proof}
Put $F=K\Q(E[3])$. Apply Chebotarev's theorem to the identity class of $F/\Q$. It supplies infinitely many completely split good primes $p\ne3$. The characteristic polynomial of arithmetic Frobenius on $E[3]$ has trace $a_p\bmod3$. At the chosen primes the action is the identity, hence $a_p\equiv2\pmod3$ and $a_p\ne0$. If $p,q$ are two distinct such primes and $a_p/\sqrt p=a_q/\sqrt q$, then $p/q=(a_p/a_q)^2$ is a square in $\Q^\times$. Its valuation at $p$ equals one. This contradiction proves that all these values are distinct. The Frobenius class in $K$ is the same at all of them.
\end{proof}

This obstruction uses neither non-CM nor an equidistribution theorem for Satake angles. It continues to hold after any fixed additional finite label is introduced, by replacing $K$ by the common Galois closure. It applies to the normalized coefficient and to functions of the finite label alone. A procedure which also uses the size of $p$, or allows the field to vary with the search interval, requires a separate analysis.
\href{https://math.mit.edu/classes/18.785/2018fa/LectureNotes28.pdf}{Chebotarev's density theorem, Theorem 28.9.}
\href{https://math.uni.lu/nt/summerschool2018/AriasdeReyna-EC-notes.pdf}{Galois representations from elliptic curves and their Frobenius polynomials.}

\subsection{Orthogonality to every fixed abelian host}

There is a stronger obstruction to a pointwise approximation on average. Fix a non-CM elliptic curve $E/\Q$. For a bounded function $h$ on the good primes define a prime logarithmic mean, when it exists, by
\[
 \langle h\rangle=\lim_{\sigma\downarrow1}
   \frac{\sum_{p\nmid N}h(p)p^{-\sigma}}
        {\log(1/(\sigma-1))}.
\]
Changing finitely many terms leaves this limit unchanged. Euler's product for $\zeta$ implies $\langle1\rangle=1$. This is the Abelian form of the prime density mean. Its use here avoids any uniform prime number theorem in a growing modulus.

\begin{theorem}\label{geo:l2}
Let $M\ge1$ and let $u:(\Z/M\Z)^\times\to\C$ be any function. In the preceding mean omit primes dividing $M$. Then
\begin{align*}
  \langle t_p\overline{u(p)}\rangle&=0,\\
  \langle|t_p-u(p)|^2\rangle
     &=1+\frac1{\varphi(M)}\sum_{b\bmod M}^{*}|u(b)|^2 .
\end{align*}
In particular the least mean-square error over all such fixed abelian weights is exactly one, attained by the zero weight.
\end{theorem}
\begin{proof}
Expand $u$ in the Dirichlet characters modulo $M$. For one character $\chi$, take a logarithm of the Euler product of the standard twist of the unitary representation attached to $E$. For real $\sigma>1$,
\[
 \log L^S(\sigma,E_u\otimes\chi)
   =\sum_{p\notin S}t_p\chi(p)p^{-\sigma}+O(1)
       \qquad(\sigma\downarrow1).
\]
Here $S$ contains the primes of bad reduction and those dividing $M$. All higher prime powers contribute a bounded quantity, since the absolute value of their trace is at most two and $\sum_p\sum_{j\ge2}p^{-j}/j<\infty$. A consistent logarithm is bounded near $1$ by holomorphy and nonvanishing there of the standard $L$-function of the cuspidal twist. The omitted local Euler factors at $1$ are finite and nonzero, using the local factors of an elliptic curve and its finite characters. Consequently the prime sum divided by $\log(1/(\sigma-1))$ tends to zero for each fixed $\chi$. This proves the first equality.

For the second equality, the trace at $p$ of the unitary symmetric-square lift is $e^{2i\theta_p}+1+e^{-2i\theta_p}=t_p^2-1$. Since $E$ has no CM, its symmetric-square lift is cuspidal by Gelbart--Jacquet. Its $L$-function is also holomorphic and nonzero at one. The same logarithm argument, now with the higher traces bounded by three, yields $\langle t_p^2-1\rangle=0$. The finite ramified local factors can again be removed at $1$; equivalently use the tempered local Weil--Deligne parameters of the unitary representation for an elliptic curve. Finally Dirichlet characters have mean zero except for the principal character, by the identical logarithm argument for $L(s,\chi)$ and Dirichlet nonvanishing. Hence $\langle|u|^2\rangle=\varphi(M)^{-1}\sum^*|u(b)|^2$. Expand the square and use $t_p\in\R$.
\end{proof}

The inputs at one are known unconditionally. The theorem gives no rate of convergence in $M$ and does not exclude cancellation in the signed difference between two prime sums. It excludes an inference from a small mean-square error on a fixed host, even if an arbitrary complex coefficient is assigned separately to every residue class.
\href{https://www.numdam.org/articles/10.24033/asens.1355/}{S. Gelbart and H. Jacquet, Ann. Sci. ENS 11 (1978), 471--542.}
H. Jacquet and J. Shalika, \emph{A non-vanishing theorem for zeta functions of $\mathrm{GL}_n$}, Invent. Math. 38 (1976), 1--16.

\subsection{A positive local weight for the correct logarithmic derivative}

There is no sign obstruction at the level of individual unitary Euler factors. For $0<\eta\le1/2$ give $\mathrm{SU}(2)$ its Haar probability measure $dg$ and put
\[
          d\nu_\eta(g)=(1+\eta\operatorname{tr}g)\,dg .
\]
The eigenvalues of $g$ are conjugates of absolute value one, so this is nonnegative. Its mass is one, since the central involution $g\mapsto-g$ changes the sign of its trace. If $\chi_j$ denotes the character of $\operatorname{Sym}^j(\C^2)$, the irreducible character orthogonality relation gives
\[
  \int\chi_j(g)\,d\nu_\eta(g)=
     \begin{cases}
         1 & j=0,\\
         \eta & j=1,\\
         0 & j\ge2 .
     \end{cases}
\]
This calculation explains exactly which average would preserve the standard representation. For an unramified prime and its $k$th power the value required of the weight is
\[
  1+\eta\operatorname{tr}(U_p^k)=1+2\eta\cos(k\theta_p)\ge0.
\]
Summing over $p,k$ gives the logarithmic derivative
\[
   -\frac{(\zeta^S)'}{\zeta^S}(s)
          -\eta\frac{(L_u^S)'(s,E)}{L_u^S(s,E)}
       \qquad(\operatorname{Re}s>1).
\]
This identity uses ordinary pointwise evaluation of the weight at $U_p^k$. Replacing $\operatorname{tr}(U_p^k)$ by $(\operatorname{tr}U_p)^k$ would produce another function. For a tempered, unitary, degree-$d$ parameter $U$ with nonreal trace, $1+\eta\operatorname{Re}\operatorname{tr}(U)$ is likewise nonnegative for $0<\eta\le1/d$. The corresponding two logarithmic derivatives have coefficients $\eta/2$ for the representation and for its contragredient.

These formulas give a target for extending the geometric count. They supply no bound for the discrepancy of this weight on the primes of an elliptic curve, and no system of geometric translates whose count evaluates it. In particular the rank bound for locally admissible auxiliary rows does not construct a Haar sample of the $U_p$.

\subsection{Exact recovery on a dyadic interval using torsion}

A finite approximation becomes exact if the field can grow with the prime interval. For an odd prime $\ell$ write $[u]_\ell$ for the integer in $[-(\ell-1)/2,(\ell-1)/2]$ reducing to $u\in\F_\ell$. Define the real class function
\[
      h_\ell(A)=\frac{[\operatorname{tr}A]_\ell}{\ell},
      \qquad A\in G_\ell=\mathrm{GL}_2(\F_\ell).
\]
It takes both signs. If a positive measure is wanted, $1+2h_\ell$ is positive everywhere. Put $K_\ell=\Q(E[\ell])$, with Galois representation $\bar\rho_\ell$. For good $p\ne\ell$ its characteristic polynomial on arithmetic Frobenius reduces to $X^2-a_p X+p\pmod\ell$.

\begin{lemma}\label{geo:decode}
For any $X\ge2$ and prime $\ell>4\sqrt X$,
\[
  t_p=\frac{\ell}{\sqrt p}h_\ell(\bar\rho_\ell(\operatorname{Frob}_p))
        \qquad(p\le X,\ p\nmid N\ell).
\]
Knowing the integer $p$ and this class determines the coefficient in \eqref{geo:newformld} at every $p^j$.
\end{lemma}
\begin{proof}
Hasse's inequality gives $|a_p|\le2\sqrt p<\ell/2$, so the centered representative of the trace residue is $a_p$. For all powers, let $b_0=2$, $b_1=a_p$, $b_{j+1}=a_pb_j-pb_{j-1}$. The characteristic polynomial of Frobenius shows by induction that $b_j$ is the sum of the $j$th powers of its two algebraic eigenvalues. Thus $b_j p^{-j/2}=\operatorname{tr}(U_p^j)$.
\end{proof}

It is important here that $\ell^{-1}a_p$ is recovered first and $p^{-1/2}$ is supplied afterwards. The finite labels at this stage do not directly give the unitary angles. For non-CM elliptic curves, Serre's open image theorem implies that $\bar\rho_\ell$ is surjective for every sufficiently large $\ell$. Indeed an open subgroup of $\mathrm{GL}_2(\widehat\Z)$ contains a principal congruence subgroup, whose reductions are the entire $G_\ell$ for $\ell$ prime to the fixed level. By Bertrand's postulate, when $P$ is large we can choose
\[
      4\sqrt{2P}<\ell< 2\lceil4\sqrt{2P}\rceil+2 .
\]
Then $\ell\asymp\sqrt P$, all the primes of $(P,2P]$ differ from $\ell$, and one surjective torsion representation recovers the whole interval. Its Galois degree is
\[
 |G_\ell|=(\ell^2-1)(\ell^2-\ell)\asymp P^2.
\]
The dependence on $P$ in this degree matters if a class-counting estimate has an implied constant depending on the extension.
\href{https://pi.math.cornell.edu/~zywina/papers/ExplicitOpenImage.pdf}{D. Zywina, Theorem 1.1, statement of Serre's adelic open image theorem.}

\subsection{The Fourier cost of reading the centered trace}

For $G=G_\ell$ with its uniform probability measure, normalize the inner product by $\langle f_1,f_2\rangle_G=|G|^{-1}\sum_A f_1(A)\overline{f_2(A)}$. Let $\widehat G$ be the complex irreducible representations, with characters $\chi_\pi$ and degrees $d_\pi$. We use the actual real-valued observable in Lemma~\ref{geo:decode}. This avoids confusing a trace in characteristic $\ell$ with the character of a complex representation.

\begin{theorem}\label{geo:fouriercost}
Let $\ell\ge3$ be prime and write the character expansion $h_\ell=\sum_{\pi\in\widehat G}c_\pi\chi_\pi$. All terms with $d_\pi=1$ have $c_\pi=0$. Furthermore
\[
  \|h_\ell\|_{L^2(G)}^2\ge\frac{\ell-1}{12\ell},
  \qquad
  \sum_{\pi\in\widehat G}d_\pi|c_\pi|
      \ge(\ell-1)\sqrt{\frac{\ell-1}{12\ell}} .
\]
More generally the first assertion holds for any odd function of the trace, where odd means $F(-u)=-F(u)$ in the ordinary complex values of the function.
\end{theorem}
\begin{proof}
For $t\in\F_\ell$ and $d\in\F_\ell^\times$, the number of matrices of trace $t$ and determinant $d$ equals
\begin{equation}
     \ell\bigl(\ell+\left(\tfrac{t^2-4d}{\ell}\right)\bigr).
 \label{geo:matrixcount}
\end{equation}
Indeed write the matrix as $\left(\begin{smallmatrix}x&y\\z&t-x\end{smallmatrix}\right)$. The equation $yz=x(t-x)-d$ has $\ell-1$ solutions in $(y,z)$ if the right side is nonzero, and $2\ell-1$ if it is zero. There are $1+(\frac{t^2-4d}{\ell})$ values of $x$ in the second case. This proves \eqref{geo:matrixcount}. The conditional distribution of $t$ at any determinant is symmetric under $t\mapsto-t$; alternatively $A\mapsto-A$ preserves its determinant. Every determinant value occurs on $|\mathrm{SL}_2(\F_\ell)|=\ell(\ell^2-1)$ matrices. Therefore its conditional mass at each $t$ is at least $1/(\ell+1)$. It follows that
\begin{align*}
 \mathbb E\bigl[h_\ell^2\mid\det=d\bigr]
  &\ge \frac1{(\ell+1)\ell^2}
             \sum_{k=-(\ell-1)/2}^{(\ell-1)/2} k^2
    =\frac{\ell-1}{12\ell}.
\end{align*}
The conditional mean itself is zero for every determinant, including for any odd function as in the statement.

We need two representation facts for this small finite group, which we prove directly. Set $U=\{I+x e_{12}:x\in\F_\ell\}$. The normal subgroup generated by $U$ contains the lower elementary transvections as well, by conjugation by a matrix interchanging the basis vectors, and these two types generate $\mathrm{SL}_2(\F_\ell)$ by elementary row elimination. For precision, diagonal matrices remaining after elimination are included because, writing $u(x)=I+xe_{12}$, $v(x)=I+xe_{21}$, the product $w(s)=u(s)v(-s^{-1})u(s)$ is $\left(\begin{smallmatrix}0&s\\-s^{-1}&0\end{smallmatrix}\right)$ and $w(s)w(-1)=\operatorname{diag}(s,s^{-1})$. Moreover each $u(x)$ is a commutator $[\operatorname{diag}(a,1),u(x/(a-1))]$ with $a\ne0,1$. Thus the commutator subgroup is $\mathrm{SL}_2(\F_\ell)$, since determinants kill every commutator. All one-dimensional complex characters factor through the determinant. Conditional mean zero kills each of those coefficients.

An irreducible representation $\pi$ that is nontrivial on $\mathrm{SL}_2$ must be nontrivial on $U$ by normal generation. Its restriction to the finite abelian $U$ splits into characters, at least one nontrivial. Conjugation by $\operatorname{diag}(a,1)$, $a\in\F_\ell^\times$, moves any nontrivial character of $U$ through all $\ell-1$ such characters. Their weight spaces have the same positive dimension. This proves $d_\pi\ge\ell-1$ for every remaining constituent. Character orthonormality and the elementary inequality between the $\ell^1$ and $\ell^2$ norms now give
\[
 \sum d_\pi|c_\pi|\ge(\ell-1)\sum|c_\pi|
 \ge(\ell-1)\left(\sum|c_\pi|^2\right)^{1/2}
 =(\ell-1)\|h_\ell\|_{L^2(G)}.
\]
Together with the conditional second moment this proves the assertion.
\end{proof}

The positivity correction $1+2h_\ell$ has constant coefficient one and exactly twice the old nonconstant coefficients. Biasing a finite group in this concrete way is possible. The modes introduced by this bias all lie outside its abelian quotient. This is already a restriction at the initial field choice in the geometric construction: a cyclotomic embedding of $\Q(E[\ell])$ would force $G_\ell$ to be abelian.

The same obstruction persists after an unrelated finite abelian host $K$ is attached to the torsion field. To see this, let $F=K\cap K_\ell$. It is abelian over $\Q$ and the fixing subgroup of $F$ in $G_\ell$ therefore contains $\mathrm{SL}_2(\F_\ell)$. Both $I$ and $-I$ act trivially on $F$, so each can be combined with the identity action on $K$ to give an element of $\operatorname{Gal}(KK_\ell/\Q)$. This follows either by the fiber product description of Galois groups of composita or by extending compatible automorphisms. Chebotarev supplies infinitely many good primes of each conjugacy class. Both classes are trivial in the chosen abelian host, whereas the raw traces are congruent to $2$ and $-2$ respectively modulo $\ell$ for these primes. Take $\ell\ge5$ for distinct residues. This argument concerns a fixed $\ell$ and infinitely many primes, without asserting where the first such primes occur.

\subsection{A sufficient estimate where uniformity can be tested}

We can formulate a concrete arithmetic input for a non-CM curve without guessing how the geometry might produce it. For each large dyadic $P$ choose $\ell(P)$ as above, and abbreviate $A_{p,\ell}=\bar\rho_\ell(\operatorname{Frob}_p)$. Suppose, for every $\varepsilon>0$, that uniformly for $P\le V\le2P$,
\begin{equation}
 \left|\sum_{P<p\le V,\ p\nmid N}
          h_{\ell(P)}(A_{p,\ell(P)})(\log p)\right|
     \ll_{E,\varepsilon} P^{1/2+\varepsilon}.
 \label{geo:targetest}
\end{equation}
The implied constant must be independent of $P$ and the chosen probe. Below we verify the precise consequence of this condition and the degree dependence of a direct attempt to prove it.

\begin{theorem}\label{geo:dyadic}
For the fixed choices of $\ell(P)\asymp\sqrt P$ made above, \eqref{geo:targetest} is equivalent to the same uniform dyadic bound with $t_p$ in place of $h_{\ell(P)}(A_{p,\ell(P)})$. Either bound implies the Riemann hypothesis for $L_u(s,E)$.
\end{theorem}
\begin{proof}
On this interval $t_p=w(p)h_{\ell(P)}(A_{p,\ell(P)})$ with $w(v)=\ell(P)/\sqrt v$. Both $w$ and $1/w$ are bounded independently of the interval. Each is monotone and has total variation $O(1)$ on $[P,2P]$. Abel summation in either direction multiplies the supremum of the partial sums by at most an absolute constant. This proves the equivalence.

Summing over dyadic intervals gives $\sum_{p\le x,p\nmid N}t_p\log p\ll_{E,\varepsilon}x^{1/2+\varepsilon}$. The absolute contribution of $p^j\le x$ with $j\ge2$ to the Chebyshev sum for \eqref{geo:newformld} is $O(\sqrt x\log^2(2x))$. Indeed each trace is bounded by two, each $p$ is at most $\sqrt x$, and for fixed $p$ the sum of the log weights is at most $\log x$. This even gives $O(\sqrt x\log(2x))$ by the trivial bound for the number of $p$. Removing the fixed bad Euler factors leaves a finite product holomorphic and nonzero for $\operatorname{Re}s>0$ in unitary normalization. In the multiplicative case the local root is $\pm p^{-1/2}$, and in the additive case the factor equals one. Thus it is enough to treat $L_u^S$.

Write $B(x)$ for the incomplete Chebyshev sum over prime powers obtained above. On $\operatorname{Re}s>1$, \eqref{geo:newformld} gives
\[
    -\frac{(L_u^S)'(s,E)}{L_u^S(s,E)}
     =s\int_1^\infty B(v)v^{-s-1}\,dv.
\]
For any compact subset of $\operatorname{Re}s>1/2$, choose $\varepsilon$ smaller than its distance from the boundary. The integral converges uniformly there and defines a holomorphic function, using the bound proved for $B$. By uniqueness of meromorphic continuation it agrees with the logarithmic derivative throughout that half-plane. A zero of $L_u^S$ would give a pole of this logarithmic derivative. Hence no zero occurs there. The weight-two functional equation, in unitary coordinates $s\leftrightarrow1-s$, excludes open-strip zeros on the other side as well. The gamma factors are regular and nonzero in this strip.
\end{proof}

To see the scope of the remaining difficulty, decompose the prime sum in \eqref{geo:targetest} into its $c_\pi$ modes. Even if one had separate estimates of the shape $B(P)d_\pi$ for each mode, a termwise absolute-value argument from those estimates would give the allowance
\[
      B(P)\sum_{\pi}d_\pi|c_\pi|
        \ \ge\ B(P)(\ell(P)-1)\sqrt{\frac{\ell(P)-1}{12\ell(P)}} .
\]
For a square-root allowance $B(P)$ and $\ell(P)\asymp\sqrt P$ this expression already has order at least $P$, before any conductor logarithms. This is a lower bound for what that \emph{particular bounding expression} costs. It gives no lower bound for the original signed sum, in which there can be cancellation across the representations. A proof of \eqref{geo:targetest} would need a direct trace estimate, cancellation between the modes, or a genuinely stronger uniform ingredient. The original abelian host arithmetic supplies none of these estimates for $h_\ell$. The finite spectral dimension bound concerns other observables and has no implication of the required strength for this one.

\subsection{A lower-cost expansion by additive trace phases}

The factor of order $\ell$ in Theorem~\ref{geo:fouriercost} arises from decomposing the observable into ordinary irreducible \emph{group} characters and applying degree-dependent estimates separately. The trace also takes values in the additive cyclic group $\F_\ell$. Fourier inversion on this latter group gives a different and smaller coefficient norm. Its basis functions are exponential sums of the trace, for which the requisite estimates are open in the present range.

\begin{lemma}\label{geo:sineexp}
Let $\ell$ be an odd prime and $m=(\ell-1)/2$. For every $t\in\F_\ell$ one has the exact identity
\[
 \frac{[t]_\ell}{\ell}
   =\sum_{j=1}^m b_j
       \sin\!\left(\frac{2\pi jt}{\ell}\right),\qquad
 b_j=\frac{(-1)^{j+1}}{\ell\sin(\pi j/\ell)} ,
\]
and, writing $H_m=\sum_{j=1}^m1/j$,
\[
   \frac{H_m}{\pi}\le\sum_{j=1}^m |b_j|\le\frac{H_m}{2}.
\]
\end{lemma}
\begin{proof}
The Dirichlet kernel $D_m(z)=\sum_{k=-m}^m e^{ikz}=\sin(\ell z/2)/\sin(z/2)$ gives
$D'_m(2\pi j/\ell)=\ell(-1)^j/(2\sin(\pi j/\ell))$. Differentiating its finite sum also gives $D'_m(z)=-\sum_{k=-m}^m k\sin(kz)$. The sine coefficient of the odd function $k\mapsto k/\ell$ is
\[
   \frac{2}{\ell^2}\sum_{k=-m}^m k\sin(2\pi jk/\ell)
     =-\frac2{\ell^2}D'_m(2\pi j/\ell)=b_j .
\]
Discrete Fourier orthogonality on the additive group proves the expansion. On $0<u\le\pi/2$ use $2u/\pi\le\sin u\le u$ and sum over $j$ for the bounds.
\end{proof}

For the dyadic prime $\ell=\ell(P)$ of Lemma~\ref{geo:decode}, put
\[
 S_j(P,V)=\sum_{P<p\le V,\ p\nmid N}
     (\log p)\sin\!\left(\frac{2\pi j a_p}{\ell}\right).
\]
Each sine observable on $G_\ell$ is conjugation invariant and has mean zero at every fixed determinant by $A\mapsto-A$. Thus there is no main term of size $P$ from the trivial group character. The exact recovery gives
\begin{equation}
 \sum_{P<p\le V,\ p\nmid N}
    h_{\ell(P)}(A_{p,\ell(P)})\log p
      =\sum_{j=1}^{(\ell-1)/2}b_jS_j(P,V).
 \label{geo:phases}
\end{equation}
Consequently, the uniform bound $|S_j(P,V)|\ll_{E,\varepsilon}P^{1/2+\varepsilon}$ for $P\le V\le2P$ and $1\le j\le(\ell-1)/2$ would imply elliptic-curve RH. Indeed the logarithm from Lemma~\ref{geo:sineexp} can be absorbed by starting with a smaller $\varepsilon$ and then applying Theorem~\ref{geo:dyadic}. More generally it suffices to bound the sum in \eqref{geo:phases} with its signs. The identity is independent of any assumption on zeros of division-field Artin functions.

There is an exact second moment revealing what averaging the additive phases can measure. Put $C_t=\sum_{P<p\le V,\,p\nmid N,\,a_p\equiv t\, (\bmod\ell)}\log p$ for $t\in\F_\ell$. With $S_j$ extended to $0\le j<\ell$ one has
\begin{equation}
 \frac1\ell\sum_{j=0}^{\ell-1}|S_j(P,V)|^2
       =\frac14\sum_{t\in\F_\ell}(C_t-C_{-t})^2 .
 \label{geo:phasemoment}
\end{equation}
Indeed the average product of the sines at $t,u$ is $(\mathbf1_{t=u}-\mathbf1_{t=-u})/2$, by finite Fourier orthogonality. Multiplying by $C_tC_u$ and summing proves the formula. Under the decoding condition different allowed integer traces are distinct residues. Thus this second moment measures the discrepancy between the total log weights at each integer trace and its negative. It retains no contribution from common increases in the counts at $t$ and $-t$. Neither Chebotarev for fixed $\ell$ nor the original spectral rank estimate provides a bound for \eqref{geo:phasemoment} uniform at $\ell\asymp\sqrt P$.

One other direct simplification can be excluded. The permutation character of the action of $G_\ell$ on the cosets of any subgroup has the same value at $A$ and $A^{-1}$, since those elements fix precisely the same cosets. More generally it is constant on generators of the same cyclic subgroup. For $\ell\ge5$ choose a generator $u$ of $\F_\ell^\times$. Then $uI$ and $u^{-1}I$ generate the same subgroup but their centered traces differ because $2u\ne2u^{-1}$. Thus no linear combination, even over $\C$, of permutation characters of subgroups of $G_\ell$ equals $h_\ell$. As the first logarithmic coefficients of Dedekind zeta functions of subfields of $K_\ell$ are exactly those permutation characters at unramified Frobenius classes, those zeta functions cannot furnish this equality at all primes. This last deduction uses surjectivity and Chebotarev for the fixed extension $K_\ell$. It leaves open bounds for the varying-interval prime sums.

\subsection{What remains valid in the local valuation calculation}

The failure of the cyclotomic readout does not affect elementary estimates on one marked Frobenius orbit. Here is the precise portion which survives a change of local labels. Fix $p\ge2$, a cycle $\Z/D\Z$ of even length, and a marked set $T$. Parity is well defined on this cycle. Write $T_e,T_o$ for the marks at even and odd indices. Fix a scale $a$ with $2a\ge p$, a loss $c\ge0$, and numbers $2c\le a_s\le2a$ at the marks. At each site $z$ let $W_z\ge1$ and $0\le U_z\le W_z+c$. For a mark $s$ put
\[
 V_s=\sum_{i\ge1,\ i\ {\rm odd}}(a_s p^{-i}-W_{s+i})_+,
 \qquad
 M_s=\sum_{i\ge1,\ i\ {\rm odd}}\min(a_s p^{-i},W_{s+i}),
 \qquad \gamma=\frac{p}{p^2-1}.
\]
The first sum has finite support. Their sum is $\gamma a_s$, by adding the two summands at each $i$. Consider the local expression
\[
 G=\gamma\sum_{s\in T}(a_s-2c)
      -\sum_{s\in T_o}\sum_{i\ge1,\ i\ {\rm odd}}
                    \min((a_s-2c)p^{-i},U_{s+i}).
\]
It is the part of the unipotent local gain before the orbit error is deducted. For $y\in T_o$ take its nearest preceding even mark $x_y$, when one exists, using strictly positive forward cyclic distance $d_y\le D$. Put $0\le b_y\le W_y$. Then
\begin{equation}
 G\ \ge\ \sum_{s\in T} V_s+
       \sum_{y\in T_o,\ x_y\ \mathrm{exists}}
           \min(b_y,a_{x_y}p^{-d_y})-
                 3c|T|\log_p(2a).
 \label{geo:localgain}
\end{equation}
To prove it, an even mark contributes $\gamma(a_s-2c)=V_s+M_s-2\gamma c$. At an odd mark its contribution is $\sum_i((a_s-2c)p^{-i}-U_{s+i})_+$ over positive odd $i$. For terms positive in $V_s$ the change costs at most $2cp^{-i}+c$. Every such $i$ satisfies $i\le\log_p(2a)$ since $W_{s+i}\ge1$. The total cost is at most $3c\log_p(2a)$, by $2p^{-i}\le1$ for $p\ge2$ (even $2c\log_p(2a)$ suffices here). This bound also absorbs $2\gamma c$ at an even mark, since $\log_p(2a)\ge1$. Finally $d_y$ is odd, and different $y$ yield different pairs $(x_y,d_y)$. The indicated term for $y$ is at most the corresponding summand in $M_{x_y}$. Summing the estimates proves \eqref{geo:localgain}.
\href{https://github.com/Deicyde/riemann-hypothesis/blob/ba2a065656ec7c638baa601994de3e0245ff6665/src/Formalism/SharpRestriction/GrowthBounds/LogarithmicSums/PartnerCountOfCentral.lean#L136}{Parity and predecessor estimate.}
\href{https://github.com/Deicyde/riemann-hypothesis/blob/ba2a065656ec7c638baa601994de3e0245ff6665/src/Formalism/SharpRestriction/GrowthBounds/LogarithmicSums/PartnerCountOfCentral.lean#L308}{Local gain comparison.}

Multiplication by any nonnegative weight depending on $p$ preserves \eqref{geo:localgain}. At the subsequent global determinant step one uses the product formula. Its multiplicities admit an elementary rigidity statement which precludes simply putting the same prime weights on a fixed determinant line.

\begin{lemma}\label{geo:rigidity}
Let $F$ be a number field. For nonzero $z\in F$, put $L_v(z)=\operatorname{ord}_v(z)\log\operatorname{N}v$ at each finite place and $L_w(z)=[F_w:\R]\log|z|_w$ at the infinite places, taking the usual absolute value at a real or complex embedding. Suppose real coefficients $c_v,d_w$ and a fixed $B$ satisfy
\[
  \left|\sum_{v\nmid\infty}c_v L_v(z)
            -\sum_{w\mid\infty}d_w L_w(z)\right|\le B
                   \qquad (z\in F^\times).
\]
Then all $d_w$ and all $c_v$ equal one common constant.
\end{lemma}
\begin{proof}
Apply the inequality to $z^n$ and divide by $n$ to obtain equality for each $z$. Applied to a unit, it says that $(d_w)_w$ is perpendicular to every vector of weighted archimedean logarithms of units. By Dirichlet's unit theorem these vectors span the hyperplane with sum of coordinates zero. Thus all $d_w$ equal one constant $d$. The same conclusion is immediate for one infinite place. For any fixed finite $v$, the ideal class group is finite, so a power $\mathfrak p_v^h$ is generated by an element $z$. Only $L_v(z)$ is nonzero at finite places, and the ordinary norm formula gives $\sum_wL_w(z)=h\log\operatorname{N}v$. The equality just proved forces $c_v=d$. This holds at every $v$.
\end{proof}

For $F=\Q$ an explicit obstruction is available even if the ordinary archimedean absolute value is kept bounded. Take primes $p,q$ with $c_p\ne c_q$ and elements $z_n=p^n/q^{m_n}$, with $m_n$ nearest to $n\log p/\log q$. Then $|\log|z_n||\le(\log q)/2$ but the reweighted finite sum is $(c_p-c_q)n\log p+O(1)$. In particular one may take $c_p=1+\eta t_p$ and $p=23,q=89$ for the curve above. A globally defined arithmetic degree has to be independent of the basis of its determinant line, and a scalar basis change already tests this product formula. Thus a successful modification must identify the desired trace weight within the distribution of geometric local gains, or change the global construction in a way that restores basis independence. This lemma addresses basis changes on a fixed line; it does not bound the local gains of a new system of maps.
\href{https://github.com/Deicyde/riemann-hypothesis/blob/ba2a065656ec7c638baa601994de3e0245ff6665/src/Development/AutomorphicArithmetic/DivisorFrames/ProductFormulaHeight.lean#L32}{The product formula for the global height in the source.}

The freedom to let the abelian modulus grow poses a different issue. Given a scale $P$, choose a prime $M>4P$ and let $S$ be the reduced classes represented by primes in $(P,2P]$. These representatives are distinct and $|S|\le P+1$. For large $P$ their complement has size at least $|S|$. Define $u$ to be $1$ on $S$, $-|S|/(\varphi(M)-|S|)$ on the other reduced classes. It has mean zero on the residue group and $|u|\le1$, but
\[
       \sum_{P<p\le2P}u(p)\log p=\sum_{P<p\le2P}\log p\sim P.
\]
Thus a square-root estimate for every varying bounded mean-zero function of a residue class would be false. Interpolating $t_p$ by an arbitrary array of congruence values on a finite list of primes does not preserve the uniformity that is required in Theorem~\ref{geo:dyadic}.

\clearpage
\section{What the retained character family sees at prime powers}\label{pp:section}

The prime sums in the geometric argument are identified with centered sums for primitive Dirichlet characters at every integer. In particular the local coefficient at $p^k$, away from the conductor, is $\chi(p)^k$. The retained terms have positive coefficients. Both features are explicit in the source. See \href{https://github.com/Deicyde/riemann-hypothesis/blob/ba2a065656ec7c638baa601994de3e0245ff6665/src/Development/AutomorphicArithmetic/PrimitiveCosets/PrimRead.lean#L48}{the prime-sum identification} and \href{https://github.com/Deicyde/riemann-hypothesis/blob/ba2a065656ec7c638baa601994de3e0245ff6665/src/Formalism/SharpRestriction/RamificationFrobenius/TotientGaloisClasses/ExplicitFormulaErrorOfHostConditions.lean#L32}{the weights and the cyclic level, lines 32--38 and 103--126}. We test this identity after averaging over $p$ separately for each $k$. In this way a test with four exponents distinguishes the elliptic symmetric square from every positive retained family, independently of the character conductors.

For any bounded prime function $b$ write
\[
             M(b)=\lim_{z\to\infty}\frac{1}{z}\sum_{p\le z}b(p)\log p
\]
when the limit exists. A fixed finite set of bad primes can be omitted. For the uncentered cyclic family of order $r$ the mean coefficient at exponent $k$ is
\begin{equation}\label{pp:cyclic}
        \frac{3r+\gcd(r,k)}{3r+1}.
\end{equation}
Indeed the family is $3r$ copies of the trivial character and all $r$ cyclic characters, divided by $3r+1$. The prime number theorem in fixed progressions gives $M(\chi^k)=1$ if the order of $\chi$ divides $k$, and zero otherwise. There are $\gcd(r,k)$ such cyclic characters. The same calculation for a finite, or absolutely summable, list of real weights $c_\chi$ gives
\begin{equation}\label{pp:divsum}
 A(k)=M\left(\sum_\chi c_\chi\operatorname{Re}(\chi(p)^k)\right)
     =\sum_{d\mid k}w_d,\qquad w_d=\sum_{\operatorname{ord}\chi=d}c_\chi .
\end{equation}
For an infinite list, the prime number theorem bounds the mean of a tail by a constant times its total variation and justifies passage to the limit. Each character may use its own conductor.

\begin{theorem}\label{pp:obstruction}
Let $E/\Q$ be non-CM, $m\ge1$, and let $U_p\in\mathrm{SU}(2)$ denote its unitary good-prime parameter. The average coefficient of the logarithmic derivative with minus sign of its $m$th symmetric power at $p^k$ is
\[
 B_m(k)=M\bigl(\operatorname{tr}((\operatorname{Sym}^m U_p)^k)\bigr)
  =\begin{cases}
       -\mathbf1_{k=2}, &m\ \text{odd},\\
       1-\mathbf1_{k=1}, &m\ \text{even}.
    \end{cases}
\]
For any $C\in\R$, $\eta>0$, no positive Dirichlet family can realize the mean profile $C+\eta B_m(k)$. For odd $m$, the error is at least $\eta/2$ on $k\in\{1,2\}$. For even $m$, it is at least $\eta/4$ on $k\in\{1,2,3,6\}$. If $A(1)=C$ exactly, these two lower bounds improve to $\eta$ and $\eta/3$, respectively.
\end{theorem}

\begin{proof}
Write the eigenvalues of $U_p$ as $e^{\pm i\theta_p}$. The Sato--Tate theorem gives density $2\sin^2\theta\,d\theta/\pi$ on $[0,\pi]$. Partial summation with the prime number theorem gives the same law for $M$. For integers $j\ge1$,
\[
 \frac2\pi\int_0^\pi 2\cos(j\theta)\sin^2\theta\,d\theta
       =-\mathbf1_{j=2}.
\]
This follows by expanding $2\sin^2\theta=1-\cos(2\theta)$. The weights of the symmetric power are $m,m-2,\ldots,-m$. For odd $m$ the list has a pair of weights $\pm1$, and its only contribution is $-1$ at $k=2$. For even $m\ge2$ there is a zero weight and a pair $\pm2$; the latter pair contributes $-1$ just at $k=1$. This computes $B_m$. See \href{https://doi.org/10.2977/PRIMS/31}{Barnet-Lamb, Geraghty, Harris and Taylor} for the Sato--Tate input.

For positive coefficients in \eqref{pp:divsum}, one has
\[
 A(2)-A(1)=w_2\ge0,\qquad
 A(6)-A(3)-A(2)+A(1)=w_6\ge0 .
\]
The corresponding differences of the target are $-\eta$ for odd $m$ and $-\eta$ for even $m$, respectively. Bounding the two or four errors by their maximum proves the claims; if the first mean is prescribed it contributes no error.
\end{proof}

For a concrete example take $E:y^2+y=x^3-x^2-10x-20$. Its $j$-invariant is $-496^3/11^5$, which excludes CM because a CM $j$-invariant is an algebraic integer. The theorem applies to its weight-two newform ($m=1$) and its symmetric square ($m=2$). For either target one may set $C=1$, $0<\eta\le1/(m+1)$, so that the target local weight $1+\eta\operatorname{tr}((\operatorname{Sym}^m U_p)^k)$ is itself nonnegative at every good prime and power. The contradiction concerns its representation by positive character weights.

\subsection{The exact cost of using signed characters}
Let $K\ge2$ be an integer. Among finite signed families satisfying $A(k)=B_m(k)$ for $1\le k\le K$, the least total variation $\sum_\chi |c_\chi|$ is
\begin{equation}\label{pp:exactcost}
  \begin{cases}
     \displaystyle\sum_{j\le K/2}|\mu(j)|,&m\ \text{odd},\\[3pt]
     \displaystyle\sum_{2\le j\le K}|\mu(j)|,&m\ \text{even}.
  \end{cases}
\end{equation}
Indeed Möbius inversion uniquely fixes $w_d$ for $d\le K$. For odd $m$ it gives $w_{2j}=-\mu(j)$ and zero at odd orders. For even $m$ it gives $w_1=0$ and $w_d=-\mu(d)$ for $d>1$. Summing the absolute values is a lower bound by the triangle inequality. Every order $d$ occurs for some Dirichlet character, as one sees by taking a prime $q\equiv1\pmod d$ and a character of order $d$ on its cyclic group of units. Assigning $w_d$ to one such character gives equality. Conjugate pairs with half the assigned weight give the same real prime means. The total costs are $3K/\pi^2+O(\sqrt K)$ and $6K/\pi^2+O(\sqrt K)$, respectively, by summing $\mu^2(n)=\sum_{h^2\mid n}\mu(h)$.

These costs remain unbounded under any fixed uniform error $e<1/2$. For even $m$, every prime $\ell\le K$ has $|w_\ell|=|A(\ell)-A(1)|\ge1-2e$. For odd $m$, every odd prime $\ell\le K/2$ has $|w_\ell+w_{2\ell}|=|A(2\ell)-A(2)|\ge1-2e$. The pairs of indices are disjoint. Consequently no signed family of bounded total variation approximates the entire profile with error less than $1/2$. This threshold is exact: constants $-1/2$ and $1/2$ realize error $1/2$ in the two cases. Exact matching here concerns only the prime means; a family attaining \eqref{pp:exactcost} need not match even a single prescribed local Euler factor.

\subsection{What happens to the original smoothing}
The geometric argument uses the ramp $(x/p^k-1)_+$. This permits a direct calculation of the scale of the discrepancy. If $k\ge2$, $v(p)$ is bounded and $M(v)=b$, then
\begin{align*}
 &\sum_{p^k\le x}v(p)\log p\,(x/p^k-1)
        -x\sum_p v(p)\log p\,p^{-k}
       =-\frac{k}{k-1}b\,x^{1/k}+o(x^{1/k}).
\end{align*}
Put $y=x^{1/k}$ and $V(y)=\sum_{p\le y}v(p)\log p=by+o(y)$. The left side is $-V(y)-x\sum_{p>y}v(p)\log p/p^k$. Partial summation makes it exactly $-kx\int_y^\infty V(t)t^{-k-1}dt$, proving the formula. The subtracted linear term converges absolutely. In particular the centered $k=2$ sector contributes $+2\sqrt x+o(\sqrt x)$ for the standard curve and $-2\sqrt x+o(\sqrt x)$ for its symmetric square. This calculation retains the weight of the actual prime sum in \href{https://www.math.uwaterloo.ca/~mrubinst/alpoge_claude_grh/zeta-paper.pdf}{Rubinstein's reconstruction, (I.2)}. The discrepancy at this scale fits within $O_\epsilon(x^{1/2+\epsilon})$. Failure of the coefficient identity leaves open a method that estimates the higher powers separately and proves a new bound on the first primes.

\clearpage
\section{The sharp character cost of recovering an elliptic trace}\label{fc:section}
An elliptic curve with surjective mod-$\ell$ representation supplies a full group $G=\mathrm{GL}_2(\F_\ell)$. On dyadic intervals one can recover $a_p$ by taking the centered representative of its trace modulo a single sufficiently large $\ell$. Section~\ref{geo:bridge} shows why bounding the irreducible modes separately incurs a cost. We now determine its order and show that changing the recovering function at unused traces cannot remove the loss. The second assertion allows the modification to depend on the whole conjugacy class.

For an odd prime $q$ write $[t]_q$ for the integer representative of $t$ between $-(q-1)/2$ and $(q-1)/2$, and $H(t)=[t]_q/q$. For a conjugacy-invariant complex function $h$ on a finite group $G$, put
\[
 c_\rho=\frac1{|G|}\sum_{g\in G}h(g)\overline{\operatorname{tr}\rho(g)},\qquad
 \mathcal C_G(h)=\sum_{\rho\in\widehat G}(\dim\rho)|c_\rho|.
\]
Thus any family of separate prime-sum estimates $B\dim\rho$ bounds the recombined observable by $B\mathcal C_G(h)$ using the triangle inequality. This is the actual norm of that procedure; a lower bound on it is not a lower bound on a signed prime sum.

\begin{theorem}\label{fc:main}
For $q\ge5$ and $h(A)=H(\operatorname{tr}A)$ on $G=\mathrm{GL}_2(\F_q)$,
\[
 \frac{(q+1)(q-3)}{4(q-1)\sin(\pi/q)}
 \ \le \mathcal C_G(h) < \frac{q^2}{2}.
\]
The lower bound already holds on irreducible principal series with two nontrivial distinct characters. In particular the cost has order $q^2$.
\end{theorem}

\begin{proof}
Let $B$ be the subgroup of upper triangular matrices. For characters $\chi,\psi$ of $\F_q^\times$ consider
$\rho_{\chi,\psi}=\operatorname{Ind}_B^G(\chi\boxtimes\psi)$, defined on $B$ by $\chi(x)\psi(y)$ on diagonal entries $x,y$. For $\chi\ne\psi$ it is irreducible of dimension $q+1$, and its isomorphism class determines the unordered pair. Indeed $G/B$ is the projective line. The two $B$-orbits there, or equivalently the two Bruhat double cosets, give by Mackey's formula the identity and the transposition as the only possible intertwining terms. Each contributes one exactly when the ordered diagonal characters agree after the indicated permutation.

Frobenius reciprocity for class functions gives
\[
 c_{\rho_{\chi,\psi}}=\frac1{(q-1)^2}
        \sum_{x,y\ne0} H(x+y)\overline{\chi(x)\psi(y)}.
\]
Set $C_\eta=(q-1)^{-1}\sum_{u\ne0}H(u)\overline{\eta(u)}$. For $\eta=\chi\psi$ odd, $H(0)=0$ removes the zero-sum terms, and substitution $x=uv$, $y=u(1-v)$ gives
\[
 c_{\rho_{\chi,\psi}}=\frac{J(\overline\chi,\overline\psi)}{q-1}C_\eta.
\]
Here a multiplicative character is extended by zero at zero when used in a Jacobi sum. If both characters and their product are nontrivial, the Jacobi sum has absolute value $\sqrt q$. To see it, multiply the two Gauss sums for the additive character $e_q(v)=\exp(2\pi i v/q)$ and group by $u=x+y$. The terms at zero vanish, giving $J(\chi,\psi)\tau(\chi\psi)=\tau(\chi)\tau(\psi)$. The Gauss sum of a nontrivial character has modulus $\sqrt q$ by additive orthogonality. For each odd $\eta$ there are precisely $(q-3)/2$ unordered pairs under consideration: choose $\chi$ other than $1,\eta$ and let $\psi=\eta/\chi$. A square character is even, precluding $\psi=\chi$.

It remains to bound the multiplicative Fourier norm of the centered trace. Since $H$ is odd, only odd $\eta$ contribute. From $H=\sum C_\eta\eta$ on $\F_q^\times$ we get
\[
 \left|\sum_{v\in\F_q} H(v)e_q(v)\right|
       \le\sqrt q\sum_{\eta\ {\rm odd}}|C_\eta|.
\]
Writing $b=(q-1)/2$, differentiation of
$\sum_{j=-b}^{b}e^{ijz}=\sin(qz/2)/\sin(z/2)$ at $z=2\pi/q$ evaluates the left side as $1/(2\sin(\pi/q))$. Multiply by $(q+1)(q-3)\sqrt q/(2(q-1))$, the dimension and multiplicity from the preceding paragraph, to obtain the lower bound. Finally Parseval for class functions and $\sum_\rho (\dim\rho)^2=|G|$ give $\mathcal C_G(h)\le |G|^{1/2}\|h\|_2<q^2/2$, since $|H|<1/2$ and $|G|<q^4$.
\end{proof}

\begin{theorem}[Recovery on part of the trace range]\label{fc:robust}
Fix $0<a<1/2$. There are constants $b_a,\epsilon_a>0$ such that the following holds for all sufficiently large primes $q$. Let $g$ be any class function on $\mathrm{GL}_2(\F_q)$. Average over $x,y\in\F_q^\times$ and $z\in\F_q$ independently and uniformly, and put $A=\left(\begin{smallmatrix}x&z\\0&y\end{smallmatrix}\right)$. If
\[
 \mathbb E\!\left[\mathbf1_{|[x+y]_q|\le aq}
              |g(A)-H(x+y)|^2\right]\le\epsilon_a^2 ,
\]
then $\mathcal C_G(g)\ge b_a q^2$. This includes exact agreement on the indicated slice, with arbitrary values outside it.
\end{theorem}
\begin{proof}
Choose a real odd smooth function $w$ on $\R/\Z$ supported in $(-a,a)$, with $c=\int_{-a}^a u w(u)\,du>0$. For example take $u$ times a nonnegative even smooth bump. Write $w_q(t)=w([t]_q/q)$ and $K_w=\sum_{j\in\Z}|\widehat w(j)|<\infty$. For a nontrivial multiplicative $\eta$,
\[
 \left|\sum_{u\ne0}w_q(u)\eta(u)\right|\le K_w\sqrt q .
\]
Expand the absolutely convergent Fourier series. For frequencies divisible by $q$ the character sum vanishes; all other complete sums have modulus $\sqrt q$.

For $t\ne0$ there are exactly $q-2$ ordered pairs $x,y\ne0$ with $x+y=t$. Consequently
\[
 \mathbb E[H(x+y)w_q(x+y)]
  =\frac{q-2}{(q-1)^2}\sum_{t\ne0}H(t)w_q(t)\longrightarrow c
\]
by an ordinary Riemann sum. Cauchy--Schwarz and the support condition compare this expression with $A_g=\mathbb E[g(A)w_q(x+y)]$ to within $\epsilon_a\|w\|_\infty$. Choose $\epsilon_a$ so that $|A_g|\ge c/2$ for all sufficiently large $q$.

Expand $F(x,y)=q^{-1}\sum_z g(A)$ in the multiplicative characters of the two diagonal variables. The coefficient of $\chi(x)\psi(y)$ is $c_{\rho_{\chi,\psi}}$ by reciprocity, whether or not the induction is irreducible. Pair one term with $w_q(x+y)$. It vanishes unless $\eta=\chi\psi$ is odd, by simultaneous negation of $x,y$. In the remaining terms the $t=0$ contribution vanishes and the sum is
\[
 \frac{J(\chi,\psi)}{(q-1)^2}\sum_{t\ne0}w_q(t)\eta(t).
\]
Here $|J|\le\sqrt q$, including the case where exactly one character is trivial (the Jacobi sum then has absolute value one). Since $\eta$ is odd one necessarily has $\chi\ne\psi$. Summing and counting the two possible orderings gives
\[
 |A_g|\le\frac{2K_w q}{(q-1)^2(q+1)}\mathcal C_G(g).
\]
Together with $|A_g|\ge c/2$ this yields the result. Only the upper triangular slice of the group was used; zero average over determinants alone could not give the same conclusion.
\end{proof}

For the mod-$\ell$ realization in Section~\ref{geo:bridge}, choose $\ell$ a fixed sufficiently large multiple in size of $\sqrt P$, so that the trace of a good Frobenius for $P<p\le2P$ belongs to the prescribed slice. Theorem~\ref{fc:robust} concerns universal recovery on matrices compatible with that trace range. It does not require that actual primes in the interval visit those matrices uniformly. In particular, any argument using the unique centered recovery on the full group pays $\asymp\ell^2\asymp P$, even before bounding individual prime sums. Starting with a square-root bound per representation and then summing in absolute value would produce an allowance of order $P^{3/2}$, up to logarithms. The low-cost sine expansion in Section~\ref{geo:bridge} avoids that norm, but exchanges it for the difficult cancellation in additive trace phases.

\subsection*{An arbitrary-rank version}
The same computation explains the scale of the obstruction for the full defining trace in higher rank. This is a statement about finite groups. Applying it to a particular higher-rank motive requires a large-image theorem and the appropriate normalizations.

\begin{theorem}\label{fc:high}
For each fixed $d\ge2$ and all sufficiently large odd primes $q$, the degree-weighted character norm of $g\mapsto H(\operatorname{tr}g)$ on $\mathrm{GL}_d(\F_q)$ is $\asymp_d q^{d^2/2}$. The lower bound holds on principal series induced from tuples of pairwise distinct nontrivial characters.
\end{theorem}
\begin{proof}
Let $B$ again be the upper triangular subgroup. For a tuple of pairwise distinct characters, induction from the diagonal extended trivially to the unipotent subgroup gives an irreducible representation. Two such inductions agree precisely up to permutation. This follows from the Bruhat decomposition indexed by permutations and Mackey's formula: the intersection in each double coset contains the diagonal torus, and agreement of the characters there is necessary and sufficient for a contribution. The dimension $[G:B]$ is the number of complete flags, $\prod_{i=1}^d\big((q^i-1)/(q-1)\big)$, hence $\asymp_d q^{d(d-1)/2}$.

Fix any odd character $\eta$. There are $\gg_d(q-1)^{d-1}$ ordered tuples of nontrivial pairwise distinct characters with product $\eta$. Choose the first $d-1$ freely. Requiring one specified entry to be trivial costs at most $(q-1)^{d-2}$ tuples. A specified coincidence costs at most $2(q-1)^{d-2}$, by choosing the remaining characters and solving an equation of degree at most two in the cyclic character group. This argument includes a coincidence with the last entry, which is determined by the product. Omit all bad choices by the union bound, then divide by $d!$ to forget the order.

The same reciprocity calculation now has $d$ diagonal variables. The generalized Jacobi sum for $x_1+\cdots+x_d=1$ equals $\prod_i\tau(\overline{\chi_i})/\tau(\overline\eta)$ by grouping the product of Gauss sums according to their additive total. Since $\eta$ is nontrivial the terms at additive total zero vanish by simultaneous scaling. Thus the modulus of the Fourier coefficient of the induced representation is
\[
      \frac{q^{(d-1)/2}}{(q-1)^{d-1}}|C_\eta| .
\]
Multiply by the number of tuples and the flag dimension, then sum over odd $\eta$. The first additive Fourier coefficient in Theorem~\ref{fc:main} gives $\sum|C_\eta|\gg\sqrt q$. The resulting lower exponent is $d(d-1)/2+(d-1)/2+1/2=d^2/2$. Parseval and $|\mathrm{GL}_d(\F_q)|<q^{d^2}$ give the upper bound.
\end{proof}

\subsection*{Why many small moduli do not give uniform decoding for free}
The preceding norm bound concerns one large residual group. Could a sum of many observables at much smaller moduli uniformly recover the same signed trace? There is an obstruction even when one allows arbitrary periodic functions. It is an assertion about deterministic reconstruction of all admissible traces, separate from the distribution of actual primes.

\begin{proposition}\label{fc:small}
Fix $D>0$, $0<\delta<1$ and $0<a<1/2$. For large $R$ consider any finite sum $F_R(n)=\sum_{j} f_j(n)$, where $f_j$ is periodic of integer period $m_j\le R^{1-\delta}$. Suppose the sum over $j$ of the sums of the absolute normalized additive Fourier coefficients of $f_j$ is $\le R^D$. Then
\[
 \liminf_{R\to\infty}\frac1R\sum_{|n|\le aR}|F_R(n)-n/R|^2\ge\frac{2a^3}{3}.
\]
The lower bound is attained by $F_R=0$.
\end{proposition}
\begin{proof}
Choose the same odd smooth compactly supported $w$ on the real interval $(-a,a)$, with $c=\int uw(u)\,du>0$. Its Fourier transform decays faster than every power. For frequency $b/m$ in lowest terms, $1\le b<m$, Poisson summation gives
\[
 \frac1R\left|\sum_{n\in\Z} w(n/R)e^{2\pi ibn/m}\right|
 \le \sum_{h\in\Z}|\widehat w(R(h-b/m))|
 \ll_{w,A}(1+R/m)^{-A}+R^{-A}
\]
for each $A>1$. The nearest integer has distance at least $1/m$ and all other distances can be summed by inverse powers. At frequency zero the sum vanishes exactly by oddness. Expanding each periodic function, then taking $A>(D+1)/\delta$, proves $R^{-1}\sum F_R(n)w(n/R)=o(1)$. Meanwhile $R^{-1}\sum (n/R)w(n/R)\to c$. By Cauchy--Schwarz the liminf in the assertion is at least $c^2/\int|w(u)|^2\,du$. Choose smooth odd cutoffs tending to $u$ in $L^2(-a,a)$. For each cutoff the limit in $R$ is taken first. The ratio on the right tends to $\int_{-a}^a u^2\,du=2a^3/3$. For $F_R=0$ the Riemann sum attains this same limit.
\end{proof}

For example fix a putative determinant $d$ invertible at the probe primes and feed each probe the companion matrix $\left(\begin{smallmatrix}0&-d\\1&n\end{smallmatrix}\right)$. Any finite-group class function evaluated on this matrix modulo $m$ is periodic in $n$ with period $m$. A polynomial number of uniformly polynomially bounded such probes automatically has polynomial total Fourier norm, by finite Parseval and Cauchy--Schwarz. Consequently their sum cannot uniformly decode even the centered trace on a fixed subinterval when all moduli are $\le R^{1-\delta}$. Products of probes may have much larger periods. Signed cancellation at the actual Frobenius classes could still give a prime-sum bound without uniform pointwise decoding. These are exactly the alternatives left open by the proposition.

\clearpage
\section{Approximating a standard function by vertical shifts of zeta}\label{vt:section}

A vertical shift can reproduce any prescribed finite set of unit Frobenius phases at distinct rational primes. What happens to the shifts beyond the region of absolute convergence determines whether this observation transfers a zero theorem. We give a necessary and sufficient condition for the transfer, a lower bound on the growth forced by a hypothetical off-line zero, and an estimate separating the matched factors from the remaining primes. Bagchi's recurrence theorem links shifts of $\zeta$ approximating $\zeta$ itself to its zero location. Here the matched phases come from a different target, and the paired shifts must also agree with its Euler factors in the absolute-convergence region. \href{https://www.kurims.kyoto-u.ac.jp/~kenkyubu/bessatsu/open/B34/pdf/B34_017.pdf}{Survey of Bagchi's recurrence theorem.}

Let $S$ be a finite set of rational primes. Work first with a unitary standard function $L$ whose incomplete Euler product is
\[
 L_S(s)=\prod_{p\notin S}[(1-\alpha_p p^{-s})(1-\bar\alpha_p p^{-s})]^{-1},
 \qquad |\alpha_p|=1,\qquad \operatorname{Re}s>1.
\]
Assume $L_S$ is holomorphic on $\mathcal H=\{s:\operatorname{Re}s>1/2\}$ and $L$ has the standard analytic properties used in Theorems~\ref{aw:criterion} and~\ref{aw:raw}. For an elliptic curve over $\Q$ take $S$ to contain the bad primes and $L(s)=L(E,s+1/2)$. The omitted factors and their inverses are nonzero on $\mathcal H$. Thus zero-freeness of $L_S$ on $\mathcal H$ is precisely the assertion about nontrivial zeros, by the functional equation.

For $X\ge\max(S\cup\{2\})$, write
\begin{align*}
 P_X(s)&=\prod_{p\le X,\,p\notin S}[(1-\alpha_p p^{-s})(1-\bar\alpha_p p^{-s})]^{-1},\\
 Z_S(s)&=\zeta(s)\prod_{p\in S}(1-p^{-s}),\qquad
 D_t(s)=Z_S(s+it)Z_S(s-it).
\end{align*}
A shift $t$ matches $X$ to precision $h$ if $|p^{-it}-\alpha_p|<h$ for all $p\le X$ outside $S$. Unique factorization makes the finitely many $\log p$ rationally linearly independent. The continuous Kronecker--Weyl theorem gives a positive asymptotic proportion of such $t$ in $[T,2T]$ as $T\to\infty$ for each fixed $X,h$. If $0<h<2$, this proportion is $(2\arcsin(h/2)/\pi)^{\#\{p\le X:p\notin S\}}$. It can be very small.

\begin{theorem}[A normal-family criterion]\label{vt:normal}
Suppose $\zeta$ has no zeros in $\mathcal H$, with its pole at $1$ excepted. Then $L_S$ has no zeros in $\mathcal H$ if and only if there exist $X_j\to\infty$, $h_j\to0$, and matched shifts $t_j\to\infty$ such that $D_{t_j}$ are eventually bounded on each compact subset of $\mathcal H$, with a bound depending on the compact set. In this event the shifts can be chosen so that $D_{t_j}\to L_S$ locally uniformly. Any prescribed positive upper bounds for the mesh $h_j$ can be respected.
\end{theorem}

The poles $1\pm it_j$ eventually leave each fixed compact set, so convergence here is convergence of functions eventually holomorphic on a neighborhood of that compact. The proof gives a positive-density selection for every finite step and estimates on the Euler tails, without a bound on the first acceptable height.

\subsection{A rectangle forced by a hypothetical zero}

For every $a>0$ the Euler products give
\begin{equation}
 \sup_{\operatorname{Re}s\ge1+a}|D_t(s)-L_S(s)|\ll_a h+X^{-a}.
 \label{vt:right}
\end{equation}
For unit $u,v$ one has $|u^k-v^k|\le k|u-v|$. The matched Euler logarithms therefore differ in total by at most $2h\sum_p p^{-1-a}/(1-p^{-1-a})$. The logarithms beyond $X$ cost at most $4\sum_{p>X}p^{-1-a}/(1-p^{-1-a})\ll_a X^{-a}$. Use the logarithms given by the absolutely convergent expansions. Exponentiation yields the claim since the real parts of these logarithms are uniformly bounded above on this half-plane. The same estimate with $h=0$ holds for $P_X-L_S$.

Suppose $\rho=\beta+i\gamma$ is a zero of $L_S$ with $1/2<\beta<1$. Choose $1/2<l<\beta$, $a>0$, and the closed rectangle
\[
 R=\{s:l\le\operatorname{Re}s\le1+a,\quad|\operatorname{Im}s-\gamma|\le H\}.
\]
Take a circle $C$ of radius $r$ about $\rho$ whose closed disk is in the interior of $R$ and whose boundary avoids zeros. Such $r$ exist since $L_S$ is holomorphic and nonzero on $\operatorname{Re}s>1$. Put $m=\min_C|L_S|>0$, $A=\max_R|L_S|$ and
\[
 \lambda=\frac{\pi}{2H},\qquad
 \omega=\frac{\sinh(\lambda(\beta-r-l))}{\sinh(\lambda(1+a-l))}
                      \cos(\lambda r),\qquad
 \kappa=\frac{a\omega}{1-\omega}.
\]
Then $0<\omega<1$.

\begin{theorem}[Quantitative obstruction]\label{vt:obstruction}
For all sufficiently large $X$, every shift matching to $h\le X^{-a}$ whose two poles are outside $R$ satisfies
\[
 \max_{s\in R}|D_t(s)|\ge c_R X^\kappa-A
\]
provided $D_t$ is zero-free on $R$. Under the zeta hypothesis in Theorem~\ref{vt:normal}, this applies to every such shift. Unconditionally $\max_R|P_X|\ge c_R X^\kappa-A$ for every sufficiently large $X$. Given any $0<k<\beta-1/2$, the rectangle and $a$ can be chosen so that $\kappa>k$.
\end{theorem}
\begin{proof}
Let $D$ stand for $D_t$ or $P_X$. Rouch\'e's theorem gives some $s\in C$ with $|D(s)-L_S(s)|\ge m$, since $D$ has no zero in the disk and $L_S$ does. Let $e=C_aX^{-a}$ bound the difference on the right edge by \eqref{vt:right}. Set $B=\max(1,e,A+\max_R|D|)$. The harmonic function
\[
 \phi(\sigma+iv)=\frac{\sinh(\lambda(\sigma-l))}{\sinh(\lambda(1+a-l))}
                           \cos(\lambda(v-\gamma))
\]
vanishes on the left, bottom and top edges, and lies in $[0,1]$ on the right. The subharmonic maximum principle gives throughout $R$
\[
 \log|D-L_S|\le(1-\phi)\log B+\phi\log e
\]
since the right side dominates on all four edges. On $C$, $\phi\ge\omega$, so $m\le B^{1-\omega}e^\omega$. Consequently $B\ge m^{1/(1-\omega)}C_a^{-\omega/(1-\omega)}X^\kappa$. For large $X$ this exceeds $1$ and $e$. Finally choose $l$ close to $1/2$ and $r$ small so that $\beta-r-l>k$. For fixed $a$, as $H\to\infty$ the exponent tends to $a(\beta-r-l)/(1+a-\beta+r)$. As $a\to\infty$ this tends to $\beta-r-l$. Choose a finite $a$ and then a finite $H$ to obtain $\kappa>k$.
\end{proof}

There is also a consequence for the heights of matching shifts. Under the starting zeta hypothesis, the classical implication from RH to the Lindel\"of bound gives $\max_R|D_t|\ll_{R,S,\nu}(1+|t|)^\nu$ for each $\nu>0$ when $|t|>2(|\gamma|+H+2)$. Indeed the bound for $\zeta$ is uniform on a fixed strip with real part bounded strictly above $1/2$, and the omitted local factors are bounded there. Combining with Theorem~\ref{vt:obstruction}, if the indicated elliptic zero exists then for every $B>0$ all sufficiently large $X$ require $|t|>X^B$ for every such shift matching to $X^{-a}$. Choose $\nu<\kappa/B$ to see this. In the cuspidal case all sufficiently accurate matching heights in fact leave every bounded interval. Otherwise a bounded matching sequence at growing cutoffs would have a limit $t_*$ with $p^{-it_*}=\alpha_p$ at every $p\notin S$. This would give $L_S=D_{t_*}$ in $\operatorname{Re}s>1$ and then by meromorphic continuation throughout $\mathcal H$. But $D_{t_*}$ has a pole at $1+it_*$: for $t_*\ne0$ its other factor is nonzero there by zeta nonvanishing on $\operatorname{Re}s=1$, and for $t_*=0$ the pole is double. This contradicts holomorphy of $L_S$. Thus the stated height consequence also excludes the bounded-height range, where proximity to an edge pole would otherwise prevent a uniform Lindel\"of estimate. This is a necessary height growth under a hypothetical zero. Kronecker equidistribution at a fixed finite set of primes gives no polynomial upper bound as the set grows. \href{https://sites.math.rutgers.edu/~zeilberg/EM18/TitchmarshZeta.pdf}{Titchmarsh, Chapter~XIV.}

For the qualitative direction of Theorem~\ref{vt:normal}, eventual local boundedness gives a locally uniformly convergent subsequence on $\mathcal H$ by Montel's theorem and a compact exhaustion. By \eqref{vt:right} its limit agrees with $L_S$ on $\operatorname{Re}s>1$, hence throughout $\mathcal H$. The sequence is eventually holomorphic and zero-free near each compact set, so Hurwitz's theorem makes the limit identically zero or zero-free. The Euler product excludes the first possibility. Every locally uniformly convergent subsequence has the same limit, which proves convergence of the given locally bounded sequence. Without the hypothesis on $\zeta$, this implication still applies to any matching sequence whose products are eventually zero-free near each compact set.

\subsection{A mean square after prescribing the prime phases}

Put $Q_X=\prod_{p\le X}p$ and $Z_{>X}(s)=\zeta(s)\prod_{p\le X}(1-p^{-s})$. Let $\mathcal B$ be a product of open arcs of positive length on the circles of some of the primes $p\le X$. Write $1_{\mathcal B}(t)$ for the event that $(p^{-it})$ lies in this box and $\mu$ for its normalized volume. Its boundary has Haar measure zero. The box may be arbitrarily small for a fixed $X$.

\begin{lemma}[Rough integers after conditioning]\label{vt:conditional}
For $s=\sigma+iv$ with $\sigma>1/2$ one has unconditionally
\begin{equation}
 \lim_{T\to\infty}\frac1{\mu T}\int_T^{2T}1_{\mathcal B}(t)
               |Z_{>X}(s+it)-1|^2\,dt
 =\sum_{\substack{n>1\\(n,Q_X)=1}}n^{-2\sigma}.
 \label{vt:meansquare}
\end{equation}
The equality persists after area integration over any fixed bounded rectangle whose closure lies in $\mathcal H$, including one crossing $\operatorname{Re}s=1$.
\end{lemma}
\begin{proof}
Fix a compact rectangle with $\operatorname{Re}s\ge\sigma_*>1/2$, decreasing $\sigma_*$ if necessary to arrange $\sigma_*<1$. For $V_N(s)=\sum_{n\le N}n^{-s}$ we need the Dirichlet approximation
\begin{equation}
 \limsup_{T\to\infty}\sup_s\frac1T\int_T^{2T}
      |\zeta(s+it)-V_N(s+it)|^2dt\ll_{\sigma_*}N^{1-2\sigma_*},
 \label{vt:zetaapprox}
\end{equation}
where the supremum is over the rectangle. Take $M=\lfloor T^{1+u}\rfloor$ with $u>0$ small enough that $(1+u)(2-2\sigma_*)<1$. Euler--Maclaurin summation at the integer $M$ gives $\zeta(w)=V_M(w)+o(1)$ uniformly for $w=s+it$, $T\le t\le2T$. Its integral term $M^{1-w}/(w-1)$ is $O(M^{1-\sigma_*}/T)=o(1)$. Its endpoint derivative terms of order $j\ge0$ are $O_j(M^{-\sigma_*}(3T/M)^j)$. After an even number $k$ of derivatives the error is $O_{k,\sigma_*}(T^kM^{1-\sigma_*-k})$, by the bounded periodic Bernoulli function and integration of $y^{-\sigma_*-k}$. Choose a fixed $k$ large enough that $(1+u)(1-\sigma_*)-ku<0$. Euler--Maclaurin first applies in the domain of absolute convergence and its integral error continues the formula to this range. This proves the uniform approximation on $[T,2T]$, an interval eventually avoiding the possible pole even at $\sigma=1$.

In the integrated square of $V_M-V_N$ the normalized diagonal is at most $\sum_{n>N}n^{-2\sigma_*}\ll_{\sigma_*} N^{1-2\sigma_*}$. For $m<n$, the absolute off-diagonal error is at most $2(mn)^{-\sigma_*}/(T\log(n/m))$ per ordered pair. For $m<n/2$ use $\log(n/m)\ge\log2$. For $n/2\le m<n$, use $\log(n/m)\ge(n-m)/n$. The total is $O_{\sigma_*}(M^{2-2\sigma_*}\log(2M)/T)=o(1)$. Combining this estimate with the uniform Euler--Maclaurin error proves \eqref{vt:zetaapprox}. This is the mean-value argument behind the classical zeta second moment. \href{https://sites.math.rutgers.edu/~zeilberg/EM18/TitchmarshZeta.pdf}{Titchmarsh, \emph{The Theory of the Riemann Zeta-Function}, Chapter~VII.}

Multiply $V_N$ by the fixed Dirichlet polynomial $\prod_{p\le X}(1-p^{-s})$. The coefficients through $N$ are $\mathbf1_{(n,Q_X)=1}$. At larger integers the coefficients are bounded by the divisor count of $Q_X$ and vanish beyond $NQ_X$. The limiting mean square of the finite excess polynomial is its sum of squared coefficient magnitudes weighted by $n^{-2\sigma}$, by integration of distinct frequencies. This is $O_X(N^{1-2\sigma_*})$, uniformly over the rectangle. Since the fixed multiplier is bounded there along all vertical shifts, \eqref{vt:zetaapprox} shows that $W_{N,X}(s)=\sum_{n\le N,(n,Q_X)=1}n^{-s}$ converges to $Z_{>X}$ in the same limiting mean-square sense, for each fixed $X$.

Equidistribute the finitely many prime phases occurring in the box and the support of $W_{N,X}$. The nonconstant terms in the latter polynomial involve only primes greater than $X$. Integrating $|W_{N,X}-1|^2$ over these independent circle coordinates kills every cross term. The limiting weighted average over $t$, divided by $\mu$, equals $\sum_{1<n\le N,(n,Q_X)=1}n^{-2\sigma}$. The indicator is Riemann integrable since its boundary is null. For the limit $N\to\infty$, use the $L^2$ triangle inequality. Multiplying by $1_{\mathcal B}$ can only decrease the unnormalized mean-square error. Dividing the norm by $\sqrt\mu$ is harmless for fixed $X$ and fixed box. This proves \eqref{vt:meansquare}. On a bounded rectangle the approximation bounds are uniform and finite polynomial limits can be integrated by bounded convergence. This proves the area assertion. The limit in the shift height is taken for each fixed cutoff before the cutoff is increased.
\end{proof}

For a compact $K\Subset\mathcal H$, take a bounded open rectangle $U$ containing $K$ and its complex conjugate $K^*$, symmetric about the real axis, with $\operatorname{Re}s>\sigma_*>1/2$ on its closure. Let $d=\operatorname{dist}(K\cup K^*,\partial U)>0$. The right side of \eqref{vt:meansquare} is at most $\sum_{n>X}n^{-2\sigma_*}\ll_{\sigma_*}X^{1-2\sigma_*}$, as a rough integer larger than one exceeds $X$. In the box there are therefore arbitrarily large $t$ such that the integrated squared error on $U$ is at most three times this upper bound times its area $|U|$. Markov's inequality leaves a positive fraction of the box's shifts. Take $t$ sufficiently large that both shifts are holomorphic near the closure of $U$. The disk mean-value inequality for the squared modulus of a holomorphic function gives
\begin{equation}
 \max_{s\in K,\,\pm}|Z_{>X}(s\pm it)-1|
       \ll_{\sigma_*}\frac{|U|^{1/2}}{d}X^{1/2-\sigma_*}.
 \label{vt:tail}
\end{equation}
For the negative shift use $Z_{>X}(s-it)=\overline{Z_{>X}(\bar s+it)}$. A single $t$ thus works for both factors even though the two shifted values are dependent. The constants in \eqref{vt:tail} are independent of the width of the box, while the first admissible height can depend on the full box.

\begin{corollary}[Recovering a finite target product by averaging]\label{vt:meanproduct}
For fixed $X$, let $\mathcal B_h$ be the boxes of chordal width $h$ centered at the phases of $\alpha_p$, with normalized volumes $\mu_h$. Locally uniformly for $\operatorname{Re}s>1/2$ one has
\[
 P_X(s)=\frac1{Z_{>X}(2s)}
 \lim_{h\downarrow0}\lim_{T\to\infty}
   \frac1{\mu_hT}\int_T^{2T}1_{\mathcal B_h}(t)D_t(s)\,dt.
\]
This formula for each finite Euler product is unconditional.
\end{corollary}
\begin{proof}
On the finite phase torus let
\[
 V(\mathbf z;s)=\prod_{p\le X,\,p\notin S}
       [(1-z_pp^{-s})(1-\bar z_p p^{-s})]^{-1}.
\]
It is continuous and uniformly bounded on the product of this torus with any fixed compact subset of $\mathcal H$. In the factorization of $D_t$, replace $Z_{>X}(s+it)$ and $Z_{>X}(s-it)$ by the finite polynomials $W_{N,X}(s+it)$ and $W_{N,X}(s-it)$ of Lemma~\ref{vt:conditional}. Cauchy--Schwarz and the mean-square approximation bound the error in the mean of the product by a quantity tending to zero uniformly on each compact as $N\to\infty$ for fixed $X,h$. For the negative shift use conjugation with $\bar s$; the two factors need not be independent. At finite $N$ the high-prime circle average of their product retains precisely the terms with equal integers in the two polynomials. Thus, writing $dm$ for probability Haar measure on the finite torus, the inner limit equals
\[
 \left(\sum_{(n,Q_X)=1}n^{-2s}\right)
     \frac1{\mu_h}\int_{\mathcal B_h}V(\mathbf z;s)\,dm(\mathbf z).
\]
The series converges absolutely and equals $Z_{>X}(2s)$, nonzero throughout $\mathcal H$. For the finite polynomial equidistribution is uniform on compact $s$ sets by compactness and uniform continuity of $V$ there; one can first cover the compact by a finite net. The area-zero boundary does not affect any average. As $h\downarrow0$, the last average tends uniformly to $P_X(s)$ on these compacts. This proves the formula.
\end{proof}

For real $s>1/2$ the extra factor is exactly the conditional second moment of the unprescribed Euler tail. In fact the series on the right of \eqref{vt:meansquare} equals
\[
 \prod_{p>X}(1-p^{-2\sigma})^{-1}-1
   \ll_{\sigma_*}\frac{X^{1-2\sigma_*}}{\log(2X)}
   \qquad(\sigma\ge\sigma_*>1/2).
\]
To verify the upper bound, take logarithms of the product and apply Chebyshev's estimate for $\pi(2^jX)$ to the dyadic sum of $p^{-2\sigma_*}$ over $p>X$. The geometric decay $2^{j(1-2\sigma_*)}$ makes the sum $O_{\sigma_*}(X^{1-2\sigma_*}/\log(2X))$. Exponentiating gives the inequality for all $X\ge2$ after adjustment of the constant. This improves the error in \eqref{vt:tail} by a factor $(\log(2X))^{-1/2}$. The unweighted integer bound is convenient for an explicit search threshold.

Indeed, suppose the finitely many phases are computable, as they are for elliptic curves. For a finite cutoff given as an integer $X$, a rational rectangular compact $K$, and prescribed phase width $h>0$ with computable data, one can search for a matching rational height satisfying the uniform tail estimate. Choose the enclosing rectangle $U$ with rational edges and a rational $1/2<\sigma_*<\inf_U\operatorname{Re}s$. One may use any positive rational $d$ below the distance from $K\cup K^*$ to $\partial U$. On any disk of radius $d$ the disk mean inequality has constant $(\pi d^2)^{-1}$. Since $\sum_{n>X}n^{-2\sigma_*}\le X^{1-2\sigma_*}/(2\sigma_*-1)$, there are arbitrarily large admissible heights with
\[
 \sup_{s\in K,\,\pm}|Z_{>X}(s\pm it)-1|
  < 2\sqrt{\frac{|U|}{\pi d^2(2\sigma_*-1)}}\,
                  X^{1/2-\sigma_*}.
\]
Choose the strict slack by first requiring an area bound with constant $3$ instead of $4$. Both this uniform inequality and the phase requirements are open in $t$ away from the poles, so the successful heights include rationals. Dovetail rational heights above the prescribed starting point and finer rectangular covers of $K$. At each height certify the inequalities by interval evaluation of the finite factors and Euler--Maclaurin for $\zeta$, using the periodic Bernoulli bound in the proof of \eqref{vt:zetaapprox}. For fixed compact input and height away from the poles, the integral remainder tends to zero with $M$ for any fixed even $k>1$; the finite-sum enclosures converge after subdivision and precision refinement. A positive strict margin on the compact guarantees a successful finite check. The search therefore terminates at every finite stage. This construction bounds the analytic error explicitly and does not provide a practical height bound.

\subsection{Recovering the target under its zero hypothesis}

Under GRH for $L$, Theorem~\ref{aw:raw} and the bounds for discarded powers give $\sum_{p\le y,\,p\notin S}(\alpha_p+\bar\alpha_p)\ll_L\sqrt y\log(2y)$. With $\sigma_K=\min_{s\in K}\operatorname{Re}s$, partial summation gives
\[
 \sup_{s\in K}\left|\sum_{p>X,\,p\notin S}(\alpha_p+\bar\alpha_p)p^{-s}\right|
       \ll_{L,K}X^{1/2-\sigma_K}\log(2X).
\]
The higher powers in the Euler logarithm have absolute tail $O_K(\sum_{n>X}n^{-2\sigma_K})=O_K(X^{1-2\sigma_K})$. The logarithm series converges locally uniformly on $\mathcal H$ and agrees with the Euler logarithm on $\operatorname{Re}s>1$. By uniqueness of continuation its exponential is $L_S$. In particular
\begin{equation}
 \sup_{s\in K}|P_X(s)-L_S(s)|
            \ll_{L,K}X^{1/2-\sigma_K}\log(2X).
 \label{vt:deterministic}
\end{equation}

Even without target GRH, selected shifts approximate the finite $P_X$ relatively. The logarithm of the ratio between the matched finite part of $D_t$ and $P_X$ is bounded on $K$ in modulus by
\[
 2h\sum_{p\le X}\frac{p^{-\sigma_K}}{1-p^{-\sigma_K}}\ll_K hX.
\]
The remaining factors of $D_t$ are $Z_{>X}(s+it)$ and $Z_{>X}(s-it)$. By \eqref{vt:tail}, for $hX\le1$ and sufficiently large $X$ the selected shifts satisfy
\begin{equation}
 \sup_{s\in K}|D_t(s)/P_X(s)-1|
       \ll_{K,U} hX+X^{1/2-\sigma_*}.
 \label{vt:relative}
\end{equation}
This bound for the finite product is unconditional. Under target GRH, \eqref{vt:deterministic} bounds $P_X$ on $K$ and identifies its limit. Choose an increasing compact exhaustion $K_j$ of $\mathcal H$, take $X_j$ large enough for the tail terms in \eqref{vt:deterministic} and \eqref{vt:relative} on $K_j$ to be less than $1/j$ after allowing for the compact's bounds. Choose the box width below the prescribed tolerance and small enough for its error term to have the same bound. Select $t_j>\max(j,t_{j-1}+1)$ by \eqref{vt:tail}. This proves the converse of Theorem~\ref{vt:normal}, independently of the zero theorem for $\zeta$.

Local boundedness of $P_{X_j}$ itself along a sequence $X_j\to\infty$ also implies zero-freeness of $L_S$, by Montel and Hurwitz applied to these nowhere-zero finite products. This deterministic statement requires no premise on $\zeta$. The starting zero hypothesis permits arbitrary matched vertical shifts to be used in the forward implication. A proof of a new fixed-curve zero theorem by this route still needs local bounds on the products across the line $\operatorname{Re}s=1$. The conditional mean square controls the residual factors in \eqref{vt:relative} and leaves the potentially large $P_X$ intact.

One elementary target illustrates why the finite match and the tail mean alone cannot force boundedness. If every $\alpha_p=-1$ and $S$ is empty, the target in $\operatorname{Re}s>1$ is
\[
 \prod_p(1+p^{-s})^{-2}=\frac{\zeta(2s)^2}{\zeta(s)^2}.
\]
Under the starting RH this ratio is holomorphic on $\mathcal H$ with a double zero at $1$. On the real interval $(1/2,1)$ its finite products tend to zero, because $\sum_p p^{-\sigma}$ diverges there, although on $(1,\infty)$ their limits are nonzero. Local boundedness of any increasing subsequence on $\mathcal H$ would give a holomorphic limit vanishing on an interval and nonzero on another, contradicting the identity theorem. The conditional tail law and arbitrarily accurate finite matching still hold for this target. Its zero at $1$ is incompatible with the standard nonvanishing theorem for an entire cuspidal function. The example tests which part of the attempted transfer actually uses the target arithmetic.

\subsection{Finite-place tempered representations over number fields}

A single vertical translate over the field $F$ assigns the same phase to all ideals of the same norm. Using rational primes and several translates eliminates this compatibility restriction. Let $F$ have degree $d$, let $\pi$ be unitary cuspidal of rank $n\ge2$ over $F$, and suppose it is tempered at all good finite primes. Remove a finite set $S$ of rational primes containing those ramified in $F$ and those under the conductor of $\pi$. For $p\notin S$, let $f_{\mathfrak p}$ be the residue degree of $\mathfrak p\mid p$. For every unit inverse root $\alpha_j(\mathfrak p)$, choose any $f_{\mathfrak p}$-th root $\eta$ of it. The polynomial identity
\[
 1-\alpha_j(\mathfrak p)z^{f_{\mathfrak p}}
       =\prod_{\omega^{f_{\mathfrak p}}=1}(1-\omega\eta z)
\]
turns the Euler factor of $L_S(s,\pi)$ over $p$ into exactly $nd$ unit linear factors in $z=p^{-s}$, since $\sum_{\mathfrak p\mid p}f_{\mathfrak p}=d$. Order their inverse roots $\xi_{p,1},\ldots,\xi_{p,nd}$ arbitrarily for each $p$. This uses neither an automorphic induction theorem nor a Galois hypothesis on $F$.

\begin{theorem}[Vertical transfer over a number field]\label{vt:field}
Under the same zeta zero hypothesis, GRH for $L(s,\pi)$ holds if and only if one can find cutoffs tending to infinity and $nd$-tuples of shifts, eventually escaping to positive infinity coordinatewise, with the following properties. At each cutoff the $j$th shift matches the chosen rational-prime phases $\xi_{p,j}$ at all $p\le X$ outside $S$, with mesh tending to zero. The products $\prod_{j=1}^{nd} Z_S(s+it_j)$ are eventually locally bounded in $\mathcal H$. Under target GRH a choice converging locally uniformly to $L_S(s,\pi)$ exists. The obstruction of Theorem~\ref{vt:obstruction} holds with the same exponent and a constant allowed to depend on $nd$. For a modular elliptic curve over $F$ only $d$ shifts are required if each is used with its negative. The same elliptic conclusion holds for a product of two entire unitary Hecke functions as in Theorem~\ref{aw:elliptic}.
\end{theorem}
\begin{proof}
At a rational prime $p$ the sum $\sum_j\xi_{p,j}$ receives no contribution from any $\mathfrak p\mid p$ with $f_{\mathfrak p}>1$, since the sum of its $f_{\mathfrak p}$-th roots of unity is zero. Thus it equals the sum of the ordinary Satake traces over the degree-one prime ideals above $p$. Under GRH the complex, unweighted sum of these traces through $p\le y$ is $O_{\pi,F}(\sqrt y\log(2y))$. To verify the complex assertion, use the nonaveraged identity \eqref{aw:explicit} at consecutive shrinking cutoffs just as in Theorem~\ref{aw:raw}: the conjugate average there was used only by the reverse one-sided implication. Taking absolute values of the complex coefficients in the sharpening argument is justified by its bound for the absolutely convergent zero sum. Division by $\log(q^k)$ gives the complex log-free sum, and temperedness bounds the removed prime powers and residue degrees exceeding one by $O_{\pi,F}(\sqrt y)$. Finite exceptional rational primes change the sharp first-prime sum by a bounded amount. Alternatively one can apply the contour computation separately to its real and imaginary parts.

Consequently partial summation bounds the first-power logarithmic tail just as in \eqref{vt:deterministic}. Its higher powers have absolutely summable tails because $|\xi_{p,j}|=1$ and the degree $nd$ is fixed. This proves local uniform convergence of the partial Euler products under GRH. Without the zero assumption they still obey the quantitative approximation on the right half-plane, with constants multiplied by the degree as necessary. Apply Lemma~\ref{vt:conditional} to the phase box of each coordinate separately and select a height for that coordinate satisfying \eqref{vt:tail}. The heights in the different coordinates may be chosen successively and arbitrarily far out. The product of their residual factors tends uniformly to one on the chosen compacts, and the sum of the finite logarithmic matching errors is $O_{K,nd}(hX)$. The Montel, Hurwitz and rectangle arguments now apply to their product. Each omitted local factor is nonzero on $\mathcal H$ by \eqref{aw:lrs}. Conjugation and the functional equation identify zero-freeness there with GRH, also for non-self-dual $\pi$.

For elliptic curves the original good roots occur in conjugate pairs. Rationalizing both polynomials belonging to one such pair yields a multiset which splits into $f_{\mathfrak p}$ conjugate pairs, even when a real root occurs with multiplicity two. Thus at each good $p$ the $2d$ phases can be recorded as $d$ pairs. Use one matched height and its negative for each pair, with the symmetric compact selection in \eqref{vt:tail}. For the stated Hecke product the sum of its two zero expansions gives the same conditional prime bound and the identical local argument applies.
\end{proof}

For a quadratic extension the local identity is especially simple. At a prime inert in $F$, an elliptic factor with roots $\alpha,\bar\alpha$ in the variable $p^{-2s}$ has the four rational-prime roots $\eta,-\eta,\bar\eta,-\bar\eta$, where $\eta^2=\alpha$. The first power sum is zero, and the second is $2(\alpha+\bar\alpha)$, as required by the residue-degree factor in the logarithmic derivative. At a split prime the four roots are the two original pairs from the two primes above $p$. This accounts for both types using the same number of shifts.

The partial Euler products and tail estimates apply to representations that are tempered outside a fixed finite set. An inverse root of modulus different from one cannot be the limit of unit vertical phases at that prime. This identifies where the present matching construction uses finite-place temperedness beyond the unconditional local estimates \eqref{aw:lrs}.

\clearpage
\section{Approximation by rank-one functions and the cost of approaching a zero}\label{slit:section}

Retain the incomplete tempered degree-two Euler product $L_S$ and the paired shifts $D_t$ from Section~\ref{vt:section}. Write $\mathcal H=\{\Re s>1/2\}$. In this section a horizontal segment from a point always extends to its right. Define
\[
 \Omega=\{s\in\mathcal H: L_S(s+x)\ne0\text{ for every real }x\ge0\}.
\]
Thus one removes the leftward horizontal ray from every zero in $\mathcal H$, stopping at its boundary. A compact subset of $\mathcal H$ meets only finitely many removed rays because the corresponding zeros have bounded ordinate and real part between the compact's left edge and 1. Consequently $\Omega$ is open. Any compact subset can be translated continuously to the absolute-convergence half-plane within $\Omega$.

\begin{theorem}\label{st:main}
There is a sequence $t_j\to+\infty$ matching the elliptic Frobenius phases on every fixed finite set of good primes to a precision tending to zero, such that
\[
       D_{t_j}(s)\longrightarrow L_S(s)
                   \quad\hbox{locally uniformly on }\Omega .
\]
No hypothesis on zeros of either function is needed for this convergence. On any given compact exhaustion of $\Omega$, the requested errors and the matching precisions may be prescribed in advance. For each finite stage the acceptable shifts have positive lower density.

The sequence can in addition be chosen eventually zero-free on every compact subset of $\mathcal H$. Its normality locus in $\mathcal H$ is then exactly $\Omega$. In particular every neighborhood of every point on the removed rays contains unbounded values of the approximants as $j\to\infty$. The domain $\Omega$ is maximal among rightward-closed domains on which an eventually zero-free matched sequence can be normal. Under zeta RH the same maximality holds for every matched sequence tending to infinite height. These assertions concern compact sets from which the poles $1\pm it_j$ have eventually escaped.
\end{theorem}

\paragraph{Positive Taylor weights.}
Let $F$ be the branch of $\log L_S$ on $\Omega$ that agrees with the Euler logarithm on $\Re s>1$. To see that this branch exists, slide any closed curve a common fixed distance to the right. The entire deformation lies in $\Omega$, and its translate contracts in the absolute-convergence half-plane. The same deformation connects points to that half-plane. Write
\[
 F(s)=\sum_{p\notin S}\sum_{k\ge1}
       \frac{\alpha_p^k+\bar\alpha_p^k}{k}p^{-ks}\qquad(\Re s>1).
\]
For $b>0$, an integer $N\ge0$ and $y\ge0$ put
\[
 w_{b,N}(y)=e^{-by}\sum_{\ell=0}^N\frac{(by)^\ell}{\ell!},
 \qquad
 A_{b,N}f(s)=\sum_{\ell=0}^N\frac{(-b)^\ell}{\ell!}f^{(\ell)}(s+b).
\]
On the exponential $q^{-s}$ the operator acts by multiplication by $w_{b,N}(\log q)$. In particular all its scalar weights lie in $[0,1]$ and increase pointwise to 1 with $N$.

\begin{lemma}\label{st:taylor}
For a compact $K\Subset\Omega$, a finite bound $Y_0$, and positive tolerances $\epsilon$ and $\nu$, there is a product $W(y)=\prod_{j=1}^J w_{b,N_j}(y)$ such that
\begin{align*}
 &\sup_{s\in K}\left|F(s)-
 \sum_{p\notin S,k\ge1}\frac{\alpha_p^k+\bar\alpha_p^k}{k}p^{-ks}W(\log p^k)\right|<\epsilon,\\
 &1-W(y)<\nu\quad(0\le y\le\log Y_0).
\end{align*}
The weighted series is absolutely convergent on a neighborhood of $K$.
\end{lemma}
\begin{proof}
Choose $B>0$ such that $\min_{s\in K}\Re s+B>1$. The compact tube $K+[0,B]$ lies in $\Omega$. Take $J$ so large that $b=B/J$ and the closed $4b$-neighborhood of the tube lies there. Increase $B$ first if necessary and decrease $b$ so that the closed $b$-neighborhood of $K+B$ lies in $\Re s>1$. Put $\eta=b/J$ and $K_j=K+jb+\{z:|z|\le j\eta\}$. Let $M\ge1$ bound $|F|$ on the closed $4b$-neighborhood.

For $s\in K_{j-1}$ the Taylor disk centered at $s+b$ of radius $2b$ is available, so $|F(s)-A_{b,N}F(s)|\le M2^{-N}$. If $|H-F|\le e$ on $K_j$ and both are holomorphic on a neighborhood, Cauchy's formula on the smaller disk of radius $\eta$ gives
\[
 |A_{b,N}H(s)-A_{b,N}F(s)|
 \le e\,C_N,\qquad C_N=\sum_{\ell=0}^N(b/\eta)^\ell .
\]
Set $e_0=\epsilon/2$. In the order $j=1,\ldots,J$, choose $N_j$ with $M2^{-N_j}\le e_{j-1}/2$, then put $e_j=e_{j-1}/(2C_{N_j})$. Impose at the same time $1-w_{b,N_j}(\log Y_0)<\nu/J$. This also controls smaller $y$, since the weight decreases in $y$. Starting with the actual $F$ on $K_J$ and applying the operators in reverse order proves that $A_{b,N_1}\cdots A_{b,N_J}F$ approximates $F$ on $K$ to error at most $e_0$. The derivatives at the last translate can all be computed by the absolute Euler series. Equivalently the composite operator uses a finite polynomial in derivatives of $F(s+B)$, which justifies the claimed weighted expansion. Each weight product is $e^{-By}$ times a polynomial in $y$. This gives absolute convergence in $\Re s+B>1$. The inequality $1-\prod_j w_j\le\sum_j(1-w_j)$ proves the finite-coordinate estimate.
\end{proof}

\paragraph{Turning the weights into local parameters.}
For a unit complex number $z$ put
\[
   H_p(z;s)=\sum_{k\ge1}(z^k+\bar z^k)p^{-ks}/k .
\]
For each $p$ choose independently a random phase $Z_p$ equal to $\alpha_p$ with probability $W(\log p)$ and otherwise uniform on the unit circle. Each positive power of a uniform phase has zero mean. Thus
\[
        \mathbb EH_p(Z_p;s)=W(\log p)H_p(\alpha_p;s).
\]
This weight uses $\log p$ even in the terms of order $k\ge2$. Their total change from the series in Lemma~\ref{st:taylor} tends uniformly to zero on any compact in $\Re s>1/2$ as the finite-coordinate tolerances tend to zero. Indeed $2\sum_{p,k\ge2}p^{-k\sigma_*}/k<\infty$ for $\sigma_*>1/2$. Discard a tail of this series first, then use the finite-coordinate bound on the remaining $\log p^k$ and $\log p$. The same argument works if the phases at primes $p\le X$ are fixed at $\alpha_p$ rather than randomized. For the terms $k=1$ at these fixed primes the additional error is at most $2\nu\sum_{p\le X}p^{-\sigma_*}$ when $Y_0\ge X$.

Choose a bounded open neighborhood $U$ of $K$ with $\Re s>\sigma_*>1/2$ and with a disk of fixed radius about every point of $K$ contained in $U$. For arbitrary finite $Y>X$, independence gives
\[
 \mathbb E\int_U\left|\sum_{X<p\le Y,\,p\notin S}
       [H_p(Z_p;s)-\mathbb EH_p(Z_p;s)]\right|^2dA(s)
 \ll_{\sigma_*}|U|\sum_{p>X}p^{-2\sigma_*}
 \ll_{\sigma_*}|U|X^{1-2\sigma_*}.
\]
Here $|H_p(z;s)|\le 2p^{-\sigma_*}/(1-2^{-\sigma_*})$. The disk mean-value inequality supplies a deterministic set of phases with uniform error $O_{K,U}(X^{1/2-\sigma_*})$ on $K$. For each fixed weight product the expected sum over $p>Y$ tends absolutely to zero on $K$, since its first powers have $W(\log p)=p^{-B}$ times a fixed polynomial in $\log p$, and the higher powers are summable.

Given a prescribed finite matching set and error, first increase $X$ beyond that set until the variance error is small enough. Next take $Y_0,\nu$ for the deterministic and higher-power errors, apply Lemma~\ref{st:taylor}, and take $Y$ large enough for the expected tail. We have proved that a finite genuine Euler product with unit roots, fixed at all $p\le X$, approximates $L_S$ on $K$ to any prescribed accuracy. Exponentiation is harmless because $F$ is bounded on $K$.

Finally prescribe narrow open arcs around these chosen phases at the good $p\le Y$. The product of the two local factors varies uniformly continuously with these finitely many phases on $K$. The conditional zeta-tail mean-square estimate \eqref{vt:meansquare}, including its disk version for both signs of the shift, bounds the two omitted zeta tails uniformly by $O_{K}(Y^{1/2-\sigma_*})$ on a positive fraction of this box. This bound is independent of the chosen finite arcs. One may enlarge $Y$ before choosing the realization to meet any required tail tolerance. Applying this construction to a compact exhaustion proves the first assertion, with each acceptable set of heights of positive lower density. All choices of the Taylor cutoffs are finite and the proof makes no uniform claim about the first acceptable shift.

\paragraph{Where normality must fail.}
We may ensure eventual zero-freeness even without information on individual zeta zeros. For any fixed compact $E\Subset\mathcal H$, surround $E\cup\bar E$ by a bounded symmetric rectangle $U_0$ with real parts at least $\sigma_0>1/2$ and with $r$-disks about these points contained inside. If $\zeta(s+it)$ vanishes on this compact, then $|Z_{>Z}(s+it)-1|=1$ somewhere there for every finite $Z$. Its integrated square over $U_0$ is at least $\pi r^2$, by the disk mean value inequality. By \eqref{vt:meansquare} with unrestricted prime phases the exceptional upper density is at most $|U_0|\sum_{n>Z}n^{-2\sigma_0}/(\pi r^2)$, which tends to zero as $Z\to\infty$. The symmetric rectangle also handles the negative shift. Remove this density-zero set at each positive-density stage, taking the compacts $E$ to exhaust $\mathcal H$. On every bounded region compactly contained in $\mathcal H$ the late chosen $D_{t_j}$ are therefore zero-free and holomorphic. Convergence on a Jordan curve surrounding zeros of $L_S$ and avoiding its zeros is impossible by Rouché's theorem if there are no poles inside. Fix $s_0=x+i\gamma_0$ on a removed ray. If $L_S(s_0)=0$ and the family were normal near $s_0$, a convergent subsequence there would agree with $L_S$ by its already known limit on a neighboring open subset of $\Omega$. Hurwitz's theorem gives a contradiction. This argument also shows that normality at any point of a ray forces local convergence to $L_S$ there, because every convergent subsequence has the same limit. For spherical normality, a meromorphic subsequential limit agreeing locally with the holomorphic $L_S$ agrees throughout the disk.

If $L_S(s_0)\ne0$, choose a rectangle with left edge crossing $s_0$, right edge in $\Re s>1$, and sufficiently short height that no other ordinate of a zero to the right of $x$ lies within it. The zeros at ordinate $\gamma_0$ to the right of $x$ are inside. Compactness and discreteness allow the horizontal edges to avoid all removed rays intersecting this rectangle. If the sequence were normal near $s_0$, shorten the rectangle's height so that its left edge lies in that neighborhood. The sequence then converges uniformly on the full boundary to $L_S$, contradicting Rouché. Thus the nonnormal locus is exactly the removed rays. Local boundedness would imply normality by Montel. For any other matched sequence that is eventually zero-free on compact sets and normal on a rightward-closed domain, the limit is fixed by its agreement with the Euler product to the right of one. Hurwitz excludes every zero in that domain, hence also every point lying immediately or arbitrarily far to its left on the same horizontal ray. Under zeta RH all matched sequences tending to infinite height possess the zero-free property.

In particular a certified elliptic height interval produces normal rank-one approximants on compact subsets of the full corridor $1/2<\Re s<\infty$ at its interior ordinates. An off-line zero outside that interval does not affect those approximants. To prove elliptic RH by this route one still must establish normality across the rays associated with any putative off-line zeros. Finite-prime matching alone permits precisely the loss of normality just described.

\subsection*{Primitive Dirichlet approximants with an explicit conductor threshold}
The vertical shifts can be replaced throughout by primitive Dirichlet characters. The change also allows general tempered targets over number fields.

\begin{theorem}\label{st:characters}
Suppose a holomorphic $L_S$ on $\mathcal H$ has, outside a finite set of rational primes, an Euler product of fixed absolute degree $d$ with unit local roots $\alpha_{p,1},\ldots,\alpha_{p,d}$. Use its horizontal zero-free domain $\Omega$ as above. There are primes $q_j\to\infty$ and tuples of primitive nonprincipal characters modulo $q_j$ such that
\[
   \prod_{v=1}^d L_S(s,\chi_{j,v})\longrightarrow L_S(s)
                \quad\hbox{locally uniformly on }\Omega,
\]
with each $\chi_{j,v}(p)$ approaching its prescribed $\alpha_{p,v}$ at every fixed good prime. For each finite approximation stage, \emph{every} sufficiently large prime conductor works, with a positive proportion of tuples. An explicit threshold is given below. For an elliptic target over $\Q$ one can use just a primitive character and its conjugate. The characters can be chosen so that every individual factor is eventually zero-free on each compact subset of $\mathcal H$. The exact normality locus of the constructed products is then $\Omega$, unconditionally.

This applies in particular to an elliptic curve over any number field for which the stated holomorphic continuation is known. At an unramified rational prime substitute the $f$th roots of each normalized Frobenius parameter at a place of residue degree $f$. Their factors multiply to the original factor at $p^{-fs}$, and all roots have modulus one. Good rational primes give total degree $2[F:\Q]$. The same conclusion holds for standard cuspidal functions known to be tempered at all but finitely many finite places, with their absolute degree after this rational-prime expansion.
\end{theorem}

\begin{proof}
We give effective details of the conditional estimate in this family. Fix a bounded rectangle $U$ with $1/2<\sigma_*<1$ and $\Re s\ge\sigma_*$ throughout its closure. It may be taken symmetric about the real axis for the conjugate-pair application. Let $Y\ge2$ be an integer, $Q_Y=\prod_{p\le Y}p$, and prescribe arcs at the primes $p\le Y$. When a prime has no constraint, one can impose any smaller arc there. Write $h_p\in(0,\pi]$ for the angular half-widths. Set
\[
 V_Y=Y^{1/2-\sigma_*}/\sqrt{2\sigma_*-1},\quad
 \mu=\prod_{p\le Y}\frac{h_p}{2\pi},\quad
 B_Y=\prod_{p\le Y}(1+p^{-\sigma_*}),
\]
\[
 C=B_Y\left[(2\sigma_*-1)^{-1/2}
        +\sup_{s\in U}|s|(1+1/\sigma_*)\right]
       +\tau(Q_Y)(2\sigma_*-1)^{-1/2}.
\]
Choose integers $N,M$ such that
\begin{align}
 &\frac{\pi[1+\log(N+1)]}{N+1}\sum_{p\le Y}h_p^{-1}\le\mu/8,
 \qquad M\ge Y,\quad C M^{1/2-\sigma_*}\le\sqrt\mu\,V_Y. \label{st:mq}
\end{align}
It is sufficient that the prime $q$ satisfy
\begin{equation}\label{st:primecut}
        q>\max\{M Q_Y,\ M Q_Y^N,\ 1+8/\mu\}.
\end{equation}
At least $(q-1)\mu/3$ primitive characters modulo $q$ then obey all the arcs and
\[
 \int_U\left|L(s,\chi)\prod_{p\le Y}(1-\chi(p)p^{-s})-1\right|^2dA(s)
      \le16|U|V_Y^2.
\]

To prove this claim, normalize character sums by $q-1$. For nonprincipal $\chi$ put
\[
 R_{a,q}(s)=\sum_{k\ge1}\big[(kq+a)^{-s}-(kq)^{-s}\big],
 \qquad 1\le a<q .
\]
Then $L(s,\chi)=\sum_{a<q}\chi(a)(a^{-s}+R_{a,q}(s))$, with $a\ge1$ understood. Group the absolutely convergent series on $\Re s>1$ and use $\sum_{a<q}\chi(a)=0$. The differences continue the identity to $\Re s>0$ by absolute locally uniform convergence. Integration of the derivative of $x^{-s}$ gives $|R_{a,q}(s)|\le \sup_U|s|(1+1/\sigma_*)q^{-\sigma_*}$. Orthogonality applied to the resulting vector of residues proves that the root mean square over nonprincipal characters of $L(s,\chi)-\sum_{n\le M}\chi(n)n^{-s}$ is at most
\[
 M^{1/2-\sigma_*}(2\sigma_*-1)^{-1/2}
       +\sup_U|s|(1+1/\sigma_*)q^{1/2-\sigma_*}.
\]
For the upper bound we may include the principal character when summing the squares of the finite residue expression. The continuation identity itself was used only for the nonprincipal characters.

Multiply the short polynomial by $\prod_{p\le Y}(1-\chi(p)p^{-s})$. Its coefficients through $M$ are exactly those of
$R_M(s,\chi)=\sum_{n\le M,(n,Q_Y)=1}\chi(n)n^{-s}$.
The excess terms stop at $M Q_Y$, have coefficient magnitudes at most $\tau(Q_Y)$ before inserting $n^{-s}$, and are again orthogonal since $q>M Q_Y$. This bounds the root mean square of the difference between the deleted $L$-function and $R_M$ by $CM^{1/2-\sigma_*}$.

For each arc use the triangular tent of height one and support that arc. The product $w$ of these tents has Haar mean $\mu$. Replace each tent by its Fejér mean of degree $N$, obtaining $0\le P\le1$ of mean $\mu$ and $\|w-P\|_\infty\le\mu/8$. Indeed a tent has Lipschitz constant $1/h_p$, while the normalized Fejér convolution has absolute first angular moment at most $\pi[1+\log(N+1)]/(N+1)$. Split the moment integral at $\pi/(N+1)$ and use the kernel bounds $N+1$ and $\pi^2/[(N+1)\theta^2]$. Since all factors are bounded by one, their errors add.

The average over all characters of $P|R_M-1|^2$ equals its torus integral \emph{exactly}. Every rational frequency has numerator and denominator at most $M Q_Y^N$. Modulo $q$ only an actual equality of these integers can therefore survive. Unique factorization and $(n,Q_Y)=1$ leave only the constant low-prime mode and diagonal high-prime terms. The average is at most $\mu V_Y^2$. The unweighted average is at most $V_Y^2$. The mass of $w$ on nonprincipal characters is at least $\mu-\mu/8-1/(q-1)\ge3\mu/4$. Its weighted square error for the deleted $L$-function is at most
\[
         2(\mu+\mu/8)V_Y^2+2C^2M^{1-2\sigma_*}\le(17/4)\mu V_Y^2
\]
at each $s$. Integrate over $U$ and apply Markov's inequality. The good mass is at least $(3/4-17/64)\mu>\mu/3$. This proves the claim. The disk mean-value bound now supplies a uniform approximation on interior compact subsets. For the conjugate character the same rectangle integral applies by conjugating $s$.

For general $d$, apply Lemma~\ref{st:taylor} to the Euler logarithm with coefficient $\sum_v\alpha_{p,v}^k/k$. Randomize each prescribed phase separately by the same weight, retaining its first $X$ prime values. The independent local variance is bounded by a constant depending on $d$ times $p^{-2\sigma_*}$. Higher powers are bounded by $d\sum_{p,k\ge2}p^{-k\sigma_*}/k$. These are exactly the two estimates used to select a finite Euler product in Theorem~\ref{st:main}. Each character coordinate has its own arcs about the resulting local phases. Choose a common $q$ past the largest threshold \eqref{st:primecut} and choose the coordinates independently from their good sets. At any one stage, a realizing prime lies below twice a suitably rounded threshold by Bertrand's postulate. The analytic Taylor bound uses an upper bound for $|\log L_S|$ on the specified zero-free tube. The conductor estimate is explicit in the resulting finite Euler data and the requested precision. To obtain eventual compact zero-freeness with an effective cutoff, fix an auxiliary compact $E\Subset\mathcal H$ with a surrounding rectangle $U_0$, disk radius $r$ and $1/2<\sigma_0<1$ as in the vertical argument. Use the definitions of $V_Z,C$ above at $\sigma_0,U_0$, writing $C_Z$ for the resulting constant. Choose integers $Z\ge2$ and $M_0\ge Z$ with
\[
 \frac{4|U_0|V_Z^2}{\pi r^2}\le\mu/6,\qquad
 C_Z M_0^{1/2-\sigma_0}\le V_Z,\qquad q>M_0Q_Z.
\]
The unweighted short rough polynomial has square error from $1$ at most $V_Z^2$ by exact character orthogonality. The root mean square of the difference between that polynomial and the deleted $L$-function is at most $C_Z M_0^{1/2-\sigma_0}$, as proved above. Thus the unweighted average integrated square error is at most $4|U_0|V_Z^2$. A zero in $E$ forces an integral of at least $\pi r^2$. Removing the exceptional characters therefore leaves at least $(q-1)\mu/6$ good primitive characters per coordinate. For conjugate pairs take $E$ symmetric. Exhausting $\mathcal H$ by such compacts gives eventual zero-freeness unconditionally. The same rectangle argument proves the exact nonnormal set.
\end{proof}

\subsection*{A conductor lower bound at an omitted zero}
Write $P_j=\prod_{v=1}^dL_S(\,\cdot\,,\chi_{j,v})$ for the primitive-character approximants. The singularity of this sequence along each removed ray has a measurable sign and a multiplicity.

\begin{proposition}\label{st:jump}
Use a sequence provided by Theorem~\ref{st:characters} with eventual compact zero-freeness. Let $1/2<x<1$ and $L_S(x+i\gamma)\ne0$. Choose $h>0$ such that $L_S$ has no zeros on the segment $[x+i(\gamma-h),x+i(\gamma+h)]$ and every zero with real part at least $x$ and ordinate in $[\gamma-h,\gamma+h]$ has ordinate $\gamma$. Let $m$ count these zeros with multiplicity. Then
\[
 \lim_{j\to\infty}\int_{\gamma-h}^{\gamma+h}
  \Re\left(\frac{P'_j}{P_j}(x+it)-\frac{L_S'}{L_S}(x+it)\right)dt
       =2\pi m.
\]
Under Dirichlet GRH, if $m>0$, any such conductor-$q$ approximation whose integrated error differs from $2\pi m$ by at most $\pi m/2$ satisfies, for sufficiently small $h$,
\[
  \log(2q)\gg_{x,\gamma,S}
              \left(\frac{m}{d h}\right)^{1/[2(1-x)]}.
\]
The lower bound applies as soon as the approximation is sufficiently good on the three other sides of a rectangle closing the segment to the right of one. The smallness threshold for $h$ may depend on $L_S$.
\end{proposition}
\begin{proof}
Choose the right edge at $\Re s=2$. Traverse the bottom, right and top edges from the lower left to the upper left. These edges lie in $\Omega$. There the logarithmic derivatives of the approximants converge uniformly to $L_S'/L_S$, by holomorphic convergence on neighborhoods of the edges. The full counterclockwise rectangular integral for $P_j'/P_j$ is eventually zero by our choice of the approximants. For $L_S'/L_S$ it is $2\pi i m$. The missing left edge points downwards in both contour integrals. Subtracting and taking imaginary parts yields the formula with its positive sign.

For completeness we give the conductor estimate uniformly at bounded heights. For primitive nonprincipal $\chi$ of conductor $q$, fixed $x\in(1/2,1)$ and $t$ in a compact interval $I$, Dirichlet GRH gives
\begin{equation}\label{st:logder}
     |(L'/L)(x+it,\chi)|\ll_{x,I} [\log(2q)]^{2(1-x)} .
\end{equation}
In the smoothed logarithmic derivative formula of
\href{https://arxiv.org/abs/2206.00819}{Chirre, Hagen and Simoni\v c, equation (2.1)}, take the two cutoffs $X$ and $2X$, where $X=\max(2,[\log(2q)]^2)$. The prime-power sum is bounded by $\sum_{n\le2X}\Lambda(n)n^{-x}\ll_x X^{1-x}$ by Chebyshev's estimate and partial summation. The nontrivial zero terms are bounded by a constant times $X^{1/2-x}\sum_\rho|x+it-\rho|^{-2}$. With $\rho=1/2+i\gamma_\rho$, each squared reciprocal is at most a constant depending on $x$ times $(3/2)/[(3/2)^2+(t-\gamma_\rho)^2]$. The sum of these Poisson kernels is $\tfrac12\log q+O_I(1)$ by the real Hadamard formula at $2+it$, its bounded Euler series and its gamma factor. The trivial zero terms are $O_x(X^{-x})$, by their spacing two and distance at least $x$ from $x+it$. This proves \eqref{st:logder}. Omitting the fixed Euler factors at $S$ adds $O_{x,S}(1)$.

For small $h$ the integral of $\Re(L_S'/L_S)$ along the segment has absolute value at most $\pi m/2$, by regularity at its midpoint. The stated accuracy therefore gives $\int\Re(P'/P)\,dt\ge\pi m$. The last integral is at most $2hd$ times the bound in \eqref{st:logder}, including the fixed omitted factors. Rearrangement proves the conductor inequality.
\end{proof}

Under Dirichlet GRH there is also an interpretation purely in terms of critical-line rank-one zeros. Put $b=x-1/2$ and let $a_\chi\in\{0,1\}$ be the parity of $\chi$. For each primitive character replace $\Re(L'/L)(x+it,\chi)$ in Proposition~\ref{st:jump} by
\[
 \sum_{\gamma_\chi}\frac{b}{b^2+(t-\gamma_\chi)^2}
 -\tfrac12\log(q/\pi)
 -\tfrac12\Re\psi\big((x+it+a_\chi)/2\big).
\]
The omitted local terms for $L_S$ are the logarithmic derivatives of $\prod_{p\in S}(1-\chi(p)p^{-s})$. Hence the excess integral of these centered Poisson sums over the target integral converges to the exact mass $2\pi m$. This statement concerns fixed intervals and their specified order of limits. It does not presuppose a uniform bound for the total variation of the signed integrands.

\begin{theorem}[Transfer using only rank-one zero ordinates]\label{st:ordinates}
The characters in Theorem~\ref{st:characters} may be chosen, unconditionally, to have the two further properties below.
\begin{enumerate}
\item For each sufficiently large selected conductor $q$, every nontrivial zero $\beta+i\gamma_\chi$ of each selected character with $|\gamma_\chi|\le(\log q)^2$ satisfies
\[
     |\beta-1/2|\le 30\log\log q/\log q.
\]
The same exponent in the conductor lower bound of Proposition~\ref{st:jump} holds for these approximants without Dirichlet GRH.
\item In the Poisson expression preceding this theorem put every nontrivial zero at $1/2+i\gamma_\chi$, retaining its actual signed ordinate and multiplicity. The resulting expression differs from $\Re(L'/L)(x+it,\chi)$ by $o(1)$ uniformly on every fixed interval of $t$ at each $x\in(1/2,1)$. The centered sums computed from rank-one ordinates alone, with the indicated gamma, conductor and deleted-factor corrections, consequently give the same limiting excess $2\pi m$.
\end{enumerate}
\end{theorem}
\begin{proof}
Montgomery's fixed-modulus density estimate, for $1/2\le v\le3/4$, is
\[
 \sum_{\chi\bmod q}N(v,H,\chi)
 \ll (qH)^{3(1-v)/(2-v)}[\log(qH)]^9
 \qquad (H\ge2).
\]
Here $N$ counts zeros with real part at least $v$ and $|\Im\rho|\le H$. See Montgomery, \emph{Topics in Multiplicative Number Theory}, Theorem 12.1, also stated in \href{https://www.impan.pl/shop/en/publication/transaction/download/product/82228}{Wolke, equation (2.2.5)}. Put $H=(\log q)^2$, $\eta=30\log\log q/\log q$, and $v=1/2+\eta$. For large $q$ the exponent is $1-2\eta/(3/2-\eta)\le1-4\eta/3$. The number of exceptional characters divided by $q-1$ is therefore $O((\log q)^{-29})$. Deleting these exceptions at any fixed positive-proportion approximation stage leaves a positive proportion at every sufficiently large prime. Horizontal reflection of the zeros of the same primitive character follows by conjugation and the functional equation, so the lower half of the asserted symmetric enclosure follows as well.

Write $s=x+it$ with $t\in I$ fixed in a bounded interval. Regardless of GRH, the real Hadamard formula at $2+iu$ implies that the number of nontrivial zeros with $|\gamma_\chi-u|\le1$ is $O(\log(q(2+|u|)))$. Each contributes a Poisson kernel bounded away from zero, and the gamma and absolute Euler terms there are bounded by these logarithms. It follows that
\[
 \sum_{|\gamma_\chi|>H}|s-\rho|^{-2}\ll_I \log(qH)/H .
\]
For $|\gamma_\chi|\le H$ we have $x-\beta\ge(x-1/2)/2$ once $q$ is sufficiently large. Comparison with the real Hadamard formula at $2+it$ now gives $\sum_{|\gamma_\chi|\le H}|s-\rho|^{-2}\ll_{x,I}\log(2q)$. In the same smoothed formula used in \eqref{st:logder}, the low-zero error is thus at most a constant times $X^{1/2+\eta-x}\log(2q)$, and the high-zero error is $O_{x,I}(X^{1-x}\log(qH)/H)$. The quantity $X^\eta$ is bounded because $X\asymp(\log q)^2$. This proves \eqref{st:logder} for every selected character with its original exponent and gives the conductor obstruction.

For the second assertion put $b=x-1/2$ and $P(u,y)=u/(u^2+y^2)$. In the general real Hadamard formula the term for a zero is $P(x-\beta,t-\gamma_\chi)$, irrespective of GRH. At $|\gamma_\chi|\le H$ pair the zeros horizontally. The two first-order Taylor terms around $b$ cancel. For $b/2\le u\le3b/2$,
\[
 |\partial_u^2P(u,y)|=|2\Re(u+iy)^{-3}|
                 \ll_b(1+y^2)^{-1}.
\]
Summation gives a total change $O_{x,I}(\eta^2\log(2q))$. For the remaining zeros either kernel separately has total absolute size $O_{x,I}(\log(qH)/H)$ by the unit-interval count. Both errors tend to zero. The finite omitted local terms and the gamma and conductor contributions are identical before and after this projection. Integrating the uniform error proves the assertion.
\end{proof}

\begin{example}[A unitary Euler product with a prescribed off-line zero]\label{st:helson}
In \href{https://arxiv.org/abs/1812.11729}{Seip, \emph{Universality and distribution of zeros and poles of some zeta functions}, Theorem 1.1}, take the one-point positive divisor $\{2/3+i\}$ and $\alpha=7/10<59/80$. Its counting function is bounded and its unit increments eventually vanish, so all the theorem's density hypotheses hold (the count vanishes whenever $\sigma\ge2/3$). Seip gives a completely multiplicative $\omega$ with $|\omega(n)|=1$ such that $Z_\omega(s)=\prod_p(1-\omega(p)p^{-s})^{-1}$ extends holomorphically to $\mathcal H$ with this single simple zero there. Then
\[
           G(s)=Z_\omega(s)\overline{Z_\omega(\bar s)}
\]
has precisely the two simple zeros $2/3\pm i$ there. Its Euler roots at every prime are a conjugate pair of modulus one and determinant one. Theorems~\ref{st:main} and \ref{st:characters} construct genuine paired rank-one approximants on the complement of the two rays. The integrated excess in Proposition~\ref{st:jump} equals $2\pi$ in the limit at every $1/2<x<2/3$ on either ray. The example makes no modularity or functional-equation assertion. It shows why holomorphic continuation on $\mathcal H$, tempered determinant-one local roots, and approximation by actual Dirichlet functions on the entire slit domain cannot by themselves imply an elliptic Riemann hypothesis.
\end{example}

\subsection*{A subconvex separation bound for the Frobenius traces themselves}

Cuspidality gives an additional obstruction before one asks about zeros. Write $\lambda_E(p)=a_p(E)/\sqrt p$ at the good primes of a fixed rational elliptic curve.

\begin{theorem}\label{st:subconvex}
Fix an elliptic curve $E/\Q$ and a finite set $S$ containing its bad primes. Suppose $\chi$ is a nonprincipal character of prime conductor $q$ and $X$ is sufficiently large. If
\[
 \left|\lambda_E(p)-\chi(p)-\overline{\chi(p)}\right|
       \le (\log X)^{-5}
       \qquad(p\le X,\ p\notin S\cup\{q\}),
\]
then for every $\varepsilon>0$ one has
\[
                q\gg_{E,S,\varepsilon} X^{4/3-\varepsilon}.
\]
This estimate holds without a hypothesis on zeros of Dirichlet or elliptic functions. It asks only that traces match. An ordering of the two Euler roots is unnecessary.
\end{theorem}
\begin{proof}
There are two distinct good primes outside $S$ at which the normalized trace avoids $-2,0,2$. For example choose primes $p>3$ splitting completely in the fixed $3$-division field of $E$, by Chebotarev. Then $a_p\equiv2\pmod3$, and $a_p=\pm2\sqrt p$ is impossible by integrality. At least one fixed witness differs from $q$. Its positive distance from $\{-2,0,2\}$ excludes $\chi^4=1$ for large $X$, since traces of fourth roots of unity belong to that set. Consequently $\bar\chi^2$ is a primitive nonquadratic character modulo $q$.

Put $D=S\cup\{q\}$ and, with the unitary normalization, compare
\[
 A(s)=L_D(s,E\otimes\bar\chi),\qquad
 B(s)=\zeta_D(s)L_D(s,\bar\chi^2).
\]
At each $p\notin D$ the two pairs of local roots both have determinant $\bar\chi(p)^2$. Their traces differ by at most $\delta=(\log X)^{-5}$ for $p\le X$. Subtracting the two reciprocal quadratic polynomials in the variable $u=p^{-s}$ gives the trace difference times $u$ divided by the product of the two quadratics. Since all four roots have modulus one, the coefficient difference at $p^k$ is at most $\delta\binom{k+2}{3}$ for $k\ge1$, and either coefficient alone has modulus at most $k+1$. Telescoping across the prime powers thus gives
\[
 |a_A(n)-a_B(n)|\le \delta\,\tau(n)
           \sum_{p^k\parallel n}k(k+2)/6
  \ll \delta\,\tau(n)[\log(2n)]^2 \qquad(n\le X).
\]
Both coefficients vanish if $(n,\prod_{p\in D}p)>1$. Fix a nonnegative smooth $w$ supported inside $(1/2,1)$ with positive integral. Summing the last display with $w(n/X)$ costs $O_w(\delta X\log^3X)$, using the elementary divisor sum.

Let $\widehat w(s)=\int_0^\infty w(v)v^{s-1}dv$. Shift the two Mellin contours to $\Re s=1/2$. The elliptic twist is cuspidal. For its critical line use
\href{https://www.renyi.hu/~gharcos/hybridbounds.pdf}{Blomer and Harcos, \emph{Hybrid bounds for twisted $L$-functions}, Theorem 2}, with their \href{https://www.renyi.hu/~gharcos/hybridbounds_addendum.pdf}{addendum}. For fixed $E$, it gives $\ll_{E,\epsilon}q^{3/8+\epsilon}(1+|t|)^{1/2+\epsilon}$. Fixed omitted local polynomials are bounded there. If $q$ is good the $q$-factor of the ramified twist is one. When $q$ is one of the fixed bad primes its deletion is covered by the same fixed bound. For $B$ the ordinary convexity estimate for the single Dirichlet factor costs $q^{1/4+\epsilon}$ times a polynomial in $|t|$, with the zeta factor contributing only a further polynomial. Deleting the prime $q$ from zeta contributes $1-q^{-s}$, uniformly bounded. The rapid vertical decay of $\widehat w$ now gives
\[
 X\widehat w(1)|\operatorname*{Res}_{s=1}B(s)|
     \ll_{E,S,w,\epsilon} X^{1/2}q^{3/8+\epsilon}
                      + \delta X\log^3 X .
\]
The residue is $L(1,\bar\chi^2)$ times the omitted local factors at one, whose product has modulus bounded below in terms of $S$. The effective lower bound for a primitive nonquadratic character is $|L(1,\bar\chi^2)|\gg1/\log(2q)$. See \href{https://arxiv.org/abs/2107.09230}{Mossinghoff, Starichkova and Trudgian, \emph{Explicit lower bounds on $|L(1,\chi)|$}}. For $q>X^2$ the conclusion already holds. Otherwise $\log(2q)\ll\log X$, and the coefficient error $O(X/\log^2 X)$ is absorbed for large $X$. Hence $X^{1/2}\ll_{E,S,\epsilon}q^{3/8+\epsilon}\log(2q)$. Taking $\epsilon$ sufficiently small in terms of the stated $\varepsilon$ proves the bound.
\end{proof}

More generally the argument gives a bound in terms of the available twist estimate. If a fixed tempered cuspidal degree-two target with trivial central character has the two trace witnesses used above and its twists satisfy a polynomial-in-height bound $q^{a+\epsilon}$ on the critical line, this proof gives $q\gg X^{1/[2\max(a,1/4)]-\varepsilon}$. If instead both the cuspidal twist and the single Dirichlet factor satisfy the conductor-aspect Lindel\"of estimate with polynomial height dependence, the argument forces $q\gg_B X^B$ for every fixed $B$. These are obstructions to simultaneous local matching, without assuming a zero of the target in $\mathcal H$.

\subsection*{How many distinct character values a number-field curve requires}

There is a separate arithmetic obstruction which applies to elliptic curves over every number field, even before continuation of their $L$-functions is available.

\begin{theorem}\label{st:packing}
Let $E/F$ be an elliptic curve over a number field. There is a positive-density set $\mathcal P$ of rational primes, splitting completely in $F$, at which every place has good ordinary reduction. At each $p\in\mathcal P$ one may choose \emph{any} normalized Frobenius root $\alpha_p$ for \emph{any} place above $p$, independently across the primes. Let $k\ge1$ be fixed and $D_k=2^{2k+1}$. Suppose a character $\chi$ of prime conductor $q$ satisfies
\[
 |\chi(p)-\alpha_p|\le\epsilon_X
 \quad(p\in\mathcal P,\ p\le X,\ p\ne q),\qquad
   2k\epsilon_X<2^{1-D_k}X^{-kD_k}.
\]
Then $\operatorname{ord}(\chi)\gg_{E,F,k} X^k/(\log X)^k$. If $E$ has complex multiplication, it suffices to assume $2k\epsilon_X<1/(2X^k)$. In particular, root matching with error smaller than every negative power of $X$ forces $q$ to exceed every fixed power of $X$. For a CM curve and any fixed $a>0$, the error bound $\epsilon_X\le X^{-a}$ gives
\[
     \operatorname{ord}(\chi)\gg_{E,F,a}
                      X^a/(\log X)^{\lceil a\rceil}.
\]
\end{theorem}

\begin{proof}
Take the primes splitting completely in the normal closure over $\Q$ of $F(E[3])$, removing finitely many primes including those below places of bad reduction and all $p\le5$. Chebotarev gives $\#\{p\le X:p\in\mathcal P\}\sim c_{E,F}X/\log X$. The group schemes of $3$-torsion at every place above such $p$ have trivial Frobenius. The integer trace $a_v$ is therefore $2$ modulo 3, so $a_v\ne0$. The Hasse bound $|a_v|\le2\sqrt p<p$ implies $p\nmid a_v$.

Write $\alpha_p=\pi_p/\sqrt p$, where $\pi_p$ is a root of $T^2-a_vT+p$ and both complex conjugates of $\pi_p$ have modulus $\sqrt p$. At any prime ideal above $p$ in a common number field containing the roots, the valuations of $\pi_p,\bar\pi_p$ are $0$ and $v(p)$ in some order. This follows from their product $p$ and their sum $a_v$, a unit. In particular $v(\alpha_p)=\pm v(p)/2$, whereas $v(\alpha_l)=0$ for every different rational prime $l$ in the construction. Thus products $z_A=\prod_{p\in A}\alpha_p$ over different $k$-element subsets are distinct even up to multiplication by a root of unity.

For two different subsets $A,B$, let $r$ be the product of the primes in their symmetric difference and $u=z_A/z_B$. There is an algebraic integer $c$ equal to the product of the appropriate $\pi_p$ or $\bar\pi_p$ such that $u=c/\sqrt r$. The number $c-\sqrt r$ is a nonzero algebraic integer in a field of degree at most $D_k$, obtained by adjoining at most $2k$ quadratic Frobenius roots and the single square root $\sqrt r$. Every conjugate has modulus at most $2\sqrt r$. Its nonzero integer norm yields
\[
      |u-1|\ge2^{1-D_k}r^{-D_k/2}
              \ge2^{1-D_k}X^{-kD_k}.
\]
In the CM case all the $\pi_p$ belong to the fixed imaginary quadratic CM field $K$. Indeed the endomorphism algebra specializes injectively, and the geometric endomorphism algebra of an ordinary elliptic curve over a finite field is imaginary quadratic. See \href{https://math.rice.edu/~av15/Files/Deuring.pdf}{\emph{Singular and supersingular moduli}, Chapter 6}. Now $u^2=c/\bar c$, so $c-\bar c$ is a nonzero algebraic integer in $K$. Its modulus is at least one, since its squared modulus is its positive integer norm. Hence $|u-1|\ge |u^2-1|/2\ge1/(2\sqrt r)\ge1/(2X^k)$.

Each product $\prod_{p\in A}\chi(p)$ lies within $k\epsilon_X$ of $z_A$. The two separation bounds therefore give a different character value for every $k$-element subset of the allowed primes other than $q$. There are $\gg_{E,F,k}(X/\log X)^k$ such subsets. For the last CM assertion take $k=\lceil a\rceil$ and apply this count at $Y=(X^a/(8k))^{1/k}\le X$. The errors satisfy the required strict inequality at $Y$, and $Y^k=X^a/(8k)$.
\end{proof}

The result applies to each coordinate of the absolute-degree character construction over $F$. At the completely split primes all local roots in the rational-prime expansion occur at degree-one places. For unordered Euler polynomials, suppose every coefficient of a monic degree-$d$ approximating polynomial of unit roots differs from that of the true local polynomial by at most $e$. Then every approximating root $z$ is within $(de)^{1/d}$ of some true root. Indeed the value of the true polynomial at $z$ is bounded by $de$ and is the product of those $d$ distances in absolute value.

For comparison, the classical Baker--W\"ustholz lower bound (\emph{Logarithmic forms and group varieties}, 1993) for nonzero linear forms in logarithms gives $\operatorname{ord}(\chi)\gg_{E,F}\epsilon_X^{-c_{E,F}}$ by using just one fixed ordinary prime different from $q$ (two candidate primes suffice). To see the implication, let $n$ be the order of the character. For a chosen nontorsion normalized root $\alpha$ at that prime, take the integer $m$ nearest $n\arg(\alpha)/(2\pi)$. The nonzero form $n\log\alpha-2m\log(-1)$ has magnitude at least a fixed negative power of $n$ and at most $(\pi/2)|\alpha^n-1|\le (\pi/2)n\epsilon_X$. There are only finitely many candidate roots at the two primes. The elementary simultaneous-product proof above gives explicit precision exponents independent of the curve and the almost first-power CM bound.

\subsection*{A bound for angular and finite Hecke characters over the base field}

\begin{theorem}\label{st:hecke-subconvex}
Let $F$ be any number field and $\pi$ a fixed cuspidal representation of $\mathrm{GL}_2(\mathbb A_F)$ with trivial central character, tempered outside a fixed finite set $S$. Let $\chi$ be a unitary Hecke character of analytic conductor $Q\ge2$. Assume that $\chi^4$ is not a pure imaginary power of the idelic norm. For fixed $a>0$ suppose
\[
  |\lambda_\pi(v)-\chi(v)-\overline{\chi(v)}|\le X^{-a}
      \quad (Nv\le X,\ v\notin S,\ \chi\text{ unramified at }v).
\]
Then, for every $\varepsilon>0$,
\[
             Q\gg_{F,\pi,S,a,\varepsilon} X^{128/103-\varepsilon}.
\]
Here $Q$ includes both finite conductor norm and archimedean conductor. In particular this applies to an automorphic elliptic curve over $F$. Over an imaginary quadratic base, for a fixed angular character of nonzero weight and no continuous norm component, its $k$th power has $Q\asymp (1+|k|)^2$, provided the finite conductor of the powers is bounded. Thus $|k|\gg X^{64/103-\varepsilon}$ is necessary for the indicated matching.
\end{theorem}
\begin{proof}
Let $D$ consist of $S$ together with the finite places ramified in $\chi$, and put $\xi=\bar\chi^2$. Compare $L_{F,D}(s,\pi\otimes\bar\chi)$ to $\zeta_{F,D}(s)L_{F,D}(s,\xi)$. The coefficient comparison of Theorem~\ref{st:subconvex} works for ideals, since its only inputs were the local quadratic polynomials and multiplicativity. The sum of the ideal divisor function for $N\mathfrak n\le X$ is $O_F(X\log(2X))$, by counting the pairs $N(\mathfrak a\mathfrak b)\le X$ and using the fixed-field estimate $\#\{\mathfrak a:N\mathfrak a\le y\}=O_F(y)$. The weighted coefficient discrepancy is therefore $O_{F,w}(X^{1-a}\log^3X)$.

We first verify the mild residue estimate needed here. For every $e>0$,
\[
       |L_{F,D}(1,\xi)|\gg_{F,S,e}Q^{-e}.
\]
Both $\xi$ and $\xi^2$ have entire Hecke functions by the norm-character exclusion. For each $\delta>0$, their incomplete values and first derivatives are $O_{F,S,\delta}(Q^\delta)$ on the segment $1\le s\le2$. Indeed their conductors are $O_F(Q)$ up to a fixed factor from doubling or quadrupling the archimedean parameters. The absolutely convergent Euler bound on $\Re s=1+b$ and the functional equation with Stirling's formula on $\Re s=-b$ give the ordinary convexity bound $Q^{O_F(b)}(2+|\Im s|)^{O_F(1)}$ on $\Re s\ge1-2b$, $\Re s\le 1+b$. Taking $b>0$ small and using Cauchy's estimate on slightly smaller disks gives the derivative bound as well. The extra deleted polynomials cost $O_{F,S,\delta}(Q^\delta)$ there, after choosing $\delta$ smaller first. For fixed $b<1/8$, all sufficiently large $Nv$ satisfy $1+(Nv)^{-1+2b}\le (Nv)^\delta$, and the product of the distinct finite ramified norms divides the conductor norm.

For real $h>0$ the nonnegative trigonometric polynomial $3+4\cos u+\cos(2u)$ in the logarithms of the incomplete Euler products gives
\[
 1\le \zeta_{F,D}(1+h)^3
          |L_{F,D}(1+h,\xi)|^4|L_{F,D}(1+h,\xi^2)|.
\]
Since $\zeta_{F,D}(1+h)\ll_F h^{-1}$ for $0<h\le1$, we get $|L_{F,D}(1+h,\xi)|\gg h^{3/4}Q^{-\delta/4}$. Choose $h=cQ^{-8\delta}$ with $c>0$ sufficiently small in terms of $F,S,\delta$. The derivative error $O(hQ^\delta)$ is then at most half the lower bound because their quotient is $O(c^{1/4}Q^{-3\delta/4})$. Taking $\delta$ small proves the claimed estimate. The residue of $\zeta_{F,D}$ is the positive fixed field residue times $\prod_{v\in D}(1-1/Nv)\gg_{F,S,e}Q^{-e}$ by the same finite-product argument.

On the line $\Re s=1/2$ use \href{https://doi.org/10.1007/s00039-014-0277-4}{Wu, \emph{Burgess-like subconvex bounds for $\mathrm{GL}_2\times\mathrm{GL}_1$}}. The exponent of $Q$ for the fixed $\pi$ twist is $3/8+\theta/4+e$, where $\theta=7/64$ is available by \href{https://doi.org/10.4007/annals.2011.174.1.18}{Blomer and Brumley}. The estimate has at most polynomial loss in the vertical variable because an imaginary translation is absorbed into the Hecke character and increases its conductor by at most a fixed power of $2+|t|$. The single rank-one factor has convexity exponent $1/4+e$. All fixed omitted Euler factors are harmless there. Additional deleted factors for the abelian terms cost $Q^e$, since $1+(Nv)^{-1/2}\le(Nv)^e$ for sufficiently large $Nv$. When $v\notin S$ and $\chi$ ramifies, the twist of the unramified $\pi_v$ has trivial local factor. We may enlarge $S$ once to contain all primes ramified in $\pi$.

Shifting the two smoothed Mellin contours as before now gives, with $e>0$ as small as desired,
\[
   XQ^{-e}\ll X^{1/2}Q^{103/256+e}+X^{1-a}\log^3X .
\]
For $Q\le X^2$, choose $e<a/4$ first and absorb the last term for large $X$. Rearrangement, choosing $e$ also sufficiently small in terms of $\varepsilon$, proves the result. Larger $Q$ already satisfies the bound. At a complex place the two real gamma factors for weight $k\ell$ give an archimedean conductor comparable to $(1+|k\ell|)^2$. The assertion for powers follows.
\end{proof}

\subsection*{Exponential conductor separation from the second prime moment}

For bounded conductors, a sum of rank-one local traces has an arithmetically determined second moment. Its diagonal already exceeds the second moment of a single cuspidal representation. This yields a stronger separation throughout the range covered by the known rank-one zero boundaries.

\begin{theorem}\label{st:marginal}
Let $\pi$ be a fixed unitary cuspidal representation of $\mathrm{GL}_d(\mathbb A_F)$, $d\ge2$, tempered outside a fixed finite set $S$ containing its ramified places. Let $\chi_1,\dots,\chi_d$ be unitary Hecke characters of $F$. Suppose no quotient of two distinct characters in this list is a pure norm power, and set $H=\sum_{i=1}^d\log(2C(\chi_i))$. Let $1/2\le\theta<1$ be an all-height upper bound for the real parts of the nontrivial zeros of the nonprincipal quotients $L(s,\chi_i\bar\chi_j)$.
If, for fixed $0\le\eta<\sqrt d-1$,
\[
 \left|\lambda_\pi(v)-\sum_{i=1}^{d}\chi_i(v)\right|\le\eta
 \qquad (Nv\le X)
\]
at all places outside $S$ and the ramification of the characters, then for all sufficiently large $X$,
\[
                 H\gg_{F,\pi,d,S,\eta} X^{1-\theta}.
\]
Repeated characters are allowed. The reviewed $7/8$ all-height rank-one input in Section~\ref{recentcm:section} yields the exponent $1/8$ for Hecke characters over every $F$. The starting Dirichlet hypothesis gives exponent $1/2$ for finite-order characters over $\Q$. For a paired angular character $\chi,\bar\chi$ of unbounded integral weight $k\ell$ over an imaginary quadratic field, with $\ell\ne0$ fixed and bounded finite conductor, the former conclusion gives $\log(2+|k|)\gg X^{1/8}$.
\end{theorem}
\begin{proof}
Fix a nonnegative smooth weight $w$ of integral one, supported in $(1/2,1)$. Use the weighted average
\[
        I_X(g)=X^{-1}\sum_v \log(Nv)w(Nv/X)g(v)
\]
over the common unramified places outside $S$. Removing the places ramified in any character changes a bounded prime trace average by $O_{d,w}(H/X)$. By the prime ideal theorem, $I_X(1)=1+o_F(1)+O(H/X)$. The fixed Rankin--Selberg prime number theorem gives
\[
 I_X(|\lambda_\pi|^2)=1+o_{F,\pi,S}(1)+O_d(H/X).
\]
Indeed the unramified logarithmic derivative of $L(s,\pi\times\tilde\pi)$ has nonnegative prime-power coefficients $|\sum_j\alpha_{v,j}^k|^2\log Nv$. Its pole at one is simple and it is nonzero on the rest of that line. Wiener--Ikehara and partial summation give the smoothed main term $X$. The powers $k\ge2$ cost $O_{F,d,w}(\sqrt X)$ by temperedness and the elementary prime-ideal Chebyshev bound. Finitely many starting Euler factors may be deleted because their logarithmic derivatives are regular on $\Re s\ge1$ by the strict standard local bound for the Rankin--Selberg roots.

For a nonprincipal quotient $\xi=\chi_i\bar\chi_j$ the smoothed prime explicit formula gives $I_X(\xi)\ll_{F,d,w} X^{\theta-1}H+H/X$. To see uniformity for angular and imaginary gamma shifts, the local zero count in a unit interval at $t$ is $O_F(\log(2C(\xi))+\log(3+|t|))$ by the real Hadamard formula at $2+it$. The Mellin transform of $w$ decays rapidly in bounded strips, so all nontrivial zeros together cost $O_{F,w}(X^\theta\log(2C(\xi)))$ before division by $X$. The quotient conductor has logarithm $O_F(H)$. In the usual contour proof one may use the left line $\Re s=-1/4$. For a unitary Hecke character the real parts of its standard real gamma shifts are nonnegative half-integers, so this line stays separated from their poles; Stirling's formula and the dual Euler series bound its integral. Trivial zeros crossed at real part zero contribute $O_{F,w}(1)$. Deleting extra ramified primes and replacing prime powers by primes has the same costs as above. This is the angular version of \eqref{ld:smooth}. The pole of a pure norm character would give a term $X^{iu}\widehat w(1+iu)$, which explains the norm-quotient condition.

If the distinct actual characters in the list have multiplicities $m_j$, expansion thus gives
\[
 I_X\left(\left|\sum_i\chi_i\right|^2\right)
   =\sum_j m_j^2+o_{F,d}(1)+O_{F,d,w}(X^{\theta-1}H),
          \qquad \sum_jm_j^2\ge d .
\]
The triangle inequality in the weighted $L^2$ norm, and the assumed pointwise bound, force the square root on the left to be at most $\sqrt{I_X(|\lambda_\pi|^2)}+\eta\sqrt{I_X(1)}$. If $H\le cX^{1-\theta}$ for a sufficiently small fixed $c>0$, the lower bound tends to at least $\sqrt{d-O(c)}$ and the upper bound to at most $1+\eta$. The fixed positive gap proves the claim.
\end{proof}

The same calculation applies to the vertical shifts of Theorem~\ref{st:main}. The two off-diagonal terms are the norm characters with exponents $\pm2t$. Their normalized pole terms are $O_w(|\widehat w(1+2it)|)=o(1)$ whenever $t\to\infty$. Under zeta RH, any such sequence matching traces through $X$ within a fixed $\eta<\sqrt2-1$ therefore satisfies $\log(2+|t|)\gg X^{1/2}$. This uses the second prime moment of the fixed cusp form and makes no zero-location assumption on that form. In fact the matching heights must escape every bounded set even at this fixed tolerance. Otherwise a subsequential limit $t_*$ would satisfy the same trace inequalities at every good prime. When $t_*=0$ the abelian square is identically four. When $t_*\ne0$, the ordinary prime number theorem at this fixed frequency gives second moment $2+2\Re[X^{2it_*}\widehat w(1+2it_*)]+o(1)$. There are arbitrarily large real scales $X$ making this real part nonnegative. In either case the $L^2$ triangle inequality contradicts $\eta<\sqrt2-1$.

\begin{corollary}[A positive proportion of distinct prime traces]\label{st:misfit}
Let $E/\Q$ be fixed and non-CM. Suppose $X\to\infty$ and $\chi$ varies over primitive Dirichlet characters of conductor $q$ with $\log(2q)=o(X^{1-\theta})$, where the same zero boundary applies to all nonprincipal powers of $\chi$. If $\eta_X\to0$, the lower limit of the normalized smoothed prime weight at which
\[
              |\lambda_E(p)-\chi(p)-\bar\chi(p)|>\eta_X
\]
is at least $1/\pi$. If the character orders tend to infinity, uniform matching at all the tested primes requires an asymptotic error of at least
\[
              2\cos(\pi/3-\sqrt3/4)-1>0.63.
\]
The conclusion with vanishing tolerance applies to every sequence of character orders.
\end{corollary}
\begin{proof}
The Sato--Tate theorem for $E/\Q$ gives the distribution of $2\cos u$, $0\le u\le\pi$, with density $2\sin^2u/\pi$ for the Frobenius traces. Partial summation gives the same limit for $I_X$. For the standard statement see \href{https://math.mit.edu/~drew/TateConferenceUnanimated.pdf}{Sutherland, \emph{Sato--Tate distributions and murmurations}, Sato--Tate theorems}. If the character orders tend to infinity, each fixed nonzero Fourier power is eventually nonprincipal, so the smoothed explicit formula above and trigonometric approximation give uniform angular distribution for $\chi(p)$. Its trace limit is therefore $2\cos u$ for uniform $u\in[0,\pi]$.

Under this latter measure the set $|2\cos u|>\sqrt2$ has mass $1/2$. Under Sato--Tate it has mass $1/2-1/\pi$. If two traces match within $\eta_X$, membership of the former in this set forces the other to exceed $\sqrt2-\eta_X$ in modulus. Both limiting distributions have no boundary atoms. Subtracting the masses gives the stated proportion. For the second assertion compare $|\chi(p)+\bar\chi(p)|>1+\eta$ to $|\lambda_E(p)|>1$. The two limiting masses are $2\arccos((1+\eta)/2)/\pi$ and $2(\pi/3-\sqrt3/4)/\pi$, respectively. Their difference is positive for every $\eta$ below the displayed cutoff. Finally, for bounded character orders the abelian traces belong to a fixed finite set, and Sato--Tate gives weight tending to zero for matching to tolerance tending to zero. Every sequence of integer orders has a bounded subsequence or a subsequence tending to infinity. Applying this to a subsequence approaching the lower limit completes the proof. The same two limiting distributions apply to paired zeta shifts when $|t|\to\infty$ and $\log(2+|t|)=o(X^{1-\theta})$, since the pole contribution for each fixed nonzero Fourier power tends to zero by the rapid Mellin decay.
\end{proof}

\begin{corollary}[Dirichlet conductors for a target over any number field]\label{st:absolute-moment}
Let $n=[F:\Q]$ and let $\pi$ satisfy the assumptions of Theorem~\ref{st:marginal}. Write $A_p$ for the first trace in its degree-$dn$ expansion at a good rational prime $p$. Suppose a product of $dn$ Dirichlet functions of characters $\chi_i$, not necessarily of the same conductor, satisfies
\[
 \left|A_p-\sum_{i=1}^{dn}\chi_i(p)\right|\le\eta
\]
at all good unramified rational primes $p\le X$, with fixed $\eta<\sqrt n(\sqrt d-1)$. Under Dirichlet GRH alone,
\[
                 \sum_{i=1}^{dn}\log(2q_i)\gg_{F,\pi,\eta}\sqrt X ,
\]
where $q_i$ are the primitive conductors. More generally a common boundary $\theta$ for the Dirichlet quotients gives exponent $1-\theta$. This applies to an automorphic elliptic curve over every number field. It requires no automorphic induction of $\pi$ to $\Q$.

If $F/\Q$ is Galois and $\pi$ has finite-order central character, the same conclusion holds for each fixed $\eta<\sqrt{dn}-\sqrt h$, where $h$ is the number of automorphisms $g$ of $F$ for which $\pi^g\simeq\pi$.
\end{corollary}
\begin{proof}
A local factor at residue degree greater than one contributes zero to the coefficient of $p^{-s}$. Thus $A_p=\sum_{v\mid p,\ f_v=1}\lambda_\pi(v)$. Cauchy's inequality bounds $|A_p|^2$ by $n\sum_{v\mid p,\ f_v=1}|\lambda_\pi(v)|^2$. With the nonnegative smoothed prime weights used above, this gives
\[
         X^{-1}\sum_p w(p/X)\log p\,|A_p|^2\le n+o(1)
\]
by the single fixed Rankin--Selberg prime number theorem over $F$. Deleting additional ramified rational primes only reduces the left side. Over $\Q$ the abelian second moment is at least $dn+o(1)+O_{F,d}(X^{\theta-1}\sum_i\log(2q_i))$ by the calculation above. The same triangle inequality proves the result.

For the refinement, in a Galois field all primes of residue degree one above an unramified $p$ occur precisely when $p$ splits completely. Its degree-one places form a free transitive orbit of $G=\mathrm{Gal}(F/\Q)$. Expanding $|A_p|^2$ therefore sums once over each pair $(v,gv)$ with $v\mid p$ and $g\in G$. Prime ideals at residue degree greater than one have total weight $O_{F,d,w}(\sqrt X\log X)$ and may be inserted into the sums over $v$. For fixed $g$ the smoothed mean of $\Re[\lambda_\pi(v)\overline{\lambda_{\pi^g}(v)}]$ tends to $1$ when $\pi^g\simeq\pi$ and to zero otherwise. One can apply Wiener--Ikehara to the nonnegative self-convolution of the isobaric sum $\pi\boxplus\pi^g$ and subtract its two diagonal means. Its pole order at one is $2+2\mathbf1_{\pi^g\simeq\pi}$. Finite-order central characters exclude any distinct imaginary norm translate on the boundary, by the Rankin--Selberg pole criterion. The mixed prime-power terms are bounded by temperedness. The imaginary parts cancel in the full sum over pairs. It follows that the left-hand prime second moment before conductor deletions is $h+o(1)$, as asserted.
\end{proof}

\subsection*{The price of approaching both sides of a removed ray}

\begin{theorem}\label{st:thin}
In Proposition~\ref{st:jump} suppose $m\ge1$, and write $s_0=x+i\gamma$ for the regular point on the ray. There are $L,h_0,B>0$ depending on the fixed local data and the number of factors with the following property. Take $0<h<h_0$ and $0<\delta\le1/16$. Let $P$ be a product of the incomplete primitive nonprincipal Dirichlet functions in Theorem~\ref{st:characters} and $H=\sum_i\log(2q_i)$. Require its individual factors to be zero-free on a fixed compact neighborhood of the horizontal segment from $s_0-L$ to $2+i\gamma$. Write
\[
   R_h=\{z: |\Re z-x|\le L,\ |\Im z-\gamma|\le h\}.
\]
Suppose $|P/L_S-1|\le\delta$ on the compact subset of $\Omega$ formed by the top, bottom and right edges of $[x,2]+i[\gamma-h,\gamma+h]$, together with the closed $h/4$-disks centered on the two horizontal edges of $R_h$. Then
\[
 BH\ \ge\ \frac{4\delta}{(1-\delta)h}
    \left(\frac{\pi m(1-\delta)-\delta}{4\delta}\right)^{
               \cos(a)\cosh(aL/h)}\qquad (0<a<\pi/2).
\]
In particular, even at fixed error $\delta=1/16$ the necessary prime conductor $q$ of a common-conductor product grows at least triple exponentially in $1/h$. For some $c>0$ and all sufficiently small $h$ it satisfies $\log\log\log(2q)\ge c/h$. At error $\delta=h$ one has $\log\log(2q)\gg e^{aL/h}\log(1/h)$ for each fixed $a<\pi/2$.

The zero-free condition on the fixed compact holds eventually for the sequence constructed above, unconditionally. Under Dirichlet GRH it is automatic for every candidate product. For paired zeta shifts tending to infinity the same conclusions hold with $H=\log(2+|t|)$, provided the individual shifts are zero-free there.
\end{theorem}

\begin{proof}
Choose a closed rectangle about $s_0$ free of target zeros and contained in $\mathcal H$. Take $L$ smaller than half its horizontal half-width, and take $h_0$ small enough that the stipulated testing sets lie in $\Omega$. This is possible by finiteness of zero ordinates in the bounded region to their right, as in the proof of Theorem~\ref{st:main}. Choose a fixed compact collar of sufficient thickness around the entire mentioned segment, still inside $\mathcal H$, for the assumed zero-freeness. Let
\[
             D=P'/P-L_S'/L_S .
\]
Then $D$ is holomorphic on a fixed neighborhood of every $R_h$, where $|D|\le BH$. We verify the uniformity of this last bound without assuming GRH. Subtract the complex Hadamard logarithmic derivative at $2+i\Im s$ from that at $s$ for each primitive character. The reciprocal differences are absolutely summable. For any zero $\rho$ the collar gives $|s-\rho|\ge c_0>0$ on the smaller rectangle, whence
\[
 \left|\frac1{s-\rho}-\frac1{2+i\Im s-\rho}\right|
     \ll_{c_0,x,L} |2+i\Im s-\rho|^{-2}.
\]
The sum on the right is $O(\log(2q_i))$ there by the unconditional real Hadamard formula at real part two, since $2-\Re\rho\ge1$. The two bounded-height gamma terms, the Euler derivative at two and the fixed omitted factors obey the same bound. The target logarithmic derivative is bounded on the local zero-free rectangle. This proves $|D|\le BH$. For the zeta shifts the same comparison applies with translated heights, gamma terms of order $O(\log(2+|t|))$, and the eventually distant poles.

On each horizontal edge of $R_h$, Cauchy's formula on the surrounding $h/4$-disk gives $|D|\le b_h=4\delta/[(1-\delta)h]$, since $D$ is the logarithmic derivative of $P/L_S$. On the three tested sides of the larger rectangle the principal branch of $\log(P/L_S)$ has modulus at most $\delta/(1-\delta)$. The argument principle, with the upward orientation on the left edge, thus gives
\[
  \max_{|u|\le h}|D(s_0+iu)|
       \ge [\pi m-\delta/(1-\delta)]/h .
\]
Indeed the difference between the upward integral of $D\,ds$ and $2\pi i m$ is the change of that principal logarithm on the other three sides. The ratio of the last lower bound to $b_h$ exceeds one.

For $z=s_0+y+iu$ introduce the harmonic function
\[
 \Phi(y,u)=\frac{\cosh(ay/h)\cos(au/h)}
                         {\cosh(aL/h)\cos a}.
\]
It is nonnegative throughout $R_h$ and at least one on both vertical sides. On the center line it is at most $[\cosh(aL/h)\cos a]^{-1}$. The subharmonic maximum principle for $\log |D|$ gives
\[
      \log(|D|/b_h)\le \Phi\log(BH/b_h).
\]
Here $BH>b_h$ by the center-line lower bound. Evaluating there and rearranging proves the inequality. The displayed asymptotics follow by taking logarithms. The Helson product of Example~\ref{st:helson} supplies actual target zeros for which the construction can meet all these finite tests at each $h$, with correspondingly large conductors.
\end{proof}

\begin{theorem}[Matching upper and lower orders]\label{st:optimal}
Use the test set and the fixed zero-free collar of Theorem~\ref{st:thin}. Let $q_{\min}(h)$ be the least prime conductor of a product as there, with the prescribed degree and relative error $1/16$ on that test set. Let $Y_{\min}(h)$ be the least cutoff for the same approximation by a finite Euler product with the prescribed degree and arbitrary unit roots at the good rational primes. For all sufficiently small $h$,
\[
 \log\log Y_{\min}(h)\asymp h^{-1},\qquad
 \log\log\log q_{\min}(h)\asymp h^{-1}.
\]
The constants may depend on the fixed target, point and collar. The same orders hold for relative error $h^b$ with any fixed $b>0$. For an elliptic target over $\Q$ the Dirichlet product may be restricted to a conjugate pair. Under Dirichlet GRH the individual zero-free collar requirement in the conductor minimum is automatic. The theorem is unconditional with that requirement as stated.
\end{theorem}

\begin{proof}
For a finite unit-root Euler product the local logarithmic derivative on the fixed rectangle is at most $O_{d,x,L}(1+Y^2)$ by absolute summation over $p\le Y$. There are no poles or zeros in $\mathcal H$. Repeating Theorem~\ref{st:thin} with this upper bound in place of $BH$ gives $\log\log Y\gg1/h$. The lower bound for the conductor was proved there. We construct products achieving both upper bounds.

Denote the tested compact set by $K_h$ and write $F=\log L_S$ for the normalized horizontal branch. Choose a fixed $B\ge3$ so that $K_h+B$ stays strictly to the right of one. For all small $h$, the closed $c_0h$-neighborhood of $K_h+[0,B]$ lies in $\Omega$ for some $c_0>0$. Indeed along the upper and lower portions the ordinate stays at least $3h/4$ from $\gamma$, and on the connecting right edge the real part is two or larger. All other zero ordinates of relevance have fixed positive separation. On that tube
\[
                         |F|\ll 1+\log(1/h).
\]
To verify this, factor out the finitely many zeros in a fixed somewhat larger rectangle inside $\mathcal H$, choosing its boundary free of zeros. The remaining nonvanishing holomorphic factor has a bounded logarithm on the smaller rectangle. Use the horizontal branch of each $\log(s-\rho)$, whose imaginary part is bounded by $\pi$ and whose real part is $O(1+\log(1/h))$ on the tube. The branches agree with $F$ up to a fixed constant by continuation from the right edge.

Set $A=A_0/h$, $c=1-e^{-A}$ and $R=1+e^{-A}/2$, with $A_0$ sufficiently large and fixed. Put
\[
       \phi(z)=-\frac{\log(1-cz)}{A}.
\]
It has nonnegative Taylor coefficients, $\phi(0)=0$ and $\phi(1)=1$. On $|z|\le R$ the inequalities $1-cR\ge e^{-A}/2$, $1+cR<2$, and $|\arg(1-cz)|\le\pi/2$ give
\[
      -\frac{\log2}{A}\le\Re\phi(z)\le1+\frac{\log2}{A},
      \qquad |\Im\phi(z)|\le\frac{\pi}{2A}.
\]
Thus $G_s(z)=F(s+B-B\phi(z))$ stays on the indicated tube for $s\in K_h$. In its Taylor series at zero, evaluation at $z=1$ followed by truncation to degree $N$ makes an error at most $\sup|G_s|R^{-N}/(R-1)$, uniformly in $s$.

Near zero all Euler logarithms for $G_s$ converge absolutely. For a frequency $y=\log p^k$, taking the first $N$ Taylor coefficients and evaluating at one multiplies the original $p^{-ks}$ by
\[
 W_N(y)=e^{-By}\sum_{j=0}^N \frac{(By/A)_j}{j!}c^j ,
 \qquad (v)_j=v(v+1)\cdots(v+j-1),\quad(v)_0=1 .
\]
Indeed $\exp(By\phi(z))=(1-cz)^{-By/A}$. All coefficients are nonnegative and their full sum at $z=1$ equals $e^{By}$, hence $0\le W_N(y)\le1$. Also
\[
       1-W_N(y)\le R^{-N}\exp(By(\phi(R)-1)).
\]
This follows by multiplying the omitted nonnegative coefficients by $R^j$. In particular $W_N$ can be made arbitrarily close to one up to the prescribed low frequencies at the same time as the approximation of $F$.

Here are the quantitative choices for a relative tolerance $\varepsilon_h$ equal to $1/16$ or $h^b$. Work in a fixed bounded neighborhood of all the $K_h$ with $\Re s\ge\sigma_*>1/2$ and fixed-radius disks for the variance estimate of Theorem~\ref{st:main}. Fix all local roots through $X_0$ with $X_0^{1/2-\sigma_*}\ll\varepsilon_h$ with a sufficiently small constant. For $0<\kappa<\sigma_*-1/2$ the sum of $p^{-k\sigma_*}/k$ over $k\ge2$, $p^k>Y_0$, is $O_{\kappa,\sigma_*}(Y_0^{-\kappa})$, by extracting $(p^k)^{-\kappa}$. Choose $Y_0\ge X_0$ so that this is sufficiently small compared with $\varepsilon_h$ and require $1-W_N(y)\le\nu$ for $y\le\log Y_0$, with $\nu\ll\varepsilon_h/(1+X_0)$. All of $X_0,Y_0,1/\nu,1/\varepsilon_h$ can be taken at most fixed powers of $1/h$. Since $\phi(R)-1\le\log2/A$, the two preceding Taylor estimates hold with errors as small as required if
\[
             N=\left\lceil C\exp(A)(A+\log(2/h))\right\rceil,
                   \qquad \log(N+2)\ll 1/h
\]
for a sufficiently large constant $C$.

For a sharper cutoff evaluate the nonnegative power series at $z=1/2$. Since $2^{-j}\ge2^{-N}$ for $j\le N$,
\[
 W_N(y)\le 2^N\exp[-By(1-\phi(1/2))]
             \le 2^N\exp(-By/2)
\]
when $A\ge2\log2$. For $y\ge C_1N$ with a sufficiently large constant the right side is at most $e^{-By/4}$. Increase $B$ at the start to at least $8$ and take $Y=\lceil\exp(C_1N)\rceil$. The mean first-power tail beyond $Y$ is then $O(\sum_{n>Y}n^{-\sigma_*-B/4})$. The higher-power tail is bounded absolutely as above. Apply the independent phase construction of Theorem~\ref{st:main} or \ref{st:characters} with probability $W_N(\log p)$ of retaining each true phase. Its variance after $X_0$ is $O(X_0^{1-2\sigma_*})$ on the fixed neighborhood. The finite phases can therefore be selected so that the Euler logarithm differs from $F$ by at most $\varepsilon_h/4$ on $K_h$, by choosing the constants in the previous estimates sufficiently small. This is relative approximation after exponentiating and gives
\[
                          \log\log Y\ll 1/h .
\]

Finally apply the quantitative phase selection of Theorem~\ref{ft:select} on a \emph{fixed} compact containing the required zero-free collar as well as every $K_h$. Use angular mesh $v=c_1\varepsilon_h/Y$ at all good $p\le Y$ and any auxiliary phases at the finitely many omitted primes. For the finite logarithmic Euler products the induced error is $O_{d,\sigma_*}(Yv)$ by summing the local Lipschitz bounds. The remaining Dirichlet tails are uniformly $1+O(Y^{1/2-\sigma_*})$ on this whole compact. In particular every individual Dirichlet factor is zero-free on the collar since its finite unit-root product never vanishes. Choose $c_1$ sufficiently small. Every sufficiently large prime with $\log q\ge C Yv^{-1}\log(2\pi(Y))$ works by that theorem, separately in each coordinate at the same $q$. Bertrand's postulate supplies such a prime with $\log q\ll Y^2\varepsilon_h^{-1}\log(2Y)$. Hence $\log\log\log q\ll1/h$, as claimed.
\end{proof}

\begin{corollary}[Polynomial cutoff on a fixed zero-free corridor]\label{st:polynomial}
Fix a compact $K\Subset\Omega$ and a second compact $E_0\Subset\mathcal H$. Let $X\ge2$, $0<\varepsilon,\nu\le1/4$. There is a finite unit-root Euler product through
\[
              Y\ll_{K,L_S,d} (2+X+\varepsilon^{-1})^{C_{K,L_S,d}}
\]
which agrees exactly with all the target roots at the good primes through $X$ and approximates $L_S$ on $K$ to relative error $\varepsilon$. A product of the specified number of actual primitive Dirichlet functions, individually zero-free on $E_0$, attains that same error on $K$ and root error $\nu$ at primes through $X$ at every prime conductor satisfying
\[
      \log q\ge C_{K,E_0,L_S,d}
                  (2+X+\varepsilon^{-1}+\nu^{-1})^{C_{K,E_0,L_S,d}} .
\]
For an elliptic target over $\Q$ a conjugate pair suffices. The exponents are fixed with the compact sets and the target, independently of all three accuracy parameters.
\end{corollary}
\begin{proof}
In the construction above, $K+[0,B]$ now has a fixed positive distance from the complement of $\Omega$, so $A$ and $R>1$ can be fixed and $\sup|G_s|$ is a constant. Choose $X_0\ge X$ with $X_0^{1/2-\sigma_*}$ sufficiently small compared with $\varepsilon$, then choose $Y_0$ for the higher powers and impose the finite-weight tolerance $O(\varepsilon/(1+X_0))$. These choices, their reciprocals when appropriate, and $\exp(B(\phi(R)-1)\log Y_0)$ are all bounded by fixed powers of $2+X+\varepsilon^{-1}$. Both Taylor errors therefore have the required size at
\[
                    N\ll_{K,L_S,d}\log(2+X+\varepsilon^{-1}).
\]
The bound $W_N(y)\le2^N e^{-By/2}$ permits $Y=\lceil\exp(C_1N)\rceil$, enlarging $N$ within the same estimate to ensure $Y\ge Y_0$ and a sufficiently small discarded tail. The mean, variance, and exponentiation arguments prove the first statement. In Theorem~\ref{ft:select} use a fixed compact containing $K\cup E_0$ and mesh a sufficiently small constant times $\min(\nu,\varepsilon/Y)$. Its conductor condition is implied by the displayed polynomial bound. Its rough tails, uniformly within $O(Y^{1/2-\sigma_*})$ of one, ensure zero-freeness on $E_0$ after a final enlargement of $N$ within the same estimate. This proves the corollary.
\end{proof}

Thus, at a regular point with a zero directly to its right, fixed relative accuracy has a minimum conductor whose third iterated logarithm grows as $1/h$. At an ordinate without a zero to the right of the test interval, all the tests for small $h$ lie on one fixed compact subset of $\Omega$. A single conductor then suffices for fixed accuracy, and for accuracy $h^b$ the logarithm of a sufficient conductor is polynomial in $1/h$. These statements quantify the analytic loss caused by an off-line zero. Its absence for a general elliptic function remains open.

For example, take the fixed unitary rectangle $3/5\le\Re s\le2$, $|\Im s|\le255$. The height-$256$ census in Theorem~\ref{cens:thousand} and Section~\ref{cens:chapter} puts this rectangle in $\Omega$ for each of the $2463$ rational elliptic isogeny classes of conductor at most $1000$, and for the $1846$ field/isogeny pairs over the $100$ fields verified there. Corollary~\ref{st:polynomial} thus gives one polynomial conductor threshold for this entire finite list by taking the largest of the finitely many constants and exponents. Above the threshold the same prime conductor works for every target, with a separate tuple of characters for each. On $|\Im s|\le63$ the same statement applies to all $38\,042$ rational classes of conductor at most $9999$ by Theorem~\ref{cens:9999}. The verification of critical zeros supplies the required horizontal corridors without an all-height elliptic assumption.

\clearpage
\section{Prime-conductor characters with prescribed Euler factors}\label{ft:main}

For every sufficiently large prime conductor we select nonreal primitive characters whose values match arbitrary prescribed phases at the small primes. The selection also controls the remaining Dirichlet function throughout a compact part of the open half-plane $\mathcal H=\{s:\Re s>1/2\}$. This supplies a conductor bound in the approximation problem of Section~\ref{vt:section}. For an elliptic curve, orthogonality of the phase errors permits the prescribed accuracy to tend to zero arbitrarily slowly.

Put $r=\pi(X)$, $P=\prod_{p\le X}p$, and for a character $\chi$ write
\[
 Z_{\chi,>X}(s)=L(s,\chi)\prod_{p\le X}(1-\chi(p)p^{-s}).
\]
Take $X\ge2$ and $0<h\le1$. For a compact set $K\subset\mathcal H$ let $K^*=\{\bar s:s\in K\}$. All conductors below are prime.

\begin{theorem}[Quantitative phase selection]\label{ft:select}
For each compact $K\subset\mathcal H$ there exists a constant $C_K$ such that, if
\begin{equation}
       \log q\ge C_K X h^{-1}\log(2\pi(X)),\label{ft:budget}
\end{equation}
then every choice $|\alpha_p|=1$ at $p\le X$ admits a nonreal primitive $\chi\pmod q$ satisfying
\begin{align}
 |\chi(p)-\alpha_p|&<h\qquad(p\le X),\label{ft:phase}\\
 |Z_{\chi,>X}(s)-1|+|Z_{\bar\chi,>X}(s)-1|
 &\ll_K X^{1/2-\Re s}/\sqrt{\log(2X)}
                         \qquad(s\in K).\label{ft:rough}
\end{align}
One may further specify any subset $J$ of the primes up to $X$. The same character may be required to satisfy
\begin{equation}
 \sup_{s\in K}\left|\sum_{p\in J}
 [\chi(p)+\bar\chi(p)-b(\alpha_p+\bar\alpha_p)]p^{-s}\right|
  \ll_K h,\qquad b\in\R,\quad 0\le1-b\ll h^2,\label{ft:centered}
\end{equation}
where $b$ depends only on $X,h$. These selections are unconditional.
\end{theorem}

\subsection{A finite measure on characters}

Write $S_D(z)=1+z+\cdots+z^{D-1}$ and set $A_j=\int_{|z|=1}|S_D(z)|^{2j}dm(z)$ for normalized Haar measure $m$. Let $K_j(z)=|S_D(z)|^{2j}/A_j$. It is a nonnegative Laurent polynomial of degree $j(D-1)$ and constant term one. Write $S_D(z)^j=\sum_{v=0}^{j(D-1)}a_vz^v$. The $a_v$ are nonnegative and $\sum a_v=D^j$. Cauchy--Schwarz, Parseval and $|S_D|\le D$ give, for $j\ge1$,
\[
 \frac{D^{2j}}{jD+1}\le\sum a_v^2=A_j\le D^{2j-2}A_1=D^{2j-1}.
\]
Choose $D=\lceil4/h\rceil$, $l=\lceil8\log(2r)\rceil$, $k=2l$, and $d=k(D-1)$. With $z_p=\chi(p)/\alpha_p$, the weight
\[
 W(\chi)=\prod_{p\le X}K_k(z_p)
 \quad\hbox{satisfies}\quad 0\le W\le (kD+1)^r=:B .
\]
Use the uniform average $\mathbb E$ over the $q-1$ characters. If $q>P^d$, then $\mathbb E W=1$. Indeed, a Laurent monomial of multi-degree at most $d$ in absolute value gives a character of the ratio of two positive integers at most $P^d$. It survives the average exactly when their residues are equal. They then agree as integers, and unique factorization forces every Laurent exponent to vanish. The same argument applies to any product of one-variable polynomials with degree at most $d$ in each variable. With $q>P^{d+2}$ it applies also when the weight is multiplied by linear or quadratic monomials in the phases and their conjugates. In these finite calculations averaging gives exactly the corresponding integral on the product torus. Write $\mathbb E_W V=\mathbb E(WV)$.

For $|z-1|\ge h$, the geometric series gives $|S_D(z)|\le2/h$. Comparing two powers yields
\[
  \mathbf1_{|z-1|\ge h}K_{2l}(z)
   \le (2/h)^{2l}\frac{A_l}{A_{2l}}K_l(z)
   \le (2l+1)\left(\frac{2}{Dh}\right)^{2l}K_l(z).
\]
Replacing a single kernel by $K_l$ leaves the mass of the product equal to one. Thus the total weighted probability of failing any phase condition is at most $r(2l+1)4^{-l}<1/8$. Comparison with the lower-power kernel is what makes $k$ independent of $h$.

\subsection{A rough Dirichlet polynomial at finite conductor}

Choose a bounded open rectangle $U$ containing $K\cup K^*$ with closure in $\mathcal H$. Write $a=\inf_U\Re s>1/2$, $A=\sup_U\Re s$ and fix $u$ with $1/(2a)<u<1$. Let $R=\lfloor q^u\rfloor$. Increasing $C_K$ in \eqref{ft:budget} assures
\[
 q>P^{d+2},\qquad RP^d<q,\qquad RP<q.
\]
In fact $\log P\ll X$, $d\ll h^{-1}\log(2r)$ and $1-u>0$ is fixed. Throughout the proof $C_K$ may be increased uniformly in $X,h$.

For every nonprincipal character and $s\in U$, P\'olya--Vinogradov and partial summation give
\[
 L(s,\chi)=\sum_{n\le R}\frac{\chi(n)}{n^s}
     +O_K(\sqrt q\log(2q)\,R^{-\Re s}).
\]
The Dirichlet series converges on $\Re s>0$ by boundedness of the character partial sums. Integrating their $O(\sqrt q\log(2q))$ bound against the derivative of $t^{-s}$ proves the estimate uniformly on $U$.
See \href{https://doi.org/10.1017/CBO9780511618314}{Montgomery and Vaughan, \emph{Multiplicative Number Theory I}, Theorem~9.18}.

Multiply the finite sum by $\prod_{p\le X}(1-\chi(p)p^{-s})$. The coefficient of $\chi(n)n^{-s}$ for $n\le R$ is $\sum_{b\mid(n,P)}\mu(b)=\mathbf1_{(n,P)=1}$. Above $R$ it is a partial sum over divisors of $P$ and has modulus at most $2^r$. Write the product as $1+V(s,\chi)+H(s,\chi)$, with $V$ the sum over rough $1<n\le R$ and $H$ supported on $R<n\le RP$. Orthogonality gives exactly
\begin{equation}
 \mathbb E_W|V(s,\chi)|^2
  =\sum_{\substack{1<n\le R\\(n,P)=1}}n^{-2\Re s}.
 \label{ft:diag}
\end{equation}
A monomial from $W$ contributes a congruence $mv\equiv nw\pmod q$ for two rough integers $m,n\le R$ and two $P$-smooth integers $v,w\le P^d$. By $RP^d<q$ there is equality of integers. The parts prime to $P$ force $m=n$ and then $v=w$. This is the constant coefficient of $W$.

For $\sigma\in[a,A]$ the sum in \eqref{ft:diag} is
$\ll_K X^{1-2\sigma}/\log(2X)$. The rough primes are bounded by summing $p^{-2\sigma}$ on successive dyadic intervals above $X$ using Chebyshev's bound for $\pi(y)$. A rough composite exceeds $X^2$, and the sum over all such integers is $O_a(X^{2-4\sigma})$, which satisfies the claimed bound since $a>1/2$. For $H$ use unweighted orthogonality first, since $RP<q$, and then $W\le B$ to obtain
\[
 \mathbb E_W|H(s,\chi)|^2\le B4^r\sum_{n>R}n^{-2\sigma}
                         \ll_K B4^r R^{1-2\sigma}.
\]
On nonprincipal characters multiplication of the P\'olya--Vinogradov error by the finite Euler factors costs at most another $2^r$ in absolute value. Both errors are negligible at the required scale. More explicitly, put $v=\min(u(2a-1),2ua-1,1)>0$. We have
$\log(B4^r)\ll r\log(2kD)\ll Xh^{-1}\log(2r)$, as also $\log(2X)\ll Xh^{-1}\log(2r)$. Increase $C_K$ so that the factors $B4^r$, $X^{2A+2}$ and $\log^2(2q)$ together are at most $q^{v/2}$. The same bound then absorbs the two real characters, whose combined weighted mass is at most $2B/(q-1)<1/8$. In particular, uniformly on $U$,
\begin{equation}
 \mathbb E_W\left[\mathbf1_{\chi\ {\rm nonreal}}
       |X^{s-1/2}(Z_{\chi,>X}(s)-1)|^2\right]\ll_K 1/\log(2X).
 \label{ft:rougharea}
\end{equation}
Only the short-polynomial orthogonality calculation ever includes the principal character's polynomial. Its meromorphic $L$-function does not enter the analytic estimate.

\subsection{Controlling the errors of many matched phases}

For $|z|=1$ with law $K_k\,dm$, put $b=\int zK_k\,dm$. It is real by symmetry. Since $(1-z)S_D(z)=1-z^D$,
\[
 2(1-b)=\int|1-z|^2K_k\,dm\le4A_{k-1}/A_k
 \le4 A_1/A_2\le 6/D^2 .
\]
The ratios $A_{j-1}/A_j$ decrease by Cauchy--Schwarz. Here $A_1=D$ and $A_2=(2D^3+D)/3$ by summing the squares of the triangular coefficients of $S_D^2$. In particular $\int|z-b|^2K_k\,dm=1-b^2\le2(1-b)$.

The real random variables
$Y_p=\chi(p)+\bar\chi(p)-b(\alpha_p+\bar\alpha_p)$ are centered in the weighted average, have second moments $O(h^2)$ and have zero covariance for distinct primes. Indeed $Y_p=2\Re(\alpha_p(z_p-b))$, and all one- and two-coordinate moments just used follow exactly from $q>P^{d+2}$. Therefore, with $F_J(s)=\sum_{p\in J}Y_p p^{-s}$,
\[
  \mathbb E_W|F_J(s)|^2\ll h^2\sum_p p^{-2a}\ll_K h^2\qquad (s\in U).
\]
Integrate this bound and \eqref{ft:rougharea} over $U$. With a fixed sufficiently large factor depending only on $U$, Markov's inequality makes the weighted measure exceeding either area bound less than $1/8$ for each bound. The exceptional sets for the phases and the real characters also have mass less than $1/8$ each. Thus the intersection of the desired events has mass at least $1/2$. A nonreal character modulo a prime is automatically primitive. In fact the number of successful characters is at least $(q-1)/(2B)$.

Now disks of a fixed radius about every $s\in K\cup K^*$ lie in $U$. The area submean inequality for the absolute square of a holomorphic function bounds its value everywhere on these compacts by a constant times its integral over $U$. Apply this both to $F_J$ and to the holomorphic function $X^{s-1/2}(Z_{\chi,>X}(s)-1)$. The latter choice preserves the exponent $1/2-\Re s$ at each target point. Finally $Z_{\bar\chi,>X}(s)=\overline{Z_{\chi,>X}(\bar s)}$ controls both rough factors of the same character. This proves Theorem~\ref{ft:select}.

\subsection{The elliptic function and local boundedness}

Let $E/\Q$ be fixed, $S$ the finite set of its bad primes, and $F_E(s)$ its unitary $L$-function with the Euler factors at $S$ omitted. At a good prime choose one normalized Frobenius eigenvalue $\alpha_p$; the other is $\bar\alpha_p$. For $X\ge\max(S\cup\{2\})$, choose arbitrary auxiliary unit phases at the primes of $S$ when applying Theorem~\ref{ft:select}, and take $J=\{p\le X:p\notin S\}$. Write
\[
 L_S(s,\chi)=L(s,\chi)\prod_{p\in S}(1-\chi(p)p^{-s}),\qquad
 D_\chi(s)=L_S(s,\chi)L_S(s,\bar\chi).
\]
The auxiliary phases at bad primes only serve to define the finite probability measure. These bad factors have been removed from both sides of the intended approximation. The missing Euler factor of each Dirichlet function at its conductor is included automatically in $Z_{\chi,>X}$.

\begin{theorem}\label{ft:normal}
Suppose the Dirichlet GRH input of this manuscript holds. Then ERH for $E$ is equivalent to the existence of prime conductors $q_j\to\infty$, nonreal characters $\chi_j\pmod {q_j}$ and cutoffs $X_j\to\infty$ matching the good phases to errors $h_j\to0$ such that $D_{\chi_j}$ is locally bounded on $\mathcal H$. In this case the products converge locally uniformly to $F_E$.

Under ERH for $E$, at every fixed compact nonempty $K\subset\mathcal H$ one may use \emph{any} prime $q$ satisfying \eqref{ft:budget} and a character as in Theorem~\ref{ft:select}, selected with the conditions of that theorem, to obtain
\begin{equation}
 \sup_K|D_\chi-F_E|\ll_{E,K}h+X^{1/2-\sigma_K}\log(2X),
 \qquad \sigma_K=\min_{s\in K}\Re s .\label{ft:ellerror}
\end{equation}
Thus if a conductor is to be chosen as well, it may be taken with $\log q\ll_K Xh^{-1}\log(2\pi(X))$.
\end{theorem}
\begin{proof}
For the reverse implication take a locally bounded matching sequence. On each half-plane $\Re s\ge1+\delta$, the Euler logarithms of the remaining factors are absolutely convergent. The difference for good $p\le X_j$ is $O_\delta(h_j)$ by $|u^m-v^m|\le m|u-v|$ on the unit circle, and the combined tails contribute $O_\delta(X_j^{-\delta})$. If a conductor prime occurred in the matched good prefix, its character value would be zero and could not meet a tolerance smaller than one. The right-hand agreement therefore holds for every sufficiently large $j$. Each $D_{\chi_j}$ is holomorphic and, by the Dirichlet input, zero-free on $\mathcal H$; the finite removed factors are zero-free on $\Re s>0$. By Montel, every subsequence has a locally uniformly convergent further subsequence. Its limit agrees with $F_E$ on $\Re s>1$, hence on $\mathcal H$. Hurwitz gives zero-freeness, since $F_E$ is not identically zero. Uniqueness of all such limits also gives convergence of the original locally bounded sequence. Modularity, the functional equation and the nonvanishing of the omitted factors on $\mathcal H$ equate this right-half-plane statement to ERH for $E$.

For the constructive direction let $P_X^E$ be the product of the target factors at good primes $p\le X$. Under ERH for $E$, Section~\ref{vt:section}, equation~\eqref{vt:deterministic}, proves
$\sup_K|P_X^E-F_E|\ll_{E,K}X^{1/2-\sigma_K}\log(2X)$ from the signed prime estimate. In particular the first-power sums $\sum_{p\in J}(\alpha_p+\bar\alpha_p)p^{-s}$ are bounded on $K$, uniformly in $X$, by the same partial summation or by subtracting the absolutely convergent higher powers from the convergent Euler logarithms.

In the difference of the two finite Euler logarithms, the first powers give the sum in \eqref{ft:centered} plus $(b-1)\sum_{p\in J}(\alpha_p+\bar\alpha_p)p^{-s}$. This is $O_{E,K}(h)$ even when the sum of the absolute first-power differences would grow with $X$. The remaining powers cost
\[
 2h\sum_p\sum_{m\ge2}p^{-m\sigma_K}\ll_K h
\]
using the actual pointwise phase matching. The ratio of the two finite products is consequently $1+O_{E,K}(h)$. Multiplication by the two rough factors in \eqref{ft:rough}, and boundedness of $P_X^E$ on $K$, prove \eqref{ft:ellerror}. Bertrand's postulate supplies a prime between the exponential of the right-hand threshold in \eqref{ft:budget} and twice that value, after rounding. On an increasing compact exhaustion, select $X_j$ so that the error at the $j$th compact tends to zero, then take $h_j$ small there and an eligible prime conductor larger than the preceding one. This gives the locally convergent sequence. The proof of this construction under elliptic ERH uses no Dirichlet zero hypothesis. That input is used to guarantee that every approximant in the reverse implication is zero-free.
\end{proof}

\subsection{Explicit failure of a uniform size bound from phase selection alone}

Take $\alpha_p=1$ at every prime, $K=\{\sigma\}$ and $1/2<\sigma<1$. For fixed $0<h<1$, \eqref{ft:phase} implies $\Re\chi(p)>1-h^2/2$. Since the prime powers beyond the first are absolutely summable at this $\sigma$, \eqref{ft:rough} implies for the selected nonreal character
\[
 \log |L(\sigma,\chi)|
 \ge (1-h^2/2)\sum_{p\le X}p^{-\sigma}-O_\sigma(1)
 \gg_{\sigma,h} \frac{X^{1-\sigma}}{\log X}
\]
for sufficiently large $X$, by the prime number theorem. For every sufficiently large prime $q$ one may take $X=c_{\sigma,h}\log q/\log\log q$ with a small positive constant. The lower bound then has order
$(\log q)^{1-\sigma}/(\log\log q)^{2-\sigma}$. With positive logarithmic growth there is no local bound for the family of all selected Dirichlet functions in this example. At these real arguments the paired product is $|L(\sigma,\chi)|^2$. These prescribed phases illustrate the role of the target's signed first-power sum and are not asserted to be elliptic Frobenius eigenvalues.

At $\sigma=1$ take $X=\log q/(\log\log q)^2$ and $h=1/\log\log X$ for sufficiently large $q$. They satisfy \eqref{ft:budget}. From $1-\Re(\chi(p)^m)\le m^2h^2/2$, the difference between the logarithms of the absolute small-prime product and the Mertens product is
\[
 \ll h^2\sum_{p\le X}\frac{p^{-1}}{(1-p^{-1})^2}
       \ll h^2\log\log X=o(1).
\]
Mertens' product theorem and \eqref{ft:rough} thus give
$|L(1,\chi)|\ge(e^\gamma-o(1))\log\log q$. Stronger extremal-value estimates for unrestricted nonprincipal characters follow from resonance; see
\href{https://arxiv.org/abs/1803.00760}{Aistleitner, Mahatab, Munsch and Peyrot, \emph{On large values of $L(\sigma,\chi)$}}. The purpose here is to check what the prescribed phases and the controlled remaining factors can force in the same construction.

\clearpage
\phantomsection\addcontentsline{toc}{part}{II. All-height bounds and arithmetic transfers}
\noindent{\small\sffamily\color{accent}II. All-height bounds and arithmetic transfers}\par
\section{A zero-free region for an entire symmetric-power family}
\label{zf:chapter}
A zero near the boundary of the critical strip contributes a large term to the logarithmic derivative just outside the strip. To turn this observation into a proof one needs a positive prime sum and a budget for the poles. For the symmetric powers of an elliptic curve, three appropriately chosen constituents give a budget of three against a zero of multiplicity at least four. This section gives the argument in both the non-CM and CM cases. Standard zero-free regions of this kind have a substantial history. In particular, J.~Thorner proves a more general theorem for twisted pairs of symmetric powers over totally real fields in \emph{Exceptional zeros of Rankin--Selberg $L$-functions and joint Sato--Tate distributions}, IMRN (2025), \href{https://arxiv.org/abs/2404.06482}{arXiv:2404.06482}. We include the particular three-constituent calculation because its costs and its limit can be read off from the same proof.

Let $E/\Q$ have conductor $N$. At a good prime, let $g_p$ be the conjugacy class in $\mathrm{SU}(2)$ with trace $a_p/\sqrt p$. Write $V_j$ for the $j$th symmetric power of the defining representation, and $u_j(g)$ for its trace. Thus $u_0=1$, $u_1=\operatorname{tr}g$, and all $u_j$ are real. The good Euler factor of the normalized $L_j$ is $\det(1-V_j(g_p)p^{-s})^{-1}$. Bad factors are the primitive local factors, with the inertia and monodromy conventions of the corresponding Weil--Deligne representation. In particular $L_1(s)=L(E,s+1/2)$. We write $L_0=\zeta$.

\begin{theorem}\label{zf:noncm}
There is an effective absolute constant $b>0$ with the following property. For every non-CM $E/\Q$, every $m\ge1$ and every real $t$, $L_m(\sigma+it)$ has no zero for
\[
 \sigma>1-\frac{b}{(m+1)^2\log\!\big(2N(m+2)(3+|t|)\big)} .
\]
There is no exceptional real zero in this region. Reflection gives a zero-free region of the same width at the other edge of the strip, for the completed function.
\end{theorem}

\paragraph{Analytic inputs.}
Newton and Thorne, \emph{Symmetric power functoriality for holomorphic modular forms, II}, Publ. Math. IH\'ES \textbf{134} (2021), \href{https://arxiv.org/abs/2009.07180}{arXiv:2009.07180}, provide cuspidal automorphy of every $L_j$, $j\ge1$, for a non-CM newform, with the local parameters of the lift. We use the standard analytic continuation, order-one bound for the completed function, and functional equation of cuspidal standard $L$-functions. Absolute convergence and nonvanishing in $\Re s>1$ follow here directly from tempered good Euler factors and the bound on the bad factors below. These inputs do not require the location of zeros in the open strip. For $j=0$ one uses the classical completed zeta function with its poles at $0$ and $1$.

\begin{lemma}[Uniform local costs]\label{zf:cost}
For each $j\ge0$, the conductor $q_j$ satisfies $\log q_j\le(j+1)\log N$. At every prime, the logarithmic derivative at $p^v$ has coefficient of absolute value at most $(j+1)\log p$ in unitary normalization. The gamma factors of $L_j$ are $j+1$ real gamma factors $\Gamma_{\R}(s+\mu)$ counted with multiplicity, with $0\le\mu\le j+2$.
\end{lemma}
\begin{proof}
For the conductor use the inertia action on $H^1$ in an $\ell$-adic realization with $\ell\ne p$. The Artin conductor is the tame codimension of inertia invariants plus the integral of codimensions of the higher ramification invariants, using upper numbering. On any of these subgroups $H$,
\[
 \operatorname{codim}(\operatorname{Sym}^j V)^H
       \le(j+1)\operatorname{codim} V^H.
\]
If the right codimension vanishes the subgroup acts trivially. Otherwise the right side is at least the entire dimension of the left space. The assertion also accounts for nontrivial unipotent tame monodromy. Apply the conductor formula term by term, and sum after multiplying by $\log p$.

At potentially good reduction the inertia image is finite. After a finite extension giving good reduction the normalized Frobenius eigenvalues have absolute value one, by Hasse and passage to a power of Frobenius. This bounds their eigenvalues on inertia invariants in every symmetric power. At potentially multiplicative reduction the local parameter is a quadratic twist of a special representation. The kernel of its monodromy in the $j$th symmetric power is a single line with normalized eigenvalue of modulus $p^{-j/2}$, unless inertia invariants remove this line. Taking traces of powers proves the bound, also for good reduction. At infinity the restriction to $\C^\times$ has angular characters $e^{i(j-2r)\theta}$, $0\le r\le j$. Conjugate pairs give $\Gamma_{\C}(s+|j-2r|/2)$, and a middle line gives $\Gamma_{\R}(s)$ or $\Gamma_{\R}(s+1)$. Splitting each complex factor into two real ones proves the stated estimates.
\end{proof}

\paragraph{The auxiliary product.}
Let $I=(0,m,2m)$, let $\tau_0=\tau_{2m}=0$, $\tau_m=t$, and put
\begin{equation}\label{zf:product}
 D(s)=\prod_{i,j\in I}\ \prod_{r=0}^{\min(i,j)}
       L_{i+j-2r}\bigl(s+i_{\C}(\tau_i-\tau_j)\bigr),
 \qquad i_{\C}=\sqrt{-1}.
\end{equation}
The inner indices are the Clebsch--Gordan decomposition
$V_i\otimes V_j=\bigoplus_{r=0}^{\min(i,j)}V_{i+j-2r}$. It follows either by multiplying the two finite Laurent sums for their characters or by successively subtracting the highest weight. Formula~\eqref{zf:product} is our definition at all primes. In particular no equality between two different prescriptions for tensoring ramified local factors is required. Its degree, the sum of the dimensions of the constituents, is
\[
 d=\left(1+(m+1)+(2m+1)\right)^2=9(m+1)^2 .
\]
Exactly the three diagonal pairs contain $V_0$, with zero shift. The only possible poles of the completed product are therefore the poles of three zeta functions at $0$ and $1$, with order at most three. At a good prime the coefficient of $p^{-v s}$ in $-D'/D$ is
\[
 \log p\,\big|1+p^{-ivt}u_m(g_p^v)+u_{2m}(g_p^v)\big|^2\ge0 .
\]
At a bad prime the sum of absolute coefficients at $p^v$ is at most $d\log p$. Hence, for $1<\sigma\le2$,
\begin{equation}\label{zf:pos}
 -\frac{D'}{D}(\sigma)\ \ge -d\sum_{p\mid N}\frac{\log p}{p^\sigma-1}
 \ \ge -d\log(2N).
\end{equation}
Conjugate factors occur with equal exponents, so the value of the logarithmic derivative here is real.

\begin{lemma}[An explicit-formula inequality]\label{zf:ef}
Suppose a product $D$ as above has total degree $d$, aggregate log conductor $\le d\log N$, at most $r$ poles at each of $0$ and $1$ in its completion, and real symmetry with respect to $1/2$. Its archimedean shifts have real part in $[0,B]$ and imaginary parts at most $H$ in absolute value. Assume its ordinary Euler product satisfies \eqref{zf:pos}. For $1<\sigma\le2$, the zeros of
$\mathcal H(s)=[s(s-1)]^r\Lambda_D(s)$, counted with multiplicity, satisfy
\[
 \sum_{\rho}\Re\frac1{\sigma-\rho}
 \le \frac r{\sigma-1}+C\{r+d\log(2N(B+3)(H+3))\}
\]
with an absolute effective constant $C$. Each term on the left is nonnegative.
\end{lemma}
\begin{proof}
The completed product has no zeros in $\Re s>1$, since its Euler product there converges absolutely and its gamma factors are free of zeros and poles. Reflection excludes zeros in $\Re s<0$. The inserted polynomial can add zeros only at the boundary. Center the order-one entire function $\mathcal H$ at $1/2$. Remove a possible central zero of order $e$. The remaining function is even by the functional equation, including a possible sign before removal. Pair its remaining zeros $z$ and $-z$ in the centered variable. Since $\sum_{z\ne0}|z|^{-2}<\infty$, the paired canonical product converges locally uniformly. Its quotient into the centered function is the exponential of a polynomial of degree at most one, by Hadamard factorization. Evenness makes the linear term vanish. Logarithmic differentiation at real $\sigma>1$ therefore gives $\mathcal H'/\mathcal H(\sigma)=\sum_{\rho}\Re(\sigma-\rho)^{-1}$. Convergence of real parts is absolute, as $0\le\Re\rho\le1$.

The logarithmic derivative on the right of the defining product for $\mathcal H$ includes $r/\sigma+r/(\sigma-1)$, half its aggregate log conductor, and the sum of the real logarithmic gamma derivatives. For $\Re z\ge1/2$, $\Re(\Gamma'/\Gamma)(z)\le C_0+\log(2+|z|)$ with an absolute effective constant. This follows from Stirling's formula with remainder in the right half-plane, treating its bounded part by the convergent digamma series. The Euler derivative is bounded above by $d\log(2N)$ by \eqref{zf:pos}. Adding the contributions proves the claim.
\end{proof}

\begin{proof}[Proof of Theorem~\ref{zf:noncm}]
In the ordered pairs $(0,m),(2m,m)$ and their reversals the constituent $V_m$ occurs once in each pair. If $L_m(\beta+it)=0$ with $1/2<\beta\le1$, its conjugate is also a zero. Thus at any $\beta<1$ the completed product has at least four zeros there, with multiplicities; the same assertion holds for $\mathcal H$ at $\beta=1$ after clearing all three zeta poles. Real gamma shifts are $\le4m+2$, and shifts in height are $\le |t|$. Set $\mathcal L=C_1d\log(2N(m+2)(|t|+3))$, choosing $C_1$ sufficiently large for the lemma. Its conclusion gives
\[
 \frac4{\sigma-\beta}\le\frac3{\sigma-1}+\mathcal L .
\]
Take $u=(8\mathcal L)^{-1}$ and $\sigma=1+u$. If $1-\beta\le u/8$, the left side is at least $32/(9u)$ and the right side is $25/(8u)$, which is impossible. Absolute convergence handles $\beta>1$. Absorb the numerical and degree constants into $b$, decreasing it so that the asserted region has left endpoint beyond $1/2$.
\end{proof}

\subsection*{CM elliptic curves and their angular powers}
Let $E/\Q$ have complex multiplication by an order in an imaginary quadratic field $K$. In unitary normalization write its inducing Hecke character as $\mu$. Its archimedean type is angular of weight one and $\mu^c=\mu^{-1}$, where $c$ denotes the nontrivial element of $\operatorname{Gal}(K/\Q)$. Put $P_j=L(s,\operatorname{AI}_{K/\Q}\mu^j)$ for $j\ge1$, writing $P_j(s)$ for the function. These are entire by Hecke's theorem, or by cuspidal quadratic induction, since no nonzero angular power is fixed by $c$. When $j=1$ this is exactly $L(E,s+1/2)$.

\begin{theorem}\label{zf:cm}
Uniformly for $j\ge1$ the angular functions $P_j(\sigma+it)$ have no zeros when
\[
 \sigma>1-\frac{b'}{\log(2N(j+2)(|t|+3))} ,
\]
where $b'>0$ is effective and absolute.
\end{theorem}
\begin{proof}
Write $W_a=\operatorname{Ind}\mu^a$, identifying $W_{-a}=W_a$, and $W_0=1\oplus\chi_K$ for the quadratic character of the field. Restricting to the index-two subgroup, then pairing the two conjugate lines (or applying the tensor-induction identity), gives $W_a\otimes W_b=W_{a+b}\oplus W_{a-b}$. Use the three constituents $1,W_j|\cdot|^{it},W_{2j}$ and their duals as in \eqref{zf:product}. The resulting product in primitive functions is
\[
 \zeta(s)^3 L(s,\chi_K)^2 P_{2j}(s)^3P_{4j}(s)
 \prod_{\eta\in\{-1,1\}}P_j(s+\eta it)^2P_{3j}(s+\eta it).
\]
One checks the degree $3+2+6+2+2(4+2)=25$. At good primes the logarithmic coefficients are squared absolute values. At bad primes their absolute values are bounded by $25\log p$ at every prime power, by unitary Hecke induction. The conductor of $\mu^a$ divides the conductor of $\mu$, because a power stays trivial on every local principal-unit subgroup on which $\mu$ was trivial. The conductor--discriminant formula therefore gives $q(W_a)\le |D_K|\operatorname{Norm}\mathfrak f_\mu=N$. It also gives $q(\chi_K)\le N$. The infinity shifts are $O(j+1)$; all height shifts are bounded by $|t|$. The four occurrences of a proposed zero at $\beta+it$ and its conjugate again defeat the three zeta poles. Lemma~\ref{zf:ef} with $d=25$ proves the claim.
\end{proof}

In an odd symmetric power of the CM curve the pairs of weights $\pm j$ all give angular inductions $W_j$. In an even power there is also a middle line. Since the original parameter has determinant one, the action of a lift of complex conjugation on the balanced monomial of degree $2h$ is $(-1)^h$. Thus the middle line is $1$ for even $h$ and $\chi_K$ for odd $h$. Every other factor satisfies Theorem~\ref{zf:cm}. Its pole, when the middle line is trivial, must be kept explicitly. The verified Dirichlet theorem determines the zeros of the middle line. The angular factors retain precisely the difficulty visible already for $W_1$.

\subsection*{What can be gained by increasing the pole budget?}
There is an exact optimization for a natural extension of the argument at $t=0$. Take a finite sum $U=\bigoplus_{a=0}^M n_a V_a$ with nonnegative integral multiplicities. The trivial constituent of $U\otimes U$ has multiplicity $r=\sum n_a^2$. The defining representation occurs precisely once in $V_a\otimes V_b$ when $|a-b|=1$, and otherwise it does not occur. Its multiplicity is $q=2\sum_{a=0}^{M-1}n_an_{a+1}$. If $r>0$, then
\[
 \frac qr\le 2\cos\frac{\pi}{M+2}<2.
\]
Indeed the left side is the Rayleigh quotient of the adjacency matrix of a path on $M+1$ vertices. For $1\le k\le M+1$, the vector with entries $\sin((a+1)k\pi/(M+2))$ has eigenvalue $2\cos(k\pi/(M+2))$. These distinct eigenvalues exhaust the spectrum. Approximating positive multiples of the first eigenvector by integer vectors shows that the supremum over multiplicities is the displayed eigenvalue.

Thus this particular device can produce a multiplicity gain over the poles, but the ratio is bounded even when arbitrarily high symmetric powers are admitted. The error from conductors and archimedean factors still occurs at the point $\sigma>1$ where the prime sum has a sign. This explains the inverse-logarithmic width in the proof. No estimate for the same Euler derivative near $\sigma=1/2$ follows from this optimization.

\clearpage
\section{An all-height zero-free region linear in the symmetric-power degree}\label{amp:chapter}
Let $E/\Q$ be non-CM, with conductor $N$. Write $L_m(s,\chi)$ for the standard unitary $L$-function of $\operatorname{Sym}^m E$ twisted by a primitive Dirichlet character of conductor $q$. In this section $\chi$ may be nonreal. Put $U_j(\cos\theta)=\sin((j+1)\theta)/\sin\theta$ for the symmetric-power character, defined by continuity at the endpoints.

\begin{theorem}\label{amp:linear}
There is an absolute effective $a>0$ such that for every $m\ge1$, every $\chi$ and all real $t$,
\[
 L_m(\sigma+it,\chi)\ne0\qquad\text{if}\qquad
 \sigma>1-\frac{a}{(m+1)\log(2Nq(m+2)(3+|t|))}.
\]
There is no exceptional real zero in this region. The functional equation gives the reflected region at the other edge for the completed function.
\end{theorem}

For comparison, the three-constituent calculation of Theorem~\ref{zf:noncm} costs $(m+1)^2$. Thorner's Corollary~1.2 in \emph{Exceptional zeros of Rankin--Selberg $L$-functions and joint Sato--Tate distributions}, IMRN (2025), gives an $M^2$ dependence in a more general totally real and two-form setting. We use the established symmetric-power automorphy and local bounds in Lemma~\ref{zf:cost}. The new ingredient is the following finite polynomial majorant. It uses symmetric powers of index as large as $O(m^2)$ with positive rational weights whose total dimension is $O(m)$.

\begin{lemma}[A pair of positive character expansions]\label{amp:kernel}
For each $m\ge1$ there are finite polynomials $S=\sum_j s_jU_j$ and $T=\sum_j b_jU_j$ with nonnegative rational coefficients for which
\begin{gather*}
 S(x)\ge |T(x)|\quad(-1\le x\le1),\qquad
 b_m=1,\quad b_0=0,\quad s_0\le3/4,\\
 S(1)+T(1)\le C(m+1),\qquad
 \max\{j:s_j+b_j>0\}\le C(m+1)^2 .
\end{gather*}
Here $C$ is absolute.
\end{lemma}
\begin{proof}
Put $M=m+1$ and $a_j=\binom{2j}{j}/4^j$. We have $a_j\le(j+1)^{-1/2}$ by induction using $a_{j+1}/a_j=(2j+1)/(2j+2)$. Define
\[
       B_M(x)=2\sum_{j=0}^{M^2} a_j x^{2j}.
\]
For $y=1-x^2\ge M^{-2}$ the binomial series for $y^{-1/2}$ has tail at most
\[
 \frac{(1-y)^{M^2+1}}{\sqrt{M^2+1}\,y}
              \le e^{-1}y^{-1/2}.
\]
For $0\le y<M^{-2}$ use monotonicity in $x^2$ and apply this inequality at $y=M^{-2}$. Consequently
\[
 \min(M,(1-x^2)^{-1/2})\le B_M(x)
 \le 2(1-x^2)^{-1/2}\quad(|x|<1),\qquad B_M(1)\le6M.
\]
The first inequality extends to the endpoints by continuity with the minimum interpreted as $M$. In particular $B_M(x)\ge|U_m(x)|$, by the sine formula and the unitary character bound $|U_m|\le M$.

Take $k=32$ and suppose first that $m>2k$. Multiplying the sine formula by $\cos^{2k}\theta$ gives
\[
  x^{2k}U_m(x)=\sum_{h=-k}^{k} a_{k,h} U_{m+2h}(x),\qquad
   a_{k,h}=4^{-k}\binom{2k}{k+h}.
\]
All these indices are positive and $a_{k,0}=a_k$. Set
\[
               T=x^{2k}U_m/a_k,\qquad S=x^{2k}B_M/a_k.
\]
The expansion of $S$ has nonnegative coefficients, since $x=U_1(x)/2$ and tensor powers decompose with nonnegative integral multiplicities. The coefficient $s_0$ is its Haar integral. Using $d\mu=(2/\pi)\sin^2\theta\,d\theta$ gives
\[
 s_0\le\frac{4}{\pi a_k}\int_0^\pi
                     \cos^{2k}\theta\sin\theta\,d\theta
      =\frac{8}{\pi a_k(2k+1)}<\frac12 .
\]
The final numerical inequality follows from $a_{32}\ge1/(2\sqrt{32})$; the bound $a_j\ge1/(2\sqrt j)$ for $j\ge1$ follows from the same recurrence by induction. The norm and degree estimates follow from $S(1)+T(1)\le7M/a_{32}$ and the degree of $B_M$.

For $1\le m\le64$ use $A=1+U_{2m}$, $B=U_m$ and
\[
       S=(A^2+B^2)/4,\qquad T=AB/2 .
\]
Clebsch--Gordan gives $s_0=3/4$, $b_m=1$ and $b_0=0$. The elementary inequality $2|AB|\le A^2+B^2$ proves pointwise domination. The cost is $O(M^2)=O(M)$ on this bounded range, completing the construction.
\end{proof}

\begin{proof}[Proof of Theorem~\ref{amp:linear}]
Use $s_j$ copies of $L_j(s)$ at shift zero, and $b_j/2$ copies each of $L_j(s+it,\chi)$ and $L_j(s-it,\bar\chi)$. This notation for rational multiplicities means that we first clear their common denominator, work with an ordinary product of integer powers, and divide the resulting logarithmic-derivative inequality by that same integer. At $p\nmid Nq$, the prime-power coefficient of the negative logarithmic derivative divided by $\log p$ is
\[
 S(\cos(v\theta_p))+
 T(\cos(v\theta_p))\Re(\chi(p)^v p^{-ivt})\ \ge0,
\]
by Lemma~\ref{amp:kernel}. The only factors that can have poles are the unshifted copies of $\zeta$, of total weight $s_0$. All other standard functions are entire, including the twists, by Newton--Thorne and ordinary character twisting. Clearing the zeta poles introduces no negative zero multiplicities.

The aggregate degree is $d=S(1)+T(1)\ll M$ because $U_j(1)=j+1$. Twisting by $\chi$ adds at most $(j+1)\log q$ to the logarithm of the conductor: on an inertia subgroup $H$ the extra codimension is at most $(j+1)\operatorname{codim}\chi^H$, and one integrates the conductor formula. At the ramified primes the local logarithmic coefficients have modulus at most $(j+1)\log p$, by the same Weil--Deligne argument as Lemma~\ref{zf:cost} applied after twisting. The nonpositive contribution of all bad primes to the real logarithmic derivative on $\sigma>1$ is therefore at worst $-d\log(2Nq)$. The real gamma shifts are $O(M^2)$ and the translated ordinates have absolute value at most $|t|$. Lemma~\ref{zf:ef} now applies after clearing denominators, with
\[
             \mathcal L=C_1 M\log(2Nq(m+2)(3+|t|)).
\]
Its absolute constant is independent of the common denominator. The product is real and self-dual because the shifts and twists occur in conjugate pairs. A zero at $\beta+it$ of the target gives a zero of total weight at least $b_m=1$ at the real point $\beta$, including when the two translated factors coincide. The lemma yields
\[
           \frac1{\sigma-\beta}\le\frac{s_0}{\sigma-1}+\mathcal L
                       \qquad(1<\sigma\le2).
\]
If $u=(16\mathcal L)^{-1}$ and $1-\beta\le u/8$, taking $\sigma=1+u$ makes the left side at least $8/(9u)$ while the right side is at most $13/(16u)$. This contradiction proves the theorem, after increasing $C_1$ to ensure $u<1$.
\end{proof}

\subsection{The degree and height costs of this argument}
A positive product that detects a standard function of degree $m+1$ with zero multiplicity $b$ has aggregate degree at least $b(m+1)$ from those factors alone. Hence a logarithmic-derivative argument that charges the aggregate degree against the detected multiplicity has degree cost at least $m+1$. The preceding construction attains this order with a pole budget strictly below $b$.

For fixed $m$ the gamma factors of the detected translates at large $|t|$ contribute $\tfrac12 b(m+1)\log|t|+O_m(b)$ to the real logarithmic derivative on $1<\sigma\le2$, by Stirling's formula. Dropping all other zero terms in the explicit formula pays this contribution without a compensating lower bound for the omitted zeros. Thus the all-height argument improves the degree dependence but still operates at distance $O(1/\log |t|)$ from the line $\Re s=1$ for fixed $m$. To reach a fixed interior line by this approach one needs additional cancellation of this gamma contribution or a lower bound for the remaining zero sum of comparable size.

\subsection{A quantitative Sato--Tate consequence}\label{amp:st}
Let $\pi_{E,I}(x)$ count primes $p\le x$ of good reduction with $\theta_p\in I\subset[0,\pi]$. Write $\pi(x)$ for the total number of rational primes through $x$.

\begin{theorem}\label{amp:discrepancy}
Uniformly for non-CM elliptic curves $E/\Q$ of conductor $N$, intervals $I$ and $x\ge3$,
\[
 \left|\frac{\pi_{E,I}(x)}{\pi(x)}-\frac2\pi\int_I\sin^2\theta\,d\theta\right|
 \ll \frac{\log(3N\log(3x))\log(3\log(3x))}{\log(3x)}
\]
with an absolute effective implied constant. For fixed $E$ this is $O_E((\log\log x)^2/\log x)$. In particular the relative Sato--Tate asymptotic holds uniformly for moving intervals whose Sato--Tate masses divided by this error tend to infinity.
\end{theorem}

\begin{proof}
Throughout, put $L=\log(3x)$, $\ell=\log(3L)$ and $J=\log(3NL)$. We give the vertical estimates to account for all constants that can grow with $m$. For $j\ge1$ write $r=j+1$, $F=L_j(\cdot,\mathbf1)$ and $\Lambda_F$ for its completion. For $|t|\le T+1$, the Hadamard identity at $2+it$ gives
\begin{equation}\label{amp:localzeros}
 \sum_{\rho=\beta+i\gamma}\frac1{1+(t-\gamma)^2}
       \ll r\log(2N(j+2)(T+4)).
\end{equation}
In fact the real logarithmic derivative of the completion is the sum over $\rho$ of $(2-\beta)/[(2-\beta)^2+(t-\gamma)^2]$. This follows by horizontal reflection in the order-one Hadamard product, as in Lemma~\ref{zf:ef}. Each term dominates one fourth of the summand in \eqref{amp:localzeros}. Lemma~\ref{zf:cost}, the digamma estimate in the right half-plane and absolute convergence of $F'/F$ at $2+it$ bound this logarithmic derivative by the right side.

Put $T=L^{16}$, $K=\log(2N(j+2)(T+4))$ and $\eta=a/(4rK)$, decreasing the absolute constant $a$ of Theorem~\ref{amp:linear} if necessary. For $|t|\le T$ and $1-\eta\le \sigma\le2$,
\[
                         |F'/F(\sigma+it)|\ll (rK)^2.
\]
Indeed, zeros with $|\gamma-t|\le1$ have $\beta\le1-4\eta$. Subtract the Hadamard logarithmic derivatives at $s=\sigma+it$ and $s_2=2+it$. The difference of a pair of zero summands is $(s_2-s)/[(s-\rho)(s_2-\rho)]$. For the nearby zeros bound this by $O(1/\eta)$, and for the rest by $O(1/[1+(t-\gamma)^2])$. Equation~\eqref{amp:localzeros} gives the required bound. The difference of the gamma logarithmic derivatives is $O(r)$ by the absolutely convergent trigamma series on $\Re z\ge1/4$. The finite-part logarithmic derivative at $s_2$ is $O(r)$.

Write $-F'/F=\sum v_j(n)n^{-s}$ and $\psi_j(x)=\sum_{n\le x}v_j(n)$. We have $|v_j(p^k)|\le r\log p$. Choose $h=xL^{-8}$ and $b=1+1/L$ and use the kernel
\[
  W(s)=\frac{(x+h)^{s+1}-x^{s+1}}{hs(s+1)}.
\]
Termwise Perron inversion on $\Re s=b$ gives $\psi_j(x)$ with a linear taper on $(x,x+h]$. Thus its difference from $\psi_j(x)$ is $O(r(h+1)\log(2x))$. At every $0<\sigma\le b$ the kernel is bounded both by $(x+h)^\sigma/|s|$ and by $2(x+h)^{\sigma+1}/(h|s(s+1)|)$. The two infinite tails on the original line contribute $O(rxL^{10}/T)$ because $\sum|v_j(n)|n^{-b}\ll rL^2$. Move the middle segment to $1-\eta$ without crossing poles or zeros. Its two horizontal segments contribute $O((rK)^2 xL^8/T^2)$, and its new vertical segment contributes
\[
                      O\big((rK)^2x e^{-\eta\log x}\log(3T)\big).
\]
All of these estimates use the finite-part logarithmic derivative.

Choose a sufficiently small absolute effective $\kappa>0$. For $1\le j\le \kappa L/(J\ell)$ we have $K\ll J$, $rK\ll L$ and $\eta\log x\ge12\ell$, after reducing $\kappa$. The preceding bounds give $\psi_j(x)\ll xL^{-4}$ in this range, adjusting an absolute constant on bounded $x$. Removing the powers $p^k$, $k\ge2$, costs $O(r\sqrt x\log^2(2x))$ by counting integers up to $x^{1/k}$. Removing the first powers at bad primes costs at most $r\log N$. In the indicated range both errors are $O(xL^{-4})$. We conclude
\[
      \left|\sum_{\substack{p\le x\\p\nmid N}}U_j(\cos\theta_p)\log p\right|
                                     \ll xL^{-4}.
\]

Use the standard degree-$B$ trigonometric majorant and minorant for an interval in the Sato--Tate measure. Their $U_0$ coefficients differ from $\mu(I)$ by $O(1/B)$ and each $U_j$ coefficient is $O(1/j)$ for $1\le j\le B$, uniformly in $I$. This is the Beurling--Selberg interval approximation in the $U_j$ basis; see the approximation used by Rouse and Thorner, \emph{The explicit Sato--Tate conjecture and densities pertaining to Lehmer-type questions}, Trans. Amer. Math. Soc. \textbf{369} (2017), Lemma 3.1. Choose $B=\lfloor\kappa L/(J\ell)\rfloor$. If $B<1$ the claimed discrepancy follows from its trivial bound. Otherwise Chebyshev's bound for $\sum_{p\le x}\log p$ gives weighted discrepancy $O(xJ\ell/L)$, including the omission of primes dividing $N$. Partial summation gives the same relative error for unweighted primes. In detail, for $\sqrt x\le y\le x$ the quantity $J_y\ell_y/L_y$ is at most twice its value at $x$; the shorter initial range contributes $O(\sqrt x\log(2x))$ to the weighted count. The lower Chebyshev bound $\pi(x)\gg x/\log x$ completes the proof.
\end{proof}

For context, the explicit result of Hoey, Iskander, Jin and Trejos Suárez, \emph{An unconditional explicit bound on the error term in the Sato--Tate conjecture}, Q. J. Math. (2022), gives an explicit constant for an error of order $\log(N\log x)/\sqrt{\log x}$ in the elliptic case. The preceding proof gives an effective absolute implied constant and keeps the dependence on the symmetric-power index throughout the contour shift.

\subsection{Hilbert modular forms and two independent factors}\label{amp:hilbert}
Let $K$ be totally real, $f=[K:\Q]$ and $D_K$ its absolute discriminant. Let $\pi$ be a unitary non-dihedral cuspidal representation of $\mathrm{GL}_2(\mathbb A_K)$ associated to a holomorphic Hilbert modular form of regular weights $k_v\ge2$, with finite-order central character $\omega$. Write $\mathfrak n$ for its conductor. When a second form $\pi'$ occurs impose the same assumptions, writing $\mathfrak n'$ for its conductor and requiring that $\pi'$ is not any character twist of $\pi$. Let $k_*$ be the largest weight of the forms in use. The required symmetric powers are cuspidal with their local parameters by Newton and Thorne, \emph{Symmetric power functoriality for Hilbert modular forms}, Ann. of Math. \textbf{203} (2026), 283--347. Equivalently one may work with the regular algebraic forms in that theorem after the indicated unitary normalization.

\begin{theorem}\label{amp:pair}
Let $m,n\ge0$, $m+n>0$, $d=(m+1)(n+1)$, and let $\chi$ be a finite-order Hecke character of conductor $\mathfrak q$. Put
\[
 \mathcal B=\log(2D_K\,\mathrm N(\mathfrak n\mathfrak n'\mathfrak q))
       +f\log\big(3k_*(m+n+2)(3+|t|)\big).
\]
For an absolute effective constant $a'>0$,
\[
 L\big(\sigma+it,\operatorname{Sym}^m\pi\times
             (\operatorname{Sym}^n\pi'\otimes\chi)\big)\ne0
               \qquad\text{whenever }\quad \sigma>1-\frac{a'}{d\mathcal B}.
\]
When $n=0$ the second form and its conductor may be omitted. In particular the assertion holds for every non-CM primitive holomorphic form over $\Q$ of integral weight at least two and arbitrary nebentype.
\end{theorem}

\begin{proof}
Use Lemma~\ref{amp:kernel} separately in the two variables and take $S=S_mS_n$, $T=T_mT_n$. For an index zero put $S_0=T_0=1$. Then $S\ge|T|$, the coefficient of $(m,n)$ in $T$ is one, and the scalar coefficient of $S$ is at most $3/4$. The scalar coefficient of $T$ is zero. Positivity of all character coefficients is preserved, $S(1,1)+T(1,1)\ll d$, and all indices are $O((m+n+2)^2)$.

At a good place write the unitary Satake class of $\pi$ as $\zeta g$ with $g\in\mathrm{SU}(2)$ and $\zeta^2=\omega(\mathfrak p)$. This uses temperedness for regular algebraic Hilbert modular forms. A choice of square root at each place suffices. Every index $i$ in the static polynomial $S_m$ is even. Use the cuspidal factor $\operatorname{Sym}^i\pi\otimes\omega^{-i/2}$ for it (and the trivial representation when $i=0$). Every index $i$ in $T_m$ has $i\equiv m\pmod2$. Use $\operatorname{Sym}^i\pi\otimes\omega^{(m-i)/2}$ there. The corresponding good traces are $U_i(g)$ and $\zeta^m U_i(g)$, respectively. Do the same for the other form. These prescriptions are independent of the choices because they are given by global integral powers of genuine Hecke characters. Tensor the relevant pairs via their Rankin--Selberg functions, and put the twisting character and the ordinate shift into the oscillating terms as in Theorem~\ref{amp:linear}. At every good prime power the coefficient is $S+T\Re z$ for $|z|=1$.

Every nontrivial auxiliary function is entire. For unequal ranks this is the Rankin--Selberg pole criterion. For equal positive indices a pole would identify the two symmetric powers up to a character twist and duality. The dual of a symmetric power of a two-dimensional parameter is itself a character twist, so their endomorphism parameters would agree at all good places. Clebsch--Gordan realizes those parameters there as the isobaric sums $\boxplus_{j=0}^i\operatorname{Sym}^{2j}\pi\otimes\omega^{-j}$ and the analogous sum for $\pi'$. All their positive-index constituents are cuspidal. Strong multiplicity one and uniqueness of the isobaric decomposition identify the unique degree-three constituents, hence $\operatorname{Ad}\pi\simeq\operatorname{Ad}\pi'$. Ramakrishnan's adjoint-fibre theorem (\emph{Modularity of the Rankin--Selberg $L$-series, and multiplicity one for $\mathrm{SL}(2)$}, Ann. of Math. \textbf{152} (2000), Theorem 4.1.2) then makes $\pi'$ a character twist of $\pi$, contrary to hypothesis. The sole polar factor is therefore $\zeta_K$ at shift zero, with weight at most $3/4$. The detected target and its dual have total weight one. Only the index $(m,n)$ in the oscillating terms is needed for this lower bound.

For clarity, the conductor estimate can be made before summing the auxiliary functions. For a local tensor of dimensions $(i+1)(j+1)=d_{ij}$, the conductor exponent is at most
\[
  C d_{ij}\big(a(\pi_\mathfrak p)+a(\pi'_\mathfrak p)+a(\chi_\mathfrak p)\big)
\]
with an absolute $C$. Indeed the actual Weil--Deligne representations, realized on inertia by the usual unipotent monodromy, have trivial action on any ramification group on which both original two-dimensional actions and $\chi$ are trivial. On any other group the new codimension is at most $d_{ij}$ and the sum of the original codimensions is at least one. Apply this to inertia and integrate on the higher ramification groups. The determinant characters used above are already trivial whenever their defining representations are trivial. This also handles arbitrary ramification of the original forms. At a ramified place the local logarithmic coefficient is bounded by $d_{ij}\log\mathrm N\mathfrak p$: tempered local parameters are bounded on the Weil factor in Weil--$\mathrm{SL}_2$ notation, and each special summand contributes on its monodromy kernel a normalized eigenvalue of absolute value at most one. This argument is stable under duals, tensor products, symmetric powers and unitary character twists. At infinity the angular exponents give nonnegative real gamma shifts $O(k_*(i+j+1))$, and the ordinate shifts have modulus at most $|t|$.

Summing these costs, the conductor including the field discriminant contributes $O(d\log(2D_K\mathrm N(\mathfrak n\mathfrak n'\mathfrak q)))$. The total archimedean contribution to the logarithmic-derivative inequality is $O(fd\log(3k_*(m+n+2)(3+|t|)))$. Excluding the ramified Euler factors costs $O(d\log(2\mathrm N(\mathfrak n\mathfrak n'\mathfrak q)))$. The proof of Lemma~\ref{zf:ef} applies to these completed functions with the stated gamma factors and discriminant. The absolute constants there come from the half-plane digamma estimate, with one term for every real gamma factor. Clearing common denominators and using the same choice of $\sigma-1$ as in Theorem~\ref{amp:linear} gives the claim.
\end{proof}

\begin{corollary}\label{amp:progression}
For $(u,q)=1$ let $\pi(x;q,u)$ count primes in that residue class and $\pi_{E,I}(x;q,u)$ count its good primes with angle in $I$. Then
\[
 \frac{|\pi_{E,I}(x;q,u)-\mu(I)\pi(x;q,u)|}{\pi(x)}
 \ll\frac{\log(3Nq\log(3x))\log(3\log(3x))}{\log(3x)} .
\]
For fixed $q$ this also bounds the relative discrepancy inside the class with an additional constant depending on $q$.
\end{corollary}
\begin{proof}
For each nonconstant angular mode apply the contour argument of Theorem~\ref{amp:discrepancy} to the twist by a primitive inducing Dirichlet character. Its dual functional equation gives the same horizontal reflection of the zeros within its own unitary strip, so that \eqref{amp:localzeros} holds with $Nq$ in place of $N$. Removing the additional primes dividing $q$ costs at most $(m+1)\log q$ for the prime sum. Character orthogonality projects onto the residue class. The division by $\varphi(q)$ cancels the number of character sums in the uniform upper bound. The constant angular mode contributes $\pi(x;q,u)$ itself. The prime number theorem in a fixed progression gives the last assertion.
\end{proof}

\subsection{One two-prime correlation equivalent to every symmetric-power ERH}\label{amp:energy}
The degree-uniform theorem also controls a quantity that detects all the symmetric powers at once. For each good prime let $x_p=a_p(E)/(2\sqrt p)$, and fix $0<r<1$. The single choice $r=1/2$ suffices. Define
\[
 K_r(x,y)=\frac{1-r^2}
 {(1-r^2)^2-4r(1+r^2)xy+4r^2(x^2+y^2)}
\]
and write $\mathcal O_r(X)$ for the signed off-diagonal sum
\[
 \mathcal O_r(X)=
  \sum_{\substack{p,q\le X,\ p\ne q\\p,q\nmid N}}
        [K_r(x_p,x_q)-1]\log p\log q .
\]
The sum is over ordered pairs. Its inputs are just finite-field point counts and the good primes.

\begin{theorem}\label{amp:energy-iff}
For any fixed non-CM elliptic curve $E/\Q$, the following two assertions are equivalent.
\begin{enumerate}
 \item Every completed standard $L$-function of $\operatorname{Sym}^j E$, $j\ge1$, has all its zeros on the unitary critical line.
 \item For one fixed $r\in(0,1)$ and every $\epsilon>0$,
 $\mathcal O_r(X)\le C_{E,r,\epsilon}X^{1+\epsilon}$ for all $X\ge2$.
\end{enumerate}
If they hold for one $r$, they hold for all such $r$. Unconditionally one has
\[
 -O_r(X\log(2X))\le\mathcal O_r(X)
       \ll_{E,r} X^2 e^{-c_{E,r}(\log(3X))^{1/3}}
\]
for effective positive constants. The diagonal omitted from $\mathcal O_r(X)$ is asymptotic to $rX\log X/(1-r)$.
\end{theorem}
\begin{proof}
Absolute uniform convergence follows from $|U_j|\le j+1$. Summing the two cosine series obtained from $\sin((j+1)\theta)\sin((j+1)\phi)$ proves
\[
                    K_r(x,y)=\sum_{j=0}^\infty r^jU_j(x)U_j(y).
\]
More explicitly the denominator before expansion is
$|1-re^{i(\theta-\phi)}|^2|1-re^{i(\theta+\phi)}|^2$, with $x=\cos\theta$, $y=\cos\phi$. Continuity treats the endpoints. Write
\[
  A_j(X)=\sum_{p\le X,\ p\nmid N} U_j(x_p)\log p,\qquad
  \mathcal E_r(X)=\sum_{j=1}^\infty r^j|A_j(X)|^2 .
\]
The exact identity $\mathcal E_r=\mathcal O_r+\mathcal D_r$ holds with
$\mathcal D_r=\sum_{p\le X,\ p\nmid N}(K_r(x_p,x_p)-1)(\log p)^2$.
Since $0\le K_r(x,x)-1\le\sum_{j\ge1}r^j(j+1)^2$,
$0\le\mathcal D_r\ll_r X\log(2X)$. Positivity of $\mathcal E_r$ gives the lower bound for the signed off-diagonal sum, for every $X$.

Assume the asserted upper bound on $\mathcal O_r$. For each fixed $j$,
$r^j|A_j|^2\le\mathcal O_r+\mathcal D_r$, and hence $A_j(X)\ll_{E,r,j,\epsilon}X^{1/2+\epsilon}$. Inserting the bad primes and the higher prime powers costs $O_{E,j}(\sqrt X\log^2(2X))$. It follows that
\[
 s\int_1^\infty\psi_j(y)y^{-s-1}\,dy
\]
is holomorphic for $\Re s>1/2$ and agrees there, by continuation from $\Re s>1$, with $-L_j'/L_j$. A zero would give the latter a pole. The functional equation excludes off-line zeros in the other half.

Conversely, suppose ERH holds for every positive $j$. In the proof of Theorem~\ref{amp:discrepancy}, take $\sigma_0=1/2+\epsilon$ for $0<\epsilon<1/4$. The nearby zeros are at distance at least $\epsilon$, and \eqref{amp:localzeros} yields
\[
       |L'_j/L_j(\sigma+it)|\ll_\epsilon
       (j+1)\log(2N(j+2)(T+4))
       \quad(\sigma_0\le\sigma\le2,\ |t|\le T).
\]
Use $h=\sqrt X$ and $T=X^2$ in the same Perron difference. Its vertical segment, horizontal segments, original tails and taper error together give
\[
 |A_j(X)|\ll_\epsilon (j+1)X^{1/2+\epsilon}
                       \log^2(2N(j+2)(X+3)).
\]
The prime-power error is included here. Thus the constants are summable after squaring against $r^j$. Reducing $\epsilon$ absorbs the logarithms and gives $\mathcal E_r\ll_{E,r,\epsilon}X^{1+\epsilon}$, proving the required upper bound for $\mathcal O_r$.

For the unconditional upper bound put $V=(\log(3X))^{1/3}$. When $1\le j\le V$, take $T=e^{4V}$ and $h=Xe^{-V}$ in the same contour argument. For fixed $E$, $K=\log(2N(j+2)(T+4))\ll_E V$, and the available $\eta$ has $\eta\log X\gg_E V$. The vertical segment and the three other errors give
$|A_j(X)|\ll_E Xe^{-c_E V}$, uniformly for $j\le V$, absorbing powers of $\log X$ by reducing $c_E$. For the other modes use $|A_j|\ll (j+1)X$ and
$\sum_{j>V}r^j(j+1)^2\ll_r e^{-c_rV}$. Summation gives the stated upper bound for $\mathcal E_r$ and hence for $\mathcal O_r$.

Finally, orthogonality of the symmetric-power characters gives
$\int(K_r(x,x)-1)\,d\mu(x)=\sum_{j\ge1}r^j=r/(1-r)$. Sato--Tate for this continuous function and partial summation, together with the prime number theorem, give the stated diagonal asymptotic.
\end{proof}

The remaining estimate concerns correlations between \emph{different} primes. One can see quantitatively how much room is left by even a small discrepancy estimate. With $d\nu_\delta=(1+\delta U_1/2)d\mu$, $0\le\delta\le1$, one has a probability measure whose only nonconstant character moment is $\delta/2$ at index one. Thus
\[
 \iint(K_r(x,y)-1)\,d\nu_\delta(x)d\nu_\delta(y)=r\delta^2/4.
\]
Taking $\delta=e^{-(\log X)^{1/3}}$ still gives $X^2r\delta^2/4$ far larger than $X^{1+\epsilon}$ for $0<\epsilon<1$. This is a statement about the strength of one-point distribution estimates, not an assertion about Frobenius classes of an elliptic curve. The rational kernel identifies the missing upper bound and the linear-degree amplification supplies the unconditional estimate stated above.

\subsection{The limitation of squares of isobaric sums}\label{amp:square}
The linear cost in Lemma~\ref{amp:kernel} improves what is possible with a single Rankin--Selberg square. We record the precise comparison, since it explains why larger isobaric sums alone stop short of the preceding theorem.

\begin{proposition}\label{amp:sqbound}
Let an isobaric sum contain $a_{j,u}\ge0$ copies of the constituent $\operatorname{Sym}^j E\otimes|\cdot|^{iu}$, finitely supported. Write
\[
 R=\sum_{j,u}(j+1)a_{j,u},\qquad P=\sum_{j,u}a_{j,u}^2 .
\]
In its Rankin--Selberg square let $B$ count the copies of the target of symmetric index $m\ge1$ at shifts $t$ or $-t$, counting each ordered constituent pair once for each target factor. If $B>P$, then $B\ll R^2/m^{3/2}$. This exponent is attained up to absolute constants. This count concerns the target multiplicity forced by the character decomposition.
\end{proposition}
\begin{proof}
Write $H_u=\sum_{j\ge m/2}a_{j,u}$ and $S_u=\sum_j a_{j,u}$. An ordered pair supplying $U_m$ has sum of indices at least $m$, so one index is at least $m/2$. Each Clebsch--Gordan multiplicity is one. The height difference must be $t$ or $-t$. Therefore
\[
 B\le2\sum_u H_u(S_{u+t}+S_{u-t})
       \le4\|H\|_2\|S\|_2,\qquad
 \|H\|_2\le\|H\|_1\le2R/m .
\]
Zero values are supplied at shifts absent from the sum. Split $S_u$ at index $J$. Cauchy--Schwarz on the first $J+1$ entries and the $\ell^1$ bound for the remaining entries imply
\[
            \|S\|_2\le\sqrt{(J+1)P}+R/(J+1).
\]
Choose $J+1$ within a factor two of $(R^2/P)^{1/3}\ge1$. Then
$B\ll R^{4/3}P^{1/3}/m$. Using $P<B$ proves $B\ll R^2/m^{3/2}$. The same proof works for nonnegative real weights. The number $P$ is the unavoidable diagonal pole budget with coincident shifts; any other shifted zeta poles do not improve this count.

For the reverse construction take one copy of each $V_0,V_2,\ldots,V_{2k}$ at shift zero and one $V_m$ at shift $t$, with $k=\lfloor\sqrt m/2\rfloor$ and $m$ sufficiently large. The total rank is $(k+1)^2+m+1\asymp m$, the pole budget is $k+2$, and each of the $2(k+1)$ ordered cross pairs contains $V_m$ at the required shift. The resulting degree cost per forced target copy is $O(m^{3/2})$, and the pole budget is bounded away from the forced multiplicity by a fixed relative margin. Lemma~\ref{zf:ef} indeed gives a zero-free region of that degree order for this construction. The unrestricted positive polynomials in Lemma~\ref{amp:kernel} remove this square constraint.
\end{proof}

\clearpage
\section{An explicit all-height exclusion for the undeformed elliptic function}\label{exstrip:chapter}
Use the unitary variable $L_E(s)=L(E,s+1/2)$, with center $1/2$. A finite certification settles the complete strip in its height range. Above it the positivity argument can be evaluated with explicit constants.

\begin{theorem}\label{exstrip:main}
For $11\mathrm a1$, all nontrivial zeros with $|\gamma|\le524288$ lie on the critical line. Every zero $\beta+i\gamma$ with $|\gamma|\ge262144$ satisfies
\[
           \frac1{14\log|\gamma|}<\beta<
                           1-\frac1{14\log|\gamma|}.
\]
For any $E/\Q$ of conductor $N$, put
\[
 P_E=\frac{4105+2113\,\mathbf1_{\rm CM}(E)}{6336},\qquad
 B_E(T)=17\log N+24+\tfrac83\log(6+|T|).
\]
Every nontrivial zero satisfies
\[
      \frac{(1-\sqrt{P_E})^2}{B_E(\gamma)}
          \le \beta\le 1-\frac{(1-\sqrt{P_E})^2}{B_E(\gamma)} .
\]
The rational elliptic curve of level eleven has $1\,924\,341$ simple positive zeros in the stated finite window. The all-height conclusion applies to its original, undeformed $L$-function.
\end{theorem}
\begin{proof}
For the $\mathrm{SU}(2)$ characters of Section~\ref{amp:chapter} take
\[
 X=\tfrac{11}8(U_1+\tfrac14U_3),\quad Y=1+U_2,\quad
 S=(X^2+Y^2)/(99/16),\qquad J=2XY/(99/16).
\]
Clebsch--Gordan gives
\[
 S=\tfrac{4105}{6336}+\tfrac{6097}{6336}U_2+
       \tfrac{2113}{6336}U_4+\tfrac{11}{576}U_6,\qquad
 J=U_1+\tfrac23 U_3+\tfrac19U_5 .
\]
Write $s_j,t_j$ for these coefficients. At each good prime power the real logarithmic-derivative coefficient $S+J\cos(\nu T\log p)$ is nonnegative, by $2|XY|\le X^2+Y^2$. Use the static factors $L(\operatorname{Sym}^jE,s)$ with weights $s_j$ and the same factors at $s\pm iT$ with weights $t_j/2$ for the oscillating part. Rational weights mean a linear combination of logarithmic derivatives; no choice of branch enters. The total degree is $D=4225/396$ and the oscillating degree is $16/3$.

For a non-CM curve all positive symmetric powers here are cuspidal; for CM the odd powers have only nonzero angular powers of the CM Hecke character after induction. In degree $2r$ the middle constituent is $\eta_K^r$, where $\eta_K$ is the quadratic character of the CM field, and the other induced constituents have nonzero infinity type. This follows by taking the monomials in the two conjugate inducing characters: the middle monomial on the index-two subgroup has trivial normalized character, and an element outside swaps the two weights, acting there by $(-1)^r$ since the normalized determinant is one. Thus among the added static factors only $\operatorname{Sym}^4$ contributes another zeta pole in the CM case. The total pole weight is precisely $P_E<1$.

At each bad prime the absolute coefficient of the weighted Euler derivative is bounded by $D$ after division by $\log p$, using the local tempered bounds for elliptic Weil--Deligne factors and their symmetric powers. Positivity therefore gives a lower bound $-D\log N$ for the sum of the negative real Euler logarithmic derivatives on the real line $\sigma=1+u>1$. Here $\sum_{p\mid N}\log p/(p^\sigma-1)\le\log N$. The conductor exponents satisfy $Q_j\le N^{j+1}$ by the local estimate of Lemma~\ref{zf:cost}. The arithmetic-conductor contribution in the Hadamard identities is at most $(D/2)\log N$. The gamma shifts for the degrees used here are real in $[0,4]$. In particular for $0<u\le1/100$ each real gamma factor contributes at most $2+\tfrac12\log(6+|T|)$ when shifted and at most $2$ when static, by $|\psi(z)-\log z|\le1/\Re z$ for $\Re z>0$. The zeta pole at zero costs at most $P_E/\sigma$ in addition to $P_E/u$ from one. Dropping all nonnegative real zero terms except the target and its conjugate, of total weight $t_1=1$, gives
\begin{equation}\label{exstrip:ineq}
           \frac1{u+1-\beta}\le\frac{P_E}{u}+B_E(T)
                              \qquad(T=\gamma).
\end{equation}
This uses $3D/2<17$, $2D+1<24$. Put $u=(\sqrt{P_E}-P_E)/B_E(T)$, which is in the required interval since $B_E(T)\ge24$ and both listed values of $P_E$ satisfy $\sqrt{P_E}-P_E<0.16$. Rearranging proves the general bound. Reflection gives the other edge.

For the sharper level-eleven estimate take a different nonnegative polynomial, now with five frequency components. Put
\[
       (P_0,P_1,P_2,P_3,P_4)=(1,3U_1,3+2U_2,3U_1,1).
\]
Expand $\left|\sum_{j=0}^4e^{ij\phi}P_j\right|^2/72$ as $\sum_{\ell=0}^4 W_\ell\cos(\ell\phi)$. Direct Clebsch--Gordan multiplication gives the table
\[
\begin{array}{c|ccccc}
 &U_0&U_1&U_2&U_3&U_4\\ \hline
 W_0&11/24&0&17/36&0&1/18\\
 W_1&0&1&0&1/3&0\\
 W_2&5/12&0&13/36&0&0\\
 W_3&0&1/6&0&0&0\\
 W_4&1/36&0&0&0&0 .
\end{array}
\]
Denote the entries by $w_{\ell j}$. The total shifted half-degree is $g=\frac12\sum_{\ell\ge1,j}(j+1)w_{\ell j}=187/72$. Positivity at good primes is built into the square. At $11$ every symmetric-power trace evaluates $U_j$ as $x^j$, with $x=11^{-\nu/2}$ at the $\nu$th prime power. Under this evaluation the Clebsch--Gordan product $U_iU_j$ becomes
\[
           \sum_{b=0}^{\min(i,j)}x^{i-b}x^{j-b}.
\]
This is the Gram matrix of the real vectors $(x^i,x^{i-1},\ldots,1,0,\ldots)$, so evaluation of this entire squared norm remains nonnegative at $11$. Each symmetric-power conductor is $11^j$ by the split Steinberg local parameter.

Set $\kappa=1/\sqrt5$ and $\sigma_1=(1+\sqrt{1+4\sigma^2})/2$.
Subtract $\kappa$ times the weighted Hadamard identities at $\sigma_1$
from the same identities at $\sigma$. The Euler coefficients acquire
the nonnegative multiplier $p^{-\nu\sigma}-\kappa p^{-\nu\sigma_1}$.
The two zeros in a functional-equation pair at ordinate $v$ have
real parts $b,1-b$. For $s=\sigma+ih$ put
\[
   K(s,b+iv)=\Re\left(\frac1{s-b-iv}+\frac1{s-1+b-iv}\right).
\]
Stechkin's pointwise inequalities give
$K(\sigma+ih,\rho)-\kappa K(\sigma_1+ih,\rho)\ge0$.
When $h=v$ and $b>1/2$, this difference is at least
$1/(\sigma-b)$. These are
\href{https://arxiv.org/abs/1106.1868}{Kadiri, Lemma~2.1,
equations (2.15)--(2.16)}. Count pairs on the line with half weight.
For a target to its right the two shifts $\pm T$ carry equal
weight $w_{11}/2=1/2$ and select the target and its conjugate.
Their total retained contribution is $1/(\sigma-\beta)$.

Let $1<\sigma=1+u\le1.02$; hence $\sigma_1<2$.
For $j=0,2,4$ write $M_j$ for the multiset consisting of
$(j/2)\bmod2$ and $l,l+1$ for $l=1,\ldots,j/2$.
These are the real gamma shifts of the even, static powers. Write $P=w_{00}$ and put
\begin{align*}
 C={}&P+\tfrac12\log11\sum_{\ell,j}j w_{\ell j}
   +\tfrac12\sum_j w_{0j}\sum_{\mu\in M_j}
                \bigl(\psi((1.02+\mu)/2)-\log\pi\bigr)\\
     &+\tfrac12\sum_{\ell\ge1,j}(j+1)w_{\ell j}\log\ell
                  <4.064 .
\end{align*}
The static digamma is increasing on the positive real axis, so its
differenced bound is $(1-\kappa)$ times the term just displayed,
including the subtraction of $\log\pi$. The conductor and $\log\ell$
terms scale by exactly $1-\kappa$. The static pole pair contributes at
most $P/u+P$ upon dropping both subtracted kernels. Thus we add
$\kappa P$ after multiplying $C$ by $1-\kappa$.

For $|T|\ge262144$, at each of the two real abscissas the shifted gamma
sum differs from $g\log(|T|/(2\pi))$ plus the $\log\ell$ terms by
less than $0.001$. Indeed $|\psi(z)-\log z|\le1/|z|$ on the right
half-plane. Integrate the Laplace identity in Section~\ref{hc:chapter} by
parts: the kernel $1/(1-e^{-v})-1/v$ increases from $1/2$ to $1$,
since $2\sinh(v/2)\ge v$. For $z=(a+\mu+i\ell T)/2$ with
$a\in\{\sigma,\sigma_1\}$ and $0\le\mu\le4$, the two errors in
$\Re\psi(z)-\log|\ell T/2|$ are at most $2/|T|$ and $18/T^2$.
Multiply their sum by $g$. The shifted zeta pole corrections at
$\pm2T$ and $\pm4T$ total at most
$\sum_{\ell=2,4}w_{\ell0}(2\sigma-1)/(\ell T)^2<0.001$,
after dropping the subtracted kernels. Thus the available
constant, including all these errors, satisfies
\[
 C'=(1-\kappa)C+\kappa P+0.004<2.456
       <2.638<(1-\kappa)g\log(2\pi).
\]
It follows that \eqref{exstrip:ineq} applies with $P=11/24$ and
$B(T)=(1-\kappa)g\log|T|$. The optimized
$u=(\sqrt P-P)/B(T)$ is less than $0.01221$. The archived \texttt{constants.py} computes the character table exactly. Its real interval evaluation verifies every numerical inequality. In particular
\[
      \frac{(1-1/\sqrt5)(187/72)}{(1-\sqrt{11/24})^2}
                <13.762<14 .
\]
This proves the strict enclosure. The lower window comes from the finite-height elliptic inventory, row of level eleven, with its complete bracket-and-count record \texttt{hybrid-11-T524288.json}.
\end{proof}

\begin{theorem}[Arbitrary elliptic reduction]\label{exstrip:wild}
For every elliptic curve $E/\Q$ put
$P^*=11/24+\mathbf1_{\rm CM}(E)/18$. At $|T|\ge262144$ every
zero $\beta+iT$ satisfies
\[
  \min(\beta,1-\beta)\ge
   \frac{(1-\sqrt{P^*})^2}
         {(1-1/\sqrt5)\bigl((187/72)\log|T|+(73/36)\log(N/11)\bigr)} .
\]
\end{theorem}
\begin{proof}
We show positivity of the five-frequency square before estimating
the conductor. At a place of potentially good reduction the
normalized Frobenius-semisimple Weil parameter $V$ has zero
monodromy, finite inertia image $I$, and image in $\mathrm{SU}(2)$
after conjugation. Indeed some positive power of normalized Frobenius
commutes with $I$ and is diagonalizable with all eigenvalues of modulus
one by potential good reduction and purity. The closure of the generated
group is then compact, and its determinant is one. At the $\nu$th
prime power the trace used in the Euler logarithm is
\[
   \operatorname{tr}(\mathrm{Fr}^{\nu}\mid(\operatorname{Sym}^j V)^I)
     =\frac1{|I|}\sum_{a\in I}
               \operatorname{tr}\operatorname{Sym}^j(\mathrm{Fr}^{\nu}a).
\]
This follows by inserting the averaging projection on invariants.
The same averaging measure and the same frequency $\phi$ occur in
every character, so the local coefficient is an average of
nonnegative squared norms.

At a potentially multiplicative place the parameter is
$\chi\otimes\operatorname{Sp}(2)$ for a quadratic or trivial
character $\chi$. For unramified $\chi$ the preceding Gram
calculation applies with $x=(\chi(\mathrm{Fr})p^{-1/2})^\nu$.
If $\chi$ is ramified then only even $j$ have a surviving
monodromy-kernel invariant. Averaging the two Gram evaluations at
$x=p^{-\nu/2}$ and $-x$ effects exactly this projection and
preserves positivity. These two types include wild additive reduction;
see \href{https://magma.maths.usyd.edu.au/~watkins/papers/ef10.pdf}
{Dummigan, Martin and Watkins, \emph{Euler factors and local root
numbers for symmetric powers of elliptic curves}} for the local parameters.

For $j\ge1$ the conductors satisfy $Q_j\le N^j$. In the potentially
good case every inertia subgroup $H$ with nontrivial image has
$V^H=0$: a finite-order element in $\mathrm{SL}_2$ fixing a
nonzero vector must be the identity. Thus
\[
 \operatorname{codim}(\operatorname{Sym}^j V)^H
   \le j+1\le 2j=j\operatorname{codim} V^H
\]
for such $H$; both codimensions vanish otherwise. Summing over
the ramification filtration gives the Artin-conductor inequality.
In the potentially multiplicative case, if $\chi$ is unramified
then the conductor exponents of $V$ and $\operatorname{Sym}^j V$
are $1$ and $j$. For ramified $\chi$ the exponent of $V$ is
$2a(\chi)$. An odd symmetric power has no inertia invariants and
has exponent $(j+1)a(\chi)$; an even power has unramified Weil
part and a single Jordan block, giving exponent $j$. Both are at
most $j\,a(V)$. Multiplication over the primes proves the claim.

Consequently the conductor term in the preceding proof increases
by at most $(1-1/\sqrt5)(73/36)\log(N/11)$. For CM the static
fourth symmetric power has one additional trivial middle constituent
of weight $1/18$. The second symmetric power, including at the
shifted frequencies, has quadratic middle constituent, and all the
odd constituents are holomorphic. The extra pole at zero adds at
most $1/18$ to $C'$; numerically $C'+1/18<2.511<2.638$. The
coefficient of $1/u$ is now $P^*$. For both values of $P^*$ the
optimized $u$ is below $0.01437$, even using only $N\ge1$.
Rearranging the optimized inequality and reflecting proves the result.
\end{proof}

\begin{corollary}\label{exstrip:primes}
For each fixed elliptic curve $E/\Q$ put $R_E=14$ if it is non-CM and $R_E=18$ if it is CM. Then
\[
 \sum_{\substack{p\le x\\p\nmid N}}\frac{a_p(E)}{\sqrt p}\log p
                 \ll_E x\exp(-\sqrt{(\log x)/R_E}).
\]
\end{corollary}
\begin{proof}
The preceding conductor adjustment is constant for fixed $E$. The coefficients of $\log|T|$ in the optimized denominator are less than $13.762$ and $17.909$ respectively. Choose $d_0>1/R_E$ slightly below the reciprocal of the applicable coefficient. At sufficiently large heights every zero has $\beta\le1-d_0/\log|\gamma|$. The finite intervening ordinates are separated from $\Re s=1$ by the general bound. In the classical truncated explicit formula for the unitary elliptic Euler derivative take $T=\exp(\sqrt{d_0\log x})$, moving it by at most one to a regular ordinate. At a half-integer cutoff its truncation error is $O_E(x\log^C(2xT)/T+\log^C(2xT))$ with a fixed $C$; this is the usual Perron estimate using $|\Lambda_E(n)|\le2\Lambda(n)$ and the unit-interval zero count (see Iwaniec--Kowalski, \emph{Analytic Number Theory}, Chapter~5). The sum over high zeros is $O_E(x\exp(-d_0\log x/\log(T+1))\log^2(2T))$ by partial summation of their counting function. Each of the finitely many lower zeros saves a fixed power of $x$. Polynomial logarithmic factors are absorbed since $\sqrt{d_0}>1/\sqrt{R_E}$. Half-integer cutoffs give all sharp sums by replacing $x$ by $\lfloor x\rfloor+1/2$. Higher prime powers cost $O(\sqrt x\log^2(2x))$ and the finitely many bad primes contribute $O_E(\log(2x))$.
\end{proof}

One further integer choice sharpens the non-CM constant. Substitute
$(7,20U_1,23+15U_2,20U_1,7)$ in the same square and divide by its
$U_1\cos\phi$ coefficient $3600$. Its values of $P,g$ and the
half-log-conductor coefficient are $413/900$, $9161/3600$ and
$1453/720$. The same interval replay gives $C'<2.400<2.585<
(1-\kappa)g\log(2\pi)$ and $u<0.0147$ for $N\ge1$.
Consequently $13.518$ can replace $14$ at level eleven and
$1/\sqrt{13.518}$ can replace $1/\sqrt{14}$ in the prime-sum
error for every fixed non-CM curve.

For comparison, \href{https://arxiv.org/abs/2509.20873}{Creech,
Hamieh, Khunger, Sinha, Streipel and Tsang, Theorem~1.1} give
the denominator $16.7053\log(kN|T|)$ for $|T|\ge1$ for newforms
of squarefree level, even weight and trivial nebentypus. They use
Stechkin differences, a quartic trigonometric polynomial and symmetric
powers through the fourth. Earlier explicit results include
\href{https://doi.org/10.1093/qmath/haac004}{Hoey, Iskander, Jin and
Trejos Su\'arez, Section~4}. The five polynomials displayed here give
a smaller denominator for $11\mathrm a1$ above the verified range
and a smaller limiting coefficient of $\log|T|$ at each fixed
non-CM elliptic curve. Theorem~\ref{exstrip:wild} also covers additive reduction.

\clearpage
\section{Explicit original zero strips over number fields}\label{rqstrip:section}
The ramified-place proof of Theorem~\ref{exstrip:wild} and its five-frequency square allow an absolute-field calculation. In this section $r=[K:\Q]$, $d_K$ is the absolute discriminant and $N=\mathrm N_{K/\Q}\mathfrak f$ is the norm of the finite conductor ideal. The constants from the square are
\[
        G=187/72,\quad A=73/36,\quad B=529/144,\quad
        \kappa=1/\sqrt5 .
\]
Its shifted half-degree is $G$, its static half-degree is $155/144$, and $B$ is their sum. The conductor-ideal weight $A$ is half the sum of $j w_{\ell j}$ over the table of Section~\ref{exstrip:chapter}.

\begin{theorem}\label{rqstrip:high}
Let $E$ be a modular elliptic curve over a totally real field $K$. This includes \emph{every} elliptic curve over \emph{every} real quadratic field. Use $P=11/24$ for a non-CM curve and $P=37/72$ for a CM curve. Every nontrivial zero $\beta+iT$ of the original unitary Hasse--Weil function at $|T|\ge262144$ satisfies
\[
 \min(\beta,1-\beta)\ \ge\
  \frac{(1-\sqrt P)^2}{(1-\kappa)\{rG\log |T|+A\log(N/11^r)+B\log d_K\}}.
\]
For a fixed curve, at sufficiently large heights the exclusion at each edge has width at least
\[
 \begin{cases}
   1/(13.762\,r\log |T|), & E\text{ without CM},\\
   1/(17.909\,r\log |T|), & E\text{ with CM}.
 \end{cases}
\] A second formula covers the original zeros at every height for all these curves and for holomorphic Hilbert newforms as follows.
\end{theorem}

\begin{theorem}\label{rqstrip:hilbert}
Let $\pi$ be a Hilbert cuspidal newform with trivial central character and even weights $k_v\ge2$ at all real places of $K$. Set $P_0=65/72$ for a non-CM representation and $P_0=69/72$ for a CM representation. Put
\[
  C_\pi(T)=A\log N+B\log d_K+(4B+2)r
                + B\sum_{v\mid\infty}\log(8+4 k_v+4|T|).
\]
Then every nontrivial unitary zero satisfies
\[
            \min(\beta,1-\beta)\ge (1-\sqrt{P_0})^2/C_\pi(T).
\]
For each fixed non-CM representation as above, the asymptotic bound of Theorem~\ref{rqstrip:high} with coefficient $13.762\,r$ holds at sufficiently large heights also at higher weights.
\end{theorem}

\begin{proof}[Proof of the two theorems]
Use the symmetric powers through the fourth. For non-CM Hilbert modular forms their cuspidal lifts with matching local parameters follow from \href{https://doi.org/10.4007/annals.2026.203.1.4}{Newton and Thorne, \emph{Symmetric power functoriality for Hilbert modular forms}}. For CM use induction of the powers of the Hecke character, with the middle finite constituent as in Section~\ref{exstrip:chapter}. \href{https://www.math.ucla.edu/~blasius/papers/hmfrc1005.pdf}{Blasius's theorem} gives temperedness at the finite places for the stated holomorphic weights. The triple of geometric inputs in the elliptic case is modularity for the specified curve, compatibility at the finite places, and weight two at each real place. For the real quadratic case the modularity theorem is \href{https://arxiv.org/abs/1310.7088}{Freitas--Le Hung--Siksek, \emph{Elliptic curves over real quadratic fields are modular}}.

Write $w_{\ell j}$ for the entries of the five-frequency table. At any finite place the proof of Theorem~\ref{exstrip:wild} applies to the unitary determinant-one two-dimensional parameter. With monodromy zero its inertia image is finite and the Frobenius action is bounded in the tempered normalization. The prime-power coefficients average the same nonnegative square over this finite inertia image. Each nontrivial ramification subgroup has no fixed vector in the original two-dimensional space: a finite-order determinant-one matrix with eigenvalue one is the identity. Hence the codimension of invariants of its $j$th symmetric power is at most $j+1\le2j$ times the indicator of such a subgroup. With nonzero monodromy the parameter is a special representation twisted by a character whose square is trivial. The positive Gram evaluation, its two-sign average, and the conductor computation are then exactly those of Theorem~\ref{exstrip:wild}. Thus the local square is nonnegative at every prime power and $\mathrm N\mathfrak f_j\le N^j$ for $j\ge1$. The absolute conductor of the $j$th power is $d_K^{j+1}\mathrm N\mathfrak f_j$. Its total half-logarithm charge before differencing is at most $A\log N+B\log d_K$.

For a weight-two curve over $K$, the archimedean factors are $r$ copies of the weight-two factors over $\Q$. Apply the differenced Hadamard argument of Section~\ref{exstrip:chapter}. Multiply its gamma and gamma-error estimates by $r$. Each Dedekind zeta pole has order one in the completed function, so its regular pole cost is no larger than $r$ times the corresponding cost already allowed in that reference estimate. Keep its singular term at the static frequency at exactly $P/u$. Its coefficient at the selected target zero is one. Relative to $r$ copies of the reference level $11$, the extra arithmetic charge is at most $(1-\kappa)\{A\log(N/11^r)+B\log d_K\}$. This gives the first denominator after optimizing in $u$. To check its range, the brace is at least $r(G\log262144-A\log11)$ and the two choices of $\sqrt P-P$ are at most $1/4$. The resulting $u<0.02$ even at $r=1$. Both denominators before rounding are computed in the accompanying interval check. Reflection gives the left-hand bound.

At arbitrary $T$ use the original square without the Stechkin subtraction. Bound every pole at one by $1/u$ and its companion at zero by $1$ at $1+u$ with $0<u\le0.02$. The static and oscillating coefficients at the trivial representation total $11/24+5/12+1/36=65/72$. The only extra pole in the CM case has weight $1/18$ from the fourth power; the positive angular characters have nonzero infinity type because every $k_v\ge2$. Thus the total pole charge is at most $P_0/u+P_0$. For $j\le4$ every real gamma shift at a real place is nonnegative and at most $2(k_v-1)+1$. The elementary digamma estimate $|\psi(z)-\log z|\le1/\Re z$ bounds one real gamma contribution at $(1+u+\mu+i\ell T)/2$ by $2+\frac12\log(8+4k_v+4|T|)$. There are $2B$ such factors per real place counted with their weights. Together with the conductor cost this gives $1/(1+u-\beta)\le P_0/u+C_\pi(T)$ for $\beta\ge1/2$. The optimal $u=(\sqrt{P_0}-P_0)/C_\pi(T)$ is in the required range by the displayed positive constants. The general formula follows. For fixed higher weights at $|T|\to\infty$, the differenced static contributions remain bounded and Stirling gives $(1-\kappa)rG\log|T|+O_\pi(1)$ for all remaining charges. The proof for the first theorem now yields the fixed-form asymptotic assertion.
\end{proof}

\begin{corollary}\label{rqstrip:generalfield}
Let $K$ instead be any number field, and let $\pi$ be a unitary cuspidal $\mathrm{GL}_2$ representation with trivial central character and cuspidal symmetric square, cube and fourth-power lifts. Then
\[
 \liminf_{|\Im\rho|\to\infty}\log|\Im\rho|\min(\Re\rho,1-\Re\rho)
                \ \ge\ \frac{(1-\sqrt{11/24})^2}{(1-\kappa)G[K:\Q]}.
\]
The assertion applies in particular to a modular non-CM elliptic curve over any number field. No Ramanujan conjecture for general Maass forms is used. Here the height threshold after rounding the asymptotic constant can depend on the fixed representation.
\end{corollary}
\begin{proof}
The symmetric square, cube and fourth-power transfers used here are the Gelbart--Jacquet and Kim--Shahidi--Kim theorems; for the non-polyhedral cuspidality criterion see \href{https://www.math.purdue.edu/~fshahidi/articles/Kim%20%26%20Shahidi%20%5B2002%2C%2021pp%5D---Cuspidality%20of%20symmetric%20powers%20with%20applications.pdf}{Kim and Shahidi, \emph{Cuspidality of symmetric powers with applications}}. For a modular non-CM elliptic curve, the solvable polyhedral exceptions in the Kim--Shahidi criterion would supply a nontrivial finite self-twist of $\pi$, $\operatorname{Sym}^2\pi$ or $\operatorname{Sym}^3\pi$. Matching good Frobenius classes and Chebotarev transfer such a self-twist to the corresponding symmetric power of the elliptic Galois representation. Serre's open-image theorem keeps that power irreducible over the finite kernel field of the twist. An intertwiner restricted there is therefore scalar; the original self-twist would be trivial. This proves the required cuspidality in this case. Remove the fixed finite set of finite places at which the unramified identities are unavailable. Their Euler logarithmic derivatives in the actual cuspidal functions are $O_\pi(1)$ uniformly at $\sigma\ge1$ by the unitary local bounds. At every other prime power the unramified parameter has determinant one and real trace, as follows from unitarity and self-duality. The identity of characters producing the square is a polynomial identity in this real trace. It therefore remains nonnegative even for a complementary-series parameter. The logarithmic-derivative sums converge absolutely at $\sigma>1$ for each cuspidal lift, by Cauchy--Schwarz with its Rankin--Selberg series at the unramified primes. All static archimedean terms and absolute conductors are bounded for these fixed forms. The shifted gamma factors number $2rG$ after expanding all complex gamma factors into real ones, and Stirling at $\ell T$ gives the same leading coefficient $(1-\kappa)rG$. The only static pole weight is $11/24$ and the shifted poles cost $O(1)$ at these heights. The differenced positivity argument and its optimizing $u$ prove the result.
\end{proof}

\subsection*{A further saving from the rank-one zero-free half-plane}
The pole at the static frequency is essential. The shifted degree-one factors can instead be estimated on their finite Euler side, producing a stricter asymptotic constant for the original degree-two zeros.

\begin{corollary}\label{rqstrip:rankone}
Under the hypotheses of Corollary~\ref{rqstrip:generalfield}, assume also that $\zeta_K$ is zero-free in some fixed half-plane $\Re s>\theta$ with $\theta<1$, apart from its pole. This holds for every abelian $K/\Q$ by the Dirichlet theorem. With $r=[K:\Q]$ one obtains
\[
 \liminf_{|\Im\rho|\to\infty}r\log|\Im\rho|\min(\Re\rho,1-\Re\rho)
   \ \ge\ \frac{(1-\sqrt{3947/8556})^2}{(1-1/\sqrt5)(6549/2852)}
    >\frac1{12.335} .
\]
It follows that every fixed non-CM rational elliptic curve, and every fixed non-CM elliptic curve over a real quadratic field, has a high-ordinate unitary zero-free width at least $1/(12.335\,r\log |T|)$ when $|T|$ is sufficiently large. The lower height for this strengthened constant is allowed to depend on the individual function.
\end{corollary}
\begin{proof}
In the trigonometric square take $(12,31U_1,36+21U_2,31U_1,12)$. The coefficient of $U_1\cos\phi$ before division is $8556$. The unshifted trivial coefficient after division is $3947/8556$. For $\ell\ge1$ leave all the $w_{\ell0}$ terms in the finite logarithmic derivative before applying the completed explicit formula to the other terms. Uniformly for $1\le\sigma\le2$ and large $|t|$, the hypothesis gives
\[
                 \frac{\zeta_K'}{\zeta_K}(\sigma+it)
                  =O_K((\log(3+|t|))^{1-\delta})=o_K(\log|t|)
\]
for some $\delta>0$. Indeed on disks strictly inside $\Re s>\theta$ at large height the logarithm from absolute convergence has modulus $O(\log |t|)$ by Borel--Carath\'eodory and the fixed-strip polynomial bound. It is bounded on a subdisk in absolute convergence. Take intermediate disks by three circles and differentiate by Cauchy with a fixed margin. The pole is outside these high disks. For an abelian field the factorization of its Dedekind function into primitive Dirichlet functions supplies $\theta=1/2$. The sum of all omitted scalar logarithmic derivatives, at $\ell T$ and both real abscissae of Stechkin's difference, is thus $o_K(\log |T|)$. This procedure estimates genuine finite Euler terms and charges no archimedean degree for them.

The half-degree cost of the shifted nontrivial terms is
\[
       \frac12\sum_{\ell\ge1,\,j\ge1}(j+1)w_{\ell j}=6549/2852.
\]
For an elementary check the dimensions of the five polynomials at the identity are $12,62,99,62,12$. The full shifted half-degree before dividing by $8556$ is the sum of the products of dimensions at distinct indices $i<k$. Subtract $\sum_{i<k}\langle P_i,P_k\rangle$, using the usual orthonormality of the $U_j$, to discard the shifted trivial terms. The static numerator is $2(12)^2+2(31)^2+36^2+21^2=3947$. The remaining completed terms obey precisely the pairwise positivity used in Corollary~\ref{rqstrip:generalfield}. Their upper bound has coefficient $(1-\kappa)r(6549/2852)\log|T|+o_\pi(\log|T|)$, together with the static pole $P/u$. Taking $u=(\sqrt{P}-P)/[(1-\kappa)r(6549/2852)\log|T|]$ gives the displayed limit. Rational interval evaluation places the reciprocal constant between $12.3343$ and $12.3344$.
\end{proof}

The alternative square at the end of Section~\ref{exstrip:chapter} also improves the unconditional non-CM constants in Theorem~\ref{rqstrip:high}. There one may simultaneously replace $(P,G,A,B)$ by
\[
             (413/900,\ 9161/3600,\ 1453/720,\ 729/200).
\]
Indeed the new total half-degree is $162^2/(2\cdot3600)$ and all the calibrated archimedean inequalities there scale by $r$ in exactly the same way. In Corollary~\ref{rqstrip:generalfield} this gives the asymptotic coefficient $13.518\,r$ over arbitrary fields, unconditionally. The $12.335\,r$ statement uses the additional scalar control just proved.

For a fixed non-CM elliptic curve over $\Q$, the same scalar saving gives the following direct prime-sum consequence of the rank-one input:
\[
 \sum_{\substack{p\le x\\p\nmid N_E}}\frac{a_p(E)}{\sqrt p}\log p
          \ll_E x\exp\!\left(-\sqrt{\frac{\log x}{12.335}}\right).
\]
Indeed the reciprocal coefficient before rounding in Corollary~\ref{rqstrip:rankone} is less than $12.3344$, leaving the strict margin required to absorb the powers of $\log x$ in the truncated explicit formula in Corollary~\ref{exstrip:primes}. The same deduction holds with denominator $12.335\,r$ in the field prime-ideal sum for every fixed modular non-CM elliptic curve in that corollary. In a fixed base field, after grouping by rational prime powers the elliptic von Mangoldt coefficients are bounded by $2r\Lambda(n)$, and higher prime-ideal powers cost $O_K(\sqrt x\log^2 x)$.

\begin{theorem}\label{rqstrip:effective}
Let $K/\Q$ be totally real and abelian. For a modular non-CM elliptic curve over $K$ put
\[
  (P,g,h,b,Z)=\left(\frac{3947}{8556},\frac{6549}{2852},\frac{16615}{8556},\frac{1841}{552},\frac{1969}{4278}\right).
\]
Write $m(L)=\min(2\sqrt L+3,\,2\log L+10)$. At every $|T|\ge262144$ define
\begin{align*}
 D_K(T)={}&(1-\kappa)\{rg\log|T|+h\log(N/11^r)+b\log d_K\}\\
          &+rZ\,m(\log(4d_K|T|)).
\end{align*}
Every original unitary zero then satisfies
$\min(\beta,1-\beta)\ge(1-\sqrt P)^2/D_K(T)$.
For the conductor-eleven curve the width is greater than $1/(13\log|T|)$ for $\log|T|\ge200$, greater than $1/(12.5\log|T|)$ for $\log|T|\ge1000$, greater than $1/(12.35\log|T|)$ for $\log|T|\ge10000$, and greater than $1/(12.335\log|T|)$ for $\log|T|\ge300000$.
\end{theorem}
\begin{proof}
Take $1<\sigma\le1.02$, $t=\ell T$ with $\ell=2,4$, and a primitive Dirichlet constituent $\chi$ of $\zeta_K$. Its conductor $q$ is at most $d_K$, by the conductor--discriminant formula. For $1\le v\le3$ the real Hadamard identity and the Dirichlet zero theorem give
\[
 -\Re\frac{L'}{L}(v+it,\chi)
 = H(v,t)-\sum_\gamma\frac{v-1/2}{(v-1/2)^2+(t-\gamma)^2},
 \qquad
 \left|H(v,t)-\tfrac12\log\frac{q|t|}{2\pi}\right|<0.001 .
\]
Indeed $H$ combines the conductor, one real gamma factor of shift zero or one, and the two real pole kernels in the trivial case. The preceding digamma bound gives errors at most $1/|t|+4/t^2$ from the gamma term and $5/t^2$ from the poles. These estimates are uniform in the character.

Set $L_0=\log(4d_K|T|)$ and $\delta=L_0^{-1/2}$. For $x=\sigma-1/2$, the kernel $K_x(y)=x/(x^2+y^2)$ satisfies $K_x(y)\ge[x/(x+\delta)]K_{x+\delta}(y)$. At $v=\sigma+\delta$, absolute convergence bounds the logarithmic derivative by $-\zeta'(1+\delta)/\zeta(1+\delta)\le1/\delta+1$. This last inequality follows from the real zeta Hadamard identity by dropping the zero terms and the negative gamma contribution. Substituting the kernel comparison in the first identity gives
\[
 -\Re(L'/L)(\sigma+it,\chi)
   \le \delta L_0+1/\delta+1.002 .
\]
Here $x\ge1/2$ and $1\le\sigma+\delta\le3$. At $\sigma_1=(1+\sqrt{1+4\sigma^2})/2$, absolute convergence similarly gives $\kappa |L'/L|\le\kappa(1+1/(\sigma_1-1))<1.172$. The full differenced Euler term of one character thus costs at most $2\sqrt{L_0}+3$. A short smoothed prime sum supplies the logarithmic alternative in $m$. Put $X=36L_0^2$ and $J(w)=((eX)^w-X^w)/w^2$. The residue formula for the integral of $-(L'/L)(s+w,\chi)J(w)$ on $\Re w=2$, with $s=\sigma+it$, bounds $|L'/L(s,\chi)|$ by
\[
  \sum_{n\le eX}\frac{\Lambda(n)}n
  +\sum_{\rho}|J(\rho-s)|+
   \sum_{\eta\ {\rm trivial}}|J(\eta-s)|+
   \mathbf1_{q=1}|J(1-s)|.
\]
The residue at zero has coefficient one since $\log(eX)-\log X=1$. Inverse Mellin gives weights $1$ up to $X$ and $\log(eX/n)$ from $X$ to $eX$. For the contour shift take horizontal lines a distance $\gg1/\log(q(|t'|+3))$ from the ordinates of zeros and then left lines between the trivial zeros. The local partial-fraction bound on each fixed strip gives $O(\log^2(q(|t'|+3)))$ on these horizontal lines (see \href{https://dms.umontreal.ca/~koukoulo/documents/publications/primes.pdf}{Koukoulopoulos, \emph{The Distribution of Prime Numbers}, Lemma 11.4}); the functional equation bounds the logarithmic derivative on the left lines. The $w^{-2}$ decay followed by exponential decay in the negative real direction removes the three added sides.

The sum over nontrivial zeros costs at most
\[
 \frac{1+e^{-1/2}}{\sqrt X}\sum_\gamma
      \frac{1}{(\sigma-1/2)^2+(t-\gamma)^2}
    \le \frac{1+e^{-1/2}}{\sqrt X}(3L_0+12.006).
\]
Indeed the kernel at the left is bounded pointwise by $6K_{3/2}$. Use the real Hadamard identity at $2+it$ and $|L'/L(2+it)|\le2$. The trivial sum is at most $(2/X)\sum_{j\ge0}(1+2j)^{-2}<4/X$, and the pole costs at most $2/t^2$. Finally, for $y\ge2$,
$\sum_{n\le y}\Lambda(n)/n\le\log y+3$. To see this, put $k=\lfloor y\rfloor$, expand $\log(k!)=\sum_{n\le k}\Lambda(n)\lfloor k/n\rfloor$ and use $\sum_{n\le k}\Lambda(n)<3k$. The latter follows by telescoping over dyadic endpoints, since every prime power in $(2^{a-1},2^a]$ contributes its prime logarithm to $\log {2^a\choose 2^{a-1}}\le2^a\log2$. Adding the earlier $1.172$ for the subtracted abscissa gives the second upper bound $2\log L_0+10$, because $L_0>13$. There are $r$ characters and the sum of their frequency weights per character is $Z$.

For clarity, normalize the five polynomials in the preceding corollary by their target coefficient after squaring. The remaining conductor weights are $h=\sum j w_{\ell j}/2$ and $b=\sum(j+1)w_{\ell j}/2-Z/2$. In the reference-level constant $C$ of Section~\ref{exstrip:chapter}, use these new weights, retaining only $j\ge1$ in the shifted sum. Its residual after differencing and removal of $g\log(2\pi)$ is bounded above by
\[
 (1-\kappa)\{C-g\log(2\pi)\}+\kappa P+0.005<-0.054 .
\]
The allowance $0.005$ exceeds $(1+\kappa)g(2/262144+18/262144^2)$. The archived interval calculation \texttt{rankone\_saving.py} verifies the inequality using the displayed real digamma constants. In $r$ copies the reference residual is consequently nonpositive. Changing the conductor contributes the two logarithmic terms in $D_K(T)$. The optimization in Theorem~\ref{rqstrip:high} now applies verbatim; even with $N=d_K=1$ the optimizing $u$ is less than $0.02$. For $K=\Q$, $N=11$, the ratio $D_K(T)/[(1-\sqrt P)^2\log|T|]$ decreases with $|T|$. Interval evaluation at the four stated logarithmic heights gives upper bounds $13$, $12.5$, $12.35$ and $12.335$, respectively, with strict inequalities.
\end{proof}

\begin{corollary}\label{rqstrip:arbitraryScalar}
For every number field $K$ the coefficient $12.335\,[K:\Q]$ in Corollary~\ref{rqstrip:rankone} applies at sufficiently large heights to each fixed $\pi$ with the stated cuspidality conditions. Furthermore, if $K$ is totally real, the effective bound for modular non-CM elliptic curves in Theorem~\ref{rqstrip:effective} holds over $K$ with $m(L)$ replaced by $8\log L+40$. The field may be non-Galois.
\end{corollary}
\begin{proof}
For the first claim it suffices to know that $\zeta_K$ has a fixed zero-free half-plane short of one. Apply \href{https://hexagonmath.org/2610.00200v1}{Cureton--Thorner, \emph{A quasi-Riemann hypothesis for Dedekind zeta functions}}, with boundary $31/32$, or \href{https://hexagonmath.org/2610.00216v2}{Hu, Theorem 1.2}, with boundary $7/8$, at the trivial character. Both are external inputs. The proof of Corollary~\ref{rqstrip:rankone} now applies.

For the effective claim use Hu's endpoint. Every nontrivial zero $\beta+i\gamma$ of $\zeta_K$ has $1/8\le\beta\le7/8$. In the previous smoothed formula take $X=(64 L_0)^8$. For $1<\sigma\le1.02$,
\[
 \frac{1}{(\sigma-\beta)^2+y^2}
  \le 72\frac{2-\beta}{(2-\beta)^2+y^2}.
\]
Indeed the ratio is bounded by the larger of $1/(2-\beta)$ and $(2-\beta)/(1-\beta)^2$, whose maximum on the stated interval is $72$. At $2+it$ the real Hadamard identity and $|\zeta_K'/\zeta_K|\le2r$ bound the sum of the kernels on the right, before multiplication by $72$, by $r(L_0/2+2.001)$. Here $\tfrac12\log d_K$ is included in $rL_0/2$ and the $r$ real gamma factors satisfy the preceding estimates. After grouping by rational prime powers, the coefficient bound $\Lambda_K(n)\le r\Lambda(n)$ gives the prime-sum bound $r(\log X+4)$, and an upper bound $1.172r$ at the subtracted abscissa. The nonpositive even zeros have multiplicity at most $r$. The simple pole costs at most $2/t^2$. Thus the scalar cost beyond $8r\log L_0$ is at most $r$ times
\[
 8\log64+4+1.172+
  \frac{1+e^{-1/8}}{64}\left(36+\frac{144.072}{L_0}\right)
       +\frac4X+\frac{2}{t^2}<40 .
\]
The application of the smoothed formula to $\zeta_K$ uses its classical functional equation and finite-order Hadamard product; on horizontal lines the corresponding local partial-fraction bound follows by that product and the fixed-field local zero count. Every remaining term of the degree-two amplifier is the same as in Theorem~\ref{rqstrip:effective}. This gives the claimed explicit bound.
\end{proof}

\subsection*{A degree-six polynomial and positive local measures}
The five-frequency square has an algebraic proof of positivity. For elliptic curves one can enlarge the search to polynomials nonnegative on the entire trace interval. There is a local reason this still works at primes of multiplicative reduction.

\begin{lemma}\label{rqstrip:measure}
At every prime power of a modular elliptic curve the unitary traces of the symmetric-power Euler factors, indexed by $j\ge0$, are the moments of $U_j(x)$ under a single probability measure on $[-1,1]$. Here $U_j(\cos\theta)=\sin((j+1)\theta)/\sin\theta$ denotes the usual Chebyshev character, with argument equal to half the two-dimensional trace.
\end{lemma}
\begin{proof}
For potential good reduction, take the distribution of the half trace of $\mathrm{Fr}^\nu a$ for uniform $a$ in the finite inertia image, in the unitary parameter of Theorem~\ref{exstrip:wild}. The average computes the trace on invariants. For the special representation with unramified quadratic twist the Euler traces are $b^j$, with $b$ the signed inverse square root of the norm of the prime power. Write $d\mu(x)=2\sqrt{1-x^2}\,dx/\pi$. Orthogonality of the $U_j$ and their generating series give
\[
 \int_{-1}^1 U_j(x)\frac{d\mu(x)}{1-2bx+b^2}=b^j .
\]
The series is uniformly convergent since $|b|<1$ and $|U_j(x)|\le j+1$. The density is positive and the zeroth moment is one. For a ramified quadratic twist take the mean of the measures at $b$ and $-b$, which retains exactly the even powers. These are all reduction types.
\end{proof}

\begin{theorem}\label{rqstrip:six}
For every fixed modular non-CM elliptic curve over a totally real number field $K$ the unitary high-ordinate zero-free width is at least $1/(12.326\,r\log|T|)$ for sufficiently large $|T|$. For the conductor-eleven curve this bound holds for $\log|T|\ge10^6$, and the width exceeds $1/(12.34\log|T|)$ already for $\log|T|\ge10000$. The same asymptotic coefficient applies over any number field to a fixed trivial-central-character cuspidal $\mathrm{GL}_2$ representation for which the symmetric lifts of indices $2,3,4,6$ are cuspidal with the unramified symmetric-power factors.
\end{theorem}
\begin{proof}
Define $A(x,y)=\sum w_{\ell j}U_j(x)T_\ell(y)$, where $T_\ell(\cos\phi)=\cos(\ell\phi)$. Here are all the integer coefficients $10^6w_{\ell j}$, with columns indexed by $j$:
\[
\begin{array}{r|rrrrrr}
 \ell&0&1&2&3&4&6\\\hline
 0&462244&0&444674&0&47038&1191\\
 1&0&1000000&0&293292&0&0\\
 2&433762&0&343537&0&0&348\\
 3&0&179116&0&1183&0&0\\
 4&33882&0&0&0&0&31\\
 6&439&0&110&0&0&0
\end{array}
\]
We verify $A(x,y)\ge0$ for all real $x$ and $-1\le y\le1$ by rational arithmetic. For the interior put $(x,y)=(2u-1,2v-1)$, with $u,v\in[0,1]$. For the exterior $x\ge1$ put $x=1/(1-u)$ and multiply the polynomial by $(1-u)^6$. Simultaneous reversal $(x,y)\mapsto(-x,-y)$ covers the other side. If $c_{ij}$ are the ordinary monomial coefficients of one of these polynomials, its bidegree-six Bernstein coefficients are
\[
 B_{ab}=\sum_{i\le a,\,j\le b}c_{ij}
 \frac{\binom ai\binom bj}{\binom6i\binom6j}.
\]
The Bernstein basis is nonnegative. Stop on any rectangle where all $B_{ab}\ge0$. Otherwise bisect in $u$ at even depths and in $v$ at odd depths. Each one-dimensional subdivision uses the endpoints of the successive rows obtained by replacing adjacent entries by their means. For the interior the numbers of accepted rectangles at depths $3,4,\dots,16$ are, in order,
\[
 2,\ 2,\ 2,\ 10,\ 14,\ 24,\ 30,\ 44,\ 68,\ 102,\ 156,\ 320,\ 152,\ 16 .
\]
For the exterior all initial coefficients are nonnegative. The exact integer recurrence and verification are in \texttt{bernstein\_amplifier.py}; it uses a common positive denominator for each array and checks the depth bound before returning. This proves global positivity of this particular polynomial.

At good primes evaluate $y=\cos(\nu T\log \mathrm N\mathfrak p)$. At bad elliptic primes apply Lemma~\ref{rqstrip:measure} to $A(\cdot,y)$. All Euler contributions are thus nonnegative. For the fixed automorphic version omit the finitely many bad primes as in Corollary~\ref{rqstrip:generalfield}; the unramified half trace is real, so the exterior check covers complementary-series places too. The nonnegative character coefficients allow precisely the same Stechkin subtraction and target-zero selection. Newton--Thorne gives the needed symmetric powers for the totally real non-CM curves. For the present weights the constants $(P,g,h,Z,b)$ in the effective scalar estimate are
\[
 \left(\frac{115561}{250000},\frac{2284863}{1000000},
       \frac{767351}{400000},\frac{468083}{1000000},
       \frac{6609519}{2000000}\right).
\]
The degree-six static shifts increase only to four. The effective proofs above therefore apply with these constants and $L_0=\log(6d_K|T|)$, allowing the sixth frequency. Interval arithmetic gives a reference residual less than $-0.050$ and the optimizing $u<0.02$ throughout $|T|\ge262144$. The asymptotic denominator $(1-\kappa)g/(1-\sqrt P)^2$ is less than $12.32555$. At conductor eleven the two stated explicit cutoffs follow from the same decreasing ratio as in Theorem~\ref{rqstrip:effective}, with upper bounds $12.32572$ and $12.33854$, respectively. This proves the assertions.
\end{proof}

\subsection*{Removing the central-character assumption}
\begin{theorem}\label{rqstrip:central}
Let $K$ be any number field and $\pi$ a fixed unitary cuspidal representation of $\mathrm{GL}_2(\mathbb A_K)$ with arbitrary central character $\omega$. Assume the symmetric-square, cube and fourth-power lifts are cuspidal. Every nontrivial zero $\beta+iT$ of its standard function, for all sufficiently large $|T|$, satisfies
\[
 \min(\beta,1-\beta)>\frac1{12.335\,[K:\Q]\log|T|}.
\]
If the sixth symmetric-power lift also exists cuspidally with the expected good Euler factors, the coefficient $12.335$ can be replaced by $12.326$. The first assertion applies in particular to non-solvable-polyhedral cuspidal $\mathrm{GL}_2$ forms with nebentypus. It requires no temperedness assumption and no square root of the global character.
\end{theorem}
\begin{proof}
At a good place with residue norm $q$, let $a,b$ be the two Satake eigenvalues raised to the prime-power exponent $\nu$. Choose a unit scalar $z$ with $z^2=ab=\omega(\mathfrak p)^\nu$. Unitarity identifies the pair with its conjugate inverse; hence $x=(a+b)/(2z)$ is real. Explicitly, $\bar a+\bar b=(a+b)/(ab)$. Write $y=\Re(zq^{-i\nu T})$. For every coefficient of the amplifiers above, $\ell\equiv j\pmod2$. The following identity therefore uses an integral power of the genuine Hecke character:
\[
 U_j(x)T_\ell(y)=
 \Re\!\left[
  \operatorname{tr}(\operatorname{Sym}^j(\operatorname{diag}(a,b)))
   (ab)^{(\ell-j)/2} q^{-i\nu\ell T}\right].
\]
Changing $z$ to $-z$ reverses both $x$ and $y$ and leaves the expression invariant. The right side is the Euler trace at shift $\ell T$ of the automorphic function
\[
    L\big(s,\operatorname{Sym}^j\pi\otimes\omega^{(\ell-j)/2}\big).
\]
In particular the target at $(\ell,j)=(1,1)$ is the original standard function with coefficient one. At the static frequency every even positive index gives a cuspidal twist; the lone static pole has coefficient $P$. All shifted scalar factors are unitary Hecke characters $\omega^{\ell/2}$. Hu's Theorem~1.2 gives a fixed zero-free half-plane for each. The disk argument of Corollary~\ref{rqstrip:rankone} bounds their logarithmic derivatives by $o_\pi(\log|T|)$, uniformly for both real abscissae. A norm-character pole has fixed ordinate before the shift and is absent at sufficiently large $|T|$.

The quartic square is nonnegative for all real $x$; the sixth-degree polynomial is nonnegative there by its exterior verification. At bad places remove the finite Euler factors with $O_\pi(1)$ cost as in Corollary~\ref{rqstrip:generalfield}. Use the real part of the Hadamard identity at the \emph{single signed shift} $+\ell T$ of each cuspidal twist. The functional equation for any unitary standard completion pairs $\rho$ and $1-\bar\rho$ in its own zero divisor. Thus Stechkin's inequality for this pair works at each ordinate individually; at the target it retains $1/(\sigma-\beta)$ when $\beta>1/2$. This operation does not require a second zero at $\bar\rho$. The resulting positive derivative and gamma costs have the same weights as in the earlier proof, with shifted coefficient $(1-\kappa)rg$. The discriminants, conductors, static gamma factors and omitted Euler factors are fixed. Optimizing $u$ proves both claims for either sign of $T$.
\end{proof}

\subsection*{Standard functions in arbitrary rank}
The same local phase argument gives a constant proportional to the absolute degree for every essentially self-dual cuspidal standard function. Here only standard and Rankin--Selberg functions are required. We use the classical analytic continuation, functional equation and pole theorem of Jacquet--Piatetski-Shapiro--Shalika. The logarithmic zero-free region for standard functions has a classical de la Vall\'ee Poussin proof, recalled in \href{https://doi.org/10.1093/imrn/rnx087}{Goldfeld--Li, introduction}. We give the constants and identify the saving from the scalar character.

\begin{theorem}\label{rqstrip:arbitraryrank}
Let $r=[K:\Q]$ and $\pi$ be a fixed unitary cuspidal representation of $\mathrm{GL}_n(\mathbb A_K)$, $n\ge2$. If $\widetilde\pi\simeq\pi\otimes\xi^{-1}$ for a unitary Hecke character $\xi$, then at every sufficiently large $|T|$ its original standard function has zero-free width at either edge greater than
\[
                   \frac1{10.919\,nr\log|T|}.
\]
For an arbitrary unitary cuspidal $\pi$ the following width holds at sufficiently large heights:
\[
             \frac{1}{(11.020\,n+3.674\,n^2)r\log|T|}.
\]
In both statements the minimum height may depend on the representation. There is no local temperedness requirement.
\end{theorem}
\begin{proof}
At an unramified prime power let $a$ be the trace of the Satake parameter raised to its prime-power exponent $\nu$, and let $t=\nu T\log q$. In the first case $\bar a=\xi(\mathfrak p)^{-\nu}a$ by unitarity and the twist relation. Choose a unit scalar $z$ with $z^2=\xi(\mathfrak p)^\nu$ and put $X=a/z\in\R$, $Y=\Re(ze^{-it})$. For $b=\sqrt2$ use
\[
 \frac{(X+bY)^2}{2b}
  = \frac{|a|^2}{2b}+\Re(ae^{-it})
     +\frac b4+\frac b4\Re(\xi(\mathfrak p)^\nu e^{-2it}).
\]
The static traces are those of $\pi\times\widetilde\pi$ and $\zeta_K$. Each function has a simple pole at one. Their total weight is $P=1/(2b)+b/4=1/\sqrt2$. The shifted standard function occurs with coefficient one. The only other shifted function is the Hecke character $\xi$ at $2T$. Its finite logarithmic derivatives cost $o_\pi(\log|T|)$ by Hu's half-plane and the disk estimate above. This removes its entire degree contribution. The positive Euler sum persists after the Stechkin difference. The fixed finite set of omitted local factors has uniformly bounded logarithmic derivative on $\sigma\ge1$: for a local Rankin--Selberg factor the two exponents from unitary generic inducing data each have magnitude strictly less than $1/2$, so its inverse roots have modulus strictly less than $q$. Absolute convergence of the unramified sums at $\sigma>1$ follows from the usual nonnegative $\pi\times\widetilde\pi$ logarithmic derivative and Cauchy--Schwarz. The fixed static archimedean contributions are bounded near the two real abscissae. Only the standard gamma factors at $T$ contribute to the leading height cost. It is $(1-\kappa)nr\log|T|/2$. The same target-pair computation gives an asymptotic reciprocal width coefficient per absolute degree
\[
  \frac{1-1/\sqrt5}{2(1-2^{-1/4})^2}=10.9186118\ldots<10.919 .
\]

For arbitrary $\pi$ put $a'=ae^{-it}$ and $b=3/4$. The nonnegative square is now
\[
 \frac{(b+\Re a')^2}{2b}
  =\frac b2+\Re a'+\frac{|a|^2}{4b}
                   +\frac{\Re(a^2e^{-2it})}{4b}.
\]
The final term is the shifted Rankin--Selberg factor for $\pi\times\pi$. Any poles of this fixed function have fixed ordinates and a bounded regular pole contribution at $2T$. The static pole weight is $b/2+1/(4b)=17/24$, while the shifted half-degree per base-field degree is $n/2+n^2/(8b)=n/2+n^2/6$. The Stechkin comparison uses the zero pairs of each completed function individually, including those of the Rankin--Selberg function. Its fixed potential pole pair is charged separately. The same calculation gives
\[
 \frac{(1-1/\sqrt5)(n/2+n^2/6)}{(1-\sqrt{17/24})^2}
           <11.020\,n+3.674\,n^2.
\]
The strict numerical margins give the stated finite thresholds for each fixed form by the definition of the limiting inequality.
\end{proof}

\begin{proposition}\label{rqstrip:optimum}
In the essentially self-dual proof, the coefficient $10.9186118\ldots$ per absolute degree is optimal among positive quadratic trace tests with constant coefficient of $X^2$, any scalar trigonometric polynomial, and nonnegative weights on the shifted trace functions. In the corresponding unrestricted complex-trace test using only $|a|^2$, linear traces and scalar functions, the static pole weight is at least the target weight.
\end{proposition}
\begin{proof}
Normalize the target coefficient to one. In the real case write a general test as $A X^2+p(\phi)X+q(\phi)$, with the first cosine coefficient of $p$ equal to one. Positivity for all real $X$ forces $A>0$ and $q\ge p^2/(4A)$. With normalized Haar measure on the circle, Parseval gives $\int p^2\ge1/2$. The static pole coefficient is consequently
\[
                P=A+\int q\ge A+1/(8A)\ge1/\sqrt2 .
\]
The shifted standard factor by itself costs half the absolute degree; nonnegative extra weights can only increase this. The square in the theorem attains both bounds together. In the complex case write the test as $A|a|^2+\Re(a p(\phi))+q(\phi)$, with the target complex Fourier coefficient of $p$ equal to one. Minimization over all $a\in\mathbb C$ gives $q\ge|p|^2/(4A)$ and Parseval gives $\int|p|^2\ge1$. Thus $P\ge A+1/(4A)\ge1$. This last conclusion concerns the specified quadratic test; the shifted $a^2$ in the general-rank proof supplies exactly the extra function outside this class.
\end{proof}

\subsection*{Elliptic curves before modular descent}
\begin{theorem}\label{rqstrip:potential}
Let $E/K$ be any non-CM elliptic curve over a totally real or CM number field; no modularity of $E$ over $K$ is assumed. There is a finite integer $D$ such that the meromorphically continued unitary Hasse--Weil function has neither zeros nor poles in either of the two regions
\[
     0<\beta\le\frac1{12.326\,D\log|T|},
   \qquad
     1-\frac1{12.326\,D\log|T|}\le\beta<1
\]
for all sufficiently large $|T|$. Here $D$ can be taken to be the largest absolute degree of the intermediate fields in the Brauer decomposition below. For every fixed $m\ge1$ the analogous assertion for the $m$th symmetric power holds with $12.326 D$ replaced by $10.919(m+1)D_m$, with a potentially different extension and degree. For a CM elliptic curve, over any number field, the Hecke-character theorem gives instead the fixed unitary strip $[1/8,7/8]$ for its nontrivial zeros.
\end{theorem}
\begin{proof}
Let $V$ be the two-dimensional elliptic Tate representation, before shifting its weight. Over a totally real field choose a finite Galois totally real modularity extension $F/K$ as in the potential modularity theorem and its solvable-descent application; see \href{https://math.stanford.edu/~conrad/modseminar/pdf/L26.pdf}{Snowden, \emph{Potential modularity and applications}}, Theorem 15 and its applications, and \href{https://webusers.imj-prg.fr/~jan.nekovar/pu/parity4app.pdf}{Wintenberger, Theorem A.1}. For every $M$ intermediate with $F/M$ solvable, $E/M$ is modular by solvable descent. Its positive symmetric powers are cuspidal by the totally real symmetric-power theorem of Newton--Thorne.

For a CM base use \href{https://doi.org/10.4007/annals.2023.197.3.2}{Allen, Calegari, Caraiani, Gee, Helm, Le Hung, Newton, Scholze, Taylor and Thorne, Theorem 7.1.11 and Corollary 7.1.13}. Their simultaneous potential automorphy theorem gives a finite Galois CM extension $F/K$ on which the finitely many required systems $\operatorname{Sym}^jV$ are all cuspidal automorphic. Take $j=1,2,3,4,6$ for the first assertion, or just $j=m$ for the second. Serre's open-image theorem gives strong irreducibility, since each symmetric power of the standard representation of $\mathrm{SL}_2$ in characteristic zero is irreducible and every finite-index restriction has the same connected algebraic monodromy.

Proposition 6.5.13(2) of Allen et al. gives the required solvable descent. In a cyclic step $L/M$ the cuspidal lift over $L$ is invariant under the cyclic Galois group: its good Euler polynomials come from a system over $M$, and strong multiplicity one applies. Cyclic base change of Arthur--Clozel supplies a descent. Its infinity type is regular algebraic in the standard arithmetic normalization. The Galois representation of this descent over the CM field $M$ matches the prescribed system on $G_L$, by good Euler factors and Chebotarev. For the existence of the Galois representation see \href{https://doi.org/10.1186/s40687-016-0078-5}{Harris--Lan--Taylor--Thorne, \emph{On the rigid cohomology of certain Shimura varieties}} and the references on Galois representations in the paper of Allen et al. Since the restriction is absolutely irreducible, its intertwining space is one-dimensional. The two extensions differ by a character of $\operatorname{Gal}(L/M)$ on this space. Twisting the automorphic descent by this finite Hecke character, or its inverse as needed, realizes the actual system over $M$. Iterate along a cyclic composition series. The fields between $K$ and $F$ here are CM since they lie in a CM field and contain the CM base. This also explains why merely enlarging a modularity field by an arbitrary nonsolvable extension would be insufficient.

Choose Brauer induction for the trivial character of $G=\operatorname{Gal}(F/K)$ in the form
\[
       1_G=\sum_i n_i\operatorname{Ind}_{H_i}^G\chi_i,
\]
with integral $n_i$, solvable elementary $H_i$ and linear characters $\chi_i$. Write $M_i=F^{H_i}$. Tensor this identity with $\operatorname{Sym}^m V$ and use the projection formula for induction. Away from a fixed finite set of places, the original Euler product is the product of the descended $L_{M_i}(s,\operatorname{Sym}^m V\otimes\chi_i)^{n_i}$, normalized by the common half-weight shift. For $m=1$, symmetric lifts through the required indices of each twisted constituent are the descended lifts twisted by the corresponding powers of $\chi_i$; good Euler factors identify them. Apply Theorem~\ref{rqstrip:central} with its sixth-power coefficient. For a general $m$, the unitary symmetric power is self-dual on good classes: the normalized two-dimensional eigenvalues have product one. Its finite twist is essentially self-dual, with scalar character $\chi_i^2$, by strong multiplicity one. Apply Theorem~\ref{rqstrip:arbitraryrank}. Take the maximum of the finitely many constituent thresholds and degrees. A product with signed integer exponents is holomorphic and nonzero wherever all its factors are holomorphic and nonzero; the signs of the Brauer coefficients cause no loss in these regions.

The omitted local corrections are fixed rational functions of $q^{-s}$. Their divisors lie on finitely many fixed real parts. Those inside the open critical strip have positive distance from its edges, hence do not affect its shrinking edge regions for large enough $|T|$. The finitely many archimedean gamma parameters have poles only at fixed ordinates. This proves the claim with the actual finite Euler factors restored. The field degree appearing here belongs to the chosen potential automorphy and induction data; an explicit bound in terms of the original conductor alone is not asserted. For CM curves the fixed-strip assertion follows directly from Deuring's factorization into unitary Hecke functions after shifting the weight and from Hu's Theorem 1.2.
\end{proof}

For example the original prime-ideal sum of a non-CM elliptic curve in Theorem~\ref{rqstrip:potential} satisfies
\[
 \sum_{\substack{\mathrm N\mathfrak p\le x\\\mathfrak p\ {\rm good}}}
       \frac{a_{\mathfrak p}(E)}{\sqrt{\mathrm N\mathfrak p}}
             \log\mathrm N\mathfrak p
   \ll_{E,K} x\exp\!\left(-\sqrt{\frac{\log x}{12.326 D}}\right).
\]
To see this apply the truncated explicit formula of Corollary~\ref{exstrip:primes} to each of the finitely many cuspidal twists in the Brauer identity. At good places their coefficients come from the elliptic representation and a unitary finite twist, so after grouping by rational norm their absolute values at prime powers are at most $2[M_i:\Q]\Lambda(n)$. One can apply this formula to the full cuspidal functions and then remove the finitely many exceptional local terms. The pre-rounding reciprocal width is less than $12.32555D$, giving room to absorb all logarithmic factors. The logarithmic derivative of the Brauer product adds with the same integer weights; the finitely many unmatched local logarithmic derivatives cost $O_{E,K}(x^{1/2}\log x)$ by the strict unitary exponent bound on the automorphic inverse roots. Higher prime powers at the matching primes cost $O_K(\sqrt x\log^2 x)$. This yields the stated estimate under exactly the field hypotheses of the theorem.

\clearpage
\section{An explicit all-height region and the cost of amplification over fields}\label{zfnum:chapter}
Write $d_F=[F:\Q]$, $D_F=|\operatorname{disc}(F)|$, and $n_E=N_{F/\Q}\mathfrak f_E$. All zeros in this section use unitary coordinates. The rational theorem of Section~\ref{zf:chapter} gave an inverse-logarithmic region for every symmetric power. Here the standard elliptic function admits constants involving only the field discriminant, the norm of the elliptic conductor and the ordinate. A separate estimate covers CM over \emph{every} number field. The coefficients from a bad prime have a positive description that removes a conductor-sized error from the explicit formula.

\subsection{The archimedean and finite-prime bounds}
For $1<\sigma\le3/2$ the first of the following estimates holds for $0\le\mu\le3$ and the second for every $\mu\ge0$:
\begin{equation}\label{zfnum:gammabound}
 \frac{\Gamma_{\R}'}{\Gamma_{\R}}(\sigma+\mu)<0,\qquad
 \Re\frac{\Gamma_{\R}'}{\Gamma_{\R}}(\sigma+\mu+ix)
 \le\frac12\log\frac{\sqrt{(3/2+\mu)^2+x^2}}2 .
\end{equation}
For the first bound the real digamma argument is at most $9/4$. The digamma function $\psi$ is increasing on positive reals and $\psi(3)=3/2-\gamma<\log\pi$. For the second, with $z=(\sigma+\mu+ix)/2$, Euler summation for the digamma series gives
\[
 \left|\psi(z)-\log z+\frac1{2z}\right|
 \le \frac12\int_0^\infty |u+z|^{-2}\,du\le\frac1{2\Re z}.
\]
Indeed the first-order Euler summation remainder has integrand $B_1(\{u\})(u+z)^{-2}$ up to sign, and $|B_1(\{u\})|\le1/2$. As $\Re z>1/2$, the last bound is below $1<\log\pi$; also $\Re(1/z)>0$. Use $\Gamma_{\R}'/\Gamma_{\R}=(\psi(z)-\log\pi)/2$. This proves both estimates.

Let $V_j=\operatorname{Sym}^j V$ for the unitary elliptic local parameter. For any ramification group $H$ in the geometric two-dimensional representation and $j\ge1$,
\[
 \operatorname{codim}(\operatorname{Sym}^jV)^H\le j\operatorname{codim}V^H.
\]
If the fixed space has dimension one, the $j$th power of a fixed vector is fixed. If it is zero, use $j+1\le2j$; the trivial action gives zero on both sides. The conductor formula on inertia and the higher ramification groups consequently gives
\begin{equation}\label{zfnum:conductor}
    a_v(V_j)\le j a_v(V).
\end{equation}
At multiplicative reduction equality holds because $V_j$ is a single special block. These conductor statements use the geometric inertia action, including the unipotent action before passing to its Weil--Deligne operator.

\begin{lemma}[Positivity at every finite prime]\label{zfnum:gram}
Let $U$ be a finite sum of the $V_j$, each with a finite-order character and a purely imaginary norm twist, and give $U\otimes U^\vee$ its geometric local factors. Every coefficient in the logarithmic derivative with minus sign of its finite Euler product is nonnegative.
\end{lemma}
\begin{proof}
At potentially good reduction the local monodromy operator vanishes and inertia has finite image $I$. For a Frobenius $g$, the trace of $g^k$ on the invariant space of $U\otimes U^\vee$ equals
\[
     \frac1{|I|}\sum_{h\in I}|\operatorname{tr} U(g^kh)|^2.
\]
Indeed averaging gives the projection onto inertia invariants; all eigenvalues of each $U(g^kh)$ have modulus one. The latter follows by passing to a finite extension giving good reduction and taking a power of $g^kh$ in that subgroup, where inertia is trivial and Weil's bound applies. This argument includes good reduction.

At potentially multiplicative reduction the elliptic parameter is an unramified special block times a quadratic character. Put $x=\pm (Nv)^{-k/2}$, where the sign is the $k$th power of the unramified sign or of the value of the twisting character on the chosen Frobenius. In the decomposition $V_i\otimes V_j=\bigoplus_{r=0}^{\min(i,j)}V_{i+j-2r}$, the $k$th coefficient in the unramified-character case is the Gram entry
\[
 M_{ij}=\sum_{r\ge0,\ r\le i,j} x^{i-r}x^{j-r}.
\]
For unramified characters collect all phases and multiplicities at index $j$ into $w_j$. Summing gives
\[
           \sum_{r\ge0}\left|\sum_{j\ge r}w_j x^{j-r}\right|^2.
\]
For arbitrary finite twists first divide the individual summands into classes by their inertia character, including the $j$th power of the curve's possible quadratic character. Only pairs in the same class retain inertia invariants. Within each class collect the unitary Frobenius phases in $w_j$ and use the same Gram identity, with $|x|=(Nv)^{-k/2}$. The full coefficient is a sum of these nonnegative contributions. Multiplication by $\log Nv$ supplies the logarithmic coefficient. These two reduction types cover every finite place.
\end{proof}

\subsection{A conductor-sensitive region for modular elliptic curves}
\begin{theorem}\label{zfnum:elliptic}
Let $E/F$ be a non-CM elliptic curve with a cuspidal automorphic representation realizing its elliptic local parameters. This includes every non-CM elliptic curve over $\Q$. Every nontrivial zero $\beta+it$ satisfies both $\beta\ge \Delta_F(t)$ and $1-\beta\ge \Delta_F(t)$, where
\begin{equation}\label{zfnum:delta}
 \Delta_F(t)=\frac{(\sqrt{70}-\sqrt{50})^2}
 {50+\tfrac{529}2\log D_F+156\log n_E+130d_F\log(\sqrt{16+t^2}/2)}.
\end{equation}
The same proof with three constituents of multiplicity one gives the simpler numerator $(2-\sqrt3)^2$ and denominator $3+18\log D_F+11\log n_E+8d_F\log(\sqrt{16+t^2}/2)$.
\end{theorem}
\begin{proof}
Write $\pi$ for the given representation, and $\pi_2$ for its Gelbart--Jacquet symmetric square. The latter is cuspidal. A dihedral parameter would become a sum of characters after a quadratic extension, whereas the geometric elliptic representation has Zariski-dense image in $\mathrm{GL}_2$ on every finite-index subgroup by Serre's open image theorem. Equivalently, a quadratic self-twist of the automorphic form would give such a self-twist of its geometric representation at almost all primes, hence everywhere by Chebotarev; the dense connected image rules it out. Use the automorphic sum
\[
   U=a\mathbf1\boxplus b(\pi\otimes |\cdot|^{it})\boxplus c\pi_2,
\]
where a coefficient denotes that many identical summands. Its Rankin--Selberg product with its contragredient has real symmetry and nonnegative logarithmic coefficients by Lemma~\ref{zfnum:gram}. It is a product of standard Rankin--Selberg completions of its cuspidal pieces. The only possible poles are at $0,1$, with order at most $p=a^2+b^2+c^2$, since the three degrees are distinct. We use the classical symmetric-square lift and the standard full local Rankin--Selberg compatibility for pairs of general linear groups. See Gelbart and Jacquet, \href{https://doi.org/10.24033/asens.1355}{\emph{A relation between automorphic representations of $\mathrm{GL}(2)$ and $\mathrm{GL}(3)$}}, and Kim and Shahidi, \href{https://arxiv.org/abs/math/0409607}{\emph{Functorial products for $\mathrm{GL}_2\times\mathrm{GL}_3$ and the symmetric cube for $\mathrm{GL}_2$}}.

At $s=\beta$ the two orientations of the degree-one and degree-two pair give $2ab$ orders of a proposed zero of $L(\pi,\beta+it)$ and its conjugate. The two orientations of the degree-two and degree-three pair supply a further $2bc$ orders. Indeed $V_1\otimes V_2=V_1\oplus V_3$, so the corresponding Rankin--Selberg factor is the product of the elliptic standard factor and its symmetric cube. The symmetric cube is holomorphic on the open strip. Kim--Shahidi's lift is cuspidal in the non-polyhedral case, and non-CM elliptic monodromy excludes finite projective type. Even almost-everywhere identification of the symmetric cube suffices here: its geometric reciprocal roots on the monodromy kernel have modulus at most one as in Lemma~\ref{zf:cost}, so the finite corrections and its nonnegative gamma shifts are holomorphic on $\Re s>0$. This also deals with a proposed zero at $\beta=1$ after removing all $p$ possible poles. Thus the entire pole-cleared auxiliary completion has at least $q=2b(a+c)$ zeros at $\beta$, counting multiplicity. For $t=0$ the paired factors coincide and their exponents add.

The total relative degree is $D=(a+2b+3c)^2$. The conductor cost at $v$ is at most $C a_v(V)$ with
\[
    C=2b^2+6c^2+4ac+2ab+8bc.
\]
This follows by adding \eqref{zfnum:conductor} over all Clebsch--Gordan constituents. The discriminant factor in the absolute conductor is $D_F^D$. Unshifted real gamma factors have shift $\mu\le3$ and have negative logarithmic derivative by \eqref{zfnum:gammabound}. The shifted factors come from $V_1,V_3$ and have shift $\mu\le5/2$ at both real and complex places. The two orientations contribute $2Jd_F$ real gamma factors with $J=2b(a+3c)$. Consequently Lemma~\ref{zf:ef}, retaining its exact Hadamard identity, yields for $1<\sigma\le 3/2$
\begin{equation}\label{zfnum:optimized}
       \frac q{\sigma-\beta}\le\frac p{\sigma-1}+B,\quad
 B=p+\frac D2\log D_F+\frac C2\log n_E
          +Jd_F\log\frac{\sqrt{16+t^2}}2.
\end{equation}
The term $p/\sigma$ was bounded by $p$. For $p<q\le2p$ take $\sigma-1=(\sqrt{pq}-p)/B$. As $B\ge p$, this lies in $(0,\sqrt2-1]$. Solving the inequality gives $1-\beta\ge(\sqrt q-\sqrt p)^2/B$. With $(a,b,c)=(4,5,3)$, the tuple $(p,q,D,C,J)$ is $(50,70,529,312,130)$, proving~\eqref{zfnum:delta}. With $(1,1,1)$ it is $(3,4,36,22,8)$. The functional equation gives the other inequality.
\end{proof}

For a finite-order Hecke character $\chi$ over $F$, the first elliptic bound in Theorem~\ref{zfnum:elliptic} extends to $L(E/F,\chi,s)$ after adding $130\log N_{F/\Q}\mathfrak f_\chi$ to the denominator. Replace $\pi\otimes|\cdot|^{it}$ in $U$ by $\pi\otimes\chi|\cdot|^{it}$. The two orientations of the cross terms contain the target and its complex conjugate with the same total zero order. Lemma~\ref{zfnum:gram} treats all ramified twists. At the real places a sign twist of each odd angular $I_j$ is isomorphic to $I_j$, and a finite-order character is trivial at complex places. Each of the $2J$ shifted degrees costs at most one additional character conductor, giving the displayed $J=130$ term after halving. The reflected statement follows by applying the same inequality for $\bar\chi$.


\subsection{Every CM elliptic curve over every number field}
Let $C/F$ have CM and let $M$ be its field of definition of all endomorphisms over $F$; thus $[M:F]\le2$. In the notation of Deuring's theorem, take either of the unitary CM characters $\nu$ over $M$ contributing to the elliptic standard function. At the complex archimedean places of $M$ its angular weights have absolute value one. If $M\ne F$ the elliptic function equals a single Hecke function over $M$. If $M=F$ it is the product for the two conjugate characters, which have the same finite conductor. We write $f_j=N_{M/\Q}\mathfrak f_{\nu^j}$, so $f_{2j}\mid f_j\mid f_1$ for $j\ge1$.

\begin{theorem}\label{zfnum:cm}
For any $j\ge1$, any nontrivial zero $\beta+it$ of $L_M(s,\nu^j)$ satisfies
\[
 1-\beta\ge \frac{(\sqrt8-\sqrt6)^2}{B_j(t)},\qquad
 \beta\ge \frac{(\sqrt8-\sqrt6)^2}{B_j(t)},
\]
where, with $d=[M:\Q]$,
\begin{align*}
 B_j(t)={}&6+8\log D_M+4\log f_j+\log f_{2j}\\
 &+4d\log\frac{\sqrt{((j+5)/2)^2+t^2}}{2}
    +d\log\frac{\sqrt{(j+5/2)^2+4t^2}}{2}.
\end{align*}
For $j=1$ this bounds every nontrivial zero of $C/F$, without an exception at $t=0$ or a modularity hypothesis. In particular, for fixed $C/F$ and $j$, the lower limit of $(1-\beta)\log |t|$ as the absolute ordinate of such zeros tends to infinity is at least $(\sqrt8-\sqrt6)^2/(5d)$.
\end{theorem}
\begin{proof}
Put $\eta=\nu^j$. Use the completed constituents
\[
  \zeta_M(s)^6\, L_M(s+it,\eta)^4L_M(s-it,\bar\eta)^4\,
  L_M(s+2it,\eta^2)L_M(s-2it,\bar\eta^2).
\]
The notation displays the finite parts; each is completed with its own conductor and archimedean factors. The product is real on the real axis and is self-dual under $s\mapsto1-s$. At a prime where $\eta$ is unramified, if $z$ is its unitary Frobenius value adjusted by the shift, the $k$th coefficient divided by $\log Nv$ is $|1+2z^k+z^{2k}|^2$. At a prime where $\eta$ is ramified, the terms linear in $\eta$ disappear; whether $\eta^2$ is ramified or not, the coefficient is at least $6-2>0$. The nonzero angular weights of both $\eta$ and $\eta^2$ exclude a pole of either completed character function or any of its norm twists. Thus the pole budget is six. A proposed zero and its complex conjugate supply eight zeros at $s=\beta$ after clearing poles, even at $t=0$.

The aggregate absolute conductor is $D_M^{16} f_j^8 f_{2j}^2$. The six unshifted zeta factors contribute negative gamma derivatives. In the eight factors translated by $t$, the largest real gamma shift is $j/2+1$. In the two translated by $2t$ it is $j+1$. Applying~\eqref{zfnum:gammabound} gives exactly the two archimedean terms in $B_j(t)$. Now optimize~\eqref{zfnum:optimized} with $p=6$, $q=8$. The dual character has the same conductor and archimedean absolute weights, hence applying its bound and the functional equation gives the reflected inequality. The lower-limit statement follows by dividing $B_j(t)$ by $\log|t|$. The same formula holds with $\nu^j$ replaced by its product with any finite-order character of $M$, using the conductors of that product and its square. Its angular weights and the positivity calculation remain valid.
\end{proof}

For an odd symmetric power $m$ of a CM elliptic curve, the normalized weights occur in pairs $\pm j$, with $1\le j\le m$ odd. Over the endomorphism field these give the characters $\nu^j,\bar\nu^j$; otherwise each pair is their quadratic induction. Thus every zero of the full odd symmetric power, over every number field, satisfies Theorem~\ref{zfnum:cm} with the single upper denominator obtained by replacing $j$ by $m$ and $f_j,f_{2j}$ by $f_1$. This bound grows only logarithmically with the symmetric-power index for fixed curve, field and ordinate. For even $m$ the same holds for all non-middle weights, while the additional weight-zero character has finite order and must be retained separately.

\subsection{All zeros over nonnormal, nonsolvable base fields}
Let $E/\Q$ be non-CM of conductor $N$. Let $K/\Q$ have degree $h$, with permutation representation $P$ on its embeddings. Assume that each irreducible constituent $\rho$ of $P$ has a cuspidal automorphic representation $\tau_\rho$ over $\Q$ matching its Frobenius parameters at all but finitely many primes. Write $P=\bigoplus m_\rho\rho$. The examples in Section~\ref{mix:towers} include every cubic and quartic field and explicitly specified nonsolvable Galois and nonnormal fields. We now estimate the entire zero set near its right edge, including zeros contributed by nontrivial permutation constituents.

\begin{theorem}\label{zfnum:artin}
There is an absolute effective $c>0$ and a constant $A_K\ge3$ determined by the finite automorphic data of $P$ such that the following holds. For every $m\ge1$ and every primitive Dirichlet character $\chi$ of conductor $q$, the geometric function of $\operatorname{Sym}^m E_K$ twisted by $\chi\circ N_{K/\Q}$ continues holomorphically to $\Re s>0$ and is nonzero when
\[
 \Re s>1-\frac{c}{h(m+1)\log(A_K Nq(m+2)(3+|\Im s|))}.
\]
For $m=1$ and $q=1$ there is also the numerical asymptotic bound
\[
 \liminf_{\substack{\beta+it\text{ a nontrivial zero}\\ |t|\to\infty}}
       (1-\beta)\log |t|\ \ge\ \frac{(\sqrt{70}-\sqrt{50})^2}{130h}.
\]
These conclusions require no automorphy theorem for $E_K$ over $K$.
\end{theorem}
\begin{proof}
Let $\pi_j$ be the cuspidal unitary symmetric power over $\Q$ for $j\ge1$, supplied by Newton--Thorne, and put $\pi_0=1$. For a Dirichlet character $\chi$ form the completed product
\[
 B_{j,\chi}(s)=\prod_\rho
       \Lambda(\pi_j\times (\tau_\rho\otimes\chi),s)^{m_\rho}.
\]
The equality of the geometric base change and this automorphic product outside finitely many local factors follows from Artin formalism. The central character of $\tau_\rho\otimes\chi$ is finite. No factor for $j\ge1$ has a Rankin--Selberg pole. Indeed this would require $\tau_\rho\otimes\chi$ to be a norm twist of $\pi_j^\vee$, and the finite determinant forces that norm twist to be trivial. Take a good rational prime completely split in the fields cut out by $\rho$ and $\chi$, outside all exceptional places. All normalized reciprocal roots of $\operatorname{Sym}^j E$ would then be one. Two consecutive eigenvalues in the symmetric power have ratio $\alpha_p/\beta_p$, the ratio of the two elliptic reciprocal roots. Thus $\alpha_p=\beta_p$ and the integral trace of $E$ would be $\pm 2\sqrt p$, an impossibility. The only pole in $B_{0,1}$ comes from the unique trivial constituent. Its completed poles are simple at zero and one.

Here is the comparison at the exceptional places. A unitary generic local representation of a general linear group has Langlands parameters formed from tempered pieces with real exponents of absolute value $<1/2$. This standard local classification and its implications for local factors are recalled in the appendix of Z.~Rudnick and P.~Sarnak, \href{https://doi.org/10.1215/S0012-7094-96-08115-6}{\emph{Zeros of principal $L$-functions and random matrix theory}}, Duke Math.~J. 81 (1996). Each $\pi_j$ is tempered locally in the Weil--Deligne sense by its elliptic construction. Tensoring a bounded Weil--Deligne parameter in the tempered $W\times\mathrm{SL}_2$ convention with another tempered parameter gives a tempered parameter, by decomposing the products of the $\mathrm{SL}_2$ blocks. Thus every reciprocal root of a finite automorphic tensor factor here has modulus $<p^{1/2}$, and every archimedean real gamma shift has real part $>-1/2$. This statement permits ramified primes and complex infinity parameters of the automorphic lifts. The corresponding geometric roots have modulus $\le1$ on the full monodromy kernel before taking inertia invariants, as in Theorem~\ref{mix:elliptic-transfer} and Lemma~\ref{zf:cost}; tensoring with a finite representation preserves this fact. The finite correction is therefore holomorphic on $\Re s>0$ and has no zero or pole on $\Re s\ge1/2$. The archimedean correction has the same properties, since the geometric gamma shifts are nonnegative. This proves holomorphic continuation to the stated half-plane and equality of zero divisors with $B_{j,\chi}$ on $\Re s\ge1/2$ for $j\ge1$.

Apply the positive character pair $S,T$ of Lemma~\ref{amp:kernel}. Use weights $s_j$ on $B_{j,1}(s)$ and $b_j/2$ on $B_{j,\chi}(s+it)$ and its conjugate mate. As in the lemma, a common denominator turns this into an ordinary product before logarithmic differentiation. The sum of absolute degrees is $h(S(1)+T(1))\ll h(m+1)$; the pole budget is $s_0\le3/4$, and a proposed zero contributes at least $b_m=1$. Conjugate constituents of $P$ have equal multiplicity, with contragredient automorphic representations by strong multiplicity one, giving the required real functional symmetry.

Fix a finite set $\mathcal B$ containing the ramification of $P$, the ramification of the $\tau_\rho$, and all their failed Frobenius comparisons. At $p\notin\mathcal B$, $p\nmid Nq$, the local coefficient is the one from the proof of Theorem~\ref{amp:linear} multiplied by $\operatorname{tr}P(\operatorname{Frob}_p^v)$, the nonnegative number of fixed embeddings. For $p\mid Nq$, $p\notin\mathcal B$, the local tensor roots have modulus at most one since the finite representation is unramified and the elliptic parameters are tempered. Their total possible loss in the prime sum at $\sigma>1$ is $O(h(m+1)\log(2Nq))$. At the primes in $\mathcal B$ the corresponding loss is at most the aggregate degree times $\sum_{p\in\mathcal B}\log p/(\sqrt p-1)$, by the preceding local bound. This term is independent of $N,q,t$.

The tensor conductor inequality $a(V\otimes W)\le(\dim W)a(V)+(\dim V)a(W)$ follows on every higher ramification subgroup from the tensor product of fixed vectors. For the tame term use the tensor product of inertia-fixed monodromy kernels. Twisting by $\chi$ costs at most the tensor dimension times $a(\chi)$. Therefore the total logarithm of the conductor is
\[
 \ll h(m+1)\left(\log(2Nq)+h^{-1}\sum_\rho m_\rho\log Q(\tau_\rho)\right),
\]
using~\eqref{zfnum:conductor}. The gamma shifts of the unshifted $B_{j,\chi}$ have absolute value $O(j+1+H_K)$, with $H_K$ a fixed archimedean parameter bound for the finitely many $\tau_\rho$. Their real parts are $>-1/2$, giving a uniform digamma bound at $\sigma>1$ by the Euler summation estimate preceding~\eqref{zfnum:gammabound}. One may take $\log A_K$ equal to a sufficiently large absolute constant times
\[
 1+\log(3+H_K)+h^{-1}\sum_\rho m_\rho\log Q(\tau_\rho)
       +\sum_{p\in\mathcal B}\frac{\log p}{\sqrt p-1}.
\]
This describes the dependence on the chosen weak lifts and their exceptional places. It makes no assertion about computing the exceptional set from the defining polynomial of $K$ alone. The completed explicit formula with these estimates gives a remaining term $O(h(m+1)\log(A_K Nq(m+2)(3+|t|)))$. The optimization used in~\eqref{zfnum:optimized}, or the final step of Theorem~\ref{amp:linear}, proves the claimed region.

For the last assertion take the fixed weights $(4,5,3)$ of Theorem~\ref{zfnum:elliptic}, now replacing each rational symmetric-power factor by $B_{j,1}$. At almost all primes positivity is again multiplied by the number of fixed embeddings. All exceptional-prime contributions and all conductors are bounded independently of $t$. The unshifted gamma derivatives are bounded on $1<\sigma\le3/2$. Exactly $260h$ real gamma factors have translations $\pm it$. Stirling's formula gives their aggregate real logarithmic derivative $130h\log|t|+O_{E,K}(1)$, uniformly on this interval as $|t|\to\infty$. Formula~\eqref{zfnum:optimized} with $p=50$, $q=70$ gives the stated lower limit.
\end{proof}

\begin{corollary}[Sato--Tate over the permutation fields]\label{zfnum:st}
Fix a field $K$ of Theorem~\ref{zfnum:artin} and a non-CM elliptic curve $E/\Q$. For prime ideals $\mathfrak p$ of $K$ of good reduction write $\theta_{\mathfrak p}$ for the Frobenius angle. Uniformly in intervals $I\subset[0,\pi]$,
\[
 \left|\frac{\#\{\mathfrak p:N\mathfrak p\le x,\ \theta_{\mathfrak p}\in I\}}
 {\pi_K(x)}-\frac2\pi\int_I\sin^2\theta\,d\theta\right|
       \ll_{E,K}\frac{(\log\log(3x))^2}{\log(3x)} .
\]
\end{corollary}
\begin{proof}
Apply the proof of Theorem~\ref{amp:discrepancy} to the entire products $B_{j,1}$. We specify the changes. Their degree is $h(j+1)$, log conductors are $O_{E,K}(j+1)$ and the gamma shifts have modulus $O_K(j+1)$ and real part $>-1/2$. The local zero count at $2+it$ and the log-derivative estimate on the zero-free strip are therefore the same as in~\eqref{amp:localzeros}, with constants depending on $E,K$. All finite roots away from the fixed exceptional places have modulus at most one. At an exceptional rational prime $p$, the total change when passing between the automorphic and geometric logarithmic prime-power sums through $x$ is $O_{E,K}((j+1)\sqrt x\log(2x))$, by the bound $p^{v/2}$ for the $v$th powers of the reciprocal roots. On the original Perron line the exceptional-prime Euler derivatives are $O_{E,K}(j+1)$, uniformly right of $1$. For the taper over $(x,x+h_0]$ use the bounded geometric coefficients and charge the exceptional-prime errors through $2x$ as above. At every rational prime power the geometric coefficients have modulus at most $h(j+1)\log p$, by local induction and the finite permutation representation. The contributions of powers of prime ideals of exponent at least two are $O_{E,K}((j+1)\sqrt x\log^2(2x))$ from the norm bound and degree $h$. The first powers at bad primes contribute $O_{E,K}(j+1)$.

It follows with $L=\log(3x)$ that the good-prime sum of $U_j(\cos\theta_{\mathfrak p})\log N\mathfrak p$ is $O_{E,K}(xL^{-4})$ for $1\le j\le c_{E,K}L/(\log(3L))^2$. Apply the same interval majorants, minorants and partial summation as in Theorem~\ref{amp:discrepancy}, now subtracting $\mu(I)$ times the actual prime-ideal count. The prime ideal theorem gives $\pi_K(x)\asymp_K x/\log x$. This proves the estimate for the indicated fields even when their normal closure is nonsolvable. For the relation between Frobenius angles and permutation or Artin data, see P.-J.~Wong, \href{https://doi.org/10.1016/j.jnt.2018.09.010}{\emph{On the Chebotarev--Sato--Tate phenomenon}}.
\end{proof}

\subsection{Potential automorphy gives a zero- and pole-free region on descent}
\begin{theorem}\label{zfnum:potential}
Let $F$ be totally real or an imaginary CM number field and $E/F$ non-CM. Fix $m\ge1$ and a finite-order Hecke character $\eta$ of $F$. The geometric unitary $L_F(\operatorname{Sym}^m E\otimes\eta,s)$ admits meromorphic continuation and has neither a zero nor a pole in a region
\[
             \Re s>1-\frac{c_{E,F,m,\eta}}{\log(3+|\Im s|)}
\]
with $c_{E,F,m,\eta}>0$. More precisely, take $S,T$ from Lemma~\ref{amp:kernel}. There is a finite Galois extension $L/F$, depending on $E,F,m$, for which every zero or pole $\beta+it$ with $\beta\ge1/2$ satisfies
\begin{equation}\label{zfnum:potential-slope}
 \liminf_{|t|\to\infty}(1-\beta)\log|t|
       \ \ge\ \frac{2(1-\sqrt{s_0})^2}{[L:\Q]T(1)}.
\end{equation}
Here the lower limit runs over the indicated divisor points. There is also $A_{E,L,m}\ge3$, independent of $\eta$, such that every such divisor point satisfies
\[
 1-\beta\ge\frac{2(1-\sqrt{s_0})^2}{[L:\Q]T(1)
 [\log(A_{E,L,m}(3+|t|))+[F:\Q]^{-1}\log N_{F/\Q}\mathfrak f_\eta]}.
\]
The proof does not bound the degree or discriminant of the potential modularity field in terms of $E,F,m$.
\end{theorem}
\begin{proof}
Use the simultaneous potential automorphy theorem of P.~Allen et al., \href{https://annals.math.princeton.edu/2023/197-3/p02}{\emph{Potential automorphy over CM fields}}, Annals of Mathematics 197 (2023), 897--1113, Theorems~1.0.1 and~7.1.11. Their convention includes totally real fields. For the finite list of all positive $j$ occurring in $S$ or $T$, together with $m$, choose a common Galois field $L/F$, totally real or CM, over which the corresponding symmetric powers are automorphic. The compatible systems are strongly irreducible by the elliptic open image theorem and the symmetric powers of the standard representation of $\mathrm{SL}_2$. The common field is obtained by taking the finite index set in Theorem~7.1.11 and its sufficient field. The soluble descent of Proposition~6.5.13(2) of the same paper then gives a cuspidal realization of each geometric symmetric power over $L^H$ whenever $H\le\operatorname{Gal}(L/F)$ is soluble. Almost-everywhere local matching suffices. The fields containing $F$ inside $L$ are again totally real or CM, so the descent proposition applies.

Brauer induction gives an identity of characters $1_G=\sum_i n_i\operatorname{Ind}_{H_i}^G\xi_i$, where $G=\operatorname{Gal}(L/F)$, the $n_i$ are integers, the $H_i$ are elementary, hence soluble, and each $\xi_i$ is one-dimensional. Put $F_i=L^{H_i}$. Local Artin formalism expresses the geometric function as the product, to powers $n_i$, of the geometric standard functions of $\operatorname{Sym}^m E_{F_i}$ twisted by $\eta_i=(\eta\circ N_{F_i/F})\xi_i$. Class field theory supplies the finite-order Hecke characters represented by $\xi_i$. For each $i$ the symmetric powers in $S,T$ have unitary cuspidal realizations $\pi_{j,i}$ over $F_i$. The geometric and automorphic local factors have the same divisor on $\Re s\ge1/2$ after their finite replacements: the geometric reciprocal roots have modulus at most one and generic unitary standard roots have modulus $<(Nv)^{1/2}$. The automorphic archimedean shifts have real part $>-1/2$, and geometric shifts are nonnegative. The corrected products furnish the stated meromorphic continuation.

Over $F_i$ use $S$ on the unshifted factors $\pi_{j,i}$, and $T/2$ on the translates of $\pi_{j,i}\otimes\eta_i$ and their conjugate mates. The untwisted realizations are self-dual by strong multiplicity one and the self-dual unitary geometric parameters. At every good unramified place the nonnegative coefficient is $S(\cos(v\theta))+T(\cos(v\theta))\Re z^v$ with $|z|=1$. Fix the exceptional places of the geometric comparison, all bad places of the finitely many automorphic forms, and the conductors of the $\xi_i$. At a new ramified place of $\eta_i$ outside this set the twisted unramified local terms disappear, leaving $S(\cos(v\theta))\ge0$. The fixed exceptional set costs $O_{E,L,m}(1)$ uniformly in $t,\eta$ by the local $<(Nv)^{1/2}$ bound. The logarithm of the additional conductor of the shifted factors contributes at most $\tfrac12 T(1)[F_i:F]\log N_{F/\Q}\mathfrak f_\eta$ to the explicit formula, up to a fixed constant. Indeed norm pullback has conductor dividing the extended conductor ideal: local elements congruent to one modulo $\mathfrak f_\eta\mathcal O_{F_i}$ have norm congruent to one modulo $\mathfrak f_\eta$. The finite-order sign choices at real places form a finite set. The only poles in this auxiliary product are contributed by the unshifted $s_0$ copies of the Dedekind zeta function. There are at least $b_m=1$ zeros at $\beta$ if the target factor has a zero at $\beta+it$. When a finite-order twist is not real, the conjugate mate supplies the other half of the weight. The two sorts of factors have total shifted archimedean degree $[F_i:\Q]T(1)$. Stirling's formula away from its poles, including bounded ordinates by continuity, and the argument of~\eqref{zfnum:optimized} consequently give
\begin{align*}
 \frac1{\sigma-\beta}&\le\frac{s_0}{\sigma-1}+\tfrac12[F_i:\Q]T(1)\mathcal H+O_{E,L,m}(1),\\
 \mathcal H&=\log(3+|t|)+[F:\Q]^{-1}\log N_{F/\Q}\mathfrak f_\eta,
                  \qquad 1<\sigma\le3/2.
\end{align*}
The constant can be increased to make the optimizer lie in this interval for all ordinates. Optimization gives the inverse-logarithmic region and leading coefficient $2(1-\sqrt{s_0})^2/([F_i:\Q]T(1))$. If $s_0=0$, the same conclusion follows by letting $\sigma$ decrease to one; in the stated construction $s_0>0$.

Intersect the finitely many regions, using $[F_i:\Q]\le[L:\Q]$. Every factor raised to a positive or negative power in the Brauer identity is holomorphic and nonzero there. The same is consequently true of the descended product, including its corrected geometric factors. This proves both the zero and pole assertions. This use of Brauer induction requires simultaneous nonvanishing; a zero of one positive factor elsewhere can be cancelled by the denominator.
\end{proof}

\begin{corollary}[A square-root-exponential prime error for a fixed symmetric power]\label{zfnum:prime}
Under the hypotheses of Theorem~\ref{zfnum:potential}, fix $E,F,m,\eta$. At good primes let $\theta_{\mathfrak p}$ denote the normalized Frobenius angle. There is $b>0$, depending on these fixed data, such that
\[
 \sum_{N\mathfrak p\le x}^{\mathrm{good,\ unram.}}
  U_m(\cos\theta_{\mathfrak p})\eta(\mathfrak p)\log N\mathfrak p
     =O_{E,F,m,\eta}\bigl(x e^{-b\sqrt{\log x}}\bigr).
\]
\end{corollary}
\begin{proof}
In the Brauer product each cuspidal factor gives a logarithmic derivative bounded by $O(\log^2(3+T))$ on $1-a/\log(3+T)\le\sigma\le2$, $|t|\le T$, for some fixed $a>0$. This follows by applying the local-zero-count and Hadamard-difference proof of~\eqref{amp:localzeros} to a factor and its conjugate mate; the product bounds the nonnegative zero count for either factor. All degrees and fields in this step are fixed. The sum of the signed logarithmic derivatives satisfies the same bound with a new constant. Finite Euler corrections at exceptional primes cost $O(\sqrt x\log x)$ in coefficient sums by the generic local bound. Apply the tapered Perron kernel in Theorem~\ref{amp:discrepancy} with $T=\exp(\sqrt{\log x})$ and $h=x/\sqrt T$. Removing the taper costs $O(h\log(2x))$ for the bounded geometric coefficients, plus the exceptional-prime allowance. The infinite tails on the original line are $O(x^2\log^2(2x)/(hT))$. On the new vertical line the bound is $O(xe^{-a\log x/\log(3+T)}\log^3(3+T))$; the two horizontal edges are $O(x^2\log^2(3+T)/(hT^2))$. Each is $O(xe^{-b\sqrt{\log x}})$ for a sufficiently small positive $b$. Grouping the prime ideals by rational prime power gives $O_{F,m}(\log p)$ coefficients at each rational prime power $p^v$ for the geometric product. Higher powers of prime ideals contribute $O_{F,m}(\sqrt{x}\log^2(2x))$. This proves the asserted first-prime sum.
\end{proof}

\begin{corollary}[Two elliptic curves under potential automorphy]\label{zfnum:pair}
Let $E_1,E_2$ be geometrically nonisogenous non-CM elliptic curves over a totally real or imaginary CM number field $F$. For $m,n\ge1$ fixed and any fixed finite-order Hecke character $\eta$, the unitary geometric tensor
\[
    L_F(\operatorname{Sym}^m E_1\otimes\operatorname{Sym}^n E_2\otimes\eta,s)
\]
is meromorphic and free of zeros and poles in $\Re s>1-c/\log(3+|\Im s|)$ for some $c>0$. With a simultaneous potential modularity field $L$ for the relevant two finite lists, the limit inferior of $(1-\beta)\log|t|$ at right-half-strip divisor points is at least
\[
 \frac{2(1-\sqrt{s_{0,m}s_{0,n}})^2}{[L:\Q]T_m(1)T_n(1)}.
\]
Here $S_j,T_j$ denote the two polynomials of Lemma~\ref{amp:kernel} for the index $j$. In particular the mixed tensor of Section~\ref{mix:chapter} has an all-height right edge after every totally real or CM base change.
\end{corollary}
\begin{proof}
Use simultaneous potential automorphy for these finite lists for both curves and descend by the same elementary-group Brauer induction. On each $F_i$ use the products $S_m(x)S_n(y)$ and $T_m(x)T_n(y)$ from Theorem~\ref{amp:pair}. Their good-prime domination, target coefficient one, shifted relative degree $T_m(1)T_n(1)$, and unshifted trivial coefficient $s_{0,m}s_{0,n}$ follow by multiplication. Each constituent is now a standard Rankin--Selberg product of the separately automorphic symmetric powers. The only poles are those supplied by the two degree-one trivial factors: the possible other equal-rank contragredient case contradicts the independence of the two geometric systems, even after a finite-order twist or restriction to a finite field extension. Indeed the connected derived monodromy projects onto $\mathrm{SL}_2$ in each coordinate. Failure to contain their product would, by Goursat's lemma for the simple adjoint groups, identify the two projective representations on a finite-index subgroup; all algebraic automorphisms of $\mathrm{PGL}_2$ are inner. The two standard representations would then differ there by a character, which has order at most two since their determinants agree. Faltings's isogeny theorem would make the curves geometrically isogenous. The full product of the two derived groups excludes the indicated equivalence of symmetric powers. Any unitary norm twist in an alleged equivalence is trivial by the finite central characters.

At finitely many exceptional places each cuspidal standard parameter is a sum of tempered pieces with real exponents strictly between $-1/2$ and $1/2$. The tensor therefore has all finite reciprocal roots of modulus strictly less than $Nv$, uniformly away from $Nv$ for this finite list, and archimedean shifts with real part strictly greater than $-1$. Thus the local corrections between the automorphic and geometric factors are nonzero and holomorphic on a fixed strip $\Re s>1-\epsilon$ for some $\epsilon>0$, and their costs in the exterior prime sum at $\sigma>1$ are bounded uniformly in the ordinate. The geometric parameters on the monodromy kernels are tempered, as follows by tensoring their tempered $W\times\mathrm{SL}_2$ parameters. Retain the smaller region if necessary at bounded ordinates. At large height the optimization from Theorem~\ref{zfnum:potential}, with the newly stated pole weight and gamma degree, gives the displayed slope. On descent every numerator and denominator is zero-free there, proving the assertion.
\end{proof}

\subsection{Exactly how often a fixed elliptic factor recurs}
An abelian extension supplies a useful converse. Suppose all finite-order twists of the elliptic function of $E/F$ are entire; this holds for modular cuspidal and CM curves. For every finite abelian extension $K/F$ Artin formalism gives
\[
          \Lambda(E_K,s)=\prod_{\xi\in\widehat{\operatorname{Gal}(K/F)}}
                             \Lambda(E_F\otimes\xi,s).
\]
Hence the Riemann hypothesis for the elliptic functions over \emph{all cyclic extensions} of $F$, including $F$ itself, is equivalent to the Riemann hypothesis for all finite-order Hecke twists of $E/F$. One implication uses that the extension cut out by a single finite-order character is cyclic and the other factors are entire. The reverse implication holds for every finite abelian extension by multiplying the twists. For a CM curve over its endomorphism field these are angular Hecke twists over $F$, as distinct from finite-order $L$-functions alone.

The amount by which restriction and tensor power force the same representation to recur can be computed without estimates. Write $V=V_\ell(E)$, with its convention chosen so its determinant is the cyclotomic character, and let $r=2k+1$. In unitary coordinates the geometric system of the same weight as $V$ is $W=V^{\otimes r}\otimes(\det V)^{-k}$. The normalization is an integral twist. For a finite extension $K/F$ form $R_{k,K}=\operatorname{Ind}_{G_K}^{G_F}(W|_{G_K})$.

\begin{proposition}\label{zfnum:multiplicity}
If $E/F$ is non-CM, the multiplicity of $V$ in $R_{k,K}$ is
\[
 u_{k,K}=\binom{2k+1}{k}-\binom{2k+1}{k-1}
        =\frac1{k+2}\binom{2k+2}{k+1}
\]
for every finite $K/F$, with the out-of-range binomial taken as zero. If $E/F$ is CM and its two endomorphism characters are interchanged over $F$, write $F'$ for the quadratic endomorphism field. Then
\[
 u_{k,K}=\begin{cases}
 \binom{2k+1}{k},& F'\not\subset K,\\
 2\binom{2k+1}{k},& F'\subseteq K.
 \end{cases}
\]
If $F$ already contains the endomorphism field, the multiplicity of either individual CM character in $R_{k,K}$ is $\binom{2k+1}{k}$. In particular the forced multiplicity of the untensored elliptic constituent under arbitrarily large base change is at most two, and it is one in the non-CM case.
\end{proposition}
\begin{proof}
Frobenius reciprocity gives $u_{k,K}=\dim\operatorname{Hom}_{G_K}(V,W)$ when $V$ is irreducible over $F$. In the non-CM case the Zariski closure of $G_K$ on $V$ is $\mathrm{GL}_2$, so the computation is unchanged on passing to any finite extension. On $\mathrm{SL}_2$ the weights $1$ and $3$ of $V^{\otimes(2k+1)}$ have multiplicities $\binom{2k+1}{k}$ and $\binom{2k+1}{k-1}$. Every odd irreducible of highest weight at least three contributes once to both weights, while the defining representation contributes only to weight one. Subtracting gives the first formula. The determinant twist makes the two central characters agree, so no multiplicity is lost in passing back to $\mathrm{GL}_2$.

Over $KF'$ write $V=\psi\oplus\psi'$. The weights in $W$ are $\psi^a(\psi')^{1-a}$ for integral $a$; they are mutually distinct because $\psi/\psi'$ has nonzero angular weight and still has infinite order after every finite extension. The weights $\psi$ and $\psi'$ each occur $\binom{2k+1}{k}$ times. When $F'\subseteq K$ the two Hom spaces contribute independently. Otherwise the nonidentity coset interchanges them, and equivariance identifies either choice with a unique choice for its conjugate. Their invariant space consequently has the dimension of one summand. When $V$ already splits over $F$, apply reciprocity separately to either character. This proves all cases.
\end{proof}

There is an exact conductor check in the non-CM semistable rational case. At any bad prime the special block in $V$ has one unipotent Jordan block. The number of irreducible $\mathrm{SL}_2$ summands in $V^{\otimes r}$, with multiplicity, is $\binom{r}{k}$, since each odd summand contains exactly one weight-one line. Each contributes one line to the kernel of the prime's monodromy. There is no wild ramification and the determinant twist is unramified there. Thus the finite conductor of $W$ over $\Q$ is precisely
\[
                    N^{2^r-\binom{r}{k}}.
\]
Its elliptic factor occurs only $u_{k,\Q}$ times. Here all complementary odd symmetric powers are entire by Newton--Thorne, so every elliptic zero actually produces at least that order in the tensor-power completion. In the CM case all odd angular factors are entire, and a CM endomorphism extension that first splits the elliptic parameter doubles the base zero orders because the two conjugate character $L$-functions over that extension agree after permuting its places. Applying Corollary~\ref{mix:cm-all} from the endomorphism field onward therefore gives the doubled analytic lower bound for towers whose remaining steps meet that corollary. When the parameter remains nonsplit, the corresponding single binomial lower bound follows by applying it from $F$. In a general setting the representation multiplicity just computed must be distinguished from possible accidental coincidences of analytic zeros of different constituents.

For a precise test of a density argument, take an analytic conductor with an archimedean factor $3+|t+\mu|$ for each real gamma factor, or any majorant thereof. The degree of $R_{k,K}$ over $\Q$ is $D=2^r[K:\Q]$, so $\log\mathcal Q\ge D\log3$. In the non-CM case Stirling's formula gives
\[
 \frac{u_{k,K}}{2^{2k+1}}\sim \frac{2}{\sqrt\pi\,k^{3/2}}.
\]
In the nonsplit CM case the corresponding quotient is asymptotic to $1/\sqrt{\pi k}$ or twice that value according as the parameter stays nonsplit or splits over $K$. Increasing field degree never increases $u_{k,K}$ beyond these two options. A density inequality of the usual shape $N(\sigma,T)\le C\mathcal Q(T)^{A(1-\sigma)}$, with $A,C>0$, when compared only to a forced multiplicity $u$ at a fixed zero, gives at most $1-\beta\ge\log(u/C)/(A\log\mathcal Q(T))$. Even with $C$ fixed and every forced copy surviving analytically, its width is $O(r/(2^r[K:\Q]))$. A linear local upper bound in $\log\mathcal Q$ similarly dominates the forced multiplicity by an increasing factor. Thus these base changes and tensor powers do not amplify a fixed off-line zero fast enough for either comparison to reach the central line. This calculation concerns these positive representation constructions and these forms of the density estimates.

\subsection{A stronger height region for CM tensors and Frobenius sectors}
The explicit numerical constants in Theorem~\ref{zfnum:cm} have a complementary fixed-data result. The relevant established input is M.~D.~Coleman, \href{https://doi.org/10.1112/S0025579300013000}{\emph{A zero-free region for the Hecke $L$-functions}}, Mathematika 37 (1990), Theorem~2. In particular a fixed Hecke character over any number field has a zero-free edge of Vinogradov--Korobov strength. See also the precise fixed-character statement in I.~Petrow, \href{https://doi.org/10.4171/JEMS/1465}{\emph{The Weyl law for algebraic tori}}, Lemma~2.12. For rational CM elliptic curves the simultaneous angular version is already stated explicitly by A.~Panidapu and J.~Thorner, \href{https://arxiv.org/abs/2105.11093}{\emph{Short-interval sector problems for CM elliptic curves}}, Lemma~3.3. We use positive and negative Brauer powers to include arbitrary finite-image twists over every ground field.

\begin{theorem}\label{zfnum:cm-artin}
Let $E/F$ be any CM elliptic curve and $\rho$ any fixed complex finite-image Galois representation of $F$. For every odd integer $m\ge1$, the geometric completion of $\operatorname{Sym}^m E\otimes\rho$ is meromorphic on the plane. For some $c=c(E,F,\rho)>0$, simultaneously for all such $m$, it has neither a zero nor a pole in either edge of width
\[
 \frac{c}{(\log(3+m+|t|))^{2/3}(\log\log(16+m+|t|))^{1/3}}.
\]
In particular, for fixed $m$ the width can be written
\[
 \frac{c(E,F,\rho,m)}{(\log(|t|+3))^{2/3}(\log\log(|t|+16))^{1/3}}.
\]
The finite Galois group of $\rho$ is unrestricted. In particular the latter fixed-data estimate holds for every CM elliptic curve over every number field.
\end{theorem}
\begin{proof}
Over the endomorphism field $M$, the nonzero weights of the odd symmetric power are $\nu^{\pm j}$, for odd $1\le j\le m$. If $M\ne F$, pair them before descending by quadratic induction. The projection formula for induction puts the tensor with $\rho$ over $M$ as the corresponding angular character tensored with $\rho|_{G_M}$. Apply Brauer induction to this last finite representation. The completed function becomes a product to signed integer powers of ordinary Hecke completions over finitely many fixed intermediate fields $M_i$, each of the form
\[
 L_{M_i}\bigl(s,(\nu\circ N_{M_i/M})^{\pm j}\xi_i\bigr)
\]
with $\xi_i$ of finite order. The fields and these finite-order characters depend on $E,F,\rho$ and do not vary with $m$. The pulled-back angular character has nonzero angular weight at each complex place. Consequently none of these factors is a norm character or has the real quadratic exceptional zero allowed in the general Hecke theorem. For fixed $m$ the minimum of the finitely many constants in Coleman's theorem proves the stronger displayed height region.

For the simultaneous assertion we use Coleman's estimate with varying angular index and ordinate as reproduced in T.~Khale, C.~O'Kuhn, A.~Panidapu, A.~Sun and S.~Zhang, \href{https://doi.org/10.1016/j.jnt.2021.04.004}{\emph{A Bombieri--Vinogradov theorem for primes in short intervals and small sectors}}, Theorem~2.5. For a fixed field the denominator is the maximum of the logarithm of the finite conductor and $(\log V)^{2/3}(\log\log V)^{1/3}$ at large $V$, where $V$ bounds the infinity-character lattice index and the absolute ordinate. The possible exceptional zero requires zero lattice index and a real finite-order character. The finite conductor of $(\nu\circ N)^j\xi_i$ divides a common modulus. Relative to the fixed lattice of unitary Hecke characters modulo finite torsion and pure norm twists, its index vector is $j$ times a fixed nonzero vector. Any extracted imaginary norm exponent also grows at most linearly with $j$. Thus take $V\asymp 3+m+|t|$ for all the factors, using the estimate for each member at its actual index and ordinate; for bounded parameters the finite lists are nonvanishing at $\Re s=1$. The possible exceptional character never occurs. The finitely many factors in each particular product are simultaneously holomorphic and nonzero in the common region, regardless of their signed exponents. Their Hecke functional equations identify the opposite edge with the right edge for the dual finite representation. Taking the smaller of the two fixed constants proves the assertion on both sides.
\end{proof}

\begin{theorem}[CM sectors with arbitrary finite Frobenius conditions]\label{zfnum:cm-sectors}
Let $P/F$ be finite Galois, $E/F$ a CM elliptic curve, and $M$ its endomorphism field. Fix a union $C$ of conjugacy classes in $\operatorname{Gal}(P/F)$. In $G=\operatorname{Gal}(PM/F)$ let $C'$ be its inverse image and $H=\operatorname{Gal}(PM/M)$. Put $\delta_+=|C'\cap H|/|G|$ and $\delta_-=|C'\setminus H|/|G|$. For some $b>0$, uniformly in intervals $I\subset[0,\pi]$,
\begin{align*}
 &\sum_{\substack{N\mathfrak p\le x,\ \operatorname{Frob}_{\mathfrak p}\in C\\
              \theta_{\mathfrak p}\in I}}^{\mathrm{good,\ unram.}}
                   \log N\mathfrak p\\
 &\qquad=x\left(\delta_+\frac{|I|}{\pi}+
             \delta_-\mathbf 1_{\pi/2\in I}\right)
       + O_{E,F,P}\!\left(x\exp\{-b(\log x)^{3/5}(\log\log(3x))^{-1/5}\}\right).
\end{align*}
The restriction to good unramified primes omits the finite bad and ramified set of the combined extension. If $M=F$, then $H=G$ and only the continuous term remains.
\end{theorem}
\begin{proof}
First consider primes of $M$ with Frobenius in $D=C'\cap H$. Its indicator is a linear combination of finitely many ordinary characters of $H$. For each component, Brauer induction expresses its twist by $\nu^j$ as a signed product of Hecke functions on fixed fields, as above. The logarithmic derivative at a prime ideal of $M$ recovers the trace times $\nu(\mathfrak p)^j$ at the first power. Put $X=\log x$, $Y=X^{3/5}(\log(3X))^{-1/5}$ and $T=\exp(Y)$. Uniformly in $1\le j\le T$, each Hecke factor is zero-free for $|t|\le T+1$, $\sigma\ge1-cY/X$ for a fixed positive $c$. On the half-width strip its finite logarithmic derivative is $O((\log(3T))^2)$. To see this, the local count of zeros in a unit interval at height $t$ is $O(\log(3T))$ by the Hadamard estimate at $2+it$, since the degrees are fixed, finite conductors are bounded and gamma parameters are $O(j+1)$. For a non-self-dual character apply the nonnegative local zero sum to the product with its conjugate; the bound holds for either member. Subtracting logarithmic derivatives at $2+it$ and $\sigma+it$ as in \eqref{amp:localzeros} now costs $O(\log(3T)X/Y)$. The Euler-factor coefficients are uniformly bounded at every finite prime, including the original ramified set.

Use the tapered Perron argument of Corollary~\ref{zfnum:prime} with $h=x/\sqrt T$. The new line contributes $O(xe^{-c'Y}\log^3(3T))$, the original tails contribute $O(x\log^2(2x)/\sqrt T)$, and the taper and horizontal edges are bounded in the same manner. Removing higher prime powers costs $O(\sqrt x\log^2(2x))$. Decreasing $b$ gives
\[
   \sum_{N\mathfrak P\le x}^{\operatorname{Frob}\in D}
          \nu(\mathfrak P)^j\log N\mathfrak P
           \ll_{E,F,P} xe^{-bY}\qquad(1\le j\le T).
\]
Here and below finitely many excluded primes do not affect the error. At $j=0$ the same fixed-data argument applies to the finite Hecke factors, with their zeta poles retained. A possible fixed exceptional real zero is bounded away from one on this finite list and can be accommodated by decreasing $c$. The pole coefficient for a character $\rho$ is $\langle\rho,1_H\rangle$ by Brauer induction and Frobenius reciprocity. The total main coefficient for the indicator of $D$ is therefore $|D|/|H|$. This proves the corresponding fixed-extension Chebotarev estimate without any Artin holomorphy assumption.

The inverse image of the angle interval under $z\mapsto\arccos(\Re z)$ in the unit circle is a union of at most two arcs of total normalized length $|I|/\pi$. Its degree-$\lfloor T\rfloor$ Beurling--Selberg majorant and minorant have constant terms within $O(1/T)$ of that length and $j$th coefficients $O(1/j)$ for $1\le j\le T$. Applying the last two estimates gives a main term $x|D||I|/(|H|\pi)$ with error $O(xe^{-b'Y})$; the factor $\sum_{j\le T}1/j$ is absorbed by decreasing $b'$. Primes of $M$ of relative residue degree two over $F$ contribute $O_{F}(\sqrt{x}\log x)$ throughout. Each remaining split prime of $F$ was counted $[M:F]$ times with the identical folded angle. Dividing by this degree gives $\delta_+$. Primes whose Frobenius is in $C'\setminus H$ have elliptic trace zero by quadratic induction. Their count has coefficient $\delta_-$ by the $j=0$ estimate for $PM/F$, and they occur exactly when $\pi/2\in I$. This proves the formula, including its endpoint convention for the atom.
\end{proof}

\subsection{Short intervals, shrinking sectors and the class field degrees}
For a conjugacy class $g^J$ of a finite group $J$, put
\[
 b_J(g)=\max\{|B|:g\in B\le J,\ B\text{ abelian}\}.
\]
Write $W=PM$. For the split classes choose representatives for the $H$-conjugacy classes in $C'\cap H$; for the nonsplit classes choose representatives for the $G$-classes in $C'\setminus H$. For this paragraph take $C'$ nonempty. Define $d_C$ to be the largest of
\[
 \frac{[W:\Q]}{b_H(g)}\quad(g\text{ split}),\qquad
 \frac{[W:\Q]}{b_G(g)}\quad(g\text{ nonsplit}).
\]
Thus $d_C\le[W:\Q]$. If $P/F$ is abelian, the split classes require only Hecke functions over $M$, and the nonsplit classes require only Hecke functions over $F$.

\begin{theorem}\label{zfnum:cm-short}
In the notation of Theorem~\ref{zfnum:cm-sectors}, put
$\mu_C(I)=\delta_+|I|/\pi+\delta_-\mathbf1_{\pi/2\in I}$.
Fix $\delta,\theta\ge0$ with
\[
                   \max(\delta,\theta)<\frac2{5d_C}.
\]
Uniformly for $x^{1-\delta}\le h\le x$ and subintervals $I$ of $[0,\pi]$ of length at least $x^{-\theta}$,
\begin{align*}
 \sum_{\substack{x<N\mathfrak p\le x+h,\ \operatorname{Frob}_{\mathfrak p}\in C\\
                          \theta_{\mathfrak p}\in I}}^{\mathrm{good,\ unram.}}
       \log N\mathfrak p
 &=h\mu_C(I)\biggl(1+O\!\left(\exp\left\{
   -b\left(\frac{\log x}{\log\log(3x)}\right)^{1/3}\right\}\right)\biggr).
\end{align*}
Here $b>0$ and the implicit constant may depend on $E,F,P,\delta,\theta$. When $F=\Q$ and $P=F$, the same assertion holds for
\[
                      \max(\delta,\theta)<5/24.
\]
\end{theorem}

For rational CM curves Panidapu--Thorner, Theorem~1.1, established the last asymptotic under $\delta+\theta<5/24$. Their Lemma~3.4 supplies the sharper density estimate used for the last range here. Short-interval and angular smoothing have a substantial history, including M.~D.~Coleman's work on primes represented by norm forms, the joint maximum range in Khale--O'Kuhn--Panidapu--Sun--Zhang for normed prime ideals, and J.~Stucky, \href{https://doi.org/10.1093/qmath/haab009}{\emph{Gaussian primes in narrow sectors}}. We give the class-dependent transfer and the details that allow the two exponents to vary separately.

\begin{proof}
For $g\in J$ and $B$ an abelian subgroup attaining $b_J(g)$, the definition of induction of class functions gives
\[
   \mathbf1_{g^J}=\frac{|B|}{|C_J(g)|}\operatorname{Ind}_B^J\mathbf1_{\{g\}},
 \qquad \mathbf1_{\{g\}}(v)=\frac1{|B|}\sum_{\xi\in\widehat B}
                     \overline{\xi(g)}\xi(v)\quad(v\in B).
\]
Indeed, at a conjugate of $g$ precisely $|C_J(g)|$ summands contribute to the induced character before division by $|B|$. Apply this identity with $J=H$ over $M$ and with $J=G$ for the nonsplit count over $F$. The logarithmic derivatives for the split angles are linear combinations of those for the ordinary Hecke characters
\[
        (\nu\circ N_{W^B/M})^{j}\xi,\qquad \xi\in\widehat B,
\]
over the indicated fixed fields $W^B$. The resulting identity holds for all powers of an unramified Frobenius by the induction law for local Euler factors. The nonsplit count uses $j=0$ over its corresponding fields. The degrees are bounded by $d_C$; the finite moduli and coefficient lists are fixed. For $j=0$, the coefficient of the trivial Hecke character over $W^B$ for a class $g^J$ equals $1/|C_J(g)|$. Summation gives the same masses as in Theorem~\ref{zfnum:cm-sectors}.

Let $N_*(\sigma,V)$ count with multiplicity the nontrivial zeros at real part at least $\sigma$ and absolute ordinate at most $V$ of these finitely many lists of Hecke functions, including each angular index $|j|\le V$ for the split lists and only index zero for the nonsplit lists. We include all zeros before class-function coefficients are applied. Coleman's density theorem (\emph{The distribution of points at which norm-forms are prime}, J. Number Theory 41 (1992), Lemma~1, also Khale et al., Theorem~2.6) gives
\begin{equation}\label{zfnum:familydensity}
 N_*(\sigma,V)\ll V^{A(1-\sigma)}(\log(3V))^B,\qquad
 A=5d_C/2,\quad \tfrac12\le\sigma\le1 .
\end{equation}
To check the use of the angular family, fix a split field $W^B$. The infinity vector is $j$ times one nonzero lattice vector and the finite ray characters range over a fixed finite list. An imaginary norm power, in coordinates for a basis of the lattice, translates the ordinate by $O(|j|)$. These characters form, up to bounded multiplicity and this translation, a subset of Coleman's full lattice box and vertical interval of size $cV$ for a fixed $c$. The same bound includes the fixed zero-index characters. In the rational unfiltered case the product of the primitive angular factors, the Riemann zeta function and the quadratic Dirichlet factor is exactly the family of Panidapu--Thorner, Lemma~3.4, up to bounded multiplicities from signs and the zero-index terms. Their estimate $(J^2+V^2)^{(12/5+\varepsilon')(1-\sigma)}(\log(3JV))^B$ supplies \eqref{zfnum:familydensity} with $A$ any fixed number greater than $24/5$ when $J\le V$. In both density estimates Coleman's zero-free estimate excludes all these zeros in $\sigma>1-\lambda_V$, where
\[
       \lambda_V=c_0(\log(3V))^{-2/3}(\log\log(16V))^{-1/3}.
\]
For the fixed finite-order characters any exceptional real zero is absorbed by decreasing $c_0$. All constants can be taken uniform over the finitely many classes.

Write $w=\max(\delta,\theta)$ and choose $\varepsilon>0$ with $A(w+3\varepsilon)<1$. Set
\[
 J=\lceil x^{\theta+\varepsilon}\rceil,\qquad V=x^{w+3\varepsilon},
 \qquad a=h x^{-\varepsilon}.
\]
Choose an integer $r\ge2$ with $2r\varepsilon>\delta+\theta+2\varepsilon$. Convolve the indicator of $[x\mp ra/2,x+h\pm ra/2]$ with $r$ copies of $a^{-1}\mathbf1_{[-a/2,a/2]}$. This gives two nonnegative weights $w^-$ and $w^+$ bounding the interval indicator below and above, with integrals $h-ra$ and $h+ra$. Endpoint choices do not affect the bounds. Their Mellin transforms $W^\pm(s)$ are entire. For $-1/4\le\sigma\le 1$, either transform satisfies
\begin{align*}
 |W^\pm(\sigma+it)|&\ll_r hx^{\sigma-1},\\
 |W^\pm(\sigma+it)|&\ll_r x^\sigma(x/a)^r |t|^{-r-1}\qquad(|t|>2r+2).
\end{align*}
The first bound uses the support. For the second, integrate by parts $r+1$ times, with the last derivative taken as a measure. Its total variation is at most $2^{r+1}a^{-r}$, from the endpoint masses of the $r+1$ interval factors.

For any one of the primitive Hecke characters $\chi$ in the calculation the smoothed prime-power sum with coefficients from $-L'/L$ equals
\[
 \mathbf1_{\chi=1}W^\pm(1)-\sum_{\rho}W^\pm(\rho)
     + O_r\big((1+h/a)x^{-1/4}\log(3x+|j|)+h/x\big).
\]
To verify the error, move Mellin inversion from $\Re s=2$ to $-1/4$. On the new line the functional equation and the convergent Euler derivative at $5/4$ give $L'/L\ll\log(3+|t|+|j|)$. The gamma shifts at a complex embedding have real shift $|j|/2$ for an angular constituent, and the shifts of the finite characters at real places are 0 or 1 in $\Gamma_{\mathbb R}$, so their poles stay separated from the new line. Split the integral at $|t|=2r+2+x/a$ and use the two preceding estimates. A possible trivial zero at $s=0$ contributes $O(h/x)$ and has uniformly bounded order. The unit-interval zero count is $O(\log(3+|j|+|t|))$ by the logarithmic derivative at $2+it$ as above. One may therefore move the contour first at heights separated from all ordinates of the particular character by $\gg 1/\log(3+|j|+|t|)$; the local logarithmic-derivative bound and the decay of $W^\pm$ make the horizontal tails vanish. The finitely many ramified Euler factors and discarded geometric bad primes have bounded coefficients at each of their prime powers, giving $O(\log(2x))$ per character, which suffices here.

For $|t|>V$, the unit-interval zero count and the second Mellin bound give a total over all $|j|\le J$ of
\[
                O_r\big(Jx(x/a)^r V^{-r}\log(3JV)\big).
\]
Dividing by $h$, this is $O(x^{\delta+\theta+\varepsilon-2r\varepsilon}\log(3x))$ since $V/(x/a)\ge x^{2\varepsilon}$. For $|t|\le V$, each Hecke functional equation pairs a zero at $\beta+it$ with one at $1-\beta+it$ of the same primitive $L$-function. It suffices to bound the sum of $x^{\beta-1}$ for the half with $\beta\ge1/2$. Partial summation of \eqref{zfnum:familydensity} bounds that sum by
\[
 O\!\left((\log(3x))^{B+1}
       \exp\{-[1-A(w+3\varepsilon)]\lambda_V\log x\}\right).
\]
Indeed $x^{\sigma-1}V^{A(1-\sigma)}
=\exp\{-(1-\sigma)[1-A(w+3\varepsilon)]\log x\}$ and no zero exceeds $1-\lambda_V$. This is $O(\exp\{-b(\log x/\log\log(3x))^{1/3}\})$ after decreasing $b$.

For a split class, the inverse image of $I$ is the union of at most two circular arcs. Its degree-$J$ trigonometric majorant and minorant have constant terms $|I|/\pi+O(1/J)$ and all nonconstant Fourier coefficients are $O(\min(|I|,1/|j|)+1/J)=O(|I|)$. Multiply the lower polynomial by $w^-$ and the upper polynomial by $w^+$. The pointwise bounds hold even where the minorant is negative, since the two time weights are nonnegative. All the zero terms are now $O(h|I|\exp\{-b(\log x/\log\log(3x))^{1/3}\})$. Replacing $W^\pm(1)$ by $h$ and the constant coefficient by $|I|/\pi$ costs $O(h|I|x^{-\varepsilon})$. Higher prime powers cost $O_F(\sqrt{x}\log^2(2x))$ for each frequency by elementary counting. Even their total cost after the Fourier sum is $O(J|I|\sqrt{x}\log^2(2x))$, a power saving relative to $h|I|$, since for a split class $d_C\ge[M:\Q]\ge2$ and hence $\delta+\theta+\varepsilon<1/2$ on reducing $\varepsilon$. The same inequality holds in the rational refined range. Dividing the split prime sum over $M$ by $[M:F]$ costs another $O_F(\sqrt{x}\log(2x))$ from relative residue degree two, also within the error. For a nonsplit class only $j=0$ occurs; $h\gg x^{1-\delta}$ with $\delta<2/5$ suffices to discard higher powers. Its indicator in the sector is exactly $\mathbf1_{\pi/2\in I}$. Summing the finitely many class estimates yields the claimed relative error. If $\mu_C(I)=0$, then the class and angle requirements exclude every good unramified prime.
\end{proof}

For instance every rational CM curve has a prime in $(x,x+x^{4/5}]$ with angle in any interval of length $x^{-1/5}$, for all sufficiently large $x$. Away from the atom the weighted number is asymptotic to $x^{3/5}/(2\pi)$. With arbitrary fixed congruence conditions the class-field version gives the same conclusion with $4/5$ increased and $1/5$ decreased by arbitrarily small fixed amounts, provided the congruence condition admits a split prime.

\subsection{Several distinct CM fields at once}
Let $E_1,\ldots,E_e$ be elliptic curves over $F$, with pairwise distinct imaginary quadratic endomorphism algebras $K_1,\ldots,K_e$. Choose their embedded copies in an algebraic closure of $F$ through the endomorphism action on the tangent spaces and put $M_i=FK_i$. For $S\subseteq\{1,\ldots,e\}$ let $L_S$ be the compositum of the $M_i$ for $i\in S$, taking $L_\varnothing=F$. The splitting signs need not be independent over $F$.

\begin{lemma}\label{zfnum:indangular}
For any nonempty $S$ and any finite extension $U/L_S$, the unitary CM characters $\nu_i\circ N_{U/M_i}$, $i\in S$, have linearly independent angular infinity vectors over $\Q$. In particular their nonzero integral combinations remain nontrivial after multiplication by any finite character or pure imaginary norm power.
\end{lemma}
\begin{proof}
The CM principal-ideal formula in Theorem~\ref{cm:factor} gives the infinity argument from $N_{U/K_i}$, up to a choice of conjugation for each $i$. On the set of all complex embeddings of $U$ its signed exponent is $1$ or $-1$ according to the restriction to $K_i$. Pass to the multiquadratic field $K$ generated over $\Q$ by the $K_i$, $i\in S$. Every embedding of $K$ has the same number $[U:K]$ of extensions to $U$. On $\operatorname{Gal}(K/\Q)$ the exponent functions are, up to independent signs, the distinct quadratic characters belonging to the distinct subfields $K_i$. Their pairwise inner products vanish and each has positive norm. A linear relation at complex places would give a relation on all embeddings, since conjugate embeddings have opposite exponents. This proves independence. Unit-circle restrictions remove both kinds of twists in the last assertion.
\end{proof}

\begin{theorem}\label{zfnum:cm-mixed}
For any fixed finite-image Galois representation $\rho$ of $F$ and all positive odd integers $m_1,\ldots,m_e$, the geometric completion of
\[
       \rho\otimes\bigotimes_{i=1}^e\operatorname{Sym}^{m_i}E_i
\]
in unitary normalization extends meromorphically to the plane. It has neither zeros nor poles in both edges of width
\[
       c\,(\log V)^{-2/3}(\log\log(5V))^{-1/3},
       \qquad V=3+|t|+\sum_{i=1}^e m_i ,
\]
with $c>0$ depending only on the curves, ground field and finite representation.
\end{theorem}
\begin{proof}
Over $L=L_{\{1,\ldots,e\}}$ the tensor's CM part decomposes into pairwise distinct character lines. They have characters $\prod_i (\nu_i\circ N_{L/M_i})^{j_i}$, where $j_i$ runs through $-m_i,-m_i+2,\ldots,m_i$. These formulas in the unitary convention follow from the Tate-module formulas with their common Tate normalization, and apply to the actual local factors. Lemma~\ref{zfnum:indangular} shows that weight tuples distinguish the lines; every coordinate $j_i$ is nonzero. The group $\operatorname{Gal}(L/F)$ acts on them by the CM conjugations. Each nonidentity element changes some sign, since the $M_i$ generate $L$. Therefore every orbit has trivial stabilizer in this group. The full parameter for one orbit is induced from its character line over $L$: its distinct translates give precisely those lines, with no choice of descent factor. Tensoring with $\rho$ first restricts $\rho$ to $L$ inside each induction.

Apply Brauer's theorem to that fixed restricted finite representation. As in Theorem~\ref{zfnum:cm-artin}, every completion in the resulting signed product is an ordinary Hecke function, now with angular type a fixed integral linear map of $(j_1,\ldots,j_e)$. At every one of the finitely many intermediate fields Lemma~\ref{zfnum:indangular} makes this map injective. The finite conductors are bounded, the lattice sizes and any translated ordinates are $O(V)$, and the real-character exception is impossible. Coleman's zero-free theorem and the functional equations prove the assertion for both edges. This proof does not require cuspidal tensor functoriality.
\end{proof}

There is a joint prime-counting statement with the same independence. Let $P/F$ be finite Galois, $C$ a union of classes, $W=P L_{\{1,\ldots,e\}}$, $G=\operatorname{Gal}(W/F)$, and $C'$ the inverse image of $C$. For $g\in C'$ let $S(g)=\{i:g|_{M_i}=1\}$, $H_g=\operatorname{Gal}(W/L_{S(g)})$ and
\[
 d_* = \max_{g\in C'}[W:\Q]/b_{H_g}(g).
\]
For a box $\mathcal I=\prod_i I_i$ in $[0,\pi]^e$ put
\[
 \mu_C(\mathcal I)=\frac1{|G|}\sum_{g\in C'}
       \prod_{i\in S(g)}\frac{|I_i|}{\pi}
       \prod_{i\notin S(g)}\mathbf1_{\pi/2\in I_i}.
\]
Take $C'$ nonempty. This formula retains every dependence among the splitting conditions of the curves and the finite extension.

\begin{corollary}\label{zfnum:cm-multisector}
Let $x^{1-\delta}\le h\le x$, $|I_i|\ge x^{-\theta_i}$ with fixed nonnegative exponents and $\max(\delta,\theta_1,\ldots,\theta_e)<2/(5d_*)$. Then
\[
 \sum_{\substack{x<N\mathfrak p\le x+h,\ \operatorname{Frob}_{\mathfrak p}\in C\\
           \theta_{i,\mathfrak p}\in I_i\ \mathrm{for\ all}\ i}}^{\mathrm{good,\ unram.}}
       \log N\mathfrak p
   = h\mu_C(\mathcal I)
      \left(1+O\!\left(e^{-b(\log x/\log\log(3x))^{1/3}}\right)\right).
\]
\end{corollary}
\begin{proof}
Partition the classes by $S=S(g)$. An index outside $S$ has angle exactly $\pi/2$, giving the atomic factors. For the other indices count over $L_S$. Each prime of $F$ split in $L_S$ has $[L_S:F]$ primes there with the same folded angles, since any conjugation only replaces an individual unit phase by its inverse. Primes of higher relative degree contribute $O_{F,L_S}(\sqrt x\log x)$ to the count. The relevant union in $H_g$ is stable under conjugation and can be treated one class at a time. Use the abelian subgroup induction in Theorem~\ref{zfnum:cm-short}. Nonzero Fourier index vectors now pull back to distinct nonzero angular lattice vectors by Lemma~\ref{zfnum:indangular}. For $|j_i|\le J_i=\lceil x^{\theta_i+\varepsilon}\rceil$ these and their possible norm shifts are bounded by a fixed multiple of $\max_i J_i$. Thus the same single full-lattice density bound \eqref{zfnum:familydensity} with $A=5d_*/2$ controls the \emph{sum} of the zero counts over this entire box. There is one constant angular mode.

For completeness, if $P_i^-,P_i^+$ are the one-coordinate minorant and majorant, the product $\prod_i P_i^+$ majorizes the box and
\[
    \prod_iP_i^+-\sum_i(P_i^+-P_i^-)\prod_{\ell\ne i}P_\ell^+
\]
minorizes it. This follows by telescoping the product with its indicators, using $P_i^+\ge0$. The products here are over the split indices. Every Fourier coefficient is $O_e(\prod_{i\in S}|I_i|)$, and both constant terms equal $\prod_{i\in S}(|I_i|/\pi)$ up to relative $O_e(x^{-\varepsilon})$. The radial convolution of Theorem~\ref{zfnum:cm-short} works with $V=x^{\max(\delta,\max_i\theta_i)+3\varepsilon}$; its high-ordinate error acquires $O(\prod_i J_i)$ in place of $O(J)$, absorbed by choosing $2r\varepsilon>\delta+\sum_i(\theta_i+\varepsilon)+\varepsilon$. The low-ordinate sum is unchanged by the preceding full-lattice estimate. To discard prime powers frequency by frequency it suffices that $\delta+\sum_{i\in S}(\theta_i+\varepsilon)<1/2$. For $S\ne\varnothing$, independence of the pure angular coordinates gives $|S|\le[L_S:\Q]/2\le d_*/2$. Since $d_*\ge2$, before adding the epsilons the left side is smaller than $1/5+2/(5d_*)\le2/5$, as required. Purely atomic classes use $\delta<2/5$ alone. Divide the class main terms by $[L_S:F]$ as above. Their coefficients add to the stated $\mu_C(\mathcal I)$, and summing the uniform relative errors proves the result.
\end{proof}

For example, take $E_1:y^2=x^3-x$ and $E_2:y^2=x^3+1$ over $\Q$. Their endomorphism fields are $\Q(i)$ and $\Q(\sqrt{-3})$. For primes splitting in both, the class field needed is biquadratic of degree four. For arbitrary intervals $I_1,I_2$, each of length $x^{-1/12}$, and $h=x^{11/12}$, the weighted number of these primes with both angle restrictions is asymptotic to $x^{3/4}/(4\pi^2)$. Here $1/12<2/(5\cdot4)$. The factor $1/4$ is the density of the common splitting condition.

\subsection{The endpoint multiplicity for arbitrary CM tensors}
The same orbit calculation determines the endpoint when some powers are even or when several curves have the same CM algebra. Drop the hypothesis that the $K_i$ are distinct and allow arbitrary integers $m_i\ge0$. Write $\kappa_i$ for the quadratic character of $M_i/F$, with value one when $M_i=F$. Form the set of formal tuples
\[
 \mathcal J=\prod_{i=1}^e\{-m_i,-m_i+2,\ldots,m_i\}.
\]
The group of the compositum of the $M_i$ over $F$ acts on $\mathcal J$ by the splitting signs. For an orbit representative $\mathbf j$, put $S=\{i:j_i\ne0\}$. Over $L_S$ define the unitary Hecke character
\[
 \xi_{\mathbf j}
  =\prod_{i\in S}(\nu_i\circ N_{L_S/M_i})^{j_i}
    \ \prod_{i\notin S}(\kappa_i\circ N_{L_S/F})^{m_i/2}.
\]
The exponent $m_i/2$ is integral whenever $j_i=0$. A character $\xi_{\mathbf j}$ with trivial angular infinity vector is finite order. Indeed it is trivial on the connected archimedean ideles by the relative-norm formula, and its ideal conductor divides a fixed modulus. It then factors through a finite ray class group.

\begin{theorem}\label{zfnum:cm-alltensors}
For any fixed finite-image representation $\rho$ of $F$ set
\[
 a(\mathbf m,\rho)=
 \sum_{\substack{\mathbf j\in\mathcal J/\operatorname{Gal}(L/F)\\
                   \xi_{\mathbf j}\text{ finite order}}}
 \dim\bigl(\rho|_{G_{L_S}}\otimes\xi_{\mathbf j}\bigr)^{G_{L_S}},
 \qquad L=L_{\{1,\ldots,e\}}.
\]
In unitary normalization the complete function of
$\rho\otimes\bigotimes_i\operatorname{Sym}^{m_i}E_i$ has a pole of order $a(\mathbf m,\rho)$ at both 0 and 1 when this integer is positive, and is regular and nonvanishing there when it is zero. Multiplying by $[s(s-1)]^{a(\mathbf m,\rho)}$ gives a meromorphic function with neither a zero nor a pole in both edges of width
\[
   \frac{c}{(\log V)^{2/3}(\log\log(5V))^{1/3}},
                 \qquad V=3+|t|+\sum_i m_i.
\]
The constant depends on $F$, the curves and $\rho$, and is independent of all the exponents. For $\rho$ trivial the pole-removed completion is entire on the whole plane.
\end{theorem}
\begin{proof}
Use the geometric CM lines in each two-dimensional parameter, and take their symmetric monomials. These are distinct vector lines even if the associated characters happen to coincide. An element outside $G_{M_i}$ exchanges the two lines of $E_i$. In a unitary basis for the Frobenius normalization its exchange matrix has off-diagonal entries $a_i,b_i$ with $a_i b_i=-1$ from determinant one. Equivalently before normalization its determinant is $-a_i b_i$ and the middle monomial is normalized by the integral determinant twist $-m_i/2$. Thus the zero-weight middle line acquires precisely $\kappa_i^{m_i/2}$. On the group fixing the two lines a monomial of signed weight $j_i$ has unitary character $\nu_i^{j_i}$. The group that fixes the labelled tensor line is exactly $G_{L_S}$: it may exchange the two lines only at indices of weight zero. The translates of each such line give the induced parameter from its own stabilizer, regardless of coincidences between the characters of different lines. Applying induction to the actual Tate-module representations first and then imposing the common analytic weight shift gives
\[
 \Lambda_F\left(s,\rho\otimes\bigotimes_i\operatorname{Sym}^{m_i}E_i\right)
  \doteq\prod_{\mathbf j\in\mathcal J/\operatorname{Gal}(L/F)}
          \Lambda_{L_S}(s,\xi_{\mathbf j}\otimes\rho|_{G_{L_S}}),
\]
where $\doteq$ permits the constant normalizations in induction of completions. In particular with $\rho$ trivial the right side uses only individual Hecke functions with positive multiplicities. Their sole possible poles are the two simple endpoint poles of a trivial Hecke character.

For general $\rho$, implement Brauer induction once for each of its restrictions to the finitely many fields $L_S$. This gives a signed product over a fixed finite collection of fields. At any field in that list, a nonzero angular vector stays nonzero after norm pullback, as may be seen on a single complex embedding above a nonzero coordinate. Such a Hecke character, even with a finite twist, has no exceptional zero. Its lattice index and any norm translation in lattice coordinates are $O(V)$ and its finite conductor is bounded. If the angular vector is zero, all the resulting Hecke characters have finite order and range over a fixed finite list because the modulus is fixed. At 1 their pole order in a particular signed product is the coefficient of the trivial Hecke character. For the Artin factor over $L_S$ this coefficient is the dimension of its invariant space, by Frobenius reciprocity applied to the Brauer decomposition. Nontrivial finite Hecke characters have nonzero value there. The same pole order occurs at 0 by the dual functional equation, since a finite representation and its dual have invariant spaces of the same dimension. This gives the formula for the exact endpoint multiplicity. Coleman's uniform zero-free theorem applies to all nonzero angular factors, while the finite list at zero angular type is covered after decreasing $c$ to accommodate any fixed exceptional zero. Their signed product, after the specified endpoint removal, is consequently regular and nonvanishing in the common two edges.
\end{proof}

For example, with $\rho$ trivial, $E_1=E_2=E$ nonsplit CM and $m_1=m_2=1$, the two orbits have representatives $(1,1)$ and $(1,-1)$. The latter gives the trivial Hecke character over the endomorphism field. The complete Rankin square has one pole at either endpoint. If all endomorphisms are already defined over $F$, the two labelled opposite tuples are separate and the order is two.

\subsection{Non-CM curves over cyclic towers of arbitrary signature}
Potential automorphy and cyclic base change can be performed in different fields. This also accommodates nonabelian Artin types without assuming their automorphy.

\begin{theorem}\label{zfnum:mixedsignature}
Let $F$ be totally real or an imaginary CM field, and let $E/F$ be non-CM. Let $A/F$ be a finite Galois extension with $A$ totally real or CM; the group $\operatorname{Gal}(A/F)$ is unrestricted. Let $K/F$ be a finite tower of cyclic Galois steps, with no condition on the signatures or on the normality of $K/F$. Fix $m\ge1$. For every complex representation $\rho$ of $\operatorname{Gal}(KA/K)$ and every finite-order Hecke character $\eta$ of $K$, the geometric unitary function
\[
             L_K(\operatorname{Sym}^m E_K\otimes\rho\otimes\eta,s)
\]
is meromorphic and has neither zeros nor poles in a right edge of inverse-logarithmic width. More precisely there are a finite extension $L/A$, Galois over $F$, and $C_0\ge3$, depending on $E,F,A,K,m$, for which every right-half-strip divisor point $\beta+it$ satisfies
\[
 1-\beta\ge\frac{2(1-\sqrt{s_0})^2}{[KL:\Q]T(1)
    [\log(C_0(3+|t|))+[K:\Q]^{-1}\log N_{K/\Q}\mathfrak f_\eta]}.
\]
Here $S,T$ are those of Lemma~\ref{amp:kernel}. Both $L$ and $C_0$ can be chosen independent of $\rho$ and $\eta$. Two geometrically nonisogenous non-CM curves over $F$ also have over $K$ the meromorphic continuation and zero- and pole-free region in Corollary~\ref{zfnum:pair}, with the same finite twists, and with $[KL:\Q]$ in its limiting coefficient.
\end{theorem}
\begin{proof}
Use Allen et al., Theorem~7.1.11, with the elliptic systems over $A$ and with its lower field equal to $F$. In the source theorem these two fields are denoted $F$ and $F_0$ respectively. For the required finite list of symmetric powers it produces a sufficient totally real or CM field $L$ containing $A$ which is Galois over $F$. Taking that sufficient field itself meets its disjointness and unramifiedness conditions. All the symmetric powers in the list are cuspidal over $L$. Their geometric systems are defined over the original $F$.

Identify $J=\operatorname{Gal}(KL/K)$ with its restriction subgroup in $\operatorname{Gal}(L/F)$. Inflating $\rho$ from $\operatorname{Gal}(KA/K)$ to $J$ is legitimate because $A\subseteq L$. Brauer induction in $J$ gives
\[
              \rho=\sum_i n_i\operatorname{Ind}_{H_i}^J\xi_i
\]
with elementary $H_i$ and one-dimensional $\xi_i$. Its inducing field is
\[
                           (KL)^{H_i}=K L^{H_i}.
\]
Indeed within $J$ the subgroup fixing $L^{H_i}$ is exactly $H_i$ by restriction, and every element of $J$ already fixes $K$. Each $L^{H_i}$ is totally real or CM: the normal closure of a CM field is CM, with central complex conjugation; any subfield is then totally real or is quadratic and totally imaginary over its conjugation-fixed field. The extension $L/L^{H_i}$ is Galois with soluble group $H_i$. Thus Allen et al., Proposition~6.5.13(2), applied to the geometric representations defined on $G_{L^{H_i}}$, descends the cuspidal symmetric powers to $L^{H_i}$. The irreducibility required in that proposition persists on $G_L$ by elliptic open image.

Each field $K L^{H_i}$ is reached from $L^{H_i}$ by adjoining the successive cyclic steps in the prescribed tower for $K/F$. Each step after compositum is again cyclic or trivial. Apply Arthur--Clozel, \href{https://www.claymath.org/library/cw/arthur/pdf/30.pdf}{\emph{Simple algebras, base change, and the advanced theory of the trace formula}}, Chapter~III, cyclic base change. The transferred symmetric powers continue to match the geometric Frobenius traces almost everywhere, since unramified base change raises each Satake root to the local residue degree. Each transfer is cuspidal after refining to prime-degree cyclic steps: loss of cuspidality would give a nontrivial finite self-twist by a character of that step. Trace comparison, Chebotarev, and semisimplicity would then give the same self-twist of the geometric symmetric power. Its restriction to the kernel of that character is irreducible, so any intertwiner is scalar on the restriction by Schur's lemma. A scalar cannot implement a nontrivial twist. This rules out the loss at every step, without requiring an automorphic descent theorem over a mixed-signature field.

For each $i$, the twist of the resulting cuspidal target over $K L^{H_i}$ is by $\xi_i(\eta\circ N_{K L^{H_i}/K})$. The local induction formula gives the signed product for the function in the theorem, with the finite Euler replacements described in Theorem~\ref{zfnum:potential}. Apply its positive-polynomial argument over each $K L^{H_i}$. That proof uses the actual good Euler parameters, cuspidality of the finite index list and the generic unitary standard bounds at the finitely many exceptional places; none requires a signature assumption. The additional conductor contribution scales with $[K L^{H_i}:K]\log N\mathfrak f_\eta$, and the height contribution with $[K L^{H_i}:\Q]T(1)$. Taking the minimum of the regions and bounding the latter degree by $[KL:\Q]$ gives the formula. There are only finitely many irreducible $\rho$ for this fixed Artin group. A general representation is a direct sum of them, so all fixed constants may be taken uniform in $\rho$. The two-curve argument over these same inducing fields uses the rank products and geometric independence in Corollary~\ref{zfnum:pair}; the necessary finite lists for both original curves are taken simultaneously over $A$. This proves the assertion for the pair as well.
\end{proof}

For example let $F=\Q(\sqrt2)$, $K=\Q(\sqrt[4]{2})$, and
\[
                        E: y^2=x^3+x+1+\sqrt2 .
\]
Then $K$ has two real places and one complex place, and $K/F$ is quadratic. The curve's $j$-invariant is $6912/(85+54\sqrt2)$. The denominator has norm $1393$, coprime to $6912$, so this $j$ is nonintegral and $E$ has no CM. It is irrational, so the curve is not a quadratic twist of a curve over $\Q$. Taking $A=F$ gives the asserted edges for this genuinely nonrational $j$ over a field of mixed signature. In a different specialization $K=F=\Q$, the theorem permits any finite Artin representation on which complex conjugation acts by a scalar. The field cut out by that representation is totally real if the scalar is $1$. If it is $-1$, conjugation is a central involution and its fixed field is totally real, giving a CM field as required.

In the situation of the theorem, fix a union $C$ of conjugacy classes of $\operatorname{Gal}(KA/K)$ and fix $m\ge1$. At good unramified prime ideals of $K$, applying Corollary~\ref{zfnum:prime} to each Brauer constituent and then expanding the indicator of $C$ in ordinary characters gives
\[
     \sum_{\substack{N\mathfrak p\le x\\ \operatorname{Frob}_{\mathfrak p}\in C}}
        U_m(\cos\theta_{\mathfrak p})\log N\mathfrak p
                 =O_{E,F,A,K,m}\bigl(xe^{-b\sqrt{\log x}}\bigr).
\]
Chebotarev supplies the constant mode. Hence every fixed positive-density Frobenius condition has the usual semicircular prime-angle distribution. The same proof applies jointly to the pair in the theorem by approximation with products of characters. The constants in the prime error are for fixed symmetric-power indices.

A concrete nonsolvable Artin type in Theorem~\ref{zfnum:mixedsignature} comes from the splitting field $A$ of
\[
                         f(X)=X^5-9X^3+16X+1
\]
over $\Q$. The signs of $f$ at $-3,-2,-1,0,1,2,3$ are $-,+,-,+,+,-,+$, giving five real roots in disjoint intervals. Modulo two it becomes $X^5+X^3+1$, with no linear factor and no factor $X^2+X+1$, and so is irreducible. Modulo seven it factors as
\[
                   (X+1)(X-3)(X-2)(X^2-3X-1).
\]
The quadratic has nonsquare discriminant $6\bmod 7$ and the factors are distinct. Frobenius therefore provides a $5$-cycle and a transposition. Conjugating that transposition by the $5$-cycle gives the edges of a connected graph on the five letters; its edge transpositions generate $S_5$. Thus $A$ is totally real and has group $S_5$. For instance take $\rho$ to be the six-dimensional exterior square of its standard four-dimensional representation. For every rational non-CM $E$ this gives an explicit degree-twelve geometric elliptic--Artin function with the stated meromorphic continuation and all-height right edge. The proof does not require Artin automorphy for this $S_5$ representation.

\subsection{Every two-dimensional Artin twist of a rational elliptic curve}
An odd two-dimensional Artin representation usually has noncentral complex conjugation. Its character therefore calls for a different nonnegative auxiliary function.

\begin{proposition}\label{zfnum:artin-character}
Let $E/\Q$ be non-CM and let $\rho$ be an irreducible complex Artin representation of dimension $r$. Suppose $\rho$ and every irreducible constituent of $\rho\otimes\rho^\vee$ admit unitary cuspidal weak Artin realizations. For fixed $E,\rho$ there are constants $C,c>0$ such that, simultaneously for all $m\ge1$ and primitive Dirichlet characters $\chi$ of conductor $q$, the geometric function $L(\operatorname{Sym}^m E\otimes\rho\otimes\chi,s)$ is meromorphic and neither vanishes nor has a pole when
\[
      \Re s>1-\frac{c}{(m+1)\log(Cq(m+2)(3+|\Im s|))}.
\]
For fixed $m,\chi$ the high-ordinate limit inferior of $(1-\beta)\log|t|$ over right-half-strip divisor points is at least $2(1-\sqrt{s_0})^2/(rT(1))$.
\end{proposition}
\begin{proof}
Write $u(g)=\mathrm{tr}\,\rho(g)$ on its finite group. The static and oscillating character expressions, where $x$ is the normalized elliptic Frobenius trace divided by two and $|z|=1$, are
\[
       S(x)\frac{1+|u(g)|^2}{2}
          \quad\hbox{and}\quad T(x)\Re(zu(g)).
\]
Their sum is nonnegative since $(1+|u|^2)/2\ge|u|$ and $S\ge|T|$. The first finite character is a nonnegative rational combination of irreducibles, coming from $(1+\rho\otimes\rho^\vee)/2$. Its trivial multiplicity is exactly one by Schur's lemma. The second uses the given irreducible $\rho$ and its conjugate mate. Use the weak Artin realizations for the finite characters and the full Newton--Thorne lifts for all elliptic symmetric powers. Each product term is a standard Rankin--Selberg function with the indicated twists. At good prime powers these give exactly the displayed formula, with $g$ itself the required power of Frobenius.

The static degree sums to $S(1)(1+r^2)/2$ and the translated degree to $rT(1)$. The only pole is the unshifted zeta weight $s_0$. To check other equal-rank possibilities with a positive symmetric index $i$, take a good rational prime splitting completely in the finite Galois field of all the Artin data. The Artin parameter has scalar Satake class there, also after any additional character twist. The $i$th elliptic symmetric-power parameter cannot be scalar: equality of its first two successive eigenvalues would give equal eigenvalues of $E$ and thus $a_p(E)=\pm2\sqrt p$, contrary to integrality. This excludes pole-producing contragredient equivalences, even with an imaginary norm shift. For zero index the static polar term is already accounted for, and the oscillating polynomial has zero constant coefficient. The target zero is counted with weight one using its dual if necessary.

For fixed weak Artin realizations the absolute aggregate degree is $O_\rho(m+1)$, and the highest elliptic index in the two polynomials is $O((m+1)^2)$. On any ramification group the conductor codimension of a tensor is bounded by its dimension times the sum of the original conductor codimensions, including the twisting character. This gives total logarithmic conductor $O_{E,\rho}((m+1)\log(3q))$. The unitary elliptic Weil--$\mathrm{SL}_2$ parameters are tempered. Tensoring with the fixed generic unitary local representations of the finite types gives reciprocal roots $<(Nv)^{1/2}$ at every exceptional place, uniformly in the elliptic index with a fixed gap in their exponents there. The same applies after finite character twisting; at a new prime of a character twist an originally unramified factor disappears. Thus finitely many replacement Euler factors preserve the divisor on $\Re s\ge1/2$ and cost $O_{E,\rho}((m+1)\log(3q))$ in the exterior positivity estimate. The archimedean shifts have real parts $>-1/2$ and magnitude $O_\rho((m+1)^2+|t|)$, so the common logarithmic-derivative proof of Theorem~\ref{amp:linear} gives the displayed region. For fixed indices only the translated factors contribute the leading $\tfrac12 rT(1)\log|t|$. Optimization as in Theorem~\ref{zfnum:potential} yields the stated limiting coefficient.
\end{proof}

\begin{corollary}\label{zfnum:allranktwo}
For every rational elliptic curve, every positive symmetric power and every two-dimensional complex Artin representation $\rho$ of $G_\Q$, the geometric unitary twisted function continues meromorphically and is holomorphic and nonvanishing in a right edge of width $c/\log(3+|t|)$ for fixed data, with the exception described below for CM even powers. For non-CM curves and irreducible odd $\rho$ the index-linear region of Proposition~\ref{zfnum:artin-character} holds, including further Dirichlet twists. For CM curves the uniform Vinogradov--Korobov region in Theorem~\ref{zfnum:cm-alltensors} holds, with its explicit endpoint pole removal when needed.
\end{corollary}
\begin{proof}
Reducible $\rho$ is a sum of Dirichlet characters. For non-CM curves and irreducible even $\rho$, the involution of complex conjugation with determinant one has two equal eigenvalues $1$ or $-1$ and so is scalar. Theorem~\ref{zfnum:mixedsignature} applies. Every irreducible odd $\rho$ comes from a cuspidal weight-one newform by C.~Khare and J.-P.~Wintenberger, \href{https://www.math.ucla.edu/~shekhar/papers/results.pdf}{\emph{Serre's modularity conjecture (I)}}, Corollary~10.2(ii). In the nondihedral case its adjoint is the cuspidal Gelbart--Jacquet lift and realizes the irreducible adjoint Artin representation. In the dihedral case write $\rho=\operatorname{Ind}_{M/\Q}\xi$. Its adjoint is the quadratic character of $M/\Q$ plus the induction of $\xi/\xi^\sigma$ for the nontrivial $\sigma\in\operatorname{Gal}(M/\Q)$. These are automorphic by degree-one class field theory and quadratic induction, also if the latter splits into two characters. Since $\rho\otimes\rho^\vee=1\oplus\operatorname{Ad}\rho$, Proposition~\ref{zfnum:artin-character} applies. The CM assertion follows from the full orbit formula without either parity restriction on $\rho$.
\end{proof}

\clearpage
\section{Complex multiplication and the angular Hecke character}\label{cm:hecke}
Let $E/k$ be an elliptic curve over a number field and suppose that $\operatorname{End}^0(E_{\overline k})$ is an imaginary quadratic field $K$. Put $H=kK$ in a fixed algebraic closure. In this section the elliptic variable is motivic: at a prime $v$ of good reduction of norm $q_v$ the local polynomial is $1-a_vT+q_vT^2$, $T=q_v^{-s}$. The critical line is $\Re s=1$. For unitary Hecke characters we use $z=s-1/2$ instead.

\begin{theorem}[The CM reduction over an arbitrary base]\label{cm:factor}
There is an algebraic Hecke character $\Psi$ of $H$ such that, for the complete set of finite Euler factors,
\[
 L(E/k,s)=
 \begin{cases}
   L_H(s,\Psi)L_H(s,\overline\Psi),&K\subset k,\\
   L_H(s,\Psi),&K\not\subset k .
 \end{cases}
\]
The same $\Psi$ describes $E/H$. At ideals coprime to its conductor, $|\Psi(\mathfrak a)|=(\mathrm N\mathfrak a)^{1/2}$. Fix the embedding of $K$ associated to the action on the tangent line, up to complex conjugation. There is an ideal $\mathfrak f$ such that $\Psi((a))=\mathrm N_{H/K}a$ for $a\equiv1\pmod{\mathfrak f}$. In particular $\eta(\mathfrak a)=\Psi(\mathfrak a)/(\mathrm N\mathfrak a)^{1/2}$ is unitary and $L_H(s,\Psi)=L_H(s-1/2,\eta)$. Finite abelian base changes do not remove its angular dependence.
\end{theorem}
\begin{proof}
We use the main theorem of complex multiplication in its Hecke-character form, including compatibility with the Tate-module action. This is the previously established arithmetic theorem of Deuring. See J.~Silverman, \emph{Advanced Topics in the Arithmetic of Elliptic Curves}, Chapter II, \S\S9--10, especially Theorem~10.5, and K.~Rubin and A.~Silverberg, \href{https://doi.org/10.1016/j.jnt.2009.01.020}{\emph{Point counting on reductions of CM elliptic curves}}, for the description by Frobenius and the relative norm. The statement applies to geometric endomorphism algebras and permits nonmaximal endomorphism orders. We explain the passage from this classical theorem to both factorizations, since replacing $k$ by $H$ without distinguishing them would introduce an extra factor.

The derivative embeds the geometric endomorphism algebra into $\overline k$ and this embedding commutes with Galois action. Thus all endomorphisms are defined over $H$, and if $H/k$ is quadratic its nonidentity element acts on the algebra by conjugation. Choose a rational prime $\ell$ and extend coefficients in $V_\ell(E)$ to $\overline\Q_\ell$. This space has one line for each of the two embeddings of $K$: it has dimension one over $K\otimes\Q_\ell$ before extending scalars. The two lines are $G_H$-stable. The CM theorem identifies the eigencharacter on one with the $\ell$-adic realization of $\Psi$, and its conjugate on the other with the realization of $\overline\Psi$. It gives the stated principal-ideal formula. The determinant is cyclotomic, in agreement with $\Psi\overline\Psi=\mathrm N$.

If $H=k$, direct sums multiply local factors. If $[H:k]=2$, every element in the other Galois coset interchanges the two lines. Choosing a vector on the first line and its translate on the second identifies the representation with the induction to $G_k$ of the character on the first. Artin induction for Euler factors gives the single $L$-function over $H$. At a place of bad reduction this equality uses the local determinant on inertia invariants. Both direct sum and local induction preserve this determinant formula, with the residue degree changing the variable to $T^{f(w/v)}$ for places above $v$. The CM curve has potentially good reduction everywhere (its $j$ is integral), so there is no missing monodromy-operator factor. To compute at a place over $\ell$ one uses any other auxiliary Tate prime. The archimedean parameters are likewise identified by the CM infinity type.

At a complex place of $H$, the local component of $\eta$ on the unit circle is $z\mapsto z$ or $\bar z$ when $|z|=1$, with inversion possible according to the idelic convention. Indeed one embedding from each conjugate pair occurs in $\mathrm N_{H/K}$, and dividing by $(\mathrm N_{H/\Q})^{1/2}$ leaves the argument. Under any finite extension $H'/H$, the character becomes $\eta\circ\mathrm N_{H'/H}$. Restrict at infinity to just one complex place of $H'$ and put unit components elsewhere. Its local norm to the complex place below is the identity up to identification, hence this circle still maps onto the unit circle. This proves the last assertion even for arbitrary finite base change.
\end{proof}

The analytic continuation and functional equation of Hecke's $L$-function apply here because this character and its conjugate have nonzero angular type. Neither can be a real power of the norm. In motivic normalization there is no pole in either of the two formulas. Therefore an off-line zero of a factor cannot disappear in the product. Complex conjugation relates the zeros of $L_H(z,\eta)$ to those of $L_H(z,\bar\eta)$, and the Hecke functional equation reflects their real parts about $1/2$. Thus for each curve the desired assertion is equivalent to the absence of zeros of this particular unitary Hecke function for $\Re z>1/2$ in the critical strip. This is a reduction to a degree-one character of the ideles of $H$. It is generally a degree greater than one as an Euler product indexed by rational primes.

\subsection*{A classical consequence for all nonzero angular powers}
One can reach the edge of the critical strip, together with angular equidistribution, without making any Riemann hypothesis. Here is an argument that also shows the size of the improvement needed in the interior. For each integer $n\ne0$, the Hecke function $L_H(z,\eta^n)$ is entire. It is nonzero on $\Re z>1$ by its absolutely convergent Euler product. On $\Re z=1$ fix a real $t$, and remove the finite set of ramified Euler factors from all the functions below. For real $\sigma>1$, the logarithm of
\[
 \zeta_H(\sigma)^3
 |L_H(\sigma+it,\eta^n)|^4
 |L_H(\sigma+2it,\eta^{2n})|
\]
is nonnegative after the common omissions. At a prime power the real coefficient in the logarithm is proportional to $3+4\cos u+\cos 2u=2(1+\cos u)^2$, with $u$ determined by the norm and the value of $\eta^n$. If the middle factor had a zero at $1+it$, then as $\sigma\downarrow1$ the product would tend to zero: the fourth-order vanishing dominates the third-order pole of the Dedekind factor, and the last factor has no pole because its angular type is $2n\ne0$. This contradicts the lower bound one. The finite removed factors are regular and nonzero on this line. In particular there is unconditional nonvanishing on the motivic edge $\Re s=3/2$, and for the completed function on the reflected edge $\Re s=1/2$.

For completeness let $\Lambda_H$ be the ideal von Mangoldt function. The functions on $\Re z>1$ with Dirichlet coefficients $\Lambda_H(\mathfrak a)(1\pm\Re(e^{i\beta}\eta(\mathfrak a)^n))$ have nonnegative coefficients, on ideals prime to the conductor. Their Dirichlet series are sums of $-\zeta'_H/\zeta_H$ and a linear combination of the two logarithmic derivatives for $\eta^{n}$ and $\bar\eta^n$. The preceding boundary result and the classical prime ideal theorem for $H$ give the hypotheses of the Wiener--Ikehara theorem with residue one at $z=1$. Consequently the two partial sums through norm $X$ are each $X+o(X)$. Subtract and use $\beta=0,\pi/2$ to obtain
\[
 \sum_{\mathrm N\mathfrak p\le X}
       \eta(\mathfrak p)^n\log\mathrm N\mathfrak p=o(X).
\]
The sum is over primes of good reduction. Higher prime powers removed from the logarithmic derivative contribute $O_H(\sqrt X\log^2(2X))$; for instance use the prime ideal theorem to majorize the sum over exponents at least two. Partial summation and Weyl's criterion give Haar equidistribution of the angles over $H$. If $H/k$ is quadratic, a prime of $k$ inert in $H$ has trace zero, because its Frobenius interchanges the two lines. Such primes have density $1/2$ by the quadratic Chebotarev theorem. Above split primes the two angles are conjugate, so the nonzero-trace distribution is half the pushforward of circle measure under $w\mapsto w+\bar w$. Primes of $H$ over inert primes of $k$ with norm at most $X$ have underlying norm at most $\sqrt X$ and do not affect this distribution.

This argument applies to every integer harmonic separately. Its error $o(X)$ still permits oscillations of order $X^{\beta}$ for every fixed $\beta<1$. For the first harmonic, a square-root error changes the information about zeros, as the concrete criterion below will show.

\section{An explicit CM curve and its zeros at real arguments}\label{cm:gaussian}
Let $E_0$ be $y^2=x^3-x$ over $\Q$, and let $K=\Q(i)$, $\mathcal O=\Z[i]$, $\lambda=1+i$. The analytic character can be constructed without first computing any special values. Every ideal $\mathfrak a$ coprime to $\lambda$ has a unique generator $\alpha\equiv1\pmod {\lambda^3}$. Indeed $\mathcal O$ is Euclidean, and the unit group $\{1,-1,i,-i\}$ maps bijectively to $(\mathcal O/\lambda^3)^\times$: the latter has order four and the four units remain distinct. Define $\Psi(\mathfrak a)=\alpha$. Uniqueness gives multiplicativity. The principal-ideal rule on residue $1$ and the angular infinity type show that this is a primitive Hecke character of conductor $\lambda^3$. Its conductor cannot divide $\lambda^2$, since $-1\equiv1\bmod\lambda^2$ and the identity ideal cannot have two different character values. Let $\eta=\Psi/\sqrt{\mathrm N}$.

\begin{lemma}[Which quadratic twist is this?]\label{cm:primary}
This $\Psi$, up to conjugation, gives $L(E_0/\Q,s)=L_K(s,\Psi)$. At odd rational primes put $T=p^{-s}$. The inverse local factors are
\[
 \begin{cases}
 (1-\pi T)(1-\bar\pi T),&p\equiv1\pmod4,\quad
          (p)=(\pi)(\bar\pi),\quad \pi\equiv1\pmod{\lambda^3},\\
 1+pT^2,&p\equiv3\pmod4 .
 \end{cases}
\]
The inverse factor is one at $2$. The conductor over $\Q$ is $32$.
\end{lemma}
\begin{proof}
The automorphism $(x,y)\mapsto(-x,iy)$ gives CM by $\Z[i]$, which is maximal. Its differential is multiplication by $i$ on the invariant differential $dx/(2y)$. The eight points of $E_0[\lambda^3]$ are rational over $K$. To see this, $E_0[\lambda]=\{O,(0,0)\}$ and $[\lambda]^2=[2i]$, so $E_0[\lambda^3]$ consists of $E_0[2]$ and the four points above $(0,0)$ under doubling. In the latter fiber the doubling formula gives $(x^2+1)^2=0$, and the solutions are $(i,\pm(1-i))$ and $(-i,\pm(1+i))$. Frobenius at each good prime $\mathfrak p\nmid2$ therefore acts trivially on this entire kernel. The relative-norm and Tate-module clauses of the CM theorem show that its scalar generates $\mathfrak p$ and is $1\bmod\lambda^3$ (or both statements conjugated). Hence it is the primary generator just constructed. The model has discriminant $64$ and good reduction at every odd prime.

At split $p$ the inverse factors over $\mathfrak p,\bar{\mathfrak p}$ give the first line. At inert $p$, the prime ideal $(p)$ has norm $p^2$ and primary generator $-p$ since $p\equiv3\bmod4$. Thus its single factor is $(1+p^{1-2s})^{-1}$. Equivalently the sum of the quadratic characters of $x^3-x$ cancels in pairs $x,-x$ in $\F_p$. At $\lambda$ the character is ramified and the local Euler factor is one. By the conductor formula for induction of a character of an imaginary quadratic field, the conductor over $\Q$ is $|D_K|\mathrm N(\lambda^3)=4\cdot8=32$; this determines the wild factor at $2$ in the conductor as well. These values agree with \href{https://www.lmfdb.org/EllipticCurve/Q/32/a/3}{LMFDB $32.\mathrm a3$}.
\end{proof}
For example $5=(-1+2i)(-1-2i)$ and $13=(3+2i)(3-2i)$ with the displayed factors primary, giving $a_5=-2$ and $a_{13}=6$. At $3$ the inverse factor is $1+3T^2$. Over the ground field $K$ both Hecke factors occur; they give the same analytic function since conjugation permutes all ideals and preserves their norm. Thus $L(E_0/K,s)=L(E_0/\Q,s)^2$. Proving the elliptic hypothesis after adjoining the endomorphisms in this example is already equivalent to proving it over $\Q$.

\subsection{A positive integral for the entire completed function}
Let $q=e^{2\pi i z}$ and write $\eta_{\mathrm D}(z)=q^{1/24}\prod_{n\ge1}(1-q^n)$ for Dedekind's eta function. The two descriptions
\begin{equation}\label{cm:theta-eta}
 f(z)=\sum_{\substack{\alpha\in\Z[i]\\\alpha\equiv1\bmod\lambda^3}}
        \alpha q^{\alpha\bar\alpha}
     =\eta_{\mathrm D}(4z)^2\eta_{\mathrm D}(8z)^2
\end{equation}
give access to different properties of the same modular form. Here are the details identifying them. Hecke's theta-series theorem, applied to the primitive character above, makes the first expression a normalized weight-two cusp form of level $32$. Its nebentypus at odd $n$ is $\chi_{-4}(n)\Psi((n))/n=1$, since the primary generator of $(n)$ is $\chi_{-4}(n)n$. The classical eta-quotient criterion of Newman and Ligozat puts the last expression in the same space. For its exponents $r_4=r_8=2$ one has $\sum\delta r_\delta=\sum(32/\delta)r_\delta=24$ and $\prod\delta^{r_\delta}$ is a square. The cusp-order formula is a sum of positive terms at each cusp. Its first Fourier coefficient is one. Finally the index of $\Gamma_0(32)$ is $48$, it has eight cusps and no elliptic points, and the modular-curve genus formula gives $\dim S_2(\Gamma_0(32))=1$. This proves \eqref{cm:theta-eta}. For the modular-form inputs see T.~Miyake, \emph{Modular Forms}, the treatment of Hecke characters, and K.~Ono, \emph{The Web of Modularity}, Theorems 1.64 and 1.65. The normalized form and its expansion are also catalogued as \href{https://www.lmfdb.org/ModularForm/GL2/Q/holomorphic/32/2/a/a/}{$32.2.\mathrm a.\mathrm a$}.

Define the completion in motivic normalization by
\[
 \Lambda_0(s)=\left(\frac{\sqrt{32}}{2\pi}\right)^s
                      \Gamma(s)L(E_0/\Q,s).
\]
The eta product is strictly positive on the positive imaginary axis. The inversion law for Dedekind's eta gives $f(i/(32y))=32y^2 f(iy)$. Set $K(u)=e^u f(ie^u/\sqrt{32})$. It follows that $K$ is a positive even function and decays faster than every exponential at both ends of the real axis. Termwise Mellin integration in $\Re s>3/2$ and the change of variable $y=e^u/\sqrt{32}$ therefore give an entire continuation with
\begin{equation}\label{cm:cosh}
 \Lambda_0(1+z)=2\int_0^\infty K(u)\cosh(zu)\,du.
\end{equation}
In particular $\Lambda_0(1+z)=\Lambda_0(1-z)$ and $\Lambda_0(s)>0$ for every real $s$. All its even Taylor coefficients at the center are strictly positive, by differentiation under the integral. The real zeros of the uncompleted function are precisely the simple zeros at the nonpositive integers, as follows by dividing by $\Gamma(s)$.

\begin{theorem}[An explicit horizontal band for a CM curve]\label{cm:band}
For $E_0\colon y^2=x^3-x$, one has $\Re\Lambda_0(1+h+it)>0$ throughout $|h|\le1/2$, $|t|\le2$. The complete $L$-function of $E_0/\Q$, and hence that of $E_0/\Q(i)$, has no zero with $|\Im s|\le2$.
\end{theorem}
\begin{proof}
Write $v=e^u$, $\kappa=2\pi/\sqrt{32}$ and $k(u)=v e^{-\kappa v}$. The eta product gives, for $u\ge0$,
\[
 K(u)=k(u)\prod_{n\ge1}(1-e^{-4\kappa nv})^2(1-e^{-8\kappa nv})^2,
 \qquad 0.97\, k(u)\le K(u)\le k(u).
\]
Indeed $r=e^{-4\kappa}<1/80$, and the product is bounded below by $1-2r/(1-r)-2r^2/(1-r^2)>0.97$, by the elementary inequality $\prod(1-a_j)\ge1-\sum a_j$ for $0\le a_j\le1$. By \eqref{cm:cosh}, half the desired real part is the integral of $K(u)\cosh(hu)\cos(tu)$ on the positive half-line.

Set $b=\pi/4$. On $[0,b]$, $\cos(tu)\ge\cos(2u)\ge1-u/b$ by monotonicity and concavity. Since $\kappa>1$, $k$ is decreasing. The last value of $k$ on each of the four successive intervals of length $b/4$ is greater than, respectively, $0.30,0.27,0.23,0.18$. For explicit comparisons, $1.10<\kappa<1.12$ and the four values $e^{jb/4}$ lie in $(1.21,1.22)$, $(1.48,1.49)$, $(1.80,1.81)$ and $(2.19,2.20)$. In each pair $(a,d)$ use $a e^{-1.12d}$. Since $\cosh(hu)\ge1$, the contribution of $[0,b]$ exceeds
\[
  \frac{0.97\cdot0.78}{32}
       (7\cdot0.30+5\cdot0.27+3\cdot0.23+0.18)>0.102.
\]
The four integer weights come from integrating the affine lower bound for the cosine. For the rest of the integral use $\cos(tu)\ge-1$ and
\[
 \cosh(hu)\le \frac{\sqrt v+1/\sqrt v}{2}\le0.9+0.1v\qquad(v\ge1).
\]
The middle member equals one at $v=1$ and its derivative $(v-1)/(4v^{3/2})$ has maximum $1/(6\sqrt3)<0.1$. Thus the possible loss is at most
\[
 \int_{2.19}^{\infty} e^{-1.1v}(0.9+0.1v)\,dv
 =\frac{e^{-2.409}}{1.1}\left(0.9+0.219+\frac1{11}\right)<0.101.
\]
All decimals in these bounds are rationals. For instance each needed exponential comparison can be checked by bounding $e^x$ between its Taylor polynomial of degree $18$ and that polynomial plus $x^{19}/(19!(1-x/20))$ for $0\le x\le5$. The positive contribution exceeds the possible loss. The Euler product excludes zeros in $\Re s>3/2$, and the functional equation handles $\Re s<1/2$. The completed square identity over $\Q(i)$ holds up to a nonzero elementary normalization factor, so every zero there comes from $\Lambda_0$.
\end{proof}
An integral such as \eqref{cm:cosh} gives an actual sign at real arguments and, with quantitative estimates, at small heights. For an explicit limitation of the qualitative positivity argument, take $0<a<1/2$ and $b>0$ and the positive even kernel
\[
 H(u)=e^{-u^2}+a\{e^{-(u-b)^2}+e^{-(u+b)^2}\}.
\]
Its Fourier transform $\sqrt\pi e^{-t^2/4}(1+2a\cos(bt))$ has nonreal zeros at $bt=(2j+1)\pi\pm i\operatorname{arcosh}(1/(2a))$, $j\in\Z$. This example isolates the additional information that an argument by Fourier transforms would have to exploit in the particular arithmetic kernel $K$.

\section{What the split Gaussian primes must supply}\label{cm:criterion}
For each $p\equiv1\pmod4$ let $e^{i\theta_p}=\pi/\sqrt p$, choosing either primary prime over $p$. Then $a_p/\sqrt p=2\cos\theta_p$. Write $F(z)=L(E_0/\Q,z+1/2)$ and define
\[
 A(X)=\sum_{\substack{p\le X\\p\equiv1\bmod4}}2\cos\theta_p\,\log p.
\]
We can isolate the precise cancellation over ordinary rational primes that is needed in this example.
\begin{theorem}[An equivalent estimate on split primes]\label{cm:angles}
The elliptic Riemann hypothesis for $E_0/\Q$ is equivalent to the assertion that for every $\epsilon>0$ one has $A(X)=O_\epsilon(X^{1/2+\epsilon})$. If the hypothesis holds, $A(X)=O(\sqrt X\log^2(2X))$.
\end{theorem}
\begin{proof}
The complete list of finite Euler factors from Lemma~\ref{cm:primary} gives
\begin{align*}
 -\frac{F'}{F}(z)
 &=\sum_{p\equiv1\bmod4}\sum_{j\ge1}
       2\cos(j\theta_p)(\log p)p^{-jz}\\
 &\quad+\sum_{p\equiv3\bmod4}\sum_{j\ge1}
       2(-1)^j(\log p)p^{-2jz} \qquad(\Re z>1).
\end{align*}
Let $P(X)$ sum these coefficients up to $p^j\le X$ in the first row and $p^{2j}\le X$ in the second. The elementary Chebyshev bound $\sum_{p\le Y}\log p\ll Y$ gives $P(X)-A(X)=O(\sqrt X)$. Indeed the contribution of the prime squares is $O(\sqrt X)$ and that of exponents at least three is $O(X^{1/3}\log(2X))$. This argument includes the inert primes and deals with every prime power, including those based at small primes.

Suppose first the asserted estimate for $A$ holds. For $\Re z>1$, the Dirichlet series just obtained equals $z\int_1^\infty P(x)x^{-z-1}\,dx$. The integral converges locally uniformly on $\Re z>1/2$ by choosing $\epsilon$ smaller than the distance from the compact set to the boundary. It extends $-F'/F$ holomorphically there. Each zero of $F$ would give its logarithmic derivative a simple pole with nonzero residue, and $F$ is entire. By the functional equation every nontrivial zero therefore has real part $1/2$.

Conversely, apply the classical explicit formula for a fixed primitive modular-form $L$-function to $P(X)$. For clarity, the form used here is the truncated Perron explicit formula, with a height $T$ between $X$ and $2X$ away from zeros, evaluated at a half-integer $X$. It expresses $P(X)$ as $-\sum_{|\Im\rho|\le T}X^\rho/\rho+O(\log^C(2X))$ for some fixed $C$, with the sum over nontrivial zeros of $F$ and their multiplicities. The hypotheses hold since $F$ is entire of fixed degree and conductor, has its primitive functional equation, and the coefficients of $-F'/F$ satisfy $|\Lambda_F(n)|\le2\Lambda(n)$. See H.~Iwaniec and E.~Kowalski, \emph{Analytic Number Theory}, Chapter~5, in particular \S\S5.5--5.7. The same zero-counting formula gives $O(\log(3+|T|))$ zeros with imaginary part in $[T,T+1]$. If $\Re\rho=1/2$, summation in unit intervals yields $\sum_{|\Im\rho|\le2X}1/|\rho|\ll\log^2(2X)$. Consequently $P(X)\ll\sqrt X\log^2(2X)$; any real $X$ gives the same sum as a half-integer with the same integer part. The $O(\sqrt X)$ prime-power correction proves the result for $A$.
\end{proof}
In particular the congruence condition on inert primes admits an explicit treatment. The term $j=1$ from those primes in $P(X)$ is $-2\sum_{p\le\sqrt X,\ p\equiv3\bmod4}\log p$. Dirichlet GRH for the two characters modulo $4$, together with the Dirichlet explicit formula, evaluates it as $-\sqrt X+O(X^{1/4}\log^2(2X))$. The remaining inert powers are $O(X^{1/4}\log(2X))$. After that evaluation the first harmonic $A(X)$ still occurs with coefficient one. The elementary $O(\sqrt X)$ treatment already suffices for the equivalence.

\subsection{A fixed finite congruence cannot detect the angular trace}
The infinity type in Theorem~\ref{cm:factor} gives a character-theoretic obstruction. There are two more concrete versions for $E_0$, one using a single prime and another using the mean square over all split primes.
\begin{theorem}[Finite twists and periodic weights]\label{cm:finite}
The Euler product of $E_0/\Q$ over finite places outside a finite exceptional set cannot agree with a finite product and quotient of $L(s-b_j,\rho_j)$ for finite-image complex Artin representations $\rho_j$ of $G_{\Q}$ and real shifts $b_j$. Moreover, for any periodic $g\colon\Z\to\C$ and $D=\operatorname{lcm}(4,\operatorname{period}(g))$,
\begin{align*}
 &\lim_{X\to\infty}
 \frac{\sum_{p\le X,\ p\equiv1\bmod4}
        |a_p/\sqrt p-g(p)|^2\log p}
      {\sum_{p\le X,\ p\equiv1\bmod4}\log p}\\
 &\hspace{20mm}=2+
 \frac{1}{\#S_D}\sum_{d\in S_D}|g(d)|^2,
 \qquad S_D=\{d\in(\Z/D\Z)^\times:d\equiv1\bmod4\}.
\end{align*}
\end{theorem}
\begin{proof}
Choose any unramified $p\equiv1\pmod4$ outside the exceptional set for all the functions. The two reciprocal roots of its elliptic polynomial are $\pi,\bar\pi$. Every available reciprocal root from a shifted Artin factor has form $p^{b_j}\zeta$, with $\zeta$ a root of unity. Cancellation inside a rational function cannot create a new such root. Equality of the Euler products on a right half-plane also gives equality of the individual local factors, by uniqueness of the Dirichlet coefficients at $p^k$. The complex absolute value forces $b_j=1/2$. If $\pi=\sqrt p\,\zeta$, then $\pi/\bar\pi=\zeta^2$, a root of unity. The fractional ideal $(\pi)/(\bar\pi)$ is nontrivial since its numerator and denominator are distinct primes. This excludes $\pi/\bar\pi$ from the roots of unity, all of which are units.

For the mean square, fix $d\in S_D$. For $n=1,2$, consider $\eta^n(\mathfrak p)\chi(\mathrm N\mathfrak p)$ for each Dirichlet character $\chi\bmod D$. The corresponding Hecke character retains the same nonzero angular type as $\eta^n$. The boundary nonvanishing and prime ideal theorem argument in Section~\ref{cm:hecke} therefore gives a weighted prime sum $o(X)$ for each twist. Character orthogonality selects $\mathrm N\mathfrak p\equiv d\bmod D$. Each $p\equiv d\bmod D$ splits in $K$, and its two prime ideals contribute the conjugate pair of angular values. Inert ideals could also satisfy the norm condition, but their total is $O(\sqrt X)$ by the elementary Chebyshev bound. We conclude that
\[
 \sum_{\substack{p\le X\\p\equiv d\bmod D}}
   2\cos(n\theta_p)\log p=o(X)\qquad(n=1,2).
\]
The prime number theorem in arithmetic progressions gives $\sum_{p\le X,\ p\equiv d\bmod D}\log p\sim X/\varphi(D)$. The conditional first moment of $2\cos\theta_p$ on each class is zero. Its conditional second moment is two, by $4\cos^2\theta=2+2\cos(2\theta)$. Expand $|2\cos\theta_p-g(d)|^2$, sum over the finitely many residues, and divide by the total mass. The finitely many primes dividing $D$ have no effect on the limit.
\end{proof}
The same local-root proof excludes such Artin products over a number field for a CM elliptic curve if equality is specified at each prime of that base. Indeed, select a rational prime split completely in a normal closure containing $H$ and omit finitely many bad primes. Each chosen place of $H$ then has norm $p$ and its relative ideal norm to $K$ is a split prime $\mathfrak p$. Its Frobenius eigenvalue and its conjugate generate the distinct ideals $\mathfrak p,\bar{\mathfrak p}$. Since the completely split primes have positive density, one can make the choice outside every finite exceptional set. The argument at infinity in Theorem~\ref{cm:factor} also distinguishes angular characters from arbitrary complex powers of the absolute idele norm, all of which are constant on each local unit circle.

The second part of Theorem~\ref{cm:finite} quantifies the loss from a fixed residue modulus. At every fixed level the best periodic least-squares predictor of the normalized split trace is zero. It leaves mean-square error exactly two. The limit is taken with $D$ fixed; a method in which $D$ grows with $X$ would need a uniform estimate in both parameters. The mean-square statement makes no assumption about which new arguments could produce such an estimate. A direct proof of the single sum bound in Theorem~\ref{cm:angles} would already establish the full elliptic hypothesis for this CM curve, at every height.

\section{Critical-line zeros over number fields}\label{cm:proportion}
The Mellin argument treats small heights. A separate theorem about modular-form zeros gives unconditional information at arbitrarily large heights and survives in certain extensions of arbitrarily large degree. For an elliptic curve $C/M$ whose completed $L$-function is entire and has the expected degree, write $N_C(T)$ for the number of its nontrivial zeros with $0<\Im s\le T$, counted with multiplicity, and $N_C^0(T)$ for the subcount with $\Re s=1$. Set $\kappa_C=\liminf_{T\to\infty}N_C^0(T)/N_C(T)$.

\begin{theorem}[A proportion after base change]\label{cm:lift}
Let $E/\Q$ have complex multiplication and let $M/\Q$ be finite and either Galois or contained in a finite solvable Galois extension. Let $B$ be the compositum of the quadratic subfields of $M$, with $B=\Q$ if there are none. Then
\[
                 \kappa_{E/M}> \frac{0.0696}{[M:B]}.
\]
The same statement holds for every elliptic curve isogenous to $E_M$ over $M$. In particular the lower proportion is greater than $6.96\%$ for all multiquadratic $M$ and greater than $2.32\%$ for all cubic $M$.
\end{theorem}
\begin{proof}
Andersen and Thorner, \href{https://doi.org/10.1515/forum-2020-0104}{\emph{Zeros of $\mathrm{GL}_2$ $L$-functions on the critical line}}, Theorem~1.1, prove $\kappa_f>0.0696$ for every holomorphic newform of arbitrary level and trivial nebentypus, in unitary coordinates. Apply this to the weight-two newform of $E/\Q$ and to that of every quadratic twist $E^\chi$ of $E$. Twisting may alter the conductor. The rationality of the coefficients and the trivial nebentypus needed here are preserved. For the elementary abelian extension $B/\Q$ the regular character is the sum of its quadratic characters, including the trivial one. Restriction and induction on the Tate module give
\[
          L(E/B,s)=\prod_{\chi\in\operatorname{Gal}(B/\Q)^\vee}
                              L(E^\chi/\Q,s)
\]
with primitive local factors, including the ramified ones. Thus zero counts add, with repetitions counted when a CM self-twist occurs. In particular $\kappa_{E/B}>0.0696$, since the set of twists is finite and each lower proportion exceeds the stated constant.

The base-change theorem of M.~R.~Murty and V.~K.~Murty, \href{https://doi.org/10.1090/conm/143/01014}{\emph{Base change and the Birch--Swinnerton-Dyer conjecture}}, Theorem~1, says that for a CM elliptic curve over any number field $k$, its $L$-function over a finite extension $U/k$ is entire, and the quotient $L(E/U,s)/L(E/k,s)$ is entire when $U/k$ is Galois or has solvable normal closure. Apply it with $k=B$ and $U=M$. If $M/\Q$ is Galois, then $M/B$ is Galois. If $M$ lies in a solvable Galois extension, its normal closure over $B$ also has solvable group. Every nontrivial zero counted for $E/B$ is consequently counted with at least the same multiplicity for $E/M$. In the open critical strip the archimedean factors do not change those multiplicities.

For $U=B,M$ the zero-counting formula gives
\[
       N_{E/U}(T)=\frac{[U:\Q]}{\pi}T\log T+O_{E,U}(T).
\]
This is the usual argument-principle formula from the completed Hecke functions of Theorem~\ref{cm:factor}. Their combined degree over $\Q$ is $2[U:\Q]$ and there are no poles. Thus $N_{E/B}(T)/N_{E/M}(T)\to1/[M:B]$. Combining the lower counts gives the claimed strict bound. Isogeny over the ground field identifies all the rational Tate modules and hence all Euler factors. For cubic $M$ the Galois closure has group contained in $S_3$, which is solvable, and $B=\Q$.
\end{proof}
For multiquadratic $M$, the first paragraph of the proof gives $\kappa_{E/M}>0.0696$ for \emph{every} rational elliptic curve, by the modularity theorem over $\Q$. The degree-adjusted bound of Theorem~\ref{cm:lift} likewise holds for every rational elliptic curve when $M/\Q$ is solvable Galois. Indeed solvable automorphic base change makes the rational modular form automorphic over $B$, and Murty--Murty, Theorem~2 and its proof by successive cyclic steps, gives holomorphy of the quotient over $M/B$. Equivalently one can lift along the cyclic tower from $B$ directly, each step retaining the old cuspidal factors in the product and all the additional twists being entire in the open critical strip. For a non-CM elliptic curve these factors remain cuspidal in solvable base change: reducibility at any finite step would contradict absolute irreducibility of the restricted Tate representation, by the open-image theorem. The CM case has already been treated in full. Here the assumption ``solvable Galois'' is sufficient, and no assertion about the quotient for an arbitrary nonnormal extension of a non-CM curve is needed.

The base-change theorem also shows precisely how all-height zero exclusion travels through these towers. If the elliptic Riemann hypothesis holds over $M$, it holds over $B$; an off-line zero downstairs would persist upstairs. Proving the hypothesis downstairs constrains only one factor of the upstairs function. Over $\Q(i)$ for $E_0$, the explicit square identity supplies the reverse implication as well. More generally when a base $U$ contains all the CM endomorphisms, the completed factors for $\Psi$ and $\bar\Psi$ are complex conjugates on the real axis (choosing the standard positive conductor normalization). Their product is nonnegative there, and any real zero has even multiplicity. Strict positivity for every real $s$ in Section~\ref{cm:gaussian} used the specific eta product.

\subsection{A quartic example with nonrational moduli}
Consider the real algebraic numbers
\[
 r=\sqrt{1+2\sqrt3/3},\qquad u=r^2,\qquad M=\Q(r),\qquad B=\Q(\sqrt3).
\]
The element $u$ has a negative conjugate in $B$, so cannot be a square there. Thus $[M:B]=2$. The field $M$ has a nonreal embedding and a real embedding, so it cannot be biquadratic: a Galois field with a real embedding has only real embeddings. Its largest multiquadratic subfield is $B$. The normal closure of this tower of quadratic extensions is solvable, since it is obtained over $B$ by adjoining square roots of $u$ and its conjugate.

On $E_0$ the $3$-division polynomial is $3x^4-6x^2-1$. Since $r$ is a root, the points with $x=r$ are the two nonzero elements of a cyclic subgroup $C$ of order three. Galois over $M$ fixes their common $x$-coordinate, so the subgroup is defined over $M$. V\'elu's isogeny formulas give the quotient
\[
  D\colon y^2=x^3+(9-30r^2)x+(42r-70r^3).
\]
Indeed for $E:y^2=x^3+ax+b$ they subtract $5\sum_{Q\in C\setminus\{O\}}(3x_Q^2+a)$ from $a$ and $7\sum_{Q\in C\setminus\{O\}}(5x_Q^3+3ax_Q+2b)$ from $b$. Both terms in each sum coincide here. See J.~V\'elu, \href{https://aghitza.github.io/publications/translation-velu.pdf}{\emph{Isog\'enies entre courbes elliptiques}}, C.~R. Acad. Sci. Paris \textbf{273} (1971), 238--241.

To verify the moduli, set $a'=9-30u$, $b'=r(42-70u)$. Reduction using $3u^2-6u-1=0$ gives
\[
 j(D)=\frac{1728\cdot4a'^3}{4a'^3+27b'^2}
 =576(115444u+17839)=76771008+44330496\sqrt3.
\]
Thus no elliptic curve over $\Q$ is isomorphic over $\overline\Q$ to $D$. The specified $M$-isogeny gives equality of $L(D/M,s)$ with $L(E_0/M,s)$, and Theorem~\ref{cm:lift} proves $\kappa_{D/M}>0.0348$. This is a curve whose isomorphism class requires a number field even though its isogeny provides a route to modular forms over the rationals.

\subsection{Unbounded degrees of the moduli fields}
For each rational prime $\ell\equiv3\pmod4$, let $C$ be any cyclic subgroup of order $\ell$ in $E_0[\ell]$. Let $M_C$ be its field of definition as a subgroup, i.e., the fixed field of its stabilizer, and form $D_C=E_0/C$ over $M_C$. The orbit of $C$ is a subset of the $\ell+1$ lines in $\F_\ell^2$, hence $[M_C:\Q]\le\ell+1$. Its normal closure is solvable. In fact the mod-$\ell$ action over $\Q(i)$ commutes with the field $\Z[i]/\ell\simeq\F_{\ell^2}$ on $E_0[\ell]$, so its image lies in the abelian group $\F_{\ell^2}^\times$. The image over $\Q(i)$ is thus an abelian normal subgroup of index at most two in the image over $\Q$. Any quotient giving the Galois closure of $M_C$ is therefore solvable.

The endomorphism order of $D_C$ over $\C$ is $\mathcal O_\ell=\Z+\ell\Z[i]$. Indeed, in a complex uniformization of $E_0$ with lattice $\mathcal O=\Z[i]$, the quotient replaces $\mathcal O$ by an overlattice $\mathcal L$ of index $\ell$ with $\ell\mathcal L\subset\mathcal O$. Multiplication by $\ell\mathcal O$ preserves $\mathcal L$. The endomorphism order lies between $\mathcal O_\ell$ and $\mathcal O$, whose index is prime, and cannot be $\mathcal O$: in that case the line $C=\mathcal L/\mathcal O$ would be $\F_{\ell^2}$-stable inside the one-dimensional $\F_{\ell^2}$-space $E_0[\ell]$. The classical ring class polynomial and class number formula now give
\[
 [\Q(j(D_C)):\Q]=h(\mathcal O_\ell)
 =\frac{h(\mathcal O)\ell}{[\mathcal O^\times:\mathcal O_\ell^\times]}
          \left(1-\frac{(-4/\ell)}{\ell}\right)=\frac{\ell+1}{2}.
\]
Here $(-4/\ell)=-1$ and the unit index is two. These are the standard CM degree and class-number formulas, for example in D.~A.~Cox, \emph{Primes of the Form $x^2+ny^2$}, Chapters~II--III. In particular both $[\Q(j(D_C)):\Q]$ and $[M_C:\Q]$ are unbounded. The preceding theorem and the ground-field isogeny give the uniform-in-form statement $\kappa_{D_C/M_C}>0.0696/(\ell+1)$ for every such $\ell$. This family proves a positive lower proportion for each individual curve. Its lower bound is allowed to decrease as the fields of moduli grow.

\subsection{Uniform improvement in anticyclotomic towers}
For the preceding family the Galois structure retains the full rational lower proportion independently of the moduli degree.
\begin{theorem}\label{cm:anticyclotomic}
Suppose $E/\Q$ has CM by an imaginary quadratic field $K$. Let $A/\Q$ be a finite Galois extension containing $K$ such that $J=\operatorname{Gal}(A/K)$ is abelian and complex conjugation acts on $J$ by inversion. For every intermediate field $\Q\subseteq M\subseteq A$ one has $\kappa_{E/M}>0.0696$. In particular the quartic curve $D$ above and every member $D_C/M_C$ of the unbounded-degree family satisfy this same strict lower bound.
\end{theorem}
\begin{proof}
Write $V=\operatorname{Ind}_{G_K}^{G_\Q}\Psi_\ell$ for the rational Tate representation after scalar extension, with $\Psi_\ell$ the realization of the character of Theorem~\ref{cm:factor}; the auxiliary Tate prime is immaterial. Put $P=\operatorname{Ind}_{G_M}^{G_\Q}1$, the permutation representation on the embeddings of $M$. Its restriction to $G_K$ factors through $J$, so is a sum of one-dimensional characters $\nu$ with nonnegative integer multiplicities $m_\nu$ totaling $[M:\Q]$. The tensor identity for induction gives
\[
 V\otimes P=\bigoplus_\nu m_\nu\operatorname{Ind}_{G_K}^{G_\Q}(\Psi_\ell\nu),
 \qquad L(E/M,s)=\prod_\nu L_K(s,\Psi\nu)^{m_\nu}.
\]
The left identity is over an algebraically closed coefficient field; the right identity includes ramified local factors by local induction.

Each finite character $\nu$, viewed by class field theory as a character of $K$, satisfies $\nu((p))=1$ for all rational primes $p$ unramified in $A$. If $p$ splits in $K$, its two Frobenius elements in $J$ are conjugate under complex conjugation and hence inverses. If $p$ is inert, the Frobenius over $K$ is the square of a member of the other coset in $\operatorname{Gal}(A/\Q)=J\rtimes\langle c\rangle$, and $(jc)^2=1$. Thus the nebentypus of the primitive theta newform of $\Psi\nu$ agrees with that of $\Psi$ at every unramified rational prime. This is the trivial Dirichlet character, since $E/\Q$ has trivial nebentypus and a Dirichlet character is determined by its values outside a finite set of primes. The angular type survives twisting and makes the theta newform cuspidal of weight two. Every factor therefore satisfies the full-level-independent Andersen--Thorner lower bound. The finite product has the same strict lower bound, with zeros counted with multiplicity.

For $D_C$, take $A$ to be the compositum of $K=\Q(i)$ and the splitting field of the action on the lines of $E_0[\ell]$. The field $M_C$ is contained in it. Over $K$ the action on the lines factors through $\F_{\ell^2}^{\times}/\F_\ell^{\times}$. Complex conjugation acts on this group by the power $\ell$, which coincides with inversion since $x^{\ell+1}$ is scalar. The natural map from $\operatorname{Gal}(A/K)$ to this projective image is injective by the definition of $A$. Thus the hypotheses hold, also for the quartic example $\ell=3$. The ground-field isogeny transfers the claim to $D_C$.
\end{proof}

\clearpage
\section{All-height bounds claimed for CM and unitary Hecke functions in 2026}\label{recentcm:section}

There are now separate results to assess for complex multiplication and for elliptic curves without complex multiplication. All exponents in the following table concern zero-freeness in the strict half-plane $\Re(s)>\theta$ in \emph{unitary} normalization. The dates identify public versions available by 23 November 2026.
\begin{center}\small
\begin{tabular}{p{.25\linewidth}p{.31\linewidth}p{.31\linewidth}}
\toprule
Work & Family and boundary & Verification reported by source\\ \midrule
Rubinstein--Restrepo Borrero, 13 October &
Angular type $-1$ and finite-order twists over $\Q(\sqrt{-3})$, $\theta=7/8$ &
Written argument and finite checks. The deposit describes Lean work as in progress.\\
Hu, version 2 &
All continuous unitary Hecke characters over all number fields, $\theta=7/8$ &
Public Lean solution to a Mathlib-only challenge. The repository reports comparator verification of the rank-one endpoint.\\
Gao, November &
Same unitary family, $\theta=(777-\sqrt{1081})/852$ &
Uses specified analytic inputs of Hu and Liu. The deposit distinguishes scalar checks from formalization of the full endpoint.\\
Liu, 20 November &
Same unitary family, $\theta=\theta_G=0.873332637445905\ldots$ &
Theorem 7.33 of a preprint. Its provenance says that a Lean formalization is in progress.\\
\bottomrule
\end{tabular}
\end{center}

The sources are the deposits of \href{https://hexagonmath.org/2610.00220v1}{Rubinstein--Restrepo Borrero}, \href{https://hexagonmath.org/2610.00216v2}{Hu}, \href{https://hexagonmath.org/2611.00334v1}{Gao} and \href{https://hexagonmath.org/2611.00741v1}{Liu}. They state these endpoints as theorems. We distinguish a consequence proved here from a recent zero-free theorem used as an input. In particular a numerical check of an optimization is a check of that optimization alone.

In the Eisenstein case write $\alpha((a))=a/|a|$ on the chosen primary generators away from the bad primes. The October paper treats $L(s,\bar\alpha\eta)$ for finite-order $\eta$. Complex conjugation reflects the corresponding character and the imaginary part. Its first part obtains the weaker elliptic arithmetic boundary $17/12$ and its second part obtains $11/8$. Its dependency table explicitly leaves higher weights and general symmetric powers as extensions to be proved. In its second argument the quoted finite-order $7/8$ result is accompanied by new CM reflection identities, inverse and ordinary moments, and an endpoint inequality. Finite tests of the reflection factors alone do not give its uniform moments. The October claim for the family $y^2=x^3+d$ has the exact normalization
\[
   L_u(s)=L(E_d,s+1/2),\qquad \Re(s)>7/8
       \quad\Longrightarrow\quad L(E_d,w)\ne0\quad(\Re w>11/8).
\]
The constants in a proof for each fixed curve may depend on $d$. For the full admissible Eisenstein angular family, a separate \href{https://hexagonmath.org/2610.00716v1}{2 November deposit by T.~Winarski} states the $7/8$ result for all angular orders, finite conductors and real norm twists. Its provenance reports a clean build and native kernel replay after correction of the public archive. Our detailed formal-source check below concerns Hu's independent arbitrary-field result.

\subsection{The scope of the formal and numerical checks}\label{recentcm:audit}
We read both source files at commit \texttt{c77efa9237c4}. The challenge imports Mathlib. Its character is a continuous homomorphism from the idèle class group of a number field to the unit circle. At an unramified prime it uses the value on a local uniformizer, and at a ramified prime it assigns zero as the ideal coefficient. Multiplicativity gives the series on integral ideals. The conclusion exhibits a function agreeing with that series for $\Re(s)>1$, holomorphic except at a subsingleton set on $\Re(s)=1$, and nonzero away from that set for $\Re(s)>7/8$. At any exceptional point it requires a simple pole with nonvanishing numerator. Thus its formulation cannot omit a zero by arbitrarily designating that zero an exceptional point. Uniqueness of meromorphic continuation identifies the exhibited function with the classical Hecke function. This quantification includes characters of infinite order.

The solution identifies the challenge's local and ideal coefficients with those in the internal theorem. Its comparison configuration lists exactly propositional extensionality, classical choice and quotient soundness. We also retrieved the archived integration job. Its final log records that the default Lean kernel accepted the solution, and its axiom output lists precisely these three axioms for the target. We compared the exported types and bodies of the four definitions for characters, ramification and coefficients. All differences are bound-variable display names. The CI checkout and the cited release have the same Git source tree. The challenge deliberately leaves the classification of poles by norm characters and the higher-rank global Weil theorem to separate arguments. These facts can be checked in the \href{https://github.com/Rogerhu12/hecke-seven-eighths/blob/c77efa9237c4635ed61933fbca98722c359c984a/src/Challenge/HeckeSevenEighthsChallenge.lean}{challenge}, the \href{https://github.com/Rogerhu12/hecke-seven-eighths/blob/c77efa9237c4635ed61933fbca98722c359c984a/src/Solution/HeckeSevenEighthsSolution.lean}{solution}, and the \href{https://github.com/Rogerhu12/hecke-seven-eighths/actions/runs/42421770943/job/128551523125}{integration job of 8 November}. The blank proof in the reference challenge is the target supplied to the comparator. The separate solution contains the proof.

Liu's precise reported enclosure is
\[
 0.8733326374459049352<\theta_G<0.8733326374459051924 .
\]
The named input is Theorem~7.33 of \href{https://hexagonmath.org/pdf/2611.00741v1}{\emph{Mixed character moments and a zero-free half-plane for unitary Hecke $L$-functions}}. Its Sections~7.3--7.6 treat a second Poisson summation and an induction on lengths in a mixed moment. Section~7.9 supplies Mellin continuation, and Section~7.18 compares the resulting exponent to zero at the stated endpoint. The contact calculation is downstream of those moment estimates. The author's public version presents a proof of the moments, together with inputs from earlier papers and companion materials. Its deposited source archive contains the single manuscript source. The Python scripts named there require a separate source. We do not have a separately checked proof of the whole chain at $\theta_G$. The rank-one Lean endpoint just described is $7/8$. The authors' stronger recent bounds and Hu's formal endpoint therefore have different verification records.

\subsection{A Mellin estimate that can be checked directly}\label{recentcm:cubes}
We give a check on an analytic step in Liu's Section~4. It identifies precisely where a prior zero-free bound for a cubed character enters a centered moment estimate. For readability we treat fixed excluded primes and state the dependence on the four lengths, which isolates the zero-sensitive step.

Fix a number field $F$ and two unitary ideal characters $\psi_1,\psi_2$. For a pure norm character $N^{it}$, first divide the sum by its scale phase $X^{it}$ (respectively $Y^{it}$), and absorb $u^{it}$ into the corresponding smooth compactly supported test $W_i(u)$. The difference below compares these dephased sums, so that the remaining pure norm characters are principal. With a fixed set of excluded primes containing the conductors, define
\[
 S(X,Y)=\sum_{\mathfrak a,\mathfrak b}^{*}
      \psi_1(\mathfrak a)\psi_2(\mathfrak b)
      W_1(N\mathfrak a/X)W_2(N\mathfrak b/Y).
\]
The star requires the product $\mathfrak a\mathfrak b$ to be cube-free and prime to the excluded primes. Suppose $XY=X'Y'=T$ and $L=\min(X,Y,X',Y')\ge1$. When $\psi_i=1$, assume a zero-free half-plane $\Re(s)>B$ for the cubic power of the opposite character, with $1/2\le B\le1$. For fixed data and every $0<\delta<1-B/3$,
\begin{equation}\label{recentcm:centered}
                   S(X,Y)-S(X',Y')=O(T L^{-\delta}).
\end{equation}
If neither character is principal, every $\delta<4/3$ is admissible without a zero-free input. The estimate is for fixed characters and test functions.

\begin{proof}
For $x=\psi_1(\mathfrak p)(N\mathfrak p)^{-s}$ and $y=\psi_2(\mathfrak p)(N\mathfrak p)^{-w}$, the good-prime generating factor is
$1+x+y+x^2+xy+y^2$. On removing the two ordinary Hecke Euler factors, the remainder is
\[
 \begin{aligned}
 H(x,y)&=(1-x)(1-y)(1+x+y+x^2+xy+y^2)\\
 &=1-x^3-x^2y-xy^2-y^3+x^3y+x^2y^2+xy^3 .
 \end{aligned}
\]
Thus $\prod H$ converges normally whenever the two real parts are at least a fixed $a>1/3$. Mellin-invert the tests, initially to the right of $1$, and shift both lines to $a\in(1/3,1/2)$. On these lines the remaining Hecke factors have at most polynomial vertical growth and the tests decay arbitrarily rapidly. The integral over the two shifted lines is $O(T^a)$. The double residue, if present, is a constant times $T\widehat W_1(1)\widehat W_2(1)$, so cancels between the two choices of lengths. Since $T\ge L^2$, the two-line integral is $O(T L^{-2(1-a)})$. This proves the second assertion by taking $a$ sufficiently close to $1/3$.

For a single residue suppose $\psi_1=1$, so $s=1$. Set $u=(N\mathfrak p)^{-1}$, $v=\psi_2(\mathfrak p)(N\mathfrak p)^{-w}$. On $\Re w\ge\sigma>0$ direct expansion gives
\[
     H(u,v)/(1-v^3)=1+
       O_\sigma\big((N\mathfrak p)^{-1-2\sigma}
               +(N\mathfrak p)^{-2-\sigma}+(N\mathfrak p)^{-3}\big).
\]
The resulting correction product converges normally there. It also has no zero: apart from $1-u$ and $1-v$, the original polynomial contains the quadratic
$v^2+(1+u)v+1+u+u^2$. Its two roots have modulus $\sqrt{1+u+u^2}>1$, whereas $|v|<1$. The remaining meromorphic factor of the residue is, up to finitely many harmless corrections,
\[
                  L_F(w,\psi_2)\,/\,L_F(3w,\psi_2^3).
\]
The zero-free assumption permits the residue line to move from $a$ to a fixed $\sigma$ satisfying $B/3<\sigma<\min(a,1-\delta)$. On this buffered line the inverse denominator has polynomial growth. For completeness, at large height take a disk centered at $2+it$, with left edge strictly between $B$ and $3\sigma$ when a shift left of $2$ is needed. A logarithm of $L_F$ exists there. The real part of this logarithm is bounded above by $O(\log(2+|t|))$ by polynomial vertical growth. Its value at the center is bounded by the convergent Euler-product logarithm. Borel--Carathéodory on a strictly smaller disk gives $|\log L_F|\ll\log(2+|t|)$ on the desired segment, hence a polynomial bound for $1/L_F$. At bounded heights the reciprocal extends across a possible principal pole. The residue is now $O(XY^\sigma)$, which is $O(TL^{-\delta})$. The exchanged single residue is identical with the two coordinates reversed. This proves~\eqref{recentcm:centered}.
\end{proof}

The factor $1/3$ in the zero-sensitive line shift comes from the argument $3w$ in the denominator. Even an exact calculation of every Euler polynomial leaves the separate mixed-moment and zero-detection arguments to be supplied. This is the scope of the preceding check on the recent preprint.

\subsection{CM curves, their powers and base changes}\label{recentcm:transfer}
Write $\mathcal H(\theta)$ for the zero-free assertion about all unitary Hecke characters of all number fields in $\Re(s)>\theta$, permitting a norm-character pole. We now give exact implications of this named hypothesis, with the reported choices $\theta=7/8$ and $\theta=\theta_G$ kept separate.

\begin{theorem}\label{recentcm:main}
Let $E/B$ have complex multiplication and let $F/B$ be any finite extension. Under $\mathcal H(\theta)$ every nontrivial zero of $L(E/F,s)$ belongs to
\[
                3/2-\theta\ \le\ \Re(s)\ \le\ 1/2+\theta.
\]
For $\operatorname{Sym}^h E$, $h\ge1$, the arithmetic strip is
\[
          1+h/2-\theta\ \le\ \Re(s)\ \le\ h/2+\theta.
\]
These statements also permit a finite-order Hecke twist at the base. In particular the elliptic arithmetic strip from the $7/8$ input is $[5/8,11/8]$. The stronger preprint input gives the interval with endpoints approximately $0.626667362554095$ and $1.373332637445905$.
\end{theorem}
\begin{proof}
Let $M$ be the quadratic CM field. Over $FM$ the unitary parameter is a sum of the Hecke characters $\eta,\eta^{-1}$ furnished by the main theorem of complex multiplication. Here a Frobenius value of its arithmetic weight-one character is divided by the positive square root of the norm. Over $F$ when $M\not\subset F$, this two-dimensional parameter is induced from $\eta$ over $FM$. These are the factorizations proved in Theorem~\ref{cm:factor}. Thus every standard factor is either a Hecke function over $F$ or the Hecke function over $FM$ identified by induction. The construction applies after arbitrary finite base extension, since the CM endomorphisms remain available over the compositum.

On the split torus the monomials in the $h$th symmetric power have characters $\eta^{h-2j}$ for $0\le j\le h$. When the quadratic extension is nontrivial, the nonmiddle terms are paired by the outside coset and form inductions of those characters. For $h=2a$ the middle monomial gives a single finite character $\epsilon_{FM/F}^{\,a}$: the normalized elliptic determinant is one, and the product of the two off-diagonal matrix entries is minus the determinant. When $M\subset F$ this character is trivial. Twisting over $F$ simply restricts the twisting character to the relevant inducing fields. This is a positive product of Hecke functions, with the usual local factors by induction at ramified places as well.

Each constituent's nontrivial unitary zeros have $\Re(s)\le\theta$, and the functional equation applied also to the conjugate gives $\Re(s)\ge1-\theta$. A possible pole of a middle norm character sits on the exterior boundary of the ordinary critical strip. Restoring arithmetic weight shifts the real parts by $h/2$. For $h=1$ this gives the displayed elliptic strip. In particular the single $\bar\alpha\eta$ family of the October curve paper can be used for its stated $j=0$ curves, while $\mathcal H(\theta)$ supplies all angular powers needed here.
\end{proof}

For a fixed CM curve $E/\Q$ the same input gives a power saving at the primes. Write $\lambda_p=a_p(E)/\sqrt p$ at good primes. The standard explicit formula for the completed elliptic function in unitary normalization gives
\[
              \sum_{\substack{p\le x\\p\nmid N_E}}\lambda_p\log p
                     =O_{E,\varepsilon}(x^{\theta+\varepsilon}).
\]
Indeed choose a truncation height between $x$ and $x+1$ away from zero ordinates, and use the unit-interval zero count $O_E(\log(2+|t|))$. The zero terms of the truncated logarithmic-derivative formula sum in absolute value to $O_E(x^\theta\log^2(2x))$, since the standard elliptic function is entire with no pole at $1$ in unitary normalization. One may first take a half-integer $x$: the Perron truncation error at that height is $O_E(\log^2(2x))$ by the uniform bound on normalized prime-power coefficients. Passing to general $x$ costs $O_E(\log(2x))$. Higher prime powers contribute $O(\sqrt x\log^2(2x))$ by the local absolute bound $2$. The finitely many bad primes contribute $O_E(\log x)$. The same proof for any fixed nontrivial angular power over the CM field gives its corresponding prime-ideal sum with constants allowed to depend on the power. For the formal $7/8$ endpoint this is an $O_{E,\varepsilon}(x^{7/8+\varepsilon})$ estimate for the normalized trace. It is weaker than the square-root estimate equivalent to elliptic GRH.

In fact all angular orders together give a uniform interval estimate. For a fixed CM curve $E/\Q$ let $M$ be its CM field and $J$ any interval in $[-2,2]$. Write $m$ for probability arc length on the unit circle and set
\[
 \mu_{\rm CM}(J)=\tfrac12\,m\{z:z+\bar z\in J\}
                             +\tfrac12\,\mathbf1_{0\in J}.
\]
For $1/2\le\theta<1$, the same $\mathcal H(\theta)$ gives
\begin{equation}\label{recentcm:st}
  \sum_{\substack{p\le x,\ p\nmid N_E\\a_p/\sqrt p\in J}}
      \log p
     = x\mu_{\rm CM}(J)+O_{E,\theta}\!\left(x^\theta\log^3(2x)\right),
\end{equation}
with the same constant for every interval $J$.

To prove this, let $\eta$ be the unitary CM character on $M$. For $1\le j\le x$ the finite conductor of $\eta^j$ divides that of $\eta$, and its archimedean weight grows linearly with $j$. Thus the usual unit-interval zero count up to height $T$ is $O_E(\log(2+j+T))$ there, uniformly in $j$. This follows from the degree-two Hecke gamma factor and conductor in the argument-principle zero count. See H.~Iwaniec and E.~Kowalski, \emph{Analytic Number Theory}, Chapter~5, pp.~101--112, for the zero count and explicit formula. The preceding explicit-formula argument therefore bounds
\[
     \left|\sum_{N\mathfrak p\le x}\eta(\mathfrak p)^j\log N\mathfrak p\right|
                \ll_{E,\theta} x^\theta\log^2(2x)\qquad(1\le j\le x)
\]
away from the fixed bad primes. The lower edge $1-\theta>0$ also bounds $1/|\rho|$ for low nontrivial zeros, uniformly in $j$. The zeroth power has sum $x+O_{E,\theta}(x^\theta\log^2(2x))$ by the same formula with the principal pole. Apply the Erd\H os--Tur\'an inequality with the positive weights $\log N\mathfrak p$ and $H=\lfloor x^{1-\theta}\rfloor$. Its truncation term is $O(x/H)$. The sum of the other terms is bounded by the displayed bound times $\sum_{j\le H}j^{-1}$. This proves a uniform $O_{E,\theta}(x^\theta\log^3(2x))$ error for arcs. The preimage of $J$ under $z\mapsto z+\bar z$ is a union of at most two arcs.

Each good split rational prime contributes two conjugate degree-one prime ideals. The inert ideals counted with norm at most $x$ have underlying rational prime $p\le\sqrt x$ and contribute $O(\sqrt x\log x)$. Hence the split rational-prime contribution is half the indicated angular main term, with the same error. At every good inert prime the elliptic trace is zero. The quadratic Hecke function for $M/\Q$ and the Riemann zeta function, both covered by the input, give their total weight $x/2+O_{E,\theta}(x^\theta\log^2(2x))$. Adding it exactly when $0\in J$ gives~\eqref{recentcm:st}. The atom at zero is retained also when zero is an endpoint of $J$.

For the real entire heat transform of a self-dual rational CM newform considered in Theorem~\ref{heat:real}, use the physical variable $s=(k/2)+iz$, where $k$ is its classical weight. The unitary strip just obtained gives $|\Im z|\le\theta-1/2$ at time zero. The strip-contraction proof of that theorem gives real zeros at every
\[
                 t\ \ge\ \frac{(\theta-1/2)^2}{2}.
\]
In particular Hu's $7/8$ input gives $t\ge9/128$. The hypotheses on the real transform in the cited theorem are part of this inference. The conclusion concerns positive heat time and supplies no all-real statement at time zero.

\subsection{Products of CM curves and a cuspidal degree-four example}\label{recentcm:tensors}
Let $E_i/F$ be elliptic curves with possibly different CM fields and choose $h_i\ge1$. Set $w=\sum_i h_i$. The arithmetic tensor product of their symmetric powers has weight $w$. Its standard function has all nontrivial zeros in
\[
              1+w/2-\theta\ \le\ \Re(s)\ \le\ w/2+\theta
\]
under $\mathcal H(\theta)$. It continues meromorphically, with no poles in the open critical strip. These assertions persist after finite extension of $F$.

Indeed, the proof of Theorem~\ref{recentcm:main} decomposes each unitary symmetric power as a sum of representations induced by characters from index-one or index-two subgroups of $W_F$. For open finite-index subgroups $U,V$ and unitary characters $\alpha,\beta$ on them, the tensor identity and Mackey decomposition give
\[
 \operatorname{Ind}_{U}^{W_F}\alpha\otimes
 \operatorname{Ind}_{V}^{W_F}\beta
 \simeq
 \bigoplus_{g\in U\backslash W_F/V}
 \operatorname{Ind}_{U\cap gVg^{-1}}^{W_F}
      \bigl(\alpha|_{U\cap gVg^{-1}}\,\beta^g|_{U\cap gVg^{-1}}\bigr).
\]
The intersection corresponds to a compositum of conjugate finite field extensions. Iteration gives a positive product of rank-one Hecke functions over number fields, including all local factors by induction. Their possible poles have real part one in unitary normalization. Applying the functional equation on the left and restoring weight gives the asserted strip. This proof gives meromorphic continuation without negative multiplicities.

For a concrete irreducible product take
\[
    E_1:y^2=x^3-x,\qquad E_2:y^2=x^3+1
\]
over $\Q$. Their quadratic CM fields are $M_1=\Q(i)$ and $M_2=\Q(\sqrt{-3})$. Write $D=M_1M_2$ and let $\eta_i$ be the first unitary CM character of $M_i$. There is just one double coset in the formula above, since the two distinct index-two subgroups generate $W_\Q$. Thus the unitary degree-four tensor is induced from the single Hecke character
\[
       \omega=(\eta_1\circ N_{D/M_1})(\eta_2\circ N_{D/M_2}).
\]
In particular $L(E_1\otimes E_2,s)=L_D(s-1,\omega)$, with the arithmetic normalization on the left. The character has four distinct Galois conjugates. To see this, choose the two complex embeddings of $D$, up to complex conjugation, to agree on $M_1$ and to be opposite on $M_2$. After independent choices of the two inducing characters, the angular exponents on the two unit circles are $(2,0)$. The independent conjugations on $M_1,M_2$ give the four distinct pairs $(2,0),(0,2),(0,-2),(-2,0)$. In particular none of these characters is a norm power. This establishes entireness of the completed degree-four function by the classical Hecke theorem.

This induction is cuspidal on $\mathrm{GL}_4(\mathbb A_\Q)$. Apply cyclic automorphic induction first from $D$ to $M_1$ and then from $M_1$ to $\Q$, using the cyclic induction and cuspidality criterion of Arthur--Clozel, \emph{Simple Algebras, Base Change, and the Advanced Theory of the Trace Formula}, \href{https://www.claymath.org/library/cw/arthur/pdf/30.pdf}{Chapter~III, Theorem~6.2}. The first conjugate pair is distinct. At the second step the two induced pairs have disjoint sets of constituents on restriction to $D$, hence the two cuspidal representations on $M_1$ are nonisomorphic. Both cyclic inductions are therefore cuspidal. Its local parameters agree with the elliptic tensor by construction.

The arithmetic Euler polynomial supplies a check on the shift by one. Writing $a=a_p(E_1)$, $b=a_p(E_2)$, multiply the four factors from the pairs of elliptic eigenvalues. The result at $p\nmid6$ is
\[
 1-abT+\{p(a^2+b^2)-2p^2\}T^2-p^2abT^3+p^4T^4 .
\]
For $p=11$, inert in both fields, both elliptic traces vanish and this is $(1-121T^2)^2$. For $p=13$, split in both fields, the traces are $6$ and $2$ and the polynomial is $1-12T+182T^2-2028T^3+28561T^4$. These traces follow by summing the quadratic character of the two displayed cubics over the prime field.

For this example the $7/8$ input puts \emph{all} zeros of the completed function in
\[
                         9/8\ \le\ \Re(s)\ \le\ 15/8.
\]
The geometric center is $3/2$. The same explicit formula as above gives the normalized prime correlation
\[
   \sum_{\substack{p\le x\\p\nmid 6}}
       \frac{a_p(E_1)a_p(E_2)}p\log p
               = O_\varepsilon(x^{\theta+\varepsilon})
\]
for either universal Hecke input. The prime trace is the product of the two normalized elliptic traces, and every normalized eigenvalue has absolute value one at a good prime. The degree-four function has no pole, so no main term is present. This yields an all-height fixed zero-free half-plane for a specific cuspidal degree-four automorphic function over a number field.

\subsection{An all-height cuspidal family in every rank}\label{recentcm:allranks}

\begin{theorem}\label{recentcm:manyranks}
For every number field $F$ and each $n\ge2$ there exist infinitely many pairwise nonisomorphic unitary cuspidal representations $\pi$ of $\mathrm{GL}_n(\mathbb A_F)$ such that $\mathcal H(\theta)$ puts every zero of their completed standard functions in
\[
                         1-\theta\ \le\ \Re(s)\ \le\ \theta.
\]
One can take inductions of infinite-order Hecke characters unramified at all finite places of a fixed cyclic extension $K/F$ of degree $n$ which splits all infinite places of $F$. In particular the $7/8$ input gives $[1/8,7/8]$ in every rank.
\end{theorem}

\begin{proof}
Choose a prime $q\equiv1\pmod{2n}$ unramified in the normal closure of $F/\Q$, by Dirichlet's theorem. The maximal real subfield of $\Q(\zeta_q)$ has a cyclic degree-$n$ subfield $D/\Q$. Every nontrivial subfield of $D$ is ramified at $q$ since $q$ is totally ramified in the cyclotomic field. Thus $D$ and the normal closure of $F$ are linearly disjoint. For $K=FD$ the extension $K/F$ is cyclic of degree $n$ and every infinite place of $F$ splits, since $D$ is totally real. Write $G=\operatorname{Gal}(K/F)$ and $A$ for the set of archimedean places of $K$. Use the standard absolute values normalized to count a complex place twice in the product formula. By the Dirichlet unit theorem the vectors
\[
       (\log|u|_v)_{v\in A},\qquad u\in\mathcal O_K^\times ,
\]
form a full lattice $\Gamma$ in $H=\{x\in\R^A:\sum_v x_v=0\}$. Let
$\Gamma^*=\{t\in H:\langle t,\Gamma\rangle\subset 2\pi\Z\}$. It is a $G$-stable full lattice in $H$. The action on the places above each infinite place of $F$ is regular, so the action on $H$ is faithful. For every $1\ne g\in G$ its fixed subspace in $H$ is proper. A finite union of proper linear subspaces cannot cover a full lattice: in the box of side $2R$ in any lattice basis there are $\gg R^{\dim H}$ lattice points, with $O(R^{\dim H-1})$ in any one proper subspace. Choose $t\in\Gamma^*$ having trivial stabilizer.

The rule
\[
       (a)\longmapsto\exp\!\left(i\sum_{v\in A}
                                     t_v\log|a|_v\right)
\]
defines a unitary character on the principal fractional ideals of $K$, since replacing $a$ by $ua$ changes the angle by a multiple of $2\pi$. It extends to the full ideal group. For an explicit extension, lift cyclic generators of the finite ideal class group and choose, for a generator of order $d$, any $d$th root of the prescribed value on the principal $d$th power of its lift. Combining this extended ideal character with the inverse archimedean character gives a continuous unitary idèle class character $\chi_t$ unramified at every finite place. Its distinct Galois conjugates already have distinct infinity types by the choice of $t$. Its type is nonconstant and thus $\chi_t$ is no norm power.

Take any tower from $K$ to $F$ with cyclic steps of prime degree. Automorphically induce $\chi_t$ one step at a time. At each step the preceding induced representation and each of its other conjugates restrict by base change to disjoint sums of conjugates of $\chi_t$ over $K$. Indeed those sums are indexed by different cosets of the corresponding subgroup of $G$, and the $G$-orbit of $\chi_t$ is free. The cyclic cuspidality criterion used in Section~\ref{recentcm:tensors} shows inductively that every step is cuspidal. The final unitary representation $\pi_t$ has completed standard function equal to the completed Hecke function of $\chi_t$, by induction of local factors. This function is entire. Its functional equation and $\mathcal H(\theta)$ applied to both $\chi_t$ and its inverse give the displayed strip.

Finally, for every positive integer $m$ the parameter $mt$ has the same stabilizer. The Euclidean lengths $\|mt\|$ distinguish the resulting $G$-orbits. Their cuspidal inductions have different archimedean parameter multisets and hence are pairwise nonisomorphic.
\end{proof}

For example take $F=\Q$, $n=2$, $K=\Q(\sqrt5)$, of class number one, and $\varepsilon=(1+\sqrt5)/2$, its fundamental unit. If $a'$ denotes the other embedding, put
\[
    t_m=\frac{\pi m}{\log\varepsilon},\qquad
    \chi_m((a))=|a/a'|^{it_m},\qquad m=1,2,\ldots .
\]
The unit relation is $2t_m\log\varepsilon=2\pi m$. Conjugation replaces $\chi_m$ by its distinct inverse, so its induction is a weight-zero dihedral Maass cusp form. Its level is $5$ by the conductor--discriminant formula for induction, its nebentypus is the quadratic character of $K$, and its Laplace eigenvalue is $1/4+t_m^2$ since its real parameters are $it_m,-it_m$. At the inert prime $2$ its normalized trace is zero. The prime $11=(4+\sqrt5)(4-\sqrt5)$ splits and its normalized Hecke eigenvalue is
\[
         2\cos\!\left(t_m\log\frac{4+\sqrt5}{4-\sqrt5}\right).
\]
These give explicit level-five Maass $L$-functions of unbounded spectral parameter to which the all-height strip applies. The real-quadratic construction of the corresponding dihedral newforms is also described in \href{https://www.aimspress.com/aimspress-data/math/2024/2/PDF/math-09-02-238.pdf}{\emph{A prime number theorem in short intervals for dihedral Maass newforms}}, §1.

\subsection{The classical global Weil group and a non-CM obstruction}\label{recentcm:obstruction}
Hu's unitary Weil statement uses the classical global Weil group $W_F$. By the classical induction theorem its finite-dimensional continuous unitary representations are integral virtual sums of inductions of unitary Hecke characters. In a decomposition
\[
       [r]=\sum_j n_j[\operatorname{Ind}_{W_{F_j}}^{W_F}\eta_j]
\]
the corresponding meromorphic function is the signed product of the lower-field functions. In $\Re(s)>\theta$ they have no zeros under $\mathcal H(\theta)$. On the line of possible poles at $s=1+it$, the pole order of the product is
\[
           \sum_j n_j\dim\operatorname{Hom}_{W_{F_j}}
                     (\mathcal N_{F_j}^{it},\eta_j)
           =\dim\operatorname{Hom}_{W_F}(\mathcal N_F^{it},r).
\]
Here $\mathcal N_F$ is the ideal norm in the convention $L(s,\mathcal N_F^{it})=\zeta_F(s-it)$. The equality is Frobenius reciprocity and restriction of norm characters. Its last term is nonnegative, so no residual zero can arise from a negative exponent at one of these points. This explains the passage from rank one to all classical unitary Weil representations in Hu, §§9.1--9.2. It applies, for example, to the full finite-image Artin family. It does not turn the local Weil--Deligne parameters of a general cuspidal representation into a single global Weil representation.

\begin{theorem}\label{recentcm:nocm}
For a non-CM elliptic curve $E/\Q$ there is no continuous two-dimensional unitary representation of the classical $W_\Q$ whose Frobenius polynomials at almost all primes equal $1-(a_p(E)/\sqrt p)T+T^2$.
\end{theorem}
\begin{proof}
Suppose such an $r$ exists. Its determinant is a unitary Hecke character of $\Q$ taking the value one on almost all prime uniformizers. Its incomplete Hecke function is then the incomplete Riemann zeta function. Comparing poles at $1$ forces the character to be trivial by the classical pole criterion for Hecke functions. The finitely many omitted factors have neither a pole nor a zero at $1$. Thus the compact closure $D$ of the image lies in $\mathrm{SU}(2)$.

A finite-dimensional continuous representation of $W_\Q$ factors through a relative Weil group $W_{M/\Q}$ of a finite Galois extension. This follows from the no-small-subgroups property of a Lie group, applied to the kernels in the inverse limit description of the global Weil group. The relative group is an extension of $\operatorname{Gal}(M/\Q)$ by the abelian idèle class group of $M$. In particular $D$ has a closed abelian normal subgroup of finite index and its identity component $D^\circ$ is a torus. See C.~Johansson, \href{https://doi.org/10.1090/tran/6847}{\emph{On the Sato--Tate conjecture for non-generic abelian surfaces}}, §3.1, proof of Proposition~7, and Hu, §9.1. This is the structural restriction on a classical global Weil representation.

If $D^\circ$ acts by scalars, it is trivial because the scalar subgroup of $\mathrm{SU}(2)$ is finite. A compact Lie group with trivial identity component is finite. The Frobenius traces would then range over a finite set, contradicting the continuous Sato--Tate law for the non-CM elliptic curve. If $D^\circ$ has two distinct characters on the representation space, their lines are permuted by $D$. If they are exchanged, the permutation character is nontrivial quadratic, hence cuts out a quadratic field by global class field theory. At every unramified prime inert in that field the matrix is off-diagonal, with trace zero. This gives density at least $1/2$ to the trace-zero primes. The Sato--Tate measure has no atom at zero, so that density is zero. The Sato--Tate theorem used in both cases is due to Barnet-Lamb, Geraghty, Harris and Taylor. See \href{https://doi.org/10.2977/PRIMS/31}{\emph{A family of Calabi--Yau varieties and potential automorphy II}}.

In the remaining case the lines are preserved. Then $r=\chi_1\oplus\chi_2$ is a sum of unitary Hecke characters. The incomplete Rankin--Selberg function of $\pi_E$ against its dual would agree with
\[
       \zeta_\Q^S(s)^2 L^S(s,\chi_1\chi_2^{-1})
                            L^S(s,\chi_2\chi_1^{-1}).
\]
The latter has pole order at least two at $1$ by nonvanishing of Hecke $L$-functions on $\Re(s)=1$ away from their poles. The Rankin--Selberg function for a cuspidal $\pi_E$ against its dual has a simple pole at $1$. Removing a common finite set of places does not change this order: the normalized elliptic parameters are tempered and the characters are unitary, so the local factors at those places are finite and nonzero there. This contradiction disposes of the reducible case.
\end{proof}

Thus the global unitary Weil theorem includes automorphic functions attached to Artin and CM-type data, whereas a fixed non-CM elliptic standard function already fails its representation-theoretic hypothesis. Consequently the unitary Hecke input supplies no fixed zero-free half-plane for the standard function of a non-CM elliptic curve.

\clearpage
\section{The all-height divisor over the inert primes of a CM curve}\label{ip:section}

There is a positive result for an infinite part of an elliptic Euler product where an inert-prime selector really does determine the Frobenius parameter. Take
\[
                 E:\ y^2=x^3-x,\qquad \chi=\chi_{-4}.
\]
Let $I_1(s)$ and $I_2(s)$ be the unitary standard and symmetric-square Euler products of $E$ taken over the rational primes $p\equiv3\pmod4$. These are precisely the inert primes in $\Q(i)$. We give the entire meromorphic divisor of suitable integer powers in $\Re s>1/4$, and of their logarithmic derivatives further left. We use the Dirichlet critical-line conclusion only for $\zeta(s)$ and $L(s,\chi)$.

\begin{theorem}\label{ip:divisor}
Assume the Dirichlet critical-line conclusion for these two functions. In $\Re s>1/4$, the following assertions hold at every height.
\begin{enumerate}
\item $I_1(s)^2$ continues holomorphically and has exactly one zero, a simple zero at $1/2$.
\item $I_2(s)^4$ continues meromorphically. Its real divisor is a double zero at $1$ and a triple pole at $1/2$. At a nonreal point its order is
\[
              2\bigl(\operatorname{ord}_s L(s,\chi)
                        -\operatorname{ord}_s\zeta(s)\bigr).
\]
In particular all its nonreal zeros and poles in this half-plane lie on $\Re s=1/2$. Shared zeros contribute their difference of multiplicities.
\end{enumerate}
\end{theorem}

\begin{proof}
For $p\equiv3\pmod4$ the Legendre symbols of $x^3-x$ and $(-x)^3-(-x)$ are negatives of one another, including the terms with value zero. Summing over $x$ shows $a_p(E)=0$. The two normalized eigenvalues of Frobenius therefore have sum zero and product one, so they are $i,-i$. The symmetric-square eigenvalues are $-1,1,-1$. Write
\[
 P(s)=\prod_{p\equiv3(4)}(1-p^{-s})^{-1},\quad
 Q(s)=\prod_{p\equiv3(4)}(1+p^{-s})^{-1},\quad
 R(s)=\frac{(1-2^{-s})\zeta(s)}{L(s,\chi)}.
\]
In $\Re s>1$ local multiplication gives $R=P/Q$, $PQ=P(2s)$ and thus
\begin{equation}\label{ip:rec}
        P(s)^2=R(s)P(2s),\qquad Q(s)^2=P(2s)/R(s).
\end{equation}
For instance at each inert prime the quotient $R$ has factor $(1+p^{-s})/(1-p^{-s})$, and all the other primes give one. The elliptic partial factors give $I_1=Q(2s)$ and $I_2=PQ^2=P^3/R^2$. Substitution in \eqref{ip:rec} yields
\begin{equation}\label{ip:two}
 I_1(s)^2=\frac{P(4s)}{R(2s)},\qquad
 I_2(s)^4=\frac{R(2s)^3P(4s)^3}{R(s)^2}.
\end{equation}
The factor $P(4s)$ is holomorphic and nowhere zero throughout $\Re s>1/4$ by absolute convergence. Thus these single-valued identities furnish the meromorphic continuations without choosing any fractional powers. In $\Re z>1/2$ the only point of the divisor of $R(z)$ is a simple pole at $1$, by the two asserted critical-line results and boundary nonvanishing at $\Re z=1$. For real $z>0$ the alternating Dirichlet beta series for $L(z,\chi)$ is positive. For $0<z<1$ the alternating eta series is positive and equals $(1-2^{1-z})\zeta(z)$, so $\zeta(z)$ is nonzero there. The correcting factor $1-2^{-z}$ never vanishes in $\Re z>0$. Now take orders on both sides of \eqref{ip:two}, including at the two real exceptional points.
\end{proof}

Removing or restoring finitely many good inert primes multiplies these functions by factors that are holomorphic and nonzero in $\Re s>0$. Thus the same divisor statements hold for the inert parts of every curve $y^2=x^3-Ax$ with $0\ne A\in\Z$, after omitting its bad primes. The same elementary cancellation of Legendre symbols works at every remaining $p\equiv3\pmod4$. On the complementary split primes the Frobenius trace varies, so the theorem on the partial product does not locate the zeros of the full elliptic function.

\subsection{Meromorphic logarithmic derivatives and their boundary}
The quadratic partial-zeta recursion \eqref{ip:rec} is classical. \href{https://arxiv.org/pdf/math/0608206}{Hashimoto, \emph{Partial zeta functions}, \S2 and Claim 3.1}, treats continuation and the natural boundary of partial zeta functions associated with a cyclic prime-degree extension. For the two elliptic products the precise poles, including multiplicity cancellations, can be read off from the two ordinary Dirichlet functions.

\begin{theorem}\label{ip:ladder}
Set $r(z)=R'(z)/R(z)$. The logarithmic derivatives $F_m(s)=I_m'(s)/I_m(s)$, initially defined on $\Re s>1$, have single-valued meromorphic continuation to $\Re s>0$, given by
\[
 \begin{split}
  F_1(s)&=-r(2s)+\sum_{j\ge2}r(2^j s),\\
  F_2(s)&=-\tfrac12 r(s)+\tfrac32\sum_{j\ge1}r(2^j s).
 \end{split}
\]
Assume the same Dirichlet critical-line conclusion. At nonreal arguments the poles in these formulas are precisely the points $s=2^{-j}\rho$ where the indicated coefficient $c_j$ is nonzero and $\operatorname{ord}_\rho R\ne0$. Their residues are $2^{-j}c_j\operatorname{ord}_\rho R$ and their real parts are $2^{-j-1}$. On the positive real axis the poles are $s=2^{-j}$ for the same indices, with residue $-2^{-j}c_j$. Every point of the imaginary axis is an accumulation point of poles of each $F_m$.
\end{theorem}

\begin{proof}
For $f=P'/P$, differentiating \eqref{ip:rec} gives $f(s)=\frac12 r(s)+f(2s)$. Hence $f(s)=\frac12\sum_{j\ge0}r(2^js)$. Initially the tail tends to zero by absolute convergence of the Euler products. On each compact set of the open right half-plane this tail converges normally: if $\sigma=\Re z\ge3$, the Dirichlet series for $r(z)$ is bounded by $2\sum_{n\ge2}(\log n)n^{-\sigma}\ll 2^{-\sigma}$. Only finitely many initial terms can have poles on a given compact set. From $I_1=Q(2s)=P(2s)/R(2s)$ one gets $F_1=2f(2s)-2r(2s)$, and from $I_2=P^3/R^2$ one gets $F_2=3f(s)-2r(s)$. These are the claimed formulas. Under the stated zero locations the distinct indices have disjoint nonreal polar lines, which gives the residues. On the real axis $r$ has only the pole at one with residue $-1$.

Here is a short way to see the accumulation at every height under the same hypothesis. The ordinary one-sided zero counts, with multiplicity, of the two real primitive Dirichlet functions obey
\[
        N_\chi^+(T)-N_\zeta^+(T)=\frac{\log 4}{2\pi}T+O(\log(2+T)).
\]
This is the fixed-conductor Riemann--von Mangoldt formula, applied at conductors $4$ and $1$. Consequently for every sufficiently large $U$, there is a zero of $L(s,\chi)$ of multiplicity exceeding that of $\zeta$ at the same point, at an ordinate $\gamma\in[U,U+C\log U]$ for an absolute $C$. Indeed subtract the count difference at the two ends and take $C$ sufficiently large. For $t>0$ use $U=2^j t$. Then $2^{-j}(1/2+i\gamma)\to it$, and for every sufficiently large $j$ each of these points is a pole of both displayed logarithmic derivatives. Conjugation gives the negative ordinates and the real dyadic poles accumulate at zero. Meromorphic continuation through any point of the imaginary axis is therefore impossible.
\end{proof}

There is also a direct consequence for the split-prime products of the same curve. Their logarithmic derivatives equal those of the full incomplete standard and symmetric-square functions minus $F_1$ and $F_2$, respectively. The full functions have global meromorphic continuation by CM theory. Their log derivatives thus have isolated singularities across the imaginary axis. If the split-prime log derivative extended through a point of that axis, subtracting would extend the corresponding inert-prime log derivative there too. The two infinite prime sets therefore each give a natural boundary, while their combined logarithmic derivative has meromorphic continuation through it.

\clearpage
\section{A power saving for concentrated Frobenius traces}\label{conc:section}
For a non-CM elliptic curve $E/\Q$ and a finite set $S\subset\Z$, write $P_E(x,S)=\#\{p\le x:p\nmid N_E,\ a_p(E)\in S\}$ and $h=|S|\ge1$. Suppose the nontrivial zeros of every finite-order Hecke function over every number field have real part at most $\theta$, with $1/2\le\theta<1$. We also use the verified Dirichlet theorem for the rational quadratic characters. The Hecke hypothesis holds with $\theta=7/8$ by the input of Section~\ref{recentcm:section}.

\begin{theorem}\label{conc:main}
For fixed $E,\theta$, uniformly in $x\ge3$ and $S$, one has
\[
 P_E(x,S)\ll_{E,\theta}
  \min\!\left\{\frac{x}{\log x},\
        \frac{h^{4/5}x^{1-2(1-\theta)/5}}{(\log x)^{3/5}}\right\}.
\]
At $\theta=7/8$ each fixed nonzero trace occurs at most $O_E(x^{19/20}/(\log x)^{3/5})$ times. Any set containing a fixed positive fraction of the good-prime traces up to $x$, counted with their prime multiplicities, contains $\gg_{E,\theta} x^{(1-\theta)/2}/\sqrt{\log x}$ distinct integers, the constant also depending on that fraction.
\end{theorem}

The splitting-prime and Borel-subgroup reduction for a single target trace is due to D.~Zywina, \href{https://arxiv.org/abs/1508.07682}{\emph{Bounds for the Lang--Trotter conjectures}}, §§2, 4 and 7. His GRH theorem gives $x^{4/5}/(\log x)^{3/5}$ for a single nonzero trace. We adapt the argument to arbitrary target sets and the two different zero inputs. For trace zero the geometric method of N.~Elkies gives the stronger unconditional estimate $O_E(x^{3/4}\log x)$. See \emph{Distribution of supersingular primes}, Ast\'erisque 198--200 (1991), 127--132.

\subsection*{The Hecke estimate in growing fields}
For a finite Galois extension $L/\Q$, put $\operatorname{rd}(L)=|D_L|^{1/[L:\Q]}$. Let $F\subseteq M\subseteq L$, $d=[F:\Q]$, and suppose $M/F$ is abelian with group $A$. Let $V$ be a positive integer divisible by the rational primes ramifying in $L$ and put $B=\log(3V\operatorname{rd}(L))$. For $C\subset A$, the number of prime ideals of $F$ of norm between $X/2$ and $X$ lying above rational primes prime to $V$ and having Frobenius in $C$ is
\begin{equation}\label{conc:analytic}
 \ll \frac{|C|}{|A|}\frac X{\log X}
      + |C|^{1/2}d B\,\frac{X^\theta}{\log X}\qquad(X\ge3).
\end{equation}
Choose a smooth nonnegative $w$ supported on $[1/4,2]$ with $w\ge1$ on $[1/2,1]$. For a character $\xi$ of $A$ put $Q_\xi=|D_F|\mathrm N\mathfrak f_\xi$. The standard smoothed explicit formula for the primitive Hecke function gives
\[
 \sum_{\mathfrak p,m}(\log\mathrm N\mathfrak p)\xi(\mathfrak p^m)
                      w(\mathrm N\mathfrak p^m/X)
   ={\bf1}_{\xi=1}\widehat w(1)X+
      O_w\big(X^\theta(\log(2Q_\xi)+d)\big),
\]
where $m\ge1$, and $\widehat w(s)=\int_0^\infty w(u)u^{s-1}\,du$. To see the uniformity, the nontrivial zeros contribute $-\sum_\rho X^\rho\widehat w(\rho)$. Integration by parts gives $|\widehat w(s)|\ll_w(1+|\Im s|)^{-3}$ on $-1/2\le\Re s\le2$. The classical count in each unit interval is $O(\log(2Q_\xi)+d\log(3+|t|))$ with an absolute constant. This follows from the logarithmic derivative at $2+it$, or the uniform local zero count in Zywina, §7. The trivial zero at zero has order $O(d+1)$. On $\Re s=-1/2$ the functional equation, absolute convergence at $3/2$ and the digamma bound give an integral $O_w(X^{-1/2}(\log(2Q_\xi)+d))$. These are the usual explicit-formula inputs. Only $X^\theta$ uses the Hecke strip. Prime powers $m\ge2$ contribute $O_w(d\sqrt X)$, since over each rational prime the sum of residue degrees is at most $d$ and $\sum_{p^m\le2X,m\ge2}\log p\ll\sqrt X$. Deleting prime ideals over $V$ now costs $O_w(d\log(2V))$.

There is a uniform individual bound $\log Q_\xi\le 2d\log\operatorname{rd}(L)$. At each nontrivial inertia subgroup $I$ in the Artin conductor formula the codimension of invariants of a one-dimensional character is at most one. For the regular representation the codimension divided by its degree is $1-1/|I|\ge1/2$. Termwise comparison bounds the local conductor exponent of $\xi$ by twice the relative discriminant exponent divided by $|A|$. The tower discriminant identity then gives $\log Q_\xi\le2d\log\operatorname{rd}(M)\le2d\log\operatorname{rd}(L)$. Finally expand $\mathbf{1}_C$ into characters of $A$. The trivial coefficient is $|C|/|A|$ and the absolute coefficient sum is at most $\sqrt{|C|}$ by Parseval and Cauchy--Schwarz. Sum the smoothed estimates, use positivity, and divide by $\log(X/2)$ to obtain \eqref{conc:analytic}.

\subsection*{Detecting a set of traces}
For every good prime $X/2<p\le X$ the integer $D_p=a_p^2-4p$ is negative and $|D_p|\le4X$ by Hasse. If $(\log X)^8\le y\le X$, at least $c\,y/\log y$ primes $y<\ell\le2y$ satisfy $(D_p/\ell)=1$, uniformly in $p$ for sufficiently large $X$. Indeed the corresponding primitive quadratic character is nonprincipal with conductor at most $4|D_p|$. Dirichlet RH bounds its log-weighted prime sum on this interval by $O(\sqrt y\log^2(8Xy))$. Removing primes dividing $2D_p$ costs $O(\log(8X))$. Combine these estimates with $\sum_{y<\ell\le2y}\log\ell\sim y$ and divide by $\log(2y)$. Discarding $\ell=p$ removes at most one candidate.

By Serre's open-image theorem, for all sufficiently large $\ell$ depending on $E$, $L=\Q(E[\ell])$ has Galois group $G=\mathrm{GL}_2(\F_\ell)$. Let $B_0$ be the upper triangular subgroup and $U$ its upper unitriangular subgroup. For $F=L^{B_0}$ and $M=L^U$, the degree $d$ is $\ell+1$ and $A=B_0/U\simeq(\F_\ell^\times)^2$. All rational ramification divides $N_E\ell$. Also $\log\operatorname{rd}(L)\ll_E\log(2\ell)$. Indeed the classical Galois discriminant inequality gives
\[
       \log\operatorname{rd}(L)
        \le \sum_{p\mid N_E\ell}\log p+\log|G|.
\]
It follows at each $p$ from the local different bound $v_{L_w}(\mathfrak D_{L_w/\Q_p})\le e_w-1+e_w v_p(e_w)$ and $e_w\mid |G|$. This controls $B$ in \eqref{conc:analytic}.

If $a_p\in S$ and $(D_p/\ell)=1$, Frobenius on $E[\ell]$ has an invariant line. The corresponding coset of $B_0$ is fixed by Frobenius and gives a degree-one prime of $F$ over $p$. Its Frobenius in $A$ is a pair $(u,v)$ satisfying $u+v\in S\bmod\ell$. Each such $p$ therefore supplies at least one prime counted by \eqref{conc:analytic}, with
\[
 C=\{(u,v)\in(\F_\ell^\times)^2:u+v\in S\bmod\ell\},
 \qquad |C|\le h(\ell-1).
\]
Each $p$ in the trace set has $\gg y/\log y$ admissible choices of $\ell$, while there are $O(y/\log y)$ primes $\ell$ in total. Summing over $\ell$ and interchanging finite sums bounds the desired number on this dyadic range by
\[
 \ll_E \frac{hX}{y\log X}
       +\frac{\sqrt h\,y^{3/2}X^\theta\log(2y)}{\log X}.
\]
Take $y=h^{1/5}X^{2(1-\theta)/5}(\log X)^{-2/5}$. If $h>X^{(1-\theta)/2}/\sqrt{\log X}$ the trivial $O(X/\log X)$ already implies the claimed bound on the dyadic range. Otherwise $y\le X$, and for fixed $\theta<1$, $y\ge(\log X)^8$ once $X$ is sufficiently large. Substitution gives the second bound of Theorem~\ref{conc:main} on each dyadic range. Summing down to $\sqrt x$ preserves its order, and the remaining primes contribute $O(\sqrt x)$. The prime number theorem gives the other member of the minimum. If $P_E(x,S)\ge c_0x/\log x$, rearrangement gives the assertion on the size of $S$.

For this trace-set test an invariant line and an upper bound for a nonnegative indicator suffice. A signed first moment of $a_p/\sqrt p$ would need cancellation between positive and negative traces. The argument above leaves that estimate open.


\clearpage
\section{Trace sets with few multiplicative quotients}\label{st:section}

Under the two zero inputs of Section~\ref{conc:section}, a set of candidate traces whose multiplicative combinations occupy few rational values admits a stronger estimate. Fix a non-CM elliptic curve $E/\Q$. Let $x\ge3$ and let $S$ be a nonempty set of nonzero integers of absolute value at most $2\sqrt x$. Put
\[
 h=|S|,\qquad K=\frac{|SS^{-1}S|}{h},\qquad
 SS^{-1}S=\{abc^{-1}: a,b,c\in S\}\subset\Q^\times .
\]
In particular $K\ge1$. The count $P_E(x,S)$ is as in Section~\ref{conc:section}.

\begin{theorem}\label{st:counts}
Uniformly in $S$ and $x$,
\[
 P_E(x,S)\ll_{E,\theta}
 \min\left\{\frac{x}{\log x},    \min(h^{4/5},K^{4/5}h^{3/5})\,
       \frac{x^{1-2(1-\theta)/5}}{(\log x)^{3/5}}\right\}.
\]
The same estimate holds with the additional allowed trace zero, with $h,K$ still computed from the nonzero part. The only-zero case is covered in Section~\ref{conc:section}. The constant is independent of $K$. For fixed $K$ the exponent of $h$ drops from $4/5$ to $3/5$. In particular a finite integer geometric progression of $h$ distinct nonzero terms has $|SS^{-1}S|\le3h-2$. Taking consecutive powers of two through $2\sqrt x$ illustrates an improvement by a factor of order $(\log x)^{1/5}$ over the unrestricted finite-set bound.
\end{theorem}

\subsection*{A nonnegative majorant on the split torus}

Let $A=\F_\ell^\times$ for an odd prime $\ell$, $n=\ell-1$, and use normalized multiplicative Fourier coefficients $\widehat g(\eta)=n^{-1}\sum_{u\in A}g(u)\overline{\eta(u)}$. Take a nonempty $T\subseteq A$, put $D=T/T$, and form
\[
 g(u)=\frac1{|T|}\sum_{v\in T}\mathbf1_D(u/v),\qquad u\in A.
\]
For $u\in T$ every $v\in T$ contributes, so $g(u)=1$ exactly. Elsewhere $0\le g(u)\le1$. Counting convolution gives $\sum_{u\in A}g(u)=|D|$. Also
\[
 \widehat g(\eta)=\frac n{|T|}
       \widehat{\mathbf1_D}(\eta)\widehat{\mathbf1_T}(\eta),
 \qquad
 \sum_\eta|\widehat g(\eta)|\le\sqrt{|D|/|T|},
\]
by Cauchy--Schwarz and Parseval. Set $g(0)=0$ and define the nonnegative function on $A^2$
\[
             H(u,v)=g(u+v)+\mathbf1_{u+v=0}.
\]
It majorizes $\mathbf1_{u+v\in T\cup\{0\}}$, and its mean on $A^2$ is
\[
      \widehat H(1,1)=\frac{(\ell-2)|D|+n}{n^2}.
\]
Indeed each nonzero sum has $\ell-2$ ordered representations and zero has $n$.

For characters $\chi,\psi$ of $A$, let $\eta=\chi\psi$ and extend the characters by zero at zero for Jacobi sums. The coefficient of $g(u+v)$ on $A^2$ is
\[
     n^{-1}J(\bar\chi,\bar\psi)\widehat g(\eta).
\]
This follows by fixing the nonzero sum $s$ and substituting $u=st$, $v=s(1-t)$. For $\eta\ne1$, the Jacobi sum has modulus $\sqrt\ell$ when both factors are nontrivial, by the Gauss-sum product formula. With exactly one factor trivial its modulus is one, since the remaining nontrivial character sums to zero over $A$ and the Jacobi sum omits one unit. For $\eta=1$ and nontrivial inverse factors the modulus is one: $t/(1-t)$ runs over the units other than $-1$ when $t\ne0,1$. For two trivial factors it is $\ell-2$. Thus over the $n$ ordered pairs with fixed product $\eta$, the sum of the Jacobi moduli divided by $n$ is at most $\sqrt\ell$ when $\eta\ne1$ and at most $2$ when $\eta=1$. Finally the zero-sum indicator is the coset $v=-u$ of the diagonal. Its Fourier coefficients have magnitude $1/n$ for the $n$ pairs with $\chi\psi=1$ and vanish otherwise. We conclude that
\begin{equation}
  \sum_{\chi,\psi}|\widehat H(\chi,\psi)|
       \le 1+2\sqrt\ell\sqrt{|D|/|T|}.
  \label{st:torus}
\end{equation}
The zero residue is included even when no original target equals zero, because an admissible nonzero target may be divisible by $\ell$.

\begin{proof}[Proof of Theorem~\ref{st:counts}]
Section~\ref{conc:section} gives the unrestricted term, also for $S\cup\{0\}$ since $h+1\le2h$. It remains to prove the $K^{4/5}h^{3/5}$ estimate. On a dyadic interval $X/2<p\le X$ with $\sqrt x\le X\le x$ write $\delta=1-\theta$ and choose
\[
     y=(\sqrt K\, h X^\delta/\log X)^{2/5}.
\]
If $K^4 h^3\ge X^{2\delta}/(\log X)^2$, the proposed bound on this interval is at least a constant multiple of $X/\log X$, so use Chebyshev's bound. Otherwise $h<X^{2\delta/3}$ and $\sqrt K\,h\le(K^4h^3)^{1/3}$ since $K\ge1$. In particular $y\le X$ for large $X$, while $K,h\ge1$ give $y\ge X^{\delta/5}$ for all sufficiently large $X$ depending only on $\theta$. These inequalities also give $(\log X)^8\le y$ and $\log(2y)\ll\log X$. Bounded $X$ is covered by enlarging the constant.

Keep the same set $S$ throughout these dyadic intervals; shrinking it at each interval need not preserve its tripling ratio. For each $s\in S$ the number of primes $y<\ell\le2y$ dividing $s$ is at most $\log(2\sqrt x)/\log y\ll_\theta1$; this also holds for $|s|=1$ with zero divisors. Thus all but $O_\theta(1)$ such primes leave at least $h/2$ elements of $S$ invertible, by summing the divisibility counts. Drop these exceptional auxiliary primes.

For each remaining $\ell$, let $S'$ consist of the invertible elements and let $T$ be their residues. Take the largest fiber $U\subseteq S'$ of the reduction map, so $|U|\ge h/(2|T|)$. For each distinct $d\in D=T/T$ choose one rational lift $b_d=s_1/s_2$ with $s_i\in S'$. The sets $b_dU$ are contained in $SS^{-1}S$, each has $|U|$ elements, and their reduction residues are distinct for distinct $d$. Consequently $|D|\,|U|\le Kh$, giving $|D|\le2K|T|\le2Kh$. This argument allows arbitrary collisions among the reduced traces.

The quadratic splitting-prime estimate of Section~\ref{conc:section} still gives $\gg y/\log y$ allowable $\ell$ for every good $p$ in the dyadic interval, after discarding the bounded exceptional set and $p$ itself. Over $\Q(E[\ell])^{B_0}$ each such $p$ with trace in $S\cup\{0\}$ supplies a degree-one prime with $H(u,v)\ge1$. In the proof of \eqref{conc:analytic}, Fourier expansion works for any nonnegative weight on the abelian Galois group. Its mean multiplies $X/\log X$ and its sum of absolute Fourier coefficients multiplies $dB X^\theta/\log X$ with the field parameters of that section. Here $d=\ell+1$, $B\ll_E\log(2\ell)$, the mean is $O(Kh/\ell)$ and \eqref{st:torus} is $O(\sqrt{K\ell})$. Sum over auxiliary primes and divide by the number of choices available to each $p$. The resulting upper bound is
\[
  \ll_{E,\theta} \frac{KhX}{y\log X}
    +\frac{\sqrt K\,y^{3/2}X^\theta\log(2y)}{\log X}
  \ll_{E,\theta}
  K^{4/5}h^{3/5}\frac{X^{1-2\delta/5}}{(\log X)^{3/5}} .
\]
Summing down to $\sqrt x$ preserves the order uniformly in $K,h$. The primes below $\sqrt x$ are absorbed by the displayed bound since $1-2\delta/5\ge4/5$. This proves the theorem.
\end{proof}

One can also read the conclusion as a restriction on an observed set of traces. If $P_E(x,S\cup\{0\})\ge c x/\log x$ for a fixed $c>0$, the second term of Theorem~\ref{st:counts}, raised to the fifth power, gives
\[
    |SS^{-1}S|\gg_{E,\theta,c}
          \frac{h^{1/4} x^{(1-\theta)/2}}{\sqrt{\log x}}.
\]
At fixed $K$ this requires $h\gg_{E,\theta,c,K}x^{2(1-\theta)/3}/(\log x)^{2/3}$. This concerns finite target sets in the Hasse range at the same value of $x$.

\clearpage
\section{Elliptic curves in nonabelian extensions of number fields}\label{ebc:chapter}

For an elliptic curve $E/\Q$ and a number field $F$, a good rational prime can give several different residue degrees in $F$. This changes the Frobenius traces seen by the Euler product. A permutation representation records all these degrees at once. Its irreducible constituents tell us exactly which additional zeros must be controlled for GRH over $F$. We use arithmetic normalization, so the proposed critical line is $\Re(s)=1$.

\subsection{Induction at every place}\label{ebc:formalism}
Take a finite Galois extension $L/\Q$ containing $F$, set $G=\operatorname{Gal}(L/\Q)$ and $H=\operatorname{Gal}(L/F)$. No normality is required of $H$. Write $V=H^1_{\mathrm{\acute et}}(E_{\overline{\Q}},\overline{\Q}_\ell)$ with geometric Frobenius, so that a good prime $p\ne\ell$ contributes $1-a_pT+pT^2$ as reciprocal factor. The local parameter of an Artin representation $\tau$ is taken in the same Frobenius convention. At a good unramified prime the parameters of $V\otimes\tau$ are $\alpha_p\lambda_j,\beta_p\lambda_j$, where $\lambda_j$ are the eigenvalues of Frobenius on $\tau$. At a general finite place use the Frobenius determinant on inertia invariants of the full tensor. Equivalently, for a Weil--Deligne parameter $(r,N)$ take $(\ker N)^I$. At a prime equal to the chosen $\ell$ use a different auxiliary prime in the compatible system.

\begin{theorem}[Base extension and the zero divisor]\label{ebc:divisor}
Let $\tau$ run through the complex irreducible representations of $G$ and put $m_\tau=\dim\tau^H$. The identity of Euler products is
\begin{equation}\label{ebc:product}
 L(E/F,s)=\prod_{\tau} L(E/\Q,\tau,s)^{m_\tau}.
\end{equation}
It holds also for the completed functions, with compatible gamma factors and absolute conductors. Suppose that each completed factor with $m_\tau>0$ is entire. Then the completed function over $F$ has all its zeros on $\Re(s)=1$ if and only if every one of those factors does. At each complex $s_0$, its order of vanishing is the sum of their orders weighted by $m_\tau$.
\end{theorem}

\begin{proof}
The tensor identity for induction is
\[
 \operatorname{Ind}_{G_F}^{G_\Q}(V|_{G_F})
       \simeq V\otimes\overline{\Q}_\ell[G/H].
\]
In the group-ring model the map from the induced representation sends $g\otimes v$ to $gv\otimes[gH]$. It respects the relation $gh\otimes v=g\otimes hv$ and is an isomorphism on each coset summand. Frobenius reciprocity gives the claimed multiplicities in the permutation representation: its character pairing with $\tau$ equals the multiplicity of the trivial character in $\tau|_H$.

At a fixed prime of $\Q$, the double cosets of the local decomposition group and $G_F$ correspond to the primes of $F$ above it. Mackey decomposition therefore gives the direct sum of inductions of the local parameters. The local induction identity for $L$-factors gives exactly the product of their factors over $F$. One way to see the identity is to let $K'/K$ be the local extension, of residue degree $d$. After passing to inertia invariants, the induced space consists of $d$ translates of the $I_{K'}$-invariants; the translates are indexed by the residue cosets. Frobenius permutes them cyclically and its $d$th power on the first translate is the Frobenius over $K'$. For a Frobenius block with product $A$, the cyclic determinant is $\det(1-T^dA)$. For a Weil--Deligne parameter the kernel of the induced monodromy operator is the induction of the kernel, so the same calculation applies before taking invariants. This proves~\eqref{ebc:product} including primes of simultaneous ramification.

The local conductor formula for induction, summed over all places, gives the absolute conductor $|D_F|^2 N_{F/\Q}\mathfrak n_{E/F}$ for the induced rank $2[F:\Q]$ system. Here $\mathfrak n_{E/F}$ is the elliptic conductor ideal. At infinity use $\Gamma_{\C}(s)=2(2\pi)^{-s}\Gamma(s)$. A real place contributes one such factor and a complex place contributes two, by the two Hodge types $(1,0),(0,1)$. Over $\Q$ twisting by a finite character at the real place acts trivially on this gamma factor, even for the sign character. Indeed the real elliptic Weil parameter is induced from $\C^\times$, where the sign character is trivial. Decomposing the restriction of $\tau$ to the order-two complex conjugation group shows that the factor for the twist is $\Gamma_{\C}(s)^{\dim\tau}$. The equality of gamma factors follows from $\sum_\tau m_\tau\dim\tau=[F:\Q]$. Conductors multiply in direct sums, giving the completed identity. The equalities also hold under the convention omitting the fixed factors $2$ everywhere.

Finally, analytic continuation preserves the identity from its common half-plane of absolute convergence. Orders of nonzero meromorphic functions add in a product. For entire factors the orders are nonnegative and all occurring $m_\tau$ are positive, so an off-line zero in any factor gives an off-line zero of the product. This proves both directions.
\end{proof}

For a finite unramified prime of good reduction, if the Frobenius cycles on $G/H$ have lengths $d_1,\ldots,d_b$, its full reciprocal factor over $F$ is
\[
        \prod_{j=1}^b(1-\alpha_p^{d_j}T^{d_j})
                        (1-\beta_p^{d_j}T^{d_j}),\qquad T=p^{-s}.
\]
For $F=L$ the representation is regular and $m_\tau=\dim\tau$. For a nonnormal field it is possible that some irreducibles of the Galois closure do not occur at all. At primes ramified in both $E$ and $L$ the full invariants in the proof are essential. For instance, a nontrivial scalar quadratic inertia action on an elliptic representation can disappear after twisting by the same quadratic inertia character.

These identities are standard Artin formalism. For elliptic curves and the standard representation of a symmetric group see R.~Lemke Oliver and F.~Thorne, \href{https://lemkeoliver.github.io/papers/21-RankGrowth.pdf}{\emph{Rank growth of elliptic curves in non-abelian extensions}}, §3. For the induction formula for the conductor see J.-P.~Serre, \emph{Local Fields}, Chapter~VI, §§1--3. Serre's formula includes the different of the extension; this accounts for the discriminant factor above.

\subsection{Which constituents are actual cuspidal functions?}\label{ebc:cuspidal}
Suppose now that $E/\Q$ has no complex multiplication over $\overline\Q$. A finite twist $\tau$ has finite monodromy, whereas the elliptic representation remains irreducible after restriction to any finite extension. This distinction gives the following control of cuspidality.

\begin{lemma}\label{ebc:irreducible}
If $K/\Q$ is finite and $\tau$ is an irreducible finite-image representation of $G_K$ in characteristic zero, then $V|_{G_K}\otimes\tau$ is absolutely irreducible. If $\theta$ is a nontrivial finite character of $G_K$, then $V|_{G_K}\otimes\theta$ and $V|_{G_K}$ are inequivalent.
\end{lemma}
\begin{proof}
Use embeddings of the algebraic coefficients of the finite representation into $\overline{\Q}_\ell$. On $G_{K'}=\ker\tau$ the elliptic image is Zariski dense in $\mathrm{GL}_2$ by Serre's open image theorem: passage to a finite extension leaves a subgroup of finite index in an open subgroup. In particular $V$ is absolutely simple on $G_{K'}$ and its endomorphism ring consists of scalars. Over $G_{K'}$ the tensor is a direct sum of copies of this simple module. Every submodule is therefore $V\otimes D$ for a subspace $D$ of the multiplicity space of these copies: apply $\operatorname{Hom}_{G_{K'}}(V,-)$ and the evaluation map. The action of $G_K$ on the multiplicity space in this description is $\tau$. Thus a $G_K$-stable submodule yields a $\tau$-stable $D$ and irreducibility follows. For the second assertion, any intertwiner commutes with all of $\ker\theta$. It is scalar on $V$. A nonzero scalar can intertwine the two representations on all of $G_K$ only if $\theta=1$. The open image theorem used here is J.-P.~Serre, \emph{Propriétés galoisiennes des points d'ordre fini des courbes elliptiques}, Invent.\ Math.\ 15 (1972), 259--331.
\end{proof}

\begin{theorem}[Finite twists induced along cyclic towers]\label{ebc:monomial}
Let $G=\operatorname{Gal}(L/\Q)$ as above. Suppose an irreducible $\tau$ of $G$ is induced by a character $\theta$ of $J\subseteq G$, and $J$ admits a chain
\[
       J=J_0\triangleleft J_1\triangleleft\cdots
                \triangleleft J_a=G,\qquad [J_{i+1}:J_i]\text{ prime}.
\]
For a non-CM elliptic curve $E/\Q$ the unitarily normalized twist
$L(E/\Q,\tau,s+1/2)$ is the standard function of a unitary cuspidal representation of $\mathrm{GL}_{2\dim\tau}(\mathbb A_\Q)$. The completed twist is entire in arithmetic normalization. Consequently Theorem~\ref{ebc:divisor} gives an equivalence with the separate cuspidal GRH statements whenever all the constituents occurring for $H$ satisfy the displayed condition.
\end{theorem}

\begin{proof}
Put $K_i=L^{J_i}$. The extension $K_i/K_{i+1}$ is cyclic of prime degree. Start from the modular representation $\pi_E$ over $\Q$ and take successive cyclic base changes down the subgroup chain to $K_0$. The result is cuspidal at each step. Indeed, a cuspidal representation whose base change to a cyclic extension of prime degree ceases to be cuspidal has a self-twist by the associated nontrivial finite character of the lower field. The local Frobenius polynomials would then give a self-twist of the attached elliptic representation, by Chebotarev and the characteristic-zero Brauer--Nesbitt theorem. This contradicts Lemma~\ref{ebc:irreducible}. At each step the arithmetic local parameters are the restrictions of the original elliptic parameters, so this argument iterates.

Twist the resulting representation over $K_0$ by the finite-order Hecke character corresponding to $\theta$. Induce back up by successive cyclic automorphic inductions. To check the intervening cuspidality, put
\[
      \tau_i=\operatorname{Ind}_{J}^{J_i}\theta,\qquad
      W_i=V|_{G_{K_i}}\otimes\tau_i.
\]
Each finite representation $\tau_i$ is irreducible. Otherwise the nonzero summands of a decomposition, induced up to $G$, would decompose $\tau$ by transitivity. Thus $W_i$ is irreducible by Lemma~\ref{ebc:irreducible}. The projection formula identifies the induction of $W_i$ to $G_{K_{i+1}}$ with $W_{i+1}$.

For a cyclic step of prime degree the only obstruction to cuspidal automorphic induction is invariance under a nonidentity Galois conjugation. If this occurs, the representations $W_i$ and its conjugate are equivalent by their good local polynomials and Chebotarev. The induced $W_{i+1}$ would then have an endomorphism algebra of dimension at least two: Frobenius reciprocity identifies that algebra with the sum of the intertwining spaces between $W_i$ and all its conjugates. This contradicts its irreducibility. Thus every induction is cuspidal. The local compatibility of cyclic induction with induction of Weil--Deligne parameters and of base change with restriction identifies the completed product everywhere, including the twists in the archimedean normalization. Cuspidal standard functions of rank at least two have entire completions. All operations preserve unitarity. This proves the theorem.

The inputs are the modularity theorem for $E/\Q$ with local compatibility, and J.~Arthur and L.~Clozel, \href{https://www.math.toronto.edu/arthur/pdf/30.pdf}{\emph{Simple Algebras, Base Change, and the Advanced Theory of the Trace Formula}}, Chapter~3, on cyclic base change, induction, and the criteria for cuspidality. For the relation to local Langlands see M.~Harris, \href{https://www.numdam.org/article/AST_2005__298__17_0.pdf}{\emph{The local Langlands correspondence: notes of (half) a course at the IHP Spring 2000}}.
\end{proof}

One class where the theorem applies uniformly consists of all finite nilpotent $G$. Every irreducible character of a finite supersolvable group is monomial, and finite nilpotent groups are supersolvable. Furthermore every proper subgroup of a finite nilpotent group is properly contained in its normalizer. This latter assertion follows by writing the group and the subgroup as products of their Sylow subgroups and using the normalizer condition in a finite $p$-group. Hence an inducing subgroup admits the required chain: enlarge it within its normalizer by the inverse image of a subgroup of prime order in the normalizer quotient, and repeat in the finite group. The monomiality result is the classical theorem for supersolvable groups; see I.~M.~Isaacs, \emph{Character Theory of Finite Groups}, Chapter~6. In particular the GRH problem for $E$ over any field inside such an $L$ is equivalent to GRH for an explicit positive list of cuspidal twists. The trivial constituent in that list recovers $E/\Q$; if there are nontrivial constituents, each brings its own zero-location requirement.

For a general solvable Galois $L/\Q$ one can still obtain a downward implication directly by choosing a tower of cyclic prime extensions from $\Q$ to $L$. At any step $M/K$ the cyclic identity is the positive product of the twists of $E/K$ by all characters of $\operatorname{Gal}(M/K)$. The elliptic representation over $K$ is cuspidally automorphic by the base-change argument in the first paragraph of the proof. Its finite character twists are cuspidal and entire. Thus GRH for $E/L$ implies GRH for $E$ over each lower field in this tower, in particular over $\Q$. This inference concerns zeros of an entire product. An ascent from $E/\Q$ would require control of the additional twists.

There is a second sufficient hypothesis for entireness which does not require a chain of inducing subgroups. Suppose an irreducible $\tau$ is already known to be strongly automorphic: a unitary cusp form $\sigma_\tau$ on $\mathrm{GL}_{d}$, $d=\dim\tau$, has its Artin local parameters at all places. The completed twist in unitary normalization is then
$\Lambda(s,\pi_E\times\sigma_\tau)$ by compatibility of local Rankin--Selberg factors with local Langlands. This product of two cuspidal data is entire unless $d=2$ and the two representations are contragredient up to a unitary norm power. The exceptional case cannot occur. At the real place the restriction of an Artin parameter to $\C^\times$ is trivial, whereas the elliptic parameter in unitary normalization restricts on the unit circle to the two nontrivial angular characters of exponents $1$ and $-1$. A norm power is trivial on this circle. Thus the analytic Rankin--Selberg theorem supplies the entireness needed in Theorem~\ref{ebc:divisor}. This deduction does not require a general tensor-product transfer to $\mathrm{GL}_{2d}$. For $d=2$ such a transfer is known by D.~Ramakrishnan, \href{https://annals.math.princeton.edu/articles/11613}{\emph{Modularity of the Rankin--Selberg $L$-series, and multiplicity one for $\mathrm{SL}(2)$}}, Theorem~M.

\subsection{A cubic field with a cuspidal degree-four factor}\label{ebc:cubic}
Let $F=\Q(\sqrt[3]{2})$, $k=\Q(\sqrt{-3})$, and $L=Fk$ the splitting field of $X^3-2$. The polynomial is Eisenstein at $2$ and its discriminant is $-108$. Its splitting group is therefore $S_3$ with quadratic resolvent $k$. This polynomial gives the full ring of integers of $F$. At $2$ the generating root is a uniformizer for the totally ramified cubic extension and hence generates its local integer ring. At $3$, the substitution $X=u-1$ gives $u^3-3u^2+3u-3$, Eisenstein again, and the same argument applies. At other primes the polynomial discriminant is a unit.

Let $\epsilon$ be the sign character and $\rho$ the two-dimensional representation on the sum-zero subspace of the permutation representation on three letters. A transposition has eigenvalues $1,-1$ on $\rho$, and a three-cycle has the two primitive cube roots of unity. The permutation and regular representations give
\begin{align}
 \Lambda(E/F,s)&=\Lambda(E/\Q,s)\Lambda(E/\Q,\rho,s),\label{ebc:cubicfactor}\\
 \Lambda(E/k,s)&=\Lambda(E/\Q,s)\Lambda(E/\Q,\epsilon,s),\notag\\
 \Lambda(E/L,s)&=\Lambda(E/\Q,s)\Lambda(E/\Q,\epsilon,s)
                               \Lambda(E/\Q,\rho,s)^2.\notag
\end{align}
These are completed identities with the conventions of Theorem~\ref{ebc:divisor}. They imply in particular that for every non-CM $E/\Q$, GRH for $E/L$ is equivalent to the conjunction of GRH for $E/F$ and for $E/k$. Over $F$ itself the sign-character twist is absent.

To identify the degree-four factor, choose a nontrivial character $\psi$ of the cyclic group $\operatorname{Gal}(L/k)=A_3$. Conjugation by the other coset sends it to $\psi^{-1}$, which is distinct from $\psi$. Hence $\operatorname{Ind}_{A_3}^{S_3}\psi=\rho$, as one also checks on characters. Class field theory gives a finite-order Hecke character of $k$. Its dihedral automorphic induction to $\Q$ is cuspidal on $\mathrm{GL}_2$ and has local parameter $\rho$. Complex conjugation is a transposition, so its determinant is odd. The conductor is $108$ by the conductor-discriminant formula for $1\oplus\rho$, using $D_F=-108$. The extra factor in~\eqref{ebc:cubicfactor} is the Rankin--Selberg product of $\pi_E$ with this dihedral form, in the arithmetic shift. For non-CM $E$, Theorem~\ref{ebc:monomial}, with the single normal inclusion $A_3\triangleleft S_3$, makes it a cuspidal degree-four standard function. Equivalently its cuspidal representation is
\[
       \operatorname{AI}_{k/\Q}
               \bigl(\operatorname{BC}_{k/\Q}(\pi_E)\otimes\psi\bigr).
\]
The base-changed elliptic representation has no finite self-twist by Lemma~\ref{ebc:irreducible}, so its twists by $\psi$ and by $\psi^{-1}$ are distinct. This explicitly checks the cyclic cuspidality criterion. The fact that the cubic $F/\Q$ itself is nonnormal causes no difficulty: the automorphic operations here use the cyclic quadratic field $k$.

For numerical local data take
\[
             E_{11}:\quad y^2+y=x^3-x^2-10x-20 .
\]
This is the conductor-$11$ optimal curve of its isogeny class. It has multiplicative reduction at $11$, good reduction elsewhere, and nonintegral $j$, hence no complex multiplication. One source for its newform and model is Y.~Martin and K.~Ono, \href{https://uva.theopenscholar.com/files/ken-ono/files/026_8.pdf}{\emph{Eta-quotients and elliptic curves}}, Theorems~1--2. The disjoint ramification gives
\[
      \mathfrak f(E_{11}\otimes\rho)
      =11^{\dim\rho}\,\mathfrak f(\rho)^2
      =11^2\,108^2=1{,}411{,}344 .
\]
Indeed, where the elliptic representation is unramified the Artin conductor of the tensor is twice that of $\rho$, and at $11$ its exponent is twice the elliptic exponent $1$. Hence the absolute conductor for $E_{11}/F$ is $11^3\,108^2$, also obtained by multiplying the absolute discriminant squared by the norm of its elliptic conductor.

For $p=2,3$ inertia acts transitively on the three embeddings of the local cubic field because its extension is totally ramified. The inertia invariants on $\rho$ therefore vanish: the invariants in the permutation representation have dimension one and the trivial line accounts for them. Since the elliptic representation is unramified, the degree-four factor at either prime is $1$. At $11$ the Frobenius is a transposition: $11\equiv2\bmod3$, the cube map on $\mathbb F_{11}$ is bijective, and the polynomial is separable there. On the one-dimensional elliptic inertia invariants Frobenius has an eigenvalue $a_{11}=\pm1$, so the twist's reciprocal factor at $11$ is $(1-a_{11}T)(1+a_{11}T)=1-T^2$.

Here is one good prime of each possible splitting type. Write $T=p^{-s}$, and let $Q_p(T)$ be the local reciprocal polynomial of $E_{11}\otimes\rho$.
\[
\begin{array}{c|c|r|l}
 p&\text{cycle lengths on the cubic roots}&a_p&Q_p(T)\\ \hline
 5&1,2&1&1+9T^2+25T^4\\
 7&3&-2&1-2T-3T^2-14T^3+49T^4\\
 31&1,1,1&7&(1-7T+31T^2)^2
\end{array}
\]
For a transposition the polynomial is
$(1-\alpha_p^2T^2)(1-\beta_p^2T^2)$.
For a three-cycle it is
$(1+\alpha_pT+\alpha_p^2T^2)(1+\beta_pT+\beta_p^2T^2)$.
The traces $a_p$ in the table come from counting the displayed elliptic equation over $\mathbb F_p$. The cubic has respectively one root, no root, and three roots at $5,7,31$. In particular a transposition prime can have zero first trace for the twist while contributing a nontrivial second trace and a nonconstant factor.

\subsection{An unconditional double central zero over the cubic field}\label{ebc:central}
Changing the elliptic curve in the same extension can force a central zero of the cuspidal degree-four twist. Take
\[
                  E_{37}:\quad y^2+y=x^3-x.
\]
Its invariants are $\Delta=37$, $c_4=48$, $c_6=-216$, so it has one prime of bad reduction, with conductor exponent one at $37$. Its reduction is nonsplit multiplicative because $-c_6=6^3$ and $6$ is a nonsquare modulo $37$. Indeed $(2/37)=-1$ and $(3/37)=1$. The $j$-invariant $48^3/37$ is nonintegral, so the curve has no complex multiplication. The real root number is $-1$, that at nonsplit multiplicative reduction is $+1$, and all good finite root numbers are $+1$. Hence $w(E_{37})=-1$.

For the standard representation $\rho_F$ of an $S_d$ field, the coprime-conductor root-number formula of V.~Dokchitser gives
\[
  w(E\otimes\rho_F)= w(E)^{d-1}\operatorname{sgn}(D_F)
                             \left(\frac{D_F}{N_E}\right)
                    \quad\text{if }(D_F,N_E)=1 .
\]
See \href{https://lemkeoliver.github.io/papers/21-RankGrowth.pdf}{Lemke Oliver--Thorne}, Lemma~3.5, citing Dokchitser, \emph{Root numbers of non-abelian twists of elliptic curves}, Corollary~2. For the present $F$, the factor $(D_F/37)=(-3/37)=1$ since $37\equiv1\bmod3$. The exponent $d-1=2$ is even and $D_F<0$, giving $w(E_{37}\otimes\rho)=-1$. The twist is self-dual in unitary normalization and entire by Theorem~\ref{ebc:monomial}; its completed arithmetic function is odd about $s=1$. It has a zero there of positive odd order. The same is true of $\Lambda(E_{37}/\Q,s)$ by its own sign. We have therefore proved
\begin{equation}\label{ebc:rank}
             \operatorname{ord}_{s=1}\Lambda(E_{37}/F,s)\ \ge\ 2 .
\end{equation}
The order is even, since it is a sum of two odd orders. The sign of the product alone would have been $+1$ and would have allowed order zero. No rank conjecture or zero-location hypothesis enters~\eqref{ebc:rank}. This identifies an actual zero on the proposed line, while other zeros of these functions remain uncontrolled by the sign calculation.

\subsection{The exact order of the degree-four zero}\label{ebc:exact}
In the cubic example the central multiplicities can be determined exactly. Write $D(s)=L(E_{37}/\Q,\rho,s)=\sum_{n\ge1} b_n n^{-s}$. Its conductor is $(37\cdot108)^2$, by disjoint ramification as in the level-$11$ computation. Put
\[
       A=\frac{3996}{4\pi^2},\qquad b=\frac{2}{\sqrt A}.
\]
Up to a positive constant the completed function is $A^s\Gamma(s)^2D(s)$. Its sign is $-1$ by the preceding calculation. We will prove
\begin{equation}\label{ebc:certorder}
      D'(1)>\frac95,\qquad
      \operatorname{ord}_{s=1}\Lambda(E_{37}/F,s)=2,\qquad
      \operatorname{ord}_{s=1}\Lambda(E_{37}/L,s)=3 .
\end{equation}

The Bessel function needed here has the real integral
$K_0(x)=\int_0^\infty e^{-x\cosh v}\,dv$. The Mellin transform of $2K_0(2\sqrt{t})$ is $\Gamma(s)^2$: take the Mellin convolution of two copies of $e^{-t}$ and substitute $u=\sqrt t\,e^v$ in $\int_0^\infty e^{-u-t/u}\,du/u$. Thus
\[
       \Theta(y)=2\sum_{n\ge1}b_nK_0(b\sqrt{ny})
\]
has Mellin transform $A^s\Gamma(s)^2D(s)$ in its half-plane of absolute convergence. Mellin inversion and the entire functional equation give $\Theta(y)=-y^{-2}\Theta(1/y)$. In detail shift the Mellin inversion line from $\Re(s)=2$ to $\Re(s)=0$ and use $s\mapsto2-s$ in the functional equation. The horizontal integrals vanish by Stirling's formula and polynomial growth of a cuspidal standard function on fixed vertical strips. Splitting the Mellin integral at one now gives the entire continuation as
\[
       \int_1^\infty \Theta(y)(y^s-y^{2-s})\,\frac{dy}{y}.
\]
Its derivative at $1$ is $2\int_1^\infty\Theta(y)\log y\,dy$. For $c>0$, integration by parts using $K_0'(u)=-K_1(u)$ and $(uK_1(u))'=-uK_0(u)$ gives
\[
  \int_1^\infty K_0(c\sqrt y)\log y\,dy
   =4\int_1^\infty zK_0(cz)\log z\,dz
   =\frac4{c^2}K_0(c).
\]
All the integrals converge exponentially and the series may be differentiated under the integral. Since $D(1)=0$, division by $A$ proves the exact identity
\begin{equation}\label{ebc:kzero}
                    D'(1)=4\sum_{n\ge1}\frac{b_n}{n}K_0(b\sqrt n).
\end{equation}

Here is a rational certificate of its sign. The simple bounds $19/100<b<1/5$, $\cosh1<2$, $\cosh2<4$ imply
\[
    K_0(b)> e^{-2/5}+e^{-4/5}>1
\]
by integrating separately over $[0,1]$ and $[1,2]$. For an upper bound use $\cosh v\ge1+v^2/2$ to get $K_0(x)\le e^{-x}\sqrt{\pi/(2x)}$. For $n\ge25$ the square root here is at most $4/3$ at $x=b\sqrt n$, and $e^{-b}<5/6$. Hence the magnitude of the possible negative contribution at such an $n$ is at most
\[
      \frac{16|b_n|}{3n}(5/6)^{\lfloor\sqrt n\rfloor}.
\]
The following are integer upper bounds in units of $1/1000$ for the sums of these last expressions restricted to $b_n<0$. All entries involve only integer point counts and rational arithmetic.
\[
\begin{array}{c|r|r}
      n&\#\{n:b_n<0\}&\text{upper bound in thousandths}\\ \hline
     2\text{--}24&0&0\\
     25\text{--}99&5&1322\\
     100\text{--}399&18&719\\
     400\text{--}1599&69&71\\
     1600\text{--}6399&242&3\\
     6400\text{--}10000&182&1
\end{array}
\]
For reproducibility the complete recurrence is as follows. At $p=2,3$ take the reciprocal polynomial $Q_p=1$, and at $37$ use $1-T+T^2$. At $37$ the elliptic sign is nonsplit and $2^{12}\not\equiv1\bmod37$, giving a three-cycle on the cubic roots. At other primes put $a=-\sum_{x\bmod p}(\frac{1+4(x^3-x)}p)$ by the quadratic formula in $y$. If $p\equiv2\bmod3$, use $Q_p=1+(2p-a^2)T^2+p^2T^4$. For $p\equiv1\bmod3$ use $(1-aT+pT^2)^2$ if $2^{(p-1)/3}=1\bmod p$, and $1+aT+(a^2-p)T^2+paT^3+p^2T^4$ otherwise. Write $q_k=[T^k]Q_p$. Compute $b_1=1$, $b_{p^j}=-\sum_{k=1}^{\min(j,\deg Q_p)}q_k b_{p^{j-k}}$, and multiply for coprime prime powers. The integer and rational loops are supplied in \texttt{checks/dihedral\_derivative.py}.

To bound the rest, $|b_n|\le d_4(n)\sqrt n$ by the local eigenvalue bound, where $d_4(n)$ counts ordered factorizations into four positive integers. Pairing factors gives $d_4(n)=\sum_{de=n}d(d)d(e)\le4\sqrt n\,d(n)\le8n$. Thus beyond $B=10000$ the tail in~\eqref{ebc:kzero} is at most $32\sum_{n>B}\sqrt n e^{-b\sqrt n}$. The summand is decreasing there. Integrating from $B$ to infinity bounds this by
\[
 64 e^{-b\sqrt B}\left(\frac{B}{b}+\frac{2\sqrt B}{b^2}+\frac2{b^3}\right)<\frac{46}{1000}.
\]
For the last inequality use $(5/6)^{100}$ for the exponential and $19/100$ for $b$ in the denominators; it is then an integer comparison. The $n=1$ term exceeds $4$. Subtract the six bounds in the table and this tail to obtain $D'(1)>9/5$. This proves that the degree-four central zero is simple.

It remains to handle the sign-character twist in the splitting field. Let $L_-(s)=L(E_{37}/\Q,\epsilon,s)=\sum c_n n^{-s}$. It has conductor $333=9\cdot37$ and root number $w(E_{37})\epsilon(-37)=+1$ by the same coprime-conductor formula. Put $A_-=\sqrt{333}/(2\pi)$, so $2<A_-<3$. The single gamma factor uses the exponential Mellin kernel in place of $2K_0$. Splitting its integral at one as above, now with sign $+1$, gives
\[
                  L_-(1)=2\sum_{n\ge1}\frac{c_n}{n}e^{-n/A_-}.
\]
Point counting at $2,5$ gives $a_2(E_{37})=a_5(E_{37})=-2$. The Euler recurrences and the trivial local factor at $3$ give $(c_1,\ldots,c_6)=(1,2,0,2,2,0)$. The first two summands inside the sum exceed $1/2+1/3=5/6$. Since $|c_n|\le d(n)\sqrt n\le2n$, the tail from $n=7$ is bounded by $2(18/25)^7/(1-18/25)<5/6$, using $e^{-1/3}<18/25$. Thus $L_-(1)>0$. Finally $\Lambda(E_{37}/\Q,s)$ has a simple central zero, proved in Lemma~\ref{ep:positive}. Inserting these three orders into~\eqref{ebc:cubicfactor} gives exactly~\eqref{ebc:certorder}.

\subsection{A zero-free region for the degree-four constituent}\label{ebc:strip}
The inverse Mellin kernel also controls zeros away from the exact center. Write
$\Phi(X)=\sum b_n K_0(b\sqrt n\,X)$, so $\Theta(y)=2\Phi(\sqrt y)$ in the notation of~\eqref{ebc:kzero}.

\begin{theorem}\label{ebc:horizontal}
The completed degree-four function of $E_{37}\otimes\rho$ has exactly one real zero. It is simple, at $s=1$. Its only zero with $|\Im(s)|\le1/3$ is that central zero.
\end{theorem}

\begin{proof}
We first certify $\Phi(X)>0$ for every $X>1$, with $\Phi(X)>1/4$ on $2\le X\le5$. All finite verifications in this paragraph are specified below. The functional equation gives $\Phi(1)=0$. Evaluation of the differentiated kernel at one gives $\Phi'(1)>217/100$. For $X\ge1$, an absolute bound on $\Phi''(X)$ is obtained by summing
\[
  |b_n|t_n^2\left(K_0(t_nX)+\frac{K_1(t_nX)}{t_nX}\right),
               \qquad t_n=b\sqrt n.
\]
Each summand decreases with $X$ by the integral representations of $K_0,K_1$. The bound at one is $531$, proving positivity of $\Phi'$ up to $1+1/256$, since $217/100-531/256>0$. For the next intervals the table gives the mesh width $h$, a bound $M$ on $|\Phi''|$ over the whole interval, and a lower bound $v$ for $\Phi$ at all mesh points. Each mesh starts at the left endpoint and the last step ends at the right endpoint even if shorter than $h$.
\[
\begin{array}{c|c|r|c}
 \text{interval}&h&M&v\\\hline
 {[257/256,17/16]}&1/128&531&1/125\\
 {[17/16,5/4]}&1/32&531&11/100\\
 {[5/4,3/2]}&1/16&180&9/25\\
 {[3/2,2]}&1/8&74&3/5\\
 {[2,5]}&1/4&18&47/100
\end{array}
\]
The maximum discrepancy from linear interpolation on a mesh cell is $Mh^2/8$, by the degree-one interpolation remainder. In each row $v-Mh^2/8>0$, and in the last it exceeds $1/4$.

Here are the details for verifying the finite bounds. Use the integer coefficients of~\eqref{ebc:kzero} through $B=65536$, calculated by its stated Euler recurrence. At rational $X$ sum $b_nK_0(t_nX)$, at $1$ also sum $-b_nt_n K_1(t_n)$, and at each required curvature endpoint sum the positive majorant above. The omitted tail in any of these calculations is less than $1/5000$. In fact for $t\ge3$ both $K_0(t)$ and $K_1(t)$ are at most $e^{-t}$. For $K_1$ insert $\cosh v$ into the integral for $K_0$ and use $(\cosh v)e^{-\cosh v}\le e^{-1}$, followed by the Gaussian bound for $K_0(t-1)$. The second derivative is bounded using $K_0(t)+K_1(t)/t$. Since $t_n<\sqrt n/5$ and $|b_n|\le8n^{3/2}$, all three tails are bounded by
\[
 8\sum_{n>B} n^{5/2}e^{-b\sqrt n}
 \ \le\ 16\int_{256}^\infty u^6 e^{-bu}\,du
 \ <\ \frac1{5000}.
\]
The sum-integral inequality uses monotonicity on the tail. The final comparison follows by integrating by parts six times, replacing $b$ in the denominators by $19/100$ and the exponential by $(5/6)^{256}$. At arguments above $25$ the same elementary bounds suffice for the finite sums. At the remaining arguments the calculation uses rigorous $128$-bit intervals for $K_0$ and $K_1$. The supplied \texttt{checks/dihedral\_kernel.py} runs all comparisons with the outward-rounded \href{https://flintlib.org/doc/arb_hypgeom/index.html}{FLINT real-ball functions}. The table bounds include the infinite tails.

For $X\ge5$ the integral representation gives
\[
 \frac{K_0(t_nX)}{K_0(bX)}\le e^{-bX(\sqrt n-1)} .
\]
The contribution in this ratio from all negative coefficients through $B$ is bounded by the exact rational sum $\sum_{n\le B,\ b_n<0}|b_n|(5/6)^{5(\lfloor\sqrt n\rfloor-1)}<1/2$, computed by the same script. For the tail, $5b(\sqrt n-1)\ge b\sqrt n$ for $n\ge2$, so the preceding integral majorant applies. The $n=1$ term alone dominates both bounds. This finishes the positivity proof for every $X>1$.

For a real $s=1+u$ the completed function is the integral of $2\Theta(y)\sinh(u\log y)$ over $1<y<\infty$. Its sign is therefore the sign of $u$. For $s=1+u+it$, $0<|t|\le1/3$, divide its imaginary part by $2t$ to get the integral of
\[
   \Theta(y)\cosh(u\log y)\frac{\sin(t\log y)}{t}.
\]
For $1\le y\le 10^4$ the integrand is nonnegative since $(\log 10^4)/3<\pi$. Its integral over $4\le y\le9$ exceeds $5/4$, by $\Theta(y)>1/2$, $\cosh(u\log y)\ge1$, $\log y>1$ and $\sin v/v>1/2$ for $|v|\le(\log9)/3$.

To bound all potentially negative contributions, take $|u|\le1/2$. For $X\ge100$ the above ratio and $bX(\sqrt n-1)\ge4\sqrt n$ for $n\ge2$ show
\[
   \sum_{n\ge2}|b_n|e^{-bX(\sqrt n-1)}
   \le 8\sum_{n\ge2}n^{3/2}e^{-4\sqrt n}
   \le16\int_1^\infty v^4e^{-4v}\,dv<1.
\]
It follows that $|\Theta(y)|\le4e^{-b\sqrt y}$ for $y\ge10^4$. Use $|\sin(t\log y)/t|\le\log y\le \sqrt y$ and $\cosh(u\log y)\le\sqrt y$. The absolute tail is at most
\[
          4\int_{10^4}^\infty y e^{-b\sqrt y}\,dy
          =8\int_{100}^\infty v^3e^{-bv}\,dv<1.
\]
As above the last inequality uses $b>19/100$ and $e^{-100b}<(5/6)^{100}$. Thus the imaginary part divided by $t$ has strict positive sign. Outside $|u|\le1/2$ the absolute Euler product on the right and the functional equation on the left exclude zeros. This proves the asserted region; simplicity at its only zero was already certified by~\eqref{ebc:certorder}.
\end{proof}

\begin{corollary}\label{ebc:fieldstrip}
For $E_{37}$ over the cubic field $F=\Q(\sqrt[3]{2})$, the sole real zero of its completed $L$-function is $s=1$, of multiplicity two. Over its $S_3$ splitting field $L$ that multiplicity is three and it is again the sole real zero. For each of the two fields the same description gives the complete list of zeros with $|\Im(s)|\le1/3$.
\end{corollary}
\begin{proof}
Theorem~\ref{ep:certificate} and Theorem~\ref{ebc:horizontal} give the required statements respectively for $E_{37}/\Q$ and its degree-four twist. For $F$ now use~\eqref{ebc:cubicfactor}. Over $L$ it remains to exclude the quadratic twist $L_-$. Its inverse Mellin kernel on $y\ge1$ is $f_-(y)=\sum c_n z^n$ with $z=e^{-y/A_-}<18/25$. Among $2\le n\le24$ the only negative coefficients are
\[
                c_7=-1,\quad c_{13}=-2,\quad c_{14}=-2,\quad
                c_{16}=-4,\quad c_{23}=-2 .
\]
These follow by the elliptic point count and Euler recurrence, twisting at $p\ne3$ by $\epsilon(p)$ and taking the trivial Euler factor at $3$. With $v=18/25$, the sum of all negative terms divided by $z$, in absolute value, is bounded by
\[
  v^6+2v^{12}+2v^{13}+4v^{15}+2v^{22}
      +2v^{24}\left(\frac{25}{1-v}+\frac{v}{(1-v)^2}\right)<\frac12 ,
\]
where $|c_n|\le2n$ was used from $n=25$ onward. Thus $f_-(y)>z/2>0$. Its even functional equation represents the completed function at $1+u+it$ as $2\int_1^\infty f_-(y)\cosh((u+it)\log y)\,dy$. In particular every real value is positive.

For $|t|\le1/3$ the cosine in the real part is nonnegative up to $y=100$. On $1\le y\le2$ we have $z>e^{-1}>1/3$, $\cos(t\log y)>1/2$, so the integral without the factor $2$ exceeds $1/12$ there. For $y\ge100$, $|f_-(y)|\le2z/(1-z)^2\le4e^{-y/3}$. With $|u|\le1/2$ the entire possible negative tail is bounded by $4\int_{100}^\infty y e^{-y/3}\,dy<1/12$. The real part is strictly positive. Outside this critical strip the Euler product and functional equation apply. This proves the assertion over $L$ by its positive factorization.
\end{proof}

\subsection{A tower with linearly growing central multiplicity}\label{ebc:kummer}
All the additional factors in a generalized dihedral extension can have negative sign at once. Let $k/\Q$ be imaginary quadratic with quadratic character $\epsilon_k$. Suppose $L/\Q$ is Galois with
\[
  \operatorname{Gal}(L/\Q)=A\rtimes\langle c\rangle,\qquad
  L^A=k,\qquad |A|\ \text{odd},
\]
where $A$ is abelian and complex conjugation $c$ inverts $A$. Put $F=L^{\langle c\rangle}$. If $E/\Q$ is non-CM, suppose that the conductor $N_E$ is prime to all primes ramified in $L/\Q$ and $\epsilon_k(N_E)=1$. We claim that
\begin{equation}\label{ebc:dihedralbound}
 \begin{split}
 \operatorname{ord}_{s=1}\Lambda(E/F,s)
  &\ge \operatorname{ord}_{s=1}\Lambda(E/\Q,s)
                                   +\frac{|A|-1}{2},\\
 \operatorname{ord}_{s=1}\Lambda(E/L,s)
  &\ge \operatorname{ord}_{s=1}\Lambda(E/k,s)+|A|-1 .
 \end{split}
\end{equation}
The two sides of each inequality have the same parity.

Indeed, the irreducibles of $G$ consist of $1$, $\epsilon_k$, and the degree-two representations $\rho_\chi=\operatorname{Ind}_A^G\chi$ for unordered pairs $\{\chi,\chi^{-1}\}$ of nontrivial characters of $A$. Odd order makes the two members distinct. Irreducibility follows by the index-two restriction test, and inequivalence by the restrictions $\chi\oplus\chi^{-1}$. Their squared dimensions sum to $2+4(|A|-1)/2=|G|$, proving completeness. Each $\rho_\chi$ is self-dual, has determinant $\epsilon_k$, and has one $c$-fixed line since $c$ interchanges its two character lines. Consequently it occurs once over $F$ and twice over $L$. The trivial character occurs in both, and $\epsilon_k$ occurs only over $L$.

Theorem~\ref{ebc:monomial} gives an entire cuspidal degree-four function for each $\rho_\chi$. The general coprime-conductor formula of \href{https://doi.org/10.1112/S0024611505015261}{Dokchitser}, Corollary~2, expresses the root number of a self-dual Artin twist as $w(E)^{\dim\rho}\det\rho(-N_E)$, interpreting its determinant as a quadratic or trivial Dirichlet character. Here this is $\epsilon_k(-N_E)=-1$ for every $\chi$, since $k$ is imaginary and $\epsilon_k(N_E)=1$. Each factor therefore contributes a positive odd central order. Summing these orders with their multiplicities proves~\eqref{ebc:dihedralbound} and the parity statements.

There is an elementary infinite family where the hypotheses can all be checked on the defining equations.

\begin{theorem}\label{ebc:tower}
Fix $E=E_{37}$ and any distinct positive primes $p_1,p_2,\ldots$ different from $3$ and $37$. Inside $\mathbb R$ choose their positive cube roots. For $m\ge1$ define
\[
 F_m=\Q(\sqrt[3]{p_1},\ldots,\sqrt[3]{p_m}),\qquad
 L_m=F_m(\sqrt{-3}).
\]
Then $[F_m:\Q]=3^m$, the two completed elliptic functions over $F_m,L_m$ are entire, and
\[
  \operatorname{ord}_{s=1}\Lambda(E/F_m,s)\ge\frac{3^m+1}{2},
  \qquad
  \operatorname{ord}_{s=1}\Lambda(E/L_m,s)\ge 3^m .
\]
Every nontrivial degree-four constituent furnishing these lower bounds is cuspidal over $\Q$.
\end{theorem}

\begin{proof}
The field $k=\Q(\sqrt{-3})$ contains the cube roots of unity. The $p_i$ have independent classes in $k^\times/k^{\times3}$. Indeed the valuation at a prime of $k$ over $p_i$ is $1$ on $p_i$ and $0$ on all other $p_j$ because $p_i\ne3$. Thus all exponents in a product representing a cube are divisible by $3$. Kummer theory gives $\operatorname{Gal}(L_m/k)\simeq(\Z/3\Z)^m$. On each chosen cube root a member of this group acts by multiplication by a cube root of unity. Complex conjugation fixes the positive roots and inverts the roots of unity, so it inverts the entire group. Also $F_m\subset\mathbb R$ and $L_m=F_mk$ show that $F_m$ is precisely the fixed field of conjugation. In particular $[F_m:\Q]=3^m$. The polynomials $X^3-p_i$ and the quadratic field $k$ are unramified outside $3\prod_i p_i$, so $L_m/\Q$ is unramified at $37$. Finally $37\equiv1\bmod3$ gives $\epsilon_k(37)=1$. The order over $\Q$ is at least one by the computed sign of $E_{37}$, and that over $k$ is at least one by the positive factorization over a quadratic extension. Formula~\eqref{ebc:dihedralbound} proves the asserted inequalities. The constituent construction proving that formula also proves entireness. This is an analytic statement about vanishing of $L$-functions and makes no assumption on the Birch--Swinnerton-Dyer conjecture.
\end{proof}

\subsection{Curves defined over the ground field}\label{ebc:ground}
Let $B$ be a number field and let $E/B$ be a curve for which modularity with local compatibility at every place is available. Its coefficients need not lie in $\Q$. Write $b=[B:\Q]$. The complete identity is
\[
       \Lambda(E/B,s+1/2)=c\,\Lambda(s,\pi_E),\qquad
       \pi_E\ \text{on } \mathrm{GL}_2(\mathbb A_B),
\]
where the nonzero constant depends on the convention for the gamma factors. This identity transfers the $T^2$ positivity bound of Theorem~\ref{li:everyheight} at every height, independently of $b$. It also transfers the positive-time result of Theorem~\ref{hg:threshold} with $d=2b$: $2b\,t\log X\ge1+\delta$ suffices at large $X$ for each fixed $\delta>0$. Indeed a real shift by $1/2$ identifies both height coordinates, and multiplication by a fixed scalar commutes with the heat operation. For $b\ge2$, Theorem~\ref{dw:windows} gives growing common height windows in which the arithmetic real parts lie in
\[
   [\,1/2+1/(2b)-\varepsilon,\;3/2-1/(2b)+\varepsilon\,]
\]
eventually for every fixed $\varepsilon>0$. For $b=1$ the shrinking interval of Theorem~\ref{common:main} gives the stronger rate.

For example the accepted theorem of Caraiani and Newton supplies this input for every elliptic curve over each of infinitely many imaginary quadratic fields, including $\Q(i)$ and $\Q(\sqrt{-d})$ for $d=2,3,5$. Their hypothesis is that $X_0(15)$ has rank zero over the field. In these examples $b=2$, so the limiting common-window interval in arithmetic coordinates is $[3/4,5/4]$. \href{https://spiral.imperial.ac.uk/handle/10044/1/147038}{A.~Caraiani and J.~Newton, \emph{On the modularity of elliptic curves over imaginary quadratic fields}.}

There are broader 2026 modularity announcements. The October 4 manuscript states the result, including local matching, for every curve over every imaginary quadratic field. A note contributed by B.~Le Hung gives extensions to totally real fields and CM fields lacking a primitive fifth root of unity. A further manuscript claims the remaining CM fields. Any application using the extra ground fields of these preprints takes the stated modularity theorem as a cited input. The implications in the preceding paragraph require only the completed identity and can be read independently of their review status. A CM \emph{ground field} here imposes no endomorphism hypothesis on the curve. Section~\ref{recentcm:section} instead uses extra endomorphisms on the curve itself.
\href{https://mathcommunity.github.io/oai-math-revised/papers/Modularity-of-elliptic-curves-over-imaginary-quadratic-fields-October-4-2026/paper.pdf}{\emph{Modularity of elliptic curves over imaginary quadratic fields}.}
\href{https://hexagonmath.org/2610.00156v1}{B.~Le Hung, October 8 note.}
\href{https://github.com/CaptainSude/elliptic-curves-cm-modularity/releases/tag/v1.0.1}{\emph{Every elliptic curve over a CM number field is modular}, v1.0.1.}

\subsection{Potential modularity and signed factorizations}\label{ebc:descent}
For comparison let $E$ now be defined over a number field $B$ and let $L/B$ be a finite Galois extension over which an automorphy theorem applies to $E$. Put $G=\operatorname{Gal}(L/B)$. A Brauer induction argument can express the trivial character as
\[
        1_G=\sum_j n_j\,\operatorname{Ind}_{H_j}^G\theta_j,
                  \qquad n_j\in\mathbb Z ,
\]
where the inducing groups may be chosen elementary and the $\theta_j$ are one-dimensional. Take $B_j=L^{H_j}$. Tensoring this identity by the elliptic representation and using induction gives
\begin{equation}\label{ebc:brauer}
     \Lambda(E/B,s)=\prod_j\Lambda(E/B_j,\theta_j,s)^{n_j}.
\end{equation}
Writing completed functions here presupposes the standard compatible archimedean factors and conductors. Suppose the particular potential-automorphy and descent results in use supply completed cuspidal functions for every twisted term on the right. It follows that the target has meromorphic continuation, and the functional equations can be multiplied with the indicated exponents. The hypothesis concerning the listed twists specifies the analytic data required; automorphy of a single function over $L$ supplies only that single function directly. Examples of potential automorphy theorems deriving meromorphic continuation in this manner are given in T.~Barnet-Lamb, T.~Gee, D.~Geraghty and R.~Taylor, \href{https://annals.math.princeton.edu/2014/179-2/p03}{\emph{Potential automorphy and change of weight}}.

At any point $s_0$ the conclusion from~\eqref{ebc:brauer} is exactly
\[
  \operatorname{ord}_{s_0}\Lambda(E/B)=
  \sum_{n_j>0} n_j\operatorname{ord}_{s_0}\Lambda(E/B_j,\theta_j)
  -\sum_{n_j<0} |n_j|\operatorname{ord}_{s_0}\Lambda(E/B_j,\theta_j).
\]
Holomorphy at $s_0$ requires that the first sum be at least the second. If every right-hand factor satisfies GRH, then both sums vanish away from $\Re(s)=1$ and the quotient is nonzero and holomorphic there. To assert an entire completion one must additionally exclude a negative order on the line itself. Conversely an entire quotient satisfying GRH need not force GRH for an individual factor: an off-line zero in a numerator can cancel against an off-line zero in a denominator. The positive product~\eqref{ebc:product} with entire factors avoids exactly this loss of information. Neither type of induction turns the Dirichlet zero-location assertion alone into a theorem for the additional elliptic cusp forms.

\clearpage
\section{Twists with prescribed small Euler factors and forced central zeros}\label{it:section}
The verified Dirichlet theorem can be used constructively even when the standard elliptic function has no Dirichlet factorization. We use it to bound the size of a prime discriminant having many prescribed local square classes. Its sign at infinity changes an elliptic root number, although all the Euler factors examined at small rational primes remain identical. Passing to a multiquadratic field makes the resulting discrepancy in central order proportional to the degree.

\begin{theorem}\label{it:main}
Let $E/\Q$ be a fixed elliptic curve of conductor $N$ with $L(E,1)\ne0$. Fix an integer $r\ge1$. For every sufficiently large $X$ there are distinct primes $q_1,\ldots,q_r$ and a multiquadratic field $K=\Q(\sqrt{-q_1},\ldots,\sqrt{-q_r})$ of degree $2^r$ with the following properties.
\begin{enumerate}
\item Every prime $p\le X$ splits completely in $K$, and all primes dividing $N$ do so as well. The Euler factor of $E/K$ over any such rational prime is the $2^r$th power of the Euler factor of $E/\Q$.
\item One has $\operatorname{ord}_{s=1}L(E/K,s)\ge2^{r-1}$. Each of the $r$ individual twists $E^{(-q_i)}$ shares every displayed small local factor with $E$ and has a central zero of odd order.
\item With $r,E$ fixed, the primes can be chosen with
\[
 \max_i \log q_i\le (\log 4+o(1))\frac{X}{\log X}\qquad(X\to\infty).
\]
\end{enumerate}
In particular one may take $E=11\mathrm a1$ by its positive real Mellin integral.
\end{theorem}

\begin{proof}
We give the uniform prime-selection argument. Put
$S=\{p\text{ odd}:p\le X\text{ or }p\mid N\}$,
$k=|S|$, and $M=8\prod_{p\in S}p$.
For a candidate prime $q\nmid M$ impose the congruence $q\equiv7\pmod8$ and for each $p\in S$ the single quadratic condition
\[
 (q/p)=(-1/p).
\]
Here the Legendre symbol is evaluated with modulus $p$; there is no requirement to choose a particular class modulo $p$. For $n$ prime to $M$ the indicator is the product of the residue-class indicator modulo $8$ and $\prod_{p\in S}(1+(-1/p)(n/p))/2$. Its expansion in characters modulo $M$ has coefficient sum in absolute value one and constant-character coefficient $\delta=1/(4\cdot2^k)$. Indeed the four terms modulo eight are pairwise distinct, and the toggled nontrivial characters at distinct odd primes are independent by the Chinese remainder theorem. Extend the indicator by zero to integers not prime to $M$.

Under the Dirichlet theorem the classical explicit formula gives, for every character of modulus $M$ and $Y\ge2$,
\[
 \sum_{n\le Y}\Lambda(n)\chi(n)
  =\mathbf1_{\chi\text{ principal}}Y+O(\sqrt Y\log^2(2MY)).
\]
The constant is absolute. This is the usual GRH uniform estimate for primitive characters with the imprimitive prime-power corrections included. Removing all proper prime powers costs $O(\sqrt Y\log^2(2Y))$ by the trivial count. Apply these estimates to the expansion and its absolute coefficient sum. The logarithmic sum over admissible primes up to $Y$ is
\begin{equation}\label{it:prime}
        \delta Y+O(\sqrt Y\log^2(2MY)).
\end{equation}
This use of an abelian theorem controls a single square class at every prime without the much smaller density $1/\varphi(M)$ of a specified simultaneous residue class.

For example choose $Y=4^k(\log(2M)+k+2)^6$. Then $\log(2MY)\ll\log(2M)+k+2$ and the main term in \eqref{it:prime} dominates its error for large $X$. Dividing by $\log Y$ shows that the number of admissible primes tends to infinity. Thus $r$ distinct ones can be chosen. Since $k\sim X/\log X$ and $\log M\sim X$, the asserted size bound follows. The prime number theorem used in this numerical simplification can equally be read from the modulus-one input.

Set $D_i=-q_i$. These are pairwise coprime odd fundamental discriminants with $D_i\equiv1\pmod8$. For each $p\in S$ we have $(D_i/p)=1$, so $p$ splits completely in $K$; the assertion at $2$ follows from the congruence modulo eight. The square classes of the $D_i$ are independent, as witnessed by the parity of the valuations at $q_i$. For each subset $J$ of $\{1,\ldots,r\}$ write $D_J=\prod_{i\in J}D_i$, with $D_\emptyset=1$. Every nonempty product is again an odd fundamental discriminant, has sign $(-1)^{|J|}$ and is prime to $N$.

The functional-equation sign for a quadratic twist of coprime conductor is $w(E^{(D)})=w(E)\chi_D(-N)$. This is the primitive quadratic-twist formula for a weight-two newform with trivial character. See M.~Radziwi\l\l\ and K.~Soundararajan, \emph{Moments and distribution of central $L$-values of quadratic twists of elliptic curves}, Invent. Math. \textbf{202} (2015), 1029--1068, introduction. Since the original central value is nonzero, $w(E)=1$. At each prime dividing $N$ we arranged $\chi_{D_J}(p)=1$. Also $\chi_{D_J}(-1)=\operatorname{sgn}(D_J)$. Every odd-sized subset therefore gives sign $-1$ and contributes at least one to the order of vanishing. Induction of the trivial character from $K$ to $\Q$ gives the positive product
\[
        L(E/K,s)=\prod_{J\subseteq\{1,\ldots,r\}} L(E^{(D_J)},s).
\]
It holds at bad primes as well by induction and tensor compatibility of the local factors; see Section~\ref{ebc:chapter}. All these factors are entire by modularity over $\Q$. There are exactly $2^{r-1}$ subsets of odd size. This proves the order assertion. At each specified prime every completion is $\Q_p$, proving the claim for the small Euler factors directly over $K$.
\end{proof}

The conductor of the degree-$2^{r+1}$ product over $\Q$ in this example is $N^{2^r}(\prod_i q_i)^{2^r}$, since a twist of coprime fundamental discriminant $D_J$ has conductor $ND_J^2$ and each $q_i$ occurs in $2^{r-1}$ subsets. In particular for the individual twists one can hide a change in central parity beyond an arbitrarily long prescribed Euler prefix with explicit conductor growth. The compared functions can both satisfy ERH: changing the order at the center makes no claim about off-axis zeros. This example singles out the global data a finite-prime transfer would still have to use.

\begin{corollary}[Growing rank and a growing Euler prefix]\label{it:family}
For the same fixed curve $E$, fix $c>0$ and let $r$ tend to infinity. One can take $X=c\log r\log\log r$ and still obtain fields from Theorem~\ref{it:main} for which, writing $\mathcal Q$ for the rational conductor of the entire product,
\[
 \operatorname{ord}_{s=1}L(E/K,s)
 \gg_{c,E}\frac{\log\mathcal Q}
 {\log\log\mathcal Q\,\log\log\log\mathcal Q}.
\]
Both the degree and the conductor vary in this assertion.
\end{corollary}
\begin{proof}
The prime number theorem gives
\[
 k=(c+o(1))\log r,\qquad \log M=O_c(\log r\log\log r).
\]
Choose any fixed $B$ larger than both $1+c\log 2$ and $2c\log 2$. Taking $Y=r^B$ in \eqref{it:prime}, the main term is $r^{B-c\log2+o(1)}$, which dominates its error $r^{B/2+o(1)}$ and also $r\log Y$. Hence there are at least $r$ distinct admissible primes below $Y$. The conductor formula gives $\log\mathcal Q\ll_{c,E}2^r r\log r$. Conversely $\log\mathcal Q\gg 2^r$ for large $r$, so $\log\log\mathcal Q\gg r$ and $\log\log\log\mathcal Q\gg\log r$. Combine with the central order lower bound $2^{r-1}$.
\end{proof}

\paragraph{A degree-eight example.}
For $E=11\mathrm a1$ and $X=19$, take $q_1=10559$, $q_2=15791$, $q_3=18191$. Trial division proves primality. For every odd prime $p\le19$ and each $i$, direct modular exponentiation gives $(-q_i)^{(p-1)/2}\equiv1\pmod p$. All three primes are $7\pmod8$. Thus every rational prime up to $19$ contributes eight identical local factors to $L(E/K,s)$, yet at the center four factors have odd order. The verifier \texttt{small\_euler\_prefix.py} reproduces these finite checks.

\clearpage
\section{A small prime distinguishing any two nonisogenous elliptic curves}\label{ld:section}

Let $E_1,E_2/\Q$ be nonisogenous over $\Q$. Let $S$ contain their primes of bad reduction and any other rational primes one wants to exclude, and put $R_S=\prod_{p\in S}p$. The $7/8$ theorem for finite-order Hecke characters over number fields implies that there is a rational prime $p\notin S$ with
\begin{equation}
 a_p(E_1)\ne a_p(E_2),\qquad p\ll [\log(2R_S)]^8 .
 \label{ld:main}
\end{equation}
The implied constant is absolute. Thus this bound holds uniformly as both conductors vary, with $R_S$ the radical of their product when only bad primes are excluded. It applies equally to curves with complex multiplication. The prime witnesses different point counts over $\F_p$. No zero-location hypothesis on either elliptic function enters the argument.

We use the finite-order, rank-one part of
\href{https://hexagonmath.org/2610.00216v2}{B.~Hu, \emph{A seven-eighths zero-free half-plane for unitary Weil $L$-functions}, version 2}; its statement and formal-source audit appear in Section~\ref{recentcm:section}. The only analytic information taken from that result is zero-freeness at all heights in $\Re s>7/8$, for every character including the trivial one away from its pole. All uniform dependence on the conductor will follow from the classical smoothed explicit formula. More generally write $\theta$ for a common boundary, $1/2\le\theta<1$. The exponent in \eqref{ld:main} is then $1/(1-\theta)$. Taking the separate November endpoint of Liu quoted in Section~\ref{recentcm:section} as input gives $7.89469346\ldots$ in place of $8$.

The finite-representation method goes back to Faltings and Serre. Under GRH, Serre's original \href{https://www.numdam.org/item/PMIHES_1981__54__123_0/}{1981 paper, \S8.3, Theorem 21}, has a further iterated-logarithm factor; the pure log-square conclusion under GRH is given by
\href{https://msp.org/pjm/2019/301-1/pjm-v301-n1-p07-p.pdf}{S.~Gun and V.~Kumar Murty, \emph{Lifting of elliptic curves}, Proposition 3.1}, citing Serre's later note on Faltings' argument. We give the reduction to a fixed-size finite quotient and show exactly which Hecke functions suffice for exponent eight. In a quadratic-twist and rational-isogeny case, Dirichlet functions alone suffice for exponent two.

\subsection{A smoothed least-prime bound in bounded degree}

Fix a nonnegative $w\in C_c^\infty((1,2))$ with $c_w=\int w(u)du>0$, and write $W(s)=\int_0^\infty w(u)u^{s-1}du$. Let $F$ be a number field of degree at most $m$, and $\psi$ a primitive finite-order Hecke character. At primes dividing its finite conductor $\mathfrak f_\psi$ its value is zero. Set $A_\psi=|d_F|\mathrm N\mathfrak f_\psi$. These functions have unitary center $1/2$, with no twist of the complex variable by the weight of an elliptic curve. If their nontrivial zeros have real part at most $\theta$, then
\begin{equation}
 \sum_{\mathfrak p,j\ge1}\psi(\mathfrak p)^j\log\mathrm N\mathfrak p\,
    w((\mathrm N\mathfrak p)^j/x)
  =\mathbf1_{\psi=1}c_w x+O_{m,w}(x^\theta\log(2A_\psi)),\quad x\ge2 .
 \label{ld:smooth}
\end{equation}
The sum ranges over prime ideals and positive powers. To check uniformity, the completed factors at infinity are products of $\Gamma_\R(s+a)$ with $a=0,1$; every complex place contributes both factors up to a constant. The conductor factor is $A_\psi^{s/2}$. The classical unit-interval zero count, including multiplicities, is
\[
  \#\{\rho:|\Im\rho-t|\le1\}\ll_m\log(2A_\psi)+m\log(3+|t|).
\]
For example, use the real part of the logarithmic derivative of the completed Hadamard product at $2+it$, after removing the principal poles; the Euler logarithmic derivative is $O(m)$ there and Stirling's formula gives the displayed conductor and height terms. The summands from zeros majorize a constant times $1/(1+(t-\Im\rho)^2)$. These are the standard Hecke explicit-formula estimates; see \href{https://bookstore.ams.org/coll-53}{Iwaniec and Kowalski, \emph{Analytic Number Theory}, Chapter~5}.

In Mellin inversion for $-L'/L$, shift from $\Re s=2$ to $-1/2$. The pole at one has residue contribution $c_wx$ for the trivial character. The nontrivial zeros contribute $-\sum_\rho x^\rho W(\rho)$, of size $O_{m,w}(x^\theta\log(2A_\psi))$ by the unit-interval count and rapid decay of $W$ in bounded strips. The possible zero of the uncompleted function at zero costs $O_{m,w}(1)$ by its gamma factors and pole criterion. On $\Re s=-1/2$ the functional equation, gamma derivatives and the absolutely convergent dual Euler series at $3/2$ give $L'/L\ll_m\log(2A_\psi)+m\log(3+|t|)$ in modulus. The integral on this line is $O_{m,w}(x^{-1/2}\log(2A_\psi))$. For the contour shift choose horizontal heights in successive unit intervals separated from all zero ordinates by $\gg_m 1/(\log(2A_\psi)+m\log(3+|t|))$, using the unit-interval count. The local partial-fraction formula gives a polynomial bound in that logarithm there, and the horizontal integrals tend to zero by rapid decay. This proves \eqref{ld:smooth} without a factor $\log x$ at the endpoint.

\begin{lemma}\label{ld:cheb}
Let $M/\Q$ be a finite Galois extension of degree $n\le m$, unramified outside a given finite set $S$. For every conjugacy class $C$ of $G=\operatorname{Gal}(M/\Q)$ there exists a $p\notin S$ with $\operatorname{Frob}_p\in C$ and
\[
        p\ll_{m,\theta} [\log(2R_S)]^{1/(1-\theta)} .
\]
It suffices to have the stated zero boundary for the finite-order Hecke functions of $M^{\langle g\rangle}$ belonging to the cyclic extension $M/M^{\langle g\rangle}$, for one $g\in C$.
\end{lemma}
\begin{proof}
Put $H=\langle g\rangle$, $F=M^H$, $z_g=|C_G(g)|$. Induction of class functions gives
\[
 \mathbf1_C=\frac{1}{z_g}\sum_{\chi\in\widehat H}\overline{\chi(g)}
                          \operatorname{Ind}_H^G\chi .
\]
At an element conjugate to $g$, the sum of induced functions before division equals $z_g$, as is seen by expanding the delta function of $g$ on the cyclic group. At all other elements it is zero. Class field theory identifies the functions of the characters $\chi$ over $F$ with primitive Hecke functions. Their conductors satisfy
$|d_F|\mathrm N\mathfrak f_\chi\le |d_M|$: the relative conductor-discriminant formula is
$|d_M|=|d_F|^{|H|}\prod_{\chi\in\widehat H}\mathrm N\mathfrak f_\chi$.

For $p\notin S$, Artin induction of the local factors identifies the coefficient at $p^j$ of the logarithmic derivative for $\chi/F$ with
$\log p\,\operatorname{tr}(\operatorname{Ind}_H^G\chi)(\operatorname{Frob}_p^j)$. Equivalently a cycle of length $f$ in the action on $G/H$ corresponds to a prime ideal of residue degree $f$, contributing $f$ times the associated character value to the trace of the $j$th power when $f\mid j$. For each $p\in S$, discarding terms over $p$ costs $O_{m,w}(\log p)$, since for each prime above it at most two successive powers of its norm fall in $[x,2x]$ and the sum of its residue degrees is at most $m$. Higher powers of rational primes away from $S$ cost $O_{m,w}(\sqrt x)$ in total: bound $\sum_{p^j\le2x,j\ge2}\log p$ by the Chebyshev estimate for squares and $O(x^{1/3}\log(2x))$ for larger powers. Applying \eqref{ld:smooth} to the displayed induced combination now gives
\[
 \sum_{\substack{p\notin S\\\operatorname{Frob}_p\in C}}\log p\,w(p/x)
  =\frac{c_w x}{z_g}
   +O_{m,w}(x^\theta\log(2|d_M|)+\sqrt x+\log(2R_S)).
\]

For this bounded-degree Galois extension,
$\log |d_M|\le n\log R_S+n\log n$.
Locally at $p$ the different exponent of a completion relative to its maximal unramified subfield is at most $e-1+e v_p(e)$. Indeed a uniformizer has an Eisenstein minimal polynomial of degree $e$, and the leading derivative term has that valuation. Other nonzero derivative terms have different valuations modulo $e$, so cancellation cannot raise the minimum valuation. Summing with residue degrees gives $v_p(d_M)\le n(1+v_p(e))$, with $e\mid n$ by Galoisness. Sum over the ramified primes to get the asserted global bound. As $z_g\le m$, take $x=C_{m,w,\theta}[\log(2R_S)]^{1/(1-\theta)}$, increasing the constant so the displayed prime sum is positive. Its support gives $p<2x$.
\end{proof}

\subsection{Detecting a trace difference across a deep congruence}

Fix a finite extension of $\Q_\ell$ with ring of integers $\mathcal O$, uniformizer $\varpi$ and residue field of size $b$. Suppose $\rho_i$ are continuous integral representations of $G_\Q$ of ranks $d_i$ unramified outside $S$. If their characteristic-zero trace functions differ, put
\[
 \mathcal A=\mathcal O\text{-span}\{(\rho_1(\gamma),\rho_2(\gamma)):
                                      \gamma\in G_\Q\}
 \ \subseteq M_{d_1}(\mathcal O)\oplus M_{d_2}(\mathcal O).
\]
It is a unital algebra since the generators are closed under multiplication. As a submodule of a finite free module over a discrete valuation ring it is free of rank $r\le d_1^2+d_2^2$. The linear functional $T(a,b')=\operatorname{tr}a-\operatorname{tr}b'$ has a nonzero image ideal $\varpi^e\mathcal O$. Divide it by $\varpi^e$ on $\mathcal A$ and reduce modulo $\varpi$. This produces a nonzero linear functional $\bar T:\mathcal A/\varpi\mathcal A\to\mathcal O/\varpi$; at least one image element of $G_\Q$ has nonzero value because those elements span $\mathcal A$. The finite image
\[
 G=\operatorname{im}\bigl(G_\Q\longrightarrow
                          (\mathcal A/\varpi\mathcal A)^\times\bigr)
                  \quad\hbox{has order at most } b^{d_1^2+d_2^2}.
\]
This map is continuous: the compact finite free module $\mathcal A$ has its induced topology from the ambient matrices, and $\varpi\mathcal A$ is open within it. The quotient uses $\varpi\mathcal A$ itself. Entrywise reduction can lose a trace difference divisible by a high power of $\varpi$ even when the normalized functional detects that difference. The extension cut out by $G$ is unramified outside $S$. Trace is conjugacy-invariant, so the entire class of the chosen element has nonzero value under $\bar T$, ensuring an actual characteristic-zero trace difference for every lift.

For a concrete algebra calculation take the closed group generated by the two scalar entries $(1,1+2^v)$, with $v\ge1$. Its $\Z_2$-algebra has basis $I=(1,1)$ and $B=(0,2^v)$ with $B^2=2^vB$. The generator has identity image under entrywise reduction, but becomes $I+B$ with $B\ne0$ and $B^2=0$ in $\mathcal A/2\mathcal A$. Its normalized trace difference is $-1$. Arbitrarily deep entrywise congruences can thus persist at a fixed quotient size.

Apply Lemma~\ref{ld:cheb}. This yields a detecting prime of size
$\ll_{b,d_1,d_2,\theta}[\log(2R_S)]^{1/(1-\theta)}$.
In characteristic zero unequal semisimple representations have unequal traces by Brauer--Nesbitt. For \eqref{ld:main} use the integral $2$-adic Tate modules. They have $d_1=d_2=2$, $b=2$, so $|G|\le256$. They are unramified at good primes other than $2$, and the arithmetic Frobenius traces there are $a_p(E_i)$. By Faltings' isogeny theorem and semisimplicity the trace functions differ when the curves are not $\Q$-isogenous. Adjoining $2$ to $S$ changes the logarithmic bound by an absolute factor. This completes the proof of \eqref{ld:main}. The same fixed-rank assertion applies to two characteristic-zero newform Galois representations with coefficients in a fixed local field, outside its residue characteristic and their levels. Their arithmetic Frobenius traces are the Hecke eigenvalues in their chosen embeddings.

For representations of $G_{K_0}$ for a number field $K_0$, form the same finite extension $M/K_0$ and its Galois closure $M'/\Q$. One may take $[M':\Q]\le ([K_0:\Q]b^{d_1^2+d_2^2})!$. In $R_S$ now include the rational primes dividing $d_{K_0}$ and those below the excluded prime ideals. The closure is unramified outside this set, since finite unramified local extensions remain unramified in a compositum. Lift the detecting element to an element of $\operatorname{Gal}(M'/K_0)$, apply Lemma~\ref{ld:cheb} to its rational conjugacy class, and choose a prime of $M'$ giving exactly that lift as Frobenius. Its restriction to $K_0$ has residue degree one: the $p$th power acts trivially on its residue field because the chosen element fixes $K_0$ pointwise. The relative Frobenius over this prime thus gives the desired trace difference. The bound on its norm has the same exponent, with a constant depending also on $[K_0:\Q]$.

\subsection{Quadratic twists seen in one fixed division field}

For a more specific detecting prime, fix $E/\Q$ of conductor $N$, put $K=\Q(E[3])$, $n=[K:\Q]\le48$, and let $D$ be a nontrivial fundamental discriminant. Suppose $\Q(\sqrt D)\not\subset K$, a condition holding for all sufficiently large $|D|$ for each fixed $E$. There exists a good $p\nmid 3ND$, split completely in $K$ and inert in $\Q(\sqrt D)$, satisfying
\[
  p\ll [\log(2\operatorname{rad}(3N|D|))]^8,
  \qquad a_p(E)\equiv2\pmod3,\quad
  a_p(E^{(D)})=-a_p(E) .
\]
Here $\operatorname{rad}(n)$ is the product of the prime divisors. The traces have opposite signs, a stronger condition than unequal traces. The bound with an absolute constant follows already by applying Lemma~\ref{ld:cheb} to the class $(1,-1)$ in
$\operatorname{Gal}(K\Q(\sqrt D)/\Q)
=\operatorname{Gal}(K/\Q)\times\{\pm1\}$, over the cyclic subgroup $\{1\}\times\{\pm1\}$.

One can see the fixed-field uniformity directly. Let $\psi_D$ be the primitive quadratic Hecke character of $K$ for $K(\sqrt D)/K$. The relative discriminant divides $D\mathcal O_K$: the order obtained from $\mathcal O_{\Q(\sqrt D)}$ by extension of scalars has relative discriminant $D\mathcal O_K$, and passage to the maximal order divides its discriminant by a squared index ideal. Thus $A_{\psi_D}\le |d_K||D|^n$. Subtract \eqref{ld:smooth} for $\psi_D$ from that formula for $\zeta_K$. Away from $3ND$, degree-one prime ideals occur precisely above rational primes split in $K$. There are $n$ such ideals over each, and the difference in the two first-power coefficients is $1-\chi_D(p)$, with $\chi_D$ the primitive Dirichlet character of $\Q(\sqrt D)$. Terms of residue degree or power greater than one in the smoothed difference cost $O_n(\sqrt x)$ by the rational-prime-power bound in Lemma~\ref{ld:cheb}. Omitting primes over $3ND$ costs $O_n(\log(2\operatorname{rad}(3N|D|)))$. Consequently
\begin{multline*}
  2n\sum_{\substack{p\nmid3ND,\ p\ {\rm split\ in}\ K\\
                     \chi_D(p)=-1}}\log p\, w(p/x)\ \ =\ \ c_wx\\
 +O_{n,w}\bigl(x^\theta\log(2|d_K||D|^n)+\sqrt x+
                     \log(2\operatorname{rad}(3N|D|))\bigr).
\end{multline*}
At a totally split prime, arithmetic Frobenius acts as the identity on $E[3]$, giving the trace congruence. At a good prime away from $D$ an elliptic quadratic twist multiplies the trace by $\chi_D(p)$. Both conclusions now follow. For two twists with discriminants $D_1,D_2$, the same argument uses the primitive character inducing $\chi_{D_1}\chi_{D_2}$, while excluding all primes dividing $3ND_1D_2$. The field-disjointness condition concerns the resulting quadratic field. Irrespective of that condition, \eqref{ld:main} still applies whenever the two twists are not $\Q$-isogenous.

\subsection{A square bound using Dirichlet functions alone}

If an extension $M/\Q$ above is abelian, expand $\mathbf1_{\{g\}}$ directly in the characters of the full group over $\Q$. The same smoothed calculation then requires only Dirichlet functions, including $\zeta$. Under the Dirichlet input of this manuscript we can take $\theta=1/2$ and obtain $p\ll_m[\log(2R_S)]^2$ for each element of its group. This does not require any Hecke zero statement over another field.

Suppose $E$ has a rational isogeny of odd prime degree $\ell$. Write $\lambda,\mu:G_\Q\to\F_\ell^\times$ for the characters on the stable line of $E[\ell]$ and its quotient, and write $\nu=\lambda/\mu$. At good $p\ne\ell$ its trace is zero modulo $\ell$ precisely when $\nu(\operatorname{Frob}_p)=-1$. Regard the signs of a nontrivial quadratic character $\chi_D$ also in $\F_\ell^\times$. Unless $\chi_D=\nu$, there is an element $g$ with $\chi_D(g)=-1$ and $\nu(g)\ne-1$. Indeed, if $\nu$ were $-1$ on this whole coset, multiplying any one such element by $\ker\chi_D$ would force $\nu$ to be trivial on that kernel, hence equal to $\chi_D$. The extension cut out by $(\nu,\chi_D)$ is abelian of degree at most $2(\ell-1)$ and is unramified outside $\ell ND$. Consequently a prime of opposite nonzero traces for $E$ and $E^{(D)}$ exists with
\[
  p\nmid\ell ND,\qquad
  p\ll_{\ell}[\log(2\operatorname{rad}(\ell N|D|))]^2.
\]
Only one quadratic character can be exceptional for this argument. For $\ell=3$ we have $\mu^{-1}=\mu$, and $\lambda\mu$ is the mod-three cyclotomic character by the Weil pairing. Thus the sole possible exception to this argument is $D=-3$. Equality $\chi_D=\nu$ is a limitation of this residual trace test, and does not assert equality of the two characteristic-zero traces.

If $E$ has complex multiplication by an imaginary quadratic field $M_0$ and $\Q(\sqrt D)\ne M_0$, choose instead the split/inert element in the biquadratic field $M_0\Q(\sqrt D)$. Deuring's reduction theorem makes a good split prime ordinary, so its trace is nonzero. The same square bound holds, excluding primes dividing $2N d_{M_0}D$, with an absolute implied constant. For the reduction theorem see S.~Lang, \emph{Elliptic Functions}, Chapter~13, Theorem~12.

The exponent supplied by smoothing and the given strip is $1/(1-\theta)$; this computation does not prove optimality for actual curves. In the opposite direction, the twists of Section~\ref{it:section} of a fixed curve with nonzero central value agree through $X$ and have different central parity, hence cannot be $\Q$-isogenous to the original curve. For an unbounded sequence of their discriminants, that construction gives
\[
  p_{\rm first}>X\ge(1/\log4-o(1))\log |D|\log\log |D|.
\]
Indeed $\log|D|\le(\log4+o(1))X/\log X$ and $\log\log|D|\le(1+o(1))\log X$. This lower example and the upper bounds concern good rational primes and the same finite-prefix notion.

\clearpage
\section{Automorphic induction and a cuspidal adjoint with GRH}\label{autx:chapter}
The transfer theorems here assume the critical-line assertion for primitive Dirichlet functions whenever a conclusion about their zeros is invoked. The factorization identities and the converses about Euler factors are unconditional. We use unitary normalization for standard automorphic $L$-functions, so the expected symmetry line is $\operatorname{Re}s=1/2$. Three questions arise even for a single newform. Does its standard function reduce to Dirichlet factors? Can its adjoint reduce to them when the standard function does not? Which Ramanujan--Petersson weights have actually been used in making the comparison? A family of holomorphic weight-one newforms answers the second question affirmatively, with an exact answer to the first. Weight two behaves differently even for CM elliptic curves.

\subsection*{Unitary degree one, isobaric sums, and a converse}
A finite-order rational idele class character gives its primitive Dirichlet function. For signs in vertical shifts, write $\omega_{\chi,t}$ for the unitary rational character whose coefficient at an unramified prime $p$ is $\chi(p)p^{it}$, $t\in\R$. Then $L(s,\omega_{\chi,t})=L(s-it,\chi)$. Every unitary character of $\Q^\times\backslash\mathbb A_\Q^\times$ is of this form. In fact this group is isomorphic to $\R_{>0}\times\widehat{\Z}^\times$: make the valuations of the finite idele zero by multiplying by a positive rational, then arrange a positive real component using $-1$. A continuous unitary character on the compact profinite factor has finite image and comes from some $U_m$, since a profinite quotient has no connected circle quotient. On the positive reals the characters are imaginary powers. The parity of the Dirichlet character also specifies the sign character of the original real idele.

Form an isobaric sum $\Pi=\boxplus_{j=1}^d\omega_{\chi_j,t_j}$ on $\mathrm{GL}_d(\mathbb A_\Q)$. At each place its standard parameter is the direct sum of the corresponding local characters. Multiplicativity of the local factors yields
\[
       L(s,\Pi)=\prod_{j=1}^d L(s-it_j,\chi_j).
\]
This is an automorphic ERH statement in every degree $d$: every zero in the open strip belongs to the central line. A trivial character in the product supplies a possible pole at $1+it_j$; such a pole is outside the open strip. If a second sum has characters $\omega_{\psi_k,u_k}$, then
\[
 L(s,\Pi\times\widetilde\Pi')
   =\prod_{j,k} L\bigl(s-i(t_j-u_k),(\chi_j\bar\psi_k)_{\rm prim}\bigr).
\]
Primitivizing after multiplication retains factors even where one or both original characters were ramified. More generally every tensor product, dual, exterior or symmetric power of these sums has a global parameter again of this form, since the weights of a tensor construction on diagonal matrices are monomials. Its automorphic realization is the isobaric sum of the corresponding rational characters. Thus these conclusions use no unproved general functorial transfer.

There is a corresponding number-field assertion. Let $K/\Q$ be abelian. Every norm pullback $\omega_{\chi,t}\circ\mathrm N_{K/\Q}$ has $L$-function $\prod_{\eta\in\widehat{\operatorname{Gal}(K/\Q)}} L(s-it,(\chi\eta)_{\rm prim})$, by the induction and reciprocity statements for Hecke characters. All the isobaric constructions just described over $K$ remain products of such rational factors. The analogous claim for an arbitrary $K$ would already include the unknown broad Dedekind hypothesis when $\chi=1$ and $t=0$.

\begin{theorem}[When a standard automorphic function can reduce in this way]\label{autx:ismo}
Let $\pi$ be an isobaric sum of unitary cuspidal automorphic representations of $\mathrm{GL}_{n_i}(\mathbb A_\Q)$, including the possibility that $\pi$ itself is cuspidal. Suppose its finite standard Euler product agrees, outside a finite set, with a finite product
$\prod_{j} L(s-it_j,\chi_j)^{b_j}$ with $t_j$ real and $b_j\in\Z$. Then every cuspidal summand of $\pi$ has degree one. After equal characters are collected, the exponents are precisely their nonnegative multiplicities. In particular a standard $L$-function of a cuspidal representation of degree at least two admits no such identity.
\end{theorem}
\begin{proof}
Move all factors with negative exponents to the left and take their isobaric sum with $\pi$. The right is the isobaric sum of the positive factors. At all unramified places outside the finite set the equality of polynomials implies equality of degrees. Those degrees are constant there, so the two representations have the same rank. The unramified Satake classes are semisimple and their characteristic polynomials determine them up to conjugacy. Thus the local representations agree almost everywhere. Equality of the products as functions in a half-plane of absolute convergence gives this same local equality by uniqueness of the ordinary Dirichlet series: the coefficients at powers of $p$ recover the local power series recursively. The strong multiplicity-one theorem for isobaric representations identifies the multisets of cuspidal summands on both sides, proving the assertion. The one-dimensional characters $\omega_{\chi,t}$ and $\omega_{\psi,u}$ agree only if $t=u$ and $\chi=\psi$, by the idele class description above.
\end{proof}
We have used the uniqueness of isobaric support and the multiplicativity of standard and pair factors. See A.~Booker and M.~Krishnamurthy, \href{https://research-information.bris.ac.uk/ws/files/68371958/gln.pdf}{\emph{A converse theorem for $\mathrm{GL}(n)$}}, Remark~1, for the precise statements and references to Jacquet--Shalika. The theorem concerns equality of functions, not an implication between their conjectural zero loci. In particular it explains why modularity of a rational elliptic curve, which supplies a cuspidal form of degree two, does not itself give a factorization proving ERH.

For finite-image families one can detect the same obstruction even if arbitrary different unitary shifts appear on both sides, without constructing an automorphic representation. Let $v_t$ be finitely many complex class functions on a fixed finite Galois group, indexed by distinct real numbers $t$. If $\sum_t p^{it}v_t(\mathrm{Frob}_p)=0$ at every unramified prime except finitely many, then each $v_t$ is zero. To prove it, for a fixed class $g$ form $P_g(y)=\sum_t e^{ity}v_t(g)$. For every fixed $\delta>0$, Chebotarev's asymptotic for this class and the prime number theorem supply a prime in the class in every interval $(e^y,e^{y+\delta})$ for all sufficiently large $y$. The derivative of $P_g$ is bounded, and the function vanishes at the logarithm of that prime. Hence $\limsup_{y\to\infty}|P_g(y)|\le\delta\sup|P_g'|$, proving convergence to zero. Averaging $P_g(y)e^{-ity}$ on $[0,Y]$ and letting $Y$ tend to infinity extracts $v_t(g)$ by orthogonality of distinct real frequencies. Thus allowing vertical shifts cannot conceal a nonlinear Artin constituent in a Dirichlet identity. This argument only uses prime traces.

\subsection*{A degree-two cusp form whose adjoint splits}
Let $K/\Q$ be quadratic, let $c$ be its nonidentity element, and let $\mu,\nu$ be unitary Hecke characters of $K$. Write $I(\mu)=\operatorname{AI}_{K/\Q}(\mu)$ for quadratic automorphic induction. We use the classical Hecke--Maass theta correspondence in this case. Its standard function is $L_K(s,\mu)$; the induction is cuspidal exactly when $\mu^c\ne\mu$. On local Weil parameters it is induction from $K$ to $\Q$, with the direct-sum convention at split places. The factors of $L(s,I(\mu)\times\widetilde {I(\nu)})$ are the pair factors for the two local parameters. For quadratic induction we use the classical theta correspondence cited at the end of this section. For the pair theory see J.~Cogdell, \href{https://people.math.osu.edu/cogdell.1/ictp-www.pdf}{\emph{Lectures on $L$-functions, converse theorems, and functoriality for $\mathrm{GL}_n$}}, for the pair theory.

\begin{lemma}[The two quadratic-induction summands]\label{autx:pair}
With $\delta_\mu=\mu/\mu^c$ and $\chi_K$ the quadratic Dirichlet character of $K$, the full finite factors satisfy
\begin{align*}
 L(s,I(\mu)\times\widetilde{I(\nu)})
      &= L_K(s,\mu/\nu)\,L_K(s,\mu/\nu^c),\\
 L(s,I(\mu)\times\widetilde{I(\mu)})
      &= \zeta(s)L(s,\chi_K)L_K(s,\delta_\mu),\\
 L(s,\operatorname{Ad} I(\mu))
      &= L(s,\chi_K)L_K(s,\delta_\mu).
\end{align*}
The identities extend to the archimedean factors when gamma factors on both sides are chosen by the same local-parameter convention.
\end{lemma}
\begin{proof}
For a group inclusion of index two, write $J(\mu)$ for its induced parameter. Restriction of $J(\nu)^\vee$ to the subgroup gives $\nu^{-1}\oplus(\nu^c)^{-1}$. The tensor identity for induction therefore gives
$J(\mu)\otimes J(\nu)^\vee=J(\mu/\nu)\oplus J(\mu/\nu^c)$. At a split prime this says explicitly that the four tensor characters are the two products at corresponding places and the two crossed products. At an inert or ramified place it follows from the same index-two identity on the Weil group; taking the determinant on inertia invariants gives the full factor. The one-dimensional Weil characters have zero monodromy operator, and induction and these sums preserve that property. Local inductivity of standard factors, and compatibility with pair factors, prove the first identity. For $\nu=\mu$ the first term is $J(1)=1\oplus\chi_K$. Removing the scalar identity inside the endomorphism representation gives the trace-zero adjoint formula. The identical decomposition on the real Weil group identifies the gamma factors by their multiplicativity and inductivity.
\end{proof}

For a finite-order character $\mu$ with $I(\mu)$ cuspidal, its entire adjoint reduces to rational Dirichlet factors if and only if $\delta_\mu$ is quadratic. Indeed $\delta_\mu^c=\delta_\mu^{-1}$. If $\delta_\mu^2\ne1$, its induction is a degree-two cusp form, so the converse in Theorem~\ref{autx:ismo} precludes a reduction for the adjoint. If $\delta_\mu^2=1$, cuspidality first excludes $\delta_\mu=1$. The remaining invariant quadratic character descends along the cyclic extension by the invariant-character argument (equivalently, its class field is a central extension of a cyclic group, hence abelian). It extends to a rational character $\theta$, with other extension $\theta\chi_K$. Then $L(s,\operatorname{Ad}I(\mu))=L(s,\chi_K)L(s,\theta)L(s,\theta\chi_K)$ with primitive conductors understood. In this setting the two extensions are in fact quadratic: if $g$ is any lift of the nontrivial coset then $\delta_\mu(g^2)=1$, since $g$ centralizes its own square. It follows that $\theta(g)^2=1$, while its values on the subgroup are already real. The three quadratic characters are distinct. Replacing $\mu$ by a rational unitary norm twist preserves the same adjoint.

For two finite-order characters with cuspidal inductions from the same $K$, the complete mixed product in the lemma reduces to rational characters exactly when $\delta_\mu=\delta_\nu$ and this common character is nontrivial quadratic. In fact the two inducing characters on its right are invariant precisely when $\delta_\mu/\delta_\nu=1$ and $\delta_\mu\delta_\nu=1$, respectively. When either fails the corresponding induction is cuspidal and survives in the isobaric sum by multiplicity-one. When both hold, the ratio $\mu/\nu$ descends, so the two cusp forms are rational twists of one another. This proves a criterion including necessity, rather than just a list of examples.

The condition changes for CM elliptic curves. For a CM weight-two newform over $\Q$, the unitarized inducing character at the complex place has angular character $(z/|z|)^{\pm1}$. Thus $\delta_\mu$ has angular character $(z/|z|)^{\pm2}$, of infinite order. Its induction in Lemma~\ref{autx:pair} remains cuspidal. In particular the accompanying quadratic Dirichlet factor of the adjoint is controlled by Dirichlet GRH, while the degree-two factor still requires a zero theorem for that Hecke character of the imaginary quadratic field. A general elliptic-curve ERH cannot be concluded by applying the finite-order argument to the algebraic Hecke character.

\subsection*{A weight-one newform of level 256}
We now give an example with its conductor, local factors, and actual Fourier coefficients. Set $N=\Q(\alpha,i)$, $\alpha^4=2$ with $\alpha>0$, and use the rotations $r(\alpha)=i\alpha$, $r(i)=i$ and a reflection $s$ given by complex conjugation. The Galois group is dihedral of order eight. The kernel of the action on $K=\Q(i)$ is the rotation subgroup $C_4$. Choose its character $\lambda(r)=i$ and let $\mu$ be the corresponding primitive finite-order Hecke character of $K$. Then $\mu^c=\mu^{-1}\ne\mu$, and the corresponding cusp form is
\[
 f(z)=\sum_{\substack{0\ne\mathfrak a\subseteq\mathcal O_K\\
                    (\mathfrak a,\mathfrak f_\mu)=1}}
              \mu(\mathfrak a)\,e^{2\pi i(\mathrm N\mathfrak a)z}.
\]
Hecke's theta-series theorem gives a holomorphic newform of weight one, normalized by its coefficient of $e^{2\pi iz}$. Its Artin representation $V=\operatorname{Ind}_{C_4}^{G}\lambda$ has eigenvalues $i,-i$ on $r$ and $1,-1$ on $s$. Its determinant is the quadratic character of $K$, so the nebentype is $\chi_{-4}$. The weight means that the classical $L(s,f)$ is already in unitary normalization.

Here is a way to check the level without computing a Gauss sum at $2$. Put $P=\Q(\alpha)$; its field discriminant is $-2^{11}$. Indeed the polynomial discriminant is $-2^{11}$, and at every odd prime the index of $\Z[\alpha]$ is therefore a unit. Over $\Q_2$ the polynomial is Eisenstein, so $\alpha$ is a uniformizer in a totally ramified degree-four extension. The valuations of nonzero $a_j\alpha^j$ for $a_j\in\Q_2$, $0\le j<4$, are distinct modulo four in the integral normalization. Hence integrality of their sum forces each $a_j$ to belong to $\Z_2$, proving the local index at $2$ is also one. The permutation representation on $G/\langle s\rangle$ of the embeddings of $P$ is $1\oplus\chi_8\oplus V$. Its linear constituents are the characters trivial at $s$ and its remaining irreducible is $V$ with multiplicity one. The permutation conductor--discriminant formula and additivity of local conductors give $\mathfrak f(V)=2^{11}/8=256$. The conductor formula for induction from $K$ consequently gives $\mathrm N\mathfrak f_\mu=256/4=64$. Thus $f$ lies in $S_1(\Gamma_0(256),\chi_{-4})$ with exact level $256$. The standard newform and finite Artin conductor agree by the theta correspondence. These invariants also identify the database objects \href{https://www.lmfdb.org/ModularForm/GL2/Q/holomorphic/256/1/c/a/}{$256.1.c.a$} and \href{https://www.lmfdb.org/ArtinRepresentation/2.256.4t3.c.a}{$2.256.4\mathrm{t}3.\mathrm{c}.\mathrm{a}$}.

For odd $p$, put $a_p=\operatorname{tr} V(\mathrm{Frob}_p)$. A Frobenius with nonzero trace must be scalar, hence central, or equivalently must split completely in $\Q(i,\sqrt2)=\Q(\zeta_8)$. Consequently $a_p=0$ unless $p\equiv1\pmod8$. In the latter case Frobenius on $\alpha$ acts by multiplication by $2^{(p-1)/4}\bmod p$, which is $+1$ or $-1$ since $2$ is then a square. Therefore $a_p=+2$ if this residue is $+1$, and $a_p=-2$ if it is $-1$. At $2$, inertia on the four embeddings of $P$ is transitive by total ramification, so its invariant subspace in the permutation representation is one-dimensional. The trivial summand already provides it. Thus $V^{I_2}=0$, the standard Euler factor at two is one and all even-index Fourier coefficients vanish. At odd $p$ the inverse Euler factor is $1-a_p p^{-s}+\chi_{-4}(p)p^{-2s}$. For example, $a_{17}=-2$ since $2^4\equiv-1\bmod17$; at $p=3,5$ the quadratic terms give $a_9=1$ and $a_{25}=-1$. The Fourier series begins $q+q^9-2q^{17}-q^{25}+O(q^{26})$ with $q=e^{2\pi iz}$.

For the cuspidal automorphic representation $\pi_f$ generated by $f$, the following functions satisfy the full critical-line assertion:
\[
 \begin{split}
 L(s,\pi_f\times\widetilde\pi_f)
      &=\zeta(s)L(s,\chi_{-4})L(s,\chi_8)L(s,\chi_{-8}),\\
 L(s,\operatorname{Ad}\pi_f)
      &=L(s,\chi_{-4})L(s,\chi_8)L(s,\chi_{-8}),\\
 L(s,\operatorname{Sym}^2\pi_f)
      &=\zeta(s)L(s,\chi_8)L(s,\chi_{-8}).
 \end{split}
\]
For the first identity the endomorphism representation is regular on $\operatorname{Gal}(\Q(\zeta_8)/\Q)$, so it holds at every prime. The third identity uses $V^\vee\simeq V$ and $V\otimes V=\operatorname{Sym}^2V\oplus\det V$. The symmetric-square parameter is itself the indicated direct sum of three rational characters, giving an isobaric automorphic realization. The local distinction at $2$ is substantial: the original standard factor there is one, whereas the tensor-square factor is $(1-2^{-s})^{-1}$ because inertia fixes the scalar endomorphisms. The tensor identity on the actual local Weil representation gives the factor at $2$ without any deletion.

More generally let $\omega_{\chi,t}$ and $\omega_{\psi,u}$ be arbitrary unitary rational characters. The representations $\pi_f\otimes\omega_{\chi,t}$ are genuine cuspidal representations of $\mathrm{GL}_2$. Their mutual Rankin--Selberg functions are
\[
  L\bigl(s,(\pi_f\otimes\omega_{\chi,t})
          \times(\pi_f\otimes\omega_{\psi,u})^\vee\bigr)
       =\prod_{\eta\in\{1,\chi_{-4},\chi_8,\chi_{-8}\}}
              L\bigl(s-i(t-u),(\eta\chi\bar\psi)_{\rm prim}\bigr).
\]
There are infinitely many inequivalent such cusps, since a self-twist must have square the trivial character by comparison of central characters, whereas rational characters of unbounded order exist. The formula gives an actual ERH assertion for a family of Rankin--Selberg products of cusps. Each individual standard $L$-function still contains an irreducible degree-two parameter.

All even tensor powers admit equally explicit automorphic realizations. If $A=\operatorname{Gal}(\Q(\zeta_8)/\Q)$, the tensor square of $V$ is $\mathrm{Reg}_A$. The tensor product of $k$ copies of $\mathrm{Reg}_A$ has character $4^k$ at the identity and zero elsewhere, so $V^{\otimes 2k}=4^{k-1}\mathrm{Reg}_A$ for $k\ge1$. On the other hand, $V\otimes\eta\simeq V$ for every $\eta\in\widehat A$, since $V$ is the unique nonlinear irreducible of this group. Thus $V^{\otimes(2k+1)}=4^k V$. For each $k$ the automorphic realization is simply the isobaric sum prescribed by this decomposition, using $f$ for each occurrence of $V$. Rational character twists can be inserted in each tensor slot and contribute their product. We have used no theorem on generic tensor power functoriality on the automorphic dual. All asserted local parameters come from these explicit decompositions, including at the ramified prime.

At an odd prime the formula for $a_p$ implies $a_p^2=4\boldsymbol{1}_{p\equiv1\bmod8}$. Its finite Fourier expansion is $1+\chi_{-4}(p)+\chi_{8}(p)+\chi_{-8}(p)$. The companion statistic $a_p$ resolves the two different fourth-power residue signs within this congruence class. Dirichlet GRH controls the four characters in the square, while deciding the square-root cancellation for the first-power twist of the signs is the genuinely degree-two problem.

\subsection*{How special is this weight-one example?}
Let $g$ be any normalized weight-one cusp eigenform over $\Q$ with associated irreducible odd finite-image representation $W$ and the standard local compatibility (the classical Artin--theta correspondence suffices for the newforms constructed here). A necessary condition for its adjoint to factor into primitive Dirichlet functions is that $[G,G]$ act by scalars on $W$, where $G$ is its finite image. Indeed, equality of prime traces, followed by Chebotarev in a finite compositum with the Dirichlet character fields, forces $\operatorname{End}(W)$ to be a sum of characters of the abelianization. Its kernel consists of the elements acting as scalars on $W$. In that case put $A=G/\ker(\operatorname{P}W)$. For any $g$ of nontrivial projective image some conjugate acts by its matrix times a nonidentity scalar, by Schur's lemma and commutativity of $A$. Thus $\operatorname{tr}W(g)=0$. The trace of $\operatorname{End}(W)$ on $A$ equals four at the identity and zero elsewhere. Its trivial multiplicity is one by irreducibility, so $|A|=4$ and the endomorphism representation is regular. It cannot be cyclic: a generator of a cyclic projective image and scalar matrices generate an abelian image, contradicting irreducibility in dimension two. Thus it is the Klein four group, and the condition is also sufficient. All three constituents of the trace-zero adjoint are then the nontrivial quadratic characters of its projective field. The statement remains necessary if a finite number of Euler factors is omitted or rational unitary shifts are allowed, by the trace argument above.

Such representations are themselves quadratic inductions. Take an order-two subgroup of the projective image and its inverse image $H$. The matrices from $H$ are generated projectively by a single nonscalar involution, so they preserve its two distinct eigenlines. Some group element outside $H$ interchanges these eigenlines by irreducibility; consequently $W$ is induced from the character of $H$ on either line. Complex conjugation acts nonscalarly because $W$ is odd: it has eigenvalues $1,-1$. Two of the three subgroups of order two in the projective image omit its class. Their fixed quadratic fields are imaginary, so this yields an imaginary-quadratic Hecke induction. This explains in general why weight-one cusps of precisely this projective type have abelian adjoint parameters.

For an elliptic curve over $\Q$, the modular form has weight two; in its classical normalization the center is $s=1$. The correct shift to the unitary convention is $L(s,\pi_E)=L(E,s+1/2)$. The degree-two nonfactorization theorem applies because $\pi_E$ is cuspidal. In the CM case the angular-character calculation also identifies exactly which additional Hecke zero theorem would be required for its adjoint. In the non-CM case a replacement of its two standard Euler roots by characters of rational primes would contradict isobaric uniqueness. None of the tensor identities above assert critical-line zeros for the individual elliptic $L(E,s)$.

\paragraph{References for the cuspidal example.}
For the classical theta realization and its conductor convention see Y.~Hu and A.~Saha, \href{https://webspace.maths.qmul.ac.uk/abhishek.saha/research/sup-norm-newforms-depth.pdf}{\emph{Sup-norms of eigenfunctions in the level aspect for compact arithmetic surfaces, II: newforms and subconvexity}}, Remark~6.5. For the general weight-one association see P.~Deligne and J.-P.~Serre, \href{https://www.numdam.org/articles/10.24033/asens.1277/}{\emph{Formes modulaires de poids $1$}}, Ann. Sci. École Norm. Sup. 7 (1974). The field, local traces and conductor in the example were computed above; the database locators specify the form independently.

\subsection*{Coprime local factors at ordinary elliptic primes}
\begin{proposition}\label{at:ordinary-obstruction}
Let $E/F$ be any elliptic curve in characteristic zero, $F$ a number field, and normalize its $h$th symmetric power to weight zero. Fix a finite list of finite-image Artin representations over arbitrary number fields, each possibly twisted by an arbitrary fixed unitary norm power $|\cdot|^{it}$. Regroup all Euler factors by rational primes. At a positive-density set of good rational primes the following holds simultaneously for all $h\ge1$.
\begin{enumerate}
\item If $h$ is odd, its local elliptic polynomial is coprime over $\mathbb C$ to every local polynomial in the list.
\item If $h$ is even and the prime splits completely in $F$, only the middle eigenvalue $1$ at each place of $F$ can be shared. In particular the total multiplicity of a shared root in the elliptic polynomial is at most $[F:\mathbb Q]$.
\end{enumerate}
For $E/\mathbb Q$ the first assertion holds at every good ordinary prime, apart from at most two rational primes per distinct nonzero $t$ and the finitely many ramified primes of the list. Such primes have density one for a non-CM curve and density one half for a CM curve.
\end{proposition}
\begin{proof}
For a fixed real $t\ne0$, at most two rational primes can have $p^{it}$ algebraic. Otherwise take the two $\mathbb Q$-independent numbers $1,it$ and the three $\mathbb Q$-independent real logarithms of distinct such primes. Independence of the latter follows by unique factorization. Three of their six exponentials are the rational primes, while the other three are algebraic by assumption, contradicting the \href{https://webusers.imj-prg.fr/~michel.waldschmidt/articles/pdf/DION2005proceedings.pdf}{six exponentials theorem, Waldschmidt, Theorem~2.1}.

At a good rational prime any regrouped reciprocal Euler root of a finite-image Artin representation is a root of unity. This remains true if a place above $p$ has higher residue degree $f$, because the roots obtained from $1-\zeta T^f$ are roots of unity. A unitary norm twist multiplies these roots by $p^{it}$, up to the sign convention for $t$. Thus outside the indicated finite exceptions all the candidate roots are either roots of unity or transcendental numbers.

For the positive-density assertion take rational primes splitting completely in the normal closure of $F(E[3])/\mathbb Q$, apart from the finitely many bad primes and the six-exponentials exceptions. This set has positive density by Chebotarev. Every residue field of $F$ there is $\mathbb F_p$. The Frobenius matrix acts trivially on $E[3]$, so its integer trace $a_v$ is $2\bmod 3$ at each $v\mid p$. For $p>4$ the Hasse bound gives $|a_v|<p$; hence $p\nmid a_v$. The polynomial $X^2-a_vX+p$ has $p$-adic root valuations $0$ and $1$ by its Newton polygon. Dividing these roots by $\sqrt p$ gives algebraic unitary roots $z,z^{-1}$ of valuations $-1/2,1/2$. Consequently $z^j$ is algebraic and cannot be a root of unity for any nonzero integer $j$. The $h$th symmetric power has exactly the eigenvalues $z^{h-2j}$ for $0\le j\le h$. For odd $h$ none survives the comparison, while for even $h$ the exponent zero appears exactly once at each split place. This proves both assertions.

Over $\mathbb Q$ only ordinarity was used in the eigenvalue comparison. For CM elliptic curves the ordinary good primes are precisely the split primes of the CM field, outside finitely many exceptions, by the CM reduction theorem. For non-CM elliptic curves a supersingular prime $p>4$ has $a_p=0$, a set of density zero by Sato--Tate and the absence of an atom at zero. The asserted densities follow.
\end{proof}

In particular an odd elliptic symmetric power, including the original Hasse--Weil function, cannot share any positive-degree Euler constituent with a fixed finite list of shifted Dirichlet functions. Passing to an arbitrary finite extension of the field of definition does not remove the obstruction. In a class-number-one CM field the nonzero angular exponents are still distinguished by the same ordinary Frobenius valuations. Only the zero exponent in an even power is available for a finite-character factor; its actual global occurrence in the CM case is the middle character of Theorem~\ref{recentcm:main}. This local result concerns inheritance through Euler factorization. It makes no claim of logical independence between conjectures on analytic zero divisors.

\begin{theorem}[The maximal finite part of elliptic tensor powers]\label{at:catalan}
Let $E/\mathbb Q$, and let $E_1,\ldots,E_{2r}$ be quadratic twists of $E$, with $r\ge1$. Write $\chi$ for the product of the twisting characters, and normalize their tensor product by its integral Tate twist to weight zero. Put $B_r=\binom{2r}{r}$ and $C_r=B_r/(r+1)$. Its largest finite-image Galois direct summand has complete standard function
\[
 \begin{cases}
   L(s,\chi)^{C_r},& E\text{ without CM},\\[2pt]
   [L(s,\chi)L(s,\chi\epsilon_K)]^{B_r/2},&
       E\text{ with CM by }K,\quad
       \epsilon_K=\epsilon_{K/\mathbb Q}.
 \end{cases}
\]
The starting Dirichlet theorem therefore places every nontrivial zero of the displayed factor on the critical line, with its indicated multiplicity, after any zeta pole is removed. No nonzero imaginary norm translate of a finite-image representation supplies an additional Euler constituent. For an odd number of twists there is no finite-image Euler constituent after unitary normalization.
\end{theorem}
\begin{proof}
Let $V$ be the two-dimensional rational $\ell$-adic representation of $E$ in the convention with Frobenius polynomial $X^2-a_pX+p$, and put $\delta=\det V$. The normalized even tensor is $V^{\otimes2r}\otimes\delta^{-r}\otimes\chi$. Removing the last finite twist preserves dimensions and maximality of the finite part.

Without CM, Serre's open-image theorem implies that the image of every finite-index subgroup is Zariski dense in $\mathrm{GL}_2$; see \href{https://msp.org/ant/2015/9-10/ant-v9-n10-p06-p.pdf}{Lombardo, introduction, citing Serre's theorem}. The kernel of any finite-image summand contains such a subgroup, so the summand is precisely made up of algebraic invariants. Over characteristic zero the tensor representation is semisimple. The multiplicity of $\operatorname{Sym}^{2j}V\otimes\delta^{-j}$, $0\le j\le r$, is
\[
             \binom{2r}{r-j}-\binom{2r}{r-j-1},
\]
where a negative lower index gives zero. Indeed on the diagonal determinant-one torus the weight-$2j$ multiplicity is the first binomial coefficient, while each symmetric power containing that weight also contains every smaller nonnegative weight of the same parity once. Successive differences give the formula, by the ordinary $\mathrm{SL}_2$ highest-weight decomposition. Only $j=0$ gives invariants. Its multiplicity is $C_r$, and the full $\mathrm{GL}_2$ action on that space is trivial after the displayed determinant normalization.

In the CM case restrict to $G_K$, extending $\ell$-adic scalars if necessary. Write $V|_{G_K}=\omega\oplus\omega^c$. The character $\omega/\omega^c$ has infinite image on every finite-index subgroup: an ordinary split Frobenius has quotient of nonzero $p$-adic valuation by the proof of Proposition~\ref{at:ordinary-obstruction}, so that its powers are all distinct. The only finite-image part therefore consists of words containing $r$ copies of each character. These $B_r$ words are all trivial on $G_K$ after tensoring by $\delta^{-r}$. The nontrivial coset permutes their word basis by complementation, possibly multiplying each basis vector by a scalar. No word equals its complement. The coset operator has trace zero, and its square acts trivially because $G_K$ acts trivially here. Thus its two eigenvalues $+1,-1$ have equal multiplicity $B_r/2$. The other words are stable in complementary pairs, so this finite part is a direct summand. Restoring $\chi$ gives the stated two characters.

These are decompositions of the actual Galois representations. Their local invariant-space Euler factors and fixed archimedean factors consequently multiply at ramified as well as good places. For a nonzero imaginary norm translate, the transcendence argument of Proposition~\ref{at:ordinary-obstruction} still separates its reciprocal roots from every algebraic geometric reciprocal root at all but finitely many good primes. For an odd tensor the exponents at a split ordinary prime are all odd, as for the odd symmetric powers, and multiplication by any finite quadratic sign cannot make a non-torsion algebraic root torsion. This gives the final assertions.
\end{proof}

\begin{corollary}[Sharp density of exact Dirichlet-pair matching]\label{at:quarter}
For a rational elliptic curve and any fixed pair of unitary Dirichlet functions, allow arbitrary imaginary norm translates of either member. Compare their complete good-prime Euler polynomials with the curve's unitary quadratic polynomial. If the curve is non-CM, the density of primes with equality is zero. If the curve has CM, the maximum of this density over all such pairs is exactly $1/4$. Any primitive quartic Dirichlet character and its conjugate attain the maximum for every rational CM curve.
\end{corollary}
\begin{proof}
A nonzero norm shift prevents equality apart from finitely many primes, by the six-exponentials argument in Proposition~\ref{at:ordinary-obstruction}. For unshifted characters all ordinary good primes are excluded. This settles the non-CM case.

For a CM curve with imaginary quadratic field $K$, an inert good prime has trace zero and unitary reciprocal roots $i,-i$. Thus if the first character of a matching pair is $\chi$, necessarily $\chi(p)\in\{i,-i\}$ and $\epsilon_K(p)=-1$. On the common finite residue group, multiplication by $-1$ interchanges the two values of $\epsilon_K$ because $K$ is imaginary. It preserves the condition $\chi\in\{i,-i\}$ because $\chi(-1)=\pm1$. By primes in arithmetic progressions the density of this necessary event is half the proportion of the group with $\chi\in\{i,-i\}$. For character order $m$ the latter proportion is zero unless $4\mid m$, and is then $2/m\le1/2$. This proves the upper bound. For quartic $\chi$ the event has density $1/4$ and its two roots $\chi(p),\bar\chi(p)$ are exactly $i,-i$. The upper bound is attained.
\end{proof}

\begin{proposition}[Two curves with the same CM field]\label{at:two-cm}
Let $E_1,E_2/\mathbb Q$ have complex multiplication by orders in the same imaginary quadratic field $K$. Choose their two $G_K$ characters $\omega_1,\omega_2$ with the same CM type, and set $\eta=\omega_1/\omega_2$. The maximal finite part of the Tate-normalized tensor of the two elliptic representations is the two-dimensional representation $\operatorname{Ind}_{G_K}^{G_{\mathbb Q}}\eta$. Its complete function is a product of two Dirichlet functions exactly when $\eta^2=1$. Otherwise it is the standard function of an irreducible odd dihedral Artin representation. Both possibilities occur even among twists of a single rational CM curve.
\end{proposition}
\begin{proof}
The ratio is finite. Indeed the two curves are isogenous over $\overline{\mathbb Q}$ by CM uniformization, since any two lattices with the same CM field are homothetic to commensurable fractional lattices in that field. An isogeny is defined over a finite extension and identifies the two characters of matching CM type on its Galois group. Write $\delta$ for the common determinant. On $G_K$ one has $\omega_i\omega_i^c=\delta$, and thus $\eta^c=\eta^{-1}$. On restricting the Tate-normalized tensor to $G_K$, the two mixed character lines have characters $\eta,\eta^{-1}$. Complex conjugation swaps the lines, so their span is the claimed induction. On each of the two remaining lines the character is $\omega_1/\omega_1^c$ or its inverse times a finite character. Its image is infinite on every finite-index subgroup by ordinary Frobenius valuations as in Theorem~\ref{at:catalan}. This establishes maximality. If $\eta=\eta^c$, its two extensions to $G_{\mathbb Q}$ assign opposite signs to complex conjugation. Their finite characters are Dirichlet and differ by $\epsilon_K$. Otherwise the distinct lines are exchanged and give an irreducible representation, with conjugation of determinant $-1$. Its finite projective image is dihedral and the quadratic induction has its usual weight-one theta realization. Local multiplicativity and Artin formalism include every ramified factor.

For example let $E_1:y^2=x^3+1$, $E_2:y^2=x^3+2$ and $K=\mathbb Q(\sqrt{-3})$. For $u^6=2$ the map $(x,y)\mapsto(u^2x,u^3y)$ identifies the first model with the second. The prime $2$ is inert in $K$, so $X^6-2$ is Eisenstein there; the extension $K(u)/K$ is cyclic of degree six because $\mu_6\subset K$. Its Galois action on the map is the full automorphism group $\mu_6$ of the cubic, acting faithfully on either CM character line. Hence $\eta$ has order six and gives the irreducible case. Taking $E_2=E_1$ or a rational quadratic twist gives the split case.
\end{proof}

\clearpage
\section{What happens beyond primitive cuspidal standard functions}\label{res:scope}
A primitive unitary cuspidal standard function on $\mathrm{GL}_n$ has a distinguished symmetry line $\operatorname{Re}s=1/2$. For a broader class there are two distinct claims one could make. The first places every nonreal zero of the completed function on that line. The second imposes that condition only inside $0<\operatorname{Re}s<1$. These claims give different answers even for a unitary representation in the discrete spectrum. In what follows all variables have unitary normalization, even when a constituent has a nonunitary real twist. Put
\[
 Z(s)=\pi^{-s/2}\Gamma(s/2)\zeta(s).
\]
This meromorphic completion has poles at 0 and 1. Hardy's theorem gives infinitely many zeros $\rho=1/2+i\gamma$ with $\gamma\ne0$. A real shift of $Z$ has poles only on the real axis. Thus every shifted Hardy zero survives in any finite product of real shifts of $Z$, whatever the multiplicities of the other zeros. This elementary observation supplies unconditional counterexamples.

\subsection*{A square-integrable residual representation}
Let $r\ge2$. The constant function on the automorphic quotient of $\mathrm{GL}_r/\mathbb Q$ modulo its split central direction generates its trivial representation. It is unitary and square-integrable there by the finite-volume theorem. Its constant term along a proper parabolic is nonzero. At each prime $p$ its spherical parameter for normalized induction is the unordered list
\[
 p^{(r-1)/2},\quad p^{(r-3)/2},\quad\ldots,\quad p^{-(r-1)/2}.
\]
Indeed the trivial representation is the spherical quotient obtained by inducing the $r$ successive half-modulus characters from its Borel. The corresponding archimedean Langlands parameter has the same real exponents, each with even parity. Consequently the completed standard function is
\begin{equation}\label{res:trivial}
 \Lambda_r(s)=\prod_{j=0}^{r-1} Z\bigl(s+(r-1)/2-j\bigr),
 \qquad \Lambda_r(1-s)=\Lambda_r(s).
\end{equation}
The reindexing $j\mapsto r-1-j$ proves the functional equation. The spectrum here is a special case of the residual classification of M\oe glin and Waldspurger. For a unitary cuspidal datum $\tau$ the general Speh representation associated with a segment of length $r$ has standard factors at unramified places given by the successive twists $L_p(s+(r-1)/2-j,\tau_p)$. Formula \eqref{res:trivial} uses the everywhere unramified trivial datum and all its archimedean factors.

Each Hardy zero gives zeros of \eqref{res:trivial} whose real parts are
\[
  1-r/2,\ 2-r/2,\ \ldots,\ r/2.
\]
For $r=2$ these are $0,1$. They refute the assertion about every nonreal completed zero. They do not refute the assertion inside the open strip. More precisely, assuming RH for $\zeta$, all nonreal zeros have the displayed real parts. For even $r$ all those parts are integers and the open strip is empty of nonreal zeros. For odd $r$ its sole possible real part there is $1/2$. The finite gamma-free product has additional trivial zeros only on the real axis. Unitarity of the global residual representation is fully consistent with these nonunitary real shifts of its standard factors.

\subsection*{A genuine Siegel cusp form and its two Euler products}
Take $k=10$ and $h=E_6\Delta$, the normalized full-level elliptic cuspidal eigenform of weight $2k-2=18$. The ring of modular forms gives $S_{18}=\Delta M_6$ of dimension one. Let $F$ be its nonzero Saito--Kurokawa lift, a scalar weight-ten holomorphic degree-two Siegel cuspidal eigenform for $\mathrm{Sp}_4(\mathbb Z)$. Its automorphic representation on $\mathrm{GSp}_4$ is cuspidal and associated to a proper parabolic in the sense of CAP (cuspidal associated to parabolic). Write $H(s)=L(s+k-3/2,h)$ and
\[
  \mathcal H(s)=(2\pi)^{-(s+k-3/2)}\Gamma(s+k-3/2)H(s).
\]
The completed function is entire and has the symmetry $s\leftrightarrow1-s$ with a root number. The two different finite Euler products of $F$ satisfy
\begin{align}
  L_{\rm spin}(s,F)&=H(s)\zeta(s-1/2)\zeta(s+1/2),\label{res:spin}\\
  L_{\rm std}(s,F)&=\zeta(s)H(s-1/2)H(s+1/2).\label{res:std}
\end{align}
The first is the degree-four spin product. It follows from the classical Saito--Kurokawa formula in Pitale, \emph{Siegel Modular Forms}, \S3.4.1, (20), by shifting the classical variable by $k-3/2$. The degree-five standard product is also computed in Exercise 3.14 there. Here is a local check. The four unitary-normalized spin parameters are $\alpha_p,\alpha_p^{-1},p^{1/2},p^{-1/2}$ with $\lvert\alpha_p\rvert=1$ by Deligne. Their second exterior power contains two copies of $1$ and the four cross products $\alpha_p^{\pm1}p^{\pm1/2}$. Removing the invariant symplectic line leaves one $1$ and the four cross products. This gives exactly \eqref{res:std}. Full level removes any exceptional finite prime.

Every nonreal Hardy zero $\rho$ gives spin zeros at $\rho\pm1/2$ with real parts $0,1$. For the standard function use Hafner's theorem that a positive proportion of the zeros of the full-level elliptic eigenform $h$ lie on its critical line. Thus $H(1/2+it)=0$ for infinitely many nonzero real $t$ and \eqref{res:std} has zeros at $it$ and $1+it$. Each asserted zero survives. The other factors are holomorphic at all such nonreal points and the gamma factors in these completions have only real shifts. A gamma function has no zeros and its poles have real arguments here. Both products therefore have nonreal zeros off the central line, even though $F$ itself is cuspidal. The group and the representation of its dual group specifying the Euler factors matter.

There is a sharper conclusion if only the open strip is counted. Assume RH for $\zeta$. For \eqref{res:spin}, neither of its shifted zeta functions has a zero or a pole inside the strip except that $\zeta(s+1/2)$ has its pole at $s=1/2$. That pole cannot alter the truth of an assertion about zeros \emph{off} the center. Thus the open-strip zero condition for the spin product is equivalent to ERH for $h$. For \eqref{res:std}, if ERH holds for $h$, all shifted nontrivial zeros of $H$ lie on the two boundaries. Its trivial zeros are to the left: those of $H$ have real parts at most $-(k-3/2)$ before the shifts. Conversely, if $\beta+it$ is a nontrivial zero of $H$ with $0<\beta<1/2$, then $(\beta+1/2)+it$ is an off-center zero of \eqref{res:std} in the open strip. If $1/2<\beta<1$, use $(\beta-1/2)+it$. The sole possible pole from the unshifted zeta factor is at $s=1$, outside this region. Hence the same equivalence holds for the standard product. We use the classical nonvanishing of the primitive elliptic eigenform on the edge $\operatorname{Re}s=1$ and reflection at the other edge. These equivalences retain the unresolved ERH for the ordinary primitive cuspidal function $H$.

\subsection*{Actual off-center zeros inside the open strip}
Enlarge the class to Hecke eigenforms of moderate growth, including non-square-integrable Eisenstein specializations. For $0<a<1/2$ put $u=1/2+a$ and specialize the classical spherical level-one Eisenstein series $E(z,u)$ there, with its standard constant-term normalization. It is regular at $u$ and has constant term
\[
       y^{u}+\frac{Z(2u-1)}{Z(2u)}y^{1-u}.
\]
Here $2u>1$ and $0<2u-1<1$ ensure the two factors are finite and the denominator is nonzero. Its normalized Hecke eigenvalue at $p$ is $p^a+p^{-a}$. Both assertions follow, for instance, by dividing the completed Fourier expansion in Zagier, \emph{Eisenstein series and the Riemann zeta function}, (14), by $Z(2u)$ and applying his Hecke formula (18). The determinant parameter equals $1$. Consequently the finite standard product and its natural completion are
\[
   L(s,E(\cdot,u))=\zeta(s-a)\zeta(s+a),\qquad
   \mathcal E_a(s)=Z(s-a)Z(s+a)=\mathcal E_a(1-s).
\]
For every Hardy zero, both $\rho-a$ and $\rho+a$ are nonreal zeros. Their real parts $1/2-a$ and $1/2+a$ lie strictly inside $(0,1)$ and differ from $1/2$. The shifts move all poles along the real axis only. For example, $u=3/4$ gives the two vertical lines $1/4$ and $3/4$ without assuming RH for any function.

The global spectral qualification is decisive. Cauchy--Schwarz in the periodic variable bounds $\int_0^1|E(x+iy,u)|^2\,dx$ below by the squared absolute value of the constant term. For large $y$ the latter is comparable with $y^{2u}$. The cusp integral with measure $dy/y^2$ diverges since $2u-2>-1$. Thus the specialization is outside the square-integrable spectrum. Its Laplace eigenvalue $u(1-u)$ is also below the range $[1/4,\infty)$ of the level-one Eisenstein unitary continuous spectrum. It is precisely the presence of the pair of opposite real twists that gives the two interior zero lines. At every place, the normalized local induction of $|\cdot|_v^a$ and $|\cdot|_v^{-a}$ lies in the unitary complementary series when $0<a<1/2$. This follows from the classical unitary dual of $\mathrm{GL}_2(\mathbb R)$ and the spherical unitary dual of $\mathrm{GL}_2(\mathbb Q_p)$ (Vogan and Tadi\'c, respectively). Twisting both inducing characters by a common unitary character preserves this property. In particular the counterexample persists even if one asks for an invariant positive-definite inner product at every local place. Its global Eisenstein realization still has the divergent norm computed above. The calculation does not extend the counterexample to primitive unitary cuspidal standard functions on $\mathrm{GL}_n$ or to the Hasse--Weil function of an elliptic curve.

\subsection*{An entire degree-two Euler product with two interior zeros}
The preceding example can be strengthened by removing all poles of the completed function. Twist the adelic spherical specialization at $u=3/4$ by $\chi(\det)$ for the primitive odd quadratic character $\chi=\chi_{-4}$. The twist is again an automorphic form of moderate growth with the same absolute value. Its two inducing characters are $\chi|\cdot|^{1/4}$ and $\chi|\cdot|^{-1/4}$. Write $\beta(s)=L(s,\chi)$ and $B_\chi(s)=(4/\pi)^{(s+1)/2}\Gamma((s+1)/2)\beta(s)$. This completed Dirichlet function is entire and satisfies $B_\chi(s)=B_\chi(1-s)$. The quadratic character has Gauss sum $2i$ and hence root number $1$. The completed standard function of the twist, up to a constant independent of $s$, is
\[
               B_\chi(s-1/4)B_\chi(s+1/4).
\]
At the sole ramified finite prime 2, both local one-dimensional parameters are nontrivial on inertia and both Euler factors equal one. The product is entire, of degree two, with an Euler product, a functional equation about $1/2$, and strictly interior noncentral zeros.

Here is a direct witness of the last assertion. Reality and the functional equation make $B_\chi(1/2+it)$ real. At $t=6$ its value is in $(0.00097,0.00098)$ and at $t=7$ it lies in $(-0.0223,-0.0222)$. The accompanying calculation of the enclosures uses the identity
\[
 \beta(s)=4^{-s}\bigl(\zeta(s,1/4)-\zeta(s,3/4)\bigr)
\]
and complex-ball evaluations of the Hurwitz zeta and gamma functions at 160 bits. Continuity gives $B_\chi(1/2+i\gamma)=0$ for some $6<\gamma<7$. Thus the degree-two completed standard function has the zeros $1/4+i\gamma$ and $3/4+i\gamma$, both in the open strip. This example also satisfies absolute convergence in a right half-plane. Its unramified parameters have moduli $p^{1/4}$ and $p^{-1/4}$, so their size records exactly the spectral feature used in the construction.

\paragraph{Sources.}
For the discrete spectrum see C.~M\oe glin and J.-L.~Waldspurger, \emph{Le spectre r\'esiduel de $\mathrm{GL}(n)$}, Ann. Sci. \!\'Ecole Norm. Sup. (4) 22 (1989), 605--674, and J.~Arthur, \emph{Eigenfamilies, Characters and Multiplicities}, \S1. For the two full-level genus-two identities see A.~Pitale, \emph{Siegel Modular Forms. A Classical and Representation-Theoretic Approach}, Lecture Notes in Mathematics 2240, Springer, 2019, \S3.4.1. For zeros on the critical line see G.~H.~Hardy, \emph{Sur les z\'eros de la fonction $\zeta(s)$ de Riemann}, C. R. Acad. Sci. Paris 158 (1914), 1012--1014, and J.~L.~Hafner, \emph{Zeros on the critical line for Dirichlet series attached to certain cusp forms}, Math. Ann. 264 (1983), 21--37. For the Fourier coefficients of the spherical specialization see D.~Zagier, \emph{Eisenstein series and the Riemann zeta function}, in \emph{Automorphic Forms, Representation Theory and Arithmetic}, Springer, 1981, 275--301.

\clearpage
\phantomsection\addcontentsline{toc}{part}{III. Positivity and heat deformation}
\noindent{\small\sffamily\color{accent}III. Positivity and heat deformation}\par
\section{Weil-form compression for a fixed elliptic curve}\label{wc:section}

Fix a rational elliptic curve, or more generally a primitive holomorphic cuspidal newform $f$ of weight $k\ge2$, trivial character and fixed level $q$. Work in unitary normalization and let $N(I)$ count its zeros with multiplicity for ordinates in $I=[T,2T)$. Let $U(I)$ count distinct zeros on the critical line and $V(I)$ count distinct unordered pairs off the line, with pairing $\rho\leftrightarrow1-\bar\rho$. Thus the number of all distinct zeros is $N_d(I)=U(I)+2V(I)$. We prove
\begin{equation}
 \liminf_{T\to\infty}\frac{U(I)+V(I)}{N(I)}
 \ \ge\ \left(\frac14+\frac1{\sqrt2}\cot\frac1{2\sqrt2}\right)^{-1}
 =0.46168822\ldots .\label{wc:orbits}
\end{equation}
In particular $N_d(I)$ has the same lower bound, with the additional term $V(I)$ before dividing by $N(I)$. Distinctness counts a zero of order $m$ once. The bound allows many off-line zeros and gives no positive proportion of simple critical-line zeros. The obstruction to getting that conclusion by the same argument is quantitative.

We start from the compression and inertia argument of
\href{https://arxiv.org/html/2608.13637v2}{Alp\"oge and Furman, \S\S2--6}. Their Theorems A and B apply to zeta and Dirichlet functions. The zero density doubles at degree two and the Rankin--Selberg prime variance stays equal to one. We track these two constants separately. The limit in \eqref{wc:orbits} follows from the fixed-curve theorem below, without a hypothesis on the zeros or a conductor average. We make no priority claim for the distinct-zero bound.

\subsection{Two normalizations and their ratio}

Write $L(f,s)$ here for the unitary function, so its completed gamma factor is a constant times $(\sqrt q/(2\pi))^s\Gamma(s+(k-1)/2)$. We use the Fourier transform $\widehat F(z)=\int F(u)e^{-izu}\,du$. Set
\[
    \ell=2\log T,\qquad X=T^\theta,\qquad L=\log X,
    \qquad \lambda=L/\ell=\theta/2,
    \qquad 0<\theta\le1 .
\]
For fixed $f$ the zero-counting formula gives
\begin{equation}
  N(I)=\frac{T\ell}{2\pi}+O_f(T),\qquad
  N([t,t+1])\ll_f\log(3+|t|).
  \label{wc:count}
\end{equation}
The first error includes the constant in the height-aspect conductor and the variation of $\log t$ over a dyadic interval. In particular $T\log X/(2\pi)$ and $N(I)$ differ at the main-term level whenever $\lambda\ne1$.

Let $\eta$ be a nonzero real even function in $C_c^\infty((-1/2,1/2))$, and put $w=\eta^2$, $A=\int w$, $B=\int w^2$. All integrals involving $w$ use its zero extension. Define
\[
 \phi(u)=\eta(u/L),\quad h=2\pi/L,\quad
 \alpha_j=T+jh,\quad J=\{j\in\Z:\alpha_j\in I\},\quad d=\#J.
\]
Take $T$ large enough that $L\ge1$. With $z_\rho=(\rho-1/2)/i$ and $v_\rho=(\widehat\phi(z_\rho-\alpha_j))_{j\in J}$, form the real symmetric matrix
\[
 G=\frac{1}{A L^2}\sum_{\Re z_\rho\in I'} m_\rho v_\rho v_\rho^{\mathsf T},
 \qquad I'=[T-\sqrt T,2T+\sqrt T).
\]
The transpose here is bilinear. An off-line pair has conjugate vectors and contributes the difference of two real positive semidefinite matrices of rank at most one. A distinct on-line zero contributes one such positive matrix.

\begin{theorem}\label{wc:moment}
For fixed $f,\theta,\eta$ in the ranges above, $\operatorname{tr}G=N(I)+o(N(I))$ and
\begin{align}
  \frac{\operatorname{tr}(G^2)}{N(I)}
  &=R_\lambda(w)+o(1),\label{wc:R}\\
  R_\lambda(w)
  &=\frac1{A^2}\left[\frac{B}{\lambda}+
          \lambda\iint |u-v|w(u)w(v)\,du\,dv\right].\notag
\end{align}
The matrices have $d/N(I)\to\lambda$.
\end{theorem}

\begin{proof}
The dimension assertion follows by counting the grid and using \eqref{wc:count}. Poisson summation of products of two translates gives the exact identity
\begin{equation}
 \sum_{j\in\Z}\widehat\phi(z-\alpha_j)\widehat\phi(z'-\alpha_j)
   =L\widehat{\phi^2}(z-z').\label{wc:poisson}
\end{equation}
Indeed the dual frequencies are the multiples of $L$ and the convolution of the two modulated windows is supported in $(-L,L)$. At frequency zero, Fourier inversion and the evenness of $\phi$ give the right hand side. This holds for complex arguments by analytic continuation or directly by Poisson summation. At $z=z'$ the sum is $AL^2$, including for an off-line zero.

Here are tail bounds that justify use of this identity at the two endpoints. For any integer $m\ge3$ integration by parts gives
\[
 |\widehat\phi(x+iy)|\ll_{\eta,m}
 L e^{L|y|/2}(1+L|x|)^{-m}.
\]
All nontrivial zeros satisfy $|\Im z_\rho|\le1/2$. When $t=\Re z_\rho$ lies inside $I$, sum the squares of this bound over grid points outside $I$. When $t$ lies outside $I$, sum them over those inside. If $s$ is the distance from $t$ to the closer endpoint, the result after division by $AL^2$ is
\[
 \ll_{\eta,m}\sqrt X(1+Ls)^{1-2m}.
\]
This follows by summing a decreasing power along two one-sided grids with $Lh=2\pi$. Using \eqref{wc:count} in successive unit intervals gives total error $O_{f,\eta}(\sqrt X\log T)$ for the trace relative to $N(I)$. For the matrix $E$ of omitted zeros with $t\notin I'$, the trace norm satisfies the same estimate restricted to $s\ge\sqrt T$. This uses $\|vv^{\mathsf T}\|_1=\|v\|_2^2$ also for complex $v$. Summing the power tail against $O_f(\log(3+T+s))$ gives $\|E\|_1=o(1)$. This proves the first moment. All subsequent calculations may use the full matrix $G+E$.

Put $b_f(p^a)=\alpha_p^a+\beta_p^a$ in terms of the normalized local parameters, allowing zero parameters at bad primes. These real numbers satisfy $|b_f(p^a)|\le2$ by the Ramanujan--Petersson theorem for holomorphic forms and the local newform factors. Write $a_n=b_f(n)\Lambda(n)/\sqrt n$ for prime powers and zero elsewhere. For real $t$ put
\[
 \mu(t)=\frac1\pi\left(\log\frac{\sqrt q}{2\pi}
    +\Re\frac{\Gamma'}{\Gamma}(k/2+it)\right),\qquad
 P_X(t)=-\frac1\pi\sum_{n\le X}a_n\cos(t\log n),\quad \nu=\mu+P_X.
\]
Weil's explicit formula for a test of Fourier support in $(-L,L)$ reads
$\sum_\rho m_\rho H(z_\rho)=\int H(t)\nu(t)dt$ with $H$ the Fourier transform of that test. There is no polar term for a cusp form. For the product of two grid translates its inverse Fourier transform has the required support. An even-test version suffices by replacing $H(t)$ with $(H(t)+H(-t))/2$, since both the full zero set and $\nu$ have this symmetry. These are the standard zero-counting and explicit formulas for a primitive automorphic function; see
\href{https://bookstore.ams.org/coll-53}{Iwaniec and Kowalski, \emph{Analytic Number Theory}, \S5.3 and \S5.5}.

Writing $K_I(t,s)=\sum_{j\in J}\widehat\phi(t-\alpha_j)\widehat\phi(s-\alpha_j)$, the formula applied entrywise gives
\[
  A^2L^4\operatorname{tr}((G+E)^2)
       =\iint_{\R^2} K_I(t,s)^2\nu(t)\nu(s)\,dt\,ds .
\]
Set $\Phi=\widehat{\phi^2}=L\widehat w(L\,\cdot)$, so the infinite-grid kernel equals $K_\infty(t,s)=L\Phi(t-s)$. For real arguments Cauchy--Schwarz and \eqref{wc:poisson} give $|K_I|,|K_\infty|\le AL^2$. We claim
\begin{equation}
 \operatorname{tr}((G+E)^2)=
  \frac{1}{A^2L^2}\iint_{I^2}\Phi(t-s)^2\nu(t)\nu(s)\,dt\,ds
  +O_{f,\eta}\big((X+\ell^2)/L^2\big).\label{wc:kernel}
\end{equation}
The following details also cover the endpoint $X=T$. Choose $m\ge4$ and put $u_j(t)=L(1+L|t-\alpha_j|)^{-m}$. Then $|\widehat\phi(t-\alpha_j)|\ll u_j(t)$.
Chebyshev's estimate and partial summation give $\sum_{n\le X}|a_n|\ll\sqrt X$.
Stirling's formula now gives $|\nu(t)|\ll B_T:=\sqrt X+\ell$ for $|t|\le4T$, and $|\nu(t)|\ll B_T+\log(1+|t|/T)$ everywhere. For $j\notin J$ one has
$\int_I u_j\ll(1+L\operatorname{dist}(\alpha_j,I))^{1-m}$, with sum of the squares of these bounds $O(1)$. On $I^2$ use
$|K_I^2-K_\infty^2|\ll L^2\sum_{j\notin J}u_j(t)u_j(s)$. Its integral with the absolute weights is $O(L^2 B_T^2)$. On the complement use symmetry and
$|K_I|^2\ll L^2\sum_{j\in J}u_j(t)u_j(s)$. For each such $j$,
$\int_\R u_j|\nu|\ll B_T$. Moreover $\sum_{j\in J}\int_{\R\setminus I} u_j|\nu|\ll B_T$. Indeed, for $|t|\le4T$ sum the convergent series in the distance of a grid point to the boundary in units of $1/L$. The remaining integral summed over $O(TL)$ points is
$\ll TL^2\int_{4T}^\infty(1+Lr/2)^{-m}(B_T+\log(1+r/T))dr\ll B_T$.
Dividing these errors by $A^2L^4$ proves \eqref{wc:kernel}.

Write $M[F,H]$ for the unnormalized double integral in \eqref{wc:kernel} with factors $F(t),H(s)$. Plancherel gives $\int\Phi^2=2\pi LB$, and $\int |x|\Phi(x)^2dx=O_\eta(1)$. Uniformly on $I$ one has $\mu(t)=\ell/(2\pi)+O_f(1)$ and $\mu'(t)\ll_f T^{-1}$. Consequently
\[
 M[\mu,\mu]=\frac{T\ell^2 L B}{2\pi}
             + O_{f,\eta}(T\ell L+\ell^2).
\]
Here replacing $\int_{I^2}\Phi(t-s)^2dt ds$ by $T\int\Phi^2$ costs $O_\eta(1)$, by its first absolute moment. With $m(s)=\int_I\Phi(t-s)^2\mu(t)dt$, one has $\sup_I|m|+\int_I|m'|\ll\ell L$. This follows by transferring the derivative to $\mu$ and including the two endpoint terms, whose $L^1$ norms are $O(\ell L)$. Integration by parts against each cosine gives $M[\mu,P_X]\ll \ell L\sum |a_n|/\log n\ll \ell L\sqrt X$.

Let $g=\phi^2\star\phi^2$ with $(w\star w)(y)=\int w(u)w(u+y)du$.
In $M[P_X,P_X]$ set $x=t-s$ and integrate first over $s\in I\cap(I-x)$. Expanding the cosines, the equal frequencies of opposite signs contribute
\[
 \frac{1}{2\pi^2}\sum_{n\le X}a_n^2\int_{|x|\le T}
  (T-|x|)\Phi(x)^2 e^{ix\log n}dx
   =\frac T\pi\sum_{n\le X}a_n^2 g(\log n)+O_\eta\left(\sum a_n^2\right).
\]
Fourier inversion supplies $2\pi g(\log n)$. The equal-sign terms have frequency $\log(nm)\ge\log4$ in $s$, so their aggregate is
$O_\eta(L(\sum|a_n|)^2)=O_\eta(LX)$.
For the unequal frequencies of opposite signs the integrated exponential has two endpoint terms with denominator $i(\log n-\log m)$. At each endpoint and each fixed $x$, the numerator factors into a function of $n$ and one of $m$ with respective absolute values $|a_n|,|a_m|$. The bilinear Montgomery--Vaughan inequality therefore bounds the off-diagonal sum by $O(\sum n|a_n|^2)$, since the minimum distance from $\log n$ to any other such frequency is $\gg 1/n$. Integrate $\Phi^2$ to obtain $O_\eta(L\sum n|a_n|^2)=O_\eta(XL^2)$, by $|b_f|\le2$ and $\sum_{n\le X}\Lambda(n)^2\ll X\log X$. In this estimate the off-diagonal error, after division by $A^2L^2N(I)$, is $O_{f,\eta}(X/(T\ell))$.

The fixed-form Rankin--Selberg prime number theorem supplies
$\sum_{p\le x}|b_f(p)|^2\log p=x+o_f(x)$; see
\href{https://www.kurims.kyoto-u.ac.jp/~kyodo/kokyuroku/contents/pdf/1511-6.pdf}{Liu and Ye, \emph{The prime number theorem for Rankin--Selberg $L$-functions}, Theorem 3.2}, with finite-place temperedness supplied by Deligne and self-duality by the trivial character. At unramified primes $b_f(p)$ is the normalized Hecke eigenvalue. The omitted prime powers in the prime number theorem contribute $O(\sqrt x\log^2 x)$ by temperedness; finitely many bad factors do not affect its main term. Partial summation, and the convergence of $\sum_p\sum_{a\ge2}(\log p)^2/p^a$, give
\[
  \sum_{n\le x}a_n^2=\tfrac12(\log x)^2+o_f((\log x)^2).
\]
Thus $L^{-2}\sum_{n\le X}a_n^2\delta_{\log n/L}$ converges weakly to $v\,dv$ on $[0,1]$, as is seen from the distribution functions at each $v$. Since $g(y)=L(w\star w)(y/L)$,
\[
 M[P_X,P_X]=\frac{TL^3}{\pi}\int_0^1 v(w\star w)(v)dv
             +o_{f,\eta}(TL^3)+O_\eta(XL^2).
\]
The integral multiplied by two is $\iint|u-v|w(u)w(v)\,du\,dv$. Divide by $A^2L^2 N(I)$, combine the estimates and use $L/\ell=\lambda$. As $X\le T$, all displayed error terms are $o(1)$ after their specified normalizations. We get \eqref{wc:R} for $G+E$, hence $\|G+E\|_{\rm HS}=O(\sqrt{N(I)})$. The trace norm estimate for $E$ gives \eqref{wc:R} for $G$ by the triangle inequality for the Hilbert--Schmidt norm.
\end{proof}

\subsection{The best window and the obstruction}

\begin{proposition}\label{wc:minimum}
For $0<\lambda<\pi/\sqrt2$, among mass-one real square-integrable windows supported on $[-1/2,1/2]$, the unique minimizer of $R_\lambda$ is
\[
 w_*(u)=\frac{\lambda\cos(\sqrt2\lambda u)}
                  {\sqrt2\sin(\lambda/\sqrt2)},\qquad
 R_*(\lambda)=\frac\lambda2+\frac1{\sqrt2}\cot(\lambda/\sqrt2).
\]
The same infimum is approached by the squares of the smooth test functions used in Theorem~\ref{wc:moment}.
\end{proposition}
\begin{proof}
On the mass-one affine hyperplane the first variation at $w$ is twice the pairing against
$F_w(u)=w(u)/\lambda+\lambda\int|u-v|w(v)dv$. For $w=w_*$, the integral of $w_*$ is one and
$F_{w_*}''=w_*''/\lambda+2\lambda w_*=0$ in the interior. By evenness $F_{w_*}$ is constant. Evaluating at $u=1/2$ gives the stated $R_*$, since $\int(1/2-v)w_*(v)dv=1/2$. For every mass-zero variation $b$, its primitive $H(u)=\int_{-1/2}^u b(v)dv$ vanishes at both endpoints. Two integrations by parts yield
\[
 \iint |u-v|b(u)b(v)du\,dv=-2\int H^2 .
\]
The difference in energies is exactly $\lambda^{-1}\int b^2-2\lambda\int H^2$, which by the Dirichlet Wirtinger inequality is at least $(\lambda^{-1}-2\lambda/\pi^2)\int b^2$. This is positive unless $b=0$. Finally $w_*>0$ on the closed interval in the stated range. Multiply its square root by even smooth cutoffs supported in the open interval and equal to one outside boundary layers shrinking to zero. Their squares converge to $w_*$ in $L^2$; the quadratic functional is continuous there at positive mass.
\end{proof}

\begin{proof}[Proof of \eqref{wc:orbits}]
Write $n_+(G)$ for the number of strictly positive eigenvalues. Each real zero contributes one positive direction before restriction, irrespective of its multiplicity. Each orbit off the axis also contributes at most one. Indeed, for an off-line vector $a+ib$ the paired contribution before normalization is $2m(aa^{\mathsf T}-bb^{\mathsf T})$. A pullback cannot increase the positive index, whence $n_+(G)\le U(I')+V(I')$. On the other hand, for $T$ large the trace is positive, and Cauchy--Schwarz on the positive spectrum gives
\[
 (\operatorname{tr}G)^2
 \le\left(\sum_{\gamma_j(G)>0}\gamma_j(G)\right)^2
 \le n_+(G)\operatorname{tr}(G^2).
\]
Here $\gamma_j(G)$ are matrix eigenvalues. By \eqref{wc:count} the enlarged interval adds $O_f(\sqrt T\log T)$ zeros. Theorem~\ref{wc:moment} thus bounds the liminf by $1/R_\lambda(w)$ from below. Take $X=T$, so $\lambda=1/2$, and then approximate the minimizing window using Proposition~\ref{wc:minimum}. Each limit in $T$ is taken for a fixed smooth window. The same lower bound holds on $[0,T)$ by dyadic summation: for any fixed loss one has the corresponding bound on every sufficiently high dyadic interval, and the remaining bounded initial segment is negligible.
\end{proof}

The simple-line inequality of Alp\"oge and Furman, \S6, has the form
\begin{equation}
  N_0^{\rm simple}(I)\ge 4\operatorname{tr}G-2N(I)-
                          \operatorname{tr}(G^2)+o(N(I)).
  \label{wc:AF}
\end{equation}
It applies to the same matrix because the on-line vector at a real ordinate satisfies $\|v_\rho\|^2\le AL^2$ by \eqref{wc:poisson} and the paired blocks have the positive index used above. With positive primes of support $n\le T$ this certificate gives no positive leading fraction at degree two. In fact, any real symmetric $d$-dimensional compression with $\operatorname{tr}G/N\to1$ and $d/N\to\lambda$ satisfies
\[
  \liminf \frac{\operatorname{tr}(G^2)}{N}\ge\lambda^{-1},
\]
since $d\operatorname{tr}(G^2)\ge(\operatorname{tr}G)^2$. When $\lambda\le1/2$ this makes the leading right hand side of \eqref{wc:AF} at most zero, independently of the prime coefficients. The rectangular limiting window has $R_{1/2}=2+1/6=13/6$; even the optimized window has $R_*(1/2)=2.16596385\ldots$.

There is an exact matrix explanation at the rectangular value. Take $N=48m$ and $d=24m$. Consider $24m$ formal off-line pairs, each a two-dimensional form of signature $(1,1)$ and mass two. A linear map may pull each back to $2e_j e_j^{\mathsf T}$ by using its positive direction only. Assign two pairs to each of $m$ coordinates, one pair to each of $22m$ coordinates and none to the final $m$. The resulting matrix has trace $48m$ and squared Hilbert--Schmidt norm $104m$. Thus its two moments give $N$ and $13N/6$, and every pair is off the line in the formal block decomposition. This example concerns the abstract inertia and moment data. The evaluation vectors of a particular modular function satisfy additional relations.

If a longer-polynomial estimate kept the same prime diagonal and controlled every error, the first possible positive fraction in this window class would start at the unique solution
\[
 \tfrac12\lambda_0+2^{-1/2}\cot(\lambda_0/\sqrt2)=2,
 \qquad \lambda_0=0.5501939647\ldots .
\]
Indeed $R_*'(\lambda)=-\tfrac12\cot^2(\lambda/\sqrt2)<0$ in the range at issue. At degree two this would require some support beyond the exponent $2\lambda_0=1.1003879294\ldots$ in $X=T^\theta$. The direct bound $O(XL^2)$ for unequal frequencies in the proof already exceeds its allowable scale there. Rankin--Selberg orthogonality evaluates the equal frequencies; a useful longer estimate would also have to bound close distinct frequencies and justify the other truncations at the enlarged support. This computation distinguishes a specific missing estimate from the separate question of the truth of elliptic-curve GRH.

\clearpage
\section{All-height heat deformation and finite Jensen tests for modular forms}\label{heat:main}

A deformation of a completed modular $L$-function satisfies a Riemann hypothesis for every sufficiently large positive heat time. The time is explicit, and the zeros are simple at every larger time. This chapter fixes the time scale, proves the strip estimate for the modular kernels, and examines its endpoint. We also give a direct positivity counterexample and rigorous finite tests for the undeformed functions of twelve elliptic-curve isogeny classes and the discriminant form.

\subsection{The physical height variable}

Let $f(z)=\sum_{n\ge1}a_n e^{2\pi inz}$ be a normalized holomorphic cuspidal newform of weight $k\ge2$ and level $N$. We allow a unitary nebentypus. In arithmetic normalization put
\[
  L(f,s)=\sum_{n\ge1}a_n n^{-s},\qquad
  \Lambda(f,s)=\left(\frac{\sqrt N}{2\pi}\right)^s\Gamma(s)L(f,s).
\]
The classical newform functional equation is $\Lambda(f,s)=\epsilon_f\Lambda(\bar f,k-s)$, where the Fourier coefficients of $\bar f$ are $\overline{a_n}$ and $|\epsilon_f|=1$. Choose a phase $c$ of modulus one with $c\epsilon_f=\bar c$. We use the variable $z$ in $s=k/2+iz$, so one real unit of $z$ equals one unit of height of an $L$-zero. Define
\begin{equation}
 K(u)=c e^{ku/2}f(i e^u/\sqrt N),\qquad
 H_t(z)=\int_{-\infty}^\infty e^{tu^2}K(u)e^{izu}\,du .
 \label{heat:kernel}
\end{equation}
The undisturbed function is $H_0(z)=c\Lambda(f,k/2+iz)$. Indeed, Mellin integration of the Fourier series gives
\[
 \int_0^\infty f(iy)y^s\,\frac{dy}{y}
       =(2\pi)^{-s}\Gamma(s)L(f,s)
\]
first sufficiently far to the right. The change of variables $y=e^u/\sqrt N$ and analytic continuation prove the identity. At both ends the integral converges absolutely for every $s$, as we now verify.

The exponential decay of a cusp form at infinity gives
$|f(iy)|\le C\exp(-2\pi y)$ for $y$ bounded below by any positive constant. One may see this from absolute convergence of its Fourier expansion and $a_1=1$. The Fricke transformation of $f$ into its dual gives
\[
 e^{-ku/2}f(i e^{-u}/\sqrt N)=
   \epsilon_f e^{ku/2}\overline{f(i e^u/\sqrt N)}
              \qquad(u\in\R).
\]
The chosen phase then gives $K(-u)=\overline{K(u)}$. In particular, there exist $b>0$, $C,B>0$ with
\[
    |K(u)|\le C\exp\{B|u|-b e^{|u|}\}.
\]
This bound dominates all derivatives of the integrand with respect to $z$ and $t$, locally uniformly even for complex $(t,z)$. Hence $H_t$ is entire in the two variables and
\begin{equation}
                  \partial_t H_t(z)=-\partial_z^2 H_t(z).
 \label{heat:pde}
\end{equation}
For real $t$ it is real on the real axis, by conjugate symmetry. It is nonzero as a function since the Fourier transform is injective and the kernel is not identically zero.

For each fixed real $t$, $H_t$ is an entire function of order at most one. To verify the growth statement, absorb $|t|u^2+B|u|$ for large $|u|$ into $b e^{|u|}/2$. For $|z|\le R$, split the integral at $u=0$ and substitute $v=e^{|u|}$. The resulting bound is a constant times $1+C_t^R\Gamma(R+1)$, allowing an additional $e^{C_t R}$ for any fixed bounded interval. Thus $\log\max_{|z|\le R}|H_t(z)|=O_t(R\log(2+R))$. On the real axis $H_t(x)\to0$ at both ends by the Riemann--Lebesgue lemma.

\subsection{Contraction of the full critical strip}

\begin{theorem}\label{heat:real}
Every zero of $H_t$ for $0\le t\le1/8$ belongs to the strip
\[
                |\operatorname{Im}z|\le\sqrt{1/4-2t}.
\]
For $t\ge1/8$ every zero is real. For $t>1/8$ all zeros are simple.
\end{theorem}

We give the contraction argument, including its constants. It is the version of de Bruijn's strip theorem for the order-one functions above.
\href{https://doi.org/10.1215/S0012-7094-50-01720-0}{N.~G. de Bruijn, Duke Math. J. 17 (1950), 197--226.}

\begin{lemma}\label{heat:oneshift}
Suppose $F$ is a nonzero real entire Fourier transform of order at most one, tending to zero at both ends of the real axis. Suppose its zeros belong to $|\operatorname{Im}z|\le D$. Then, for $h>0$, the function
\[
             \tfrac12\bigl(F(z+ih)+F(z-ih)\bigr)
\]
has all its zeros in $|\operatorname{Im}z|\le\sqrt{\max(D^2-h^2,0)}$.
\end{lemma}
\begin{proof}
Hadamard factorization writes $F$ as a real constant times $z^m e^{az}$ and a genus-one canonical product over its nonzero zeros, with $a$ real. To see the reality of the exponential coefficient, divide by the canonical product formed from the real zeros and the conjugate pairs; the quotient is zero-free and real on the real axis up to a constant sign. There is at least one zero, since otherwise the exponential factor alone cannot tend to zero at both ends.

Write $z=x+iy$ with $y>0$ and $y^2+h^2>D^2$. A real zero $v$ gives $|z+ih-v|>|z-ih-v|$. For a pair $v\pm ib$, comparison of the products of the two distances follows from
\begin{align*}
 &|((z+ih-v)^2+b^2)|^2
       -|((z-ih-v)^2+b^2)|^2\\
 &\hspace{20mm}=8yh\bigl((x-v)^2+y^2+h^2-b^2\bigr)>0.
\end{align*}
For a real zero, the exponential in its genus-one factor has the same modulus at $z+ih$ and $z-ih$. For a conjugate pair, the product of the two exponential corrections has this property because the sum of the two reciprocal zeros is real. The same is true of $e^{az}$. When the denominator $F(z-ih)$ is nonzero, the modulus ratio of the two values consequently exceeds one. Each zero or conjugate pair supplies a factor greater than one, and any fixed such factor supplies a strict lower bound for the ratios of all subsequent partial products. If the denominator vanishes, the numerator cannot vanish: a zero in the numerator would have imaginary part $y+h>D$. Thus the displayed average cannot vanish at $z$. Reflection deals with the lower half-plane.
\end{proof}

A single shift-average corresponds inside Fourier integration to multiplication by $\cosh(hu)$. For $n\ge1$ put $h=\sqrt{2t/n}$ and apply the lemma $n$ times. All intermediate kernels satisfy the same type of two-sided decay and conjugate symmetry, giving the order-one bound and the decay on the real axis required at each application. At the $n$th step the strip width is at most $\sqrt{\max(D^2-2t,0)}$. As $n$ increases,
\[
       \cosh(\sqrt{2t/n}\,u)^n\longrightarrow e^{tu^2},
       \qquad \cosh(\sqrt{2t/n}\,u)^n\le e^{tu^2}.
\]
The inequality follows by integrating $\tanh v\le v$ for $v\ge0$. Dominated convergence is uniform for $z$ in compact sets. Hurwitz's theorem excludes zeros of the limit outside the stated strip, since the limiting function is not identically zero. This argument can start from any time $t_0$, and a forward increment $t$ costs $2t$ in the squared strip width.

At time zero the normalized Euler product gives nonvanishing for $\operatorname{Re}s>(k+1)/2$. Here we use the established Ramanujan bounds for holomorphic newforms, or directly absolute convergence and nonvanishing of the unitary Euler product in its right half-plane. The functional equation gives nonvanishing on the opposite side and the completed cuspidal function is entire. In the $z$ variable this puts all zeros in the closed strip $D=1/2$. The contraction just proved gives the first two conclusions of Theorem~\ref{heat:real}.

\subsection{Splitting a multiple zero}

Suppose a real zero $x_0$ at time $t_0$ has multiplicity $m$. Write $A=H_{t_0}^{(m)}(x_0)/m!\ne0$. For $q>0$ tending to zero, the power series in the two variables and \eqref{heat:pde} give
\begin{align}
 q^{-m}H_{t_0+q^2}(x_0+qw)&\longrightarrow A P_m(w),\notag\\
 q^{-m}H_{t_0-q^2}(x_0+qw)&\longrightarrow A i^{-m}P_m(iw),\qquad
 P_m(w)=\sum_{r=0}^{\lfloor m/2\rfloor}
       \frac{(-1)^r m!}{r!(m-2r)!}w^{m-2r}.
 \label{heat:splitting}
\end{align}
For the first limit, the coefficient of a term with $r$ time derivatives and $j$ spatial derivatives contains $H_{t_0}^{(2r+j)}(x_0)$. It consequently vanishes whenever $2r+j<m$. At degree $2r+j=m$ it gives exactly the displayed polynomial. For the second limit replace $q^2$ by $-q^2$. The power series converge uniformly on compact $w$-sets, justifying these limits there.

The polynomials $P_m$ are the monic Hermite polynomials with generating function $\exp(wx-x^2)=\sum_mP_m(w)x^m/m!$. They have $m$ distinct real roots. For a short proof, the recurrence $P_{m+1}=wP_m-2mP_{m-1}$ and $P_0=1$, $P_1=w$ show strict interlacing by induction: evaluate $P_{m+1}$ at consecutive roots of $P_m$, and use its leading sign at both infinities to obtain the remaining two roots. The second limiting polynomial in \eqref{heat:splitting} has purely imaginary distinct roots. If $m\ge2$ at least two of these are nonreal. By Rouché's theorem at least two zeros of $H_{t_0-q^2}$ are nonreal for every sufficiently small $q>0$. In the forward direction one similarly obtains exactly one zero in each small disk around a real rescaled root. Conjugation preserves each disk and forces its single zero to be real and simple.

A multiple zero at any $t_0>1/8$ would now contradict reality for nearby times on its left, completing Theorem~\ref{heat:real}. This is a statement at every height. The argument does not require a numerical lower bound on gaps between consecutive zeros.

\subsection{The Newman constant in this normalization}

Let $\lambda_f$ be the infimum of all real times when $H_t$ has only real zeros. Contraction shows that once one such time is reached all later times have the same property. Local uniform convergence as $t$ varies and Hurwitz show that this set is closed. It is nonempty by Theorem~\ref{heat:real}. Dobner proved that for any nonconstant $L$-function in the extended Selberg class with its prescribed completion, every negative deformation time has noncritical zeros. The normalized series $F(s)=L(f,s+(k-1)/2)$ belongs to that class. Its gamma factor has coefficient of $s$ equal to one, its shift $(k-1)/2$ is nonnegative, its completed function has finite order, and the usual functional equation with conjugation has center $1/2$. There is no polar polynomial for this cuspidal series. Its completion differs by a nonzero constant from $c\Lambda(f,s+(k-1)/2)$ with the phase chosen above.
\href{https://doi.org/10.4064/aa200603-23-7}{A. Dobner, Acta Arith. 201 (2021), 29--62, Theorems 1--3.}

Dobner's critical-line variable $x$ occurs at $s=(1+ix)/2$ for $F$. At time zero his Fourier transform is therefore a scalar multiple of $H_0(x/2)$. Changing $u$ to $2u$ in \eqref{heat:kernel} shows that his time $T$ gives a scalar multiple of $H_{T/4}(x/2)$. In particular the signs of negative and positive time agree, while his numerical constant is four times $\lambda_f$. In physical height coordinates we obtain
\[
  0\le\lambda_f\le1/8,\qquad
  \text{all nontrivial zeros of }L(f,s)\text{ lie on }\operatorname{Re}s=k/2
           \quad\Longleftrightarrow\quad \lambda_f=0 .
\]
The equivalence follows from closedness of the interval of real-zero times, Dobner's lower bound, and the equality $H_0(z)=c\Lambda(f,k/2+iz)$. The gamma factor has no zeros; its poles cancel the trivial zeros of $L(f,s)$. For $t>0$ the theorem concerns the transformed entire functions defined by the Fourier integral. It gives no estimate reducing the upper endpoint $1/8$ to zero for the original modular form.

\subsection{A positive modular-type decay bound does not force real zeros}

A positive even kernel gives positive even moments. The following example determines exactly how far these conditions go, even for double-exponential decay and a transform whose zeros are all confined to a prescribed narrow strip. For $a>0$ let
\[
 q(u)=e^{-\cosh u},\qquad K_a(u)=3q(u)+q(u-a)+q(u+a).
\]
It is strictly positive, real analytic, even, and $K_a(u)\ll_a\exp(-c_a e^{|u|})$. Denote its transform by $T_a$. A change of variables gives
\[
      T_a(z)=(3+2\cos(az))\int_\R q(u)e^{izu}\,du .
\]
Hence $T_a$ has the explicit nonreal zeros
\begin{equation}
      z=\frac{(2j+1)\pi\ \pm i\operatorname{arcosh}(3/2)}{a},
                \qquad j\in\Z .
 \label{heat:falsepositive}
\end{equation}

In fact these are all its nonreal zeros. For completeness set
\[
       Y_z(x)=\int_0^\infty e^{-e^x\cosh u}\cos(zu)\,du .
\]
The other factor is $2Y_z(0)$. For the kernel in this last integral, $\partial_x^2-e^{2x}=\partial_u^2$ as operators applied to that kernel. Two integrations by parts in $u$ give
\[
          -Y_z''(x)+e^{2x}Y_z(x)=z^2Y_z(x).
\]
The derivatives at $u=0$ which occur in the boundary terms vanish, and double-exponential decay controls infinity. Both $Y_z$ and $Y_z'$ decrease fast enough at positive infinity to have zero boundary term after integration in $x$. For example $e^{-r\cosh u}$ with $r=e^x\ge2$ can be bounded by $e^{-r/2}e^{-(r/2)\cosh u}$, absorbing the factor $r\cosh u$ from the derivative in a smaller exponential. This function of $x$ is not identically zero. Indeed, after changing $u=v/\sqrt r$, the integral $e^r\sqrt r\,Y_z(\log r)$ tends to $\int_0^\infty e^{-v^2/2}\,dv$. The bound $\cosh u\ge1+u^2/2$ and $|\cos(zv/\sqrt r)|\le e^{|z|v}$ justify dominated convergence.

If $Y_z(0)=0$, multiply the differential equation by $\overline{Y_z}$ and integrate over $x\ge0$. Integration by parts now yields
\[
        z^2\int_0^\infty |Y_z|^2\,dx=
             \int_0^\infty(|Y_z'|^2+e^{2x}|Y_z|^2)\,dx>0 .
\]
Thus $z$ is real. For $a\ge2\operatorname{arcosh}(3/2)$ the counterexample has every zero in $|\operatorname{Im}z|\le1/2$, while infinitely many of its zeros remain off the real line. The missing property in an argument relying only on strip confinement, kernel positivity, and decay must use further arithmetic or analytic structure.

\subsection{What the central derivatives must satisfy}

For a real entire function $H(z)=\sum_{j\ge0}h_jz^j/j!$, define its Jensen polynomial by
\[
   J_h^{d,n}(X)=\sum_{j=0}^d\binom dj h_{n+j}X^j,
                \qquad d\ge1,\quad n\ge0.
\]
A real polynomial is hyperbolic when all of its complex zeros are real, allowing multiplicities; the zero polynomial is allowed in closure statements. For a nonzero real entire $H$ of order at most one, all its zeros are real if and only if $J_h^{d,0}$ is hyperbolic for every $d$. This gives a criterion for all modular $H_0$, including a complex nebentypus after its phase is fixed.

Here are details of the equivalence. If $H$ has only real zeros, its Hadamard product has real zeros in all its finite partial products; the real exponentials in such products can each be approximated locally uniformly by $(1+bz/M)^M$. A diagonal sequence of real-rooted polynomials therefore tends to $H$ locally uniformly. If $P$ is such a polynomial, write $P(D)$ for the polynomial in $D=d/dX$. Then $P(D)X^d=X^d J_P^{d,0}(1/X)$ where the $J$ on the right is formed using $P^{(j)}(0)$. Each real linear factor $D-b$ preserves real-rootedness. Indeed if $Q$ has simple real zeros $x_j$, then $Q'/Q=\sum_j(X-x_j)^{-1}$ decreases from $+\infty$ to $-\infty$ between successive zeros. For $b\ne0$ there is exactly one additional solution on the outside interval with the appropriate sign, and for $b=0$ Rolle's theorem suffices. Repeated zeros follow by approximation or factoring their multiplicities. This proves hyperbolicity of the Jensen polynomial for $P$, hence for $H$ by local uniform convergence and its resulting derivative convergence. The same reasoning applies to $H^{(n)}$ using derivatives of the approximating polynomials.

Conversely, as $d\to\infty$, $J_h^{d,0}(z/d)\to H(z)$ locally uniformly. Indeed the coefficient multiplier $d(d-1)\cdots(d-j+1)/d^j$ is at most one and tends to one for each fixed $j$. The power series for $H$ converges absolutely on compact disks. Hurwitz on the upper and lower half-planes gives the required reality of every zero.

For a newform with positive root number and positive $K(u)$ one can make a stronger, simpler test by using only its even moments. In this case take $c=1$ and put
\[
  M_{2j}=\int_\R u^{2j}K(u)\,du>0,\qquad
  \gamma_j=\frac{j!}{(2j)!}M_{2j},\qquad
  F(w)=\sum_{j\ge0}\gamma_j w^j/j!.
\]
We assume here also that $K$ is even, as happens for the real eta products below. Then $H_0(z)=F(-z^2)$ and $F$ has order at most $1/2$. Since all its coefficients are positive, $F$ has no nonnegative real zero. Thus all zeros of $H_0$ are real exactly when every $J_\gamma^{d,0}$ has only negative real roots. Pointwise positivity of $K$ alone does not imply this. For example the discriminant test at degree two already requires $\gamma_1^2\ge\gamma_0\gamma_2$, i.e. $3M_2^2\ge M_0 M_4$. Ordinary moment Cauchy--Schwarz supplies the different inequality $M_2^2\le M_0M_4$.

For these even forms the widely used fixed-degree asymptotic is known: Wagner's theorem for modular $L$-functions shows that for each fixed $d$ the polynomial $J_\gamma^{d,n}$ is hyperbolic for every sufficiently large shift $n$. It does not give hyperbolicity of $J_\gamma^{d,0}$ as $d$ grows without bound. The two quantifiers correspond to different entries in the array of polynomials. \href{https://doi.org/10.1515/forum-2019-0104}{I. Wagner, Theorem 1.3 and Corollary 1.5.}
\href{https://doi.org/10.1073/pnas.1902572116}{M. Griffin, K. Ono, L. Rolen and D. Zagier, the Jensen--Pólya criterion and Hermite limits.}

\begin{lemma}\label{heat:triangle}
Suppose $\gamma_j>0$ for $0\le j\le D$. If $J_\gamma^{D,0}$ has $D$ distinct negative roots, then $J_\gamma^{d,n}$ has $d$ distinct negative roots for every $d\ge1$, $n\ge0$ with $d+n\le D$.
\end{lemma}
\begin{proof}
From $J=J_\gamma^{m,n}$ its two immediate descendants are given by
\[
       J_\gamma^{m-1,n+1}=\tfrac1m J',\qquad
       J_\gamma^{m-1,n}=J-\tfrac Xm J'.
\]
The first has one root strictly between consecutive roots of $J$. To see the second, reverse $J$ to $R(X)=X^mJ(1/X)$, differentiate, then reverse in degree $m-1$. This gives $mJ-XJ'$. Reversal and differentiation preserve distinct negative roots in these operations since all endpoint and leading coefficients are positive. Each step lowers the degree and either increases the shift by one or preserves it. Starting from $(D,0)$ produces exactly the stated triangle.
\end{proof}

\subsection{Interval certificates at the central point}

We apply Lemma~\ref{heat:triangle} to the discriminant form $\Delta=\eta^{24}$ and to the twelve weight-two eta-quotient newforms classified by Martin and Ono. For precision, the table lists the arguments and exponents of $f(z)=\prod_d\eta(dz)^{r_d}$. A bare argument has exponent one. The table uses $d^{[r]}$ for an argument $d$ with exponent $r$.
\href{https://doi.org/10.1090/S0002-9939-97-03928-2}{Y. Martin and K. Ono, Proc. Amer. Math. Soc. 125 (1997), 3169--3176, Theorems 1 and 2.}
\[
\begin{array}{r|l@{\qquad}r|l}
 N&d^{[r_d]}&N&d^{[r_d]}\\ \hline
 1&1^{[24]}&11&1^{[2]},11^{[2]}\\
 14&1,2,7,14&15&1,3,5,15\\
 20&2^{[2]},10^{[2]}&24&2,4,6,12\\
 27&3^{[2]},9^{[2]}&32&4^{[2]},8^{[2]}\\
 36&6^{[4]}&48&4^{[4]},12^{[4]},2^{[-1]},6^{[-1]},8^{[-1]},24^{[-1]}\\
 64&8^{[8]},4^{[-2]},16^{[-2]}&80&4^{[6]},20^{[6]},2^{[-2]},8^{[-2]},10^{[-2]},40^{[-2]}\\
 144&12^{[12]},6^{[-4]},24^{[-4]}&&
\end{array}
\]
All factors are positive on the imaginary axis. In each row the exponents are invariant under $d\mapsto N/d$. The eta transformation $\eta(i/y)=\sqrt y\,\eta(iy)$ for $y>0$ therefore gives $K(-u)=K(u)$ with $c=1$. For detail, $\sum r_d=2k$ and $\prod d^{r_d}=N^{k}$ by this symmetry. Multiplying the square roots from the eta transformation precisely cancels the weight $e^{-ku/2}$ appearing at $-u$. Thus the positive central moment criterion applies to each listed form, including those written as genuine quotients.

We give a reproducible analytic error bound for all moments used in the certificate. Let $D$ denote its maximum degree. For $0\le n\le D$ put
\[
   b=6+2n,\qquad F_n(w)=w^{2n}e^{kw/2}f(i e^w/\sqrt N).
\]
This function is holomorphic for $|\operatorname{Im}w|<\pi/2$, because every eta argument then lies in the upper half-plane and the eta function has no zero there. At a point $w=x+iv$ with $x\ge0$, $|v|\le1$, the nome of $\eta(i d e^w/\sqrt N)$ has absolute value
$\rho_d=\exp(-2\pi d e^x\cos(v)/\sqrt N)\le e^{-1/2}$. Indeed $d/\sqrt N\ge1/6$ throughout the table and $\cos(1)>1/2$. The absolute value of the logarithm of the modulus of the infinite-product part of one eta factor is at most
\[
 \sum_{j\ge1}\frac{\rho_d^j}{1-\rho_d^j}
      \le \frac{\rho_d}{(1-\rho_d)^2}<4.
\]
This absolute bound handles negative eta exponents as well. Now $\sum_d|r_d|\le24$, $\sum_d d r_d=24$, and $N\le144$. The leading exponential of the eta product consequently has modulus
$\exp(-2\pi e^x\cos(v)/\sqrt N)\le\exp(-e^x/4)$. Since $k/2\le6$ and $|x+iv|\le x+1\le e^x$, we obtain
\begin{equation}
     |F_n(x+iv)|\le \exp(96+b|x|-e^{|x|}/4)
             \qquad (x\in\R,\ |v|\le1).
 \label{heat:stripmaj}
\end{equation}
For the left half-strip this follows by the already proved evenness, extended by the identity theorem.

Use mesh $h=1/L$ and cutoff $U$ with $UL$ integral. The full even trapezoidal sum truncated at the endpoints is
\[
   A_n=h F_n(0)+2h\sum_{j=1}^{UL}F_n(jh).
\]
If $e^U>4b$ then
\begin{equation}
 |M_{2n}-A_n|\le
   \frac{4e^{96}4^b(b-1)!}{e^{2\pi L}-1}
       +2e^{96-e^U/8}8^b(b-1)! .
 \label{heat:error}
\end{equation}
To prove the first term, use Poisson summation on a mesh $h$ for the full real-line trapezoidal sum. The restriction of $F_n$ to the real line is Schwartz, by Cauchy's estimates in a smaller strip and \eqref{heat:stripmaj}. For the Fourier frequency $2\pi j/h$ move the line of integration to imaginary part $-1$ if $j>0$, and to $1$ if $j<0$, with the transform convention $e^{-i\xi x}$. The frequency factor has modulus $e^{-2\pi |j|L}$. By \eqref{heat:stripmaj}, the remaining integral is bounded by
$2e^{96}\int_0^\infty e^{bx-e^x/4}dx\le2e^{96}4^b(b-1)!$, using $v=e^x/4$. Sum the geometric series in both frequency signs. For the second term, the majorant decreases for $x\ge U$, since $e^U/4>b$. Hence $h\sum_{j>UL}e^{bjh-e^{jh}/4}$ is at most the integral of that majorant over $x\ge U$. Separate $e^{-e^x/4}$ as two equal factors, bound one by $e^{-e^U/8}$ and integrate the other using $v=e^x/8$. This proves \eqref{heat:error}.

\begin{theorem}\label{heat:certificate}
For all thirteen rows, every central Jensen polynomial $J_\gamma^{d,n}$ with $d\ge1$, $n\ge0$, $d+n\le256$ has $d$ distinct negative roots. This is a total of $427{,}648$ certified polynomials, including $394{,}752$ for the twelve elliptic-curve rows.
\end{theorem}
\begin{proof}
For the degree-$256$ calculation, take $L=1024$, $U=11$, and 3072-bit ball arithmetic. At each node the implementation evaluates the eta quotient in complex balls using Arb, projects onto the real line, forms the power $u^{2n}$ recursively, and sums with nonnegative real weights. Widening the result by both terms of \eqref{heat:error} gives an enclosure for $M_{2n}$. Multiplication by the exact rational $\binom{256}{n}n!/(2n)!$ gives the enclosure for the coefficient of $X^n$.

For each form the certificate supplies $257$ increasing dyadic rationals $x_0<\cdots<x_{256}=0$ with denominator $2^{32}$. Interval Horner evaluation at these points proves that the sign at $x_j$ is $(-1)^{256-j}$. There is at least one zero in each of the $256$ disjoint negative intervals. As the degree is exactly $256$, all are simple and these are all the roots. Lemma~\ref{heat:triangle} gives $256\cdot257/2=32{,}896$ strict hyperbolicity conclusions for each form. The exact dyadic numerators and the checker are included with the manuscript. Root approximations serve only to propose dyadic test points. The replay uses the recorded rational points and interval sign evaluations alone.
\end{proof}

These finite conclusions concern central derivative polynomials. They assert neither $427{,}648$ zeros of the modular $L$-functions nor a strip-height bound for their zeros. The obstruction that remains at $t=0$ is control as the polynomial degree increases without bound. Each degree introduces a new moment and the sufficient range of shifts in the known fixed-degree theorem may change with that degree.

\subsection{A simultaneous limitation of all finite tests}

\begin{theorem}\label{heat:jointlimit}
Fix $D\ge1$, $T>0$ and $0<b\le1/2$. There exists a strictly positive, even, real analytic kernel with two-sided double-exponential decay whose Fourier transform has all the following properties. Every zero is in $|\operatorname{Im}z|\le b$, all zeros with $|\operatorname{Re}z|\le T$ are real, and all the folded Jensen polynomials with $1\le d$, $d+n\le D$ have distinct negative roots. Nevertheless the transform has infinitely many nonreal zeros.
\end{theorem}
\begin{proof}
We first verify strict Jensen hyperbolicity for the unshifted base $q(u)=e^{-\cosh u}$. Its transform $H$ is real and even and has only real zeros by the energy identity above. It has infinitely many zeros. Otherwise its order at most one and Hadamard factorization would give $H(z)=P(z)e^{az+c}$ with $P$ a polynomial and $a$ real. No such nonzero function tends to zero on the real axis at both ends. Since $H(0)>0$, the folded function $F$ defined by $H(z)=F(-z^2)$ has infinitely many negative zeros, order at most $1/2$, and positive Taylor coefficients. Its Hadamard product has genus zero and is a constant times the product of its normalized linear factors.

For any $D$ remove $D$ such factors, writing $F(z)=R(z)G(z)$, where $R$ has $D$ negative roots counted with multiplicity. Then $G$ is also a locally uniform limit of polynomials with negative real roots and $G(0)>0$. The operator $G(\partial_X)$ is well defined on polynomials by a finite Taylor sum. Acting on $X^D$ gives a real-rooted polynomial of degree $D$, by the Jensen principle proved above. Now apply the $D$ factors of $R(\partial_X)$ to this polynomial. Each is a nonzero scalar times $\partial_X+r$ for some $r>0$. At an old root of multiplicity $h$ it leaves multiplicity exactly $h-1$. Every new root is simple as follows from strict decrease of the logarithmic derivative between consecutive distinct roots, including the appropriate outer interval. No other roots occur, by degree. Thus after all the factors have acted every root of $F(\partial_X)X^D$ is simple. Its reciprocal polynomial is $J_\gamma^{D,0}$, with positive constant and leading coefficient. Its zeros are therefore distinct and negative.

Now set
\[
 K_{a,b}(u)=2\cosh(ab)q(u)+q(u-a)+q(u+a).
\]
Its transform equals $(2\cosh(ab)+2\cos(az))H(z)$, whose complete set of extra zeros is $(2j+1)\pi/a\pm ib$, $j\in\Z$. Taking $a<\pi/T$ places every extra zero beyond the prescribed finite height. As $a\downarrow0$, every one of the first $2D$ moments converges to four times the corresponding base moment, by changing variables in the two shifted integrals or by dominated convergence. The coefficients of $J^{D,0}$ consequently converge to those of a polynomial with $D$ distinct negative roots. Choose rational test points separating those roots with strictly alternating signs; continuity preserves the signs and hence all $D$ strict negative roots for sufficiently small $a>0$. Lemma~\ref{heat:triangle} supplies every other test in the finite triangle. All these conditions can be satisfied by the same $a$.
\end{proof}

This example has no asserted Euler product or modular transformation law. It specifies exactly why finite zero counts, finite Jensen tests and a global strip result do not by themselves establish the arithmetic statement. A proof for the undeformed modular $L$-functions must also use a property which distinguishes their kernels from this family, or control an infinite limiting process.

\clearpage
\section{Critical zeros after a heat time decreasing with height}\label{hs:section}
For the Riemann \(\xi\)-function, Ki, Kim and Lee proved that at each positive heat time all sufficiently high zeros are real and simple. D.~H.~J. Polymath obtained effective estimates for the same deformation. Here the first Dirichlet coefficient of a holomorphic newform gives a direct proof with a heat time of size \(1/\log x\) at physical height \(x\). Combined with strict motion of an outermost zero, it improves the non-strict \(1/8\) of Section~\ref{heat:main} for every fixed form. The endpoint at time zero is precisely the open question.
\href{https://doi.org/10.1016/j.aim.2009.04.003}{H.~Ki, Y.~Kim and J.~Lee, Adv. Math. 222 (2009), 281--306.}
\href{https://doi.org/10.1007/s40687-019-0193-1}{D.~H.~J. Polymath, Res. Math. Sci. 6 (2019), 31.}

\begin{theorem}\label{hs:large}
For every normalized holomorphic newform \(f\) of weight \(k\ge2\) there are constants \(K_f,X_f>0\) with the following property. If \(0<\tau\le1/8\), then every zero \(z\) of \(H_\tau\) with \(x=|\Re z|\ge X_f\) and \(\tau\log x\ge K_f\) is real and simple. At unbounded positive heights and each fixed \(\tau>0\) up to \(1/8\), two consecutive zeros have spacing asymptotic to \(\pi/\log x\). The physical Newman constant satisfies
\[
                       0\le\lambda_f<1/8 .
\]
More generally, if all initial zeros satisfy \(|\Im z|\le D\) with \(0<D\le1/2\), then \(\lambda_f<D^2/2\).
\end{theorem}

\subsection*{The first Dirichlet coefficient dominates two half integrals}
Put \(c=k/2\), \(A=\sqrt N/(2\pi)\) and \(s=c+iz\). Retain the phase \(\omega\) of Section~\ref{heat:main}, chosen so that \(\omega\Lambda(c+iz)\) is real on the real \(z\)-axis. A Gaussian integration of the Mellin kernel gives
\[
 H_\tau(z)=B_\tau(z)+\overline{B_\tau(\bar z)},\qquad
 B_\tau(z)=\frac{\omega}{\sqrt{4\pi\tau}}\int_0^\infty
                      e^{-u^2/(4\tau)}\Lambda(s+u)\,du.
\]
Indeed Gaussian averaging over the whole real \(u\)-axis multiplies the Fourier kernel at \(v\) by \(e^{\tau v^2}\). Its double-exponential decay justifies Fubini for complex \(z\). The negative half becomes the reflected term by the functional equation.

For positive \(x=\Re z\) large and \(|\Im z|\le3/5\), write \(L=\log|As|\), \(\alpha=\arg s\). Put
\[
        M_\tau(z)=\omega A^s\Gamma(s)
                    \exp\{\tau(\log A+\log s)^2\}.
\]
We show that given any \(\epsilon>0\) one can choose \(K_f,X_f\), also depending on \(\epsilon\), such that
\begin{equation}\label{hs:approx}
       |B_\tau(z)-M_\tau(z)|\le\epsilon |M_\tau(z)|
       \qquad (x\ge X_f,\ \tau\log x\ge K_f,\ 0<\tau\le1/8).
\end{equation}
The same proof on the negative-height half-plane uses the principal argument near \(-\pi/2\).

All ensuing estimates are uniform on these horizontal strips. Stirling's formula, or its differentiated form integrated along a segment, gives
\[
 \frac{\Gamma(s+u)}{\Gamma(s)}=
 e^{u\log s}\left(1+O_f((1+\log x)^2/x)\right),
                    \qquad 0\le u\le 8\log x.
\]
For the remaining \(u\), integrate \(\Re\psi(s+h)\le\log(|s|+h)+O_f(1)\), valid uniformly for \(h\ge0\) by Stirling since \(\Re s\ge2/5\) and \(|\Im s|=x\). Thus
\[
 \log\left|A^u\Gamma(s+u)/\Gamma(s)\right|
       \le u(L+C_f)+u^2/x .
\]
For \(u\ge8\log x\) and \(x\) large this and the Gaussian give a bound \(e^{-u^2/(8\tau)}\). The Dirichlet series is bounded there by absolute convergence. The integrated, normalized tail is \(O(e^{-c_1(\log x)^2})\) for some \(c_1>0\); this estimate is relative to \(|\omega A^s\Gamma(s)|\), hence also to \(|M_\tau|\) up to a bounded factor.

Choose \(h_0=h_0(f)\) so large that the normalized first coefficient \(a_1=1\) and absolute convergence give
\[
       L(f,s+u)=1+O_f(2^{-u/2}) \qquad (u\ge h_0).
\]
For instance bound the sum over \(n\ge2\) by \(2^{-u/2}\sum_{n\ge2}|a_n| n^{-\Re s-u/2}\). For \(0\le u\le h_0\), the polynomial growth of the newform \(L\)-function on a fixed vertical strip gives \(|L(f,s+u)|\ll_f x^{C_f}\). This is the usual convexity consequence of its functional equation and absolute convergence.

In units of \(|M_\tau|\), the contribution of \(0\le u\le h_0\), both with and without \(L(f,s+u)\), is
\[
   O_f\!\left(\tau^{-1/2}
         \exp\{(h_0+C_f+1)\log x-\tau L^2+\tau\alpha^2\}\right).
\]
Since \(L=\log x+O_f(1)\) and \(\tau\le1/8\), this tends to zero uniformly on increasing \(K_f>h_0+C_f+3\) if necessary. For the part containing \(2^{-u/2}\), completion of the square gives a relative bound
\[
   O_f\bigl(\exp\{-\tau L\log 2+O(\tau)\}\bigr).
\]
Here one bounds the absolute integral by the whole-line Gaussian with linear exponent \(u(L-\tfrac12\log 2)\). Integrating the Stirling error against the Gaussian with linear exponent \(Lu\) costs \(O_f((1+\log x)^2/x)\), because the ratio between the corresponding absolute Gaussian mass and \(|e^{\tau(\log A+\log s)^2}|\) is \(e^{\tau\alpha^2}=O(1)\). Finally one can extend the modeled integral to negative \(u\): its omitted mass relative to \(|M_\tau|\) is \(O(e^{-\tau L^2+\tau\alpha^2})\), since \(L>0\). The whole-line Gaussian with complex exponent \(u(\log A+\log s)\) is exactly the multiplier in \(M_\tau\). This proves \eqref{hs:approx}. For fixed \(\tau>0\) its relative error tends to zero.

\subsection*{One zero in each phase interval}
On the positive real axis choose a continuous argument \(\phi(x)=\arg M_\tau(x)\) and put \(\ell_x=\log(Ax)\). Uniformly as \(x\to\infty\),
\[
 \phi'(x)=\Re\{\log A+\psi(c+ix)+2\tau\log(A(c+ix))/(c+ix)\}
            =\ell_x+O_f((1+\log x)/x).
\]
In particular \(\phi\) is eventually increasing without bound. For bounded complex \(w\), logarithmic differentiation and Stirling give
\[
 \frac{M_\tau(x+w)}{M_\tau(x)}
     = \exp\{(-\pi/2+i\ell_x)w\}\,
            \left(1+O_f((1+\log x)/x)\right).
\]
In a phase interval with endpoints \(\phi^{-1}(j\pi)\), \(\phi^{-1}((j+1)\pi)\), let \(x_j=\phi^{-1}((j+1/2)\pi)\). The two endpoints are \(x_j+(\pm\pi/2+o(1))/\ell_{x_j}\), uniformly in \(\tau\). Compare on the enclosing rectangle with these vertical edges and \(|\Im z|\le3/5\) to
\[
  G(z)=2(-1)^{j+1}|M_\tau(x_j)|\,e^{-\pi(z-x_j)/2}
                          \sin\bigl(\ell_{x_j}(z-x_j)\bigr).
\]
On the horizontal edges the quotient of \(|\sin(a+ib)|\) by \(\cosh b\) tends uniformly to one since \(|b|=(3/5)\ell_{x_j}\). On the vertical edges the same quotient tends uniformly to one since \(a=\pm\pi/2+o(1)\) and \(|\sin(a+ib)|^2=\cosh^2 b-\cos^2 a\). Thus \(|G|\) exceeds, for example, three quarters of the sum of the two local exponential moduli, for large \(x_j\). Take \(\epsilon=1/20\) in \eqref{hs:approx}. Its error and the two local linearization errors total less than \(|G|\) on the boundary. Rouché's theorem gives exactly one zero with multiplicity, as \(G\) has one. Conjugation preserves the rectangle, so the zero is real and simple. Raise \(K_f,X_f\) by a fixed amount to ensure the first and last neighboring rectangles satisfy \eqref{hs:approx} whenever a requested ordinate does. Section~\ref{heat:main} confines all zeros to \(|\Im z|\le1/2\), so the argument exhausts that end of the zero set. At the negative end apply the same proof to the dual form, whose deformed function reflects the real ordinates. For fixed \(\tau>0\) the relative error tends to zero; at the real zero the sine consequently tends to zero. Its location is \(x_j+o(1/\ell_{x_j})\), proving the asserted spacing.

\subsection*{Strict contraction at a nonreal outer zero}
Put \(t_*=D^2/2\). The strip theorem gives \(|\Im z|^2\le D^2-2t\) for \(0<t<t_*\). If all zeros are real for any such \(t\), the strict bound on \(\lambda_f\) follows. Otherwise, in any closed interval of positive times inside \((0,t_*)\), the preceding theorem confines all nonreal zeros to one compact rectangle, uniformly in time.

One can choose an interior time \(t_0\) at which all those zeros are simple. For details fix an interior time and a slightly larger rectangle with boundary free of zeros, persisting for nearby times. There are only finitely many initial zeros therein. At any zero of multiplicity \(m\), real or complex, the local rescaling in \eqref{heat:splitting} at time increment \(q^2>0\) gives the Hermite polynomial with \(m\) distinct roots. On small disjoint circles about these rescaled roots Rouché gives exactly one zero each. Simultaneously at the finitely many centers this yields a nearby time with simple zeros. Compactness and the boundary check exclude other zeros entering the smaller rectangle.

At \(t_0\) let \(z_j\) be a zero of maximum positive imaginary part \(v\). The heat equation and the implicit-function theorem give \(z_j'=H''(z_j)/H'(z_j)\). The real genus-one product used in Section~\ref{heat:main} now gives
\[
    \Im z_j'=-2\sum_{h\ne j}\frac{v-\Im z_h}{|z_j-z_h|^2}<-\frac1v .
\]
In the logarithmic derivative at the removed simple zero the real exponential coefficient drops out. The genus corrections \(1/z_h\), including the one for \(j\), sum with zero imaginary part in conjugate pairs; a zero at the origin contributes through its ordinary polynomial factor. The displayed imaginary-part series converges absolutely by strip confinement and order one. Every summand before multiplication by \(-2\) is nonnegative. The conjugate of \(z_j\) gives exactly \(-1/v\), and any of the real zeros established above gives a further strictly negative term.

There are only finitely many maximizing zeros. On a small time interval they have analytic branches, and compactness keeps all lower branches initially separated from their maximum. Therefore the right derivative of the squared maximum at \(t_0\) is the maximum of \(2v\Im z_j'\) over the maximizing branches, strictly less than \(-2\). For a sufficiently small \(h>0\) the actual squared strip width at \(t_0+h\) is strictly less than \(D^2-2(t_0+h)\). Starting the strip contraction afresh there proves all zeros real at a time strictly less than \(t_*\). Dobner's lower bound in Section~\ref{heat:main} gives the other endpoint.

\clearpage
\section{Positive heat time for general automorphic completions}\label{hg:section}

Let $\pi$ be a fixed unitary cuspidal representation of $\mathrm{GL}_n$ over a number field $F$, and suppose its completed standard function is entire. Put $d=n[F:\Q]$. In unitary normalization write
\[
 \Lambda(s)=G(s)L(s,\pi),\qquad G(s)=Q^{s/2}\prod_{\nu=1}^d\Gamma_{\R}(s+\mu_\nu),
 \qquad\Gamma_{\R}(w)=\pi^{-w/2}\Gamma(w/2).
\]
Constants independent of $s$ in the completion can be included in $G$. The gamma duplication formula accounts for the complex places. Choose $|c|=1$ so that $H_0(z)=c\Lambda(1/2+iz)$ is real on the real axis. This follows from the functional equation with the conjugate contragredient, including when $\pi$ is not self-dual. In physical height coordinates define
\begin{equation}\label{hg:heat}
 H_t(z)=\frac{1}{\sqrt{4\pi t}}\int_{\R} e^{-u^2/(4t)} H_0(z-iu)\,du
 \qquad(t>0).
\end{equation}
A Fourier kernel containing $e^{izu}$ is multiplied by $e^{tu^2}$ in this convention.

\begin{theorem}\label{hg:main}
Fix $t_{\max}>0$. There are constants $C,X_0$ depending on the completion and $t_{\max}$ such that, for $0<t\le t_{\max}$, every zero $z$ of $H_t$ satisfying
\[
                |\Re z|\ge X_0,\qquad t\log|\Re z|\ge C
\]
is real and simple. For fixed $t>0$ the spacing of consecutive zeros at either end of the real axis, measured at height $|x|$, is $(2\pi+o(1))/(d\log|x|)$. Moreover there is a number $t_\pi<1/8$ such that every zero of $H_t$ is real and simple for all $t\ge t_\pi$.
\end{theorem}

For the Riemann function the high-zero reality and the strict improvement of the elementary upper bound are due to \href{https://doi.org/10.1016/j.aim.2009.04.003}{Ki, Kim and Lee}. \href{https://arxiv.org/abs/1904.12438}{Polymath} gives effective zeta approximations, and \href{https://arxiv.org/abs/2610.27343}{Freire et al.} treat primitive real Dirichlet characters. \href{https://doi.org/10.4064/aa200603-23-7}{Dobner} treats general extended-Selberg-class deformations, including the opposite, negative-time regime. The proof below allows complex archimedean shifts, keeps the height-dependent time threshold, and uses only analytic properties of the individual completed standard function. It gives no numerical improvement uniform in the conductor.

\subsection{Confinement by real Gaussian averages}
The standard continuation, functional equation and order-one bound on the completion, absolute convergence of its normalized Dirichlet series in $\Re s>1$, and polynomial growth of $L$ in any fixed vertical strip are the automorphic inputs. See \href{https://arxiv.org/pdf/2103.05634}{Humphries--Thorner, \S3} and \href{https://bookstore.ams.org/coll-53}{Iwaniec--Kowalski, Chapter 5}. Write $z=x+iy$. The bound
\begin{equation}\label{hg:horizontal}
 \sup_{x\in\R}|H_0(x+iy)|
       \le \exp\{C_\pi(1+|y|)\log(2+|y|)\}
\end{equation}
follows on the half-plane $1/2-y$ sufficiently large by $|\Gamma(a+ib)|\le\Gamma(a)$ for $a>0$ and absolute convergence, on the opposite half-plane by the functional equation, and in the intervening strip by Stirling's formula and polynomial growth. On each fixed horizontal line $H_0(x+iy)\to0$ as $x\to\pm\infty$. Formula~\eqref{hg:horizontal} justifies \eqref{hg:heat} locally uniformly in $z$ and implies the same decay on each fixed horizontal line for $H_t$, by dominated convergence. Convolution in \eqref{hg:heat} has the semigroup property. Gaussian integration by parts gives $\partial_t H_t=-\partial_z^2H_t$; the integral is holomorphic in $t$ near every positive real time. For every $0<\epsilon<1$, the order-one bound $|H_0(z)|\le\exp(C_{\pi,\epsilon}(1+|z|)^{1+\epsilon})$ gives the corresponding bound for $H_t$, by separating $|z|^{1+\epsilon}$ and $|u|^{1+\epsilon}$ under the integral. Thus $H_t$ is real entire of order at most one.

All zeros of $H_0$ lie in $|\Im z|\le1/2$, by Euler-product nonvanishing and reflection. The strip contraction of Lemma~\ref{heat:oneshift} applies here without an explicit Fourier kernel. Indeed its proof uses the conjugate-paired Hadamard product, the order-one bound, and decay on the real axis. For a starting time $v\ge0$, average $H_v(z+ih)$ and $H_v(z-ih)$ successively $N$ times with $h=\sqrt{2t/N}$. Each finite average is real entire of order at most one and tends to zero on the real axis. If the starting strip has half-width $D$, iteration gives half-width $\sqrt{\max(D^2-2t,0)}$. The limiting average is $H_{v+t}$: the sum of $N$ independent signs times $h$ converges in distribution to a real Gaussian of variance $2t$, and the tail probabilities are bounded by $2e^{-R^2/(4t)}$ uniformly in $N$. This follows from $(\cosh(hs))^N\le e^{ts^2}$ and exponential Markov. The order bound with any $1+\epsilon<2$ gives uniform integrability and local uniform convergence on compact $z$-sets. The finite shift averages are nonzero if their predecessor is: the strict modulus comparison outside its zero strip in Lemma~\ref{heat:oneshift} gives this as well. The saddle estimate below guarantees that each limiting function is nonzero, independently of strip contraction. Hurwitz's theorem therefore yields
\begin{equation}\label{hg:basicstrip}
        |\Im z|\le D(t)=\sqrt{\max(1/4-2t,0)}
\end{equation}
at every zero of $H_t$, and gives the same contraction from any later starting strip.

\subsection{The contribution from positive real translates}
Set $s=1/2+iz$, $X=|\Re z|$, and restrict temporarily to $|\Im z|\le3/5$. Put
\[
 b(s)=G'(s)/G(s),\quad M_t(z)=cG(s)e^{t b(s)^2},\quad
 B_t(z)=\frac{c}{\sqrt{4\pi t}}\int_0^\infty e^{-u^2/(4t)}\Lambda(s+u)\,du .
\]
For $J^*(z)=\overline{J(\bar z)}$ the functional equation gives $H_t=B_t+B_t^*$. We show that, for $t\log X$ sufficiently large and $X\to\infty$,
\begin{equation}\label{hg:one}
  B_t(z)=M_t(z)\left(1+
        O_{\pi,t_{\max}}\left(e^{-c_1t\log X}+\frac{(\log X)^2}{X}\right)\right)
\end{equation}
with $c_1>0$, uniformly on the indicated strip. The integral of $G(s)e^{b(s)u}$ on the full Gaussian line is exactly $G(s)e^{t b(s)^2}$ even for complex $b$. The estimates below bound the three errors in replacing its integrand by the positive-translate half.

Stirling's formula, uniformly for both signs of $\Re z$, gives
\begin{equation}\label{hg:stir}
 \begin{gathered}
  A:=\Re b(s)=\tfrac d2\log X+O_\pi(1),\qquad |\Im b(s)|\ll_\pi1,\\
  b'(s)=O_\pi(X^{-1}),\qquad b''(s)=O_\pi(X^{-2}).
 \end{gathered}
\end{equation}
For $0\le u\le U=C_0\log X$ with $C_0$ fixed, the same derivative estimates on the segment to $s+u$ yield
$G(s+u)=G(s)e^{u b(s)}(1+O_\pi(U^2/X))$. After taking absolute values in the linear model and dividing by $|e^{t b(s)^2}|$, the Gaussian becomes a fixed bounded factor $e^{t(\Im b(s))^2}$ times the normal probability measure with mean $2tA$ and variance $2t$. Denote the corresponding real normal variable by $V$.

Choose a fixed $u_0$ so large that for $u\ge u_0$, absolute convergence at $\Re s+u/2$ gives $L(s+u)=1+O_\pi(2^{-u/2})$. Its error under the tilted Gaussian is therefore
\[
   O_{\pi,t_{\max}}\bigl(\mathbb E\,2^{-V/2}\bigr)
    =O_{\pi,t_{\max}}\bigl(e^{-tA\log2+t(\log2)^2/4}\bigr).
\]
On $0\le u\le u_0$ we have $|L(s+u)|\ll_\pi X^{C_2}$ by fixed-strip growth. Choose the lower bound on $t\log X$ sufficiently large that $u_0\le tA$ and $X^{C_2}\mathbb P(V\le u_0)\ll X^{-2}$. Indeed this probability is at most $e^{-tA^2/4}$, and \eqref{hg:stir} makes the required bound uniform down to $t=C/\log X$. This bound also accounts for the negative portion of the linear model. On $[u_0,U]$ the relative gamma Taylor error costs $O_{\pi,t_{\max}}(U^2/X)$ since $L$ is bounded there. On $[0,u_0]$ the small probability already absorbs the possible polynomial growth.

For completeness the remaining positive tail is negligible for both integrals. Stirling's formula for the digamma function in the sector containing all $s+u+\mu_\nu$, $u\ge0$, gives
\[
       |G(s+u)/G(s)|\le
       \exp\{u(C_3+(d/2)\log(X+u+2))\}.
\]
For $C_0\log X\le u\le X$, take $C_0$ sufficiently large in terms of $t_{\max},d,C_3$ so that the exponent here is at most $u^2/(8t)$. For $u>X$ the same holds for sufficiently large $X$ since $\log(2u+2)=o(u)$. The Gaussian tail after normalization is $O_{\pi,t_{\max}}(X^{-2})$; the linear model satisfies the same estimate by \eqref{hg:stir}. These inequalities prove \eqref{hg:one}. In particular its relative error is at most $1/100$ throughout the strip after increasing the constants.

\subsection{One zero per phase interval}
On either high half-strip choose a holomorphic branch $\ell_t(z)=\log M_t(z)$, and define
\[
       \phi_t(z)=\frac{\ell_t(z)-\ell_t^*(z)}{2i},\qquad
       h_t(z)=\frac{\ell_t(z)+\ell_t^*(z)}2 .
\]
At real $x$, $\phi_t(x)$ is the phase of $M_t(x)$ and
\[
  M_t+M_t^*=2e^{h_t}\cos\phi_t,\qquad
  \phi_t'(x)=\Re\{b(s)+2t b(s)b'(s)\}=\tfrac d2\log|x|+O_{\pi,t_{\max}}(1).
\]
Moreover $\ell_t''(z)=O_{\pi,t_{\max}}(1/X)$ from \eqref{hg:stir}, hence $\phi_t(x+iy)=\phi_t(x)+iy\phi_t'(x)+O_{\pi,t_{\max}}(1/X)$ on $|y|\le3/5$. In particular $\Re\phi_t'(x+iy)>0$ for large $X$. The identity $\Im\phi_t(x+iy)=\int_0^y\Re\phi_t'(x+iv)\,dv$ shows that $\phi_t(z)$ can be real only on the real axis in these regions.

Use successive real points $x_j,x_{j+1}$ where $\phi_t(x_j)=j\pi$, and form the rectangle with these vertical sides and horizontal sides at $\Im z=\pm 11/20$. On its vertical sides $\Re\phi_t$ is within $O(1/X)$ of an integral multiple of $\pi$. On its horizontal sides $|\Im\phi_t|\gg_d\log X$. In each case
\[
   |M_t+M_t^*|\ge\tfrac12(|M_t|+|M_t^*|)
\]
for sufficiently large $X$. Indeed the ratio on the left to the sum on the right before multiplying by $1/2$ is $|\cos\phi_t|/\cosh(\Im\phi_t)$. By \eqref{hg:one}, Rouché's theorem gives $H_t$ the same number of zeros as $\cos\phi_t$ in the rectangle. The latter has one simple real zero, by strict growth on the real line and the preceding sign of $\Im\phi_t$. The single zero of the real entire $H_t$ must itself be real and simple. In particular $H_t$ is nonzero as a function. The rectangles have width $O_d(1/\log X)$. Taking the lower bound on $t\log X$ in the theorem twice as large as in \eqref{hg:one} puts both sides of each relevant rectangle in the admissible range. The confinement \eqref{hg:basicstrip} covers all zeros in their height range.

For fixed $t$, the relative error in \eqref{hg:one} tends to zero. The zero in the $j$th rectangle thus satisfies $\phi_t(x)=\pi(j+1/2)+o(1)$ by the cosine expression on the real axis. Subtract the phases of successive zeros and apply the mean value theorem and \eqref{hg:stir}. This proves the spacing assertion.

\subsection{The heat threshold from a critical-line bound}
The first Dirichlet coefficient alone need not approximate the saddle with vanishing relative error when $t\log X$ stays fixed. Keeping the shifted $L$-factor gives a coefficient depending only on absolute degree and the available pointwise bound on the critical line.

\begin{theorem}\label{hg:threshold}
Suppose $0\le\alpha<1$ and, for each $\epsilon>0$,
\[
  |L(1/2+iu,\pi)|\ll_{\pi,\epsilon}
           (2+|u|)^{(d/4)(1-\alpha)+\epsilon}.
\]
The choice $\alpha=0$ is always valid. Fix $\delta>0$ and $t_{\max}>0$. For $0<t\le t_{\max}$ every zero with $|\Re z|=X$ sufficiently large in terms of the fixed data is real and simple as soon as
\[
                       d\,t\log X\ge\frac1{1+\alpha}+\delta .
\]
For every elliptic curve over $\Q$, in particular, $t\log X\ge3/8+\delta$ suffices. For all fixed ranks over $\Q$ a positive $\alpha$ is available.
\end{theorem}

The elliptic assertion uses $\alpha=1/3$ from the Weyl bound for primitive holomorphic cusp forms of arbitrary level of \href{https://arxiv.org/abs/1707.01576}{Booker, Milinovich and Ng}. The general rational result uses \href{https://arxiv.org/abs/2109.15230}{Nelson, Theorem 1.1}. No pointwise subconvexity input over a general number field is needed for the choice $\alpha=0$.

\begin{proof}
The high-zero assertion already proved for large $t\log X$ lets us restrict to
$A_0\le dC\le dC_0$, where $C=t\log X$, $A_0=1/(1+\alpha)+\delta$ and $C_0$ is fixed. Take $|\Im z|\le1/2$, write $\sigma=\Re s=1/2+v_0$ and $\sigma_*=1/2+|v_0|$. The larger of the two gamma models has size comparable to
\[
             m_X=|G(\sigma_*+ix)|\exp\big(t\Re(b(\sigma_*+ix)^2)\big),
             \qquad x=\Re z .
\]
In fact $|M_t(z)|/m_X\asymp_{\pi,C_0} X^{(d/2)(\sigma-\sigma_*)}$, by Stirling and $b'(s)=O(1/X)$. Under the tilted Gaussian the real mean is $H=2t\Re b(s)=dC+O(1/\log X)$ and its variance is $2t$.

We first bound the contribution where $0\le u$ and $0\le\sigma+u\le1$. By Phragmén--Lindelöf between $0,1/2,1$, the functional equation and the assumed critical-line bound yield, for $0\le v\le1$,
\[
 |L(v+ix)|\ll_{\pi,\epsilon} X^{g(v)+\epsilon},\qquad
       g(v)=\tfrac d2(1-v)-\tfrac{d\alpha}{2}\min(v,1-v).
\]
The usual $O(X^\epsilon)$ bound holds just to the right of one. After dividing our half integral by $m_X$, the Gaussian and the gamma ratio show that on the bounded range in question the integrand is at most $O(\sqrt{\log X})$ times $X^{E+\epsilon}$, with
\[
       E=\tfrac d2(\sigma-\sigma_*)+g(\sigma+u)
                           -(u-dC)^2/(4C).
\]
Replacing $H$ by $dC$ costs only a bounded factor here. Use
$\min(\sigma+u,1-\sigma-u)=1/2-|v_0+u|$ and $|v_0+u|\le |v_0|+u$. Expanding the square and completing the remaining quadratic in $u$ gives
\begin{align*}
 E&\le\tfrac d4(1-dC-\alpha)-\tfrac d2(1-\alpha)|v_0|
                   +\tfrac{d\alpha}{2}u-u^2/(4C)\\
  &\le\tfrac d4(1-\alpha)\{1-(1+\alpha)dC\}.
\end{align*}
This is strictly negative with a fixed margin. With $\epsilon$ chosen inside that margin the whole contribution is $o(m_X)$.

Fix $0<b<\min(1/8,(A_0-1/2)/8)$. For the dominant half ($\sigma=\sigma_*$), and for either half when $|v_0|\le2b$, we have $\sigma+H>1+4b$ for sufficiently large $X$. The shifted model
\[
              \mathcal M_t(z)=M_t(z)L(s+2tb(s))
\]
therefore has size comparable to $|M_t(z)|$: both $L$ and $1/L$ are bounded in the relevant half-plane by their absolutely convergent Dirichlet series. In the range $1<\sigma+u\le1+b$ the real Gaussian center is at least $3b$ away. Since $t\le C_0/\log X$, the saving is a fixed negative power of $X$, sufficient against $X^\epsilon$. The contribution with the constant $L(s+2tb(s))$ is negligible on the whole range $\sigma+u\le1+b$, including the omitted negative tail. When $\sigma+u\ge1+b$, absolute convergence also bounds the derivative of $L$ uniformly, giving
\[
 |L(s+u)-L(s+2tb(s))|\ll_{\pi,b}|u-H|+O_{\pi,b}(t).
\]
The expectation of the right side is $O_{\pi,b}(\sqrt t+t)=o(1)$. The previous gamma Taylor estimate up to $U=C_1\log X$ and its far-tail estimate give $|B_t-\mathcal M_t|=o(m_X)$ on these halves. All error terms are uniform for the stated compact range of $C$.

For the smaller half with $|v_0|\ge b$, it suffices to bound $B_t$ directly. Its gamma size is $O(m_X X^{-db})$. On $\sigma+u>1$ the $O(X^\epsilon)$ bound for $L$ in a fixed strip, followed by absolute convergence farther right, gives $o(m_X)$ after Gaussian integration. On $\sigma+u\le1$ we already have the strictly negative exponent $E$. Thus the larger half dominates whenever $b\le |\Im z|\le1/2$ and $H_t$ has no zero there.

On $|\Im z|\le b$, both shifted models apply, with error $o(|\mathcal M_t|+|\mathcal M_t^*|)$. Repeat the phase construction using $\mathcal M_t$. Its logarithm exists on a neighborhood of the rectangles. Its first and second logarithmic derivatives differ from those of $M_t$ by $O_{\pi,b,C_0}(1)$ from the logarithmic derivative of $L$ and its derivative in the half-plane of absolute convergence. The new phase derivative thus has real part $\frac d2\log X+O(1)$ everywhere in the narrow strip, and the second phase derivative is $O(1)$. On the vertical sides with $|\Im z|\le(\log X)^{-1/2}$ the real part of the phase differs from $j\pi$ by $O(1/\log X)$. On the rest of those sides and on the horizontal edges $\Im z=\pm b$, its imaginary part has magnitude $\gg\sqrt{\log X}$. These estimates give the same lower bound for the cosine on the boundary. To allow for the small change in $X$ between neighboring vertical sides, use the estimates first with $\delta/2$ and $C_0+1$. Rouché supplies exactly one real simple zero in each rectangle. Combined with the dominance outside the narrow strip and \eqref{hg:basicstrip}, this proves the theorem.
\end{proof}

\subsection{The global time and the polar case}
Apply the outermost-zero argument of Section~\ref{hs:section} at an interior time of $(0,1/8)$. Its inputs are the real-entire order-one flow, the heat equation, strip contraction, and uniform confinement of nonreal zeros to a compact rectangle on each closed positive time interval. All four have just been proved for $\pi$. There are infinitely many real zeros, giving a strict term in the velocity inequality there. Consequently strip contraction reaches width zero at a time strictly less than $1/8$, and every later time has only simple real zeros by the backward splitting in \eqref{heat:splitting}. Choosing an intermediate upper time $t_\pi<1/8$ proves the stated inclusive assertion.

One can also include the meromorphic rank-one exceptions with the ordinary pole removal. In the convention $L(s)=\zeta_F(s-i\omega)$, the completed function has simple poles at $i\omega$ and $1+i\omega$. Multiply it by
\[
                  P_\omega(s)=(s-i\omega)(s-1-i\omega)
\]
\emph{before} applying the Gaussian in \eqref{hg:heat}. Its restriction to $s=1/2+ix$ is $-1/4-(x-\omega)^2$, so this polynomial preserves the reality normalization. Its values at the removed poles after multiplication by the meromorphic completion are nonzero. Use $P_\omega G$ throughout the saddle estimates. At high ordinates its extra logarithmic derivative and its first two derivatives are $O_\omega(X^{-1})$, $O_\omega(X^{-2})$, and $O_\omega(X^{-3})$. The global translate bound \eqref{hg:horizontal} persists: sufficiently far to the right the two powers of $s$ can be absorbed into a single gamma factor by two applications of $\Gamma(w+1)=w\Gamma(w)$, and the other side follows by reflection. The far positive tail changes by a polynomial factor already covered by the Gaussian estimate. Thus every conclusion of the theorem also holds for the pole-removed completion of a norm translate of a Dedekind zeta function.

\subsection{The lower endpoint with general archimedean shifts}
For clarity we also give the fixed-strip form of Dobner's negative-time argument. His definition of the extended Selberg class requires nonnegative real parts for the shifts in $\Gamma(\omega s+\mu)$. Unitary automorphic parameters allow negative real shifts. The calculation below only uses Stirling's formula far from the poles and applies to these shifts as well. The displaced argument, damped Dirichlet series and almost-periodicity argument are due to \href{https://arxiv.org/pdf/2005.05142}{Dobner, Theorems 4--5 and Lemma 3}.

\begin{theorem}\label{hg:negative}
For each fixed $h>0$, there is $B=B(\pi,h)>0$ such that every interval $[T,T+B]$ at sufficiently large positive height contains the imaginary part of a zero $w$ of the time $-h$ deformation in the $s$ variable satisfying
\[
                      \Re w=hd\log T+O_{\pi,h}(1).
\]
In particular there are $\gg_{\pi,h} T$ such zeros up to height $T$. This holds also for the pole-removed norm twists above. For all the completions considered here the set of real-zero times is a closed interval $[\lambda_\pi,\infty)$ with
\[
                  0\le\lambda_\pi<1/8 .
\]
The critical-line assertion for the undeformed standard function is equivalent to $\lambda_\pi=0$.
\end{theorem}

\begin{proof}
At negative time $-h$ the physical $z$ variable is averaged against a real Gaussian of variance $2h$. Equivalently, in the $s$ variable the deformed completion is
\[
  X_{-h}(s)=\frac1{\sqrt{4\pi h}}\int_\R e^{-v^2/(4h)}\Lambda(s+iv)\,dv .
\]
This is the same heat deformation. For example its common Taylor series in physical coordinates is $\sum_{j\ge0}(-t)^jH_0^{(2j)}(z)/j!$. The order-one bound and Cauchy's estimate on circles of radius $j^{1/(1+\epsilon)}$, with any fixed $0<\epsilon<1$, make this series locally uniformly convergent for all complex $t,z$. Termwise Gaussian integration of the Taylor series of $H_0$ gives both the positive- and the negative-time integrals. Absolute interchange follows from the same derivative bound and the absolute Gaussian moments, whose growth at order $k$ is $O_h((C_h\sqrt k)^k)$. In the polar case $\Lambda$ here means the pole-removed completion and $G$ includes $P_\omega$.

Write $L(s)=\sum_{n\ge1}a(n)n^{-s}$ and put
\[
  J_h(s)=s+2h b(s),\qquad D_h(s)=\sum_{n\ge1}a(n)n^{-s}e^{-h(\log n)^2}.
\]
For $s$ in a fixed bounded range of real parts and $T=\Im s\to+\infty$ we claim
\begin{equation}\label{hg:negapprox}
    \frac{X_{-h}(J_h(s))}{G(s)e^{h b(s)^2}}=D_h(s)+O_{\pi,h,\mathrm{strip}}(T^{-1}).
\end{equation}
Indeed choose a fixed real $a$ sufficiently far right for absolute convergence and free of gamma poles. In the integral defining $X_{-h}(J_h(s))$ shift the vertical contour to $\Re w=a$. The integrand is the entire completion times $\exp((w-J_h(s))^2/(4h))$, with normalizing factor $1/(i\sqrt{4\pi h})$. For each fixed $s$ the horizontal ends tend to zero by Gaussian decay in any fixed vertical strip. The identity
\[
  b(s)(w-s)+\frac{(w-J_h(s))^2}{4h}
       = h b(s)^2+\frac{(w-s)^2}{4h}
\]
transforms the normalized integrand into
\[
   L(w)e^{(w-s)^2/(4h)}
        \frac{G(w)}{G(s)e^{b(s)(w-s)}}.
\]
On $|\Im(w-s)|\le T^{1/4}$ the last quotient is $1+O_{\pi,\mathrm{strip}}((1+|w-s|)^2/T)$ by the gamma Taylor estimate. The integral of this error is $O_{\pi,h,\mathrm{strip}}(T^{-1})$ since $L(w)$ is bounded and $\Re(w-s)$ is bounded. On the remaining part of the line, with $v=\Im(w-s)$, Stirling on the fixed vertical lines and \eqref{hg:stir} bound the quotient by $O(T^C(1+|v|)^C e^{C|v|})$. To see this also for $v$ near $-T$, use the exponential gamma bound $e^{-\pi d|T+v|/4}$ for the numerator, allowing a fixed polynomial, and the corresponding lower bound at $s$; the ratio of exponential factors is at most $e^{\pi d|v|/4}$. The Gaussian makes the whole tail negligible. With the quotient replaced by one, integrate the absolutely convergent Dirichlet series termwise. The elementary vertical Gaussian integral of $n^{-w}$ equals $n^{-s}e^{-h(\log n)^2}$, proving \eqref{hg:negapprox}.

The series $D_h$ is entire of order at most two: for fixed $a>1$ bound its absolute value by $\sum |a(n)|n^{-a}$ times $\sup_{u\ge0}\exp((a-\Re s)u-hu^2)$. It has a zero. In fact a zero-free entire function of order at most two must be the exponential of a polynomial of degree at most two. Its limit $D_h(\sigma)\to1$ on the positive real axis forces that polynomial to be constant: its real part must tend to zero, and the exponential of the imaginary part of a nonconstant real polynomial has no limit there. But $D_h$ is nonconstant by Dirichlet series uniqueness. Some $a(n)$ with $n>1$ is nonzero, since otherwise the supposed entire completion would consist of a gamma product with infinitely many poles, even after the possible quadratic pole removal.

Take a fixed disk centered at a zero of $D_h$ with boundary free of zeros, and let $e>0$ be the minimum of $|D_h|$ on the boundary. The set of vertical shifts $v$ for which $D_h(s+iv)$ differs from $D_h(s)$ by less than $e/3$ throughout the disk has bounded gaps. This is Bohr almost-periodicity: truncate the absolutely convergent series uniformly on the disk's real strip, and use a neighborhood of the identity in the compact closure of the rotation $v\mapsto(n^{-iv})_{n\le N}$. Finitely many translates of this neighborhood cover the closure, giving a return in every sufficiently long interval. At every sufficiently large such shift, \eqref{hg:negapprox} and Rouché produce a zero $s_v$ of $X_{-h}\circ J_h$ in the shifted disk. The normalizing factor is nonzero there. Equation~\eqref{hg:stir} gives
\[
        \Re J_h(s_v)=hd\log v+O_{\pi,h}(1),\qquad
        \Im J_h(s_v)=v+O_{\pi,h}(1).
\]
Bounded gaps for the shifts, enlarged by these bounded displacements, yield the asserted intervals. Taking disjoint intervals with a fixed gap gives the lower count even if a zero is at an endpoint.

Thus every negative time has nonreal zeros in physical coordinates. The semigroup contraction makes the set of nonnegative real-zero times an upper interval. It is closed by local uniform convergence and Hurwitz, since none of the limiting functions is identically zero. Theorem~\ref{hg:main} gives its strict upper endpoint bound, and the identity at time zero gives the equivalence.
\end{proof}

\subsection{Rankin--Selberg convolutions and tensor motives}\label{hg:rs}

Let $\pi,\pi'$ be unitary cuspidal representations of $\mathrm{GL}_m(\mathbb A_F)$ and $\mathrm{GL}_n(\mathbb A_F)$ over the same number field. Write $D=mn[F:\Q]$. In the convention $L(s,\pi\otimes|\det|^{ia})=L(s+ia,\pi)$, the polar case is $m=n$ and $\pi'\simeq\widetilde\pi\otimes|\det|^{-i\omega}$ for some real $\omega$. Here set $P(s)=(s-i\omega)(s-1-i\omega)$, and otherwise set $P=1$. Denote the meromorphic completion prior to removal by $\Lambda_{\rm RS}$, and choose $|c|=1$ so that $cP(s)\Lambda_{\rm RS}(s)$ is real at $s=1/2+ix$ for real $x$.

\begin{corollary}\label{hg:rscor}
Deform this pole-removed completion in physical coordinates by \eqref{hg:heat}. Its real-zero times are $[\lambda_{\pi,\pi'},\infty)$, where $0\le\lambda_{\pi,\pi'}<1/8$. Equality $\lambda_{\pi,\pi'}=0$ is equivalent to the critical-line assertion for the nontrivial zeros of the Rankin--Selberg function. At every fixed $t>0$ all sufficiently high zeros are real and simple with spacing $(2\pi+o(1))/(D\log |x|)$. Uniformly for $0<t\le t_{\max}$ and fixed $\delta>0$, reality and simplicity hold at all sufficiently large $|x|$ when $Dt\log |x|\ge1+\delta$. The improvement to $1/(1+\alpha)+\delta$ holds under the critical-line bound of Theorem~\ref{hg:threshold} with $d=D$. At each negative time $-h$, zeros with real part $hD\log T+O_{\pi,\pi',h}(1)$ in the unitary $s$ variable occur with bounded gaps in their positive ordinates.
\end{corollary}

\begin{proof}
The Rankin--Selberg theory of \href{https://arxiv.org/pdf/2103.05634}{Humphries--Thorner, \S3.2, (3.5)--(3.13)} gives the finite local factors, the product of $D$ real gamma factors after duplication, the functional equation, order one after pole removal, and the fixed-pair convexity bound $|L(1/2+ix,\pi\times\pi')|\ll_{\pi,\pi',\epsilon}(2+|x|)^{D/4+\epsilon}$. Their central-character normalization absorbs continuous norm twists; restoring the shifts gives the polar case just specified. Each genuine pole is simple. On the center line $P(1/2+ix)=-1/4-(x-\omega)^2$ in that case, and it satisfies $P(s)=\overline{P(1-\bar s)}$. The dual pair has conjugate coefficients and gamma parameters, so the required phase exists for every pair. Polynomial bounds on the fixed strips at high ordinates follow from the same convexity estimate and the functional equation. The archimedean shifts have real parts at least $-\vartheta_m-\vartheta_n>-1$; they need not all be nonnegative.

To verify the absolute convergence used in the saddle, write $A_v,B_v$ for the two Satake matrices at a common unramified place of norm $q$. For real $\sigma>1$, Cauchy--Schwarz gives
\[
 \sum_{v,k}\frac{|\operatorname{tr}A_v^k\operatorname{tr}B_v^k|}{kq^{k\sigma}}
 \le
 \left(\sum_{v,k}\frac{|\operatorname{tr}A_v^k|^2}{kq^{k\sigma}}\right)^{1/2}
 \left(\sum_{v,k}\frac{|\operatorname{tr}B_v^k|^2}{kq^{k\sigma}}\right)^{1/2},
\]
where all sums here use common unramified places. The two sums on the right are convergent logarithms of incomplete self Rankin--Selberg products. Indeed their local logarithmic coefficients are nonnegative, while the full self products have nonnegative Dirichlet coefficients and converge absolutely on this half-plane. At the finitely many remaining places the actual pair parameters in (3.7) have modulus at most $q^{\vartheta_m+\vartheta_n}<q$. Their logarithmic series therefore converge there as well, even when the local factors have lower degree. Exponentiation gives absolute Dirichlet convergence of both the mixed product and its reciprocal, including after grouping by rational norm. Its logarithmic derivatives are bounded on closed half-planes strictly to the right of one by shrinking the half-plane slightly. These facts give the zero strip by reflection and all of the right-hand estimates in Theorems~\ref{hg:main} and~\ref{hg:threshold}. At high ordinates the contribution $P'/P$ to $b(s)$ is $O_{\omega}(1/|x|)$, with the differentiated bounds already checked in the polar case. Those proofs now apply with $d=D$.

For the negative-time assertion one must ensure that grouping by rational norm leaves a nonconstant Dirichlet series. At every rational prime $p$ unramified in $F$ and avoiding both conductors, the reciprocal factor over $p$ is
\[
           \prod_{v\mid p}\det\left(I-(A_v\otimes B_v)T^{f_v}\right),
                  \qquad T=p^{-s},\qquad f_v=[k_v:\F_p].
\]
It has constant coefficient one and degree $mn\sum_{v\mid p} f_v=D$ with nonzero leading coefficient, since both unramified matrices are invertible. Its reciprocal Taylor series must thus have a nonzero coefficient of degree between $1$ and $D$, inclusive. That coefficient is the rationally grouped Dirichlet coefficient at the corresponding power of $p$. Gaussian damping multiplies it by a positive number and preserves its nonvanishing. This supplies the nonconstancy needed in Theorem~\ref{hg:negative}; the remaining negative-time estimates use exactly the gamma and absolute-convergence properties just verified.
\end{proof}

The same proof permits any nonempty fixed finite product of such functions with positive integral multiplicities. Use the sum of their absolute degrees and the product of their pole-removal polynomials, with multiplicities, before taking the heat flow. Order at most one, reality after a constant phase, absolute convergence and the convexity bound pass to this product. The logarithmic derivatives of the finitely many polynomial factors give the same lower-order bounds, and the common good-prime determinant has nonzero leading coefficient in the summed degree. This statement concerns the heat flow of the whole product.

For two rational elliptic curves $E,E'$, their unitarily normalized functions are weight-two $\mathrm{GL}_2$ functions. The variable $S$ of the tensor motive $H^1(E)\otimes H^1(E')$ satisfies $S=s+1$. Its arithmetic critical line is $\Re S=3/2$, its degree is four and its archimedean factor is $\Gamma_{\C}(S)\Gamma_{\C}(S-1)$ up to a constant. The polar case occurs exactly when $E,E'$ are $\Q$-isogenous. Indeed their central characters are trivial, so a norm twist relating the pair to a contragredient has $\omega=0$. Equality of the two unitary elliptic representations is equivalent to rational isogeny by the isogeny theorem of Faltings. In that case the pole-removal polynomial in the arithmetic variable is $(S-1)(S-2)$. Thus the preceding corollary applies also to the self-pair at its true tensor-motive center.

More generally, for non-CM rational elliptic curves the tensor
$\operatorname{Sym}^a H^1(E)\otimes\operatorname{Sym}^b H^1(E')$ with $a,b\ge1$ gives a pair of cuspidal representations by \href{https://doi.org/10.1007/s10240-021-00126-4}{Newton--Thorne, Theorem A}. Its absolute degree is $(a+1)(b+1)$ and its arithmetic center is $(a+b+1)/2$, since each good eigenvalue has absolute value $p^{(a+b)/2}$ before conversion to unitary coordinates. For CM elliptic curves the symmetric powers decompose into unitary Hecke characters and their quadratic inductions after normalization: on the CM field the characters are the monomials $\psi^{a-j}\bar\psi^j$, paired by conjugation, with the middle line when $a$ is even. Thus the finite-product statement covers mixed and CM pairs as well. Finally, the tensor of four rational elliptic curves can be grouped into two pairs. \href{https://arxiv.org/pdf/math/0007203}{Ramakrishnan, Theorem M and \S4.2--4.3} identifies the fourfold function with the Rankin--Selberg convolution of the two $\mathrm{GL}_4$ isobaric tensor products, including the local tensor factors. Their cuspidal constituents are unitary in the present normalization: at common good primes all eigenvalues have modulus one, forcing the real norm exponent of each constituent to vanish by the modulus of its determinant. This gives the same heat conclusions in degree sixteen, at the arithmetic center $5/2$.

For two rational elliptic curves the preceding $\mathrm{GL}_4$ tensor transfer also improves the coefficient in the sufficient heat threshold. \href{https://ultronozm.github.io/math/standard.html}{Nelson, Theorem 1.1 and (1.2)} gives every relative subconvex saving $\alpha<2/(3r^5-2r^4-r^2)$ for rank $r\ge2$ over $\Q$, and the fixed rank-one Weyl bound gives $\alpha<1/3$. All constituent ranks in the elliptic pair are at most four and their sum is four. Multiplying their critical-line bounds therefore permits any $\alpha<1/1272$, since $2/(3\cdot4^5-2\cdot4^4-4^2)=1/1272$. In particular, for every $\delta>0$ one may require
\[
                 t\log|x|\ge\frac{318}{1273}+\delta
\]
for the high real simple zeros of every fixed rational elliptic pair, with the starting height depending on the pair and $\delta$. For two fixed level-one $\mathrm{GL}_2/\Q$ cusp forms, \href{https://arxiv.org/pdf/2101.12106}{Blomer--Jana--Nelson, Theorem 2} gives the stronger bound $|L(1/2+ix,\pi\times\pi')|\ll (1+|x|)^{2/3+\epsilon}$. Here $D=4$ and $\alpha=1/3$ give $t\log|x|\ge3/16+\delta$.

\clearpage
\section{A uniform all-height heat bound in the level and weight}\label{hc:chapter}
Use the physical height and time of Section~\ref{heat:main}. For a holomorphic cuspidal newform $f$ of weight $k\ge2$ and level $N$, put
\[
                   c=k/2,\qquad A=\sqrt N/(2\pi),\qquad
            F_0(s)=\omega A^s\Gamma(s)L(f,s).
\]
Here $|\omega|=1$ makes $F_0(c+ix)$ real on the real $x$-axis, by the functional equation. In the $s$ variable the heat integral averages $F_0(s+u)$ against the real Gaussian of variance $2t$. Write $F_t$ for this average and $\lambda_f$ for the least time with every zero on $\Re s=c$. In particular $\lambda_f=0$ is the undeformed ERH.

\begin{theorem}\label{hc:theorem}
Uniformly over all these newforms and all heights,
\[
             0\le\lambda_f<0.123585,\qquad
                 \lambda_f<0.121481\quad(N\ge11).
\]
The second bound includes the curve $11\mathrm a1$. Moreover, with $D=\log(A(c+5))-1$, one has
\[
                       \lambda_f\le 8/D\qquad(D\ge64).
\]
In particular, $\lambda_f\ll1/\log(3+k\sqrt N)$ with an absolute constant. For every fixed $t>0$, all the zeros at time $t$ are on the critical line and simple for every sufficiently large $k\sqrt N$.
\end{theorem}

The mechanism giving decay with the conductor was outlined for Dirichlet characters in an October 2026 \href{https://mathoverflow.net/q/515667}{MathOverflow question}, with the bound of order $\log\log q/\log q$ suggested there. We give a direct estimate for the modulus and phase of a Gaussian average. It proves the numerical bounds by controlling the motion of nonreal zeros and the stronger large-parameter estimate by comparing the two reflected halves of the integral. For the finite-conductor numerical approach to primitive real Dirichlet characters see \href{https://arxiv.org/abs/2610.27343}{Freire, Gonz\'alez-Riquelme, Mu\~noz, Poblete and Salinas}.

\subsection*{Coherent integration to the right of the critical strip}
Fix $T\in\R$ and abbreviate $F(v)=F_0(c+v+iT)$. For $B\ge6$,
\begin{equation}\label{hc:initial}
      |F(v)|\le |F(B)|\quad(-1/2\le v\le B),\qquad
                 |F(v)|\ \text{is nondecreasing on }[B,\infty).
\end{equation}
Indeed the centered order-one product may be paired under $a+i\gamma\longleftrightarrow -a+i\gamma$, where $|a|\le1/2$. The centered genus corrections of each pair have zero real sum. The linear exponential has purely imaginary coefficient, by reflection in the critical line, also after factoring out a possible central zero. For each zero, $|v-a|\le|B-a|$ throughout the indicated interval, because $v\ge-1/2\ge2a-B$. The moduli of the ordinary linear factors satisfy the comparison individually. Zeros on the symmetry line give the same conclusion. Taking limits of the paired product proves \eqref{hc:initial}. The same pairing gives for $r>1/2$, in the heat-deformed function,
\begin{equation}\label{hc:poisson}
 \partial_r\log|F_t(c+r+iT)|=
               \sum_\rho\frac{r-a_\rho}{(r-a_\rho)^2+(T-\gamma_\rho)^2},
        \qquad \rho=c+a_\rho+i\gamma_\rho.
\end{equation}
The sum converges absolutely. Section~\ref{heat:main} gives order at most one, reflection and $|a_\rho|\le1/2$ at these positive times.

The two normalized local roots of $f$ have modulus at most one, including the surviving factors at the bad primes. This is Deligne's bound and local newform theory in the conventions of Section~\ref{heat:main}. On $v\ge B$ the normalized Dirichlet variable has real part at least $6+1/2$. The Euler logarithm and its first two derivatives therefore obey
\[
  \max\{|\log L|,|(\log L)'|\}\le 2\sum_{n\ge2}n^{-6}<0.044,\qquad
  |(\log L)''|\le 2\sum_{n\ge2}n^{-5}<0.094 .
\]
Use $\log n\le\sqrt n$ and bound the decreasing sums by their first term and the integral from two onward. For $a\ge7$, the digamma series gives
\[
 |\Im\psi(a+iT)|\le\frac\pi2+\frac1{2a},\qquad
 |\psi'(a+iT)|\le a^{-2}+a^{-1}.
\]
The first estimate follows by bounding $\sum_{j\ge0}|T|/((a+j)^2+T^2)$ by its first term and the integral, the second from the corresponding series of reciprocal squares. Consequently on $v\ge6$ the phase $\phi(v)=\arg F(v)$ can be chosen $1.7$-Lipschitz, and $(\log|F(v)|)''<1/3$.

For real $r$ and $0<t\le1/8$ let $P_r$ be the mass of the density
\[
       w_r(v)=e^{-(v-r)^2/(4t)}|F(v)|\quad(v\ge B),
\]
and write $\mu$ and $V$ for its mean and variance after division by $P_r$. Then $V\le3/11$. In fact $W=-\log w_r$ satisfies $W''\ge1/(2t)-1/3\ge11/3$. Integrate $(v-\mu)(W'(v)-W'(\mu))\ge(11/3)(v-\mu)^2$ by parts. The integral of $(v-\mu)W'(v)$ against the probability density equals $1+(B-\mu)w_r(B)/P_r\le1$. The Gaussian dominates at infinity. The phase bound and $\cos x\ge1-x^2/2$ give
\begin{equation}\label{hc:coherence}
 \left|\int_B^\infty e^{-(v-r)^2/(4t)}F(v)\,dv\right|
        \ge P_r(1-1.7^2 V/2)>0.60P_r.
\end{equation}
Here the complex integral was rotated by $-\phi(\mu)$.

\subsection*{A lower bound from every height}
Suppose $R\ge B+3/2$. Reflect the negative half of the full integral for $F_t(c+R+iT)$ by the functional equation. Its modulus density at $v\ge0$ is $e^{-(v+R)^2/(4t)}|F(v)|$. For $v\ge B$ this is at most $e^{-RB/t}w_R(v)$. The two intervals $0\le v\le B$ together contribute at most
\[
  2|F(B)|\sqrt{\pi t}\,e^{-(R-B)^2/(4t)}
                  \le 2P_R e^{-(R-B)^2/(4t)} .
\]
Use \eqref{hc:initial}, the elementary Gaussian tail bound, and integrate $w_R$ from $R$ to infinity for the last inequality. Since $2e^{-9/2}+e^{-48(15/2)}<0.023$, the modulus of $\sqrt{4\pi t}\,F_t(c+R+iT)$ lies between $(0.60-0.023)P_R$ and $(1+0.023)P_R$. This applies also at $R+h$ for an integer $h\ge1$.

Shift the integration variable by $h$ in the part $v\ge B+h$ of $P_{R+h}$. Gamma recurrence gives $A^h\prod_{j=0}^{h-1}|c+v+j+iT|$ for the ratio of gamma and conductor factors. The two Euler logarithm bounds give the remaining modulus ratio at least $e^{-0.088}$. Together the preceding constants give
\begin{equation}\label{hc:ratio}
   \frac{|F_t(c+R+h+iT)|}{|F_t(c+R+iT)|}
          \ge \frac{A^h}{2}\prod_{j=0}^{h-1}|c+B+j+iT|.
\end{equation}

Let $y>0$ be the largest imaginary part of a zero in physical $z$-coordinates and $T$ the real part of one such zero. Equivalently the corresponding $s$-zero lies at $c-y+iT$. Remove it and its reflected zero from the sum in \eqref{hc:poisson}, at a time when these zeros are simple. Their combined integral over $R\le r\le R+h$ is at most $2\log(1+h/(R-1/2))$. For any remaining zero, the integrated contribution is at most $h(R+h/2+1/2)/(1+(T-\gamma_\rho)^2)$, since $(r-a_\rho)^2\ge1$. Thus
\[
 \sum_{\rho}^{*}\frac1{1+(T-\gamma_\rho)^2}\ge K,\qquad
 K=\frac{\log(\tfrac12 A^h\prod_{j=0}^{h-1}(c+B+j))
                   -2\log(1+h/(R-1/2))}{h(R+h/2+1/2)} .
\]
The star removes the indicated pair. Pair the remaining physical zeros $x_j\pm i y_j$. Their terms in the outermost velocity formula of Section~\ref{hs:section} have numerators $y\mp y_j\ge0$. Each squared denominator is at most $(T-x_j)^2+1$, since $y+|y_j|\le1$. The numerators of a pair sum to $2y$, while a real zero contributes numerator $y$. For any $0<\kappa\le K$ we obtain
\[
                 y'\le -1/y-2\kappa y,\qquad
                        (y^2)'\le-2-4\kappa y^2.
\]
At any compact subinterval of positive times all nonreal zeros are confined to a fixed compact set by Theorem~\ref{hs:large}. Locally enclose them in finitely many contours avoiding zeros on the boundaries. The Weierstrass polynomials inside those contours have analytic discriminants; these are not identically zero because forward Hermite splitting separates the roots at every collision as in Section~\ref{hs:section}. Multiple-root times are locally isolated. Between them the maximum of finitely many squared heights satisfies the differential inequality at almost every time and is locally Lipschitz. Integrate on those intervals and pass continuously across collision times; the maximum is continuous there by the local root theorem. Once the maximum reaches zero the strip theorem keeps it there. Starting at a positive time tending down to zero, with initial squared height at most $1/4$, proves
\[
             \lambda_f\le \frac{\log(1+\kappa/2)}{4\kappa}.
\]
For all weights and levels take $(B,R,h)=(16,35/2,15)$, lower-bound $c$ by one and $A$ by $1/(2\pi)$, and obtain $K>0.046$. For $N\ge11$, $(6,15/2,9)$ gives $K>0.117$. Substitution proves the decimal bounds in Theorem~\ref{hc:theorem}. The short replay \texttt{checks/heat-conductor/constants.py} evaluates only these elementary real numbers.

\subsection*{Every height after time of order the inverse log conductor}
Here is the estimate removing the iterated logarithm that arises from integrating the outer-zero inequality. It separates the two reflected half integrals at every height. For $-1/2\le r\le1/2$ write
\[
             J(r)=\int_r^\infty e^{-(v-r)^2/(4t)}F(v)\,dv .
\]
Up to the positive Gaussian normalization,
\[
                 F_t(c+r+iT)=J(r)+\overline{J(-r)}.
\]
Indeed conjugation and $v\mapsto-v$ turn the latter integral into the part below $r$. The choice of phase in $F_0$ removes the root number here.

Take $B=6$ and suppose $D\ge64$ and $t=8/D$. On $v\ge B$,
\begin{equation}\label{hc:slope}
               (\log|F(v)|)'\ge D .
\end{equation}
For $a>1$ the digamma series gives $\Re\psi(a+iT)\ge\psi(a)\ge\log(a-1)$. For the second inequality one may use
\[
 \log z-\psi(z)=\int_0^\infty e^{-zu}
             \left(\frac1{1-e^{-u}}-\frac1u\right)du,\qquad \Re z>0 .
\]
The integrand in parentheses lies between zero and one, yielding $\psi(a)\ge\log a-1/a\ge\log(a-1)$. This identity follows by subtracting the elementary Laplace formulas for $\log z$ and $\psi(z)$. Together with the Euler derivative bound it proves \eqref{hc:slope}.

Split $J(r)=I+K$ at $B$, with $I$ the high part. Write $P=P_r$. Integrating the lower bound $|F(B)|e^{D(v-B)}$ for $v\ge r+2tD\ge B$ gives
\[
 P\ge |F(B)|\sqrt{\pi t}\exp\{D(tD-B+r)\}.
\]
Since $tD=8$, the exponent is at least $3D/2$. Using \eqref{hc:initial} on $[r,B]$, we obtain
\[
  |K|/P\le (B+1/2)e^{-3D/2}/\sqrt{\pi t}<10^{-3},\qquad
  2t|F(r)|/P<10^{-3}.
\]
These inequalities hold for every $D\ge64$ by $t=8/D$. In particular $|J(r)|>0.599P$ by \eqref{hc:coherence}. Differentiate the integral at its moving lower endpoint. With $\mu,V$ from its high-part absolute density,
\begin{align*}
 2t\,J'(r)={}&-2tF(r)+(\mu-r)I\\
 &+\int_B^\infty(v-\mu)e^{-(v-r)^2/(4t)}F(v)\,dv
      +\int_r^B(v-r)e^{-(v-r)^2/(4t)}F(v)\,dv .
\end{align*}
For the last integral bound the absolute integral on $[r,B]$ just as for $K$, before any cancellation, and multiply by $B+1/2$. The preceding centered integral has modulus at most $P\sqrt V$ by Cauchy--Schwarz. Finally $\Re(I/J)\ge1-10^{-3}/0.599$. Thus
\[
 2t\Re(J'/J)\ge
  (\mu-r)(1-10^{-3}/0.599)
         -\frac{10^{-3}+\sqrt{3/11}+(B+1/2)10^{-3}}{0.599}>0 ,
\]
since $\mu-r\ge B-1/2$. The function $|J(r)|$ is strictly increasing on the whole interval. At $r\ne0$ the summands $J(r)$ and $\overline{J(-r)}$ consequently have unequal moduli. No off-line zero exists in the critical strip; outside it the strip theorem applies. This proves the asserted $8/D$. The later simplicity and lower bound $\lambda_f\ge0$ follow from Section~\ref{heat:main}.

\begin{theorem}[All cuspidal standard functions at fixed degree]\label{hc:automorphic}
Let $\pi$ have absolute degree $d$ and an entire cuspidal standard completion in the conventions of Section~\ref{hg:section}. Set
\[
 \mathcal Q_* = Q\inf_{u\in\R}\prod_{j=1}^d(3+|\mu_j+iu|),\qquad
 B=8(d+1),\qquad D_*=\tfrac12\log\mathcal Q_*-3(d+1).
\]
Here $Q$ is the arithmetic conductor including the discriminant factor, and $\mu_j$ are the unitary archimedean shifts. Whenever $D_*\ge64(d+1)^2(B+2)$, the least physical all-real time satisfies
\[
                         \lambda_\pi\le(B+2)/D_* .
\]
In particular $\lambda_\pi\ll_d 1/\log(3+\mathcal Q_*)$. The minimum over norm translates makes the bound invariant under shifts of all ordinates. In degree one, for every primitive nonprincipal Dirichlet character, this gives a physical upper bound $O(1/\log(3+q))$. In the usual $\xi((1+ix)/2)$ convention multiply the physical upper bound by four.
\end{theorem}

\begin{proof}
Let $c=1/2$, and take $F(v)$ as the phased completion at $c+v+iT$. The paired-product comparison \eqref{hc:initial} and the reflection formula for $J$ use only the initial strip, reality on the critical line and order one, all available in Section~\ref{hg:section}. We give the uniform right-half estimates. The elementary local unitary bounds $|\alpha_p|\le p^{1/2}$ and $\Re\mu_j\ge-1/2$ suffice; see \href{https://arxiv.org/pdf/2103.05634}{Humphries--Thorner, \S3} for the stronger standard bounds. If the base field differs from $\Q$, split a residue-degree-$f$ linear factor as $f$ rational-prime factors with roots of modulus at most $p^{1/2}$ after passing from $p^f$ to $p$. There are at most $d$ rational roots at any prime.

On $v\ge B$ the Euler logarithm and each of its first two derivatives have modulus at most $d\sum_{n\ge2} n^{1-v}<1/4$. For an archimedean factor put $z=(c+v+\mu_j+iT)/2$, so $\Re z\ge v/2$. Its contribution to the phase derivative has magnitude at most $\tfrac12(\pi/2+1/(2\Re z))$. The absolute value of its log second derivative is at most $\tfrac14((\Re z)^{-2}+(\Re z)^{-1})$. Thus the phase of $F(v)$ is $2d$-Lipschitz, and $(\log|F(v)|)''\le1$. Also $3+|\mu_j+iT|\le3|c+v+\mu_j+iT|$, by the triangle inequality and $v\ge16$, $\Re\mu_j\ge-1/2$. The preceding digamma integral shows $\Re\psi(z)\ge\log|z|-1/\Re z$. Summing
\[
 \Re\frac{d}{dv}\log\Gamma_{\R}(c+v+\mu_j+iT)
       =\tfrac12\Re\psi(z)-\tfrac12\log\pi
\]
over $j$, and inserting the conductor and Euler derivative, proves $(\log|F(v)|)'\ge D_*$ for every $T$.

Take $t=(B+2)/D_*\le1/(64(d+1)^2)$. The integration-by-parts variance bound now gives $V\le4t$. The phase estimate yields
\[
                  |I|\ge(1-8d^2t)P\ge(7/8)P .
\]
The same completed square gives $|F(B)|/P\le e^{-3D_*/2}/\sqrt{\pi t}$ uniformly in $|r|\le1/2$. In particular the absolute low-part integral divided by $P$ is at most
\[
             (B+1/2)e^{-3D_*/2}\sqrt{D_*/(\pi(B+2))}
                       < 10^{-3}/(B+1),
\]
and $2t|F(r)|/P<10^{-3}$. These bounds follow already from $D_*\ge B+2$, $B\ge16$, by monotonicity of $e^{-3x/2}\sqrt x$. In the displayed inequality for $2t\Re(J'/J)$ one can therefore use $|J|/P>0.8$, $\sqrt V\le1$, $\mu-r\ge B-1/2$, and a total numerator for the three remainder terms at most $1.002$. The result is strictly positive. Reflection and strip confinement give the theorem. When the large-parameter hypothesis fails the unconditional $\lambda_\pi\le1/8$ from Section~\ref{hg:section} supplies the claimed order bound with a constant depending only on $d$. For Dirichlet characters the archimedean shift is zero or one, so $\mathcal Q_*$ is $3q$ or $4q$.
\end{proof}

\subsection*{The leading constant in the hybrid conductor bound}
At small time the right-hand integral in \eqref{hc:coherence} keeps an arbitrarily large proportion of its absolute mass. Splitting the two halves at the fixed critical line, instead of the moving center of each Gaussian, improves the leading coefficient.

\begin{theorem}\label{hc:leading}
Fix $d$, $0\le\theta\le1/2$ and $\epsilon>0$. Suppose that the normalized rational-prime roots have modulus at most $p^\theta$ and the archimedean shifts have real part at least $-1/2$. Write
\[
                \mathcal Q(T)=Q\prod_{j=1}^d(3+|\mu_j+iT|).
\]
There is $t_0=t_0(d,\theta,\epsilon)>0$ such that at every $0<t\le t_0$, every zero of $F_t$ at height $T$ lies on the critical line whenever
\[
                      t\log\mathcal Q(T)\ge 1+2\theta+\epsilon .
\]
Consequently along every such family with $\mathcal Q_*\to\infty$,
\[
            \limsup\lambda_\pi\log\mathcal Q_*\le1+2\theta .
\]
For holomorphic newforms this gives $\limsup\lambda_f\log(k\sqrt N)\le1/2$ as $k\sqrt N\to\infty$. For primitive nonprincipal Dirichlet characters the upper limit of the physical $\lambda_\chi\log q$ is at most one.
\end{theorem}
\begin{proof}
We can suppose $0<\epsilon\le1$. Put $B=1/2+\theta+\epsilon/8$, and write $F(v)$ as before. On $v\ge B$ the Euler product and its first two logarithmic derivatives are bounded in modulus by constants depending only on $d,\theta,\epsilon$, since $c+v-\theta\ge1+\epsilon/8$. Here and below the bound on the Euler product includes its reciprocal. The elementary digamma estimates used above, now with $\Re z\ge B/2$, give a uniform Lipschitz bound for the phase and a uniform upper bound for $(\log|F|)''$ on this right ray. They also give
\[
        (\log|F(v)|)'\ge D_T:=\tfrac12\log\mathcal Q(T)-C_{d,\theta,\epsilon}.
\]
Indeed $3+|\mu_j+iT|\le(2+3.5/B)|c+v+\mu_j+iT|$ by the triangle inequality and $\Re\mu_j\ge-1/2$.

We need a uniform upper bound on $0\le v\le B$, where the earlier monotonicity comparison need not apply. In the strip $|\Re(s-c)|\le B$ multiply the completed function by $\exp((s-c-iT)^2)$. On the right edge at ordinate $T+w$ its modulus is at most $C|F(B)|\exp(B^2+C|w|-w^2)$. The Euler factors have bounded modulus and reciprocal on that edge, and the vertical logarithmic derivative of the modulus of the gamma-conductor factor is bounded there by the uniform phase-derivative estimate for the gamma factors. The same modulus bound holds on the left edge by reflection. The maximum principle on expanding rectangles, using order at most one against the negative square on the horizontal sides, gives
\[
                    |F(v)|\le C'|F(B)|\qquad(0\le v\le B)
\]
uniformly in the form and in $T$.

With the splitting point fixed at zero, put
\[
 A_r=\int_0^\infty e^{-v^2/(4t)+vr/(2t)}F(v)\,dv,\qquad
 \sqrt{4\pi t}\,F_t(c+r+iT)=e^{-r^2/(4t)}(A_r+\overline{A_{-r}}).
\]
Take $r_0=\min(1/4,\epsilon/16)$. For sufficiently small $t$ the hypothesis gives $tD_T\ge B+\epsilon/4$, uniformly in $T$. For $-r_0\le r\le1/2$ the absolute mass $P$ on $v\ge B$ in $A_r$ dominates its absolute mass on $[0,B]$ by an arbitrarily large factor as $t_0\downarrow0$, uniformly in the parameters. To check this, first cancel the common factor $e^{r^2/(4t)}$ so that the Gaussian is centered at $r$. The high mass is at least $|F(B)|\sqrt{\pi t}\exp(D_T(tD_T-B+r))$ by completion of the square, since $r+2tD_T\ge B$. The low mass is at most $BC'|F(B)|$. Their ratio tends to zero for low over high since $D_T\ge(B+\epsilon/4)/t$ and $tD_T-B+r\ge3\epsilon/16$.

For the high part the integration-by-parts estimate now gives variance $V\le4t$ once $t_0$ is small; its phase is uniformly Lipschitz. Thus its complex modulus is at least $(1-O(t))P$. In particular $|A_r|\ge(1-o(1))P$ uniformly on the indicated range, where the error is taken as $t_0\downarrow0$. For $|r|\le r_0$, differentiate $A_r$ at its fixed endpoint: $2tA_r'$ is the same integral with an extra factor $v$. If $\mu\ge B$ is the mean of the high absolute mass, its centered part has modulus at most $P\sqrt V$, and the low contribution has modulus at most $B$ times the low absolute mass. The ratio of the high complex integral to $A_r$ has real part at least $1-o(1)$. Consequently
\[
                   2t\Re(A_r'/A_r)\ge \mu(1-o(1))-O(\sqrt t)-o(1)>0 .
\]
This proves strict comparison $|A_r|>|A_{-r}|$ for $0<r\le r_0$. For $r_0\le r\le1/2$ the high absolute mass of $A_{-r}$ is at most $e^{-Br/t}P$, and its low absolute mass is at most the negligible low mass of $A_r$. Hence the same strict comparison holds there. Reflection and strip confinement finish the pointwise assertion. Take $t=(1+2\theta+\epsilon)/\log\mathcal Q_*$ and then let $\epsilon$ decrease for the second assertion. The two real shifts of a holomorphic newform are $(k-1)/2,(k+1)/2$ in unitary normalization, so $\log\mathcal Q_*=2\log(k\sqrt N)+O(1)$, and Deligne permits $\theta=0$. The analogous degree-one calculation gives the last assertion.
\end{proof}

\clearpage
\section{Subconvexity and all-height heat bounds for elliptic twists}\label{hsub:chapter}
Use the physical time of Section~\ref{hc:chapter} and the unitary coordinate $c=1/2$. In this section a weight-$k$ holomorphic newform is put in that coordinate by translating its arithmetic variable by $(k-1)/2$. The physical time and heights are invariant under this translation. Write $\mathcal Q(T)$ for the analytic conductor in Theorem~\ref{hc:leading}, $\mathcal Q_*=\inf_T\mathcal Q(T)$, and $U=\log\mathcal Q(T)$. All the gamma shifts in the present applications have nonnegative real part, and the finite Euler roots are tempered.

\begin{theorem}\label{hsub:global}
There is an absolute $\delta>0$ such that for holomorphic newforms varying in both weight and level,
\[
  \limsup_{k\sqrt N\to\infty}\lambda_f\log(k\sqrt N)
                         \le\frac1{2(1+4\delta)}.
\]
One may use any positive saving $\delta<1/4$ in the uniform central subconvexity bound $L(1/2+iT,\pi)\ll_\nu\mathcal Q(T)^{1/4-\delta+\nu}$ for all $\nu>0$. Such a saving for $\mathrm{GL}_2/\Q$ is the theorem of P.~Michel and A.~Venkatesh, \href{https://doi.org/10.1007/s10240-010-0025-8}{\emph{The subconvexity problem for $\mathrm{GL}_2$}}, Theorem~1.1, including unitary norm twists.
\end{theorem}

\begin{theorem}\label{hsub:elliptic}
Fix a rational elliptic curve $E$. For every $\epsilon>0$ there is $b>0$ such that, with $O_{E,\epsilon}(D^{1-b})$ exceptions among squarefree $a$ with $D\le |a|\le2D$, one has
\[
              \lambda_{E^{(a)}}\log |a|\le\tfrac14+\epsilon .
\]
All zeros are real in physical height coordinates at the indicated time and all larger times, and simple at each strictly larger time. The assertion holds for a proportion tending to one of prime parameters in any fixed reduced progression. A fixed finite collection of original rational elliptic curves can be treated simultaneously.
\end{theorem}

\subsection*{Comparing the two Gaussian exponents}
We first give the common estimate. Fix $B>1/2$ close to $1/2$. At each height write $F(v)$ for the phased completion at $c+v+iT$, and $A_r$ for its positive-half integral in the proof of Theorem~\ref{hc:leading}, so the two terms at real displacement $r$ are $A_r$ and $\overline{A_{-r}}$ up to the common real normalization. Suppose, with uniform constants, that
\begin{equation}\label{hsub:strip}
   |F(v)|\le C|F(B)|e^{-aU(1-v/B)}\quad(0\le v\le B),
       \qquad 0\le a<B/2 .
\end{equation}
Write $s=tU$ for the product of time and logarithmic conductor in this proof; it is a real number. If $s$ ranges in a bounded interval, a sufficient condition for every zero at this height to be on the critical line, as $U\to\infty$, is
\begin{equation}\label{hsub:threshold}
              s>\frac{2B^2}{B+2a}
\end{equation}
with a fixed positive margin. This is uniform for $a$ in a fixed compact subset of $[0,B/2)$.

Indeed put $D_T=U/2-C_B$, where the slope and phase estimates on $v\ge B$ in Theorem~\ref{hc:leading} apply. The high absolute mass $P$ in $A_r$ is at least
\[
 |F(B)|\sqrt{\pi t}\exp\left(t(D_T+r/(2t))^2-D_TB\right)
\]
provided its completed-square center $2tD_T+r$ is at least $B$. The low absolute mass $P_0$ satisfies by \eqref{hsub:strip}
\[
  P_0\le C B |F(B)|
       \exp\left(-aU+t(aU/B+r/(2t))^2\right).
\]
The second inequality maximizes a real quadratic over the full line, which is an upper bound for its maximum on $[0,B]$. At $r=0$ the difference of the high and low exponents, divided by $U$, is
\[
      s\left(\tfrac14-a^2/B^2\right)-B/2+a+O_B(U^{-1}) .
\]
Its sign changes at exactly \eqref{hsub:threshold}. For general $r$ the additional difference divided by $U$ is $r(1/2-a/B)+O_B(U^{-1})$. The threshold exceeds $B$ on the stated compact parameter sets. Hence a fixed $r_0>0$ suffices simultaneously for $-r_0\le r\le 1/2$, with the center above $B$ and $P_0/P\ll e^{-b_0U}$ for some $b_0>0$. The phase Lipschitz bound and variance $O(t)$ show that the high complex integral has modulus at least $(1-O(t))P$. For $|r|\le r_0$ differentiate $A_r$ at its fixed endpoint zero. If $\mu\ge B$ is the mean of the high mass, the numerator $2tA_r'$ is its high complex integral times $\mu$, plus terms of modulus at most $P O(\sqrt t)+BP_0$. Its ratio with $A_r$ has real part at least $\mu(1-o(1))-O(\sqrt t)>0$. For $r\ge r_0$ the high absolute mass at $-r$ is at most $e^{-Br/t}P$, and its low absolute mass is no larger than $P_0$. In both cases the two reflected terms have different moduli for $r>0$. Reflection and strip confinement exclude all off-line zeros. This proves the common estimate.

\begin{proof}[Proof of Theorem~\ref{hsub:global}]
We verify \eqref{hsub:strip}. Let $\gamma=1/4-\delta+\nu$. At both of the real gamma factors in a holomorphic newform the unitary shifts are nonnegative. Stirling on $\Re z\ge1/4$ gives, uniformly in the weight, level and ordinate,
\[
 \left|\frac{G(c+iy)}{G(c+B+iy)}\right|\ll_B\mathcal Q(y)^{-B/2}
\]
for the gamma-conductor part $G$. This follows also on bounded arguments by compactness. On $\Re s=c+B$ the Euler product and its reciprocal are bounded uniformly for fixed $B$, and the gamma modulus at ordinates $T+w$ and $T$ differs by at most a factor $e^{C_B|w|}$ using the uniform imaginary-digamma bound on that line. Since $\mathcal Q(T+w)/\mathcal Q(T)$ and its reciprocal are at most $(1+|w|)^d$, the subconvex bound gives
\[
  |F_0(c+i(T+w))|\le C|F(B)|\mathcal Q(T)^{-a}
                      e^{C_B|w|}(1+|w|)^{C_B},\qquad
                       a=B/2-\gamma .
\]
Here $F_0$ denotes the phased global completion. Multiply it in $c\le\Re z\le c+B$ by $\exp((z-c-iT)^2-(aU/B)(z-c))$. The two vertical edges are bounded by $C'|F(B)|e^{-aU}$, and the horizontal ends vanish by order at most one. The maximum principle gives \eqref{hsub:strip}. Take $B\downarrow1/2$ and $\nu\downarrow0$ after fixing a positive margin. The threshold tends to $1/(1+4\delta)$. Products $s\ge 1+\epsilon$ for arbitrarily small times are already covered by Theorem~\ref{hc:leading}, so only bounded $s$ enter the comparison. At $t=(1/(1+4\delta)+\epsilon)/\log\mathcal Q_*$ every height is covered for sufficiently large conductor. Finally $\log\mathcal Q_*=2\log(k\sqrt N)+O(1)$, proving the result.
\end{proof}

The same argument for primitive nonprincipal Dirichlet characters with I.~Petrow and M.~Young, \href{https://arxiv.org/abs/1908.10346}{\emph{The fourth moment of Dirichlet $L$-functions along a coset and the Weyl bound}}, Theorem~1.1, gives the physical upper limit $\limsup\lambda_\chi\log q\le3/4$ by using $\delta=1/12$. This particular argument uses no prior Dirichlet critical-line theorem. In degree one the prior Dirichlet result treated earlier in this manuscript already gives the exact value $\lambda_\chi=0$; the strict elliptic improvements above concern a different degree.

\begin{proof}[Proof of Theorem~\ref{hsub:elliptic}]
Write $Q_a(T)$ for the analytic conductor of the weight-two newform of $E^{(a)}$. Uniformly on $D\le|a|\le2D$,
\[
         \log Q_a(T)=2\log D+2\log(3+|T|)+O_E(1).
\]
In fact after extracting the bounded twisting factor supported on $2N_E$ as in Theorem~\ref{ang:quadratic}, the remaining fundamental character is coprime to the level of one of finitely many fixed base forms. Its conductor is comparable to $D$, and the conductor of the twist is the base level times the square of this conductor.

For any small fixed $\nu>0$, outside $O_{E,\nu}(D^{1-b_\nu})$ twisting parameters we claim that
\begin{equation}\label{hsub:sup}
   |L(1/2+iy,E^{(a)})|\ll_{E,\nu}D^\nu(1+|y|)
                         \qquad (|y|\le 2D^2)
\end{equation}
in unitary coordinates. The quadratic large-sieve and central approximate-functional-equation estimate in the proof of Theorem~\ref{ang:quadratic} gives
\[
 \sum_a|L(1/2+iy,E^{(a)})|^2
          \ll_{E,\nu'}D(2+|y|)[D(2+|y|)]^{\nu'}
\]
for every $\nu'>0$. The same bound holds for $|L'|^2$, on enlarging $\nu'$ slightly. To check it, take a circle of radius $1/\log(3D(2+|y|))$ around $1/2+iy$. On that circle the two sums in the same approximate functional equation have the same lengths $O_E(D(2+|y|))$ and gamma-ratio modulus $O_E(1)$. Moving the real part by the stated radius costs a bounded factor on the length scale, and on the tail shells it is absorbed by rapid decay. The maximal-prefix large sieve thus gives the same square moment on the circle; Cauchy's integral formula costs $O(\log^2(3D(2+|y|)))$ for the squared derivative. On every unit interval the supremum of the square is at most twice the integral of $|L|^2+|L'|^2$. Integrating through $|y|\le H+1$ and summing intervals consequently gives
\[
  \sum_a \sup_{|y|\le H}|L(1/2+iy,E^{(a)})|^2
                    \ll_{E,\nu'} D^{1+4\nu'}(1+H)^2
             \qquad(1\le H\le4D^2),
\]
after taking the initial moment epsilon smaller if necessary. Markov's inequality at threshold $D^\nu(1+H)$, and a union over dyadic $H$, prove \eqref{hsub:sup} for $\nu'<\nu/4$.

Put $U_0=2\log D$, $h_T=\log(3+|T|)/\log D$, and take heat time $t=(1/2+\epsilon_0)/U_0$, with $\epsilon_0>0$ fixed. Then $U=\log Q_a(T)=(1+h_T)U_0+O_E(1)$. Theorem~\ref{hc:leading} covers all $h_T\ge1-\epsilon_1$ if $\epsilon_1>0$ is small enough in terms of $\epsilon_0$. For the other heights, \eqref{hsub:sup} supplies at the central edge in the maximum-principle argument the exponent saving
\[
              a_0(h)=\frac{1-h}{4(1+h)},\qquad 0\le h\le1-\epsilon_1.
\]
More precisely, one can use $a=a_0(h_T)-2\nu$ for all sufficiently large $D$. Indeed the upper bound $D^\nu(1+|T+w|)$ divided by $Q_a(T+w)^{B/2}$ is at most
$C Q_a(T)^{-a} (1+|w|)^C$ when $|T+w|\le2D^2$ and $B\ge1/2$. The gamma comparison contributes $e^{C_B|w|}$. Outside this ordinate range the usual convexity bound replaces \eqref{hsub:sup}; since the relevant $|T|\le D$ there, $e^{-w^2}$ against the fixed polynomial conductor bounds absorbs this tail uniformly. Thus the identical maximum-principle argument proves \eqref{hsub:strip}. Choose $\nu$ small enough that $a>0$ throughout this compact height range.

At $B=1/2$ the limiting threshold from \eqref{hsub:threshold}, divided by $1+h$, equals
\[
              \frac{1}{(1+h)(1+4a_0(h))}=\frac12
                    \qquad(0\le h\le1-\epsilon_1).
\]
This expression is continuous also at $h=0$. First take $\nu>0$ small and $B>1/2$ sufficiently close to $1/2$ that the threshold divided by $1+h$ is less than $1/2+\epsilon_0/2$ throughout the compact range. With these choices $a$ stays a fixed positive distance below $B/2$, as required for the Gaussian exponent comparison. The chosen time gives $s=tU>(1+h)(1/2+\epsilon_0/2)$ for large $D$. Every height is consequently free of off-line zeros at this time, and heat monotonicity and strict forward splitting give the assertions about later times. Since $\log|a|=\log D+O(1)$, reducing $\epsilon_0$ proves the displayed bound of the theorem. The prime number theorem in each fixed reduced progression and a finite union of exceptional sets prove the final two assertions.
\end{proof}

\subsection*{Using the zero-density strip inside the Euler-product boundary}
The scale $1/4$ in Theorem~\ref{hsub:elliptic} can itself be improved. The preceding explicit coefficient is retained because the improvement depends on the saving in hybrid subconvexity.

\begin{theorem}\label{hsub:strict}
For every rational elliptic curve $E$ there exist $\kappa,b>0$ such that, outside $O_E(D^{1-b})$ squarefree parameters $D\le |a|\le2D$,
\[
                      \lambda_{E^{(a)}}\log |a|\le\tfrac14-\kappa .
\]
The same conclusion holds for a proportion tending to one in each fixed reduced progression of prime twisting parameters.
\end{theorem}
\begin{proof}
Write $w_0=(1+4\delta)^{-1}<1$ for a uniform newform constant from Theorem~\ref{hsub:global}. In its proof this is also the pointwise threshold at every height: $t\log Q_a(T)\ge w_0+\epsilon$ suffices for all small $t$, with fixed $\epsilon>0$. Choose
\[
            \max(0,2w_0-1)<\tau_0<\tau_1<1 .
\]
Take $0<u<1/2$ so close to $1/2$ that $2u/(3-4u)>\tau_1$. Theorem~\ref{ang:quadratic}, with a height exponent strictly between these two numbers, shows that outside a power-saving set of twists the finite $L$-function has no zeros in $\sigma\ge1/2+u$ through $2D^{\tau_1}$. Include the exceptional set of \eqref{hsub:sup}. All ensuing estimates are uniform in the remaining twists.

We first control the logarithm just to the left of $\Re s=1$. Fix $\sigma_L>1/2+u$, $\sigma_R>1$ so close to $1/2+u$ and $1$ that $(\sigma_L+\sigma_R)/2<1$. For $|y|\le D^{\tau_1}$ and $\sigma_L\le\sigma\le\sigma_R$, the standard local logarithmic-derivative formula gives $|L'/L(\sigma+iy)|\ll\log D$: within one unit the $O(\log D)$ zeros are a fixed horizontal distance from this segment, and subtracting the logarithmic derivative at $2+iy$ bounds the remaining terms by $O(\log D)$ as in \eqref{amp:localzeros}. On the right edge take the absolutely convergent Euler logarithm. Horizontal integration fixes a holomorphic branch of $\log L$ on the rectangle with $|\log L|\ll\log D$ throughout and $|\log L|\ll1$ on its right edge.

For $|y|\le 2D^{\tau_0}$ interpolate the modulus of this analytic logarithm between the two edges. One may multiply it by $\exp((z-\sigma_R-iy)^2)$ and by the exponential of a linear function interpolating the logarithms of the two edge bounds. The horizontal sides at heights $\pm D^{\tau_1}$ are negligible by the just-proved logarithmic bound. The maximum principle gives
\[
   |\log L(\sigma+iy)|
      \ll (\log D)^{(\sigma_R-\sigma)/(\sigma_R-\sigma_L)} .
\]
Choose any $B<1/2$ sufficiently close to $1/2$ that $c+B>(\sigma_L+\sigma_R)/2$. Applying the last inequality on small circles about points of the right ray (and using absolute convergence beyond the right edge) gives, uniformly for $v\ge B$ and $|y|\le(3/2)D^{\tau_0}$,
\[
  |\log L(c+v+iy)|+|(\log L)'(c+v+iy)|+|(\log L)''(c+v+iy)|
                         \ll (\log D)^{1/2-b_1}
\]
for some $b_1>0$. Thus the right-half integral remains coherent even with its cutoff inside the critical strip. More explicitly its phase Lipschitz bound is $O(U^{1/2-b_1})$ for $h_T\le\tau_0$, its logarithmic modulus curvature is $o(U)$ from above, and its slope is at least $U/2-O(U^{1/2-b_1})$. At $t\asymp U^{-1}$ the variance bound remains $O(t)$ and the phase error is $O(U^{-2b_1})$. The logarithmic error in the high Gaussian exponent in \eqref{hsub:threshold} is $o(U)$, too small to change any strict comparison. On the new right edge the Euler factor and its reciprocal are bounded by $\exp(o(U))$; this contributes $o(U)$ to both vertical boundary logarithms in the proof of \eqref{hsub:strip}. For the Gaussian-damped comparison at $T+w$ first restrict to $|w|\le D^{\tau_0}/4$, where the new estimate holds for all the present $T$. At larger $|w|$ ordinary convexity on both vertical edges, the gamma-factor comparison, and $e^{-w^2}$ absorb the tails uniformly. Therefore the proof of the common estimate still applies with a multiplicative $\exp(o(U))$ in \eqref{hsub:strip}.

For $h=h_T\le\tau_0$, the actual normalized central saving, up to an arbitrarily small fixed loss allowed by \eqref{hsub:sup}, is now
\[
                a_B(h)= B/2-\frac{h}{2(1+h)} .
\]
Indeed $\log(1+|T|)=h\log D+O(1)$ and $U=2(1+h)\log D+O_E(1)$. We may take $B$ sufficiently close to $1/2$ that $a_B$ is positive on this compact range. As $a$ increases to $a_B(h)$, the required lower bound for $tU_0$ from \eqref{hsub:threshold}, where $U_0=2\log D$, tends to
\[
                f(B,h)=\frac{2B^2}{2B(1+h)-h}.
\]
Now $f(1/2,h)=1/2$, while $\partial_B f(1/2,h)=1-h\ge1-\tau_0>0$. Lowering $B$ by a sufficiently small fixed amount consequently makes $\max_{0\le h\le\tau_0} f(B,h)<1/2$. First fix such a $B$, then take the loss in $a$ positive and sufficiently small so the strict margin persists. In particular the actual $a$ stays below $B/2$ even when $h=0$, guaranteeing uniformity of the Gaussian comparison.

For $h_T\ge\tau_0$, the universal threshold gives the sufficient lower bound $(w_0+\epsilon)/(1+\tau_0)$ on $tU_0$, up to $o(1)$. This is strictly below $1/2$ for sufficiently small $\epsilon>0$ by choice of $\tau_0$. Choose one constant $v_*<1/2$ strictly above the sufficient bounds in both height ranges. At $t=v_*/(2\log D)$ every zero is on the critical line. The difference $\log|a|-\log D=O(1)$ can be absorbed into a fixed strictly positive margin, proving the theorem and its prime specialization.
\end{proof}

\begin{corollary}\label{hsub:numeric}
For each fixed rational elliptic curve $E$ and all sufficiently large $D$, one may impose the explicit inequality
\[
              \lambda_{E^{(a)}}\log |a|\le 247/1000
\]
with a power-saving exceptional count as in Theorem~\ref{hsub:strict}. Without any exceptional set the same fixed-twist family has $\limsup_{|a|\to\infty}\lambda_{E^{(a)}}\log|a|\le64/153$.
\end{corollary}
\begin{proof}
H.~Wu, \href{https://doi.org/10.1007/s00039-014-0277-4}{\emph{Burgess-like subconvex bounds for $\mathrm{GL}_2\times\mathrm{GL}_1$}}, bounds $L(1/2,\pi\otimes\chi)$ for fixed $\pi$ by
$\ll_{\pi,\nu} C(\chi)^{1/2-(1-2\vartheta)/8+\nu}$ in the full analytic conductor of $\chi$, including its archimedean norm twist. Use the admissible Kim--Sarnak exponent $\vartheta=7/64$ over $\Q$. Here $C(\chi)\asymp_E D(3+|T|)$, so the degree-two saving in $Q_a(T)\asymp_E C(\chi)^2$ is $\delta=25/512$, allowing an arbitrarily small epsilon loss as above. Hence $w_0=128/153$. The universal assertion follows by dividing this number by two. Only finitely many fixed auxiliary base curves are used when twisting parameters share a prime factor with the original conductor.

For the density assertion take
\[
 \tau_0=696/1000,\qquad \tau_1=7/10,\qquad u=44/100,
 \qquad B=472/1000 .
\]
Here $2u/(3-4u)=22/31>7/10$ and $B>(1/2+u)/2=47/100$. Thus one may for instance take $\sigma_L=941/1000$ and $\sigma_R=1001/1000$ for the logarithm interpolation. As $2B<1$, the rational function $f(B,h)$ increases with $h$. Its maximum in the low range and the high-range bound are respectively
\[
        f(B,\tau_0)=6962/14141,\qquad
        w_0/(1+\tau_0)=4000/8109 .
\]
Both are strictly below $4937/10000$. All auxiliary losses in the savings may be taken small enough for this strict inequality to persist. Therefore $t=(4937/10000)/(2\log D)$ suffices on the good twists. Its product with $\log|a|$ is at most $4937/20000+o(1)<247/1000$.
\end{proof}

\subsection*{Times smaller than every fixed inverse-log coefficient in centered windows}
The all-height bound and a count through a growing centered window answer different questions. In a centered window the coefficient of the inverse logarithm may be made arbitrarily small.

\begin{theorem}\label{hsub:local}
Fix a rational elliptic curve $E$, $0<\alpha<1$, and $c_0>0$. Put $L=\log D$, $H=\exp(L^\alpha)$. Except for
\[
             O_{E,\alpha,c_0}\!\left(D\exp\{-bL\log\log L/\log L\}\right)
\]
squarefree $a$ with $D\le|a|\le2D$, all the zeros through physical height $|T|\le H$ of the heat deformation of $E^{(a)}$ are real at \emph{every} time $t\ge c_0/L$. They are simple for $t>c_0/L$. The conclusion holds in relative density one among primes in each fixed reduced progression. For each fixed $\alpha$ one may replace $c_0$ by some function tending to zero sufficiently slowly and retain the density-one conclusions.
\end{theorem}
\begin{proof}
We show that a short Euler sum controls the phase on any fixed ray starting at unitary real part $1/2+B$, for every $B>0$. It suffices to take $0<B<1/4$. From Theorem~\ref{ang:quadratic}, with $u=B/4$, there are $b_0,\tau>0$ such that, outside $O(D^{1-b_0})$ parameters, the right-hand half of the strip $\sigma\ge1/2+B/4$ contains no zero through $2D^\tau$. Write $M_a(z)=-L'_a/L_a(z)$ for the finite unitary function of the elliptic twist. On $1/2+B/2\le\sigma\le3$ and $|y|\le D^\tau$ the local zero-count argument used above gives $|M_a(\sigma+iy)|\ll_B L$.

Choose a smooth $W$ equal to one on $[0,1]$ and zero on $[2,\infty)$, and set $X=L^{C_B}$ with $C_B$ sufficiently large. On $1/2+3B/4\le\sigma\le 5/2$ and $|y|\le4H$, Mellin inversion gives
\[
 M_a(z)=\sum_{n\le2X}\frac{\Lambda_a(n)}{n^z}W(n/X)+O_B(1),
                \qquad z=\sigma+iy ,
\]
where $\Lambda_a$ are the logarithmic-derivative coefficients. To verify this, the Mellin transform of $W$ has a simple pole of residue one at zero and is rapidly decreasing on fixed vertical lines, by integration by parts on the support of $W'$. Move its inversion line from $2$ to $-B/8$, up to cutoffs $|\Im w|=D^\tau/2$. The new line costs $O_B(LX^{-B/8})$; the horizontal sides and the original tails are $O_B(1)$ using rapid decrease, the just-proved log-derivative bound and absolute convergence on the initial line. The enlargement $5/2+2>3$ causes no issue, since to the right of $2$ the Euler derivative is bounded absolutely. On smaller rectangles Cauchy's formula on circles of fixed small radius gives the differentiated version with $O_B(1)$ as well.

For the first derivative and the sum itself the contributions of powers $\ell^k$, $k\ge2$, are $O_B(1)$ uniformly, by temperedness and absolute convergence on each of these rays. The first powers, up to bounded contributions from fixed primes, have the form
\[
       P_a(z)=\sum_{\ell\le2X} b_\ell(z)\chi_a(\ell),\qquad
       |b_\ell(z)|\ll_B (\log(2\ell))^2\ell^{-1/2-B/2}.
\]
Here $\chi_a$ is the primitive quadratic character after fixing one of the finitely many bounded twisting factors in the proof of Theorem~\ref{ang:quadratic}. In particular $\sum_\ell|b_\ell|^2\ll_B1$. We need a bound on the maximum over the full window, so we give the moment estimate explicitly.

After removing the bounded local factors, the varying parameter is an odd fundamental discriminant $d$ in a fixed sign and progression, with $|d|\asymp_E D$. At primes outside the fixed progression modulus $Q$ one has $\chi_d(\ell)=(d/\ell)$. Drop squarefreeness to majorize the nonnegative $2k$th moment by the sum over \emph{all} integers in that progression. In the expansion, frequencies $m$ which are squares have coefficient mass at most $(2k-1)!!(\sum|b_\ell|^2)^k$ by pairing equal primes. For a nonsquare $m$ the character $(d/m)$ has sum zero over a complete set of residues modulo $m$, also when imprimitive, as a prime of its squarefree kernel lies outside $Q$. Coprimality with $Q$ thus bounds the incomplete progression sum by $O_Q(m)$. Here $m\le(2X)^{2k}$ and the total absolute coefficient mass is at most $(C_B X)^{2k}$. Therefore
\[
  D^{-1}\sum_a |P_a(z)|^{2k}
          \ll_{E,B} (C_B k)^k+D^{-1}(C_B X)^{4k}.
\]
Taking $k=\lfloor c_B L/\log L\rfloor$ with $c_B>0$ suitably small bounds the last term by $D^{-1/2}$. Markov's inequality at threshold $\sqrt L/(\log L)^{1/4}$ now gives an exceptional fraction $O(\exp\{-bL\log\log L/\log L\})$. This bound survives the union over a grid in $1/2+B\le\Re z\le2$ and $|\Im z|\le3H$. Indeed the polynomial coefficient derivatives in either coordinate cost at most a fixed power of $X$, so a mesh of spacing $L^{-C'_B}$ for a sufficiently large constant $C'_B$ extends the estimate to the full rectangle with at most a fixed factor in the threshold. There are $O(HL^{2C'_B})$ grid points, and $L^\alpha=o(L\log\log L/\log L)$. Both the sum and its first derivative have the same moment and modulus estimates. Beyond real part two the Euler series is absolutely bounded. Consequently on $v\ge B$, $|y|\le 2H$,
\[
  |\log L_a(c+v+iy)|+|(\log L_a)'(c+v+iy)|
      +|(\log L_a)''(c+v+iy)|=o_B(\sqrt L)
\]
on the good twists, since the logarithm itself follows by horizontal integration from the absolutely convergent edge.

Finally use \eqref{hsub:sup} with an arbitrarily small fixed $\nu>0$. Here $\log(3+|T|)=o(L)$, so its central bound is $e^{\nu L+o(L)}$ throughout $|T|\le2H$. Exactly as in Theorem~\ref{hsub:strict}, with $U=2L+o(L)$, the coherent estimates and \eqref{hsub:strip} are available on this interior cutoff ray with effective $a$ as close as desired to $B/2$ from below, and an additional $e^{o(U)}$ harmless at a fixed strict margin. In the Gaussian-damped three-lines proof restrict first to $|w|\le H/2$ for $|T|\le H$; outside, convexity and the factor $e^{-w^2}$ remove the tail. All estimates can be established on a window enlarged by any fixed factor at the outset. The right-hand side of \eqref{hsub:threshold} tends to $B$ as $a\uparrow B/2$. Choose $B<\min(1/4,c_0)$, then fix the losses in $a$ so that this threshold is strictly less than $2c_0$. For every $c_0/L\le t\le 1/L$ the product $s=tU$ exceeds that threshold uniformly and stays bounded. The complex-half comparison excludes off-line zeros throughout the enlarged window. If $c_0>1$ the intermediate interval is empty. For $t\ge1/L$ the universal all-height bound for this fixed-twist family in Corollary~\ref{hsub:numeric} already suffices. At a multiple real zero, the backward local Hermite expansion produces nonreal zeros at arbitrarily close smaller times. For $t>c_0/L$ these would lie within the enlarged window, a contradiction. This proves simplicity.

The exceptional count is $o(D/\log D)$, giving the prime specialization. For the final assertion take $c_0=1/n$ on successive, sufficiently late ranges of $D$ to make the exceptional fraction at most $1/n$ among squarefree parameters and in each of the first $n$ fixed reduced progressions. This produces the claimed slowly decreasing coefficient.
\end{proof}

\begin{theorem}[Polynomial centered height at arbitrarily small inverse-log time]\label{hsub:poly}
Fix $c_0>0$ and put $\tau=\min(c_0,1/4)^2/32$. Aside from $O_{E,c_0}(D^{1-b})$ squarefree twisting parameters $D\le|a|\le2D$, for some $b>0$, all zeros through physical height $|T|\le D^\tau$ of the heat-deformed elliptic function of $E^{(a)}$ are real simultaneously at every $t\ge c_0/\log D$, and simple at every strictly larger time. This remains true for a proportion tending to one of the prime twisting parameters in any fixed reduced progression.
\end{theorem}
\begin{proof}
The gain over Theorem~\ref{hsub:local} comes from separating the small primes before taking the high moment. Put $B=\min(c_0,1/4)$, $L=\log D$, $Y=L$, $X=L^{2/B}$, and $k=\lfloor BL/(16\log L)\rfloor$. Then $\tau=B^2/32$. For large $D$ the nonsquare error in the preceding moment proof is $O(D^{-2/5})$ because $(C_B X)^{4k}\le D^{3/5}$. For the short Euler approximation take $u=B/8$. The height exponent in Theorem~\ref{ang:quadratic} can be taken larger than $2\tau$ since $2u/(3-4u)> B/12>2\tau$. Start at $\Re z\ge c+15B/16$ and move the Mellin line to $-3B/4$, which leaves $\Re(z+w)\ge c+3B/16>c+u$. Its error is $O_B(L X^{-3B/4})=O_B(L^{-1/2})$. Cauchy differentiation for $\Re z\ge c+B$ can use circles of radius $B/32$.

For both prime polynomials on $v\ge B$, the primes no larger than $Y$ cost, without cancellation,
\[
              O_B\big(Y^{1/2-B}(\log(2Y))^2\big).
\]
For the remaining primes the squared coefficient sum, now using the actual real part $\sigma\ge1/2+B$, is
\[
 \sum_{\ell>Y}|b_\ell|^2
       \ll_B\sum_{n>Y}(\log(2n))^4 n^{-1-2B}
       \ll_B Y^{-2B}(\log(2Y))^4.
\]
Divide the $2k$th moment over these large primes by $L^{(1-B/2)k}$ to test the threshold $L^{1/2-B/4}$. The square frequencies cost at most
\[
      \left(C_B k L^{-2B}(\log(2L))^4 L^{-(1-B/2)}\right)^k
          \le L^{-Bk}\le D^{-B^2/17}
\]
for all sufficiently large $D$. The nonsquare error remains $O_B(D^{-2/5})$. The small-prime contribution is $o_B(L^{1/2-B/4})$. As in the preceding proof, the coefficient derivatives have polynomial bounds in $L$ and a fine polynomial mesh suffices in both coordinates. The grid has $O_B(D^\tau L^{C'_B})$ points, including a fixed enlargement in height. The resulting exceptional proportion is bounded by a negative power of $D$. On every good twist the logarithm and its first two derivatives are $O_B(L^{1/2-B/4})=o_B(\sqrt L)$ on the right ray through that height.

Use \eqref{hsub:sup} with $\nu=B/16$. At $h_T\le\tau+o(1)$ the available boundary saving is at least $B/2-h_T/(2(1+h_T))-\nu$, exceeding $B/4$. Thus \eqref{hsub:strip} can be used with $a=B/4$ and the harmless $e^{o(U)}$ from the interior right edge. The sufficient threshold for $tU$ is now $4B/3$. For $c_0/\log D\le t\le1/\log D$ one has $tU\ge2c_0+o(1)$, uniformly and with a strict margin. The preceding phase, reflection, and simplicity arguments apply to every such time. The global bound handles $t\ge1/\log D$. All exceptional sets in this proof have power saving, proving also the prime assertion.
\end{proof}

\subsection*{An explicit coefficient tending to zero for prime elliptic twists}
The preceding short Euler proof can be made uniform in its cutoff when the logarithmic-loss mean estimate for prime twists replaces estimates with an epsilon power of the conductor.

\begin{theorem}\label{hsub:explicitlocal}
Let $E/\Q$ have odd conductor, let $8N_E\mid Q$, and sample $D\le p\le2D$, $p\equiv1\pmod Q$. Write $L=\log D$, $S=\log L$. Uniformly for any prescribed
\[
 100\frac{\log S}{S}\le B\le1/8,\qquad H=\exp(B^2L/65536),
\]
all but a proportion $O_{E,Q}(e^{-bB^2L})$ of the sampled twists have only real heat zeros throughout $|T|\le H$, at every physical time $t\ge B/L$. They are simple at every strictly larger time. In particular one may take
\[
 t\ge \frac{100\log\log\log D}{\log D\,\log\log D},
 \qquad
 |T|\le\exp\!\left(\frac{10000}{65536}\,
           \frac{\log D\,(\log\log\log D)^2}{(\log\log D)^2}\right).
\]
\end{theorem}
\begin{proof}
We track the losses as $B$ decreases. The prime-twist estimate in Theorem~\ref{etr:strip} applied with $u=B/8$ and height $H^{10}$ bounds the exceptional proportion for a zero to its right by
\[
        O_{E,Q}\left((1+B^{-1})H^{25}e^{-BL/16} L^C\right)
                                   \ll_{E,Q} e^{-BL/30}.
\]
The analytic conductors at these and all smaller heights have logarithm $O_E(L)$. On the interior line $\sigma\ge c+3B/16$ the local-zero estimate used above is now $O_E(L/B)$ since the gap is at least $B/16$. For the Euler approximation choose $X=L^{2/B}$ and move the Mellin line by $-3B/4$ as in Theorem~\ref{hsub:poly}, with horizontal cutoffs at $H^{10}/2$. The Mellin transform of the fixed $W$ is bounded by $C_j(1+|y|)^{-j}(1+B^{-1})$ on this line, since its only nearby singularity is its simple pole at zero. Thus the new vertical error is $O_E(B^{-2} L^{-1/2})$ and on differentiation it costs at most an additional $O(B^{-1})$ by Cauchy's formula. The original tails and horizontal segments are negligible: $\log X=2S/B=o(B^2L)$ uniformly, and the Mellin transform has arbitrary fixed polynomial decay. Higher prime powers, with up to one extra logarithmic derivative, cost a fixed power of $1/B$ by summing $\sum n^{-1-B}(\log(2n))^2$ and the resulting geometric tails. These costs are negligible compared with the moment threshold below.

Again set $Y=L$, $k=\lfloor BL/(16S)\rfloor$. The small-prime bound is $O(L^{1/2-B}S^2)$. The remaining squared coefficient mass is $O(B^{-5}L^{-2B}S^4)$; this follows by integrating $(\log x)^4 x^{-1-2B}$ above $L$. Both implied constants are absolute up to the fixed local data. The prime sums here can be majorized by the moments over all integers in the fixed progression at a cost $O_Q(L)$ after normalization on primes. This does not affect the exponential saving. For the small primes, $L^{-3B/4}S^2=o(1)$, and for the large-prime paired terms the $k$th root of the moment after testing $L^{1/2-B/4}$ is at most
\[
              C k B^{-5}S^4 L^{-2B}L^{-(1-B/2)}
                      \le C S^8L^{-3B/2}\le L^{-B}.
\]
These inequalities are uniform since $BS\ge100\log S$. The total nonsquare error is $O(D^{-2/5})$ before the prime normalization: the coefficient mass is bounded by $(C X)^{2k}$ and the incomplete-character cost by $(2X)^{2k}$, with $\log X=2S/B$. Finally $L^{-Bk}\le e^{-B^2 L/17}$ for large $D$. The time and real-part grid has $O(H X^{C'})$ points with an absolute fixed exponent, since differentiating these two polynomials on this region costs only powers of $X$. Its density cost is more than absorbed by this last saving. These estimates hold also on a fixed enlargement of the time window. We obtain a bound $O_E(L^{1/2-B/4})=o(\sqrt L)$ for each of the Euler logarithm and its first two derivatives on $v\ge B$ through $2H$. Horizontal integration from 2 gives the claim for the logarithm itself.

For the central boundary use \eqref{etr:central}. Its squared-derivative analogue costs only another power of $L$ for $|y|\le 4H$. Indeed the proof there uses a smooth dyadic partition and Fourier inversion on each shell. On the Cauchy circle of radius $1/\log(3D(2+|y|))$, the additional norm weight $n^{-(\sigma-1/2)}$ and its first four derivatives in the dyadic log variable are bounded, apart from a bounded factor on the normal length and a small power on the rapidly decaying outer shells. The dual gamma ratio gains a bounded factor. Thus the same proof holds on that circle, and Cauchy's formula costs $O(L^2)$ for the squared derivative. Integrating the two moments over unit intervals and using Sobolev as in \eqref{hsub:sup} now bounds the prime mean squared supremum through $3H$ by $O_{E,Q}(H^4 L^{C'})$. At threshold $e^{BL/32}$ Markov's inequality excludes a proportion $O_{E,Q}(e^{-BL/30})$. This step uses the stated log-power estimate for the sampled primes.

Use $U=\log Q_p(T)=2L+O_E(B^2L+1)$. The new central bound and the gamma ratio between real parts $c$ and $c+B$ give an available saving larger than $B/4$ in \eqref{hsub:strip}. The error from the interior right boundary is $O_E(L^{1/2-B/4})=o(BL)$ in the logarithm. The gamma-ratio and vertical gamma estimates are uniform for $0\le B\le1/8$ since the fixed weight-two gamma argument stays in $\Re z\ge1$. The parts of the Gaussian-damped vertical comparison with $|w|>H/2$ again follow from polynomial growth and $e^{-w^2}$. These statements hold on an enlarged window by starting with the preceding estimates on fixed multiples of $H$.

For clarity the shrinking-$B$ comparison has a uniform positive margin. With $a=B/4$ and $B/L\le t\le 1/L$, one has $s=tU\ge2B-O_E(L^{-1})$. At $r=0$ the high-minus-low logarithmic exponent divided by $U$ is at least $B/8-o(B)$. Lowering $r$ to $-B/8$ loses at most $B/32+o(B)$; increasing it only improves the estimate. The completed-square center stays above $B$. The slope error $O_E(L^{1/2-B/4})$ contributes $o(BL)$ even at $t=1/L$. Variance is $O(t)$ and the phase error is $O(L^{-B/2})=o(1)$. In the derivative estimate $\sqrt t=o(B)$, so $|A_r|$ is increasing for $|r|\le B/8$. For $r\ge B/8$, the reflected high mass is at most $e^{-B^2 L/8}$ times the original mass. This proves the result for the intermediate times. The all-height bound for $t\ge1/L$, and simplicity at strictly larger times, are as before.
\end{proof}

\clearpage
\section{Local strip contraction at the square of the original width}\label{hlc:section}

A local bound for the horizontal displacements of the ordinary zeros gives a bound at positive heat time on the same physical height interval. The time is essentially half the square of the initial half-width. The possible zeros outside the larger initial interval enter through the second logarithmic derivative of an outer factor.

\begin{theorem}\label{hlc:main}
Fix a rational elliptic curve $E$ and $\eta>0$. Let $D\le |d|\le2D$ be squarefree, $L=\log D$, $1\le H\le\exp(\sqrt L)$, $H_1=4H+L(1+\log L)^3$, and $0<u\le1/2$. Suppose that every ordinary nontrivial zero of $L(s,E^{(d)})$ with ordinate at most $H_1$ in absolute value has distance at most $u$ from the critical line. For sufficiently large $D$ depending only on $E,\eta$, the physically normalized heat deformation has only real simple zeros through $|T|\le H$ at \emph{every} time
\[
                         t\ge(1+\eta)u^2/2 .
\]
\end{theorem}

\begin{corollary}\label{hlc:prime}
Fix an elliptic curve of odd conductor, $8N_E\mid Q$, and $M,\eta,\delta>0$. There is $K=K(E,Q,M,\delta)$ such that, for any prescribed $1\le H\le\exp(\sqrt{\log D})$, all but an $O_{E,Q,M,\delta}((\log D)^{-M})$ proportion of primes $p\in[D,2D]$, $p\equiv1\pmod Q$, have only real simple heat zeros through $|T|\le H$ at every time
\[
 t\ \ge\ \frac{1+\eta}{2}
       \left(\frac{(5+\delta)\log(H+2)+K\log\log D}{\log D}\right)^2 .
\]
In particular the time is $O_{E,Q,A,M}((\log\log D/\log D)^2)$ through $H=(\log D)^A$ for every fixed $A>0$.
\end{corollary}
\begin{proof}[Proof of the corollary]
Apply Theorem~\ref{etr:strip} with height $H_1$ and with $u$ equal to the positive fraction in the display. Its exceptional proportion is $\ll_{E,Q}(1+u^{-1})(H_1+1)^{5/2}\exp(-uL/2)L^C$. The displayed choice cancels the height factor and permits any specified logarithmic saving by increasing $K$. Horizontal reflection supplies the other half of the strip. Theorem~\ref{hlc:main} gives the assertion.
\end{proof}

\subsection*{A uniform count for the heat evolution of the finite zero polynomial}
Write $F(z)$ for the real phased ordinary completion in physical coordinates $z=x+iy$. Its real part is physical height and $|y|$ is distance from the critical line. Let $P$ be the monic polynomial formed from \emph{all} zeros of $F$ with $|x_j|\le H_1$, including multiplicity. Its coefficients are real. Let $n=\deg P$ and $S=1+\log((H+2)L)$. The usual unit-interval estimate in the explicit formula gives at most $C_E\log(3D(2+|v|))$ ordinary zeros with $v\le x_j\le v+1$. In particular the initial polynomial has at most $C_E L$ zeros in any unit interval, and $\log(2+n)\ll_E S$. For an entire function $J$ occurring here put $J_t=\exp(-t\partial_z^2)J$, with the vertical Gaussian definition of Section~\ref{hg:section}.

\begin{lemma}\label{hlc:count}
For $0\le t\le L^{-1}$ and every real $x$, the disk $|z-x|\le1$ contains at most $C_E L$ zeros of $P_t$, counted with multiplicity.
\end{lemma}
\begin{proof}
Use a large constant $R$, depending only on $E$. At each $z_*=w+iR$ write $A=P'/P(z_*)$ and
\[
 M_*(h)=\frac{P(z_*+h)}{P(z_*)}e^{-Ah},\qquad
 B_2=\sum_j |z_*-w_j|^{-2}\ll_E L/R ,
\]
where $w_j$ range over the roots of $P$. The bound follows by summing $CL/(R^2+j^2)$ over the horizontal unit intervals. In addition $|\Im A|\le(R+1)B_2\ll_E L$. The continuous real function $w+2t\Re A(w+iR)$ takes every real value, since $P$ is a polynomial and its logarithmic derivative tends to zero as $w\to\pm\infty$. Choose $w$ so that this value is the prescribed center $x$. Then $z_0=z_*+2tA$ has real part $x$ and imaginary part within $O_E(1)$ of $R$.

The elementary inequality $|(1+v)e^{-v}|\le e^{|v|^2/2}$ gives $|M_*(h)|\le e^{B_2|h|^2/2}$ for every complex $h$. On $|h|\le R/3$, Taylor's formula for $\log(1+h/(z_*-w_j))$ also gives
$|M_*(h)-1|\le C B_2 |h|^2e^{C B_2 |h|^2}$. Translation of the polynomial Gaussian integral now gives the exact formula
\[
 J(h):=\frac{P_t(z_0+h)e^{-Ah}}{P(z_*)e^{tA^2}}
    =\int_{\R} \frac{e^{-y^2/(4t)}}{\sqrt{4\pi t}} M_*(h+iy)\,dy .
\]
The translated line is permissible by polynomial Gaussian decay. Choose $R$ large so $tB_2\ll_E1/R$ is sufficiently small. The local Taylor bound then gives $|J(0)-1|\le O(tB_2)+O(e^{-R^2/(80t)})<1/2$, the tail using the global bound for $M_*$. On $|h|\le5R$ the global bound and Gaussian integration give $|J(h)|\le2e^{25B_2 R^2}$. As $|z_0-x|\le R+O_E(1)$, Jensen's inequality between radii $2R$ and $5R$ bounds the count in $|z-x|\le1$ by $O_E(1+B_2 R^2)=O_E(L)$. At $t=0$ the original unit-interval bound suffices.
\end{proof}

\subsection*{Spacing after the polynomial zeros reach the real line}
\begin{lemma}\label{hlc:gap}
Write $M=C_E L\ge2$ with constant large enough for Lemma~\ref{hlc:count}. If $(1+\eta)u^2/2\le t\le L^{-1}$ and $n\ge2$, all roots of $P_t$ are real and simple, and their minimum gap $g$ satisfies
\[
                    g\gg_\eta \min\{\sqrt{t/M},1/M\}.
\]
\end{lemma}
\begin{proof}
Polynomial strip contraction gives real roots from $t_1=u^2/2$ onwards. Indeed the proof of Lemma~\ref{heat:oneshift} applies directly to the finite product of the real factors and conjugate pairs of any nonconstant real polynomial; no decay condition is needed for that product. Iterating its shift average $m$ times with $h=\sqrt{2t/m}$ converges coefficientwise to $P_t$, by the Taylor series of the finite derivative operator. Its leading coefficient stays fixed, so Hurwitz gives the claimed polynomial strip contraction. For $v>t_1$ they are simple by the backward Hermite splitting at a multiple root, as in Section~\ref{heat:main}. Write $r_j$ for the ordered roots at these times. The heat equation gives $r_j'=2\sum_{k\ne j}(r_j-r_k)^{-1}$. At a smallest adjacent gap $g=r_{j+1}-r_j$, its derivative is $4/g$ minus twice the contributions $g/[d(d+g)]$ from the other roots, where $d$ is the distance to the nearer endpoint.

On either side the roots with $d<1$ number at most $M$, increasing the constant if necessary. Their ordered distances are at least $g,2g,\ldots$. Telescoping bounds their contributions on each side by $(1-1/(M+1))/g$. More distant roots contribute $O(Mg)$ altogether, by the same unit-interval count and $\sum k^{-2}<\infty$. Consequently $(g^2)'\ge8/(M+1)-CM g^2$ almost everywhere. The minimum of finitely many smooth gaps is locally Lipschitz at simple times, also when several gaps tie. Integrate this inequality from just after $t_1$ with nonnegative starting value. Since $t-t_1\ge\eta t/(1+\eta)$ it gives the claim.
\end{proof}

\subsection*{The outer factor and the differential product identity}
\begin{proof}[Proof of Theorem~\ref{hlc:main}]
The all-height bound of Corollary~\ref{hsub:numeric} handles $t\ge L^{-1}$ with strict simplicity, so assume $(1+\eta)u^2/2\le t\le L^{-1}$. Fix a physical ordinate $x$ with $|x|\le H$ at which we test for nonreal zeros. The ratio $Q_0=F/P$ is an entire real function of order at most one, with precisely the distant zeros left. Throughout $|\Re z-x|\le H_1/4$ its logarithm has second derivative bounded by
\[
           |(\log Q_0)''(z)|\le
           \sum_{|x_j|>H_1}|z-w_j|^{-2}\ll_E L/H_1 .
\]
Indeed the second derivative of the genus-at-most-one Hadamard product has precisely this absolutely convergent series up to sign. The ordinary unit-interval bound above gives the estimate after summing
$\sum_{j\ge H_1/2}\log(3D(2+H_1+j))/j^2$. It holds for \emph{every} imaginary part. The real number $a=Q_0'(x)/Q_0(x)$ satisfies $|a|\ll_E LS$. To see this evaluate the logarithmic derivative at $x+2i$. The original completion there has logarithmic derivative $O_E(L)$ by its functional equation and the Euler and gamma derivatives on the opposite side of the strip. For $P$ the absolute sum of the reciprocal distances at $x+2i$ is $O_E(LS)$, by the ordinary unit-interval bound. The second-derivative estimate moves the difference back to $x$.

Define $K(z)=Q_0(z)e^{-a(z-x)}/Q_0(x)$. Its logarithm in this vertical strip can be chosen to vanish together with its derivative at $x$. Thus
\[
             |\log K(z)|\le C_E (L/H_1)|z-x|^2 .
\]
Choose $R'=C_1 S$ with $C_1$ sufficiently large in terms of $E,\eta$. For large $D$ it is less than $H_1/16$, $2t|a|+4<R'/2$, and $L(R')^2/H_1=o(1)$, all uniformly in $H$. Indeed $S=O(\log L)$ if $H\le L^3$, and $L(1+\log((H+2)L))^2/H=o(1)$ uniformly if $H>L^3$. On $|z-x|\le 3R'$ vertical Gaussian integration therefore gives
\[
                 |K_t(z)-1|\ll_{E,\eta} L(R')^2/H_1=:\epsilon_D=o(1).
\]
In detail use $|e^w-1|\le |w| e^{|w|}$ with $|w|\ll_E (L/H_1)((3R')^2+y^2)$ inside the Gaussian integral. The Gaussian second moment remains $O(t)$ after this quadratic perturbation. Hence for $|z-x|\le R'$ and $k\ge1$ Cauchy's inequality gives
$|K_t^{(k)}(z)|/k!\le\epsilon_D/(R')^k$, on adjusting the implied constant once.

For every polynomial $P$ of degree $n$ and every entire $K$ of order strictly below two one has the exact identity
\begin{equation}\label{hlc:leibniz}
               (PK)_t(z)=\sum_{k=0}^n\frac{(-2t)^k}{k!}
                                P_t^{(k)}(z) K_t^{(k)}(z).
\end{equation}
Indeed $(e^{vz}K)_t=e^{vz-tv^2}K_t(z-2tv)$ by Gaussian translation. Differentiate at $v=0$ to obtain each monomial case; the order bound justifies the contour translation and differentiation. In particular, up to a nonzero constant,
$F_t(z)=e^{az-ta^2}(PK)_t(z-2ta)$.

Let $r_j$ be the roots of $P_t$ and put $g_*=\min(1,g)$ using Lemma~\ref{hlc:gap}; use $g_*=1$ for $n\le1$. Outside the open disks of radii $g_*/4$ centered at the $r_j$, one has
\[
 \Sigma(z)=\sum_j |z-r_j|^{-1}\ll \frac{1+\log(2+n)}{g_*},
 \qquad |P_t^{(k)}(z)/P_t(z)|\le\Sigma(z)^k .
\]
The first estimate orders the real roots by distance and uses their spacing. For the second, the derivative is a sum over ordered choices of distinct factors of the polynomial. Lemma~\ref{hlc:gap} implies there that
\[
   2t\Sigma(z)\ll_{E,\eta} S(t+\sqrt{tM}+tM)\ll_{E,\eta} S\le R'/2
\]
for $C_1$ sufficiently large, also if $n\le1$. Inserting the preceding Cauchy estimates into \eqref{hlc:leibniz} gives, outside these disks and for $|z-x|\le R'$,
\[
               |(PK)_t(z)/P_t(z)-1|\le 2\epsilon_D<1 .
\]
We now test a potential zero of $F_t$ at ordinate $x$ and imaginary part $|y|\le1/2$; the global strip bound excludes all larger displacements. Its translated position $x+iy-2ta$ is within $R'/2-2$ of $x$. If it lies outside the disks, the last inequality excludes it. If it lies in a disk, the entire boundary of that disk is within $R'$ of $x$. Rouch\'e's theorem gives exactly one zero of $(PK)_t$ inside, counted with multiplicity. Both the disk and function are real under conjugation, hence this unique zero is real and simple. The multiplier and displacement are real in the stated sense, so the same conclusion holds for $F_t$. The argument holds for each $x$ and each $t$ in the range, proving the theorem.
\end{proof}

\begin{corollary}[A lower bound for local gaps]\label{hlc:fullgap}
In Theorem~\ref{hlc:main}, for $(1+\eta)u^2/2\le t\le L^{-1}$, two distinct zeros of the full heat deformation with physical ordinates in $[-H,H]$ are separated by at least
\[
             c_{E,\eta}\min\!\left\{\sqrt{\frac{t}{L}},\frac1{L}\right\}.
\]
\end{corollary}
\begin{proof}
It suffices to consider two physical zeros at distance at most one, since the displayed lower bound tends to zero. Freeze the real slope $a$ at the first zero's ordinate $x$. Both translated points and both associated disks then lie within $R'$ of $x$ and use exactly the same model $(PK)_t$. They must belong to different disks since each disk contains one simple zero. Consequently the spacing is at least $g_*/2$, and Lemma~\ref{hlc:gap} applies. If $n\le1$ two such zeros cannot occur.
\end{proof}

The exponent two above comes directly from the finite polynomial strip contraction. The argument transfers it to a local interval even though the full ordinary completed function could have more distant off-line zeros outside the original larger interval.

\subsection*{A density-one result for each individual ordinary function}
\begin{theorem}\label{hlc:fixed}
Fix a cuspidal standard function of $\mathrm{GL}_2/\Q$ with unitary central character. Write $L=\log X$. For every $M,\eta>0$, any prescribed $1\le H\le\exp(\sqrt L)$, and a relative measure $1-O_{\pi,M}(L^{-M})$ of $v\in[1,2]$, all its physical heat zeros through $|T-vX|\le H$ are real and simple simultaneously at every time
\[
 t\ge \frac{1+\eta}{2}
        \left(\frac{8\log(4H+L(1+\log L)^3)+8(M+2)\log L}{\log X}\right)^2
\]
for sufficiently large $X$ depending on $\pi,M,\eta$. In particular, for almost every $v\in[1,2]$, at all sufficiently large dyadic scales $X=2^j$ and throughout $|T-vX|\le(\log X)^A$, the time $C_{A}(\log\log X/\log X)^2$ suffices. The latter assertion holds simultaneously for all rational elliptic curves and indeed all the indicated cuspidal standard functions, with the starting scale allowed to depend on the function and on $A$. One may take any fixed $C_A>32(\max\{A,1\}+4)^2$.
\end{theorem}
\begin{proof}
For a fixed $\pi$ the estimate in Corollary~\ref{ep:density} bounds the number of ordinary zeros of displacement greater than $u$ at positive ordinate through $3X$ by $O_\pi(XL\exp(-uL/8))$. Each such zero rules out a set of $v\in[1,2]$ of measure at most $2H_1/X$, where $H_1=4H+L(1+\log L)^3=o(X)$. Use the fraction appearing in the theorem for $u$. The resulting excluded measure is $O_\pi(L^{-M})$.

We spell out why the preceding deterministic transfer applies when the window moves through the zeros of one fixed form. Replace the original function $F(z)$ by $F(z+vX)$, and $P$ by the polynomial of all zeros with $|\Re z|\le H_1$ in the translated coordinates. Its coefficients are real also for nonselfdual forms, since the phased completion is real for real physical arguments. Its ordinary count per unit is $O_\pi(L)$ within this window. For the distant zeros use
\[
 \log(3+|vX+j|)\ll L+\log(3+|j|)
\]
in the convergent sum for the second logarithmic derivative. The bound $O_\pi(L/H_1)$ thus holds throughout the same vertical strip. At physical displacement $2i$, the completed log derivative of the original translated form is $O_\pi(L)$ by Stirling, the functional equation and the absolutely convergent Euler derivative. These are exactly the ordinary-function inputs in the preceding three arguments.

Moreover their short-time range $t\le 1/L$ may throughout be replaced by $t\le C_\pi/L$ for any fixed $C_\pi$. Increase the fixed constant $R$ in the centered Jensen argument so that $tB_2\ll_\pi C_\pi/R$ is sufficiently small and $|z_0-x|+1<2R$; the count remains $O_\pi(L)$. The gap estimate and $t\Sigma\ll_{\pi,\eta,C_\pi} S$ likewise keep the same powers, so the finite Leibniz argument applies verbatim with constants depending on $C_\pi$. Choose $C_\pi$ sufficiently large for Theorem~\ref{hg:main} to handle all $t\ge C_\pi/L$ in $|T-vX|\le H$, first up to $1/8$ and then at all larger times. This completes the single-scale assertion.

For the simultaneous version take $M=2$ and sum over dyadic $X$. The first Borel--Cantelli lemma gives one full-measure set for each fixed function. As in Theorem~\ref{common:main}, there are countably many such functions up to imaginary norm twists; the latter translate the physical ordinate by a constant and commute with the heat flow. Apply the single-scale result to the countable representatives with enlarged windows $2(\log X)^A$. To handle every real $A>0$ at once take the intersection for all positive rational exponents and use exponents decreasing to $A$ to retain any prescribed strict constant $C_A>32(\max\{A,1\}+4)^2$. The freedom in $\eta$ and in enlarging the eventual starting scale absorbs the strict constant margin. This proves the assertion.
\end{proof}

\begin{corollary}[Actual completions of elliptic curves over number fields]\label{hlc:fields}
The full single-scale and simultaneous fixed-function assertions of Theorem~\ref{hlc:fixed} hold for the physical heat deformation of each of the following Hasse--Weil completions:
\begin{enumerate}
\item $E/K$, where $K/\Q$ is finite abelian and $E$ is the base change of a rational elliptic curve;
\item every elliptic curve $E/K$ over an imaginary quadratic field with geometric complex multiplication by $K$.
\end{enumerate}
The conclusion applies to the heat deformation of the \emph{whole} number-field completion. In the second case the curve can be outside all rational base-change isogeny classes.
\end{corollary}
\begin{proof}
Each ordinary completed function in this list is a finite product of degree-two cuspidal completed standard functions over $\Q$, up to a nonzero scalar and the common shift to unitary normalization. In the abelian case Artin factorization gives the product over the Dirichlet characters of $\operatorname{Gal}(K/\Q)$ of the twists of the elliptic newform, including the full ramified factors as in Section~\ref{ebc:chapter}. In the second case the CM theorem (Theorem~\ref{cm:factor}) gives two angular Hecke characters over $K$. The Hecke theta series, equivalently quadratic automorphic induction, attaches a rational-base cuspidal standard function of degree two to each character. Cuspidality follows because the nonzero angular infinity type changes under conjugation of $K$. These identifications include the gamma and conductor factors up to a scalar.

The ordinary zero count and exceptional-interval estimate add over this fixed finite list, so they obey the very same bounds with constants depending on the original curve and field. The complete product in physical coordinates is phased real, entire of order one, has its zeros in the initial strip, and has logarithmic derivative $O_{E,K}(\log X)$ on the exterior line used in the proof. Thus the finite polynomial and the outer factor must now be formed from the zeros of the entire product. Their estimates prove the short-time conclusion for the heat of that product directly. At $t\ge C_{E,K}/\log X$ the proof of Theorem~\ref{hg:main} applies to this entire finite product: its standard analytic inputs are the single product functional equation, order-one bound, absolute convergence of its Dirichlet series and its reciprocal sufficiently far to the right, the fixed list of gamma shifts and polynomial vertical growth in fixed strips. All follow by multiplication from the fixed factors; strip contraction gives the larger-time step as before. Countable intersection gives simultaneity also in these number-field cases.

For a concrete instance take $K=\Q(i)$ and $E:y^2=x^3-(1+2i)x$. Its $j$-invariant is $1728$. At the prime $(1+2i)$ the integral discriminant $64(1+2i)^3$ has valuation three. The displayed model is minimal there since any integral model whose discriminant valuation is smaller would reduce it by a positive multiple of twelve. Thus reduction at this prime is bad. At the conjugate prime $(1-2i)$ the same discriminant is a unit, so reduction is good. Rational base change would have the same good-reduction status at conjugate primes, and the good-reduction criterion on the rational Tate module makes this status isogeny invariant. This curve is therefore not isogenous over $K$ to a rational base change.
\end{proof}

\subsection*{Counting the nonreal zeros themselves}
\begin{theorem}\label{hlc:defect}
Fix a completion covered by Theorem~\ref{hlc:fixed} or Corollary~\ref{hlc:fields}, and $0<c<\sqrt2/8$. At any physical heat time $t\ge0$ let $D(t;X)$ count, with multiplicity, its nonreal zeros and its multiple real zeros with physical ordinate in $[X,2X]$. For all sufficiently large $X$, uniformly in $t$,
\[
 D(t;X)\ll_{F,c}
 X(\log X)^3(1+\log\log X)^3
                  \exp\bigl(-c\sqrt t\log X\bigr).
\]
For every fixed $M>0$ there is $C_M>0$ such that
\[
 t\ge C_M(\log\log X/\log X)^2
       \quad\Longrightarrow\quad D(t;X)\ll_{F,M} X/(\log X)^M .
\]
For sufficiently large $X$ the count vanishes when $t\ge C_F/\log X$.
\end{theorem}
\begin{proof}
Write $L=\log X$. The last assertion is Theorem~\ref{hg:main} and its fixed-product extension above. For smaller $0\le t\le C_F/L$ we first show the unconditional bound $O_F(L)$ for the number of all deformed zeros in any unit interval of physical ordinates inside $[X-2,2X+2]$. Use the centered proof of Lemma~\ref{hlc:count} directly on the complete function, regardless of the locations of its ordinary zeros within $|\Im z|\le1/2$. At $z_*=x+iR$ the ratio
\[
       M_*(h)=\frac{F(z_*+h)e^{-hF'/F(z_*)}}{F(z_*)}
       =\prod_j\left(1+\frac{h}{z_*-w_j}\right)
                           \exp\!\left(-\frac{h}{z_*-w_j}\right)
\]
is its centered genus-one Hadamard product, with the affine exponential canceled by taking the logarithmic derivative at $z_*$. The ordinary unit count gives $B_2\ll_F L/R$ for every fixed sufficiently large $R$ and all sufficiently large $X$. For instance, group the horizontal distances in unit intervals, use $O_F(L+\log(3+j))$ zeros in the $j$th pair, and sum against $(R^2+j^2)^{-1}$. For fixed $R$ the extra $O_F(\log(3+R)/R)$ is absorbed for large $X$. The implied coefficient of $L/R$ can be taken independent of $R$. At the exterior point $|F'/F(z_*)|\ll_F L$ by Stirling, functional equation and the absolutely convergent Euler product, with the coefficient of $L$ again independent of fixed $R$ for large $X$. Take $R$ sufficiently large in terms of $F,C_F$. Then $z_0=z_*+2tF'/F(z_*)$ differs from $z_*$ by a bounded distance less than $R/4$, and the Gaussian translation, centered product inequalities and Jensen argument in the lemma apply without change. Contour translation uses order one. The bound obtained is $O_F(L)$ per unit, since all zeros are within $|\Im z|\le1/2$.

Choose $\eta>0$ so that $1/(8)\sqrt{2/(1+\eta)}>c$ and set $u=\sqrt{2t/(1+\eta)}$. For positive $t$ in the remaining range we have $0<u<1/2$ when $X$ is large. Partition $[X,2X]$ into intervals of length at most one, using integer endpoints in the interior. If one contains a nonreal zero or a multiple real zero, the translated deterministic transfer in the proof of Theorem~\ref{hlc:fixed}, with $H=2$, forces an \emph{ordinary} zero of displacement greater than $u$ within $O(L(1+\log L)^3)$ of its midpoint. Each ordinary zero can affect at most $O(L(1+\log L)^3)$ of these intervals. The ordinary density estimate bounds the number of such zeros by $O_F(XL e^{-uL/8})$, also for the finite products by summation. Multiply by the bound per affected interval and the $O_F(L)$ deformed multiplicity there to obtain the result. At $t=0$ the ordinary unit bound suffices. Taking $c\sqrt{C_M}>M+4$ gives the second assertion.
\end{proof}

\begin{corollary}[Comparison of the complete zero counts]\label{hlc:total}
For any fixed completion in Theorem~\ref{hlc:defect}, write $N_t(I)$ for the count of all its heat zeros with physical ordinates in $I$, with multiplicity. For any fixed $C>0$, uniformly over all intervals $I\subset [X,2X]$ and $0\le t\le C/\log X$,
\[
                 N_t(I)=N_0(I)+O_{F,C}(\log X).
\]
In particular, for $0<c<\sqrt2/8$ the fraction of nonreal or multiple real zeros in the full dyadic block is
\[
  \frac{D(t;X)}{N_t([X,2X])}
       \ll_{F,c} (\log X)^2(1+\log\log X)^3
                              e^{-c\sqrt t\log X}
\]
whenever the denominator is nonzero. At the times in Theorem~\ref{hlc:fixed} lying below $C/\log X$, each good moving interval of length $2H\to\infty$ contains $(d/\pi+o(1))H\log X$ simple real heat zeros, where $d$ is the full absolute degree.
\end{corollary}
\begin{proof}
Again put $L=\log X$, and use the full-function centered estimate of Theorem~\ref{hlc:defect}. With fixed large $R$, $z_*=w+iR$, $A(w)=F'/F(z_*)$ and $z_0(w)=z_*+2tA(w)$, it gives uniformly for $w\in[X-O(1),2X+O(1)]$
\[
 \left|\frac{F_t(z_0(w))}{F(w+iR)e^{tA(w)^2}}-1\right|<1/2,
 \qquad
 \sup_{|h|\le10R}
       \big|F_t(z_0(w)+h)/F_t(z_0(w))\big|\le e^{C_{F,R}L}.
\]
The second inequality uses the same centered Gaussian bound on the larger fixed disk, putting back $e^{A(w)h}$; $|A(w)|\ll_F L$. On this exterior line also $|A'(w)|\ll_{F,R}1$ by the differentiated gamma and Euler series. Thus $\Re z_0(w)$ is strictly increasing for large $X$, differs from $w$ by $O_{F,R,C}(1)$, and $\Im z_0(w)$ stays between $R/2$ and $3R/2$ if $R$ was taken large enough.

Temporarily move the endpoints of $I$ by distances at most one, within a bounded enlargement of $[X,2X]$, so that their vertical lines miss all ordinary and deformed zeros. The unit counts cost $O_{F,C}(L)$ for either function. At each endpoint $b$ choose $w_b$ such that $\Re z_0(w_b)=b$. Take the contour with top side the corresponding segment of $z_0(w)$, bottom side its complex conjugate and vertical ends. The strip confinement puts all the desired zeros inside. On the top, the change in argument of $F_t$ equals the change in argument of $F(w+iR)$ up to $O_{F,C}(L)$. Indeed the factor $e^{tA(w)^2}$ contributes endpoint phases of magnitude $O(tL^2)=O_C(L)$, and the quotient in the first display has a single-valued argument of magnitude less than $\pi/2$. The shifts of at most $O_{F,R,C}(1)$ between the $w_b$ and the physical endpoints cost $O_{F,R,C}(L)$ by the exterior log derivative bound. Conjugation gives the same comparison on the lower side.

The change in argument on each vertical end is $O_{F,C}(L)$. To verify this take its top $z_0$, write $h_b=\Im z_0$, and form the entire function of $v$
\[
 B(v)=F_t(z_0-iv)/F_t(z_0),\qquad
 B_{\R}(v)=\tfrac12(B(v)+\overline{B(\bar v)}).
\]
We have $B_{\R}(0)=1$ and $|B_{\R}(v)|\le e^{C_{F,R}L}$ on $|v|\le10R$. Jensen bounds its real zeros on $[0,2h_b]$ by $O_{F,R}(L)$. Between consecutive zeros of the real part the nonvanishing value of $B$ lies in one open half-plane, where its continuous argument changes by at most $\pi$. This proves the asserted end bound. The same argument at $t=0$ bounds the ends of the ordinary rectangle with horizontal sides at $\pm iR$. The argument principle now compares the two full counts with error $O_{F,C}(L)$.

The ordinary Riemann--von Mangoldt formula gives $N_0([X,2X])\sim dXL/(2\pi)$ and $N_0(I)=d|I|L/(2\pi)+O_F(|I|+L)$ inside a dyadic block. It holds for the fixed products by addition. Use $C=C_F$ from Theorem~\ref{hlc:defect}; above its time range the defect count vanishes. These estimates prove the fraction and the assertion for growing intervals. If an interval extends slightly beyond the dyadic block, the same proof works on a fixed enlargement such as $[X/2,3X]$.
\end{proof}

\clearpage
\section{Two elliptic heat constants from finite walls and all-height repulsion}\label{hcfinite:chapter}

\begin{theorem}\label{hcfinite:bound}
In physical time the two elliptic de Bruijn--Newman constants satisfy
\[
             0\le\lambda_{11\mathrm a1}<0.084678,\qquad
             0\le\lambda_{234446.\mathrm a1}<0.069420 .
\]
The second assertion holds for both possible analytic ranks of the curve.
\end{theorem}

\begin{proof}
We first certify that every nonreal zero, if one exists, is at large height. For $11\mathrm a1$ take $Y=65395/16$. For every $0\le t\le1/8$ the vertical physical boundary $\Re z=Y$, $|\Im z|\le3/4$, is free of zeros. The check uses the even weight-two Mellin kernel for $[0,-1,1,-10,-20]$, with $B=2\pi/\sqrt{11}$. Here are the data for four such walls.
\[
\begin{array}{c|c|c|c|c}
 Y & (D,W,M) & (J,K) & \mathrm{retained\ degrees} &\mathrm{zeros\ inside}\\ \hline
 65395/16 &(1683,10,3456)&(1451,1501)&(28,17)&17374\\
 65441/32 &(846,9,1758)&(798,846)&(25,15)&7792\\
 32613/32 &(427,9,909)&(464,495)&(22,14)&3432\\
 8149/8   &(427,9,909)&(464,495)&(22,14)&3430
\end{array}
\]
For the four rows, precision is respectively $10320,5200,2816,2816$ bits. They use $19,32,15,15$ Bernstein rectangles, with maximum subdivision depths $14,16,10,10$, throughout the full heat and transverse intervals. At $65395/16$ the already scaled signed moments are converted to polynomial coefficients by outward rounding at $256$ bits; the whole calculation through formation and scaling of the moments uses the indicated $10320$ bits. Repeating the polynomial enclosure at the full $10320$ bits gives the identical degrees, subdivisions and count. The proof is the elliptic quadrature and bivariate Bernstein enclosure of Section~\ref{r4:heat}. Its error bounds apply unchanged: $u^2/8\le18+\tfrac12 B\cos(1)e^u$ for $u\ge0$ holds with the larger value of $B$ at level eleven, and the Fourier bound $|a_n|\le2n$ holds for the level-eleven form. The normalized joint error is below $2^{-22}$. The finite replay is \texttt{checks/heat-conductor/elliptic\_wall.py}, with grid-height parameters $4096,2048,1024,1024$. At $32613/32$ a second mesh at $3072$ bits, $(D,W,M)=(433,9,922)$ with grid-height parameter $1040$, gives the same enclosures and total count. At each positive wall the input brackets at height $4096$ contain respectively $8687,3896,1716,1715$ roots below it, with none at the boundary. At time zero these roots and their reflections are simple and there is no central zero. The strip theorem excludes zeros at the two remaining edges of the symmetric rectangle. Its number of nonreal zeros cannot increase with time by the local Hermite splitting argument of Theorem~\ref{r4:phase}; a multiple real zero at positive time would require a previous nonreal pair. Thus every zero inside the first wall is real and simple throughout. For $234446.\mathrm a1$, Corollary~\ref{r4:upperphase} supplies the corresponding statement with $Y=65257/32$ for $t\ge t_0=2^{-8231}$. Use $t_0=0$ at level eleven.

For all remaining possible outermost zeros put $|T|\ge Y$, $A=\sqrt N/(2\pi)$ and $R=3/2$. We sharpen \eqref{hc:ratio} between $R$ and $R+h$, with $h=2$. Use the notation of Section~\ref{hc:chapter}, here with $c=1$. On the right ray $v\ge B$ put $\sigma_B=B+1/2$ and
\[
 Z_B=\zeta(\sigma_B),\qquad
 j_{1,B}=-2(\log\zeta)'(\sigma_B),\qquad j_{2,B}=2(\log\zeta)''(\sigma_B).
\]
The Euler logarithm gives $|L|,|L|^{-1}\le Z_B^2$ and first and second log derivatives bounded by $j_{1,B},j_{2,B}$ respectively. The integral digamma identity gives
\[
       (\log|F|)'\ge D_B=\log(AY)-(1+B)^{-1}-j_{1,B}.
\]
Use this estimate with $B=1,3$, for which $D_B>0$. In the right-line comparison at $R_B=B+1/2$ the two low Gaussian tails relative to the absolute mass $P$ on $v\ge B$ are at most
\[
       2\exp\{-(R_B-B+2tD_B)^2/(4t)\}\le 2e^{-D_B}.
\]
Indeed bound the numerator by $2|F(B)|\sqrt{\pi t}\exp(-(R_B-B)^2/(4t))$ as before; $|F(v)|\le|F(B)|$ on $0\le v\le B$ follows by the initial product since $B\ge1$. Bound $P$ from below by $|F(B)|\sqrt{\pi t}\exp(D_B(R_B-B)+tD_B^2)$, integrating above the shifted Gaussian mean. The reflected high part costs at most $e^{-8R_BB}P$. Write $\delta_B=2e^{-D_B}+e^{-8R_BB}$.

For the numerator at $B=3$, the phase derivative is bounded by $l_B=\pi/2+1/Y+j_{1,B}$, and the logarithmic modulus second derivative by $u_B=\pi/(2Y)+Y^{-2}+j_{2,B}$. The first gamma estimate uses the first term and integral in the digamma sum with $|T|\ge Y$. For the second, sum $((c+v+j)^2+T^2)^{-1}$ by its first term and its integral. The variance bound in the previous proof gives complex coherence at least $w_B=1-l_B^2/(2(4-u_B))$ for this numerator. For the denominator only the upper bound $(1+\delta_1)P$ is needed. Two gamma recurrences and the Euler bounds now give
\[
 \frac{|F_t(c+R+h+iT)|}{|F_t(c+R+iT)|}\ge C(AY)^h,\qquad
 C=\frac{w_3-\delta_3}{(1+\delta_1)(Z_1Z_3)^2}.
\]

The removed conjugate pair at $T$ contributes at most
\[
       \log\frac{(R+h)^2-1/4}{R^2-1/4}
\]
to the integrated Poisson formula: its exact expression has $y^2$ in place of $1/4$ and increases with $y^2$. In the remaining terms $r-a_\rho\ge1$, and the numerators of reflected pairs sum to $2r$, giving the common integral coefficient $h(R+h/2)$ against the kernel $(1+(T-\gamma_\rho)^2)^{-1}$. Thus the additional repulsion coefficient can be taken to be
\[
 \kappa=
 \frac{\log(C(AY)^h)-\log(((R+h)^2-1/4)/(R^2-1/4))}{h(R+h/2)}.
\]
Starting the contraction at $t_0$, with squared height at most $1/4$, gives the all-height upper bound $t_0+\log(1+\kappa/2)/(4\kappa)$ as in Section~\ref{hc:chapter}. The two $\kappa$ values are greater than $2.1676$ and $3.8894$, respectively. The outward-rounded gamma, zeta and elementary evaluations used for the displayed bounds are in \texttt{checks/heat-conductor/fixed\_constants.py}.
\end{proof}

\clearpage
\section{The exact zero boundary as a limit of heat thresholds of oldforms}\label{oldheat:section}
Let $f(z)=\sum_{n\ge1}a_n e^{2\pi i n z}$ be a normalized holomorphic newform of weight $k\ge2$ and level $N$. Write $c=k/2$ in its arithmetic variable, and put
\[
 b_f=\sup_{\Lambda_f(\rho)=0}|\Re\rho-c|\in[0,1/2].
\]
For each integer $M\ge2$ take the classical oldform
\[
                  g_M(z)= f(z)+M^{k/2}f(Mz).
\]
The two degeneracy maps place it in the cusp space of level $NM$, with induced nebentype. Its completed Mellin transform at that level is phased to be real on the critical line. For this possibly nonprimitive function let $\lambda^+(g_M)$ be the infimum of the \emph{nonnegative} physical heat times after which all its zeros lie on that line. Using the nonnegative threshold avoids any prerequisite about negative-time deformations of an oldform.

\begin{theorem}\label{oldheat:main}
The undeformed rightmost zero boundary of the newform is recovered exactly by the ordinary limit
\[
                 \lim_{M\to\infty}(\log M)\lambda^+(g_M)=b_f .
\]
At every individual level,
\[
 0\le\lambda^+(g_M)\le
    \frac{\sqrt{1+b_f^2(\log M)^2}-1}{(\log M)^2}.
\]
In particular ERH for one rational elliptic curve $E$ is equivalent to $\lambda^+(f_E+M f_E(M\,\cdot))=o(1/\log M)$; it is enough that the scaled thresholds tend to zero along \emph{one} unbounded subsequence. The result holds also when no zero attains $b_f$.
\end{theorem}

\begin{proof}
The Mellin identity for the degeneracy map is $L(g_M,s)=(1+M^{c-s})L(f,s)$. Changing the conductor factor from $N^{s/2}$ to $(NM)^{s/2}$ gives, up to a nonzero constant,
\[
   G_{R,0}(s)=\cosh(R(s-c))F_0(s),\qquad R=\tfrac12\log M,
\]
where $F_0$ is the phased completion of $f$. This also proves the required reflection at the old level. The new member introduces an explicit lattice of critical-line zeros. The Euler product and zero divisor of the underlying primitive standard function are unchanged as representation-theoretic data. Here we deform the actual completed Mellin transform of $g_M$.

Write $F_t$ and $G_{R,t}$ for the respective Gaussian deformations, with variance $2t$ in the real $s$ direction. Completing the square in the two exponential terms gives the exact identity
\begin{align*}
  2e^{-tR^2}G_{R,t}(s)={}& e^{R(s-c)}F_t(s+2Rt)\\
                         &+e^{-R(s-c)}F_t(s-2Rt).
\end{align*}
All integrals converge locally uniformly by the fixed-completion bounds. De Bruijn's strip contraction puts the zeros of $F_t$ in $|\Re s-c|\le d_t=\sqrt{\max(0,b_f^2-2t)}$. For $r=\Re s-c>0$, the paired Hadamard product gives
\[
              |F_t(s+h)|\ge |F_t(s-h)|\qquad(h\ge d_t,\ h>0).
\]
Indeed a conjugate-reflected pair of real-part displacements $\pm b$ at ordinate $\gamma$ contributes squared-distance products whose difference has sign $8rh((\Im s-\gamma)^2+r^2+h^2-b^2)>0$. A critical-line zero gives the same strict orientation, and exponential corrections have equal moduli. Also $F_t(s+h)\ne0$ there, since $r+h>d_t$. If $h=2Rt\ge d_t$ the unequal moduli of the two exponentials in the exact identity preclude a zero in $r>0$. Reflection handles $r<0$. Solving $(2Rt)^2+2t\ge b_f^2$ and using the all-real persistence of de Bruijn proves the finite-level upper bound, including the case $b_f=0$ by inspection at $t=0$.

If $b_f>0$, fix $0<h<b_f$. Some zero $\rho$ has $\Re\rho>c+h$, by reflection. Choose a small closed disk about $\rho-h$ strictly inside $\Re s>c$ whose boundary contains no zero of $F_0(s+h)$. Take $t=h/(2R)$, so $t\downarrow0$. On this disk
\[
  2e^{-tR^2-R(s-c)}G_{R,t}(s)
         = F_t(s+h)+e^{-2R(s-c)}F_t(s-h)
           \longrightarrow F_0(s+h)
\]
locally uniformly, by Gaussian convergence on compacts and exponential decay of the second coefficient. Rouch\'e's theorem places the same positive number of zeros inside for every sufficiently large $R$. Hence $2R\lambda^+(g_M)\ge h$ eventually, by persistence of the all-real property at later times. Let $h$ increase to $b_f$ and combine with the upper bound. No maximal zero was selected.
\end{proof}

More generally, for every entire cuspidal standard completion the same formula recovers its right-half width from the product with $\cosh(R(s-1/2))$. The holomorphic weight-$k$ statement identifies every member with a classical cusp form by the two degeneracy maps. The compact convergence also recovers multiplicities: for any fixed $h>0$ and any closed disk in $\Re s>c$ with boundary avoiding the zeros of $F_0(s+h)$, the deformed oldform at physical time $h/\log M$ eventually has exactly the same number of zeros in that disk. In each such disk zeros therefore approach the original divisor translated left by $h$.

\subsection*{A single automorphic place in any number field}
There is an unconditional same-field realization of the symmetric multiplier in every rank $n\ge2$. Fix any one unramified finite place $v$ of a unitary cuspidal $\pi$ on $\mathrm{GL}_n$ over a number field $K$. Write $q=\mathrm N v$, $c=1/2$ and $F=\Lambda_\pi$. Let $g_j=\operatorname{diag}(\varpi^j,1,\ldots,1)$, where $\varpi$ is a uniformizer, and let $W^\circ$ be the unramified Whittaker function normalized at the identity, using an additive character of conductor zero. For $k\ge1$ form the actual local Whittaker vector
\[
 W_k(g)=\tfrac12 q^{k(n-2)/2}W^\circ(gg_k)
                +\tfrac12 q^{-k(n-2)/2}W^\circ(gg_{-k}).
\]
Let $F_k$ be the completed standard zeta product in which the $v$ factor is evaluated on $W_k$. The other finite standard factors and the archimedean completion take their usual values. All these vector data belong to the same cuspidal representation. The intrinsic primitive standard function of that representation is $F$.

\begin{corollary}\label{oldheat:localvector}
For $R=k\log q$ one has $F_k(s)=F(s)\cosh(R(s-c))$. In particular
\[
          \lim_{k\to\infty}2k\log q\,\lambda^+(F_k)=\Delta_\pi,
\qquad \Delta_\pi=\sup_{F(\rho)=0}|\Re\rho-c| .
\]
There is also an assertion entirely at time zero. For any nonnegative integers $j_k$ with $j_k/R\to a\ge0$,
\[
        2R^{-j_k}e^{-R(s-c)}F_k^{(j_k)}(s)
                     \longrightarrow F(s+a)\qquad(\Re s>c)
\]
locally uniformly, with the resulting equality of multiplicities in fixed disks whose boundaries avoid the limiting zeros. In particular, choose once and for all $j_k=\lfloor\sqrt k\rfloor$. ERH for $\pi$ is equivalent to every zero of $F_k^{(j_k)}$ lying on the critical line for infinitely many integers $k$. Under ERH it holds for every $k$.
\end{corollary}
\begin{proof}
Normalize multiplicative Haar measure to give the local units volume one. The unramified standard local integral is
\[
 I(s,W)=\int_{K_v^\times}W(\operatorname{diag}(y,1,\ldots,1))
                      |y|^{s-(n-1)/2}\,d^\times y.
\]
For $j<0$ the spherical Whittaker function vanishes on $g_j$. At the other diagonal elements its value is
\[
 W^\circ(g_j)=q^{-j(n-1)/2}h_j(\alpha_1,\ldots,\alpha_n)\qquad(j\ge0),
\]
where $h_j$ is the complete symmetric polynomial in the Satake parameters. Thus $I(s,W^\circ)=L_v(s,\pi)$ by summing a geometric series in each parameter. See \href{https://www.numdam.org/article/CM_1980__41_2_207_0.pdf}{Casselman and Shalika (1980)} or \href{https://doi.org/10.2969/jmsj/06610017}{Miyauchi (2014)}. Changing $y$ to $y\varpi^j$ shows that right translation by $g_j$ multiplies this integral by $q^{j(s-(n-1)/2)}$. Inserting the displayed $W_k$ gives the asserted cosine factor on the right, hence everywhere by continuation. Theorem~\ref{oldheat:main} and its proof for a general completion give the threshold limit.

For the derivative formula write $D=d/ds$. Expanding the two exponentials by Leibniz gives the exact normalized expression
\[
 (1+D/R)^{j_k}F(s)+(-1)^{j_k}e^{-2R(s-c)}(1-D/R)^{j_k}F(s).
\]
On a fixed compact set choose a Cauchy radius $B$ larger than $1+\sup_k(j_k/R)$. If $M_B$ bounds the modulus of $F$ on the enlarged compact set, the absolute $l$th term of either binomial is at most $M_B(j_k/(RB))^l$. Each coefficient of $F^{(l)}$ tends to $(\pm a)^l/l!$, so the two binomials converge uniformly to $F(s\pm a)$ by the Taylor series. The exponential multiplying the second tends to zero on the indicated open half-plane. The multiplicity assertion follows by Rouch\'e. If $F$ has ERH, then $F_k(c+iz)$ is real entire up to phase, of order at most one with only real zeros. Its Hadamard product is in the Laguerre--P\'olya class, whose derivatives have only real zeros. Conversely if $F$ has an off-line zero take its reflected representative on the right and a small disk around it. Here $a=0$ and the uniform limit gives a zero in the same disk for every sufficiently large $k$.
\end{proof}

\clearpage
\section{A heat-conductor criterion for the undeformed ERH}\label{hsc:chapter}
The bound in Theorem~\ref{hc:automorphic} is of order $1/\log\mathcal Q_*$. Could one replace the implied bounded constant by a quantity tending to zero? For products with a separated ordinate this improvement, even in a lower mean, contains the original ERH. We prove a quantitative lower bound from the horizontal displacement of any hypothetical zero.

Use unitary normalization. Let $\pi$ have an entire cuspidal standard completion, and fix the nonprincipal Dirichlet character $\chi$ modulo four. More generally the auxiliary $\nu$ below may be this character, of degree $m=1$, or a fixed holomorphic newform associated to a rational elliptic curve, of degree $m=2$. Use the entire product
\[
               P_{\tau,0}(s)=\Lambda_\pi(s)\Lambda_\nu(s+i\tau)
                      \qquad(\tau\ge1)
\]
and its physical heat deformation $P_{\tau,t}$ by the full real Gaussian in $s$. A constant phase makes the product real on $\Re s=1/2$. Use $\lambda(\tau)$ for its least \emph{nonnegative} time at which every zero lies on this line. For representations over $\Q$ this is the standard function of the isobaric sum $\pi\boxplus(\nu\otimes|\cdot|^{i\tau})$ with the norm-twist sign convention corresponding to $s+i\tau$. For different base fields the product is the displayed analytic product, or the direct-sum $L$-function on the product of the two groups.

\begin{theorem}\label{hsc:criterion}
If $\Lambda_\pi(\rho)=0$ for $\rho=\beta+i\gamma$ with $\beta>1/2$, then for each $0<a<(\beta-1/2)/m$ the set of $\tau$ for which
\[
                       \lambda(\tau)\log\tau >a
\]
has positive lower Lebesgue density. For each of the specified pairs, $\lambda(\tau)\log(3+\tau)$ is bounded unconditionally.

In particular, for any $E/\Q$, use $\pi=\nu=\pi_E$. Then ERH for $E$ is equivalent to
\[
          \liminf_{X\to\infty}\frac1X\int_1^X
                      \lambda(\tau)\log(3+\tau)\,d\tau=0 .
\]
For arbitrary entire cuspidal standard functions, ERH for all of them is equivalent to this condition for every $\pi$ with the same auxiliary $\nu=\chi$.
\end{theorem}

The fixed parameter $a$ has a geometric meaning. At physical heat time $a/\log\tau$, a zero near $\rho$ is detected near $\rho-ma$ in the unitary $s$-plane. The proof only uses the assertion while $\Re(\rho-ma)>1/2$.

\subsection*{A positive-density choice of ordinates}
Write $\Lambda_\nu(s)=G_\nu(s)L_\nu(s)$, including its conductor in $G_\nu$. Choose a closed disk $D$ about $\rho$ entirely inside $1/2<\Re s<1$. Its radius can be taken small enough that the translate of the disk by $-ma$ is still strictly to the right of $1/2$. The completed standard function is nonzero on the line $\Re s=1$ by its classical boundary theorem; this justifies the upper endpoint in choosing $D$.

By vertical universality the shifts satisfying
\[
                    \sup_{s\in D}|L_\nu(s+i\tau)-1|<1/4
\]
have positive lower density. For the character this is the Dirichlet universality theorem of Voronin, Bagchi and Gonek, as reviewed in \href{https://arxiv.org/abs/1407.4216}{K.~Matsumoto's survey}. For the elliptic form use \href{https://doi.org/10.1070/IM2003v067n01ABEH000419}{Laurin\v cikas, Matsumoto and Steuding, Izv. Math. \textbf{67} (2003), 77--90}; their arithmetic strip $k/2<\Re s<(k+1)/2$ becomes this same unitary strip. Cauchy's inequalities give uniformly bounded derivatives on every smaller concentric disk.

We also use a density-one estimate. For every fixed $M>2$ and $\epsilon>0$, except for a set of $\tau$ of density zero,
\begin{equation}\label{hsc:typical}
 \sup_{\substack{1/2\le\sigma\le M\\ |v|\le M}}
                  |L_\nu(\sigma+i\tau+iv)|\le\tau^\epsilon .
\end{equation}
Here is a proof to keep the height averaging separate from the statement about an individual $\pi$. For fixed degree one or holomorphic degree two,
\[
  \int_{X-M-2}^{2X+M+2}|L_\nu(\sigma+iy)|^2dy
                         \ll_{\nu,M,\eta}X^{1+\eta}
             \quad(1/2\le\sigma\le M+1)
\]
for every $\eta>0$, uniformly in $\sigma$. For $\sigma=1/2$ this is the classical central mean square: see \eqref{ang:means} for fixed degree two, and the ordinary Dirichlet-polynomial mean value theorem and approximate functional equation of length $O(X^{1/2+\eta})$ for the fixed character. For the remaining $\sigma$ interpolate the $L^2$ norms on $1/2$ and $M+2$ of $L_\nu(s)\exp((s-3iX/2)^2/X^2)$ by the three-lines theorem in $L^2$. On the left, the central mean square on dyadic intervals and the Gaussian give $O_{\nu,\eta}(X^{1+\eta})$ for the squared norm. On the right absolute convergence gives $O_{\nu,M}(X)$. The cutoff is bounded below on the ordinate interval in the display, uniformly in the bounded real parts.

By the functional equation and Stirling, extension to $\sigma\ge1/2-\delta$ costs at most $X^{2m\delta}$ in the squared estimate for each fixed small $\delta>0$. Apply the disk mean-value bound for $|L_\nu|^2$ with radius $\delta$ at every candidate for the supremum in \eqref{hsc:typical}. A large value forces the area integral on the surrounding vertical rectangle of fixed height to be $\gg_{\delta}X^{2\epsilon}$. Integration over $\tau\in[X,2X]$ and Fubini bound the exceptional measure by $O_{\nu,M,\delta,\eta}(X^{1+\eta+2m\delta-2\epsilon})$. Take $\eta+2m\delta<\epsilon$ and sum the dyadic estimates. Intersecting \eqref{hsc:typical} with the universality set leaves positive lower density. Only one choice of $M$ and $\epsilon$ will be required.

\subsection*{Two possible saddles}
Let $K$ be a small closed disk about $s_0=\rho-ma$ whose translate $K+ma$ lies strictly inside $D$. Choose its radius so that $\Lambda_\pi$ has no zeros on $\partial K+ma$. Fix $t=a/\log\tau$ and put $b_\tau=G_\nu'(s_0+i\tau)/G_\nu(s_0+i\tau)$. Stirling gives, with $u$ real,
\begin{gather*}
 \Re b_\tau=(m/2)\log\tau+O(1),\qquad \Im b_\tau=O(1),\\
 \frac{G_\nu(s+u+i\tau)}{G_\nu(s_0+i\tau)}
       =e^{(s+u-s_0)b_\tau}(1+O_U(\tau^{-1}))
                  \qquad(s\in K,\ |u|\le U).
\end{gather*}
The second bound follows by integrating the second logarithmic derivative, $O_U(\tau^{-1})$, twice on this fixed neighborhood. The modulus of the normalized Gaussian with the linear exponent $b_\tau u-t b_\tau^2$ is $e^{t(\Im b_\tau)^2}$ times the Gaussian of mean $2t\Re b_\tau=ma+O(1/\log\tau)$ and variance $2t$.

We claim that along the indicated positive-density set,
\begin{equation}\label{hsc:saddle}
 \frac{P_{\tau,t}(s)}
 {G_\nu(s_0+i\tau)e^{(s-s_0)b_\tau+t b_\tau^2}}
      =\Lambda_\pi(s+ma)L_\nu(s+ma+i\tau)+o(1)
              \quad\hbox{uniformly for $s\in K$.}
\end{equation}
Constant phases in the product are taken consistently on both sides. Choose a fixed $\eta_0>0$ so that $K+[ma-\eta_0,ma+\eta_0]$ lies inside $D$. On $|u-ma|<\eta_0$ the derivatives and values of the product $\Lambda_\pi(s+u)L_\nu(s+u+i\tau)$ are bounded uniformly. The Gaussian integral in this part equals its value at $u=ma$ plus $o(1)$, by the preceding mean and variance and the derivative bound. The entire linear Gaussian has complex integral exactly one after the indicated normalization; its omitted absolute tail tends to zero.

On a fixed $|u|\le U$ the auxiliary functional equation and the typical bound on the reflected half give
\[
       |L_\nu(s+u+i\tau)|\ll_U
             \tau^{\,m(1/2-\Re s-u)_++\epsilon},
             \qquad (x)_+=\max(x,0).
\]
Choose $M$ large enough for these translates and their reflections. The normalized remaining exponent to base $\tau$, apart from $\epsilon+o(1)$, is
\[
  -\frac{(u-ma)^2}{4a}+m(1/2-\Re s-u)_+
   =\begin{cases}
       -(u-ma)^2/(4a),&\Re s+u\ge1/2,\\
       -(u+ma)^2/(4a)-m(\Re s-1/2),&\Re s+u<1/2 .
     \end{cases}
\]
Away from $|u-ma|<\eta_0$, it is bounded above by a fixed negative number since $\min_{s\in K}\Re s>1/2$. Choose $\epsilon$ smaller than half its absolute value. The prefactor $1/\sqrt{4\pi t}$ and the compact integration interval do not affect the resulting decay.

For clarity the infinite tails can be bounded before choosing $M$. For all $s\in K$, all real $u$ and sufficiently large $\tau$,
\[
 \left|\frac{\Lambda_\pi(s+u)\Lambda_\nu(s+u+i\tau)}
               {G_\nu(s_0+i\tau)}\right|
 \le \exp\{C(1+|u|)(\log(3+\tau)+\log(2+|u|))\}.
\]
For bounded $u$ this follows from polynomial growth and Stirling on a fixed strip. Far to the right use absolutely convergent Euler products and integrate the gamma logarithmic derivatives along positive real translates. Their real parts are bounded by $C\log(3+\tau+|u|)$ on those rays; fixed imaginary translations at the starting point cost a bounded factor. Far to the left reflect both completed factors, which preserves their respective imaginary ordinates and reduces to the preceding bound. The fixed first factor also satisfies the entire order-one gamma bound on those rays. The further scalar normalization in \eqref{hsc:saddle} costs at most $\exp(C_a\log(3+\tau))$ on $K$. Choosing $U$ sufficiently large in terms of these constants and $a$, the Gaussian $e^{-u^2\log\tau/(4a)}$ dominates the whole estimate on $|u|>U$ by $O(e^{-u^2\log\tau/(8a)})$. This proves \eqref{hsc:saddle}.

\subsection*{Persistence of an off-line zero}
On $\partial K$ the modulus of the main term in \eqref{hsc:saddle} is at least three quarters of $\min_{\partial K}|\Lambda_\pi(s+ma)|>0$. The auxiliary factor is zero-free throughout $K+ma$ by the same universality event. Rouch\'e's theorem gives at least one zero of $P_{\tau,t}$ inside the off-line disk $K$, indeed exactly the multiplicity sum of the fixed zeros in $K+ma$. The all-real property, once attained, persists by the heat strip theorem. Hence $\lambda(\tau)>a/\log\tau$ on this positive-lower-density set for all sufficiently large $\tau$.

The inverse-log upper bound of Theorem~\ref{hc:automorphic} also applies to these products: the archimedean factors concatenate, local roots concatenate, the order-one entire products have the same reflection symmetry, and all initial zeros lie in the strip. These are exactly the hypotheses used in its reflected-half proof. For fixed factors of absolute degrees $d,m$,
\[
                 \mathcal Q_*(P_{\tau,0})\asymp_{\pi,\nu}
                            (3+\tau)^{\min(d,m)} .
\]
To see the lower bound, either the common shift has magnitude less than $\tau/2$, and the $m$ shifted auxiliary parameters have magnitude $\gg\tau$, or the $d$ fixed parameters do. Testing the common shift at $0$ and $-\tau$ gives the matching upper bound. Thus $\lambda(\tau)\log(3+\tau)$ is bounded.

The integral in the theorem is well-defined: for each fixed nonnegative $t$ the property that $P_{\tau,t}$ has all its zeros on the line is closed in $\tau$, by local uniform convergence and Hurwitz, and the functions are nonzero. Nonvanishing as an entire function also follows from Gaussian heat inversion for these order-one data, or from the large-height saddle of Section~\ref{hg:section} at fixed positive time. The threshold is therefore measurable. Under ERH for both factors the product already has all zeros on the line at time zero, so its nonnegative threshold vanishes identically. Conversely any off-line zero reflects, if needed, to a zero with $\beta>1/2$ and yields a strictly positive lower mean by the density bound. For the elliptic self-product the hypothesis on both factors is the single ERH for $E$. Using the fixed character for all $\pi$ gives the asserted general equivalence.

\subsection*{Recovery of the exact horizontal boundary with a Dirichlet auxiliary}
The positive-density lower bound has a matching all-shifts upper bound if the fixed auxiliary character satisfies its own ERH. The target may have off-line zeros. Define
\[
           \Delta_\pi=\sup_{\Lambda_\pi(\rho)=0}(\Re\rho-1/2)
                    \quad(0\le\Delta_\pi\le1/2).
\]

\begin{theorem}\label{hsc:exact}
Suppose the fixed primitive nonprincipal Dirichlet auxiliary $\chi$ satisfies ERH, and that the target $\pi$ has absolute degree $d\ge2$. For its product with the translate of this one character,
\[
             \limsup_{\tau\to\infty}\lambda(\tau)\log\tau=\Delta_\pi .
\]
More precisely, for every $a>\Delta_\pi$ all sufficiently large $\tau$ satisfy $\lambda(\tau)\log\tau\le a$, and for $0<a<\Delta_\pi$ the reverse strict inequality holds on a set of positive lower density.
\end{theorem}

\begin{proof}
The lower bound is Theorem~\ref{hsc:criterion}. For the upper bound fix $\Delta_\pi<D<a$, set $L=\log\tau$, $t=a/L$, and at height $T$ denote
\[
 p=L^{-1}\log\mathcal Q_\pi(T),\qquad q=L^{-1}\log\mathcal Q_\chi(T+\tau),
 \qquad W=a(p+q).
\]
At heights with $W\ge3$, Theorem~\ref{hc:leading} with $\theta=1/2$ already excludes off-line zeros, for all sufficiently large $\tau$. It applies to this product with its concatenated local and archimedean parameters. In the remaining argument $p,q$ belong to a fixed compact interval, and $p,q\ge0$. Since the two factors are fixed,
\begin{equation}\label{hsc:dichotomy}
          q\ge1-o(1)\quad\text{or}\quad p\ge d-o(1)\ge2-o(1).
\end{equation}
Indeed either $|T+\tau|\ge\tau/2$ or $|T|\ge\tau/2$. All $o(1)$ bounds in the proof are uniform over the remaining heights as $\tau\to\infty$.

We spell out the estimates for the finite Euler factors. If an entire fixed standard function has right zero boundary $\Theta$, the Euler logarithm, with its branch obtained from the right, is $o(\log(3+|y|))$ uniformly on fixed compact substrips of $\Re s>\Theta$ as $|y|\to\infty$. The same holds for any fixed number of its derivatives. For a proof take disks centered at $b+iy$ with $b>2$ fixed. On an outer disk contained in the zero-free half-plane the real part of the logarithm is $O(\log(3+|y|))$ by fixed-strip polynomial growth. Borel--Carath\'eodory gives this bound for its modulus on a slightly smaller disk, since the center has bounded Euler logarithm. On a fixed innermost disk in absolute convergence it is bounded. The three-circles theorem bounds its modulus on intermediate disks by $O((\log(3+|y|))^{1-\delta})$ with $\delta>0$ on each fixed margin. Choose $b$ and the radii to cover the desired compact substrips, using absolute convergence to the right; Cauchy's inequalities give the derivative assertion. Bounded heights contribute $O(1)$ on compact sets away from zeros and poles.

Consequently for $v\ge D$ the logarithms of both finite parts at $1/2+v+iT$ and $1/2+v+i(T+\tau)$, and their first two derivatives, are $o(L)$ on every bounded complex neighborhood with a fixed margin from the respective zero boundaries. On $0\le v\le D$, the two modulus estimates we need are
\begin{align}
 \log|L_\pi(1/2+v+iT)|&\le \tfrac12p(D-v)L+o(L),\label{hsc:lowtarget}\\
 \log|L_\chi(1/2+v+i(T+\tau))|&\le o(L).\label{hsc:lowaux}
\end{align}
For the first estimate use the preceding logarithm bound on $v=D$, the functional equation on $v=-D$, and the Phragm\'en--Lindel\"of convexity bound between. Concretely the two edge upper bounds on the finite part are $(3+|T|)^{o(1)}$ and $(3+|T|)^{dD+o(1)}$, giving $(3+|T|)^{d(D-v)/2+o(1)}$. Fixed bounded heights are harmless. This interpolation is the standard convexity argument for a fixed GL standard function (Section~\ref{hg:section}); one may first take an arbitrarily small fixed power on both edges and then let the power decrease. For the auxiliary, ERH gives the analogous zero-free logarithm bound for every positive $v$. Interpolate across $|v|\le\eta$ using the functional equation and let $\eta\downarrow0$, obtaining the Lindel\"of bound at and to the right of its critical line. Since $p,q$ are bounded, these fixed-function estimates give the stated uniform $o(L)$.

Let $F(v)$ denote the phased product completion at $1/2+v+iT$ before heat evolution, and let $C_T$ denote the modulus of the product of its gamma-conductor factors at $v=0$. Every shift of these two fixed functions has real part strictly greater than $-1/2$ by the local unitary bounds cited in Theorem~\ref{hc:automorphic}, so the gamma factors here are regular. On every fixed interval $0\le v\le U$ their combined modulus lies between constant multiples of
\[
                     C_T\exp((p+q)Lv/2).
\]
This follows by integrating the digamma estimates, with real derivative $(p+q)L/2+O_U(1)$ and imaginary derivative bounded in modulus on the real integration ray. Their second logarithmic derivatives are bounded there and on small complex neighborhoods. Work with the fixed-splitting-point integral $A_r$ from Theorem~\ref{hc:leading}. At $r=0$ the main squared-Gaussian exponent, in units of $L$ relative to $C_T$, is $W^2/(4a)$ at $v=W$. The upper exponent on $0\le v\le D$ from \eqref{hsc:lowtarget}--\eqref{hsc:lowaux} is
\[
                         -v^2/(4a)+qv/2+pD/2+o(1).
\]
There is a fixed gap between its maximum and the main exponent. If $aq\ge D$ its maximum without the $o(1)$ is at $D$ and the gap is $(W-D)^2/(4a)$. If $aq<D$, the maximum is at $aq$ and the gap is $p(a(p+2q)-2D)/4$. In this second case the first alternative of \eqref{hsc:dichotomy} fails for sufficiently large $\tau$, so $p\ge2-o(1)$. In the first case $W\ge a-o(1)>D$ by the same dichotomy. Both gaps are uniformly positive. This is the explicit use of the degree hypothesis.

For a small fixed $r_0>0$ this gap persists for $-r_0\le r\le1/2$, and $\mu=W+r$ stays a fixed distance to the right of $D$. Indeed the exponent gains $rv/(2a)$ on the low interval, whereas its high maximum becomes $\mu^2/(4a)$; for $r\ge0$ this only enlarges the gap since $W>D$. For negative $r$ use compactness and then decrease $r_0$. On $D\le v\le U$, \eqref{hsc:lowtarget} is replaced by the $o(L)$ logarithm, and the absolute exponent is at most $\mu^2/(4a)-(v-\mu)^2/(4a)+o(1)$. The integral for sufficiently large fixed $v>U$ is exponentially negligible on this scale: integrating the gamma logarithmic derivatives to the right and using bounded Euler products gives
\[
 |F(v)|/C_T
       \le \exp\{(p+q)Lv/2+C(1+v)\log(2+v)+C\}
          \quad(v\ge U).
\]
For example, on these rays $\Re\psi(z+x)\le\log(3+|\Im z|)+\log(2+x)+O(1)$, uniformly for the fixed initial gamma parameters and all relevant heights. The negative quadratic dominates this bound when $v>U$, since $p+q$ is bounded. These estimates leave an exponentially small absolute error, relative to $C_T\exp(L\mu^2/(4a))$, outside any fixed small real neighborhood of $\mu$. They also apply with an extra factor $v$.

The phase inside that neighborhood requires attention; a bound for $|\log L|$ by $o(L)$ alone does not give a relative approximation to the integral along the initial real segment. Write the integrand there as
\[
           C_T e^{L\mu^2/(4a)}\exp\{-(v-\mu)^2/(4t)+h(v)\}.
\]
The analytic $h$ combines the finite Euler logarithms and the gamma logarithms after subtracting their real linear slope $(p+q)L/2$. Choose the constant arguments modulo $2\pi$ at the real center. On a fixed smaller complex disk about $\mu$ one has $h,h',h''=o(L)$ uniformly, by the zero-free margin and the earlier digamma bounds. The map $v\mapsto\mu+2t h'(v)$ is a strict contraction of a fixed closed disk of smaller radius into itself for large $\tau$. Its fixed point $z=\mu+o(1)$ is the complex stationary point. Deform the short real segment to the parallel horizontal segment through $z$, with endpoints matched by short linear segments. On these connectors the Gaussian loses a fixed multiple of $L$, whereas $h=o(L)$. For real $x$ on the new line, Taylor's estimate gives, after factoring out the value at $z$,
\[
        -(z+x-\mu)^2/(4t)+h(z+x)
        +(z-\mu)^2/(4t)-h(z)=-x^2/(4t)+o(L)x^2 .
\]
The linear term vanishes by the stationary equation. On rescaling by $\sqrt t$, dominated convergence, uniform because the coefficient of $x^2$ in the error is uniformly $o(L)$, gives an integral $\sqrt{4\pi t}(1+o(1))$ times the factored value. The same reasoning with the factor $v$ gives $z+o(1)$ times this integral. The modulus of the factored value relative to $C_T e^{L\mu^2/(4a)}$ is $\exp(o(L))$, since $|z-\mu|=o(1)$ and $h(z)=o(L)$. The exponential tail gaps therefore justify the relative approximation for the full integrals.

It follows that $2t\Re(A_r'/A_r)=\mu+o(1)>0$ for $|r|\le r_0$. For $r_0\le r\le1/2$, the absolute integral of $A_{-r}$ on $v\ge D$ is at most $e^{-r_0 D/t}$ times the absolute integral of $A_r$ there, and its low part is bounded by the exponentially negligible low part at positive $r$. The preceding estimates bound the total absolute integral of $A_r$ by $C_T\exp(L\mu^2/(4a)+o(L))$, and give the same expression for its actual modulus. Hence $|A_{-r}|<|A_r|$ for these $r$, and the derivative gives it for $0<r\le r_0$ as well. The reflected halves cannot cancel off the critical line at any height. This proves the upper bound and the theorem.
\end{proof}

\clearpage
\section{An ordinary limit for the whole zero strip of a translated product}\label{hol:chapter}
Fix an entire cuspidal standard completion $F$ of absolute degree $d$, in the unitary normalization $c=1/2$. Let $V$ be a second fixed entire cuspidal standard completion of absolute degree $m$, all of whose zeros lie on the line $\Re s=c$. For the application to every rational elliptic curve one may take $V=\Lambda(\,\cdot\,,\chi_{-4})$. Only the Dirichlet input is required for this choice. Set
\[
 P_{\tau,0}(s)=F(s)V(s+i\tau),\qquad
 \Delta=\sup_{F(\rho)=0}(\Re\rho-c),\qquad L=\log\tau.
\]
Write $P_{\tau,t}$ for its physical Gaussian deformation, $V_t$ for that of $V$, $b_\tau(t)$ for the supremum of $|\Re\rho-c|$ over its deformed zeros, and $\lambda_\tau^+$ for the least nonnegative time with $b_\tau(t)=0$. The phase of the product is chosen as the product of the two constituent phases. The factors are fixed throughout the limit in $\tau$.

\begin{theorem}\label{hol:profile}
For every fixed $a>0$ there is local uniform convergence on $\Re s>c$,
\[
          \frac{P_{\tau,a/L}(s)}{V_{a/L}(s+i\tau)}
                           \longrightarrow F(s+ma).
\]
It holds for every sequence $\tau\to\infty$ and preserves the multiplicity in every closed disk in this half-plane whose boundary avoids the limiting divisor. If $d\ge2m$, then for all $a\ge0$,
\begin{equation}\label{hol:width}
        \lim_{\tau\to\infty} b_\tau(a/L)=(\Delta-ma)_+,
    \qquad
        \lim_{\tau\to\infty} L\lambda_\tau^+=\Delta/m .
\end{equation}
For $a>\Delta/m$ the width is exactly zero at all sufficiently large $\tau$. In particular elliptic ERH is equivalent to $\lambda^+_\tau=o(1/\log\tau)$ for the product of that curve's unitary completed function with translates of the single fixed character $\chi_{-4}$. Any one sequence of ordinates tending to infinity, including the positive integers, suffices for this test. For two constituents over $\Q$ these are standard functions of actual isobaric representations. Their thresholds measure the original divisor via the stated limit.
\end{theorem}

\begin{proof}
Fix a compact disk $K$ in $\Re s>c$. Write $V=G L_V$ with $G$ its archimedean and conductor factor. For an ERH auxiliary, the logarithm bound proved in Theorem~\ref{hsc:exact}, applied to this fixed function, bounds $\log L_V$ and its first two derivatives by $O_K(L^{1-\delta})$ on a fixed neighborhood of $K+ma+i\tau$, for some $\delta>0$. Its proof uses only absolute convergence on a right subdisk, polynomial growth, three circles and the zero-free right half-plane. Interpolating across the critical line with the functional equation gives the Lindel\"of bound uniformly on each fixed closed right strip, as in \eqref{hsc:lowaux}. These two estimates concern every ordinate. Use
\[
 t=a/L,\qquad B(s)=G'(s+i\tau)/G(s+i\tau),\qquad \mu(s)=2tB(s).
\]
Stirling gives $B(s)=mL/2+O_K(1)$, hence $\mu(s)=ma+O_K(L^{-1})$. On a small disk $|u-ma|<3e$ choose the analytic logarithm
\[
 h_s(u)=\log L_V(s+u+i\tau)
     +\log\frac{G(s+u+i\tau)}{G(s+i\tau)}-u B(s).
\]
For the gamma logarithm use the branch obtained by horizontal integration from zero and its analytic continuation. Its remainder and its first two derivatives are $O_{K,a}(1)$ by Stirling. Uniformly for $s\in K$ we have $h_s,h_s',h_s''=O(L^{1-\delta})$, decreasing $\delta$ if necessary. After dividing by $G(s+i\tau)e^{tB(s)^2}$ the local Gaussian integrand, with the usual factor $1/\sqrt{4\pi t}$, is
\[
           e^{-(u-\mu(s))^2/(4t)+h_s(u)}F(s+u).
\]
The expression without $F$ gives the auxiliary-only heat integral with exactly the same normalization.

Here are uniform tail bounds for both integrals. On $|u|\le U$ the gamma ratio has logarithmic modulus $mLu/2+O_{K,U}(1)$ before completing the square. The fixed factor $F(s+u)$ is bounded there. By the functional equation and the Lindel\"of estimate just proved, the normalized absolute exponent, divided by $L$, is at most
\[
 -\frac{(u-ma)^2}{4a}+m(c-\Re s-u)_++o(1).
\]
On $\Re s+u\ge c$ this is bounded away from zero on the negative side outside $|u-ma|<e$. On its complement the two leading terms equal $-(u+ma)^2/(4a)-m(\Re s-c)$, which is also uniformly negative. For $|u|>U$ use the bound $\exp(C(1+|u|)(L+\log(2+|u|)))$ before the Gaussian, proved by gamma integration and reflection in Theorem~\ref{hsc:criterion}. An adequately large fixed $U$ makes the Gaussian dominate it there. Thus the absolute error after removing the interval $[ma-e,ma+e]$ is $O(e^{-c_KL})$ with some $c_K>0$, for both integrals.

The map $u\mapsto\mu(s)+2t h_s'(u)$ is a strict contraction into a closed disk about $ma$ of radius $e/2$ for all sufficiently large $\tau$. Its unique fixed point $z_s$ satisfies $z_s=ma+O(L^{-\delta})$. Deform the retained segment to $z_s+[-e,e]$ and connect its ends to the original ends. The short connectors still have an exponentially small normalized integral, since their real displacements from $ma$ stay at least $e/2$ and their imaginary parts are $o(1)$. Divide out the common nonzero saddle value
\[
                \exp\bigl(-(z_s-\mu(s))^2/(4t)+h_s(z_s)\bigr).
\]
Its modulus is $\exp(o(L))$, so the tails remain negligible in relative size. At $u=z_s+x$ for real $x$ the remaining exponent is $-x^2/(4t)+O(L^{1-\delta}x^2)$, by cancellation of the linear term. After $x=\sqrt{t}\,y$ it is bounded in real part by $-y^2/8$ for large $L$. Dominated convergence gives $1+o(1)$ for the auxiliary integral after division by both normalizations. The fixed amplitude in the other integral differs from $F(s+ma)$ by $O_K(|x|+L^{-\delta})$. Its Gaussian average, after division by the auxiliary integral, differs by $O_K(L^{-1/2}+L^{-\delta})$. This proves the quotient limit even at zeros. Also $V_t$ has no zeros off the line by de Bruijn. Rouch\'e's theorem now gives the asserted divisor convergence. At a zero of multiplicity $j$ on a fixed disk containing no other zero, the displaced zeros lie within $O_K(L^{-\min(1/2,\delta)/j})$ of that limiting zero, by its local factorization.

For the all-height assertion fix $a>0$ and $D>\Delta$, and consider $r=\Re s-c$ with
\[
               (D-ma)_++\eta\le r\le1/2
\]
for a fixed $\eta>0$. The case of an empty range is immediate. Use the notation $p=L^{-1}\log\mathcal Q_F(T)$, $q=L^{-1}\log\mathcal Q_V(T+\tau)$ and $W=a(p+q)$ from Theorem~\ref{hsc:exact}. Large $W\ge3$ is covered by Theorem~\ref{hc:leading}. For the remaining compact conductor range, at least one of $q\ge m-o(1)$ or $p\ge d-o(1)$ holds. The logarithm bounds on $v\ge D$ and the low-interval estimates \eqref{hsc:lowtarget}--\eqref{hsc:lowaux} hold unchanged for these degrees, using ERH only for $V$. For the positive-half integral $A_r$, put $w=r+W$. Its high saddle exponent is $w^2/(4a)$ in units of $L$. The upper exponent on $0\le v\le D$ is
\[
             -v^2/(4a)+(r+aq)v/(2a)+pD/2+o(1).
\]
If $r+aq\ge D$, its gap below the high maximum is $(w-D)^2/(4a)+o(1)$. Here $w-D$ is bounded below by a positive constant since $r+ma\ge D+\eta$ and either $q\ge m-o(1)$ or $p\ge d-o(1)$. If $r+aq<D$, then $q\ge m-o(1)$ is impossible, so $p\ge d-o(1)\ge2m-o(1)$. The gap is now
\[
              p\{2r+a(p+2q)-2D\}/4+o(1),
\]
which has a uniform positive lower bound for the same reason. The high saddle stays a fixed distance to the right of $D$. The complex stationary-point calculation in Theorem~\ref{hsc:exact} therefore applies verbatim to $A_r$ at all these heights. It gives its modulus as the gamma-conductor prefactor times $\exp(Lw^2/(4a)+o(L))$. In the reflected half $A_{-r}$ the absolute high integral is smaller than the positive absolute high integral by a factor at most $e^{-rD/t}$; its absolute low integral is bounded by the exponentially negligible positive low integral. Since $r\ge\eta$, the two halves cannot cancel. Reflection gives the symmetric exclusion on the left. Sending $D$ down to $\Delta$ and $\eta$ down to zero proves the upper limit in \eqref{hol:width}. The quotient limit at a disk about any zero with $\Re\rho>c+ma$ gives the matching lower limit, including an unattained supremum by approximation. At $a=0$ the divisor is exactly the union of the two initial divisors.

If $a>\Delta/m$, choose $D$ strictly between $\Delta$ and $ma$. The same two gap bounds hold already at $r=0$ and hence uniformly also for $-r_0\le r\le1/2$ with some fixed $r_0>0$. The differentiated high-saddle estimate of Theorem~\ref{hsc:exact} gives $2t\Re(A'_r/A_r)=r+W+o(1)>0$ on $|r|\le r_0$. It excludes a zero at every $0<r\le r_0$. This proves exact purity at these times. If $a<\Delta/m$ the positive limiting width forbids purity for all sufficiently large shifts. Persistence of purity at later times and then $a\to\Delta/m$ give the threshold formula, including $\Delta=0$.
\end{proof}

\begin{corollary}\label{hol:abelianfield}
Let $K/\Q$ be any finite abelian extension of degree $m$, and let $\pi$ be any unitary cuspidal representation of $\mathrm{GL}_n(\mathbb A_K)$, $n\ge2$. There is a single finite-order Hecke character $\eta$ of $K$, independent of $n$ and $\pi$, such that all the comparisons in Theorem~\ref{hol:profile} are standard functions of isobaric representations over the very same field $K$. For every such $\pi$,
\[
   m\lim_{\tau\to\infty}(\log\tau)\lambda^+
             (\pi\boxplus(\eta\otimes|\cdot|^{i\tau}))=\Delta_\pi.
\]
At time $a/(m\log\tau)$ the limiting width of this product is $(\Delta_\pi-a)_+$. In particular this statement applies to the standard function of a modular elliptic curve over $K$, with a shift by $1/2$ for its arithmetic variable.
\end{corollary}
\begin{proof}
Choose a primitive Dirichlet character $\chi$ ramified at a new prime where $K$ is unramified. Let $\eta$ be its pullback by the idele norm $K\to\Q$. By induction of the trivial representation and the projection formula, the full completed Hecke function of $\eta$ factors as the product of the completed Dirichlet functions for $\chi\psi$, where $\psi$ runs over the characters of $\operatorname{Gal}(K/\Q)$. Take each product character at its primitive conductor. This is Artin formalism, including inertia invariants and the archimedean and conductor factors. The choice of new prime makes every $\chi\psi$ nonprincipal. Hence the completion is entire and the Dirichlet theorem puts all of its zeros on the line. Its absolute degree is $m$ while that of the target is $nm\ge2m$. Apply Theorem~\ref{hol:profile} to the two constituents over $K$.
\end{proof}

\clearpage
\section{Two unknown zero boundaries under a separated heat shift}\label{htb:section}

Let $\pi,\nu$ be fixed entire cuspidal standard completions of absolute degrees $d,m$, in the unitary normalization and with the local bounds of Section~\ref{hc:chapter}. They may belong to different number fields. Use the phased analytic product $P_{\tau,0}(s)=\Lambda_\pi(s)\Lambda_\nu(s+i\tau)$ and its physical heat deformation from Section~\ref{hsc:chapter}. Let $B_\pi,B_\nu$ be their initial rightmost zero displacements from $1/2$, and $B(\tau,a)$ the supremum of those displacements for $P_{\tau,a/\log\tau}$. Reflection makes each of these the half-width of the whole divisor. Put
\[
     c_\pi=\min(m,d/2),\qquad c_\nu=\min(d,m/2).
\]

\begin{theorem}\label{htb:upper}
For each $a>0$ the following bounds hold without an ERH assumption on either function:
\begin{align*}
 \limsup_{\tau\to\infty}B(\tau,a)
   &\le \max(0,B_\pi-ac_\pi, B_\nu-ac_\nu),\\
 \limsup_{\tau\to\infty}\lambda(\tau)\log\tau
   &\le \max(B_\pi/c_\pi,B_\nu/c_\nu).
\end{align*}
If $d\ge2m$ and $B_\pi\ge2B_\nu$, both estimates have exact ordinary limits:
\[
   \lim_{\tau\to\infty} B(\tau,a)=(B_\pi-ma)_+,
       \qquad
   \lim_{\tau\to\infty}\lambda(\tau)\log\tau=B_\pi/m .
\]
The same conclusion with the roles reversed holds when $m\ge2d$ and $B_\nu\ge2B_\pi$. In particular a fixed ERH auxiliary of degree $m$ yields both exact limits for every target of absolute degree at least $2m$.
\end{theorem}

\begin{proof}
Write $L=\log\tau$, $t=a/L$, and use the normalized conductor logarithms $p=\log\mathcal Q_\pi(T)/L$, $q=\log\mathcal Q_\nu(T+\tau)/L$ at an arbitrary height $T$. When $a(p+q)\ge3$ the pointwise bound of Theorem~\ref{hc:leading} excludes all off-line zeros uniformly as $\tau\to\infty$. Otherwise $p,q$ lie in a fixed compact set. At every subsequential limit one has either $p\ge d$ or $q\ge m$, since either $|T|\ge\tau/2$ or $|T+\tau|\ge\tau/2$.

Choose fixed $D_\pi>B_\pi$, $D_\nu>B_\nu$ and set $D=\max(D_\pi,D_\nu)$. The right-hand zero-free logarithm bound and strip interpolation in the proof of Theorem~\ref{hsc:exact}, now applied separately to each function, give on $v\ge0$ in bounded intervals
\[
 \log|L_j(1/2+v+iT_j)|
   \le \tfrac L2 p_j(D_j-v)_++o(L),
 \quad (p_\pi,p_\nu)=(p,q),\quad(T_\pi,T_\nu)=(T,T+\tau).
\]
On fixed neighborhoods to the right of $D$ the logarithms and their first two derivatives are $o(L)$. The estimates are uniform, including bounded values of either $T_j$.

Let $r$ be a real displacement, $P=p+q$ and $\mu=r+aP$. In the positive half-integral $A_r$ of Theorem~\ref{hc:leading}, subtract the high saddle exponent $L\mu^2/(4a)$ after division by the central gamma modulus. On $0\le v\le D$ the remaining upper exponent in units of $L$ is
\[
     -\frac{(v-\mu)^2}{4a}
       +\frac12\sum_{j\in\{\pi,\nu\}}p_j(D_j-v)_++o(1).
\]
The positive-part sum is the maximum over subsets $I\subset\{\pi,\nu\}$ of $\sum_{j\in I}p_j(D_j-v)$. For $P_I=\sum_{j\in I}p_j$, completion of the square bounds the corresponding exponent, even on the entire real $v$-axis, by
\[
       \frac12\sum_{j\in I}p_jD_j
              -\frac{P_I}{4}\big(2r+a(2P-P_I)\big).
\]

Take $r\ge0$ above the asserted width bound by a fixed positive margin. Take the $D_j$ close enough to $B_j$ that $D_j\le r+ac_j-\epsilon$ for both $j$ with a fixed $\epsilon>0$. At any limiting pair $p,q$ with the stated alternative, $p+2q\ge2c_\pi$ and $2p+q\ge2c_\nu$. These inequalities give strict negativity for each singleton of positive weight. For the set of both indices use
\[
                    (p+q)^2\ge 2(pc_\pi+qc_\nu).
\]
If $q\ge m$ then $2pq\ge2pm$ and $q^2\ge qm$. If $p\ge d$ then $p^2\ge pd$ and $2pq\ge2qd$. The definitions of $c_\pi,c_\nu$ prove the display in both cases. For a subset of zero limiting weight the original exponent on $[0,D]$ is $-(v-\mu)^2/(4a)$, strictly negative there because $P\ge\min(d,m)\ge\max(c_\pi,c_\nu)$ and $\mu\ge D+\epsilon$. Thus strict negativity holds at every point of the compact limiting parameter set, including $p=0$ or $q=0$. Continuity and subsequence compactness give a uniform negative gap for all large $\tau$.

The high interval $v\ge D$ has the same complex-saddle estimate as in Theorem~\ref{hsc:exact}: at $\mu$ the residual logarithm and its first two derivatives are uniformly $o(L)$ on a fixed disk. The shifted stationary point is $\mu+o(1)$, giving $|A_r|=C_T\exp(L\mu^2/(4a)+o(L))$ with $C_T$ the central gamma modulus. The corresponding absolute integral obeys the same upper bound, and the low part loses a fixed exponential factor. At positive $r$ bounded away from zero the high absolute integrand of $A_{-r}$ is smaller by at least $e^{-rD/t}$; its low absolute mass is bounded by that of $A_r$. Thus the reflected terms have unequal moduli. This gives the asserted width bound since the initial strip and its heat contraction already exclude $r>1/2$.

If $a>\max(B_\pi/c_\pi,B_\nu/c_\nu)$, both strict inequalities for $D_j$ hold even for all $r$ in a neighborhood of zero, of either sign. The same saddle asymptotic with the extra factor $v$ gives $2t\Re(A'_r/A_r)=\mu+o(1)>0$ there. Strict monotonicity of $|A_r|$ gives the reflected comparison for every $r>0$ in that neighborhood. This proves the upper time bound, without inferring exact reality merely from a vanishing limiting width.

For the lower bounds we state precisely how the unknown auxiliary width enters. At each compact subset of
\[
              \Re s-1/2>\max(0,B_\nu-ma/2)
\]
one has the local uniform limit
\begin{equation}\label{htb:quotient}
   \frac{P_{\tau,t}(s)}
        {(\Lambda_\nu)_t(s+i\tau)}\longrightarrow \Lambda_\pi(s+ma),
                      \qquad t=a/\log\tau .
\end{equation}
Use consistent constant phases. The denominator is nonzero on the chosen compact for large $\tau$ by the shared-saddle proof below; no ERH for $\nu$ is used.

Write $r=\Re s-1/2$ and choose $D_\nu>B_\nu$ still satisfying $r>D_\nu-ma/2$ there. Normalize the numerator and denominator Gaussian integrals by the auxiliary gamma factor and its linear tangent, as in \eqref{hsc:saddle}. On bounded real $u$ the normalized Gaussian exponent is $-(u-ma)^2/(4a)+o(1)$ in units of $L$. Put $w=r+u$. The bound for $L_\nu$ adds respectively $0$, $m(D_\nu-w)/2$, or $-mw$ on $w\ge D_\nu$, $|w|\le D_\nu$, or $w\le-D_\nu$, up to $o(1)$. This follows from the zero-free right logarithm, interpolation across the two edges, and the functional equation on the left. In the middle region the sum of the two exponents is
\[
           -u^2/(4a)-m^2a/4+m(D_\nu-r)/2<0.
\]
On the far left it is $-(u+ma)^2/(4a)-mr<0$. Both inequalities have a fixed margin on the compact. The only remaining saddle is $u=ma$, where $r+ma>D_\nu$. The fixed target and its derivatives are bounded on compact translates, and at unbounded $u$ the Gaussian dominates the gamma and Euler growth as in the proof of \eqref{hsc:saddle}.

On a fixed disk about that saddle take the logarithm of the normalized auxiliary alone. Its residual logarithm after subtraction of the gamma slope and its first two derivatives are $o(L)$. The complex stationary point is $ma+o(1)$ and is the same for the two integrals. Move to the horizontal segment through it. Its nonzero common leading factor has modulus $\exp(o(L))$ in the chosen normalization. Rescaling the Gaussian by $\sqrt t$ leaves $1+o(1)$ in the denominator. In the numerator it leaves $\Lambda_\pi(s+ma)+o(1)$ by the uniform continuity of the target in a larger compact neighborhood. This is an absolute error after division by the common factor, so it remains valid at zeros of the target. Negative exponential gaps dispose of the connectors and the rest of the integral, proving \eqref{htb:quotient}.

Finally assume $d\ge2m$ and $B_\pi\ge2B_\nu$, so $c_\pi=m$ and $c_\nu=m/2$. The upper width simplifies to $(B_\pi-ma)_+$. If $ma<B_\pi$, take original target zeros with displacement arbitrarily close to $B_\pi$. Their translates by $-ma$ satisfy $r>0$ and $r>B_\nu-ma/2$ once sufficiently close, because $B_\pi>B_\nu+ma/2$. Rouché applied to \eqref{htb:quotient} on small disks gives the same multiplicities to the right of the line at all sufficiently large shifts. This gives the matching lower width and excludes an earlier all-real time by persistence. If $ma\ge B_\pi$ the lower width zero is automatic. The symmetric assertion follows by translating the product's ordinate and allowing $\tau\to-\infty$ in the same bounds.
\end{proof}

\clearpage
\section{Elliptic self-shifts and zeros of high ordinary derivatives}\label{hdp:section}
For a rational elliptic curve $E$, use its unitary completion $F=\Lambda_E$.
Let $\Delta=\sup_{F(\rho)=0}(\Re\rho-1/2)$. Denote the nonnegative
threshold for $F(s)F(s+i\tau)$ from Section~\ref{hsc:chapter} by
$\lambda_E(\tau)$. Lebesgue density refers to positive real shifts.

\begin{theorem}\label{hdp:main}
Put $C(b)=\max\{b/2,b(1-2b)\}$. For every $\epsilon>0$, some $c>0$
satisfies, on $[X,2X]$ outside a set of measure $O_{E,\epsilon}(X^{1-c})$,
\[
 \frac{\Delta}{2}-\epsilon\ \le       \lambda_E(\tau)\log\tau\ \le\ C(\Delta)+\epsilon .
\]
In particular, if $\Delta\ge1/4$, the scaled threshold converges in
density, and in every finite positive mean power of its absolute difference,
to $\Delta/2$. If $\Delta=0$ the threshold vanishes identically.
For every $E$, its scaled threshold is at most $1/4+\epsilon$
outside the stated exceptional set. These assertions require no ERH
assumption on an auxiliary function.
\end{theorem}

\subsection*{A local quotient and the lower bound}
\begin{theorem}[Recovery outside a power-saving exceptional set]\label{htb:density}
Take the fixed degree-one Dirichlet auxiliary or a fixed degree-two rational elliptic auxiliary from Theorem~\ref{hsc:criterion}. Fix $a>0$ and a compact $K\Subset\{\Re s>1/2\}$. Without an ERH assumption on either function, for some $c,e>0$ and outside a set of shifts of measure $O(X^{1-c})$ in $[X,2X]$, one has uniformly for $s\in K$
\[
      \frac{P_{\tau,a/\log\tau}(s)}
           {(\Lambda_\nu)_{a/\log\tau}(s+i\tau)}
        = \Lambda_\pi(s+ma)+O((\log\tau)^{-e}).
\]
The denominator is nonzero there. For a fixed target zero $\rho$ with $\Re\rho-ma>1/2$ and multiplicity $k$, the heated product has $k$ zeros within distance $O((\log\tau)^{-e/k})$ of $\rho-ma$ on these shifts, after choosing $K$ to include a fixed neighborhood. The constants can depend on the fixed compact, functions, and $a$.

Consequently, for every $a<B_\pi/m$ with $a>0$, the shifts satisfying $\lambda(\tau)\log\tau>a$ have natural Lebesgue density one. In particular
\[
 \liminf_{X\to\infty}\frac1X\int_1^X
              \lambda(\tau)\log(3+\tau)\,d\tau\ge B_\pi/m.
\]
For a rational elliptic self-product with width $B_E$, the lower mean is at least $B_E/2$ and the upper limit of the individual scaled times is at most $B_E$.
\end{theorem}
\begin{proof}
The saddle neighborhoods of $K+ma$ lie a fixed positive distance inside $\Re s>1/2$. Cover them by finitely many disks in that half-plane to which the right-hand Borel--Carathéodory and three-circles argument in Theorem~\ref{hsc:exact} applies, using centers sufficiently far into absolute convergence. More explicitly, one may place a right-hand center at the ordinate of each of finitely many points and move it far enough right that its disk from the center to the required neighborhood stays to the right of $1/2+\delta$ for a fixed $\delta>0$. Choose a slightly larger disk with this property and a smaller concentric disk entirely in absolute convergence.

Except for $O(X^{1-c_1})$ shifts in $[X,2X]$, all the translated outer disks are zero-free for $L_\nu$. Indeed each zero of real part at least $1/2+\delta$ excludes a set of shifts of bounded length and only heights $O(X)$ are involved. For the fixed Dirichlet character use Montgomery's classical zero-density bound, and for the holomorphic form use \href{https://arxiv.org/abs/2103.01956}{Beckwith, Liu, Thorner and Zaharescu, Theorem 1.1}. Both bound the relevant number by $O(X^{1-c_1})$ for some $c_1>0$. Finite-part trivial zeros do not enter these disks. On each remaining disk take the logarithm agreeing with the bounded Euler logarithm at the right-hand center. Standard polynomial growth bounds its real part above by $O(\log X)$. Borel--Carathéodory shrinks the disk once to bound the entire logarithm by $O(\log X)$ there. On the small concentric Euler disk it is $O(1)$. Three circles and another fixed shrinking for Cauchy derivatives give bounds $O((\log X)^{1-e_0})$ for the logarithm and its first two derivatives on the required saddle neighborhoods with a fixed $e_0>0$.

Away from the saddle use \eqref{hsc:typical}. Its proof gives a further exceptional measure $O(X^{1-c_2})$ for any one required fixed small exponent and sufficiently large fixed window. The Dirichlet character here and the rational elliptic auxiliary are self-dual, so the corresponding reflected estimate is available at negative ordinates as well. For $r=\Re s-1/2>0$ the upper Gaussian exponents on the two sides are now $-(u-ma)^2/(4a)$ and $-(u+ma)^2/(4a)-mr$, with the arbitrarily small fixed error from \eqref{hsc:typical}. On the fixed compact choose that error smaller than the positive off-saddle gaps. This gives exponentially small omitted tails as in \eqref{hsc:saddle}, uniformly for all the remaining shifts.

The common auxiliary stationary point then satisfies $z=ma+O((\log\tau)^{-e_0})$, from its defining contraction and the gamma slope error $O(1/\log\tau)$. On the horizontal segment through $z$ the fixed amplitude differs from $\Lambda_\pi(s+ma)$ by $O((\log\tau)^{-e_0}+|u-z|)$. The normalized absolute Gaussian integral of $|u-z|$ is $O((\log\tau)^{-1/2})$; the residual second derivative is $o(\log\tau)$, so the same bound holds with its local exponential weight included. The common denominator has magnitude at least half its Gaussian main amplitude for large $\tau$. Subtracting the target value times the denominator before division gives the claimed uniform error with $e=\min(e_0,1/2)$.

At an isolated zero of order $k$ the modulus of $\Lambda_\pi(s+ma)$ on a small centered circle of radius $r'$ is at least a fixed multiple of $(r')^k$. Take $r'$ a sufficiently large constant times $(\log\tau)^{-e/k}$ and apply Rouché. Given $a<B_\pi/m$, choose a zero with $\Re\rho-1/2>ma$. Its shifted neighborhood stays off the line, forcing $\lambda(\tau)\log\tau>a$ except on the indicated power-saving set. Summing the dyadic exception bounds gives natural density one. The nonnegative lower mean follows first for each such $a$, and then by taking its supremum. The self-product upper bound is Theorem~\ref{htb:upper} with $d=m=2$.
\end{proof}

\subsection*{A simultaneous bound for the two finite factors}
Write $f=L_E$ for the finite function. For a fixed $0\le k\le1$,
$R=X^k$, and fixed $\eta>0$, one has
\begin{equation}\label{hdp:sup}
 \sup_{|u|\le 3R+2(\log X)^2}
  |f(1/2+iu) f(1/2+i(u\pm\tau))|\le X^{k/2+\eta}
\end{equation}
for all $\tau\in[X,2X]$ outside a set of measure $O(X^{1-c_\eta})$.
Both signs and any fixed finite list of $k$ can be imposed together.

Indeed put $S=3R+2(\log X)^2+2$.
For $|\sigma-1/2|\le\delta$ and every fixed $\xi>0$, the
central mean square \eqref{ang:means}, interpolation to the right
and the functional equation on the left, as in \eqref{hsc:typical},
give
\[
 \int_{|u|\le S}|f(\sigma+iu)|^2du\ll S^{1+\xi+4\delta},
 \qquad
 \sup_{|u|\le S}\int_X^{2X}|f(\sigma+i(u\pm\tau))|^2d\tau
          \ll X^{1+\xi+4\delta}.
\]
The second interval is contained in $[-7X,7X]$ for large $X$.
Multiply the integrands before applying Fubini: the integrated
squared product over $\tau,u,\sigma$ is $\ll X^{1+k+3\xi+8\delta}$
for large $X$, absorbing the fixed logarithms by the extra $\xi$.
The disk submean inequality for the holomorphic product at each
possible central ordinate bounds its squared supremum by
$O_\delta(1)$ times the area integral over the enlarged rectangle
for that $\tau$. Markov's inequality, with $3\xi+8\delta<\eta$,
proves \eqref{hdp:sup}.

Fix $a>C(\Delta)$ and choose $D>\Delta$ so close to $\Delta$
that $a>C(D)$. Write $L=\log\tau$, $t=a/L$, and at physical
height $T$ put
\[
 M=\frac{\log\mathcal Q_E(T)+\log\mathcal Q_E(T+\tau)}{2L},
          \qquad W=2aM .
\]
Theorem~\ref{hc:leading} disposes of $W\ge3$. In the remaining
range $M$ is bounded and both physical ordinates are polynomially
bounded in $X$. The single-function estimate \eqref{hsc:lowtarget}
gives globally, uniformly on $0\le v\le D$,
\begin{equation}\label{hdp:convex}
 \log |f(1/2+v+iT) f(1/2+v+i(T+\tau))|
        \le M(D-v)L+o(L).
\end{equation}

Suppose first $b=\min(|T|,|T+\tau|)\le2X$. Define
$k=\min(1,\log(3+b)/L)$; the larger absolute ordinate is
between $X/2$ and $4X$. Since $\mathcal Q_E(y)\asymp_E(3+|y|)^2$,
$M=1+k+O_E(1/L)$. Fix a mesh on $[0,1]$ of sufficiently small
maximum spacing $\eta$, including both endpoints. Apply
\eqref{hdp:sup} at all mesh points, rounding $k$ upward.
The factor $3$ in that estimate covers $b\le2X$, since $\tau\le2X$.
It therefore also covers all $T+w$ with $|w|\le(\log X)^2$,
using $u=T+w$ or $u=T+\tau+w$ according to the smaller ordinate.
On the central line the product has log modulus at most
$(k/2+O(\eta))L$. Combining with \eqref{hdp:convex} at $v=0$
and interpolating to its right edge yields
\begin{equation}\label{hdp:low}
 \begin{aligned}
  &\log |f(1/2+v+iT) f(1/2+v+i(T+\tau))|\\
  &\hspace{1em}\le [H(1-v/D)+O(\eta)+o(1)]L,\qquad
       H=\min\{k/2,DM\}.
 \end{aligned}
\end{equation}
For this interpolation apply three lines to the two finite factors
translated by $iT$ and $i(T+\tau)$, multiplied by $e^{z^2}$,
where $0\le\Re z\le D$ denotes displacement from $1/2$.
On the portion $|\Im z|\le(\log X)^2$, changing the ordinates
by $\Im z$ changes $M$ by $O(\log\log X/L)$.
On the rest of the edges fixed-strip polynomial growth and the
Gaussian damping make the contribution negligible, including
when $W$ is bounded and the ordinates exceed $X$ by a fixed power.
The right edge has $o(L)$ Euler logarithms by $D>\Delta$,
as proved in Theorem~\ref{hsc:exact}. If $b>2X$, we use only
\eqref{hdp:convex}, or \eqref{hdp:low} with $H=DM$; here
$M\ge2-o(1)$.

\subsection*{The height at which the bound is largest}
In the positive half-integral at $r=0$ the high Gaussian has its
maximum $aM^2$ at $W=2aM$, in units of $L$ after removing the
central gamma modulus. The low side $0\le v\le D$ from
\eqref{hdp:low} has upper exponent
\[
        -v^2/(4a)+Mv+H(1-v/D)+O(\eta)+o(1).
\]
For $0\le H\le DM$ this is strictly below $aM^2$, uniformly in
$v$ up to the displayed errors, when
\[
                        a>\frac{D^2}{2DM-H}.
\]
Indeed define $a_0=D^2/(2DM-H)$. At $a=a_0$, the difference
$(2aM-v)^2/(4a)-H(1-v/D)$ is the square
$(v-(2D-2a_0M))^2/(4a_0)$, with its zero in $[0,D]$.
At $a>a_0$ it increases strictly for each $v\in[0,D]$ since
$2aM>D$ and its derivative with respect to $a$ is
$M^2-v^2/(4a^2)$. For $M\ge2-o(1)$ and $H=DM$ the threshold
is at most $D/2+o(1)$. In the other range set $M=1+k$ in the
limit. If $D\ge1/4$, $2D(1+k)-H\ge2D$ for $0\le k\le1$.
If $0<D<1/4$, this denominator decreases to
$D/(1-2D)$ at $k=2D/(1-2D)$ and increases thereafter.
Thus its inverse times $D^2$ has maximum exactly $C(D)$ over
the two ranges. In particular the worst intermediate heights
for this estimate have $b=\tau^{2D/(1-2D)+o(1)}$ when $D<1/4$.

Because $a>C(D)$, compactness gives a positive uniform gap before
the $O(\eta)+o(1)$ errors and gives $W>D$ with a fixed margin.
Choose the mesh and $\eta$ sufficiently small that the gap persists.
At displacement $r\ge0$ the high maximum becomes
$(W+r)^2/(4a)$, while the low exponent gains $rv/(2a)$, which
only enlarges the gap relative to the high maximum. By compactness
the gap and $W+r>D$ persist for $-r_0\le r\le1/2$ with some
$r_0>0$. At the high saddle the combined analytic remainder
after subtracting the gamma slope, with its first two derivatives,
is $o(L)$ on a fixed disk: both finite factors now have a margin
from their actual zero boundary, independently of the exceptional
shifts. The complex-saddle deformation and tail comparison in
Theorem~\ref{hsc:exact} apply verbatim with slope $ML$ and
$\mu=W+r$. In detail they give $2t\Re(A'_r/A_r)=\mu+o(1)>0$
near $r=0$ and give, for $r_0\le r\le1/2$, exponential dominance
of the positive half over the reflected half. Indeed the latter's
high absolute integrand is at most $e^{-rD/t}$ times the former's,
and its low part has the uniform gap just proved. This rules out
every off-line zero at $t=a/L$, including arbitrarily small
positive displacements. The initial heat strip excludes $r>1/2$.

Apply the upper conclusion with $a=C(\Delta)+\epsilon$ and the
lower one with $a=\Delta/2-\epsilon$ when positive. Under ERH both
factors already have all zeros on the line. For $\Delta\ge1/4$
the two boundaries agree; uniform boundedness of the scaled times
from Theorem~\ref{hsc:criterion} upgrades convergence in density
to mean convergence of every fixed positive power of the absolute
error. Finally $C(b)\le1/4$ for $0\le b\le1/2$.

\begin{corollary}\label{hdp:width}
Let $B(\tau,a)$ be the supremum of the rightward zero displacements
of the heated elliptic self-product at $t=a/\log\tau$. Set
$s(b)=1/(1-2b)$ for $0\le b<1/4$ and $s(b)=2$ for
$1/4\le b\le1/2$. For every fixed $a,\epsilon>0$,
\[
 (\Delta-2a)_+-\epsilon\le B(\tau,a)
            \le(\Delta-a s(\Delta))_++\epsilon
\]
outside a set of measure $O_{E,a,\epsilon}(X^{1-c})$ on $[X,2X]$
for some $c>0$. For $\Delta\ge1/4$ both bounding profiles coincide
and $B(\tau,a)\to(\Delta-2a)_+$ in density, also in all finite
positive mean powers of the absolute error.
\end{corollary}
\begin{proof}
The lower bound follows by applying Theorem~\ref{htb:density} to
one fixed zero sufficiently close to the right boundary when
$\Delta>2a$. At any nonnegative displacement $r$, dominance over
the low part follows provided $r+2aM$ strictly exceeds
\[
  \max_{0\le v\le D}\big\{v+2\sqrt{aH(1-v/D)}\big\}
           \le D+aH/D .
\]
The inequality is $2\sqrt{(D-v)aH/D}\le D-v+aH/D$. The conditions
of the saddle comparison, including $r+2aM>D$, thus follow from
$r>D-a(2M-H/D)$ by a fixed margin. The optimization above gives
$2M-H/D\ge s(D)-o(1)$ (extend $s(D)=2$ when $D>1/2$).
Take $D>\Delta$ close enough, the mesh sufficiently fine, and
$r\ge(\Delta-a s(\Delta))_++\epsilon$. The high saddle dominates
at each such displacement. Here $r$ is already bounded away from
zero, so the direct absolute comparison of the reflected half suffices.
The initial width is bounded by $1/2$ for every heat time, giving
the stated mean-power convergence.
\end{proof}

\begin{corollary}\label{hdp:lindelof}
Write $\ell=\limsup_{y\to+\infty}\log^+|L_E(1/2+iy)|/\log y$.
Then $0\le\ell\le\min(\Delta,1/3)$. If $\Delta>0$, at every
sufficiently large shift one has the limiting bound
\[
 \limsup_{\tau\to\infty}\lambda_E(\tau)\log\tau
       \le\frac{\Delta^2}{2\Delta-\ell}\le \frac38 .
\]
If $0<\Delta<1/4$, the density upper bound $C(\Delta)$ in
Theorem~\ref{hdp:main} can be decreased to
$\Delta^2(1-2\ell)/(2\Delta-\ell)$. In particular Lindel\"of
for the fixed elliptic function implies, for every $\Delta$,
that the scaled thresholds converge in density to $\Delta/2$
and that their ordinary limsup equals $\Delta/2$.
\end{corollary}
\begin{proof}
The single-factor bound \eqref{hsc:lowtarget} at the center gives
$\ell\le\Delta$ by taking $D\downarrow\Delta$. The arbitrary-level
Weyl bound for the fixed holomorphic form gives $\ell\le1/3$;
see \href{https://arxiv.org/abs/1707.01576}{Booker, Milinovich
and Ng}. Self-duality covers negative ordinates. Uniformly when
$W<3$, the central bound for the pair is now $H=\ell M+o(1)$
in units of $L$: apply the definition with $\ell+\epsilon'$
at each ordinate, absorb bounded heights and then send
$\epsilon'\downarrow0$. Use the same strip localization and
positive-half comparison. Since $M\ge1-o(1)$ at all heights,
the deterministic sufficient threshold in the limit is
$\Delta^2/(2\Delta-\ell)$. Its denominator is at least $\Delta>0$.
For fixed $\ell\le\Delta$ the quotient is increasing in $\Delta$,
and for fixed $\Delta$ it increases in $\ell$; the constraints
$\Delta\le1/2$ and $\ell\le1/3$ thus give $3/8$.

In density take $H=\min(k/2,\ell(1+k))$ for $M=1+k$.
For $\ell\le\Delta<1/4$ the two expressions meet at
$k=2\ell/(1-2\ell)\in[0,1]$. As in the preceding proof,
$2D(1+k)-H$ decreases and then increases when $D$ is
sufficiently close to $\Delta$. The critical denominator tends
to $(2\Delta-\ell)/(1-2\ell)>0$. For $M\ge2-o(1)$ the
deterministic threshold is at most half its global maximum and
is covered by the displayed density bound, since $1-2\ell>1/2$.
If $\ell=0$, the deterministic upper and Theorem~\ref{htb:density}
match at $\Delta/2$, and the case $\Delta=0$ is immediate.
\end{proof}

\subsection*{Simultaneous ordinary derivatives before heat evolution}
The original product itself admits an inverse result using only ordinary
differentiation. Its differentiation order can grow proportionally to
the logarithm of the shift.

\begin{theorem}\label{hdp:derivative}
Let $F=\Lambda_E$, $P_\tau(s)=F(s)F(s+i\tau)$, $L=\log\tau$.
For fixed $A>0$ and $K\Subset\{\Re s>1/2\}$ there are $c,e>0$
such that, outside a set of measure $O_{E,A,K}(X^{1-c})$ in
$[X,2X]$, simultaneously for every integer $0\le j\le A L$,
\begin{equation}\label{hdp:diffquot}
 \frac{P_\tau^{(j)}(s)}{F^{(j)}(s+i\tau)}
      =F(s+j/L)+O_{E,A,K}(L^{-e})\qquad (s\in K).
\end{equation}
The denominators do not vanish there. Suppose $\Delta>0$. For each
fixed $0<\epsilon<\Delta$, on all but a power-saving exceptional
set, every derivative of order $0\le j\le(\Delta-\epsilon)L$
has an off-line zero. Consequently, for any prescribed integer orders
$j(\tau)=o(\log\tau)$, ERH for $E$ is equivalent to positive
upper density of shifts for which $P_\tau^{(j(\tau))}$ has all
its zeros on the critical line (for this assertion the orders are
assumed measurable and nonnegative).
\end{theorem}

\begin{proof}
Use the same exceptional-disk argument as for Theorem~\ref{htb:density},
now on a fixed complex neighborhood of the compact set $K+[0,A]$
contained strictly to the right of $1/2$. On this neighborhood the
Euler logarithm of $L_E(s+i\tau)$ and its first three derivatives
are $O(L^{1-e_0})$ for some $e_0>0$. Include also a fixed smaller
neighborhood of $K$, and impose the typical finite-part bound
\eqref{hsc:typical} on the fixed region needed for the other arcs
below. The exceptional set is independent of $j$. All differentiations
are with respect to $s$.

For $j\ge1$ put $b=j/L$. By Cauchy's formula, up to the same
nonzero factor $j!/(2\pi b^j)$ the numerator and denominator are
the integrals on $-\pi\le\theta\le\pi$ of
\[
         F(s+b e^{i\theta}+i\tau)e^{-ij\theta},
\]
with and without the amplitude $F(s+b e^{i\theta})$.
Divide both integrands by the nonzero $F(s+b+i\tau)$.
On a fixed complex disk about $\theta=0$ the logarithm of the
common integrand, with constant value zero there at $\theta=0$, is
\[
 \Phi(\theta)= j(e^{i\theta}-1-i\theta)+h(\theta)-h(0).
\]
Indeed subtract the real slope $L$ from the logarithm of the auxiliary
completion in its $s$-variable. Stirling leaves a bounded gamma
slope and bounded higher derivatives on the fixed translates; the
Euler bounds above and the chain rule then give
$|h^{(n)}|\ll_A b L^{1-e_0}=j O_A(L^{-e_0})$ for $1\le n\le3$,
decreasing $e_0$ if needed. These bounds are uniform down to very
small $b$.

The rest of the real circle loses $e^{-c_1j}$ in modulus beyond
any fixed small arc about zero, for some $c_1>0$. For $b$ smaller
than a fixed positive fraction of $\operatorname{dist}(K,\{\Re s=1/2\})$,
the whole circle remains in the chosen zero-free neighborhood of
$K$, reducing that fraction if needed. The residual logarithm
therefore varies by $O(bL^{1-e_0})=o(j)$ around the circle, and
$\Re(e^{i\theta}-1)=\cos\theta-1$ gives the gap. Otherwise
$b$ stays bounded below. After extracting the tangent to the gamma
factor, the absolute exponent in units of $L$ is at most
\[
 b(\cos\theta-1)+2(1/2-\Re s-b\cos\theta)_+
                    +\eta+o(1),
\]
by the functional equation and \eqref{hsc:typical}. In the reflected
part this is $-b(1+\cos\theta)-2(\Re s-1/2)+\eta+o(1)$; elsewhere
the first term suffices. The removed auxiliary finite value at $s+b$
has logarithmic modulus $o(L)$. Choose $\eta$ smaller than the fixed
gaps. This proves the claimed tail estimate in both ranges.

For completeness, on a smaller fixed disk the equation $\Phi'(z)=0$
becomes
\[
                  z=-i\log(1+i h'(z)/j).
\]
The branch at one defines a contraction for large $X$, since
$|h'|/j+|h''|/j=O(L^{-e_0})$ there. It has a unique fixed point
$z=O(L^{-e_0})$. Deform the central arc in the $\theta$ variable
to a parallel segment through $z$. The short connectors have
relative exponential loss $e^{-c_2j}$, using the same real endpoints
and $z=o(1)$. At the stationary point $\Phi(z)=o(j)$,
$\Phi''(z)=-j(1+O(L^{-e_0}))$ and $\Phi'''=O(j)$ on the disk.
On the segment $\theta=z+x$, $x$ real and sufficiently small,
the modulus after extracting $e^{\Phi(z)}$ is at most $e^{-c_3jx^2}$.
Rescaling by $\sqrt j$ therefore gives a nonzero denominator integral
$e^{\Phi(z)}\sqrt{2\pi/j}(1+o(1))$ as $j,X\to\infty$; all estimates
are uniform for $j\le AL$. This lower bound, with a fixed constant
factor, holds for $j$ above a sufficiently large fixed threshold
and $X$ sufficiently large. One may equivalently obtain it by bounding
the rescaled tails first and using the cubic Taylor remainder on a
fixed rescaled interval.

On the new segment subtract the constant amplitude $F(s+b)$.
The amplitude difference is $O_{E,A,K}(b(|z|+|x|))$; on the omitted
arcs and connectors it is $O_{E,A,K}(b)$ by bounded derivatives of
the fixed entire $F$ on a compact set. Dividing the error integral
by the denominator accordingly costs at most
\[
 O_{E,A,K}\big(b L^{-e_0}+b/\sqrt j+b\sqrt j\,e^{-c_4j}\big).
\]
Since $b=j/L\le A$, this is $O(L^{-e})$ with
$e=\min(e_0,1/2)$. For the finitely many smaller positive $j$, on
the entire circle the normalized common integrand differs uniformly
by $o(1)$ from $e^{j(e^{i\theta}-1-i\theta)}$. Its latter integral
is exactly $2\pi e^{-j} j^j/j!>0$; the same subtraction costs
$O(b)=O(1/L)$. For $j=0$ the quotient is exact. This proves
\eqref{hdp:diffquot}.

Choose a fixed zero $\rho$ of displacement $>\Delta-\epsilon/2$.
For all $0\le j/L\le\Delta-\epsilon$ the translates of a sufficiently
small disk about $\rho$ remain in one compact subset to the right
of the critical line. Rouch\'e on each translated boundary recovers
the fixed positive multiplicity, simultaneously for all these orders.
With a zero of order $m$, a radius $C L^{-e/m}$ suffices by the
local factorization at $\rho$. Conversely under ERH the original
product has only critical-line zeros, and all its ordinary derivatives
retain this property by Jensen's disk theorem in the real Hardy
coordinate, as in Theorem~\ref{disk:paths}. This proves the equivalence.
\end{proof}

\begin{corollary}\label{hdp:laguerre}
Let $\rho$ be a fixed zero of $F$ of order $m$ with $\Re\rho>1/2$,
and fix $j\ge1$. Put $c_m=F^{(m)}(\rho)/m!$. On each fixed compact
set of complex $x$, on all but a power-saving exceptional set of shifts,
\[
 \frac{L^{m-j}P_\tau^{(j)}(\rho+x/L)}{c_m F(\rho+i\tau)}
   = e^x (1+d/dx)^j x^m+O(L^{-e}).
\]
In this scale $\max(m-j,0)$ zeros remain at $x=0$ exactly. The
other $\min(m,j)$ zeros are within $O(L^{-e})$ of fixed, simple
negative real numbers. In particular $P'_\tau$ has a zero at
$\rho-m/L+O(L^{-1-e})$.
\end{corollary}
\begin{proof}
On a fixed enlarged compact in $x$, Taylor expansion at $\rho$ and
the auxiliary slope $L+O(L^{1-e})$ give
$L^m P_\tau(\rho+x/L)/(c_mF(\rho+i\tau))=x^m e^x+O(L^{-e})$.
Cauchy's estimate gives the same error after $j$ ordinary
$x$-derivatives. With $n=\min(m,j)$ the polynomial on the right is
$n! x^{(m-j)_+} L_n^{(|m-j|)}(-x)$ by the Rodrigues formula.
The positive simple Laguerre roots give the claimed simple negative
locations by Rouch\'e, and ordinary differentiation preserves the
indicated exact multiplicity at $\rho$.
\end{proof}

\begin{theorem}[Recovery from derivatives on the critical line]\label{hdp:online}
Let $I\subset\mathbb R$ be compact and $0<b_0<A$. At
$s=1/2+iy$, $y\in I$, set $b=j/L$ and
\[
     D_{j,b}U(s)=U^{(j)}(s)+\frac{b}{j+1}U^{(j+1)}(s).
\]
Outside a set of shifts of measure $O(X^{1-c})$ in $[X,2X]$,
simultaneously for $y\in I$ and integers $b_0L\le j\le AL$,
\[
     \frac{D_{j,b}P_\tau(s)}{D_{j,b}(F(\cdot+i\tau))(s)}
                   = F(s+b)+O(L^{-e}) .
\]
The denominator is nonzero. In particular, for each fixed rectangle
strictly to the right of the critical line on whose boundary $F$ is
nonzero, the winding number of the interpolated derivative ratios
computed on the critical line gives its exact zero count with
multiplicity, for all sufficiently large good shifts.
\end{theorem}
\begin{proof}
Use the same Cauchy circle of radius $b$ for both derivative orders.
Their combination inserts the multiplier $1+e^{-i\theta}$ in its
integral. If $U(s)=\omega\overline{U(1-\bar s)}$, $|\omega|=1$,
the full circle integral for $D_{j,b}U(s)$, apart from
$j!/(2\pi b^j)$, is
\[
 J_+(U)+\omega(-1)^j\overline{J_-(U)},\qquad
 J_\pm(U)=\int_{-\pi/2}^{\pi/2}
  (1\pm e^{-i\theta})U(s+b e^{i\theta})e^{-ij\theta}\,d\theta .
\]
This follows by sending the left arc to $\theta\mapsto\pi-\theta$
and using that $s$ lies on the critical line. Both $P_\tau$ and
$F(\cdot+i\tau)$ have such a reflection phase.

The positive saddle neighborhoods now lie a fixed distance to the
right of the line because $b\ge b_0$. Use the logarithm bounds of
Theorem~\ref{htb:density} there and its typical bound on the rest
of the right semicircle, entirely in $\Re s\ge1/2$. Off a fixed
small central arc its exponent $b(\cos\theta-1)+\eta+o(1)$
is strictly negative for sufficiently small $\eta>0$.
Normalize both right semicircle integrals by the same auxiliary
stationary contribution
\[
       Z=F(s+b+i\tau)e^{\Phi(z)}\sqrt{2\pi/j},
\]
using $z$ and $\Phi$ from the preceding saddle calculation. For
$U=F(\cdot+i\tau)$ put $V=1$, and for $U=P_\tau$ put $V=F$.
Then
\[
  J_+(U)= Z(2V(s+b)+O(L^{-e})),\qquad
  J_-(U)= Z O(L^{-e}).
\]
For the normalized estimates with a rate one may use the exact
constant-amplitude right semicircle integral in place of $Z$:
it is $Z(1+o(1))$ and nonzero. Relative to that integral the
claimed power errors follow directly by subtracting the amplitudes
at $\theta=0$, as in \eqref{hdp:diffquot}; on the recentered
segment this costs $O(|z|+j^{-1/2})$. Equivalently, relative to
the displayed $Z$ itself the bounded cubic remainder and the
$O(L^{-e_0})$ relative Hessian error in that calculation give
$O(L^{-e})$ after decreasing $e>0$. The multiplier of $J_-$
vanishes at zero. Hence its conjugate is $O(L^{-e}|Z|)$ regardless
of phase, and the nonzero denominator and quotient follow.
Only ordinary derivatives on the critical line occur.

Choose $b_0,A$ with fixed margins around the horizontal displacements
of the rectangle. Interpolate the displayed ratios linearly in $b$
between consecutive $j/L$; each is continuous in $y$. Bounded
derivatives of $F$ on the fixed rectangle make the interpolation
error relative to $F$ at most $O(L^{-e}+1/L)$. Its minimum modulus
on the boundary is positive, so straight-line homotopy preserves
the winding number there. The argument principle for $F$ gives
the asserted integer without any holomorphy requirement on the
interpolant.
\end{proof}

\begin{theorem}[Ordinary derivative zeros in bounded and growing windows]\label{hdp:clock}
Fix a nondegenerate compact interval $I$. Let $\Delta_I$ be the maximum
of zero and the right displacements of the zeros of $F$ with ordinate
in $I$. For any fixed $\Delta_I<b_0<A$, except on a set of shifts of
measure $O(X^{1-c})$ in $[X,2X]$, every zero of $P_\tau^{(j)}$ with
ordinate in $I$ is simple and lies on the critical line, simultaneously
for $b_0L\le j\le AL$. Consecutive ordinates in $I$ have spacing
\[
                  (\pi/L)(1+O(L^{-e}))
\]
for some $e>0$. Suppose neither endpoint is the ordinate of an off-line
zero of $F$. Let $J_I(\tau)$ be the least $j\ge0$ having all these
zeros on the line, with value $+\infty$ if no such order exists. Then
for every $\epsilon>0$, outside $O_{E,I,\epsilon}(X^{1-c})$ shifts,
\[
               |J_I(\tau)/\log\tau-\Delta_I|<\epsilon .
\]
In particular this recovers the horizontal extent of the original
zeros in $I$ using only ordinary derivatives at heat time zero.

For the single completion $F$ there is the following unconditional
high-ordinate conclusion. For every fixed $0<b_0<A$, there are
$\kappa,c,e>0$ such that, outside $O(X^{1-c})$ shifts, throughout
$[\tau-X^\kappa,\tau+X^\kappa]$ all zeros of every
$F^{(j)}$ for $b_0L\le j\le AL$ are simple and lie on the line.
Their gaps obey the same spacing estimate. For either assertion the
number in every subinterval of the indicated window of length $h$ is
$hL/\pi+O(1+hL^{1-e})$.
\end{theorem}
\begin{proof}
First use a fixed window. Write $U=VF(\cdot+i\tau)$ with $V=F$
or $V=1$. Choose a fixed neighborhood of
$\{1/2+b+iy:y\in I,\ b_0\le b\le A\}$ on which $V$ is bounded
away from zero. This exists by compactness and the hypothesis
$b_0>\Delta_I$ when $V=F$. Enlarge $I$ slightly if necessary for
endpoint estimates. Define the exact analytic right-half integral
\[
 Z_j(s)=\int_{-\pi/2}^{\pi/2}U(s+b e^{i\theta})e^{-ij\theta}\,d\theta,
         \qquad b=j/L.
\]
On a fixed strip $|\Re s-1/2|\le r_0$, $\Im s\in I$, with $r_0>0$
sufficiently small, its unique dominant saddle is still the right one.
Indeed the right endpoint of the Cauchy circle has a fixed zero-free
collar for the auxiliary finite function on the exceptional-set
complement, as before. On the other portions of the semicircle
the exponent bound is
$b(\cos\theta-1)+2(1/2-\Re s-b\cos\theta)_++\eta+o(1)$.
At the critical line it has a uniform negative gap outside the saddle
neighborhood; perturbing $\Re s$ by at most $r_0$ costs at most $2r_0$.
Thus the gap persists, including for centers slightly to the left.
The stationary calculation gives $Z_j(s)\ne0$ throughout this strip.

Furthermore, uniformly on the strip one has
\[
                       Z'_j(s)/Z_j(s)=L+O(L^{1-e}).
\]
To see this, differentiate under the integral and subtract $LZ_j$.
Near the saddle the extra amplitude after removing the common auxiliary
is $V'(w)+V(w)((F'/F)(w+i\tau)-L)$ with $w=s+b e^{i\theta}$.
It is $O(L^{1-e})$ on the deformed complex segment by the Euler
and gamma logarithmic estimates. Its normalized Gaussian absolute
integral has the same upper scale as the original denominator, which
has nonzero leading amplitude $V(s+b)$ bounded below. On the tails
Cauchy's estimate for the differentiated integrand in the $s$ variable
uses a small fixed radius; moving $s$ by that radius costs at most
$e^{C\delta L}$ in the previous tail bound, absorbed by its strict gap
upon choosing the radius $\delta$ small. This proves the display for
both choices of $V$.

The full Cauchy numerator for $U^{(j)}(s)$ is exactly
$Z_j(s)+\omega(-1)^j\overline{Z_j(1-\bar s)}$, by reflection of
the other half-circle. The positive real part of $Z'_j/Z_j$ makes
$|Z_j|$ strictly increasing horizontally throughout the strip.
Consequently there are no off-line zeros there. On its center the
zeros occur each time $\arg Z_j(1/2+iy)$ passes a fixed value
modulo $\pi$; its derivative in $y$ is
$\Re(Z'_j/Z_j)=L+O(L^{1-e})$. This gives the claimed simplicity,
spacing and subinterval counts, including an $O(1)$ endpoint error.

For $r_0\le\Re s-1/2\le1/2$ the denominator of \eqref{hdp:diffquot}
does not vanish. When $V=F$, its numerator also does not vanish since
$F(s+b)$ has a uniform positive minimum modulus on the compact region.
Reflection excludes the other side. Jensen's disk theorem in the real
coordinate excludes zeros of the derivatives beyond the original
critical strip. Thus the assertion covers every real part. If
$\Delta_I>0$ the maximum is attained at a zero with ordinate interior
to $I$ under the endpoint convention. The recovery by Rouch\'e in
Theorem~\ref{hdp:derivative} then rules out purity of all orders up to
$(\Delta_I-\epsilon/2)L$ in $I$ itself. The upper bound follows by
testing an order between $(\Delta_I+\epsilon/3)L$ and
$(\Delta_I+2\epsilon/3)L$. For $\Delta_I=0$ positivity gives the
lower bound automatically.

Finally for $V=1$ take a unit window, a slightly broader range of
orders $b_0/2\le j/\log\tau\le 2A$, and exceptional exponent $c_0>0$
there. Use this unit-window assertion at $\tau+n$ for every integer
$|n|\le X^\kappa+2$, taking $0<\kappa<\min(c_0/2,1/2)$.
The enlarged dyadic range $[X/2,3X]$ costs a fixed factor. Translation
preserves measure, so the union costs $O(X^{1-c_0+\kappa})$.
Use overlapping windows of fixed length, for example $[-1,1]$, so
each consecutive gap locally lies in one of them. The local inequalities
also control the phase derivative on each such window; subinterval
counts over long intervals follow from the uniform individual gap
estimate with one $O(1)$ endpoint error. Since
$\log(\tau+n)=L+O(X^{\kappa-1})$, the same $L$ applies to the full
growing window. This proves the single-function assertion.
\end{proof}

\begin{theorem}[An exponential band for standard completions]\label{hdp:expband}
Let $F$ be a fixed entire unitary cuspidal standard completion of
absolute degree $d$ over a number field, with the normalization and
local bounds of Section~\ref{hc:chapter}. The theorem also holds
for a fixed norm translate of a Dedekind zeta completion after
multiplying by the quadratic which clears its two poles. Write
$B=\sup_{F(\rho)=0}(\Re\rho-1/2)$. Put $L_d(T)=(d/2)\log|T|$.
For any fixed $B<b_0<A$, for all sufficiently large $|T|$ and
every integer $b_0L_d(T)\le j\le A L_d(T)$, the zeros of $F^{(j)}$
with ordinates in $[T-1,T+1]$ are simple and lie on the critical
line. All consecutive spacings there are
\[
            \frac{\pi}{L_d(T)}\{1+O((\log|T|)^{-e})\}
\]
for some $e=e(F,b_0,A)>0$. In particular, for fixed $0<\eta<2/d$,
all sufficiently large $j$ have this all-zero conclusion throughout
the two bands
\[
           \exp(\eta j)\le |\Im s|\le\exp((4/d-\eta)j).
\]
There is no exceptional set of heights in this statement.
\end{theorem}
\begin{proof}
First take $T>0$. In the proof of Theorem~\ref{hdp:clock} take
$V=1$, auxiliary $F$, $b=j/L_d(T)$ and divide exponents by $L_d(T)$.
Stirling for the fixed $d$ real gamma factors, with complex gamma
factors counted twice, gives the real logarithmic slope $L_d(T)+O(1)$,
a bounded imaginary slope, and bounded higher logarithmic derivatives
on the fixed complex translates in use. A fixed polynomial clearing
poles adds $O(1/T)$ to the slope. Completed conjugate reflection
$F(s)=\omega\overline{F(1-\bar s)}$ holds also for a non-self-dual
unitary representation.

Choose $B<D<b_0$. On bounded strips a fixed distance to the right
of $1/2+B$ the finite Euler logarithm and its first three derivatives
are $O((\log T)^{1-e})$ for some $e>0$, unconditionally at all large
heights, by the right-hand disk argument in Theorem~\ref{hsc:exact}.
For $v$ in a bounded interval, the logarithmic modulus of the finite
part at real displacement $v$ and height $T+O(1)$ has the following
upper bounds in units of $L_d(T)$:
\[
 \begin{cases}
    o(1),&v\ge D,\\
    D-v+o(1),&-D\le v\le D,\\
    -2v+o(1),&v\le-D.
 \end{cases}
\]
The middle line is strip convexity between $-D$ and $D$ and the
leftmost real parts use the functional equation, as in
\eqref{hsc:lowtarget}. Fixed poles are far from the ordinates
under consideration. All bounds are uniform in the bounded real
and imaginary translates here.

At a circle centered at displacement $r$ let $v=r+b\cos\theta$.
The normalized modulus exponent adds $b\cos\theta-b$ to the
three respective bounds. In the middle range this gives
$D-r-b+o(1)$. For the right semicircle and $|r|\le r_0$ choose
$0<r_0<\min(D/2,(b_0-D)/3)$; the middle-range bound is strictly
negative and the range $v<-D$ is absent. For $v\ge D$, outside a
fixed small arc about zero there is the usual strictly negative gap
$b(\cos\theta-1)$. On a disk about the right saddle the actual
logarithm is available with a fixed margin. Consequently the
right-half calculation for $Z_j$, including
$Z'_j/Z_j=L_d(T)+O(L_d(T)^{1-e})$, now holds at every height.
For $r_0\le r\le1/2$ use the full circle. In the range $v<-D$
the exponent is $-b(1+\cos\theta)-2r+o(1)$, uniformly negative.
The other off-saddle ranges have just been bounded. The unique
right saddle gives a nonzero full derivative there. The reflection
and phase argument of Theorem~\ref{hdp:clock} handles $|r|\le r_0$,
and Jensen's disk theorem excludes real parts beyond the initial
critical strip. This also proves the spacing estimate.

The negative heights obey the same argument, since the real gamma
slope is $(d/2)\log|T|+O(1)$ there as well, and reflection in
$\Re s=1/2$ preserves the chosen physical ordinate. Finally in the
stated exponential bands the ratio $j/L_d(T)$ stays in a fixed
compact subinterval of $(1/2,\infty)$, whereas $B\le1/2$.
The local result covers the bands, with the usual convention for
gaps between consecutive zeros in the same connected interval.
\end{proof}

For a related classical high-derivative question see
\href{https://doi.org/10.1070/IM2005v069n03ABEH000538}{Rezvyakova's
fixed-order lower bound for the critical-line proportion of the zeros
of derivatives of the Riemann $\xi$-function}.

\begin{theorem}[All ordinates through exponential height]\label{hdp:fulldisk}
Under the hypotheses of Theorem~\ref{hdp:expband}, fix $H>0$
with $dBH<2$; if $B=0$, any fixed $H>0$ is allowed. For all
sufficiently large derivative orders $j$, every zero of the ordinary
derivative $F^{(j)}$ through $|\Im s|\le e^{Hj}$ is simple and lies
on the critical line. Uniformly for consecutive ordinates $T,T'$
there,
\[
          T'-T=\frac{2\pi(1+o(1))}{d\log(2+j+|T|)} .
\]
Without any information beyond the critical strip one may choose
every $H<4/d$. Thus a rational elliptic completion permits every
$H<2$, and the pole-cleared Riemann completion permits every $H<4$.
\end{theorem}
\begin{proof}
Fix a center height $T_0$, $|T_0|\le e^{Hj}+2$, and write
\[
 Y=2+j+|T_0|,\qquad S=(d/2)\log Y,\qquad b=j/S .
\]
All estimates below are uniform for centers $s$ of real displacement
$|\Re s-1/2|\le1/2$ and $|\Im s-T_0|\le2$.
There is a fixed lower margin $\liminf b\ge 2/(dH)>B$.
It suffices by subsequences to handle bounded $b$ and $b\to\infty$.
In the bounded case $|T_0|$ grows exponentially with $j$ and
$S/L_d(T_0)\to1$; Theorem~\ref{hdp:expband} supplies all the
conclusions on bounded neighborhoods (enlarge its fixed unit window
or use overlaps).

Suppose $b\to\infty$. Apply Cauchy's semicircle splitting with
radius $b$ to $F$ itself. On a fixed small complex disk
$|\theta|<c$ around the right saddle the points $w=s+b e^{i\theta}$
satisfy $\Re w\gg b$ and $|w|\asymp b+|T_0|$. In particular the
finite factor is uniformly close to one, and the logarithm of the
whole completion is analytic there. Stirling and absolute convergence
give, uniformly for $1\le n\le3$,
\begin{align*}
  (\log F)'(w)&=S+O_F(1+\log\log Y),\\
  (\log F)^{(n)}(w)&\ll_F (b+|T_0|)^{1-n}+e^{-c' b}
                      \qquad(n=2,3).
\end{align*}
Indeed $\log(2+b+|T_0|)=\log Y+O(1+\log\log Y)$, the gamma
arguments have bounded arguments, and an absolutely convergent
Dirichlet series with constant term one has exponentially decaying
logarithmic derivatives at these real parts. A fixed pole-clearing
polynomial contributes $O_F((b+|T_0|)^{-n})$ to the $n$th
logarithmic derivative. Subtracting the slope $S$ in the $w$
variable therefore gives in the $\theta$ variable a residual $h$
with $|h'|+|h''|+|h'''|\ll_F b(1+\log\log Y)=o(j)$.
The stationary calculation with large parameter $j=bS$ is valid
on the right arc.

We give an absolute estimate on the rest of that semicircle because
the imaginary part of a gamma argument can cross zero. Put
$V_0=2+b+|T_0|$. Uniformly for real $|\theta|\le\pi/2$,
\begin{equation}\label{hdp:globalgamma}
 \log\frac{|F(s+b e^{i\theta})|}{|F(s+b)|}
    \le \frac d2 b(\cos\theta-1)\log V_0
                             +O_F(b+\log V_0).
\end{equation}
To prove this, start with $\Re(s+b e^{i\theta})$ above a fixed
large constant. For each real gamma factor use Stirling on its right
half-plane and the elementary bound
$-y\arg(x+iy)=-\pi|y|/2+O(x)$ for $x>0$. The logarithm of the
full absolute value is at most
\[
 (d/2)b\cos\theta\log V_0
       - (d\pi/4)|\Im s+b\sin\theta|+O_F(b+\log V_0).
\]
The coefficient of $\log|x+iy|$ in Stirling can be taken positive
by the choice of the large constant. Replacing it by its upper bound
with $\log V_0$ incurs only the error shown. At $\theta=0$ the
corresponding expression is a two-sided estimate, since
$|s+b|\asymp V_0$. At the remaining bounded real parts,
fixed-strip polynomial growth, Stirling at large imaginary parts,
and boundedness of the pole-cleared entire function at bounded
heights give the same upper bound. There $b\cos\theta=O_F(1)$.
The two imaginary absolute values differ by at most $b$, proving
\eqref{hdp:globalgamma}. The same inequalities allow a small fixed
complex displacement of $s$ before normalization, at a cost
$O_F(b+\log V_0)$.

Now $b+\log V_0=o(j)$ and $(d/2)\log V_0/S\to1$. Thus the
entire real right semicircle away from any fixed neighborhood of
$\theta=0$ loses $e^{-c_1j}$ in the normalization by $F(s+b)$.
The differentiated tails have the same property by Cauchy's estimate
on a small fixed disk in the center variable. Deform to the complex
stationary point exactly as in Theorem~\ref{hdp:derivative}. It lies
at $o(1)$ and the normalized leading modulus is $e^{o(j)}/\sqrt j$.
Hence the exact right-half integral $Z_j(s)$ never vanishes on
the whole initial critical strip of centers in the unit window.
Inserting $(F'/F)(w)-S=O_F(1+\log\log Y)$ on the right saddle
and absorbing the differentiated tails as above gives
\[
             Z'_j(s)/Z_j(s)=S+O_F(1+\log\log Y)=S(1+o(1)).
\]
The right absolute integral on the deformed central segment and the
actual right integral have comparable magnitudes by the same
Gaussian estimate. This justifies the relative error even when the
original real-arc integral has significant phase cancellation.

For every $r>0$ in the initial strip, strict monotonicity in $r$
prevents cancellation with the conjugate reflected half. On the
line the phase increases with derivative $S(1+o(1))$, so the zeros
are simple, occur with a gap $\pi(1+o(1))/S$, and meet every unit
subinterval for large $j$. The logarithm defining $S$ changes by
$o(1)$ across any such subinterval. The two cases cover every
sequence of centers by subsequence compactness, giving the uniform
conclusion through the indicated height, including the origin.
\end{proof}

The local cosine limits under repeated differentiation proved by
\href{https://doi.org/10.1016/j.jnt.2005.11.004}{H.~Ki for Riemann
$\Xi$} and by \href{https://eprints.whiterose.ac.uk/id/eprint/134951/1/ManyDerivsXi_Resubmitted.pdf}{J.~Gunns and C.~Hughes for Selberg
completions} concern fixed compact sets in their rescaled coordinate.
Here the varying center $T_0$ ranges from the origin through exponential
height and the comparison controls all real parts in the original strip.

\begin{corollary}[Counting the zeros of a high derivative]\label{hdp:count}
Fix $0<H<4/d$. At each real $u$ let $S_j(u)>0$ be the unique
solution of
\[
                     S_j(u)=\frac d2\log(2+|u|+j/S_j(u)).
\]
For all sufficiently large $j$, uniformly on every subinterval
$[v,w]\subset[-e^{Hj},e^{Hj}]$,
\[
 \#\{\rho:F^{(j)}(\rho)=0,\ v\le\Im\rho\le w\}
       =\frac1\pi\int_v^w S_j(u)\,du+O_{F,H}(1+w-v).
\]
In particular the consecutive spacings at ordinate $u$ equal
$\pi/S_j(u)+O_{F,H}(S_j(u)^{-2})$, and
$2S_j(0)/d=\log j-\log\log j+O_d(1)$.
\end{corollary}
\begin{proof}
The right-hand side of the defining equation decreases from infinity
to a finite limit as its positive argument increases, giving
existence and uniqueness. Uniformly through the stated height put
$S=S_j(T_0)$ and $b=j/S$. Then
$S=(d/2)(\log(2+j+|T_0|)+O(1+\log\log(2+j+|T_0|)))$.
This follows directly by bounding $S$ above by $O_d(\log(2+j+|T_0|))$
and below by a positive constant, then substituting these bounds
into the defining logarithm. Thus $b$ has a fixed lower margin
above $1/2$. In the bounded-$b$ case the finite Euler logarithmic
derivative stays bounded on a fixed disk about the right saddle of
the critical-line center, by absolute convergence there. In arbitrary
rank this last assertion follows by Cauchy--Schwarz on the unramified
Euler logarithms against the absolutely convergent nonnegative prime-power
series of $\log L(\sigma,\pi\times\widetilde\pi)$ and $\log\zeta_K(\sigma)$
at a slightly smaller $\sigma>1$; the finitely many bad factors are
regular and nonzero there, and $K$ denotes the base field. The gamma
derivative, including its imaginary part, differs from $S$ by $O(1)$.
When $b\to\infty$, $|s+b e^{i\theta}|\asymp 2+b+|T_0|$ on
the saddle disk, and the new definition of $S$ likewise makes
$(F'/F)(s+b e^{i\theta})-S=O_F(1)$ there. All other saddle
and gap estimates persist for this radius. Equivalently, separate
$b$ above a sufficiently large fixed constant from the complementary
compact range; both the constants in this slope bound and the
comparison of absolute and actual integrals are uniform.
The insertion proof for $Z'_j$ now yields $Z'_j/Z_j=S+O_{F,H}(1)$
at the center and for its imaginary translates of size at most one.
For consecutive zeros in that window phase integration proves the
gap estimate. Implicit differentiation gives $|S'_j(u)|\le d/4$
away from $u=0$, and the same Lipschitz bound holds through zero.
Integrating $S_j$ over each complete gap thus gives $\pi$ with
error bounded by a constant times the gap length. Summing and
allowing the two endpoint gaps (or no complete gap) gives the
displayed count. At $u=0$ substituting $S_j(0)\asymp\log j$
in its equation gives the stated second term.
\end{proof}

In Theorems~\ref{hdp:expband} and~\ref{hdp:fulldisk} one may
replace $B$ by the limiting right displacement
$\limsup_{F(\rho)=0,\ |\Im\rho|\to\infty}(\Re\rho-1/2)$.
Indeed each fixed line strictly to its right admits the same
zero-free logarithm disks at all sufficiently large ordinates.
Only the bounded-radius case uses these disks, and its ordinates
tend to infinity. The growing-radius proof is already in absolute
convergence at its saddle.

\begin{corollary}[A fixed auxiliary for ordinary derivatives]\label{hdp:fixedaux}
Let $G$ be any fixed entire unitary cuspidal standard completion as in
Theorem~\ref{hdp:expband}, and retain one fixed rational elliptic
completion $F$ as auxiliary. In Theorem~\ref{hdp:online} replace
$P_\tau$ by $G(s)F(s+i\tau)$; the reconstructed function is $G(s+b)$.
The local conclusion for $J_I(\tau)$ in Theorem~\ref{hdp:clock}
then holds with $\Delta_I$ the right displacement for $G$.
In particular ERH for $G$ is equivalent to the following test:
for every bounded interval with rational endpoints and every fixed
$\epsilon>0$, at a density-one set of shifts all zeros in that
interval of the ordinary derivative of $G(s)F(s+i\tau)$ of order
$\lfloor\epsilon\log\tau\rfloor$ are simple and on the critical
line. This criterion uses no zero hypothesis on $F$.
\end{corollary}
\begin{proof}
In \eqref{hdp:diffquot} and its on-line version only the translated
auxiliary supplies the density input. The target and its derivatives
are bounded on the fixed compact sets throughout the stationary
argument. Its functional equation supplies the numerator's own
reflection phase, and the original product still has its zeros in
the critical strip. The right-half comparison and the lower translated
zero count therefore go through for $V=G$. Under ERH all local
displacements vanish and every fixed positive scaled order is above
the threshold. Conversely, if a right-hand zero exists, choose it
in the interior of a rational interval, with its endpoints avoiding
the finitely many nearby off-line ordinates, and take $\epsilon$
smaller than its displacement. Simultaneous translated-zero recovery
precludes the proposed purity on all but a power-saving exceptional
set. Reflection treats a zero on the left.
\end{proof}

\clearpage
\section{Ordinary zeta and supersingular arithmetic auxiliaries}\label{hzp:chapter}

\subsection*{A translated standard factor under ordinary RH}
Write $c=1/2$ and use the entire version of the completed zeta function
\[
               V(s)=s(s-1)\Gamma_{\R}(s)\zeta(s),
\]
which differs from Riemann's $\xi$ by a fixed nonzero scalar. Its only zeros are the nontrivial zeta zeros. The two integers $0,1$ are regular nonzero values of $V$.

\begin{corollary}\label{hzp:zeta}
Assume the ordinary Riemann hypothesis. Let $F$ be any fixed entire cuspidal standard completion of absolute degree $d\ge2$, and $\Delta$ its rightmost zero displacement in unitary normalization. For the physical deformation of $F(s)V(s+i\tau)$ the full conclusions of Theorem~\ref{hol:profile} hold with $m=1$. In particular, if $\lambda_\tau^+$ is its least nonnegative all-real time and $b_\tau(t)$ its strip half-width,
\[
 \lim_{\tau\to\infty}(\log\tau)\lambda_\tau^+=\Delta,\qquad
 \lim_{\tau\to\infty} b_\tau(a/\log\tau)=(\Delta-a)_+
                    \quad(a\ge0).
\]
For $F$ over $\Q$ the underlying meromorphic product is the standard completion of $\pi\boxplus|\cdot|^{i\tau}$. Its two potential poles are removed by $(s+i\tau)(s+i\tau-1)$ \emph{before} heat evolution. If $F$ already has a zero at either pole, the entire product has exactly the corresponding order from $F$ there. For a different base field this is the specified product of completions over two fields.
\end{corollary}
\begin{proof}
Factor the auxiliary on the right half-plane as
\[
       G^\dagger(s)=(s+2)^2\Gamma_{\R}(s),\qquad
       L^\dagger(s)=\frac{s(s-1)}{(s+2)^2}\zeta(s).
\]
The artificial point $-2$ lies to the left of the half-plane where the logarithm of $L^\dagger$ is used. The singularity at $1$ cancels, leaving a nonzero value. Under ordinary RH the resulting $L^\dagger$ is holomorphic and zero-free on $\Re s>1/2$, and its analytic logarithm and first two derivatives satisfy the fixed-margin subpower estimates of Theorem~\ref{hsc:exact}. Indeed its additional rational factor and reciprocal are bounded on any sufficiently far right strip, and the entire pole-removed auxiliary has the same functional equation. The three-circles and convexity arguments use $\log L^\dagger$ on the zero-free half-plane, and yield the central Lindel\"of bound for its modulus as well. On a fixed bounded-height region the cancellation at $1$ must be made in this combined factor.

The two logarithmic derivatives of the extra factor in $G^\dagger$ are $2/(s+2)$ and $-2/(s+2)^2$. They are bounded on all right rays used in the proof. At the high shifted ordinate they are $O(1/\tau)$ and $O(1/\tau^2)$ on compact sets. The large-height gamma slope, the lower/upper saddle dichotomy, and every term of the proof of Theorem~\ref{hol:profile} therefore apply with $m=1$. In the unbounded negative tail reflect the entire completed product first, as in that proof. On the very large hybrid-conductor range of Theorem~\ref{hc:leading} the same right-ray argument applies: $L^\dagger$ and its reciprocal are uniformly bounded on every $\Re s\ge1+\epsilon$, and all extra gamma slopes and curvatures cost bounded constants there. This checks the use of the pole-removed factor throughout. The stated order at a canceled pole follows directly from the regular nonzero values of $V$ at zero and one.
\end{proof}

\subsection*{An unconditional quadratic character in fixed positive characteristic}
For odd $q\equiv3\bmod4$ let $\xi_q$ be the quadratic idele class character of $\F_q(x)$ belonging to $y^2=x^3-x$. Its infinity place is ramified, as are the three roots of the cubic. Its full function-field standard $L$-function is
\[
             L(s,\xi_q)=1+q^{1-2s},\qquad
           q^{s-1/2}L(s,\xi_q)=2\cosh((s-1/2)\log q).
\]
To verify it, Artin factorization gives the quotient of the smooth elliptic curve zeta function by that of the rational base. The two factors $1-q^{-s}$ and $1-q^{1-s}$ cancel, leaving the $H^1$ numerator $1-a_q q^{-s}+q^{1-2s}$. The quadratic character of $\F_q^\times$ takes $-1$ to $-1$. Pairing $x$ and $-x$ in the character sum for $a_q$ therefore gives $a_q=0$, including the zero contributions at the branch points. The extension has exact constant field $\F_q$ because its smooth curve has a rational point. This also proves the identity for $q=3^{2n+1}$, giving an unbounded sequence in one fixed characteristic.

For any fixed entire standard completion $F$ as above put $H_{q,0}(s)=F(s)q^{s-1/2}L(s,\xi_q)$. Along this sequence the least nonnegative physical all-real times satisfy, unconditionally,
\[
     \lim_{q\to\infty}2\log q\,\lambda^+(H_q)=\Delta,\qquad
 \lambda^+(H_q)\le
       \frac{\sqrt{1+4(\log q)^2\Delta^2}-1}{4(\log q)^2}.
\]
Indeed Theorem~\ref{oldheat:main} proves the comparison for $F(s)\cosh(R(s-c))$ by horizontal Gaussian translation and the centered Hadamard product. Its proof of the comparison uses only the initial zero strip, reflection, order-one bounds and de Bruijn contraction, available for this general $F$ by Section~\ref{hg:section}. Set $R=\log q$. These products use an actual positive-characteristic degree-one automorphic constituent and a characteristic-zero one. The common analytic $s$ variable makes their heat deformation meaningful.

For a rational elliptic curve $E$ one can take the polynomial from its own reduction. By \href{https://people.math.harvard.edu/~elkies/math_pubs.html}{N.~Elkies, Invent. Math. \textbf{89} (1987), 561--567}, $E$ has infinitely many good supersingular primes $p>3$. At these $p$, $p\mid a_p(E)$ and $|a_p(E)|\le2\sqrt p<p$ for $p\ge5$, so $a_p(E)=0$. The $H^1$ numerator is again $1+p^{1-2s}$. Therefore the preceding ordinary limit and finite-$p$ upper bound recover $\Delta_E$ on the sequence of supersingular reductions of the same curve, using $F(s)=\Lambda(E,s+1/2)$. The reduction polynomial here is evaluated at $p^{-s}$; the good local denominator in the global \emph{unitary} Euler product is evaluated at $p^{-(s+1/2)}$. This shift of half a unit places the reduction polynomial's zeros on the common symmetry line.

\clearpage
\section{How far a nonreal zero can travel under differentiation}\label{disk:section}
Write $F(z)=c_\pi\Lambda(1/2+iz,\pi)$ with the real normalization of Section~\ref{hg:section}. In this coordinate the initial zero corresponding to $\rho=\beta+i\gamma$ has horizontal position $\gamma$ and absolute vertical displacement $|\beta-1/2|$. Local information on the original divisor gives a finite-height theorem for billions of orders of differentiation at once. We use the classical Jensen disks for a real polynomial, and the shifted-product comparison behind de Bruijn's theorem. For the classical results see \href{https://msp.org/pjm/1973/44-1/pjm-v44-n1-p21-p.pdf}{the discussion of Jensen's disk theorem in \emph{Pacific J. Math.} \textbf{44} (1973)}, and \href{https://arxiv.org/abs/math/0601653}{Lagarias, \emph{Zero spacing distributions for differenced $L$-functions}}.

At one step take either $F\mapsto F'$, with cost parameter $h=0$, or choose $h>0$ and $\alpha\in\R$ and take
\[
 F(z)\longmapsto \tfrac12\bigl(e^{i\alpha}F(z+ih)+e^{-i\alpha}F(z-ih)\bigr).
\]
Let $F_m$ result after $m$ such steps with parameters $h_j$. All shifts are in physical height coordinates.

\begin{theorem}\label{disk:paths}
For each nonreal zero $x+iy$ of $F_m$, there is a nonreal zero $\gamma+ib$ of $F$ such that
\[
              \frac{(x-\gamma)^2}{m}+y^2+\sum_{j=1}^m h_j^2\le b^2 .
\]
In particular, if all initial zeros through $|\gamma|\le T$ lie on the line, all zeros of $F^{(m)}$ with $|\Re z|\le T-\sqrt m/2$ are real. If the verified initial zeros are simple, these derivative zeros are simple as well. On every interval wholly inside this band, the consecutive zeros of each derivative have exactly one zero of the next derivative between them, whenever both orders fall within the stated bounds.
\end{theorem}
\begin{proof}
All iterates are real entire of order at most one. For this completion all iterates decay exponentially along both directions in each fixed horizontal strip, by Stirling, the fixed-strip polynomial bound on $L$, and Cauchy's formula for derivatives. They are not identically zero: on the Fourier-transform side the two operations multiply the original continuous Fourier transform by a nonzero linear function or by a nonidentically zero sum of two real exponentials with nonzero complex constants. Such a multiplier has a discrete set of real zeros. Translation of the integration line is justified by the same exponential decay. In particular a nonzero iterate has infinitely many zeros, since a finite-zero real order-one function is a polynomial times an exponential and cannot decay at both ends of the real line.

For a derivative step, if $z=x+iy$ is an inherited nonreal zero choose the identical parent zero. Otherwise use the logarithmic derivative of the real genus-one product. At $y>0$ a real zero contributes a strictly negative imaginary part. A conjugate pair $v\pm ib$ contributes the negative of
\[
 \frac{2y((x-v)^2+y^2-b^2)}{|(z-v)^2+b^2|^2}.
\]
The exponential corrections in the combined pair are real in the logarithmic derivative, and the series of imaginary parts converges absolutely. Absence of a pair with $(x-v)^2+y^2\le b^2$ makes the derivative nonzero. Reflection handles $y<0$. At a shifted step, if the two shifted values both vanish choose the parent with larger absolute imaginary part. In the remaining case the modulus comparison in Lemma~\ref{heat:oneshift} is valid locally for each pair, and the difference of squared moduli has the sign of $(x-v)^2+y^2+h^2-b^2$. If all pairs have positive sign, the two aggregate moduli are unequal for $y\ne0$ (real-zero factors have the same strict orientation). Thus a nonreal zero must have a nonreal parent for which $(x-v)^2+y^2+h^2\le b^2$. This applies for every $\alpha$.

Follow such parents backwards. The inequalities telescope after summing the squared single-step displacements. Cauchy--Schwarz bounds the squared total horizontal displacement by $m$ times their sum, proving the first assertion. For the finite-height consequence an off-line parent must have $|\gamma|>T$ and $|b|\le1/2$, which contradict the first assertion within the claimed band.

To check simplicity for ordinary derivatives, suppose $F^{(j)}$ has a multiple real zero at $x$ in the band, with $1\le j\le m$. Put $V=F^{(j-1)}$. If $V(x)=0$, its multiplicity is one greater, so descend at the same point. Otherwise $V'(x)=V''(x)=0$, forcing $(V'/V)'(x)=0$. Each real zero of $V$ contributes strictly negatively to this logarithmic second derivative. Each conjugate pair $v\pm ib$ contributes $-2((x-v)^2-b^2)/((x-v)^2+b^2)^2$. Therefore some pair satisfies $(x-v)^2\le b^2$, or else the sum is strictly negative. Follow this pair backwards using the nonreal-zero inequalities already proved. This gives an original off-line zero within the same forbidden horizontal distance $\sqrt j/2$ of $x$. The other possible end of descent is a multiple real zero of the initial function there, contrary to the simplicity assumption. For interlacing, the logarithmic derivative of $V$ is strictly decreasing away from its real zeros throughout the interior where no such disk reaches: any offending pair would give the same excluded ancestry. Its limits between successive simple roots are $+\infty$ and $-\infty$. This gives the asserted exact interlacing.
\end{proof}

For the conductor-$11$ elliptic curve take $T=524\,288$. Throughout $|\Re z|\le262\,144$, the zeros of every derivative of orders $0\le m\le274\,877\,906\,944$ are real and simple. The conclusion concerns derivatives of the \emph{completed} function in physical height and uses only the original rectangle certificate. For the uniform conductor census, nontrivial central multiplicities can propagate to lower derivatives; the reality assertion still applies to every isogeny class.

\begin{corollary}\label{disk:measure}
For a fixed cuspidal $\mathrm{GL}_2/\Q$ completion take $m\ge1$ with $\sqrt m\le T/2$. In $[T,2T]$ all real parts of nonreal zeros of all the above iterates with $\sum h_j^2\ge h^2$ lie in one set of length
\[
            \ll_\pi \sqrt m\,T(1+h\log T)e^{-h\log T/8}
                           \qquad(0<h\le1/2).
\]
In particular for one shift, if $h(T)\log T\to\infty$ arbitrarily slowly, all nonreal zeros of all shifts of size at least $h(T)$ are confined to real parts of total length $o(T)$ on each such interval. This last measure concerns positions on the real axis, not a zero count.
\end{corollary}
\begin{proof}
The first theorem projects every admissible nonreal zero into the intervals of radii $\sqrt m\sqrt{b^2-h^2}$ about the original $\gamma$ with $|b|>h$. Only $|\gamma|\le3T$ need be considered. The number of such points with $|b|\ge v$ is $\ll_\pi T^{1-v/8}\log(2T)$ by the fixed-form density theorem of Section~\ref{common:section} and functional-equation symmetry. Bound each radius by $\sqrt m |b|$, sum the tails by the layer-cake identity, and integrate the density estimate from $h$ to $1/2$. The term at $h$ and the tail give exactly the claimed bound, uniformly in $h$.
\end{proof}

\clearpage
\section{Every shifted quadratic Jensen polynomial of thirteen modular forms}
\label{jq:section}
The finite Jensen triangle in Section~\ref{heat:main} leaves infinitely many shifts even in degree two. For the twelve elliptic eta quotients and $\Delta$, we can settle that entire row. A difference between two of the kernel shapes also makes it possible to test exactly what this conclusion implies for the zeros of the Fourier transform.

Use the weight $k=2c$, the center $c$, the even Mellin kernel $K$ and the time parameter of Sections~\ref{num:chapter} and~\ref{heat:main}. Put $H_t(z)=\int_\R K(u)e^{tu^2}e^{izu}\,du$ and define $\gamma_n(t)$ by
\[
 H_t(z)=\sum_{n\ge0}\frac{\gamma_n(t)}{n!}(-z^2)^n.
\]
All $\gamma_n(t)>0$ for real $t$. Its shifted quadratic Jensen polynomial is
\[
 J^{2,n}_t(X)=\gamma_n(t)+2\gamma_{n+1}(t)X+\gamma_{n+2}(t)X^2 .
\]

\begin{theorem}\label{jq:all}
For all thirteen forms, every $J^{2,n}_t$ has two distinct negative roots for every $n\ge0$ and every $t\ge0$. For the twelve forms other than the level-$80$ elliptic newform this holds for every real $t$, including negative time. For the level-$80$ form and each fixed $n$, the quadratic has a conjugate pair of nonreal roots for all sufficiently negative $t$.
\end{theorem}

In particular, in twelve of the actual arithmetic examples the test consisting of \emph{all} shifted quadratics remains strictly satisfied even in the negative-time range where Dobner's theorem supplies nonreal zeros of the deformed transform. At $t=0$ the theorem gives a necessary condition for ERH at every shift.

\begin{lemma}[A covariance test]\label{jq:cov}
Put $f(x)=K(\sqrt x)$ for $x\ge0$, $g(x)=-(\log f)'(x)$, and let $T$ be real. For $r=n+1/2$ put $m_r=\int_0^\infty x^{r-1} f(x)e^{Tx}\,dx$. The strict quadratic Jensen inequality at $(n,T)$ is equivalent to
\[
 \iint_{(0,\infty)^2}(xy)^r f(x)f(y)e^{T(x+y)}
                  (x-y)(g(x)-g(y))\,dx\,dy>0 .
\]
In particular, strictly increasing $g$ suffices for every $n,T$.
\end{lemma}
\begin{proof}
Evenness makes $f$ smooth at zero and our kernels are strictly positive. Double exponential decay in $\sqrt x$ justifies the integrations by parts for every fixed $T$. The duplication formula gives
$\gamma_n(T)=\sqrt\pi\,4^{-n} m_r/\Gamma(r)$.
Writing $b_n=m_r/\Gamma(r)$, the desired inequality is $b_{n+1}^2>b_nb_{n+2}$. Integration by parts gives
\[
 \frac{b_n}{b_{n+1}}=\frac{r m_r}{m_{r+1}}
    =\mathbb E[g(X)-T],
\]
where $X$ has density proportional to $x^r f(x)e^{Tx}$. At the next index the distribution is size-biased by $X$. Thus $b_{n+1}/b_{n+2}-b_n/b_{n+1}$ equals $\operatorname{Cov}(X,g(X))/\mathbb E X$. For independent copies $X,Y$, twice the covariance is $\mathbb E[(X-Y)(g(X)-g(Y))]$. This proves the criterion.
\end{proof}

\paragraph{The derivative to be bounded.}
Let $\lambda=2\pi/\sqrt N$, and use the integer eta exponents $b_d$ from the table of Section~\ref{num:chapter}. With $L(u)=\log K(u)$ on the real axis, its Lambert series is
\[
 L(u)=cu-\lambda e^u
       -\sum_d b_d\sum_{m\ge1}\sigma_{-1}(m)e^{-\lambda dm e^u},
 \qquad \sigma_{-1}(m)=\sum_{v\mid m}\frac1v .
\]
This follows by expanding each $\log(1-q^{dn})$ absolutely and putting $m=nv$. For real $u$ all differentiations below are locally normally convergent. Define
\[
 M(u)=e^{-u}L'''(u),\quad R(u)=L'(u)-uL''(u).
\]
Then $R'(u)=-uL'''(u)$, $R(0)=M(0)=0$ by evenness, and for $u>0$,
\[
       g(u^2)=-\frac{L'(u)}{2u},\qquad g'(u^2)=\frac{R(u)}{4u^3}.
\]
In particular, $M<0$ on the positive axis implies the strict monotonicity required in Lemma~\ref{jq:cov}.

For a term with $y=\lambda d m e^u$, the contribution to $M$ after the leading $-\lambda$ is
$b_d\sigma_{-1}(m)\lambda dm (y^2-3y+1)e^{-y}$. Its derivative is
$b_d\sigma_{-1}(m)\lambda dm\,y(-y^2+5y-4)e^{-y}$. The terms beyond the leading term in $L'$ and $L''$ are respectively $b_d\sigma_{-1}(m)y e^{-y}$ and $b_d\sigma_{-1}(m)y(1-y)e^{-y}$.

\paragraph{Certified kernel bounds.}
The following strict inequalities are sufficient. They can be reproduced with the short interval verifier in \texttt{checks/kernel-inequalities}. We give the entire-axis tail here so that the finitely checked domain is explicit.
\begin{enumerate}
\item For $N\ne80$ in the table, $M'<0$ on $[0,1/8]$ and $M<0$ on $[1/8,3]$.
\item For $N=80$, $M'(0)>1$, $R(3/4)>0$, and $M<0$ on $[3/4,3]$. On $0\le u\le3/4$,
$0<g(u^2)<13/25<g(1)<1$. On $0\le u\le1$ one has $K(u)<1$. On $11/10\le u\le12/10$ one has $K(u)>1/4$, and on $16/10\le u\le17/10$ one has $K(u)>1/10$. Finally $g((12/10)^2)<3/5$ and $g((16/10)^2)>3/4$.
\item For every row, $M(u)<0$ for $u\ge3$.
\end{enumerate}

For the last claim note that $\lambda\ge1/2$, $\lambda d\ge1$ whenever $b_d\ne0$, and $\sum_d |b_d|\le24$. If $u\ge3$ then $y\ge20m$. Since $y^3e^{-y}$ is decreasing there and $\lambda dm\le y$,
\[
 |M(u)+\lambda|\le24\cdot 2\cdot20^3\sum_{m\ge1}m^4 e^{-20m}
 \le \frac{384000e^{-20}}{1-16e^{-20}}<\frac14 .
\]
For the compact checks keep $m\le72$. In every one of $L,L',L'',M,M'$ the entire error is at most
\[
   10^7(2\cdot72+10)e^{-72/2}.
\]
Indeed $\sigma_{-1}(m)\le m$, $y\ge m$, and $y^4e^{-y/2}\le(8/e)^4$; each polynomial factor is bounded by $10y^4$ for $y\ge1$. The bound follows from $\sum_{m>P}m e^{-m/2}\le(2P+10)e^{-P/2}$. For $R$ multiply this error by $1+u$.

The program groups equal $dm$ with exact rational coefficients, evaluates the truncated expressions on a whole interval with outward-rounded balls, adds the preceding errors, and bisects until the required inequality has strictly positive margin. For $0\le u\le1/8$ the quotient defining $g$ is controlled by
$g(u^2)=-\tfrac12\int_0^1L''(vu)dv$, which also covers $u=0$. The largest finite subdivision here is the level-$80$ check. The accompanying JSON gives the number of boxes and a rational lower bound on the certified margin for every inequality. The compact-interval tests and the analytic infinite tail together establish assertions 1--3.

\begin{proof}[Proof of Theorem~\ref{jq:all}]
For the twelve rows with $N\ne80$ we have $M(0)=0$ and $M'<0$ initially, then $M<0$ for the rest of the positive axis. Hence $R>0$ on the positive axis and Lemma~\ref{jq:cov} proves all quadratic inequalities for all times.

At level $80$, $R(3/4)>0$ and $M<0$ thereafter, so $g$ is increasing on $[(3/4)^2,\infty)$. Since $g(1)>13/25$ and $g(x)<13/25$ before $(3/4)^2$, any negative contribution to the double integral of Lemma~\ref{jq:cov} comes from $[0,1]^2$. On that square $0<g<1$ and $0<f<1$ by the certified bounds and monotonicity where applicable. Its negative contribution has magnitude at most $e^{2T}$ when $T\ge0$.

Take the rectangle $(11/10)^2\le x\le(12/10)^2$ and $(16/10)^2\le y\le(17/10)^2$. Its sides have lengths greater than $1/5$ and $3/10$. There $xy>3$, $y-x>1$, $g(y)-g(x)>1/8$, $f(x)>1/4$ and $f(y)>1/10$. Also $x+y>2$. Its positive contribution alone is greater than $e^{2T}3^r/6000$. This dominates the possible negative contribution for $r\ge8$. Thus for every $T\ge0$ the strict quadratic inequality holds at indices $n\ge8$.

For $n<8$, strictness at $T=0$ is included in the independently established finite Jensen triangle of Theorem~\ref{heat:certificate}. To propagate it for $T\ge0$, write $D_n=b_{n+1}^2-b_nb_{n+2}$ with the $b_n$ of the lemma. As $b_n'=r b_{n+1}$, with prime denoting time differentiation, one has
\[
 D_n'=(r+2)(b_{n+1}b_{n+2}-b_nb_{n+3}).
\]
Suppose $D_{n+1}>0$ for all $T\ge0$. At a first time with $D_n=0$, the ratio $b_{n+1}/b_n=b_{n+2}/b_{n+1}$ would be strictly greater than $b_{n+3}/b_{n+2}$. Hence $D_n'>0$, impossible at a first zero approached from positive values. Downward induction from $n=7$ proves the assertion.

It remains to see what happens at large negative time for $N=80$. Put $a=-T$ and $l(x)=\log f(x)$. Evenness and the Taylor series at zero give $l''(0)=L^{(4)}(0)/12=M'(0)/12>0$. For fixed $r>0$, Taylor's formula at zero integrated against $x^{r-1}e^{-ax}$ gives
\[
 \frac{m_r}{\Gamma(r)}=f(0)a^{-r}\!
 \left(1+\frac{r l'(0)}a
       +\frac{r(r+1)[l'(0)^2+l''(0)]}{2a^2}
       +O_r(a^{-3})\right).
\]
On a fixed neighborhood the cubic Taylor error integrates to $O_r(a^{-r-3})$; on its complement exponential decay in $a$ dominates the fixed moment. Substitution into the logarithmic second difference yields
\[
 2\log b_{n+1}-\log b_n-\log b_{n+2}
       =-l''(0)a^{-2}+O_n(a^{-3})<0
\]
when $a$ is sufficiently large. Its discriminant then has strictly negative sign. This proves the final assertion.
\end{proof}

\subsection*{A propagation principle for positive heat times}
The preceding proof used a downward induction from a uniform bound at the higher indices. There is also a way to propagate all shifts together from one time. The variance comparison below is an instance of convex ordering for probability sequences with a log-concave density ratio. See A.~Marsiglietti and J.~Melbourne, \href{https://doi.org/10.1093/imrn/rnag023}{\emph{Concentration inequalities for log-concave sequences}}, Theorem~2.6. We include its short two-crossing specialization to identify the strict case needed here.

\begin{proposition}\label{jq:propagate}
Let $a>0$, and let $f$ be a nonnegative measurable function that is positive on a set of positive measure on $(0,\infty)$. Suppose its moments $\int_0^\infty x^{a+n-1}f(x)e^{tx}\,dx$ are finite for every $t\ge0$ and integer $n\ge0$. Put
\[
  B_n(t)=\frac{1}{\Gamma(a+n)}
             \int_0^\infty x^{a+n-1} f(x)e^{tx}\,dx .
\]
If $(B_n(0))_{n\ge0}$ is log-concave, then $(B_n(t))_{n\ge0}$ is log-concave for every $t\ge0$. Every inequality is strict for every $t>0$.
\end{proposition}
\begin{proof}
Write $r=a+n$ and fix $t>0$. Expansion of the exponential and monotone convergence give a probability distribution on $j=0,1,\ldots$,
\[
 p_j=\frac{(r)_j}{j!}\,\frac{B_{n+j}(0)t^j}{B_n(t)} .
\]
Here $(r)_j=\Gamma(r+j)/\Gamma(r)$. Its first two factorial moments, obtained by shifting the summation index, are
\[
 \mu=\mathbb E J=rt\frac{B_{n+1}(t)}{B_n(t)},\qquad
 \mathbb E[J(J-1)]=r(r+1)t^2\frac{B_{n+2}(t)}{B_n(t)}.
\]
Consequently the desired inequality is equivalent to $\operatorname{Var}J\le \mu+\mu^2/r$.

Compare $p_j$ with the negative-binomial probabilities
$q_j=(1-v)^r(r)_jv^j/j!$ at $v=\mu/(r+\mu)$. Differentiating their generating function $((1-v)/(1-vz))^r$ gives the same mean $\mu$ and variance $\mu+\mu^2/r$. The ratios $p_j/q_j$ are proportional to $B_{n+j}(0)(t/v)^j$ and are log-concave. Thus the indices where $p_j>q_j$ form an interval of integers. If the distributions differ, the two equalities $\sum_j(p_j-q_j)=\sum_jj(p_j-q_j)=0$ require negative terms on both sides of this interval. Indeed, a single change of sign would give a nonzero signed first moment after subtracting a constant times the total mass. The positive interval therefore has finite endpoints $h,k$ with $h>0$. The quadratic
$(j-h+1/2)(j-k-1/2)$ is negative on the interval and positive on its complement. Multiplying by $p_j-q_j$ and summing gives a strictly negative number. The affine terms drop out, so $\sum j^2p_j<\sum j^2q_j$. If $p=q$, their variances agree and the required weak inequality still holds. In fact $p=q$ would force $B_{n+j}(0)=C d^j$ for all $j\ge0$ and some $C,d>0$. The exponential expansion then gives $B_n(1/d)=C\sum_{j\ge0}(r)_j/j!=\infty$, by the binomial series and monotone convergence as its argument increases to one. This contradicts the moment assumption. This proves the proposition.
\end{proof}
For $a=1/2$ and $f(x)=K(\sqrt x)$ the proposition shows that the infinite row of strict quadratic inequalities at time zero is already sufficient for every positive time. It imposes no condition on higher-degree Jensen polynomials.

\clearpage
\section{A finite Jensen degree that detects nonreal zeros at deep negative times}
\label{jl:section}
Fix a holomorphic cuspidal newform $F$ of level $N$ and weight $k$ with real Fourier coefficients. Its root number is $(-1)^e$, $e\in\{0,1\}$. In this section take the real Mellin kernel
\[
 K(u)=e^{ku/2}F(ie^u/\sqrt N),\qquad
 H_T(z)=\int_{\R} K(u)e^{Tu^2}e^{izu}\,du.
\]
The Fricke symmetry gives $K(-u)=(-1)^eK(u)$. For the odd sign the unphased $H_T$ is purely imaginary on the real axis, so we remove $iz$ at the origin. Define real coefficients by
\[
 \frac{H_T(z)}{(iz)^e}
   =\sum_{n\ge0}\frac{\gamma_n(T)}{n!}(-z^2)^n,
 \qquad
 J_T^{d,n}(X)=\sum_{j=0}^d\binom dj\gamma_{n+j}(T)X^j.
\]

\begin{theorem}[A common detecting degree on growing ranges of shifts]
\label{jl:allforms}
There are a finite degree $d_F\ge2$ and a positive constant $C_F$, depending only on the newform, such that $J_{-a}^{d_F,n}$ has a nonreal zero whenever $n\ge0$ and $a\ge C_F(n+1)$. This holds for either sign of the functional equation. Thus a single sufficiently negative time admits a consecutive block of shifts of length proportional to its absolute value, at one and the same degree.
\end{theorem}

\begin{lemma}[An exact gamma integral]\label{jl:asymp}
Put $f(x)=K(\sqrt x)/x^{e/2}$, with its analytic continuation at zero, $r=n+e+1/2$ and $c_e=\sqrt\pi/2^e$. Define
\[
 P_d(f;x,z)=\sum_{j=0}^d\binom dj(-1)^j f^{(j)}(x)z^{d-j}.
\]
For $a>0$ let $X$ have gamma distribution with shape $r+d$ and rate $a$. Then one has the exact identity
\begin{equation}
 \frac{4^n a^{r+d}}{c_e}\,J_{-a}^{d,n}(-4a+4z)
      =\mathbb E P_d(f;X,z).
 \label{jl:gamma}
\end{equation}
For $d$ fixed, $a\to\infty$ and $n/a\to x_0\ge0$, this converges coefficientwise to $P_d(f;x_0,z)$. At $x_0=0$ a bound on the pointwise error on $|z|\le Z$ is
\[
 \frac{r+d}{a}\sum_{j=0}^d\binom dj Z^{d-j}
                    \sup_{x\ge0}|f^{(j+1)}(x)|.
\]
\end{lemma}
\begin{proof}
For $s>0$ write $F_s(h)=a^s\int_0^\infty x^{s-1}e^{-ax}h(x)\,dx/\Gamma(s)$. By parity, $x=u^2$ and the duplication formula give $\gamma_n(-a)=c_e4^{-n}a^{-r}F_r(f)$. The function $f$ is analytic near zero. All its fixed derivatives are bounded on the positive axis and decay at infinity, by differentiating the exponentially convergent newform Fourier series and using $d/dx=(2\sqrt x)^{-1}d/du$. Integration by parts has no boundary terms and yields
\[
 F_{s+1}(h)-F_s(h)=a^{-1}F_{s+1}(h')
\]
for $h=f$ and each of its derivatives. If $E$ translates $s$ by one, the rescaled left side of \eqref{jl:gamma} equals $a^d(1-E+(z/a)E)^d F_r(f)$. The operators $E$ and $E-1$ commute. In the binomial expansion the term using $(E-1)^j$ is $\binom dj(-1)^j z^{d-j}F_{r+d}(f^{(j)})$ by the integration by parts identity. This proves \eqref{jl:gamma}.

Since $\mathbb E X=(r+d)/a$, the mean-value theorem gives the displayed error at zero. For arbitrary fixed $x_0$, use
\[
 \mathbb E|X-x_0|\le |(r+d)/a-x_0|+\sqrt{r+d}/a
\]
from the gamma variance. The same derivative bounds imply coefficientwise convergence. Equivalently the unscaled formula
\[
 J_T^{d,n}(X)=\frac{c_e4^{-n}}{\Gamma(r+d)}
      \int_0^\infty x^{r+d-1}e^{Tx}
           P_d(f;x,X/4-T)\,dx
\]
holds at every real or complex $T$: it follows directly from \eqref{jl:gamma} for negative $T$, and the two entire functions agree by analytic continuation. The double exponential decay in $\sqrt x$ dominates every fixed exponential in $x$.
\end{proof}

\begin{proof}[Proof of Theorem~\ref{jl:allforms}]
Write $f(x)=x^h g(x)$ near zero, $h\ge0$, $g(0)\ne0$. It is an analytic germ because $K$ has definite parity and is analytic near zero. It is nonzero because $F$ is a nonzero holomorphic form. In the remainder of the proof write $P_d(f;z)=P_d(f;0,z)$. For any $D\ge0$ the derivative identity at zero gives
\[
 P_{D+h}(f;z)=(-1)^h\frac{(D+h)!}{D!}\,P_D(g;z).
\]
We prove that some $P_D(g)$ has a nonreal zero. Suppose that they all have only real zeros. Then for every positive $D$ the reversed and rescaled polynomial
\[
 B_D(z)=\frac1{g(0)}\sum_{j=0}^D\binom Dj g^{(j)}(0)(z/D)^j
       =\prod_i(1+\alpha_{D,i}z)
\]
has only real inverse-root parameters $\alpha_{D,i}$. Zeros at zero of the unreversed polynomial reduce the degree after reversal. Comparison of the first two coefficients gives
\[
 \sum_i\alpha_{D,i}=g'(0)/g(0),\qquad
 \sum_i\alpha_{D,i}^2=(g'(0)/g(0))^2-(1-1/D)g''(0)/g(0).
\]
These power sums are uniformly bounded. Since $|1+w|\le\exp(\operatorname{Re}w+|w|^2/2)$, the sequence is locally uniformly bounded on $\C$. A normal subsequence converges to an entire function. For every $j$ its $j$th derivative at zero equals $g^{(j)}(0)/g(0)$, as follows from the coefficients of $B_D$. Thus the germ $g$ extends to an entire function.

This is impossible at the cusp $-1$. Indeed $K(u)=u^{e+2h}g(u^2)$ near zero would give an entire extension of $K$, agreeing with the original Mellin kernel on the strip $|\operatorname{Im}u|<\pi/2$. The map $u\mapsto\tau=ie^u/\sqrt N$ is locally invertible at $u_*=\log\sqrt N+i\pi/2$, where $\tau=-1$. Multiplication by $e^{-ku/2}$ would extend $F$ holomorphically across $-1$. Cuspidality gives $F(-1+iy)=O(y^{-k}e^{-c/y})$ as $y\downarrow0$ for some $c>0$. In fact the scaling matrix at that rational cusp makes the imaginary part equal to $1/y$ up to its fixed cusp-width normalization, and the Fourier development there has strictly positive first exponent. Such infinite-order decay is impossible for a nonzero holomorphic germ. This proves the desired existence of a nonreal zero of $P_D(g)$. Take a closed disk containing that zero in its interior, disjoint from the real axis and with no zeros on its boundary. Write $d_F=D+h$ and let $m_0>0$ be the minimum modulus of the limit on the boundary. At degree $d_F$ the error bound of Lemma~\ref{jl:asymp} is at most $B_F(n+e+1/2+d_F)/a$ there. A constant $C_F$ can be chosen so that $a\ge C_F(n+1)$ makes this less than $m_0$ simultaneously for all $n$. Rouch\'e's theorem proves the claim. Degree zero and degree one polynomials with real coefficients do not have nonreal zeros.
\end{proof}

\begin{corollary}[Effective search]\label{jl:effective}
Given the level, weight and an algorithm for the Fourier coefficients of the normalized newform, one can compute a witnessing degree and a constant $C_F$ with the property in Theorem~\ref{jl:allforms}.
\end{corollary}
\begin{proof}
All derivatives of $f$ at zero are computable with enclosures. On $|u|=1/2$, the newform argument has imaginary part at least $e^{-1/2}\cos(1/2)/\sqrt N$. The Deligne bound, or the coarser $|a_m|\le m^{k+2}$ for these normalized newforms, bounds the Fourier tails geometrically with any fixed number of derivatives. The parity quotient is regular at zero. Enumerate degrees and rational disks disjoint from the real axis, increasing the coefficient precision. A nonzero winding number of the limiting polynomial on such a disk, with boundary image separated from zero, can be certified by subdivision and interval bounds for the polynomial and its derivative. The proof of Theorem~\ref{jl:allforms} guarantees eventual success. This procedure requires no decision about whether any particular central derivative vanishes exactly.

To bound all the derivatives in Lemma~\ref{jl:asymp}, fix $R=1/4$. A bound for $f$ on $|x|\le R$ follows from the preceding estimate at $|u|=1/2$ and the maximum modulus principle. For any real $x\ge R/2$ and $|w-x|\le R/8$, take the principal square root $u=\sqrt w$. One has $|w|\ge3R/8$, $v=\operatorname{Re}u\ge0$ and $|\operatorname{Im}u|<1/2$. Consequently
\[
 |f(w)|\le (3R/8)^{-e/2}\,
     e^{kv/2-c e^v/2}\sum_{m\ge1}m^{k+2}e^{-cm/2},
         \qquad c=2\pi\cos(1/2)/\sqrt N .
\]
The real exponential has an elementary finite maximum, and the convergent power-geometric sum is computable. For $0\le x\le R/2$ the disk $|w-x|\le R/8$ lies inside $|w|\le R$. These two estimates give one computable uniform bound $G$ on all such disks. Thus $\|f^{(j)}\|_\infty\le j! G(8/R)^j$ by Cauchy's estimate. Substitution into Lemma~\ref{jl:asymp} gives a computable $C_F$ from the disk and its certified minimum modulus.
\end{proof}

\subsection{Explicit detecting degrees}

For the thirteen eta products in Section~\ref{num:chapter}, $e=h=0$ and $f(0)>0$. The table gives the first degree at which the limiting polynomial has a nonreal zero and the exact number of its real roots at that degree. All its roots there, and at the preceding degrees, are simple. Consequently these are also the first failing degrees for the original Jensen polynomials whenever $a=-T\ge C_F(n+1)$ with a sufficiently large row-dependent constant. Apply the same bound to the finitely many lower degrees and disks around their simple roots. Real conjugation and a disk around each simple real root preserve its reality in the approximating polynomials.
\[
\begin{array}{c|rrrrrrrrrrrrr}
 N&1&11&14&15&20&24&27&32&36&48&64&80&144\\ \hline
 \text{first degree}&10&7&5&6&5&6&4&4&4&3&5&2&4\\
 \text{real roots}&8&5&3&2&3&4&2&2&2&1&3&0&2\\
 q_F&38&22&18&20&18&20&16&16&16&14&19&12&17
\end{array}
\]
The last row supplies certified constants $C_F=10^{q_F}$ valid both for the failing degree and all its lower degrees. These are sufficient bounds on the heat times. At level $48$ this gives a direct distinction between two successive degrees. Every quadratic is strictly hyperbolic at every shift and time by Theorem~\ref{jq:all}. At each fixed shift the cubic loses two real roots sufficiently far in negative time. At level $1$ the first such witness occurs in degree ten.

Here are the analytic bounds and exact-arithmetic steps to reproduce the counts. Put $l=\log f$ and write $L(u)=\log K(u)$ as in Section~\ref{jq:section}. Replacing $f(x)$ by $e^{-l'(0)x}f(x)/f(0)$ translates the roots of $P_d$ and does not change their number. This follows either from Leibniz's rule or from the exponential generating function $e^{zt}f(-t)$ for $P_d(f;z)$. Thus use $l_1=0$ and
\[
 l_j=\frac{j!}{(2j)!}L^{(2j)}(0)\qquad (2\le j\le12).
\]
For instance the centered cubic is $z^3+3 l_2 z-l_3$. Its three roots are real and distinct precisely when $4(-l_2)^3>l_3^2$.

Let $S(q,i)$ denote a Stirling number of the second kind. For $q=2j$ and $y=\lambda dm$, the Lambert term satisfies
\[
 \left.\frac{d^q}{du^q}\exp(-ye^u)\right|_{u=0}
       = e^{-y}\sum_{i=0}^q S(q,i)(-y)^i.
\] This follows by $q$ applications of $y\,d/dy$. The leading term $-\lambda e^u$ contributes $-\lambda$. Keep $m\le P=1024$ and put $B_q=\sum_i S(q,i)$. Since $y\ge m$, $\sigma_{-1}(m)\le m$, $\sum_d |b_d|\le24$ and $y^q e^{-y/2}\le(2q)^q$, a uniform absolute error for $L^{(q)}(0)$ is
\[
 24 B_q(2q)^q(2P+10)e^{-P/2}.
\]
Here $y\ge1$ on the discarded range; the final factor bounds $\sum_{m>P}m e^{-m/2}$. This bound applies to all required derivatives.

The resulting enclosure for each $l_j$ uses 768-bit outward-rounded real balls. Recover the centered Taylor derivatives $c_j$ of $e^{-l'(0)x}f(x)/f(0)$ by $c_0=1$ and
\[
 c_j=\sum_{i=1}^j\binom{j-1}{i-1} l_i c_{j-i}.
\]
Form the polynomials with coefficients $(-1)^j\binom dj c_j$ at $z^{d-j}$. Take each polynomial, its derivative and the successive negative Euclidean remainders. Every leading coefficient in the sequence has definite nonzero interval sign and every last remainder is a nonzero constant. Thus Sturm's theorem counts simple real roots by the difference between its sign variations at negative and positive infinity. In polynomial division the highest coefficient cancels identically, so the verifier discards that coefficient and encloses each lower coefficient. A second check uses disjoint complex root enclosures. Every box meeting the real axis has a conjugate box disjoint from every other isolating box, forcing the enclosed root to be real. All other enclosures are disjoint from the real axis. The full count for every degree two through twelve is in \texttt{negative-limits.json}; \texttt{negative\_limits.py} regenerates it from the eta exponents.


For the constants in the last row, the program uses the root enclosures of the centered monic limit. Around each simple real root at a lower degree it chooses a real-centered disk; at the first failing degree it chooses a nonreal-centered disk disjoint from the axis. The radius $b$ is smaller than one quarter of the lower distance from its center $w$ to every other root enclosure. The enclosed root is less than $b/2$ from $w$. A nonreal disk also has $b<|\operatorname{Im}w|/2$. The product factorization of the monic polynomial thus bounds its minimum modulus on the circle from below by
\[
 (b-|w-\rho_i|)\prod_{j\ne i} (|w-\rho_j|-b).
\]
The original uncentered limit has the additional factor $f(0)$ and its circle satisfies $|z|\le |w|+|l'(0)|+b$. The values $f(0)$ and $l'(0)$ are enclosed by the same Lambert series, at derivative orders $0$ and $2$, respectively. The latter includes the factor $1/2$ from $u^2=x$.

For the uniform Cauchy bound of Corollary~\ref{jl:effective} on these even rows, put $s=k+2$, $S(t)=s!/(1-e^{-t})^{s+1}$ and $c=2\pi\cos(1/2)/\sqrt N$. The geometric series majorant $\sum_{m\ge1}m^s e^{-tm}\le S(t)$ gives the admissible common bound
\[
 G=e^{k/4}S(ce^{-1/2})+(k/(c\exp(1)))^{k/2}S(c/2).
\]
Let $Z=|w|+|l'(0)|+b$ and let $m_0$ be the certified uncentered minimum modulus on the circle. The program verifies by directed rounding that $10^{q_F}$ strictly exceeds
\[
 (d+1/2)\,G\,m_0^{-1}
       \sum_{j=0}^d\binom dj Z^{d-j}(j+1)!32^{j+1}
\]
on every circle needed for the claim. The factor $d+1/2$ bounds $(n+1/2+d)/(n+1)$ for $d\ge1$. Thus the constants certify growing blocks of actual Jensen polynomials directly.

\subsection{The unique quadratic transition at level eighty}

\begin{theorem}\label{jl:transition}
For the level-$80$ elliptic curve there is a unique $T_n$ for each $n\ge0$ with
\[
       -32(n+1)<T_n<0 .
\]
Its shifted quadratic has nonreal roots for $T<T_n$, one double root for $T=T_n$ and two distinct negative roots for $T>T_n$. There is a number $x_*>0$, determined by the kernel, such that
\[
   \lim_{n\to\infty}\frac{-T_n}{n}=\frac1{x_*},
   \qquad 0.5988363970<\sqrt{x_*}<0.5988363971 .
\]
In particular the asymptotic slope is approximately $2.78858330$.
\end{theorem}

\begin{proof}
Continue to write $l=\log f$ and $g=-l'$. We first record the additional certified bounds,
\[
 f(x)l''(x)>1/100 \quad(0\le x\le1/4),\qquad
 f(x)l''(x)>-1/16 \quad(x\ge0).
\]
The same Lambert series as in Section~\ref{jq:section} give $M'(u)>1/3$ and $K(u)>3/5$ on $[0,1/8]$. There $M(u)>u/3$ for $u>0$ and $-R(u)=\int_0^u t e^t M(t)\,dt\ge u^3/9$. Thus $f(u^2)l''(u^2)=-K(u)R(u)/(4u^3)\ge(3/5)/36$, including zero by continuity. On $[1/8,1/2]$ and $[1/2,3]$ directly enclose this expression on subdivided intervals to verify the two claimed lower bounds. For $u\ge3$ write $v=e^u\ge20$. The Lambert tails in $L,L',L''$ are individually less than $1/4$ by the majorant for $M$ in Section~\ref{jq:section}, since their polynomial factors are bounded by $y^3$ for $y\ge20$. From $c=1$, $1/2<\lambda<1$ one obtains
\[
 |K(u)l''(u^2)|
 \le e^{u+1/4-\lambda v}
        \frac{2+v+u(1+v)}{4u^3}
 \le \frac{e^{1/4}}{18}v^2 e^{-v/2}<1/16 .
\]
The last inequality uses monotonicity at $v\ge20$. The script \texttt{level80\_negative.py} encloses the extra compact ranges by the same outward-rounded method as the earlier certificate.

For $d=2$ the exact identity of Lemma~\ref{jl:asymp} gives a quadratic in $z$ with coefficients $\mathbb E f,-2\mathbb E f',\mathbb E f''$. Here $X$ has gamma shape $n+5/2$ and rate $a>0$. If $\mathbb E_f$ denotes expectation additionally weighted by $f(X)>0$ and renormalized, the discriminant divided by $4(\mathbb E f)^2$ equals
\[
         -\mathbb E_f l''-\operatorname{Var}_f(l').
\]
For $a\ge32(n+1)$ the gamma exponential moment yields
\[
 \mathbb P(X>1/4)\le2^{n+5/2}e^{-a/8}
       \le 2^{5/2}e^{-4}(2e^{-4})^n<1/9.
\]
The displayed curvature bounds give $\mathbb E[f(X)l''(X)]>(8/9)/100-(1/9)/16>0$, so the quadratic has negative discriminant throughout this region.

The function $g$ has exactly one critical point $x_*$, a strict minimum. In fact $g'(u^2)=R(u)/(4u^3)<0$ through $u=1/2$ by the positive curvature. The compact verifier proves $M(u)<0$ from $1/2$ through $3/4$. The earlier bounds continue this sign to infinity. Thus $R'(u)=-u e^u M(u)>0$ for $u\ge1/2$ and $R(3/4)>0$. It crosses zero exactly once. Evaluating $R$ with the Lambert cutoff raised to $128$ at the two rational endpoints in the theorem gives strictly opposite signs. Also $g(x)\to\infty$ as $x\to\infty$, since $L'(u)\sim-\lambda e^u$.

For a fixed $n$, put $r=n+1/2$, $w(x)=x^r f(x)$ and
\[
    A(T)=\frac{\int_0^\infty g(x) w(x)e^{Tx}\,dx}
                {\int_0^\infty w(x)e^{Tx}\,dx}.
\]
Then $A'(T)=\operatorname{Cov}_T(X,g(X))$ in the probability law proportional to $w(x)e^{Tx}$, so its sign and the quadratic discriminant agree by Lemma~\ref{jq:cov}. At negative infinity $A(T)\to g(0)$ by Laplace concentration at zero, for example by scaling $x$ by $1/|T|$ and dominated convergence. At positive infinity the probability of every bounded interval tends to zero: compare its integral to any unit interval strictly to its right. The lower boundedness of $g$ and $g(x)\to\infty$ then imply $A(T)\to\infty$. The negative-discriminant estimate already proved gives $A'(T)<0$ far to the left. Thus $A$ attains a finite global minimum $m<g(0)$.

For completeness, if a continuous real $h$ changes sign $k$ times on $(0,\infty)$ and the needed exponentially tilted integrals are finite, then $\int h(x)w(x)e^{Tx}\,dx$ has at most $k$ real zeros, counted with multiplicity, provided $h$ is nonzero on some interval and $w>0$. For the functions used here the integrals are real analytic in $T$. When $k=0$ the integral has constant strict sign. Otherwise choose a sign-changing point $s$. Multiplication by $x-s$ removes one sign change. The integral of $(x-s)h(x)$ is $e^{sT}$ times the derivative of $e^{-sT}\int h(x)w(x)e^{Tx}\,dx$. Rolle's theorem, with multiplicities, proves the assertion by induction. Apply this argument to $h(x)=g(x)-c$. For $c\ge g(0)$ it changes sign at most once; for all $c$ it changes sign at most twice.

It follows that $A(T)-c$ has at most two real zeros counted with multiplicity, and at most one when $c\ge g(0)$. At $c=m$ the minimum is thus unique and its multiplicity is exactly two. For $m<c<g(0)$ at least one crossing on each side of that minimum follows from the two limits. Both crossings must then be simple. For $c\ge g(0)$ every crossing is likewise simple. Hence the minimum is the only point where $A'$ vanishes. The discriminant sign on its two sides now follows, as does the stated interval for its location, since all shifted quadratics are strictly hyperbolic at zero. Their coefficients are positive at every real time.

Finally fix $y>0$ and take $a=n/y$ for $n\to\infty$. The gamma limit in Lemma~\ref{jl:asymp} gives a quadratic whose discriminant divided by four is
\[
    f'(y)^2-f(y)f''(y)=f(y)^2 g'(y).
\]
It is strictly negative for $0<y<x_*$ and strictly positive for $y>x_*$. The uniqueness of $T_n$ therefore bounds the limiting ratios $-T_n/n$ above by $1/y$ whenever $y<x_*$ and below by $1/y$ whenever $y>x_*$. Let $y$ approach $x_*$ from either side.
\end{proof}

There is also an expansion of the transition time itself. Put $s=n+5/2$ and $m=s/a$. A gamma variable with this mean and shape has variance $m^2/s$ and third central moment $2m^3/s^2$. Its fourth central moment is $3m^4/s^2+6m^4/s^3$. Taylor expansion of $h=f,f',f''$ around $m$ therefore gives, uniformly for $m$ near $x_*$,
\[
 \mathbb E h(X)=h(m)+\frac{m^2}{2s}h''(m)+O(s^{-2}).
\]
All the necessary derivatives are bounded. Put $D_0=f'^2-ff''$ and $D_2=2f'f'''-(f'')^2-ff^{(4)}$. At the transition the quadratic discriminant equation is
\[
 0=D_0(m)+\frac{m^2}{2s}D_2(m)+O(s^{-2}).
\]
We know $m\to x_*$, and $D_0'(x_*)=f(x_*)^2 g''(x_*)>0$ by $R'(\sqrt{x_*})>0$. It follows first that $m-x_*=O(1/s)$ and then, by one more Taylor expansion, that
\[
 -T_n=\frac{s}{x_*}+\frac{D_2(x_*)}{2D_0'(x_*)}+O(1/s)
      =\frac{n}{x_*}+c_*+O(1/n).
\]
Writing $u=\sqrt{x_*}$, the intercept is
\[
 c_* =\frac1{u^2}+\frac{L'(u)}{2u}
                   +\frac{1+M'(u)/M(u)}{4u},\qquad
       6.016800<c_*<6.016801 .
\]
Indeed $l''(x_*)=0$ gives $D_2/(2D_0')=l'+l^{(4)}/(2l''')$. At this point $uL''(u)=L'(u)$, hence $l'''(x_*)=L'''(u)/(8u^3)$ and $l^{(4)}(x_*)=L^{(4)}(u)/(16u^4)-6L'''(u)/(16u^5)$. Substitute $L'''(u)=e^uM(u)$. The displayed enclosure is evaluated on the whole certified interval for $u$, without replacing the root by a floating-point estimate.

The argument gives successive terms to any fixed order. Indeed the moment generating function of $X-m$ is
\[
 \exp\left(\sum_{j\ge2}\frac{m^j t^j}{j s^{j-1}}\right)
\]
near $t=0$. Its derivatives at zero give all central moments explicitly. Taylor expansion through degree $2b+1$ of each $h$ has error $O(s^{-b-1})$ for $m$ in the indicated compact interval, since the even moment of order $2b+2$ has that bound. The nonzero coefficient $D_0'(x_*)$ determines each successive power of $1/s$ in the expansion of the root. These asymptotic statements do not give an exact linear formula for any finite shift.

\begin{corollary}\label{jl:block}
The transition times satisfy $T_0>T_1>T_2>\cdots$. For the level-$80$ newform all its shifted quadratics are strictly hyperbolic exactly at $T>T_0$, with
\[
       \frac{100358703}{2^{24}}<-T_0<
              \frac{100358704}{2^{24}} .
\]
At any real time the shifts with a nonreal pair form an initial segment of the nonnegative integers. Its length is $x_*|T|+O(1)$ as $T\to-\infty$. If $T=T_n$, exactly that index has a double root.
\end{corollary}
\begin{proof}
Use the positive normalized moments $b_n$ and $D_n=b_{n+1}^2-b_nb_{n+2}$ from Section~\ref{jq:section}. The unique minimum in the preceding proof has multiplicity two. Hence its first derivative vanishes simply with positive slope; the positive differentiable factors relating this derivative to $D_n$ give $D_n'(T_n)>0$. At this time the differentiation formula used in Section~\ref{jq:section} simplifies to
\[
 D_n'(T_n)=(r+2)\frac{b_{n+1}}{b_{n+2}}D_{n+1}(T_n)>0 .
\]
Thus $T_n>T_{n+1}$. The two-term expansion gives $-T_n=n/x_*+O(1)$, so the strictly decreasing sequence goes to negative infinity. Inverting this estimate and using the signs supplies the counting assertion.

For the finite numerical bracket define
\[
 M_j(a)=\int_{\R} u^{2j}K(u)e^{-au^2}\,du .
\]
The sign at shift $n$ is that of $(n+3/2) M_{n+1}(a)^2-(n+1/2)M_n(a) M_{n+2}(a)$. Evaluate the trapezoidal sum on the real grid of step $1/256$ through $|u|\le9$ using the level-$80$ eta product. The first error term in \eqref{heat:error} acquires only the factor $e^a$, since on the two displaced lines $|e^{-a(u\pm i)^2}|\le e^a$. Its real tail error is unchanged for $a\ge0$. In that formula use $b=6+2j$, so $e^9>4b$ for all $j\le12$ here. After widening each 128-bit moment ball by both errors, the algorithm retains two adjacent dyadic numerators with the indicated opposite signs. Uniqueness identifies the enclosed transition. The same test gives the following sample of numerators $v_n$, where $v_n/2^{24}<-T_n<(v_n+1)/2^{24}$.
\[
 \begin{array}{c|rrrrrr}
 n&0&1&2&3&5&10\\ \hline
 v_n&100358703&147315867&194194744&241038735&334678586&568678602
 \end{array}
\]
All eleven enclosures for $0\le n\le10$ are in \texttt{level80-negative.json}. In particular the entire shifted quadratic row remains strictly hyperbolic even throughout $T_0<T<0$, although the deformed completed function already has nonreal zeros there.
\end{proof}

\clearpage
\phantomsection\addcontentsline{toc}{part}{IV. Finite zero divisors and arithmetic reconstruction}
\noindent{\small\sffamily\color{accent}IV. Finite zero divisors and arithmetic reconstruction}\par
\section{Counting every zero in a finite rectangle}\label{num:chapter}
Consider the normalized weight-two newforms listed below in Theorems 1 and 2 of Y.~Martin and K.~Ono, \emph{Eta-quotients and elliptic curves}, Proc. Amer. Math. Soc. \textbf{125} (1997), 3169--3176. Each is the Mellin form of a specified rational elliptic curve. We also treat Ramanujan's weight-twelve form $\Delta$ of level one. Our method counts all complex zeros in a rectangle by an argument change and produces as many disjoint sign changes on its symmetry line. Both counts include a boundary check. The actual intervals and their independent replay belong to the accompanying archive. Booker's \emph{Artin's conjecture, Turing's method and the Riemann hypothesis}, Experiment. Math. \textbf{15} (2006), 385--407, provides background on rigorous Fourier methods for counting zeros. Here the eta formulas allow a direct contour count with no Turing remainder.

For the formulas below $d^b$ in the second column indicates the factor $\eta(d\tau)^b$. A plain $d$ indicates exponent one. The level-one form is included in the same notation. The last four lines are the cases in which taking an absolute value of the product requires care.
\begin{center}\small
\begin{tabular}{rll}
\toprule Level & Eta factors & Weight\\\midrule
1   & $1^{24}$ &12\\
11  & $1^2\,11^2$ &2\\
14  & $1\,2\,7\,14$ &2\\
15  & $1\,3\,5\,15$ &2\\
20  & $2^2\,10^2$ &2\\
24  & $2\,4\,6\,12$ &2\\
27  & $3^2\,9^2$ &2\\
32  & $4^2\,8^2$ &2\\
36  & $6^4$ &2\\
48  & $4^4\,12^4/(2\,6\,8\,24)$ &2\\
64  & $8^8/(4^2\,16^2)$ &2\\
80  & $4^6\,20^6/(2^2\,8^2\,10^2\,40^2)$ &2\\
144 & $12^{12}/(6^4\,24^4)$ &2\\\bottomrule
\end{tabular}
\end{center}

\paragraph{A theorem with an effective witness.}
For the following finite-height result the ordinate is in the arithmetic variable $s$, with no scaling of its imaginary part. Denote by $M_f(T)$ the number of zeros of the completed $L$-function with $0<\Im s\le T$, counted with multiplicity. The table following the proof records the heights and counts certified by the published programs.

\begin{theorem}\label{num:thm}
For each form and each height in the certificate table, every nontrivial zero with $|\Im s|\le T$ lies on $\Re s=k/2$ and is simple. There is no real zero of the completed function, even outside this rectangle. For positive ordinates the listed intervals, one per zero, exhaust the zeros without overlap in their interiors.
\end{theorem}

\subsection*{Mellin symmetry and absence of real zeros}
Use $q=e^{2\pi i\tau}$ and $\eta(\tau)=q^{1/24}\prod_{n\ge1}(1-q^n)$. In each row $f(\tau)=\prod_d\eta(d\tau)^{b_d}$ with $\sum d b_d=24$, $\sum b_d=2k$, and $b_{N/d}=b_d$. Set $c=k/2$, $A=\sqrt N/(2\pi)$, and
\[
 K(u)=e^{cu}f(ie^u/\sqrt N),\qquad
 \Lambda(s)=A^s\Gamma(s)L(f,s).
\]
The eta transformation $\eta(-1/\tau)=(-i\tau)^{1/2}\eta(\tau)$, using its branch positive on the imaginary axis, exchanges $d$ and $N/d$. Direct multiplication gives $K(-u)=K(u)$. This identity continues to $|\Im u|<\pi/2$, where the eta arguments stay in the upper half-plane. For real $u$, $K(u)>0$; all eta factors on the positive imaginary axis are positive. The leading $q$ gives double exponential decay at positive infinity and evenness gives the same at the other end. Mellin inversion first in the absolutely convergent half-plane, followed by analytic continuation, gives the entire function
\begin{equation}\label{num:mellin}
 \Lambda(c+z)=\int_{\R}K(u)e^{zu}\,du.
\end{equation}
It is positive for every real $z$. At good primes the roots of the arithmetic Euler polynomial have modulus $p^{c-1/2}$. This is Hasse for elliptic curves and Deligne for $\Delta$. The normalized bad roots are bounded by one. The Euler product and reflection therefore confine all zeros to $|\Re s-c|\le1/2$. We count in the wider strip $|\Re s-c|\le1$ so that the right-hand argument is particularly simple.

\subsection*{First enclosure, using the real axis}
For the first nine rows $N\le36$ and $b_d\ge0$. Given an integer $T\ge1$, take $S=\lceil T/2\rceil+24$, $h=1/S$, and the least integer $U\ge3$ with $e^U>6T+360$. Write $J=US$, $w_j=hK(jh)$, and form
\[
 P(z)=w_0+\sum_{j=1}^J w_j(e^{jhz}+e^{-jhz}).
\]
For $|\Re z|\le1$, $|\Im z|\le T$ we have
\begin{equation}\label{num:bound}
 |\Lambda(c+z)-P(z)|\le
 \frac{e^{T+62}}{e^{2\pi S}-1}+184320\exp(-e^U/2).
\end{equation}
Here is a derivation using only bounds on the eta product. On $u=x\pm i$ put $v=e^{|x|}$. By evenness it suffices to use the positive real part of $u$. Then $|q|\le e^{-v/2}$ since $\sqrt N\le6$ and $\cos 1>1/2$. The product after its leading $q$ has absolute value at most $\exp(24|q|/(1-|q|))\le e^{48}$. This follows by summing $\log|1-q^{dn}|\le |q|^{dn}$ and using $\sum b_d\le24$. On either displaced line the integral of the absolute value of $K(u)e^{zu}$ is at most
\[
 2e^{48+T}\int_0^\infty e^{7x}e^{-e^x/2}\,dx
 \le 2e^{48+T}2^7\Gamma(7).
\]
Apply Poisson summation with step $h$ to $u\mapsto K(u)e^{zu}$ on the real axis. Its Fourier transform at frequency $2\pi l/h$ can be bounded by moving the integral a unit vertically in the direction of decay. Analyticity and the double-exponential estimate justify closing the finite rectangles and letting their sides tend to infinity. Summing the nonzero frequencies costs $2/(e^{2\pi S}-1)$, giving the first term in \eqref{num:bound}. On the real axis the remaining eta factors are at most one, and $|q|\le e^{-v}$. Beyond $U$, $e^{(c+1)x-e^x}$ decreases. The two omitted tails of the sum are bounded by
\[
 2\int_{e^U}^{\infty}v^ce^{-v}\,dv
 \le 2e^{-e^U/2}2^{c+1}\Gamma(c+1)
 \le184320 e^{-e^U/2}.
\]
This proves the bound without approximating a completed function by a floating-point point value.

\subsection*{Second enclosure, moving the integration line}
At height $T$ the gamma factor on the line is of size roughly $T^{c-1/2}e^{-\pi T/2}$. The first method subtracts quantities of ordinary size to find a result that small. A contour displacement recovers much of the exponential factor before computing. This method handles every row, including the eta quotients with negative exponents.

Take $\delta=\min(1/2,64/T)$ and $\alpha=\pi/2-2\delta$. For $z=x+it$, $|x|\le1$, $0\le t\le T$, set
\[
 G(z)=e^{-i\alpha z}\Lambda(c+z)
   =\int_{\R}K(u+i\alpha)e^{zu}\,du .
\]
Indeed this is \eqref{num:mellin} shifted up by $i\alpha$. On the critical line $G(it)=e^{\alpha t}\Lambda(c+it)$ has the same sign as the original function. Let $\lambda=\delta/4$ and
\[
 B=\frac{2c!}{(1-e^{-\lambda})^{c+1}},
 \qquad I=2B c!\lambda^{-(c+1)}.
\]
For $|b-\alpha|\le\delta$ one has $2\pi\cos b/\sqrt N\ge\lambda$. All the forms here are normalized Hecke eigenforms and the coefficient bound $|a_n|\le d(n)n^{c-1/2}\le2n^c$ follows from their local roots. Hence for $v=e^{|u|}$, evenness and the $q$-series imply
\[
 |K(u+ib)|\le B v^c e^{-\lambda v}.
\]
For this inequality use $\sum_{n\ge1}n^c r^n\le c!r/(1-r)^{c+1}$, obtained from the binomial series; the reflected height $-b$ has the same cosine. In particular the integral of the right side times $e^{|u|}$ is at most $I$. This proof controls denominators in the eta quotients without differentiating a reciprocal near a cusp.

Given integers $S,U>0$ with $\lambda e^U>c+1$, put $h=1/S$, $J=US$, and
\[
 v_j=hK(jh+i\alpha),\quad
 P_{\alpha}(z)=v_0+\sum_{j=1}^J(v_je^{jhz}+\overline{v_j}e^{-jhz}).
\]
Evenness and real Taylor coefficients give the conjugate weights at negative indices. Poisson summation on lines at distance $\delta$ on either side of the displaced line gives
\begin{equation}\label{num:tilt}
 |G(z)-P_{\alpha}(z)|\le
 E=\frac{2I e^{\delta T}}{e^{2\pi\delta S}-1}
        +2^{c+1}I e^{-\lambda e^U/2}.
\end{equation}
The second term follows by monotone comparison with the two tail integrals and by retaining half the exponential decay for the factor at $e^U$. For the first term the change in $e^{zu}$ under the small shift has modulus at most $e^{\delta T}$. These same estimates justify the initial displacement from zero to $\alpha$. The implementation chooses integers making each term smaller than $e^{-205}(T+3)^{-5}$ and explicitly checks the inequalities. Their validity is independent of how these integers were proposed. For large $T$ the cost in precision now comes from a fixed loss $e^{-128}$ and powers of $T$, instead of $e^{-\pi T/2}$.

\subsection*{Winding along the upper edge}
The two enclosures have the same form, with $\alpha=0$ for the first. Write $z=x+iT$ and expand the finite exponential polynomial at $x=1/2$. Set $M=|v_0|+2\sum_{j=1}^J|v_j|$. On $0\le x\le1$ the Taylor remainder of degree $D$ is at most
\[
 M e^U (U/2)^{D+1}/(D+1)! .
\]
Apply the real integral remainder to each exponential in $x$, bounding its derivative anywhere on the interval by $(jh)^{D+1}e^{jh}$ times the modulus of its weight. The coefficients, including the contributions from the negative indices, are explicitly
\[
 [y^n]P_\alpha(1/2+y+iT)
 =\mathbf1_{n=0}v_0+
 \sum_{j=1}^J \frac{(jh)^n}{n!}
 \left(v_je^{jh(1/2+iT)}+(-1)^n\overline{v_j}e^{-jh(1/2+iT)}\right).
\]
The verifier evaluates this polynomial by interval Horner arithmetic at dyadic subintervals of $[-1/2,1/2]$. It adds both the quadrature bound and the Taylor bound. If the enclosing rectangle lies in one of the four open half-planes cut out by the axes, the curve avoids zero and its argument change on that segment equals the principal argument of the quotient of its endpoints, taken in the direction of traversal. The segment is subdivided until both this test and the strict endpoint argument bound $<\pi$ succeed. An unresolved segment aborts the proof. Let $H$ be the sum of the certified changes for $G$ from $x=1$ to $x=0$. The corresponding change for $\Lambda$ is $H-\alpha$.

On the right edge $\Re s=c+1$ the continuous argument beginning with zero at the bottom equals
\[
 \theta(T)=\Im\{s\log A+\log\Gamma(s)\}+\arg L(f,s),\qquad s=c+1+iT.
\]
The logarithm of gamma is the analytic branch real on the positive axis. The principal argument of the Euler factor product is the correct lifted argument here: the absolutely convergent logarithm satisfies
$|\Im\log L(f,s)|\le2\log\zeta(3/2)<\pi$, using $\zeta(3/2)<3$ by integral comparison. The top right endpoint enclosure divided by the appropriate gamma and contour factors supplies that principal argument without needing an approximate Euler product.

Traverse the upper rectangle counterclockwise, first from its real-axis midpoint to its lower right corner and then along the right edge and half the top. Reflection in the vertical symmetry line conjugates the image; reversing the path on the reflected half gives the same argument change. The entire bottom edge has positive image. The argument principle thus gives
\begin{equation}\label{num:wind}
 M_f(T)=\frac{\theta(T)+H-\alpha}{\pi}.
\end{equation}
All sides have been proved nonzero, including the top midpoint. The computed interval on the right must contain exactly one integer. The proof requires this integer test; estimating the nearest integer from digits would lose the guarantee.

\subsection*{Exhaustion on the symmetry line}
For an even grid length $m>2J$, the discrete Fourier transform of the weights gives $P_{\alpha}(it_l)$ at $t_l=2\pi S l/m$. The convention of the software uses a negative exponential; the positive-index weight is placed at index $-j\pmod m$. In the real-axis method the two weights agree. Compare the real part at each mesh point with the rigorous bound $E$ or \eqref{num:bound}. Save a sign only when its enclosing interval and error are strictly on one side of zero. Signs begin at $t=0$, positive by the Mellin integral. Include the separately certified endpoint $T$. Each successive opposite pair supplies a distinct zero in its open interval by continuity. Endpoints can be shared by two closed brackets, but no endpoint is a zero. Equality with \eqref{num:wind} proves that every zero in the upper rectangle is on the line and occurs with multiplicity one. Complex conjugation gives the same assertions below the real axis. This proves Theorem~\ref{num:thm} upon checking the reported interval runs.

\paragraph{Arithmetic behind a run.}
The scripts use FLINT/Arb real and complex balls. The errors in \eqref{num:bound} and \eqref{num:tilt} are added to the outward-rounded arithmetic enclosures. Elementary functions, gamma, the eta function, and the discrete Fourier transform are evaluated with enclosures. The input consists of the displayed finite list of integers and the chosen height. Intermediate precision and the proposed transform length affect whether a run succeeds, not whether a successful inequality can be false. The JSON records store the integer zero count, top sign, transform length, grid scale, and bracketing indices. They can be regenerated by replaying the program. For an independent check at low height, the archive also evaluates sign endpoints by a paired incomplete-gamma expansion with integer Fourier coefficients. See F.~Johansson, \emph{Arb: efficient arbitrary-precision midpoint-radius interval arithmetic}, IEEE Trans. Comput. \textbf{66} (2017), 1281--1292.


\subsection*{Certifying a narrow pair of zeros between Fourier mesh points}
The larger-height runs can use several smaller transforms or a local interpolation. Both changes preserve the proof above. For a transform of total period $m=bl$ in grid indices, restrict first to indices $k=r+lj$. Rotating the coefficient at signed frequency $a$ by $e^{2\pi iar/m}$ leaves the factor $e^{2\pi iaj/b}$. Coefficients may now be added modulo $b$, and a length-$b$ transform evaluates all ordinates in that residue class. Repeating over $r$ gives exactly the old transform. This reordering bounds the memory for one transform.

For interpolation let $\Delta=2\pi S/m$ and $P(t)=P_\alpha(it)$. This is real for real $t$. Its $n$th derivative on the real axis has absolute value at most $MU^n$ with $M$ as above. Suppose a point has index $k+\theta$, $0<\theta<1$. Interpolate from the $2l$ consecutive values with indices $k+1-l,\ldots,k+l$. Since the product of the $l$ node distances on either side of $\theta$ is at most $l!$ in grid units, the classical real interpolation remainder gives
\[
 \left|P((k+\theta)\Delta)-
   \sum_{j=1-l}^{l}P((k+j)\Delta)
        \prod_{\substack{i=1-l\\i\ne j}}^{l}\frac{\theta-i}{j-i}\right|
 \le M(U\Delta)^{2l}\frac{(l!)^2}{(2l)!}.
\]
The remainder is for the finite trigonometric polynomial itself. One adds the quadrature error $E$ once after the interpolation. Mesh-value enclosures can be stored as pairs of binary floating-point endpoints once their containment has been checked against the original balls. The interpolator multiplies the outward endpoints by interval enclosures of the signed rational Lagrange weights. Bounds are used only when they exclude zero strictly.

A numerical search proposes a dyadic $\theta$ near a possible extremum when two zeros may lie between coarse mesh points. The theorem places no reliability requirement on this search. A proposed sign is used only following the rigorous test, with its nearest known neighbors retained, and the final ordered sign sequence must exhaust the independent winding count. If the search overlooks a pair, the equality fails and no certificate is issued for that run. For the interpolated records the reported integer mesh incorporates the dyadic subgrid denominator. The same bracket-to-height conversion continues to apply.

\paragraph{An independently formalized counting step.}
The archive includes a Lean proof of the following finite implication: if every disjoint bracket contains at least one positive zero multiplicity, the complement contains a list of positive multiplicities, and the sum of all multiplicities equals the number of brackets, then each bracket contains exactly one simple zero and the complement is empty. It also checks the integer bracket order for one concrete run. The analytic enclosures and the argument principle provide the hypotheses of that lemma.

At the common height $65\,536$ the thirteen forms have $2\,672\,374$ positive zeros in total. Every one is simple and on its own critical line. The kernels exclude real zeros and the coefficients are real, so the total for both signs, counting one divisor for each form, is $5\,344\,748$. The Gaussian contour totals at this common height and the line sign brackets are independent pieces of the proof.

\paragraph{Certificate table.} The two fixed columns report exact counts through the indicated height. The last two columns report the strongest completed run for that row. A dash means that the indicated height was not run as a separate certificate. A certificate at greater height covers the entire smaller closed rectangle as well.
\begin{center}\small\begin{tabular}{rrrrr}\toprule
Level & $M_f(256)$ & $M_f(1024)$ & Highest $T$ & $M_f(T)$\\\midrule
1 & 224 & 1337 & 131072 & 373228\\
11 & 319 & 1725 & 524288 & 1924341\\
14 & 328 & 1765 & 65536 & 199679\\
15 & 331 & 1776 & 65536 & 200399\\
20 & 342 & 1823 & 65536 & 203400\\
24 & 350 & 1853 & 65536 & 205302\\
27 & 355 & 1872 & 65536 & 206530\\
32 & 362 & 1899 & 131072 & 445523\\
36 & 367 & 1919 & 65536 & 209531\\
48 & 379 & 1965 & 65536 & 212531\\
64 & --- & 2012 & 65536 & 215532\\
80 & --- & 2049 & 65536 & 217859\\
144 & --- & 2144 & 65536 & 223990\\
\bottomrule\end{tabular}\end{center}
For example, the total number through both signs of height is $2M_f(T)$ since the center and the whole real axis are zero-free for these forms. Each integer in the table is extracted directly from the replayable record. A bracket $[a,b]$ in a JSON file means endpoints $2\pi S a/m$ and $2\pi S b/m$; the string \texttt{T} specifies the exact top height.

\clearpage
\section{An independent zero count from the arithmetic coefficients}\label{gauss:chapter}
At large height, a direct computation of the Mellin integral on the positive real axis loses the exponential factor in the gamma function. We give a second high-ordinate count that starts with the integer Fourier coefficients. It uses a smoothed approximate functional equation and retains a polynomial enclosure for a whole horizontal edge, or for a short two-dimensional rectangle. Thus a numerical zero total can be checked against a contour calculation whose modular kernel and sampling variable differ from those in Section~\ref{num:chapter}.

Let $k=2c$, $\lambda=2\pi/\sqrt N$ and
\[
       \Lambda(s)=\lambda^{-s}\Gamma(s)\sum_{n\ge1}a_n n^{-s}.
\]
In this section the newform has real coefficients and positive root number, and is one of the thirteen forms of Section~\ref{num:chapter}. The normalization makes $\Lambda(c+it)$ real. At positive height $t$ let $D_t=\lambda^{-c}|\Gamma(c+it)|$. This positive divisor converts an exponentially small number into an ordinary-size real number without changing its sign.

\subsection*{Two shifts of one integral}
Put $s_0=c+it$, choose $a=3$, $R=12$, and set $g(z)=\exp(z^2/R^2-i\pi z/2)$. For $|\Re w|<a$ define
\[
 U(w)=\frac{1}{2\pi D_t}\int_{\R}
       \Lambda(s_0+a+iv)\frac{g(a+iv-w)}{a+iv-w}\,dv .
\]
The smoothed approximate functional equation in this normalization is
\begin{equation}\label{gauss:afe}
        \Lambda(s_0+w)/D_t= U(w)+\overline{U(-\bar w)} .
\end{equation}
Shift the right contour to its reflection about $\Re s=c$; the sole residue is the value at $s_0+w$. On the reflected integral use $\Lambda(u)=\Lambda(2c-u)$, real Dirichlet coefficients, and $\overline{g(-\bar z)}=g(z)$. The horizontal integrals disappear because of the Gaussian. This gives \eqref{gauss:afe}. This is the classical smoothed approximate functional equation; see M.~Rubinstein, \href{https://arxiv.org/abs/math/0412181}{\emph{Computational methods and experiments in analytic number theory}}. What requires verification here is the size of its error at the advertised ordinates.

Write $z=b+iv$ with real $b$ and strip off the factor
$e^{w^2/R^2+i\pi w/2}$. In the resulting first integral the $w$-dependent factor relative to $w=0$ is
\[
                 \frac{e^{-2zw/R^2}}{1-w/z}.
\]
The reflected integral has the analytically conjugate coefficients, with alternating signs. On the imaginary $w$-axis the stripped factor is positive and on the real $w$-axis its phase is $\pi w/2$.

We first truncate the series after $P$ terms. Although the retained integral is computed on $b=3$, the integral of the discarded tail can be shifted to $b=32$: its Dirichlet series is absolutely convergent between these lines, $|w|\le1$, and there are no poles. The Hasse--Deligne bound gives $|a_n|\le d(n)n^{c-1/2}\le2n^c$ and hence
\[
    \sum_{n>P}|a_n|n^{-c-b}\le \frac{2P^{1-b}}{b-1}.
\]
This independent shift allows a manageable truncation even when $t$ is large.

\subsection*{Explicit bounds for the quadrature}
For an integer $b\ge2$, put $u=b+c$ and
\[
 C_b=\frac{(t/\lambda)^b}{b-1}
 \exp\!\left(\frac{b^2+2b}{R^2}
           +\frac{u(u+1)}{t}
           +\frac12\left(\frac{uR}{t}+\frac2R\right)^2\right).
\]
After removing the series and the factor $1/(2\pi)$, the normalized, stripped integrand is bounded by $C_b e^{-v^2/(2R^2)}$ on $|w|\le1$. To verify this, for $y=t+v$ use
\[
 e^{\pi|y|/2}|\Gamma(1+iy)|
     =\left(\frac{2\pi|y|}{1-e^{-2\pi|y|}}\right)^{1/2}
     \le (1+2\pi|y|)^{1/2},
\]
with the continuous value at zero. The denominator satisfies
$e^{\pi t/2}|\Gamma(c+it)|\ge t^{c-1}\sqrt{2\pi t}$ by recurrence. Recurrence in the numerator gives
\[
  e^{\pi v/2}\frac{|\Gamma(c+b+i(t+v))|}{|\Gamma(c+it)|}
      \le t^b\exp\!\left(\frac{u(u+1+|v|)}{t}\right).
\]
Indeed the $u-1$ linear factors and the remaining square-root ratio are bounded after division by their powers of $t$ by $(1+(u+1+|v|)/t)^u$. Use $|z-w|\ge b-1$, $|e^{-2zw/R^2}|\le e^{2(b+|v|)/R^2}$, and complete the square in $|v|$ to obtain the bound above. The Gaussian integral is at most $I=2R\sqrt\pi$. Including both integrals in \eqref{gauss:afe}, the truncation allowance is
\[
             E_1=\frac{2P^{1-b}I C_b}{(b-1)\pi},\qquad b=32.
\]

For the retained integrals use $v=jh$, $|j|\le J$, $h=1/16$, $Jh=180$. The strip trapezoidal theorem on the two boundary lines $b=2,4$ gives the error
\[
 E_2=\frac{9I(C_2+C_4)}{\pi(e^{2\pi/h}-1)}.
\]
Here the sum of the absolute values of the finite Dirichlet series is bounded by $\zeta(b+1/2)^2<9$ throughout the strip. The omitted samples cost at most
\[
 E_3=\frac{18h C_3 e^{-(Jh)^2/(2R^2)}}
               {\pi(1-e^{-(2J+1)h^2/(2R^2)})}.
\]
This follows by comparison with a geometric series beyond $J$. We may take $P=\lceil4(T+32)/\lambda\rceil+32$ where $T=t$ for a contour center and $T$ is the enclosing certificate height for a sampled endpoint at the ordinates reported here. Both the sufficiency of this choice and the three error bounds are tested anew for each run. For signs at $w=0$ alone the program uses $h=1/12$ on boundary lines $b=1,5$. In the corresponding constants $b-1$ is replaced by $b$, the term $2b/R^2$ is deleted, and $2/R$ is deleted inside the square. The strip half-width for this calculation is two.

\subsection*{The two finite polynomials}
Let $z_j=3+ijh$. Put
\[
 W_j=\frac{h}{2\pi}\,
     \frac{\Gamma(s_0+z_j)}{|\Gamma(s_0)|}\lambda^{-it-z_j}
                    \frac{g(z_j)}{z_j}.
\]
Reconstruct $a_n$ as exact integers by Euler's pentagonal formula for the positive and negative eta exponents separately, and divide the resulting power series in $\Z[[q]]$. An independent logarithmic-derivative recurrence checks the first 96 coefficients. For the finite sum $F_j=\sum_{n\le P}a_n n^{-c-3-it-ijh}$, choose $M=32768$ and round $h\log n$ to a multiple $2\pi m/M$. The program verifies every residual $|\delta_n|=J|h\log n-2\pi m/M|<G=1.001\pi J/M$ before using it. Expand $e^{-ij\delta_n/J}$ to degree $K$. The relative phase approximation has absolute error at most $e^G G^{K+1}/(K+1)!$. Grouping by $m$ makes the finite trigonometric sums discrete Fourier transforms. Since $\sum |a_n|n^{-c-3}<2$, and
\[
 \left|\frac{e^{-2z_jw/R^2}}{1-w/z_j}\right|
 \le \frac32 e^{2(3+Jh)/R^2}\qquad (|w|\le1),
\]
the total Fourier approximation error for both integrals is at most
\[
  E_4=6 e^{2(3+Jh)/R^2}
             \left(\sum_j |W_j|\right)e^G G^{K+1}/(K+1)!.
\]
Choose $K$ so that this is less than $10^{-13}$. The simpler sign checker can first transform the $W_j$ and then apply the same phase expansion term by term in $n$. This reverses two finite sums.

Finally expand $e^{-2z_jw/R^2}/(1-w/z_j)$ at $w=0$. Its modulus on $|w|=2$ is at most $3e^{4(3+Jh)/R^2}$. For the approximations $\widetilde F_j$ just computed, Cauchy's estimate bounds the combined tail beyond degree $D$ on $|w|\le1$ by
\[
 E_5=6e^{4(3+Jh)/R^2}\!
                  \left(\sum_j|W_j\widetilde F_j|\right)2^{-D}.
\]
Take $D$ so that $E_5<10^{-13}$. If the coefficients of the first retained polynomial are $b_m$, those of the total are $b_m+(-1)^m\bar b_m$. This gives one interval polynomial for the entire disk, with absolute error $E_1+\cdots+E_5$. All gamma evaluations, arithmetic of coefficients, Fourier transforms and interval comparisons in this second calculation are performed with directed enclosures. A floating-point bin choice has no effect on correctness after the interval test for its residual.

For a full top-edge count put $w=x$, $0\le x\le1$, and apply the argument computation of \eqref{num:wind} with phase correction $\pi/2$. For a localized count use the rectangle $|\Re w|\le17/32$, $|\Im w|\le3/4$. Its corners lie strictly inside the unit disk. At each segment the interval polynomial is enclosed in an open coordinate half-plane before adding the principal change of argument. Along $w=i v$ the same polynomial has real coefficients when written in $v$. Its sign brackets must exhaust the closed-rectangle count for the localized assertion of simplicity. The strip of all possible nontrivial zeros is contained between the vertical edges by the Euler product and functional equation.

The level-$11$ check through $4096$ recovers exactly $8\,708$ positive zeros. Its last bracket has normalized endpoint values in $(-0.142,-0.141)$ and $(0.198,0.199)$. At level $32$ the last bracket below $65\,536$ has endpoints $t_-,t_+$ satisfying
\[
 65535.8080<t_-<65535.8081,\qquad
 65535.8519<t_+<65535.8520 .
\]
Their normalized values lie in $(-0.164,-0.163)$ and $(0.055,0.056)$ respectively. These signs check the original brackets. Section~\ref{num:chapter} supplies the full sign exhaustion.

\begin{theorem}[Seven simple zeros near one million]\label{gauss:million}
For the conductor-$11$ curve, every nontrivial zero with
$|\Im s-10^6|\le3/4$ is simple and has real part one. There are exactly seven. Their ordinates belong, one each, to the intervals
\[
             10^6+(j/64,(j+1)/64),\qquad
             j=-42,-25,-16,11,15,44,46 .
\]
\end{theorem}
\begin{proof}
Apply the disk construction at $t=10^6$, using $2\,111\,529$ integer Fourier coefficients. The disk error is less than $6.33\cdot10^{-12}$. The directed contour winding lies in $(7-10^{-7},7+10^{-7})$. Strictly opposite signs at the endpoints give seven different zeros on the symmetry line. The argument principle counts seven with multiplicity in the entire rectangle. Thus each is simple and there are no other zeros in this band. The output and all dyadic endpoints are in \texttt{block-11-1000000.json}. The gap from lower certified windows to this particular band is outside this theorem.
\end{proof}

\begin{proposition}[Ramanujan's cusp form at the same ordinate]\label{gauss:milliondelta}
Every nontrivial zero of \(L(\Delta,s)=\sum_{n\ge1}\tau(n)n^{-s}\) with
\(|\Im s-10^6|\le3/4\) lies on \(\Re s=6\) and is simple. The six ordinate brackets have the form in Theorem~\ref{gauss:million}, with
\[
                 j=-46,-34,-7,13,30,37 .
\]
\end{proposition}
\begin{proof}
Use the weight-twelve gamma factor and \(636\,673\) exact coefficients. On the unit disk the error is less than \(1.96\cdot10^{-12}\). For this run the polynomial enclosure on each boundary segment is recentered: at its midpoint \(m\), on a radius \(h\), the variation is at most
\[
  \sum_{j=1}^{5}\frac{|P^{(j)}(m)|h^j}{j!}
      +h^6\sum_{n\ge6}\binom n6|[z^n]P|.
\]
Taylor's integral remainder gives the last term because the entire segment is inside the unit disk. All seventy segment enclosures exclude the origin. The winding lies between \(6-4\cdot10^{-11}\) and \(6+4\cdot10^{-11}\). Each of the six real axis intervals gives a strict sign change. Their number equals the winding. The complete record is \texttt{block-1-1000000.json}.
\end{proof}

At small nonnegative center heights one can avoid negative powers of $t$ in the preceding quadrature bound. Put $u=b+c$ and $X=t+u+1+R\sqrt u$. The same gamma recurrence gives
\[
 C_b^{\rm low}=\frac{3e^{-\pi t/2}\lambda^{-b}X^{u-1/2}}
 {(b-1)|\Gamma(c+it)|}
 \exp\!\left(\frac{b^2+2b}{R^2}
             +\frac12\left(\frac{uR}{X}+\frac2R\right)^2\right)
\]
in place of $C_b$. Indeed $\sqrt{1+2\pi|t+v|}\le3\sqrt{1+t+|v|}$ and each of the $u-1$ recurrence factors is at most $1+t+u+|v|\le Xe^{|v|/X}$. This bound works even at $t=0$. The consecutive-unit-strip replay uses these low-height constants for $t<128$ and the original ones thereafter. Neighboring closed bands meet only at their certified zero-free boundary. Each sign-change count is checked against a full local winding, independently of any lower-height accumulated count.

\subsection*{A cyclotomic field at high ordinate}
Write \(E=11.\mathrm{a1}\) and \(K=\Q(\zeta_5)\). Its four rational factors are the base newform and its twists by the quadratic and two quartic characters modulo \(5\), as proved in the chapter on non-self-dual twists. For a factor with complex coefficients write \(r\) for its unit-modulus phase, so that \(r\Lambda(s)=\overline{r\Lambda(2-\bar s)}\). For the character with \(\chi(2)=i\) one has \(r=\overline{\tau(\chi)}/\sqrt5\), for its dual take \(\bar r\), and for the other two \(r=1\). In \eqref{gauss:afe} replace the first integral by the integral for \(r\Lambda\). Its reflected integral is then its conjugate at \(-\bar w\). Thus multiply each Dirichlet coefficient by its character and the first polynomial by \(r\) before forming \(b_m+(-1)^m\bar b_m\). Every absolute bound in this chapter still holds; twisting only decreases the available coefficient majorant.

\begin{theorem}\label{gauss:field}
Every zero of \(L(E/K,s)\) in the band \(|\Im s-100\,000|\le3/4\) is simple and has real part one. There are \(22\) zeros. The four factors, in the order base, quadratic, quartic with \(\chi(2)=i\), and dual, contribute \(6,5,5,6\), respectively. By conjugation the band about \(-100\,000\) has the same total.
\end{theorem}
\begin{proof}
Perform the disk winding separately at conductor \(11\) for the base and conductor \(275\) for the twists. In each case strict sign brackets exhaust the winding, as in Theorem~\ref{gauss:million}. Bisect each proposed bracket, retaining a half only if its endpoint signs are certified. At a common denominator \(2^{24}\) the distance between successive intervals, including those from different factors, is at least \(74147/2^{24}\). The product factorization therefore counts \(22\) simple zeros. The arithmetic and all intervals are recorded in \texttt{field-100000.json}. At the lower test ordinate \(192\), all eleven signed brackets intersected their independent Mellin brackets in the height-\(4096\) lists. Where a coarse bracket straddled the test boundary, an independent upper incomplete gamma evaluation established the side of the cut.
\end{proof}

\begin{theorem}[Cyclotomic base change at height one million]\label{gauss:fieldmillion}
For $E=11.\mathrm{a1}$ over $K=\Q(\zeta_5)$ the band $|\Im s-10^6|\le3/4$ contains exactly $26$ zeros. Every one is simple and has real part one. The base, quadratic, quartic and conjugate quartic factors contribute $7,6,7,6$ of them. The genuinely quadratic isogeny class over $\Q(\sqrt5)$ in Corollary~\ref{qt:genuine} therefore has exactly $13$ simple zeros in this band. Reflection gives the same totals at negative height.
\end{theorem}
\begin{proof}
Use the same four polynomial enclosures as in Theorem~\ref{gauss:field}. The base requires $2\,111\,529$ coefficients and the three twists each require $10\,557\,515$. To reconstruct the level-eleven coefficients exactly, put $P(q)=\prod_{j\ge1}(1-q^j)$ and write the newform as $qP(q)^2P(q^{11})^2$. Euler's pentagonal identity gives $P$ modulo $q^M$ for $M=10\,557\,515$. Work first modulo the prime $p=2\,013\,265\,921$. Zero-pad the coefficient list to length $2^{25}\ge2M$ and square by a number-theoretic Fourier transform. Truncate back to $M$, dilate a copy by $q\mapsto q^{11}$ and multiply with the same padding. Thus reduction modulo $x^{2^{25}}-1$ causes no wraparound in either truncated product. The program trial-divides the prime and checks that the successive roots constructed from $31$ have the asserted power-of-two orders and that the inverse scale multiplies to one. The divisor and Hasse bounds give $|a_n|\le d(n)\sqrt n\le2n<p/2$ through the largest index. Each signed integer coefficient is therefore the unique lift in the allowed range.
 
An independent check certifies every proposed coefficient by finite-field arithmetic. For $p>29$ the change $X=36x-12$, $Y=216y+108$ carries the minimal model $y^2+y=x^3-x^2-10x-20$ to $Y^2=X^3-13392X-1080432$. Write $n=p+1-a_p$ and check $|a_p|\le2\sqrt p$. If a point $P$ on the latter curve satisfies $[n]P=O$ and $[n/q]P\ne O$ for a prime $q\mid n$, then $q^{v_q(n)}$ divides its order. Multiply the prime-power divisors certified by any such points, using each prime once, to obtain $D\mid n,\#E(\mathbb F_p)$. As soon as $D^2>16p$, these multiples of $D$ in the Hasse interval of length $4\sqrt p$ must coincide. The same argument on a quadratic twist tests $2(p+1)-n$, since its trace is the negative of the original trace. Direct point counts handle the rest. Of the $699\,164$ primes, $678\,684$ were certified on the curve, $20\,215$ on its twist, and $265$ by complete enumeration, including $p=11$. The Hecke recurrence at every composite, with the degree-one factor at $11$, reconstructs all $9\,858\,350$ composite coefficients. The two complete integer files coincide byte for byte. As a third check, PARI/GP computes the entire list directly from the original integral model by \texttt{ellan} and independently verifies its prime entries by \texttt{ellap}. All $10\,557\,515$ of its coefficients agree. The code, counts and checksum are in the coefficient audit accompanying \texttt{field-1000000.json}.

Each factor has a boundary winding equal to its number of strict line brackets. In the indicated factor order the four contours required $82,118,114,128$ certified boundary segments. On refinement to denominator $2^{24}$ the least gap between adjacent brackets, including different factors, is at least $14631/2^{24}$. This proves simplicity of every product zero. All brackets are recorded in \texttt{field-1000000.json}. The two quartic factors give precisely the indicated quadratic-field curve by the complete factorization in Corollary~\ref{qt:genuine}.
\end{proof}

\paragraph{A consecutive-block comparison.}
For $11.\mathrm{a1}$ we also tiled the full interval of positive heights up to $1024.5$ by the independent local Gaussian calculation. The $1025$ closed rectangles have centers $j=0,1,\ldots,1024$ and half-height $1/2$; the initial rectangle covers heights of both signs around zero and contains no zero. Each local winding separately equals its number of strict sign brackets. The rectangle boundaries were enclosed using $52\,168$ contour segments in total. Their $1727$ positive brackets, in order, each have nonempty interior intersection with the matching tilted-Mellin bracket from the height-$4096$ list. The next Mellin bracket lies above the top boundary. This tests every individual zero of the low range by two representations of the completed function. The two consecutive-strip files and the exact comparison are identified in \texttt{match-11-1024.json}.

\paragraph{A CM factor at height one million.}
For the conductor-$32$ curve $E_0:y^2=x^3-x$, Gaussian ideals give an independent arithmetic input for an isolated high band. A primary Gaussian integer of odd norm has real part $a$ odd, imaginary part $b$ even and $a+b\equiv1\pmod4$. By Lemma~\ref{cm:primary}, each odd-norm ideal has exactly one such generator, and its contribution to the weight-two series is that generator. To compute the $n$th coefficient enumerate $a_0>0$ odd, $b\ge0$ even with $a_0^2+b^2=n$, choosing $a=a_0$ or $-a_0$ by the congruence. Add $2a$ if $b>0$ and $a$ if $b=0$. For even $n$ the coefficient is zero. Conjugation explains the factor two.

All $3\,601\,413$ coefficients used by the Gaussian contour at $H=10^6$ coincide with those of $\eta(4z)^2\eta(8z)^2$. For a further arithmetic check, at a split prime the prime decomposition $p=a^2+b^2$ was obtained by the Euclidean algorithm from a square root of $-1$ modulo $p$. Each returned decomposition and its primary congruence were tested by exact integer arithmetic. The split trace $2a$ and inert trace zero then give every odd coefficient by the Hecke recurrence. The resulting list agrees through the full range, including all $128\,318$ split and $128\,508$ inert primes. A separate calculation from the integral equation $y^2=x^3-x$ agrees at all $3\,601\,413$ indices by PARI/GP \texttt{ellan} and at all $256\,827$ primes, including $2$, by \texttt{ellap}. Exhaustive point counts at fourteen primes also agree. For example $\#E_0(\mathbb F_{3\,601\,369})=3\,598\,096$, giving $a_{3\,601\,369}=3274$.

The seven simple critical zeros at $|\Im s-10^6|\le3/4$ for $E_0$ exhaust an integer winding on $98$ contour segments. Their brackets are disjoint from all 26 earlier brackets for $11.\mathrm{a1}/\Q(\zeta_5)$. Therefore the five-dimensional abelian variety over $\Q$
\[
 E_0\ \times\ \operatorname{Res}_{\Q(\zeta_5)/\Q}
    (11.\mathrm{a1}_{/\Q(\zeta_5)})
\]
has exactly $33$ zeros of its $H^1$-function in this band. All are simple and on $\Re s=1$, with every pair separated by more than $14631/2^{24}$. Indeed restriction of scalars induces the corresponding Galois parameter, hence preserves its completed $L$-function. The product adds the fifth factor. Sorting all 33 disjoint intervals on the same integer lattice proves the separation.

\begin{theorem}[Eight simple CM zeros near ten million]\label{gauss:tenmillion}
For $E_0:y^2=x^3-x$ over $\Q$, exactly eight nontrivial zeros lie in $|\Im s-10^7|\le3/4$. Each is simple and lies on $\Re s=1$. The eight ordinates belong respectively to
\[
        10^7+(j/64,(j+1)/64),\qquad
        j=-44,-30,-20,-8,6,20,31,42.
\]
\end{theorem}
\begin{proof}
The first disk polynomial uses the Gaussian ideal sum for $36\,012\,800$ coefficients. The compressed eta product in $q^4$ independently agrees at every index through $4\,000\,000$. With $R=12$, $h=1/16$, $J=2880$ and transform size $32768$, the absolute disk error is below $1.068\cdot10^{-10}$. The winding interval lies in $(8-7\cdot10^{-10},8+7\cdot10^{-10})$. The eight disjoint signed brackets listed in the theorem exhaust this count.

PARI/GP independently computed the $2\,205\,005$ prime traces in this range from the integral curve. A segmented sieve checked the prime list and reconstructed every one of the $36\,012\,800$ coefficients by the local Hecke recurrences. The reconstructed integers agreed term by term with the Gaussian-ideal sum. Direct exhaustive point counts at the four largest primes in the odd classes modulo eight agreed as well.

For a second calculation construct $38\,263\,598$ integer coefficients from prime decompositions and Hecke recurrences. The signed 32-bit coefficient fingerprint over the entire overlap, including the initial zero at index zero, agrees with the Gaussian-ideal list. Use $224$-bit interval arithmetic, $R=14$, $h=1/20$, $J=4480$, and transform size $65536$. The disk, quadrature and truncation allowances follow the displayed bounds above at these parameters. The new disk error is below $2.673\cdot10^{-12}$, and the second winding is within $1.76\cdot10^{-11}$ of eight. Both calculations refine each of their eight strict sign brackets to the same pair of consecutive multiples of $2^{-24}$. This completes the zero count in this band. Conjugation gives the reflected band around $-10^7$.
\end{proof}

\clearpage
\section{From a finite rectangle to positive Li coefficients}\label{li:section}
A Li coefficient depends on zeros at every height. How can a finite list of zeros prove its sign? Theorem~\ref{li:everyheight} gives a uniform answer in every degree. Above the verified rectangle the few ordinates that can contribute negatively are outweighed by a positive contribution from distant zeros. Proposition~\ref{li:uncertain-center} also allows an unresolved pair near the center. For the elliptic curve of conductor $234446$ it proves strict positivity through index $96\,251\,863$.

Use unitary normalization. Let \(Z\) be the nontrivial zeros with multiplicities and put
\[
     \lambda_m=\lim_{U\to\infty}\sum_{\substack{\rho\in Z\\|\Im\rho|\le U}}
                   \bigl[1-(1-1/\rho)^m\bigr]\qquad(m=1,2,\dots).
\]
We use the real part, which converges absolutely for each fixed \(m\): at a zero with large ordinate \(\gamma\), the linear term has real part \(O_m(\gamma^{-2})\), as do the higher terms together. The equivalence of ERH and nonnegativity of all these real parts is due to Li, Bombieri--Lagarias and, for automorphic functions, Lagarias. Partial Li criteria have also been studied by Chaudhary, Khmiri and Mazhouda. Our estimates below specify the constants and the forms to which the numerical conclusions apply.
\href{https://doi.org/10.5802/aif.2311}{J.~Lagarias, Ann. Inst. Fourier 57 (2007), 1689--1740.}
\href{https://doi.org/10.7153/jca-2026-28-07}{M.~P. Chaudhary, J.~Khmiri and K.~Mazhouda, J. Class. Anal. 28 (2026).}

\begin{theorem}\label{li:explicit}
Suppose a completed cuspidal standard function has all its nontrivial zeros with \(|\Im\rho|\le T\) on its critical line, where \(T\ge1\). Then
\[
                  \Re\lambda_m>0\qquad (1\le m\le 2\pi T-\pi).
\]
For any of the thirteen forms in Section~\ref{num:chapter}, if \(T\ge4096\) is an integer satisfying the same hypothesis, then
\[
                  \Re\lambda_m>0\qquad (1\le m\le T^2\log T).
\]
In particular all thirteen have \(139\,548\,959\) strictly positive initial coefficients. For each row whose certified height is \(65\,536\), this increases to \(47\,632\,711\,549\). For every rational elliptic curve of conductor at most \(1000\) the first conclusion and the census alone give \(1605\) strictly positive coefficients.
\end{theorem}

\subsection*{The effect of one reflected pair}
For \(\rho=\beta+i\gamma\) with \(g=|\gamma|>0\), write \(1-1/\rho=e^{u+i\vartheta}\) with principal argument and set \(\theta=|\vartheta|\). Reflection in the critical line changes \(u\) to \(-u\) and leaves \(\vartheta\) fixed. For \(0\le\beta\le1\),
\[
 |u|=\frac12\left|\log\frac{g^2+(1-\beta)^2}{g^2+\beta^2}\right|
       \le\frac{|1-2\beta|}{2g^2},\qquad
 \theta=\int_{-\beta}^{1-\beta}\frac{g\,dx}{g^2+x^2},
 \qquad\frac{g}{1+g^2}\le\theta\le\frac1g .
\]
Thus the contribution per zero, averaged with its reflection, is \(1-\cos(m\theta)\cosh(mu)\). On the line the expression is \(1-\cos(m\theta)\), also at the central zero with \(\theta=\pi\).

If \(g>T\) and \(m\le2\pi T-\pi\), then \(0<m\theta<2\pi\). Whenever \(m\theta\le\pi/2\), we have \(m|u|\le m\theta\), since \(g\ge1\). On \([0,\pi/2]\) the product \(\cos(x)\cosh(x)\) is at most one, by differentiation or by \(\tan x\ge\tanh x\). In \([\pi/2,3\pi/2]\) the cosine is nonpositive. In the remaining interval put \(y=2\pi-m\theta\). Now
\[
         y\ge 2\pi-m/g\ge\pi/g,\qquad m|u|\le m/(2g^2)<\pi/g.
\]
Use the same product inequality at \(y\). All high terms are therefore nonnegative. At sufficiently large \(g\) each is strictly positive because \(0<m\theta<\pi/2\) and \(\cos(x)\cosh(x)<1\) for \(x>0\). There are infinitely many such zeros by the zero-counting formula. This proves the first assertion.

\subsection*{Two bounds on the ordinates}
Write \(v=(k-1)/2\), \(A=\sqrt N/(2\pi)\). For the thirteen forms, the completed function in the unitary variable has gamma and power factors \(\Gamma(s+v)A^{s+v}\). For real \(t\) define
\[
 S(t)=\sum_{\rho}\frac{2-\beta}{(2-\beta)^2+(t-\gamma)^2}.
\]
The logarithmic derivative of the Hadamard product gives
\[
     S(t)=\log A+\Re\psi(2+v+it)+\Re\frac{L'}{L}(2+it)
\]
with \(L\) in unitary normalization. To check the real constant, evaluate the same logarithmic derivative on the critical line at a nonzero of the completed function. The real part there vanishes by the functional equation; the reflected pairs cancel in the sum of real parts, which is absolutely convergent. Thus no real constant remains at \(2+it\).

We have \(|\log A|<2\). The normalized local roots have modulus at most one by the weight bound; thus \(|L'/L(2+it)|\le2\sum_{\ell\ge2}\log(\ell)/\ell^2<3\), by an integral comparison starting at \(2\). For \(\Re z>0\), comparison of \(\sum_{j=0}^{M-1}(z+j)^{-1}\) with its integral on \([0,M]\) gives \(|\psi(z)-\log z|\le\int_0^\infty |z+x|^{-2}\,dx\le1/\Re z\). Since \(2+v\le15/2\),
\begin{equation}\label{li:poisson}
 \log|t|-5.5\ \le S(t)\le\log(|t|+8)+5.5 .
\end{equation}
The lower estimate is used only for \(t\ne0\). A zero in any unit interval contributes at least \(8/17\) to \(S\) at its midpoint: \(1\le2-\beta\le2\) and \(|t-\gamma|\le1/2\). Consequently the number of zeros in any such interval of midpoint \(w\) is at most \(3[\log(|w|+9)+6]\).

For an integer \(T\ge4096\), summing the last bound on the two signs of the height gives
\begin{equation}\label{li:tail}
   \sum_{|\gamma|>T}|\gamma|^{-4}
   \le6\sum_{j=T}^\infty\frac{\log(j+10)+6}{j^4}
   \le\frac{3(\log T+8)}{T^3}.
\end{equation}
Indeed \(\log(j+10)\le\log(2j)\); integrate the decreasing majorant from \(T-1\), using
\(\int_y^\infty(\log(2x)+6)x^{-4}dx=(\log(2y)+6+1/3)/(3y^3)\).

Next fix an integer \(m\ge25\,000\). On the range \(m/4\le g\le m/2\) the angle satisfies \(1.99\le m\theta\le4\); hence \(1-\cos(m\theta)\ge5/4\). To lower-bound the number of such ordinates, integrate \(S(t)\) on \(9m/32\le t\le15m/32\), and on its reflection. At a fixed positive \(t\) the zeros outside \([m/4,m/2]\) have \(|\gamma-t|\ge m/32\). Put \(h=\lfloor m/32\rfloor\). In each of the two unit intervals at distances in \([j,j+1]\), \(j\ge h\), the number of zeros is at most \(3(\log(40j)+6)\), since \(t\le15m/32\). Each adds at most \(2/j^2\). Their total is at most
\[
  12\sum_{j=h}^{\infty}\frac{\log(40j)+6}{j^2}
       \le 12\,\frac{\log(40(h-1))+7}{h-1}<1 .
\]
The inequality follows at \(h\ge780\) by monotonicity. The same calculation works for negative \(t\). On each interval the integral of the summand for a single zero is at most its whole-line integral \(\pi\). Equation~\eqref{li:poisson} and \(t\ge m/4\) on the positive interval now give
\begin{equation}\label{li:base}
 \sum_{\rho}[1-\cos(m\theta_\rho)]
  \ \ge\ \frac54\,\frac{3m}{8\pi}\,(\log m-8)
  \ >\ \frac{m}{8}(\log m-8).
\end{equation}

\subsection*{Allowing all unverified heights}
If \(|\gamma|\le T\) then \(u=0\). For the remaining zeros use
\(\cosh y-1\le \tfrac12 y^2e^{|y|}\) and \eqref{li:tail}. The loss upon replacing the sum in \eqref{li:base} by \(\Re\lambda_m\) is at most
\[
        \frac38(\log T+8)\frac{m^2}{T^3}e^{m/(2T^2)}.
\]
We need only consider \(2\pi T-\pi<m\le T^2\log T\), so \(m>25\,000\). The ratio of this loss to the last expression in \eqref{li:base} increases with \(m\) on this range, because \(m/(\log m-8)\) does. At the upper endpoint it is at most
\[
    \frac{3(\log T+8)\log T}
         {\sqrt T\,[\,2\log T+\log\log T-8\,]}
      \ \le\ \frac{3(\log T+8)}{\sqrt T}<1.
\]
The last bound decreases for \(T\ge4096\) and already holds there. This proves Theorem~\ref{li:explicit}. The small directed-rounding script \texttt{li\_prefix.py} checks the numerical inequalities and integer parts.

\begin{proposition}[A fixed-function asymptotic form]\label{li:asympt}
Fix a completed unitary cuspidal standard function for which every zero has \(|\Re\rho-1/2|\le D\), with \(0<D\le1/2\). For each \(0<\varepsilon<1\) and every sufficiently large \(T\) at which all zeros through height \(T\) lie on the line, one has
\[
   \Re\lambda_m>0\qquad
      \left(1\le m\le \frac{1-\varepsilon}{D}\, T^2\log T\right).
\]
The threshold depends on the fixed function. For any of the all-height Hecke-induced examples of Section~\ref{recentcm:section}, \(D=3/8\) is admissible.
\end{proposition}
\begin{proof}
The unit-interval bound \(O(\log(2+|t|))\) and the Riemann--von Mangoldt formula give \(\sum_{|\gamma|>T}|\gamma|^{-4}\ll(\log T)/T^3\) and \(\gg m\log m\) ordinates in \(m/4\le|\gamma|\le m/2\). Constants may depend on the fixed function. The latter ordinates give \(\sum(1-\cos(m\theta))\gg m\log m\). The loss is \(O_D(m^2T^{-3}\log T\,e^{Dm/T^2})\), since \(|u|\le D/\gamma^2\). For \(2\pi T-\pi<m\le(1-\varepsilon)T^2\log T/D\), its ratio to the positive lower bound is \(O_{D,\varepsilon}(T^{-\varepsilon}\log T)\), tending to zero. Smaller indices are covered individually by the first assertion of Theorem~\ref{li:explicit}.
\end{proof}

\begin{corollary}\label{li:gl2}
For a fixed cuspidal unitary standard \(L\)-function on \(\mathrm{GL}_2/\Q\), let \(0<A<17/8\). For all sufficiently large \(T\), a critical-line certificate through \(T\) implies \(\Re\lambda_m>0\) for every \(m\le A T^2\log T\).
\end{corollary}
\begin{proof}
Put \(a_\rho=|\beta-1/2|\). The zero-density input in Section~\ref{common:section} bounds the number of zeros of height at most \(U\) with \(a_\rho\ge b\) by \(O_\pi(U^{1-b/8}\log U)\) for \(0<b<1/2\). For \(b=0\) use the zero-counting formula. Subdivide \([0,1/2]\) into finitely many intervals of length at most \(\eta\), with left endpoints \(b\). In the bin \(b\le a_\rho\le b+\eta\), the loss from zeros above \(T\) is
\[
 O_\pi\!\left(m^2\,e^{m(b+\eta)/T^2}
            \sum_{\substack{|\gamma|>T\\ a_\rho\ge b}}|\gamma|^{-4}\right)
 \ll_\pi m^2 e^{m(b+\eta)/T^2} T^{-3-b/8}\log T.
\]
The last inequality follows by summing on the dyadic intervals above \(T\). Against the lower bound \(\gg_\pi m\log m\), the ratio for \(2\pi T-\pi<m\le A T^2\log T\) is
\[
           O_{\pi,A}\bigl((\log T) T^{-1+(A-1/8)b+A\eta}\bigr).
\]
If \(A\ge1/8\), its exponent is at most \(-1+(A-1/8)/2+A\eta\), strictly negative for a sufficiently small fixed \(\eta\). If \(A<1/8\), choose \(A\eta<1\). The finitely many bins therefore give a total loss smaller than the positive contribution for all sufficiently large \(T\). The preceding geometric estimate covers the other integers \(m\).
\end{proof}

\subsection*{Summing the high-zero loss before taking a maximum}
\begin{theorem}\label{li:integral-tail}
For the thirteen forms of Theorem~\ref{li:explicit}, let $T\ge4096$ be an integer through which every zero is on the line. Put $y=T-1$, $L=\log T$, and $C=\log(2y)+19/3$. Then $\Re\lambda_m>0$ throughout
\[
  1\le m\le 2y^2\bigl(L+2\log L-\log 17\bigr).
\]
A sharper numerical endpoint follows from a single inequality. If $M\ge25\,000$ and $A=M/(2y^2)$ satisfy
\begin{equation}\label{li:single-test}
 6yC\left(\frac{\sinh A}{A}-1\right)<\frac M8(\log M-8),
\end{equation}
then all integer indices $1\le m\le M$ have $\Re\lambda_m>0$. At $T=4096$ this certifies the first $342\,205\,229$ indices. At $T=65\,536$ it certifies the first $118\,740\,061\,700$.
\end{theorem}
\begin{proof}
Let $E_m$ denote the sum of $\cosh(m u_\rho)-1$ over $|\Im\rho|>T$. Each unit interval above a positive integer $j\ge T$, together with its reflection, contains at most $6(\log(2j)+6)$ ordinates. Expanding the nonnegative terms of $\cosh-1$, set $a=m/(2y^2)$. For $k\ge1$ the function $(\log(2t)+6)t^{-4k}$ decreases on $[y,\infty)$. Its integral is
\[
 y^{1-4k}\left[\frac{\log(2y)+6}{4k-1}+\frac1{(4k-1)^2}\right].
\]
Thus, since $4k-1\ge2k+1$ and $4k-1\ge3$,
\begin{equation}\label{li:sinh}
 E_m\le6y C\sum_{k\ge1}\frac{a^{2k}}{(2k+1)!}
      =6y C\left(\frac{\sinh a}{a}-1\right).
\end{equation}
The positive part of the sum has lower bound \eqref{li:base}. For $m\ge25\,000$ the ratio of \eqref{li:sinh} to $m(\log m-8)/8$ increases with $m$, since each $m^{2k-1}/(\log m-8)$ does. Smaller indices are covered by the termwise result. This proves \eqref{li:single-test} as a sufficient test.

For a closed formula set $A=L+2\log L-\log17$ and $M=2y^2A$. Here $A\ge L$, $C\le37L/20$ and $\log M-8\ge4L/3$. The last bound follows from $2L/3+\log(2L)+2\log(1-1/T)>8$ at $T=4096$ and monotonicity. Using $\sinh A< e^A/2$ and $e^A=TL^2/17$, the ratio in \eqref{li:single-test} is strictly below
\[
  \frac{12C e^A}{yA^2(\log M-8)}
      \le \frac{12}{17}\frac{37}{20}\frac34\frac{4096}{4095}<1.
\]
The two numerical endpoints use \eqref{li:single-test} directly. Directed $160$-bit interval evaluations put the left side divided by the right side between $0.9999999780$ and $0.9999999781$ at the first endpoint, and between $0.9999999999701$ and $0.9999999999702$ at the second. The ratio exceeds one at each succeeding index. The exact integer inputs and checks are in \texttt{li\_tail.py}.
\end{proof}

For the fixed function in Proposition~\ref{li:asympt}, the same integral gives
\begin{equation}\label{li:secondlog}
 \Re\lambda_m>0\quad\left(1\le m\le
 \frac{(T-1)^2}{D}\bigl[\log T+2\log\log T-C_{F,D}\bigr]\right)
\end{equation}
whenever $T$ is sufficiently large and the line has been verified through $T$. Indeed, the unit-interval count for a fixed function is $O_F(\log(2+|t|))$. Replace $m/2$ by $mD$ in \eqref{li:sinh}. With $a=mD/(T-1)^2$, division by the positive lower bound $\gg_F m\log m$ gives $O_{F,D}(e^a/(T a^2))$ at the displayed endpoint. Increasing $C_{F,D}$ makes this quantity smaller than one. The ratio of the full series bound to $m\log m$ is increasing for $m\ge3$, so the same comparison applies down to the large-index threshold of the fixed-function lower bound. Fixed small indices are covered by Theorem~\ref{li:explicit}. The argument for a real $T$ uses $y=\lfloor T\rfloor-1$ in the integral and an increase of the constant, since the resulting change of $a$ at the endpoint is $O((\log T)/T)$.

\subsection*{A height independent of the representation and the field}
Chaudhary, Khmiri and Mazhouda prove nonnegativity for $m\le T^2$ once $T\ge T_1(\pi)$ for an irreducible unitary cuspidal representation of $\mathrm{GL}_n(\mathbb A_\Q)$. Their footnote on p.~113 asks whether one can choose $T_1$ independently of $\pi$. The following bound gives such a choice for standard functions over every number field. The constants also allow arbitrary degree and unitary norm twists. \href{https://files.ele-math.com/articles/jca-28-07.pdf}{Chaudhary--Khmiri--Mazhouda, Theorem~1, J. Class. Anal. 28 (2026), 107--119.}

\begin{theorem}\label{li:uniform}
Let $F$ be a number field, let $\pi$ be any irreducible unitary cuspidal representation of $\mathrm{GL}_n(\mathbb A_F)$, and use the zero divisor of its completed standard function in the unitary variable, after removing any poles. Put $L=\log T$. For every integer $T\ge4096$ at which all nontrivial zeros of ordinates $|\gamma|\le T$ are on the critical line,
\begin{equation}\label{li:all-conductors}
 \Re\lambda_m(\pi)>0\qquad
  \left(1\le m\le 2(T-1)^2 [L+2\log L-\log 576]\right).
\end{equation}
\end{theorem}
\begin{proof}
Write $d=n[F:\Q]$ and use $\Gamma_\C(s)=\Gamma_\R(s)\Gamma_\R(s+1)$ to express the gamma factor as $\prod_{j=1}^d\Gamma_\R(s+\mu_j)$. At finite places $|\alpha_j(\mathfrak p)|\le (\mathrm N\mathfrak p)^{1/2}$ and $\Re\mu_j\ge-1/2$ by \eqref{aw:lrs}. Set
\[
   H(x)=\log Q_\pi+\sum_{j=1}^d\log(8+|\mu_j|+x),\qquad x\ge0.
\]
The absolute conductor $Q_\pi=|D_F|^n\mathrm N\mathfrak q_\pi$ is at least one. The conductor and the gamma parameters enter both sides of the comparison through $H$.

Use the same positive sum $S(t)$ as in \eqref{li:poisson}. At $2+it$ the Euler product satisfies $|\log L(2+it,\pi)|\le d\log\zeta(3/2)<d\log3$ on the branch from the Euler product, and $|(L'/L)(2+it,\pi)|<5d$. To see this, group prime ideal powers with norm $p^a$. The sum of residue degrees of the contributing prime ideals is at most $[F:\Q]$. For the logarithmic derivative this gives $d\sum_{k\ge2}(\log k)k^{-3/2}<5d$ by integration. For the logarithm itself, division by $a\log p$ gives $d\log\zeta(3/2)$.

The Hadamard logarithmic derivative at $2+it$ has gamma contribution
\[
 \tfrac12\sum_j\log|2+\mu_j+it|-\tfrac d2\log(2\pi)+e(t),
       \qquad |e(t)|\le 2d/3.
\]
Here the real part of each digamma argument is at least $3/4$. If $L$ is a norm twist of the Dedekind zeta function, multiply the completion by $(s-i\tau)(s-1-i\tau)$ for the appropriate $\tau\in\R$. The real logarithmic derivative of these factors at $2+it$ belongs to $[0,3/2]$. As before the functional equation determines the real Hadamard constant. Consequently
\begin{equation}\label{li:uniform-unit}
 S(t)\le\tfrac12 H(|t|)+8d,
 \qquad \#\{\gamma:|\gamma-w|\le1/2\}\le (17/8)S(w).
\end{equation}

We first prove a lower bound for the positive sum, uniformly in the archimedean parameters. Assume $m\ge25\,000$ and let $I$ be either $[9m/32,15m/32]$ or its negative. Its length is $\ell=3m/16$. For every $\mu_j$,
\[
  \ell^{-1}\int_I\log|2+\mu_j+it|\,dt
      \ge \log(8+|\mu_j|+m)-5.
\]
If $|\mu_j|\ge2m$, then $|2+\mu_j+it|\ge |\mu_j|/2\ge(8+|\mu_j|+m)/4$. Otherwise, replace the modulus by $|t+\Im\mu_j|$. Among intervals of length $\ell$ the integral of $\log|t|$ is smallest for the one centered at zero. Its average there is $\log(\ell/2)-1$, while $8+|\mu_j|+m\le4m$. This proves the bound since $1+\log(128/3)<5$. Integrating $L'/L$ over $I$ uses two endpoint values of $\log L$, with absolute mean at most $2d\log3/\ell$. In particular,
\[
   \ell^{-1}\int_I S(t)\,dt\ \ge\ \tfrac12 H(m)-4.1d .
\]

For $t\in I$, all zeros outside the band with the corresponding sign and with $m/4\le |\gamma|\le m/2$ have $|t-\gamma|\ge m/32$. Put $h=\lfloor m/32\rfloor$. In each unit interval at distance between $j$ and $j+1$ from $t$, $j\ge h$, the last bound in \eqref{li:uniform-unit} is at most $3H(m+j+1)$ because $H(m+j+1)\ge d\log 25000$. Also
\[
 H(m+j+1)\le H(m)+d\log(1+(j+1)/m)
            \le H(m)+d\log(3j/(h-1)).
\]
The kernel of each zero is at most $2/j^2$. Integrating a decreasing majorant from $h-1$, the total outside the band at a fixed $t$ is at most
\[
  \frac{12}{h-1}\bigl[H(m)+d(1+\log 3)\bigr]<0.019 H(m).
\]
Since $H(m)\ge d\log 25000$, the retained mean exceeds $H(m)/16$. Integrate on both signed windows, use the kernel integral $\pi$ and $1-\cos(m\theta)\ge5/4$ on both bands. This gives
\begin{equation}\label{li:uniform-base}
 \sum_\rho [1-\cos(m\theta_\rho)]
   \ge\frac54\frac{3m}{8\pi}\frac{H(m)}{16}
   >\frac{mH(m)}{128}.
\end{equation}

Put $y=T-1$ and $a=m/(2y^2)$. By \eqref{li:uniform-unit}, the two unit intervals with $j\le |\gamma|\le j+1$ contain at most $8H(j+1)$ zeros altogether for $j\ge T$, since $H(j+1)\ge d\log4095$. For these $j$,
\[
     H(j+1)\le H(y)+d\log(2j/y).
\]
Integrate $(H(y)+d\log(2t/y))t^{-4k}$ from $y$ and sum as in \eqref{li:sinh}. The high-zero loss is at most
\[
 8y\bigl[H(y)+d(\log2+1/3)\bigr]\left(\frac{\sinh a}{a}-1\right)
 \le9yH(y)\left(\frac{\sinh a}{a}-1\right).
\]
For $m>2\pi T-\pi$ we have $m\ge y$ and $m\ge25\,000$. Hence the positive sum is at least $m H(y)/128$. The quotient of the last loss bound by this quantity increases with $m$, term by term in the series for $\sinh(a)/a-1$. At the claimed endpoint, $a=L+2\log L-\log576\ge8L/11$ for $T\ge4096$. The quotient there is strictly below
\[
    \frac{288 e^a}{y a^2}
    \le \frac{288}{576}\frac{4096}{4095}\left(\frac{11}{8}\right)^2<1.
\]
All smaller indices in the termwise range are covered by Theorem~\ref{li:explicit}.
\end{proof}

In particular all unitary cuspidal standard functions in this theorem have $\Re\lambda_m>0$ for $m\le T^2$ whenever the line has been verified through any real $T\ge4097$. Indeed apply \eqref{li:all-conductors} at $\lfloor T\rfloor$ and use $a\ge8\log(4096)/11>6$. For the Riemann zeta-function, the Platt--Trudgian verification through $3\cdot 10^{12}$ and \eqref{li:all-conductors} yield strict positivity for more than $5.2\cdot10^{26}$ initial Li coefficients. \href{https://doi.org/10.1112/blms.12460}{D.~Platt and T.~Trudgian, Bull. London Math. Soc. 53 (2021), 792--797.}

\subsection*{Rankin--Selberg convolutions without a tensor-product lift}
\begin{theorem}\label{li:pairs}
Let $\pi$ and $\pi'$ be unitary cuspidal representations of $\mathrm{GL}_n$ and $\mathrm{GL}_{n'}$ over any number field $F$. In Theorem~\ref{li:uniform} one may replace the standard function by the completed Rankin--Selberg function $L(s,\pi\times\pi')$. The index bound, the threshold and strict positivity are unchanged. When a pole occurs, remove the two reflected poles from the completion before taking its Li coefficients.
\end{theorem}
\begin{proof}
Use the actual local factors for the pair. With $d=nn'[F:\Q]$ and complex gamma factors split, the finite reciprocal roots have modulus at most $\mathrm N\mathfrak p$ and the $d$ archimedean shifts satisfy $\Re\mu_j\ge-1$. These estimates include ramified primes. The absolute conductor is $Q=|D_F|^{nn'}\mathrm N\mathfrak q_{\pi\times\pi'}$. Rankin--Selberg theory gives the functional equation, nonvanishing on the one-line away from poles, and order one after pole removal. There is at most one reflected pair of simple poles: upon removing the norm twists of the two cuspidal representations they occur at $0,1$ exactly in the contragredient case. Undoing the translation puts them at $i\tau,1+i\tau$ for some real $\tau$. Multiplication by $(s-i\tau)(s-1-i\tau)$ preserves horizontal reflection of the zero divisor. For the input on the actual local factors and the completed product see \href{https://arxiv.org/pdf/2103.05634}{Humphries--Thorner, \S3, (3.5)--(3.11)}. For nonvanishing use \href{https://doi.org/10.1090/S0273-0979-1980-14769-2}{F.~Shahidi, \emph{On nonvanishing of $L$-functions}, Bull. AMS 2 (1980)}. On their Li coefficients see \href{https://www.math.bas.bg/infres/MathBalk/MB-24/MB-24-217-228.pdf}{A.~Od\v{z}ak, \S3}. No assertion about a functorial lift of the pair enters the argument.

Move the Poisson kernel to $\sigma=5/2$. Write $S_\sigma(t)$ for the sum of
$(\sigma-\beta)/((\sigma-\beta)^2+(t-\gamma)^2)$ over the zeros. Use
$H(x)=\log Q+\sum_j\log(8+|\mu_j|+x)$ as before. Since $\sigma-1=3/2$, the bounds $|\log L|\le d\log3$ and $|L'/L|<5d$ still hold. The real parts of the digamma arguments are at least $3/4$. The real contribution of the two pole-removing factors is in $[0,16/15]$, because it is at most $1/\sigma+1/(\sigma-1)$. The same estimates therefore give
\[
 S_\sigma(t)\le H(|t|)/2+8d,\qquad
 \frac1{|I|}\int_I S_\sigma(t)\,dt\ge H(m)/2-4.1d
\]
on the two intervals $I$ of the preceding proof. For the second inequality, when $|\mu_j|\ge2m$, the reverse triangle inequality still gives $|\sigma+\mu_j+it|\ge|\mu_j|/2$ there. Otherwise the same integral of $\log|t+\Im\mu_j|$ applies.

For $|t-\gamma|\le1/2$, the kernel is at least $5/13$, its minimum at $\sigma-\beta=5/2$. When $m\ge25\,000$, a unit interval at distance $j$ as above consequently contains at most $(7/2)H(m+j+1)$ zeros, since
\[
 \frac{13}{5}\left(\frac12+\frac8{\log25000}\right)<\frac72 .
\]
For $j\ge h=\lfloor m/32\rfloor$, each such zero contributes at most $(5/2)j^{-2}$. Accounting for both directions from $t$, the exterior contribution is bounded above by
\[
 \frac{35}{2(h-1)}[H(m)+d(1+\log3)]<0.0273H(m).
\]
The retained mean still exceeds $H(m)/16$, since $1/2-4.1/\log25000-0.0273>1/16$. The integral of an individual kernel is $\pi$ at this abscissa also. The lower phase mass $mH(m)/128$ of \eqref{li:uniform-base} follows.

For $j\ge T\ge4096$, the two unit intervals with $j\le|\gamma|\le j+1$ contain at most $8H(j+1)$ zeros, because $(13/5)(1/2+8/\log4095)<4$. The hyperbolic-sine tail estimate and its comparison with the phase mass now have exactly the same constants as in Theorem~\ref{li:uniform}. The reflected-pair argument covers all the smaller indices. This proves the assertion.
\end{proof}

For elliptic curves $E,E'/\Q$, modularity identifies the tensor of their two unitary standard parameters with the Rankin--Selberg local parameters. Each elliptic shift back to arithmetic coordinates contributes $1/2$. Thus the theorem applies to verification about the line $\Re s=3/2$ for $L(E\otimes E',s)$ in arithmetic coordinates. In the isogenous case the translated poles are removed first. The resulting Li coefficients always refer to the shift to the unitary center.

For a finite product of these completions with positive integer multiplicities, the same criterion holds independently of the number of factors. Indeed the completions with their poles removed are nonzero at every point of the two boundary lines. Clearing the poles of the product consequently gives the product of the individually regularized completions, up to a nonzero constant. Its zero divisor is their sum with positive multiplicities. A verified height for the product is therefore a verified height for each factor. Absolute convergence of the real Li sum at each fixed index then adds their strict inequalities.

\subsection*{The quadratic criterion at every height}
\begin{theorem}\label{li:everyheight}
Let $L$ be the standard function of a unitary cuspidal representation over any number field, or the Rankin--Selberg function of any two such representations over the same field. Use unitary normalization and the zero divisor of the completion after pole removal. For every real $T\ge0$,
\[
  [\,|\Im\rho|\le T\implies\Re\rho=1/2\,]\quad\Longrightarrow\quad
        \Re\lambda_m(L)>0\qquad (m\in\Z,\ 1\le m\le T^2).
\]
In particular, if $\Re\lambda_m(L)\le0$, there is an off-line zero of absolute ordinate at most $\sqrt m$.
\end{theorem}
\begin{proof}
It suffices to use $T=\sqrt m$. For $1\le m\le32$ we have $m\le2\pi\sqrt m-\pi$, by concavity and evaluation at the two endpoints. The reflected-pair lemma applies. Assume $m\ge33$.

Write $d$ for the absolute degree, $Q\ge1$ for the absolute conductor and $\mu_j$ for the shifts in the $d$ real gamma factors. The standard analytic inputs for both kinds of function are given in \href{https://arxiv.org/pdf/2103.05634}{Humphries--Thorner, \S3}. We use only $\Re\mu_j\ge-1$ and $|\alpha_j(\mathfrak p)|\le\mathrm N\mathfrak p$, including ramified places. With
\[
 P_\sigma(t;\rho)=
 \frac{\sigma-\beta}{(\sigma-\beta)^2+(t-\gamma)^2},
 \qquad S_\sigma(t)=\sum_{\rho=\beta+i\gamma}P_\sigma(t;\rho),
\]
the Hadamard logarithmic derivative gives, for $\sigma\ge3$,
\[
 S_\sigma(t)=\frac12\log Q+
  \frac12\sum_{j=1}^d\log\frac{|\sigma+\mu_j+it|}{2\pi}
        +e_\sigma(t)+\Re\frac{L'}L(\sigma+it)+p_\sigma(t).
\]
Here $|e_\sigma(t)|\le d/(\sigma-1)$ and $0\le p_\sigma(t)\le 1/\sigma+1/(\sigma-1)$, the latter term arising only from the reflected pair of simple poles. Indeed $|\psi(z)-\log z|\le1/\Re z$ for $\Re z>0$, from the integral with kernel $1/(1-e^{-x})-1/x$ bounded between zero and one. The reflected pair $\rho,1-\bar\rho$ makes the real constant in the Hadamard identity vanish on the critical line. Grouping prime ideals by the rational prime below them gives $|L'/L(\sigma+it)|\le d[-\zeta'(\sigma-1)/\zeta(\sigma-1)]$. At $\sigma=3$ this is below $0.65d$. At $\sigma\ge16.5$ it is below $0.001d$, by integration of $(\log x)x^{-15}$ beyond the first term.

\emph{A lower bound from distant zeros.} Set $r=m/2$, $R=2m=4r$ and
\[
 D=\log Q+\sum_{j=1}^d\log\frac{|R+\mu_j|}{2\pi}
       \ \ge\ d\log\frac{65}{2\pi}.
\]
For $g=|\gamma|\le r$, the ratio $P_R(0;\rho)/P_r(0;\rho)$ increases with $g$ at fixed $\beta$. At $g=r$, the numerator in the expression
\[
 \frac{(r-\beta)+r^2/(r-\beta)}{(R-\beta)+r^2/(R-\beta)}
\]
is $2r+\beta^2/(r-\beta)\le2r+1/(r-1)$. The denominator is at least $R-1+r^2/R$. Their quotient is less than $0.481$ for $r\ge16.5$, by division by $r$. For $g>r$ one has $P_R(0;\rho)\le R/g^2$. Consequently
\[
 R\sum_{g>r}\frac1{g^2}\ge S_R(0)-0.481 S_r(0)>0.21D .
\]
For the strict inequality, $|\mu_j+r|\le|\mu_j+R|$. The lower bound for the difference is
\[
 \tfrac12(1-0.481)D-
 d\big[1/65+0.001+0.481(2/31+0.001+4/31)\big].
\]
Use the displayed lower bound on $D/d$.

With $\theta=|\arg(1-1/\rho)|$ and $u=\log|1-1/\rho|$, the real contribution per zero after horizontal pairing is $F_m=1-\cos(m\theta)\cosh(mu)$. For $g>r$ we have $0<m\theta\le2$ and $|mu|\le m/(2g^2)<1$. On $0\le x\le2$, $1-\cos x\ge(1-\cos2)x^2/4$ by monotonicity of $\sin(x/2)/(x/2)$. Also $\cosh y-1\le y^2$ for $|y|\le1$. Using $\theta\ge g/(1+g^2)$ gives
\[
 F_m\ \ge\ \frac{m^2}{g^2}
 \left[\frac{1-\cos2}{4}(1+g^{-2})^{-2}-\frac1{4g^2}\right]
       >0.34\,\frac{m^2}{g^2}.
\]
The sum over these distant zeros is therefore greater than $0.0357mD$.

\emph{An upper bound on all negative contributions.} Zeros with $g\le\sqrt m$ are on the line and contribute nonnegatively. If $F_m<0$ for any other zero, the distance from $m\theta$ to a multiple $2\pi k$ is less than $m|u|\le1/2$. Indeed a negative contribution requires a positive cosine. At angular distance $w\le\pi/2$, the inequality $\cos w\cosh w\le1$ shows that $w\ge m|u|$ precludes negativity. Also $\theta>|u|$ for $g>1$, so $k\ge1$. The positive integer $k$ satisfies $2\pi k\le\sqrt m+1/2$. Write $c=m/(2\pi k)$. From $g\theta\in[g^2/(1+g^2),1]$ one obtains
\[
 |g-c|\le\frac{c}{g}\left(\frac12+\frac{g}{1+g^2}\right),\quad
 \frac{g}{c}\ge 1-\frac1{1+g^2}-\frac1{2g}.
\]
At $g\ge\sqrt{33}$ these imply $|g-c|<1$.

For $|t-\gamma|\le1$, $P_3(t;\rho)\ge3/10$. Thus both signed windows for any one $k$ contain at most $(10/3)(S_3(c)+S_3(-c))$ zeros. Since $c\le m/(2\pi)$ and $|R+\mu_j|\ge2m-1$, each gamma modulus at $3\pm ic$ is at most $|R+\mu_j|(1+33/(130\pi))$. The two sums therefore total at most
\[
 D+\big[\log(1+33/(130\pi))+1+1.30\big]d+5/3.
\]
Each negative summand has magnitude at most $\cosh(1/2)-1$. Substituting $D/d\ge\log(65/(2\pi))$, $d\ge1$, gives total loss less than $1.17D$ per possible $k$. The estimates permit counting a zero in more than one window.

Finally
\[
  0.0357m > \frac{1.17}{2\pi}(\sqrt m+1/2)\qquad(m\ge33),
\]
because $m/(\sqrt m+1/2)$ increases and the inequality is strict at $33$. Distant zeros alone outweigh all possible negative summands. Absolute convergence of the real Li sum now proves the theorem.
\end{proof}

\begin{theorem}[An unconditional positive tail]\label{li:positive-tail}
For every integer $m\ge1$ and every $T\ge\sqrt m$, the real Li sum over the zeros of absolute ordinate strictly greater than $T$ is positive, regardless of the locations of the zeros below $T$. For $m\ge100$, put
\[
 v=\max(T,m/2),\qquad D_v=\log Q+\sum_{j=1}^d\log\frac{|4v+\mu_j|}{2\pi}.
\]
Then the tail satisfies the two-sided estimate
\[
 \frac1{128}\frac{m^2}{v}D_v
  < \sum_{|\gamma|>T}\Re\left[1-\left(1-\frac1\rho\right)^m\right]
  <3\frac{m^2}{v}D_v .
\]
Both constants are absolute. In particular the low-zero sum, even if some low zeros lie off the line, determines a rigorous interval for the complete real coefficient. Lagarias's Theorem~1.1 gives an asymptotic expansion involving the incomplete coefficient for a fixed standard function over $\Q$. \href{https://aif.centre-mersenne.org/item/10.5802/aif.2311.pdf}{J.~Lagarias, Ann. Inst. Fourier 57 (2007), 1689--1740.}
\end{theorem}
\begin{proof}
The cut preserves both members of every horizontal pair. If $m\le32$ the termwise proof applies to every retained pair, and gives strict positivity at arbitrarily large ordinates. If $m\ge33$ and $\sqrt m\le T\le m/2$, all the distant zeros used in the preceding proof stay in the tail. Its negative part can only decrease when the comparison starts at $T$ instead of $\sqrt m$. In particular its sum exceeds $[0.0357m-1.17(\sqrt m+1/2)/(2\pi)]D_v$. For $m\ge100$ this lower bound exceeds $mD_v/64$, since $(\sqrt m+1/2)/m$ decreases.

If $T\ge m/2$ and $m\ge33$, every retained pair satisfies $F_m>0.34m^2/g^2$. The comparison of the two Poisson kernels in the preceding proof works for every real $r\ge16.5$ and $R=4r$, with $D$ defined using that $R$. Use $r=T$. It gives a lower bound $0.34(0.21)m^2 D_v/(4T)$, which is larger than the one stated.

For the upper bounds note first that $S_{4v}(0)<0.522D_v$. Indeed the three positive error bounds at $4v\ge66$ total at most $d(1/65+0.001+2/65)$ and $D_v/d\ge\log(65/(2\pi))$. If $T\le m/2$, every retained pair satisfies
\[
 F_m\le 2.13\min(1,m^2/g^2)\le10.65mP_{2m}(0;\rho).
\]
For the first inequality use $1+\cosh(1/2)<2.13$ when $g\le m$, and $1-\cos x\le x^2/2$, $\cosh y-1\le y^2$ when $g>m$. The second follows directly from $m\le2m-\beta\le2m$. Summation gives less than $6mD_v$. If $T\ge m/2$ then $F_m\le m^2/g^2$ for $g>T$, by the same cosine and hyperbolic cosine inequalities. On these zeros $P_{4T}(0;\rho)\ge3T/(17g^2)$ because $3T\le4T-\beta\le4T$. The upper bound follows from $(17/3)(0.522)<3$.
\end{proof}

\paragraph{Formal verification.}
The declaration \texttt{from\_local\_factors} in \texttt{LiFullDivisorTransfer} proves Theorem~\ref{li:everyheight} for the original complex Li powers, simultaneously for every real $T\ge0$ and positive integer $m\le T^2$. Its inputs are the strip and horizontal reflection of the divisor with multiplicities, the finite-height assertion, the conductor and gamma parameters, the local Euler roots with their degree and size bounds, and the usual real Hadamard identity for the completed function. That identity admits a translated pair of poles. Lean derives the required convergence of the real zero sums, unboundedness of the ordinates, the bounds for the Euler series and the actual complex digamma function, and the Poisson kernel estimates used in the positivity proof. The cited automorphic theorems supply the interpretation of these inputs for each function. The archive contains the replay of all fifteen modules and the axiom reports for 73 declarations. They use only propositional extensionality, classical choice and quotient soundness. The two-sided inequalities for the high-zero tail have the human proof above.

\begin{corollary}[Unitary isobaric sums]\label{li:isobaric}
Theorem~\ref{li:everyheight} and both bounds of Theorem~\ref{li:positive-tail} also hold for the standard function of a nonempty unitary isobaric sum and for the convolution of two such sums. In the bounds, $Q$ and the list of gamma shifts are those of the whole product.
\end{corollary}
\begin{proof}
Write each isobaric representation as a finite sum of unitary cuspidal constituents, allowing repetitions and imaginary norm twists. The standard function is the product of their standard functions. Multiplicativity of the Rankin--Selberg factors gives the product over ordered pairs of constituents in the convolution case. Each factor is holomorphic and nonzero on the two vertical boundary lines, apart from its indicated poles, by the one-line nonvanishing theorem and reflection. Consequently pole removal gives an entire function whose zero divisor is the sum of the individual pole-removed divisors, including coincident zeros with their full multiplicities. The rectangle hypothesis therefore passes to every constituent. The real Li sums and their tails add by absolute convergence. Conductors multiply and gamma parameters concatenate, so $D_v$ for the product equals the sum of the constituent values. Summing the strict inequalities proves the result. For the product identity see \href{https://people.math.ethz.ch/~kowalski/moment-theory.pdf}{Blomer, Fouvry, Kowalski, Michel, Mili\'cevi\'c and Sawin, \emph{The second moment theory of families of $L$-functions}, \S2.3}.
\end{proof}

\begin{corollary}[How far a low zero must leave the line]\label{li:displacement}
For either a cuspidal standard function or a cuspidal Rankin--Selberg pair, suppose $m\ge100$ and $\Re\lambda_m\le0$. Put
\[
        a_m=\operatorname{arcosh}(1+\sqrt m/96).
\]
There is a zero $\rho=1/2+x+i\gamma$ with $|\gamma|\le\sqrt m$ satisfying
\[
 |x|>(\gamma^2+1/4+x^2)\tanh(a_m/m),
 \qquad |\gamma|<\sqrt{m/(2a_m)} .
\]
In particular an enclosure with positive horizontal width can certify strict Li positivity even when it has yet to decide whether every low zero is exactly on the line.
\end{corollary}
\begin{proof}
Write $t=\sqrt m$ and $D=D_{m/2}$. By Theorem~\ref{li:positive-tail} the high zeros contribute more than $mD/64$. The reflected contribution per low zero is at least $-(\cosh(m u)-1)$, where $u=\log|1-1/\rho|$. Thus the sum of $\cosh(m u)-1$ over $|\gamma|\le t$ exceeds $mD/64$.

We need an upper bound on the number of such zeros, counted with multiplicity. Each contributes at least $(2t-1)/(5t^2)\ge0.38/t$ to $S_{2t}(0)$. In the logarithmic derivative formula $|2t+\mu_j|\le|2m+\mu_j|$. The sum of the upper gamma and Euler errors and the possible pole terms at $2t\ge20$ is at most
\[
       d(1/19+0.001+1/20+1/19)<0.157d .
\]
Since $D/d\ge\log(199/(2\pi))>3.45$, we have $S_{2t}(0)<0.546D$. The number of low zeros is therefore less than $3tD/2$. Some low zero must satisfy $\cosh(m u)-1>\sqrt m/96$, so $|u|>a_m/m$. From the two squared distances to $0$ and $1$ one obtains the exact identity
\[
       \tanh |u|=\frac{|x|}{\gamma^2+1/4+x^2}.
\]
This proves the first inequality. For $\gamma\ne0$ the bound $|u|\le1/(2\gamma^2)$ gives the second one, which also holds at $\gamma=0$.
For a positive isobaric product a nonpositive real coefficient forces one constituent coefficient to be nonpositive, so the same displacement statement applies to the product.
\end{proof}

\subsection*{Classical groups and their twists}
Let $G$ be one of the split groups $\mathrm{SO}_{2n+1}$, $\mathrm{Sp}_{2n}$ or $\mathrm{SO}_{2n}$, taking $n\ge2$ in the even orthogonal case. Its standard transfer has degree $N=2n,2n+1,2n$, respectively. We use the standard representation of the dual group in these degrees. This choice specifies the $L$-function of the group.

\begin{corollary}\label{li:classical}
Let $\pi$ be a globally generic cuspidal automorphic representation of $G$ over a number field. Theorems~\ref{li:everyheight} and~\ref{li:positive-tail} hold for its completed standard function and for its completed Langlands--Shahidi twist by any unitary cuspidal $\tau$ on $\mathrm{GL}_m$ over the same field. The quantitative displacement conclusion of Corollary~\ref{li:displacement} holds as well. Use the pole-removed divisor and the unitary normalization in each case.
\end{corollary}
\begin{proof}
Cogdell, Kim, Piatetski-Shapiro and Shahidi identify the standard transfer with an isobaric sum $\Pi=\boxplus_i\Pi_i$ of unitary cuspidal representations. Their local lift identifies the twisted $L$- and epsilon factors also at ramified places. The archimedean correspondence gives the gamma factors. In detail, use Theorems~7.1--7.2 for the isobaric image, Proposition~7.5 and Lemma~7.2 for all generic twists locally, and the proof of Theorem~7.4 to identify the local construction with the component of the global lift. Strong multiplicity one identifies the global lifts when the auxiliary place in that proof is changed. See \href{https://people.math.osu.edu/cogdell.1/functoriality-www.pdf}{\emph{Functoriality for the classical groups}, Publ. Math. IH\'ES 99 (2004), 163--233}. Hence the desired completion is the one for $\Pi\times\tau$; taking $\tau$ trivial gives the standard function. Corollary~\ref{li:isobaric} proves the bounds and the assertion about the divisor. The proof of Corollary~\ref{li:displacement} supplies the claimed witness if a real coefficient is nonpositive.
\end{proof}

\subsection*{A million coefficients for the conductor census}
For a holomorphic cuspidal newform, the gamma shifts have a fixed imaginary part equal to zero before any norm twist. This permits a shorter verified height. The following numerical consequences use the same unitary normalization as Theorem~\ref{li:explicit}.
\begin{corollary}\label{li:weight-census}
Every rational elliptic curve of conductor at most $1000$ has strictly positive Li coefficients through index $1\,021\,171$. This covers all $2463$ isogeny classes in the census. For all thirteen forms through height $4096$, the prefix reaches $389\,268\,779$. For each of those forms verified through $65\,536$, it reaches $128\,635\,042\,312$, and for the level-$32$ form individually it reaches $128\,892\,238\,059$.
\end{corollary}
\begin{proof}
Suppose a holomorphic cuspidal newform has weight $2\le k\le12$, level $N$ and unitary variable $s$. Write $b=\frac12\log N$, $v=(k-1)/2$, and assume the line has been verified through an integer $T\ge256$. At $2+it$ the two finite roots per prime have modulus at most one, so $|L'/L|\le 2[-\zeta'(2)]/\zeta(2)<1.4$. Indeed $-\zeta'(2)\le (\log2)/4+(1+\log2)/2$ by integration, and $\zeta(2)>3/2$. Since $|\psi(2+v+it)-\log(2+v+it)|\le0.4$, the Hadamard sum satisfies $S(t)\le b+\log(|t|+8)$.

Take $m\ge1605$ and the two signed windows of length $\ell=3m/16$ used above. The bound $|\log L(2+it)|\le 2\log\zeta(2)<2\log2$ shows that the mean of the real logarithmic derivative over either window has magnitude at most $4\log2/\ell$. Hence the mean of $S$ there is at least $b+\log m-3.65$, using $|t|\ge m/4$ and the constant $-\log(2\pi)$ in the completion. Here $h=\lfloor m/32\rfloor\ge50$. For $j\ge h$, either unit interval at distances $[j,j+1]$ has midpoint $w$ with $|w|+8\le17j$, since $|t|\le15m/32$. The Poisson kernel estimate gives total leakage outside the signed ordinate band at most
\[
  \frac{17}{2}\sum_{j\ge h}\frac{b+\log(17j)}{j^2}
  \le\frac{17}{2(h-1)}[b+\log(17(h-1))+1]
  <0.174 b+1.35.
\]
The angle bound $1-\cos(m\theta)\ge 5/4$ is valid for $m\ge1605$ as well, since $m\theta\ge2/(1+4/m^2)>1.99$ and $m\theta\le4$ on the band. Multiplying the retained mean by $5/4$ and the total window length divided by $\pi$ yields a positive base at least
\[
   \frac{15m}{32\pi}\left(\frac45 b+\log m-5\right).
\]
For the complementary high-zero loss, the two unit intervals with $j\le|\gamma|\le j+1$, $j\ge T$, together contain at most $(17/4)(b+\log(2j))$ zeros. Termwise integration as in \eqref{li:sinh} shows that the loss is at most $(17/4)y[b+\log(2y)+1/3](\sinh a/a-1)$, with $y=T-1$ and $a=m/(2y^2)$. Therefore it suffices that the following quotient at the desired final index $M$ be less than one, where $A=M/(2y^2)$:
\[
  \frac{136\pi}{15}\,\frac{y[b+\log(2y)+1/3]}{M[4b/5+\log M-5]}
           \left(\frac{\sinh A}{A}-1\right).
\]
For each $b$ the quotient increases with $m$ on $m\ge1605$ by the positive series argument. At a fixed endpoint its dependence on $b$ is a quotient of two affine functions with positive denominator. Consequently testing $b=0$ and $b=\frac12\log N_{\max}$ suffices for every $1\le N\le N_{\max}$. The termwise range $m\le2\pi T-\pi$ covers all indices below $1605$. The four inputs $(T,N_{\max},M)$ are
\[
\begin{array}{r|r|r}
256&1000&1\,021\,171\\
4096&144&389\,268\,779\\
65\,536&144&128\,635\,042\,312\\
65\,536&32&128\,892\,238\,059.
\end{array}
\]
The interval checks in \texttt{li\_tail.py} establish the strict inequality at both conductor endpoints in each row. Apply the height certificates in Sections~\ref{num:chapter} and~\ref{cens:chapter}.
\end{proof}

\subsection*{A height of sixteen for holomorphic newforms}
Let $f$ be a primitive holomorphic cuspidal eigenform over $\Q$ of any positive integral weight $k$, level $N$, and nebentypus. Use its unitary normalization, with gamma factor $(2\pi)^{-s}\Gamma(s+v)$ up to a constant, $v=(k-1)/2$. All finite reciprocal roots have absolute value at most one. At good primes this is Deligne's theorem for weight at least two, and the finite-image theorem for weight one. At bad primes use the newform Euler factors and temperedness. See \href{https://www.numdam.org/articles/10.24033/asens.1277/}{Deligne--Serre, \emph{Formes modulaires de poids 1}}, and \href{https://bookstore.ams.org/coll-53}{Iwaniec--Kowalski, Chapters 5 and 14}.

\begin{theorem}\label{li:holodyadic}
If every zero of the completion through real height $T\ge16$ lies on its critical line, then $\Re\lambda_m(f)>0$ for $1\le m\le T^2$. The threshold is independent of $k,N$ and the nebentypus. For $2\le k\le12$ the same assertion holds for $T\ge12$.
\end{theorem}
\begin{proof}
Write $b=(\log N)/2$ and $H(n)=b+\log(v+2+n)$. We first count the phase mass at any integer scale $n\ge1605$, restricting to the open bands $n/4<|\gamma|<n/2$. Put $S(t)=\Re(\Lambda'/\Lambda)(2+it)$ as in \eqref{li:poisson}. The real digamma error is at most $1/(2+v)\le1/2$. Also $|L'/L|\le-2\zeta'(2)/\zeta(2)<1.3$. For the last inequality bound $\sum_{j\le10}(\log j)/j^2$ directly, the remaining terms by $(1+\log 10)/10$, and use $\zeta(2)>1.6$. Thus $S(w)\le b+\log(v+2+|w|)$.

Average over either signed interval $9n/32\le|t|\le15n/32$. Their individual length is $3n/16$. For $r=9/32$, Cauchy--Schwarz gives $|v+2+it|\ge r(v+2+n)/\sqrt{1+r^2}$. The integral of the finite logarithmic derivative costs at most $4\log2$ at the two endpoints. Together these estimates give mean at least $H(n)-3.67$, using
\[
 \log(\sqrt{1+r^2}/r)+\log(2\pi)+1/2+
                 4\log2/(3\cdot1605/16)<3.67 .
\]
Put $h=\lfloor n/32\rfloor$. Every zero outside the corresponding open band is at least $n/32$ from $t$. For the two unit intervals of distances between $j$ and $j+1$, $j\ge h$, each midpoint $w$ satisfies $|w|+2\le17j$. Also
\[
 b+\log(v+17j)\le H(n)+\log\big(j/(h-1)\big)
\]
since $17(h-1)\le n$ and $j/(h-1)\ge1$. A unit interval contains at most $(17/8)S(w)$ zeros and a distant zero contributes at most $2/j^2$ to $S(t)$. Integration of a decreasing majorant bounds total leakage by
\[
 \frac{17}{2(h-1)}[H(n)+1].
\]
Here $h\ge50$. Thus the retained mean is at least $0.825H(n)-3.85$. Each retained kernel integrates to at most $\pi$. The two intervals and the phase bound $1-\cos(n\theta)\ge 5/4$ give
\[
 \sum_{n/4<|\gamma|<n/2}(1-\cos(n\theta))\ \ge\
          \frac{15n}{32\pi}\,[0.825H(n)-3.85].
\]

Take $u=\max(m,1605)$ and $n=2^j u$ for $j=0,1,\ldots$. The open bands are disjoint. Since $n\theta/2\le2$ there, concavity of $\sin$ on $[0,\pi]$ gives $1-\cos(m\theta)\ge(m/n)^2(1-\cos(n\theta))$. Hence the full phase mass satisfies
\[
      B_m\ \ge\ \frac{15}{16\pi}\,\frac{m^2}{u}
                         [0.825H(u)-3.85].
\]
This bound includes zeros at arbitrarily high ordinates.

At an integer verified height $J$ put $y=J-1$, $a=m/(2y^2)$. The two unit shells of absolute height $j$ through $j+1$ contain at most $(17/4)[b+\log(v+2+j+1/2)]$ zeros. The same moment integral as in \eqref{li:sinh} therefore gives
\[
 E_m\le\frac{17}{4}y[b+\log(v+2y+4)+1/3]
                 \left(\frac{\sinh a}{a}-1\right).
\]
Indeed $v+2+j+1/2\le(v+2y+4)j/y$ for $j\ge J$. The moment integral of $\log(t/y)$ contributes at most $1/3$ after division by the common factor. These bounds imply strict positivity if $E_m/B_m<1$. The resulting majorant for this ratio increases with $m$: the series for $(\sinh a/a-1)/m^2$ has nonnegative coefficients, and $u/(0.825H(u)-3.85)$ increases for $u\ge1605$, by differentiation.

It suffices to test $m=(J+1)^2$, since $J=\lfloor T\rfloor$. Here $v+2y+4\le v+u+2$. Dividing the two bounds, use the decreasing function $(H+1/3)/(0.825H-3.85)$ at $H=\log u$. For $16\le J\le40$ direct interval arithmetic bounds every one of the resulting 25 ratios by $0.993$. For $J\ge41$, $u=m$, $a\le(42/40)^2/2$ and $\sinh(a)/a-1<0.053$. The decreasing factor is at most $7/2$ since $u\ge1605$. The ratio is then at most $(68\pi/15)(7/2)(0.053)/J<1$.

For $2\le k\le12$, the preceding subsection gives phase mass at least $(15n/(32\pi))(4b/5+\log n-5)$ on the union $n/4<|\gamma|<n/2$ of the two signed bands, at every scale $n\ge1605$. Its proof bounds the complementary contribution even when a zero is at a band endpoint. Summing the geometric series, using $\sum_{j\ge0}j2^{-j}=2$, gives
\[
 B_m\ge\frac{15}{16\pi}\frac{m^2}{u}
                         [4b/5+\log(2u)-5].
\]
The corresponding upper bound for $E_m$ is that in Corollary~\ref{li:weight-census} with $b+\log(2y)+1/3$. Its derivation works at $J\ge12$: for $j\ge12$ the midpoint satisfies $|w|+8\le 2j$. The ratio again increases with $m$. For $12\le J\le15$ apply the last two bounds at $m=(J+1)^2$. Inspect $b=0$ and the limit $b\to\infty$, the two extreme values of a quotient of affine functions. All eight values are less than $0.92$. The larger heights already follow from the first part.
\end{proof}

This dyadic estimate gives longer certified prefixes on the actual curves. For $2\le k\le12$, $N\le N_{\max}$, test the ratio at $m=M$, both for $b=0$ and for $b=\frac12\log N_{\max}$. The tests are strict and cover smaller indices by the monotonicity in the proof.
\[
\begin{array}{r|r|r}
 J&N_{\max}&\text{positive indices through }M\\\hline
 32&9999&10\,781\\
 64&9999&53\,664\\
 256&1000&1\,150\,720\\
 4096&275&418\,442\,908\\
 4096&144&418\,937\,857\\
 65\,536&144&135\,799\,925\,615\\
 65\,536&32&136\,063\,625\,162\\
 131\,072&1&574\,565\,494\,060\\
 131\,072&32&571\,912\,188\,515\\
 131\,072&11&572\,678\,194\,104\\
 262\,144&11&2\,400\,255\,893\,995\\
 524\,288&11&10\,035\,901\,614\,218
\end{array}
\]
The interval tests, including the independent-weight estimates and their integer ranges, are in \texttt{li\_dyadic\_holomorphic.py}. Each numerical consequence uses a zero-height certificate for the named conductor range or form. Increasing the conductor range requires checking its own critical-line hypothesis.

In particular the genuinely quadratic isogeny class of Corollary~\ref{qt:genuine} and the base change to $\Q(\zeta_5)$ both have strictly positive real Li coefficients through $418\,442\,908$. Their factors over $\Q$ are among the newforms of levels $11$ and $275$ whose zeros are certified in Section~\ref{qt:chapter}. Their weights are two and the real parts of the Li coefficients add across those factors.

\subsection*{Li positivity with an unresolved central divisor}
The lower bound for $B_m$ also applies to approximate enclosures of low zeros. This matters for $234446.\mathrm{a1}$, where a pair of zeros could coincide at the center or split in either direction.

\begin{proposition}\label{li:uncertain-center}
For the curve $234446.\mathrm{a1}$ the first $96\,251\,863$ Li coefficients are strictly positive, in all three central configurations of Theorem~\ref{r4:cluster}.
\end{proposition}
\begin{proof}
Use unitary coordinates. Put $M=96\,251\,863$, $J=2048$, $b=(\log234446)/2$, and $\delta=2^{-4108}$. At most two zeros of absolute ordinate at most $J$ can lie off the line. Each of these lies in $|\rho-1/2|<\delta$. By the triangle inequality,
\[
 \big|\log|1-1/\rho|\big|
 \le\log\frac{1+2\delta}{1-2\delta}
 \le \frac{4\delta}{1-2\delta}<5\delta .
\]
In the phase decomposition the two together can cost at most $2(\cosh(5m\delta)-1)\le50m^2\delta^2$ for $1\le m\le M$, since $5M\delta<1$.

For $u=\max(m,1605)\le M$, the lower bound just proved gives
\[
 B_m\ge\frac{15}{16\pi}\frac{m^2}{u}
       [4b/5+\log(2u)-5]>\frac{m^2}{4M}.
\]
Here the bracket exceeds one already at $u=1605$. The estimate for the error from $|\gamma|>J$ used only the strip and the unit-shell counts. Its ratio to $B_m$ is bounded by the increasing majorant of the preceding proof. Evaluating that majorant at $M$ gives a value below $0.99999995$. The additional central loss divided by $B_m$ is less than $200M\delta^2<10^{-8}$. Their sum is below one, as required. The final inequalities are checked by outward-rounded interval arithmetic in the same supplementary program.
\end{proof}

\subsection*{Real archimedean parameters and symmetric powers}
The argument with distant bands applies to other weights and degrees. Consider a cuspidal standard function or a cuspidal Rankin--Selberg pair over a number field. Suppose the actual finite reciprocal roots have absolute value at most one, also at ramified primes, and that all the shifts $\mu_j$ in the $d$ real gamma factors are real and nonnegative. Put $Q$ for the absolute conductor. No upper bound on $Q,d$ or the $\mu_j$ is required.

\begin{theorem}\label{li:hodge}
For such an entire completion a verification through any real $T\ge18$ gives $\Re\lambda_m>0$ for $1\le m\le T^2$. If the completion has a reflected pair of simple poles, the same statement for its pole-removed Li coefficients holds for $T\ge23$. For the symmetric square, cube and fourth power of $43.\mathrm{a1}$ the verified prefixes reach $1\,106\,109$, $1\,105\,385$ and $9\,885$, respectively.
\end{theorem}
\begin{proof}
Use $u=\max(m,1605)$ and
\[
 H(x)=d^{-1}\left[\log Q+\sum_{j=1}^d\log(2+\mu_j+x)\right].
\]
Rescale $S(t)$ to $2S(t)/d$. The Euler contribution to this last quantity has magnitude less than $1.3$ by the prime ideal norm grouping and the bound for $-2\zeta'(2)/\zeta(2)$ in the preceding proof. In the logarithmic derivative of $\Gamma_{\R}(s+\mu_j)$ the digamma argument at $2+it$ has real part at least one. After rescaling, the mean of the digamma errors is therefore bounded by one. The constant term from the real gamma factors is $-\log(2\pi)$. Since $1.3+1-\log(2\pi)<0.5$, the new upper bound is
\[
               2S(w)/d\le H(|w|)+U ,
\]
where $U=0.5$ for an entire completion and $U=3.5$ after removal of two simple poles. Indeed their rescaled real contribution is between zero and $3/d$. On the two signed intervals the lower mean is at least $H(n)-4.17$: the Euler integral after normalization costs at most $4\log2$ at the endpoints, and each of the real shifts satisfies the same Cauchy--Schwarz lower bound as $v$ in the preceding proof. The increase from the previous constant $3.67$ is the increase of the digamma error from $1/2$ to $1$.

For $j\ge h=\lfloor n/32\rfloor$, the numerator in each remote shell is bounded by $H(n)+U+\log(j/(h-1))$. Hence after removing the complement of the open bands the retained mean exceeds $0.825H(n)-V$, where $V=4.44$ for an entire completion and $V=4.96$ with poles: subtract at most $17[H(n)+U+1]/98$. Undo the normalization $2/d$ when integrating the kernel. The dyadic sum and the moment bound become
\[
\begin{split}
 B_m&\ge \frac d2\,\frac{15}{16\pi}\frac{m^2}{u}
                 [0.825H(u)-V],\\
 E_m&\le \frac d2\,\frac{17}{4}y[H(2y+2)+U+1/3]
                        \left(\frac{\sinh a}{a}-1\right),
\end{split}
\]
where $J=\lfloor T\rfloor$, $y=J-1$, $a=m/(2y^2)$. The second estimate uses $2+\mu_j+k+1/2\le(2+\mu_j+2y+2)k/y$ for $k\ge J$. The factors $d/2$ cancel. As before the ratio of the majorants increases with $m$, since $uH'(u)\le1$ and $0.825\log1605-V>0.825$ in both cases.

At $m=(J+1)^2$ one has $2y+2\le u$ and $H(u)\ge\log u$. Thus the conductor-dependent ratio $(H(2y+2)+U+1/3)/(0.825H(u)-V)$ is at most $(\log u+U+1/3)/(0.825\log u-V)$. For $18\le J\le40$ in the entire case the resulting loss-to-mass ratio is below $0.986$. With poles it is below $0.91$ for $23\le J\le40$. For $J\ge41$ the logarithmic ratio is less than $10$, $\sinh(a)/a-1<0.053$, and $y/u\le1/J$. The ratio of the bounds for $E_m$ and $B_m$ is therefore at most $(68\pi/15)(10)(0.053)/J<1$. This proves both uniform statements.

Finally retain the actual $H$ in those two bounds. The respective degrees are $3,4,5$, the conductors are $43^2,43^3,43^4$, and their real gamma shifts are
\[
 (1,1,2),\qquad(1/2,3/2,3/2,5/2),\qquad(0,1,2,2,3).
\]
Their finite local bounds follow from Hasse at the good places and the single Steinberg block at $43$. The heights already certified in the symmetric-power sections are $256,256,32$. Direct interval evaluation verifies $E_m<B_m$ at each of the three stated terminal indices. All numerical inequalities in this proof are checked by \texttt{li\_dyadic\_parameters.py}.
\end{proof}

\subsection*{Larger prefixes from dyadic ordinate windows}
\begin{theorem}\label{li:small-universal}
Let $L$ be a standard unitary cuspidal function over a number field or the Rankin--Selberg function of two such representations. Use unitary normalization and the zero divisor of the pole-cleared completion. If every zero with $|\Im\rho|\le T$ lies on the critical line, then
\[
                 \Re\lambda_m>0\qquad(1\le m\le bT^2,\quad T\ge T_0)
\]
for every row of the following table.
\[
\begin{array}{c|rrrrrrrrrr}
T_0&10&12&16&32&64&128&256&1024&4096&65536\\
\hline
b&1.7&3&4.5&7.4&9.7&11.6&13.5&16.8&20&26.2
\end{array}
\]
Each threshold is independent of the ranks, field, conductor and archimedean parameters. Here $T$ may be real.
\end{theorem}
\begin{proof}
Write $d$ for the number of real gamma factors and let $H(x)=\log Q+\sum_{j=1}^d\log(8+|\mu_j|+x)$ as in Theorem~\ref{li:uniform}. Throughout the proof put $z=d/H(T)$, $q=13/3$, and $P=\pi/2+5/3$. At the abscissa $3$, the Poisson sum is
\[
 S(t)=\sum_\rho\frac{3-\beta}{(3-\beta)^2+(t-\gamma)^2}.
\]
Its logarithmic derivative identity includes the positive contribution $p(t)$ from the pole-clearing polynomial. In the case of a pole pair, $p(t)$ is the sum of two translates of $b/(b^2+t^2)$ with $b=2,3$; in other cases it is zero. The local parameters at finite places have modulus at most the residue-field size, and $\Re\mu_j\ge-1$. Hence $|\log L(3+it)|\le d\log\zeta(2)$ on the Euler-product branch, and $|L'/L(3+it)|\le d[-\zeta'(2)/\zeta(2)]<0.58d$. For the last numerical bound sum $\Lambda(n)/n^2$ up to $2048$ and majorize the rest by $(1+\log 2048)/2048$.

For $\Re w>0$ one has
\[
 \psi(w)=\log w-\frac1{2w}+R(w),\qquad |R(w)|\le\frac1{12(\Re w)^2}.
\]
Indeed in the Laplace integral for $\log w-\psi(w)$, the integrand left after subtracting $1/2$ is $\tfrac12\coth(u/2)-1/u$, between zero and $u/12$ by the partial fractions of $\coth$ and $\sum n^{-2}=\pi^2/6$. In this proof $w=(3+\mu_j+it)/2$ has real part at least one. Define
\[
 c=-\tfrac12\log(2\pi)+1/24+0.58 .
\]
The sign of $-\Re(1/w)$ gives $S(t)\le H(|t|)/2+cd+p(t)$. In a four-unit interval every zero contributes at least $1/q$ to $S$ evaluated at its midpoint, since $2\le3-\beta\le3$ and $\max_{2\le b\le3}(b+4/b)=13/3$. For any set of sample points separated by at least four,
\[
                    \sum p(t)\le P .
\]
For a single kernel split the sample points at its center. On each side its sum is at most $1/b+\tfrac14\int_0^\infty b/(b^2+u^2)\,du$. Summing the two poles proves the claim, including every imaginary translation.

\paragraph{The high-zero loss.} For the selected table entry let $C_b=[\cosh(b/2)-1]/b^2$ and $S_s(T)=\sum_{h\ge0}(T+4h)^{-s}$ for $s>1$. For $m\le bT^2$, $|u_\rho|\le1/(2\gamma^2)$ above $T$. The positive series of $\cosh$ therefore gives $\cosh(m u_\rho)-1\le C_bm^2/\gamma^4$. On the positive and negative shells with lower absolute height $v=T+4h$, bound counts at the two midpoints. Their collective set is four-separated. Since
\[
 H(v+2)\le H(T)+d\log(1+2/T)+d\log(v/T),\qquad c+\tfrac12\log(1+2/10)<0,
\]
using $\log u\le u/e$ for $u\ge1$ gives the following bound for the entire loss $E_m$:
\begin{equation}\label{li:dyadic-loss}
 \frac{E_m}{m^2H(T)}\le C_bq\left[S_4(T)+z\left(\frac{S_3(T)}{eT}+\frac{P}{T^4}\right)\right].
\end{equation}
For $s>1$ the quantity $T^{s-1}S_s(T)$ decreases in $T$. The gamma integral expresses it as an integral of $v^{s-1}e^{-v}/[T(1-e^{-4v/T})]$ divided by $\Gamma(s)$. Its denominator increases, by $1+x\le e^x$. In applications to \eqref{li:dyadic-loss}, fix $G_0=\max(144,bT_0^2)$ and put $G=\max(G_0,m)\le\max(G_0,bT^2)$. Both $\max(G_0,bT^2)/T^3$ and $\max(G_0,bT^2)/T^4$ decrease for $T\ge T_0$.

\paragraph{A pair of spectral windows.} Indices $m\le2\pi T-\pi$ are already covered. Otherwise $G\ge m>T$. For $X=2^jG$, use the disjoint open ordinate bands $X<|\gamma|<2X$, with interior windows $I=[10X/9,17X/9]$ and $-I$, of length $\ell=7X/9$ each. For $r=|3+\mu|$,
\[
 \tfrac12\log|3+\mu+it|+\tfrac12\log|3+\mu-it|
          \ge \tfrac12\log(t+r)+\tfrac12\log|t-r|\quad(t>0),
\]
away from the single point $t=r$, by the reverse triangle inequality for $(3+\mu)^2+t^2$. If $r\le2X$, the first term is at least $\tfrac12\log(X+r)$ and the second averages at least $\tfrac12[\log(\ell/2)-1]$. The latter assertion follows by centering the interval of length $\ell$ at the singularity, which minimizes the integral of $\log|t-r|$. If $r\ge2X$, use $r^2-t^2\ge r^2-(17X/9)^2$. The relative cost against $\log(r+X)$ is at most $\log(3/\sqrt{4-(17/9)^2})$. The first case costs at most $(1+\log(54/7))/2$. Both costs are smaller than $1.53$. Finally $8+|\mu|+X\le(1+11/X)(r+X)$.

For a single gamma parameter the reciprocal term in the digamma estimate has integrated absolute cost at most $\pi/2$ on either window, since $\int_{\mathbb R}\Re(1/w)\,dt=2\pi$. The real Euler logarithmic derivative costs at most $2d\log\zeta(2)/\ell$ in average by its logarithmic primitive. The average of $S$ on the union is thus at least $H(X)/2-d A_X$, where
\[
 A_X=\tfrac12[1.53+\log(1+11/X)+\log(2\pi)]+\tfrac1{24}
          +\frac{2\log\zeta(2)+\pi/2}{\ell} .
\]

At $t\in I$ remove the zeros outside $(X,2X)$. On the left and right, the respective distances to the excluded ordinates are at least $D_-=t-X$ and $D_+=2X-t$. For either side use consecutive four-unit intervals at distances $[4h,4h+4]$ with $4h\ge4\lfloor D_\pm/4\rfloor$. Put $Y_\pm=D_\pm-8>0$. For each such interval the kernel bound per zero can use $5/(2(4h)^2)$ after averaging horizontal reflections, since the numerators $3-\beta$ average to $5/2$. The pole part is bounded using four-separation. For the regular part one has at each midpoint
\[
 H(|t\pm(4h+2)|)\le H(X)+d\log(C_\pm 4h/Y_\pm),\qquad C_-=3,\ C_+=2.
\]
Indeed $|t\pm(4h+2)|\le t+4h+2$. At $4h=Y_\pm$ the last expression is at most $C_\pm X$, and the right-hand majorant $C_\pm X(4h)/Y_\pm$ grows with slope at least one. The same estimates apply on $-I$ with the two sides interchanged. Comparing the resulting decreasing logarithmic densities divided by $u^2$ with their integrals from $Y_\pm$ bounds the average leakage on either window by
\[
 B_X\left[\tfrac12H(X)+d\left(c+\tfrac12+\tfrac14\log6\right)\right]
       +\frac{5qP}{\ell (X/9-8)},\qquad
 B_X=\frac{5q}{4\ell}\log\frac{8X/9-8}{X/9-8}.
\]
For detail, the integrals on either side use $\int_Y^\infty [D+\log(Cu/Y)]u^{-2}\,du=[D+1+\log C]/Y$. The two averages of $Y_\pm^{-1}$ both equal $\ell^{-1}\log[(8X/9-8)/(X/9-8)]$. The pole costs are at most $(5/2)qP(Y_-^{-2}+Y_+^{-2})$, whose mean is at most the last term displayed. The regular majorants decrease on $[Y_\pm,\infty)$: differentiating gives this directly from $H(X)\ge d\log(8+X)$, $X\ge144$ and the stated $c$.

Put
\[
 U_X=(1-B_X)/2,\qquad V_X=A_X+B_X(c+\tfrac12+\tfrac14\log6)
                          +\frac{5qP}{\ell(X/9-8)}.
\]
The retained average divided by $H(T)$ is at least
\[
 R_X(z)=\max\left(0,\ U_X\max(1,z\log(8+X))-zV_X\right)
\]
when $U_X>0$; take $R_X=0$ otherwise. Here $H(X)\ge\max(H(T),d\log(8+X))$ and $1/H(T)\le z$ since $d\ge1$. Each retained zero integrates to at most $\pi$ on its signed window. Therefore the two signed open bands together contain at least $2\ell H(T)R_X(z)/\pi$ zeros. This bound remains valid at band endpoints, since endpoint zeros were counted with the discarded part.

\paragraph{A positive mass across dyadic bands.} Write $X_j=G_0\cdot2^j$. The functions $A_X$, $B_X$ and $V_X$ decrease for $X\ge144$ and $R_X(z)$ increases. The reference angle at index $G$ on the $j$th band is at most one and at least $2^{-j-1}/[1+(2X_j)^{-2}]$. This follows from $|\gamma|/(1+\gamma^2)\le\theta_\rho\le1/|\gamma|$. On $[0,\pi]$ one has $\sin(vx)\ge v\sin x$ for $0\le v\le1$. Hence the actual $1-\cos(m\theta_\rho)$ is at least $(m/G)^2 w_j$, where
\[
 w_j=1-\cos\left(\frac{2^{-j-1}}{1+(2X_j)^{-2}}\right).
\]
Summing the first fourteen disjoint bands gives
\[
 \sum_\rho [1-\cos(m\theta_\rho)]\ge\frac{m^2H(T)}{G}\,C(z),\qquad
 C(z)=\frac{14}{9\pi}\sum_{j=0}^{13}2^j w_j R_{X_j}(z).
\]
All comparisons retain $z$ until this step; $0\le z\le1/\log(8+T_0)$. By \eqref{li:dyadic-loss} and its monotonicity calculation it suffices to verify on this compact interval
\[
 G_0 C_b q\left[S_4(T_0)+z\left(\frac{S_3(T_0)}{T_0e}+\frac P{T_0^4}\right)\right]< C(z).
\]
The directed check in \texttt{li\_dyadic.py} bounds the ratio by $0.9999$ for each row on every one of $4096$ equal cells. On $[z_0,z_1]$ it evaluates the numerator at $z_1$ and bounds each positive part from below by using $z_0$ in the maximum and $z_1$ in the subtraction. To bound $S_s(T_0)$ take the first $128$ terms and integrate the remainder from $T_0+4\cdot127$. Thus the unverified high zeros cost strictly less than the positive mass. This proves the theorem.
\end{proof}

\subsection*{Recovering the best all-height strip from square indices}
There is a converse question for a sparse sequence of tests. A negative Li coefficient detects an off-line zero. Can a function with an off-line zero evade detection along all square indices? The answer also determines the optimal strip of the zero divisor by applying the tests to norm translates. We use the power-sum method of \href{https://websites.umich.edu/~lagarias/doc/bombieri.pdf}{Bombieri--Lagarias} and \href{https://www.epfl.ch/labs/pde/wp-content/uploads/2020/03/Lecture2.1.pdf}{Weyl's polynomial equidistribution theorem}.

\begin{theorem}\label{li:sparse}
For a standard or Rankin--Selberg completion as in Theorem~\ref{li:everyheight}, let $P\in\Z[X]$ be nonconstant, with $P(0)=0$ and positive leading coefficient. Then
\[
 \limsup_{k\to\infty}
 \big(1+\max(0,-\Re\lambda_{P(k)})\big)^{1/P(k)}
     =\sup_\rho\left|1-\frac1\rho\right|.
\]
ERH is equivalent to $\Re\lambda_{P(k)}\ge0$ for all sufficiently large $k$. If ERH fails, the set of $k$ witnessing a strictly negative value has positive lower natural density.
\end{theorem}
\begin{proof}
Write $z_\rho=1-1/\rho$. When all zeros lie on the line, $|z_\rho|=1$ and the real summands are nonnegative. Otherwise reflection gives at least one $|z_\rho|>1$. Since $|z_\rho|\to1$ as $|\Im\rho|\to\infty$ and there are only finitely many zeros in a bounded interval of height, the supremum $R>1$ is attained by a nonempty finite multiset $A$. There is $r$ with $1<r<R$ for which every other $|z_\rho|\le r$. Writing $z_\rho=R\exp(2\pi i\alpha_\rho)$ for $\rho\in A$, we obtain
\[
 \Re\lambda_j=-R^j\sum_{\rho\in A}\cos(2\pi j\alpha_\rho)
                         +O_L(j\log(2j)\,r^j).
\]
Indeed all zeros in $A$ have absolute height at most $j$ for large $j$. The other zeros of height at most $j$ each contribute at most $1+r^j$ in absolute value and there are $O_L(j\log(2j))$ of them. The remaining real tail is $O_L(j\log(2j))$ by Theorem~\ref{li:positive-tail} with $m=T=j$. This also proves the desired upper bound for the limsup.

Let $H$ be the closure of the cyclic group generated by $\boldsymbol\alpha=(\alpha_\rho)_{\rho\in A}$ in the corresponding finite-dimensional torus. This closed subgroup has a connected identity component $H^0$ of finite index. Let $b$ be the order of $H/H^0$. The vector $b\boldsymbol\alpha$ generates a dense subgroup of $H^0$: multiplication by $b$ maps $H$ onto $H^0$, since it kills the component group and is surjective on the torus $H^0$. For $k=b\ell$ we have
\[
 P(k)\boldsymbol\alpha=\big(P(b\ell)/b\big)\, b\boldsymbol\alpha.
\]
The scalar polynomial in parentheses has integer coefficients and is nonconstant. Every nontrivial character of $H^0$ takes an irrational argument at $b\boldsymbol\alpha$; otherwise its image on a dense subgroup would be finite. Weyl's theorem and the character criterion give equidistribution in $H^0$ as $\ell$ runs over the positive integers. The conclusion includes the case $H^0=\{0\}$. A neighborhood of zero where every indicated cosine exceeds $1/2$ contains an open set of positive Haar measure with null boundary. Hence a set of $k$ of positive lower natural density has sum of cosines at least $|A|/2$. The error above is $o_L(R^j)$, so on this set and for all sufficiently large $k$,
\[
                   \Re\lambda_{P(k)}\le-(|A|/3)R^{P(k)} .
\]
This proves the lower bound, density assertion and equivalence.
\end{proof}

For each $\tau\in\Q$ use the translated function $L(s+i\tau)$ with the corresponding translated pole divisor. This is a unitary norm twist of the standard representation, or of one factor in a Rankin--Selberg pair. Let $B$ be the supremum of the displayed limsups over rational $\tau$, taking even just $P(k)=k^2$. In unitary normalization the least possible upper edge of the all-height strip is
\[
           \sup_\rho\Re\rho=\frac{B}{1+B},
\]
with value $1$ when $B=\infty$. Indeed a zero $\beta+i\gamma$ contributes the radius
\[
 \left(\frac{(1-\beta)^2+(\gamma-\tau)^2}
                  {\beta^2+(\gamma-\tau)^2}\right)^{1/2}.
\]
For $\beta\le1/2$ its supremum over rational $\tau$ is $(1-\beta)/\beta$, by density. For $\beta>1/2$ its supremum is one. Horizontal reflection completes the proof. Thus for the all-height unitary Hecke strip of Section~\ref{recentcm:section} these growth factors are at most seven. ERH corresponds exactly to one.

\begin{theorem}[Prime-indexed tests]\label{li:prime-sparse}
Let $P\in\Z[X]$ be a nonconstant polynomial with positive leading coefficient. Suppose that for each positive integer $b$ there is a residue $a$ coprime to $b$ with $P(a)\equiv0\pmod b$. For every completion in Theorem~\ref{li:sparse}, ERH is equivalent to
\[
           \Re\lambda_{P(p)}\ge0
           \quad\text{for every sufficiently large prime }p .
\]
If ERH fails, a positive lower proportion of the primes have $\Re\lambda_{P(p)}<0$. The limsup of $(1+\max(0,-\Re\lambda_{P(p)}))^{1/P(p)}$ along primes equals $\sup_\rho|1-1/\rho|$. In particular one may use $P(p)=p-1$, $p+1$, or $(p-1)^k$ for any fixed $k\ge1$.
\end{theorem}
\begin{proof}
Only the lower bound in the off-line case requires a change. Use the torus $H$, its identity component $H^0$, the integer $b=|H/H^0|$ and the dominant angle vector $\boldsymbol\alpha$ from the preceding proof. Choose $a$ as in the hypothesis. For every prime $p\equiv a\pmod b$, write $P(p)\boldsymbol\alpha=(P(p)/b)b\boldsymbol\alpha\in H^0$. These points are equidistributed in $H^0$ along this progression. Indeed for any nontrivial character of $H^0$, write its value at the dense generator $b\boldsymbol\alpha$ as $e^{2\pi i\xi}$. Then $\xi$ is irrational. The character at $P(p)\boldsymbol\alpha$ is $e^{2\pi i\xi P(p)/b}$. In the finite Fourier expansion of the progression indicator every resulting polynomial phase has the form
\[
                   \xi P(x)/b+\ell x/b+\mathrm{constant}.
\]
It has an irrational nonconstant coefficient. The prime polynomial equidistribution theorem gives mean zero for each such phase over the primes. Divide by the prime number theorem density $1/\varphi(b)$ to get mean zero along the progression. The character criterion proves the assertion, also for the trivial torus. The cited prime equidistribution theorem is due to Vinogradov and Rhin. \href{https://doi.org/10.4064/aa-23-3-217-248}{G.~Rhin, Acta Arith. 23 (1973), 217--248.}

Select a positive-measure continuity set around zero on which all dominant cosines exceed $1/2$. Its frequency among all primes in the progression, measured relative to all primes, is its Haar measure divided by $\varphi(b)$. On these primes the dominant power sum makes the coefficient at most $-(|A|/3)R^{P(p)}$ for all sufficiently large $p$. This proves the assertion and the exponential limsup. The on-line case follows termwise. For the three displayed examples use $a=1,-1,1$, respectively.
\end{proof}

Consequently the optimal all-height strip in the preceding formula can be recovered using only the ordinary Li coefficients at indices $p-1$ of rational vertical translates. A positive-density assertion here counts the prime parameters, even when the indices themselves have zero density among the integers.

\begin{theorem}[Tests with arbitrarily slowly growing counting function]\label{li:thin}
Let $h$ be any nondecreasing unbounded function on the positive integers. There is an infinite set $S$ of indices, depending only on $h$, such that $\#(S\cap[1,X])\le h(X)$ for every sufficiently large $X$. For every standard or Rankin--Selberg function in Theorem~\ref{li:sparse}, ERH is equivalent to eventual nonnegativity of $\Re\lambda_m$ on $m\in S$.

More precisely, suppose ERH fails and use $R>1$ and the multiset $A$ of dominant transformed zeros from Theorem~\ref{li:sparse}. For every sufficiently large positive integer $M$, some $1\le j\le 8^{|A|}$ satisfies
\[
                  \Re\lambda_{jM}<-\tfrac12 |A|R^{jM}.
\]
The set $S$ can be chosen so that
\[
 \limsup_{\substack{m\to\infty\\m\in S}}
   \frac{\max(0,-\Re\lambda_m)}{R^m}=|A|.
\]
\end{theorem}
\begin{proof}
Write $d=|A|$ and use the dominant angle vector $\boldsymbol\alpha\in(\R/\Z)^d$ from Theorem~\ref{li:sparse}. For each positive integer $q\ge2$, sort the $q^d+1$ points $i M\boldsymbol\alpha$ for $0\le i\le q^d$ into the $q^d$ half-open cubes of side $1/q$ in $[0,1)^d$. Two occupy one cube. Subtract their indices to obtain $1\le j\le q^d$ with every coordinate of $jM\boldsymbol\alpha$ within $1/q$ of an integer. Take $q=8$. The $d$ dominant cosines are all at least $1/\sqrt2$. The normalized error in the dominant power sum tends to zero uniformly over indices $n\ge M$ as $M\to\infty$, since $n\log(2n)(r/R)^n\to0$. This proves the strictly negative estimate for every sufficiently large $M$.

Choose integers $M_k$ recursively with $M_{k+1}>kM_k$ and $h(M_k)\ge k(k+1)/2$. Take
\[
             S=\bigcup_{k\ge1}\{M_k,2M_k,\ldots,kM_k\}.
\]
For $M_k\le X<M_{k+1}$ its counting function is bounded by $k(k+1)/2\le h(X)$. For any one fixed function failing ERH let $q_k=\lfloor k^{1/d}\rfloor$. For all large $k$ this tends to infinity and $q_k^d\le k$. The cube argument finds $1\le j_k\le k$ with all dominant cosines at $j_kM_k$ tending to one. Consequently $-\Re\lambda_{j_kM_k}/R^{j_kM_k}\to d$. The absolute upper bound $d+o(1)$ for the normalized negative part follows from the same power-sum expansion. A function satisfying ERH has nonnegative real coefficients at every index. These observations prove the conclusions with a single $S$ for all functions.
\end{proof}

The short blocks of multiples are a substantive part of the construction. For example, in the class of abstract horizontally reflected zero divisors take the four transforms $z=R e^{2\pi i/3},\bar z,1/z,1/\bar z$, with $R>1$ and corresponding zeros $\rho=1/(1-z)$. They lie strictly in the unitary strip. In fact the first has real part $(1+R/2)/(1+R+R^2)<1/2$, and reflection accounts for its partner. The real Li sum is
\[
        4-2(R^m+R^{-m})\cos(2\pi m/3).
\]
It is positive at every index $m=2^k$, although the divisor contains off-line zeros. This example concerns the power-sum argument for a reflected divisor; it asserts no realization by an automorphic function.

\clearpage
\section{A symmetric zero count for the rank-two curve of conductor 389}\label{r2:chapter}
The completed $L$-function of an elliptic curve of positive even analytic rank vanishes at the midpoint of the usual counting rectangle. Counting only a half rectangle with the real axis as bottom edge would place the central zero on the integration path. For the following curve we count on both sides of the real axis instead. We keep the arithmetic variable, so the line of symmetry is $\operatorname{Re}s=1$.
\[
          E:\ y^2+y=x^3+x^2-2x,\qquad N=389.
\]
Our computation certifies all the zeros in the indicated window, including multiplicities, from this equation. Its interval for the leading central coefficient supplies an independent check on the multiplicity at the point that defeats a sign-change detector.

\begin{theorem}\label{r2:window}
For $E$ every zero of its completed Hasse--Weil function with $|\operatorname{Im}s|\le2048$ lies on $\operatorname{Re}s=1$. The zero at $s=1$ is double and every other zero in this range is simple. There are $5064$ positive-ordinate zeros in this window. Moreover
\[
    0.7593165002 < L''(E,1)/2 < 0.7593165004.
\]
\end{theorem}
The form, its analytic rank, and high-precision central derivatives have long been tabulated. We use no numerical rank or zero as input. What follows establishes the finite zero count with reproducible enclosures and makes no assertion of a first computation. The following proof and the independently runnable \texttt{checks/elliptic-389/certify.py} are two parts of the certificate.

\subsection*{Exact vanishing and parity}
The Weierstrass invariants are $c_4=112$, $c_6=-856$ and $\Delta=(112^3-856^2)/1728=389$. The equation is minimal, since the discriminant is squarefree. The only bad prime is $389$, which is multiplicative as $389\nmid112$. It is split since $-c_6\equiv95^2\pmod {389}$. The conductor is therefore $389$ and the two negative local root numbers, at $389$ and at infinity, give global sign $+1$. Modularity identifies the $L$-series with that of a normalized weight-two newform $f$. Its completion is
\[
       F(z)=B^{-(1+z)}\Gamma(1+z)L(E,1+z),
       \qquad B=2\pi/\sqrt{389},\qquad F(-z)=F(z).
\]
It is entire, real under complex conjugation, and its zeros correspond to the nontrivial zeros of the Hasse--Weil function. Euler convergence and the functional equation already exclude $|\operatorname{Re}z|\ge1$.

For completeness two independent rational points can be certified in small finite fields. At the good primes $5,7,11$ the group orders are $9,13,16$ by direct counting. Prime-to-$p$ torsion injects at good $p$, so the first two orders exclude all rational torsion. Indeed a possible $5$-primary part is excluded at 7 and a $7$-primary part at 5. For all other primes use the coprimality. Set $P=(0,0)$ and $Q=(1,0)$. The line $y=0$ gives $P+Q=(-2,-1)$. The doubles in $E(\mathbb F_{11})$ are
\[
             O,\quad(3,5),\quad(6,4),\quad(6,6).
\]
To reproduce the list, for each affine point $(x,y)$ with $2y+1\ne0$ compute $m=(3x^2+2x-2)/(2y+1)$ and put $x'=m^2-1-2x$ and $y'=m(x-x')-y-1$. The remaining points double to $O$. None of $P,Q,P+Q$ reduces to a double. Hence the reductions of $P,Q$ are independent modulo $2E(\mathbb F_{11})$. In any integer relation $aP+bQ=O$, torsion-freeness allows us to divide out a common power of two. Reducing leaves a nonzero vector in $\mathbb F_2^2$, a contradiction. Thus $E(\mathbb Q)$ has rank at least two. The rank-zero theorem of Gross--Zagier and Kolyvagin gives $L(E,1)=0$ because nonvanishing would force $E(\mathbb Q)$ finite. Evenness of $F$ then gives a zero of order at least two. No implication of the general Birch--Swinnerton-Dyer conjecture in rank two is assumed.

\subsection*{The coefficients and a signed Mellin kernel}
Let $a_p=p+1-\#E(\mathbb F_p)$, counting also the point at infinity on the special fibre at $p=389$. There $a_{389}=1$. At odd $p$ the count is one quadratic-symbol sum
\[
   a_p=-\sum_{x\bmod p}\left(\frac{1+4(x^3+x^2-2x)}p\right).
\]
The computer builds the exact integer $a_n$ using multiplicativity and $a_{p^j}=a_pa_{p^{j-1}}-pa_{p^{j-2}}$ at good primes. At 389 it uses the degree-one Euler recurrence. Hasse's bound on the two good Euler roots, with multiplicativity, gives $|a_n|\le d(n)\sqrt n\le 2n$. Define
\[
  K(u)=e^u f(ie^u/\sqrt N),\qquad
  F(z)=\int_{-\infty}^{\infty}K(u)e^{zu}\,du.
\]
The Mellin formula holds by termwise integration first on the right and then by analytic continuation. The Fricke transformation with sign $+1$ gives $K(-u)=K(u)$ throughout $|\operatorname{Im}u|<\pi/2$. The identity $F(0)=0$ rules out an everywhere strictly positive kernel.

Take $u=x+iy$ with $x\ge0$, $|y|\le1$. Setting $v=B e^x\cos y$ gives $v>e^x/6$. Hence
\[
 |f(ie^u/\sqrt N)|\le\frac{2e^{-v}}{(1-e^{-v})^2}
                <90\exp(-e^x/6).
\]
For $|\operatorname{Re}z|\le1$, $|\operatorname{Im}z|\le T$ each horizontal boundary integrand is bounded by $90e^T\exp(2x-e^x/6)$ and the reflection supplies the other half-line. The sum of the two boundary $L^1$ norms is at most $4\cdot90\cdot36e^T<20000e^T$, since substitution $r=e^x/6$ bounds the half-line integral by $36\Gamma(2)$. The trapezoid rule on the full line at spacing $h$ therefore has error at most $20000 e^T/(e^{2\pi/h}-1)$. Indeed each nonzero Fourier alias at frequency $2\pi m/h$ gains $e^{-2\pi|m|/h}$ by moving the contour to the appropriate horizontal line. The sum for $m\ne0$ is a geometric series with the two boundary integrals as constants. Exponential decay on the vertical ends justifies that displacement.

On the real line $B>3/10$, so the analogous denominator constant is $2/(1-e^{-3/10})^2<30$. Truncating at $|u|\le W$ costs at most
\[
       60\int_W^\infty e^{2x-3e^x/10}\,dx
         <3000\exp(-3e^W/20).
\]
Here the majorant is decreasing beyond $W$ and splitting the exponential in half bounds the residual integral by $(20/3)^2\Gamma(2)$. Write $\epsilon$ for the sum of the displayed discretization and truncation bounds. It is uniform on the whole rectangle. At a positive grid point $u=jh$ compute $q=\exp(-Be^u)$ and enclose $f$ by its exact series through $n=M$ with radius
\[
                2q^{M+1}\frac{(M+1)-Mq}{(1-q)^2}.
\]
Writing $w_j=he^{jh}f(ie^{jh}/\sqrt N)$, the enclosed finite approximation to $F$ is
\[
        A(z)=w_0+\sum_{j=1}^{W/h}w_j(e^{jhz}+e^{-jhz}),
                 \qquad |F(z)-A(z)|\le\epsilon.
\]
Interval uncertainty in the coefficients stays inside the interval evaluations of $A$. Every numerical operation after the integer point counts uses directed real or complex ball arithmetic in FLINT/Arb.

\subsection*{A full contour and the signed count}
Parametrize the upper right half-edge by $z=1/2+x+iT$, $-1/2\le x\le1/2$. Expand the exponentials of $A$ in the real variable $x$ through degree $d$. Putting $S=|w_0|+2\sum_{j>0}|w_j|$, the discarded Taylor terms together with the previous approximation have absolute value at most
\[
             \epsilon+Se^W\frac{(W/2)^{d+1}}{(d+1)!}.
\]
Indeed the factors at $x=0$ cost at most $e^{W/2}$ and the exponential Taylor bound costs a second $e^{W/2}$. Signed weights require the absolute-value mass $S$. On an interval in $x$, substitute its midpoint-radius ball in the polynomial by Horner's rule and add this remainder to both complex coordinates. Bisect until the images of all subintervals lie in some open coordinate half-plane. This certifies nonvanishing of the edge. For a subinterval with right endpoint value $b$ and left endpoint value $a$, the variation oriented to the left is $\arg(a/b)$ on the principal branch, since every path within that half-plane admits arguments in a common interval of length $\pi$. The code checks that the quotient argument stays strictly within $(-\pi,\pi)$. Their sum is $\theta_{\rm top}$.

At $s=2+it$, each good Euler root has modulus $\sqrt p$, so the absolutely convergent Euler logarithm satisfies
\[
 |\log L(E,2+it)|\le 2\log\zeta(3/2)\le2\log 3<\pi.
\]
The bad Euler factor is smaller than this majorant as well. Thus the continuous Euler argument from $t=0$ equals the principal argument at $T$. Compute the completed argument $\theta_{\rm right}$ for the edge from $s=2$ to $s=2+iT$ as the sum of this Euler argument and the imaginary part of $\log\Gamma(2+iT)-(2+iT)\log B$, using the branch real at $2$. The Euler factor at $T$ is evaluated by dividing the enclosed completed value by this explicit gamma factor. The initial value at $s=2$ is strictly positive by its convergent real Euler product. Symmetry in the real and imaginary axes relates the other three quarters of the boundary $|\operatorname{Re}z|\le1$, $|\operatorname{Im}z|\le T$ to this quarter. Reflection of a path reverses orientation and conjugation reverses angle, so their oriented variations are equal. The total number of zeros, with multiplicity, is consequently
\begin{equation}\label{r2:full-count}
              C(T)=\frac{2}{\pi}(\theta_{\rm right}+\theta_{\rm top}).
\end{equation}
No condition on values of $F$ on the real axis inside the contour has been used. The entire completion has no poles.

Two independent devices account for these multiplicities. First $F(it)$ is real. Sampling the finite exponential sum $A(it)$ by a discrete Fourier transform at $t_k=2\pi k/(hJ)$, $0<t_k<T$, and including the already enclosed top value gives intervals of opposite certified signs. A candidate sign is used only if its enclosure, enlarged by $\epsilon$, avoids zero. The first sampling index is positive, so the known central zero supplies no sign by assumption. These successive intervals have disjoint interiors. A change contributes at least one positive zero and, by conjugation, a second negative zero. Second, Cauchy's inequality on $|z|=1/2$ gives $|F''(0)-A''(0)|\le8\epsilon$. Hence $2\sum_{j>0} w_j(jh)^2$ with that extra radius encloses $F''(0)$. The guaranteed central zero and evenness give $L''(E,1)/2=BF''(0)/2$. Its stated positive interval proves that the central multiplicity equals two. It is obtained directly from the equation, without division by a small value of $L$.

The arithmetic in the certificate isolates a unique integer in \eqref{r2:full-count}. Its parameter choices and numerical bounds are summarized below. $D=h^{-1}$, $M$ is the highest Dirichlet coefficient, $d$ the horizontal Taylor degree, and $J$ the number of Fourier samples on one period. The $b$ entries count sign-changing positive brackets, so $C=2+2b$ accounts for every zero in the rectangle.
\[
\begin{array}{r|rrrrrrr|rr}
 T&D&W&M&d&\hbox{arcs}&J&\hbox{bits}&b&C\\\hline
 32&40&8&948&97&6&2048&320&36&74\\
 256&152&10&4084&241&10&16384&992&463&928\\
 1024&536&11&14836&626&343&131072&3296&2306&4614\\
 2048&1048&12&29172&1081&230&262144&6368&5064&10130
\end{array}
\]
An independent evaluation of the first 36 positive brackets, using an incomplete-gamma expansion at every endpoint, also gives the rational enclosure
\[
             2.876099071260 < \gamma_1 < 2.876099071261
\]
for the least positive zero ordinate. Every bracket in the table contains exactly one simple zero. No zero elsewhere on or off the center line is left unaccounted for. Nonvanishing of the upper and lower edges makes the height bound inclusive. This proves Theorem~\ref{r2:window}.

\subsection*{A finite-window criterion for arbitrary rational elliptic curves}
The treatment of the real axis extends to every rational elliptic curve. Let $E'/\mathbb Q$ have conductor $N'$, root number $w$, and completed function $F_{E'}(z)$ centered at zero. Suppose a rigorous argument gives a lower bound $m\ge0$ for the central multiplicity and that $m$ has parity $w=(-1)^m$. For $T>0$ with no zero on the two horizontal boundaries, let $C$ be the full contour count and $b$ the number of disjoint positive-ordinate sign-change intervals for $i^{-m}F_{E'}(it)$ on $(0,T)$. Both $C$ and the signs can be obtained to arbitrary precision by the construction below. Then
\[
                       C\ge m+2b.
\]
Equality certifies the exact central multiplicity, simplicity in every bracket, and the complete zero assertion throughout the window. A positive certified central derivative of the requisite order is convenient, but the equality of counts itself suffices. The inequality counts an actual zero in each sign-changing interval, its conjugate at negative ordinate, and the zeros already forced at the origin. The phase $i^{-m}$ is necessary for odd parity, where the completed function on the symmetry line is purely imaginary.

To see why the analytic evaluation can be made uniform in this generality, set $B'=2\pi/\sqrt{N'}$, $b_0=B'\cos(1)$ and $C_0=2/(1-e^{-b_0})^2$. The coefficient bound $|a_n|\le2n$ holds at all good primes, and at bad primes the actual Euler roots preserve it. The same signed kernel has parity $w$ under $u\mapsto-u$. On the complex strip used above its newform factor is bounded by $C_0\exp(-b_0e^x)$ for $x\ge0$. The sum of the horizontal boundary integrals is at most $4C_0e^T/b_0^2$. On the real line set $C_1=2/(1-e^{-B'})^2$ instead. For $B'e^W\ge2$, truncation costs at most $8C_1(B')^{-2}\exp(-B'e^W/2)$. The geometric Fourier-coefficient remainder works at every $q=\exp(-B'e^x)$. Thus taking $h$ small, $W$ large and sufficiently many coefficients and bits resolves any strict sign and any compact nonvanishing contour segment. Rotation by $w$ changes no increment of argument in the four-quarter contour proof. A compact nonvanishing horizontal edge admits a finite cover by subarcs in open coordinate half-planes. Finally, if every positive zero in the window is simple and has real part one, a sufficiently fine mesh yields one sign change per zero. This gives a terminating certificate in such a window when the exact central multiplicity is known as a lower bound. Its proof supplies a criterion for further finite windows without implying anything about arbitrarily large zeros of a fixed curve.

\paragraph{Sources and reproduction.}
For modularity of semistable elliptic curves use Wiles and Taylor--Wiles. For the analytic rank-zero implication use H.~Darmon, \emph{Rational Points on Modular Elliptic Curves}, CBMS 101, Theorem 1.14, which cites Gross--Zagier and Kolyvagin. Tate's algorithm or the multiplicative-reduction criterion in J.~H.~Silverman, \emph{The Arithmetic of Elliptic Curves}, gives the stated reduction and conductor. The curve is also LMFDB 389.a1 and its newform is 389.2.a.a. For each row, run \texttt{certify.py} with the corresponding integer \texttt{--height}, using Python and \texttt{python-flint} 0.9.0. The JSON files include every bracket as integer grid endpoints. Reproduction commands, exact output hashes and an independent incomplete-gamma calculation for the smallest window accompany the chapter.

\clearpage
\section{Four zeros in a disk about the center of a rank-four elliptic curve}\label{r4:chapter}

The curve
\[
 E:\quad y^2+xy=x^3-x^2-79x+289
\]
has Mordell--Weil rank four. Its rational invariant $j=3801^3/468892$ is nonintegral, so the curve has no complex multiplication. Its analytic rank remains either two or four. Kezuka and Zagier identify it as the exceptional curve in the conductor range through $500\,000$ for the unconditional comparison of the two ranks. The LMFDB reliability statement makes the same distinction for 234446.a1. We give a finite theorem that preserves this unresolved alternative. All variables in this section have arithmetic center $s=1$.

\begin{theorem}\label{r4:cluster}
The completed function of $E$ has exactly four zeros, counted with multiplicity, in
\[
                          |s-1|<2^{-4108}.
\]
The boundary is zero-free. In the full horizontal slab $|\Im s|\le2048$, every other zero is simple and lies on $\Re s=1$. There are $7149$ such zeros of positive ordinate and their $7149$ conjugates of negative ordinate.

At least two of the four zeros are at $1$. The other two are either both at $1$, or form a simple pair $1\pm it$ with $t>0$, or form a simple pair $1\pm x$ with $x>0$. The last alternative is equivalent to $L''(E,1)<0$. In particular all zeros in this slab satisfy elliptic GRH if and only if $L''(E,1)\ge0$.
\end{theorem}

For the associated heat deformation, Corollary~\ref{r4:upperphase} identifies the last possible nonreal pair through height $65257/32$ by the unique positive solution, if one exists, of an explicit scalar equation.

The arithmetic database supplies the four rational points used below. Their linear independence has a short finite-field certificate. Only one point of infinite order is needed for the analytic argument.

\subsection{The exact central zero and the first two possible coefficients}
The usual integral invariants are $c_4=3801$, $c_6=-232605$ and $\Delta=468892=2^2\cdot117223$. The integer $117223$ is prime and both bad primes give multiplicative reduction because $c_4$ is a unit there. This integral model is minimal at both primes and its conductor is $N=234446$. At $2$ counting the affine points on the reduced equation gives $a_2=-1$. At $117223$ the reduction splits, since $32289^2\equiv-c_6\pmod{117223}$. The product of the archimedean sign and the two local signs $-a_p$ is $+1$.

The point $P=(0,17)$ reduces to a nonidentity point of a group of order seven modulo $3$. Modulo $5$ the whole group has order ten. Both assertions follow by counting solutions to the displayed cubic. Torsion divisible in order by seven injects on its seven-primary part modulo the good prime $5$. Consequently $P$ has infinite order. The analytic-rank-zero implication of Gross--Zagier and Kolyvagin gives $L(E,1)=0$. Modularity and the even functional equation make the order at least two.

For a direct check of all four independent directions, order the points on this model as
\[
               (6,-1),\quad(4,3),\quad(5,-2),\quad(8,7).
\]
For each prime in the table write $m=\#E(\mathbb F_p)$. Its two-primary subgroup has order two. Multiplication by $m/2$ therefore gives a homomorphism to that subgroup. The row records which reductions of the four points map to its nonidentity element.
\begin{center}\small
\begin{tabular}{rrcc}
\toprule
$p$ & $m$ & nonidentity element & four bits\\ \midrule
5 & 10 & $(4,3)$ & $1100$\\
11 & 18 & $(8,7)$ & $0111$\\
23 & 30 & $(19,2)$ & $1000$\\
73 & 82 & $(13,30)$ & $1001$\\
\bottomrule
\end{tabular}
\end{center}
Direct counting and the chord-and-tangent law reproduce the table. These four row vectors have rank four over $\mathbb F_2$. Every integral relation between the points would thus have even coefficients. The odd group order at $3$ excludes rational two-torsion, so one can divide any such relation by two, contradicting primitivity. This proves rank at least four. The upper bound also follows from the nonzero $5$-adic quartic and the rank inequality in Theorem~\ref{r4:sha}, using these four points for the lower bound. An independent PARI two-descent returns lower and upper ranks $[4,4]$.

Put $B=2\pi/\sqrt N$ and
\[
 F(z)=B^{-(1+z)}\Gamma(1+z)L(E,1+z)
          =b_2z^2+b_4z^4+b_6z^6+\cdots .
\]
The even entire function $F$ is real on the real axis. The comparison obtained by outward arithmetic was
\begin{equation}\label{r4:numbers}
      |b_2|<2^{-8219},\qquad
      689.2327823095<b_4<689.2327823096 .
\end{equation}
In the original Hasse--Weil function the same bounds give
\[
 |L''(E,1)|<2^{-8224},\qquad
 8.94384739589<\frac{L^{(4)}(E,1)}{24}<8.94384739591.
\]
Indeed $L''(E,1)=2Bb_2$ and $2B<1/32$. The fourth coefficient equals
\[
 B\biggl(b_4+b_2\left(\frac{(\gamma+\log B)^2}{2}
                              -\frac{\pi^2}{12}\right)\biggr),
\]
by the expansion of $e^{z\log B}/\Gamma(1+z)$. This formula includes the possible contribution from a nonzero second derivative.

\subsection{An exponentially accurate second-coefficient calculation}\label{r4:quadrature}
Let $f$ be the normalized weight-two newform of $E$ and set
\[
  K(u)=e^u f(ie^u/\sqrt N)
      =e^u\sum_{n\ge1}a_n e^{-Bne^u}.
\]
The Fricke involution gives $K(-u)=K(u)$, and the Mellin transform gives $F(z)=\int_{\R}K(u)e^{zu}\,du$. Each prime coefficient is obtained by counting solutions of this equation modulo $p$. Multiplicativity and the good-prime recursion $a_{p^k}=a_pa_{p^{k-1}}-pa_{p^{k-2}}$ complete the integer list, with $a_{p^k}=a_p^k$ at multiplicative primes. In particular $|a_n|\le d(n)\sqrt n\le2n$.

Here are the error estimates for the central moments. Write $h=1/D$, let $W$ be a positive integer, put $b=B\cos 1$, and set
\[
 C_0=\frac{2}{(1-e^{-b})^2},\qquad C_1=\frac{2}{(1-e^{-B})^2}.
\]
The finite trapezoidal approximation is the even exponential sum
\[
 A(z)=w_0+\sum_{j=1}^{DW}w_j(e^{jhz}+e^{-jhz}),\qquad
 w_j=hK(jh).
\]
Here $A$ uses the exact weights. Their numerical intervals are handled separately. Throughout $|\Re z|,|\Im z|\le1$ one has
\begin{equation}\label{r4:trapezoid}
 |F(z)-A(z)|\le\epsilon:=
 \frac{4C_0 e}{b^2(e^{2\pi D}-1)}
           +\frac{8C_1}{B^2}\exp(-Be^W/2)
\end{equation}
when $Be^W\ge2$. For the first term, on the lines $u=x\pm i$ the absolute value of the integrand for $x\ge0$ is at most $eC_0 e^{2x-b e^x}$. Reflection gives the same bound for the other half. The four half-line integrals have total at most $4eC_0/b^2$. Shifting in the Fourier transform and summing the two geometric alias tails gives the denominator $e^{2\pi D}-1$, as in Section~\ref{r2:chapter}. For the second term, $e^{2x-Be^x}$ decreases beyond $W$, so bound each omitted real-axis mesh tail by an integral. With $v=Be^x$, its size is at most $C_1 B^{-2}(Be^W+1)e^{-Be^W}$ per side. The displayed weaker second term follows.

For $y=e^{jh}$, $q=e^{-By}$, the truncation after coefficient $m$ in $w_j$ has radius at most
\[
                  2hy\,q^{m+1}\frac{m+1-mq}{(1-q)^2}.
\]
From a global maximum $M$ the implementation halves the number of terms successively when $y$ doubles, rounding each length upward. The displayed tail is applied to the actual $m$ at every point. Arb polynomial arithmetic encloses the whole finite sum. Finally Cauchy's estimate on $|z|=1/2$ bounds the error in the coefficient of $z^k$ in $F-A$ by $2^k\epsilon$. Thus the numerical sums for even $k\ge2$ are
\[
     \frac{2}{k!}\sum_{j=1}^{DW}w_j(jh)^k
                    \quad\text{with an additional radius }2^k\epsilon.
\]
This is a direct evaluation of the Taylor coefficient. No difference of two nearly equal values of $L$ is used for the primary enclosure.

\begin{center}\small
\begin{tabular}{rrrrrr}
\toprule
Working bits & $D$ & $W$ & $M$ & $\epsilon<$ & $|b_2|<$\\ \midrule
512 & 63 & 12 & 30148 & $2^{-538}$ & $2^{-494}$\\
2048 & 232 & 13 & 112194 & $2^{-2069}$ & $2^{-2028}$\\
4096 & 458 & 14 & 221589 & $2^{-4118}$ & $2^{-4075}$\\
8240 & 915 & 14 & 442943 & $2^{-8261}$ & $2^{-8219}$\\
\bottomrule
\end{tabular}
\end{center}
All the listed runs give the same ten-place enclosure for $b_4$ in~\eqref{r4:numbers}.

We made a second calculation using four evaluations of real upper incomplete gamma functions per Fourier coefficient. For $0<x<1$ define
\[
 I(x)=\sum_{n\ge1}a_n\left\{
 \frac{\Gamma(1+x,Bn)}{(Bn)^{1+x}}+
 \frac{\Gamma(1-x,Bn)}{(Bn)^{1-x}}\right\}.
\]
Splitting the Mellin integral at zero, reflecting, and integrating each term gives $I(x)=F(x)$. The series is absolutely convergent. Since $y^x,y^{-x}\le y$ for $y\ge1$, the omitted terms past $m$ have total absolute value at most
\[
  \frac{4}{B}\left(1+\frac1{B(m+1)}\right)
          \frac{e^{-B(m+1)}}{1-e^{-B}} .
\]
Take $h_0=2^{-66}$ and $m=29436$. Compute $I(h_0)$ and $I(2h_0)$ at 512-bit precision with Fourier coefficients obtained independently from PARI/GP. Then
\[
  \frac{16I(h_0)-I(2h_0)}{12h_0^2},\qquad
  \frac{I(2h_0)-4I(h_0)}{12h_0^4}
\]
approximate $b_2,b_4$ respectively. The unit-disk bound $2^{30}$ proved below gives respective errors at most $7\cdot 2^{30}h_0^4$ and $6\cdot2^{30}h_0^2$ after the displayed integration tails and ball-rounding errors are included separately. Indeed the Taylor remainder above degree four at either argument $y$ is bounded by $2^{30}y^6/(1-y^2)$. The resulting independent bounds give $|b_2|<2^{-230}$ and the same narrow interval for $b_4$. This calculation uses positive real gamma arguments and no trapezoidal mesh. The two coefficient-construction algorithms were compared entry by entry through all $442943$ coefficients of the longest central calculation and all $639771$ of the largest contour calculation.

\subsection{From analytic inequalities to exactly four zeros}
On $|z|\le1$ and $u\ge0$ the integrand is bounded by $C_1 e^{2u-Be^u}$. Integrating and reflecting as above yields the uniform majorant
\[
                  |F(z)|\le 2C_1/B^2<2^{30}.
\]
Cauchy's estimate now gives $|b_{2j}|<2^{30}$ for every $j$. In particular for $r=|z|<1$,
\[
         |F(z)-b_2z^2-b_4z^4|
                 \le 2^{30}\frac{r^6}{1-r^2}.
\]
At $r=2^{-4108}$, division by $r^4$ bounds the combined quadratic and higher terms by
\[
                 2^{-8219}r^{-2}+2^{30}\frac{r^2}{1-r^2}
                      <1<b_4.
\]
Rouché's theorem compares $F$ with $b_4 z^4$ and certifies that no zero belongs to the circumference.

For the position of the two remaining zeros, write
\[
    F(z)=z^2 H(z^2),\qquad H(v)=b_2+b_4v+b_6v^2+\cdots.
\]
The four-zero count puts precisely one zero of $H$ inside $|v|<r^2$, counted with multiplicity. Complex conjugation makes that zero real, since a nonreal zero of $H$ would bring a second distinct zero. The derivative on the real interval $|v|\le r^2$ differs from $b_4$ by at most
\[
                    2^{30}\bigl((1-r^2)^{-2}-1\bigr)<1.
\]
It is strictly positive there. Hence $b_2=0$ puts both remaining zeros at the center, $b_2>0$ puts them on the imaginary $z$-axis, and $b_2<0$ puts them on the real $z$-axis. This proves the asserted classification, including simplicity in the last two cases. If $z_0\ne0$ is a member of the residual pair, the same Taylor estimate gives a relative location formula
\[
       \left|\frac{-b_2}{b_4 z_0^2}-1\right|
                 \le \frac{2^{30}r^2}{b_4(1-r^2)} .
\]
Thus a future nonzero determination of the quadratic coefficient also determines the pair's separation to the displayed relative precision. Since zeros of a nonzero entire function are isolated, the analytic rank is four if and only if every sufficiently small dyadic disk about zero contains four zeros. Rank two leaves a positive distance from zero to the residual pair, beyond which smaller disks count only the central double zero. The finite bounds above give an explicit upper bound on that distance in the rank-two case.

The same deduction applies to any self-dual automorphic completion with real coefficients. Suppose its centered function has the parity of $k$ and an exact zero of order at least $k$, and write
\[
        G(z)=z^k(a+bz^2+O(z^4)),\qquad b\ne0.
\]
If $|G|\le U$ on $|z|=R$, then the sufficient condition for exactly $k+2$ zeros in the open disk of radius $r<R$ is
\[
       |a|r^{-2}
       +\frac{U r^2}{R^{k+4}(1-r^2/R^2)}<|b|.
\]
Indeed Cauchy's coefficient estimates give the second term as the entire higher-order tail divided by $r^{k+2}$. Rouché supplies the count. There is one zero, including multiplicity, of $G(z)/z^k$ as a function of $v=z^2$, hence its $v$-coordinate is real. At $v=\pm r^2$ the sign of this real function is the sign of $bv$, by the same strict inequality. Its value $a$ at zero determines on which half interval the residual zero falls. This gives an explicit criterion for local critical-line verification with one unresolved pair for either parity of the functional equation.

\subsection{Exhausting the punctured horizontal slab}\label{r4:slab}
Use the full symmetric counting rectangle from Sections~\ref{r2:chapter} and~\ref{cens:chapter}. The argument-principle count is independent of any central rank. The Euler product controls the vertical edge $\Re(s)=2$. The heights below use the unshifted real Mellin integral. The horizontal edge is bounded by complex interval Taylor evaluation of the same signed Mellin sum, with interval subdivision into half-planes. On the critical line, strict opposite signs of $F(it)$ at two mesh endpoints certify a zero in the open interval between them.

\begin{center}\small
\begin{tabular}{rrrrrrrr}
\toprule
 $T$ & bits & $D$ & $W$ & $M$ & mesh $J$ & total & positive brackets\\ \midrule
 64  & 448  & 56  & 11 & 28205 & 4096  & 306  & 151\\
 256 & 1024 & 152 & 12 & 87389 & 16384 & 1448 & 722\\
 1024 & 3328 & 536 & 14 & 324124 & 262144 & 6700 & 3348\\
 2048 & 6400 & 1048 & 14 & 639771 & 524288 & 14302 & 7149\\
\bottomrule
\end{tabular}
\end{center}

The bracket numbers are indices in units $2\pi D/J$. The first positive brackets are $(20,21)$, $(29,30)$, $(134,135)$ and $(138,139)$ in the four rows. All lie above $1.69$, so none intersects the cluster. Distinct intervals certify distinct zeros, as common endpoints always have strict sign. The equality between each total and four plus twice the number of sign changes leaves multiplicity one for every point in these brackets and excludes any other zero, including on the horizontal edges, on the real axis outside the cluster, and anywhere off the critical line outside the cluster. The Euler product and functional equation exclude zeros beyond the vertical edges in the same slab. On the real axis the completed function is therefore positive whenever $|s-1|$ is at least the certified cluster radius. Its sign is positive in the right Euler-product region and the functional equation gives the left side. The first bracket and the next lower endpoint in the $T=64$ row show that there is exactly one positive zero outside the cluster in $[1,2]$. Bisection with independently recomputed signs of the real Mellin sum gives it a dyadic interval of length $2^{-200}$. Its ordinate starts
\[
                1.73353224715782688313137173222659311839091305\ldots .
\]
The exact consecutive numerators are in the companion record.

\subsection{A Heegner base change and the remaining sign}\label{r4:heegner}

Take $K=\Q(\sqrt{-7})$. Both primes of $N$ split in $K$, since $-7\equiv1\pmod8$ and $117223\equiv1\pmod7$. Let $E_-$ denote the quadratic twist by its primitive character $\chi_{-7}$. Its conductor is $49N=11\,487\,854$ and its root number is $-1$, by the coprime quadratic twist formula $w(E_-)=w(E)\chi_{-7}(-N)$. At all integers its coefficients are $a_n\chi_{-7}(n)$. A minimal model has invariants $[1,-1,0,-3880,-91372]$.

\begin{theorem}\label{r4:heegner-thm}
The curve $E/K$ has Mordell--Weil rank five and analytic rank three or five. Its completed function has exactly five zeros in the disk of Theorem~\ref{r4:cluster}, counted with multiplicity. The residual pair relative to the triple central zero is the pair in that theorem. In $|\Im s|\le32$ every other zero is simple and on the critical line. There are $157$ positive and $157$ negative such zeros. A necessary and sufficient condition for this pair to satisfy elliptic GRH is
\[
                     L^{(3)}(E/K,1)\ge0.
\]
\end{theorem}
\begin{proof}
For the twist use the same Mellin quadrature with $B_-=B/7$, sign $-1$, $192$ working bits and $(D,W,M)=(24,12,60000)$. Its odd Mellin kernel gives the first derivative of the completed function as $2\sum_{j>0} w_j jh$ with additional radius $2\epsilon$ by Cauchy's formula. Multiplication by $B_-$ gives
\[
             34.1720344609<L'(E_-,1)<34.1720344610 .
\]
The unit-disk bound for its completion $F_-$ is $2C_{1,-}/B_-^2<2^{42}$ by the same integral as above. On $|z|=2^{-24}$ the part beyond the linear coefficient, divided by $|z|$, is bounded by $2^{42}2^{-48}/(1-2^{-48})<1<F_-'(0)$. Hence its simple central zero is its only zero in this larger disk.

The quadratic base-change identity gives $L(E/K,s)=L(E,s)L(E_-,s)$, and the same product for the completions up to a constant. The twist has analytic rank one, hence arithmetic rank one by Gross--Zagier and Kolyvagin. The two eigenspaces for $\operatorname{Gal}(K/\Q)$ on $E(K)\otimes\Q$ identify with the rational spaces of $E$ and $E_-$. This gives arithmetic rank $4+1=5$. The zero counts add. As $L(E,1)=L'(E,1)=L(E_-,1)=0$, the product rule yields the exact identity
\begin{equation}\label{r4:third}
                L^{(3)}(E/K,1)=3L''(E,1)L'(E_-,1).
\end{equation}
The certified positive last factor and Theorem~\ref{r4:cluster} give the sign criterion.

For the finite slab count the two factors separately at $T=32$ with $384$ working bits. For $E$ the rectangle contains $140=4+2\cdot68$ zeros. For $E_-$ it contains $179=1+2\cdot89$. In both cases the horizontal profiles use $D=40$. The widths are $11,13$ and the Fourier cutoffs are $18341,128386$, respectively. Bisection of each strict central-line sign change isolates the $68+89$ positive roots in pairwise disjoint open dyadic intervals of length $2^{-18}$. Hence the completed product has $319$ zeros in the slab and all $314$ outside the cluster are simple.
\end{proof}

Equation~\eqref{r4:third} singles out a geometric target. At this Heegner discriminant the classical Gross--Zagier formula expresses $L'(E/K,1)$ as a positive constant times a Néron--Tate height. The product rule already gives $L'(E/K,1)=0$ for either sign of $L''(E,1)$, so the associated Heegner point is torsion. A nonnegative third-derivative pairing for this specific complex base-change function would settle the residual GRH question in the entire certified slab for $E/\Q$. A pairing allowed to vanish would suffice. Yun and Zhang established higher derivative identities with Heegner--Drinfeld cycles over global function fields, including ramified settings. Realizing such a third derivative over this imaginary quadratic number field, together with a proof that the resulting pairing is nonnegative, is the missing geometric assertion. See Gross and Zagier, \href{https://doi.org/10.1007/BF01388809}{\emph{Heegner points and derivatives of $L$-series}}, and Yun and Zhang, \href{https://annals.math.princeton.edu/2019/189-2/p02}{\emph{Shtukas and the Taylor expansion of $L$-functions} (II)}. The distinction between first and higher derivatives over number fields is discussed by C.~Li, \href{https://www.math.columbia.edu/~chaoli/ASWSurvey.pdf}{\emph{From sum of two squares to arithmetic Siegel--Weil formulas}}.

For this curve the $5$-adic analogue of exact fourth-order vanishing is known. We reproduce the decisive modular-symbol computation. Write $\omega_1$ for the least positive real period and
\[
    S_m(a)=\omega_1^{-1}\Re\left(2\pi i\int_{i\infty}^{a/m}f(z)\,dz\right).
\]
These are integers here by the period-lattice calculation in C.~Wuthrich, \href{https://doi.org/10.1090/mcom/3274}{\emph{Numerical modular symbols for elliptic curves}}, §9. In the ordinary $p$-adic $L$-function the symbols are divided by the full real volume, which is $2\omega_1$ because the real locus has two components.

For $m=5,25$ choose $a'$ with $Naa'\equiv-1\pmod m$ and put $q=e^{-B/m}$. Splitting the cusp path at imaginary height $1/(m\sqrt N)$ and applying the Fricke transformation as in Wuthrich, §4.1, gives
\[
 \omega_1 S_m(a)=\sum_{n\ge1}\frac{a_nq^n}{n}
     \left(\cos(2\pi na/m)+\cos(2\pi na'/m)\right).
\]
The plus sign is the root number of $E$. The tail after $40000$ coefficients is bounded by $4q^{40001}/(1-q)$. For the period, if $e_1>e_2>e_3$ are the three roots of $4x^3-3x^2-316x+1156$, the real elliptic integral gives
\[
 \omega_1=\frac{\pi}{\operatorname{AGM}(\sqrt{e_1-e_3},\sqrt{e_1-e_2})}
             \in(1.4863,1.4864).
\]
The root intervals and the series isolate unique integers.

Let $\alpha$ be the $5$-adic unit root of $X^2-a_5X+5$. As $a_5=-4$, its residue is $1$. The ordinary distribution at the residue $a+25\mathbb Z_5$, reduced modulo $5$, is consequently $(S_{25}(a)-S_5(a))/2$. In the cyclotomic coordinate with generator $1+5$, put
\[
 k(a)=\frac{\log\langle a\rangle}{\log(1+5)}
        \equiv 4\frac{a^4-1}{5}\pmod5 .
\]
This congruence follows by raising $a$ to the fourth power modulo $25$. The fourth Mahler coefficient $\binom{k(a)}4$ is $1$ modulo $5$ when $k(a)\equiv4$, and zero otherwise. The four contributing classes give
\[
\begin{array}{c|rrrr}
a & 3&4&21&22\\ \hline
 S_{25}(a)&8&-7&-7&8\\
 S_5(a)&-8&8&8&-8
\end{array}
\]
and hence $[T^4]\mathcal L_5(E,T)\equiv(16-15-15+16)/2=1\pmod5$. This recovers Wuthrich's coefficient with the full real period. The Schneider--Kato inequality $\operatorname{ord}_{T=0}\mathcal L_p(E,T)\ge\operatorname{rank}E(\Q)$ holds for the odd good ordinary primes used here, as in the chain preceding Proposition~10.1 of W.~Stein and C.~Wuthrich, \href{https://wstein.org/papers/mcom2649-iwasawa-alg.pdf}{\emph{Algorithms for the arithmetic of elliptic curves using Iwasawa theory}}. Thus the $5$-adic order is \emph{exactly four}. Wuthrich, §9, also records triviality of the $5$-primary Tate--Shafarevich group for this curve.

This makes the missing comparison explicit. The $5$-adic quadratic coefficient vanishes. The ordinary measure interpolates normalized complex central values for finite-order twists, while $b_2$ is a logarithmic second moment for the single untwisted complex function. If one could prove for this curve the single implication $[T^2]\mathcal L_5(E,T)=0\Rightarrow L''(E,1)=0$, Theorem~\ref{r4:cluster} would give analytic rank four. Classical interpolation supplies central values of the twists and contains no such derivative identity. The proposed rank--defect formula of Kezuka and Zagier likewise implies $\operatorname{ord}_{s=1}L(E,s)\ge4$ here and is explicitly conjectural.

A simple positivity attempt with modular integrals can be tested on this curve. Put
\[
 A(y)=\int_{\log y}^{\infty}K(u)\,du
     =\frac1B\sum_{n\ge1}\frac{a_n}{n}e^{-Bny},\qquad y\ge1 .
\]
Integration by parts gives $b_2=2\int_0^{\infty} uA(e^u)\,du$. The exact central value gives $A(1)=0$. Evaluating the first $1200$ terms gives both signs for the primitive
\[
                  -29<A(2)<-28,\qquad 8<A(32)<9 .
\]
For $q=e^{-By}$ the omitted terms have absolute value at most $2q^{1201}/(B(1-q))$. Consequently pointwise nonnegativity of this primitive cannot supply the desired integral inequality. The ordinary weight-two cusp-to-cusp modular symbol controls the central value. The logarithmic second moment above retains the cancellation measured by~\eqref{r4:numbers}.

\subsection{Nine ordinary primes and the Tate--Shafarevich group}\label{r4:ordinary}

\begin{theorem}\label{r4:sha}
For every prime $5\le p\le31$ the ordinary cyclotomic $p$-adic $L$-function of $E$ has order exactly four at the trivial character. The cyclotomic $p$-adic height pairing on $E(\Q)\otimes\Q_p$ is nondegenerate and the $p$-primary part of $\operatorname{Sha}(E/\Q)$ is trivial.
\end{theorem}

Here are the numerical and local inputs. We use the generator $1+p$ and the full real period throughout. The column $A_4$ gives $[T^4]\mathcal L_p$ modulo $p$, except at $7$ where the modulus is $49$. The column $v_p(R)$ is the valuation of the ordinary cyclotomic height regulator before the division by $\log_p(1+p)^4$. The prime $\ell$ in each row supplies an irreducible Frobenius, and the last four primes test saturation of the four rational points used above.

\begin{center}\small
\begin{tabular}{rrrrrl}
\toprule
$p$ & $a_p$ & $A_4$ & $v_p(R)$ & $\ell$ & saturation primes\\ \midrule
5 & $-4$ & 1 & 2 & 3 & 13,23,59,71\\
7 & $-5$ & $21\bmod49$ & 5 & 5 & 3,19,47,79\\
11 & $-6$ & 2 & 4 & 3 & 127,167,199,241\\
13 & $-6$ & 11 & 4 & 11 & 7,97,127,181\\
17 & $-6$ & 4 & 4 & 3 & 41,71,101,151\\
19 & $-8$ & 1 & 4 & 5 & 149,449,631,829\\
23 & $-6$ & 12 & 4 & 3 & 61,163,419,773\\
29 & $-6$ & 2 & 4 & 3 & 83,701,727,971\\
31 & $-4$ & 23 & 4 & 5 & 109,223,457,601\\
\bottomrule
\end{tabular}
\end{center}
\begin{proof}
Each prime is good ordinary by its displayed trace. We first detail the uniform modular-symbol calculation and the additional precision at $7$. For precision modulus $p^h$, $h=1$ or $2$, set $m=p^{h+1}$. Lift the unit root $\alpha$ of $X^2-a_pX+p$ modulo $p^h$. For a unit residue $a\bmod m$ recover the unique $j\bmod p^h$ such that
\[
             (1+p)^{(p-1)j}\equiv a^{p-1}\pmod m .
\]
By the ordinary measure formula of Stein and Wuthrich, §3, this residue contributes to $[T^k]\mathcal L_p$ modulo $p^h$ the quantity
\[
 \binom jk\frac{\alpha^{-(h+1)} S_{m}(a)
                       -\alpha^{-(h+2)} S_{m/p}(a)}{2}
               \qquad(0\le k\le4).
\]
Here $p\ge5$, so $k!$ is a unit and the Mahler coefficient depends only on $j\bmod p^h$. To calculate all the same-denominator symbols efficiently, group the coefficients $a_nq^n/n$ by $n\bmod m$ and take a discrete Fourier transform in complex balls. Its real entries give the needed cosine sums. Use $q=e^{-B/m}$, an integer cutoff $M$ with $4q^{M+1}/(1-q)<1/16$, and the real period enclosed above. All symbol intervals contain a unique integer by the previously cited integrality theorem. The longest integer Fourier vector, of length $1138363$, was independently reconstructed as follows. The change $X=4x-1$, $U=4(2y+x)$ gives $U^2=X^3-1267X+17230$. At a good prime $l$ and a reduced rational point $P$, the trace $a_l$ is one of the integers $a$ with $|a|\le2\sqrt l$ and $(l+1-a)P=0$. Write $H=\lfloor2\sqrt l\rfloor$, $d=\lfloor\sqrt{2H}\rfloor+1$ and $a=-H+jd+i$ with $0\le i<d$. Matching the list of $iP$ against $(l+1+H)P-jdP$ supplies every candidate. Testing the remaining three rational points leaves a unique answer except at $17,53,181$, where direct quadratic-character sums give the trace. The resulting exact integer vector agrees with PARI at every entry. At $p=7$ the quartic coefficient vanishes modulo $7$. Taking $m=343$, $M=379073$ and $\alpha\equiv16\pmod{49}$ gives $21\pmod{49}$, of valuation exactly one. The other eight quartics are units. In every case the first four coefficients vanish modulo the working precision, while the exact order lower bound already follows from the rank inequality.

Next we check integral divisibility, which is needed to control leading valuations. At $q=117223$ the discriminant valuation is one and reduction is split multiplicative. For $p\ne q$ as in the table, Tate uniformization describes the $p$-torsion by a primitive $p$th root of unity and a $p$th root of the Tate parameter. The first lies in an unramified extension, while the parameter has valuation one, so inertia acts by a nontrivial transvection on $E[p]$. At the displayed auxiliary $\ell$ the discriminant $a_\ell^2-4\ell$ is a nonsquare modulo $p$. In successive rows its residues are $2,3,8,5,14,15,20,26,27$. Thus the image has no invariant line over $\mathbb F_p$. Conjugate the transvection by an element moving its fixed line. In the basis of the two fixed lines the resulting two transvections and their powers give both elementary unipotent subgroups, hence generate $\mathrm{SL}_2(\mathbb F_p)$. The Weil pairing supplies the full determinant. The image is $\mathrm{GL}_2(\mathbb F_p)$. For $p\ge5$ this gives surjectivity on the Tate module by Stein--Wuthrich, Proposition~7.2. In particular $E(\Q)[p]=0$.

Let $f_E(T)$ be a characteristic series of the cyclotomic dual Selmer group. Kato's integral divisibility in Stein--Wuthrich, Theorem~7.3, reads $\mathcal L_p=f_E d$ with $d\in\mathbb Z_p[[T]]$. Schneider and Perrin--Riou, in Theorem~6.1 there, give
\[
  4\le\operatorname{rank}E(\Q)\le\operatorname{ord}_T f_E
                     \le\operatorname{ord}_T\mathcal L_p\le4 .
\]
Equality implies finiteness of $\operatorname{Sha}(E/\Q)[p^\infty]$ and nondegeneracy of the height pairing. The same theorem gives the valuation of the leading coefficient $f_E^*$,
\[
 v_p(f_E^*)=v_p\!\left((1-\alpha^{-1})^2
     \frac{\prod_\ell c_\ell}{\# E(\Q)[p^\infty]^2}
     \#\operatorname{Sha}(E/\Q)[p^\infty]\,
     \frac{R}{\log_p(1+p)^4}\right)
                  \le v_p([T^4]\mathcal L_p).
\]
The inequality uses integral divisibility in the displayed direction. The local Tamagawa numbers at the bad primes are $c_2=2$ and $c_{117223}=1$ from the multiplicative reduction types $I_2,I_1$. All Tamagawa and rational torsion valuations therefore vanish in these rows.

For $p\ge7$ in the table, $p\nmid\#E(\mathbb F_p)$ by the traces. In particular the multiplier is a unit, and $R/\log_p(1+p)^4$ is $p$-integral. Here is a direct justification of the latter integrality, independent of the regulator calculation. Multiply any rational point $P$ by $n=2\#E(\mathbb F_p)$, a $p$-adic unit. Then $nP$ has nonsingular reduction at every prime and lies in the formal group at $p$. The ordinary canonical sigma-function formula, Stein--Wuthrich, equation (4.1), gives its height as twice the $p$-adic logarithm of an element of $\Q_p^\times$. This logarithm belongs to $p\mathbb Z_p$ for odd $p$. Division by $n^2$ preserves this ideal. Polarization and division by $\log_p(1+p)$ now give integral pairing entries on the entire Mordell--Weil lattice. Thus the valuation of the finite group order is at most zero for $p\ge11$ in the table and at most one for $p=7$. A finite odd-primary Tate--Shafarevich group of an elliptic curve has square order by the nondegenerate alternating Cassels--Tate pairing. Hence the group is trivial also at $7$. At the anomalous prime $5$ its triviality was already proved by Wuthrich in the cited §9 example. See also B.~Poonen and M.~Stoll, \href{https://math.stanford.edu/~conrad/BSDseminar/refs/PoonenStoll.pdf}{\emph{The Cassels--Tate pairing on polarized abelian varieties}}, for the alternating pairing in the elliptic case.

As an independent check including $5$, the regulator column yields zero as the upper bound on the group valuation in every row, without the square-order argument. To verify the lattice input, at each listed saturation prime $l$ the order $s_l=\#E(\mathbb F_l)$ is exactly divisible by $p$ to the first power. Multiply the reductions of the four rational points by $s_l/p$ and express them in the cyclic group of order $p$. The four resulting rows have rank four over $\mathbb F_p$. The individual row vectors and their generators are supplied in the arithmetic record. This proves that the span of the points has index prime to $p$ in the rank-four Mordell--Weil lattice. Computing its canonical ordinary height determinant with \texttt{ellpadicregulator} at precisions nine and fourteen gives the displayed valuations at both precisions. A convention multiplying the pairing by two has no effect on these odd-prime valuations. The logarithm denominator has valuation four. Finally $v_p(1-\alpha^{-1})=v_p(1-a_p+p)$ since the second root of the Frobenius polynomial is divisible by $p$. Thus the multiplier valuation is two at $5$ and zero elsewhere. The resulting sums of regulator and multiplier valuations are $0,1,0,0,0,0,0,0,0$, exactly the respective analytic quartic valuations. This proves the asserted bounds by the independent regulator calculation as well.
\end{proof}

For the Heegner field $K=\Q(\sqrt{-7})$ the same nine primary groups of $E/K$ are consequently finite, and there is an explicit identification
\[
     \operatorname{Sha}(E/K)[p^\infty]
       \simeq\operatorname{Sha}(E_-/\Q)[p^\infty]
              \qquad(5\le p\le31).
\]
Indeed the sum and difference of the two Galois conjugates give an isogeny between the Weil restriction of $E/K$ and $E\times E_-$. The inverse up to multiplication by two is given by the same two-character decomposition. At an odd prime this gives an isomorphism on $p$-primary global and local cohomology, hence on the kernels of localization. Shapiro's lemma identifies the Weil-restriction kernel with $\operatorname{Sha}(E/K)$. The other eigenspace is zero by Theorem~\ref{r4:sha}. Gross--Zagier and Kolyvagin give finiteness for the analytic-rank-one twist over $\Q$.

These nine arithmetic comparisons include no complex second-derivative identity. In particular the higher-rank obstruction remains even for a curve with nondegenerate height pairings and vanishing Tate--Shafarevich primary parts at every odd ordinary prime through $31$.

\subsection{The sign of every higher even Taylor coefficient}\label{r4:kernel}

The single undecided Taylor coefficient has a further distinction. Write $M_j=F^{(2j)}(0)=2\int_0^\infty u^{2j}K(u)\,du$.

\begin{theorem}\label{r4:positive-derivatives}
There are exactly two zeros $0<\alpha<\beta$ of $K$ on the half-line $u\ge0$. Both are simple, $1<\alpha<2$ and $4<\beta<5$. For every $j\ge2$ one has
\[
             M_j>0,\qquad M_{j+1}>\beta^2 M_j .
\]
For every real $x$ and every $j\ge2$, the derivative $F^{(2j)}(x)$ is strictly positive.
\end{theorem}
\begin{proof}
Write $q=\exp(-Be^u)$ and $P(q)=\sum_{k\ge0}a_{k+1}q^k$, so $K(u)=e^u qP(q)$. As $u$ increases, $q$ decreases strictly from $e^{-B}\in(0.945,0.99)$ to zero. The following table gives uniform interval signs. The last column gives the number of closed subintervals in the certificate.
\begin{center}\small
\begin{tabular}{cccc}
\toprule
$q$ interval & function & sign & subintervals\\ \midrule
$[0,0.42]$ & $P$ & $+$ & 5\\
$[0.42,0.43]$ & $P'$ & $-$ & 1\\
$[0.43,0.94]$ & $P$ & $-$ & 24\\
$[0.94,0.945]$ & $P'$ & $+$ & 32\\
$[0.945,0.99]$ & $P$ & $+$ & 41\\
\bottomrule
\end{tabular}
\end{center}
Here are details of the finite sign test. At each midpoint use the first $m=4096$ terms of $P$ and $P'$. With $Q=0.99$ and $c=1-Q$, the uniform omitted radii are
\[
 E_0=2Q^m((m+1)-mQ)/c^2,\qquad
 E_1=2Q^{m-1}\left(\frac{m(m+1)}c+
     \frac{(2m+1)Q}{c^2}+\frac{Q(1+Q)}{c^3}\right).
\]
These follow by summing $2(k+1)Q^k$ and $2k(k+1)Q^{k-1}$ for $k\ge m$. On a subinterval of radius $d$ and right endpoint $r$ the bound $|P''|\le C=12/(1-r)^4$ follows by twice differentiating $2/(1-q)^2$. If the midpoint values of the two polynomials are $v,w$, use radii $E_0+d(|w|+E_1)+Cd^2/2$ for $v$ and $E_1+Cd$ for $w$. Bisection gives the displayed strict signs at 128-bit precision. The endpoint checks $q(1)>0.945$, $q(2)<0.94$, $q(4)>0.43$, $q(5)<0.42$ place the two zeros in the asserted $u$ intervals.

Set $A=\alpha^2$ and $C=\beta^2$. The kernel is positive, negative and positive on the three successive intervals. Hence $(u^2-A)(u^2-C)K(u)$ is nonnegative on the half-line and positive off the two roots. If $D_j=M_{j+1}-CM_j$, then
\[
        D_{j+1}-A D_j
        =2\int_0^\infty u^{2j}(u^2-A)(u^2-C)K(u)\,du>0 .
\]
Since $C<25$, the bounds for $M_1=2b_2$ and $M_2=24b_4$ give $D_1>0$. Consequently $D_j>0$ for every $j\ge1$, and induction from $M_2>0$ proves the inequalities. For $j\ge2$ the even Taylor expansion of $F^{(2j)}(x)$ has strictly positive coefficients at every even power and positive constant term. Entireness permits its evaluation at every real $x$.
\end{proof}

\subsection{Resolving the quartet at two signs of heat time}\label{r4:heat}

Let $H_t(w)=\int_{\R}K(u)e^{tu^2+iwu}\,du$ in the physical height variable of Section~\ref{heat:main}, so $H_0(w)=F(iw)$. Take $T$ to be the largest height in the table in Section~\ref{r4:slab}.

\begin{theorem}\label{r4:heat-thm}
For every
\[
                     2^{-8230}\le t\le 2^{-6700},
\]
all the zeros of $H_t$ in $|\Re w|\le T$ are real and simple. Four belong to $|w|<5\sqrt t$. The two positive ones in that disk belong to $(0,2\sqrt t)$ and $(2\sqrt t,4\sqrt t)$. For $H_{-t}$ there are instead four simple purely imaginary zeros in the small disk, one in each reflection of $i(0,2\sqrt t)$ and $i(2\sqrt t,4\sqrt t)$. Inside $|\Re w|\le T$, $|\Im w|\le1$ these are the only nonreal zeros of $H_{-t}$.
\end{theorem}
\begin{proof}
For $\eta=\pm1$, $|a|\le5$, consider the even entire function of $x$
\[
       G(x)=\int_{\R} K(u)e^{\eta x^2u^2}\cos(axu)\,du.
\]
For $|x|\le R=1/32$ it has modulus bounded by $8C_1/B^2<2^{32}$. In fact $u^2\le e^u$ for $u\ge0$, $R^2<B/2$, $5R<1$, and reflection reduces the bound to $2C_1\int_0^\infty e^{2u-(B/2)e^u}\,du$. Inserting $x=\sqrt t$, the two leading coefficients and Cauchy's tail bound give
\[
 H_{\eta t}(a\sqrt t)
  =(2\eta-a^2)b_2 t+(a^4-12\eta a^2+12)b_4t^2+E,\quad
 |E|\le\frac{2^{32}(\sqrt t/R)^6}{1-t/R^2}.
\]
The exact constant term is zero. Over the whole stated interval, $|b_2|/(b_4t)<3$ and $|E|/(b_4t^2)<1$. At $\eta=1$ the quartic polynomial at $a=0,2,4$ takes the respective values $12,-20,76$. The possible quadratic contribution at those points, in these units, is at most $6,6,42$. Thus the three strict signs are $+,-,+$. At $\eta=-1$ take $a=0,2i,4i$. The same signs follow, since the quartic has the same values and the quadratic changes sign. Both functions are real along the test axes by evenness and real Taylor coefficients. On $|a|=5$ the leading monomial $a^4$ dominates the combined errors, as
\[
                 5^4-12\cdot5^2-12-3(2+5^2)>1.
\]
Rouché gives four zeros in each small disk. The sign changes and reflection already supply four distinct zeros, proving simplicity and the given locations.

To retain the outer counts, work in the centered arithmetic coordinate $z=iw$. For every real $h$ with $|h|\le2^{-6700}$ and $|\Re z|\le1$,
\[
 \left|H_h(-iz)-F(z)\right|
 \le 2C_1|h|\int_0^\infty e^{3u-(B/2)e^u}\,du
 \le 32C_1|h|/B^3 < 2^{40}|h|.
\]
Indeed $|e^{hu^2}-1|\le |h|u^2 e^{|h|u^2}$ and $u^2\le e^u$ on $u\ge0$. The bound is independent of $\Im z$.

The existing horizontal edge and critical-line intervals enclose these perturbed values too. More precisely the boundary-integral estimate underlying their trapezoid error discarded the factor
\[
 b^2\int_0^\infty e^{2x-be^x}\,dx=(1+b)e^{-b}<1
\]
when it bounded each half-line by $1/b^2$. Here $b=B\cos1$. The stated error budget for the exact finite exponential sum therefore exceeds the required undeformed budget by at least
\[
 \bigl(1-(1+b)e^{-b}\bigr)
       \frac{4C_0e^T}{b^2(e^{2\pi D}-1)}.
\]
For the rows at heights $64,256,1024$ its base-two logarithm has floor $-399,-993,-3366$, respectively. At $(T,D)=(2048,1048)$ the floor is $-6529$. Each exceeds $-6659$, while the perturbation is at most $2^{-6660}$. This leaves every horizontal subdivision in its original open half-plane and preserves each certified sign at a line endpoint.

On the right vertical edge $s=2+iv$, the Euler product gives $|L(E,s)|\ge\zeta(3/2)^{-2}>1/9$. Indeed each normalized local reciprocal is bounded by $(1+p^{-3/2})^2$, and $\zeta(3/2)<3$ by integral comparison. The identity $|\Gamma(2+iv)|^2=(1+v^2)\pi v/\sinh(\pi v)$ gives $|\Gamma(2+iv)|\ge e^{-\pi|v|/2}$, including the limit at zero. As $B<1$ and $T\le2048$, the completed function thus has modulus greater than $2^{-4650}$ everywhere on this edge. Reflection gives the other edge. The continuous homotopy from $0$ to $h$ avoids all four boundary edges, so its rectangle count is constant.

The existing positive brackets, their reflections and the four small-disk zeros now exhaust that count for both signs. For positive $t$, Theorem~\ref{heat:real} already confines every zero worldwide to $|\Im w|\le1/2$. This gives the asserted whole slab. For the negative times the rectangle assertion uses only the preceding count.
\end{proof}

\begin{corollary}\label{r4:heat-ray}
Throughout the physical height window $|\Re w|\le1$, every zero of $H_t$ is real and simple for every $t\ge 2^{-8230}$. There are exactly four such zeros for $2^{-8230}\le t\le1/8$.
\end{corollary}
\begin{proof}
Use the fixed boundary $\Re w=1$, $|\Im w|\le3/4$. At $0\le t\le1/8$ one has uniformly
\[
                    \Re H_t(1+iv)>1\qquad (|v|\le3/4).
\]
To check the inequality take $D=24$, $W=12$, $M=14000$ and the weights of Section~\ref{r4:quadrature}. The real part of the finite exponential sum becomes
\[
       w_0+2\sum_{j=1}^{DW} w_j
             \cos(j/D)e^{t(j/D)^2}\cosh(vj/D).
\]
For $x\ge0$ one has $x^2/8\le18+(b/2)e^x$. For $x\le12$ the constant suffices, and beyond $12$ the ratio $x^2e^{-x}$ is decreasing with $144e^{-12}/8<b/2$. Applying the previous contour shift by one to the heat integrand consequently bounds the discretization error by
\[
   \frac{16 C_0 e^{19}}{b^2(e^{2\pi D}-1)}
        +\frac{32 C_1e^{18}}{B^2}\exp(-Be^W/4)<2^{-100},
\]
including the real-mesh tail. In the real sum the multiplier $e^{t(j/D)^2}\cosh(vj/D)$ increases in $t,v$ on the positive quadrant. Use its endpoint bounds on $0\le v\le3/4$, with $t$ in $[0,1/16]$ and $[1/16,1/8]$ successively. Direct interval sums give the claimed strict lower bound on both rectangles. Reflection covers the negative heights and the opposite edge.

The remaining two edges $\Im w=\pm3/4$ contain no zeros at any of these times by Theorem~\ref{heat:real}. The number of zeros in this fixed rectangle stays four throughout $[0,1/8]$, starting with the cluster count at time zero. At time $2^{-8230}$ all four are distinct and real. This property persists until $1/8$. Indeed simple real zeros persist locally by conjugation and the implicit function theorem. At a first limiting time the four zeros would still be real by compactness and the boundary bounds. A multiple one is impossible by the backward splitting formula in~\eqref{heat:splitting}: immediately before that time it would have had nonreal zeros in every sufficiently small fixed disk around it. This gives a contradiction. For $t>1/8$ Theorem~\ref{heat:real} gives simplicity and reality in the entire plane.
\end{proof}

\subsection{A finite-height transition time determined by one coefficient}\label{r4:transition}

The same positive-time conclusion persists to substantially greater heights. For a positive height $Y$, call the boundary certified if
\[
       H_t(Y+iv)\ne0\qquad
                (0\le t\le1/8,\ |v|\le3/4).
\]
There is one unresolved scalar in the phase diagram, independent of the certified boundary chosen.

\begin{theorem}\label{r4:phase}
Let $Y$ be any of the heights in the following table. If $L''(E,1)\ge0$, all zeros in $|\Re w|\le Y$ are real for every $t\ge0$ and all are simple for every $t>0$. If $L''(E,1)<0$, there is a unique $t_*>0$ with $H_{t_*}(0)=0$, and
\[
          0<t_*<\frac{-b_2}{6b_4}<2^{-8231}.
\]
For $0<t<t_*$ precisely two zeros in this slab are nonreal. They are simple and purely imaginary, and all other zeros are real and simple. At $t=t_*$ the pair becomes a double zero at the origin. At every $t>t_*$ all zeros in the slab are real and simple. The counts in the table hold for all $0\le t\le1/8$, including multiplicities.
\end{theorem}

\begin{center}\small
\begin{tabular}{rrrrrrrr}
\toprule
$Y$ & bits & $D$ & $W$ & $M$ & degrees & boxes & zeros\\ \midrule
$477/8$ & 492 & 35 & 11 & 15863 & $(33,19)$ & 7 & 284\\
$253$ & 1268 & 114 & 12 & 39271 & $(40,22)$ & 15 & 1432\\
$15411/16$ & 4108 & 405 & 14 & 125238 & $(48,24)$ & 18 & 6266\\
\bottomrule
\end{tabular}
\end{center}

To certify each boundary, use the preceding heat quadrature with $e^{19}$ replaced by $e^{18+Y}$ in the error bound. Put $p=16t-1$, $q=4v/3$, $z=q^2$. It suffices to treat $|p|\le1$, $0\le q\le1$. In the finite exponential sum expand $e^{pu^2/16}$ through degree $J$ in $p$, and $e^{-3qu/4}$ through degree $K$ in $q$. If $P=W^2/16$, $Q=3W/4$, and $S=|w_0|+2\sum_{j>0}|w_j|$, the combined Taylor remainder is at most
\[
  S e^{2P+Q}\left(\frac{P^{J+1}}{(J+1)!}
                         +\frac{Q^{K+1}}{(K+1)!}\right).
\]
This follows by truncating the two exponentials successively with $u$ and $-u$ in the finite mesh. Reflection makes the real polynomial even in $q$ and the imaginary polynomial odd. Thus the approximation after multiplication by $e^{\pi Y/2}$ is $R(p,z)+iq I(p,z)$. All coefficients are enclosed by summing signed moments of $w_j e^{u^2/16+iYu}$. The initial pairs $(J,K)$ in successive rows were $(93,96)$, $(209,209)$ and $(557,535)$. In the remaining high degrees the sum of the absolute normalized coefficients supplies a uniform bound. Discarding them gives the degrees in $(p,z)$ tabulated above and a total normalized error less than $2^{-22}$ in each coordinate.

For completeness, the finite range proof uses the convex hull property of the Bernstein basis. For degree $d$ on $[0,1]$ this basis consists of $\binom di x^i(1-x)^{d-i}$, nonnegative functions summing to one. The degree-$d$ Bernstein coefficients of $x^k$ are $\binom{i}{k}/\binom{d}{k}$. For $(2x-1)^k$ they are
\[
 \frac{1}{\binom dk}
       \sum_{r=0}^k(-1)^{k-r}\binom ir\binom{d-i}{k-r}.
\]
These follow by expanding the product of $k$ factors, with $i$ positive and $d-i$ negative choices when restricted to a subset of size $k$. The tensor product of these formulas converts $p^j z^k$ using $x=(p+1)/2$. Bernstein subdivision at $1/2$ takes repeated pairwise arithmetic averages and produces exact coefficient enclosures on both half intervals. On every final parameter rectangle either all real coefficients of $R$ have one sign beyond the error, or the same is true of the coefficients of $I$ after multiplying the sign margin by the minimum $q$. This proves boundary nonvanishing over the continuum of parameter values.

At time zero all positive non-cluster roots come with the disjoint endpoint intervals of Section~\ref{r4:slab}. Counting their right endpoints below $Y$ gives $140$, $714$ and $3131$ in the respective rows. The next interval in each case starts above $Y$. The four cluster zeros and the negative reflections give the stated total. The global strip and the new vertical bound protect this rectangular count through all $t\in[0,1/8]$.

\begin{proof}[Proof of Theorem~\ref{r4:phase}]
Write $\Phi(t)=H_t(0)$. The even moments and entireness in $t$ give
\[
   \frac{\Phi(t)}t=2b_2+12b_4 t+
                      \sum_{j\ge3}\frac{M_j}{j!}t^{j-1}.
\]
By Theorem~\ref{r4:positive-derivatives}, every coefficient after the constant is positive. The right-hand side is therefore strictly increasing for $t\ge0$. It has a positive zero exactly when $b_2<0$, it has at most one such zero, and in that case it is already positive at $t=-b_2/(6b_4)$. Its derivative there and at its zero is positive. The inequality $6b_4>2^{12}$ gives the displayed bound for $t_*$.

Here is the zero-count argument that identifies this scalar zero with the last nonreal pair. Within any fixed bounded rectangle whose boundary avoids zeros, the number of nonreal zeros cannot increase in forward heat time. To see this locally at a time $s$, enclose the finitely many distinct zeros in disjoint small disks. Disks about nonreal zeros can avoid the real axis. Rouché preserves each disk's total multiplicity for nearby times. Simple real zeros stay real, and every multiple real zero of order $m$ has the Hermite splitting of~\eqref{heat:splitting}. All $m$ descendants are simple and real immediately afterwards, whereas immediately beforehand at least two were nonreal for $m\ge2$. This proves local monotonicity, hence monotonicity on the entire compact time interval. It also shows that every multiple real zero strictly decreases the count from its left.

If $b_2\ge0$, the initial count of nonreal roots is zero. All roots are simple for positive times by the same backward splitting. Suppose $b_2<0$. There are two nonreal roots at time zero, and their initial nearby descendants stay nonreal. Whenever there are just two, they are simple and purely imaginary, since conjugation and evenness would otherwise supply either four distinct roots or total multiplicity at least four. Their only possible meeting with the real axis is at the origin. Every such meeting requires $\Phi(t)=0$, so it can occur only at $t_*$. At that time $\Phi'(t_*)>0$, while the heat equation gives $H_{t_*}''(0)=-\Phi'(t_*)$. The multiplicity is exactly two. The Hermite splitting makes these roots real just afterwards. Monotonicity excludes any further nonreal root or multiple zero at later times within the certified interval. Finally Theorem~\ref{heat:real} gives reality and simplicity for all $t>1/8$.
\end{proof}

\begin{corollary}\label{r4:upperphase}
The phase description of Theorem~\ref{r4:phase} holds throughout the fixed physical height slab
\[
                    |\Re w|\le \frac{65257}{32}=2039.28125 .
\]
In particular every zero in the slab is real and simple for every $t\ge2^{-8231}$.
\end{corollary}
\begin{proof}
Use nested protected rectangles during two successive time intervals. The undeformed count in this slab is $14236$. The first row below uses the indicated wider rectangle. The counts include multiplicities.

\begin{center}\small
\begin{tabular}{rrrrrr}
\toprule
$Y$ & heat interval & initial degrees & retained degrees & boxes & zeros\\ \midrule
$2045$ & $[0,1/16]$ & $(832,934)$ & $(35,24)$ & 7 & 14280\\
$65257/32$ & $[1/16,1/8]$ & $(833,934)$ & $(39,26)$ & 7 & 14238\\
\bottomrule
\end{tabular}
\end{center}
Both calculations use $5504$-bit balls and the common grid $D=847$, $W=14$, $M=256191$. In each row write the time as the midpoint of its interval plus its half-length times $p\in[-1,1]$. The same Taylor remainder and Bernstein argument apply with $W^2/16$ in the time-variable tail replaced by $W^2/32$, and with the midpoint heat weight in each signed moment. In the first interval, the original height-$2048$ inventory contains $7138$ positive non-cluster intervals below $2045$, with the next interval entirely above it. Thus the count is $4+2\cdot7138$. At $1/16$ all these roots are real and simple. Thereafter the second, narrower rectangle inherits this property by nonvanishing on its own boundary and backward Hermite splitting. At $t>1/8$ the global heat theorem applies again.

For the exact second count we also measured the positive real roots between the boundaries at the common time $t=1/16$. Put $c=(2045+65257/32)/2$, $a=(2045-65257/32)/2$, and $x=(w-c)/a$. Approximate $e^{\pi c/2}H_{1/16}(c+ax)$ on $-1\le x\le1$. At distance $1/2$ from the real segment the same quadrature bound applies with the horizontal factor $e^{18+2045+1/2}$. Cauchy's estimate multiplies that error by at most $2a$ for the $x$-derivative. For the finite sum expand $e^{iaxu}$ through degree $n=1316$. With $A=aW$ and $S$ the mass of the heat-weighted mesh after the same scaling, the value remainder is at most $S e^A A^{n+1}/(n+1)!$. Its derivative remainder is at most $n+1$ times that bound. The sum of $k$ times the absolute $k$th coefficients also bounds the errors from discarding the highest positive degrees, simultaneously for the value and derivative. The retained degree is 123 and the joint error is below $2^{-18}$. On each dyadic subinterval use the Bernstein coefficient range for the value, or a strict derivative sign from consecutive coefficient differences together with endpoint signs. This certifies exactly 21 positive simple roots in the connecting band, with every interval recorded. The $14280-2\cdot 21=14238$ remaining zeros stay inside the narrower rectangle through $t=1/8$. Since the undeformed count at that narrower height was $14236$, real zeros crossed its boundary before the switch. The wider first rectangle controls this change of count.
\end{proof}

In the second case the transition time is itself fixed to high relative precision in terms of the exact $b_2$. On $|t|=1/1024$ the bound $u^2/1024<(B/2)e^u$ for $u\ge0$ gives $|\Phi(t)|\le4C_1/B<2^{23}$. Cauchy's estimate at $t=t_*$ consequently yields
\[
 0<\frac{-b_2}{6b_4t_*}-1
 \le\frac{2^{23}\,1024^3\,t_*}{12b_4(1-1024t_*)}<2^{-8170}.
\]
Thus determining the missing archimedean sign would also decide whether the finite-height transition is at zero or at this explicit positive root.

The center admits an exact description even at arbitrarily large negative times. There are exactly two real zeros of $t\mapsto\Phi(t)$ counted with multiplicity. If $b_2=0$, they coincide at zero. If $b_2\ne0$, one is zero and the other is a simple real number $s_*$ of sign opposite to $b_2$, with $|s_*|<2^{-8230}$. In the latter case these are the only two heat times at which the center is a zero, and the spatial multiplicity at each is two. In the former case the center is a zero only at time zero and has spatial multiplicity four there.

Indeed, for $A=\alpha^2$ and $C=\beta^2$ from Theorem~\ref{r4:positive-derivatives}, differentiation under the heat integral gives for every real $t$
\[
 (\partial_t-C)(\partial_t-A)\Phi(t)
  =2\int_0^\infty (u^2-A)(u^2-C)K(u)e^{tu^2}\,du>0.
\]
If $\Phi$ had three real zeros counted with multiplicity, two applications of Rolle's theorem, after the respective integrating factors $e^{-At}$ and $e^{-Ct}$, would contradict this strict inequality. At zero we have $\Phi'(0)=2b_2$ and $\Phi''(0)=24b_4>0$. If $b_2=0$ these already give the double zero in time. If $b_2\ne0$, $\Phi$ is strictly negative just to the side of zero opposite to $b_2$. The sign tests at $a=0$ in Theorem~\ref{r4:heat-thm} show that $\Phi(2^{-8230})$ and $\Phi(-2^{-8230})$ are both strictly positive. This produces the other zero with the stated location. The at-most-two result ensures it is simple. Evenness in $w$ and $H_t''(0)=-\Phi'(t)$ then yield the spatial multiplicities. In particular the unresolved assertion of analytic rank four is equivalent to the strict positivity of $H_t(0)$ at every nonzero real heat time.

\begin{theorem}\label{r4:moment-envelope}
For every $k\ge1$ the real function $\Phi^{(k)}(t)$ has exactly two simple real zeros $a_k<b_k$. For $k\ge1$ they interlace according to
\[
               a_{k+1}<a_k<b_{k+1}<b_k .
\]
Both zeros are negative for $k\ge2$. Put $r_1=\alpha^2$, $r_2=\beta^2$, $J(v)=K(\sqrt v)$ and $m=k+1/2$. Then as $k\to\infty$,
\[
  a_k=-\frac m{r_1}-\frac{J''(r_1)}{2J'(r_1)}+o(1),\qquad
  b_k=-\frac m{r_2}-\frac{J''(r_2)}{2J'(r_2)}+o(1).
\]
The following open intervals enclose the two slopes and additive constants after multiplication by $10^{10}$.
\[
\begin{array}{c|rr}
 &10^{10}\cdot\text{\rm slope}&10^{10}\cdot\text{\rm constant}\\ \hline
a_k &(-4400219988,-4400219987)&(1364938545,1364938546)\\
b_k &(-569165708,-569165707)&(475391432,475391433)
\end{array}
\]
\end{theorem}
\begin{proof}
Each additional time derivative multiplies the kernel by $u^2$. Thus $(\partial_t-r_2)(\partial_t-r_1)\Phi^{(k)}(t)>0$ holds without changing the two crossing points. Each derivative has at most two real zeros counted with multiplicity. At $t=-h$ with $h$ large and fixed $k$, substituting $v=hu^2$ gives
\[
   h^{k+1/2}\Phi^{(k)}(-h)
       =\int_0^\infty v^{k-1/2}e^{-v}K(\sqrt{v/h})\,dv
       \longrightarrow \Gamma(k+1/2)K(0)>0 .
\]
The kernel is bounded, so dominated convergence applies. When $\Phi$ has two simple zeros $a<b$, its derivative is negative at $a$ and positive at $b$. Its far negative sign just proved therefore puts one zero of $\Phi'$ before $a$ and one between $a,b$. When $\Phi$ has a double zero at zero, $\Phi'$ has a simple zero there with positive slope and is strictly negative just to its left. Its far negative sign gives the other root. This starts the induction. The same argument at every following pair of simple zeros proves the strict interlacing and exhausts the possible multiplicities. The positivity of all $M_j$ for $j\ge2$ shows directly that $\Phi^{(k)}(t)>0$ for $t\ge0$, $k\ge2$.

For the limiting constants substitute $v=u^2$ without rescaling. The sign of $\Phi^{(k)}(-h)$ for $h>0$ is that of $\mathbb E J(V)$, where $V$ has the gamma density $h^m v^{m-1}e^{-hv}/\Gamma(m)$. Fix $r=r_1$ or $r_2$ and take $h=m/r-s$ with fixed real $s$. The mean and variance give
\[
  \mathbb E(V-r)=\frac{r^2s}{m}+O(m^{-2}),\qquad
  \mathbb E[(V-r)^2]=\frac{r^2}{m}+O(m^{-2}).
\]
The fourth centered moment is $O(m^{-2})$ by the gamma moment formula, so the third absolute centered moment is $O(m^{-3/2})$. On a fixed neighborhood of $r>0$, Taylor-expand $J$ through degree two. The probability outside that neighborhood is $O(m^{-2})$ by the fourth moment. The same fourth-moment bound controls the expectations of the discarded polynomial terms there, since $r$ and the neighborhood are fixed and $J$ is bounded globally. As $J(r)=0$, this yields
\[
 m\mathbb E J(V)=r^2\bigl(sJ'(r)+\tfrac12 J''(r)\bigr)+O(m^{-1/2}).
\]
The two roots of $K$ are simple, hence $J'(r)\ne0$. Test $s$ on either side of $-J''(r)/(2J'(r))$, arbitrarily closely. For large $m$ the signs are opposite. The two resulting windows are disjoint for $r_1,r_2$. As there are precisely two roots, this proves both formulas.

For the supplied numerical constants, bisection of the same $P(q)$ intervals gives $1.5075190374<\alpha<1.5075190375$ and $4.1916077742<\beta<4.1916077743$. At either root write $x=u$, $c=Be^x$ and $q=e^{-c}$. Logarithmic differentiation at the simple root gives
\[
 -\frac{J''(x^2)}{2J'(x^2)}
   =\frac{1}{4x^2}-\frac{1}{4x}
                \left(3(1-c)-cq\frac{P''(q)}{P'(q)}\right).
\]
The code bisects the $q$ brackets 92 times and uses the geometric series tail on complex disks of radius $1/100$, followed by Cauchy's estimate for its first two derivatives. Interval evaluation then gives the stated rational endpoints.
\end{proof}

\subsection{The natural heat transform of the degree-four base change}\label{r4:degreefour}

The rank-five example over $K=\Q(\sqrt{-7})$ also gives a heat statement for a degree-four automorphic function over $\Q$, with its natural single physical height. Let $K_+$ and $K_-$ be the original even and twisted odd Mellin kernels. Write
\[
    G_t(w)=\frac1i\int_{\R}(K_+*K_-)(u)e^{tu^2+iwu}\,du .
\]
At time zero this is $F(iw)F_-(iw)/i$. The integral deforms the completed product itself, including the cross term in $e^{t(x+y)^2}$.

\begin{theorem}\label{r4:degreefour-thm}
For $2^{-8230}\le t\le2^{-6700}$ the $319$ zeros of $G_t$ in the full physical slab $|\Re w|\le32$ are simple and real. One is at zero. Its next two positive zeros lie in $(\sqrt t,3\sqrt t)$ and $(3\sqrt t,5\sqrt t)$, with their two negatives. For $G_{-t}$ these four neighbors lie instead on the imaginary axis at the corresponding absolute heights, and are the only nonreal zeros in $|\Re w|\le32$, $|\Im w|\le1$.
\end{theorem}
\begin{proof}
The centered product has Taylor series $F(z)F_-(z)=a z^3+d z^5+\cdots$. If $c_1,c_3$ are the coefficients of $z,z^3$ in $F_-$, then $a=b_2c_1$ and $d=b_4c_1+b_2c_3$. The twist derivative in Theorem~\ref{r4:heegner-thm} gives $c_1>1$ and its unit-disk bound gives $|c_3|<2^{42}$. Hence $d>688c_1$ and $|a|/(dt)<3$ on the stated interval.

As in the quartic test, put $w=yx$, $t=x^2$ and take $|y|\le6$. The function of $x$ with signed time $\eta x^2$, $\eta=\pm1$, is odd. At $|x|=1/64$ its modulus is less than $2^{76}$. Indeed, in the double integral use $(u+v)^2\le2(u^2+v^2)$, $6/64<1$ and $2/64^2<B_-/2$. Absolute integration then gives the product of the two bounds $8C_{1,+}/B^2$ and $8C_{1,-}/B_-^2$. Cauchy's estimate yields
\[
 G_{\eta t}(y\sqrt t)
 =a t^{3/2}(6\eta y-y^3)+d t^{5/2}(y^5-20\eta y^3+60y)+E,
 \qquad |E|\le\frac{2^{76}(64\sqrt t)^7}{1-64^2t}.
\]
The identity follows by expanding the heat exponential and the Fourier exponential, or by the heat equation. The remainder divided by $dt^{5/2}$ is less than one. At positive time the quintic at $y=1,3,5$ is $41,-117,925$, respectively. The cubic errors, in the same units, are bounded by $15,27,285$. Thus there are two strictly positive zeros between successive test points, two negatives by oddness, and the exact zero at the origin. At negative time take $y=i,3i,5i$ and divide values by $i$, giving the same quintic signs and the opposite cubic. The five zeros exhaust the disk $|w|<6\sqrt t$ by Rouché, since $6^5-20\cdot6^3-60\cdot6-3(6^3+6\cdot6)>1$.

For $|\Im w|\le1$, $|h|\le2^{-6700}$, the product changes by less than $2^{90}|h|$. In the preceding double integral use $|e^{h(u+v)^2}-1|\le2|h|(u^2+v^2)e^{2|h|(u^2+v^2)}$. The integral with one additional squared variable is bounded by $32C_{1,\pm}/B_\pm^3$ for that factor. The other factor has bound $8C_{1,\mp}/B_\mp^2$. Summing the two products and doubling gives the stated factor $2^{90}$. On each horizontal side of the undeformed degree-four rectangle, multiplication gives modulus greater than $2^{-620}$. For each factor separately, the available alias reserve computed in the proof of Theorem~\ref{r4:heat-thm} is greater than $2^{-310}$ at $T=32$, $D=40$. The open half-plane test for its larger error intervals consequently gives this strict lower bound for its exact modulus. On the right or left vertical side each factor has modulus greater than $2^{-80}$ by the same Euler and gamma estimates. The rectangle therefore keeps all $319$ zeros.

It remains to preserve the real brackets. At all $314$ endpoints of the disjoint positive dyadic intervals in Theorem~\ref{r4:heegner-thm}, re-evaluation of both factors in $384$-bit balls gives product modulus greater than $2^{-310}$. The same signs survive the perturbation. Thus $157$ positive and $157$ negative real brackets, together with the central quintet, exhaust the rectangle for both signs of time. Finally the convolution kernel has an absolute bound $C\exp(-c e^{|u|/2})$ for some $c,C>0$. Indeed at least one of $|v|,|u-v|$ is at least $|u|/2$ and one may reserve half the double-exponential decay of both factors to obtain an integrable majorant. This gives the order-one growth and Fourier decay required in the proof of Theorem~\ref{heat:real}. At time zero every zero has $|\Im w|\le1/2$ by the two Euler products and functional equations. Strip contraction confines all positive-time zeros to this same strip. This proves the full-slab assertion.
\end{proof}

\subsection{Reproducibility and prior results}
The integer recursions, outward ball bounds, and complete endpoints for the $7149$ positive brackets are supplied with the companion records. The \href{https://www.lmfdb.org/EllipticCurve/Q/234446/a/1}{elliptic-curve record} identifies the rational Mordell--Weil rank. Consult the LMFDB \href{https://www.lmfdb.org/L/2/234446/1.1/c1/0/2/Reliability}{analytic reliability note} for the rank alternative, and Y.~Kezuka and D.~Zagier, \href{https://people.mpim-bonn.mpg.de/zagier/files/preprints/BSD_Dichotomy.pdf}{\emph{A Birch--Swinnerton-Dyer dichotomy}}, §4, pp.~9--10, for the status of exact vanishing of the second derivative. The arithmetic-rank-zero implication used above is the theorem of Gross--Zagier and Kolyvagin in the form cited in Section~\ref{r2:chapter}.

\clearpage
\section{A conductor census for rational elliptic curves}\label{cens:chapter}
The signed-kernel method gives a uniform finite statement about every rational elliptic curve in an entire conductor range. Because isogenous curves have the same $L$-function at every place, the objects counted here are isogeny classes. The first conductor range contains 690 different primitive automorphic functions, including forms with additive reduction at 2 or 3. Larger conductor ranges also contain exact double and triple central zeros. The Hasse--Weil variable has center $s=1$.

\begin{theorem}\label{cens:thm}
For every elliptic curve $E/\mathbb Q$ with conductor $N\le388$, every zero of its completed Hasse--Weil function with $|\operatorname{Im}s|\le256$ lies on $\operatorname{Re}s=1$ and is simple. Across the 690 isogeny classes there are $300\,605$ positive-ordinate zeros and $254$ central zeros in this window, counting each function separately. Exactly $436$ classes have analytic rank zero and $254$ have analytic rank one.
\end{theorem}
For each class of rational elliptic curves use the normalized weight-two newform furnished by modularity. Conversely every normalized weight-two cuspidal newform of trivial nebentype with rational Hecke eigenvalues defines an elliptic curve over $\mathbb Q$ by the Eichler--Shimura construction, with conductor equal to its primitive level. The theorem therefore covers precisely that rational-coefficient family of cuspidal automorphic functions in this conductor range. In unitary automorphic normalization it places all these zeros on $\operatorname{Re}s=1/2$. A shared ordinate belonging to two distinct members is counted once for each member. The positive and negative ordinates occur in conjugate pairs.

\subsection*{Which curves enter the theorem}
For completeness of the enumeration we use J.~E.~Cremona's tables, derived from his modular-symbol computation. The book explicitly enumerates all modular elliptic curves of conductor below 1000. The corrected electronic tables give minimal equations and all isogeny classes in the stated range. Modularity identifies every rational elliptic curve in that range with one of the enumerated classes. Select the class representative whose Cremona curve number is one. Its optimality is unnecessary, since its $L$-function is an isogeny invariant. The source extractor queries only the conductor, the class label and the curve's Weierstrass coefficients. It checks that one record is returned for every class in the conductor range. Its 690 rows are included in the research package. The tables contain $1612$ rational isomorphism classes through this bound. The number of isogeny classes in each prefix and the zero counts at height 256 are
\[
\begin{array}{r|r|rr|rr}
 N\text{ at most}&\text{classes}&\text{order 0}&\text{order 1}&
   \text{positive zeros}&\text{full multiplicity}\\\hline
 100&93&76&17&35721&71459\\
 200&281&206&75&115629&231333\\
 300&475&315&160&202005&404170\\
 388&690&436&254&300605&601464
\end{array}
\]
For example, the first row covers $306$ curves up to rational isomorphism. Counting one copy per curve there would repeat contributions from isogenous curves.

\subsection*{How the equation determines the analytic inputs}
Suppose the model is $y^2+a_1xy+a_3y=x^3+a_2x^2+a_4x+a_6$ and set
\[
 b_2=a_1^2+4a_2,\qquad b_4=a_1a_3+2a_4,\qquad b_6=a_3^2+4a_6.
\]
For odd $p$, completing the square and summing the quadratic character gives the exact coefficient
\[
     a_p(E)=-\sum_{x\bmod p}\left(\frac{4x^3+b_2x^2+2b_4x+b_6}{p}\right).
\]
For $p=2$ directly test the four pairs $(x,y)$. At a multiplicative prime the singular fibre gives $a_p=1$ or $-1$ and at an additive prime it gives $a_p=0$, as in the local Hasse--Weil definition. The degree-two prime-power recurrence is used when $p\nmid N$ and the degree-one recurrence $a_{p^j}=a_p^j$ otherwise. These exact coefficients determine $f$ to any requested truncation. With $q\in(0,1)$, the tail beyond coefficient $M$ has radius at most $2q^{M+1}((M+1)-Mq)/(1-q)^2$ by the Hasse bound $|a_n|\le2n$.

The sign $w$ at infinity and all bad finite places combined could be read from local root numbers. Here it is determined by a different test, which works uniformly even at $p=2$. Writing $B=2\pi/\sqrt N$ and $f_y=f(iy/\sqrt N)$, the exact Fricke identity is
\begin{equation}\label{cens:fricke}
                      f_{1/y}=w y^2 f_y,\qquad w\in\{1,-1\}.
\end{equation}
This follows because a normalized newform of weight two and trivial character is a Fricke eigenform, and the sign of the completed functional equation includes the weight-two factor $i^2$. Evaluate both sides from the $q$-series at $y=3/2$ and $y=2$ with the preceding tail enclosures. At each pair both values have strict determined signs, so $\operatorname{sgn}(f_{1/y})\operatorname{sgn}(f_y)$ fixes $w$ uniquely. Both pairs give the same sign. The residual interval of \eqref{cens:fricke} also contains zero at each pair. No tabulated rank, sign, $L$-value or zero is used.

\subsection*{Central multiplicity and the closed top edge}
Write $K(u)=e^uf(ie^u/\sqrt N)$. Then $K(-u)=wK(u)$ and
\[
 F(z)=\int_{-\infty}^{\infty} K(u)e^{zu}\,du
     =B^{-(1+z)}\Gamma(1+z)L(E,1+z).
\]
The general rectangle bound from Section~\ref{r2:chapter} applies with $b=B\cos(1)$, $C_0=2/(1-e^{-b})^2$ and $C_1=2/(1-e^{-B})^2$. For spacing $h=1/D$ and width $W$ such that $Be^W\ge2$, it gives the uniform error
\[
 \epsilon=\frac{4C_0e^T}{b^2(e^{2\pi D}-1)}+
              \frac{8C_1}{B^2}\exp(-Be^W/2).
\]
Writing $w_j=hK(jh)$ for $j\ge0$, the signed finite transform is
\[
 A(z)=w_0+\sum_{j=1}^{DW}w_j\{e^{jhz}+w e^{-jhz}\}.
\]
When $w=-1$, the fixed-point value $w_0$ is exactly zero. Each individual weight is computed with its own interval radius. For $w=1$, $A(0)$ enlarged by $\epsilon$ certifies the nonvanishing central value. For $w=-1$, parity forces $F(0)=0$, while $2\sum_{j>0} jhw_j$ enlarged by $2\epsilon$ encloses $F^{(1)}(0)$ by Cauchy's derivative estimate on $|z|=1/2$. Every one of these parity-compatible leading quantities is nonzero. The saved interval in every record has rational endpoints and excludes zero. This proves the orders appearing in the table independently of the contour.

For either sign, conjugation and the functional equation identify the four oriented quarter-contour variations. A constant factor $w$ rotates values without changing their increments of argument. Hence the symmetric count equals $2(\theta_{\rm right}+\theta_{\rm top})/\pi$ as in \eqref{r2:full-count}. Its interval isolates a single integer in every case. At the center axis the real scalar to test is $F(it)$ when $w=1$ and $F(it)/i$ when $w=-1$. This change matters in the finite Fourier transform. With the convention $\widehat v_k=\sum_jv_j e^{-2\pi ijk/J}$, use $v_{-j}=w_j$ and $v_j=w w_j$ for positive $j$ so that $\widehat v_k=A(2\pi i kD/J)$. Upper-edge interval subdivision certifies that there is no zero on $t=T$. The lower edge follows by conjugation. On the right edge the absolute Euler logarithm is bounded by $2\log\zeta(3/2)<\pi$, and outside this edge Euler convergence and reflection exclude other completed zeros at the same heights.

Each sign-changing interval with positive endpoints provides a positive zero, and its reflection provides a negative one. None overlaps the known central zero. For every one of the 690 rows the certified contour integer is exactly the recorded central order plus twice the number of such intervals. Every bracketed zero therefore has multiplicity one, and the list exhausts the rectangle. This proves Theorem~\ref{cens:thm}.

\subsection*{Certifying a double center without a rank table}
In the continuation to larger conductors a positive Fricke sign can accompany an exact double central zero. An interval containing zero cannot certify this equality. We use the following independent arithmetic witness. Find a point
\[
     P=(X/d^2,Y/d^3)\in E(\mathbb Q)
\]
by integral arithmetic in the Weierstrass equation. At a prime $p$ of good reduction coprime to $d$, suppose its reduction has order divisible by $\ell^e$. At a second good prime $q\ne\ell$, suppose $\ell^e\nmid\#E(\mathbb F_q)$. Then $P$ has infinite order. If its order were the finite integer $n$, reduction at $p$ forces $\ell^e\mid n$. The $\ell$-primary subgroup generated by the $\ell$-primary part of $P$ injects at $q$, contradicting the second order test. The witness consists solely of $P,p,q,\ell^e$, the exact order of the first reduction, and the exact group order at $q$. Indeed the first order is checked by the usual finite group law and the second is given by the same character sum as the Fourier coefficients.

By the Gross--Zagier--Kolyvagin rank-zero implication an infinite-order rational point forces $L(E,1)=0$. With even parity there is then a central zero of order at least two. The moment
\[
       2\sum_{j>0} w_j(jh)^2
\]
with extra radius $8\epsilon$ encloses $F^{(2)}(0)$ by Cauchy's inequality. A nonzero interval proves that the multiplicity is exactly two. The complete count in the resulting rectangle is $2+2b$ for $b$ positive brackets, exactly as in the single-curve theorem. For example on 433a the integral point $(0,1)$ reduces to a point of order six at $p=3$, whereas $\#E(\mathbb F_7)=11$. The factor two of the first order is excluded at 7. This method treats algebraic order of points, exact central vanishing and numerical nonvanishing of the second derivative for each double center.

\begin{theorem}\label{cens:thousand}
For all rational elliptic curves of conductor at most $1000$, every nonreal zero of the completed Hasse--Weil function with $|\operatorname{Im}s|\le256$ is simple and lies on $\operatorname{Re}s=1$. The sole real zero allowed in this range is the center. The $2463$ isogeny classes comprise $1321$ nonvanishing centers, $1124$ simple centers and $18$ double centers. There are $1\,163\,668$ positive-ordinate zeros and $2\,328\,496$ zeros with multiplicity in the symmetric windows, counted per isogeny class.
\end{theorem}
Among these $2463$ functions the smallest positive zero belongs to $974\mathrm b$ and has ordinate in $(0.406780166836,0.406780166837)$. Indeed, compare the first brackets in exact rational units of $2\pi$. Only $730\mathrm b,974\mathrm b,978\mathrm c$ have a lower endpoint below the smallest of the upper endpoints. Incomplete-gamma sign bisection within each of these three unique-root intervals isolates the stated minimum. Consequently in the entire closed range $|\operatorname{Im}s|\le0.406780166836$ the only possible zero for any curve of conductor at most $1000$ is its center.

The independently certified height-$64$ subwindow has $221\,277$ positive zeros and a symmetric count of $443\,714$. The central-order and derivative enclosures agree at the two heights. The model list represents $5113$ rational isomorphism classes. A curve of conductor 1000 would require the prime 5 to have conductor exponent three, whereas at primes $p\ge5$ the local conductor exponent of an elliptic curve is at most two. Thus the printed modular-symbol enumeration through conductor $999$, with its corrected electronic data, covers this endpoint. The same contour, central-order test and arithmetic witness above prove the theorem. In every case the count equals the central order plus twice the number of positive sign intervals. The exceptional 18 point witnesses are recorded below. The remaining classes are settled by the zeroth or first central derivative according to parity. With $r=\operatorname{ord}(P\bmod p)$, the column $b$ gives a prime power dividing $r$ which does not divide $\#E(\mathbb F_q)$ and is coprime to $q$.
\[
\begin{array}{c|r|rrrrr}
 \text{class}&(x,y)&p&r&b&q&\#E(\mathbb F_q)\\\hline
\mathrm{389a}&(0,0)&3&3&3&7&13\\
\mathrm{433a}&(0,1)&3&6&2&7&11\\
\mathrm{446d}&(0,2)&3&7&7&5&10\\
\mathrm{563a}&(1,1)&3&5&5&7&13\\
\mathrm{571b}&(0,1)&3&2&2&11&15\\
\mathrm{643a}&(-1,3)&3&6&2&7&11\\
\mathrm{655a}&(-1,5)&3&7&7&11&16\\
\mathrm{664a}&(-1,4)&3&7&7&5&10\\
\mathrm{681c}&(0,1)&5&10&2&7&11\\
\mathrm{707a}&(0,3)&3&6&2&5&9\\
\mathrm{709a}&(0,0)&3&5&5&7&12\\
\mathrm{718b}&(0,0)&3&6&2&5&9\\
\mathrm{794a}&(0,1)&3&2&2&7&11\\
\mathrm{817a}&(0,2)&3&3&3&5&8\\
\mathrm{916c}&(0,1)&3&7&7&5&9\\
\mathrm{944e}&(1,4)&3&7&7&5&9\\
\mathrm{997b}&(-1,0)&3&5&5&11&14\\
\mathrm{997c}&(1,5)&3&5&5&7&12
\end{array}
\]

\subsection*{A second conductor range}
Through height $64$, every one of the $5854$ rational elliptic isogeny classes of conductor at most $2000$ has all nonreal zeros simple and critical. There are $564\,357$ positive zeros. The central orders $0,1,2$ occur on $2896,2861,97$ classes, respectively, giving the symmetric count $1\,131\,769$. Classes of conductors from $1001$ to $2000$ use the full contour anew. When a preselected Fricke test yields zero-containing intervals, choose a different positive rational $y$; the record retains the two nonzero test pairs. This was required on three classes, where $5/4$ replaced $2$.

At height 32, a separate enumeration yields $241\,078$ positive zeros and $485\,211$ total zeros through conductor $2000$. For conductors at most $1000$ its exact lower-height counts can be recovered from the height-64 certificate without another contour. Each positive zero is unique in its bracket, so brackets lying entirely on one side of 32 settle their counts. To split one containing 32, evaluate the real central-line scalar there by the incomplete-gamma expansion. The sign immediately to the right of zero is fixed by the certified central order and derivative, and each preceding bracket reverses it. Therefore the straddling bracket's root is below 32 precisely when its left-side sign and the independently enclosed sign at 32 are opposite. There were 599 such sign determinations, all strict. This proves the subordinate integer counts and excludes a boundary zero. Every one of the 97 double centers has the arithmetic and second-derivative witnesses described above.

\subsection*{An exact triple central zero}
A negative sign only forces one vanishing. If its first derivative interval contains zero, the arithmetic implication for simple central zeros provides a second test. Suppose $P,Q\in E(\mathbb Q)$ and a prime $\ell$ satisfy the following conditions. One good reduction at $q\ne\ell$ has order prime to $\ell$. For each projective direction $[a:b]\in\mathbb P^1(\mathbb F_\ell)$ find a good prime $p$ at which $aP+bQ$ is not an $\ell$-multiple. The first condition implies $E(\mathbb Q)[\ell]=0$. If $uP+vQ=O$ were an integer relation, its coefficients could be divided by any common power of $\ell$. Reducing the resulting pair modulo $\ell$ gives one of the projective directions. A nonzero scalar multiple modulo $\ell$ preserves membership in $\ell E(\mathbb F_p)$, so the relation contradicts its designated reduction test. Hence $P,Q$ are independent.

The Gross--Zagier--Kolyvagin simple-zero theorem implies that $L(E,1)=0$ and $L^{(1)}(E,1)\ne0$ would force $E(\mathbb Q)$ to have rank one. Two independent rational points exclude this possibility. The odd Fricke sign now forces central order at least three. The finite approximation to $F^{(3)}(0)$ is $2\sum_{j>0} w_j(jh)^3$ with radius $48\epsilon$ by Cauchy's inequality. A nonzero enclosure proves order exactly three. The full contour must then account for this central multiplicity along with the positive brackets and their conjugates.

For an explicit case take the curve $5077\mathrm a1$, given by $y^2+y=x^3-7x+6$. Use $P=(0,2)$ and $Q=(-1,3)$. The group $E(\mathbb F_3)$ has order seven, so the rational group has no 2-torsion. At 5 the nonidentity doubles are $(1,0),(1,4),(4,1),(4,3)$. The reductions of $P$ and $P+Q=(2,-1)$ are absent from this list. At 7 the only affine doubles are $(0,2),(0,4)$, which exclude $Q$. These three direction tests certify independence. The two Fricke evaluations give the negative sign and the moment proves
\[
  117.837959237939 < F^{(3)}(0)<117.837959237940.
\]
At $T=2048$ the contour counts $11803$ zeros. The $5900$ positive sign intervals and their conjugates together with the triple zero account for every one, including multiplicities. Thus every nonreal zero of this completed function with $|\operatorname{Im}s|\le2048$ is simple and has real part one, with a triple real zero at $s=1$. For a proposed central order of four or higher this arithmetic argument has a genuine limit. The rank-zero and simple-zero theorems exclude analytic orders zero and one. Excluding order two requires an additional exact vanishing result that the two rational-point tests do not provide.

\subsection*{Every rational isogeny class below conductor ten thousand}
\begin{theorem}\label{cens:9999}
For every rational elliptic curve of conductor at most $9999$, every zero of its completed Hasse--Weil function with $|\operatorname{Im}s|\le64$ lies on $\operatorname{Re}s=1$. All nonreal zeros in this window are simple. Across the $38\,042$ rational isogeny classes, counted separately, there are $4\,256\,834$ positive zeros and $8\,537\,231$ zeros with multiplicity in the full symmetric windows. The unique triple center is $5077\mathrm a$. No higher central order occurs.
\end{theorem}
\begin{proof}
Cremona's complete conductor enumeration gives $38\,042$ isogeny classes represented by $64\,687$ rational isomorphism classes in this range. The model extraction uses only equations, conductors and class labels. The normalized contour certificate for each selected equation is checked as above. A separate implementation of finite group addition checks all $1969$ rational nontorsion witnesses for double centers and the independent-point witness for the single triple center. In each row the contour integer equals its central order plus twice the number of positive sign intervals. Exact summation gives
\[
\begin{array}{r|r|rrrr|r|r}
 N\le&\text{classes}&m=0&m=1&m=2&m=3&\text{positive}&\text{full}\\\hline
 3000&9515&4559&4713&243&0&954177&1913553\\
 5000&17314&7914&8811&589&0&1821604&3653197\\
 7500&27749&12122&14428&1198&1&3029686&6076199\\
 9999&38042&16450&19622&1969&1&4256834&8537231
\end{array}
\]
where $m$ is the central order. The complete ordered roster is checked against the pinned model list. For the $2620$ classes also treated at height $256$, all $238\,319$ corresponding positive brackets in the shorter window intersect the higher-height brackets of the same ordinal, and all central derivative enclosures are compatible.
\end{proof}

For earlier numerical investigations, see S.~Fermigier, \emph{Z\'eros des fonctions $L$ de courbes elliptiques}, Experiment. Math. 1 (1992), 167--173, and S.~Akiyama and Y.~Tanigawa, \emph{Calculation of values of $L$-functions associated to elliptic curves}, Math. Comp. 68 (1999), 1201--1231. The latter paper includes the conductor-$5077$ example among four curves tested to height $400$. S.~J.~Miller's \emph{Investigations of zeros near the central point of elliptic curve $L$-functions}, Experiment. Math. 15 (2006), studies the dependence of low-lying zeros on rank. Rubinstein's larger numerical enumeration is cited below.

In unitary normalization every curve in this conductor-$9999$ census has strictly positive Li coefficients through index $53\,664$ by the holomorphic dyadic bound in Section~\ref{li:section}. This uses its complete height-$64$ rectangle.

The $16\,450$ classes with nonzero central value and the $19\,622$ with a simple central zero have rational Mordell--Weil ranks zero and one, respectively, and finite Tate--Shafarevich groups by the cited rank-zero and rank-one theorems.

\subsection*{A catalogue of multiquadratic base changes}
A multiquadratic field $K/\mathbb Q$ determines a finite vector space $S_K\subset\mathbb Q^\times/\mathbb Q^{\times2}$, consisting of 1 and the squareclasses whose square roots lie in $K$. Conversely $S$ determines $K=\mathbb Q(\sqrt d:d\in S)$, with $[K:\mathbb Q]=|S|$. Write $E^d$ for the quadratic twist of $E$ by the class $d$. Consider all $E/\mathbb Q$ and all nontrivial multiquadratic $K$ for which every $E^d$, $d\in S_K$, has conductor at most $1000$. The condition includes $E$ itself. Their $L$-functions over $K$ satisfy the critical-line assertion for $|\operatorname{Im}s|\le256$. The complete catalogue, with the field and the $K$-isogeny class both retained, is
\[
\begin{array}{c|rrr|r}
[K:\mathbb Q]&\text{fields}&\text{pairs before descent}&
  \text{distinct pairs over }K&\text{largest central order}\\\hline
2&49&3164&1604&2\\
4&21&648&182&4\\
8&1&28&5&2
\end{array}
\]
There are no further degrees under the stated conductor condition. A pair before descent retains a rational isogeny class and the field. The next column identifies rational classes when they become isogenous over the specified field. It gives 1791 pairs across 71 distinct number fields.

For the zero theorem use the regular character of the elementary abelian group $\operatorname{Gal}(K/\mathbb Q)$. Induction of the trivial representation decomposes as the sum of its quadratic characters $\chi_d$, each once. Tensor with the rational Tate module of $E$ and use $V_\ell(E)\otimes\chi_d\simeq V_\ell(E^d)$. Artin formalism, including inertia at the bad primes as in Section~\ref{ebc:chapter}, gives
\[
        \Lambda_K(s,E)=\prod_{d\in S_K}\Lambda_{\mathbb Q}(s,E^d)
\]
with consistent completions, up to an immaterial nonzero normalization constant. This identifies the full zero divisor, including multiplicities, in the certified window. It also explains the orbit count. If two rational classes $E_1,E_2$ are isogenous over $K$, then $\operatorname{Hom}_{G_K}(V_\ell E_1,V_\ell E_2)\ne0$. Frobenius reciprocity identifies this space with the direct sum of the $\operatorname{Hom}_{G_{\mathbb Q}}(V_\ell E_1,V_\ell E_2\otimes\chi_d)$ over $d\in S_K$. A nonzero summand, by the isogeny theorem of Faltings, is equivalent to $E_1$ being rationally isogenous to $E_2^d$. Conversely that condition supplies a $K$-isogeny. Thus the quotient taken in the table has exactly the asserted meaning.

To reproduce the finite enumeration, keep the whole list of $5113$ rational isomorphism classes. For one chosen model $E$ in each isogeny class, search all models $E'$ with $j(E')=j(E)$. If $c_4c_6\ne0$, put $q=(c'_6/c_6)/(c'_4/c_4)$. Equal $j$ gives $c'_4/c_4=q^2$, $c'_6/c_6=q^3$. Since the quadratic twist scales the two invariants by $d^2$ and $d^3$, and a rational change of scale acts by fourth and sixth powers, $E'$ is rationally isomorphic to $E^d$ exactly for the squareclass $d=[q]$. For $j=0$ require $c'_6/c_6=q^3$ with $q\in\mathbb Q$ and take $d=[q]$. For $j=1728$ require $c'_4/c_4=q^2$ with $q\in\mathbb Q$ and use both $[q]$ and $[-q]$. These criteria follow equally from the short Weierstrass models in characteristic zero. Searching all target models accounts for twists that are isogenous to a chosen representative without being isomorphic to it. The resulting finite squareclass set for each $E$ contains precisely the twists that remain in the census.

Starting with $\{1\}$, enlarge each current space $U$ to $U\cup dU$ whenever this union is contained in the available twist set. Every admissible subspace occurs by adding a basis one element at a time. Twisting partitions each set of rational classes into the $K$-isogeny orbits proved above. The exact calculation and every orbit are included in \texttt{twist\_fields.py} and its output. The unique possible degree-eight field is $\mathbb Q(i,\sqrt2,\sqrt3)=\mathbb Q(\zeta_{24})$. One of its five classes comes from the rational classes $36\mathrm a,144\mathrm a,576\mathrm a,576\mathrm e$. Its function is the product of the squares of their four functions. In the window through 256 it has 3500 positive zeros with multiplicity, a double center and 7002 total zeros. Over $\mathbb Q(i,\sqrt{30})$ the rational classes $576\mathrm h$ and $800\mathrm a$ yield a product of their two squares with central order four. The factorization forces these repetitions. The separate interval comparison of Theorem~\ref{csep:joint256} shows that distinct rational constituents have disjoint nonreal zeros throughout this window.

\paragraph{Non-Galois radical fields.}
The extra automorphisms at $j=0,1728$ give another finite family for which the rational constituents can be found exactly. Use $E_c:y^2=x^3+c$ with $m=6$ at $j=0$ and $E_c:y^2=x^3+cx$ with $m=4$ at $j=1728$. Take $n=3,6$ in the first case and $n=4$ in the second. Let $d>0$ have exact order $n$ in $\mathbb Q_{>0}^{\times}/\mathbb Q_{>0}^{\times n}$ and set $K=\mathbb Q(d^{1/n})$ using the positive real root. Then
\begin{equation}\label{cens:radicals}
 \operatorname{Res}_{K/\mathbb Q}(E_c\otimes K)
       \ \sim_{\mathbb Q}\prod_{k=0}^{n-1}E_{c d^{km/n}}.
\end{equation}
In particular the base-change $L$-function is exactly the product of the indicated $n$ rational elliptic functions, with conductor and archimedean normalizations as above.

To prove this put $F=\mathbb Q(\mu_m)$. All geometric endomorphisms of $E_c$ are defined over $F$, an imaginary quadratic field. A positive rational number which is an $n$th power in $F$ is an $n$th power in $\mathbb Q$: the positive real $n$th root differs from any root by an element of $\mu_n\subset F$, and $F\cap\mathbb R=\mathbb Q$. The class of $d$ therefore still has order $n$ over $F$. Kummer theory gives $[F(d^{1/n}):F]=n$, so $[K:\mathbb Q]=n$. The normal closure is $FK$, with dihedral Galois group of order $2n$, since complex conjugation fixes the real generator and inverts the $n$th roots of unity.

The curves in \eqref{cens:radicals} become isomorphic over $K$ by scaling with powers of $d^{1/n}$. Distinct curves in that list are pairwise nonisogenous over $\mathbb Q$. Indeed fix a scaling isomorphism between the $i$th and $j$th curve over $FK$. Under $G_F$ it transforms by the Kummer character of exponent $i-j$, embedded in the automorphisms $\mu_m$. Every geometric homomorphism between these two curves is this isomorphism composed with a CM endomorphism, and the latter is fixed by $G_F$. A nonzero invariant homomorphism would force the character to be trivial. Its order is $n/\gcd(n,i-j)>1$. By the adjunction for restriction of scalars, each of the $n$ pairwise nonisogenous rational elliptic curves thus maps nontrivially into the left side of \eqref{cens:radicals}. Their sum has dimension $n$, which is the full dimension. This proves the isogeny. The Tate-module argument used for the Jacobians gives the identity of local factors at every place. The general theory of twists of restrictions of scalars is treated by Mazur, Rubin and Silverberg, \emph{Twisting commutative algebraic groups}. A related cyclic cubic formula over a base containing $\mu_3$ is given in \emph{Restriction of scalars and cubic twists of elliptic curves}, J. Korean Math. Soc. 58 (2021), 123--132.

Requiring all the rational constituents in \eqref{cens:radicals} to have conductor at most $1000$ gives the following catalogue for these positive-radical constructions. The representatives $d$ in each row give distinct fields.
\[
\begin{array}{c|c|c|l}
n&\text{fields}&\text{field/isogeny pairs}&d\\\hline
3&6&17&2,3,5,6,10,12\\
4&5&7&2,3,5,12,18\\
6&3&3&5,80,108
\end{array}
\]
To check completeness within the construction, compare all same-$j$ model coefficients $c,c'$. Keep positive ratios $c'/c=q^{m/n}$ with $q>0$ rational, and reduce the prime valuations of $q$ modulo $n$. Keep a generated cyclic subgroup precisely when its full orbit consists of available models. Two generators of the same subgroup define the same pure field. Conversely, equal normal closures over $F$ give the same Kummer subgroup, and the injection of positive rational power classes into the corresponding classes over $F$ proved above gives equality already over $\mathbb Q$. The distinct $K$-isogeny classes are exactly the displayed twist orbits by adjunction and \eqref{cens:radicals}. For $K=\mathbb Q(5^{1/6})$, the base extension of $27\mathrm a1$ has constituents $27\mathrm a,225\mathrm a,225\mathrm b,675\mathrm a,675\mathrm c,675\mathrm e$. It has 2696 positive zeros with multiplicity and a double center through height 256. None of these 14 radical fields is Galois over $\mathbb Q$. Their normal closures furnish genuine nonabelian Galois examples. Since $F$ and $K$ are linearly disjoint, the permutation representation for $FK$ is that for $K$ tensored with $1\oplus\chi_F$. Every elliptic constituent in \eqref{cens:radicals} has CM by $F$, hence is a quadratic self-twist by $\chi_F$ (its Tate module is induced from $G_F$). Consequently $L(E_c/FK,s)=L(E_c/K,s)^2$, with the multiplicities doubled. This also gives examples with dihedral Galois groups of orders 6, 8 and 12 without a new analytic input.

For a degree-16 example take $L=\mathbb Q(3^{1/4},i,\sqrt2)$ and the curve $64\mathrm a1$. The splitting field of $x^4-3$ has dihedral group of order eight and quadratic subfields $\mathbb Q(i),\mathbb Q(\sqrt3),\mathbb Q(\sqrt{-3})$. Thus $\sqrt2$ enlarges its degree by two and $\operatorname{Gal}(L/\mathbb Q)$ is the direct product of the dihedral group and a group of order two. The pure quartic factors are $64\mathrm a,288\mathrm a,288\mathrm e,576\mathrm h$. Passage to the normal closure squares their product. Twisting by 2 supplies the additional factors $32\mathrm a,576\mathrm f,576\mathrm g,288\mathrm d$, again each squared. The direct product of the field groups justifies tensoring the permutation representation by $1\oplus\chi_2$. The resulting degree-32 Hasse--Weil function over $L$ has $7092$ positive zeros with multiplicity, central order four, and $14188$ total zeros throughout $|\operatorname{Im}s|\le256$. All lie on the center line.


Taken together, these constructions give $1846$ different field/isogeny pairs over $100$ different fields. There are $1791$ multiquadratic pairs, $27$ pairs over pure radical fields, $27$ over their dihedral normal closures, and the degree-$16$ example. Distinct pure fields of the same listed degree have distinct normal closures by the Kummer-subgroup criterion. The three possible nonabelian normal-closure orders differ, and distinguish those fields from the multiquadratic fields and from the nonnormal pure fields. For a fixed radical field, distinct twist orbits still give distinct isogeny classes over its normal closure. Indeed, the restriction of scalars to $\mathbb Q$ there has the original orbit constituents, each twice, so a shared isogeny class after base change would put the two rational curves in the same orbit by adjunction. Across the $1846$ functions the positive-with-multiplicity count is $1\,952\,547$ and the two-sided count is $3\,906\,569$ through height 256. The central orders $0,1,2,3,4$ occur with frequencies $683,880,266,5,12$. The rational-factor exponents are all one on $1731$ products and all two on $115$ products.

\paragraph{Data and prior results.}
M.~O.~Rubinstein, \emph{Computational methods and experiments in analytic number theory}, arXiv:math/0412181, p.~43, reported the first 1000 zeros for elliptic curves of conductor below 1000 and the first 100 below conductor 8000. He explains which truncation and roundoff bounds were omitted in those computations. The present finite statements come with new reproducible interval evaluations and exact counts in closed rectangles. We claim no priority for the numerical observation of agreement with the critical-line conjecture in these ranges.
For the equations and enumeration use J.~E.~Cremona, \emph{Algorithms for Modular Elliptic Curves}, second ed., Cambridge University Press, 1997, and the corrected electronic \emph{Elliptic Curve Data}. For modularity and the weight-two correspondence use the modularity theorem and the Eichler--Shimura construction. The accompanying program \texttt{checks/elliptic-census/certify.py} can recompute the theorem directly from the supplied equations. Its source ledger gives the archived database checksum, extraction query and precisions. The per-class output records the two-sided count, both central derivative bounds, the subdivision and Fourier parameters, and every positive bracket. At height 64 an independent upper-incomplete-gamma evaluation checked the first positive bracket of every one of the 2463 classes, for a total of 4926 enclosed endpoint signs. The theorem concerns bounded ordinates and supplies no uniform limit in height.

\clearpage
\section{Elliptic functions over positive multiradical fields}\label{cmr:chapter}
For the non-CM example of Section~\ref{ebc:kummer}, cubic-radical extensions introduce new degree-four cusp forms. For curves with extra automorphisms the same base fields admit a complete descent to rational elliptic factors. The result also covers fourth and sixth roots and allows several independent radicands. It turns the elliptic critical-line hypothesis over $\mathbb Q$ for the specified CM twists into its counterpart over infinitely many nonnormal fields and their nonabelian normal closures.

Let $m=6$ and $E_c:y^2=x^3+c$ for $j=0$, or $m=4$ and $E_c:y^2=x^3+cx$ for $j=1728$, with $c\in\mathbb Q^\times$. Let $S$ be any finite subgroup of
\[
   \mathbb Q_{>0}^\times/(\mathbb Q_{>0}^\times)^m .
\]
For $d\in S$ use any positive representative, and let $K=\mathbb Q(d^{1/m}:d\in S)$, choosing every positive root. Both the field and the rational isomorphism class of $E_{cd}$ are independent of these choices. Put $F=\mathbb Q(\mu_m)$ and $L=FK$.

\begin{theorem}\label{cmr:tower}
One has $[K:\mathbb Q]=|S|$ and
\[
 \operatorname{Res}_{K/\mathbb Q}(E_c\otimes K)
       \ \sim_{\mathbb Q}\prod_{d\in S} E_{cd}.
\]
The factors on the right are pairwise nonisogenous over $\mathbb Q$. If $S$ has exponent greater than two, $K$ is nonnormal and $L$ is its normal closure. The Galois group of $L/\mathbb Q$ is $\operatorname{Hom}(S,\mu_m)\rtimes C_2$, where the involution acts by inversion. Moreover
\begin{align*}
 L(E_c/K,s)&=\prod_{d\in S}L(E_{cd}/\mathbb Q,s),\\
 L(E_c/L,s)&=\left(\prod_{d\in S}L(E_{cd}/\mathbb Q,s)\right)^2.
\end{align*}
The identities include every ramified finite Euler factor. They hold for compatible completions up to a nonzero normalization constant. In particular the full critical-line assertion for either field is equivalent to that assertion for each listed rational elliptic function.
\end{theorem}
\begin{proof}
If a positive rational $a$ has an $m$th root in $F$, multiplying that root by a suitable element of $\mu_m$ gives its positive real root. Since $F$ is imaginary quadratic, this root is rational. Consequently positive rational power classes inject into $F^\times/F^{\times m}$. The Kummer pairing identifies $\operatorname{Gal}(L/F)$ with $A=\operatorname{Hom}(S,\mu_m)$. As $K$ is a subfield of $\mathbb R$, $K\cap F=\mathbb Q$ and $[L:K]=2$, so $[K:\mathbb Q]=|S|$. Complex conjugation fixes all the chosen roots and inverts the roots of unity. This proves the semidirect product description. The conjugates of these positive roots contain their multiples by all the values of the Kummer characters. When the exponent exceeds two, these values generate $F$ over $\mathbb Q$: their orders include 3 in the sixth-root case or 4 in the fourth-root case. Thus the normal closure is $L$ and $K$ is nonnormal.

A change of Weierstrass scale with multiplier $d^{1/m}$ makes $E_{cd}$ isomorphic to $E_c$ over $K$. For distinct $d,e\in S$, the comparison isomorphism of $E_{cd}$ and $E_{ce}$ transforms under $G_F$ by the Kummer character paired with $d/e$. This is a nontrivial character by the perfect pairing; its inverse would serve equally upon changing the direction of the isomorphism. Every geometric endomorphism of $E_c$ is defined over $F$. Hence every geometric homomorphism between the two twists is the comparison isomorphism composed with an $F$-fixed CM endomorphism, and no nonzero such homomorphism is fixed by $G_F$. In particular the curves are pairwise nonisogenous over $\mathbb Q$.

The adjunction defining restriction of scalars gives a nonzero rational homomorphism from each $E_{cd}$ into $\operatorname{Res}_{K/\mathbb Q}E_c$. Pairwise nonisogenous elliptic curves are distinct simple factors up to isogeny, so their sum has dimension $|S|$, the entire dimension of the restriction of scalars. This proves the first isogeny. Taking rational Tate modules and inertia invariants gives the identity of the finite Euler factors at every prime. Induction of the archimedean parameters and the conductor formula give the statement on completions, as in Section~\ref{ebc:chapter}.

Finally $K$ and $F$ are linearly disjoint, so the permutation representation for $FK$ is the tensor product of the permutation representations for $K$ and $F$. Each rational CM factor is a self-twist by the quadratic character of $F$, since its Tate module is induced from $G_F$. Tensoring the first product by $1\oplus\chi_F$ therefore repeats every factor twice. The functions are entire by modularity, so the divisor of either product is the sum of its factor divisors with the indicated positive multiplicities. This proves both directions of the last assertion.
\end{proof}

For example take distinct primes $p_1,\ldots,p_r$ and the subgroup generated by their classes modulo $m$th powers. Then $K=\mathbb Q(p_1^{1/m},\ldots,p_r^{1/m})$ and $[K:\mathbb Q]=m^r$. The rational functions in the theorem are indexed by all $m^r$ exponent vectors. Thus the transfer allows fields of unbounded degree with only degree-two rational constituents. It does not furnish the unbounded-height critical-line assertion for those constituents.

\subsection*{Enumerating the finite catalogues}
The construction can be classified completely under a bound $X$ on every rational constituent conductor. For each input rational isogeny class, use \emph{all} its Weierstrass models up to rational isomorphism in the complete list, and retain the two indicated $j$-invariants. The short-model parameter $c$ is $-c_6/864$ when $j=0$, and $-c_4/48$ when $j=1728$. Its class in $\mathbb Q^\times/\mathbb Q^{\times m}$ is independent of the rational short-model scaling. For each fixed starting model the allowable positive power classes are precisely the ratios $c'/c>0$ among the same-$j$ models in the list. Reduce every prime exponent modulo $m$. Enumerate all subgroups of these allowable classes by adjoining one element $d$ at a time, testing every element of $U\langle d\rangle$ before retaining it. Any admissible group is reached by adjoining its generators.

We restrict to subgroups of exponent exceeding two in the counts below. The square-root fields in the complementary case are already covered by the multiquadratic calculation. Two positive fields in the present list with the same $m$ are isomorphic precisely when their subgroups agree. Indeed an isomorphism of number fields extends to an automorphism of $\overline{\mathbb Q}$, so their normal closures agree. Kummer theory recovers the subgroup over $F$, and the injection of positive rational power classes recovers $S$. The converse is immediate. A field from the fourth-root family has normal closure containing an element of order four in its Galois group. A field from the sixth-root family has only elements of orders dividing six, or involutions, in that group. This separates the two families. The rational twist orbit of the starting model determines its $K$-isogeny class, and also its $L$-isogeny class, by adjunction and Theorem~\ref{cmr:tower}. Thus the resulting orbit counts count actual pairs of a field and an isogeny class there.

In the table, the column labelled $S$ gives its invariant cyclic factors. The columns count nonnormal fields $K$ and pairs over them before taking their normal closures; the latter give the same numbers again.
\[
\begin{array}{c|r|rr|rr}
 S&[K:\mathbb Q]&\multicolumn{2}{c|}{X=1000}&\multicolumn{2}{c}{X=9999}\\
 &&\text{fields}&\text{pairs}&\text{fields}&\text{pairs}\\\hline
 C_3&3&6&17&22&117\\
 C_4&4&5&7&36&97\\
 C_6&6&3&3&35&76\\
 C_2\times C_4&8&1&1&14&23\\
 C_3\times C_3&9&1&1&4&8\\
 C_2\times C_6&12&0&0&3&3\\
 C_4\times C_4&16&0&0&3&3\\
 C_3\times C_6&18&0&0&4&4\\\hline
 \text{total}&&16&29&121&331
\end{array}
\]
For $X=1000$ the two noncyclic entries add three new pairs to the $1846$ already described in Section~\ref{cens:chapter}, because the normal closure of the quartic--quadratic compositum is precisely its earlier degree-16 example. In all, the two constructions yield $1849$ pairs over $103$ fields through height 256. Their $1\,968\,171$ positive zeros counted with multiplicity correspond to $3\,937\,828$ total zeros across the symmetric windows. The degree-nine field has a short description. For $K=\mathbb Q(\sqrt[3]2,\sqrt[3]3)$ the curve $27\mathrm a1$ has nine rational constituents, from labels $27\mathrm a,36\mathrm a,108\mathrm a,243\mathrm a,243\mathrm b,972\mathrm a,972\mathrm b,972\mathrm c,972\mathrm d$. Its central order is three and its positive count is $4026$. Its normal closure $K(\sqrt{-3})$ has order-$18$ generalized dihedral group, central order six and positive count $8052$.

\begin{theorem}\label{cmr:finite}
For the $331$ positive-radical pairs in the $X=9999$ column, and for all their normal closures, every nontrivial zero through $|\operatorname{Im}s|\le256$ lies on $\operatorname{Re}s=1$. On these $662$ functions there are $2\,446\,389$ positive and $4\,894\,917$ total zeros with multiplicity. Together with the $1791$ multiquadratic pairs for $X=1000$, this gives $2453$ distinct pairs over $313$ fields, with $4\,262\,133$ positive and $8\,527\,783$ total zeros.
\end{theorem}
\begin{proof}
The factor orbits involve $155$ different rational elliptic functions. The conductor-$1000$ theorem covers $37$ of them and an additional full symmetric-contour calculation covers the remaining $118$. Arithmetic point witnesses certify the two double centers in the supplementary list. The resulting weighted sums give the stated counts by Theorem~\ref{cmr:tower}. As an independent finite-place check, for each orbit choose an independent basis of $S$ with cyclic orders $e_j$. Write the corresponding $m$th-root generators as positive $e_j$th roots of rational integers $r_j$. Their independent binomial equations have product of degrees $[K:\mathbb Q]$. At a prime $q$ away from the binomial discriminants and all conductors, count the roots of each binomial in $\mathbb F_{q^k}$ and multiply to count all embeddings of the residue algebra. If there are $c_f$ primes of degree $f$, the result equals $\sum_{f\mid k}f c_f$, recovering the residue degrees by induction on $k$. From these degrees and the ordinary elliptic Frobenius traces, $2317$ complete Euler polynomials were compared with the proposed rational products at two completely split and five other good primes for each of the $331$ orbits.
\end{proof}

Theorem~\ref{csep:joint256} gives exact nonreal multiplicities throughout this catalogue. In $2035$ of the $2453$ factorizations every primitive exponent is one. In the other $418$ every exponent is two. The primitive ordinates are pairwise disjoint, so every nonreal zero through $256$ has exactly the respective multiplicity one or two.

For a degree-$36$ example put $K=\mathbb Q(189^{1/6},2^{1/3})$, $L=K(\sqrt{-3})$ and $E=27\mathrm a1$, with equation $y^2+y=x^3-7$. The $18$ curves contributing over $K$ are
\[
\begin{gathered}
27\mathrm a,36\mathrm a,108\mathrm a,441\mathrm a,441\mathrm b,
1323\mathrm a,1323\mathrm b,1323\mathrm m,\quad
1764\mathrm a,1764\mathrm b,1764\mathrm c,\\
5292\mathrm a,5292\mathrm b,5292\mathrm c,5292\mathrm h,5292\mathrm i,5292\mathrm j,5292\mathrm k.
\end{gathered}
\]
The central order over $K$ is nine, there are $9178$ positive zeros with multiplicity, and the symmetric count is $18365$. The normal closure $L$ has central order eighteen, $18356$ positive zeros and total count $36730$. These central orders are exact and do not require a Birch--Swinnerton-Dyer assumption over $K$ or $L$.

\begin{corollary}\label{cmr:rank}
For the preceding $E,K,L$ one has $\operatorname{rank}E(K)=9$, $\operatorname{rank}E(L)=18$, and the Tate--Shafarevich groups over both fields are finite. More generally, for $2446$ of the $2453$ cases of Theorem~\ref{cmr:finite}, the stated central order equals the Mordell--Weil rank and the Tate--Shafarevich group is finite.
\end{corollary}
\begin{proof}
On the displayed orbit, nine rational factors have nonzero central value and nine have a simple central zero. The latter are $441\mathrm b$, $1323\mathrm a$, $1323\mathrm m$, $1764\mathrm b$, $1764\mathrm c$, $5292\mathrm b$, $5292\mathrm c$, $5292\mathrm h$, $5292\mathrm i$. For each individual factor, Gross--Zagier and Kolyvagin give equality of analytic and algebraic ranks and finiteness of its Tate--Shafarevich group. Weil restriction identifies the rational points of the restricted abelian variety with the points over the original field. The isogenies in Theorem~\ref{cmr:tower} add ranks with multiplicity and preserve finiteness of Tate--Shafarevich groups. Shapiro's lemma, applied globally and at every completion, identifies the Tate--Shafarevich group for the restriction of scalars with that over the original field. In the full catalogue exactly $2446$ entries have all rational factor orders at most one, so the same proof applies to each. This argument makes no leading-coefficient claim.
\end{proof}

\subsection*{A larger catalogue at height sixty-four}
Taking every rational elliptic constituent of conductor at most $9999$ in the multiquadratic classification of Section~\ref{cens:chapter} gives the following counts. The pairs are taken up to isogeny over the indicated embedded field.
\[
\begin{array}{r|r|r|r}
 [F:\mathbb Q]&\text{fields}&(F,E)\text{ pairs}&\max\operatorname{ord}_{s=1}\\\hline
 2&340&49755&4\\
 4&320&11409&6\\
 8&20&528&7
\end{array}
\]
The quadratic-twist pairings also passed $298\,530$ separate good-prime trace comparisons, three split and three inert primes for each of the $49\,755$ pairs. Combining these $61\,692$ multiquadratic entries with the $662$ nonquadratic radical and normal-closure entries of Theorem~\ref{cmr:finite} gives $62\,354$ distinct pairs over $922$ fields. Theorem~\ref{cens:9999} and the full local factorizations give $16\,413\,564$ positive zeros and $32\,903\,342$ symmetric-window zeros with multiplicity through height $64$, all on the critical line. Exactly $61\,390$ factorizations have primitive exponents one throughout and $964$ have primitive exponents two throughout. Theorem~\ref{csep:joint9999} separates every pair of distinct primitive positive ordinates in these products. Hence in the first $61\,390$ cases every nonreal zero is simple through $64$, and in the other $964$ every nonreal zero has multiplicity exactly two.

In $60\,225$ of these entries every rational constituent has analytic order zero or one. Applying the proof of Corollary~\ref{cmr:rank} gives equality of the algebraic rank with the verified central order and finiteness of the Tate--Shafarevich group in all these cases. The maximal rank in this proved subset is $18$. Every number here comes from summing the listed individual primitive orders and zero counts with their isogeny multiplicities.

\begin{corollary}\label{cmr:li}
All $62\,354$ field functions in the height-$64$ catalogue have strictly positive real Li coefficients at every index $1\le n\le53\,664$, using unitary normalization. For the $1849$ entries already certified through $256$ with all rational constituents of conductor at most $1000$, this holds through $n=1\,150\,720$.
\end{corollary}
\begin{proof}
Each rational constituent of the first catalogue has weight two, conductor at most $9999$ and a complete critical-line window through $64$. The holomorphic dyadic coefficient bound of Section~\ref{li:section} thus gives the first $53\,664$ strictly positive coefficients on every constituent. At height $256$ with conductor bound $1000$, the corresponding value is $1\,150\,720$. The field zero divisors are positive integral sums of the rational divisors and their real Li coefficients add with those same positive multiplicities.
\end{proof}

\subsection*{Standard CM functions of unbounded degree}
A different cyclotomic curve supplies higher-degree cuspidal constituents. Let $p$ be any odd prime, $F_p=\mathbb Q(\zeta_p)$ and
\[
 C_d:\quad y^2=d x^p-1,\qquad A_d=\operatorname{Jac}(C_d).
\]
Use the smooth projective curve, so $\dim A_d=(p-1)/2$. A change of rational $p$th-power representative of $d$ preserves the rational isomorphism class. For any finite subgroup $S\subset\mathbb Q_{>0}^\times/(\mathbb Q_{>0}^\times)^p$ put
\[
   K=\mathbb Q(d^{1/p}:d\in S),\qquad L=KF_p.
\]

\begin{theorem}\label{cmr:prime}
The standard degree-$(p-1)$ functions of the $A_d$, $d\in S$, arise from pairwise inequivalent cuspidal automorphic inductions of algebraic Hecke characters of $F_p$, with the arithmetic variable for pure cohomology of weight one (or equivalently the corresponding unitary normalization after shifting the center to $1/2$). At all finite places one has
\[
 L(A_1/K,s)=\prod_{d\in S}L(A_d/\mathbb Q,s),\qquad
 L(A_1/L,s)=\left(\prod_{d\in S}L(A_d/\mathbb Q,s)\right)^{p-1}.
\]
The complete zero divisors obey the same identities, with central line $\operatorname{Re}s=1$. Thus critical-line assertions for every rational constituent transfer to fields with Galois closure group $(C_p)^r\rtimes C_{p-1}$, where $r=\dim_{\mathbb F_p}S$ is unrestricted. Conversely the assertion over either of these base fields implies it for each rational constituent.
\end{theorem}
\begin{proof}
If a rational $d$ becomes a $p$th power in $F_p$, its norm $d^{p-1}$ is a rational $p$th power; coprimality shows that $d$ was already a $p$th power. Thus the Kummer classes inject, and $[L:F_p]=p^r$. Adjoining the $r$ real $p$th roots to $\mathbb Q$ has degree at most $p^r$ and at least $[L:F_p]$. It follows that $[K:\mathbb Q]=p^r$, $K$ and $F_p$ are linearly disjoint, and $\operatorname{Gal}(L/K)$ acts on $\operatorname{Gal}(L/F_p)$ by cyclotomic scalar multiplication.

Fix an auxiliary prime $v\ne p$ and extend $v$-adic coefficient scalars to $\overline{\mathbb Q}_v$. Write $V$ for the first cohomology of $A_1$. The automorphism $x\mapsto\zeta_p x$ has every nontrivial $p$th-root eigenvalue exactly once on $V$. Indeed on holomorphic differentials $x^{j-1}dx/y$, for $1\le j\le (p-1)/2$, the eigenvalues are $\zeta_p^j$ and the conjugate differentials give the remaining eigenvalues. Comparison with singular cohomology gives the same multiplicities on $V$. The group $G_{F_p}$ preserves the eigenlines and $G_{\mathbb Q}$ permutes them transitively. Consequently $V=\operatorname{Ind}_{G_{F_p}}^{G_{\mathbb Q}}\psi$ for the character on any one chosen eigenline; choose the one with eigenvalue $\zeta_p$. The CM theorem identifies it with the $v$-adic character of an algebraic Hecke character, equivalently the Jacobi-sum character in this cyclotomic family.

For $d\in S$ the coordinate isomorphism $C_d\to C_1$ over $K$ is $x\mapsto d^{1/p}x$. Accordingly its comparison on the chosen line twists $\psi$ by the Kummer character paired with $d$, up to a uniform inversion depending on the cohomology convention. On $G_{F_p}$ the permutation representation $P_K$ on the embeddings of $K$ is the regular representation of $\operatorname{Gal}(L/F_p)$. It is the direct sum of these Kummer characters as $d$ runs through $S$, and inversion permutes that same set. The tensor identity for induction now gives
\[
 V\otimes P_K
   =\operatorname{Ind}_{G_{F_p}}^{G_{\mathbb Q}}
        \bigl(\psi\otimes P_K|_{G_{F_p}}\bigr)
   \simeq \bigoplus_{d\in S} H^1(A_{d,\overline{\mathbb Q}},\overline{\mathbb Q}_v).
\]
Every such induced representation is unchanged by twisting with any character of $\operatorname{Gal}(F_p/\mathbb Q)$, because those characters restrict trivially to the inducing group. Linear disjointness gives $P_L=P_K\otimes P_{F_p}$, and $P_{F_p}$ is the sum of its $p-1$ characters. This proves the product and power, including local factors at bad primes. The same induction at the archimedean place and the conductor formula give the completed assertions.

For cuspidality, the CM type of $C_1$ is $\Phi=\{1,\ldots,(p-1)/2\}\subset (\mathbb Z/p\mathbb Z)^\times$. Its stabilizer under multiplication is trivial: the sum of its elements is $(p^2-1)/8$, nonzero modulo $p$. Any multiplier preserving $\Phi$ must preserve that nonzero sum. The conjugates of the Hecke character thus have distinct infinity types, a property unchanged by finite-order Kummer twisting. The cuspidality criterion for cyclic automorphic induction, applied through the cyclic tower $F_p/\mathbb Q$, gives the claimed cuspidal representations. Two different $d$'s cannot produce equivalent inductions: equality would conjugate one inducing character to the other, equality of infinity types forces the conjugation to be trivial, and then the Kummer pairing forces the two power classes to agree. The positive factorization of entire completions gives the two implications about zero location.
\end{proof}

For example $p=5$ and $K=\mathbb Q(2^{1/5},3^{1/5})$ give $25$ distinct rational cuspidal standard functions of degree four. The Galois closure has degree $100$ and the exponent of each constituent there is four. This example uses the CM statement for the genus-two Jacobians and makes no appeal to a checked zero window for them. For $p=3$ the rational change $X=dx,Y=dy$ gives $Y^2=X^3-d^2$, recovering the cubic elliptic twists. For the cyclotomic CM type and Jacobi-sum characters see M.~Emory and H.~Goodson, \emph{Sato--Tate distributions of $y^2=x^p-1$ and $y^2=x^{2p}-1$}, J. Algebra 597 (2022), 241--265, \S\S2,5.1. The analytic induction theorem of Arthur and Clozel is the automorphic input.

\begin{corollary}\label{cmr:subfields}
In either positive-radical family, let $A$ be the Kummer group over the CM field and $G=\operatorname{Gal}(L/\mathbb Q)$. For \emph{every} intermediate field $M=L^H$, put $B_H=A\cap H$, $e=[G:AH]$, and let $S_H\subset S$ be the annihilator of $B_H$ under the Kummer pairing. The function over $M$ is the product of the rational twists indexed by $S_H$, each with exponent $e$. This applies both to $A_1$ in Theorem~\ref{cmr:prime} and to $E_c$ in Theorem~\ref{cmr:tower}.
\end{corollary}
\begin{proof}
Conjugation on the abelian group $A$ is by scalar multiplication, either by inversion or by an element of $(\mathbb Z/p\mathbb Z)^\times$. Hence every subgroup of $A$ is stable under $G$. An $A$-orbit on $G/H$ is thus always a copy of $A/B_H$. There are $[G:AH]$ such orbits. The permutation representation restricted to $A$ consequently contains each character in $S_H$ exactly $e$ times. Tensoring this restriction with the chosen CM eigenline and inducing gives the formula, exactly as in the proof of Theorem~\ref{cmr:prime}. For the elliptic family the two eigenlines of $\mu_m$ have primitive exponents $1,-1$, and are interchanged by complex conjugation. Its representation is therefore induced from either one, giving precisely the same tensor argument there. In particular the rank equality and finiteness in Corollary~\ref{cmr:rank} extend to \emph{every} intermediate field of its degree-$36$ example, since only a subset of the same rank-zero and rank-one rational factors occurs.
\end{proof}

\subsection*{A certified primitive genus-two rectangle}
The Jacobian $J$ of $y^2+y=x^5$ is geometrically simple, with conductor $3125$ and bad Euler factor $1$ at $5$. These local data appear in the genus-two database at class $3125.\mathrm a$, and the curve is $3125.\mathrm a.3125.1$. The substitution $Y=2y+1$ gives $Y^2=4x^5+1$. Thus $J$ is a quadratic twist by $-1$ of the $p=5,d=4$ case of Theorem~\ref{cmr:prime}, after replacing $x$ by $-x$. Its degree-four standard function has an unconditional continuation and functional equation. This example has a geometrically simple abelian surface rather than a product of elliptic factors.

The signed rectangle calculation through height $8$ gives $14$ zeros, all nonreal, simple and on $\operatorname{Re}s=1$. To prove this set $B=(2\pi)^2/\sqrt{3125}$, $c=2\sqrt B$ and
\[
 \Lambda(s)=B^{-s}\Gamma(s)^2 L(J,s),\qquad
 G(u)=2e^u\sum_{n\ge1}a_n K_0(c\sqrt n e^{u/2}).
\]
The modified Bessel function $K_0$ gives $\Lambda(1+z)=\int_{\mathbb R}G(u)e^{zu}\,du$. Indeed the Mellin transform of $2K_0(2\sqrt x)$ is $\Gamma(s)^2$. Absolute convergence first gives this identity on the right of the critical strip. Inverse Mellin transformation and the completed functional equation give $G(-u)=wG(u)$ and the bilateral continuation. One may shift the contour by the vertical gamma decay and the usual strip growth of the entire CM function.

At odd good $p$, let $t_1=p+1-\#C(\mathbb F_p)$ and $t_2=p^2+1-\#C(\mathbb F_{p^2})$. The polynomial is
\[
   1-t_1X+\tfrac12(t_1^2-t_2)X^2-pt_1X^3+p^2X^4 .
\]
The program counts the two traces directly by the quadratic character of $1+4x^5$ over the corresponding fields. Over $\mathbb F_{p^2}$ this character equals the ordinary Legendre symbol of the norm. For $p\equiv2,3\bmod5$, taking fifth powers permutes both finite fields, so both traces are zero. The original equation also gives $\#C(\mathbb F_q)=q+1$ at $q=2,4$, by summing over $y$ and inverting the fifth power map. The finite prime factor at $5$ is one. These polynomials and multiplicativity compute all needed coefficients. Purity gives $|a_n|\le d_4(n)\sqrt n\le n^4$.

Here are explicit analytic errors for the finite transform. Use mesh $h=1/D$, log cutoff $W$, and in the sample at $u=jh\ge0$ retain $n\le\lceil(V/(ce^{u/2}))^2\rceil$. Put $S_0=4\cdot2^{10}9!$. For $|\operatorname{Re}z|\le2$, $|\operatorname{Im}z|\le T$ a uniform bound for its error is
\[
 \eta=\frac{10^{15}e^T}{e^{2\pi D}-1}
 +\frac{8h S_0 e^{3W-e^{W/2}/2}}{e^h-1}
 +\frac{16h c^{-10}\Gamma(10,V)}{1-e^{-2h}} .
\]
For the first term work on the two parallel lines $|\operatorname{Im}u|=1$. When $u=x+iy$, $x\ge0$, the Bessel argument has real part at least $c\cos(1/2)\sqrt n e^{x/2}>\sqrt n e^{x/2}$. The integral $K_0(q)=\int_0^\infty e^{-q\cosh v}dv$ and $\cosh v\ge1+v^2/2$ give $|K_0(q)|\le e^{-\operatorname{Re}q}\sqrt{\pi/(2\operatorname{Re}q)}\le2e^{-\operatorname{Re}q}$ here. Since $\sqrt n e^{x/2}\ge(\sqrt n+e^{x/2})/2$, the integrand is bounded by $4 S_0 e^T e^{3x-e^{x/2}/2}$, using $\sum n^4e^{-\sqrt n/2}\le2\int_0^\infty t^4e^{-\sqrt t/2}\,dt=S_0$. This last comparison follows by integrating from $n$ to $n+1$ and using $e^{1/2}<2$. The integral in $x$ is at most $2\cdot2^6 5!$. Reflection bounds the negative half-line, so the sum of the boundary integrals is below $10^{15}e^T$. The analytic-strip trapezoidal rule yields the first term. On the real axis the factor $e^T$ disappears, and the logarithmic derivative of $e^{3x-e^{x/2}/2}$ is at most $-1$ beyond the stated $W$. Summing the two omitted geometric tails gives the second term. Finally, for $b=ce^{u/2}$ and $b\sqrt M\ge V\ge8$,
\[
 \sum_{n>M}n^4e^{-b\sqrt n}\le\int_M^\infty t^4e^{-b\sqrt t}\,dt
       \le2 b^{-10}\Gamma(10,V).
\]
Multiplying by the sample weights, bounding the two $e^{\pm zu}$ factors, and summing $e^{-2jh}$ proves the last term.

The same local coefficient estimate gives strict signs and the equality interval at both tests $u=\pm\log(3/2)$ and $u=\pm\log2$, proving $w=+1$. For real positive Bessel arguments beyond the working threshold the program uses the expansion
\begin{align*}
 K_0(x)&=\sqrt{\frac{\pi}{2x}}e^{-x}
       \left(\sum_{j=0}^{k} b_jx^{-j}+R_k(x)\right),\\
 b_0&=1,\quad b_j=-\frac{(2j-1)^2}{8j}b_{j-1},
 \qquad |R_k(x)|\le |b_{k+1}|x^{-k-1}.
\end{align*}
To prove the remainder, replace $\cosh v$ in the integral representation by $1+t/x$. After the change of variables, the normalized integral is the average of $(1+t/(2x))^{-1/2}$ against $e^{-t}t^{-1/2}dt/\sqrt\pi$. Taylor's theorem on $t\ge0$ bounds the absolute remainder by its first omitted term. Integrating the powers of $t$ gives the displayed recurrence. Smaller arguments use their interval Bessel evaluation directly.

For the final winding calculation use $s=3+it$ on the right edge. Its continuous gamma contribution is $\operatorname{Im}(2\log\Gamma(s)-s\log B)$, and the Euler-product argument is in $(-\pi,\pi)$ by $4\log\zeta(5/2)<\pi$. The opposite side follows by reflection. On the top edge write $z=1+iT+x$, $-1\le x\le1$, and expand the finite weight polynomial in $x$. With mass $M_0=|w_0|+2\sum|w_j|$, degree $k$ gives uniform Taylor error at most $M_0 e^{2W}W^{k+1}/(k+1)!$. Add $\eta$, enclose subarcs in open half-planes, and use the same quadrant orientation and integer winding argument as in Section~\ref{r2:chapter}. At $T=8$ one obtains $D=25$, $W=13$, top Taylor degree $109$, and $260$ accepted top subarcs. The centered transform gives a strictly positive interval at zero. Its seven positive sign brackets and their seven reflections exhaust the integer winding $14$. The upper and lower boundaries are zero-free, and the Euler product and its reflection exclude further completed zeros outside the vertical sides.

The independent height-$32$ contour gives $54$ positive roots and $108$ total. At height $64$ the counts are $136$ and $272$, again with no central zero. For the latter computation the top path in its real parameter $x$ is multiplied by the positive scalar $e^{-a(1+x)}$, $a=\log((1+T^2)/B)$. This does not change any increment of argument or the sign at its left endpoint. In the rescaled finite series the horizontal exponents are $\pm u-a$; summing the Taylor remainders bounds their total by $M_0 e^{2W}(W+a)^{k+1}/(k+1)!$. The transform error is still bounded by $\eta$ because the multiplier is at most one on this path. With $D=61,W=14$, top degree $232$ and $13$ verified subarcs, the integer $272$ is exhausted by the signed intervals. At $T=8$ both the original top calculation and a separate rescaled calculation give the integer $14$. Over $\mathbb Q(\sqrt5)$ and $\mathbb Q(\zeta_5)$ the Hasse--Weil function of this Jacobian is its rational function to the second and fourth powers respectively, by the cyclotomic self-twists in Theorem~\ref{cmr:prime}. Their nonreal zeros consequently have exact multiplicities two and four throughout this window.

\subsection*{A nonabelian quintic base change of the genus-two Jacobian}
Let $K=\mathbb Q(2^{1/5})$, $F=\mathbb Q(\zeta_5)$, $L=KF$, and write $J_d$ for the Jacobian of $y^2+y=d x^5$ at $d=1,2,4,8,16$. The same coordinate scaling as in Theorem~\ref{cmr:prime} gives
\[
 L(J_1/K,s)=\prod_{d=1,2,4,8,16}L(J_d/\mathbb Q,s),\qquad
 L(J_1/L,s)=L(J_1/K,s)^4 .
\]
The two completed functions have degrees $20$ and $80$ as functions over $\mathbb Q$. The second field has Galois group $C_5\rtimes C_4$ acting faithfully on the order-five group.

We detail the arithmetic identification of the four additional rational functions. At primes away from $2$ and $5$ replace $1+4x^5$ in the two point counts by $1+4d x^5$. At 2 the field $F$ is unramified and inert. Adjoining $2^{1/5}$ there gives a totally ramified tame extension of degree five, by the Eisenstein polynomial $x^5-2$. The base $J_1$ has good reduction there and each nontrivial Kummer twist acts nontrivially on all four inertia lines. Thus the new local factor is one and the conductor exponent is four.

At 5 put $F_5^{\mathrm{loc}}=\mathbb Q_5(\zeta_5)$, with uniformizer $\lambda=\zeta_5-1$. The base Hecke character has local conductor exponent two because induction multiplies its conductor norm by $|D_F|=5^3$, while the base conductor is $5^5$. Locally the order-five Kummer character of 2 also has conductor exponent two. To see this write $\alpha^5=16$. The Newton polygon for $(1+X)^5-16$ is the single segment from $(0,4)$ to $(5,0)$ in the $\lambda$-adic valuation. This gives total ramification degree five, with $v(\alpha-1)=4$ and $v(\lambda)=5$ upstairs. Thus $\pi=\lambda/(\alpha-1)$ is a uniformizer. A nonidentity $\sigma$ sends $\alpha$ to $\zeta_5^j\alpha$, and
\[
    v(\sigma\pi-\pi)=5+v((\zeta_5^j-1)\alpha)-4-4=2 .
\]
The sole lower ramification break is one, hence every nontrivial character in this cyclic extension has conductor exponent two. Taking a fourth power does not change the generated quintic Kummer field, so this treats 2 and its nontrivial powers.

The product character at 5 has conductor exponent at most two and cannot be unramified. If it were unramified, its induction restricted to inertia would be the regular representation of the cyclic tame quotient of order four, because $F_5^{\mathrm{loc}}/\mathbb Q_5$ is totally ramified and cyclic of degree four. A generating four-cycle has determinant $-1$. The first cohomology of a polarized abelian surface has determinant the square of a cyclotomic character up to the dual convention; that character is unramified at 5 for auxiliary prime different from 5. Its determinant is therefore one on inertia, a contradiction. The local factor at 5 is consequently one, and the possible global conductors for all four new twists are $2^4 5^4$ and $2^4 5^5$.

There are now two discrete possibilities and an unknown sign. For each conductor $N$ use $B=(2\pi)^2/\sqrt N$ and enclose both sides of $G(-u)=wG(u)$ for $u=\log(3/2),\log 2$ and $w=\pm1$. Even the candidate that fails has a well-defined rapidly convergent Bessel series for these tests. The true conductor and sign satisfy both identities, so excluding all other pairs identifies the true pair. The checker retains one pair in every case. For the uniform trapezoid estimate of the preceding subsection replace the constant $2$ in the separated exponential $e^{-(\sqrt n+e^{x/2})/2}$ by $3$ when necessary. Indeed $c\cos(1/2)>2/3$ even at the largest candidate. The coefficient-sum bound then becomes $4\cdot3^{10}9!$, and the two edge integrals together are at most $10^{18}e^T$. In the sample tail use $e^{-e^{W/2}/3}$. The explicit coefficient-tail term with $c^{-10}$ needs no alteration. For the negative-sign case the odd symmetry forces the central value to vanish. The first derivative of the centered transform, enlarged by $2\eta$ by Cauchy's estimate at radius $1/2$, excludes zero.

At height $16$ the resulting five complete rectangles are
\[
\begin{array}{r|r|r|r|r|r}
 d&N&w&\text{central order}&\text{positive}&\text{total}\\\hline
 1&3125&+1&0&21&42\\
 2&50000&-1&1&27&55\\
 4&50000&+1&0&27&54\\
 8&10000&+1&0&23&46\\
 16&50000&+1&0&28&56
\end{array}
\]
Thus $L(J_1/K,s)$ has central order one, $126$ positive zeros and $253$ total zeros with multiplicity; $L(J_1/L,s)$ has central order four, $504$ positive zeros and $1012$ total zeros. Every zero in these closed rectangles lies on $\operatorname{Re}s=1$. The global factorization includes the ramified primes. Separately, counting residue degrees of $X^5-2$ and taking powers of Frobenius in the base degree-four polynomial verifies ten full degree-$20$ polynomials at good primes of the four splitting types $(1,4),(5),(1,2,2),(1,1,1,1,1)$.

Corollary~\ref{cmr:subfields} gives the complete result for all intermediate fields of $L/\mathbb Q$. Those inside $F$ have functions $L(J_1/\mathbb Q,s)^{[M:\mathbb Q]}$. For those outside $F$ the Hasse--Weil function is $L(J_1/K,s)^{[M:\mathbb Q]/5}$. Indeed in the latter case $H\cap C_5$ is trivial. The three subfields of the cyclotomic quartic, the five conjugate quintic fields, the five conjugate degree-ten fields and $L$ give $14$ embedded fields. To count them, write $G$ as the affine transformations of $\mathbb F_5$. A subgroup disjoint from translations injects into $C_4$. If nontrivial it fixes one common point: a generator $x\mapsto ax+b$, $a\ne1$, has a unique fixed point. There are consequently five possible stabilizing points for the order-four and order-two subgroups, and the trivial subgroup occurs once. For the degree-ten field $K(\sqrt5)$ the height-$16$ central order is two and the total is $506$.

At height $64$ the five twists have respective positive counts $136$, $164$, $165$, $148$ and $164$, in order of increasing $d$. The real-root orders remain $0,1,0,0,0$. Each symmetric contour uses the positive top rescaling above; the required numbers of verified horizontal subarcs are $13,9,24,20,14$. To compare rational constituents, bisect the signed brackets with the finite transform until their endpoints are consecutive multiples of $10^{-10}$. For numerical stability, if evaluation of the weight polynomial at $Z=e^{it/D}$ has a zero-containing scalar interval, write it as $\sum_h e^{24iht/D} P_h(Z)$ with $\deg P_h\le23$. Summing the individual sine or cosine balls supplies a second fallback. These are equal evaluations of the same finite transform and carry the same proved uniform error. In the integer sort of all $777$ resulting intervals the smallest gap between successive endpoints is $571202$ units of $10^{-10}$. All primitive positive ordinates are therefore pairwise distinct.

It follows that over the quintic field $K$ the center has order one and every nonreal zero through $64$ is simple. The full count is $1555$. Over each degree-ten field the center has order two and every nonreal zero has multiplicity two, giving $3110$ zeros; over the full degree-twenty normal closure the center has order four and each nonreal zero has multiplicity four, giving $6220$ zeros. The fields inside $F$ have no central zero; their $136$ distinct positive ordinates have constant multiplicity $[M:\mathbb Q]$. These conclusions cover every embedded intermediate field described above.

Finally, the $T_0=64$ row of Theorem~\ref{li:small-universal} gives strictly positive real Li coefficients through $\lfloor9.7\cdot64^2\rfloor=39\,731$ on each of the five primitive degree-four CM functions. Their positive product decompositions prove the same assertion for $J_1$ over each of the fourteen embedded intermediate fields of the quintic splitting field.

\clearpage
\section{Separating ordinates of distinct elliptic functions}\label{csep:chapter}
A sign interval isolating a zero of one function can overlap a sign interval for a different function. The two argument-principle counts then allow either a common ordinate or two distinct ordinates. To decide between these possibilities in a bounded computation, we give a fast independent sign enclosure on the same axis. All precise cutoffs and separation assertions below refer to the individual rational isogeny classes, so that the equal functions belonging to isogenous curves have already been identified.

\subsection*{An explicit bound for the normalized Hardy function}
Let $E/\mathbb Q$ have conductor $N$, write $B=2\pi/\sqrt N$, and set
\[
 F(z)=B^{-(1+z)}\Gamma(1+z)L(E,1+z).
\]
Its integral and functional equation are those of Section~\ref{cens:chapter}. If $w$ is its sign, put $c=1$ for $w=1$ and $c=1/i$ for $w=-1$. The function
\[
              H(z)=c\exp(\pi z/2)F(iz)
\]
is entire, takes real values on the real axis, and has the same real zeros, with the same orders, as $F(iz)$. The exponential removes the dominant gamma-factor decay in the positive direction. This makes local interpolation at large positive height possible with bounds that grow only polynomially with height.

\begin{lemma}\label{csep:strip}
For every integer $A\ge2$, put $R=2A+1$. Throughout $|\operatorname{Im}z|\le A$,
\[
            |H(z)|\le 12 B^{-(A+1)}|R-iz|^{A+1}.
\]
\end{lemma}
\begin{proof}
On the two horizontal boundaries $z=x\pm iA$, the functional equation $F(v)=wF(-v)$ lets us evaluate $F$ on $\operatorname{Re}v=A$, at one of the imaginary parts $x$ and $-x$. The absolute coefficient bound $|a_n|\le d(n)\sqrt n$ gives
\[
 |L(E,1+A\pm ix)|\le\zeta(A+1/2)^2<4.
\]
The identity $|\Gamma(1+ix)|^2=\pi|x|/\sinh(\pi|x|)$ gives
\[
 e^{\pi x/2}|\Gamma(1+ix)|\le\sqrt{1+2\pi |x|}.
\]
For $x\ge0$, this follows from $u/(1-e^{-u})\le1+u$, applied to $u=2\pi x$ and interpreted by continuity at zero. Negative $x$ gives a smaller left side. The factors in $\Gamma(1+A+ix)/\Gamma(1+ix)$ are $j+ix$, $1\le j\le A$. On either boundary let $d=|R-iz|$. Then $d^2=(R\pm A)^2+x^2$, so $|j+ix|\le d$. Also $d\ge\max(1,|x|)$ and $\sqrt{1+2\pi |x|}\le3d$. The claimed estimate follows there.

To pass to the entire strip, divide by $(R-iz)^{A+1}$. Its zero is outside the strip. In the strip the absolute Mellin integral bounds $F(iz)$ uniformly as a function of $x$, since the exponentially decaying Fourier kernel can be integrated against $e^{A|u|}$. Thus the quotient grows at most exponentially in $|x|$. Multiply by $e^{-\eta z^2}$ with $\eta>0$. The product tends uniformly to zero at either infinite end of the strip, and on the horizontal edges its additional modulus factor is at most $e^{\eta A^2}$. The maximum principle on larger and larger rectangles gives the stated constant times $e^{\eta A^2}$. Letting $\eta$ decrease to zero proves the estimate.
\end{proof}

\subsection*{A 96-point interpolation with a proved remainder}
Take $A=12$, with lattice spacing $\delta=2\pi D/J$ as in the finite Fourier certificates. For an integer $a$ choose 96 real sample nodes $t_j=(a-47+j)\delta$, $0\le j\le95$, and let $P$ denote their Lagrange interpolant to the real function $H$. Put $S=\max_j |t_j|$. Then for every $a\delta\le t\le(a+1)\delta$,
\begin{equation}\label{csep:remainder}
 |H(t)-P(t)|\le
 12\left(\frac{S+3A+1}{B}\right)^{A+1}
 \left(\frac{\delta}{A}\right)^{96}(48!)^2 .
\end{equation}
Indeed, for every real $\xi$ between the sample nodes, the circle of radius $A$ about $\xi$ belongs to the closed strip in Lemma~\ref{csep:strip}. On that circle, $|R-iz|\le S+3A+1$. The Cauchy derivative estimate therefore bounds $|H^{(96)}(\xi)|/96!$ by the first factor of \eqref{csep:remainder} times $A^{-96}$. In the real Lagrange remainder, each of the 48 node distances to the left and each of the 48 distances to the right can be bounded by the corresponding factor in $\delta^{48}48!$. This proves the displayed inequality. The fact that $H$ is real on this axis permits the usual real interpolation remainder with a single intermediate $\xi$.

Each required $H(t_j)$ already has an interval evaluator from the discrete Fourier transform. The few nodes close to the top boundary can exceed the original ordinate bound $T$. Running the same trapezoid estimate out to $T+49\delta$ increases its aliasing term by $e^{49\delta}$ and leaves its real-axis truncation estimate unchanged. Hence multiplying the original total error by $e^{49\delta}$ is sufficient at all 96 nodes, including negative nodes near the origin. Multiply each signed Fourier interval by $e^{\pi t_j/2}$. This produces the enclosures used for the interpolation.

For $u=t/\delta-(a-47)$ off the integer nodes, the exact barycentric formula is
\[
 P(t)=\frac{\displaystyle\sum_{j=0}^{95}
       \frac{(-1)^j\binom{95}{j}H(t_j)}{u-j}}
       {\displaystyle\sum_{j=0}^{95}\frac{(-1)^j\binom{95}{j}}{u-j}}.
\]
The denominator is verified to exclude zero; adding \eqref{csep:remainder} yields a rigorous sign whenever the final interval excludes zero. There is also a direct alternative using the same trapezoid transform. Let $P_0(Z)=\sum_{j=0}^{DW}w_j Z^j$ be the finite weight polynomial and $Z=e^{it/D}$. The unnormalized real scalar is enclosed by $2\operatorname{Re}P_0(Z)-w_0$ for $w=1$, or by $2\operatorname{Im}P_0(Z)$ for $w=-1$, in either case enlarged by the original uniform error. The implementation uses this direct bound whenever the interpolation does not certify the requested sign.

\subsection*{From narrow brackets to products}
Each original interval contains exactly one simple positive zero of its assigned function. A floating approximation proposes two nearby decimal rationals strictly inside that interval. Opposite signs at these endpoints, proved by the preceding interval operations, isolate that same zero. For a class with central order $m$, the sign of the real scalar immediately to the right of zero is
\[
                (-1)^{\lfloor m/2\rfloor}\operatorname{sgn} F^{(m)}(0).
\]
It reverses at each successive simple zero. The checker verifies both endpoint signs in this orientation and records the two numerators over a common decimal denominator. The proposal does not enter the proof.

There are two exact comparisons of the resulting intervals. For a specified set of products it suffices to compare distinct primitive factors that actually occur together. Intervals with disjoint interiors on the original lattices already certify different zeros. The lattice units are rational multiples of a common $2\pi D$, so testing their ordering requires only integer comparisons. The comparison for the refined intervals likewise uses integers. Alternatively, sorting the intervals for \emph{every} positive zero across an entire family gives pairwise distinctness when every lower endpoint is at least the preceding largest upper endpoint. Open endpoints suffice because all endpoint test signs are strict. By conjugation the negative ordinates obey the same separation. The exact multiplicity in a product at each nonreal zero is then simply the exponent of its unique primitive factor. At the center the multiplicities add.

\begin{theorem}\label{csep:joint64}
For the $2463$ rational elliptic isogeny classes of conductor at most $1000$, all $221\,277$ positive ordinates in $(0,64]$ are pairwise different, including across different classes. The smallest distance between any two equals the gap between the $56$th positive ordinates of $582\mathrm d$ and $618\mathrm f$ and lies in $(1.0630831165684567\cdot10^{-9},1.0630831165684569\cdot10^{-9})$.
\end{theorem}
\begin{proof}
Refining the complete first set of contour brackets by the two-sign procedure above gives $221\,277$ rational intervals. In their integer sort, the smallest neighboring lower gap is $104380$ in units of $10^{-14}$, its upper bound is $108380$, and every other neighboring gap is at least $287951$ in those units. The two identified intervals contain the respective $56$th zeros. A nonneighboring pair cannot have a smaller distance than both of the intervening neighboring distances. This identifies the minimizing pair. Bisection of its two original brackets to consecutive multiples of $10^{-25}$, as in Section~\ref{ehc:chapter}, supplies the narrower stated gap. The full list and its integer comparison are supplied in the research package.
\end{proof}


\begin{theorem}\label{csep:joint256}
For the $2463$ rational elliptic isogeny classes with conductor at most $1000$, the $1\,163\,668$ positive ordinates through $256$ are pairwise distinct across the entire family. The complete set of $2620$ primitive functions obtained by adjoining the rational constituents from Theorems~\ref{jzero:thm} and~\ref{cmr:finite} has $1\,249\,609$ distinct positive ordinates through the same height. The unique closest pair in both sets consists of the $212$th positive zero of $726\mathrm e$ and the $181$st of $153\mathrm c$. Its gap is between $2.7010365347071\cdot10^{-10}$ and $2.7010365347073\cdot10^{-10}$.
\end{theorem}
\begin{proof}
All $157$ added functions have complete contour certificates. After strict two-sign refinement, sort the $1\,249\,609$ intervals by lower endpoint. The smallest neighboring lower gap is $25010$ in units of $10^{-14}$ and its upper gap is $29010$. The second-smallest lower gap is $72730$. Thus every neighboring pair is disjoint and only the first pair can minimize the distance. Its two intervals, written as offsets in units of $10^{-14}$ from $125.585873835$, are $(71540,73540)$ and $(98550,100550)$ respectively. Direct upper-incomplete-gamma sums verify all four endpoint signs at $760$ bits. Both members belong to the original $2463$. The independent incomplete-gamma bisection in Section~\ref{ehc:chapter} narrows their ordinates to intervals of width $10^{-23}$ and yields the stated gap by subtraction.
\end{proof}

\begin{theorem}\label{csep:joint9999}
All $4\,256\,834$ positive ordinates through $64$ of all $38\,042$ rational elliptic isogeny classes of conductor at most $9999$ are pairwise distinct. The unique closest pair comprises the $45$th positive zero of $5331\mathrm b$ and the $37$th of $848\mathrm b$, and their gap satisfies
\[
       3.9189115762142\cdot10^{-12}<\Delta<3.9189115762144\cdot10^{-12}.
\]
The closest pair belonging to one primitive function comprises the $100$th and $101$st positive zeros of $3520\mathrm v$, with gap between $0.0015203336312973447145595$ and $0.0015203336312973447145597$.
\end{theorem}
\begin{proof}
For each zero of Theorem~\ref{cens:9999}, two strict finite-transform signs enclose its ordinate in an interval with denominator $10^{14}$. On the added functions the midpoint proposals are obtained simultaneously by ordinary floating interpolation; every interval must pass the direct signed-polynomial check on both endpoints, strictly inside its original bracket. Sorting the $4\,256\,834$ intervals leaves one overlapping pair. Bisection on those two functions to consecutive multiples of $10^{-14}$ resolves both and the full sort has no overlaps.

To identify the minimum, compare lower and upper possible gaps on the sorted intervals. Only two pairs can minimize the distance. Further strict signed bisection of their four endpoints to denominator $10^{16}$ gives a lower gap of $39188$ and an upper gap of $39190$ units for the identified pair. Every other neighboring pair then has lower gap at least $83900$ units. Independent upper-incomplete-gamma sums check all four signs of the two minimizing intervals. An additional signed bisection to denominator $10^{25}$ gives the displayed gap. When only consecutive ordinates of a common primitive function are compared, the smallest lower gap is $152033362930$ and the next lower gap is $267357843898$ in units of $10^{-14}$. The first upper gap is $152033363330$ units. Again bisection to denominator $10^{25}$ gives the stated interval for this unique within-function minimum. The two sets of refined ordinates are also inputs of Section~\ref{ehc:chapter}.
\end{proof}

The single isobaric sum $\boxplus_{[E],\,N_E\le9999}\pi_E$ on $\mathrm{GL}_{76084}/\mathbb Q$ therefore has a completed standard function with $4\,256\,834$ distinct simple positive zeros, their conjugates and a center of order $23\,563$ through height $64$. Its total multiplicity is $8\,537\,231$. This follows by multiplication of the $38\,042$ primitive completions and addition of the central orders. In any product of powers of these primitives, the multiplicity at a positive ordinate through $64$ is exactly the exponent of its unique primitive factor.

Even within the shorter interval $(0,7]$ there are $276\,535$ distinct positive ordinates. The unique closest pair in that interval comprises zero $4$ of $8688\mathrm g$ and zero $3$ of $1856\mathrm m$. Their separation exceeds $2.459\cdot10^{-11}$. Consequently a measurement of any one of those ordinates with absolute error below $1.2295\cdot10^{-11}$ identifies it uniquely within this finite list.

\subsection*{A sharp positive-height isogeny criterion}
The following elementary inverse statement uses only first ordinates. It is stronger than merely detecting whether a single $L$-function vanishes somewhere in a window.

\begin{proposition}\label{csep:inverse}
Let $E_1,\dots,E_r$ be pairwise nonisogenous elliptic curves over $\mathbb Q$. Suppose their completed $L$-functions have distinct first positive critical-line ordinates $\gamma_1,\dots,\gamma_r$. Put $G=\max_i\gamma_i$. Among all abelian varieties rationally isogenous to products of these curves, of arbitrary dimension, the positive critical-line zero divisor through height $G$ determines the rational isogeny class. This is the smallest such height when only positive ordinates are retained.
\end{proposition}
\begin{proof}
The category up to rational isogeny gives a unique vector of nonnegative multiplicities $n_i$ to every such product. Its completed $L$-function is the product of the corresponding powers at every place. Suppose two products have different multiplicity vectors and select the smallest first ordinate $\gamma_i$ at which their corresponding exponents differ. The factors with earlier first ordinates have equal exponents in the two products, so they contribute equally at $\gamma_i$, even if a later zero of one of them happens to equal $\gamma_i$. The factors with larger first ordinate do not vanish there. The difference at $\gamma_i$ is therefore precisely $(n_i-n'_i)\operatorname{ord}_{1+i\gamma_i}\Lambda(E_i,s)\ne0$. This proves injectivity. If $T<G$, the curve with first ordinate $G$ and its square have the same empty positive zero divisor through $T$ and are not isogenous. This proves optimality.
\end{proof}

Here is an exact way to check the proposition's premise in a conductor census without resolving pairs of higher zeros. Start with the first positive interval supplied by each complete contour certificate. Run the signed-kernel polynomial evaluator at a small local ordinate bound exceeding its upper endpoint. Its Fricke sign agrees with the original computation. Strict signs just inside the two ends and subsequent integer bisection give a new interval between consecutive multiples of $10^{-14}$. Sorting these first-marker intervals uses only integer comparisons. The class $11\mathrm a$ has first positive ordinate
\[
  6.362613894713088701<\gamma_{11\mathrm a}<6.362613894713088702;
\]
both decimal endpoints have independent strict signs from the upper-incomplete-gamma representation. The finite support and sharp endpoint of the resulting criterion are recorded below.

For the $2463$ rational elliptic isogeny classes of conductor at most $1000$, every pair of first-marker intervals is disjoint. The least first ordinate belongs to $974\mathrm b$ and lies in $(0.40678016683626,0.40678016683627)$. The second-largest belongs to $14\mathrm a$ and lies in $(5.57928681742950,5.57928681742951)$. Thus the optimal height of Proposition~\ref{csep:inverse} for this family is exactly $\gamma_{11\mathrm a}$. The closest pair of first ordinates in this family belongs to $874\mathrm c$ and $318\mathrm b$. Their certified gap exceeds $0.00000069512540$.

\begin{corollary}\label{csep:inverse9999}
The positive critical-line zero divisor through the first positive ordinate of $11\mathrm a$ determines the rational isogeny class of every abelian variety isogenous to a product of rational elliptic curves of conductors at most $9999$. This height is optimal over products of arbitrary dimension.
\end{corollary}
\begin{proof}
The strict signed first-marker procedure gives $38\,042$ intervals between consecutive multiples of $10^{-14}$. Every neighboring pair in their sorted list is disjoint. The closest first-marker pair belongs to $6405\mathrm l$ and $7056\mathrm v$ and its lower gap is $80304$ such units. The next lower gap is $86544$ units. The two largest first ordinates still belong to $11\mathrm a$ and $14\mathrm a$, so Proposition~\ref{csep:inverse} applies.
\end{proof}

There is a sharp central zero-free height for each observed central order. In the table the actual least positive ordinate is strictly between $k/10^{14}$ and $(k+1)/10^{14}$. Thus the listed open neighborhood of the center is uniform over every curve having the indicated order, and its exact limiting ordinate is attained by the specified class.
\[
 \begin{array}{r|r|c|r}
  m&\text{classes}&\text{minimizer}&k\\\hline
  0&16450&8415\mathrm b&9872560751564\\
  1&19622&9269\mathrm b&51830817685672\\
  2&1969&8789\mathrm d&102506721060797\\
  3&1&5077\mathrm a&205247285847993
 \end{array}
\]
For the overall minimum a separate upper-incomplete-gamma test reduced the original $38\,042$ lower-bracket comparisons to just two candidate functions. The smaller ordinate belongs to $8415\mathrm b$ and lies in $(0.098725607515642906,0.098725607515642907)$, while the other belongs to $6045\mathrm a$ and exceeds $0.116777543062593182$.

\subsection*{Additive relations among the first ordinates}
For every rational isogeny class $E$ of conductor below $10\,000$, write $\gamma_E$ for its first positive ordinate. Let $A$ consist of these $38\,042$ real numbers and zero. An unordered pair in the next statement allows both entries to be equal.

\begin{theorem}\label{csep:sidon}
The $723\,653\,946$ unordered pair sums from $A$ are all different. Their least positive separation $\delta$ satisfies
\[
                 1.73819\cdot10^{-19}<\delta<1.73823\cdot10^{-19}.
\]
The two sums attaining this gap are
\[
 \gamma_{6171\mathrm f}+\gamma_{2310\mathrm h}\quad\hbox{and}\quad
 \gamma_{6195\mathrm c}+\gamma_{5727\mathrm c}.
\]
In particular, the $723\,615\,903$ positive differences of elements of $A$ are distinct.
\end{theorem}
\begin{proof}
The first-ordinate certificates put each nonzero $a\in A$ strictly between $L_a/10^{14}$ and $(L_a+1)/10^{14}$, with integer $L_a$. Set $L_0=0$ and retain $0$ itself as an exact interval. The integer lower endpoints of the pairs are $L_a+L_b$ for $a\le b$. Two integer endpoints differing by at least five guarantee that their true sums differ by more than $3\cdot10^{-14}$. It suffices to enumerate every other comparison.

Sort the integers $L_a$. Partition the integer value range of the sums into half-open blocks. For a block $[v,w)$, and each fixed first index $i$, binary searches in the sorted list give the contiguous set of indices $j\ge i$ such that $v-L_i\le L_j<w-L_i$. This produces each unordered pair once. Sort each block by its sum. When passing to the next block retain every previous integer sum within four units of the previous maximum. A backward scan from each entry now gives all other entries within four units, including across block boundaries. Repeated comparisons are removed by the unordered pair of pair indices. The exact count of entries is $38\,043\cdot38\,044/2$.

The scan gives $10\,383$ comparisons involving $25\,018$ nonzero ordinates. Inside each previously isolated first interval, bisection with strict signed Mellin bounds at each tested endpoint gives a new bracket between consecutive multiples of $10^{-24}$. Adding and subtracting these fine brackets uses integer arithmetic. At the least gap the resulting open bounds are $173819$ and $173823$ such units. Every other candidate lower bound is at least $8209520$ units. The upper-incomplete-gamma expansion independently verifies both endpoint signs for each of the four functions attaining the minimum.

Finally, equality $a-b=c-d>0$ implies $a+d=c+b$. Equality of the unordered pairs gives $a=c$ and $b=d$, since the other possibility would make the difference zero. This proves the difference assertion.
\end{proof}

For the $2463$ classes of conductor at most $1000$, the same proof gives $3\,036\,880$ distinct unordered sums after adjoining zero. The closest sums use the pairs $(906\mathrm a,988\mathrm b)$ and $(427\mathrm a,441\mathrm d)$. The gap is between $8.0307571320\cdot10^{-14}$ and $8.0307571324\cdot10^{-14}$. The fine endpoints here also have an independent incomplete-gamma check.

\begin{corollary}\label{csep:fourth}
For arbitrary complex coefficients $(c_a)_{a\in A}$ put $Z(t)=\sum_{a\in A}c_a e^{iat}$. Then
\[
 \lim_{X\to\infty}\frac1X\int_0^X |Z(t)|^4\,dt
       =2\left(\sum_{a\in A}|c_a|^2\right)^2-\sum_{a\in A}|c_a|^4 .
\]
At finite $X>0$ the absolute error is at most $2(\sum_a|c_a|)^4/(X\delta)$.
\end{corollary}
\begin{proof}
Expand the fourth power. A zero frequency is given by $a+b=c+d$. Theorem~\ref{csep:sidon} says that the two unordered pairs coincide, giving the two orderings of each unequal pair and the single ordering of a repeated pair. A nonzero frequency $\omega$ has $|\omega|\ge\delta$, and its average exponential has modulus at most $2/(X|\omega|)$. Summing absolute values of the remaining coefficient products proves the finite bound.\end{proof}

The argument decides all relations having at most two summands on each side after zero entries are inserted. It leaves higher-order relations and the rational linear independence of the full elliptic zero sets open. The program and fine signed endpoints are supplied with the elliptic census.

\subsection*{Finite determination for an arbitrary fixed family}
There is also an unconditional existence result without a conductor cutoff. In the following statement a zero divisor records a complex position and its multiplicity. This convention matters when a zero may be off the critical line.

\begin{theorem}\label{csep:exist}
Fix pairwise nonisogenous elliptic curves $E_1,\ldots,E_r$ over $\mathbb Q$ and a rational number $H\ge0$. There exist $r$ disjoint rectangles with rational vertices, all above height $H$ and contained in a bounded upper half-strip, such that the numbers of completed zeros in those rectangles determine the multiplicity vector of every product $\prod_i E_i^{n_i}$. A search with rigorous termination finds such rectangles from the equations of the curves. For the general existence assertion one may replace the elliptic curves by pairwise nonisogenous $\mathbb Q$-simple abelian varieties whose completed Hasse--Weil functions are entire of order at most one with the usual weight-one functional equations. Effectivity in this last extension requires effective evaluation of their completions on compact sets.
\end{theorem}
\begin{proof}
Write $A_i$ for the varieties, $g_i=\dim A_i$, $L_i$ for their finite Hasse--Weil functions and
\[
  \Lambda_i(s)=c_i Q_i^{s}\Gamma(s)^{g_i}L_i(s),\qquad
       Q_i=\sqrt{N_i}(2\pi)^{-g_i},
\]
where $c_i\ne0$ is constant. The zeros lie in a fixed vertical strip by the absolutely convergent Euler product on the right and the functional equation on the left. Complex conjugation preserves the divisor.

For every complex $z$ with $\operatorname{Im}z>H$, form the integer row
$v_z=(\operatorname{ord}_z\Lambda_i)_{1\le i\le r}$. If the rows do not span $\mathbb Q^r$, choose a nonzero integer vector $m$ in their common nullspace. The meromorphic product $P=\prod_i\Lambda_i^{m_i}$ then has no zeros or poles above $H$ or below $-H$. Only finitely many are left in the intervening compact part of the zero strip. Applying Hadamard factorization to the two entire products on opposite sides of the quotient gives
\[
           P(s)=e^{a+bs}\prod_{z\in D}(s-z)^{k_z}
\]
for a finite set $D$ and nonzero integers $k_z$. Common canonical-product factors cancel; the finitely many remaining factors of genus one contribute only the displayed polynomial powers and a linear exponential.

On taking logarithmic derivatives at large real $\sigma$, one obtains
\[
 \sum_i m_i\log Q_i+
       \left(\sum_i m_i g_i\right)\frac{\Gamma'}{\Gamma}(\sigma)
       +\sum_i m_i\frac{L'_i}{L_i}(\sigma)
   = b+\sum_{z\in D}\frac{k_z}{\sigma-z}.
\]
The logarithmic gamma derivative grows as $\log\sigma$, whereas each finite Euler logarithmic derivative is $O(2^{-\sigma})$, by absolute convergence at one fixed sufficiently large abscissa. First $\sum_i m_i g_i=0$, and then $b=\sum_i m_i\log Q_i$. The rational function that remains on the right is exponentially small along the positive real axis, hence identically zero. Its residues force $D$ to be empty. Integrating the remaining equality and taking $\sigma\to+\infty$ shows $\prod_i L_i^{m_i}=1$.

At every common good prime $p$, uniqueness of coefficients in the logarithmic Dirichlet series gives
\[
                \sum_i m_i\operatorname{Tr}(\mathrm{Frob}_p\mid V_\ell A_i)=0
                \qquad(p\ne\ell).
\]
Separate the positive and negative exponents into two products. Their rational Tate modules have equal dimensions, and their traces agree on a Chebotarev-dense set. Faltings's semisimplicity and the character criterion give isomorphic Tate modules; his isogeny theorem then gives a rational isogeny of the two products. Uniqueness of multiplicities in the category of abelian varieties up to isogeny now forces every $m_i=0$, a contradiction.

Select $r$ linearly independent rows $v_z$. The union of the zero sets of the finitely many input functions is discrete. Small disjoint rational rectangles surrounding the selected points can therefore be chosen with precisely those zeros in their interiors and no boundary zeros. Their matrix of argument-principle counts has full rank. For a product with exponent column $n$, the rectangle counts are this matrix times $n$, which proves the assertion.

For effectivity, enumerate rational rectangles above $H$ and interleave increasing precisions and finer boundary subdivisions on all of them. Effective holomorphic enclosures on compact sets certify any boundary avoiding the finite union of zero sets, and the argument principle gives the integer count for each factor. The Mellin expansion and its exponential tails give such evaluations for a modular elliptic curve on any fixed compact rectangle. Enumeration of rational boxes and subdivision can be done in parallel rounds, so a box whose boundary passes through a zero cannot stall the search. Test finite sets of $r$ disjoint successful boxes for a nonzero determinant. The existence argument guarantees termination.
\end{proof}

The established literature gives distinct-zero results for different $L$-functions; see E.~Bombieri and A.~Perelli, \emph{Distinct zeros of $L$-functions}, Acta Arith. 83 (1998), 271--281, and Q.~Gazda, \emph{A theorem for distinct zeros of $L$-functions}, arXiv:1504.06556. The preceding argument gives its finite-family form using the Tate-module isogeny theorem. It gives no numerical height bound for an arbitrary family. In the certified conductor range, the sharp height in Proposition~\ref{csep:inverse} uses the positive critical-line divisor alone.

\clearpage
\section{An explicit collision under backward heat flow}\label{ehc:chapter}

The closest pair in Theorem~\ref{csep:joint256} consists of simple zeros of two different elliptic functions. Their product has two exceptionally close simple zeros. Its heat deformation lets us follow both in a fixed disk, including the precise time at which they leave the real axis.

For $E_1=726\mathrm e$ and $E_2=153\mathrm c$ use the physical ordinate of Section~\ref{heat:main} and put
\[
 P_0(z)=(-i)\Lambda(E_1,1+iz)\Lambda(E_2,1+iz).
\]
The two Fricke signs are $-1,+1$, so $P_0$ is real on the real axis. If $K$ is the convolution of their phased elliptic Mellin kernels, define $P_t(z)=\int_{\mathbb R}e^{t y^2}K(y)e^{izy}\,dy$. At $t=0$ this is the completed standard function of an isobaric sum of degree four. Write
\[
 m=125.58587383586045,\qquad h=10^{-10},\qquad u=-t,\qquad v=u/h^2 .
\]

\begin{theorem}\label{ehc:collision}
For $0\le v\le1$, the disk $|z-m|<5h$ contains exactly two zeros of $P_{-h^2v}$, with multiplicity, and has no boundary zeros. There is a unique collision value
\[
 0.9119497952278610 < v_* < 0.9119497952278611.
\]
For $v<v_*$ both zeros are real and simple. At $v_*$ they coalesce into a double zero $z_*$ satisfying
\[
125.58587383586047230844558 < z_* < 125.58587383586047230844559 .
\]
For $v>v_*$ they form a simple nonreal conjugate pair. At $v=0$ these are the $212$th positive zero of $726\mathrm e$ and the $181$st positive zero of $153\mathrm c$.
\end{theorem}

\begin{proof}
The two kernels have two-sided double-exponential decay, so convolution, Fourier transformation and Gaussian smoothing can be interchanged by absolute convergence. For $u>0$,
\[
       P_{-u}(z)=\frac1{\sqrt{4\pi u}}\int_{\mathbb R}
                  e^{-y^2/(4u)}P_0(z+y)\,dy .
\]
Indeed Gaussian integration multiplies the frequency-$\omega$ integrand by $e^{-u\omega^2}$. This identity holds for complex $z$ and also gives the heat equation in $u$. The value at zero time follows by continuity.

Use the two elliptic signed transforms of Section~\ref{cens:chapter} at height $256$ with common denominator $D=152$, both widths $W_i=9$, and respective analytic errors $\epsilon_i$. Their phased product is a trigonometric polynomial
\[
                   G(z)=\sum_k b_k e^{i\omega_k z},\qquad
             \omega_k=k/D,\qquad |\omega_k|\le\Omega=W_1+W_2.
\]
Its coefficients are obtained by finite convolution. Let $M_i$ be the sum of the absolute values of the individual finite coefficients, and let $B_i=2\pi/\sqrt{N_i}$. On $|\operatorname{Im}z|\le1$ a bound for the true individual transform is
\[
                  U_i=\frac4{B_i^2(1-e^{-B_i})^2}.
\]
Indeed the Fourier integrand on each half-line has absolute integral bounded by integrating $2e^{2y}e^{-B_i e^y}/(1-e^{-B_i})^2$ over $y\ge0$. On $|\operatorname{Im}z|\le8h$ the individual finite transform has modulus at most $2M_i$, since $e^{8hW_i}<2$. Whenever $|z-m|\le8h$ and $|y|\le1$, the original elliptic approximation thus gives
\[
          |P_0(z+y)-G(z+y)|\le U_1\epsilon_2+2M_2\epsilon_1 .
\]
For all real $y$ a corresponding bound is $U_1U_2+4M_1M_2$. Gaussian smoothing and $\operatorname{erfc}(a)\le e^{-a^2}$ for $a\ge0$ show that the smoothed product differs from the smoothed finite polynomial by at most
\[
 E=U_1\epsilon_2+2M_2\epsilon_1+
        (U_1U_2+4M_1M_2)e^{-1/(4h^2)}
\]
throughout the disk and throughout $0\le u\le h^2$. At $u=0$ the first two terms alone suffice.

Rescale $z=m+hx$, and multiply the transform by $S=e^{\pi m}/h^2$ and by a constant sign chosen to make the leading quadratic coefficient positive. These multipliers preserve all zeros. Set $J=30$ and $K=15$. Up to the same constant sign, form the real polynomial
\[
 Q(x,v)=S\sum_{j=0}^{J}\sum_{k=0}^{K}
            \frac{G^{(j+2k)}(m)h^{j+2k}}{j!k!}x^jv^k .
\]
For each Fourier term this is the rectangular Taylor expansion of
$e^{i\omega hx}e^{-\omega^2h^2v}$. Writing $X=8h\Omega$ and $V=h^2\Omega^2$, the normalized true transform differs from $Q$ on $|x|\le8$, $0\le v\le1$ by at most
\[
 S\left[E+M_1M_2e^{X+V}
          \left(\frac{X^{J+1}}{(J+1)!}+\frac{V^{K+1}}{(K+1)!}\right)\right]
       <10^{-30}.
\]
The exponential remainder follows by truncating the two factors successively and applying the absolute Taylor bound. Cauchy's inequalities at radius one give derivative errors $10^{-30}$ and $2\cdot10^{-30}$ on $|x|\le5$.

Write $Q(x,v)=\sum B_j(v)x^j$ and $A=B_2(0)>0$. The interval calculation over the entire real parameter interval $[0,1]$ gives $25A>5.0765\cdot10^7$ and
\[
 |B_0|+5|B_1|+25|B_2-A|+\sum_{j=3}^{30}5^j|B_j|+10^{-30}
       <3.709\cdot10^6 .
\]
Thus $Ax^2$ dominates on $|x|=5$, and Rouch\'e's theorem gives two zeros. At each of $x=-5,5$ the oriented function is positive. On the real interval $[-5,5]$ its second derivative exceeds $4.0612\cdot10^6$, by the lower bound
\[
          2B_2-\sum_{j=3}^{30}j(j-1)5^{j-2}|B_j|-2\cdot10^{-30}.
\]
Its first derivative has opposite strict signs at $x=\pm10^{-3}$. There is consequently a unique interior minimum. The exact rescaled heat equation is $\partial_v=\partial_x^2$ on this function. The lower curvature bound shows that the minimum increases strictly with $v$. At $v=4/5$ its value is below $-454650$, and at $v=1$ it is above $357590$.

For clarity, the final parameter bisection needs no approximate double-root solver. Since the minimum occurs in $|x|<10^{-3}$, set
\[
 q=B_0-B_1^2/(4B_2),\qquad
 R=10^{-30}+\sum_{j=3}^{30}10^{-3j}|B_j|.
\]
The stationary point of the quadratic part also belongs to this interval. The true minimum is in $[q-R,q+R]$, by testing this point for the upper bound and minimizing the quadratic for the lower bound. Strict signs at the two decimal endpoints give the stated interval for $v_*$. Substituting that full parameter interval and bisecting the strictly increasing first derivative gives
\[
      0.0002230844558 < (z_*-m)/h < 0.0002230844559 .
\]
Before the crossing, strict convexity and the two positive boundary values give two distinct real roots. At the crossing the positive second derivative gives multiplicity two. Afterward there are no real roots, so the two zeros counted by Rouch\'e are a simple conjugate pair. The original two primitive intervals from Theorem~\ref{csep:joint256} lie inside this disk at $v=0$.
\end{proof}

At the last time $u=10^{-20}$ the upper nonreal root is isolated more closely. In the coordinate $x=(z-m)/h$ it lies within $10^{-15}$ of
\[
                  0.0002230844866064+0.4196431931346889\,i .
\]
For this rational complex center $c$ and $r=10^{-15}$ the verifier checks
\[
 |Q(c,1)|+10^{-30}+
 \sum_{j=2}^{30}\frac{|\partial_x^jQ(c,1)|\,r^j}{j!}
            < r\,|\partial_xQ(c,1)|.
\]
Rouch\'e's theorem against the linear term on $|x-c|=r$ gives a single zero in this disk. Complex conjugation gives the other disk.

An independent incomplete-gamma bisection places the two original primitive roots between consecutive multiples of $10^{-23}$, with respective lower numerators
\[
 12558587383572542048167834,\qquad 12558587383599552413514906 .
\]
If $\Delta$ denotes their original separation and $u_*=h^2v_*$ the physical backward collision time, direct rational comparison gives
\[
                        \left|\Delta^2/8-u_*\right|<10^{-33}.
\]

The same local test applies to the two $56$th positive zeros identified in Theorem~\ref{csep:joint64}. Direct signed refinement puts their difference between $1.0630831165684567\cdot10^{-9}$ and $1.0630831165684569\cdot10^{-9}$. For the product belonging to $582\mathrm d$ and $618\mathrm f$ the unique local double collision occurs at a backward time between $1.4126821409161288\cdot10^{-19}$ and $1.4126821409161292\cdot10^{-19}$. Here the rescaling center is the midpoint of the two refined bracket midpoints and the spatial unit is half their separation. The certificate records these rational inputs exactly.

The arithmetic certificate uses $1152$-bit real and complex balls, exact convolution indices and interval Horner evaluation of the real parameter polynomials. Its source is \texttt{checks/elliptic-heat-collision/verify.py}. The collision is localized in both time and ordinate. Global bounds for the de Bruijn--Newman constant are discussed in Section~\ref{heat:main}.

\subsection*{A collision in one primitive cuspidal function}
For the single elliptic curve $E=3520\mathrm v$, Theorem~\ref{csep:joint9999} identifies its $100$th and $101$st positive zeros as the closest two consecutive ordinates of any one curve in the full height-$64$ census. Their original isolating intervals are refined to consecutive multiples of $10^{-25}$. Let $m$ be the midpoint of their two bracket midpoints and let $h$ be half the difference of the bracket midpoints. The exact inputs are
\[
       m=59.05258189002039225512126855,\qquad
       h=0.0007601668156486723572798 .
\]

\begin{theorem}\label{ehc:primitive}
The phased primitive completion of $3520\mathrm v$ and its heat deformation have exactly two zeros in $|z-m|<5h$ at each backward time $0\le u\le h^2$. They undergo a unique real double collision at
\[
 \begin{aligned}
 2.8892816288482551&<10^7u_*<2.8892816288482557,\\
 59.0525835039724727&<z_*<59.0525835039724728 .
 \end{aligned}
\]
Both are real and simple for $u<u_*$. For $u>u_*$ they form a simple nonreal conjugate pair. At $u=h^2$ the upper zero has imaginary part in $(0.0007601668156254201,0.0007601668156254217)$.
\end{theorem}
\begin{proof}
Use one signed elliptic transform in the proof of Theorem~\ref{ehc:collision}. The local Gaussian approximation error is now $\epsilon_1$ and the global mass bound is $U_1+2M_1$. The normalizing multiplier is $e^{\pi m/2}/h^2$ up to its constant sign. With $J=40$, $K=20$, the uniform approximation error in $|x|\le8$, $0\le v\le1$, is below $10^{-30}$. On the radius-five circle the main quadratic modulus exceeds $118383$, and the remaining terms have modulus below $6244$. On the real diameter the true oriented second derivative exceeds $9168$, and the first derivative changes sign between $-0.01$ and $0.01$. The same count and strict monotonicity thus hold.

Here the parameter test uses the positive curvature more directly. If $F$ denotes the exact oriented transform on the real diameter, $F''\ge c>0$, and $x_0$ lies there, convexity gives
\[
 F(x_0)-\frac{|F'(x_0)|^2}{2c}\ \le\ \min_{[-5,5]}F\ \le\ F(x_0).
\]
A Newton iterate for $Q_x$ supplies only a candidate $x_0$. The polynomial and its derivative, with the proved errors, enclose both sides of this inequality. The minimum is negative at $v=1/4$ and positive at $v=1$. Bisection gives
\[
 \begin{aligned}
 0.5000023692473899&<u_*/h^2<0.5000023692473900,\\
 0.0021231551381&<(z_*-m)/h<0.0021231551382 .
 \end{aligned}
\]
At $u=h^2$ the same affine Rouch\'e check as above gives a disk of radius $10^{-15}$ in the $x$-plane about
$0.0042462936242975+0.9999999999694127\,i$ for the upper root. Rescaling gives the displayed conclusions.
\end{proof}

The most closely spaced pair across distinct functions in the full census gives a still smaller collision time for an isobaric product. Use the two constituents $5331\mathrm b$ and $848\mathrm b$ of Theorem~\ref{csep:joint9999}. In the midpoint and half-gap normalization from their $10^{-25}$ initial intervals, $m=28.9682960274461888587998548$ and $h=1.95945578810715\cdot10^{-12}$. The same strict quadratic domination, curvature and endpoint checks give
\[
 1.9197334927733315\cdot10^{-24}<u_*<1.9197334927733319\cdot10^{-24}.
\]
For $0\le u\le h^2$ these two zeros are simple and real before the unique collision, double at the collision, and a simple nonreal conjugate pair afterward. Each of these calculations uses the undeformed automorphic function specified by its constituent list.

\subsection*{A tunable self-dual automorphic family}
The local mechanism has an unconditional application to infinitely many unitary isobaric parameters. Let $H(z)$ be a phased completed rational elliptic function in the physical ordinate, and let $\gamma$ be any simple real zero. For real $a$ form
\[
                   H(z-a)H(z+a).
\]
Up to constant phase this is the completed standard function of the self-dual isobaric sum $\pi_E|\det|^{-ia}\boxplus\pi_E|\det|^{ia}$. The shifts of the finite argument and of the two archimedean factors agree; their opposite conductor phases cancel.

\begin{proposition}\label{ehc:tuning}
For all sufficiently small $a>0$, the two zeros $\gamma-a,\gamma+a$ undergo a unique local collision at a backward heat time $u_*(a)$. Put $g(x)=H(\gamma+x)/x$ and use the local logarithm of $g(x)/g(0)$. Then
\[
 \begin{aligned}
 u_*(a)&=\frac{a^2}{2}-(\log g)''(0)a^4+O(a^6),\\
 z_*(a)&=\gamma-2(\log g)'(0)a^2+O(a^4).
 \end{aligned}
\]
Every sufficiently small prescribed positive backward time occurs at such a collision. If the completed elliptic function has only critical-line zeros through height $T$ and $|\gamma|\le T-\tfrac12$, then $(\log g)''(0)<0$. In particular $u_*(a)>a^2/2$ for all sufficiently small $a>0$ at each of these zeros.
\end{proposition}
\begin{proof}
Write $A=a^2$ and $x=z-\gamma$. The normalized heat deformation $F(x,A,u)$ at backward time $u$ is analytic near $(0,0,0)$. Indeed the modulated individual kernels have double-exponential decay uniformly when $a$ is bounded in a complex neighborhood, and their convolution can be multiplied by $e^{-u y^2}$ for complex $u$ near zero. Interchanging the two shifts shows that the function is even in $a$, hence analytic in $A$. Normalize by dividing by $g(0)^2$. The starting function and heat equation are
\[
 f(x,A)=F(x,A,0)=(x^2-A)\frac{g(x-a)g(x+a)}{g(0)^2},
 \qquad F_u=F_{xx}.
\]
At the origin $F=F_x=0$ and $F_u=F_{xx}=2$. The Jacobian of $(F,F_x)$ in $(u,x)$ has determinant four. The implicit-function theorem gives unique analytic solutions $u(A),x(A)$ near zero.

Put $\ell=(\log g)'(0)$ and $d=(\log g)''(0)$. The starting derivatives at $x=0$ are
\[
\begin{aligned}
f&=-A-dA^2+O(A^3), & f_x&=-2\ell A+O(A^2),\\
f_{xx}&=2-4\ell^2 A+O(A^2), & f_{xxx}&=12\ell+O(A),\\
f_{xxxx}&=24(2\ell^2+d)+O(A).
\end{aligned}
\]
For instance $\log(g(x-a)g(x+a)/g(0)^2)$ equals $2\log(g(x)/g(0))+A(\log g)''(x)+O(A^2)$. The heat Taylor expansion starts with $F=f+u f_{xx}+u^2f_{xxxx}/2$. The degree-one coefficient of $F=0$ gives $u=A/2+u_2A^2+O(A^3)$. Setting $x=x_1 A+O(A^2)$, the degree-one coefficient of $F_x=0$ and the degree-two coefficient of $F=0$ give
\[
            2x_1+4\ell=0,\qquad (x_1+2\ell)^2+2u_2+2d=0 .
\]
This proves the expansions. Since $u'(0)=1/2$ in $A$, all sufficiently small positive times occur. On a sufficiently small real neighborhood $F_{xx}>0$ and hence $F_u>0$. On a fixed small complex circle the original double zero is isolated. Continuity and Rouch\'e preserve its total multiplicity two. Positive boundary values on the real diameter, strict convexity and the monotonic minimum identify this double zero as the real-to-nonreal collision of the two specified roots.

For the final sign, Hadamard factorization of the order-one elliptic completion gives the absolutely convergent sum
\[
           (\log g)''(0)=-\sum_{\rho\ne\gamma}\frac{1}{(\gamma-\rho)^2}.
\]
Here the zeros are taken in the physical ordinate and counted with multiplicity. Each real zero contributes a strictly positive term before the minus sign. All nonreal zeros, if any, have $|\operatorname{Im}\rho|\le1/2$ by the Euler product and the functional equation. By the hypothesis on the window they also satisfy $|\operatorname{Re}\rho|>T$, and therefore $|\operatorname{Re}\rho-\gamma|>|\operatorname{Im}\rho|$. The real part of $(\gamma-\rho)^{-2}$ is then strictly positive as well. Conjugate terms pair to give a real sum, so $d<0$.
\end{proof}

In particular, degree, finite conductor and an upper bound on the archimedean parameters cannot furnish a positive universal lower bound for the spacing of noncentral critical zeros of unitary self-dual isobaric functions. Fix one elliptic curve with a certified simple positive zero. For all sufficiently small $a>0$ its two distinct norm-twisted constituents give two simple local zeros of spacing exactly $2a$. Their product has degree four and the same finite conductor $N^2$ for every $a$, with archimedean parameters in a fixed compact interval.

For example, at least $4\,216\,266$ positive zeros across the $38\,042$ classes of Theorem~\ref{cens:9999} satisfy this sign criterion. Count those initial brackets whose entire upper end is at most $63.5$ using $\pi<355/113$. At height $256$ the same sufficient integer test includes $1\,160\,931$ of the positive zeros in the conductor-$1000$ census. The cuspidal input $\pi_E$ contributes its two unitary shifts to the indicated degree-four isobaric sum. For the historical relation between close zeros and heat flow see G.~Csordas, W.~Smith and R.~S.~Varga, \emph{Lehmer pairs of zeros, the de Bruijn--Newman constant $\Lambda$, and the Riemann Hypothesis}, Proc. Sympos. Appl. Math. 48 (1994).

\subsection*{Polynomial heat limits in all even automorphic degrees}
Varying more than two character shifts allows any monic polynomial with real roots as a local model. Let $p(x)=\prod_{j=1}^d(x-\lambda_j)$ with $\lambda_j\in\mathbb R$, and let $H$ and $\gamma$ be as in Proposition~\ref{ehc:tuning}. Denote by $P_t^{(\varepsilon)}$ the heat deformation of $\prod_j H(z-\varepsilon\lambda_j)$ in the physical ordinate.

\begin{theorem}\label{ehc:polynomial}
Uniformly on compact subsets of $\mathbb C^2$ in $(x,v)$,
\[
 \frac{P_{-\varepsilon^2v}^{(\varepsilon)}(\gamma+\varepsilon x)}
      {(\varepsilon H'(\gamma))^d}
  =\sum_{k=0}^{\lfloor d/2\rfloor}
       \frac{v^k}{k!}\,p^{(2k)}(x)+O(|\varepsilon|)
       \qquad(\varepsilon\longrightarrow0).
\]
The convergence holds with all fixed partial derivatives. For real $\varepsilon$ the undeformed inputs are completed standard functions of unitary isobaric sums on $\mathrm{GL}_{2d}/\mathbb Q$. A multiset of roots $\{\lambda_j\}$ stable under negation gives a self-dual sum. Every real collision of multiplicity exactly two in the polynomial heat equation persists at nearby scaled time and position in the corresponding automorphic heat family.
\end{theorem}
\begin{proof}
Before division, the left numerator is an integral over $\mathbb R^d$. Its integrand is the product of the original cuspidal kernels, multiplied by
\[
 \exp\left(i\gamma Y+i\varepsilon(xY-\sum_j\lambda_j y_j)
                       -\varepsilon^2vY^2\right),
                \qquad Y=\sum_j y_j .
\]
The individual two-sided double-exponential bounds dominate this expression and its derivatives on compact subsets of all the complex parameters. The integral is analytic there. Its Taylor coefficient $C_n(x,v)$ at degree $n$ in $\varepsilon$ is a polynomial in $(x,v)$, and differentiating the integral gives
\[
               \partial_v C_n=\partial_x^2 C_n .
\]
At $v=0$ the product of the $d$ simple-zero expansions starts with $\varepsilon^dH'(\gamma)^d p(x)$. Thus $C_n(x,0)=0$ for $n<d$, and $C_d(x,0)=H'(\gamma)^d p(x)$. For a polynomial solution of the heat equation, its coefficients in powers of $v$ are successively determined by its value at $v=0$. This gives $C_n=0$ for $n<d$ and the stated finite heat polynomial for $C_d$. Holomorphy on larger compact neighborhoods gives the uniform remainder and its derivative bounds.

The automorphic assertion follows by the norm-twist formula and multiplication of standard factors in an isobaric sum. If the heat polynomial $Q(x,v)$ has a real double root at $(x_0,v_0)$ with $Q_{xx}(x_0,v_0)\ne0$, the Jacobian determinant of $(Q,Q_x)$ in $(v,x)$ is $Q_{xx}^2$ there. The divided transform extends analytically to $\varepsilon=0$ by the preceding argument. The implicit-function theorem therefore gives a unique nearby double-root position and time for all sufficiently small $\varepsilon$, real when $\varepsilon$ is real. A small surrounding circle and Rouch\'e's theorem preserve the count of two. The exact heat equation determines which side has two real roots, as in Proposition~\ref{ehc:tuning}.
\end{proof}

\clearpage
\section{Every rationally split modular Jacobian}\label{jzero:chapter}
The conductor census gives a single height bound for a geometric family specified without a level cutoff. Write $J_0(N)$ for the Jacobian of the modular curve $X_0(N)$ over $\mathbb Q$. A positive-dimensional $J_0(N)$ \emph{splits over $\mathbb Q$ into elliptic curves} if it is $\mathbb Q$-isogenous to a product of elliptic curves defined over $\mathbb Q$. We use the Hasse--Weil function of $H^1$, with central line $\operatorname{Re}s=1$.

\begin{theorem}\label{jzero:thm}
Let $N$ be any positive integer for which $J_0(N)$ splits over $\mathbb Q$ into elliptic curves. Every zero of its completed $H^1$ $L$-function with $|\operatorname{Im}s|\le256$ lies on $\operatorname{Re}s=1$. This includes all 71 positive-dimensional examples in the classification, up to dimension 205 and $L$-degree 410.
\end{theorem}

The level classification is prior work. Elkies, Howe and Ritzenthaler, \emph{Genus bounds for curves with fixed Frobenius eigenvalues}, Theorem 5.1, give its corrected complete form. If $N=2^en$ with $n$ odd, their theorem gives an explicit finite list of allowed $n$ and bounds on $e$. There are 86 levels counting the genus-zero curves. The last two are 1152 and 1200. They corrected the earlier classification by adding 28 and deleting 672. Our use of the theorem is to ensure that the geometric statement covers every $N$, including any that could lie above the dyadic and conductor cutoffs used to find the factors of a given Jacobian.

Let $\mathcal E(M)$ denote the set of rational isogeny classes of elliptic curves of conductor $M$. Put $d(k)=\sum_{a\mid k}1$. For a split level $N$ there is a rational isogeny
\begin{equation}\label{jzero:isog}
        J_0(N)\ \sim_{\mathbb Q}\
        \prod_{M\mid N}\ \prod_{E\in\mathcal E(M)} E^{d(N/M)}.
\end{equation}
Indeed newform theory decomposes $S_2(\Gamma_0(N))$ into its primitive systems and their degeneracy images. A primitive normalized form $f$ of level $M$ contributes the linearly independent forms $f(az)$ for $a\mid N/M$. The Eichler--Shimura construction gives a simple modular abelian variety $A_f$ per Galois orbit, and its dimension is the degree of the coefficient field. The split condition implies that each such $A_f$ is an elliptic curve. Conversely every $E\in\mathcal E(M)$ occurs by modularity and the equality of conductor and primitive level. This proves \eqref{jzero:isog} with the indicated multiplicities. In particular rationality of the Fourier coefficients and oldform multiplicity are checked at the primitive level, not on each degeneracy image as if it were another primitive automorphic function. The multiplicities also follow from the standard modular Jacobian old/new decomposition.

An isogeny over $\mathbb Q$ induces an isomorphism on the rational $\ell$-adic Tate modules compatible with Galois action. At a finite prime $p$, use $\ell\ne p$ and take inertia invariants. Both the invariant space and its Frobenius characteristic polynomial respect the direct sums in \eqref{jzero:isog}. Thus the full Hasse--Weil Euler product, including every bad prime, satisfies
\[
      \Lambda(s,J_0(N))=\prod_{M\mid N}\prod_{E\in\mathcal E(M)}
             \Lambda(s,E)^{d(N/M)}
\]
with the customary multiplicative gamma and conductor factors. These follow from direct sums on the real Weil group and additivity of the conductor exponent. The constituent functions are entire. At a point $s_0$ the order for the Jacobian is the sum of the constituent orders times $d(N/M)$. This keeps repeated elliptic constituents and coincident ordinates of distinct curves with their true multiplicities.

Here is an independent integer check on the enumeration. The genus formula for $X_0(N)$ reads
\[
 g_0(N)=1+\frac{\mu(N)}{12}-\frac{e_2(N)}4-\frac{e_3(N)}3
               -\frac12\sum_{d\mid N}\varphi(\gcd(d,N/d)),
 \qquad \mu(N)=N\prod_{p\mid N}(1+1/p).
\]
If $4\nmid N$, the order-two elliptic point count is $e_2(N)=\prod_{p\mid N}(1+(\frac{-4}{p}))$, and it vanishes otherwise. If $9\nmid N$, the corresponding order-three count is $e_3(N)=\prod_{p\mid N}(1+(\frac{-3}{p}))$, and it vanishes otherwise. Use Kronecker symbols at 2. For each of the 86 levels the script checks
\[
        g_0(N)=\sum_{M\mid N}d(N/M)\,\#\mathcal E(M).
\]
The union of the needed classes contains just 222 distinct elliptic functions. Of these, 136 have conductor at most 388. The remaining 86 have respective conductors $396,400,432,576,600,720,1152,1200$, with class counts $3,8,8,9,9,10,20,19$. For these last two levels the models are selected by the same database extraction as for the main census. Each of the 222 functions has its own symmetric-window certificate. Their central orders are zero in 171 cases and one in 51. Hence every primitive weight-two system occurring at any of the split levels has only simple critical-line zeros in the certified window, even when the product for its Jacobian has multiple zeros. The completed product equality now proves Theorem~\ref{jzero:thm}. A genus-zero Jacobian gives the empty product and contributes no zeros.

For the three largest genera the actual multiplicities at height 256 are
\[
\begin{array}{r|rr|rrr}
 N&g_0(N)&\text{degree}&\operatorname{ord}_{s=1}&\text{positive}&\text{full}\\\hline
720&121&242&8&48086&96180\\
1152&161&322&32&69529&139090\\
1200&205&410&28&86122&172272
\end{array}
\]
For example the central zero for $J_0(1200)$ has multiplicity $28$ because its elliptic factors with simple central zeros are repeated with their degeneracy multiplicities. Simplicity for individual primitive functions would give the wrong conclusion about this product.

Theorem~\ref{csep:joint256} identifies the exact nonreal multiplicities at every level. Distinct primitive elliptic constituents have disjoint positive ordinates through $256$. A zero belonging to a constituent of conductor $M$ therefore has multiplicity exactly $d(N/M)$. The numbers of distinct positive ordinates for $N=720,1152,1200$ are $16361,27753,30902$ respectively. For $J_0(1200)$ precisely $342$ positive ordinates attain the largest multiplicity $12$, all from the primitive class $20\mathrm a$. The negative ordinates have the same multiplicities by conjugation.

There is an arithmetic consequence at every level of the classification. All $222$ rational elliptic constituents have analytic order zero or one. Gross--Zagier and Kolyvagin identify their Mordell--Weil ranks with these orders and give finiteness of their Tate--Shafarevich groups. The rational isogeny decomposition then gives
\[
 \operatorname{rank}J_0(N)(\mathbb Q)=\operatorname{ord}_{s=1}\Lambda(J_0(N),s),
 \qquad |\mathop{\rm Sha}(J_0(N)/\mathbb Q)|<\infty
\]
for every one of the classified split levels. The three displayed examples have Mordell--Weil ranks $8,32,28$ respectively.

\paragraph{Sources and reproducibility.}
For the global classification see N.~D.~Elkies, E.~W.~Howe and C.~Ritzenthaler, \emph{Genus bounds for curves with fixed Frobenius eigenvalues}, Proc. Amer. Math. Soc. 142 (2014), 71--84, Theorem 5.1. For degeneracy multiplicities use Atkin--Lehner--Li newform theory and the Eichler--Shimura decomposition. The script \texttt{checks/elliptic-census/split\_jacobians.py} reproduces the genus equalities and records the complete weighted factor list at every positive genus. The genus formula and the classification are cited inputs. The zero-window consequence uses the independently certified elliptic factors.

\clearpage
\section{An asymmetric zero count for a quartic twist}\label{qt:chapter}

Take the elliptic newform $f(\tau)=\eta(\tau)^2\eta(11\tau)^2=\sum_{n\ge1}a_n e^{2\pi i n\tau}$ and the primitive character $\chi$ of conductor $5$ determined by $\chi(2)=i$. Set $g=f\otimes\chi$ and $M=275$. In arithmetic coordinates its completed function is
\[
 \Lambda_\chi(s)=\left(\frac{\sqrt M}{2\pi}\right)^s
             \Gamma(s)\sum_{n\ge1}a_n\chi(n)n^{-s}.
\]
The series defines the function on its initial half-plane. The Mellin integral below continues it to an entire function. The coefficients are genuinely complex: $a_2=-2$, so the coefficient of $2^{-s}$ is $-2i$ for $\chi$ and $2i$ for $\bar\chi$. In particular the two standard functions and the two automorphic representations are distinct. The four-dimensional product $L_\chi L_{\bar\chi}$ would conceal which factor contributes each zero. We shall count zeros of $L_\chi$ individually.

\begin{theorem}\label{qt:count}
At every height certified for $\chi$ in Table~\ref{qt:table}, every zero of $\Lambda_\chi(s)$ with $|\operatorname{Im}s|\le T$ is simple and belongs to $\operatorname{Re}s=1$. The function has no zero on the real axis. For $T=256$, there are $451$ zeros of positive ordinate and $449$ of negative ordinate. At $T=4096$ the corresponding counts are $10807$ positive and $10806$ negative. For its dual $f\otimes\bar\chi$, the two counts are reversed.
\end{theorem}

\subsection*{Ramification, root number, and the sign on the critical line}
The level-$11$ form corresponds to the rational isogeny class $11.\mathrm a$. At each $p\nmid55$ the primitive local factor after twisting is
\[
 (1-a_p\chi(p)p^{-s}+\chi(p)^2p^{1-2s})^{-1}.
\]
The original curve has split multiplicative reduction at $11$, with $a_{11}=1$. The character is unramified there and $\chi(11)=1$, so this local factor remains $(1-11^{-s})^{-1}$ and the conductor exponent remains $1$. At $5$ an unramified two-dimensional Weil representation is twisted by a ramified character of conductor exponent one. Its inertia invariants vanish, its local factor is one, and its conductor exponent is two. The twist is thus cuspidal of exact conductor $11\cdot5^2$ with nebentypus $\chi^2$. These facts also follow from Atkin and Li, \href{https://doi.org/10.1007/BF01390245}{\emph{Twists of newforms and pseudo-eigenvalues of $W$-operators}}. The two embeddings belong to newform orbit $275.2.\mathrm b.\mathrm a$; the quadratic companion used below is $275.2.\mathrm a.\mathrm b$. See the \href{https://www.lmfdb.org/ModularForm/GL2/Q/holomorphic/11/2/a/a/}{twist table for $11.2.\mathrm a.\mathrm a$}.

We derive the root number to fix the sign conventions in the computation. Write $N=11$, $q=5$, $e(v)=e^{2\pi iv}$, $\tau(\chi)=\sum_{a=1}^{q-1}\chi(a)e(a/q)$, and
\[
 T_x=\begin{pmatrix}1&x\\0&1\end{pmatrix},\qquad
 W_N=\begin{pmatrix}0&-1\\N&0\end{pmatrix}.
\]
The Gauss sum identity for primitive characters gives
\[
 g=\frac1{\tau(\bar\chi)}\sum_{a=1}^{4}\bar\chi(a) f|T_{a/q}.
\]
Indeed its $n$th Fourier coefficient is $a_n\chi(n)$, including zero for $5\mid n$. The eta transformation gives $f|W_N=-f$ in the weight-two, determinant-normalized slash action. For $1\le a<q$ choose $b$ with $Nab\equiv-1\pmod q$. Then
\[
 q^{-1}T_{a/q}W_{Nq^2}
 =\begin{pmatrix}(1+Nab)/q&a\\Nb&q\end{pmatrix}
   W_N T_{b/q},
\]
whose middle matrix belongs to $\Gamma_0(N)$. A scalar matrix acts trivially. The change of variable $a\mapsto b$ gives
$\bar\chi(a)=\bar\chi(-1)\chi(N)\chi(b)$. Using
$\tau(\chi)\tau(\bar\chi)=\chi(-1)q$ therefore yields
\[
 g|W_M=-\varepsilon_\chi(f\otimes\bar\chi),\qquad
 \varepsilon_\chi=\chi(N)\tau(\chi)^2/q=\tau(\chi)^2/5 .
\]
Consequently $\Lambda_\chi(s)=\varepsilon_\chi\Lambda_{\bar\chi}(2-s)$. The conjugate coefficient relation also gives $\Lambda_{\bar\chi}(s)=\overline{\Lambda_\chi(\bar s)}$. Thus reflection across the critical line acts on the zeros of this single function, whereas conjugation of ordinates passes to the distinct dual function.

Put $r=\overline{\tau(\chi)}/\sqrt5$. It has modulus one and $r\varepsilon_\chi=\bar r$. The rotated function $F(1+z)=r\Lambda_\chi(1+z)$ satisfies $F(1-\bar z)=\overline{F(1+z)}$. In particular $F(1+it)$ is real for real $t$. Set
\[
 K(u)=r e^u g(ie^u/\sqrt M),\qquad
 F(1+z)=\int_\R K(u)e^{zu}\,du .
\]
The Fricke identity gives $K(-\bar u)=\overline{K(u)}$ in $|\operatorname{Im}u|<\pi/2$. The integral follows first from the Dirichlet series for $\operatorname{Re}z$ large and then from double exponential decay at both ends. Since $\chi(-1)=-1$ one has $r/\tau(\bar\chi)=-1/\sqrt5$. In particular the checker evaluates the kernel by the exact formula
\[
 K(u)=-\frac{e^u}{\sqrt5}\sum_{a=1}^4\bar\chi(a)
       f(a/5+ie^u/\sqrt M).
\]
The midpoint of a horizontally displaced integration line has real kernel by the displayed symmetry. No positivity of that kernel is assumed.

\subsection*{Uniform quadrature for a complex kernel}
For a signed height $t$, choose a real tilt $\alpha$ and a margin $d>0$. In the tilted computations take $d=\min(1/2,32/\max(1,|t|))$ and $\alpha=\operatorname{sgn}(t)(\pi/2-2d)$ with $\alpha=0$ at $t=0$. The independent unshifted computations use $\alpha=0$, $d=1$. We then have $\cos(\alpha+v)\ge d/2$ for $|v|\le d$. In the tilted case this follows from $\cos(\alpha+v)\ge\sin d\ge d/2$ away from $t=0$. Put $\lambda=d/6$. At $u=x+i(\alpha+v)$ with $x\ge0$ use
$2\pi/\sqrt{275}>1/3$ and $|a_n\chi(n)|\le d(n)\sqrt n\le2n$ to obtain
\[
 |K(u)|\le B e^x\exp(-\lambda e^x),\qquad
 B=\frac2{(1-e^{-\lambda})^2}.
\]
Here $d(n)$ denotes the divisor count. Conjugate reflection handles $x<0$ with $x$ replaced by $|x|$.

Write $G_\alpha(z)=e^{-i\alpha z}F(1+z)$. Contour translation gives
$G_\alpha(z)=\int_\R K(x+i\alpha)e^{zx}dx$. With $D$ a positive integer, $h=1/D$, integer $U\ge3$, $J=UD$, and $w_j=hK(jh+i\alpha)$, put
\[
 P(z)=w_0+\sum_{j=1}^J\left(w_j e^{jhz}+\bar w_j e^{-jhz}\right).
\]
For $|\operatorname{Re}z|\le1$ and $|\operatorname{Im}z|\le |t|$ set $M_0=2B/\lambda^2$ and $v_U=\lambda e^U$. When $v_U>2$, the explicit error is
\begin{equation}\label{qt:err}
 |G_\alpha(z)-P(z)|\le
 E=\frac{2M_0 e^{d|t|}}{e^{2\pi dD}-1}
       + M_0(1+v_U)e^{-v_U}.
\end{equation}
In fact the $L^1$ norm on either line displaced by $d$ in the $x$ variable is at most
$2B e^{d|t|}\int_0^\infty e^{2x-\lambda e^x}dx
\le M_0 e^{d|t|}$. Poisson summation and contour translation of each nonzero Fourier coefficient give the first term in \eqref{qt:err}. The majorant $e^{2x-\lambda e^x}$ decreases for $x\ge U$; both omitted trapezoidal tails together cost at most $2B\int_U^\infty e^{2x-\lambda e^x}dx=M_0(1+v_U)e^{-v_U}$. The checker selects $D,U$ so each error term is below $e^{-144}(|t|+3)^{-4}$ in the tilted method or below $e^{-2|t|-64}$ in the unshifted method, and checks both inequalities by outward rounding. Distinct tilts at positive and negative heights keep the respective values above the absolute error.

\subsection*{An argument count without ordinate symmetry}
We count on the rectangle $0\le\operatorname{Re}s\le2$, whose right vertical edge is in the open half-plane of absolute convergence and whose left edge is its reflection. This contains all completed zeros, by the Euler product and reflection. Here containment refers to the real-coordinate range; the ordinate range is imposed separately. On the right edge, $s=2+it$, define
\[
 \theta(t)=\arg r+\operatorname{Im}\{s\log(\sqrt M/(2\pi))+\log\Gamma(s)\}
             +\operatorname{Im}\log L_\chi(s),
\]
with the Euler branch for $\log L$ and the gamma branch real on the positive axis. The Euler term has modulus at most $2\log\zeta(3/2)<\pi$ by Hasse and the local factors. It equals the principal argument of the quotient of the right endpoint by its known gamma, conductor and rotation factors. This step fixes the winding on the right edge even when $|t|$ is large.

Let $\Delta_G(t)$ be the change of argument of $G_\alpha$ along $z=x+it$ from $x=1$ to $x=0$. Its value for $F$ is $\Delta_F(t)=\Delta_G(t)-\alpha$. When the half-edge is certified nonzero, set
\[
 h(t)=\frac{\theta(t)+\Delta_F(t)}{\pi}.
\]
This is an integer because the value on the central axis is real and nonzero. Equivalently the checker encloses the displayed quantity and checks that its unique integer has the parity of the certified axis sign. For $a<b$, with the edges at $a$ and $b$ certified, the right edge contributes $\theta(b)-\theta(a)$. Reflection by $s\mapsto2-\bar s$ shows that the left edge traversed downwards gives the same change, while the full top and bottom edges contribute $2\Delta_F(b)$ and $-2\Delta_F(a)$. The argument principle therefore gives
\begin{equation}\label{qt:winding}
 N(a,b)=h(b)-h(a).
\end{equation}
In particular both $h(0)-h(-T)$ and $h(T)-h(0)$ must be evaluated. A zero on the real axis away from the symmetry line is excluded only after the entire horizontal edge at height zero has been certified; knowing the value at its midpoint would not suffice.

To certify each half-edge, expand $P(1/2+x+it)$ in powers of real $x$, $|x|\le1/2$. The coefficient of degree $n$ for the $j$th pair is
\[
 \frac{(jh)^n}{n!}
 \left(w_j e^{(1/2+it)jh}+(-1)^n\bar w_j e^{-(1/2+it)jh}\right).
\]
The constant term also receives $w_0$. If $H=2\sum_{j=0}^J|w_j|$, truncation at $m$ costs at most $H e^U (U/2)^{m+1}/(m+1)!$ by Taylor's formula. The checker takes $m$ so this is below $E$ and adds both errors. On a real subinterval it evaluates the resulting complex polynomial by interval Horner arithmetic. If its enclosure lies in an open coordinate half-plane, the entire path avoids zero there and its argument change is the principal argument of the endpoint ratio, using left over right for a traverse from right to left. Otherwise subdivide. The sum of all certified argument increments gives $\Delta_G(t)$.

Finally evaluate $P(i t_k)$ on the grid $t_k=\pm2\pi D k/m_0$, $k\ge0$, using a discrete Fourier transform of size $m_0>2J$. For positive ordinates and the negative-exponent transform convention the slots $-j$ and $j$ contain $w_j$ and $\bar w_j$, respectively. For negative ordinates exchange the two slots. The result is real, as is $G_\alpha(it_k)=e^{\alpha t_k}F(1+it_k)$. Values exceeding the error in absolute value certify the sign of $F$ at that grid point. Use the separately certified axis signs at the endpoints as well. Intervals between successive certified samples of opposite signs have disjoint interiors and each contains a line zero of odd multiplicity. When their number equals \eqref{qt:winding}, each contains precisely one simple zero and the full rectangle contains no others. This proves Theorem~\ref{qt:count} once the integer computations in the table are made.

\begin{table}[htbp]\centering\small
\begin{tabular}{@{}rllrrrr@{}}
\toprule
$T$&Character&Contour&$h(-T)$&$h(T)$&$N(-T,0)$&$N(0,T)$\\
\midrule
16 & $\chi$ & real & $-15$ & $14$ & 14 & 15\\
256 & $\chi$ & real and tilted & $-450$ & $450$ & 449 & 451\\
1024 & $\chi$ & tilted & $-2251$ & $2249$ & 2250 & 2250\\
4096 & $\chi$ & tilted & $-10807$ & $10806$ & 10806 & 10807\\
% Quartic maximum-height row.
256 & $\chi^2$ & tilted & $-449$ & $449$ & 449 & 449\\
1024 & $\chi^2$ & tilted & $-2250$ & $2250$ & 2250 & 2250\\
4096 & $\chi^2$ & tilted & $-10807$ & $10807$ & 10807 & 10807\\
% Quadratic maximum-height row.
\bottomrule
\end{tabular}
\caption{Individual contour counts. For the phase choices in the proof, $h(0)=-1$ for $\chi$ and $h(0)=0$ for $\chi^2$.}\label{qt:table}
\end{table}

The program \texttt{quartic.py} produces a JSON certificate for each row. Each certificate records the signed lifts, the subdivisions certifying the horizontal path, the mesh and truncation parameters, and two lists of intervals for the absolute ordinates distinguished by sign. The repeated count at $256$ uses two different contours, precisions and Fourier meshes. On large meshes the checker factors the transform over subgrids with residue classes of the sample index; this gives the same Fourier sum term by term and bounds memory use. In detail, for mesh $m_0=bq$ and index $k=r+qv$, its Fourier multiplier factors as $e^{2\pi i jr/m_0}e^{2\pi i jv/b}$. Rotate the coefficient by the first factor, sum over its residue modulo $b$, then take a size-$b$ transform. The real-axis signs at height zero are checked directly and agree across the two signed transforms.

An independent check avoids translated eta values and Fourier interpolation. Let $a=2\pi/\sqrt{275}$ and $s=1+it$. Splitting the Mellin integral at $u=0$ and using its conjugate symmetry gives
\[
 F(1+it)=2\operatorname{Re}\sum_{n\ge1}
     r a_n\chi(n)(an)^{-s}\Gamma(s,an).
\]
For the terms beyond $n=L$, substitute the integral
$(an)^{-s}\Gamma(s,an)=\int_1^\infty e^{-any}y^{it}\,dy$. The absolute error is at most $4e^{-a(L+1)}/(a(1-e^{-a}))$, uniformly in real $t$. The program \texttt{audit\_quartic.py} uses $L=150$ and integer coefficients computed by multiplying the defining eta-product factors as polynomials. It encloses the incomplete gamma values and verifies both endpoint signs of every one of the 29 intervals through height $16$. Bisection, followed by fresh sign evaluations at the stated rational decimal endpoints, yields the following isolating intervals for the two zeros nearest the real axis on their respective sides:
\[
 -2.0532709<\gamma_-<-2.0532708,\qquad
  1.1447753<\gamma_+<1.1447754 .
\]
No equality between the two absolute ordinates is possible.

\subsection*{An elliptic curve over the quartic cyclotomic field}
Let $K=\Q(\zeta_5)$ and let $E$ be the conductor-$11$ elliptic curve above. The regular representation of $\operatorname{Gal}(K/\Q)$ is the sum of $1,\chi,\chi^2,\bar\chi$. Artin formalism for the Tate module, or cyclic base change, gives in arithmetic normalization
\[
 L(E/K,s)=L(f,s)L(f\otimes\chi^2,s)
               L(f\otimes\chi,s)L(f\otimes\bar\chi,s).
\]
This is an identity of complete finite Euler products by induction of the Tate module. At the totally ramified prime $5$ the curve has good reduction, the residue field stays $\F_5$, and only the untwisted local factor survives on the right. The prime $11\equiv1\pmod5$ splits completely; all four multiplicative factors are $(1-11^{-s})^{-1}$. The conductor ideal has norm $11^4$. Since $|d_K|=5^3$, the rational conductor of the degree-eight function is $|d_K|^2 11^4=11\cdot275^3$, the product of its four rational conductors. Its four complex gamma factors agree as well. The product of the four root numbers is $+1$. For the companion $\chi^2=(\,\cdot\,/5)$ one has $\chi^2(11)=1$, $\tau(\chi^2)=\sqrt5$, root number $1$, and rotation $1$. The Fourier expansion and conductor are obtained by exactly the same character twist. Its kernel is now $e^u\sum_a\chi(a)^2 f(a/5+ie^u/\sqrt M)/\sqrt5$. These substitutions are the program's quadratic-companion option. The two signed counts are computed separately as a check even though the function has real coefficients.

At $T=256$ the four counts in the positive half-strip are $319,449,451,449$, respectively. The first is the level-$11$ row of Section~\ref{num:chapter}, and the remaining ones are in Table~\ref{qt:table} or its conjugate. At $T=1024$ the respective counts are $1725,2250,2250,2250$. Thus $L(E/K,s)$ has exactly $3336$ zeros with $|\operatorname{Im}s|\le256$ and $16950$ with $|\operatorname{Im}s|\le1024$, counted with multiplicity, all on $\operatorname{Re}s=1$. At $4096$ the four positive counts are $8708,10807,10807,10806$, giving $82256$ zeros in the full strip through that height.

\begin{corollary}\label{qt:base}
Over $K=\Q(\zeta_5)$, all $82256$ zeros through $|\operatorname{Im}s|\le4096$ are simple.
\end{corollary}
\begin{proof}
It remains to exclude coincidence between zeros from different character factors. Take the individually certified positive-ordinate brackets for $f$, $f\otimes\chi^2$, $f\otimes\chi$, and $f\otimes\bar\chi$, reflecting the negative brackets of $\chi$ for the last list. Replace each endpoint by a rational multiple of $10^{-12}$ rounded toward the interior; the terminal height stays exact. Reevaluate the two signs of every interval. Their signs still oppose, so the smaller interval retains the unique root of its original bracket. The sign tester writes $q_t=e^{it/D}$ and evaluates $P(it)=2\operatorname{Re}\sum_{j=0}^J w_jq_t^j-w_0$. All rounding is outward and the analytic bound \eqref{qt:err} is deducted before deciding a sign. For the base form the corresponding bound in Section~\ref{num:chapter} applies. Evaluating in polynomial blocks prevents excessive growth of interval widths in a long Horner recurrence.

Sort the four lists together by their left endpoints. When two intervals overlap, bisect them and keep the subinterval with opposite endpoint signs; the uniqueness already proved for the original interval justifies the update without a monotonicity assumption. Repeat until successive intervals are strictly separated. The archive contains the final integer endpoints and their factor labels. For the height $4096$ express each of the four original positive-ordinate lists in the common unit $2\pi/2^{50}$. Every nonterminal endpoint is represented exactly, since its Fourier mesh is a power of two dividing $2^{50}$ and its spacing denominator is an integer. A terminal endpoint $4096$ is rounded down and its sign is certified again. The original endpoint signs are otherwise inherited from the individually certified lists. Bisect the wider of two colliding intervals at an integer midpoint in these units, using the block-polynomial evaluator as above. After $41$ rounds the $41128$ positive intervals are strictly disjoint, with least reported separation $1241344\cdot 2\pi/2^{50}$. For the decimal grid the same program \texttt{quartic\_base\_change.py} certifies $1668$ pairwise disjoint positive intervals through $256$ and $8475$ through $1024$. The smallest separations between reported endpoints are $23968449\cdot10^{-12}$ and $3306\cdot10^{-12}$, respectively, lower bounds for all adjacent positive-ordinate zero gaps in the indicated ranges. Every zero in them is simple for the product, since only one factor vanishes there and its order is one. An auxiliary audit checks every factor label, bracket containment and separation directly against the original tables. For the pair with the smallest reported interval separation it also recalculates the four endpoint signs by a staggered rule. Its positive abscissae are $(j+1/2)/D$, paired with their negatives. The Poisson aliasing coefficients acquire factors $(-1)^k$ of absolute value one. The first omitted half-grid tail is bounded by integrating from $U-1$ instead of $U$ in \eqref{qt:err}, because the majorant is decreasing there. This independent test uses margin $d=\min(2/5,18/(H+1))$, tilt $\operatorname{sgn}(t)(\pi/2-3d)$ and separate choices of $U,D$ for the maximum ordinate $H$ of the tested pair, rounded up. Reality of the product and absence of central zeros prove the result.
\end{proof}

\begin{corollary}[A curve over a real quadratic field with no rational model]\label{qt:genuine}
Let $F_0=\Q(\sqrt5)$ and $d=-(5+\sqrt5)/2$. The quadratic twist $E^{(d)}/F_0$ admits no model over $\Q$. All $43226$ zeros of its completed Hasse--Weil function up to absolute ordinate $4096$ are simple and lie on $\operatorname{Re}s=1$.
\end{corollary}
\begin{proof}
The square of $\zeta_5-\zeta_5^{-1}$ is $(\zeta_5+\zeta_5^{-1})^2-4=d$. Hence $K=F_0(\sqrt d)$. The nontrivial character of $\operatorname{Gal}(K/F_0)$ has exactly two extensions to the cyclic group $\operatorname{Gal}(K/\Q)$, namely $\chi$ and $\bar\chi$. Its induced representation is their direct sum. Twisting the Tate module of $E/F_0$ by this quadratic character and applying induction gives
\[
 L(E^{(d)}/F_0,s)=L(f\otimes\chi,s)L(f\otimes\bar\chi,s).
\]
At the single prime over $5$, twisting the good Tate module by the ramified tame quadratic character gives conductor exponent two and local factor one. At both primes over $11$ the quadratic extension splits, so each multiplicative factor is $(1-11^{-s})^{-1}$. The curve has good reduction elsewhere. Its conductor ideal has norm $5^2 11^2$ and $|d_{F_0}|=5$, so the absolute conductor is $275^2$, with two complex gamma factors and root number $+1$. At $256$ the positive-height count is $451+449=900$. The negative count is the same and the center is nonzero. The certified disjoint brackets in Corollary~\ref{qt:base} show that the two factors have no zero in common through absolute height $1024$. At $1024$ the positive and negative counts are each $2250+2250=4500$. At $4096$ they are each $10807+10806=21613$. Applying exactly the same dyadic refinement to the two quartic factors alone gives an independent check. The resulting $21613$ positive brackets have smallest separation $158892032\cdot2\pi/2^{50}>0$. Reflection gives separation on the negative side as well.

Finally $j(E)$ is rational and different from $0$ and $1728$; for example its multiplicative reduction at $11$ gives negative $11$-adic valuation. Every rational model with this $j$ would be a quadratic twist $E^{(b)}$ by some $b\in\Q^\times$. An isomorphism over $F_0$ between this model and $E^{(d)}$ would imply $b/d\in F_0^{\times2}$. This would give $K=F_0(\sqrt b)$, an extension of $\Q$ with Galois group of exponent at most two. But $\operatorname{Gal}(K/\Q)$ is cyclic of order four. This contradiction proves the first assertion.
\end{proof}

\subsection*{When descent up to isogeny already gives descent}
The quadratic example has a stronger property: its isogeny class over $F_0$ contains no curve from $\Q$. The obstruction can be read from any finite abelian Galois group.

\begin{proposition}\label{qt:descent}
Let $L/\Q$ be a finite abelian extension with Galois group $G$, let $H\le G$, and put $F=L^H$. For a rational elliptic curve $A$ and a quadratic or trivial character $\psi:H\to\{1,-1\}$, write $B/F$ for the corresponding twist of $A/F$. The following statements are equivalent.
\begin{enumerate}
\item $B$ is $F$-isogenous to an elliptic curve defined over $\Q$.
\item $B$ is $F$-isomorphic to a quadratic or trivial twist of $A$ defined over $\Q$.
\item $\psi$ is trivial on $H\cap G^2$, where $G^2=\{g^2:g\in G\}$.
\end{enumerate}
In all cases the completed $L$-function of $B/F$ is the product of the completed elliptic twists $L(A,\chi,s)$ for the characters $\chi$ of $G$ with $\chi|_H=\psi$, with the same gamma-factor convention.
\end{proposition}
\begin{proof}
By linear algebra over $\mathbf F_2$, a character $H\to\{1,-1\}$ extends to such a character of $G$ exactly if it vanishes, in additive notation, on $H\cap G^2$. Indeed the induced functional on the image of $H$ in the vector space $G/G^2$ extends to the whole space. A quadratic extension of the character recovers $B$ over $F$ by twisting $A$ over $\Q$. Thus (3) implies (2), and (2) implies (1).

Write $V_A$ for the rational Tate module with coefficients extended to $\overline{\Q}_\ell$. If $C/\Q$ is as in (1), induction to $G_\Q$ and the projection formula give
\[
  V_C\otimes\operatorname{Ind}_H^G\mathbf 1
   \simeq\bigoplus_{\chi|_H=\psi}(V_A\otimes\chi).
\]
Here induction on the finite quotient is inflated to $G_\Q$. Finite abelian character orthogonality gives the displayed summands and their multiplicity one. Local Artin formalism for this latter decomposition also gives the stated completed identity, including ramified factors, as in Theorem~\ref{ebc:divisor}. Each elliptic Tate representation over $\Q$ is absolutely irreducible: Faltings gives semisimplicity and the scalar endomorphism ring, since a rational elliptic endomorphism acts by a rational algebraic integer on invariant differentials. The trivial character occurs in $\operatorname{Ind}_H^G\mathbf 1$. Therefore $V_C\simeq V_A\otimes\chi$ for at least one $\chi$ in the sum. Comparison of their determinants, both elliptic determinants being cyclotomic, gives $\chi^2=1$. Its restriction to $H$ is $\psi$, proving (3).
\end{proof}

For $G=C_4$ and $H=G^2$, the nontrivial $\psi$ violates (3). Thus the $43\,226$ zeros in Corollary~\ref{qt:genuine} concern an elliptic isogeny class over a real quadratic field that cannot be obtained by base extension from $\Q$. More generally, for $G$ cyclic of order $2r$ and $|H|=2$, the nontrivial twist descends up to isogeny exactly when $r$ is odd. When $r$ is even all its rational factors come from nonquadratic characters. The obstruction is independent of the conductor and of complex multiplication.

\subsection*{Two independently resolved near coincidences}
Let $t_1<t_2$ be the closest positive ordinates in the ranges indicated below. Each printed endpoint is an integer multiple of $10^{-18}$. The two roots in the first row come from $\chi^2$ and $\chi$. The second row uses $\chi$ and $\bar\chi$.
\[
\begin{array}{c|r|r|r}
\text{field}&T&\lfloor 10^{18}t_1\rfloor&\lfloor 10^{18}t_2\rfloor\\ \hline
\Q(\zeta_5)&4096&997519238744324364166&997519238754827295989\\
\Q(\sqrt5)&4096&869192507874627891774&869192509007710574939
\end{array}
\]
Thus the first gap lies strictly between $10502931822\cdot10^{-18}$ and $10502931824\cdot10^{-18}$. The second lies between $1133082683164\cdot10^{-18}$ and $1133082683166\cdot10^{-18}$. To certify uniqueness of each minimum, remove the gap between those two brackets from its corresponding separation file. Every other positive bracket pair in the first file is separated by at least $79446016\cdot2\pi/2^{50}$, and every other adjacent pair in the second by at least $5084545024\cdot2\pi/2^{50}$. Both lower bounds exceed the respective upper bounds above.

We verified the refined endpoints in two ways. The first evaluates the eta products on the staggered contour. The second counts points over finite fields on $y^2+y=x^3-x^2-10x-20$. At good $p$ this gives $a_p=p+1-\#E(\F_p)$, using direct enumeration of the numbers of $y$ for each cubic value; at $11$ use $a_{11}=1$. The Hecke recurrence $a_{pn}=a_pa_n-pa_{n/p}$, with $a_{n/p}=0$ for $p\nmid n$, gives the composite coefficients at good $p$. At $11$ use $a_{11n}=a_n$. With these integers the staggered integrator evaluates $K(u)=r e^u\sum a_n\chi(n)\exp(-2\pi n e^u/\sqrt M)$ directly. For a sample $u=x+i\alpha$ in $k\le x\le k+1$ set $b=(2\pi/\sqrt M)e^k\cos\alpha$ and $q=e^{-b}$. The contribution to the absolute kernel error beyond $n=L$ is bounded by
\[
 2e^{k+1}q^{L+1}\left(\frac{L+1}{1-q}+\frac{q}{(1-q)^2}\right),
\]
by $|a_n|\le2n$. The cutoffs make this less than $e^{-160}$ at every node, adding at most $2Ue^{-160}$ to the integral error. The direct series and staggered eta methods certify opposing signs at identical rational endpoint pairs to $10^{-18}$. The root counts from the argument principle guarantee that each bracket contains exactly the one claimed root. In particular the close pair near $997.51923875$ consists of two distinct zeros with verified factor labels.

\subsection*{The same count for an arbitrary coprime primitive twist}
The signed winding formula works for any primitive Dirichlet character $\xi$ of conductor $q$ prime to $11$. The same matrix identity with $N=11$ and variable $q$ yields root number $\xi(11)\tau(\xi)^2/q$. At a prime dividing $q$, both local roots of the original unramified representation acquire the same ramified twist, so its conductor exponent is twice the character's local conductor exponent and its local Euler polynomial is one. At $11$ the original linear factor acquires $\xi(11)$, and the new conductor is exactly $11q^2$. Choose either unit phase $r_\xi$ satisfying $r_\xi^2=\overline{\xi(11)\tau(\xi)^2/q}$. The rotated kernel is
\[
 \frac{r_\xi e^u}{\tau(\bar\xi)}
 \sum_{a\bmod q}\bar\xi(a)
       f(a/q+ie^u/(q\sqrt{11})).
\]
It obeys the same conjugate reflection and double-exponential bounds, now taking any positive $\lambda$ below $(2\pi/(q\sqrt{11}))\min_{|v|\le d}\cos(\alpha+v)$. All horizontal bounds and the right-edge Euler logarithm bound apply. A second rotation choice changes every lifted integer by a common odd integer and leaves their differences intact.

For any two ordinates whose horizontal boundary segments are zero-free, the completed function has a strictly positive minimum modulus on those segments. Increasing the precision, the cutoff and the sample density, and subdividing, will therefore certify the two argument variations to any prescribed accuracy. If every interior zero lies simply on the central line, finitely many disjoint pairs of nearby ordinates of opposite signs exist by simplicity and compactness. A sufficiently fine mesh with sufficiently small absolute error detects at least that many sign changes, and the argument count forces equality. Conversely, equality certifies exactly this zero statement. This termination assertion assumes the indicated simple zero configuration for the converse direction. It does not give its truth for untested twists or unbounded heights.

\subsection*{Asymmetry at arbitrarily high ordinates}
For $\gamma>0$, let $m_+(\gamma)$ and $m_-(\gamma)$ denote the total multiplicities of the zeros of $\Lambda_\chi$ with respective ordinates $\gamma$ and $-\gamma$, including any off-line zeros. We prove two unconditional statements. The first determines a particular oscillatory sum of the signed differences. The second shows that each sign supplies a positive proportion of surplus multiplicity in every sufficiently high dyadic band.

\begin{theorem}\label{qt:asymmetry}
As $T\to\infty$,
\[
 \sum_{0<\gamma\le T}(m_+(\gamma)-m_-(\gamma))e^{i\gamma\log2}
      =\frac{i\log2}{\pi}T+O\!\left(\frac T{\log T}\right).
\]
Moreover, with $u_+=\max(u,0)$, both
\[
 \sum_{T<\gamma\le2T}(m_+(\gamma)-m_-(\gamma))_+\gg T\log T,
 \qquad
 \sum_{T<\gamma\le2T}(m_-(\gamma)-m_+(\gamma))_+\gg T\log T .
\]
The quantities count multiplicities. They do not require simplicity of a zero or uniqueness of its real part at a given ordinate.
\end{theorem}
\begin{proof}
For the first assertion put $D(s)=L_\chi'(s)/L_\chi(s)-L_{\bar\chi}'(s)/L_{\bar\chi}(s)$. The completions have identical gamma and conductor factors, so $D$ is also their logarithmic derivative difference. Differentiating the two functional equations gives $D(2-s)=D(s)$. At a zero $\rho$ its residue is $d(\rho)=\operatorname{ord}_\rho\Lambda_\chi-\operatorname{ord}_\rho\Lambda_{\bar\chi}$. On the right edge $\operatorname{Re}s=2$, write $D(s)=-\sum_{n\ge2} b(n)n^{-s}$, absolutely convergently. The arithmetic coefficient is $b(2)=(\chi(2)-\bar\chi(2))a_2\log2=-4i\log2$. The logarithm of a local factor introduces a $1/j$ at a $j$th power; its logarithmic derivative multiplies by $j\log p$, giving this coefficient at $p=2$.

Integrate $2^sD(s)/(2\pi i)$ counterclockwise over the rectangle of real parts $0,2$ and heights $0,T$, initially excluding zero ordinates at the top. The bottom contains no poles by Theorem~\ref{qt:count}. On the upward right edge only $n=2$ has zero frequency, and the others integrate to $O(1)$ since $|\log(2/n)|$ is bounded below for integers $n\ne2$. The downward left edge also costs $O(1)$ by $D(s)=D(2-s)$ and the nonzero frequencies $\log(2n)$. The full upper-edge integral is $O(\log T)$ uniformly even when $T$ approaches a zero ordinate. Indeed, on a fixed strip the logarithmic derivative is the sum of $1/(s-\rho)$ over $|\operatorname{Im}\rho-T|<1$ with error $O(\log T)$, and there are $O(\log T)$ such terms. Each term integrates against $2^s$ to $O(1)$. To see this, subtract $2^\beta$ from $2^\sigma$ inside the integral for $\rho=\beta+i\gamma$. The resulting numerator is $O(|\sigma-\beta|)$. The remaining integral of $1/(\sigma-\beta+i(T-\gamma))$ is bounded since $1/2\le\beta\le3/2$, the two endpoints stay away from zero, and their continuous argument difference has magnitude at most $\pi$. Consequently
\[
 \sum_{0<\operatorname{Im}\rho\le T}d(\rho)2^\rho
       =-\frac{b(2)}{2\pi}T+O(\log T)
       =\frac{2i\log2}{\pi}T+O(\log T).
\]
For an exceptional ordinate take a decreasing sequence of admissible heights tending to it. Local finiteness of the zero set includes exactly the zeros at that ordinate. This is Landau's formula evaluated at the first differing coefficient. For the general formula and the two local estimates used on the horizontal edge, see \href{https://www.cs.uleth.ca/~akbary/uniform.pdf}{Akbary--Murty, equation (5) and Lemmas 3--4}, following M. R. Murty and V. K. Murty.

To pass from zero locations to ordinates, take $c=1/8$ in \href{https://arxiv.org/pdf/2103.01956}{Beckwith--Liu--Thorner--Zaharescu, Theorem 1.1}. For either of these fixed unitary-normalized cuspidal representations it bounds the number of zeros of arithmetic real part at least $1+v$ and absolute ordinate at most $T$ by $O(T^{1-cv}\log T)$ for $v\ge0$. Reflection across the central line and integration with weight $2v$ give
\[
 \sum_{|\gamma|\le T}(\beta-1)^2
   \ll T\log T\int_0^{1/2}v\,e^{-cv\log T}\,dv
   \ll T/\log T .
\]
Paired real parts $1+v,1-v$ contribute the factor $2\cosh(v\log2)$ after averaging their two $2^\rho$ weights and removing $e^{i\gamma\log2}$. Its difference from $2$ is $O(v^2)$. The dual has, at $+\gamma$, the same multiplicity as the original at $-\gamma$. Dividing the preceding Landau formula by $2$ proves the first assertion.

For the positive-proportion assertions we apply the short-interval joint distribution theorem of \href{https://matwbn.icm.edu.pl/ksiazki/aa/aa83/aa8336.pdf}{Bombieri--Perelli, \emph{Distinct zeros of $L$-functions}, Theorem 2}. We give the hypothesis checks and the deduction at the level of ordinates. Their analytic hypotheses B(I)--B(III) hold for the two unitary twists by the Euler factors, Hasse bound, entireness of finite order, and functional equations above. Their density hypothesis D follows from the stated bound with $c=1/8$. Under B(IV) the two self-correlations of prime coefficients divided by $p$ must equal $\log\log x+C+O(1/\log x)$ through $x$, and the cross-correlation must equal $C'+O(1/\log x)$. Away from $5,11$, the diagonal coefficient is $\lambda_f(p)^2$ and the cross-coefficient is $\lambda_f(p)^2\chi^2(p)$, where $\lambda_f(p)=a_p/\sqrt p$. Apply the Rankin--Selberg prime number theorem of \href{https://www.intlpress.com/site/pub/files/_fulltext/journals/pamq/2007/0003/0002/PAMQ-2007-0003-0002-a004.pdf}{Liu--Ye, Theorem 2.3} to the self-dual pairs $(f,f)$ and $(f,f\otimes\chi^2)$. In the second pair the central characters forbid a nontrivial imaginary norm twist, and the distinct coefficients at $2$ rule out equality. The respective logarithmically weighted prime sums are thus $x+O(xe^{-c'\sqrt{\log x}})$ and $O(xe^{-c'\sqrt{\log x}})$ for some $c'>0$. Removing the powers of degree at least two costs $O(\sqrt x\log^2 x)$ by the finite-place bounds. Partial summation with $1/(t\log t)$ gives precisely B(IV) in both cases. All hypotheses here are unconditional.

Let $N_1,N_2$ count positive ordinates with multiplicity for the quartic twist and its dual. Write $B(T)=\sum_{T<\gamma\le2T}(m_+(\gamma)-m_-(\gamma))_+$. Suppose $B(T)=o(T\log T)$ along a sequence. Along it take
\[
 M=\min\!\left(\sqrt{\log T},\sqrt{\frac{T\log T}{1+B(T)}}\right),
 \qquad h=M/\log T .
\]
Then $M\to\infty$, $M\le\log T/\log\log T$, and $h B(T)=o(T)$. Bombieri--Perelli's Theorem 2 gives independent nondegenerate Gaussian limits for the two argument increments on $(t,t+h]$ divided by their respective fixed positive multiples of $\sqrt{\log M}$, with $t$ uniform in $[T,2T]$. The two counting formulas have the same gamma term and conductor. Thus on a set of $t$ of measure $\gg T$, the first count increment exceeds the second by at least $c_0\sqrt{\log M}$ for some $c_0>0$. Indeed take a bounded open rectangle of positive Gaussian mass with the first standardized argument increment above one and the second below minus one. The argument-principle counting formula converts their difference to the difference of increments of $N_j$, up to $O(h+T^{-1})$, as in equations (14)--(16) of the cited paper. Discard $t>2T-h$ to keep both endpoints inside the same band. On the remaining interval,
\[
 \int_T^{2T-h}\!\big(N_1(t+h)-N_1(t)-N_2(t+h)+N_2(t)\big)_+\,dt
       \le h B(T).
\]
Each positive surplus at a single ordinate is counted over a set of $t$ of length at most $h$. The Gaussian lower bound contradicts $hB(T)=o(T)$. This proves $B(T)\gg T\log T$, and exchanging the two functions proves the other inequality. This is the short-interval argument of Bombieri--Perelli applied after grouping the zero multiplicities at each ordinate.
\end{proof}

The cumulative counts themselves have unbounded discrepancies. For every fixed $A>0$, a set of heights of measure $\gg_A T$ in $[T,2T]$ satisfies
\[
 |N_1(t)-N_2(t)|>A\sqrt{\log\log T}.
\]
Indeed take $M=\sqrt{\log T}$ in the same joint limit and choose the two bounded Gaussian rectangles sufficiently far apart in the imaginary coordinates. At a set of $t$ of measure $\gg_A T$ the difference of count increments then exceeds $2A\sqrt{\log\log T}$. At least one of $t,t+h$ has the stated cumulative discrepancy. Translating the second set preserves its measure, and discarding the final interval of length $h$ keeps the heights in the band.


\begin{corollary}[Two fixed degree-two cusp forms]\label{qt:allpairs}
Let $\pi_1,\pi_2$ be distinct fixed unitary cuspidal representations of $\mathrm{GL}_2/\Q$. Suppose each standard gamma factor can be written as a product of $\Gamma(\lambda s+\mu)$ with $\lambda>0$ and $\operatorname{Re}\mu\ge0$, up to a conductor factor. This includes holomorphic newforms of every weight at least one and arbitrary nebentypus, and Maa\ss\ newforms tempered at infinity. Let $m_j(\gamma)$ be the total multiplicity at ordinate $\gamma$. In every sufficiently high dyadic interval, both $\sum_{T<\gamma\le2T}(m_1(\gamma)-m_2(\gamma))_+$ and its reverse are $\gg_{\pi_1,\pi_2} T\log T$. For every fixed good prime $p$ their normalized prime coefficients determine the oscillatory discrepancy
\[
 \sum_{0<\gamma\le T}\big(m_1(\gamma)-m_2(\gamma)\big)e^{i\gamma\log p}
  =-\frac{(\lambda_{\pi_1}(p)-\lambda_{\pi_2}(p))\log p}{2\pi\sqrt p}\,T
     +O_{p,\pi_1,\pi_2}(T/\log T).
\]
In particular, for every genuinely non-self-dual form in this class, its positive ordinates and its reflected negative ordinates have the stated two positive-proportion discrepancies, with multiplicities.
\end{corollary}
\begin{proof}
For the weighted identity subtract the two unitary Landau formulas at $p$, divide by $\sqrt p$, and pair real parts $1/2+v$ and $1/2-v$. The same second-moment bound on $v$ gives the displayed error. The density estimate applies to both forms. The gamma hypothesis and the standard analytic properties give B(III). At good finite primes the Satake parameters satisfy $|\alpha|\le p^{7/64}$ by the Kim--Sarnak bound, with the stronger bound one for holomorphic forms. At each of the finitely many other primes the standard parameters have modulus $<p^{1/2}$ by local unitarity and genericity. Increasing the fixed exponent at these primes if needed still gives B(I) with an exponent strictly below $1/2$.

For B(IV) consider all four transfers $\Pi_{jk}=\pi_j\boxtimes\widetilde\pi_k$. These are isobaric automorphic representations on $\mathrm{GL}_4/\Q$ by \href{https://arxiv.org/abs/math/0007203}{Ramakrishnan, \emph{Modularity of the Rankin--Selberg $L$-series}, Theorem M}. The constituents of Ramakrishnan's isobaric sum are unitary cuspidal representations in his conventions, Sections 2.1--2.2. A one-dimensional norm character produces its actual simple pole on $\operatorname{Re}s=1$, since every other cuspidal standard factor is nonvanishing there. The Rankin--Selberg pole classification then identifies a trivial constituent exactly once on the diagonal. Off the diagonal there are no norm characters unless the two forms are related by a nonzero imaginary norm twist, in which case exactly one corresponding norm character occurs. A self-twist by a nonzero imaginary norm power is impossible by comparing central characters.

Apply the prime number theorem of Liu--Ye separately to each constituent paired with the trivial self-dual representation of $\mathrm{GL}_1$. At every good prime the combined first coefficient is $\lambda_{\pi_j}(p)\overline{\lambda_{\pi_k}(p)}$. Good higher prime powers cost $O(x^{1/2+7/32}\log^2 x)$ since each transfer parameter is bounded by $p^{7/32}$. All powers of the finitely many excluded primes contribute $O(\sqrt x\log^2 x)$ by the local bound $|\alpha|\le p^{1/2}$ for a unitary cuspidal constituent. Except in the nonzero norm-twist case, the weighted prime sums are $\mathbf1_{j=k}x+O_{\pi_1,\pi_2}(xe^{-c\sqrt{\log x}})$ with $c>0$. In that case the off-diagonal main term is $x^{1+i\tau}/(1+i\tau)$ for a fixed $\tau\ne0$. Its contribution after partial summation with $1/(t\log t)$ is a constant plus $O_\tau(1/\log x)$, by integrating $e^{i\tau v}/v$ by parts. This verifies B(IV) in all cases. For B(II), at a good prime the two parameters are either on the unit circle or, by local unitarity, have the form $z,1/\bar z$, hence the same complex argument. In both cases the sum of their squared moduli is bounded by $2+|\lambda_\pi(p)|^2$. The diagonal prime estimate just established therefore gives B(II). Bombieri--Perelli's joint theorem and the short-interval inequality in the preceding proof now apply. The Riemann--von Mangoldt main terms have the same leading coefficient because both degrees are two. Differences in conductors or gamma parameters add at most $O(h+T^{-1})$ to the difference of increments. Finally the ordinates for the dual are the negatives of the original ordinates, giving the last assertion.
\end{proof}

\begin{corollary}[Stability under perturbing the ordinates]\label{qt:matching}
Under the hypotheses of Corollary~\ref{qt:allpairs}, form the two positive-ordinate multisets in $(T,2T]$. Fix $A>0$. For all sufficiently large $T$, every one-to-one partial matching which moves its matched entries by at most $A/\log T$ leaves $\gg_{A,\pi_1,\pi_2}T\log T$ unmatched entries in each multiset. More generally, in any matching which covers the smaller multiset completely, the sum of its absolute displacements, each truncated at $1/\sqrt{\log T}$, is $\gg_{\pi_1,\pi_2} T\sqrt{\log\log T}$.
\end{corollary}
\begin{proof}
Take $h=M/\log T$ as in the proof of Theorem~\ref{qt:asymmetry}. If $U_1$ entries of the first list are unmatched and $d_j$ is the distance moved by pair $j$, the integral of the positive difference of the two window counts over $T<t<2T-h$ is at most
\[
                         hU_1+\sum_j\min(h,d_j).
\]
Indeed the sets of starting heights putting the two members of one pair in $(t,t+h]$ are intervals of length $h$ translated by $d_j$. The portion counted only by the first has length $\min(h,d_j)$. The Gaussian limit lower-bounds the same integral by $cT\sqrt{\log M}$ for $M\to\infty$ in the admissible range. If a matching at distance $A/\log T$ had $U_1=o(T\log T)$ along a sequence, choose $M$ in the earlier proof with $U_1$ in place of $B(T)$. Then $hU_1=o(T)$ while the sum over the $O(T\log T)$ matched pairs is $O_A(T)$, a contradiction. Interchange the two functions for the opposite direction. Finally, for a matching covering the smaller list take that list first, so $U_1=0$, and set $M=\sqrt{\log T}$. This proves the second assertion. The comparison uses multiplicity and one-to-one matching; several zeros can be close to the same counterpart.
\end{proof}

In particular, over any fixed finite family of distinct eligible forms, a list on a high dyadic interval remains distinguishable even after its entries move by $A/\log T$ and a sufficiently small fixed fraction are removed or inserted. A list compatible with two candidates would match most entries of those two candidates within $2A/\log T$ by composing the two matchings. The family is finite, so the discrepancy constants have a positive minimum.

For the quartic twist, both positive and reflected negative lists contain exactly $5855$ ordinates between $2048$ and $4096$. The certificate brackets rigorously bound the maximum cardinality of a matching with tolerance $A/\log 2048$ as follows.
\[
\begin{array}{r|r|r|r}
A&\text{lower bound}&\text{upper bound}&\text{necessarily unmatched per list}\\\hline
0.1&254&849&5006\\
0.5&1788&2616&3239\\
1&3563&4083&1772\\
2&5185&5322&533\\
4&5735&5764&91\\
8&5855&5855&0
\end{array}
\]
To compute the upper bound allow any pair whose two intervals approach within the tolerance. For the lower bound accept a pair only when all its possible distances are within the tolerance. Process the two sorted lists from the bottom, discarding an interval too early for all remaining partners; when both first intervals qualify, match them. This is optimal for the indicated eligibility rule: a later match to either first interval can be replaced by the first-first pair, and if both first intervals were matched, the two former partners qualify for one another by the ordering of endpoints. Rounding the tolerance outwards or inwards, respectively, preserves the bounds. The minimum possible mean displacement in a complete matching lies in $[0.1859,0.2115]$ and its minimum possible maximum displacement in $[0.7772,0.8100]$. Both minima are attained by matching in order, by the elementary exchange inequality for absolute distance; applying the corresponding endpoint bounds to each sorted pair gives the displayed intervals.

The same argument recovers the finite Euler factors from the positive ordinates alone. Write $A_p(k)=\alpha_{p,1}^k+\alpha_{p,2}^k$ in unitary normalization, padding a local factor of degree less than two with zero parameters. For each fixed prime $p$ and each fixed integer $k\ge1$, the one-form version of the preceding computation gives
\[
 A_p(k)=-\frac{2\pi p^{k/2}}{\log p}\,\frac1T
     \sum_{0<\gamma\le T}e^{ik\gamma\log p}
     +O_{\pi,p,k}(1/\log T).
\]
Indeed the coefficient of $p^{-ks}$ in $-L'/L$ is $A_p(k)\log p$, also at a ramified prime. Apply Landau's formula at $x=p^k$, divide by $p^{k/2}$, and use the quadratic bound for the horizontal displacement. Taking the two limits for $k=1,2$ recovers the local polynomial
\[
 1-A_p(1)U+\tfrac12\big(A_p(1)^2-A_p(2)\big)U^2.
\]
The strong multiplicity-one theorem therefore recovers the global automorphic representation from its positive-ordinate multiset. For $E/\Q$, this determines $a_p=\sqrt p\,A_p(1)$ and hence $\#E(\F_p)$ at every good prime. Faltings's isogeny theorem gives the rational isogeny class. These are statements about the full multiset. The estimate keeps $p,k$ fixed and does not turn a single finite list of approximate ordinates into an exact recovery procedure for every prime.

\begin{corollary}[Exact point counts from finitely many approximate ordinates]\label{qt:inverse}
Fix $E/\Q$. For every sufficiently large $T$ depending on $E$, suppose the list of its positive completed-zero ordinates through $T$ is given with multiplicities, and each entry $\gamma$ is replaced by $\widehat\gamma$ with $|\widehat\gamma-\gamma|\le(\log T)^{-2}$. Simultaneously for every good prime $p\le\log T/(\log\log T)^2$, the nearest integer to
\[
 -\frac{2\pi p}{T\log p}\sum_{0<\gamma\le T}
            \cos(\widehat\gamma\log p)
\]
is $a_p=p+1-\#E(\F_p)$. The same assertion holds up to $p\le c_E\log T/\log\log T$ for a sufficiently small fixed $c_E>0$. Completeness of the list, including the multiplicities, is part of the hypothesis. No location hypothesis on the real parts is imposed.
\end{corollary}
\begin{proof}
Use unitary coordinates and put $v=|\operatorname{Re}\rho-1/2|$, $b=\log p$ and $c=1/8$. The uniform one-sided Landau formula of \href{https://www.cs.uleth.ca/~akbary/uniform.pdf}{Akbary--Murty, Proposition 15}, with exponent $\theta=1/4$ admissible for this tempered form, gives
\[
 \sum_{0<\gamma\le T}p^{\rho-1/2}
   =-\frac{a_p\log p}{2\pi p}T+O_E(p^{3/4}\log T).
\]
The finitely many real zeros, if included in the source convention, cost $O_E(\sqrt p)$ after normalization. For $b\le c\log T/2$, the density estimate yields
\[
 \sum_{0<\gamma\le T}v^2e^{bv}
   \ll_E T\log T\int_0^{1/2}(2u+bu^2)e^{-(c\log T-b)u}\,du
   \ll_E T/\log T.
\]
Since $\cosh(bv)-1\ll b^2v^2e^{bv}$, horizontal reflection replaces the Landau sum by $\sum e^{i\gamma\log p}$ at a cost $O_E(Tb^2/\log T)$. Taking real parts, multiplying by $-2\pi p/(T\log p)$ and perturbing the ordinates, the total error in $a_p$ is
\[
 O_E\!\left(\frac{p^{7/4}\log T}{T\log p}
     +\frac{p\log p}{\log T}+\frac{p}{\log T}\right).
\]
For the last term use $|\cos(b\widehat\gamma)-\cos(b\gamma)|\le b(\log T)^{-2}$ and $N_E^+(T)\ll_E T\log T$. The error tends uniformly to zero in the first stated range. In the second range its limit superior is at most a form-dependent constant times $c_E$, which can be chosen below $1/2$ after multiplication by that constant. The exact point count then follows by integrality.
\end{proof}

For a holomorphic newform $f$ of weight $k\ge1$ with known coefficient field $K$ of degree $d$, the rounded statistic has a corresponding algebraic version. Replace the cosine by $e^{i\widehat\gamma\log p}$ and $p$ in the prefactor by $p^{k/2}$. This approximates the algebraic integer $a_f(p)=p^{(k-1)/2}\lambda_f(p)$. The proof multiplies the unitary error by $p^{(k-1)/2}$ in lieu of $\sqrt p$. In the polylogarithmic prime ranges below this costs $O_f(p^{k/2}\log(2p)/\log T)$ including the perturbation. Restrict the candidates $a\in\mathcal O_K$ to the finite set satisfying $|\sigma(a)|\le2p^{(k-1)/2}$ for every embedding $\sigma$ of $K$. The true coefficient satisfies that condition by Ramanujan for all conjugate forms. If $a\ne a'$ are two candidates, the absolute value of the nonzero integer norm of $a-a'$ is at least one. Bounding the other $d-1$ factors gives
\[
 |a-a'|\ge (4p^{(k-1)/2})^{1-d}
\]
in the observed embedding. Hence a single positive-ordinate list recovers all good-prime coefficients by nearest-candidate selection throughout
\[
 p\le\left(\frac{\log T}{(\log\log T)^2}\right)^{2/(d(k-1)+1)}
\]
for all sufficiently large $T$ depending on $f$. The error divided by the separation tends uniformly to zero in this range. If the complete lists are supplied for every coefficient-field embedding, simultaneous recovery holds up to the exponent $2/k$ in the same display. The error in every embedding then tends to zero, and two different algebraic integers cannot differ by less than one in all embeddings. This argument includes rational-coefficient forms of higher weight and makes explicit the input needed for the coefficient-field improvement.

\begin{theorem}[Prime recovery from a bounded rectangle]\label{qt:nearlinear}
Fix a primitive holomorphic cuspidal newform $f$ of weight $k\ge1$, conductor $N$, and known coefficient field $K$ of degree $d$ in a specified complex embedding. For all sufficiently large $T$, the complete multiset of its unitary zeros $\rho$ with $0<\operatorname{Im}\rho\le T$ determines $a_f(p)$ for every good prime
\[
                         p\le c_{k,d}\,T/\log T.
\]
The same data determine every Fourier coefficient $a_f(n)$ for integers $1\le n\le c_{k,d}T/\log T$, including terms at ramified primes. The data may have complex coordinate errors of at most $T^{-2d(k+1)}$ for each listed zero. The procedure uses $k,N,K$ and this list alone. When the entire listed rectangle has been certified on the critical line, its ordinates supply the required coordinates. General lists of ordinates with unknown real parts instead satisfy the preceding polylogarithmic conclusion.
\end{theorem}
\begin{proof}
Write $b(p^j)=(\alpha_{p,1}^j+\alpha_{p,2}^j)\log p$ in unitary normalization and set $b(n)=0$ elsewhere. Ramanujan gives $|b(n)|\le 2\Lambda(n)$, including the factors of lower degree. Put $z_\rho=(\rho-1/2)/i$, $L=\log p$, $t_0=T/2$, $H=T/(C\sqrt{\log T})$. Use
\[
 h(z)=e^{-((z-t_0)/H)^2+iLz},\qquad
 \widehat h(u)=\sqrt\pi H\,e^{i(L-u)t_0-H^2(L-u)^2/4},
\]
with $\widehat h(u)=\int_{\mathbb R} h(t)e^{-iut}\,dt$. The Weil explicit formula in this convention reads
\begin{align*}
 \sum_\rho h(z_\rho)&=G(h)-\frac{1}{2\pi}\sum_{n\ge2}
 \frac{b(n)\widehat h(\log n)+\overline{b(n)}\widehat h(-\log n)}{\sqrt n},\\
 G(h)&=\frac1\pi\int_{\mathbb R} h(t)
  \left(\log\frac{\sqrt N}{2\pi}+\operatorname{Re}\psi(k/2+it)\right)dt ,
\end{align*}
where $\psi=\Gamma'/\Gamma$. Compare \href{https://www.numdam.org/item/PMIHES_2000__91__55_0.pdf}{Iwaniec--Luo--Sarnak, equation (4.6)}. For the asymmetric complex test here, integrate $h((s-1/2)/i)\Lambda'_u(s)/\Lambda_u(s)$ around the strip between $1-c$ and $c>1$. Take heights avoiding zeros by $\gg_f1/\log(|t|+2)$ on each unit interval. The Gaussian kills the horizontal integrals, using the local partial fraction expansion there. The functional equation converts the left integral to the conjugate form on the right. Moving the two gamma integrals to the central axis gives $G(h)$, since their sum has the displayed real gamma kernel and crosses no gamma poles. Inserting the two absolutely convergent Dirichlet series on the right yields the displayed Fourier terms. This derivation neither pairs positive with negative ordinates nor requires $h$ even.

Compute the complex number
\[
 B_p=-\frac{2\sqrt\pi\,p^{k/2}}{H\log p}
     \left(\sum_{0<\operatorname{Im}\rho\le T}h(z_\rho)-G(h)\right).
\]
The term $n=p$ from the first Fourier transform is exactly $a_f(p)$. With $T$ large, all omitted roots contribute to $B_p$ at most
\[
 \ll_f\frac{p^{(k+1)/2}\log T}{\log p}
                    \exp\!\left(-\frac{T^2}{8H^2}\right).
\]
Indeed $|h(z_\rho)|\le\sqrt p\,e^{1/(4H^2)}e^{-(\operatorname{Im}\rho-t_0)^2/H^2}$, and the number of roots in each unit interval centered at $v$ is $O_f(\log(2+|v|))$. Factor off $e^{-T^2/(8H^2)}$ on the omitted range and sum the remaining half Gaussian. This includes any roots of ordinate zero.

For $p/2\le n\le 2p$, $n\ne p$, we have $|\log(n/p)|\ge |n-p|/(2p)$. The total error from these $n$ in the first Fourier transform is
\[
 \ll p^{(k-1)/2}e^{-H^2/(16p^2)}
\]
when $H/p\ge 4$, by summing the resulting Gaussian over nonzero integers $n-p$. The remaining $n$ and the entire mirrored transform together cost
$O(p^{(k+1)/2}\log(2p)e^{-c_0 H^2})$
with an absolute $c_0>0$. For example, on $2^jp<n\le2^{j+1}p$, $j\ge1$, the sum before the $p^{k/2}/L$ normalization is
$O(\sqrt{2^jp}\log(2^{j+1}p)e^{-H^2j^2(\log2)^2/4})$. On $n<p/2$ the elementary sum $\sum_{n<p/2}\log n/\sqrt n\ll\sqrt p\log(2p)$ suffices. For the mirror use $\log(np)\ge\log 4$ and bound the second half of the Gaussian by $(np)^{-2}$ for large $H$.

Choose $C$ so that $C^2/8>dk+3$, and then $c_{k,d}>0$ so that $1/(16C^2c_{k,d}^2)>dk+3$. The combined error is $o((4p^{(k-1)/2})^{1-d})$, uniformly throughout the asserted range. Each allowed coordinate perturbation changes $h$ by $O(\sqrt p(1+\log p)T^{-2d(k+1)})$: differentiate on the straight segment using $|\operatorname{Im}z|\le1/2+T^{-2d(k+1)}$ and $|\operatorname{Re}z-t_0|\le T$. There are $O_f(T\log T)$ entries, so the normalized perturbation is $O_f(p^{(k+1)/2}(\log T)^{3/2}T^{-2d(k+1)})$, also negligible. The algebraic norm separation proved above identifies the unique nearest candidate. For $k=2$, $K=\Q$, this is ordinary integer rounding for the curve's Frobenius trace.
\end{proof}

To verify the assertion for all indices, center the Fourier transform at $\log m$ for $m=p^j$ and use $m^{k/2}/\log p$ in the prefactor. The main term is the arithmetic power sum $m^{(k-1)/2}A_p(j)$ at good and bad primes alike. The estimates just established with $p$ replaced by $m$ gain at most a factor $\log m/\log p\le (\log T)/\log 2$, which is absorbed by the exponent margins in the choices of $C,c_{k,d}$. Each power sum is an algebraic integer all of whose conjugates are bounded by $2m^{(k-1)/2}$, so the same separation identifies it exactly. Expanding each reciprocal Euler polynomial and using multiplicativity gives every $a_f(n)$ through the stated bound.

\begin{theorem}[Determination of the conductor from a positive list]\label{qt:conductor}
Fix the holomorphic weight $k\ge1$. Suppose the complete positive ordinates of a primitive newform of conductor $N$ through $T$ are supplied with individual errors at most $0\le\delta\le1$. Their real parts are unrestricted. Set $t_0=T/2$, $H=T/(8\sqrt{\log T})$, $w(t)=e^{-((t-t_0)/H)^2}$, and
\[
 D=\frac{2\sqrt\pi}{H}\sum_{0<\gamma\le T}w(\widehat\gamma)
       +2\log(2\pi)-\frac{2}{\sqrt\pi H}
                    \int_{\mathbb R}w(t)\operatorname{Re}\psi(k/2+it)\,dt.
\]
Uniformly for the level and nebentypus,
\[
       |D-\log N|\ll_k \log(2NT)(H^{-2}+\delta/H+T^{-8}).
\]
When $N\le c_k T^2/(\log T)^2$ and $\delta\le1/H$, its value is the nearest integer to $e^D$, for sufficiently large $T$ and sufficiently small $c_k>0$. With the coarser errors $\delta\le(\log T)^{-2}$ this holds uniformly for $N\le c_k T\sqrt{\log T}$.
\end{theorem}
\begin{proof}
The local zero bound needed uniformly in $N$ is
\[
       \#\{\rho:|\operatorname{Im}\rho-v|\le1\}
                  \ll_k \log(2N(2+|v|)).
\]
Indeed the logarithmic derivative of the order-one Hadamard product gives
$\sum_\rho\operatorname{Re}(2+iv-\rho)^{-1}=\operatorname{Re}(\Lambda'_u/\Lambda_u)(2+iv)$. The real constant is fixed by taking any nonzero point on the central line: the logarithmic derivative there has real part zero, while the reciprocal distances cancel in real part by horizontal reflection. The series of these real parts converges absolutely. Every zero in the indicated band contributes at least $1/5$ on the line $\operatorname{Re}s=2$. The gamma term there is $O_k(\log(2+|v|))$ apart from $\log(\sqrt N/(2\pi))$, and the Euler logarithmic derivative is absolutely bounded by $2\sum_{n\ge2}(\log n)n^{-2}$.

Apply the preceding explicit formula to $w$, whose transform is $\sqrt\pi H e^{-iut_0-H^2u^2/4}$. The normalized Euler terms are $O(e^{-cH^2})$, by $|b(n)|\le2\log n$, $|\log n|\ge\log 2$ and Gaussian decay. The omitted roots have normalized sum $O_k(T^{-8}\log(2NT))$ by the unit-interval bound and the same half-Gaussian tail split as in Theorem~\ref{qt:nearlinear}.

At each included ordinate write $u=\operatorname{Re}\rho-1/2$. By reflection the real offsets $u$ and $-u$ have the same multiplicity. Their average contribution differs from $w(\gamma)$ by
\[
 \ll H^{-2}\bigl(1+((\gamma-t_0)/H)^2\bigr)
              e^{-((\gamma-t_0)/H)^2}
\]
by the second derivative bound on $w(\gamma+iv)$ for $|v|\le1/2$ and $H\ge1$. Replacing $\gamma$ by $\widehat\gamma$ adds
$O(\delta H^{-1}(1+|\gamma-t_0|/H)e^{-(\gamma-t_0)^2/(2H^2)})$
at each root by the first derivative bound. Summing both by unit intervals and dividing by $H$ gives the two claimed errors. The conductor factor in the explicit formula is $(H/\sqrt\pi)\log(\sqrt N/(2\pi))$, which gives the displayed sign and normalization. In each stated range, $N|D-\log N|$ is bounded by a $k$-dependent constant times $c_k+o(1)$. Choose $c_k$ small enough that $|e^D-N|<1/2$.
\end{proof}

\begin{theorem}[A uniform level range and recognition without level or weight]\label{qt:recognition}
Let $T$ be sufficiently large. At known weight $k$ the direct conductor recovery from positive ordinates is valid uniformly for
\[
 \begin{array}{c|c}
 \text{error in each ordinate}&\text{range for }N\\\hline
 \delta\le 1/T & N\le c_k T^2/\log T\\
 \delta\le (\log T)^{-2} & N\le c_k T\log T .
 \end{array}
\]
If complete complex zero coordinates are supplied with error at most $e^{-T/100}$, the range includes $N\le e^{T/200}$, for sufficiently large $T$ depending on $k$. Certified critical-line coordinates are included here.

There is also a single reconstruction procedure, with no input weight, level, nebentypus or coefficient field, which eventually returns each fixed primitive holomorphic cuspidal newform from a complete positive-ordinate list with errors at most $(\log T)^{-2}$. Its threshold may depend on the input form. For elliptic curves over $\Q$ it returns the rational isogeny class.
\end{theorem}
\begin{proof}
Take a fixed nonnegative $C^2$ function $v$, supported in $[1/3,2/3]$, of mass one. For instance normalize $(x-1/3)^3(2/3-x)^3$ on that interval, extending by zero. Put $v_T(x)=v(x/T)$, $g_T(x)=(\pi T)^{-1/2}e^{-x^2/T}$ and $W=v_T*g_T$, using its entire continuation. Now $\int W=T$ and $|\widehat W(u)|\le T e^{-Tu^2/4}$. Use the explicit formula again, with the level statistic
\[
 D_k=\frac{2\pi}{T}\sum_{0<\gamma\le T} W(\widehat\gamma)
       +2\log(2\pi)-\frac2T\int_{\mathbb R}W(t)\operatorname{Re}\psi(k/2+it)\,dt .
\]
The Euler sum after normalization is $O(e^{-T(\log2)^2/8})$. For $|\operatorname{Im}z|\le1$ and $0\le j\le2$, differentiating the compactly supported factor gives $|W^{(j)}(z)|\ll_j T^{-j}$. Outside $0<\operatorname{Re}z\le T$ one also has
\[
 |W(z)|\ll \sqrt T\, e^{-T/18}
  e^{-\operatorname{dist}(\operatorname{Re}z,[T/3,2T/3])^2/(2T)} .
\]
This follows by bounding the absolute value of the complex Gaussian kernel, noting that the distance to its source interval is at least $T/3$. Sum this bound over omitted zeros by the preceding unit-interval estimate. Their contribution after dividing by $T$ is $O_k(e^{-T/100}\log(2NT))$. Inside the interval, the second derivative bound, horizontal pairing and the first derivative bound for perturbations now give
\[
  |D_k-\log N|\ll_k \log(2NT)
            (T^{-2}+\delta/T+e^{-T/100}).
\]
This proves both ranges by rounding. With complex coordinates replace $W(\widehat\gamma)$ by $W((\widehat\rho-1/2)/i)$ and take the real part of the resulting statistic. The contribution $T^{-2}$ disappears entirely and the other bounds hold unchanged. At $N\le e^{T/200}$, multiplying the remaining error by $N$ still gives $o(1)$.

For the assertion without prescribed weight, fix the actual form, put $U=\sqrt{\log T}$ and construct $W$ now at scale $U$. Use those entries with $0<\widehat\gamma\le U$. A mistaken inclusion or exclusion at an endpoint by at most $(\log T)^{-2}$ costs exponentially little in $U$, since $W$ is exponentially small near both endpoints and the local count is $O_f(\log(2U))$. Stirling's formula, for each fixed $k$, gives uniformly on $U/4\le t\le3U/4$
\[
 \operatorname{Re}\psi(k/2+it)-\operatorname{Re}\psi(1+it)
              =\frac{k(k-2)}{8t^2}+O_k(t^{-4}).
\]
The weighted integral outside that subinterval costs $O_k(e^{-cU})$ after normalization, because $[U/3,2U/3]$ is at distance at least $U/12$ from the exterior of $[U/4,3U/4]$ and the gamma kernel grows logarithmically. Thus $D_2=\log N+O_f((\log U)/U^2)$ at this scale. First recover $N$ by rounding $e^{D_2}$.

To identify the smaller weight-dependent correction, use the unconditional degree-two density estimate in Corollary~\ref{ep:density}. Its second-moment consequence is
\[
 \sum_{0<\gamma\le U}|\operatorname{Re}\rho-1/2|^2\ll_f U/\log U .
\]
In the second-derivative pairing bound keep this squared displacement for each zero. After dividing by $U$, the resulting error is $O_f(1/(U^2\log U))$. The errors from perturbed ordinates, $O_f((\log T)^{-2}\log U/U)$, are also $o(U^{-2})$. Put $J_U=U^{-1}\int_{U/4}^{3U/4}W(t)t^{-2}\,dt\asymp U^{-2}$. With $N$ already recovered, the statistic
\[
                    \frac{D_2-\log N}{2J_U}
                   =\frac{k(k-2)}8+o_f(1)
\]
uses only the positive list. The values $k(k-2)/8$ increase strictly through integers $k\ge1$ and have minimum spacing $1/8$. The unique nearest value therefore gives the actual integer weight for all sufficiently large $T$. The finite modular-space identification below now proves the asserted global recovery. For elliptic curves use the modularity and isogeny theorems.
\end{proof}

The uniform conductor conclusions also apply to a fixed archimedean gamma type of a standard unitary cuspidal representation $\pi$ of $\mathrm{GL}_n$ over a fixed number field $F$. Write
$\Lambda(s)=Q_\pi^{s/2}\prod_j\Gamma_{\mathbb R}(s+\mu_j)L(s)$
and suppose the pole divisor is known, including any norm twist of the Dedekind zeta function. Subtract the test values at its poles from the zero sum. The archimedean integrand in the explicit formula is then
\[
 \tfrac12\log Q_\pi+
   \sum_j\operatorname{Re}\frac{\Gamma_{\mathbb R}'}{\Gamma_{\mathbb R}}
                           (1/2+\mu_j+it).
\]
This follows by the identical residue computation, since $\operatorname{Re}\mu_j>-1/2$ for cuspidal local parameters, and the dual integrand supplies the conjugate gamma terms. The finite unitary root bound $|\alpha_{\mathfrak p,j}|\le(\mathrm N\mathfrak p)^{1/2}$ gives $|b(m)|\le n[F:\mathbb Q]\Lambda(m)m^{1/2}$ after collecting absolute norms: at $m=p^a$ the relevant residue degrees sum to at most $[F:\mathbb Q]$. Thus the Gaussian suppresses every nonzero Euler frequency exponentially. The real Hadamard derivative at $2+it$, with the finite pole corrections, gives $O_{F,n,\mu}(\log(2Q_\pi(2+|t|)))$ zeros per unit interval. These were all the estimates used for the uniform conductor ranges. Accordingly the same rounding recovers the absolute integer $Q_\pi=|D_F|^n\mathrm N\mathfrak q_\pi$, and hence the ideal norm $\mathrm N\mathfrak q_\pi$. Distinct ideals can have that same norm. The constants here may depend on $F,n$ and the fixed gamma parameters.

The direct prime-by-prime range can be distinguished from global recognition in a fixed modular space. For fixed known weight $k$, even the positive-ordinate data of Corollary~\ref{qt:inverse}, with real parts unspecified, eventually recover the \emph{whole} newform. For any fixed input form, Theorem~\ref{qt:conductor} first recovers its previously unknown conductor $N$. Set
\[
 B=\left\lfloor \frac{k}{12}[\mathrm{SL}_2(\mathbb Z):\Gamma_1(N)]\right\rfloor.
\]
For each of the finitely many $p^j\le B$, including bad primes, the fixed-frequency one-sided formula above recovers the arithmetic power sum $u_j=p^{j(k-1)/2}A_p(j)$ to an error tending to zero. The prescribed ordinate perturbations still cost $O_{p,j,f}(1/\log T)$. The finitely many normalized newforms of weight $k$ on $\Gamma_1(N)$ supply finitely many candidate power sums in the observed embedding. Their distinct values have positive separation, so the sufficiently accurate approximation identifies each one without knowing the coefficient field. Writing $c_j=a_f(p^j)$, the local identity between $L$ and its logarithmic derivative yields $jc_j=\sum_{r=1}^j u_r c_{j-r}$, $c_0=1$. Multiplicativity now determines all coefficients through $B$. The characteristic-zero Sturm bound on $\Gamma_1(N)$, or directly its valence formula, makes these coefficients an injective coordinate system on the space: a nonzero form cannot vanish at the cusp to order greater than $k[\mathrm{SL}_2(\mathbb Z):\Gamma_1(N)]/12$. Thus the form and all its remaining coefficients are determined as an element of the known space. This deduction makes a modular-space computation after determining the level; the direct local estimators use neither enumeration nor recognition of eigenforms. See \href{https://www.williamstein.org/books/modform/stein-modform.pdf}{Stein, \emph{Modular Forms, a Computational Approach}, Section 9.4.1, on Sturm's theorem}.

For a weight-zero cuspidal Hecke--Maa\ss\ newform the same limiting statistic extracts its Laplace eigenvalue. Suppose its real gamma factor is
$\Gamma_{\mathbb R}(s+e+ir)\Gamma_{\mathbb R}(s+e-ir)$, $e\in\{0,1\}$, $r\in\mathbb R$, and put $\lambda=1/4+r^2$. These are the usual tempered-at-infinity factors in weight zero. Write $a=1/4+e/2$. Its gamma kernel, including the $\pi$-normalization, is
\[
 -\log\pi+\tfrac12\sum_{\pm}\operatorname{Re}\psi(a+i(t\pm r)/2).
\]
The average of the logarithmic terms is $\log(t/2)-r^2/(2t^2)+O_r(t^{-4})$. The next coefficient is $4(a(a-1)/2+1/12)=-1/24$ at either parity. Subtracting the weight-two kernel $-\log(2\pi)+\operatorname{Re}\psi(1+it)$ leaves
$-\lambda/(2t^2)+O_r(t^{-4})$. The same zero-density and bump estimates apply. Thus the recovered quantity $(D_2-\log N)/(2J_U)$ tends to $-\lambda/2$. For a holomorphic form it likewise equals minus half the Casimir eigenvalue $k(2-k)/4$ in the limit. See \href{https://doi.org/10.1007/s40993-022-00393-y}{A. Seymour--Howell, \emph{Rigorous computation of Maass cusp forms of squarefree level}, Section 5.4}, for the weight-zero conventions.

Once the integer level has been determined, the cuspidal Laplacian on $\Gamma_1(N)\backslash\mathbb H$ has discrete spectrum of finite multiplicity and finitely many eigenvalues in any bounded range. Hence one can use the finitely many distinct spectral values near the limiting estimate, including the finitely many holomorphic weights in that range, as candidates. For each fixed form the closest candidate is eventually its exact Casimir eigenvalue. Its candidate Hecke eigenlines form a finite set, with the primitive normalization. Strong multiplicity one distinguishes any two representations in this set at some prime, so finitely many prime coefficients distinguish all pairs. The fixed-prime positive-ordinate Landau estimates above recover their values to arbitrarily small error, which identifies the line by finite separation. This Maa\ss\ conclusion is identification relative to the discrete cuspidal spectrum; the finite numerical enumeration given here is for the holomorphic level-$11$ example. See \href{https://sites.math.rutgers.edu/~sdmiller/sl3/sl3spectrum.pdf}{S.~Miller, \emph{On the existence and temperedness of cusp forms}, Theorem 1.1, on Selberg's Weyl law}.

An unrestricted unitary norm twist introduces a continuous unknown. In the convention $L_{f_\tau}(s)=L_f(s+i\tau)$ its ordinates shift by $-\tau$. For an endpoint-free finite interval and any positive input tolerance, sufficiently small changes of $\tau$ preserve the interval membership and fit the very same approximate data. Exact tolerant recovery of that real parameter is consequently impossible. Its limiting value can nevertheless be read from the two smoothed spectral averages. In the compact test at scale $U$, define
$J_{1,U}=U^{-1}\int_{U/4}^{3U/4}W(t)t^{-1}\,dt$ and $J_{2,U}=J_U$. Replacing $t$ by $t+\tau$ in the archimedean expansion gives
\[
 (D_2-\log N)/2=\tau J_{1,U}+(b-\tau^2/2)J_{2,U}+o_f(U^{-2}),
\]
where $b=k(k-2)/8$ or $-\lambda/2$ for the untwisted form. The norm twist has the same conductor and the same horizontal second-moment bound. Initially $D_2\to\log N$ still recovers the integer level. Form the two displayed equations at $U$ and $2U$. Since $J_{j,U}\sim U^{-j}\int v(x)x^{-j}dx$ for $j=1,2$, the $2$-by-$2$ determinant is asymptotic to a nonzero constant times $U^{-3}$. Solving gives $\tau+o_f(U^{-1})$ and $b-\tau^2/2+o_f(1)$, respectively. The underlying discrete spectral value is consequently recovered as before. At each fixed prime dephasing by $p^{i\widehat\tau}$ with this consistent estimate preserves convergence to the untwisted coefficient. This identifies the discrete base representation while retaining the necessary distinction between convergence and exact finite identification of the continuous twist.

\subsection*{An effective inverse computation at height 4096}
For the level-$11$ curve, put $T=4096$, $H=640$, $L=\log p$, $Q=\sqrt{11}/(2\pi)$, and take $h(t)=e^{-t^2/H^2}\cos(Lt)$. By self-duality and the certified nonzero center, the finite spectral sum is $S=2\sum_{0<\gamma\le T}h(\gamma)$, with $8708$ terms before the factor two. Write $C_p=-2\sqrt\pi p S/(HL)$. The following explicit bounds use no Frobenius trace from the unknown coefficient list.

The completed function is even about $1/2$. Pairing opposite centered zeros in its order-one Hadamard product gives
\[
 \sum_\rho\frac1{4+\gamma^2}
 \le\sum_\rho\operatorname{Re}\frac1{2-\rho}
 =\frac{\Lambda'_u}{\Lambda_u}(2)
 \le \log Q+\psi(5/2)+2\sum_{n\ge2}\frac{\log n}{n^2}<5.
\]
The first bound uses $1\le 2-\operatorname{Re}\rho\le2$. For the last inequality use $Q<1$, $\psi(5/2)<1$, $\log n\le\sqrt n$, and $\sum_{n\ge2}n^{-3/2}<2$. The omitted spectral sum has modulus at most
\[
 Z_p=5\sqrt p(4+T^2)\exp\{-T^2/H^2+1/(4H^2)\},
\]
since $(4+t^2)e^{-t^2/H^2}$ is decreasing for $t\ge T$ and
$|\cos(Lz_\rho)|\le \sqrt p$. This bounds every unknown higher zero, including any off the line.

Now $\widehat h(x)=\tfrac12\sqrt\pi H(e^{-H^2(x-L)^2/4}+e^{-H^2(x+L)^2/4})$ is nonnegative, even, and has mass $\pi$ on $x\ge0$. The elementary digamma integral gives
\[
 G(h)=\frac1\pi\left[(\log Q-\gamma_{\!E})\widehat h(0)
 +\int_0^\infty\frac{\widehat h(0)-\widehat h(x)}{e^x-1}\,dx\right].
\]
Split at $x=L/2$. On $0\le x\le L/2$ the derivative is at most $(3/4)\sqrt\pi H^3L e^{-H^2L^2/16}$ in modulus and $x/(e^x-1)\le1$. On the other part use the nonnegative mass and integrate $(e^x-1)^{-1}$ exactly for the constant term. Consequently
\[
 |G(h)|\le D_p=\frac1{\sqrt p-1}+
  \frac{\widehat h(0)(2-\log(1-p^{-1/2}))
           +\sqrt\pi H^3L^2 e^{-H^2L^2/16}}{\pi}.
\]
The uniformly small mirror and remote prime-power contributions, before multiplication by $p/L$, have absolute sum at most
\[
 F_p=4p^3 e^{-H^2(\log4)^2/8}
        +4p^{-3} e^{-H^2(\log(2p))^2/8}.
\]
For $n>4p$, half of the Gaussian is bounded by $(n/p)^{-3}$. For the mirror use $(np)^{-3}$ on $n\ge2$. Multiply by $2\log n/\sqrt n$ and sum, using $\sum_{n\ge2}(\log n)n^{-7/2}<2$.

By the explicit formula, a rigorous radius about $C_p$ for $a_p$ is therefore
\[
 \frac{2\sqrt\pi p}{HL}(Z_p+D_p)
 +\frac pL\left[F_p+\sum_{\substack{2\le n\le4p\\n\ne p}}
  \frac{2\Lambda(n)}{\sqrt n}e^{-H^2\log^2(n/p)/4}\right].
\]
Refine each certified root bracket to at most eight units on the certificate's fine grid, rechecking both endpoint signs by the interpolation bound in Section~\ref{num:chapter}. In computing $S$, the sum of the ordinate radii incurs at most $2(L+1/H)$ times that sum by differentiation. The resulting enclosures contain exactly one integer for $108$ of the $167$ good primes below $1000$, including every good prime through $251$. At each of these $108$ primes this integer also equals the unique nearest integer to $C_p$. Some representative outward-rounded enclosures are
\[
\begin{array}{r|r|r|r}
p&\text{lower bound}&\text{upper bound}&a_p\\\hline
2&-2.0227&-1.9453&-2\\
97&-7.0150&-6.9825&-7\\
251&-23.0236&-22.9744&-23\\
797&52.9525&53.0491&53\\
919&9.9277&10.0832&10\\
929&-30.0815&-29.9083&-30
\end{array}
\]
The two bounds $Q<1$ and $|\log Q-\gamma_{\!E}|<2$ hold for every integer conductor from $4$ through $39$. Thus this single-frequency arithmetic inference even works with that range in place of knowledge of the exact level. The finite inversion uses a zero-only input file. The program computes independent finite-field point counts after making its determinations and checks all $167$ enclosures, including the wider ones.

More traces can be identified by solving for closely spaced frequencies together. For a finite block $B$ of tested primes put
\[
 M_{p,q}=\frac{p\log q}{q\log p}\,
      e^{-H^2\log^2(q/p)/4}\,\mathbf1_{q\le4p},\qquad p,q\in B.
\]
With the same $C_p$, subtract from it every contribution of a known prime power $n=q^j\le4p$, apart from the first powers of the primes in $B$. Extend the uncertainty radius by the Hasse bound for each remaining unknown term. Once $a_q$ has been recovered, the numerator of its $j$th power sum satisfies $u_0=2$, $u_1=a_q$, $u_j=a_q u_{j-1}-q u_{j-2}$ at good primes. Its contribution is $(p/\log p)(u_j\log q/q^j)e^{-H^2\log^2(q^j/p)/4}$. At the multiplicative prime $11$ the numerator is $a_{11}^j$. Its first coefficient is inferred by including $11$ in the same procedure. Thus the resulting interval vector contains $M(a_q)_{q\in B}$. Interval matrix inversion, when successful, encloses that integer vector coordinate by coordinate. This argument accounts also for all unlisted primes larger than the maximum tested prime.

Taking all primes through $1032$ in increasing order, start a new block whenever $H\log(q/p)>5$ for consecutive tested primes $p<q$. The largest block has four members. With the identical zero list, the interval calculation now determines \emph{all $167$ good-prime traces through $1000$}. For instance at $p=997$ the block enclosure is $[37.9431,38.0584]$, giving $a_{997}=38$ and $\#E(\mathbb F_{997})=960$. These coupled enclosures also pass the separate point-count check after reconstruction.

For the larger cutoff $4096$, refine the Fourier grid a further $16$ binary places, again enclosing all $8708$ roots in intervals of at most eight fine units. In computing the Gaussian sum one can use the derivative bound $(L+2b/H^2)e^{-a^2/H^2}$ on each positive ordinate interval $[a,b]$. The archimedean term admits a more precise elementary correction. Put $g(x)=(e^x-1)^{-1}$. The term $G(h)+g(L)$ has absolute value at most
\begin{align*}
 &\frac1\pi\bigg[\widehat h(0)\{2-\log(1-p^{-1/2})\}
       +\sqrt\pi H^3L^2e^{-H^2L^2/16} \\
 &\qquad +(g(L)+L|g'(L)|)\tfrac L2\sqrt\pi H e^{-H^2L^2/16}
       +\frac{2\sqrt\pi}{H}|g'(L)| e^{-H^2 L^2/4}\bigg]
       +\frac{\sqrt p(\sqrt p+1)}{H^2(\sqrt p-1)^3}.
\end{align*}
To see this, approximate $g(x)$ on $x\ge L/2$ by its tangent at $L$. Its second derivative is positive and at most $\sqrt p(\sqrt p+1)/(\sqrt p-1)^3$ there. The normalized density $\widehat h(x)/\pi$ on $x\ge0$ is the law of $|L+Z|$, where $Z$ has density $He^{-H^2z^2/4}/(2\sqrt\pi)$. Thus its second centered moment about $L$ is at most $2/H^2$. The first centered moment is at most $2e^{-H^2 L^2/4}/(\sqrt\pi H)$ by integrating the Gaussian tail that folds across zero. The interval $x<L/2$ contributes the derivative bound used above and at most the displayed exponential term for the tangent. These estimates give the stated correction. Finally, any uniquely determined integer from one interval solve can be moved to the known side of the equations in the next round. At every round the right-hand sides still enclose the exact residuals. Using this procedure, \emph{all $563$ good-prime traces through $4096$} are identified from ordinates through that same height. It gives $a_{4091}=92$, $a_{4093}=94$.

An additional direct computation reaches $8192$ from the same spectral height $4096$. Average $2J+1$ equal-width Gaussians centered at $jD$, $-J\le j\le J$, writing
\[
 w(t)=\frac1{2J+1}\sum_{j=-J}^J e^{-((t-jD)/H)^2}.
\]
Use $w(t)\cos(t\log p)$ as test. Each Fourier Gaussian acquires the real multiplier
\[
  A(u)=\frac1{2J+1}\sum_{j=-J}^J e^{-ijDu},\quad
   u=\log(n/p)\ \hbox{or}\ \log(np),\qquad |A(u)|\le1 .
\]
The block matrices above acquire this multiplier as well. All absolute Euler tails still apply. The omitted-zero envelope replaces $T^2$ in $e^{-T^2/H^2}$ by $(T-JD)^2$; for the parameters below, $(4+t^2)e^{-(t-JD)^2/H^2}$ is decreasing for $t\ge4096$. For the gamma term the tangent correction becomes $-w(0)g(L)$ because the signed half-line Fourier mass is $\pi w(0)$. Its second-moment and folded first-moment error bounds still hold: $A$ is real and even, and its modulus is at most one. In the small-$x$ derivative bound include the additional term $\sqrt\pi H JD L e^{-H^2L^2/16}$, from $|A'|\le JD$. These modifications justify exactly the same interval calculation for the translated tests.

Use $(H,J,D)=(640,0,520),(200,5,520),(80,21,160)$ in succession. Each pass keeps every uniquely identified integer from preceding systems. Intersections of rigorous intervals also give valid enclosures. The prime buffer extends to $9728$; its coefficients are unknowns throughout unless independently inferred by these equations. This identifies $1024$ of the $1027$ good-prime traces through $8192$. The only unresolved primes are $8161,8167,8171$, with respective integer candidate sets
\[
       \{-39,-38,-37\},\qquad
       \{121,122,123,124,125\},\qquad
       \{-29,-28,-27\}.
\]
The explicit formula also holds for a Gaussian centered in frequency at $\log x$ when $x$ is an arbitrary positive real. In the unshifted $H=640$ formula, subtract every proved arithmetic contribution, and bound the remaining frequencies outside these three by the same Hasse majorants. The resulting interval contains the exact linear combination of the three unknowns with weights
$(x\log p)/(p\log x)\,e^{-H^2\log^2(p/x)/4}$. At $x=8064,8066,8068,8070,8072,8074,8076$ the respective numbers of surviving triples are $43,21,15,8,5,3,1$. The sole triple is
\[
                      (a_{8161},a_{8167},a_{8171})=(-38,123,-28).
\]
The last equation alone gives a short proof of uniqueness. Divide it by the positive coefficient of $a_{8161}$. For $(a,b,c)=(a_{8161},a_{8167},a_{8171})$ the directed enclosures at $x=8076$ give
\[
 -15.7249\le a+\nu b+\xi c\le-15.6378,\qquad
 .1955\le\nu\le.1956,\qquad .0620\le\xi\le.0621 .
\]
Write $(a,b,c)=(-38,123,-28)+(u,v,w)$, where $|u|,|w|\le1$ and $|v|\le2$. As $-38+123\nu-28\xi\in[-15.6923,-15.6772]$, it follows that $|u+\nu v+\xi w|\le .0545$. If $u\ne0$ the left side is at least $1-2(.1956)-.0621=.5467$. If $u=0$ and $v\ne0$ it is at least $.1955-.0621=.1334$. Finally $u=v=0$, $w\ne0$ gives at least $.0620$. Hence $u=v=w=0$.

Thus all $1027$ good-prime traces through twice the certified zero height are directly determined, independently of the modular-space identification below. Every trace was checked against a finite-field point count after inference. The single-frequency test at $11$ gives $0.9915\le a_{11}\le1.0136$, hence $a_{11}=1$. Consequently these zero data recover every Fourier coefficient $a_1,\ldots,a_{8192}$. Indeed $a_1=1$ and, at each good prime, $a_{pm}=a_pa_m-p\,\mathbf1_{p\mid m}a_{m/p}$, interpreting the second term as zero for $p\nmid m$. At the bad prime $a_{11m}=a_{11}a_m$. As an independent arithmetic check, expand $q\{\prod_{j\ge1}(1-q^j)(1-q^{11j})\}^2$ through $q^{8192}$ by Euler's pentagonal number identity and exact integer polynomial multiplication. All $8192$ coefficients agree, including $a_{4096}=-64$ and $a_{8192}=128$.

The signed quartic example gives a direct test with genuinely complex coefficients. Use the $10807$ positive and $10806$ negative zeros of $f\otimes\chi$ through height $4096$, with
\[
       h(t)=e^{-t^2/100^2+i t\log p},\qquad
       C_p=-\frac{2\sqrt\pi p}{100\log p}
                   \sum_{|\gamma|\le4096}h(\gamma).
\]
The gamma kernel is even, so its integral against the imaginary part is zero, and the preceding uncorrected archimedean bound applies with $H=100$. Horizontal reflection implies $\sum_\rho\operatorname{Re}(2-\rho)^{-1}=\operatorname{Re}(\Lambda'_u/\Lambda_u)(2)$ for this twist too. Indeed the logarithmic derivative has real part zero at any nonzero point on the critical line by its constant phase there. In the order-one Hadamard expansion the corresponding sum of real parts of the reciprocal distances there is zero, since opposite horizontal distances cancel. This identifies the real constant in that expansion and yields the claimed identity at $2$. The bound for its right-hand side is now $7$, using conductor $275$ in $Q$. The same omitted-root and Euler estimates therefore apply, with $7$ replacing $5$ in $Z_p$. The root brackets from the separation file retain the original sign of the ordinate for $\chi$; only the negative side was reflected in plotting positive spacings. The two real coordinate intervals isolate the Gaussian integer $a_{f\otimes\chi}(p)$. At $p=2$ they are
\[
          [-0.2831,0.4836]+i[-2.1350,-1.8621],
\]
which contains exactly $-2i$. The same computation recovers $i$ at $3$, $-2i$ at $7$ and $-4i$ at $13$. The conjugate twist interchanges the two ordinate lists and conjugates every inferred coefficient. Independent point counts for $E$ followed by the character twist agree with the four recovered values.

The zero-frequency calculation also determines the conductor effectively at this height. Use $w$ as in Theorem~\ref{qt:conductor}, now with $t_0=2048$ and $H=300$. All listed zeros have their real parts certified, so no second-derivative replacement is needed. At the test point $2+it_0$, the Hadamard identity bounds the sum of $(4+(\gamma-t_0)^2)^{-1}$ by $b+\tfrac12\log N$, where
$b=-\log(2\pi)+\operatorname{Re}\psi(5/2+it_0)+4$. The omitted normalized spectral contribution is consequently bounded by
\[
 \frac{2\sqrt\pi}{H}(4+t_0^2)
 e^{-(t_0/H)^2+1/(4H^2)}(b+\tfrac12\log N).
\]
The normalized Euler sum is bounded by $16e^{-H^2(\log2)^2/8}$, by the same $n^{-3}$ majorant. Integrate the gamma average on $128\le t\le3968$ on a grid of spacing $1/4$, multiplying each midpoint digamma enclosure by the exact Gaussian interval mass. On $[a,b]$ there, the derivative bound $|\psi'(1+it)|\le\sum_{n\ge1}(n^2+t^2)^{-1}\le\pi/(2a)$ controls the midpoint error. The complementary tails are bounded by $2(4+\log(2+t_0)+2H/(2+t_0))e^{-1920^2/(2H^2)}$ in the normalized gamma average. This follows by the digamma series bound $|\operatorname{Re}\psi(1+it)|\le4+\log(2+|t|)$, taking half of the Gaussian outside and integrating the rest. At each positive zero bracket use the supremum bound for $|w'|$ to enclose its contribution. The error in $D$ is now bounded by $A+B\log N$ with explicit $A>0$, $0<B<1$, giving $e^{(D-A)/(1+B)}\le N\le e^{(D+A)/(1-B)}$. No incoming upper bound on the conductor is needed. The four positive lists separately give
\[
\begin{array}{c|c|r}
\text{list}& \text{enclosure for }N&\text{unique integer}\\\hline
f &[10.9894,11.0109]&11\\
f\otimes\chi^2 &[274.8355,275.1673]&275\\
f\otimes\chi &[274.7503,275.2541]&275\\
f\otimes\bar\chi &[274.9156,275.0841]&275
\end{array}
\]

\subsection*{Determining both weight and conductor from one finite spectrum}
For the finest positive list of $8708$ critical-line coordinates, there is an effective identification among primitive holomorphic newforms which requires neither weight nor conductor as input. Take $H=100$, $c_1=1024$, $c_2=3072$. Write
\[
 s(c)=\frac{\sqrt\pi}{H}\sum_{0<\gamma\le4096}e^{-(\gamma-c)^2/H^2},
 \qquad
 P_a(c)=\frac1{\sqrt\pi H}\int_{\mathbb R}
       e^{-(t-c)^2/H^2}\operatorname{Re}\psi(a+it)\,dt .
\]
In the explicit formula $s(c)$ equals $\log(\sqrt N/(2\pi))+P_{k/2}(c)$ up to the two tails. To bound these without prior knowledge of $N,k$, put $X=\tfrac12\log N+\log(1+k)$ and
\[
 q=\frac{\sqrt\pi}{H}(4+1024^2)e^{-1024^2/H^2+1/(4H^2)},
 \qquad e_0=8e^{-H^2(\log 2)^2/8}.
\]
The omitted part in each equation has modulus at most $q(X+\log 4100+5)+e_0$. For the spectral tail, evaluate the Hadamard sum at $2+ic$. The elementary bound $|\psi(z)-\log z|\le1/\operatorname{Re}z$ gives an upper bound $X+\log 4100+5$ for that sum, using Hasse at $\operatorname{Re}s=2$. The lower bound $P_{k/2}(c)\ge\log(1+k)-\log4-2$ follows from the same digamma estimate since $k\ge1$. Thus the first equation gives
\[
 X\le\frac{s(c_1)+\log(2\pi)+\log4+2+q(\log4100+5)+e_0}{1-q}<20.
\]
This makes both error bounds explicit absolute numbers.

For $a\ge1/2$, $P_a(c_2)-P_a(c_1)$ is strictly decreasing in $a$. Indeed trigamma reflection gives $\operatorname{Re}\psi'(1/2+it)=\pi^2/(2\cosh^2(\pi t))$. For $a>1/2$ its harmonic extension is convolution with the Poisson kernel at width $a-1/2$: the real part of the trigamma is bounded on the half-plane and tends to zero at infinity by Stirling. Both the boundary values and the Poisson kernel are positive, even and strictly decreasing on $t>0$. Their convolution shares these properties, as does the further convolution with the Gaussian. This follows, for example, by writing positive even decreasing functions as superpositions of centered interval indicators and integrating their intersection lengths; strict decrease follows from the strictly decreasing positive densities here. Hence $\partial_a P_a(c_2)<\partial_a P_a(c_1)$ for $c_2>c_1>0$.

The following intervals are outward-rounded, in units of $10^{-12}$. The middle row includes the errors from the data and the two unknown tails.
\[
\begin{array}{c|r|r}
 &\text{lower endpoint}&\text{upper endpoint}\\\hline
 P_{1/2}(c_2)-P_{1/2}(c_1)&1100748712408&1100748729642\\
 s(c_2)-s(c_1)\text{ with error}&1100748613334&1100748613336\\
 P_{3/2}(c_2)-P_{3/2}(c_1)&1100748281647&1100748298882
\end{array}
\]
For reproducibility the two gamma averages at each weight use composite Simpson integration on $[c-768,c+768]$ with spacing $1/8$. For a midpoint of any such panel the analytic function
$e^{-(z-c)^2/H^2}\{\psi(a+iz)+\psi(a-iz)\}/2$
has modulus below $16$ on the circle of radius $32$, for $a=1/2,3/2$. Indeed shift the digamma argument by $34$, bound the reciprocal correction by $34/224$, then use $|\psi(w)-\log w|\le1/\operatorname{Re}w$ and $|e^{-(z-c)^2/H^2}|\le e^{32^2/H^2}$. These bounds apply at every real point of the integration interval. Cauchy's formula bounds the fourth derivative by $4!\,16/32^4$ in the composite Simpson error. The discarded normalized gamma tail is at most $2(5+\log(2+c)+2H/(2+c))e^{-768^2/(2H^2)}$. This proves the comparisons to the indicated precision.

The strict inequalities and monotonicity leave exactly the integer weight $k=2$. Returning to the conductor equation with this weight and the same fine positive list gives $10.9978\le N\le11.0022$, hence $N=11$. Since $X_1(11)$ has genus one, $S_2(\Gamma_1(11))$ is one-dimensional, and normalization determines the newform $\eta(\tau)^2\eta(11\tau)^2$ uniquely. This also identifies the elliptic isogeny class. The finite input here includes prior certification that all its positive zeros up to $4096$ lie on the critical line. The limiting recognition theorem above handles lists where those real parts are unspecified. See \href{https://wstein.org/msri06/refs/Diamond-Im-Modular_forms_and_modular_curves.pdf}{Diamond--Im, \emph{Modular forms and modular curves}, the example of $X_1(11)$}.

\clearpage
\section{A Cauchy-integral certificate for a primitive degree-three function}\label{sq:cauchy}
Use the non-CM curve $E\colon y^2+y=x^3+x^2$ from Section~\ref{sp:twists}. Its symmetric square is a cuspidal standard representation of $\mathrm{GL}_3/\Q$ by the symmetric square theorem of Gelbart and Jacquet, or by Newton and Thorne's theorem used there. We give a different finite-height verification which uses the completed function on a line of absolute convergence. It does not require evaluating a degree-three inverse Mellin kernel.

\begin{theorem}\label{sq:height8}
The standard $L$-function of $\operatorname{Sym}^2(43.\mathrm a1)$ has exactly fourteen nontrivial zeros with $|\Im s|\le8$. Every one is simple and has real part $1/2$ in unitary coordinates. The positive ordinates lie one each in $(j/64,(j+1)/64)$ for
\[
                  j=113,183,248,318,375,392,447.
\]
The motivic central line in this example is $\Re s=3/2$.
\end{theorem}
For a good rational prime $p$, write $t=a_p(E)$ and take the motivic variable $T$. The inverse Euler polynomial is
\[
                   (1-pT)\bigl(1-(t^2-2p)T+p^2T^2\bigr).
\]
Its roots in the reciprocal convention are $p,\alpha^2,\bar\alpha^2$. At $43$ the inverse polynomial is $1-T$, the conductor exponent is two, and the local root number is one by Dummigan--Martin--Watkins, Proposition~4.1. At the real place the unitary parameter is $I_2\oplus\operatorname{sgn}$ because the angular exponents $2,-2$ are paired and the middle monomial changes sign under the other Weil-group component. The two local signs multiply to $i^3i=1$. This gives conductor $43^2$, sign $+1$, and gamma factor $\Gamma_{\mathbb C}(s+1)\Gamma_{\mathbb R}(s+1)$, with $\Gamma_{\mathbb R}(s)=\pi^{-s/2}\Gamma(s/2)$. Thus, up to a positive constant independent of $s$, the completion in unitary coordinates is
\begin{equation}\label{sq:completion}
 \Lambda_2(s)=\delta^{-s}\Gamma((s+1)/2)\Gamma(s+1) F_2(s),
 \qquad \delta=2\pi^{3/2}/43,\quad F_2(s)=\sum\lambda_n n^{-s}.
\end{equation}
The integer coefficients in motivic normalization are $n\lambda_n$ and satisfy $|\lambda_n|\le d_3(n)$ in unitary normalization. We have $\Lambda_2(s)=\Lambda_2(1-s)$, and $\Lambda_2$ is entire.

\subsection{Recovering a rectangle from the right line}
Put $\sigma=15$, $r=\sigma-1/2=29/2$, and let $w=s-1/2$. The Cauchy integral formula, with the left side reflected onto the right side, yields
\begin{equation}\label{sq:two-kernels}
 \Lambda_2(1/2+w)=\frac{1}{2\pi}\int_{\mathbb R}\Lambda_2(\sigma+iy)
 \left(\frac1{r+iy-w}+\frac1{r+iy+w}\right)dy.
\end{equation}
To verify it, integrate the completion divided by its Cauchy denominator around the rectangle with vertical sides $\Re s=\sigma$ and $1-\sigma$ and let the horizontal sides tend to infinity. The gamma factors decay exponentially at the two ends, and the finite part has at most polynomial growth in every fixed strip. On the left side change variables using the functional equation. Its negative orientation cancels the minus in the reflected Cauchy denominator. No residues occur except that at the evaluation point. The same formula with a minus between the kernels applies to an odd functional equation.

Let $b=1/12$, $J=1080$, $N=4096$, $K=150$, $s_j=\sigma+ijb$, $u_j=r+ijb$. Define explicit finite quantities
\[
 Z_j=\delta^{-s_j}\Gamma((s_j+1)/2)\Gamma(s_j+1)
        \sum_{n\le N}\lambda_n n^{-s_j},\qquad
 P(w)=\frac b\pi\sum_{k=0}^K w^{2k}\sum_{j=-J}^{J}\frac{Z_j}{u_j^{2k+1}}.
\]
The coefficients are real since $Z_{-j}=\overline{Z_j}$. Uniformly for $|\Re w|\le1/2$ and $|\Im w|\le8$ we claim
\begin{equation}\label{sq:uniform}
                     |\Lambda_2(1/2+w)-P(w)|<10^{-17}.
\end{equation}
This is the analytic truncation allowance. Arithmetic rounding in the evaluation of $P$ is accounted for separately by interval arithmetic.

First, rotation of Euler's integral for the gamma function through angle $\operatorname{sgn}(t)$ radians gives
\[
 |\Gamma(x+it)|\le(\cos 1)^{-x}\Gamma(x)e^{-|t|}
                      \le2^x\Gamma(x)e^{-|t|},\qquad x>0.
\]
Absolute convergence gives $|F_2(x+iy)|\le\zeta(x)^3<8$ for $x\ge2$. Since $\delta>1/4$, uniformly on $13\le x\le17$ we obtain
\[
 |\Lambda_2(x+iy)|\le8\cdot4^{17}2^{27}\Gamma(9)\Gamma(18)e^{-3|y|/2}
                       < 3\cdot10^{38}e^{-3|y|/2}.
\]
On $x=15$ the same expression specialized to $x$ is less than $10^{33}e^{-3|y|/2}$. Shift the $y$ variable in \eqref{sq:two-kernels} into the strip $|\Im y|\le2$. Both denominators have modulus at least $12$ there for the entire target rectangle. Integration of the last bound shows that the sum of the two boundary norms, including the factor $1/(2\pi)$, is less than $10^{39}$. The analytic trapezoidal theorem therefore bounds the discretization by $10^{39}/(e^{4\pi/b}-1)$. The omitted samples $|j|>J$ contribute at most
\[
        b\cdot10^{33} e^{-135}/(e^{3b/2}-1).
\]
The constant absorbs both symmetric tails and both Cauchy denominators, since on the real $y$-axis those denominators have modulus at least $14$.

On the retained right line the Dirichlet tail is bounded by $\zeta(2)^3 N^{-13}<8N^{-13}$, using $|\lambda_n|\le d_3(n)$. Put $C_0=\delta^{-15}\,7!\,15!$, a bound for the gamma prefactor on that line without any decay restriction. Its use in the finite sum gives a coefficient-truncation allowance $b(2J+1)C_0\,8N^{-13}$. Finally $|w|<9$ and $|u_j|\ge r$. The sum of the two geometric expansions of the Cauchy denominators contains only even powers. Truncation at $w^{2K}$ costs at most
\[
   b(2J+1)C_0\cdot 8\frac{(9/r)^{2K+2}}{1-(9/r)^2}.
\]
Here $8$ bounds the absolutely convergent finite Dirichlet polynomial on $\Re s=15$, and dropping the factor $1/(\pi r)$ only increases the estimate. The four displayed error allowances sum to less than $10^{-17}$. This proves \eqref{sq:uniform}.

\subsection{Argument change and signs}
The integer coefficients in $\sum_{n\le N}\lambda_n n^{-s_j}$ are computed from the Weierstrass equation by direct quadratic-residue counts. Interval arithmetic evaluates the real coefficients of $P$. For clarity, the enclosing algorithm uses a polynomial only about its displayed center. Each side of the counterclockwise contour with vertices $\pm1/2\pm8i$ is first partitioned into subsegments of length at most $1/2$. On one such segment let $z_m$ be its midpoint and $v$ its displacement, and express its image by the exactly composed polynomial $Q(x)=P(z_m+vx)=\sum_k c_kx^k$ on $|x|\le1/2$. An explicit curvature bound is
\[
                  |Q''(x)|\le \sum_{k\ge2}k(k-1)|c_k|2^{-(k-2)}.
\]
On any subinterval of length $l$ in this parameter, the linear interpolation error is at most $l^2/8$ times that curvature bound. Just as in the degree-four contour, accept a subinterval when its endpoints have a common signed component beyond the sum of the interpolation and analytic errors. Bisect any unresolved subinterval. Accepted images have their argument increments enclosed by the principal argument of the quotient of endpoint evaluations, and the polynomial and the exact completion are homotopic away from zero on the contour.

The interval computation accepts $136$ subintervals and encloses the winding around zero by the unique integer $14$. On the line $w=it$ the exact completion is real, and \eqref{sq:uniform} certifies the seven changes of sign in the positive intervals in Theorem~\ref{sq:height8}. The even functional equation supplies seven reflected zeros. All fourteen zeros allowed by the winding are exhausted by these distinct roots, so each has multiplicity one. The center is positive in normalization \eqref{sq:completion}, with an enclosure in the accompanying result. Outside the closed critical strip the Euler product and the functional equation exclude any further zeros at these heights. This proves Theorem~\ref{sq:height8}.

\subsection{Translated polynomials through height 32}\label{sq:height32}
Keeping the Cauchy expansion local in the ordinate gives a stronger finite theorem for the same function. All $96$ zeros with $|\Im s|\le32$ are simple and lie on the unitary critical line. The corresponding counts above the real axis are as follows.
\begin{center}
\begin{tabular}{rrr}
\toprule
Ordinate interval & Winding & Accepted contour intervals\\\midrule
$[0,8]$ & 7 & 74\\
$[8,16]$ & 12 & 96\\
$[16,24]$ & 14 & 122\\
$[24,32]$ & 15 & 140\\\bottomrule
\end{tabular}
\end{center}
The included height-32 result supplies all $48$ disjoint positive brackets, each of width $1/64$. Here are the changes needed to verify the stronger assertion. In \eqref{sq:two-kernels} use $\sigma=27$, $r=53/2$, $b=1/14$, $Y=220$, $J=Y/b$, $N=16384$. On a band with center $it_0$, $t_0=4,12,20,28$, write $w=it_0+\xi$ with $|\xi|<5$. Expand the first denominator at $u_j-it_0$ and the second at $u_j+it_0$ through degree $150$. Their moduli before adding $\xi$ are each at least $r$. If $B_k=\sum_{j=-J}^J Z_j/(u_j-it_0)^{k+1}$, then the actual polynomial coefficient is $b(B_k+(-1)^k\overline{B_k})/(2\pi)$. This follows by changing $j$ to $-j$ in the second sum. The translated polynomial takes real values along $\xi\in i\mathbb R$.

In the $y$-strip $|\Im y|\le4$ use the explicit gamma contour rotation above with $x\le31$ to bound the completed right-line values by $10^{83} e^{-3|\Re y|/2}$. On the sampling line they are bounded by $10^{68} e^{-3|y|/2}$. The normalizing factors from the two Cauchy denominators and their two strip-integrals are less than one. With $C_0=\delta^{-27}13!27!$, the four error allowances can be taken as
\begin{align*}
 &10^{83}/(e^{8\pi/b}-1),\qquad b\cdot10^{68}e^{-3Y/2}/(e^{3b/2}-1),\\
 &8b(2J+1)C_0 N^{-25},\qquad
 8b(2J+1)C_0(5/r)^{151}/(1-5/r).
\end{align*}
Their sum is less than $10^{-46}$. The checker evaluates the slightly sharper explicit gamma upper bound for $10^{68}$, and the displayed rounded bound also suffices. In each band the composed-polynomial verification is carried out on initial segments of length at most $1/4$ and certifies that no horizontal boundary contains a zero. Every winding equals the number of strictly certified critical-line sign changes, with simple roots consequently forced. Conjugation and the absence of zeros on the real segment give the full count. This second computation also recovers the seven positive zeros of the original height-eight calculation using different quadrature and truncation parameters.

For the larger height $64$ in the same translated-polynomial method the parameters are $\sigma=55$, $r=109/2$, $b=1/20$, $Y=450$, $J=9000$ and $N=16384$. Expand about the successive centers $4i,12i,\ldots,60i$ through degree $224$, and use strip half-width $6$ about the frequency contour. For odd positive $x$ write
\[
 C(x)=8\cdot4^x 2^{3(x+1)/2}((x-1)/2)!x!,\qquad
 C_0=\delta^{-55}27!55!.
\]
The same four error allowances, with the full strip and line prefactors, are
\begin{align*}
 &C(61)/(e^{12\pi/b}-1),\qquad b C(55)e^{-3Y/2}/(e^{3b/2}-1),\\
 &8b(2J+1)C_0N^{-53},\qquad
       8b(2J+1)C_0(5/r)^{225}/(1-5/r).
\end{align*}
The four bounds are respectively below $2.02\cdot10^{-146}$, $1.27\cdot10^{-133}$, $8.00\cdot10^{-87}$, and $7.73\cdot10^{-97}$. Their sum is less than $10^{-80}$. The right-line Dirichlet phases are recomputed directly at every thirty-second step, and intervening steps use complex-ball multiplication at $896$-bit precision. The tightened strip norm takes account of both Cauchy denominators, bounded below in modulus by $48$ throughout this strip.

\begin{theorem}[Symmetric square through height 64]\label{sq:height64}
The primitive degree-three standard $L$-function of $\operatorname{Sym}^2(43.\mathrm a1)$ has precisely $234$ nontrivial zeros with $|\Im s|\le64$. They are simple, lie on the unitary critical line, and form $117$ conjugate pairs.
\end{theorem}
\begin{proof}
The eight translated polynomials give the following verified counts. Every band starts at the previous upper endpoint, with the first starting at zero.
\begin{center}
\begin{tabular}{r|rrrrrrrr}
\toprule
Upper ordinate & 8&16&24&32&40&48&56&64\\
Winding & 7&12&14&15&16&17&18&18\\
Panels &74&96&122&140&140&144&148&148\\\bottomrule
\end{tabular}
\end{center}
The winding counts a total of $117$ zeros in the upper rectangles with multiplicity. The $117$ disjoint open intervals in the accompanying integer report each give a strict critical-line sign change and have width $1/64$. Thus each contains exactly one zero of multiplicity one. The contours are zero-free. In particular the lower boundary of the first rectangle excludes real zeros. Conjugation gives the other $117$ zeros, and the Euler product and functional equation exclude additional zeros outside the critical strip. All $48$ positive brackets at heights through $32$ agree with the differently truncated height-32 calculation.
\end{proof}

An audit with $b=1/21$, $Y=454$, $N=16352$, $928$-bit arithmetic, and phase resets every $24$ samples produced the same complete output. This includes every winding, every contour subdivision count, and each of the $117$ sign brackets. All four analytic error allowances were recomputed with these parameters.

Repeating the audit grid at $960$ bits and bisecting with certified signs puts each of the $117$ positive zeros into an open interval of width $2^{-24}$. The first and last lower endpoint numerators are $29867679$ and $1070267979$, respectively. The full list accompanies the certificate.

\subsection{Cauchy contours at larger ordinates}\label{sq:large}
A much larger ordinate makes the completed function exponentially small. There are two ways to control the arithmetic error in Cauchy reconstruction. The first retains the original completion and moves its integration line farther right. The second multiplies the completed function by a nowhere-vanishing exponential before applying the Cauchy integral formula. Both use the certified frequency gridding of Section~\ref{sf:grid}.

For the bands $64\le\Im s\le128$, retain the original completion and take
\[
 (\sigma,b,Y,N,d,K)=(97,1/25,860,32768,8,380),
\]
where $d$ denotes the Cauchy strip half-width and $K$ the local polynomial degree. With $A(x)=\delta^{-x}((x-1)/2)!\,x!$ and $C(x)=8\cdot2^{3(x+1)/2}A(x)$ for odd $x$, the four allowances are
\begin{align*}
 &C(105)/(e^{16\pi/b}-1),\qquad bC(97)e^{-3Y/2}/(e^{3b/2}-1),\\
 &8b(2J+1)A(97)N^{-95},\qquad
 8b(2J+1)A(97)(5/r)^{381}/(1-5/r).
\end{align*}
Here $r=\sigma-1/2$ and $J=Y/b$. These expressions are respectively below $10^{-199},10^{-245},2\cdot10^{-155},3\cdot10^{-216}$. Their sum is less than half of the uniform allowance $10^{-144}$. Each band has width eight and each polynomial is centered at its midpoint. The Fourier gridding encloses the exact retained finite Dirichlet sum, with the absolute gridding remainder included in the returned arithmetic intervals.

For bands above $128$ set $\theta=3/2$ and $H(w)=e^{-i\theta w}\Lambda_2(1/2+w)$. Its zeros and multiplicities are unchanged, and $H(it)$ has the same real sign as the original completion. In the Cauchy rectangle the exponential growth in the imaginary direction is strictly slower than the vertical decay of the gammas, so the horizontal integrals still tend to zero. Reflection gives the precise replacement for \eqref{sq:two-kernels}, with $u=r+iy$,
\[
 H(w)=\frac1{2\pi}\int_{\mathbb R}\Lambda_2(\sigma+iy)
 \left(\frac{e^{-i\theta u}}{u-w}+\frac{e^{i\theta u}}{u+w}\right)dy.
\]
If $Z_j$ is the earlier finite right-line approximation before tilting, the first local denominator uses the signed weights $W_j=e^{-i\theta r}e^{\theta jb}Z_j$. The second weight is $\overline{W_{-j}}$. Thus exactly the same real and imaginary parity projection works after summing over the signed first weights. The positive and negative entries of this first sum must be weighted separately.

Use $(\sigma,b,Y,N,d,K)=(99,1/24,2200,49152,12,365)$. Rotating the gamma integrals through one radian bounds the gamma prefactor times $e^{\theta |y|}$ on the sampling line by $A_*=2^{150}A(99)$. To retain exponential decay in the quadrature bounds, rotate through $7/5$ radians. Because $\cos(7/5)>1/6$, the weighted completed integrand on a line of real part $x$ has majorant $C_*(x)e^{-3|y|/5}$ with
\[
                 C_*(x)=8\cdot6^{3(x+1)/2}A(x).
\]
Use $C_*(111)$ throughout the widened strip. Its two Cauchy denominators have modulus at least $86$, and after their normalization the integrals of the decaying exponentials contribute less than one. The four allowances become
\begin{align*}
 &C_*(111)/(e^{24\pi/b}-1),\qquad bC_*(99)e^{-3Y/5}/(e^{3b/5}-1),\\
 &8b(2J+1)A_*N^{-97},\qquad 8b(2J+1)A_*(5/r)^{366}/(1-5/r).
\end{align*}
Their sum is below $10^{-110}/2$. The code may further discard an entry when its entire complex-ball modulus satisfies $|W_j|<10^{-110}/[64b(2J+1)]$. It simultaneously discards the conjugate second-kernel entry. On all target rectangles the sum of these omissions, including the factor $b/(2\pi)$ and both denominators, is smaller than $10^{-110}/64$. All verification uses the larger allowance $10^{-110}$.

\begin{theorem}[Symmetric square through height 128]\label{sq:height128}
The symmetric square of $43.\mathrm a1$ has exactly $554$ nontrivial zeros of absolute ordinate at most $128$. Each is simple and critical in unitary coordinates.
\end{theorem}
\begin{proof}
For clarity, the new positive-band windings and accepted panel counts from the unweighted line at $97$ are
\begin{center}
\begin{tabular}{r|rrrrrrrr}
\toprule
Upper ordinate & 72&80&88&96&104&112&120&128\\
Winding &20&19&19&20&19&21&21&21\\
Panels &164&160&172&168&184&182&178&180\\\bottomrule
\end{tabular}
\end{center}
Each winding exactly equals the number of disjoint strict positive sign brackets recorded for the band. No horizontal or vertical boundary contains a zero. The new bands add $160$ positive roots to the $117$ from Theorem~\ref{sq:height64}. Reflection gives all $554$ and multiplicity one is forced by the equality of counts.
\end{proof}

\paragraph{An independent smaller right line.}
Lemma~\ref{sp:taillemma} also permits a degree-three check of every positive band with the same fixed parameters. Tilt by $\theta=3/2$, evaluate the retained polynomial at $\sigma=31$, and move its omitted series alone to $\sigma'=99$. Use $b=1/18$, $Y=1100$, $N=40960$, strip width $d=8$, local degree $K=280$, and $896$-bit arithmetic. Taking $\phi=51/50$ in the tail lemma gives decay $c=3/100$, with $\cos\phi>1/2$. With the earlier $A(x)$ and $C_*(x)$, the coefficient-tail contribution becomes
\[
 \frac{16N^{-97}2^{150}A(99)}{\pi(3/100)98}.
\]
The other three bounds use $C_*(39)$ on the enlarged strip, $C_*(31)$ on the sampling line, and $2^{48}A(31)$ for the finite weighted Dirichlet polynomial. The logarithms to base ten of the trapezoidal, frequency-tail, coefficient-tail and Taylor errors are respectively below $-259.07,-183.94,-125.16,-137.68$. The total remains less than $10^{-110}/2$, with the same small-weight omission allowance as above. This run obtains prime traces by finite-group orders and exact fallback, as described in Section~\ref{sp:tailshift}. Thus its real line, prime computation, coefficient truncation, frequency spacing, and precision differ from the main upper-band sweep.

\begin{theorem}[Symmetric square through height 256]\label{sq:height256}
The primitive standard degree-three $L$-function of $\operatorname{Sym}^2(43.\mathrm a1)$ has exactly $1276$ nontrivial zeros with $|\Im s|\le256$. They are simple and form $638$ conjugate pairs on the critical line.
\end{theorem}
\begin{proof}
The smaller right-line Cauchy computation checks every eight-unit band from zero to $256$. Its windings separately agree with the number of strict critical-line sign changes in each band. Above $128$ the results are
\begin{center}
\begin{tabular}{r|rr|r}
\toprule
Range & \multicolumn{2}{c|}{Eight-unit positive windings} & Cumulative total\\\midrule
$[128,144]$&22&21&640\\
$[144,160]$&22&22&728\\
$[160,176]$&21&22&814\\
$[176,192]$&22&23&904\\
$[192,208]$&23&23&996\\
$[208,224]$&23&23&1088\\
$[224,240]$&23&24&1182\\
$[240,256]$&23&24&1276\\\bottomrule
\end{tabular}
\end{center}
All positive sign intervals have width $1/64$. The highest has endpoints $16381/64$ and $16382/64$. The new bands add $361$ distinct positive zeros to the $277$ positive zeros below $128$. Winding forces multiplicity one, and the contours exclude boundary zeros. Reflection gives the full count.
\end{proof}
The alternate far-right sweep also reaches $256$. The two complete computations agree on all $32$ positive-band windings and every one of the $638$ brackets. Their total numbers of accepted contour panels are $5546$ and $5422$, differing in fifteen bands.

For the signed bisections used in Section~\ref{sf:separation}, an additional square calculation uses working line $\sigma=33$, omitted-series line $\sigma'=101$, $N=43008$, spacing $b=1/19$, cutoff $Y=1150$, strip half-width $8$ and degree $K=290$. It obtains the prime traces by exact residue sums and evaluates the analytic part at $960$-bit precision. The tail bound from Lemma~\ref{sp:taillemma} is now
\[
           \frac{16N^{-99}2^{153}A(101)}{\pi(3/100)100}.
\]
The other bounds use $C_*(41)$ on the enlarged strip, $C_*(33)$ on the sampling line, and $2^{51}A(33)$ for the finite weighted polynomial. The base-ten logarithms of the four errors in the same order are below $-272.87,-189.23,-128.71,-147.24$. This again permits uniform allowance $10^{-110}$ and the preceding small-weight threshold. Starting at each strictly changing coarse bracket, retain the half on whose endpoints the real parts have opposite certified signs. This proves every accepted refined enclosure for the completed function itself.
The $\sigma=31$ bisections enclose all $638$ positive ordinates through $256$ in intervals of width $2^{-24}$. The highest has successive endpoint numerators $4294303699$ and $4294303700$. On the first $117$ positive zeros the endpoints agree with the independent low-height calculation.
The independent line-$33$ calculation returns precisely the same $638$ integer endpoint pairs. Its cumulative files at heights $32$, $64$, $128$ and $256$ reproduce all $48$, $117$, $277$ and $638$ positive enclosures, respectively.

\clearpage
\section{Forced central zeros in every even degree}\label{sp:twists}
Consider $E:y^2+y=x^3+x^2$ over $\Q$. Its invariants are $c_4=16$, $c_6=-280$ and $\Delta=-43$. In particular its conductor is $43$ and its only bad prime is multiplicative. The reduction there is nonsplit because $-c_6\equiv22\bmod43$ is a nonsquare, as one verifies by Euler's criterion. The $j$-invariant $-4096/43$ is not an algebraic integer, so $E$ has no complex multiplication. These data can also be compared with \href{https://www.lmfdb.org/EllipticCurve/Q/43/a/1}{LMFDB $43.\mathrm a1$}.

\begin{theorem}\label{sp:family}
Fix any even $d=2m\ge4$. For every prime $q\equiv1\bmod4$, $q\ne43$, let $\chi_q$ be its quadratic character and set $\Pi_q=\operatorname{Sym}^{d-1}(E)\otimes\chi_q$ in unitary normalization. This is a cuspidal standard $L$-function on $\mathrm{GL}_{d}/\Q$. Its arithmetic conductor $Q_q$ and its global sign $W_q$ are
\[
 Q_q=43^{d-1}q^d,\qquad W_q=(-1)^{m(m+1)/2}\left(\frac q{43}\right).
\]
Among all rational primes a set of natural density $1/4$ gives an odd-order zero at $s=1/2$. For fixed $d$ the number of these functions with $Q_q\le X$ is asymptotic to
\[
 \frac14\operatorname{Li}\left((X/43^{d-1})^{1/d}\right).
\]
In degree four the untwisted symmetric cube already has sign $-1$ and conductor $79507$.
\end{theorem}
\begin{proof}
Modularity over $\Q$ attaches a weight-two newform with trivial nebentypus to $E$. Newton and Thorne's \href{https://doi.org/10.1007/s10240-021-00126-4}{symmetric power theorem}, Theorem~A, gives cuspidality, regularity and the symmetric-power parameters at every place. Twisting by the unitary quadratic character preserves cuspidality. Since the symmetric exponent is odd, $\Pi_q$ is also the $(d-1)$st symmetric power of the quadratic twist of $E$. Its self-duality is of symplectic type. Its completed standard function is entire and satisfies its self-dual functional equation. We compute the sign in this functional equation.

At $43$ a unitary Steinberg parameter twisted by the nontrivial unramified quadratic character describes the Tate module. Its $(d-1)$st symmetric power is the single special block of dimension $d$, with the same unramified quadratic sign. Multiplication by $\chi_q$ changes the unramified character value to $-(q/43)$. A special block with unramified character value $a=\pm1$ has conductor exponent $d-1$ (the kernel of monodromy has dimension one) and root number $(-a)^{d-1}$. These formulas also follow from Dummigan, Martin and Watkins, \href{https://magma.maths.usyd.edu.au/~watkins/papers/ef10.pdf}{\emph{Euler factors and local root numbers for symmetric powers of elliptic curves}}, Proposition~4.1. The local sign is therefore $(q/43)$.

At $q$, before twisting, the representation is unramified of determinant one and has zero monodromy. Upon twisting, every one of the $d$ lines in its semisimplified Weil parameter acquires the tame quadratic inertia character. Thus the conductor exponent is $d$ and its local Euler factor is one. The product formula for epsilon factors of unramified twists gives the local root number $W(\chi_{q,v})^d$, since all unramified character values multiply to one. The quadratic Gauss sum identity gives $W(\chi_{q,v})^2=\chi_{q,v}(-1)=1$ as $q\equiv1\bmod4$. Hence the local root number is one. All other finite places are unramified and contribute sign one. We use the standard additive character of conductor zero at finite primes.

At the real place write $I_j=\operatorname{Ind}_{W_\C}^{W_\R}(z/|z|)^j$. The parameter $I_1$ of $E$ has, on $W_\C$, angular exponents $1$ and $-1$. In the odd symmetric power the positive exponents are $1,3,\ldots,2m-1$, with their negatives paired by complex conjugation. This gives one $I_{2j-1}$ for each $1\le j\le m$. A sign twist of such an induced representation is isomorphic to itself. The local root number of $I_{2j-1}$ is $i^{2j}=(-1)^j$, and its gamma factor in unitary coordinates is $\Gamma_\C(s+(2j-1)/2)$ with $\Gamma_\C(s)=2(2\pi)^{-s}\Gamma(s)$. Multiplying gives the stated archimedean sign. The character $\chi_q$ is even. This establishes the completed conductor and sign, including the archimedean contribution.

When $W_q=-1$ the holomorphic function $z\mapsto\Lambda(\Pi_q,1/2+z)$ is odd. It is not identically zero by the Euler product. Thus it has a zero of finite odd order at the center. Among the $\varphi(172)$ reduced residues modulo $172$, exactly one quarter have $q\equiv1\bmod4$ and the required Legendre sign modulo $43$. The prime number theorem in these fixed progressions gives the asymptotic. Distinct $q$ give distinct arithmetic conductors. Omitting the twist gives finite sign $+1$ at $43$ and the stated archimedean sign for $m=2$.
\end{proof}

For example $q=5$ has $(5/43)=-1$, so the degree-six function $L(\operatorname{Sym}^5 E\otimes\chi_5,s)$ has a central zero with conductor $43^5 5^6$. A fixed $q\equiv1\bmod4$ with $(q/43)=1$ yields central zeros simultaneously in every degree $d\equiv2,4\bmod8$ in this family, and the other Legendre sign yields them in every degree $d\equiv0,6\bmod8$.

\section{A Bessel certificate for the symmetric cube of 43.a1}\label{sp:cube}
Martin and Watkins studied central values and functional equations of symmetric-power $L$-functions numerically, and Booker developed rigorous methods for certifying zeros of higher-degree $L$-functions. We give the finite coefficient, special-function and contour bounds for this curve. See \href{https://magma.maths.usyd.edu.au/~watkins/papers/antsVII.pdf}{Martin and Watkins, \emph{Symmetric powers of elliptic curve $L$-functions}}, and \href{https://arxiv.org/abs/math/0507502}{Booker, \emph{Artin's conjecture, Turing's method and the Riemann hypothesis}}.

\begin{theorem}[A finite-height theorem in degree four]\label{sp:height8}
For the primitive standard $L$-function of $\operatorname{Sym}^{3}(43.\mathrm a1)$, every nontrivial zero $\rho$ with $|\Im\rho|\le8$ is simple and satisfies $\Re\rho=1/2$ in unitary coordinates. There are exactly $23$ such zeros. The noncentral ordinates occur in opposite pairs. The positive ordinates lie one each in the intervals
\[
 \frac{j}{64}<\Im\rho<\frac{j+1}{64},\qquad
 j\in\{76,162,209,224,279,309,336,414,437,480,500\}.
\]
\end{theorem}

Write $F(s)=L(\operatorname{Sym}^3 E,s)$ in unitary coordinates and normalize its completion by
\[
 \Lambda(s)=\delta^{-s}\Gamma(s+1/2)\Gamma(s+3/2)F(s),
 \qquad\delta=\frac{(2\pi)^2}{\sqrt{43^3}},\qquad c=2\sqrt\delta.
\]
Thus $\Lambda(s)=-\Lambda(1-s)$. The center in motivic coordinates is $2$. Write $F(s)=\sum_{n\ge1}\lambda_n n^{-s}$ on $\Re s>1$. At good $p$ let $a_p=p+1-\#E(\F_p)$, and put $A_n=n^{3/2}\lambda_n$. The inverse Euler polynomial in motivic coordinates is
\begin{equation}\label{sp:polynomial}
 (1-pa_pT+p^3T^2)\{1-(a_p^3-3pa_p)T+p^3T^2\}.
\end{equation}
Indeed its reciprocal roots are $p\alpha,p\bar\alpha,\alpha^3,\bar\alpha^3$ if $\alpha+\bar\alpha=a_p$ and $\alpha\bar\alpha=p$. At $43$ its inverse polynomial is $1+T$. These formulas determine all $A_n$ by inversion of power series and multiplicativity. The normalized good roots have modulus one by Hasse's theorem, and the bad root has absolute value at most one, so $|\lambda_n|\le d_4(n)\le n^3$. The last inequality bounds the number of possible choices of the first three entries in an ordered factorization of $n$.

\subsection{Central simplicity from two positive Bessel weights}
Let $K_j$ denote the modified Bessel function of the second kind, and put $k(y)=2yK_1(2\sqrt y)$, $v_n=c\sqrt n$. The Bessel Mellin integral gives
\[
  \int_0^\infty k(y)y^s\frac{dy}{y}
       =\Gamma(s+1/2)\Gamma(s+3/2),\qquad \Re s>-1/2.
\]
Put $G(u)=e^{u/2}\sum_n\lambda_n k(\delta n e^u)$. Inverse Mellin transformation yields $G(-u)=-G(u)$. In detail, its inverse transform on a positive vertical line is that of $\Lambda(1/2+w)$. Move it to the opposite vertical line, where the functional equation supplies the negative sign, and substitute $w\mapsto-w$. No poles are crossed. Stirling's estimates for the two gamma factors and the fixed-strip polynomial growth of the automorphic $L$-function justify the shift. It follows that
\begin{equation}\label{sp:oddmellin}
 \Lambda(1/2+z)=2\int_0^\infty G(u)\sinh(zu)\,du.
\end{equation}
The rapid decay at both ends follows from the positive-argument Bessel decay and oddness.

\begin{lemma}\label{sp:moments}
In this normalization the central derivative and a useful absolute majorant are
\[
 D=\Lambda'(1/2)=4\sum_{n\ge1}\lambda_n
       \left(K_1(v_n)+\frac{2K_0(v_n)}{v_n}\right),\qquad
 H=16\sum_{n\ge1}|\lambda_n|\frac{K_1(v_n)}{v_n^2}.
\]
If $|h|\le1/2$, then
\[
 \Re\frac{\Lambda(1/2+h+it)}{h+it}
       \ge D-\max\{h^2,3t^2/2\}\,H,
\]
with the quotient at zero interpreted by continuation.
\end{lemma}
\begin{proof}
For one $n$, denote its positive-kernel weight in \eqref{sp:oddmellin} by $W_n(z)$, leaving out $\lambda_n$. Substitute $t=v_n e^{u/2}$ in $W_n'(0)$. Temporarily writing $v=v_n$, this gives
\[
 W_n'(0)=\frac4v\int_v^\infty t^2 K_1(t)\log(t/v)\,dt
        =\frac4v\int_v^\infty tK_2(t)\,dt
        =4K_1(v)+\frac8vK_0(v).
\]
The first integration by parts uses $(t^2K_2)'=-t^2K_1$. The second computation uses $K_2=K_0+2K_1/t$, $(tK_1)'=-tK_0$ and $K_0'=-K_1$. The same substitution at $z=1$ gives
\[
 W_n(1)=\int_v^\infty(t^4/v^3-v)K_1(t)\,dt
       =4K_1(v)+8K_2(v)/v.
\]
Indeed $\int_v^\infty t^4K_1(t)dt=v^4K_0(v)+4v^3K_1(v)+8v^2K_2(v)$. Subtract to get $W_n(1)-W'_n(0)=16K_1(v)/v^2$. All special-function identities here are the standard Bessel recurrences and integrals recorded by NIST in \href{https://dlmf.nist.gov/10.29}{DLMF §§10.29, 10.32 and 10.43}.

Now $\sinh(zu)/z=\int_0^u\cosh(zx)\,dx$ and $\Re\cosh((h+it)x)=\cosh(hx)\cos(tx)$. Its deviation above one is at most $h^2(\cosh x-1)$ using power series and $|h|\le1$. Its deviation below one is at most $(t^2x^2/2)\cosh(x/2)\le(3t^2/2)(\cosh x-1)$. The last inequality follows by comparing coefficients using $(2j)(2j-1)/4^{j-1}\le3$ for $j\ge1$. Integrating these bounds against the nonnegative kernel for $W_n$ shows that the error in the real part of $W_n(z)/z$ from $W'_n(0)$ has absolute value at most $\max\{h^2,3t^2/2\}(W_n(1)-W'_n(0))$. Summation proves the lemma.
\end{proof}

\begin{corollary}\label{sp:deriv}
One has $D>9.287$ and $H<33.551$. The simple central zero is the only zero of the completion with $|\Im s|\le3/7$.
\end{corollary}
\begin{proof}
The interval calculation appears in the accompanying symmetric-cube verifier. The option \texttt{--moments-only} reproduces this part alone. It computes the integer coefficients in \eqref{sp:polynomial} by counting the solutions at each prime. At odd good $p$ the trace is minus the sum of the quadratic character at $1+4x^2(x+1)$ over $\F_p$. At $2$ the trace is $-2$. It evaluates the displayed Bessel series through $N=4096$ using Arb ball arithmetic. An explicit tail bound applies to both. For $v\ge1$,
\[
 0<K_0(v)\le K_1(v)\le e^{-v}\sqrt{\pi/(2v)}e^{1/(2v)}\le 3e^{-v}.
\]
To obtain the middle inequality, start from the integral representation of $K_1$. Replace the exponent $-v\cosh x$ by $-v(1+x^2/2)$ and evaluate the resulting Gaussian integral against $\cosh x$. For $n>N$, $v_n>40$, and the factor multiplying $e^{-v_n}$ is less than two. Since $c>7/10$, both tails are bounded by
\[
 32\sum_{n>N}n^3e^{-0.7\sqrt n}
 \le64(10/7)^8\Gamma(8,0.7\sqrt N)<2\cdot10^{-5}.
\]
The summand decreases here, and the incomplete gamma of positive integral order is just a finite exponential-polynomial expression. Together with the interval sums this gives the stated bounds. The central derivative excludes multiplicity greater than one. In the critical strip use $\max\{h^2,3t^2/2\}\le27/98$ and $9.287-(27/98)33.551>0$. Absolute convergence of the Euler product and reflection deal with points outside the closed strip.
\end{proof}

\subsection{A uniform enclosure for the critical rectangle}
One can extend the computation to noncentral zeros without assuming that sign changes on the line exhaust all zeros. Here we derive the approximation and the winding test used by the verifier. Put $b=1/16$, $U=12$, $J=U/b=192$ and take $|\Re z|\le1/2$, $|\Im z|\le8$. At each $u=jb$, $1\le j\le J$, put $a=ce^{u/2}$ and $N_j=\max(1,\lceil(72/a)^2\rceil)$. Increasing this ceiling to accommodate interval arithmetic is permissible. Truncate the sum defining $G(jb)$ to $n\le N_j$ and call the result $G_j$. Define
\begin{equation}\label{sp:sinc}
 S(z)=2b\sum_{j=1}^J G_j\sinh(jbz).
\end{equation}
Then throughout the rectangle
\begin{equation}\label{sp:uniform}
             |\Lambda(1/2+z)-S(z)|<10^{-11}.
\end{equation}

To prove this, first consider trapezoidal quadrature on the whole line in \eqref{sp:oddmellin}, before truncating either sum. The function $G(u)$ is analytic in $|\Im u|<\pi$. On its right half it is given by a normally convergent Bessel series, and the identity $G(-u)=-G(u)$ extends by analyticity. On overlapping compact regions these series and identities agree. For $u\ge0$ and $|y|\le1$, the Bessel argument is $v e^{iy/2}$ with $v=c\sqrt n e^{u/2}$. Its real part is at least $0.85v$. The integral for $K_1$ gives
\[
 |K_1(ve^{iy/2})|\le K_1(0.85v)\le 4e^{-0.85v},
\]
since $0.85v\ge0.85\cdot0.7$ and the preceding Gaussian bound is then less than four times the exponential. Write $z=h+it$. The inequality $\sqrt n e^{u/2}\ge(\sqrt n+e^{u/2})/2$, together with $0.85c>1/2$, gives
\[
 |G(u+iy)e^{z(u+iy)}|
 \le2e^{|t|} e^{2u-e^{u/2}/4}
                 \sum_{n\ge1}n^4e^{-\sqrt n/4}.
\]
The bound $B=4\cdot4^{10}9!$ for the sum follows by comparing each summand with twice the integral of $x^4e^{-\sqrt x/4}$ from $n$ to $n+1$. The integral of $e^{2u-e^{u/2}/4}$ for $u>0$ is at most $2\cdot4^4\,3!$. By odd reflection the same majorization holds on the other half-line with $u$ replaced by $|u|$. The sum of the two boundary integrals over $\Im u=\pm1$ is therefore less than $10^{17}e^{|t|}$. The analytic-strip trapezoidal theorem, for example Trefethen and Weideman, \href{https://doi.org/10.1137/130932132}{\emph{The exponentially convergent trapezoidal rule}}, Theorem~5.1, bounds the discretization error by $10^{17}e^{8}/(e^{2\pi/b}-1)$. The same estimate follows by shifting the Fourier transform to the two boundaries and summing its nonzero Poisson frequencies.

On the real axis use the same bound with $e^{|t|}$ omitted. Beyond $U$, the logarithmic derivative of $e^{2u-e^{u/2}/4}$ is at most $-2$. The omitted sample tails at both ends are therefore at most
\[
       4bB\, e^{2U-e^{U/2}/4}/(e^{2b}-1).
\]
For the truncations in $n$, the positive Bessel argument is $a\sqrt n>72$ when $n>N_j$. Use $K_1(v)\le2e^{-v}$ and $|\lambda_n|\le n^3$ to bound the omitted $G$ terms by $e^{u/2}a^2\sum_{n>N_j}n^4e^{-a\sqrt n}$. The summand decreases for $a\sqrt n\ge8$. Its integral from $N_j$ to infinity then yields $2c^{-8}e^{-7u/2}\Gamma(10,72)$. Since $|2\sinh(zu)|\le 2e^{u/2}$, the total omitted contribution is at most
\[
         4b c^{-8}\Gamma(10,72)/(1-e^{-3b}).
\]
The three displayed error bounds add to less than $10^{-11}$, as can be verified with the finite-sum formula for $\Gamma(10,72)$. The computer uses a ball enclosure for each retained Bessel value in addition to this analytic error. For real $v>24$ it accelerates these evaluations using
\[
 K_1(v)=\sqrt{\frac{\pi}{2v}}e^{-v}\left(\sum_{r=0}^{24}a_r v^{-r}+R\right),
 \qquad |R|\le |a_{25}|v^{-25}.
\]
Here $a_0=1$ and $a_r=a_{r-1}(4-(2r-1)^2)/(8r)$. For an explicit justification apply Taylor's theorem to $(1+x)^{1/2}$ on $x\ge0$. Its first omitted term bounds the absolute remainder there, since the relevant derivative decreases in absolute value. Integrating at $x=t/(2v)$ against $e^{-t}t^{1/2}\,dt/\Gamma(3/2)$ in the Laplace integral for $K_1$ gives precisely these coefficients and bound. Direct Arb evaluation at all arguments gives an alternative.

\subsection{The contour verification}
Use the counterclockwise rectangle with vertices $\pm1/2\pm iT$ in the centered $z$-plane. Set $P_k(Y)=b\sum_{j=1}^J G_j(jb)^kY^j$ for $k\in\{0,2\}$. Then $S^{(k)}(z)=P_k(e^{bz})-P_k(e^{-bz})$. The verifier computes ball enclosures at exact dyadic endpoints using these two polynomials. A global bound for the third derivative is
\[
         C_3=2b\sum_{j=1}^J|G_j|(jb)^3\cosh(jb/2).
\]
For each side interval of Euclidean length $l$ and midpoint $z_m$, the image under $S$ lies within $l^2(|S''(z_m)|+l C_3/2)/8$ of the linear interpolation of its endpoint values. This is the absolute-valued linear interpolation error, by integrating twice along the segment. Add the error from \eqref{sp:uniform}. Accept the interval only when both endpoint values have some common signed coordinate, real or imaginary, beyond this bound. Bisect each unresolved interval. On every accepted interval the image of $S$ stays in that open coordinate half-plane and its distance from zero exceeds the uniform approximation error. The original analytic function and $S$ consequently have the same boundary winding. The argument change for $S$ on each accepted interval is the principal argument of the quotient of its endpoint values, since the entire segment image lies in one open half-plane. Summing the enclosures and dividing by $2\pi$ must isolate a unique integer. The argument principle gives the number of zeros counted with multiplicity, including any that were off the critical line.

On $z=it$, $S(it)$ and $\Lambda(1/2+it)$ are purely imaginary. The verifier tests their signs using \eqref{sp:uniform} at successive positive ordinates $j/64$. A strict change of sign gives at least one critical-line zero and a reflected negative one. The forced central zero supplies a further zero. When this lower count equals the winding, each of those zeros is simple and all zeros in the rectangle have been accounted for. This last equality is checked as an assertion independently of the expected list of ordinates.

The output at $T=4$ is winding number $9$ on $180$ accepted intervals, with four positive sign changes. The output at $T=8$ is winding number $23$ on $5868$ intervals, with the eleven positive sign changes listed in Theorem~\ref{sp:height8}. Both counts include the central zero and the negative ordinates. The Euler product and the functional equation exclude further zeros outside the closed critical strip. This completes the finite-height proof. The reproducible source and full integer output are included as the symmetric-cube verifier and its height-eight result. The arithmetic construction at every even degree in Theorem~\ref{sp:family} establishes central vanishing there without providing corresponding noncentral counts in those other degrees.

\subsection{Rotation of the Mellin contour at larger heights}\label{sp:rotated}
At positive ordinate $t$ the unitary completion contains two gamma factors and has small values. It is advantageous to evaluate the holomorphic function
\[
            H(w)=e^{-2iw}\Lambda(1/2+w),\qquad w=h+it.
\]
This multiplier never vanishes, and its value on $w=it$ is the positive real $e^{2t}$. Moving the $u$ contour in the bilateral Mellin integral to $\Im u=2$ yields
\[
              H(w)=\int_{\mathbb R}G(u+2i)e^{wu}\,du.
\]
The strip of analyticity and the Bessel decay used above justify the shift. In particular the integral is again absolutely convergent throughout every fixed compact set of $w$.

For $|t|\le64$, $|h|\le1/2$, use spacing $b=1/48$, cutoff $U=14$ and strip half-width $d=4/5$ about the new contour. At $u\ge0$ put $v=c\sqrt n e^{u/2}$. In the whole strip, the real part of $ve^{i(2+y)/2}$ is at least $0.16v$ for $|y|\le d$. The real-integral bound for $K_1$ gives $|K_1(ve^{i(2+y)/2})|\le400e^{-0.16v}$, since $v\ge c>0.7$. Splitting the exponential as before gives a majorant for $|G(u+i(2+y))e^{w(u+iy)}|$ of
\[
 200 e^{d|t|} e^{2u-e^{u/2}/20}
                         \sum_{n\ge1}n^4 e^{-\sqrt n/20}.
\]
The sum is bounded by $4\cdot20^{10}9!$, and integration on the positive half-line costs at most $2\cdot20^4\,3!$. The negative half-line is treated by $G(-z)=-G(z)$, using the conjugate complex angles. Both horizontal strip integrals together are therefore less than $10^{29}e^{d|t|}$. The resulting trapezoid error is less than $10^{29}e^{64d}/(e^{2\pi d/b}-1)$.

On the sampling line itself, $\cos1>0.54$, and $|K_1(ve^{i})|\le10e^{-v\cos1}$. Now the separated majorant uses the exponent $-e^{u/2}/6$, and the coefficient sum is bounded by $B=4\cdot6^{10}9!$. The omitted sample tail at $|u|>U$ is at most
\[
             10bB e^{2U-e^{U/2}/6}/(e^{2b}-1).
\]
At a nonnegative sample write $s_0=\cos1$. The truncation index is
\[
             N_j=\max\{1,\lceil(150/(s_0ce^{u/2}))^2\rceil\}.
\]
On the discarded tail $s_0v>150$ the Bessel estimate is $2e^{-s_0v}$. Comparing the decreasing function $x^4e^{-s_0ce^{u/2}\sqrt x}$ with its integral as before gives total coefficient-truncation allowance
\[
              4b c^{-8}s_0^{-10}\Gamma(10,150)/(1-e^{-3b}).
\]
At the stated mesh the three allowances are respectively below $2.84\cdot10^{-54}$, $2.62\cdot10^{-53}$ and $1.94\cdot10^{-42}$. Their sum is less than $10^{-40}$. They account for both tails, including the contribution of the sample at $u=0$.

The negative samples are determined by $G(-u+2i)=-\overline{G(u+2i)}$. Thus approximate $H(w)$ by
\[
  S_2(w)=\sum_{j=0}^{U/b}
       \bigl(B_j e^{jbw}-\overline{B_j} e^{-jbw}\bigr),
\]
where $B_j=bG(jb+2i)$ with its indicated Dirichlet cutoff for $j>0$. For $j=0$ use half of the imaginary part times $i$ of the truncated $bG(2i)$, so that this single sample is counted once. At $z=ve^{i}$, the preceding large-argument expansion of $K_1$ remains valid with the same absolute bound for the remainder within the parentheses calculated at $v=|z|$. In its Laplace integral the Taylor remainder now uses $|1+xt|\ge1$ for $\Re x\ge0$, $t\ge0$, since the argument of $z$ is one radian. The checker takes degree $64$ at $45<v\le70$ and $40$ at $v>70$.

For $v\le45$ put $x=z^2/4$, $a_k=1/(k!(k+1)!)$ and let $H_k$ denote the harmonic numbers. The checker uses the convergent integer-order series
\[
 K_1(z)=\frac1z+\frac z2(\log(z/2)+\gamma)\sum_{k\ge0}a_kx^k
               -\frac z4\sum_{k\ge0}(H_k+H_{k+1})a_kx^k.
\]
This is the principal-branch series in DLMF \S10.31. The arguments used here have phase one radian and stay in the open right half-plane. For $m=120$ and $x_0=45^2/4$ the ratio $x_0/((m+1)(m+2))<1/20$, so the two absolute remainders after terms through $m-1$ are respectively at most $2a_m x_0^m$ and $4(m+2)a_m x_0^m$. The second follows from $H_k+H_{k+1}\le2(k+1)$ and the geometric tail. The code encloses both complex coordinates for every special-function remainder. For the enlarged point-counting range it computes the same Legendre sum as before with exact signed 64-bit integers, explicitly checking $4N^2(N+1)+1<2^{63}$ first.

Finally partition $[0,64]$ into bands of width four with centers $t_0=2,6,\ldots,62$. On each band expand $S_2(it_0+\xi)$ through degree $192$. Since $|\xi|<2.1$, its uniform Taylor remainder is bounded by
\[
 2\Big(\sum_j |B_j|\Big)e^{2.1U}\frac{(2.1U)^{193}}{193!}.
\]
Added to $10^{-40}$, this is checked to be less than $10^{-39}$. The moment coefficient at order $k$ is the sum of $B_j e^{it_0jb}(jb)^k/k!$ minus $(-1)^k$ times its conjugate. Each contour is now treated with local polynomial composition, a curvature bound, certified half-planes, and winding, as in the preceding argument. The first band starts its lower boundary at ordinate $1/2$ to avoid placing the central zero on the contour. The earlier height-four certificate counts all zeros below that boundary. Each critical-line sign test uses the imaginary part, because $e^{2t}$ is positive there.

\begin{theorem}[Symmetric cube through height 64]\label{sp:height64}
The primitive degree-four standard $L$-function of $\operatorname{Sym}^3(43.\mathrm a1)$ has precisely $341$ nontrivial zeros with $|\Im s|\le64$. All are simple and lie on the unitary critical line. They comprise the central zero and $170$ conjugate pairs.
\end{theorem}
\begin{proof}
The rotated Mellin enclosure and the verified argument changes of Section~\ref{sp:rotated} give the following counts in consecutive positive bands. The first band begins at $1/2$. Every other band begins at the previous upper endpoint. Each winding independently equals the number of critical-line sign changes in that band.
\begin{center}
\begin{tabular}{r|rrrrrrrr}
\toprule
Upper ordinate & 4&8&12&16&20&24&28&32\\
Winding & 4&7&9&9&9&11&11&11\\
Panels &36&64&76&88&100&102&110&116\\\midrule
Upper ordinate &36&40&44&48&52&56&60&64\\
Winding &12&11&13&12&11&14&12&14\\
Panels &114&118&126&128&132&136&130&146\\\bottomrule
\end{tabular}
\end{center}
Every positive zero lies in a recorded open interval $(j/128,(j+1)/128)$. The unrotated height-eight calculation puts the eleven lowest positive zeros in wider intervals, each containing its rotated counterpart. In particular there is no positive zero below $1/2$. All rectangle sides are certified nonvanishing, including the common band boundaries and the two vertical boundaries. Conjugation, the earlier central simplicity certificate, and the zero-free exterior of the critical strip complete the count.
\end{proof}

The rotated samples give an independent numerical check at the central point. The reflected sum satisfies $S_2(0)=2i\sum_j\Im B_j$ and $S'_2(0)=2\sum_j jb\Re B_j$. Its value at zero encloses zero after adding the analytic allowance. Cauchy's estimate on $|w|=1/4$ bounds the derivative error by $4\cdot10^{-40}$. Since $H'(0)=\Lambda'(1/2)$, the resulting interval for the central derivative is
\[
 \left|\Lambda'(1/2)-9.287125361262821099805089842096070325498\right|
              <8.21\cdot10^{-40}.
\]
It lies in the real Bessel moment enclosure from Corollary~\ref{sp:deriv}.

Repeating the entire winding count with spacing $1/52$ and $448$-bit arithmetic in place of $1/48$ and $384$ bits gives the same sixteen winding numbers, contour subdivision counts and $170$ sign brackets. The analytic error bounds above still hold when the three mesh-dependent expressions are evaluated at the new spacing. Both estimates of the central derivative also agree.

A further calculation moves the integration line to $\Im u=19/10$ and uses mesh $1/50$, $480$-bit arithmetic, and the same cutoffs $14,150$. The integrand is then the transform for $e^{-(19/10)iw}\Lambda(1/2+w)$. All three strip and tail bounds stay valid with $\cos(19/20)$ in the coefficient bound, since the minimum real part of the rotated Bessel argument has increased. The Taylor allowance is checked anew. Bisection using strict interval signs puts all $170$ positive ordinates in open intervals of width $2^{-24}$. In particular, the lowest and highest are enclosed by
\[
 \frac{20100112}{2^{24}}<\gamma_1<\frac{20100113}{2^{24}},\qquad
 \frac{1073621973}{2^{24}}<\gamma_{170}<\frac{1073621974}{2^{24}}.
\]
This second rotation verifies the critical-line locations with a different contour. The full integer intervals are recorded in the accompanying ordinate file.

\subsection{Higher rotated Bessel enclosures}\label{sp:large}
The estimates in Section~\ref{sp:rotated} retain all their constants for a higher target $T$ after replacing the discretization factor $e^{64d}$ by $e^{Td}$, the cutoff $150$ by $V$, and the physical truncation $14$ by $U$. The following choices give two further uniform enclosures for $H(w)=e^{-2iw}\Lambda(1/2+w)$, using $|\Re w|\le1/2$ and $|\Im w|\le T$.
\begin{center}
\begin{tabular}{rrrrrrrr}
\toprule
$T$ & $1/b$ & $U$ & $V$ & Dirichlet cap & Local degree & Bits & Error\\\midrule
128 &72&16&240&352317&270&640&$10^{-75}$\\
256 &124&17&397&964031&360&1024&$10^{-139}$\\\bottomrule
\end{tabular}
\end{center}
For the first row the base-ten logarithms of the three pre-polynomial bounds are less than $-83.70$, $-187.23$, and $-78.97$, respectively. For the second row the corresponding upper bounds are $-152.74$, $-326.33$, and $-145.20$. The corresponding sums are less than $10^{-76}$ and $10^{-140}$. The degree-$K$ polynomial on each band of width four introduces at most
\[
 2\Big(\sum_j|B_j|\Big)e^{2.1U}\frac{(2.1U)^{K+1}}{(K+1)!},
\]
which is checked from the actual ball enclosures before any winding is accepted. The entire calculation is done using the locally counted Euler factors, with the signed 64-bit discriminant check valid up to both displayed caps.

For the height-$128$ row, evaluate the convergent $K_1$ series through terms $k<250$ when $v\le93$. For $93<v<155$ use the asymptotic polynomial of degree $110$, and use degree $78$ at larger $v$. For the height-$256$ row use the convergent series through $k<435$ when $v\le165$, then the asymptotic degrees $190$ for $v<254$ and $132$ thereafter. For a convergent cutoff $m$ and argument ceiling $v_0$, the program first checks $(v_0^2/4)/[(m+1)(m+2)]<1/10$. Consequently the two power-series tails are bounded by $2a_m(v_0^2/4)^m$ and $4(m+2)a_m(v_0^2/4)^m$ as above. The complex asymptotic remainder is bounded by the first omitted coefficient in absolute value times $v^{-k-1}$ within its normalized parentheses. These special-function error enclosures accompany every retained value independently of the analytic truncation allowance. An additional symmetry check encloses zero both for the truncated real part of $bG(2i)$ with its coefficient tail and for the complete approximation to $H(0)$ with its uniform error.

\subsection{Shifting only the omitted series}\label{sp:tailshift}
A second contour calculation for this degree-four example uses its Dirichlet series and the standard gamma function directly. It gives an independent way to control the Bessel calculation. The analytic observation is that the abscissa for proving a tail estimate need not be the abscissa for evaluating the retained coefficients.

\begin{lemma}[A separate contour for a Dirichlet tail]\label{sp:taillemma}
Write $D_N(s)=\sum_{n>N}\lambda_n n^{-s}$, assume $|\lambda_n|\le d_D(n)$, and let the gamma prefactor be $Q^s\prod_{j}\Gamma(a_j s+b_j)$ with $Q>0$, $a_j>0$, and real $b_j$. Start the tilted two-kernel Cauchy formula on a line $\sigma>1$ where all gamma arguments have positive real part. Its exponential tilt is $e^{-i\theta w}$, with $w=s-1/2$. Let $A=\sum_j a_j$, take $0<\phi<\pi/2$ with $c=A\phi-|\theta|>0$, and choose $1<v\le\sigma'$ with $\sigma'\ge\sigma$. Uniformly on $|\Re w|\le1/2$, the part of the two continuous kernels coming from $D_N$ is at most
\[
 \frac{2\zeta(v)^D N^{v-\sigma'}}{\pi c(\sigma'-1)}
   Q^{\sigma'}\prod_j (\cos\phi)^{-(a_j\sigma'+b_j)}\Gamma(a_j\sigma'+b_j).
\]
The same estimate holds for either sign of the functional equation.
\end{lemma}
\begin{proof}
Separate the finite and infinite Dirichlet series before discretizing. In each tail integral shift the real part of its right-line variable from $\sigma$ to $\sigma'$. The two Cauchy poles lie to its left, the series converges normally throughout, and the gamma arguments have positive real parts. Rotating the Euler integrals by $\phi$ gives exponential decay $e^{-A\phi|y|}$ along each of the bounded family of vertical lines. Both tilts have modulus at most $e^{|\theta||y|}$. The horizontal sides therefore disappear. On the new line absolute convergence and the divisor bound give $|D_N(\sigma'+iy)|\le\zeta(v)^D N^{v-\sigma'}$. Each Cauchy denominator has modulus at least $\sigma'-1$. The two kernels, the factor $1/(2\pi)$, and the integral of $e^{-c|y|}$ give the prefactor in the assertion.
\end{proof}

For a complex gamma shift $b_j$ whose real part meets the same positivity hypothesis, replace $b_j$ on the right by $\Re b_j$ and multiply by $e^{\phi|\Im b_j|}$. This follows from $|a_jy+\Im b_j|\ge a_j|y|-|\Im b_j|$. For a nonselfdual pair the two right-line tails can be treated separately by the same one-kernel argument, retaining the individual gamma constants.

For the symmetric cube take $\theta=12/5$, so that $H(w)=e^{-(12/5)iw}\Lambda(1/2+w)$ is purely imaginary on the critical line with its original signed imaginary part multiplied by $e^{12t/5}$. The working right line is $\sigma=65/2$, whereas the omitted coefficients are bounded on $\sigma'=257/2$. Write $m=32$, $m'=128$, and
\[
 A(x)=\delta^{-(x+1/2)} x!(x+1)!,\qquad
 C(x)=16\cdot9^{2x+3}A(x),\qquad A_*=3^{2m+3}A(m).
\]
Here $A(x)$ is the exact real gamma prefactor at abscissa $x+1/2$. Use $b=1/18$, $Y=1400$, $J=Y/b$, coefficient cutoff $N=1100000$, integration half-width eight and local polynomial degree $220$ on four-unit bands. The right line and enlarged strip admit decay $e^{-|y|/2}$ after the tilt, by taking the Euler rotation angle $29/20$ and using $\cos(29/20)>1/9$. For the nondecaying arithmetic norm use angle $6/5$ and $\cos(6/5)>1/3$. For the tail alone use $\phi=121/100$, so that $c=1/50$ and $\cos\phi>7/20$. The resulting four bounds are
\begin{align*}
 &\frac{C(m+8)}{e^{16\pi/b}-1},\qquad \frac{bC(m)e^{-Y/2}}{e^{b/2}-1},\\
 &\frac{32N^{2-\sigma'}A(m')(20/7)^{2m'+3}}{\pi(1/50)(\sigma'-1)},\qquad
 16b(2J+1)A_*\frac{(2.1/m)^{221}}{1-2.1/m}.
\end{align*}
Their base-ten logarithms are respectively below $-180.51,-138.46,-102.54,-124.67$. The total is less than $10^{-98}/2$. As before the finite exponential sums are enclosed by certified frequency gridding. Individual signed finite weights below $10^{-98}/[64b(2J+1)]$ may be omitted together with their paired terms. The total additional uniform error is less than $10^{-98}/64$.

In the first denominator use $W_j=e^{-i\theta m}e^{\theta jb}Z_j$. The second denominator now has the minus sign from the functional equation, and its signed weight is $-\overline{W_{-j}}$. Therefore if $B_k=\sum_j W_j/(m+ijb-it_0)^{k+1}$, the local coefficient is $b(B_k-(-1)^k\overline{B_k})/(2\pi)$. This gives purely imaginary values along the line, which supplies the sign test. The right-line enclosure uses $896$-bit arithmetic and the same rigorous contour interpolation test as above.

\paragraph{Independent recovery of the prime traces.}
The second checker can obtain the many good traces using Hasse's theorem and the rational point $P=(0,0)$. At an odd good $p$ complete the square to write $Y^2=x^3+x^2+1/4$, so that $P=(0,1/2)$. Every possible group order lies in $[L,R]$, where $L=p+1-\lfloor2\sqrt p\rfloor$ and $R=p+1+\lfloor2\sqrt p\rfloor$, and the true order annihilates $P$. Take $s=\lfloor\sqrt{R-L}\rfloor+1$. Tabulate $jP$ for $0\le j<s$, retaining all collisions. Compare $-(L+ks)P$ with the table for every $k$ with $L+ks\le R$. This enumerates every $n\in[L,R]$ for which $nP=0$. A unique candidate is necessarily the group order. When there is more than one, the checker performs the exact quadratic-character sum and verifies that the resulting order is among the candidates. It treats $2$ and $43$ by their explicit local factors. For an independent implementation check, all $2762$ prime traces at primes up to $25000$ agreed with the direct residue-counting implementation. This arithmetic procedure uses no assumed equidistribution or order statistic about $P$ modulo $p$.

\begin{theorem}[Symmetric cube through height 128]\label{sp:height128}
The symmetric cube of $43.\mathrm a1$ has exactly $791$ nontrivial zeros with $|\Im s|\le128$. All are simple and lie on the critical line. There are $395$ conjugate pairs and one central zero.
\end{theorem}
\begin{proof}
The separate-tail Cauchy calculation first recovers every bracket of the rotated Mellin calculation on $60<\Im s<64$. Its subsequent positive windings, in successive four-unit intervals, are
\begin{center}
\begin{tabular}{r|rrrr|r}
\toprule
Range & \multicolumn{4}{c|}{Successive positive windings} & Total zeros through upper endpoint\\\midrule
$[64,80]$ &13&13&13&14&447\\
$[80,96]$ &14&13&14&14&557\\
$[96,112]$&15&14&14&15&673\\
$[112,128]$&14&14&15&16&791\\\bottomrule
\end{tabular}
\end{center}
At each band the checker separately certifies that the corresponding number of intervals of strict critical-line sign change equals the winding. The contours exclude boundary zeros. There are therefore $225$ new simple positive zeros above $64$. The lower range, reflection, and exterior of the critical strip are already covered by Theorem~\ref{sp:height64}.
\end{proof}

\begin{theorem}[Symmetric cube through height 256]\label{sp:height256}
The primitive standard $L$-function of $\operatorname{Sym}^3(43.\mathrm a1)$ has precisely $1803$ nontrivial zeros with $|\Im s|\le256$. All are simple and critical, comprising $901$ conjugate pairs and the central zero.
\end{theorem}
\begin{proof}
The same separated-tail bounds of Section~\ref{sp:tailshift} apply uniformly on all the remaining four-unit bands. The two certified counts on each band, from winding and from critical-line signs, agree. The finite outputs give
\begin{center}
\begin{tabular}{r|rrrr|r}
\toprule
Range & \multicolumn{4}{c|}{Successive positive windings} & Cumulative total\\\midrule
$[128,144]$&14&15&15&16&911\\
$[144,160]$&15&15&15&16&1033\\
$[160,176]$&15&16&15&16&1157\\
$[176,192]$&16&15&17&15&1283\\
$[192,208]$&16&17&15&17&1413\\
$[208,224]$&15&17&15&17&1541\\
$[224,240]$&16&17&16&16&1671\\
$[240,256]$&17&16&17&16&1803\\\bottomrule
\end{tabular}
\end{center}
Every positive sign bracket is an open interval of width $1/128$. The last has endpoints $32732/128$ and $32733/128$. All boundaries of the rectangles have been certified nonzero. The $506$ new positive roots above $128$ account, with their reflections, for the increase from $791$ to $1803$ zeros.
\end{proof}

For an additional change of contour one may use $\theta=11/5$, $\sigma=63/2$, $\sigma'=289/2$, $N=940000$, $b=1/20$, $Y=1550$, strip half-width $8$, and local degree $260$. The raw gamma rotation is $11/10$ with reciprocal cosine bounded by $9/4$. The strip angle is $7/5$, with reciprocal cosine bounded by $6$ and remaining decay $e^{-3|y|/5}$. The omitted-series angle is $111/100$ with reciprocal cosine below $9/4$ and remaining decay $e^{-|y|/50}$. Substitution into Lemma~\ref{sp:taillemma} and the three finite-polynomial bounds gives base-ten error logarithms below $-244.41,-255.66,-123.10,-181.29$, respectively in the same order as before. These values permit an analytic allowance $10^{-119}$. The corresponding polynomial arithmetic uses $1024$ bits. This tilt computes $e^{-(11/5)iw}\Lambda(1/2+w)$ and permits a separate signed bisection of the critical roots.

The high Bessel calculation of Section~\ref{sp:large} independently recovers every one of the $225$ positive Cauchy brackets on $(64,128)$. The Bessel verification uses $2514$ accepted panels there and the Cauchy verification uses $2506$. Seven of the sixteen panel counts differ, while all sixteen windings and all $225$ strict root brackets agree.

\subsection{Grouping the rotated Bessel sum in $\sqrt n$}\label{sp:groups}
The high Bessel sum contains many terms at nearby arguments. Its rotated samples can be enclosed without evaluating a separate special function for each such term. For $a=ce^{u/2}$ and $\omega=e^{i}$, extract the factor $a^2/2$ from the sample. The remaining sum is
\[
         \sum_{n\le N_u}\frac{A_n}{\sqrt n}K_1(a\omega\sqrt n),
\]
where $A_n$ is the integer motivic symmetric-cube coefficient. Round the original ceiling $N_u$ upwards to the end of the unit interval in $\sqrt n$ containing it. This retains more terms and leaves the preceding discarded-tail upper bound valid. Use direct Bessel enclosures for $n<256$. For each integer $j\ge16$, the bin contains $j^2\le n<(j+1)^2$. Put $x=j+1/2$, $y_n=(\sqrt n-x)/x$ and precompute signed moments
\[
        M_k=\sum_{j^2\le n<(j+1)^2}\frac{A_n}{\sqrt n}y_n^k.
\]
The moments depend on the curve and bin, so a single precomputation serves every sample $u$.

For complex $h$ with $|h|\le x\cos(1)/2$, the real part of $a\omega(x+h)$ is at least $ax\cos(1)/2>3$. The integral representation and the real large-argument bound show that $|K_1(a\omega(x+h))|<4$ on this entire disk. With $\eta=1/(x\cos1)$, the absolute Taylor remainder at every point of the bin after degree $m$ is at most
\[
                       4\eta^{m+1}/(1-\eta).
\]
The program selects the bin degree so that this is below $10^{-175}$. It adds the error times $\sum |A_n|/\sqrt n$ to the complex enclosure of the computed bin sum. Thus signed cancellation in the moments plays no role in the tail bound.

Write $z=a\omega x$ and expand $K_1(z(1+y))=\sum_{k\ge0}g_k y^k$. The two starting values are $g_0=K_1(z)$ and $g_1=-zK_0(z)-K_1(z)$. Substitution into the modified Bessel equation gives, for $k\ge0$ with $g_{-1}=g_{-2}=0$,
\[
 g_{k+2}=\frac{z^2(g_k+2g_{k-1}+g_{k-2})+(1-k^2)g_k-(k+1)(2k+1)g_{k+1}}
                   {(k+1)(k+2)}.
\]
Consequently $\sum_{k=0}^m g_kM_k$ is the entire retained Taylor approximation in the bin. Every recurrence operation and moment is enclosed separately. To enclose $K_0(z)$ at $|z|\le165$, use
\[
 K_0(z)=-(\log(z/2)+\gamma)\sum_{i\ge0}\frac{(z^2/4)^i}{(i!)^2}
                  +\sum_{i\ge0}H_i\frac{(z^2/4)^i}{(i!)^2}.
\]
With $m_0=435$ and $v_0=165^2/4$, the two tails after terms $i<m_0$ are at most $2v_0^{m_0}/(m_0!)^2$ and $2(m_0+2)v_0^{m_0}/(m_0!)^2$. These follow from $v_0/(m_0+1)^2<1/10$ and $H_i\le i+1$. For larger arguments use the same asymptotic degrees as for $K_1$, with coefficient multiplier $-(2i-1)^2/(8i)$ appropriate to $K_0$. The integral remainder follows by Taylor expansion of $(1+t/(2z))^{-1/2}$ in the same right half-plane. This bounds both starting values used in the recursion.

The final prefactor per sample is $b e^{u/2}\omega^3 a^2/2$. The completed bin sum replaces only the finite arithmetic evaluation of that sample. The strip quadrature, omitted-frequency tails, analytic coefficient tails, half sample at zero, local-polynomial Taylor bound, and contour winding remain those of the height-$256$ rotated Bessel verifier. The grouped calculation uses $1152$ bits.
The change in the number of special-function evaluations is explicit. Before discarding zero coefficients, the individual cutoffs request $120023459$ coefficient--sample pairs. The grouped calculation visits $224200$ bins across the sampling points, with two Bessel starting values per visit, and requires at most $297244$ additional evaluations at indices $n<256$. Thus at most $745644$ starting special-function values suffice, followed by the finite recurrences. The global integer cap after rounding complete bins is $964323$ and the largest bin degree is $189$.

Using these grouped Bessel enclosures, the verifier certifies the complete upper range of Theorem~\ref{sp:height256} independently of the Cauchy calculation. On the $48$ four-unit bands above $64$, all $731$ strict positive root brackets and all winding numbers agree. The grouped Bessel calculation requires $8162$ accepted contour panels there, while the Cauchy calculation uses $8134$, with different counts in $28$ bands. The lower Bessel theorem supplies the center and the first $170$ pairs. Thus the full $1803$-zero result is certified by both analytic constructions.

The $\theta=11/5$ Cauchy check gives the same $48$ upper winding numbers with $8160$ accepted panels. Its strict-sign bisections put all $731$ positive ordinates above $64$ in intervals of width $2^{-24}$. Combining them with the $170$ enclosures from the $19/10$ Bessel rotation gives this width for all $901$ positive ordinates through $256$. The last lies strictly between $4290265591/2^{24}$ and $4290265592/2^{24}$.
Repeating the $170$ low bisections with Bessel rotation $2i$, mesh $1/48$ and $512$-bit enclosures gives exactly the same integer endpoints as rotation $19i/10$ at mesh $1/50$ and $480$ bits.
Recomputing the upper $731$ bisections with the $12/5$ Cauchy tilt gives the same endpoints as the $11/5$ tilt. Thus the two complete degree-four lists agree on all $901$ width-$2^{-24}$ intervals, with separate Bessel rotations below height $64$ and separate right-line Cauchy approximations above it.

\clearpage
\section{The symmetric fourth of the conductor-43 curve}\label{sf:chapter}
Let $E\colon y^2+y=x^3+x^2$. Its fourth symmetric power gives a cuspidal standard representation of $\mathrm{GL}_5/\Q$ by Newton and Thorne's symmetric power theorem, since $E$ is non-CM as verified in Section~\ref{sp:twists}. We compute its complete local data before evaluating any values. The discussion uses unitary coordinates, so the central line corresponds to motivic $\Re s=5/2$.

\subsection{Arithmetic and archimedean parameters}
At a good prime put $t=a_p(E)$, $q=t^2-2p$, and take the motivic variable $X$. The inverse Euler polynomial is
\begin{equation}\label{sf:euler}
 (1-p^2X)(1-pqX+p^4X^2)(1-(q^2-2p^2)X+p^4X^2).
\end{equation}
Indeed, for reciprocal roots $\alpha,\beta$ of $E$, the five weights are $p^2$, $p\alpha^2,p\beta^2$ and $\alpha^4,\beta^4$. The displayed traces follow from $\alpha+\beta=t$ and $\alpha\beta=p$. At the single bad prime $43$ the curve has nonsplit multiplicative reduction. Proposition~4.1 of Dummigan, Martin and Watkins gives the inverse polynomial $1-X$, conductor exponent four, and local sign $+1$ for the fourth symmetric power. At every other finite prime the representation is unramified. Thus its conductor is $43^4=3418801$. The normalized Dirichlet coefficients are $\lambda_n=A_n/n^2$, with integral $A_n$ obtained by local power-series inversion and multiplicativity. Hasse's bound on the good reciprocal roots, together with the bad normalized factor $1-43^{-2}43^{-s}$, gives $|\lambda_n|\le d_5(n)$ for every $n$. Here $d_5(n)$ counts ordered factorizations of $n$ into five factors. This estimate includes all prime powers and all multiples of the conductor.

At the real place the symmetric fourth of the weight-one angular parameter decomposes into $I_4\oplus I_2\oplus\mathbf1$. The middle monomial $x^2y^2$ is fixed under the nonidentity component of the Weil group. The local root numbers multiply to $i^5i^3=1$ and its gamma factor is $\Gamma_{\mathbb R}(s)\Gamma_{\mathbb C}(s+1)\Gamma_{\mathbb C}(s+2)$. In a normalization differing from the standard completion by a positive constant let
\begin{equation}\label{sf:completed}
 \Lambda_4(s)=\delta^{-s}\Gamma(s/2)\Gamma(s+1)\Gamma(s+2)
       \sum_{n\ge1}\lambda_n n^{-s},\qquad \delta=4\pi^{5/2}/43^2,
\end{equation}
where the series is initially taken on $\Re s>1$. The completed cuspidal function is entire and satisfies $\Lambda_4(s)=\Lambda_4(1-s)$. The sign and the real coefficients imply $\Lambda_4(1/2+it)\in\mathbb R$ for real $t$.

\begin{proposition}[Persistence under finite extension]\label{sf:strongirred}
For every number field $K$ and every integer $m\ge1$, the restriction of the $\ell$-adic representation $\operatorname{Sym}^m V_\ell(E)$ to $G_K$ is absolutely irreducible. It cannot be induced from a representation of a proper finite extension of $K$. In particular, after any finite base change the good Euler factors of the degree-three, degree-four and degree-five systems cannot agree prime by prime with a product of compatible character factors.
\end{proposition}
\begin{proof}
By Serre's open image theorem for this non-CM curve, the image of $G_{\mathbb Q}$ contains an open subgroup of $\mathrm{GL}_2(\mathbb Z_\ell)$. The image of every finite-index $G_K$ has finite index in the original image, so is also open. It is Zariski dense in $\mathrm{GL}_2$ over $\overline{\mathbb Q}_\ell$, because a polynomial vanishing on a product of four open $\ell$-adic coordinate balls vanishes identically, by the polynomial identity in one variable applied to each coordinate. The symmetric power of the standard representation is irreducible for $\mathrm{SL}_2$ in characteristic zero, so a subspace invariant under the Galois image must be either zero or the full symmetric power.

An induced representation from a proper finite extension restricts to a direct sum over its distinct cosets on the Galois group of its normal closure. Each summand has smaller dimension. This contradicts the irreducibility just proved after passing to that normal closure. Finally, equality of the good Euler factors of semisimple compatible representations identifies their semisimplifications by Chebotarev density and the character criterion. This excludes an equality of the good factors with those of a sum of characters. The open image input is Serre, \href{https://www.college-de-france.fr/media/jean-pierre-serre/UPL5874918517843398173_Serre_proprie_te_s_galoisiennes_des_courbes_elliptiques.pdf}{\emph{Propri\'et\'es galoisiennes des points d'ordre fini des courbes elliptiques}}, Inventiones 15 (1972), 259--331.
\end{proof}

\subsection{A rectangle from two reflected Cauchy integrals}
The two-kernel formula of \eqref{sq:two-kernels} holds here with $\Lambda_4$ in place of $\Lambda_2$. Use the right line $\sigma=32$, $r=63/2$, frequency spacing $b=1/20$, cutoff $Y=195$, $J=Y/b=3900$, and $N=147456$ Dirichlet coefficients. In the first positive band write $w=s-1/2=4i+\xi$, so that $|\xi|<5$ whenever $|\Re w|\le1/2$ and $0\le\Im w\le8$. Expand each Cauchy denominator through degree $K=232$ as in Section~\ref{sq:height32}. We give all constants needed for the resulting uniform approximation.

For even positive integer $x$ put
\[
 A(x)=\delta^{-x}(x/2-1)!\,x!\,(x+1)!,\qquad
 C(x)=2^{5+5x/2+3} A(x).
\]
Rotating each Euler gamma integral by one radian in the sign of its imaginary part gives $|\Gamma(a+iv)|\le 2^a\Gamma(a)e^{-|v|}$. Since $|F_4(s)|\le\zeta(2)^5<2^5$ on $\Re s\ge2$, the completed function in $|\Re s-32|\le4$ is bounded in modulus by $C(36)e^{-5|\Im s|/2}$. On the sampling line use $C(32)e^{-5|\Im s|/2}$. The two poles of the Cauchy kernels stay at least $27$ from the integration variables throughout the widened strip. The sum of the two boundary norms, including $1/(2\pi)$ and the integral of the exponential bound, is consequently smaller than $C(36)$. The omitted Dirichlet tail is at most $2^5N^{2-32}$ by the global divisor bound just proved. The four allowances for infinite trapezoidal summation, discarded samples, Dirichlet truncation and expansion in $\xi$ are therefore bounded by
\begin{align*}
 &\frac{C(36)}{e^{8\pi/b}-1},\qquad
 \frac{bC(32)e^{-5Y/2}}{e^{5b/2}-1},\\
 &2^5b(2J+1)A(32)N^{-30},\qquad
 2^5b(2J+1)A(32)\frac{(5/r)^{233}}{1-5/r}.
\end{align*}
These are respectively below $10^{-38},10^{-55},2\cdot10^{-21},10^{-52}$. Their sum is less than $10^{-18}/2$. The interval computations use the larger allowance $10^{-18}$. Apart from these analytic truncations, every numerical operation is enclosed by ball arithmetic.

\subsection{Gridding a finite Dirichlet polynomial}\label{sf:grid}
At these conductors the same Dirichlet polynomial must be enclosed at many equally spaced frequencies. We describe the acceleration used here and for the larger degree-three rectangles. It applies when $d\in\{3,5\}$ and the unitary coefficients satisfy $|\lambda_n|\le d_d(n)$. For $n\le127$ one may compute the contribution directly. For the others partition the integers into doubled intervals $[L,2L)$, starting at $L=128$ and truncating the last interval at $N+1$. Let $M$ be the smallest power of two with $M\ge2J$. For each $n$ select an approximate nearest integer $h_n$ to $M b\log n/(2\pi)$ and put
\[
        u_n=J(b\log n-2\pi h_n/M),\qquad
        g=(1001/1000)\pi J/M.
\]
The program certifies $|u_n|<g$ after selecting the index. For $0\le j\le J$ we have the exact factorization
\[
 e^{-ijb\log n}=e^{-2\pi ijh_n/M}e^{-i(j/J)u_n}.
\]
Expand the final exponential through degree $k_L$ on each doubled interval. The error in that entire group, including the coefficients $\lambda_nn^{-\sigma}$, is at most
\[
    2^dL^{2-\sigma}e^g\frac{g^{k_L+1}}{(k_L+1)!}.
\]
This uses $\sum_{n\ge L}d_d(n)n^{-\sigma}\le L^{2-\sigma}\zeta(2)^d$. Increase the group orders until the sum of these explicit remainders is below $\tau$, where $\tau=\varepsilon/[100b(2J+1)A(\sigma)]$ and $\varepsilon$ is the analytic error allowance. On the degree-three line the raw gamma constant from that calculation replaces $A(\sigma)$.

For each exponent $k$ accumulate $\lambda_nn^{-\sigma}u_n^k/k!$ in bin $h_n$ over the groups with $k\le k_L$. A length-$M$ discrete Fourier transform evaluates that coefficient sum at every $j$, including the first exponential in the factorization. Horner's rule with argument $-ij/J$ then sums the powers. Adding the explicit absolute truncation error to both complex coordinates encloses the original finite Dirichlet polynomial. This gridding affects computational speed only. The Cauchy estimates and winding argument already allow arbitrary rigorous enclosures of those polynomial values. As an independent check, the program recomputes the full finite Dirichlet sum term by term at $j=0,\lfloor J/3\rfloor,J$ and requires overlap with each Fourier enclosure.

\begin{theorem}[A primitive degree-five rectangle]\label{sf:height8}
The fourth symmetric power of $43.\mathrm a1$ has exactly $32$ nontrivial zeros with $|\Im s|\le8$. They are all simple and have $\Re s=1/2$ in unitary coordinates. The sixteen positive ordinates lie one each in the open intervals $(j/64,(j+1)/64)$ with
\[
 j\in\{45,76,126,152,178,229,283,315,342,352,376,405,431,463,481,504\}.
\]
\end{theorem}
\begin{proof}
The verifier uses the exact Euler polynomials and $640$-bit arithmetic in the construction above. Its counterclockwise contour with vertices $0,1,1+8i,8i$ in the unitary $s$-plane has winding number $16$ on $134$ accepted panels. The same signed-component test and interpolation error estimate as in the degree-three proof certify that no boundary or subdivision ambiguity can change the winding. The sixteen displayed intervals all have strict opposite signs of the real completion at their endpoints, so supply distinct critical-line zeros. Their number equals the winding counted with multiplicity. The lower boundary proves absence of a real zero, and conjugation accounts for the negative half. The Euler product and functional equation cover the exterior of the critical strip.
\end{proof}

At $s=1/2$ the same enclosing polynomial gives
\[
  |\Lambda_4(1/2)-28.42555674994986462|<5.37\cdot10^{-18}.
\]
Bisection with strict endpoint signs puts all sixteen positive roots in intervals of width $2^{-24}$ recorded in the machine report. The first lower endpoint numerator is $11947086$. The completed function is positive everywhere on the real axis. Indeed the Euler factors and gamma factors are positive for real $s>1$, and the contour proves nonvanishing of the real segment $[0,1]$. The even functional equation gives positivity on the rest of the real axis.

\begin{theorem}[Symmetric fourth through height 16]\label{sf:height16}
Exactly $76$ nontrivial zeros of $\operatorname{Sym}^4(43.\mathrm a1)$ satisfy $|\Im s|\le16$. Every one is simple and lies on the critical line.
\end{theorem}
\begin{proof}
Repeat the Cauchy integral with $(\sigma,b,Y,N,d,K)=(46,1/24,330,440000,6,330)$ and $1024$-bit arithmetic. Here $d$ is the half-width of the integration strip. In the four displayed allowances above replace $32$ by $46$, $36$ by $52$, $8\pi$ by $12\pi$, the Dirichlet exponent $-30$ by $-44$, and the geometric exponent $233$ by $331$. The four resulting bounds are respectively smaller than $2\cdot10^{-115},10^{-118},5\cdot10^{-41},4\cdot10^{-110}$. Use an overall allowance $10^{-34}$ on each translated polynomial. The independent windings on $0\le\Im s\le8$ and $8\le\Im s\le16$ are $16$ and $22$, on $134$ and $240$ accepted panels, respectively. All of these are matched by distinct strictly signed positive brackets. The lower band also reproduces every height-eight bracket computed with the smaller right line and coefficient cutoff. Multiplicity, reflection, and the exterior of the strip are handled as in Theorem~\ref{sf:height8}.
\end{proof}
Bisection in the second computation gives width $2^{-24}$ for all $38$ positive roots. Its independent central value satisfies
\[
 |\Lambda_4(1/2)-28.425556749949864615647206307011527|<1.60\cdot10^{-34}.
\]

\paragraph{Real zeros of the three finite $L$-functions.}
The three certificates also determine all real zeros of the finite Euler products after continuation, including the trivial zeros. The positive Euler factors show that every completed value is positive for real $s>1$. The functional equations give nonvanishing for $s<0$. On $[0,1]$ the contour certificates exclude real zeros of the symmetric square and symmetric fourth, and leave only the simple central zero of the symmetric cube. At a pole of the gamma product the real completed value is nonzero, so the order of vanishing of the finite part is exactly the pole order. In unitary coordinates this gives the following complete list.
\begin{center}
\begin{tabular}{lll}
\toprule
Power & Real zero of the finite part & Multiplicity\\\midrule
$2$ & negative odd integer & 2\\
    & negative even integer & 1\\
$3$ & $1/2$ or $-1/2$ & 1\\
    & $-3/2,-5/2,\ldots$ & 2\\
$4$ & $0,-1$ &1\\
    & negative even integer &3\\
    & negative odd integer $\le-3$ &2\\\bottomrule
\end{tabular}
\end{center}
The simple zero at $1/2$ for the symmetric cube is the only real nontrivial zero in this list.

\subsection{A larger degree-five contour with an arithmetic audit}\label{sf:large}
For four consecutive eight-unit bands the raw gamma constant $A(\sigma)$ and the rotated bound $C(\sigma+d)$ remain the ones above. The two sets of parameters used for a common height of $32$ are
\begin{center}
\begin{tabular}{lrrrrrrr}
\toprule
Method & $\sigma$ & $1/b$ & $Y$ & $N$ & $d$ & $K$ & Bits\\\midrule
Residue sum &52&28&440&1200000&6&410&1280\\
Group order &54&30&450&1180000&6&430&1408\\\bottomrule
\end{tabular}
\end{center}
The first row obtains each good trace by the full finite-field character sum. The second enumerates the multiples of the rational point in Hasse's interval as in Section~\ref{sp:tailshift}, with exhaustive character counting whenever more than one candidate remains. Both rows use the good factors of \eqref{sf:euler}, the inverse factor $1-X$ at $43$, and the archimedean parameters of \eqref{sf:completed}. The four uniform errors are bounded by
\begin{align*}
 &C(\sigma+d)/(e^{2\pi d/b}-1),\qquad
 bC(\sigma)e^{-5Y/2}/(e^{5b/2}-1),\\
 &32b(2J+1)A(\sigma)N^{2-\sigma},\qquad
 32b(2J+1)A(\sigma)(5/r)^{K+1}/(1-5/r).
\end{align*}
For both rows their sum is smaller than $10^{-60}/2$, and the checker uses $10^{-60}$ as the uniform analytic allowance. Three independent full Dirichlet evaluations also check the Fourier evaluations on each row. The finite-field sieve and the Fourier frequency mesh therefore both change between these two computations.

\begin{theorem}[Symmetric fourth through height 32]\label{sf:height32}
The cuspidal degree-five standard $L$-function of $\operatorname{Sym}^4(43.\mathrm a1)$ has precisely $188$ nontrivial zeros with $|\Im s|\le32$. They are simple and form $94$ conjugate pairs on the critical line. Each positive ordinate is enclosed in an open interval of width $2^{-24}$.
\end{theorem}
\begin{proof}
With the first row above the four band counts are
\begin{center}
\begin{tabular}{r|rrrr}
\toprule
Upper ordinate &8&16&24&32\\
Winding &16&22&27&29\\
Panels &134&240&292&302\\\bottomrule
\end{tabular}
\end{center}
All $94$ lower bounds from strictly changing critical-line signs are realized in disjoint intervals, and all upper bounds from the contour counts agree. Each positive zero therefore has multiplicity one. The $38$ refined brackets below $16$ agree with those already obtained from the different abscissa $46$. The final upper-band zero lies in $(534490526/2^{24},534490527/2^{24})$. As in Theorem~\ref{sf:height8}, reflection and the exterior of the strip complete the count.
\end{proof}
The height-$32$ polynomial gives $|\Lambda_4(1/2)-a|<2.06\cdot10^{-60}$ for the decimal
\[
       a=28.42555674994986461564720630701152705828881114887031479314912.
\]

The second row of the parameter table gives byte-identical numerical output, including the four windings and panel counts, both sets of $94$ brackets, and the central enclosure. The two runs recover the same roots from different finite-field algorithms and different analytic sampling parameters.

\paragraph{An independent check on the first 25000 coefficients.}
At each of the $2762$ primes $p\le25000$, count the fibers of the map $y\mapsto y^2+y$ on $\mathbb F_p$, and add the sizes of the fibers above $x^2(x+1)$ for all $x$. This counts the original affine cubic directly, even at $p=2$ and $43$. The motivic trace is $p$ minus the affine count, giving $-2$ and $-1$ at these two primes. At a good prime let $t_0=2$, $t_1=a_p$ and $t_i=a_pt_{i-1}-pt_{i-2}$. Thus $t_i=\alpha^i+\beta^i$. Define $U_0(x,y)=1$, $U_1(x,y)=x$, $U_j(x,y)=xU_{j-1}(x,y)-yU_{j-2}(x,y)$. For symmetric power $m$ the sum of the $i$th powers of its reciprocal roots is $S_i=U_m(t_i,p^i)$. Taking the logarithmic derivative of the local generating series gives
\[
     H_0=1,\qquad kH_k=\sum_{i=1}^k S_i H_{k-i},\qquad H_k=A_{p^k}.
\]
At $43$ the single reciprocal root instead gives $H_k=(-1)^{mk}$. Exact multiplication over the prime factors of each $n$ recovers the three sets of integral coefficients from these recurrences. Each of the $75000$ values so obtained agrees with the reciprocal-polynomial implementation used for the zeros.

\subsection{Separating zeros of different symmetric powers}\label{sf:separation}
A common zero of two simple factors would have multiplicity at least two in their product. We compare their ordinates after putting \emph{all} three functions in unitary coordinates. Write $D=2^{24}$, and let $\gamma_{d,j}$ be the $j$th positive ordinate of the degree-$d$ completion. Thus $d=3,4,5$ refer to symmetric powers $2,3,4$, respectively.

\begin{theorem}[Joint simplicity through height 32]\label{sf:joint32}
The normalized symmetric square, cube and fourth of $43.\mathrm a1$ have pairwise disjoint nontrivial zeros at $|\Im s|\le32$. Their product has exactly $427$ such zeros. All are simple and lie on the critical line. Among all positive zeros of the product the closest pair is $(\gamma_{3,39},\gamma_{5,74})$, with
\[
             8972/D<\gamma_{3,39}-\gamma_{5,74}<8974/D.
\]
\end{theorem}
\begin{proof}
The individual counts give $48,71,94$ positive ordinates, respectively. Only the symmetric cube vanishes at the center, and it does so simply. Each positive ordinate has a signed enclosure $(l/D,(l+1)/D)$. For two disjoint enclosures with lower numerators $l$ and $m$, the absolute distance is strictly between $(|l-m|-1)/D$ and $(|l-m|+1)/D$. Comparing the integer endpoints gives the table. The last column bounds $D$ times the distance of every other positive pair of the indicated degrees strictly from below.
\begin{center}
\begin{tabular}{ccrrrc}
\toprule
Degrees & Closest indices & First $l$ & Second $m$ & $|l-m|-1$ & Other pairs\\\midrule
$3,4$ & $20,30$ &272666915&272627243&39671&105274\\
$3,5$ & $39,74$ &441502638&441493665& 8972& 58496\\
$4,5$ & $37,50$ &321949464&322006395&56930& 68412\\\bottomrule
\end{tabular}
\end{center}
For example the second row checks $48\cdot94=4512$ pairs against the square and fourth-power lists. For each row even the upper bound at the closest pair is below the last column. Within the three individual positive lists the successive gaps, multiplied by $D$, are strictly greater than $745300$, $1162431$ and $872155$, respectively. The first ordinate in every list exceeds $1/2$. Thus the smallest gap among all ordinates of the product is attained by the displayed pair and its reflection. Conjugation gives disjointness at negative ordinates. Opposite signs and the single central zero cannot give a coincidence. The orders in the product add, which yields $2(48+71+94)+1=427$ simple zeros.
\end{proof}

\begin{theorem}[Joint simplicity of the square and cube through height 256]\label{sf:joint256}
For $|\Im s|\le256$, the unitary symmetric square and cube of $43.\mathrm a1$ have disjoint nontrivial zeros. Their product has exactly $3079$ such zeros, all simple and critical. Its smallest gap is attained at the $585$th positive degree-three ordinate and the $827$th positive degree-four ordinate, and at their negatives. More precisely,
\[
     3200/D<\gamma_{3,585}-\gamma_{4,827}<3202/D.
\]
\end{theorem}
\begin{proof}
Theorems~\ref{sq:height256} and~\ref{sp:height256} give $638$ and $901$ simple positive roots. Signed bisection on the tilted Cauchy polynomials gives the strict width-$1/D$ enclosures in the preceding chapters, with the low cube brackets independently supplied by its two Bessel rotations. The two displayed ordinates have lower integer endpoints $3991966902$ and $3991963701$. The remaining $638\cdot901-1$ pairs from different lists all have lower gap bound at least $4373/D$. Within the two positive lists the least gap is strictly greater than $598680/D$ and $393926/D$, respectively. Thus the closest cross-factor pair also realizes the smallest gap in the entire product. The first positive root in both lists lies above one, so including the simple central cube zero and the reflected roots preserves this minimum. Adding the multiplicities gives $2(638+901)+1=3079$.
\end{proof}
For comparison, the same square/cube products have $575$ and $1345$ simple zeros through heights $64$ and $128$. The minimal positive cross-factor gaps multiplied by $D$ are in $(39671,39673)$ and $(22657,22659)$, respectively.

\begin{corollary}[Rankin--Selberg zeros through height 32]\label{sf:rankin32}
Let $\pi$ be the unitary representation of $\mathrm{GL}_2$ attached to $43.\mathrm a1$. At absolute ordinate at most $32$, the degree-eight Rankin--Selberg function $L(s,\pi\times\operatorname{Sym}^{3}\pi)$ has exactly $284$ nontrivial zeros and the degree-nine function $L(s,\operatorname{Sym}^2\pi\times\operatorname{Sym}^2\pi)$ has exactly $292$. Each zero is simple and critical. In the same range $L(s,\pi\times\pi)$ has $104$ simple critical zeros.
\end{corollary}
\begin{proof}
Put $\Lambda_\zeta(s)=\pi^{-s/2}\Gamma(s/2)\zeta(s)$. The unitary two-dimensional local parameter $V$ of $\pi$ has trivial determinant. The Clebsch--Gordan identity for $\operatorname{Sym}^a V\otimes\operatorname{Sym}^b V$ gives the summands $\operatorname{Sym}^{a+b-2j} V$ for $0\le j\le\min(a,b)$. At the bad prime the identity is an isomorphism of Weil--Deligne parameters including monodromy. Indeed it is equivariant for the Weil group and the algebraic $\mathrm{SL}_2$ of the parameter. At the real place it is the same tensor decomposition for the Weil representation. The local Langlands compatibility with factors of pairs therefore identifies the three completions, up to the fixed positive normalization constants used above, with
\[
\begin{aligned}
 \Lambda(\pi\times\operatorname{Sym}^3\pi,s)&=\Lambda_2(s)\Lambda_4(s),\\
 \Lambda(\operatorname{Sym}^2\pi\times\operatorname{Sym}^2\pi,s)&=\Lambda_\zeta(s)\Lambda_2(s)\Lambda_4(s),\\
 \Lambda(\pi\times\pi,s)&=\Lambda_\zeta(s)\Lambda_2(s).
\end{aligned}
\]
The trace-zero endomorphism representation of $\operatorname{Sym}^2 V$ has the two summands $\operatorname{Sym}^2 V$ and $\operatorname{Sym}^4 V$. Thus the degree-eight product also gives the adjoint function of $\operatorname{Sym}^2\pi$. The tensor-factor identities use the local correspondence of Harris--Taylor and Henniart, as stated in \href{https://nms.kcl.ac.uk/martin.dickson/bb/notes/cogdell_1.pdf}{Cogdell, \emph{Langlands correspondences for $\mathrm{GL}_n$}, Lecture~1}. The last two completions have simple poles at $0$ and $1$ because the other displayed factors are nonvanishing there. The three arithmetic conductors are $43^6,43^6,43^2$ by additivity of the conductor exponent under direct sums.

There are four positive zeta zeros below $32$. Their strict width-$1/D$ brackets have lower numerators $237141336$, $352691299$, $419612559$ and $510444718$. The count and brackets follow from the rigorous zeta-zero and counting functions of \href{https://flintlib.org/doc/acb_dirichlet/zeta-zeros.html}{FLINT}, evaluated at $128$ and $192$ bits. In particular, the fifth ordinate is strictly above $32$. The separations of the zeta zeros from the square and fourth-power lists, multiplied by $D$, are strictly greater than $1441791$ and $1434482$, respectively. Theorem~\ref{sf:joint32} separates the two symmetric powers. There are no central zeros in these three factors. The three positive counts are therefore $48+94$, $4+48+94$, and $4+48$. Conjugation gives the stated totals.
\end{proof}

\begin{corollary}[The Rankin square through height 256]\label{sf:rankin256}
The degree-four Rankin--Selberg function $L(s,\pi\times\pi)$ has exactly $1498$ nontrivial zeros at $|\Im s|\le256$. They are all simple and critical.
\end{corollary}
\begin{proof}
The same rigorous zeta count gives $111$ positive simple zeta zeros below $256$, with the $112$th above $256$. All $111$ have strict enclosures of width $1/D$ from the zeta isolation algorithm. Compare them with the $638$ positive intervals for $\Lambda_2$. Among the resulting $70818$ pairs, the closest is the $96$th zeta and $560$th symmetric-square ordinate, with lower interval numerators $3847643319$ and $3847665877$. Their distance is strictly between $22557/D$ and $22559/D$. All other cross-factor gaps exceed $45793/D$. Thus the two factors in Corollary~\ref{sf:rankin32} have disjoint zeros throughout this range. Their multiplicities add to $2(111+638)=1498$ and all are simple. The same count gives $262$ and $638$ at the intermediate heights $64$ and $128$.
\end{proof}

The final integer comparison repeats the entire separation calculation on a second set of brackets. In degree three the two Cauchy sampling lines are $31$ and $33$. In degree four the two computations change the Bessel rotation below height $64$ and the Cauchy tilt above it. In degree five the two methods in Section~\ref{sf:large} change both the coefficient algorithm and the right line. All $638$, $901$, and $94$ positive interval endpoint pairs agree in the respective proved ranges, and both collections give identical comparisons at every reported height.

\clearpage
\section{A cuspidal symmetric fifth power on \texorpdfstring{$\mathrm{GL}_6$}{GL(6)}}\label{fifth:chapter}
Let $E$ be the curve with Cremona label $11\mathrm a1$,
\[
                y^2+y=x^3-x^2-10x-20.
\]
Use $F(s)=L(\operatorname{Sym}^5 E,s+5/2)$ in unitary coordinates. The motivic center is~$3$.

\begin{theorem}\label{fifth:low}
For $|\Im s|\le28$, the zeros of the completed $F(s)$ comprise its simple central zero and $68$ simple conjugate pairs. All lie on $\Re s=1/2$. In motivic normalization,
\[
  4.2972319356146513 < L'(\operatorname{Sym}^5E,3) < 4.2972319356146514 .
\]
\end{theorem}

\subsection{The degree-six local factors}
The minimal model has $c_4=496$ and discriminant $-11^5$. Its rational $j$-invariant has negative $11$-adic valuation, which rules out complex multiplication. The modularity theorem and the symmetric-power theorem of Newton and Thorne identify $F$ with a cuspidal standard $L$-function of $\mathrm{GL}_6/\Q$. We use its analytic continuation, finite order, functional equation and local compatibility. The general inputs are Newton and Thorne, \emph{Symmetric power functoriality for holomorphic modular forms, II}, Publ. Math. IHÉS \textbf{134} (2021), Theorem A, and Martin and Watkins, \emph{Symmetric powers of elliptic curve $L$-functions}, ANTS VII (2006), §§2--4 for the local conventions.

For $p\ne11$, write $t=a_p(E)$ and let $\alpha+\beta=t$, $\alpha\beta=p$. Pair the six weights as
\[
       p^2(\alpha,\beta),\qquad p(\alpha^3,\beta^3),\qquad (\alpha^5,\beta^5).
\]
Put $u=t^2-2p$, $v=t^3-3pt$, $w=uv-p^2t$. The pairs have traces $p^2t,pv,w$ and determinant $p^5$. The motivic inverse Euler polynomial is
\begin{equation}\label{fifth:euler}
   (1-p^2tX+p^5X^2)(1-pvX+p^5X^2)(1-wX+p^5X^2).
\end{equation}
At $11$ the reduction is split multiplicative, since $-c_6=-20008\equiv1\pmod{11}$. The level-eleven newform has $a_{11}=1$. The special representation has one monodromy block of length two, whose fifth symmetric power has one block of length six. The conductor exponent is $6-1=5$. The one-dimensional kernel of monodromy has motivic Frobenius eigenvalue $1$, and the local sign is $(-1)^5=-1$. Thus the motivic inverse Euler factor at $11$ is $1-X$ and the unitary inverse factor is $1-11^{-5/2}X$.

At the real place the three induction characters have angular exponents $1,3,5$. Their gamma shifts are $1/2,3/2,5/2$ and their combined epsilon factor is $i^{2+4+6}=1$. Removing a positive constant independent of $s$ gives
\begin{equation}\label{fifth:arch}
 \Lambda(s)=Q^s\Gamma(s+1/2)\Gamma(s+3/2)\Gamma(s+5/2)F(s),
 \qquad Q=\frac{\sqrt{11^5}}{(2\pi)^3}.
\end{equation}
The completed function is entire, obeys $\Lambda(s)=-\Lambda(1-s)$ and takes conjugate values at conjugate points. Its zeros lie in $0\le\Re s\le1$ by absolute convergence on the right and reflection on the left.

Let $A_n$ be the motivic coefficients. The eigenvalues above have magnitude $p^{5/2}$ at good primes. Expanding the six geometric series and including the smaller factor at $11$ gives $|A_n|/n^{5/2}\le d_6(n)$. Each $A_n$ is computed from \eqref{fifth:euler} and multiplicativity. The input $t$ can be read from the normalized newform $\eta(z)^2\eta(11z)^2$. Euler's pentagonal formula computes its truncated series over the integers. Independently, let $T_k=\alpha^k+\beta^k$ and define $U_0=1$, $U_1=x$, $U_{j+1}=xU_j-yU_{j-1}$. The $k$th power sum of the six eigenvalues is $U_5(T_k,p^k)$. Thus
\[
   r A_{p^r}=\sum_{k=1}^r U_5(T_k,p^k) A_{p^{r-k}}.
\]
The second computation checks this integer identity at every retained good prime power. For example $a_2(E)=-2$ and $a_3(E)=-1$ give $A_2=8$ and $A_3=-16$.

\subsection{The central strip with an odd functional equation}
Put $\sigma=m+1/2$ and $w=s-1/2$. The reflected Cauchy integrals give
\[
 \Lambda(1/2+w)=\frac1{2\pi}\int_{\mathbb R}\Lambda(\sigma+iy)
       \left(\frac1{m+iy-w}-\frac1{m+iy+w}\right)dy.
\]
As in \eqref{sq:two-kernels}, the integrals on the ends of the rectangle vanish by the gamma decay and polynomial vertical growth. Define $G(w)=-i\Lambda(1/2+w)$, which is real on $w\in i\mathbb R$. Center a band at $it_0$, writing $w=it_0+\xi$. On the contour $|\Re\xi|\le1/2$, $|\Im\xi|\le4$, one has $|\xi|<5$. We establish the same uniform approximation for the disk $|\xi|\le5$.

Let $P_N(s)=\sum_{n\le N}A_n n^{-s-5/2}$, $b=1/d$, $J=dY$ and $s_j=\sigma+ijb$. Write $V_j$ for \eqref{fifth:arch} with $F(s_j)$ replaced by $P_N(s_j)$. Define
\[
 M_k=\sum_{j=-J}^J\frac{V_j}{[m+i(jb-t_0)]^{k+1}},
 \qquad
 P(\xi)=\frac b\pi\sum_{k=0}^K c_k\xi^k,\qquad
 c_k=\begin{cases}\Im M_k,&k\ \mathrm{even},\\-i\Re M_k,&k\ \mathrm{odd}.\end{cases}
\]
Indeed, reflection in the critical line changes $\xi$ to $-\overline\xi$. The reflected summand therefore contributes $-(-1)^k\overline{M_k}$ to the $k$th unscaled coefficient of $\Lambda$. Multiplication by $-i$ gives the stated projection.

Put $A(m)=Q^{m+1/2}m!(m+1)!(m+2)!$ and $C(m)=64\,2^{3m+6}A(m)$. Euler's gamma integral rotated through one radian bounds each $\Gamma(a+iy)$ by $2^a\Gamma(a)e^{-|y|}$, and $\zeta(2)^6<64$. In the strip of half-width $W$ about the numerical integration line one may use $C(m+W)e^{-3|y|}$. Here $Q>1$ and all real gamma arguments exceed two, so the right endpoint bounds the entire strip. The Cauchy kernels have denominators of modulus at least $m-W-5\ge2$. Their boundary norms, including $1/(2\pi)$, sum to less than $C(m+W)$. The four absolute allowances for the infinite trapezoidal formula, discarded samples, omitted coefficients and the local Taylor expansion are respectively
\begin{align*}
 &\frac{C(m+W)}{e^{2\pi W/b}-1},\qquad
    \frac{b C(m)e^{-3Y}}{e^{3b}-1},\\
 &64b(2J+1)A(m)N^{3/2-m},\qquad
  64b(2J+1)A(m)\frac{(5/m)^{K+1}}{1-5/m}.
\end{align*}
For the third bound use $\sum_{n>N}d_6(n)n^{-\sigma}\le\zeta(2)^6N^{2-\sigma}$. For the fourth expand each Cauchy denominator geometrically and bound $|V_j|$ on the unrotated line. Section~\ref{sf:grid} applies with six in the divisor bound. The exponent $\sigma+5/2=m+3$ is integral. The Fourier gridding remainder is added to the output balls. Direct sums at $j=0,\lfloor J/3\rfloor,J$ overlap the Fourier enclosures.

\begin{center}\small
\begin{tabular}{r|rrrrrrr}
\toprule
$|t|$ & $m$ & $N$ & $d$ & $Y$ & $W$ & $K$ & $-\log_{10}\varepsilon$\\
\midrule
4&20&36\,000&14&85&4&154&14\\
20&32&190\,000&21&155&4&220&36\\
20&30&205\,000&19&152&5&230&35\\
28&40&670\,000&29&210&4&258&62\\
28&39&770\,000&25&218&5&263&63\\
\bottomrule
\end{tabular}
\end{center}
In every row the four allowances sum to less than $\varepsilon/2$. In the last two rows their base-ten logarithms are less than $(-95,-77,-62,-72)$ and $(-120,-93,-64,-79)$ respectively. The two final settings use $1024$ and $1088$ bits.

\begin{proof}[Completion of the proof of Theorem~\ref{fifth:low}]
The first band is centered at zero, with the forced zero in its interior. On each contour the enclosed polynomial and the true completion have the same winding. Compose $P$ with each linear segment. Its maximum second derivative times one eighth of the squared segment length bounds the deviation from its chord, as in Section~\ref{sq:height32}. Subdivide until one coordinate of both chord endpoints exceeds the deviation plus $\varepsilon$ with the same strict sign. The true image and chord then lie in a common open half-plane through the origin. Summing the principal arguments of the successive polynomial endpoint ratios and dividing by $2\pi$ encloses exactly one integer.

\begin{center}
\begin{tabular}{r|rrrr}
\toprule
Band&$[-4,4]$&$[4,12]$&$[12,20]$&$[20,28]$\\
\midrule
Winding&7&17&23&25\\
Disjoint signed intervals&7&17&23&25\\
Accepted boundary segments&84&164&238&278\\
\bottomrule
\end{tabular}
\end{center}
The signed intervals come from $G(it)$ on the mesh $j/64$. Each sign change gives a distinct zero in the open interval. In the first band the positive brackets have lower numerators $125,163,223$ and width $1/64$. The interval $(-1/64,1/64)$ contains the exact zero at zero. The argument principle counts multiplicity, so equality with the number of disjoint intervals forces every zero to be simple. Reflection supplies the negative bands, giving $7+2(17+23+25)=137$.

For the derivative, Cauchy's inequality on $|\xi|=1$ bounds $|G'(0)-P'(0)|$ by $\varepsilon$. Equation~\eqref{fifth:arch} and $F(1/2)=0$ give $L'(\operatorname{Sym}^5 E,3)=\Lambda'(1/2)/(2\sqrt{Q})$. The final row encloses this value in the displayed interval.
\end{proof}

\begin{corollary}\label{fifth:sign}
The real even entire function $w\mapsto\Lambda(1/2+w)/w$, extended at zero by its derivative, is strictly positive on the whole real axis.
\end{corollary}
\begin{proof}
The rectangle in Theorem~\ref{fifth:low} excludes other real zeros in the critical strip. Outside this strip the functional equation and the Euler product exclude zeros. On the positive real ray sufficiently far to the right each paired factor in \eqref{fifth:euler} and every gamma factor is positive. Continuity gives strict positivity on both rays and the derivative certificate treats zero.
\end{proof}

\subsection{Moving the Dirichlet tail to a separate contour}\label{fifth:tilted}
For higher ordinates the exponential decay of the completed function makes each sign smaller. We give the error constants for an exponential factor on the entire Cauchy integrand. This also treats the companion symmetric cube needed to study tensor products.

Write $n=3$ or $5$, $g=(n+1)/2$, $D=2^{2g}$ and $Q_n=\sqrt{11^n}/(2\pi)^g$. Let
\[
 \Lambda_n(s)=Q_n^s\prod_{j=0}^{g-1}\Gamma(s+1/2+j)
                           L(\operatorname{Sym}^n E,s+n/2).
\]
The sign is $+1$ for the cube and $-1$ for the fifth power. For the cube the archimedean factors have exponents $1,3$ and root number $i^6=-1$, while the split Steinberg parameter at $11$ has sign $(-1)^3=-1$ and conductor exponent~$3$. The two good-prime cubical pairs have traces $pt,t^3-3pt$ and determinant $p^3$, with coefficients bounded by $d_4$ after unitary normalization.

Fix $0<\theta<g\pi/2$. Put $\mathcal H_n(w)=e^{-i\theta w}\Lambda_n(1/2+w)$ and use $-i\mathcal H_5$ or $\mathcal H_3$ on the critical line. There the new multiplier is positive. Both normalized functions $B$ satisfy $B(-\overline w)=\overline{B(w)}$. Orienting the two Cauchy lines as before gives the integral of the right-hand value of $B$ divided by $m+iy-w$, plus the complex conjugate right-hand value divided by $m-iy+w$, with common measure $dy/(2\pi)$. In terms of the unscaled samples of $\mathcal H_n$, the fifth-power projections are those above. The cubical projections are $\Re M_k$ for even $k$ and $i\Im M_k$ for odd $k$.

Write $m_t+1/2$ for a second line farther to the right. With $0\le u<\pi/2$, define
\[
 A_n(x,u)=Q_n^{x+1/2}
       \frac{\prod_{j=0}^{g-1}(x+j)!}
            {(\cos u)^{gx+g(g+1)/2}}.
\]
The rotated gamma estimate yields the bound
$D A_n(x,u)\exp[-(gu-\theta)|y|]$ for the weighted completed Dirichlet polynomial on the line $\Re s=x+1/2$ whenever $u\ge\theta/g$. Set
\[
 u_0=\theta/g,\quad u_1=u_0+2/5,\quad u_t=u_0+1/250,
 \qquad c=2g/5,\quad c_t=g/250 .
\]
Every setting used below has $u_1<\pi/2$. The four allowances are
\begin{align}
 E_1&=\frac{D\max(1,Q_n^{-2W})A_n(m+W,u_1)}{e^{2\pi W/b}-1},&
 E_2&=\frac{bD A_n(m,u_1)e^{-cY}}{e^{cb}-1},\notag\\
 E_3&=\frac{2D A_n(m_t,u_t)N^{3/2-m_t}}{\pi c_t(m_t-5)},&
 E_4&=Db(2J+1)A_n(m,u_0)
                \frac{(5/m)^{K+1}}{1-5/m}.\label{fifth:errors}
\end{align}
For $E_3$, split the exact Cauchy integral into the polynomial and the omitted Dirichlet tail \emph{before} discretization. Each tail integral can be shifted from $m+1/2$ to $m_t+1/2$ since its Dirichlet series converges absolutely there and exponential decay makes the horizontal integrals vanish. Apply
$\sum_{k>N}d_{2g}(k)k^{-m_t-1/2}\le D N^{3/2-m_t}$.
On the shifted line the two denominators are at least $m_t-5$ throughout $|\xi|\le5$, and integrating $e^{-c_t|y|}$ gives $2/c_t$. Counting both integrals gives $E_3$.

The remaining integral contains only the polynomial. In its quadrature strip $\max(1,Q_n^{-2W})$ bounds the change of the exponential base relative to the right endpoint; the gamma norms are bounded there by their right-endpoint values. For the settings used here $4/[\pi c(m-W-5)]<1$. This bounds both boundary norms for both kernels by the numerator of $E_1$. The omitted samples give $E_2$, and geometric expansion about $it_0$ gives $E_4$. Each retained Fourier value is enclosed using Section~\ref{sf:grid}. At each of three frequencies the computed interval also overlaps the direct sum of its first $\min(N,32768)$ terms enlarged by $D\min(N,32768)^{3/2-m}$ for the rest.

Finally discard weighted samples whose absolute values are strictly less than $\varepsilon/[64b(2J+1)]$. Both Cauchy kernels together contribute less than $\varepsilon/64$ from these samples, since $m-5>1$. Every row checks $\sum E_j<\varepsilon/2$, so the entire disk estimate remains strictly within $\varepsilon$. The winding and axis signs are then obtained from this scaled polynomial exactly as for Theorem~\ref{fifth:low}.

\begin{theorem}[The companion symmetric cube]\label{fifth:cube}
The completed symmetric cube of $11\mathrm a1$ has exactly $648$ zeros through $|\Im s|\le132$. All are simple, with $324$ positive ordinates on the critical line and their negatives. Its central value is strictly positive.
\end{theorem}
\begin{proof}
For the first computation use
\[
 (m,m_t,N,d,Y,W,K,\theta)=(31,82,43000,15,586,6,234,9/5)
\]
with $\varepsilon=10^{-88}$ and $768$ bits. The second uses
\[
 (32,80,44000,16,600,6,232,19/10)
\]
with $\varepsilon=10^{-83}$ and $800$ bits. In order, the base-ten logarithms of the four terms of \eqref{fifth:errors} are less than $(-113,-96,-92,-100)$ and $(-118,-91,-87,-96)$. The two central enclosures give $1.09611832799977<\Lambda_3(1/2)<1.09611832799979$. The band centered at zero has winding $4$. The successive positive bands starting with $[4,12]$ have windings
\[
 10,\ 14,\ 17,\ 16,\ 20,\ 18,\ 21,\ 21,\quad
 21,\ 22,\ 23,\ 23,\ 24,\ 23,\ 24,\ 25.
\]
In each band these integers equal the number of disjoint certified sign changes. The independent dyadic brackets for every zero have a nonempty open intersection. The central bracket is absent in each computation.
\end{proof}

\subsection{Rankin--Selberg products and the fifth tensor power}\label{fifth:tensors}
At each place, let $V$ denote the two-dimensional unitary local parameter of $E$. The central character is trivial. Clebsch--Gordan gives
\begin{align*}
 V\otimes\operatorname{Sym}^4V
   &\simeq\operatorname{Sym}^5V\oplus\operatorname{Sym}^3V,\\
 \operatorname{Sym}^2V\otimes\operatorname{Sym}^3V
   &\simeq\operatorname{Sym}^5V\oplus\operatorname{Sym}^3V\oplus V,\\
 V^{\otimes5}
   &\simeq\operatorname{Sym}^5V\oplus
        (\operatorname{Sym}^3V)^{\oplus4}\oplus V^{\oplus5}.
\end{align*}
For the last identity, successively tensor the symmetric powers by $V$ and use $V\otimes\operatorname{Sym}^kV=\operatorname{Sym}^{k+1}V\oplus\operatorname{Sym}^{k-1}V$ for $k\ge1$. In motivic coordinates each lower-weight term carries the determinant power required to restore weight five. For example the middle line has motivic factors $L(\operatorname{Sym}^5 E,s)$, $L(\operatorname{Sym}^3E,s-1)$ and $L(E,s-2)$. The decompositions hold for the entire Weil--Deligne parameters, including monodromy at $11$. Hence all Euler, conductor and gamma factors agree, and these products have conductors $11^8$, $11^9$ and $11^{22}$, respectively. The first two are the Rankin--Selberg functions on $\mathrm{GL}_2\times\mathrm{GL}_5$ and $\mathrm{GL}_3\times\mathrm{GL}_4$.

A common zero of two constituents would increase the multiplicity of the associated product. To exclude this, intersect the two signed brackets of each symmetric-power zero, using exact rational arithmetic. For the base curve use its positive intervals from the degree-two certificate through $4096$, whose unit is $2\pi\cdot3412/2^{54}$. Round both endpoints outward onto $2^{-48}\mathbb Z$, including the uncertainty in $\pi$, and retain those below the chosen height. The program checks strictly that the height bisects no interval. The real central value of $L(E,s)$ in motivic coordinates is positive by the Mellin integral of $\eta(iy)^2\eta(11iy)^2$, positive for every $y>0$. Sorting all three positive lists and checking that consecutive enclosing intervals are disjoint proves the asserted multiplicities. Conjugation treats the negative ordinates.

\begin{theorem}[The common window through height one hundred]\label{fifth:high}
The standard cuspidal $\mathrm{GL}_6$ function of $\operatorname{Sym}^5 E$ has $725$ zeros through $|\Im s|\le100$, all simple on the unitary critical line, including a simple central zero. The two Rankin--Selberg products above have $1181$ and $1369$ zeros respectively, all simple on that line. The fifth tensor power has $1369$ distinct zeros and $3489$ with multiplicity. The multiplicities are $1$ on the fifth-power spectrum, $4$ on the cubic spectrum and $5$ on the elliptic spectrum.
\end{theorem}
\begin{proof}
In \eqref{fifth:errors} take $n=5$. The first computation uses
\[
(m,m_t,N,d,Y,W,K)=(40,104,5060000,21,638,7,338)
\]
with $\theta=23/10$, $\varepsilon=10^{-111}$ and $1216$ bits. The second uses
\[
 (42,107,4950000,21,660,7,343)
\]
with $\theta=9/4$, $\varepsilon=10^{-112}$ and $1280$ bits. The respective base-ten logarithms of the allowances are less than $(-146,-122,-113,-126)$ and $(-136,-123,-115,-127)$. One run obtains the elliptic prime traces from the Hasse-interval group-order certificate through $10557515$, the other from the eta product. The latter also checks both symmetric-power coefficient recurrences. Each computation gives winding $7$ in the central band. The successive positive bands have windings
\[
 17,\ 23,\ 25,\ 28,\ 29,\ 31,\ 33,\ 32,\ 34,\ 34,\ 36,\ 37.
\]
Each winding is exhausted by sign changes. The two computations intersect strictly for all $362$ positive brackets, each of width $2^{-32}$. The three total counts through this height, ordered by symmetric exponents $1,3,5$, are $(188,456,725)$. Among distinct positive ordinates of the indicated pairs of constituents, the certified minimum separations exceed $0.00595$ for exponents $(1,3)$, $0.000354$ for $(1,5)$ and $0.000327$ for $(3,5)$. Adding the counts with their Clebsch--Gordan multiplicities proves the result.
\end{proof}

\subsection{Recovering Frobenius traces from symmetric-power zeros}\label{fifth:inverse}
Consider the inverse question in which the individual elliptic traces are withheld. The available inputs are a finite list of symmetric-power zero intervals, the conductor, the gamma shifts, the central multiplicity, the reduction type at $11$ and the Hasse bounds. For $n=3,5$ let $\ell_n(q^j)$ be the normalized $j$th power sum of the local eigenvalues at $q$. For good $q$ and a proposed elliptic trace $t$, it equals
\[
    \ell_n(q^j;t)
       =q^{-jn/2}U_n(T_j,q^j),\qquad
       T_0=2,\quad T_1=t,\quad T_{j+1}=tT_j-qT_{j-1}.
\]
At $11$ we use only $|\ell_n(11^j)|\le11^{-jn/2}$. Write $K_n(z)=\log Q_n+\tfrac12\sum_{k=1}^g[\psi(k+iz)+\psi(k-iz)]$. For a positive integer $p$, put $L=\log p$. The entire even test functions used here are
\[
 h(t)=\tfrac12\left(e^{-(t-c)^2/H^2}+e^{-(t+c)^2/H^2}\right)\cos(tL),
 \qquad c\ge0,\quad H>0.
\]
Define $\Phi(v)=e^{-H^2v^2/4}\cos(cv)$, and for each prime power $m=q^j$ put
\[
   W_{p,H,c}(m)=\frac{\log q}{\sqrt m}
       [\Phi(\log (m/p))+\Phi(\log(mp))].
\]
The explicit formula, with the Fourier transform $\widehat h(v)=\int h(t)e^{-ivt}\,dt$, gives the exact identity
\begin{equation}\label{fifth:inverseformula}
 \sum_{q,j}\ell_n(q^j)W_{p,H,c}(q^j)
   =\frac{2\sqrt\pi}{H}
     \left[\frac1\pi\int_\R h(t)K_n(t)\,dt
          -\sum_{\rho}h\big((\rho-1/2)/i\big)\right].
\end{equation}
Indeed $\widehat h(v)=H\sqrt{\pi}[\Phi(v-L)+\Phi(v+L)]/2$. The absence of poles eliminates a polar term. The Gaussian estimates below apply also to zeros outside the critical line.

For both powers one has
\[
 \sum_{\rho=\beta+i\gamma}\frac1{4+\gamma^2}
 \ \le\ \Re\frac{\Lambda_n'}{\Lambda_n}(2)<20 .
\]
To see the first inequality, pair the zeros about $1/2$ in the Hadamard product. The paired factors have summable quadratic tails. Evenness or oddness about the center rules out an extra linear exponential. Each zero contributes
$(2-\beta)/[(2-\beta)^2+\gamma^2]\ge1/(4+\gamma^2)$ to the real logarithmic derivative. For the second inequality, $|\log Q_n|<1$, each gamma logarithmic derivative at $2$ is less than $2$, and the finite-part logarithmic derivative is bounded by $6\sum_{k\ge2}(\log k)/k^2<12$. Thus the contribution of all zeros beyond $|\gamma|=T$ is bounded by
\[
 20\sqrt p\,(4+T^2)
       \exp\!\left(-\frac{(T-c)^2}{H^2}+\frac{1}{4H^2}\right),
 \qquad T-c>H .
\]
Here $|\cos((x+iy)L)|\le\sqrt p$ for $|y|\le1/2$, and $(4+x^2)e^{-(x-c)^2/H^2}$ decreases for $x\ge T$. The central zero, when present, contributes $e^{-c^2/H^2}$.

For the integral in \eqref{fifth:inverseformula}, the series for the digamma function gives $|\psi(a+ix)-\psi(a)|\le6|x|$ for $a\ge1/2$, by comparison with $\sum_{k\ge0}(k+1/2)^{-2}<6$. The real values $|\psi(a)|$ are below $3$ on $[1/2,7/2]$. Thus on the two lines $\Im z=\pm1/2$ the modulus of $K_n(z)$ is at most $10+18|\Re z|$. Use trapezoidal spacing $b=1/18$ and cutoff $M=c+10H$. For the unnormalized full-line integral the discretization and tail bounds are
\[
 \frac{2\sqrt p\,e^{1/(4H^2)}
   [(10+18c)H\sqrt\pi+18H^2]}{e^{\pi/b}-1},\qquad
 2H^2\left(\frac{5+9c}{M-c}+9\right)e^{-(M-c)^2/H^2}.
\]
The first follows by integrating the two boundary norms. Beyond $M$ the positive majorant decreases, so its integral bounds the omitted sampled tail. These two bounds include the negative half-axis. Finally the combined arithmetic sum over $m>4p$ is at most
\[
                     24p^2e^{-H^2\log^2(4)/8}.
\]
In fact $|\ell_n(m)|\le6$ and each Gaussian is bounded by $e^{-H^2\log^2(4)/8}(m/p)^{-2}$ on this tail whenever $H^2\log(4)/8\ge2$. The sum of $(\log m)m^{-5/2}$ is less than two.

For each remaining prime $q$ start with all integers $|t|\le2\sqrt q$. Group its powers $q^j\le4p$ before taking the interval hull of the possible contributions in \eqref{fifth:inverseformula}. Subtract the enclosures for all other primes from the computed spectral side. Delete a candidate for $q$ if its contribution misses this resulting interval. The true trace survives each deletion by the explicit formula; iteration can only reduce the lists. The final point-count comparison uses no information during the deletion process.

\begin{proposition}\label{fifth:inverse100}
The height-$100$ fifth-power spectrum determines the first motivic coefficient of every good prime through $89$ under these constraints. It also determines the underlying elliptic trace at all these primes except possibly $31$. Adjoining the height-$132$ cubical spectrum determines every good elliptic trace through $97$ using centered Gaussians, and through $181$ using translated Gaussians.
\end{proposition}
\begin{proof}
Use $c=0$, $H=15,18,21,24,27,30$ and test centers $\log p$ for good primes $p\le113$. Interval elimination gives $22$ singleton elliptic traces and $23$ singleton fifth-power coefficients from the degree-six spectrum. At $31$ the surviving traces are $-3$ and $7$, both with $U_5(t,31)=-5544$. The independent degree-four equations leave $7$ alone and give $24$ singleton elliptic traces in total. For example $a_2=-2$ and $a_3=-1$ are recovered from the fifth-power spectrum alone. For the translated calculation use all good prime test centers through $193$ and adjoin the following $(H,c)$ pairs to the centered ones:
\[
 \begin{gathered}
 (15,15),(15,28),(13,39),(12,48),(10,57),(9,64),(8,71),\\
 (15,36),(13,52),(12,64),(10,76),(9,85),(8,95).
 \end{gathered}
\]
Keep only shapes satisfying $T-c>H$ for the spectrum being used. These $1591$ constraints remove respectively $642,292,144,156,86,35,14$ candidates on the successive passes before stabilization. The $41$ singletons include every good prime through $181$. The ledgers give the surviving candidate sets and recovered integers.
\end{proof}

\begin{theorem}[The final common rectangle]\label{fifth:final}
Through $|\Im s|\le132$ the completed fifth symmetric power has $1023$ simple critical-line zeros, its central zero and $511$ conjugate pairs. The $\mathrm{GL}_2\times\mathrm{GL}_5$ product has $1671$ simple zeros and the $\mathrm{GL}_3\times\mathrm{GL}_4$ product has $1945$ simple zeros, all on the critical line. The fifth tensor power has $1945$ distinct zeros and $4985$ with multiplicity.
\end{theorem}
\begin{proof}
The two runs of \eqref{fifth:errors} use respectively
\[
 (m,m_t,N,d,Y,W,K)=(42,120,9800000,21,720,8,375)
\]
with $\theta=5/2$, $\varepsilon=10^{-132}$ and $1344$ bits, and
\[
 (43,122,9050000,22,730,8,368)
\]
with $\theta=49/20$, $\varepsilon=10^{-131}$ and $1376$ bits. The four base-ten logarithms are below $(-172,-142,-135,-152)$ and $(-190,-144,-134,-145)$ respectively. The first input is constructed from the exact eta product. The second uses the independently verified group orders and Newton recurrence. Both repeat the full contour calculation from zero, recovering all the earlier band windings. The four successive bands from $100$ to $132$ have windings $36,38,37,38$, each equal to its number of signed intervals. Every positive bracket from the two runs intersects in the same open dyadic interval of width $2^{-32}$. The factor counts are $274,648,1023$. The minimum gaps between positive zeros of different factors exceed $0.00479$, $0.000193$ and $0.0000111$ for exponent pairs $(1,3),(1,5),(3,5)$, respectively. The identities in Section~\ref{fifth:tensors} give the reported totals with the same multiplicities as before.
\end{proof}

\clearpage
\section{A primitive mixed elliptic tensor over the rationals}\label{mix:chapter}
Consider the two curves
\[
 E\colon y^2+y=x^3-x^2,\qquad F\colon y^2+xy+y=x^3+x^2.
\]
Their invariants $(c_4,\Delta)$ are $(16,-11)$ and $(1,-15)$. Both models have good reduction away from the indicated discriminant and multiplicative reduction at its prime divisors. Their conductors are $11$ and $15$. These small equations lie in the usual isogeny classes $11\mathrm a$ and $15\mathrm a$, respectively. The database labels are \href{https://www.lmfdb.org/EllipticCurve/Q/11/a/3}{11.a3} and \href{https://www.lmfdb.org/EllipticCurve/Q/15/a/7}{15.a7}. We use unitary coordinates for $L(E\times F,s)$ throughout. The arithmetic central line is $\Re s=3/2$.

\subsection{The tensor is cuspidal and stays irreducible over number fields}
Let $\pi_E,\pi_F$ be the unitary forms attached to the two curves by modularity. Their nonintegral $j$-invariants, $-4096/11$ and $-1/15$, exclude CM, so both forms are nondihedral. At $11$ the local parameter of $E$ has nonzero monodromy and that of $F$ has zero monodromy. A twist by a local character preserves the rank of the monodromy operator. The forms are therefore inequivalent under every Hecke-character twist. Ramakrishnan's Theorem M gives a cuspidal tensor transfer $\Pi=\pi_E\boxtimes\pi_F$ on $\mathrm{GL}_4/\Q$, with the local tensor parameters. Its standard completion is entire and self-dual. See \href{https://arxiv.org/abs/math/0007203}{D.~Ramakrishnan, \emph{Modularity of the Rankin--Selberg $L$-series, and multiplicity one for $\mathrm{SL}(2)$}}, Annals of Mathematics 152 (2000), 45--111.

\begin{proposition}\label{mix:strong}
For every number field $K$ and every auxiliary prime $\ell\notin\{3,5,11\}$, the restriction of $V_\ell(E)\otimes V_\ell(F)$ to $G_K$ is absolutely irreducible. The joint algebraic monodromy group of the two Tate modules is the full subgroup of $\mathrm{GL}_2\times\mathrm{GL}_2$ on which the two determinants agree. In particular the tensor cannot arise by induction from a proper finite extension of $K$.
\end{proposition}
\begin{proof}
Let $H$ be the Zariski closure over $\overline{\Q}_\ell$ of the joint image in $\mathrm{GL}_2\times\mathrm{GL}_2$. Each projection is all of $\mathrm{GL}_2$, by Serre's open image theorem and the finite index of $G_K$ in $G_\Q$. At a prime of $K$ over $11$, inertia acts trivially on $F$. Its action on $E$ contains a nontrivial unipotent element, by the Tate-curve description of multiplicative reduction. Passing to a finite extension only replaces the unipotent inertia group by a finite-index subgroup. Thus the subgroup $\{g:(g,1)\in H\}$ is normal in $\mathrm{GL}_2$ and contains a nonidentity unipotent. The powers of such an element are Zariski dense in its one-dimensional unipotent group. Conjugating gives both upper and lower unipotent groups, which generate $\mathrm{SL}_2$. Interchanging the two curves at a prime above $3$ gives the other copy of $\mathrm{SL}_2$. Hence $H$ contains $\mathrm{SL}_2\times\mathrm{SL}_2$. Both determinants are the cyclotomic character. For any $(A,B)$ of equal determinant, lift $A$ to some $(A,B')\in H$ using the first projection, then multiply by $(1,(B')^{-1}B)\in H$. This proves the asserted equality of algebraic groups. The external tensor of the two standard representations is irreducible on this product, proving the assertion. Induction from a proper extension would split into smaller summands on the Galois group of its normal closure, contradicting the same assertion for that finite extension.
\end{proof}

The two alternating polarization forms give a symmetric nondegenerate form on the external tensor by multiplying their pairings on pure tensors. After unitary normalization the tensor is of orthogonal type. Its motivic Hodge numbers in bidegrees $(2,0),(1,1),(0,2)$ are $1,2,1$, including the repeated middle weight.

\subsection{All local factors and two constructions of the coefficients}
Write $a=a_p(E)$, $b=a_p(F)$ and take the motivic Euler variable $X$. Multiplying the four products of the two sets of Frobenius eigenvalues gives the good reciprocal polynomial
\begin{equation}\label{mix:good}
   1-abX+\{p(a^2+b^2)-2p^2\}X^2-p^2abX^3+p^4X^4.
\end{equation}
At any of $3,5,11$ exactly one curve has multiplicative reduction. If its trace is $e=\pm1$ and the other has eigenvalues $u,v$, the retained eigenvalues are $eu,ev$. The local reciprocal polynomial is $1-abX+pX^2$. More explicitly, $(a,b)=(-1,-1),(1,1),(1,-4)$ at $3,5,11$, respectively. Each of these local parameters consists of two special blocks of dimension two, giving conductor exponent two. In unitary normalization the two unramified characters on these blocks have values $e u/\sqrt p$ and $e v/\sqrt p$. The two Steinberg signs multiply to $e^2uv/p=1$. The conductor is consequently $165^2=27225$.

At infinity the tensor decomposes as $I_1\otimes I_1=I_2\oplus\mathbf1\oplus\operatorname{sgn}$. The factors are $\Gamma_{\mathbb C}(s+1)\Gamma_{\mathbb R}(s)\Gamma_{\mathbb R}(s+1)$ and the sign is $i^3\cdot1\cdot i=1$. Equivalently the four real gamma shifts are $0,1,1,2$. Put $\delta=4\pi^2/165$. Up to a fixed positive constant the even completion is
\begin{equation}\label{mix:complete}
 \Lambda(s)=\delta^{-s}\Gamma(s)\Gamma(s+1)\sum_{n\ge1}\lambda_n n^{-s},
 \qquad \lambda_n=A_n/n,\qquad \Lambda(s)=\Lambda(1-s).
\end{equation}
Here $A_n$ are integers and the series converges absolutely for $\Re s>1$. Hasse's bound and the explicit bad factors show that the normalized reciprocal roots have modulus at most one everywhere. Thus $|\lambda_n|\le d_4(n)$ for all $n$.

For the first arithmetic construction, count the two Weierstrass equations over $\mathbb F_p$. At odd primes the discriminants in $y$ are $1+4x^2(x-1)$ and $(x+1)^2+4x^2(x+1)$. Summing their quadratic characters, with a minus sign, gives the two traces, also at multiplicative reduction. At $2$ enumeration of the two affine equations gives $a_2=-2$ and $b_2=-1$. Inverting the local reciprocal polynomials as geometric series and multiplying across distinct primes yields $A_n$.

A second construction obtains every $A_n$ without a finite-field point count. The normalized weight-two newforms are
\[
 f(z)=\eta(z)^2\eta(11z)^2=\sum_{n\ge1}a_nq^n,\qquad
 g(z)=\eta(z)\eta(3z)\eta(5z)\eta(15z)=\sum_{n\ge1}b_nq^n.
\]
For both products the two eta congruence sums equal $24$, the character is trivial since the product of the level divisors with their exponents is a square, and the cusp orders are positive. Both products start with $q$. The genus formula gives $\dim S_2(\Gamma_0(11))=\dim S_2(\Gamma_0(15))=1$, so modularity of the curves identifies them uniquely. This uses the eta-product modularity and cusp-order criteria, for example K.~Ono, \emph{The Web of Modularity}, Theorems~1.64--1.65.

The exact convolution giving the tensor coefficients is
\begin{equation}\label{mix:squares}
 A_n=\sum_{\substack{h^2\mid n\\(h,165)=1}}h^2\,a_{n/h^2}b_{n/h^2}.
\end{equation}
Indeed, at a good prime let $U_j,V_j$ be the degree-$j$ complete symmetric polynomials in the two pairs of eigenvalues. Multiplication of geometric series gives
\[
 \sum_{j\ge0}U_j V_j X^j
   =\frac{1-p^2X^2}{\prod_{e\in\{\alpha,\beta\},\ f\in\{u,v\}}(1-efX)}.
\]
This identity also follows by expanding $U_j=(\alpha^{j+1}-\beta^{j+1})/(\alpha-\beta)$ and its counterpart, then using polynomial continuity when eigenvalues coincide. At a bad prime one eigenvalue of its retained factor is zero if padded to a pair; the numerator is then one. Multiplying the identities proves \eqref{mix:squares}. The second implementation builds the eta series by Euler's pentagonal number theorem and exact truncated integer polynomial multiplication, then sums the square divisors. This construction changes both the source of the prime data and the reconstruction of coefficients at composite indices.
Term-by-term comparison through $n=350000$ gives exact equality of all $350000$ motivic integers $A_n$ obtained in these two ways. The analytic error bounds below use the all-$n$ divisor bound proved from the Euler factors.

\subsection{Two contours for the mixed gamma factor}\label{mix:contours}
Write $w=s-1/2$ and take $H(w)=e^{-i\theta w}\Lambda(1/2+w)$. The exponential has no zeros and its value on $w=it$ is $e^{\theta t}>0$. Thus the signed real values of $H(it)$ locate the zeros of the original completion there. Apply the two-kernel Cauchy formula of \eqref{sq:two-kernels} with this tilt. The finite Dirichlet sum is evaluated at $\sigma$ and the omitted infinite tail is bounded on a separate line $\sigma'$ by Lemma~\ref{sp:taillemma}. The parameters in the two computations are as follows. Here $b$ is the frequency spacing, $Y$ the frequency cutoff, $d$ the integration-strip half-width and $K$ the degree for each translated polynomial.
\begin{center}
\begin{tabular}{lrrrrrrrrr}
\toprule
Arithmetic & $\theta$ & $\sigma$ & $\sigma'$ & $N$ & $1/b$ & $Y$ & $d$ & $K$ & Bits\\\midrule
Point counts & $7/4$&27&99&350000&22&910&7&300&1024\\
Eta products & $2$&29&109&430000&24&1000&8&320&1088\\\bottomrule
\end{tabular}
\end{center}
Let $r=\sigma-1/2$, $J=Y/b$, and $A(x)=\delta^{-x}(x-1)!x!$ at the integral abscissae used in the bounds. Choose gamma rotation angles $(\phi,\phi')$ on the enlarged integration strip and the omitted-series line, respectively, and certified upper bounds $(B,B')$ for their reciprocal cosines. The first row uses $(5/4,1;4,2)$ and the second row $(27/20,11/10;5,9/4)$. Set $c=2\phi-\theta>0$, $c'=2\phi'-\theta>0$ and $C(x)=16 B^{2x+1}A(x)$. On $\Re s=x$, Euler's integral rotation yields $C(x)e^{-c|\Im s|}$ after multiplication by either tilted weight and the finite Dirichlet sum, since $\zeta(2)^4<16$. On the two continuous strip boundaries, the same bound holds with $C(\sigma+d)$, which dominates both gamma norms there. The absolute values of the two Cauchy denominators exceed $\sigma-d-1$. After both integrals and their normalizing factors are included, the summed boundary norm is bounded by $C(\sigma+d)$ since $4/[\pi c(\sigma-d-1)]<1$. On the sampling line the analogous constant in front of the two omitted geometric frequency tails is also below one. For the nondecaying bound on the retained samples, rotating by angle one gives $A_* = 2^{2\sigma+1}A(\sigma)$ because $\theta\le2$ and $1/\cos(1)<2$.

For each eight-unit positive rectangle use the translated variable $w=it_0+z$, $|z|<5$. The four absolute error bounds, covering the infinite trapezoidal rule, omitted frequencies, omitted coefficients and Taylor expansion of both kernels, are therefore
\begin{align*}
 &\frac{C(\sigma+d)}{e^{2\pi d/b}-1},\qquad
 \frac{bC(\sigma)e^{-cY}}{e^{cb}-1},\\
 &\frac{32 N^{2-\sigma'} A(\sigma') (B')^{2\sigma'+1}}{\pi c'(\sigma'-1)},\qquad
 16b(2J+1) A_*\frac{(5/r)^{K+1}}{1-5/r}.
\end{align*}
Their base-ten logarithms are strictly below $(-280.95,-190.55,-106.81,-125.48)$ and $(-362.60,-182.97,-107.99,-141.84)$, respectively. The two computations use the same uniform allowance $\varepsilon=10^{-102}$. Each checks that the four bounds sum to less than $\varepsilon/2$.

For precision about the polynomial actually evaluated, form the weighted finite samples
\[
 W_j=e^{-i\theta r+\theta jb}\delta^{-(\sigma+ijb)}\Gamma(\sigma+ijb)
       \Gamma(\sigma+1+ijb)\sum_{n\le N} A_n n^{-\sigma-1-ijb}.
\]
The finite Fourier computation and its Taylor remainder are those of Section~\ref{sf:grid} with coefficient normalization $A_n/n$. Three further direct sums, at $j=0,\lfloor J/3\rfloor,J$, must overlap the Fourier enclosures. A sample can be discarded when its entire weighted modulus is below $\varepsilon/[64b(2J+1)]$, together with its associated reflected term. The resulting absolute loss in the two Cauchy kernels is less than $\varepsilon/64$. Put
\[
 B_k=\sum_{j\text{ retained}}\frac{W_j}{(r+ijb-it_0)^{k+1}}.
\]
The coefficient of $z^k$ is $b\{B_k+(-1)^k\overline{B_k}\}/(2\pi)$. The four analytic bounds, omission allowance and interval evaluation of the remaining finite terms give the required uniform enclosure. The half-plane subdivision and winding computation in Section~\ref{sq:cauchy} now apply unchanged.

\begin{theorem}[A mixed cuspidal tensor through height 128]\label{mix:128}
The primitive degree-four tensor of $11.\mathrm a3$ and $15.\mathrm a7$ has exactly $744$ nontrivial zeros of absolute ordinate at most $128$. They form $372$ conjugate pairs, each is simple, and their real part is $1/2$ in unitary coordinates. Each positive ordinate lies in an open interval of width $2^{-24}$.
\end{theorem}
\begin{proof}
Partition $[0,128]$ into consecutive eight-unit ordinate ranges. On each closed critical rectangle the boundary of the local polynomial and the corresponding completion lie in a common open coordinate half-plane on every accepted subdivision. Adding their certified argument increments gives the winding. The table gives the winding $W$ and the numbers $P,Q$ of accepted panels for the point-count and eta computations. The central positive value is outside all zero intervals.
\begin{center}
\begin{tabular}{rrrr|rrrr}
\toprule
Height &$W$&$P$&$Q$& Height&$W$&$P$&$Q$\\\midrule
8&10&86&86&72&26&246&242\\
16&16&128&128&80&26&256&258\\
24&19&152&146&88&25&268&268\\
32&20&162&176&96&28&270&276\\
40&22&184&182&104&26&280&284\\
48&23&196&190&112&28&288&288\\
56&24&220&220&120&27&290&292\\
64&24&234&226&128&28&288&286\\\bottomrule
\end{tabular}
\end{center}
Each row supplies exactly $W$ disjoint intervals where the real value of the tilted completion strictly changes sign. Testing their midpoint signs and retaining the half with opposite endpoint signs narrows every positive root to width $2^{-24}$. The two computations give the same $372$ endpoint pairs. The first and last have lower endpoint numerators $30545744$ and $2142850567$ on denominator $2^{24}$. In the very first band the ten coarse lower numerators on denominator $64$ are $116,168,225,268,312,343,381,437,464,495$, each with upper numerator one greater. These bounds are strict. Every band winding already counts multiplicity and possible off-line zeros, so equality with the number of signed intervals forces simplicity and exhausts the band. The nonzero boundaries and conjugation complete the count in the closed critical strip. The Euler product is nonzero for $\Re s>1$, and the functional equation covers $\Re s<0$.
\end{proof}

In the normalization of \eqref{mix:complete} the two independent central enclosures give $|\Lambda(1/2)-a|<2\cdot10^{-68}$ for
\[
 a=5.7421496838267976050759825618437234180494853235861402292952152884882.
\]
The completion is positive on the entire real axis. On $s>1$ all local factors and the gamma factors are positive; the functional equation gives positivity for $s<0$. Real continuity and the nonvanishing on the real segment $[0,1]$ supplied by the lower contour fill the remaining interval. The finite part has one simple trivial zero at $s=0$ and double zeros at each negative integer. These are all its real zeros, since their multiplicities are exactly the pole orders of $\Gamma(s)\Gamma(s+1)$.

\subsection{An even Bessel--Mellin inversion}\label{mix:bessel}
The two right-line computations use different exponential weights but share the Cauchy formula. The mixed archimedean parameters also permit an independent inversion in terms of the Bessel function $K_1$. Define
\[
 k(y)=2\sqrt y\, K_1(2\sqrt y),\qquad c=2\sqrt\delta=4\pi/\sqrt{165},
 \qquad G(u)=e^{u/2}\sum_{n\ge1}\lambda_n k(\delta n e^u).
\]
The Mellin transform of $k$ is $\Gamma(s)\Gamma(s+1)$ for $\Re s>0$. Absolute convergence on a sufficiently far right vertical line and Mellin inversion therefore identify $G$ with the inverse transform of \eqref{mix:complete} in its centered variable. Shifting that inversion line through the entire completion, and using the even functional equation, gives $G(-u)=G(u)$. The identity extends to the common strip of analyticity. The inverse transform consequently decays at both real ends and
\[
 e^{-2iw}\Lambda(1/2+w)=\int_{\mathbb R}G(u+2i)e^{wu}\,du.
\]
The shift is legitimate because the scaled Bessel arguments stay in the right half-plane on the shifted strips, and reflection gives the bound at the negative end.

Here $v=c\sqrt n e^{u/2}$ for $u\ge0$. With complex argument $z$ of modulus $v$, the mixed kernel is $zK_1(z)$, whereas the term estimated in Section~\ref{sp:rotated} is $(z^2/2)K_1(z)$. The absolute ratio $2/v$ is bounded by $2/c<3$. Since $0.7<c<1$, all majorants there based on $|\lambda_n|\le n^3$ apply to the present coefficient sum after multiplication by three, using the present value of $c$. The even reflection has the same absolute bounds at the other end. Set $T=64$, $U=15$, $V=180$, $b=1/56$, $d=4/5$ and $s_0=\cos(1)$. The respective bounds for quadrature, the omitted sample tail and the truncated coefficient tail are
\begin{align*}
 &3\cdot10^{29}e^{64d}/(e^{2\pi d/b}-1),\\
 &30b(4\cdot6^{10}9!)e^{2U-e^{U/2}/6}/(e^{2b}-1),\\
 &12b c^{-8}s_0^{-10}\Gamma(10,V)/(1-e^{-3b}).
\end{align*}
Their base-ten logarithms are below $-70.53,-102.72,-54.48$, respectively, with total below $10^{-51}$. For $u=jb$ retain $n\le\max(1,\lceil(V/(s_0ce^{u/2}))^2\rceil)$. The finite coefficients come from \eqref{mix:squares}, and the complex values of $K_1$ are enclosed by the convergent and asymptotic series with remainders from Section~\ref{sp:rotated}.

Let $C_j=bG(jb+2i)$ with this finite coefficient cutoff. Evenness and conjugation make the exact zero-frequency sample real. Take half its real part at $j=0$, checking its missing imaginary part against the coefficient-tail enclosure. The finite approximation is then $\sum_j(C_je^{jbw}+\overline{C_j}e^{-jbw})$. At each center $it_0$ expand the exponentials through degree $400$. On $|w-it_0|<5$ the sum of the absolute Taylor remainders is bounded by
\[
        2\sum_j |C_j| e^{5U}(5U)^{401}/401!.
\]
The verifier checks that this, added to the three analytic allowances, is below $10^{-49}$. Rounding and kernel-evaluation errors are carried in the finite interval coefficients using $640$ bits. This gives a separate polynomial approximation constructed from a Mellin kernel on a displaced integration line.

The resulting Bessel calculation independently counts $316$ simple critical zeros through height $64$. All eight windings and all $158$ coarse positive brackets coincide with the Cauchy results. At denominator $2^{24}$, $157$ Bessel brackets have exactly the Cauchy endpoints. The last Bessel bracket has numerators $(1070346298,1070346300)$ and contains the sharper Cauchy bracket $(1070346299,1070346300)$. This Bessel contour accepts $1354$ panels; the $7/4$ Cauchy tilt accepts $1362$, with five different individual subdivision counts. The Bessel central enclosure also overlaps the independently evaluated Cauchy value.

\subsection{The mixed tensor through height 256}\label{mix:high}
For the upper half of the calculation we retain the two-kernel formula of Section~\ref{mix:contours} and replace each row of parameters. The successive columns are the tilt, the two real abscissae, the last retained coefficient, reciprocal frequency spacing, cutoff, quadrature half-width, Taylor degree and working precision.
\begin{center}
\begin{tabular}{lrrrrrrrrr}
\toprule
Arithmetic & $\theta$ & $\sigma$ & $\sigma'$ & $N$ & $1/b$ & $Y$ & $d$ & $K$ & Bits\\\midrule
Point counts & $12/5$&31&136&720000&22&1900&8&330&1088\\
Eta products & $49/20$&33&140&850000&23&2200&8&340&1152\\\bottomrule
\end{tabular}
\end{center}
On the strip the gamma rotation is $29/20$ with reciprocal cosine below $9$ in both cases. On the omitted-coefficient line the rotation and reciprocal-cosine bounds are $(121/100,20/7)$ and $(247/200,40/13)$. In the raw Taylor estimate use $A_*=B_0^{2\sigma+1}A(\sigma)$, where the rotation and bound $(\phi_0,B_0)$ are $(121/100,3)$ and $(123/100,31/10)$. Each has $2\phi_0\ge\theta$, so the modulus incurs no increase from the frequency after rotating the two gamma factors. The four bounds in Section~\ref{mix:contours}, in the same order, now have base-ten logarithms below $(-288.41,-265.36,-112.37,-139.43)$ and $(-298.78,-271.61,-113.46,-146.50)$. Their sums are below $10^{-108}/2$. The geometric Taylor estimate uses $|z|<5$; the omitted semi-infinite Dirichlet tail is needed only on $|\Re w|\le1/2$, where Lemma~\ref{sp:taillemma} applies throughout each rectangle. Rounded finite weights use this new precision. The two tables of integer coefficients agree term by term through $720000$.

\begin{theorem}[A full rectangle for the mixed elliptic tensor]\label{mix:256}
The primitive standard $L$-function of $11.\mathrm a3\boxtimes15.\mathrm a7$ has exactly $1716$ nontrivial zeros with $|\Im s|\le256$. All are simple. They have real part $1/2$ in unitary coordinates, or $3/2$ in motivic coordinates. The $858$ positive ordinates have disjoint certified open brackets of width $2^{-24}$.
\end{theorem}
\begin{proof}
Each of the two new parameter sets first rechecks the winding and strict coarse brackets in the common band $[120,128]$, before extending its corresponding proof of Theorem~\ref{mix:128}. The newly computed eight-unit windings and accepted panel counts are as follows, in the notation of the earlier table.
\begin{center}
\begin{tabular}{rrrr|rrrr}
\toprule
Height &$W$&$P$&$Q$& Height&$W$&$P$&$Q$\\\midrule
136&30&282&282&200&30&298&300\\
144&28&284&284&208&31&298&298\\
152&29&290&290&216&32&298&298\\
160&30&292&292&224&31&300&300\\
168&29&292&292&232&31&294&294\\
176&30&294&294&240&32&306&310\\
184&30&308&308&248&31&310&310\\
192&30&306&306&256&32&308&310\\\bottomrule
\end{tabular}
\end{center}
There are $486$ further positive roots. In all $32$ bands the count of disjoint strict sign changes equals the winding with multiplicity. Thus every zero is simple and lies on the central line. The two computations agree on all $858$ dyadic numerator pairs and certify all rectangle boundaries zero-free. The highest pair is $(4291265578,4291265579)$, again on denominator $2^{24}$. The accepted panel totals are $8308$ and $8316$. Fourteen individual band counts differ. Conjugation supplies the other $858$ zeros.
\end{proof}

\paragraph{An independent complete Bessel count.}
Keep the displaced even Bessel inversion of Section~\ref{mix:bessel} and take $T=256$, $U=17$, $V=397$, $b=1/124$, $d=4/5$, $K=550$, with $1152$-bit arithmetic. In the three estimates of that section replace $e^{64d}$ by $e^{256d}$ and insert these new cutoffs. The resulting base-ten logarithms are below $-152.27,-325.86,-145.65$, with sum of the errors below $10^{-144}$. The local exponential Taylor error on $|z|<5$ is at most $2\sum_j|C_j| e^{85}85^{551}/551!$. Including the three errors, the interval test gives a total below $10^{-140}$.

The evaluation of each retained finite coefficient sum also changes. Put $a=ce^{u/2}$, $\omega=e^i$. The mixed sample is
\[
 b e^{u/2}\omega^2a\sum_{n\le N_u}\frac{A_n}{\sqrt n}K_1(a\omega\sqrt n).
\]
Round $N_u$ up to the right end of the same unit bin in $\sqrt n$. This only decreases the positive omitted-tail majorant. Evaluate $n<256$ directly. For the higher terms, use the signed moments, Taylor disks and explicit $K_0,K_1$ recurrences of Section~\ref{sp:groups} with the mixed integers~\eqref{mix:squares}. The disk real-part bound required there holds already at the smallest center since $c\cos(1)(33/4)>3$. The per-argument Taylor remainder is below $10^{-175}$, multiplied by the unsigned sum $\sum |A_n|/\sqrt n$ in each bin. The $736$ possible bins extend the initial coefficient cutoff to $565503$. All $2109$ nonnegative frequencies together require $167875$ bin evaluations. The symmetry check, real half-weight at zero, translated polynomial and final winding use the even parity of the mixed tensor.

This calculation certifies all $32$ band windings, all $858$ strict fine intervals with precisely the same endpoints as the two Cauchy computations, and overlapping central-value enclosures. Its total of $8312$ accepted panels differs from the point-count Cauchy total in $18$ individual bands. Thus the same full rectangle and all its isolated zeros have been certified by the two-kernel Cauchy formula and, independently, by the displaced even Bessel transform.

\subsection{The same critical zeros over cyclic towers}\label{mix:towers}
\begin{theorem}\label{mix:persist}
Let $\Q=K_0\subset K_1\subset\cdots\subset K_m=K$ be any finite tower in which each step is cyclic. The tensor of $E_K$ and $F_K$ is the standard function of a cuspidal representation of $\mathrm{GL}_4(\mathbb A_K)$. Its conductor ideal is $\prod_{\mathfrak p\mid165}\mathfrak p^2$, and its functional-equation sign is $+1$. Every nontrivial zero of the rational mixed tensor occurs at the same complex point over $K$, with at least its original multiplicity. In particular its degree-$4[K:\Q]$ completion contains at least the $1716$ distinct critical zeros of Theorem~\ref{mix:256} in the same height range.
\end{theorem}
\begin{proof}
A cyclic step can be refined into steps of prime degree. Cyclic automorphic base change successively supplies both degree-two forms over each $K_j$, with the local parameters obtained by restriction. See J.~Arthur and L.~Clozel, \href{https://press.princeton.edu/books/paperback/9780691085180/simple-algebras-base-change-and-the-advanced-theory-of-the-trace}{\emph{Simple Algebras, Base Change, and the Advanced Theory of the Trace Formula}}. The possible noncuspidal base change of a cuspidal $\mathrm{GL}_2$ form is a sum of two characters. Its local monodromy vanishes everywhere, whereas the respective multiplicative reduction above $11$ or $3$ preserves a nonzero monodromy operator. Each base change is therefore cuspidal. For the same reason it is nondihedral: local automorphic induction of a character has zero monodromy. At a place above $11$ the two ranks of monodromy still differ, ruling out a character twist between the two forms. Ramakrishnan's Theorem M now produces the cuspidal mixed tensor over each $K_j$.

At a prime above $3,5$ or $11$ one restricted curve remains multiplicative, possibly changing split type, and the other remains good. There are still two unramified special blocks with signs whose product is one. Their conductor exponents sum to two. All other finite places are good. At a real place the earlier factor gives sign one. At a complex place the angular exponents of the tensor are $2,0,0,-2$. A character $(z/|z|)^k$ has local unit root number $i^{|k|}$ in the standard convention, so the product is $i^4=1$. Changing to the trace additive character at a finite place has no effect on the normalized sign because the representation has trivial determinant. This proves the stated arithmetic conductor and sign; the absolute conductor in the completion over $\Q$ is $|D_K|^4\prod_{\mathfrak p\mid165}(N\mathfrak p)^2$.

For one cyclic step $L/M$ let $\Pi_M$ denote the preceding cuspidal mixed tensor. Local induction and restriction with its finite permutation representation give, in standard normalization of the completed functions,
\[
       \Lambda(E_L\times F_L,s)=
       \prod_{\chi\in\widehat{\operatorname{Gal}(L/M)}}
                \Lambda(\Pi_M\otimes\chi,s).
\]
This is Artin formalism, using the sum of the characters in the regular representation and the compatibility of local tensor, induction and restriction factors, including the discriminant contribution. Every twisted factor is cuspidal and entire. The trivial character contributes exactly the preceding completion. Induction along the tower proves the divisor inclusion. Distinct complex points remain distinct, giving the numerical lower bound with their original ordinates. Additional factors can contribute further zeros and higher multiplicities.
\end{proof}

\begin{corollary}[Nonnormal fields]\label{mix:dihedral}
Suppose $L/\Q$ is Galois with group $A\rtimes C_2$, where $A$ is abelian and $C_2$ acts by inversion. For every intermediate $K$, the mixed tensor completion over $K$ is entire and divisible by its rational completion with entire quotient. This includes every non-Galois cubic field. Together with Theorem~\ref{mix:persist}, every cubic number field therefore has all the same certified rational critical zeros.
\end{corollary}
\begin{proof}
Put $M=L^A$, a quadratic field, and $H=\operatorname{Gal}(L/K)$. The permutation representation $P=\operatorname{Ind}_H^{\operatorname{Gal}(L/\Q)}1$ contains the trivial representation once. Every irreducible summand of $P$ has dimension one or is induced from a character of $A$. Indeed, restriction to $A$ decomposes into eigenspaces, and the other coset interchanges the spaces for $\nu$ and $\nu^{-1}$. A non-fixed pair gives the two-dimensional induced irreducible; for an inversion-fixed character the involution diagonalizes on its eigenspace and gives one-dimensional irreducibles.

Tensor $P$ with the rational four-dimensional parameter and use local Artin formalism. A character summand gives an entire finite-character twist of $\Pi_\Q$. A two-dimensional summand $\operatorname{Ind}_{M/\Q}\nu$ gives the completed standard function of the twist $\Pi_M\otimes\nu$ over the quadratic field, entire by Theorem~\ref{mix:persist}. These are positive integer multiplicities in a direct sum, so after removing the trivial summand the remaining product is entire. This proves the claim for the geometric completion over $K$ without requiring automorphic descent of the tensor to a nonnormal $K$. A non-Galois cubic field has normal closure with group $S_3=C_3\rtimes C_2$, as required.
\end{proof}
The number-field degrees covered by the corollary are unbounded. For instance, take the ring class field of $\Z+\ell\Z[i]$ with $\ell\equiv3\pmod4$ prime. It is generalized dihedral over $\Q$; its subfield fixed by complex conjugation has degree $(\ell+1)/2$ by the class-number formula used in Section~\ref{cm:proportion}. The two elliptic curves in the tensor remain non-CM over these fields.

\begin{corollary}[All quartic fields]\label{mix:quartic}
For every quartic number field $K$, the quotient of the geometric mixed completion over $K$ by its rational completion continues holomorphically to $\Re s>0$. Thus every rational critical-strip zero is present over $K$, with at least its rational multiplicity.
\end{corollary}
\begin{proof}
The transitive quartic groups are $C_4,C_2^2,D_4,A_4,S_4$, with $D_4$ here of order eight. The first three were treated above. In the other cases the permutation representation is $1\oplus\rho$, where the augmentation representation $\rho$ is irreducible of dimension three. The strong Artin theorem for $A_4$ and $S_4$ supplies a cuspidal $\tau$ on $\mathrm{GL}_3/\Q$ with the Frobenius parameters of $\rho$ at almost all primes. See P.~J.~Cho, \href{https://doi.org/10.1093/qmath/has046}{\emph{The strong Artin conjecture and large class numbers}}, Quarterly Journal of Mathematics 65 (2014), 101--111. The completed Rankin--Selberg function $\Lambda(\Pi\times\tau,s)$ is entire by the standard Rankin--Selberg theorem of Jacquet, Piatetski-Shapiro and Shalika for cuspidal representations of different degrees.

Artin formalism identifies the desired quotient with the completion attached to the tensor of the rational mixed parameter and $\rho$. It agrees with $\Lambda(\Pi\times\tau,s)$ outside a finite set of places. Each missing automorphic finite factor is removed by its inverse Euler polynomial, an entire function of $s$. Each desired geometric local factor can be inserted holomorphically on $\Re s>0$: its reciprocal roots have modulus at most one. Indeed, the Artin parameter has finite image and zero monodromy, so the kernel of tensor monodromy is the tensor of its space with the original kernel. On this kernel the mixed local reciprocal roots have modulus at most one by the good and multiplicative calculations above. Passing to inertia invariants preserves the bound. At infinity all geometric gamma shifts have nonnegative real part, and the inverse automorphic gamma functions are entire. The conductor quotient is an everywhere nonzero exponential. These finitely many replacements prove holomorphy on the stated half-plane using only the almost-everywhere Artin matching, which suffices for the claimed zero inclusion.
\end{proof}
For an explicit example $x^4-x-1$ is irreducible modulo two, since it has neither a root nor the unique irreducible quadratic factor. Its cubic resolvent $y^3+4y-1$ has no rational root, and its discriminant is $-283$. The Galois group is therefore $S_4$. Its quartic root field is one of the nonnormal fields in Corollary~\ref{mix:quartic}.

\begin{proposition}[A finite-image transfer for the mixed tensor]\label{mix:artin}
Let $L/\Q$ be finite Galois and $K$ an intermediate field. Suppose each irreducible constituent of the permutation representation on the embeddings of $K$ is associated at almost all places with a cuspidal automorphic representation over $\Q$. Then the geometric mixed completion over $K$, divided by its rational completion, is holomorphic on $\Re s>0$.
\end{proposition}
\begin{proof}
The trivial constituent has multiplicity one. For any other constituent $\rho$ let $\tau$ denote its unitary automorphic representation. The correction of the finitely many exceptional local factors for $\Pi\times\tau$ is holomorphic on $\Re s>0$ by the proof of Corollary~\ref{mix:quartic}. Its completed Rankin--Selberg function is entire unless $\tau$ is a unitary norm twist of $\Pi^\vee$. This exception requires degree four. The central character of $\tau$ has finite order: it agrees almost everywhere with $\det\rho$, so this follows from multiplicity one for Hecke characters. Since $\Pi$ has trivial determinant, a norm twist by $|\cdot|^{it}$ could therefore occur only at $t=0$. If $\tau=\Pi^\vee$, the normalized four-dimensional Tate representation, dualized, would have the same good Frobenius polynomials as $\rho$. Its normalization uses the integral Tate twist dividing the motivic reciprocal roots by $p$. Chebotarev density and the semisimple character criterion would identify this representation with a finite-image representation. Its restriction to the finite kernel field of $\rho$ remains irreducible of dimension four by Proposition~\ref{mix:strong}, also after the Tate twist. This contradiction excludes the exception. Taking the product over the nontrivial constituents with their multiplicities proves the result.
\end{proof}

\begin{corollary}[Nonsolvable base fields of unbounded degree]\label{mix:nonsolvable}
Let $L$ be the splitting field of $x^5-x^3-x^2+x+1$. For every prime $p\ne1609$ and every intermediate field of $L\Q(\zeta_p)/\Q$, every critical-strip zero of the rational mixed tensor occurs over that field, with at least its original multiplicity. In particular this holds for the full nonsolvable extensions of degrees $120(p-1)$.
\end{corollary}
\begin{proof}
Calegari, \href{https://doi.org/10.1007/s00208-012-0839-4}{\emph{The Artin conjecture for some $S_5$-extensions}}, Theorem~5.1, proves automorphy for all irreducible complex representations of $\operatorname{Gal}(L/\Q)\simeq S_5$. The discriminant of the quintic field is the prime $1609$. The only nontrivial abelian subextension of $L$ is $\Q(\sqrt{1609})$. It ramifies at $1609$, whereas $\Q(\zeta_p)$ with $p\ne1609$ does not. The compositum has Galois group $S_5\times(\Z/p\Z)^\times$. Each irreducible character of this product is the tensor of an $S_5$ irreducible and a finite abelian character. Its automorphic representation is the corresponding finite-character twist, still cuspidal. Apply Proposition~\ref{mix:artin} to the permutation representation for any intermediate field.
\end{proof}

There are also infinitely many pairwise distinct quintic fields with nonsolvable normal closures to which this zero inclusion applies. For a monic quintic impose the following squarefree reductions:
\begin{align*}
 f\bmod 2&=x^5+x^2+1,\\
 f\bmod 3&=x(x-1)(x+1)(x^2+1),\\
 f\bmod 5&=x(x^2+2)(x^2+3).
\end{align*}
The first polynomial is irreducible: it has no linear root and is divisible by neither $x^2+x+1$ nor either irreducible cubic over $\mathbb F_2$, as direct remainders show. The last two factorizations have degrees $(1,1,1,2)$ and $(1,2,2)$. The Galois group contains a five-cycle and a transposition, so is $S_5$. Indeed the five-cycle conjugates the transposition through the edges of a connected graph on five vertices, and its edge transpositions generate the full symmetric group. Require also that the polynomial have exactly one real root. Complex conjugation and the Frobenius at $5$ then both have cycle type $(2,2,1)$. Calegari's Theorem~1.1 applies to the irreducible four-dimensional augmentation representation. The quintic permutation representation is $1\oplus\rho_4$, so Proposition~\ref{mix:artin} gives the asserted inclusion for its root field.

For example $f=x^5+15x^2+56x+15$ has these reductions. Its derivative has minimum $56-(45/2)(3/2)^{1/3}>0$. To obtain distinct fields recursively, choose a prime $q>5$ unramified in all preceding fields. Choose the coefficient $a$ in $f=x^5+ax^2+bx+c$ so it is divisible by $q$ and meets the three reductions. Choose $b$ in its required classes modulo $30q$, divisible by $q$ and large enough that $5x^4+2ax+b>0$ on $\mathbb R$. Finally impose $c\equiv q\pmod{q^2}$ along with the three original reductions. The Chinese remainder theorem gives all these choices. This polynomial is Eisenstein at $q$, hence its quintic field is ramified at $q$ and is different from every preceding one. Its positive derivative preserves the required signature. The argument calls only on the four-dimensional part of Calegari's Theorem~1.1; the stronger all-representation input for the single discriminant-$1609$ field comes from Theorem~5.1.

\subsection{Transfer of elliptic zeros over number fields}
\begin{theorem}\label{mix:elliptic-transfer}
Let $C/F$ be an elliptic curve over a number field. Assume that $C$ is CM or has a cuspidal automorphic representation on $\mathrm{GL}_2/F$ agreeing with its unitary parameters at almost all places. Let $K/F$ be finite, and suppose each nontrivial irreducible constituent of the permutation representation on the embeddings of $K$ agrees almost everywhere with a cuspidal automorphic representation over $F$. Then in unitary coordinates the quotient
\[
            \Lambda(C_K,s)/\Lambda(C_F,s)
\]
has a holomorphic continuation to $\Re s>0$. Every critical-strip zero over $F$ occurs over $K$ at the same complex point, with at least its original order.

For $F=\Q$ this applies to every elliptic curve, every field of degree at most four, every generalized dihedral intermediate field of Corollary~\ref{mix:dihedral}, every subfield in Corollary~\ref{mix:nonsolvable}, and the infinite family of nonsolvable quintic root fields constructed above. Over arbitrary $F$, every zero of every CM elliptic curve persists in every extension of degree at most three, in every generalized dihedral intermediate field, and in every finite tower of these extensions.
\end{theorem}
\begin{proof}
Suppose first that $C$ has a unitary cuspidal parameter $\pi_C$ over $F$. Let the permutation representation be $1\oplus\bigoplus_j m_j\rho_j$. The automorphic representation $\tau_j$ associated to $\rho_j$ has finite central character $\det\rho_j$ by multiplicity one for Hecke characters. Use its unitary normalization. Artin formalism writes the quotient on a far right half-plane as the product of the completed factors of the geometric tensors of $C$ with $\rho_j$, with powers $m_j$. At almost every place these agree with the factors of $\Lambda(\pi_C\times\tau_j,s)$.

We show that each Rankin--Selberg completion is entire. Its only possible exception would be $\tau_j=\pi_C^\vee\otimes |\cdot|^{it}$ in degree two. The central character of $\pi_C$ is trivial by the determinant of the normalized elliptic Frobenius parameters. Comparison of central characters gives $t=0$. The normalized traces at good primes for $\pi_C^\vee$ equal those of $\pi_C$, since the two normalized eigenvalues have product one. Were these traces the traces of a finite group representation, they would belong to a single number field $D$. Take rational primes $p$ splitting completely in a normal closure over $\Q$ of $F(C[3])$, avoiding the finite set of bad or exceptional primes. There are infinitely many by Chebotarev. At every $v\mid p$ the residue field has size $p$, and the integral trace $a_v$ is $2$ modulo $3$, since Frobenius acts trivially on $C[3]$. Thus $a_v\ne0$. Having $a_v/\sqrt p\in D$ would force $\sqrt p\in D$ for infinitely many distinct $p$. A fixed number field contains only finitely many quadratic subfields, by passing to its finite Galois closure. This contradiction excludes the exception.

To replace a finite exceptional factor its automorphic inverse is an entire Euler polynomial in $(Nv)^{-s}$. All normalized geometric reciprocal roots on the kernel of monodromy of $C$ have modulus at most one \emph{before} taking inertia invariants. After a finite Galois local extension $C$ has good or split multiplicative reduction. For good reduction the claim there is the Weil bound with unitary normalization. For split multiplicative reduction the one-dimensional kernel has normalized Frobenius eigenvalue of modulus $q^{-1/2}$ for residue field size $q$. A positive power of an original Frobenius belongs to the subgroup of the semistable extension. On that subgroup inertia acts trivially on the monodromy kernel, and the power acts with the stated semistable moduli raised to the corresponding relative power. The eigenvalues of the original Frobenius on the same kernel therefore have modulus at most one. The Artin parameter has finite image and zero monodromy, so tensoring preserves this bound on the whole kernel. Taking inertia invariants afterwards preserves it as well. The inserted factors are consequently holomorphic on $\Re s>0$. At a real place a finite Artin parameter consists of signs and trivial characters, and either twist of the elliptic parameter $I_1$ has factor $\Gamma_{\mathbb C}(s+1/2)$. At a complex place a finite Artin parameter is trivial, while the elliptic angular exponents $\pm1$ give factors with shift $1/2$. All inserted geometric gamma factors are holomorphic on this half-plane. Dividing out any automorphic gamma factors uses their entire reciprocals. Conductors contribute nonvanishing exponentials. The corrected product proves the assertion for cuspidal $\pi_C$.

It remains to cover CM over every $F$. If the endomorphisms become defined only over a quadratic extension, the two conjugate CM characters are distinct by their opposite angular types. Their quadratic automorphic induction is cuspidal and the preceding proof applies. When all endomorphisms are defined over $F$, the CM factorization from Deuring gives two unitary Hecke characters $\psi$ and $\bar\psi$. Tensoring either one with each finite $\rho_j$ gives at almost all places the standard function of $\tau_j\otimes\psi$ or $\tau_j\otimes\bar\psi$. These standard completions are entire for $\dim\rho_j>1$. In dimension one the character associated to $\rho_j$ has finite order, whereas each CM factor has nonzero angular exponent at a complex place. Its twist cannot be a pure norm character, so its Hecke completion is entire as well. The same correction of the finite exceptional factors and gamma functions applies to the product. This proves the assertion for all CM curves.

For rational $C$ modularity is known. The required permutation constituents in the rational examples were identified in Corollaries~\ref{mix:dihedral}--\ref{mix:nonsolvable} and the subsequent quintic construction. Over any base field a generalized dihedral group has one-dimensional constituents and quadratic inductions of Hecke characters, as in the proof of Corollary~\ref{mix:dihedral}; when two-dimensional, these inductions are cuspidal. Every non-Galois cubic has normal group $S_3$, while cyclic extensions use characters. Applying the proved assertion successively is legitimate for a CM curve because it retains CM over every finite extension. Thus orders can only increase along any of the stated towers.
\end{proof}

\begin{corollary}[Entire CM quotients]\label{mix:cm-all}
For a CM elliptic curve over any number field $F$, the completed quotient in Theorem~\ref{mix:elliptic-transfer} is entire on the whole plane. The same conclusion holds for every Galois extension $K/F$, of arbitrary group, and every field $K$ contained in a solvable Galois extension of $F$. These conditions may be applied successively in a tower. In particular every certified rational CM zero in this manuscript persists with at least its original multiplicity over these fields and over the nonnormal nonsolvable quintic fields constructed above.
\end{corollary}
\begin{proof}
For the Galois and solvable-closure cases the stronger field coverage is already Theorem~1 of M.~R.~Murty and V.~K.~Murty, \href{https://doi.org/10.1090/conm/143/01014}{\emph{Base change and the Birch--Swinnerton-Dyer conjecture}}, Contemporary Mathematics 143 (1993), 481--494. Their theorem gives an entire quotient of the finite parts. Their proof uses positive character induction and Deuring's theorem. The geometric gamma factors in unitary normalization all have positive half-integer shift and are holomorphic on $\Re s>0$. The finite quotient thus gives a holomorphic completed quotient there. Under the automorphic-constituent hypothesis of Theorem~\ref{mix:elliptic-transfer}, holomorphy on the same half-plane was established directly, including for the nonsolvable quintics.

In either case both elliptic completions are entire and satisfy their self-dual functional equations by CM theory. Their meromorphic quotient $Q$ therefore satisfies $Q(s)=(\epsilon_K/\epsilon_F)Q(1-s)$ in conductor-normalized coordinates. For every point with $\Re s\le0$ the point $1-s$ belongs to $\Re s>0$, so $Q$ is holomorphic there as well. Multiplication of the resulting entire quotients proves the assertion for towers. The orders of every zero can only increase at each step.
\end{proof}

\clearpage
\section{Cohomology of powers of elliptic curves}\label{coh:section}
The tensor and symmetric-power results give zero theorems for the individual cohomology groups of abelian varieties of unbounded dimension. We first treat a pair of arbitrary elliptic curves over $\Q$, and then obtain exact finite zero divisors for powers of a particular curve. Use the unitary coordinate for a motive of weight $j$, so that the arithmetic variable is $s+j/2$. Its proposed critical line is $\Re s_{\rm arith}=(j+1)/2$.

\begin{theorem}\label{coh:tensors}
Let $E,F/\Q$ be elliptic curves and $A=E^g\times F^h$, with $g,h\ge0$, $g+h>0$. For each $0\le j\le2(g+h)$ the unitary completed $H^j$-function of $A$ is a product, with nonnegative integral multiplicities, of completed Rankin--Selberg functions of cuspidal constituents of symmetric powers of $E$ and $F$. Exponent zero denotes the trivial character. The conclusion holds also for the tensor of any number of copies of $H^1(E)$ with any number of copies of $H^1(F)$. After removing poles, the uniform partial Li criterion and the two-sided unconditional tail estimates of Theorems~\ref{li:everyheight} and~\ref{li:positive-tail} apply to each of these geometric functions.
\end{theorem}
\begin{proof}
Normalize the two two-dimensional parameters $V,W$ to have weight zero. Their determinants are then trivial. If $a\ge0$, repeated application of the two-dimensional Clebsch--Gordan formula gives
\[
 V^{\otimes a}\simeq\bigoplus_{0\le i\le a/2}
 \bigl(\operatorname{Sym}^{a-2i}V\bigr)^{\oplus c_{a,i}},
 \qquad c_{a,i}=\binom ai-\binom a{i-1}.
\]
Indeed, tensoring a positive symmetric power by $V$ gives the two adjacent powers, and tensoring $\operatorname{Sym}^0V$ by $V$ gives just $V$. Pascal's identity gives the stated multiplicities, starting with $c_{0,0}=1$. We recover the tensor of $a$ copies of $V$ and $b$ copies of $W$ by tensoring the two displayed decompositions. Each summand is $\operatorname{Sym}^{a-2i}V\otimes\operatorname{Sym}^{b-2k}W$. This identity is a representation identity before taking inertia invariants or the kernel of monodromy, so it identifies the full local factors at bad primes as well. Newton and Thorne give cuspidal symmetric powers of each non-CM curve at every positive exponent, with the matching local parameters. Local Rankin--Selberg theory identifies the tensor factor. See \href{https://doi.org/10.1007/s10240-021-00126-4}{Newton--Thorne, \emph{Symmetric power functoriality for holomorphic modular forms, II}, Theorem A} and the local correspondence cited in Section~\ref{sf:chapter}.

For a CM curve use the index-two induction of its unitary Hecke character over the CM field $K$, as in Section~\ref{cm:hecke}. Upon restriction to $K$ the two basis lines carry $\chi$ and $\chi^c$, and the outer coset interchanges these lines. In the $k$th symmetric power the monomials of bidegrees $(k-i,i)$ and $(i,k-i)$ for $i<k/2$ therefore form the induction of $\chi^{k-i}(\chi^c)^i$. Their two characters have distinct infinity types because $k-2i>0$. The induction is consequently cuspidal on $\mathrm{GL}_2$ by quadratic automorphic induction. For $k=2i$ the middle monomial $(xy)^i$ gives the character $\eta_{K/\Q}^i$. Its restriction to $K$ is trivial, and an element outside $K$ interchanges the basis lines by a determinant-one matrix, so it sends $xy$ to $-xy$. This proves the requisite positive isobaric decomposition in the CM case, also when the two curves have different CM fields.

By the K\"unneth formula a summand of $H^j(A)$ chooses from each curve an $H^0$, an $H^1$ or an $H^2$, with the subscripts adding to $j$. The $H^0$ is trivial and $H^2$ is the determinant. The Tate twists disappear upon making the weight-$j$ unitary normalization. This gives exactly positive sums of the tensors already considered. At infinity the same decomposition identifies the gamma factors. Conductors multiply across the resulting direct sum. One-line nonvanishing and reflection allow removal of boundary poles factor by factor, including when $E$ and $F$ are isogenous. The Li sums and their tails are additive, and $D_v$ is additive as in Corollary~\ref{li:isobaric}. This proves the assertions.
\end{proof}

\begin{corollary}\label{coh:counts}
Put $E=43.\mathrm{a1}$. For every $g\ge1$, each cohomology group listed below has all nontrivial zeros on its arithmetic critical line throughout the stated band. Omit any group of degree greater than $2g$. Write $b_k=\binom gk$, setting it to zero for $k>g$, and $c=g\binom{g-1}{2}$. The counts include both signs and the center with multiplicity.
\[
\begin{array}{c|r|l}
 \text{Group} & |\Im s|\text{ at most} &\text{Number of zeros} \\ \hline
 H^2(E^g)&256&1276b_2+222(g+b_2)\\
 H^3(E^g)&256&1803b_3+749[2b_3+g(g-1)]\\
 H^4(E^g)&32&188b_4+96(3b_4+c)+8(2b_4+c+b_2).
\end{array}
\]
Poincar\'e duality gives the same counts for the complementary degrees with their own arithmetic centers. The first $884\,736$ unitary Li coefficients in the first two rows and the first $7577$ in the third row are positive.
\end{corollary}
\begin{proof}
Let $L_i$ mean the unitary completed symmetric $i$th power and $L_0$ the completed zeta function. The possible choices of positive cohomological degrees in the three K\"unneth decompositions are respectively $(2)$ or $(1,1)$, then $(2,1)$ or $(1,1,1)$, then $(2,2)$ or $(2,1,1)$ or $(1,1,1,1)$. They give the products
\[
\begin{aligned}
 H^2 &: L_2^{b_2}L_0^{g+b_2},\\
 H^3 &: L_3^{b_3}L_1^{2b_3+g(g-1)},\\
 H^4 &: L_4^{b_4}L_2^{3b_4+c}L_0^{2b_4+c+b_2}.
\end{aligned}
\]
Here the notation denotes the completed function of each cohomology group. At height $256$ the input counts for $L_0,L_1,L_2,L_3$ are $222,749,1276,1803$ from the conductor census and the audited higher-rank counts in Section~\ref{sf:chapter}. At height $32$ the counts for $L_0,L_2,L_4$ are $8,96,188$. Every factor has all its zeros on the line in its indicated band. Multiplicities add even at common zeros, proving the table. The normalized factors are self-dual, so the K\"unneth duality twist in the complementary degree gives the same unitary zero divisor. For the Li coefficients, Theorem~\ref{li:small-universal} gives the respective index bounds $13.5\cdot256^2=884\,736$ and $\lfloor7.4\cdot32^2\rfloor=7577$ on every constituent. The real Li coefficients add with the displayed nonnegative multiplicities.
\end{proof}

\begin{corollary}[Totally real base fields]\label{coh:totally-real}
Let $k$ be a totally real number field and let $E_1,E_2/k$ be modular elliptic curves. For $g,h\ge0$ with $g+h>0$, every completed cohomological function of $E_1^g\times E_2^h$, and every tensor function built from copies of $H^1(E_1)$ and $H^1(E_2)$, has the positive factorization of Theorem~\ref{coh:tensors} over $k$. Its pole-removed zero divisor satisfies both the uniform finite-height Li criterion and the two-sided tail bounds there. For each fixed totally real $k$ there is a finite set $J_k\subset k$ such that the conclusion applies whenever both $j(E_1),j(E_2)$ lie outside $J_k$. The conclusion holds for \emph{every} pair of elliptic curves when $[k:\Q]=2$ or $3$, when $[k:\Q]=4$ and $\sqrt5\notin k$, or when $k/\Q$ is abelian and unramified at $3,5,7$.
\end{corollary}
\begin{proof}
The parameter of a modular elliptic curve is attached to a parallel-weight-two Hilbert eigenform, hence is essentially square integrable at every real place. At a finite place $v$, choose an auxiliary residue characteristic $\ell$ different from that of $v$. Modularity identifies the global semisimple Galois representations by the good Frobenius polynomials. Local--global compatibility for the Hilbert form identifies the Frobenius-semisimplified Weil--Deligne parameter at $v$, with its monodromy, as in \href{https://websites.math.leidenuniv.nl/edixhoven/ics/AG_Fall_2005/notes/edix3.pdf}{Edixhoven, Theorem~6.1}. In the non-CM case \href{https://annals.math.princeton.edu/2026/203-1/p04}{Newton--Thorne, \emph{Symmetric power functoriality for Hilbert modular forms}, Theorem~A} supplies the cuspidal symmetric power with the specified Weil--Deligne parameter at every place. In the CM case the parameter is induced from the quadratic CM extension of $k$. The two conjugate Hecke characters have distinct weights at each real place of $k$, so all the noncentral paired monomials induce cuspidally, and the even middle monomial gives the indicated power of the quadratic character. At a ramified place these decompositions are made before taking inertia invariants and the kernel of monodromy. Tensoring, taking K\"unneth summands, and adding the resulting zero divisors now give the first assertion. The absolute conductors multiply and the archimedean factors include all real places of $k$, exactly as required in the two-sided estimate. Finiteness of the possibly nonmodular $j$-invariants is \href{https://arxiv.org/pdf/1310.7088}{Freitas--Le Hung--Siksek, Theorem~5}, applied to the fixed base field. The stated field classes follow respectively from the modularity theorems of \href{https://arxiv.org/abs/1310.7088}{Freitas--Le Hung--Siksek}, \href{https://arxiv.org/abs/1901.03436}{Derickx--Najman--Siksek}, \href{https://arxiv.org/abs/2103.13975}{Box}, and \href{https://doi.org/10.5802/jtnb.1047}{Yoshikawa, Theorem~1.2}.
\end{proof}

\begin{corollary}[Adding arbitrarily many CM factors]\label{coh:mixed-cm}
Let $k$ be totally real. Consider a product of powers of elliptic curves over $k$. Suppose its non-CM curves, up to $k$-isogeny and quadratic twist, belong to at most two classes, each represented by a modular elliptic curve over $k$. There is no limit on the number of CM classes. Each individual weight-normalized cohomological function of the product satisfies Theorems~\ref{li:everyheight} and~\ref{li:positive-tail}, including the same absolute constants. Over $k=\Q$ the modularity hypothesis is automatic.
\end{corollary}
\begin{proof}
In a K\"unneth summand the two nondihedral kinds give, after Clebsch--Gordan decomposition, a tensor $V_a\otimes W_b$ of two unitary symmetric-power parameters, allowing degree-zero powers. A quadratic twist adds a unitary character. Each remaining $H^1$ from a CM curve is a sum of characters or a quadratic induction of one character, and the normalized $H^0,H^2$ terms are trivial. The tensor of any finite list of these CM terms is a positive sum of parameters $\operatorname{Ind}_{K/k}\chi$, with $K/k$ multiquadratic or trivial and $\chi$ a unitary Hecke character on $K$. Indeed the bases for the quadratic inductions are indexed by the cosets of normal index-two subgroups. The diagonal group permutes the tuples of basis lines; on the stabilizer of any tuple the action is the product of characters. The stabilizer is the intersection of the index-two groups for that summand. This is the Mackey decomposition in Section~\ref{recentcm:tensors}. A twist from $k$ is absorbed by restricting it to $K$.

For each orbit the local projection formula identifies its contribution with
\[
 L_K\bigl(s,(\mathrm{BC}_{K/k}V_a\otimes\chi)
                   \times\mathrm{BC}_{K/k}W_b\bigr).
\]
Here base change denotes automorphic base change of the corresponding symmetric-power forms. It exists by successive cyclic base change along the quadratic tower. It is a unitary isobaric sum, allowing a single cuspidal summand. Artin formalism and local compatibility at all places identify the completed functions, including the change of absolute conductor. These are the cyclic base-change inputs cited in Section~\ref{ebc:chapter}. We may apply Corollary~\ref{li:isobaric} over $K$. Since all multiplicities are positive, summing the real Li coefficients and the tail bounds proves the result over $k$. In low cohomological degree there is a separate way to handle additional non-CM types.
\end{proof}

\begin{corollary}[Low cohomological degrees with unrestricted isogeny classes]\label{coh:four}
Let $k$ be a number field and $A=E_1\times\cdots\times E_g$ a product of elliptic curves modular over $k$, in the sense of cuspidal or unitary isobaric parameters compatible at all places. There is no restriction on how many distinct isogeny or twist classes occur. If $\min(j,2g-j)\le4$, the completed unitary $H^j$-function satisfies Theorems~\ref{li:everyheight} and~\ref{li:positive-tail}. In particular every cohomological degree is included when $\dim A\le4$. More generally any tensor containing at most four of the non-CM $H^1$ parameters can be multiplied by arbitrarily many CM $H^1$ parameters and retains these conclusions.
\end{corollary}
\begin{proof}
First suppose $j\le4$. In a K\"unneth summand there are at most $j$ factors of degree one. The normalized determinants from degree two are trivial. Put the degree-one factors into two groups of at most two. For a single member use its modular parameter, for an empty group the trivial character, and for two members the degree-four tensor transfer of \href{https://arxiv.org/abs/math/0007203}{Ramakrishnan, Theorem M}. Distribute first over any isobaric linear factors. The tensor transfer is isobaric and has the indicated tensor parameters locally. The all-place formulation is given in \href{https://www.mathunion.org/fileadmin/ICM/Proceedings/ICM2002.2/ICM2002.2.ocr.pdf}{Cogdell--Piatetski-Shapiro, \emph{Converse theorems, functoriality, and applications to number theory}, pp.~125--126}. For the three- and fourfold local $L$- and $\varepsilon$-identities at simultaneous supercuspidal places see also Ramakrishnan, Proposition~4.3.1. Its constituents are unitary in this normalization: at unramified good places all the constituent roots are among products of eigenvalues of modulus one, so their determinant characters have trivial absolute-value exponent. Consequently the tensor of the two groups gives a Rankin--Selberg function of unitary isobaric sums. The projection to invariants at a ramified place takes place after the decomposition of the full parameter. This proves the assertion by Corollary~\ref{li:isobaric}. Poincar\'e duality and the weight normalization give the complementary degrees.

For the last assertion decompose the CM tensor into quadratic-tower inductions as above. On the multiquadratic extension use cyclic base change of the two groups just constructed and twist one group by the inducing character. Their Rankin--Selberg function gives exactly the contribution from that orbit, and the positive sum over orbits concludes the proof.
\end{proof}

\clearpage
\section{Critical zeros in real cyclotomic fields of increasing degree}\label{cyt:fields}
A recently obtained family mean value theorem gives information at heights which increase together with the base field. The finite characters occurring in the mean value theorem are primitive. The base change to an actual number field also contains the characters induced from smaller conductors. The following count includes their full contribution in the denominator.

Let $E/\Q$ be any elliptic curve. In rank $d=2$ let $L(s,\Pi)$ be its unitarily normalized Hasse--Weil function, and in rank $d=3$ assume $E$ has no CM and let $\Pi=\operatorname{Sym}^2E$ in the same normalization. Write $N$ for the arithmetic conductor of $\Pi$. Both are cuspidal with trivial central character, and the rank-three representation is self-dual. For any $q>2$ put $M_q=\Q(\zeta_q)^+$, the maximal totally real cyclotomic subfield. Write $Z_q(T)$ for the number of all nontrivial zeros with $|\Im s|\le T$ of the base-changed standard function $L(s,\Pi/M_q)$, and $Z_q^0(T)$ for the count on the critical line. The counts include multiplicities. In rank two this is the $L$-function of $E$ over the number field $M_q$ itself.

\begin{theorem}\label{cyt:main}
Fix $0<\omega<1$ for $d=2$, or $0<\omega<1/3$ for $d=3$. Let
\[
 c_2(\omega)=\frac{1-\omega}{3(1+\omega)},\qquad
 c_3(\omega)=\frac{1-3\omega}{9(1-\omega)}.
\]
Define $e_2(N)=1$ if $2\mid N$ and $e_2(N)=3/4$ otherwise, and set
\[
 D_N=e_2(N)\prod_{\substack{p>2\\p\nmid N}}
 \left(1+\frac1{p(p-2)}-\frac1{p^2(p-1)(p-2)}\right)^{-1}.
\]
For every sufficiently large $Q$, there exists $q\in(Q,2Q)$ coprime to $N$ for which
\[
                   \frac{Z_q^0(Q^\omega)}{Z_q(Q^\omega)}
                           >D_N c_d(\omega).
\]
Consequently, at the fixed exponent $\omega=1/100000$, one obtains the following strict lower proportions over fields of unbounded degree and at heights tending to infinity:
\begin{center}
\begin{tabular}{lr}
\toprule
Base-changed function & Proportion on its critical line\\
\midrule
$L(E/M_q,s)$ for any rational $E$ & $17.28\%$\\
$L(E/M_q,s)$ if $E/\Q$ has even conductor & $23.05\%$\\
$L(\operatorname{Sym}^2E/M_q,s)$ for rational non-CM $E$ & $5.76\%$\\
\bottomrule
\end{tabular}
\end{center}
The fields selected for different rows or different curves can differ.
\end{theorem}
\begin{proof}
For a primitive Dirichlet character $\chi$ of conductor $m$ prime to $N$, let $N_\chi(T)$ and $N^0_\chi(T)$ be the corresponding full and critical zero counts of $L(s,\Pi\otimes\chi)$. Write $\phi^*(q)$ for the number of primitive characters modulo $q$, counting both parities. Use only eligible moduli $(q,N)=1$ and $\phi^*(q)>0$. Fix $T=Q^\omega$ and put $B=(d/\pi)(1+\omega)T\log Q$.

The previous analytic input is Corollary~1.13 and (1.41) of B.~Conrey, C.-H.~Kwan, Y.~Lin and C.~L.~Turnage-Butterbaugh, \href{https://arxiv.org/abs/2607.00282}{\emph{Critical zeros and unconditional mean value theorems for twisted $\mathrm{PGL}(2)$ and $\mathrm{PGL}(3)$ $L$-functions}}, July 2026 preprint. For some $k_d(\omega)>c_d(\omega)$ and all sufficiently large $Q$, their theorem gives
\begin{equation}\label{cyt:input}
 \sum_q \frac{2}{\phi^*(q)}\sum_{\chi}^{\flat}N^0_\chi(T)
 > k_d(\omega)\sum_q\frac{2}{\phi^*(q)}\sum_{\chi}^{\flat}N_\chi(T).
\end{equation}
Here $Q<q<2Q$ runs over the eligible moduli and the inner sums run over even primitive characters of conductor exactly $q$. The stated rank-three result is unconditional for self-dual $\Pi$, which holds for the symmetric square with trivial normalized determinant. The result for $d=2$ applies at arbitrary conductor.

The degree-$d$ zero-counting formula, uniformly in these finite twists, gives
\[
 N_\chi(T)=\frac{T}{\pi}\log(Nm^d)+\frac{dT}{\pi}\log T
                   +O_\Pi(T+\log(m(T+2))).
\]
Indeed the conductor is $Nm^d$ since $(m,N)=1$, the gamma parameters are fixed for even twists, and the coefficients satisfy the divisor bound from the elliptic good-prime parameters. Apply the usual argument-principle zero count as in Iwaniec--Kowalski, \emph{Analytic Number Theory}, Theorem~5.8; the factor of two relative to the one-sided count comes from including both ordinates. Thus $N_\chi(T)=(1+o(1))B$ uniformly at $m=q\in(Q,2Q)$ and $N_\chi(T)\le(1+o(1))B$ uniformly at every conductor $m\mid q$.

The number of even primitive characters of an eligible conductor $q$ is $(\phi^*(q)+O(1))/2$. Indeed the sum of $\chi(-1)$ over primitive characters is, by M\"obius inversion, supported in the inner ordinary character sum only at the divisors $1$ and $2$, hence is bounded by two. Furthermore $\phi^*(q)\gg\sqrt q$ uniformly on eligible moduli, as follows from $\phi^*(p)=p-2$ and $\phi^*(p^a)=p^{a-2}(p-1)^2$ for $a\ge2$. Consequently each outer term of the right-hand sum in \eqref{cyt:input}, before the factor $k_d$, is $(1+o(1))B$ uniformly.

Now every even character modulo $q$, including one induced from a proper divisor, occurs once in the regular representation of $\operatorname{Gal}(M_q/\Q)$. Artin formalism applied to the elliptic Tate representation, or its symmetric square, yields the finite-place identity
\[
     L(s,\Pi/M_q)=\prod_{\substack{\chi\bmod q\\\chi(-1)=1}}
                                      L(s,\Pi\otimes\chi^{\mathrm{prim}}).
\]
Each factor here is primitive with its true local factors. All are entire and the archimedean factors are regular and nonzero in the critical strip. It follows that
\[
 Z_q^0(T)\ge\sum_\chi^{\flat}N^0_\chi(T),\qquad
 Z_q(T)\le (1+o(1))\frac{\phi(q)}2 B.
\]
If different factors share a zero, the multiplicities add, exactly as required. In particular this argument does not infer simplicity after base change.

Let $A_N$ be the mean of the indicator of eligible integers and $B_N$ the mean of $g(q)=\boldsymbol{1}_{q\ \mathrm{eligible}}\phi(q)/\phi^*(q)$, with value zero on ineligible integers. We next prove that $A_N/B_N=D_N$. For odd $p\nmid N$ the factor of $g$ is $(p-1)/(p-2)$ at exponent one and $p/(p-1)$ at all larger exponents. At $p\mid N$ positive exponents are excluded. At $p=2$ not dividing $N$, exponent one is excluded while all exponents at least two have factor two. The indicator is identically one at each odd prime outside $N$. The mean of $g$ at such $p$ is
\[
 b_p=1+\frac1{p(p-2)}-\frac1{p^2(p-1)(p-2)}.
\]
For instance, if $g=1*h$ is Dirichlet convolution with the constant-one function, then $h(p)=1/(p-2)$, $h(p^2)=-1/((p-1)(p-2))$, and higher powers vanish. This gives $b_p=1+h(p)/p+h(p^2)/p^2$. The series $\sum |h(n)|/n$ converges because its Euler factors are $1+O(1/p^2)$, apart from finitely many primes. Therefore the mean $B_N$ exists and equals the product of these local means: sum the convolution using $\lfloor X/d\rfloor/X\to1/d$ and dominated convergence. At primes dividing $N$ both means include $1-1/p$, which cancels. At $2\nmid N$ the indicator has mean $3/4$, whereas $g$ has mean $1/2+2/4=1$. These observations give $A_N/B_N=D_N$. The same means govern dyadic intervals by subtraction.

Multiply the inequalities for full field counts by $2/\phi^*(q)$ and sum over eligible $Q<q<2Q$. The ratio of these two weighted sums, critical to total, has lower limit at least $k_d(\omega)A_N/B_N$. Every summand in the denominator is positive for large $Q$. Some individual field therefore has ratio strictly larger than $c_d(\omega)D_N$ for every sufficiently large $Q$. The strict gap comes from $k_d(\omega)>c_d(\omega)$. Its degree $\phi(q)/2$ tends to infinity.

To verify the numerical claims one may keep the rational Euler factors $b_p$ through the odd integer $P=10001$ and use $b_p\le\exp(1/(p(p-2)))$ for the rest. The sum of $1/(n(n-2))$ over all odd $n>P$ telescopes to $1/(2P)$. Hence the inverse tail exceeds $1-1/(2P)$. Omitting primes dividing $N$ can only improve the lower bound. The supplied rational checker obtains respectively $0.172888$, $0.230518$, and $0.057629$ as downward-rounded lower bounds at $\omega=1/100000$. This proves the table.
\end{proof}
The conductor and the height both grow in this theorem. For $E_0:y^2=x^3-x$ the even-conductor line applies. The independence of the choice of $q$ from the twisting character is furnished by summing over the complete character group of the selected field. The theorem does not assign these proportions to the zeros at unbounded heights over any one fixed $M_q$.

\clearpage
\phantomsection\addcontentsline{toc}{part}{V. The global function field theorem}
\noindent{\small\sffamily\color{accent}V. The global function field theorem}\par
\section{The global function field theorem}\label{ff:global}
Let $C$ be a smooth, projective, geometrically connected curve of genus $g$ over $k=\mathbb F_q$, and set $K=k(C)$. The field $k$ is the full constant field of $K$. A closed point $x$ of $C$ has residue field of size $q_x=q^{d_x}$, where $d_x=[k(x):k]$. Put $T=q^{-s}$. For a reciprocal root $\gamma$ of a polynomial in $T$ we have
\begin{equation}\label{ff:dictionary}
  1-\gamma q^{-s}=0
\quad\Longleftrightarrow\quad
  \Re(s)=\frac{\log|\gamma|}{\log q},\qquad
  \Im(s)\log q\equiv \arg\gamma\pmod{2\pi}.
\end{equation}
Thus one reciprocal root describes an entire arithmetic progression of zeros, all of the same order as the corresponding root of the polynomial. We take all Euler factors, including those at ramified places. The word \emph{completed} in this chapter has that meaning. There are no archimedean places.

\begin{theorem}[Automorphic RH over a global function field]\label{ff:main}
Let $\pi$ be a cuspidal automorphic representation of $\mathrm{GL}_r(\mathbb A_K)$ with unitary central character, in unitary normalization. The completed standard $L(s,\pi)$ is a rational function of $T=q^{-s}$. Every zero has real part $1/2$. It is a polynomial with constant coefficient $1$, except that for a degree character on $\mathrm{GL}_1$ it is a rotation of the zeta function of $C$. In this exception its poles are on the lines of real parts $0$ and $1$.

For two such representations, the completed Rankin--Selberg function $L(s,\pi\times\pi')$ is rational and every zero has real part $1/2$. Its poles occur precisely when $\pi'$ is the contragredient of $\pi$ twisted by a unitary degree character. In that case the two lines of poles have real parts $0$ and $1$.
\end{theorem}
The deep inputs are Grothendieck's trace formula, Deligne's theorem on weights, and Lafforgue's correspondence with compatibility at every place. We give the full deduction, including the numerator degree, the locations of all exceptional poles, and an explicit constant in the prime-sum bound.

\subsection{The determinant on the complete curve}\label{ff:determinant}
Choose a prime $\ell\nmid q$ and a complex embedding of the coefficient field for measuring weights. On a dense open $j:U\hookrightarrow C$ take a lisse $\overline{\mathbb Q}_\ell$-sheaf $\mathcal V$ of rank $r$, with geometric generic fibre $V$. Equivalently it is a continuous representation unramified on $U$. Write $I_x$ for inertia and $F_x$ for geometric Frobenius on $V^{I_x}$. The geometric Frobenius on the cohomology of $C_{\bar k}$ will be denoted by $F$. Its eigenvalue on $\overline{\mathbb Q}_\ell(-1)$ over $k$ is $q$. The Euler product is
\begin{equation}\label{ff:euler}
 L_C(T,V)=\prod_{x\in |C|}\det(1-T^{d_x}F_x\mid V^{I_x})^{-1}.
\end{equation}
Each coefficient at $T=0$ involves finitely many places. The stalk of the ordinary direct image $j_*\mathcal V$ over $x$ identifies with $V^{I_x}$. In Weil--Deligne notation $(r_x,N_x)$, the vector space in this determinant is $(\ker N_x)^{I_x(r_x)}$. Indeed, on a sufficiently small open subgroup of inertia the original action is the exponential of the tame character times $N_x$. Invariance under this subgroup gives $\ker N_x$. On this kernel the remaining action of inertia is given by $r_x$. Frobenius semisimplification leaves the determinant unchanged. This distinction matters at every place of multiplicative reduction of an elliptic curve.

Say that $\mathcal V$ is pointwise pure of weight $w$ if for each $x\in |U|$ every Frobenius eigenvalue on $V$ has absolute value $q_x^{w/2}$ in the chosen complex embedding. For the geometric sheaves below the assertion holds in every embedding. Write
\[
 H^i_{\rm mid}(V)=H^i(C_{\bar k},j_*\mathcal V),\quad
 h_i=\dim H^i_{\rm mid}(V),\quad
 P_i(T)=\det(1-TF\mid H^i_{\rm mid}(V)).
\]
On a smooth curve the associated perverse middle extension is $(j_*\mathcal V)[1]$. Our notation uses ordinary cohomological degrees $0,1,2$.

\begin{theorem}[Rationality, purity and the degree]\label{ff:cohomology}
For a sheaf pointwise pure of weight $w$,
\begin{align}
 L_C(T,V)&=\frac{P_1(T)}{P_0(T)P_2(T)},\label{ff:ratio}\\
 |\gamma|&=q^{(w+i)/2}\qquad
   (\gamma\text{ an eigenvalue on } H^i_{\rm mid}(V)).\label{ff:weight}
\end{align}
The three polynomials have pairwise disjoint sets of zeros. Let $G_{\rm geom}=\pi_1(U_{\bar k})$, acting on $V$. The outside terms are
\[
 H^0_{\rm mid}(V)=V^{G_{\rm geom}},\qquad
 H^2_{\rm mid}(V)=V_{G_{\rm geom}}(-1).
\]
A subscript denotes coinvariants. With $a_x(V)=r-\dim V^{I_x}+\operatorname{Swan}_x(V)$,
\begin{equation}\label{ff:degree}
 h_1=h_0+h_2+r(2g-2)+\sum_{x\in|C|} d_x a_x(V).
\end{equation}
\end{theorem}
\begin{proof}
Apply the Grothendieck trace formula on the proper curve $C$ to $j_*\mathcal V$. The coefficient of $T^n$ in $T\,d\log L_C/dT$ is
\[
 S_n(V)=\sum_{d_x\mid n}d_x\,\operatorname{tr}
 (F_x^{\,n/d_x}\mid V^{I_x}).
\]
A closed point of degree $d_x$ contributes $d_x$ points over $\mathbb F_{q^n}$ exactly when $d_x\mid n$. The trace formula equates this coefficient to
$\sum_{i=0}^2(-1)^i\operatorname{tr}(F^n\mid H^i_{\rm mid}(V))$.
The determinant identity follows since both formal series in~\eqref{ff:ratio} have constant coefficient $1$.

A section over $C_{\bar k}$ of $j_*\mathcal V$ is a section over $U_{\bar k}$, giving the invariants. The quotient of $j_*\mathcal V$ by $j_!\mathcal V$ is supported on finitely many points. Over $\bar k$ its positive-degree cohomology vanishes. Hence $H^2_{\rm mid}(V)=H^2_c(U_{\bar k},\mathcal V)$, which is $V_{G_{\rm geom}}(-1)$ by Poincaré duality.

The assertion about weights is Deligne's theorem for a complete nonsingular curve. In middle degree both weight inequalities can be seen separately. The point-supported quotient makes $H^1_c(U_{\bar k},\mathcal V)\to H^1_{\rm mid}(V)$ surjective. The low-degree sequence for $Rj_*$ makes $H^1_{\rm mid}(V)\to H^1(U_{\bar k},\mathcal V)$ injective. Thus the middle term is the image of the map from compactly supported to ordinary cohomology. Deligne's bound on compactly supported cohomology gives weights at most $w+1$ for the source. Apply that bound to the dual sheaf of weight $-w$ and use Poincaré duality to get weights at least $w+1$ for the target. Their common subquotient has weight $w+1$. For the theorem, including the outside degrees, see P.~Deligne, \href{https://doi.org/10.1007/BF02684780}{\emph{La conjecture de Weil II}}, §3.2.3. Distinct $i$ give distinct absolute values.

Finally the Grothendieck--Ogg--Shafarevich formula on $U$ gives
\[
 \chi_c(U_{\bar k},\mathcal V)
 =r\left(2-2g-\sum_{x\in C\setminus U}d_x\right)
   -\sum_{x\in C\setminus U}d_x\operatorname{Swan}_x(V).
\]
Adding the Euler characteristic of the point-supported quotient adds
$\sum_{x\in C\setminus U}d_x\dim V^{I_x}$. Equating the result with
$h_0-h_1+h_2$ gives~\eqref{ff:degree}. For the curve formulas see N.~Katz,
\href{https://web.math.princeton.edu/~nmk/arizona34.pdf}{\emph{$L$-functions and monodromy: four lectures on Weil II}}, Lecture~1.
\end{proof}

\begin{corollary}\label{ff:sum}
Every zero of $L_C(q^{-s},V)$ has real part $(w+1)/2$. Every pole has real part $w/2$ or $(w+2)/2$. For each $n\ge1$,
\[
 \left| S_n(V)-\operatorname{tr}(F^n\mid H^0_{\rm mid}(V))
                  -\operatorname{tr}(F^n\mid H^2_{\rm mid}(V))\right|
 \le h_1q^{n(w+1)/2}.
\]
For a finite set $B$ of places, deleting its local factors may add zeros with real part at most $w/2$.
\end{corollary}
\begin{proof}
Apply~\eqref{ff:dictionary} to the three disjoint weight circles. The bound follows by the triangle inequality for the $h_1$ middle eigenvalues counted algebraically. Deleting a factor multiplies by
$\det(1-T^{d_x}F_x\mid V^{I_x})$. The middle-extension stalk at $x$ has weights at most $w$ by Deligne's local bound (\emph{Weil II}, §1.8). For an eigenvalue $\eta$ on this stalk the reciprocal roots introduced have absolute value $|\eta|^{1/d_x}\le q^{w/2}$. At a pole some factors may cancel.
\end{proof}

Conversely, given the determinant identity, a uniform bound
$|\operatorname{tr}(F^n\mid H^1_{\rm mid})|\le A R^n$ forces every middle eigenvalue to have modulus at most $R$. Indeed, the generating series $\sum_{n\ge1}\operatorname{tr}(F^n)z^n$ is analytic for $|z|<R^{-1}$. Its rational expression is $\sum_\gamma m_\gamma\gamma z/(1-\gamma z)$ over distinct eigenvalues with positive integral multiplicities. Every nonzero $\gamma$ gives an actual pole at $1/\gamma$. This argument uses the uniform estimate for all $n$ and permits cancellations among traces for individual $n$. For the matching lower bound, cohomological duality pairs the middle group for $V$ with that for $V^\vee$, with a Tate twist. Deligne's theorem supplies the upper bound for both.

\subsection{Cuspidal representations and arbitrary unitary twists}\label{ff:cuspidal}
Use the unitary Satake normalization, in which the unramified parameters of a cuspidal representation with finite-order central character have absolute value one. Write $\omega_z(a)=z^{\deg(a)}$ for the character of $\mathbb A_K^\times/K^\times$ taking the value $z^{d_x}$ on a local uniformizer of degree $d_x$. A twist on $\mathrm{GL}_r$ is through its determinant. The identity of local factors gives
\begin{equation}\label{ff:twist}
       L_C(T,\pi\otimes\omega_z)=L_C(zT,\pi).
\end{equation}
At a ramified place this character scales Frobenius on the same inertia-fixed space.

First suppose that the central character has finite order. Lafforgue gives an irreducible sheaf $V_\pi$ of rank $r$, unramified off finitely many places, with finite-order determinant and pure of weight $0$. Global-local compatibility identifies its Frobenius-semisimplified Weil--Deligne representation with the parameter of $\pi_x$ at every place. Hence $L_C(T,\pi)=L_C(T,V_\pi)$. We use L.~Lafforgue, \href{https://doi.org/10.1007/s002220100174}{\emph{Chtoucas de Drinfeld et correspondance de Langlands}}, Theorem~VI.9, Corollary~VII.5 and Theorem~VII.6. At unramified places the identity uses the full representation. At the other places it uses the kernel of monodromy followed by inertia invariants.

The $G_{\rm geom}$-invariants of an irreducible arithmetic representation are stable under the arithmetic group. Thus they vanish unless all of $V_\pi$ is geometrically trivial. In that case the representation factors through $\operatorname{Gal}(\bar k/k)$, which is procyclic. An irreducible representation of this group over an algebraically closed coefficient field has dimension one: a Frobenius eigenline is stable under Frobenius and under its closure. Applying the same argument to the dual handles coinvariants. In the geometrically nontrivial case, $P_0=P_2=1$ and every reciprocal root of $P_1$ has modulus $q^{1/2}$. In rank one with a geometrically trivial character, let $z$ be the Frobenius eigenvalue. Then the product is
\[
 Z_C(zT)=\frac{\det(1-zTF\mid H^1(C_{\bar k},\overline{\mathbb Q}_\ell))}
                    {(1-zT)(1-qzT)}.
\]
Rank one with nontrivial geometric monodromy belongs to the polynomial case. The argument uses arithmetic irreducibility. Higher-rank sheaves induced across a constant extension can be geometrically reducible.

The degrees of closed points have greatest common divisor $1$. For instance the Weil bound gives points over $\mathbb F_{q^n}$ for every sufficiently large $n$, and every such $n$ is divisible by this greatest common divisor. Thus $\deg$ on the idèle class group maps onto $\mathbb Z$. Its degree-zero subgroup is compact: the image of the product of the local units is compact, and its quotient in the degree-zero subgroup is a subgroup of the finite group $\mathrm{Pic}^0(C)(k)$. A smooth central character $\nu$ has finite image on this compact subgroup, since its kernel there is open. Choose an idèle class $a$ of degree one. For unitary $\nu$, choose $z$ on the unit circle with $z^r\nu(a)=1$. Then the central character of $\pi\otimes\omega_z$ is finite on the degree-zero group and trivial on $a$, hence finite on the whole idèle class group. Formula~\eqref{ff:twist} rotates the $T$-plane without changing any modulus. This proves the standard-function assertion of Theorem~\ref{ff:main}, including potentially transcendental $z$.

If $\pi_0$ is unitary and $u\in\mathbb C$, then $\pi=\pi_0\otimes |\det|^{u}$ has $L(s,\pi)=L(s+u,\pi_0)$, and its zeros have real part $1/2-\Re(u)$. An isobaric sum of unitary cuspidal data has standard function the product of its summands' standard functions. Their numerator zeros have real part $1/2$ and their possible poles have real part $0$ or $1$. The multiplicities of zeros therefore add. For arbitrary residual representations the conclusion changes. The trivial representation of $\mathrm{GL}_2(\mathbb A_K)$ has parameters $q_x^{1/2},q_x^{-1/2}$ and function
\[
                 Z_C(q^{1/2}T)Z_C(q^{-1/2}T).
\]
For $g>0$ it has zeros on the two lines of real parts $1$ and $0$. Its poles have real parts $-1/2$, $1/2$ and $3/2$, so they cannot remove these zeros.

\subsection{Products, exceptional poles, and the field of constants}\label{ff:products}
Let $V$ and $W$ denote the two irreducible sheaves for cuspidal data of finite-order central character, using a common $U$. Lafforgue's Theorem~VI.9(ii) identifies the local Rankin--Selberg factor at every $x$ with the factor on $V\otimes W$. This includes places where both representations are ramified. The local monodromy operator is $N_V\otimes1+1\otimes N_W$, and the kernel and inertia invariants are taken after forming the tensor product. For example, a nontrivial ramified character paired with its inverse gives a trivial character. Both of its original invariant spaces vanish at that place, but the tensor-product factor is $(1-T^{d_x})^{-1}$. The tensor product on $U$ is pure of weight zero, since its unramified Frobenius eigenvalues are products of eigenvalues of the two factors. Theorem~\ref{ff:cohomology} proves the zero assertion. Twisting either cusp form by a unitary degree character gives the same rotation as in~\eqref{ff:twist}.

We determine the denominator explicitly. Let $H$ be geometric monodromy in the sense that all intertwining maps below commute with $G_{\rm geom}$. Then
\begin{equation}\label{ff:hom}
           (V\otimes W)^H=\operatorname{Hom}_H(V^\vee,W).
\end{equation}
The right-hand side is stable under arithmetic Frobenius. If nonzero it contains a Frobenius eigenvector $f$ of eigenvalue $\alpha$. The relation is
$F_W f=\alpha f F_{V^\vee}$.
Consequently $f$ intertwines the actions of both $H$ and $F$ on $V^\vee\otimes\omega_\alpha$ and $W$. The parameter $\alpha$ is an $\ell$-adic unit because the Frobenius action and its inverse preserve a lattice over a finite coefficient extension. The character of integral powers of Frobenius with parameter $\alpha$ therefore extends continuously to $\operatorname{Gal}(\bar k/k)$. The group generated by $H$ and a lift of $F$ is dense in the arithmetic fundamental group, so $f$ is an arithmetic intertwiner. It is an isomorphism by irreducibility. Equation~\eqref{ff:weight} on $H^0$ gives $|\alpha|=1$ in the complex embedding. Conversely, every such isomorphism gives an invariant vector in~\eqref{ff:hom}. Applying the argument to the two duals classifies $H^2$ as well. Local compatibility and uniqueness in the global correspondence identify this twist condition with $\pi'\simeq\pi^\vee\otimes\omega_\alpha$. After undoing the preliminary finite-central-character twists one obtains the condition for arbitrary unitary cuspidal data. This proves the rest of Theorem~\ref{ff:main}.

There can be several poles in a single imaginary period. To compute their number, restrict $V^\vee$ to $H$ and choose a nonzero irreducible $H$-submodule. The sum of its arithmetic conjugates is a nonzero arithmetic-stable subspace and so exhausts $V^\vee$. A sum of simple submodules is semisimple. There are only finitely many isomorphism classes of simple constituents. The action of $F$ on these is transitive, since a union of orbits gives an arithmetic-stable sum of isotypic components. Let $m$ be the size of the orbit. Each isotypic component has multiplicity one. Here are the details. Write a chosen component as $A\otimes D$, with $H$ acting only on the irreducible $A$. The operator $F^m$ preserves it. Since $F^m$ conjugates the action on $A$ to an equivalent one, there is an intertwining automorphism $B$ of the underlying space of $A$ implementing this conjugation. By Schur's lemma $(B^{-1}\otimes1)F^m$ has the form $1\otimes D_0$ with $D_0$ invertible. Choose an eigenline $D_1$ of $D_0$. The sum of the $m$ translates by powers of $F$ of $A\otimes D_1$ is stable under $H$ and $F$. Irreducibility forces it to exhaust $V^\vee$. Its intersection with the chosen component has multiplicity one, proving the assertion.

Now $\operatorname{End}_H(V^\vee)$ has a basis of $m$ projectors onto these pairwise inequivalent simple summands, cyclically permuted by $F$. For $W=V^\vee\otimes\omega_\alpha$ the action on~\eqref{ff:hom} is $\alpha$ times that cyclic permutation. Applying this computation to the two duals gives the dual action on coinvariants, which is then twisted by $(-1)$ in degree two. We obtain
\begin{equation}\label{ff:poles}
              P_0(T)=1-(\alpha T)^m,\qquad
              P_2(T)=1-(q\alpha T)^m.
\end{equation}
Each factor has simple zeros. Changing the choice of $\alpha$ by a degree self-twist changes it by an $m$th root of unity and leaves these polynomials unchanged. Indeed, by the same Hom calculation the allowable scalars form exactly the eigenvalues of the cycle on the $m$ projectors, translated by $\alpha$. In the common case of geometric irreducibility $m=1$. Allowing $m>1$ accounts for poles from representations induced across constant extensions.

This reasoning extends to any finite list of pure sheaves and their duals. Tensoring sheaves of weights $w_1,\ldots,w_b$ gives weight $\sum_iw_i$ on a common open, with dualization reversing the weight. A Schur functor homogeneous of degree $h$ sends a pure sheaf of weight $w$ to one of weight $hw$, provided its fibre is nonzero. One can see this at each unramified place by triangularizing Frobenius and using that each resulting eigenvalue is a product of $h$ eigenvalues of the source. Apply the middle-extension construction to this resulting sheaf. All zeros of its completed function have real part $(hw+1)/2$. The argument proves the weight assertion even in instances where one does not use an automorphic lift for the particular functor.

\subsection{Constant extension and counts of primes}\label{ff:constantbase}
For $b\ge1$, let $C_b=C\otimes_k\mathbb F_{q^b}$ and let $V_b$ be the pullback of $V$. The Frobenius on every cohomology group over the new constant field is $F^b$. The geometric cohomology groups agree since extending the constants is étale and preserves the inertia groups used in the stalks of $j_*\mathcal V$. If $\gamma$ runs over the eigenvalues for $C$ in degree $i$, then
\[
                     P_{i,b}(X)=\prod_\gamma(1-\gamma^bX).
\]
Consequently the closed Euler products satisfy the exact factorization
\begin{equation}\label{ff:constantidentity}
       L_{C_b}(T^b,V_b)=\prod_{\zeta^b=1}L_C(\zeta T,V).
\end{equation}
In fact $\prod_{\zeta^b=1}(1-\zeta\gamma T)=1-(\gamma T)^b$ separately in each cohomological degree. Locally a point of degree $d$ splits into $\gcd(b,d)$ points of degree $d/\gcd(b,d)$ over the new constants. The new Frobenius at each of them is the old one to the power $b/\gcd(b,d)$ on the same inertia invariants. This gives the same product directly. In particular the weight in~\eqref{ff:weight} is independent of $b$ when absolute values on $C_b$ are measured using $q^b$. The determinant identity describes possible increases of multiplicity when distinct eigenvalues acquire the same $b$th power.

The trace bound also gives a statement for prime places with no prime powers in the sum. Write $N_d=\#\{x\in |C|:d_x=d\}$ and
\[
 B_n(V)=\sum_{\deg x=n}\operatorname{tr}(F_x\mid V^{I_x}),\qquad
 M_n(V)=\operatorname{tr}(F^n\mid H^0_{\rm mid}(V))
               +\operatorname{tr}(F^n\mid H^2_{\rm mid}(V)).
\]
For $n\ge1$, explicit bounds holding for every genus and every rank are
\begin{equation}\label{ff:primes}
 |nB_n(V)-M_n(V)|
 \ \le\ \left(h_1+2(1+g)r\,\frac{q}{q-1}\right)q^{n(w+1)/2}.
\end{equation}
Indeed the difference $S_n(V)-nB_n(V)$ is supported on proper divisors $d$ of $n$. Each eigenvalue in a term of degree $d$ has, after raising to the power $n/d$, absolute value at most $q^{nw/2}$. The Weil bound for the curve gives
$dN_d\le \#C(\mathbb F_{q^d})\le q^d+1+2gq^{d/2}\le2(1+g)q^d$.
Every proper divisor is at most $\lfloor n/2\rfloor$ and the sum of $q^d$ over this whole range is at most $q^{n/2}q/(q-1)$. Adding Corollary~\ref{ff:sum} proves~\eqref{ff:primes}. For a geometrically nontrivial cuspidal standard sheaf the two main terms vanish. For an exceptional Rankin--Selberg pair,~\eqref{ff:poles} gives
\[
              M_n=
              \begin{cases}m\alpha^n(1+q^n),&m\mid n,\\
                            0,&m\nmid n.
              \end{cases}
\]
Thus even the main term for prime degrees can depend on a congruence class of the degree. A bound for all $n$ has to retain that dependence.

\section{Elliptic curves over global function fields}\label{ff:elliptic}
Let $E$ be an elliptic curve over $K=k(C)$. On the nonempty open of good reduction its smooth proper model $f:\mathcal E\to U$ yields the rank-two sheaf $V=R^1 f_*\overline{\mathbb Q}_\ell$. At a place $x$ of good reduction its factor is
\[
             (1-a_xT^{d_x}+q_xT^{2d_x})^{-1},
       \qquad a_x=q_x+1-\#E_x(k(x)).
\]
The two eigenvalues have absolute value $q_x^{1/2}$ by the Riemann hypothesis for an elliptic curve over a finite field. Hence $w=1$. At split and nonsplit multiplicative reduction the factor is respectively $(1-T^{d_x})^{-1}$ or $(1+T^{d_x})^{-1}$. At additive reduction it is $1$. These are the standard local computations on the inertia invariants of elliptic cohomology. The conductor still records wild ramification at additive places. For these local results and the global function see D.~Ulmer, \href{https://arxiv.org/abs/1101.1939}{\emph{Elliptic curves over function fields}}, Lecture~1, §§8--9.

\begin{theorem}\label{ff:ec}
For every elliptic curve over $k(C)$, every zero of the arithmetic Hasse--Weil function $L(E/K,s)$ has real part $1$. Its possible poles have real parts $1/2$ and $3/2$. For a nonisotrivial curve it is a polynomial in $T=q^{-s}$ of degree
\[
                         4g-4+\deg\mathfrak n_E,
     \qquad \deg\mathfrak n_E=\sum_xd_xa_x(V).
\]
The unitarily normalized function is $L(E/K,s+1/2)$, with its zeros on the line of real part $1/2$. More generally, the function of the middle extension of $\operatorname{Sym}^h V$ has its zeros on $\Re(s)=(h+1)/2$ for every integer $h\ge0$.
\end{theorem}
\begin{proof}
The product just described agrees everywhere with~\eqref{ff:euler}. Apply Theorem~\ref{ff:cohomology} with $w=1$. The shift $s\mapsto s+1/2$ divides each local eigenvalue by $q_x^{1/2}$. The symmetric-power result follows by applying the tensor construction above to the good fibres before taking the middle extension.

For nonconstant $j(E)$, the $K\bar k/\bar k$-trace of $E$ is zero. A nonzero trace would make $E$ isogenous to a constant elliptic curve, and elliptic curves isogenous to a constant curve have constant $j$. For clarity, after passing to an algebraic closure, the inseparable part of an isogeny is an iterate of relative Frobenius; the separable part has finite kernel of geometric torsion points. All torsion points on a constant elliptic curve over an algebraically closed extension already come from the algebraically closed constant field. Quotienting by this kernel therefore again gives a constant curve. The Lang--Néron theorem now makes $E(K\bar k)$ finitely generated, so its $\ell$-power torsion is bounded. In particular the geometric invariants of $T_\ell E$ vanish. The Weil pairing identifies its dual over the geometric group with itself (the Tate twist is geometrically trivial). Invariants and dual invariants in $R^1f_*$ therefore vanish as well. Formula~\eqref{ff:degree} gives the degree. For the trace and finite generation statements see Ulmer, Lecture~1, §§5--7, and S.~Lang and A.~Néron, \emph{Rational points of abelian varieties over function fields}, Amer.\ J.\ Math.\ 81 (1959), 95--118.
\end{proof}

A constant curve $E_0/k$ illustrates why the denominator is present. If the eigenvalues of its Frobenius are $\alpha,\beta$, then $\alpha\beta=q$ and $|\alpha|=|\beta|=q^{1/2}$. At $x$ the two eigenvalues are $\alpha^{d_x}$ and $\beta^{d_x}$. Multiplication over all places gives
\begin{equation}\label{ff:constantE}
                  L(E_0/K,T)=Z_C(\alpha T) Z_C(\beta T).
\end{equation}
For the geometrically constant sheaf, $h_0=h_2=2$ and $h_1=4g$. If an isotrivial curve is a ramified twist of a constant curve, the geometric group can act nontrivially and those dimensions must be computed for the twist. For example twisting by a geometrically nontrivial quadratic character in odd characteristic makes $-1$ act on the constant two-dimensional fibre, which kills all geometric invariants and coinvariants.

\subsection{A constant example with a double central zero}\label{ff:constant-example}
Take $q=5$, $C=E_0$ and $E_0:y^2=x^3-x$. The values of $x^3-x$ at $x=0,1,2,3,4$ are $0,0,1,4,0$. Counting square roots and including infinity gives $\#E_0(\mathbb F_5)=8$. Thus $\alpha+\beta=-2$, $\alpha\beta=5$, and
\[
              Z_C(T)=\frac{1+2T+5T^2}{(1-T)(1-5T)}.
\]
In~\eqref{ff:constantE}, the numerator has reciprocal roots
$\alpha^2,\alpha\beta,\beta\alpha,\beta^2$. Since
$\alpha^2+\beta^2=-6$, direct multiplication gives
\begin{equation}\label{ff:double}
 L(E_0/\mathbb F_5(C),T)
 =\frac{(1+6T+25T^2)(1-5T)^2}
 {(1+2T+5T^2)(1+10T+125T^2)}.
\end{equation}
The numerator has the predicted degree four. Its root $T=1/5$ has multiplicity two, giving arithmetic central order two. Its remaining reciprocal roots also have modulus five. The reciprocal roots of the two denominator factors have moduli $\sqrt5$ and $5\sqrt5$. This example distinguishes poles arising from the constant elliptic curve from the zeros arising from degree-one cohomology on the base.

\subsection{A nonisotrivial example including the place at infinity}\label{ff:legendre}
Let $K=\mathbb F_5(t)$ and start with the Legendre curve
\[
               E_t:\quad y^2=x(x-1)(x-t).
\]
Its discriminant is $16t^2(t-1)^2$ and its $c_4$ is $16(t^2-t+1)$, so the finite places of bad reduction are exactly $0$ and $1$. At both places the reduction is multiplicative. At $0$ the tangent slopes at the node have squares $-1$, which is a square in $\mathbb F_5$. At $1$ their squares are $1$. Both are split. With $u=t^{-1}$, substituting $x=tX$, $y=tY$ gives
\[
                     Y^2=tX(X-u)(X-1).
\]
Thus at infinity $E_t$ is the ramified quadratic twist of the split multiplicative Legendre curve $E_u$ at $u=0$. The two eigenvalues of inertia on the twisted representation for an element mapping nontrivially to the quadratic quotient are both $-1$, so its invariant space is zero. All these ramification actions are tame in characteristic $5$: the multiplicative action is tame unipotent, and the quadratic character has order two. The conductor exponents are $1,1,2$. Since $j(E_t)$ has a pole at $t=0$, the curve is nonisotrivial. Formula~\eqref{ff:degree} gives numerator degree $-4+4=0$, so the completed $L$-function is $1$.

Now take the quadratic twist, specified by
\[
               E':\quad y^2=(t-2)x(x-1)(x-t).
\]
The places $0$ and $1$ remain multiplicative, with signs multiplied by the quadratic characters of $-2$ and $-1$ in $\mathbb F_5$. These signs are respectively $-1$ and $+1$. The formerly good place $t=2$ acquires ramified quadratic inertia acting by $-1$, so its conductor exponent is $2$ and its Euler factor is $1$. At infinity the combined twist of $E_u$ is $t(t-2)=u^{-2}(1-2u)$. This is a square in $\mathbb F_5((u))$, since $1-2u$ is a unit congruent to $1$ and the residue characteristic is odd. Infinity is now split multiplicative with exponent $1$. All other places have good reduction. The conductor degree is $1+1+2+1=5$, so the polynomial has degree exactly one. Its linear coefficient is the sum of the traces over the six rational places, including infinity.

For a good finite fibre at $t=a$ this trace is
\[
   -\chi(a-2)\sum_{x\in\mathbb F_5}\chi(x(x-1)(x-a)),
\]
where $\chi$ is the quadratic character extended by $\chi(0)=0$. The following table includes the local computations just made at the four bad places.
\begin{center}
\begin{tabular}{c|rrrrrr}
place $a$ & $0$ & $1$ & $2$ & $3$ & $4$ & $\infty$\\\hline
$\operatorname{tr}(F_a\mid V^{I_a})$ & $-1$ & $1$ & $0$ & $2$ & $2$ & $1$
\end{tabular}
\end{center}
Their sum is $5$. Since $\log L$ and $L-1$ have the same first coefficient,
\begin{equation}\label{ff:linear}
                    L(E'/K,T)=1+5T.
\end{equation}
The reciprocal root is $-5$. Thus the zeros have real part $1$ and imaginary parts $(2j+1)\pi/\log 5$ for $j\in\mathbb Z$. The central point itself is nonzero. The normalized function has the factor $1+\sqrt5\,5^{-s}$. For every $b$, restriction to $\mathbb F_{5^b}(t)$ gives $1-(-5)^bX$ in the variable $X=(5^b)^{-s}$. In particular an even extension degree moves a zero to the actual central point without changing its real part.

The theorems in this section apply to global fields with a finite field of constants. Their proof uses a single Frobenius acting on the cohomology of the complete base curve. For an elliptic curve over a number field, the Frobenius at each good prime still controls the corresponding local factors, with the same elliptic weight. Deligne's theorem on the cohomology of $C_{\bar k}$ supplies no single global Frobenius determinant for that number-field Euler product. Consequently the theorem over $k(C)$ does not settle the number-field elliptic or automorphic RH.

\clearpage
\section{What an all-height proof would still have to establish}\label{close:main}
For the number-field problem, consider a single non-CM curve such as $11\mathrm a1$. One can build its exact Frobenius trace from successively larger torsion fields, detect every zero in a specified rectangle, and place zeros in an unconditional region at every height. To prove that the entire zero divisor lies on $\Re s=1$, the quantifiers have to meet. In Section~\ref{aw:section}, a bound on the trace sums for one fixed object on arbitrarily distant windows suffices. The family argument gives many twists on each window. Passing from those statements to a single prescribed twist would require a new uniform estimate. A counting argument cannot exchange these quantifiers.

Section~\ref{apw:section} closes part of that gap for each individual curve in the conductor census. For every $0<\eta<1$, there are arbitrarily large $Y$ on which the unsmoothed sum $\sum_{p\le x}a_p(E)$ has the required one-sided bound simultaneously for $Y\le x\le Y+Y^{1-\eta}$. The proof uses the absence of a real zero to the right of the line, which the finite census can establish at every real point. Theorem~\ref{rawfilter:main} reduces the sufficient interval to $[Y,YH(Y)]$ for any chosen $H(Y)\to\infty$. In $\log x$, the length of the unconditional interval tends to zero. Even arbitrarily slow divergence would suffice. If a rightmost nonreal zero exists off the line and no real zero is as far right, some fixed positive lengths are already possible by Theorem~\ref{cinv:attained}. This identifies a specific estimate for a fixed curve whose improvement would settle its ERH.

There are two other precise ways to express what remains. For standard functions and arbitrary Rankin--Selberg pairs, Theorem~\ref{li:everyheight} controls the contribution of every unverified zero to the first $T^2$ Li coefficients at any verified height $T$. ERH requires nonnegativity at every positive integer index. For the heat deformation the unknown is a single real number: Theorem~\ref{hg:negative} identifies ERH with $\lambda_\pi=0$, and Theorem~\ref{hc:automorphic} bounds it by $O_d(1/\log(3+\mathcal Q_*))$ uniformly at every height. Both theorems retain the contribution from the whole divisor. Neither estimate reaches the endpoint needed to solve the all-height problem.

The translated-product comparison now gives the exact rate required in the heat approach. For a character of degree one, Theorem~\ref{hol:profile} identifies the ordinary limit of $(\log\tau)\lambda^+_\tau$ with the undeformed displacement $\Delta$. A bound with a universal constant times $1/\log\tau$ permits a positive value for $\Delta$; a little-oh estimate even on a single sequence would force it to vanish. The comparison proves the existence and exact value of this limit in terms of the original divisor. Determining that it vanishes for one prescribed general elliptic curve remains the essential step.

The direct character average supplies a different numerical target for a proof. At real part $1/2+u$ its corrected projective two-scale difference is $O(N^{-2u+\epsilon})$. A bound at $N^{-4u}$ would exclude zeros to the right of unitary $1-u$ for $0<u\le1/4$. The full-strip formulation of Section~\ref{pc:section} also tracks which coincident zeta zeros could mask an elliptic zero after meromorphic continuation. The geometric classification in Theorem~\ref{albs:hyper} shows what happens if additional rational parameters are introduced: any such quadratic family carrying weight-one Frobenius traces factors through a hyperelliptic pencil. The analytic distribution of its parameter values is still available for a stronger estimate.

The finite computation has another instructive feature. At any height where the argument calculation works, the strip count is exact before the verifier starts finding sign changes. Equality with the number of brackets is a separate arithmetic fact. A nonreal pair or a multiple critical-line zero could obstruct equality. Extending the verified data to a sequence of unbounded heights for one fixed curve would imply its ERH. Our algorithms terminate with a theorem at each recorded height; the proofs provide no theorem that this exhaustion succeeds for unbounded heights. A stopping test remains a stopping test even when every completed test has passed.

At a real point, symmetry alone determines the parity of the central order. If the sign is negative a zero is forced. It takes additional arguments to prove simplicity or to force order two or higher for positive sign. The conductor-389 proof uses rational points to obtain the necessary exact vanishing, where numerical approximation at the midpoint cannot certify an exact zero. The base-extension arguments use a different mechanism: several entire factors can carry independently forced odd orders. The resulting order can grow proportionally to the field degree. This establishes increasingly large \emph{multiplicity} at a single point already known to lie on the line; it places no new bound on the real parts of other zeros.

For a general primitive standard $L$-function over a number field, Dirichlet factorization usually fails at the very first Euler coefficient. The strip bound, central parity, and the prime-sum criterion survive well beyond the range of an abelian argument. In positive characteristic, the representation on middle cohomology has a weight fixed by purity, and that weight gives the required absolute value for every reciprocal root. For elliptic curves over number fields we know the corresponding absolute values of the two \emph{local} good-prime eigenvalues. No analogous theorem in this work controls a global set of reciprocal roots of the completed function. This is where the complete function-field proof ends and the characteristic-zero problem remains.

\end{document}
